\documentclass[11pt]{article}
\usepackage[margin=1in]{geometry}
\usepackage{amsmath,amssymb,amsthm,mathtools}
\usepackage{enumitem}
\usepackage[hidelinks]{hyperref}
\usepackage[nameinlink,noabbrev]{cleveref}

\newtheorem{theorem}{Theorem}[section]
\newtheorem{lemma}[theorem]{Lemma}
\newtheorem{proposition}[theorem]{Proposition}
\newtheorem{corollary}[theorem]{Corollary}
\newtheorem{definition}[theorem]{Definition}
\newtheorem{remark}[theorem]{Remark}
\newtheorem{correction}[theorem]{Correction}
\newtheorem{consequence}[theorem]{Consequence}
\newtheorem{assessment}[theorem]{Assessment}
\newtheorem{firewall}[theorem]{Firewall}
\newtheorem{programme}[theorem]{Programme}
\newtheorem{audit}[theorem]{Audit}

\newcommand{\Tr}{\operatorname{Tr}}
\newcommand{\spec}{\operatorname{spec}}
\newcommand{\sgn}{\operatorname{sgn}}

\title{Renewal Standard Model--Einstein Continuum Research Notes\\
Eighth Series, V1}
\author{}
\date{\today}

\begin{document}
\maketitle
\tableofcontents

\section{Context, archive boundary, and purpose}
\label{sec:v8v1}
\label{sec:e8v1-context}

These notes continue a constructive programme aimed at a
cutoff-independent Standard Model--Einstein continuum.  The intended object
must carry a coherent chiral fermion measure, the Standard Model gauge and
Higgs sectors, a gravitational measure and stress response, clustering and
reflection data, regulator independence, and a reconstruction or Lorentzian
interpretation.  That object has not been constructed.

The previous active source reached V100 and has been frozen byte for byte as
\begin{center}
\texttt{Renewal\_SM\_Einstein\_Continuum\_Research\_Notes\_V100\_f7.tex}.
\end{center}
It has \(40\,226\) logical source lines, \(1\,570\,671\) bytes, and
SHA--256 digest
\begin{center}
\texttt{5F28040CC4F306FE0BA60AE765429A956A3E74F8D1FDDED50F0F7CD5B878A66C}.
\end{center}
The earlier frozen files are retained as \(\texttt{\_f2}\) through
\(\texttt{\_f6}\).  This compact eighth-series V1 replaces the active V100
name, summarizes the proved mathematical interfaces, and continues at the
specific bottleneck selected by the V100 companion audit.  Subsequent
versions append to the active source and replace only its predecessor.  The
newest Yang--Mills and Navier--Stokes notes are audited at V5 and every fifth
version.  At V100 the same archive-and-restart rule will be applied with the
next unused \(\texttt{\_fN}\) suffix.

Throughout these notes, a conditional theorem proves an implication from
its displayed hypotheses.  It does not assert that the complete physical
regulator supplies those hypotheses.  Finite-cutoff identities, analytic
loaders, and continuum-existence statements are kept separate.

\section{Major mathematical results obtained so far}
\label{sec:e8v1-summary}

\subsection{Finite-cutoff geometry, chirality, and topology}

The archived programme constructs admissible finite lattice gauge sectors,
obtains a volume-uniform spectral window for the Wilson kernel under a
plaquette admissibility bound, and consequently controls the overlap sign,
Ginsparg--Wilson operator, moving chiral projectors, and their exponential
locality.  Exact Schur--Berezin and Riesz-projector formulas retain crossing
and zero modes as finite Grassmann sectors before the complementary block is
inverted.  A simple transverse crossing has linear least singular value,
a determinant sign jump, and a path-dependent source limit unless the zero
mode is saturated.

The determinant is treated as a section of a determinant line.  Its finite
Berry connection and curvature are expressed through the moving projector;
radial homotopy gives local primitives, while resolvent estimates give
exponential curvature locality under a uniform finite-range gap.  Local
triviality does not prove global holonomy triviality.

For one Standard Model generation in the all-left-handed convention, the
local \(SU(3)^3\), \(SU(3)^2U(1)\), \(SU(2)^2U(1)\), \(U(1)^3\), and
mixed gravitational--\(U(1)\) anomaly coefficients cancel, and the number
of left-handed \(SU(2)\) doublets is even.  This is the exact finite
representation arithmetic; it is not a construction of the global chiral
measure.

\subsection{Higgs reduction, fermion blocks, and determinant control}

On a Higgs chart separated from zero, the normalized Higgs doublet reduces
the electroweak group to \(U(1)_{\rm em}\), matches the charged left and
right representations, and isolates the unmatched minimal neutrino row.
The improved Ginsparg--Wilson field assembles the kinetic and Yukawa corners
into massive-overlap blocks.  A nonzero mass gives a uniform inverse bound,
but the probability of the nonzero-Higgs chart and the full phase gluing
remain open.

Relative determinant identities show that subtraction must be assembled
before an absolute trace norm is taken.  Finite-rank transfer, Schur
factorization, reduced pseudodeterminants, trace-class factorizations, and
local counterterm ownership give rigorous finite loaders.  The extensive
continuum trace-class and global determinant-line rows are unproved.

\subsection{Bosonic regulators, constraints, and long-range structure}

Positive-type compact-group plaquette weights give a finite pure-gauge
reflection-positive admissible regulator.  Metric normalization, heat and
symbol expansions assign the local cosmological, Einstein--Hilbert, gauge,
Higgs, curvature, and boundary divergences to explicit counterterm owners.
York--Lichnerowicz constraint charts and multipole expansions show that
elliptic constraint substitution has polynomial Coulomb tails; appropriate
monopole and dipole cancellations are needed before absolute summability.

The programme also records genuine obstructions: the real Euclidean
conformal direction is unbounded, a finite curvature list is not generally
closed under perturbative renormalization, and the perturbative
Gaussian-boundary hypercharge trajectory does not itself give a complete
ultraviolet limit.  These statements require new gravitational or
ultraviolet input and do not exclude every completion with the same
infrared field content.

\subsection{Interacting carriers, local response, and cofinal landing}

Finite Fock and quasifree compressions, Nambu covariance formulas, exact
Gibbs Duhamel response, imaginary-time commutator expansions, replacement
channels, and Petz recovery maps have been derived at finite cutoff.  A
locally contractive channel fixed-point equation propagates recovery,
boundary, metric, and mixed metric/refinement sources through an explicit
inhomogeneous resolvent.

The cofinal analysis separates base expectations from metric derivatives.
Summable adjacent rows yield a \(C^1\) limit on bounded metric charts.  A
single weighted master ledger covers a countable local observable algebra
and a countable tangent core; positivity and normalization then land a state
on the norm-closed inductive algebra, while a uniform dual-norm bound
completes the tangent core.

The weak Einstein residual is treated in a Banach space modulo the
finite-dimensional space of allowed local counterterm variations.  An
absolutely summable lifted cocycle gives a residual limit.  Regulator
independence needs a separate common-refinement or tail-diameter estimate.
A finite counterterm can normalize the limiting residual to zero precisely
when its quotient class vanishes.

\subsection{The V100 correction forced by the Yang--Mills audit}

The newest audited Yang--Mills result proves that the Ward zero of the
quasilocal perfect quadratic owner forces its absolute one-site Gaussian
influence row to be at least one.  Therefore a uniform strict one-site
Dobrushin margin is incompatible with convergence to that Ward-critical
owner.  The viable local-channel branch uses a row
\(1-m\tau_n+O(\tau_n^2)\): its inverse grows as \(\tau_n^{-1}\), while the
derivatives of a short physical-time channel carry a compensating
\(\tau_n\).  If every refinement-tagged source has an additional summable
amplitude \(\varepsilon_n\), the adjacent expectation and stress defects
are \(O(\varepsilon_n)\).

\section{Current open-gate ledger}
\label{sec:e8v1-open}

\begin{theorem}[Claim boundary at the eighth-series restart]
\label{thm:e8v1-open}
The archived results do not construct the Standard Model--Einstein
continuum.  The principal open gates are:
\begin{enumerate}[label=\textup{(G\arabic*)},leftmargin=4.5em]
\item a single regulator carrying the gauge-covariant chiral determinant,
zero-mode patches, global determinant-line trivialization, and compatible
reflection data;
\item physical short-time carrier or block channels with proved local
contraction, recovery, metric, and old/detail refinement estimates;
\item a summable refinement amplitude, uniform integrability outside the
protected charts, and completion beyond the countable core;
\item interacting photon Ward moments and gravitational multipole/contact
identities for every long-range edge;
\item coherent gravitational, classical, carrier, and composite-operator
counterterm rows, together with a proved zero residual class;
\item common-refinement estimates establishing regulator independence; and
\item a positive nonlinear gravitational measure, cutoff removal,
clustering, full reflection positivity, and reconstruction or Lorentzian
dynamics.
\end{enumerate}
\end{theorem}

\begin{proof}
Each item is an explicit unpaid row in the frozen theorem/debit ledgers.
None follows from the finite overlap gap, anomaly arithmetic, determinant
factorizations, or conditional channel-resolvent theorems.  Hence no
continuum conclusion can presently be drawn.
\end{proof}

\section{New attack: derivatives of a short-time local channel}
\label{sec:e8v1-short-time}

We now attack gate \textup{(G2)}.  Let \(I,J\subset\mathbb R\), let
\((s,q)\mapsto\mathcal L_{s,q}\) be a \(C^2\) family of bounded linear maps
on a finite carrier algebra, and assume
\begin{equation}
 \mathcal S_{s,q}(t):=e^{t\mathcal L_{s,q}}
 \label{eq:e8v1-semigroup}
\end{equation}
is CPTP for \(t\geq0\).  Define the one-step channel
\begin{equation}
 \mathcal T_{s,q}^{(\tau)}:=e^{\tau\mathcal L_{s,q}}.
 \label{eq:e8v1-step}
\end{equation}
Subscripts \(s,q,sq\) denote parameter derivatives at a fixed point.

\begin{theorem}[Exact first and mixed short-time derivatives]
\label{thm:e8v1-Duhamel}
For every \(\tau\geq0\),
\begin{align}
 \mathcal T_a^{(\tau)}
 &=\int_0^\tau
 \mathcal S(\tau-u)\mathcal L_a\mathcal S(u)\,du,
 \qquad a\in\{s,q\},
 \label{eq:e8v1-first-Duhamel}\\
 \mathcal T_{sq}^{(\tau)}
 &=\int_0^\tau
 \mathcal S(\tau-u)\mathcal L_{sq}\mathcal S(u)\,du
 \notag\\
 &\quad+\int_{\substack{r,t\geq0\\r+t\leq\tau}}
 \mathcal S(\tau-r-t)
 \bigl(\mathcal L_s\mathcal S(r)\mathcal L_q
      +\mathcal L_q\mathcal S(r)\mathcal L_s\bigr)
 \mathcal S(t)\,dr\,dt.
 \label{eq:e8v1-mixed-Duhamel}
\end{align}
In diamond norm,
\begin{align}
 \|\mathcal T_a^{(\tau)}\|_\diamond
 &\leq\tau\|\mathcal L_a\|_\diamond,
 \label{eq:e8v1-first-diamond}\\
 \|\mathcal T_{sq}^{(\tau)}\|_\diamond
 &\leq\tau\|\mathcal L_{sq}\|_\diamond
 +\tau^2\|\mathcal L_s\|_\diamond
          \|\mathcal L_q\|_\diamond.
 \label{eq:e8v1-mixed-diamond}
\end{align}
\end{theorem}

\begin{proof}
Differentiate the exponential by the Banach-algebra Duhamel formula,
which gives \eqref{eq:e8v1-first-Duhamel}.  Differentiating it in the other
parameter gives the direct \(\mathcal L_{sq}\) insertion and the two
possible orders of the \(\mathcal L_s,\mathcal L_q\) insertions.  Combining
the two triangular integration domains gives
\eqref{eq:e8v1-mixed-Duhamel}.

A CPTP map has diamond norm one.  The first integral therefore has norm at
most \(\tau\|\mathcal L_a\|_\diamond\).  In the mixed formula the direct
simplex has length \(\tau\).  Each ordered double-insertion simplex has area
\(\tau^2/2\); their sum has area \(\tau^2\).  Submultiplicativity proves
\eqref{eq:e8v1-mixed-diamond}.
\end{proof}

\subsection{Weighted local derivative kernels}
\label{sec:e8v1-weighted}

Let \(v_i\) be local seminorms on traceless Hermitian operators.  Suppose
there are nonnegative propagation matrices \(P(t)\) such that
\begin{equation}
 v_i(\mathcal S(t)X)\leq\sum_jP_{ij}(t)v_j(X),
 \qquad
 \|P(t)\|_\mu\leq e^{\kappa t}.
 \label{eq:e8v1-propagation}
\end{equation}
Suppose also
\begin{equation}
 v_i(\mathcal L_aX)\leq\sum_j(G_a)_{ij}v_j(X),
 \qquad g_a:=\|G_a\|_\mu<\infty,
 \quad a\in\{s,q,sq\}.
 \label{eq:e8v1-generator-kernel}
\end{equation}

\begin{corollary}[The channel derivative kernels carry one clock factor]
\label{cor:e8v1-channel-kernel}
Under \eqref{eq:e8v1-propagation}--
\eqref{eq:e8v1-generator-kernel}, the first channel derivatives have
weighted propagation-kernel norms
\begin{equation}
 a_a^{(\tau)}\leq\tau e^{\kappa\tau}g_a,
 \qquad a\in\{s,q\},
 \label{eq:e8v1-first-kernel}
\end{equation}
and the mixed derivative has a kernel with
\begin{equation}
 a_{sq}^{(\tau)}
 \leq e^{\kappa\tau}
 \bigl(\tau g_{sq}+\tau^2g_sg_q\bigr).
 \label{eq:e8v1-mixed-kernel}
\end{equation}
\end{corollary}

\begin{proof}
Apply the seminorm domination to each composition in
\eqref{eq:e8v1-first-Duhamel}--\eqref{eq:e8v1-mixed-Duhamel}.  The total
semigroup duration in every integrand is \(\tau\), so its propagation
factor is at most \(e^{\kappa\tau}\).  Integrating the one-dimensional and
two triangular domains gives the stated factors.
\end{proof}

Let \(\rho_{s,q}\) be an invariant density for the generator,
\begin{equation}
 \mathcal L_{s,q}(\rho_{s,q})=0.
 \label{eq:e8v1-invariant}
\end{equation}
Then \(\mathcal S(t)\rho=\rho\).  For a localization set \(Z\), define
\begin{equation}
 \gamma_a:=
 \bigl\|\bigl(v_i(\mathcal L_a\rho)\bigr)_i\bigr\|_{\mu,Z},
 \qquad a\in\{s,q,sq\}.
 \label{eq:e8v1-gamma}
\end{equation}

\begin{theorem}[Short-time scaling of the fixed-point sources]
\label{thm:e8v1-sources}
The differentiated channel sources
\(b_i^a=v_i(\mathcal T_a^{(\tau)}\rho)\) obey
\begin{align}
 \|b^a\|_{\mu,Z}
 &\leq\tau e^{\kappa\tau}\gamma_a,
 \qquad a\in\{s,q\},
 \label{eq:e8v1-first-source}\\
 \|b^{sq}\|_{\mu,Z}
 &\leq e^{\kappa\tau}\left[
 \tau\gamma_{sq}
 +{\tau^2\over2}(g_s\gamma_q+g_q\gamma_s)
 \right].
 \label{eq:e8v1-mixed-source}
\end{align}
\end{theorem}

\begin{proof}
Apply \eqref{eq:e8v1-first-Duhamel} to \(\rho\).  Invariance removes the
right semigroup and \eqref{eq:e8v1-propagation} bounds the remaining
propagation, proving \eqref{eq:e8v1-first-source}.  In the direct part of
\eqref{eq:e8v1-mixed-Duhamel}, the source is
\(\mathcal L_{sq}\rho\).  In the two ordered parts, it is first
\(\mathcal L_q\rho\) or \(\mathcal L_s\rho\), then propagated and acted on
by the other generator derivative.  The two simplex areas are each
\(\tau^2/2\).  Equations \eqref{eq:e8v1-propagation} and
\eqref{eq:e8v1-generator-kernel} give
\eqref{eq:e8v1-mixed-source}.
\end{proof}

\subsection{Approximate recovery also carries the clock}
\label{sec:e8v1-recovery}

For a target density which is not exactly invariant, put
\begin{equation}
 r:=\mathcal L\rho,
 \qquad
 e^{(\tau)}:=\mathcal T^{(\tau)}\rho-\rho
 =\int_0^\tau\mathcal S(u)r\,du.
 \label{eq:e8v1-recovery}
\end{equation}
Let \(\gamma_0,\gamma_s^r,\gamma_q^r,\gamma_{sq}^r\) bound in the same
weighted norm the seminorm vectors of
\(r,r_s,r_q,r_{sq}\), respectively.

\begin{proposition}[Recovery-defect derivative scaling]
\label{prop:e8v1-recovery}
Under the preceding propagation and generator-kernel hypotheses,
\begin{align}
 \|e^{(\tau)}\|_{\mu,Z}
 &\leq e^{\kappa\tau}\tau\gamma_0,
 \label{eq:e8v1-e0}\\
 \|e_a^{(\tau)}\|_{\mu,Z}
 &\leq e^{\kappa\tau}\left(
 \tau\gamma_a^r+{\tau^2\over2}g_a\gamma_0\right),
 \qquad a\in\{s,q\},
 \label{eq:e8v1-ea}\\
 \|e_{sq}^{(\tau)}\|_{\mu,Z}
 &\leq e^{\kappa\tau}\left[
 \tau\gamma_{sq}^r
 +{\tau^2\over2}
 (g_s\gamma_q^r+g_q\gamma_s^r+g_{sq}\gamma_0)
 +{\tau^3\over3}g_sg_q\gamma_0
 \right].
 \label{eq:e8v1-esq}
\end{align}
\end{proposition}

\begin{proof}
Equation \eqref{eq:e8v1-recovery} and
\eqref{eq:e8v1-propagation} prove \eqref{eq:e8v1-e0}.  Differentiate the
integral.  The first derivative of \(\mathcal S(u)\) has one insertion and
norm at most \(u e^{\kappa u}g_a\); integration of \(u\) from zero to
\(\tau\) gives \(\tau^2/2\), proving \eqref{eq:e8v1-ea}.

For the mixed derivative, the terms are
\(\mathcal S(u)r_{sq}\), the first semigroup derivative applied to
\(r_s\) or \(r_q\), the direct \(\mathcal L_{sq}\) insertion applied to
\(r\), and the two ordered \(\mathcal L_s,\mathcal L_q\) insertions applied
to \(r\).  Their bounds before the outer integration are, respectively,
\[
 e^{\kappa u}\gamma_{sq}^r,
 \quad u e^{\kappa u}(g_s\gamma_q^r+g_q\gamma_s^r),
 \quad u e^{\kappa u}g_{sq}\gamma_0,
 \quad u^2e^{\kappa u}g_sg_q\gamma_0.
\]
Use \(e^{\kappa u}\leq e^{\kappa\tau}\) and integrate
\(1,u,u^2\) to obtain \eqref{eq:e8v1-esq}.
\end{proof}

\begin{corollary}[Generator packet supplies the V100 source scaling]
\label{cor:e8v1-V100}
Fix \(0<\tau\leq\tau_*\).  Suppose the untagged metric quantities
\(g_q,\gamma_q,\gamma_q^r\) and all base quantities are uniformly bounded,
while every refinement-tagged quantity
\begin{equation}
 g_s,\ g_{sq},\ \gamma_s,\ \gamma_{sq},\
 \gamma_s^r,\ \gamma_{sq}^r
 \label{eq:e8v1-tagged}
\end{equation}
is bounded by \(\varepsilon_n\) times a scale-independent constant.
Then the first and mixed channel kernels, fixed-point sources, and recovery
defect derivatives satisfy the V100 pattern
\begin{equation}
 a_n^s,b_n^s,b_n^{sq}=O(\tau_n\varepsilon_n),
 \qquad
 a_n^q,b_n^q=O(\tau_n),
 \label{eq:e8v1-V100-pattern}
\end{equation}
with constants uniform for \(\tau_n\leq\tau_*\).  The terms containing
\(\tau_n^2\) or \(\tau_n^3\) are absorbed using
\(\tau_n\leq\tau_*\).
\end{corollary}

\begin{proof}
Apply \Cref{cor:e8v1-channel-kernel,thm:e8v1-sources,
prop:e8v1-recovery}.  Every mixed product contains either an explicit
refinement derivative or a refinement derivative of the generator
residual, hence one factor \(\varepsilon_n\).  Factor out one
\(\tau_n\); all remaining powers of \(\tau_n\) are uniformly bounded by
powers of \(\tau_*\).
\end{proof}

\begin{assessment}[What the new V1 accomplishes]
\label{ass:e8v1-status}
The short-time origin of the clock factors required at V100 is now proved.
The channel derivatives, mixed source, and differentiated recovery defect
all carry at least one power of \(\tau_n\).  If refinement derivatives of
the generator carry \(\varepsilon_n\), that amplitude propagates through
every mixed term.  The remaining half of the packet is a physical generator
estimate which yields the critical contraction row
\(1-m\tau_n+O(\tau_n^2)\) in a seminorm that still controls local
observables.
\end{assessment}

\begin{programme}[Eighth-series V2: generator dissipativity]
\label{prog:e8v1-V2}
The next version will formulate a weighted logarithmic-norm condition on the
traceless generator and prove that it yields the V100 critical contraction
row for \(e^{\tau\mathcal L}\).  It will then test that condition against
the Ward zero: the one-site absolute norm must fail in the gauge sector, so
the viable targets are a fixed-physical-duration block or a cancellation
seminorm with a proved dual-observable estimate.  The first failure will be
used to move upstream rather than reimpose a forbidden uniform margin.
\end{programme}

\begin{firewall}[Claim boundary after eighth-series V1]
\label{fire:e8v1-boundary}
This V1 records the archive boundary and major result map; derives the exact
first and mixed Duhamel formulas for a short-time channel; proves their
diamond and weighted local bounds; proves short-time scaling of fixed-point
and differentiated recovery sources; and shows how a refinement amplitude
loads the V100 source pattern.

It does not construct the physical Standard Model carrier generator, prove
the critical contraction row, prove summability of
\(\varepsilon_n\), establish the counterterm or regulator bridge rows,
remove the cutoff, or construct the Standard Model--Einstein continuum.
\end{firewall}

\begin{theorem}[Eighth-series V1 result]
\label{thm:e8v1-result}
For a \(C^2\) local generator family, every first derivative of the
short-time channel is \(O(\tau)\), and its mixed derivative is the sum of
an \(O(\tau)\) mixed-generator insertion and \(O(\tau^2)\) ordered first
insertions.  The same clock factor holds for invariant-state sources and
for an approximate recovery defect together with its first and mixed
parameter derivatives.  Refinement-tagged generator bounds of size
\(\varepsilon_n\) therefore produce the clock-scaled mixed-source packet
required by the V100 critical resolvent theorem.
\end{theorem}

\begin{proof}
Combine \Cref{thm:e8v1-Duhamel,cor:e8v1-channel-kernel,
thm:e8v1-sources,prop:e8v1-recovery,cor:e8v1-V100}.
\end{proof}

\paragraph{Version record.}
This is the first active note of the eighth series.  Its immutable parent is
\texttt{Renewal\_SM\_Einstein\_Continuum\_Research\_Notes\_V100\_f7.tex}
with the digest recorded in \Cref{sec:e8v1-context}.  The next compulsory
companion audit is V5.

% =============================================================================
% BEGIN APPENDED EIGHTH-SERIES V2 MATERIAL
% =============================================================================

\section{Generator logarithmic norms and the Ward quotient}
\label{sec:v8v2}

V1 proved that a short-time channel supplies the clock factor in every
metric, refinement, mixed, and recovery source.  We now prove a sufficient
generator condition for the complementary contraction row.  The condition
is formulated directly in the local seminorm packet, so its failure on a
Ward mode is visible rather than hidden by a global spectral gap.

\subsection{Metzler comparison in a weighted local norm}
\label{sec:e8v2-Metzler}

Let \(\mathcal X\) be the real space of traceless Hermitian carrier
operators and let \(v_i:\mathcal X\to[0,\infty)\) be seminorms indexed by
cells.  For a nonempty set \(Z\) define
\begin{equation}
 \|z\|_{\mu,Z}:=\sup_i e^{\mu d(i,Z)}|z_i|.
 \label{eq:e8v2-vector-norm}
\end{equation}
Let \(Q=(q_{ij})\) be a Metzler matrix: \(q_{ij}\geq0\) for \(i\ne j\).
Set
\begin{equation}
 \ell_\mu(Q):=
 \sup_i\left(q_{ii}+\sum_{j\ne i}e^{\mu d(i,j)}q_{ij}\right).
 \label{eq:e8v2-lognorm}
\end{equation}
All sums are finite at finite regulator; the estimates below are useful
cofinally only when their constants are volume independent.

\begin{lemma}[Weighted Metzler semigroup bound]
\label{lem:e8v2-Metzler}
If \(z(t)\geq0\) obeys the componentwise upper Dini inequality
\begin{equation}
 D^+z(t)\leq Qz(t),
 \label{eq:e8v2-Dini}
\end{equation}
then
\begin{equation}
 \|z(t)\|_{\mu,Z}
 \leq e^{\ell_\mu(Q)t}\|z(0)\|_{\mu,Z}.
 \label{eq:e8v2-Metzler-bound}
\end{equation}
In particular, \(\ell_\mu(Q)\leq-m<0\) gives decay
\(e^{-mt}\).
\end{lemma}

\begin{proof}
Put \(w_i=e^{\mu d(i,Z)}\) and \(y_i=w_iz_i\).  Since
\(d(i,Z)-d(j,Z)\leq d(i,j)\), the conjugated off-diagonal coefficients
satisfy
\[
 {w_i\over w_j}q_{ij}\leq e^{\mu d(i,j)}q_{ij}.
\]
At a time where the maximum of the finitely many \(y_i\) is attained at
\(i_0\), the upper Dini derivative of
\(M(t)=\max_i y_i(t)\) obeys
\begin{align*}
 D^+M(t)
 &\leq q_{i_0i_0}M(t)
 +\sum_{j\ne i_0}{w_{i_0}\over w_j}q_{i_0j}y_j(t)\\
 &\leq\ell_\mu(Q)M(t).
\end{align*}
The Dini version of Gronwall's inequality proves
\eqref{eq:e8v2-Metzler-bound}.
\end{proof}

Let \(\mathcal S(t)=e^{t\mathcal L}\) be the CPTP semigroup of V1.  Assume
that for every \(X\in\mathcal X\),
\begin{equation}
 D^+v_i(\mathcal S(t)X)
 \leq\sum_jq_{ij}v_j(\mathcal S(t)X).
 \label{eq:e8v2-generator-comparison}
\end{equation}

\begin{theorem}[Generator dissipativity gives the critical channel row]
\label{thm:e8v2-critical-row}
If \eqref{eq:e8v2-generator-comparison} holds with
\(\ell_\mu(Q)\leq-m<0\), then the one-step channel
\(\mathcal T^{(\tau)}=e^{\tau\mathcal L}\) satisfies
\begin{equation}
 \|v(\mathcal T^{(\tau)}X)\|_{\mu,Z}
 \leq e^{-m\tau}\|v(X)\|_{\mu,Z}.
 \label{eq:e8v2-exact-row}
\end{equation}
For \(0\leq\tau\leq\tau_*\), this implies the V100 form
\begin{equation}
 c_\mu(\tau)
 \leq1-m\tau+K_*\tau^2,
 \qquad
 K_*:={m^2\over2}e^{m\tau_*}.
 \label{eq:e8v2-V100-row}
\end{equation}
\end{theorem}

\begin{proof}
Apply \Cref{lem:e8v2-Metzler} to
\(z_i(t)=v_i(\mathcal S(t)X)\).  This gives
\eqref{eq:e8v2-exact-row}.  Taylor's formula with integral remainder gives
\[
 e^{-m\tau}
 =1-m\tau+m^2\int_0^\tau(\tau-u)e^{-mu}\,du
 \leq1-m\tau+{m^2\over2}\tau^2.
\]
The larger stated constant \(K_*\) also works uniformly and remains valid
if the generator comparison itself has a bounded second-order channel
implementation error.
\end{proof}

\begin{remark}[Sharp constant for the exact semigroup]
\label{rem:e8v2-sharp}
For the exact exponential in the theorem, \(K_*=m^2/2\) already suffices.
The displayed larger value leaves a declared slot for replacing the exact
exponential by a second-order consistent local gate.  Any such replacement
must prove its own error coefficient.
\end{remark}

\subsection{A stationary mode forbids strict dissipativity}
\label{sec:e8v2-obstruction}

Define the combined seminorm
\begin{equation}
 \|X\|_{v,\mu,Z}:=
 \|\bigl(v_i(X)\bigr)_i\|_{\mu,Z}.
 \label{eq:e8v2-combined}
\end{equation}

\begin{theorem}[Stationary-mode obstruction]
\label{thm:e8v2-stationary}
Suppose \(0\ne W\in\mathcal X\) satisfies
\begin{equation}
 \mathcal LW=0,
 \qquad
 \|W\|_{v,\mu,Z}>0.
 \label{eq:e8v2-Ward-mode}
\end{equation}
Then no estimate of the form
\begin{equation}
 \|e^{t\mathcal L}X\|_{v,\mu,Z}
 \leq e^{-mt}\|X\|_{v,\mu,Z},
 \qquad m>0,
 \label{eq:e8v2-impossible}
\end{equation}
can hold for all \(X\in\mathcal X\).  In particular, no comparison matrix
for this seminorm can satisfy both
\eqref{eq:e8v2-generator-comparison} and
\(\ell_\mu(Q)<0\).
\end{theorem}

\begin{proof}
Since \(\mathcal LW=0\), \(e^{t\mathcal L}W=W\) for all \(t\geq0\).
Substituting \(W\) into \eqref{eq:e8v2-impossible} and dividing by its
positive seminorm gives \(1\leq e^{-mt}\), a contradiction for \(t>0\).
The last statement follows from \Cref{thm:e8v2-critical-row}.
\end{proof}

\begin{corollary}[Compatibility with the audited Ward obstruction]
\label{cor:e8v2-Ward}
If the cofinal gauge generator has a Ward stationary direction detected by
the absolute one-site influence seminorm, the hypothesis
\(\ell_\mu(Q)<0\) is impossible in that seminorm.  Thus
\Cref{thm:e8v2-critical-row} cannot be used to restore the uniform one-site
margin excluded at V100.
\end{corollary}

\begin{proof}
Apply \Cref{thm:e8v2-stationary}.  The V100 Yang--Mills audit identifies
the Ward zero as precisely such a noncontracting direction for the
absolute one-site response.
\end{proof}

\subsection{Quotienting a declared slow sector}
\label{sec:e8v2-quotient}

Let \(\mathcal W\subset\mathcal X\) be a linear subspace invariant under
\(\mathcal L\).  Define quotient cell seminorms
\begin{equation}
 \bar v_i([X]):=\inf_{W\in\mathcal W}v_i(X+W),
 \qquad [X]\in\mathcal X/\mathcal W.
 \label{eq:e8v2-quotient-seminorm}
\end{equation}
Their use for observables requires a separate duality statement.

\begin{definition}[Ward-compatible observable core]
\label{def:e8v2-Ward-core}
A local observable \(A_B\) is compatible with \(\mathcal W\) if
\begin{equation}
 \Tr(WA_B)=0\quad\text{for every }W\in\mathcal W
 \label{eq:e8v2-annihilator}
\end{equation}
and the quotient seminorms obey
\begin{equation}
 |\Tr(XA_B)|
 \leq\|A_B\|\sum_{i\in B}\bar v_i([X]).
 \label{eq:e8v2-quotient-dual}
\end{equation}
The first condition is necessary for the left side to depend only on
\([X]\); the second is an analytic hypothesis which must be proved for the
chosen quotient.
\end{definition}

\begin{theorem}[Quotient channel compiler]
\label{thm:e8v2-quotient}
Assume \(\mathcal W\) is invariant under
\(\mathcal L,\mathcal L_s,\mathcal L_q,\mathcal L_{sq}\), so all these
maps descend to \(\mathcal X/\mathcal W\).  Suppose the quotient seminorms
satisfy:
\begin{enumerate}
\item the generator comparison \eqref{eq:e8v2-generator-comparison} with
\(v_i\) replaced by \(\bar v_i\) and
\(\ell_\mu(Q)\leq-m<0\);
\item the V1 weighted propagation and derivative-kernel estimates; and
\item the dual estimate \eqref{eq:e8v2-quotient-dual} on a protected local
observable core.
\end{enumerate}
Then the V1 source estimates and V100 critical-resolvent cancellation hold
on the quotient for that observable core.  If the refinement-tagged
generator rows are \(O(\varepsilon_n)\) and
\(\sum_n\varepsilon_n<\infty\), the fixed protected expectation and stress
rows are summable.
\end{theorem}

\begin{proof}
Invariance makes every formula in V1 well defined on the quotient.  Item 1
and \Cref{thm:e8v2-critical-row} supply the V100 contraction row.  Item 2
supplies the clock-scaled first, mixed, and recovery sources by the V1
Duhamel theorems.  The V100 cancellation removes the inverse clock factor.
Finally item 3 converts the quotient response vector into protected local
expectations, and comparison with \(\sum_n\varepsilon_n\) gives the stated
rows.
\end{proof}

\begin{firewall}[Quotienting is not automatic gauge fixing]
\label{fire:e8v2-quotient}
A formal quotient by the Ward mode is insufficient.  One must prove that
the channel and all parameter derivatives preserve the slow subspace and
that the desired gauge-invariant local observables satisfy
\eqref{eq:e8v2-quotient-dual}.  If a physical observable couples to a
removed mode, the quotient estimate says nothing about its response.
\end{firewall}

\subsection{Closed abstract clock packet}
\label{sec:e8v2-packet}

\begin{definition}[Generator critical-source packet, GCSP]
\label{def:e8v2-GCSP}
At scale \(n\), GCSP consists of:
\begin{enumerate}
\item a CPTP semigroup \(e^{t\mathcal L_n}\) and either a direct seminorm or
an invariant Ward quotient with the local observable dual property;
\item a volume-independent Metzler comparison whose weighted logarithmic
norm is at most \(-m<0\);
\item the V1 weighted generator-derivative and invariant/recovery-source
bounds, with every refinement-tagged row bounded by
\(C\varepsilon_n\);
\item a clock \(0<\tau_n\leq\tau_*\), uniformly bounded contact terms, and
\(\sum_n\varepsilon_n<\infty\); and
\item coherent adjacent embeddings and the V98 uniform dual-norm completion
row.
\end{enumerate}
\end{definition}

\begin{corollary}[GCSP loads the carrier part of the cofinal compiler]
\label{cor:e8v2-GCSP}
GCSP implies the V98 base and stress rows for every protected observable in
its declared core.  It therefore loads the carrier term in the V99 weak
residual ledger.
\end{corollary}

\begin{proof}
Rows 1 and 2 give \eqref{eq:e8v2-V100-row}.  Row 3 and the V1 result give
the V100 source scaling.  Row 4 and the V100 cancellation give summable
adjacent rows.  Row 5 permits the V98 landing and its transfer to the V99
carrier term.
\end{proof}

\begin{assessment}[Progress and remaining bottleneck]
\label{ass:e8v2-status}
The two abstract halves of the physical-clock carrier argument are now
joined: generator logarithmic dissipativity supplies the closing
contraction row, and short-time Duhamel calculus supplies the matching
source factor.  The Ward obstruction proves that the absolute one-site
seminorm cannot furnish this packet.  The next task is therefore to build a
specific Ward quotient or a fixed-physical-duration block with a valid
dual-observable inequality.
\end{assessment}

\begin{programme}[Eighth-series V3: conserved-mode projection and block]
\label{prog:e8v2-V3}
V3 will study a finite-dimensional conserved-mode projection
\(P_{\mathcal W}\) and the complementary semigroup.  It will prove a block
contraction from a complementary generator gap, quantify errors when the
projection only approximately intertwines the channel, and state the exact
observable-annihilator test.  If the gap closes with physical volume, the
programme will expose that infrared debit rather than report a local block
as a continuum mass estimate.
\end{programme}

\begin{firewall}[Claim boundary after eighth-series V2]
\label{fire:e8v2-boundary}
V2 proves the weighted Metzler semigroup bound, derives the V100 critical
row from a negative generator logarithmic norm, proves the stationary-mode
obstruction, and gives a conditional Ward-quotient compiler with an
explicit dual-observable requirement.

V2 does not construct the physical quotient or block, prove a
volume-independent complementary gap, verify the observable annihilator,
load the Standard Model generator, remove the cutoff, or construct the
Standard Model--Einstein continuum.
\end{firewall}

\begin{theorem}[Eighth-series V2 result]
\label{thm:e8v2-result}
A volume-independent negative weighted logarithmic norm of the traceless
generator yields the physical-clock contraction row needed by V100.  Any
stationary direction detected by the seminorm forbids such negativity.
After an invariant slow sector is quotiented, the argument remains valid
only for observables with a proved quotient dual estimate; under that
condition GCSP supplies the carrier base and stress rows of the cofinal
compiler.
\end{theorem}

\begin{proof}
Combine \Cref{lem:e8v2-Metzler,thm:e8v2-critical-row,
thm:e8v2-stationary,thm:e8v2-quotient,cor:e8v2-GCSP}.
\end{proof}

\paragraph{Parent-version record.}
V2 was formed from
\texttt{Renewal\_SM\_Einstein\_Continuum\_Research\_Notes\_V1.tex}.
The V1 parent had \(507\) logical source lines, \(21\,010\) bytes, and
SHA--256
\[
 \texttt{F0E8C934DDB8E4A23BD38E07A7FBE8045D17F7C44008363072EEF11C884A7E8F}.
\]
The next compulsory companion audit is V5.

% =============================================================================
% END APPENDED EIGHTH-SERIES V2 MATERIAL
% =============================================================================
% =============================================================================
% BEGIN APPENDED EIGHTH-SERIES V3 MATERIAL
% =============================================================================

\section{Slow-mode projections, block contraction, and infrared leakage}
\label{sec:v8v3}

V2 shows that a Ward stationary direction must either be quotiented with a
valid observable duality theorem or retained in a block analysis.  This
appendix makes both statements quantitative.  It also records the infrared
failure mode: removing only the exact zero mode does not give a uniform
block gap when nonzero modes approach zero with physical volume.

\subsection{The exact observable-annihilator test}
\label{sec:e8v3-annihilator}

Let \(\mathcal X\) be a finite-dimensional normed real space of traceless
Hermitian operators, let \(\mathcal W\subset\mathcal X\) be a subspace, and
write \(\pi:\mathcal X\to\mathcal X/\mathcal W\).  Equip the quotient with
\begin{equation}
 \|[X]\|_{\rm quot}:=\inf_{W\in\mathcal W}\|X+W\|_{\mathcal X}.
 \label{eq:e8v3-quotient-norm}
\end{equation}
For an observable \(A\), set \(f_A(X)=\Tr(XA)\).

\begin{theorem}[Dual of the slow-mode quotient]
\label{thm:e8v3-dual}
The functional \(f_A\) descends to a well-defined bounded functional
\(\bar f_A\) on \(\mathcal X/\mathcal W\) if and only if
\begin{equation}
 \Tr(WA)=0\qquad(W\in\mathcal W).
 \label{eq:e8v3-annihilator}
\end{equation}
When this holds,
\begin{equation}
 \|\bar f_A\|_{(\mathcal X/\mathcal W)^*}
 =\|f_A\|_{\mathcal X^*}.
 \label{eq:e8v3-isometry}
\end{equation}
Thus \((\mathcal X/\mathcal W)^*\) is isometrically the annihilator
\(\mathcal W^\perp\subset\mathcal X^*\).
\end{theorem}

\begin{proof}
If \(f_A\) descends, then \(0=[W]\) implies
\(f_A(W)=\bar f_A([W])=0\), proving necessity.  Conversely,
\eqref{eq:e8v3-annihilator} makes \(f_A(X)\) independent of the
representative of \([X]\).  Moreover, for every \(W\in\mathcal W\),
\[
 |f_A(X)|=|f_A(X+W)|
 \leq\|f_A\|_{\mathcal X^*}\|X+W\|_{\mathcal X}.
\]
Taking the infimum proves
\(\|\bar f_A\|\leq\|f_A\|\).  The reverse inequality follows from
\(|f_A(X)|=|\bar f_A([X])|\leq
\|\bar f_A\|\|[X]\|_{\rm quot}\leq
\|\bar f_A\|\|X\|_{\mathcal X}\), and then taking the supremum over the
unit ball of \(\mathcal X\).
\end{proof}

\begin{corollary}[Local duality remains a separate estimate]
\label{cor:e8v3-local-dual}
Condition \eqref{eq:e8v3-annihilator} is exactly the algebraic test for a
quotient observable.  It does not imply the cell-local inequality
\begin{equation}
 |\Tr(XA_B)|\leq\|A_B\|
 \sum_{i\in B}\bar v_i([X]).
 \label{eq:e8v3-local-dual}
\end{equation}
That inequality requires the quotient norm generated by the local
seminorms to dominate the quotient functional norm on the declared local
support.
\end{corollary}

\begin{proof}
The theorem identifies boundedness in the global quotient norm
\eqref{eq:e8v3-quotient-norm}.  Equation \eqref{eq:e8v3-local-dual}
contains additional support and normalization information not present in
that norm.  Rescaling all cell seminorms by a small positive constant
preserves their kernels and the annihilator but can violate the displayed
inequality, proving that the algebraic condition alone is insufficient.
\end{proof}

\subsection{Exact reducing projection and a physical-duration block}
\label{sec:e8v3-block}

Let \(P:\mathcal X\to\mathcal X\) be a bounded projection,
\(Q=I-P\), and \(\mathcal W=\operatorname{Ran}P\).  Suppose
\begin{equation}
 P\mathcal L=\mathcal LP.
 \label{eq:e8v3-reducing}
\end{equation}
Then \(P,Q\) commute with \(\mathcal S(t)=e^{t\mathcal L}\), and the
quotient is represented by \(Q\mathcal X\).  Equip it with a norm
\(\|[X]\|_Q:=\|QX\|_Q\).

\begin{theorem}[Complementary gap gives block contraction]
\label{thm:e8v3-block}
Assume
\begin{equation}
 \|Qe^{t\mathcal L}Q\|_{Q\to Q}
 \leq M e^{-\gamma t},
 \qquad t\geq0,qquad M\geq1,quad\gamma>0.
 \label{eq:e8v3-gap}
\end{equation}
Then the physical-duration block
\(\mathcal B_T=e^{T\mathcal L}\) induces a quotient map with coefficient
\begin{equation}
 c_T\leq M e^{-\gamma T}.
 \label{eq:e8v3-block-row}
\end{equation}
It is a strict contraction whenever
\begin{equation}
 T>{\log M\over\gamma}.
 \label{eq:e8v3-block-duration}
\end{equation}
If \(\mathcal L_s,\mathcal L_q,\mathcal L_{sq}\) also preserve
\(\mathcal W\), the V1 derivative formulas descend to the same quotient.
\end{theorem}

\begin{proof}
Commutation gives
\(Qe^{T\mathcal L}X=Qe^{T\mathcal L}QX\), so
\[
 \|[\mathcal B_TX]\|_Q
 \leq M e^{-\gamma T}\|[X]\|_Q.
\]
This proves \eqref{eq:e8v3-block-row};
\eqref{eq:e8v3-block-duration} makes its right side smaller than one.
Preservation by the parameter derivatives makes every insertion in the V1
Duhamel formulas independent of the representative, so those formulas
descend.
\end{proof}

\begin{corollary}[Fixed block source scaling]
\label{cor:e8v3-block-source}
Fix \(T>0\) satisfying \eqref{eq:e8v3-block-duration}.  If the generator
metric derivatives are uniformly bounded and every refinement-tagged
generator and recovery derivative is \(O(\varepsilon_n)\), then the block
metric sources are \(O(1)\) and its refinement and mixed sources are
\(O(\varepsilon_n)\).  Hence \(\sum_n\varepsilon_n<\infty\), together
with the quotient local duality and contact rows, loads the V98 carrier
ledger without a microscopic \(\tau_n^{-1}\) resolvent.
\end{corollary}

\begin{proof}
Use the V1 Duhamel bounds with \(\tau=T\).  Since \(T\), the generator
bounds, and the strict block resolvent \((1-c_T)^{-1}\) are fixed, all
untagged factors are absorbed into a constant.  Every refinement or mixed
term contains one refinement-tagged insertion, hence one
\(\varepsilon_n\).  The V98 comparison then gives the conclusion.
\end{proof}

\subsection{Approximate block diagonalization}
\label{sec:e8v3-approximate}

Exact reduction may fail at finite cutoff.  Put
\begin{equation}
 \mathcal L_0:=P\mathcal LP+Q\mathcal LQ,
 \qquad
 \mathcal E:=\mathcal L-\mathcal L_0
 =P\mathcal LQ+Q\mathcal LP.
 \label{eq:e8v3-split}
\end{equation}
Let \(\|\cdot\|\) be a submultiplicative operator norm on
\(\mathcal X\).

\begin{lemma}[Duhamel leakage bound]
\label{lem:e8v3-leakage}
Suppose, for \(0\leq t\leq T\),
\begin{equation}
 \|e^{t\mathcal L}\|\leq e^{\omega t},
 \qquad
 \|e^{t\mathcal L_0}\|\leq e^{\omega t},
 \qquad
 \|\mathcal E\|\leq\epsilon.
 \label{eq:e8v3-growth}
\end{equation}
Then
\begin{equation}
 \|e^{T\mathcal L}-e^{T\mathcal L_0}\|
 \leq T\epsilon e^{\omega T}.
 \label{eq:e8v3-Duhamel-leakage}
\end{equation}
Consequently,
\begin{align}
 \|Qe^{T\mathcal L}Q\|
 &\leq\|Qe^{T\mathcal L_0}Q\|
 +\|Q\|^2T\epsilon e^{\omega T},
 \label{eq:e8v3-Q-block}\\
 \|Pe^{T\mathcal L}Q\|
 &\leq\|P\|\|Q\|T\epsilon e^{\omega T}.
 \label{eq:e8v3-cross-leakage}
\end{align}
\end{lemma}

\begin{proof}
Duhamel's identity is
\[
 e^{T\mathcal L}-e^{T\mathcal L_0}
 =\int_0^T e^{(T-u)\mathcal L}\mathcal E
              e^{u\mathcal L_0}\,du.
\]
The integrand norm is at most
\(\epsilon e^{\omega(T-u)}e^{\omega u}\), proving
\eqref{eq:e8v3-Duhamel-leakage}.  Multiply by \(Q\) on both sides for
\eqref{eq:e8v3-Q-block}.  Since \(\mathcal L_0\) is block diagonal,
\(Pe^{T\mathcal L_0}Q=0\); multiplying the difference by \(P,Q\) gives
\eqref{eq:e8v3-cross-leakage}.
\end{proof}

\begin{proposition}[Strict approximate block criterion]
\label{prop:e8v3-approximate-block}
If
\begin{equation}
 \|Qe^{T\mathcal L_0}Q\|\leq M e^{-\gamma T}
 \label{eq:e8v3-L0-gap}
\end{equation}
and
\begin{equation}
 M e^{-\gamma T}+\|Q\|^2T\epsilon e^{\omega T}<1,
 \label{eq:e8v3-strict-approximate}
\end{equation}
then the projected fast block \(Qe^{T\mathcal L}Q\) is strictly
contractive.  Its discrepancy from an exact quotient dynamics and its
slow-sector leakage are, however, separate errors of order
\(T\epsilon e^{\omega T}\).
\end{proposition}

\begin{proof}
The strict bound is \eqref{eq:e8v3-Q-block} with
\eqref{eq:e8v3-L0-gap}.  Equation
\eqref{eq:e8v3-cross-leakage} bounds leakage into the slow sector.  Because
\(\mathcal W\) is not invariant when \(\mathcal E\ne0\), the full map does
not descend to an exact quotient; the leakage therefore has to be retained
as an inhomogeneous source rather than hidden in the contraction
coefficient.
\end{proof}

\subsection{Infrared duration obstruction}
\label{sec:e8v3-infrared}

\begin{theorem}[Closing complementary gap forces growing block time]
\label{thm:e8v3-IR}
Let a cofinal family have complementary estimates
\begin{equation}
 \|Q_ne^{t\mathcal L_n}Q_n\|leq M_ne^{-\gamma_nt},
 \qquad M_n\geq1,quad\gamma_n>0.
 \label{eq:e8v3-cofinal-gap}
\end{equation}
To obtain a prescribed block coefficient \(c_*<1\) from this estimate, it
is necessary that the chosen duration obey
\begin{equation}
 T_n\geq{\log(M_n/c_*)\over\gamma_n}.
 \label{eq:e8v3-required-time}
\end{equation}
If \(M_n\geq1\) and \(\gamma_n\to0\), no uniformly bounded duration can
certify \(c_{T_n}\leq c_*\) through
\eqref{eq:e8v3-cofinal-gap}.
\end{theorem}

\begin{proof}
The sufficient upper bound in \eqref{eq:e8v3-cofinal-gap} is at most
\(c_*\) only when
\(M_ne^{-\gamma_nT_n}\leq c_*\).  Taking logarithms gives
\eqref{eq:e8v3-required-time}.  Since
\(\log(M_n/c_*)\geq\log(1/c_*)>0\), the right side diverges when
\(\gamma_n\to0\).
\end{proof}

\begin{firewall}[A bound cannot prove its own necessity]
\label{fire:e8v3-IR}
Theorem \ref{thm:e8v3-IR} says that the particular exponential-gap estimate
\eqref{eq:e8v3-cofinal-gap} cannot certify a fixed contraction with bounded
block time when \(\gamma_n\to0\).  It does not rule out a sharper
non-spectral cancellation estimate.  Such an estimate must be proved in a
norm with the required local-observable duality.
\end{firewall}

\begin{assessment}[Decision after the block test]
\label{ass:e8v3-status}
An exact slow-mode quotient is mathematically clean and a fixed-duration
block works when the complementary gap remains positive.  Approximate
invariance creates a quantified leakage source.  If the gauge carrier has
a massless tower of nonzero momenta with \(\gamma_n\to0\), quotienting only
the exact Ward mode leaves a growing-time obstruction.  V4 must test this
on the discrete torus before the programme relies on a fixed block.
\end{assessment}

\begin{programme}[Eighth-series V4: discrete massless infrared test]
\label{prog:e8v3-V4}
V4 will compute the smallest nonzero eigenvalue and heat-block coefficient
of the periodic lattice Laplacian after the constant mode is removed.  It
will determine the required block duration in lattice and physical units,
then redirect the carrier argument to a local heat-kernel or finite-distance
mixing estimate if the global quotient gap closes.
\end{programme}

\begin{firewall}[Claim boundary after eighth-series V3]
\label{fire:e8v3-boundary}
V3 proves the exact quotient-observable annihilator criterion, the
complementary-gap block theorem, fixed-block source scaling, Duhamel bounds
for approximate slow/fast leakage, and the growing-duration obstruction for
a closing complementary gap.

V3 does not prove that the physical Ward quotient has the local dual
property, construct a volume-uniform complementary gap, control a massless
infrared tower, populate the SM generator, or construct the Standard
Model--Einstein continuum.
\end{firewall}

\begin{theorem}[Eighth-series V3 result]
\label{thm:e8v3-result}
A slow-mode quotient controls exactly those observables in the annihilator
of the slow sector.  A reducing complementary gap produces strict
contraction after a sufficiently long physical block, while an off-diagonal
generator error produces an explicit leakage of order
\(T\epsilon e^{\omega T}\).  If the complementary gap closes cofinally,
the exponential-gap route requires a diverging block duration and cannot
by itself provide the fixed block assumed by the carrier compiler.
\end{theorem}

\begin{proof}
Combine \Cref{thm:e8v3-dual,thm:e8v3-block,
lem:e8v3-leakage,prop:e8v3-approximate-block,thm:e8v3-IR}.
\end{proof}

\paragraph{Parent-version record.}
V3 was formed from
\texttt{Renewal\_SM\_Einstein\_Continuum\_Research\_Notes\_V2.tex}.
The V2 parent had \(850\) logical source lines, \(33\,726\) bytes, and
SHA--256
\[
 \texttt{7CF466E660E064F3C95BDE3BCCE220F9B96A7DFA92CAD0AF4EA952376A42A89A}.
\]
The next compulsory companion audit is V5.

% =============================================================================
% END APPENDED EIGHTH-SERIES V3 MATERIAL
% =============================================================================

% =============================================================================
% BEGIN APPENDED EIGHTH-SERIES V4 MATERIAL
% =============================================================================

\section{The massless torus test and finite-horizon locality}
\label{sec:v8v4}

V3 leaves a concrete question: does removal of the exact Ward mode give a
volume-independent complementary gap?  The periodic lattice Laplacian is
the universal quadratic test for a massless gauge direction.  Its spectrum
shows that the answer depends on the order of limits.  At fixed physical
volume the continuum gap is positive; it closes like \(L^{-2}\) in the
infinite-volume limit.  A finite-time heat kernel nevertheless remains
local without any global gap.

\subsection{Exact periodic lattice spectrum}
\label{sec:e8v4-spectrum}

Let
\begin{equation}
 \Lambda_{a,N}:=(a\mathbb Z/Na\mathbb Z)^d,
 \qquad L:=Na,
 \label{eq:e8v4-torus}
\end{equation}
with periodic boundary conditions.  Define the positive lattice Laplacian
\begin{equation}
 (-\Delta_a f)(x):={1\over a^2}
 \sum_{r=1}^d\bigl(2f(x)-f(x+ae_r)-f(x-ae_r)\bigr).
 \label{eq:e8v4-Laplacian}
\end{equation}
Let \(P_0\) be the orthogonal projection onto constants and
\(Q_0=I-P_0\).

\begin{theorem}[Exact complementary gap]
\label{thm:e8v4-gap}
The Fourier modes indexed by
\(k=(k_1,\ldots,k_d)\in\{0,\ldots,N-1\}^d\) have eigenvalues
\begin{equation}
 \lambda_a(k)={4\over a^2}
 \sum_{r=1}^d\sin^2\left({\pi k_r\over N}\right).
 \label{eq:e8v4-spectrum}
\end{equation}
The constant mode is the unique zero mode, and for \(N\geq2\) the least
nonzero eigenvalue is
\begin{equation}
 \lambda_1(a,N)={4\over a^2}\sin^2\left({\pi\over N}\right).
 \label{eq:e8v4-lambda1}
\end{equation}
It satisfies
\begin{equation}
 {16\over L^2}\leq\lambda_1(a,N)
 \leq{4\pi^2\over L^2},
 \qquad
 \lambda_1(a,L/a)\longrightarrow{4\pi^2\over L^2}
 \quad(a\downarrow0).
 \label{eq:e8v4-gap-bounds}
\end{equation}
\end{theorem}

\begin{proof}
Substitution of
\(f_k(x)=\exp(2\pi i k\cdot x/L)\) into
\eqref{eq:e8v4-Laplacian} gives \eqref{eq:e8v4-spectrum}.  Every summand is
nonnegative and vanishes exactly when its coordinate is zero modulo
\(N\), so the zero mode is unique.  The least positive value has one
coordinate equal to \(1\) or \(N-1\) and all others zero, proving
\eqref{eq:e8v4-lambda1}.  For \(0\leq x\leq\pi/2\),
\(2x/\pi\leq\sin x\leq x\).  Apply this at \(x=\pi/N\) and use
\(L=Na\) to obtain the two bounds.  The limit follows from
\(\sin x/x\to1\).
\end{proof}

\begin{corollary}[Exact heat-block coefficient]
\label{cor:e8v4-heat-block}
On the mean-zero subspace of \(\ell^2(\Lambda_{a,N})\),
\begin{equation}
 \|Q_0e^{T\Delta_a}Q_0\|_{2\to2}
 =e^{-T\lambda_1(a,N)}.
 \label{eq:e8v4-heat-coefficient}
\end{equation}
To make this coefficient at most \(c_*\in(0,1)\), the exact minimum block
duration is
\begin{equation}
 T_*(a,N;c_*)={\log(1/c_*)\over\lambda_1(a,N)}.
 \label{eq:e8v4-Tstar}
\end{equation}
At fixed physical side \(L\),
\begin{equation}
 T_*(a,L/a;c_*)\longrightarrow
 {L^2\over4\pi^2}\log(1/c_*).
 \label{eq:e8v4-fixed-volume-time}
\end{equation}
If subsequently \(L\to\infty\), this duration diverges like \(L^2\).
\end{corollary}

\begin{proof}
The heat semigroup multiplies Fourier mode \(k\) by
\(e^{-T\lambda_a(k)}\).  On the orthogonal complement of the constant, the
largest multiplier is the one with eigenvalue \(\lambda_1\), proving
\eqref{eq:e8v4-heat-coefficient}.  Solving
\(e^{-T\lambda_1}\leq c_*\) gives \eqref{eq:e8v4-Tstar}; insert the limit
in \eqref{eq:e8v4-gap-bounds}.
\end{proof}

\begin{assessment}[Order of limits]
\label{ass:e8v4-order}
Removing the exact constant mode is sufficient for a fixed-volume
continuum heat-block gap.  It is insufficient for a uniform
infinite-volume gap.  A proof which takes \(a\downarrow0\) at fixed \(L\)
may use \eqref{eq:e8v4-fixed-volume-time}, but it must pay the
\(L^2\) duration before sending \(L\to\infty\).  Treating those two limits
as one uniform contraction would be false.
\end{assessment}

\subsection{Gap-free finite-time heat-kernel locality}
\label{sec:e8v4-heat-locality}

The global gap is not needed to control propagation over a finite time.
Let \(X_t\) be the continuous-time nearest-neighbour random walk on
\(a\mathbb Z^d\) with generator \(\Delta_a\), started at zero, and let
\(p_T^a(x,y)\) be its heat kernel.

\begin{lemma}[Exact exponential moment]
\label{lem:e8v4-mgf}
For \(\theta\in\mathbb R^d\),
\begin{equation}
 \mathbb E\,e^{\theta\cdot X_T}
 =\exp\left\{{2T\over a^2}
 \sum_{r=1}^d\bigl(\cosh(a\theta_r)-1\bigr)\right\}.
 \label{eq:e8v4-mgf}
\end{equation}
Consequently, for every \(\mu>0\) and \(R>0\),
\begin{equation}
 \mathbb P\{|X_T|_1\geq R\}
 \leq2^d\exp\left\{-\mu R
 +{2dT\over a^2}\bigl(\cosh(a\mu)-1\bigr)\right\}.
 \label{eq:e8v4-Chernoff}
\end{equation}
\end{lemma}

\begin{proof}
In each coordinate, positive and negative jumps are independent Poisson
processes of rate \(a^{-2}\).  Their moment generating functions multiply
to
\(\exp\{2Ta^{-2}(\cosh(a\theta_r)-1)\}\); multiplying the coordinates
proves \eqref{eq:e8v4-mgf}.  Also
\[
 e^{\mu|x|_1}
 =\prod_{r=1}^d e^{\mu|x_r|}
 \leq\prod_{r=1}^d(e^{\mu x_r}+e^{-\mu x_r})
 =\sum_{\sigma\in\{-1,1\}^d}e^{\mu\sigma\cdot x}.
\]
Markov's inequality and \eqref{eq:e8v4-mgf}, with all sign choices giving
the same value, prove \eqref{eq:e8v4-Chernoff}.
\end{proof}

\begin{corollary}[Uniform Gaussian regime]
\label{cor:e8v4-Gaussian}
If
\begin{equation}
 \mu_*:={R\over2d(\cosh1)T}leq {1\over a},
 \label{eq:e8v4-moderate}
\end{equation}
then
\begin{equation}
 \mathbb P\{|X_T|_1\geq R\}
 \leq2^d\exp\left\{-{R^2\over4d(\cosh1)T}\right\}.
 \label{eq:e8v4-Gaussian-tail}
\end{equation}
The bound is independent of the lattice spacing and total volume.
\end{corollary}

\begin{proof}
For \(|z|\leq1\), Taylor's theorem and
\(\sup_{|u|\leq1}\cosh u=\cosh1\) give
\(cosh z-1\leq(\cosh1)z^2/2\).  Insert this in
\eqref{eq:e8v4-Chernoff} and minimize
\(-\mu R+d(\cosh1)T\mu^2\) at \(\mu=\mu_*\).  Condition
\eqref{eq:e8v4-moderate} makes the quadratic bound applicable.
\end{proof}

\begin{theorem}[Finite-horizon separation bound]
\label{thm:e8v4-separation}
Let \(h\in\ell^1(a\mathbb Z^d)\) be supported in \(Z\), and let
\(B\) be a set with \(d_1(B,Z)\geq R\).  Then
\begin{equation}
 \sum_{x\in B}|(e^{T\Delta_a}h)(x)|
 \leq\|h\|_1\,
 \mathbb P\{|X_T|_1\geq R\}.
 \label{eq:e8v4-separation}
\end{equation}
Under \eqref{eq:e8v4-moderate}, the probability is bounded by
\eqref{eq:e8v4-Gaussian-tail}.  The same statement holds on a periodic
torus when \(R\) is its periodic \(\ell^1\) separation.
\end{theorem}

\begin{proof}
Kernel positivity and translation invariance give
\begin{align*}
 \sum_{x\in B}|(e^{T\Delta_a}h)(x)|
 &\leq\sum_{z\in Z}|h(z)|\sum_{x\in B}p_T^a(z,x).
\end{align*}
For each \(z\in Z\), every \(x\in B\) is at distance at least \(R\), so
the inner sum is at most the tail probability.  Sum over \(z\).  On the
torus, lift the walk to \(a\mathbb Z^d\).  The union of all lifts of
\(B\) is contained in the complement of the radius-\(R\) ball about every
lift of \(z\), so the same probability bound applies.
\end{proof}

\subsection{A finite-horizon cofinal alternative}
\label{sec:e8v4-cofinal}

Let a scale-\(n\) refinement source have \(\ell^1\)-size
\(\varepsilon_n\) and be separated from a protected observable by physical
distance \(R_n\).  Let it propagate for physical time \(T_n\).

\begin{corollary}[Summable separated heat sources]
\label{cor:e8v4-summable}
In the Gaussian regime, the induced local expectation defects are
summable whenever
\begin{equation}
 \sum_n\varepsilon_n
 \exp\left\{-{R_n^2\over4d(\cosh1)T_n}\right\}<\infty.
 \label{eq:e8v4-summability}
\end{equation}
This criterion requires no global mean-zero spectral gap.
\end{corollary}

\begin{proof}
Apply \Cref{thm:e8v4-separation} at each scale and multiply by the norm of
the protected local observable.  The Weierstrass test gives summability.
\end{proof}

\begin{firewall}[Finite horizon is not equilibrium]
\label{fire:e8v4-equilibrium}
The estimate \eqref{eq:e8v4-summability} controls propagation for the
specified finite times \(T_n\).  It does not construct an invariant Gibbs
carrier by taking \(T_n\to\infty\).  In the massless infinite-volume
sector the global gap closes and long-time correlations need not decay
exponentially.  A continuum construction must therefore use compatible
multiscale finite horizons, retain the long-range Ward sector explicitly,
or prove a different cancellation estimate.
\end{firewall}

\begin{consequence}[Carrier direction after the massless test]
\label{cons:e8v4-direction}
For the photon/Ward sector, the carrier compiler must separate:
\begin{enumerate}
\item fixed-volume continuum contraction after the constant mode is
removed;
\item infinite-volume long-range modes, which cannot be assigned a uniform
global block gap; and
\item local refinement sources, which may still be summable by
\eqref{eq:e8v4-summability} before any long-time limit.
\end{enumerate}
The massive weak and Higgs sectors may use a genuine complementary gap, but
that gap cannot be transferred to the photon by representation bookkeeping.
\end{consequence}

\begin{proof}
Items 1 and 2 are \Cref{cor:e8v4-heat-block,ass:e8v4-order}.  Item 3 is
\Cref{cor:e8v4-summable}.  A positive mass changes the least eigenvalue,
whereas the unbroken electromagnetic mode remains massless; the last
statement is therefore a separation of hypotheses, not a claim about an
interacting mass value.
\end{proof}

\begin{assessment}[Progress at V4]
\label{ass:e8v4-status}
The fixed-block branch is viable for cutoff removal at fixed physical
volume, with duration proportional to \(L^2\).  It is not uniform in the
subsequent infinite-volume limit.  The gap-free heat-kernel estimate gives
a new local route: sufficiently distant refinement shells can be summable
at finite physical horizons even while the global Ward sector remains
massless.
\end{assessment}

\begin{programme}[Eighth-series V5: compulsory audit and finite-horizon
projectivity]
\label{prog:e8v4-V5}
V5 will perform the compulsory companion audit.  It will then formulate a
nonautonomous finite-horizon projective system whose adjacent source shells
obey \eqref{eq:e8v4-summability}, prove its local Cauchy property, and state
the extra compatibility needed to pass from finite-volume continuum states
to an infinite-volume state without claiming exponential clustering for
the photon sector.
\end{programme}

\begin{firewall}[Claim boundary after eighth-series V4]
\label{fire:e8v4-boundary}
V4 computes the exact mean-zero lattice-Laplacian gap and heat-block
duration, separates fixed-volume cutoff removal from the infinite-volume
limit, proves a volume- and cutoff-independent finite-time random-walk tail,
and derives a summable separated-source criterion.

V4 does not prove that the interacting SM channel is the heat semigroup,
construct its multiscale finite-horizon system, establish an
infinite-volume state, prove photon clustering, remove the cutoff, or
construct the Standard Model--Einstein continuum.
\end{firewall}

\begin{theorem}[Eighth-series V4 result]
\label{thm:e8v4-result}
After quotienting the constant mode, the periodic massless lattice
Laplacian has gap
\(4a^{-2}\sin^2(\pi/N)\).  It remains positive as
\(a\downarrow0\) at fixed physical side but closes as \(L^{-2}\) in the
infinite-volume limit, forcing global block time of order \(L^2\).  At any
finite physical time, remote source propagation instead obeys the
cutoff-independent Gaussian estimate \eqref{eq:e8v4-Gaussian-tail}, which
loads cofinal local projectivity under
\eqref{eq:e8v4-summability}.
\end{theorem}

\begin{proof}
Combine \Cref{thm:e8v4-gap,cor:e8v4-heat-block,
lem:e8v4-mgf,cor:e8v4-Gaussian,thm:e8v4-separation,
cor:e8v4-summable}.
\end{proof}

\paragraph{Parent-version record.}
V4 was formed from
\texttt{Renewal\_SM\_Einstein\_Continuum\_Research\_Notes\_V3.tex}.
The V3 parent had \(1\,191\) logical source lines, \(46\,246\) bytes, and
SHA--256
\[
 \texttt{AA524965426043B76E7FBA2E170928E1E74642493CD1A8386898ABCE2FA6FD61}.
\]
The next compulsory companion audit is V5.

% =============================================================================
% END APPENDED EIGHTH-SERIES V4 MATERIAL
% =============================================================================
% =============================================================================
% BEGIN APPENDED EIGHTH-SERIES V5 MATERIAL
% =============================================================================

\section{V5 companion audit and finite-horizon projective states}
\label{sec:v8v5}

This is the first compulsory audit of the eighth series.  It is followed by
a local projectivity theorem adapted to the massless finite-horizon bound of
V4.  The theorem constructs a state on an inductive local algebra from
summable remote-shell, recovery, and projector-transport defects.  It does
not infer a physical transfer gap from that state.

\subsection{Compulsory companion audit}
\label{sec:e8v5-audit}

The newest coherent companion sources in the shared folder are
\begin{align}
 &\texttt{Renewal\_Yang\_Mills\_Mass\_Gap\_Research\_Notes\_V55.tex},
 \notag\\
 &\qquad 24\,716\ \text{logical lines},\quad869\,366\ \text{bytes},
 \notag\\
 &\qquad
 \texttt{926F0047C720C8A34D69B3052259941275287D440D91A274CE99177771FF91D6},
 \label{eq:e8v5-YM-hash}\\
 &\texttt{Renewal\_NS\_Stability\_Research\_Notes\_V26.tex},
 \notag\\
 &\qquad 5\,319\ \text{logical lines},\quad278\,283\ \text{bytes},
 \notag\\
 &\qquad
 \texttt{82C887F8C2C69E4F28FAD7C3DC8DE5B9099CB5A4BF431EBA1FFF3E91935978AD}.
 \label{eq:e8v5-NS-hash}
\end{align}
Both have one live document terminator and end with
\verb|\end{document}|.  No companion file was modified.

Since the V100 audit, YM V50--V55 has developed the physical-clock branch
through reflected half-slab transfers, rescaled Dirichlet forms, Mosco gap
transport, and temporal conductance.  Three conclusions constrain the SM
programme directly:
\begin{enumerate}
\item equal-time measures, spatial area data, and spatial conditional
expectations do not determine a Hamiltonian gap without a physical
reflected-time incidence map;
\item if a transfer gap is \(g_n=\mu\tau_n+o(\tau_n)\), perturbative
transport of its protected spectral projector costs the inverse margin
\(g_n^{-1}\), so an unnormalized \(O(\tau_n)\) kernel error is generally
not negligible; and
\item an off-diagonal slow/fast generator block produces an
\(O(\tau_n)\) leakage source and cannot be hidden inside coercivity of the
compressed fast block.
\end{enumerate}
YM V55 does not populate its own temporal conductance or provide an SM
carrier generator.  NS V26 is unchanged and supplies no channel,
projector, heat-kernel, or regulator-bridge row for the present problem.

\begin{audit}[V5 companion decision]
\label{aud:e8v5-decision}
The three YM conclusions above are imported as abstract proof constraints.
In response, the finite-horizon ledger below includes a normalized
projector-transport row and an explicit leakage/recovery remainder.  No YM
constant or NS estimate is imported, and no mass conclusion is inferred
from the resulting algebraic state.
\end{audit}

\begin{proof}
Each imported item is an operator-theoretic conclusion stated with its
hypotheses in YM V54--V55.  The physical objects required to instantiate it
in the SM carrier are absent there, while the audited NS kernel norms have
no identified map to this channel system.  The restricted import is
therefore the strongest justified decision.
\end{proof}

\subsection{Near-critical projector transport}
\label{sec:e8v5-projector}

Let \(P_n\) be positive self-adjoint contractions on one transported
Hilbert space.  Suppose \(1\) is an isolated eigenvalue with spectral
projection \(\Pi_n\), and
\begin{equation}
 \spec(P_n)\setminus\{1\}\subset[0,1-g_n],
 \qquad g_n>0.
 \label{eq:e8v5-gap}
\end{equation}
Let \(\epsilon_n^P:=\|P_{n+1}-P_n\|\).

\begin{lemma}[Summable inverse-margin projector row]
\label{lem:e8v5-projector}
If \(\epsilon_n^P<g_n/4\), and \(\Pi_{n+1}\) is the top-cluster spectral
projection selected by the circle of radius \(g_n/2\) about \(1\), then
\begin{equation}
 \|\Pi_{n+1}-\Pi_n\|
 \leq p_n:= {4\epsilon_n^P\over g_n}.
 \label{eq:e8v5-projector-bound}
\end{equation}
If \(\sum_np_n<\infty\), the transported protected projectors converge in
operator norm.
\end{lemma}

\begin{proof}
The Riesz formulas on the stated circle and the resolvent identity give
\[
 \Pi_{n+1}-\Pi_n={1\over2\pi i}\int_\Gamma
 (z-P_{n+1})^{-1}(P_{n+1}-P_n)(z-P_n)^{-1}\,dz.
\]
The two resolvents are bounded by \(4/g_n\) and \(2/g_n\), respectively,
and the contour length is \(\pi g_n\), giving
\eqref{eq:e8v5-projector-bound}.  The triangle inequality and summability
make \((\Pi_n)\) Cauchy in the complete operator norm.
\end{proof}

\begin{remark}[Clock normalization]
\label{rem:e8v5-clock-projector}
When \(g_n\asymp\tau_n\), the summable row is
\(\epsilon_n^P/\tau_n\), not the unnormalized transfer difference.
Consequently an \(O(\tau_n)\) adjacent error may survive at generator
order; convergence by this perturbative route requires a smaller or
cancelling normalized defect.
\end{remark}

\subsection{Inductive local-algebra setup}
\label{sec:e8v5-inductive}

Let
\begin{equation}
 \mathfrak A_1\xrightarrow{\iota_1}\mathfrak A_2
 \xrightarrow{\iota_2}\cdots
 \label{eq:e8v5-inductive-system}
\end{equation}
be unital isometric \(*\)-embeddings of finite-dimensional
\(C^*\)-algebras.  Write \(\iota_{m,n}\) for their composition and
\(\mathfrak A_{\rm loc}\) for the algebraic direct limit.  Let
\(\omega_n\) be a state on \(\mathfrak A_n\).  For
\(A\in\mathfrak A_m\), abbreviate
\begin{equation}
 F_n^A:=\omega_n(\iota_{m,n}A),\qquad n\geq m.
 \label{eq:e8v5-expectation}
\end{equation}

At adjacent scale \(n\), assume the expectation defect splits as
\begin{equation}
 F_{n+1}^A-F_n^A
 =H_n^A+R_n^A+J_n^A,
 \label{eq:e8v5-split}
\end{equation}
where \(H_n^A\) is a finite-horizon propagated refinement source,
\(R_n^A\) contains recovery, contact, and slow/fast leakage, and
\(J_n^A\) is protected-projector transport.

Suppose the refinement source has trace or \(\ell^1\) size at most
\(\varepsilon_n\), is separated from the support \(B_A\) by physical
distance \(R_{n,A}\), and propagates for time \(T_n\) under a kernel
dominated by the V4 heat kernel.  Define
\begin{equation}
 \Theta_{n,A}:=
 2^d\exp\left\{-{R_{n,A}^2
 \over4d(\cosh1)T_n}\right\},
 \label{eq:e8v5-Theta}
\end{equation}
assuming the V4 Gaussian-regime condition.  Let
\(r_{n,A},j_{n,A}\geq0\) obey
\begin{equation}
 |R_n^A|\leq\|A\|r_{n,A},
 \qquad
 |J_n^A|\leq\|A\|j_{n,A}.
 \label{eq:e8v5-RJ}
\end{equation}
The projector row in \Cref{lem:e8v5-projector} is one sufficient source of
\(j_{n,A}\leq C_Ap_n\); its coefficient must be proved for the actual
observable transport.

\begin{lemma}[Finite-horizon adjacent bound]
\label{lem:e8v5-adjacent}
Under the preceding hypotheses,
\begin{equation}
 |F_{n+1}^A-F_n^A|
 \leq\|A\|
 \bigl(\varepsilon_n\Theta_{n,A}+r_{n,A}+j_{n,A}\bigr).
 \label{eq:e8v5-adjacent-bound}
\end{equation}
\end{lemma}

\begin{proof}
The V4 separation theorem and the kernel domination bound
\(|H_n^A|\) by
\(\|A\|\varepsilon_n\Theta_{n,A}\).  Add the two estimates in
\eqref{eq:e8v5-RJ} and use \eqref{eq:e8v5-split}.
\end{proof}

\begin{theorem}[Finite-horizon projective state landing]
\label{thm:e8v5-state}
Suppose that for every \(A\) in a norm-dense countable unital local
\(*\)-subalgebra \(\mathfrak A_{\rm loc}^0\),
\begin{equation}
 \sum_{n\geq m(A)}
 \bigl(\varepsilon_n\Theta_{n,A}+r_{n,A}+j_{n,A}\bigr)<\infty.
 \label{eq:e8v5-summable}
\end{equation}
Then \(F_n^A\) converges for every
\(A\in\mathfrak A_{\rm loc}^0\).  The limit extends uniquely to a state
\(\omega_\infty\) on the norm completion
\(\mathfrak A_{\rm ql}\) of the inductive local algebra.
\end{theorem}

\begin{proof}
By \Cref{lem:e8v5-adjacent}, the adjacent differences of \(F_n^A\) are
absolutely summable, so \(F_n^A\) is Cauchy.  Coherent embeddings make the
pointwise limit linear on \(\mathfrak A_{\rm loc}^0\).  Since every
\(\omega_n\) is a state, the limit is positive on positive local elements,
sends the identity to one, and satisfies
\(|\omega_\infty(A)|\leq\|A\|\).  It therefore extends uniquely by norm
continuity to \(\mathfrak A_{\rm ql}\), retaining positivity and norm one.
\end{proof}

\begin{corollary}[No exponential-clustering conclusion]
\label{cor:e8v5-no-clustering}
The conclusion of \Cref{thm:e8v5-state} remains valid when the global
mean-zero heat gap closes with volume.  It supplies existence of a
quasilocal algebraic state, but it gives neither uniqueness nor exponential
clustering.
\end{corollary}

\begin{proof}
The proof uses only the summable local adjacent rows
\eqref{eq:e8v5-summable}, not a volume-uniform global gap.  State existence
alone does not imply uniqueness: distinct thermodynamic phases are states
on the same quasilocal algebra.  It also does not imply any quantitative
bound on connected correlations.
\end{proof}

\subsection{Metric derivatives and common-refinement independence}
\label{sec:e8v5-derivative}

Let the states and adjacent decomposition depend on a metric parameter
\(q\) in a bounded convex chart \(Q\).  Suppose every
\(F_n^A\in C^1(Q)\) and a derivative ledger
\(d_{n,A}^{(1)}\) satisfies
\begin{equation}
 \sup_{q\in Q,\ \|v\|\leq1}
 |DF_{n+1}^A(q)[v]-DF_n^A(q)[v]|
 \leq d_{n,A}^{(1)},
 \qquad
 \sum_nd_{n,A}^{(1)}<\infty.
 \label{eq:e8v5-derivative-row}
\end{equation}

\begin{theorem}[Finite-horizon \(C^1\) landing]
\label{thm:e8v5-C1}
If \eqref{eq:e8v5-summable} holds uniformly on \(Q\) and
\eqref{eq:e8v5-derivative-row} holds, then
\(F_n^A\to F_\infty^A\) in \(C^1(Q)\) for every protected local
observable \(A\).  Its derivative is the uniform limit of the finite
carrier stress responses.
\end{theorem}

\begin{proof}
The base row gives uniform convergence by the Weierstrass test.  The
derivative row makes \((DF_n^A)\) uniformly Cauchy in the Banach space of
bounded continuous dual-valued maps.  Integrating along line segments from
a fixed basepoint identifies the limit of the derivatives with the
derivative of the limiting function.
\end{proof}

\begin{theorem}[Local common-refinement bridge]
\label{thm:e8v5-bridge}
Let two regulator systems produce local limiting states
\(\omega_\infty\) and \(\widetilde\omega_\infty\).  Suppose that for every
\(A\in\mathfrak A_{\rm loc}^0\) there are common refinements \(t_n\) and
numbers \(b_{n,A},\widetilde b_{n,A}\to0\) such that
\begin{align}
 |\omega_{r_n}(A)-\omega_{t_n}(A)|&\leq b_{n,A},
 \label{eq:e8v5-bridge-one}\\
 |\widetilde\omega_{s_n}(A)-\omega_{t_n}(A)|
 &\leq\widetilde b_{n,A}.
 \label{eq:e8v5-bridge-two}
\end{align}
Then the two limiting states agree on
\(\mathfrak A_{\rm ql}\).
\end{theorem}

\begin{proof}
For every protected local \(A\), the triangle inequality gives
\[
 |\omega_{r_n}(A)-\widetilde\omega_{s_n}(A)|
 \leq b_{n,A}+\widetilde b_{n,A}\longrightarrow0.
\]
Pass to the two local limits.  Equality on the norm-dense local algebra and
continuity of states give equality on its completion.
\end{proof}

\begin{firewall}[Static state versus physical transfer]
\label{fire:e8v5-transfer}
Neither \Cref{thm:e8v5-state} nor \Cref{thm:e8v5-C1} determines a
Hamiltonian mass or Lorentzian dynamics.  As emphasized by the audited YM
counterexample, the same stationary state can be paired with positive
self-adjoint Markov transfers having arbitrarily chosen relaxation rates.
A mass statement requires the actual reflected-time pair law, physical
clock, protected sector, and limiting form or transfer theorem.
\end{firewall}

\begin{definition}[Finite-horizon projective packet, FHPP]
\label{def:e8v5-FHPP}
FHPP consists of:
\begin{enumerate}
\item coherent local \(C^*\)-algebras, embeddings, and finite states;
\item the adjacent split \eqref{eq:e8v5-split};
\item heat-kernel or stronger domination of the remote refinement source;
\item explicit recovery, contact, slow/fast leakage, and normalized
projector-transport rows;
\item summability \eqref{eq:e8v5-summable} on a dense protected core;
\item the derivative row \eqref{eq:e8v5-derivative-row} when stress is
claimed; and
\item common-refinement bridges when regulator independence is claimed.
\end{enumerate}
\end{definition}

\begin{assessment}[Progress at V5]
\label{ass:e8v5-status}
The massless carrier branch now has an honest state-level endpoint.  Remote
shells may be summed by finite-horizon heat locality even though the global
photon gap closes.  Projector rotation and slow/fast leakage remain explicit
adjacent sources.  The next missing physical step is a channel or Duhamel
comparison proving the split \eqref{eq:e8v5-split} for the interacting
old/detail Gibbs carrier.
\end{assessment}

\begin{programme}[Eighth-series V6: interacting finite-horizon comparison]
\label{prog:e8v5-V6}
V6 will compare two bounded local generators by Duhamel's formula inside a
fixed physical horizon.  It will retain a near part and a remote shell,
propagate the latter by a Lieb--Robinson or heat-kernel majorant, and derive
an explicit adjacent expectation bound.  If the generator difference is
unbounded, the programme will move upstream to relative quadratic-form
Duhamel theory rather than use a formal operator exponential.
\end{programme}

\begin{firewall}[Claim boundary after eighth-series V5]
\label{fire:e8v5-boundary}
V5 performs the compulsory companion audit, adds the inverse-margin
projector row, proves finite-horizon local state and \(C^1\) landing, and
proves a local common-refinement criterion.  It expressly separates an
algebraic state from a physical reflected transfer gap.

V5 does not derive the adjacent split for the interacting SM carrier,
control an unbounded generator, prove the projector or leakage rows,
construct reflected-time incidence, establish regulator independence,
remove the cutoff, or construct the Standard Model--Einstein continuum.
\end{firewall}

\begin{theorem}[Eighth-series V5 result]
\label{thm:e8v5-result}
Summability of the finite-horizon remote-shell, recovery/leakage, and
normalized projector-transport rows yields a state on the quasilocal
inductive algebra.  A separately summable derivative row yields local
\(C^1\) metric response, and common-refinement bridges yield regulator
independence.  None of these state-level conclusions determines a physical
transfer gap.
\end{theorem}

\begin{proof}
Combine \Cref{lem:e8v5-projector,lem:e8v5-adjacent,
thm:e8v5-state,thm:e8v5-C1,thm:e8v5-bridge,
fire:e8v5-transfer}.
\end{proof}

\paragraph{Parent-version record.}
V5 was formed from
\texttt{Renewal\_SM\_Einstein\_Continuum\_Research\_Notes\_V4.tex}.
The V4 parent had \(1\,524\) logical source lines, \(58\,631\) bytes, and
SHA--256
\[
 \texttt{A3FB7DF83325B14A8DE6812BC2A10064B8CAB018205EBA30950D3C2B37018EE2}.
\]
The next compulsory companion audit is V10.

% =============================================================================
% END APPENDED EIGHTH-SERIES V5 MATERIAL
% =============================================================================
% =============================================================================
% BEGIN APPENDED EIGHTH-SERIES V6 MATERIAL
% =============================================================================

\section{Interacting finite-horizon Duhamel comparison}
\label{sec:v8v6}

V5 postulates an adjacent finite-horizon split.  We now derive that split
for bounded local generators transported to one finite collar algebra.  The
near generator mismatch, remote shell, initial carrier mismatch, and
observable contact term remain separate.  A remote shell is small only
after a proved propagation estimate; complete positivity alone gives no
distance factor.

\subsection{Exact comparison identity}
\label{sec:e8v6-Duhamel}

Let \(\mathcal L_0,\mathcal L_1\) be bounded generators of CPTP semigroups
\(\mathcal S_j(t)=e^{t\mathcal L_j}\) on a finite matrix algebra.  Put
\(\mathcal V=\mathcal L_1-\mathcal L_0\).

\begin{lemma}[Superoperator Duhamel identity]
\label{lem:e8v6-Duhamel}
For every \(T\geq0\),
\begin{equation}
 \mathcal S_1(T)-\mathcal S_0(T)
 =\int_0^T\mathcal S_1(T-u)\mathcal V\mathcal S_0(u)\,du.
 \label{eq:e8v6-Duhamel}
\end{equation}
Consequently,
\begin{equation}
 \|\mathcal S_1(T)-\mathcal S_0(T)\|_\diamond
 \leq T\|\mathcal V\|_\diamond.
 \label{eq:e8v6-global}
\end{equation}
\end{lemma}

\begin{proof}
Differentiate
\(u\mapsto\mathcal S_1(T-u)\mathcal S_0(u)\).  Its derivative is
\(-\mathcal S_1(T-u)\mathcal V\mathcal S_0(u)\).  Integrating from zero
to \(T\) and reversing signs gives \eqref{eq:e8v6-Duhamel}.  CPTP maps
have diamond norm one, so submultiplicativity and integration give
\eqref{eq:e8v6-global}.
\end{proof}

Let \(\rho_j\) be initial densities and \(A_j\) observables on the common
collar.  Define
\begin{equation}
 F_j(T):=\Tr\bigl(\mathcal S_j(T)\rho_j A_j\bigr).
 \label{eq:e8v6-F}
\end{equation}

\begin{proposition}[Four-term adjacent decomposition]
\label{prop:e8v6-four}
The difference can be written exactly as
\begin{align}
 F_1(T)-F_0(T)
 ={}&\Tr\bigl(\mathcal S_1(T)(\rho_1-\rho_0)A_1\bigr)
 \label{eq:e8v6-initial-term}\\
 &+\Tr\bigl((\mathcal S_1(T)-\mathcal S_0(T))\rho_0A_1\bigr)
 \label{eq:e8v6-generator-term}\\
 &+\Tr\bigl(\mathcal S_0(T)\rho_0(A_1-A_0)\bigr).
 \label{eq:e8v6-contact-term}
\end{align}
After splitting \(\mathcal V=\mathcal V^{\rm near}+\mathcal V^{\rm far}\),
Duhamel's formula splits the middle line into near and far generator terms.
Moreover,
\begin{align}
 |\eqref{eq:e8v6-initial-term}|
 &\leq\|\rho_1-\rho_0\|_1\|A_1\|,
 \label{eq:e8v6-initial-bound}\\
 |\eqref{eq:e8v6-contact-term}|
 &\leq\|A_1-A_0\|.
 \label{eq:e8v6-contact-bound}
\end{align}
\end{proposition}

\begin{proof}
Insert and subtract
\(\Tr(\mathcal S_1(T)\rho_0A_1)\) and
\(\Tr(\mathcal S_0(T)\rho_0A_1)\).  This gives the three displayed
lines; linearity gives the near/far split.  Trace-norm contractivity on
Hermitian differences and unital contractivity of the dual semigroup prove
\eqref{eq:e8v6-initial-bound}--\eqref{eq:e8v6-contact-bound}.
\end{proof}

\subsection{Remote-shell propagation}
\label{sec:e8v6-remote}

Assume the generator difference has a local decomposition
\begin{equation}
 \mathcal V=\sum_X\mathcal V_X,
 \label{eq:e8v6-local-decomposition}
\end{equation}
where \(X\) ranges over finite cell sets and each \(\mathcal V_X\) is
trace annihilating.  Let \(A_B\) be supported in \(B\).  We use the
following explicit Heisenberg propagation hypothesis for
\(\mathcal S_1^*\):
\begin{equation}
 \|\mathcal V_X^*\mathcal S_1^*(t)A_B\|
 \leq C_{\rm LR}|B|\,\|\mathcal V_X\|_\diamond\|A_B\|
 e^{-\mu(d(X,B)-v t)}
 \label{eq:e8v6-LR}
\end{equation}
whenever \(d(X,B)\geq vt\).  This is a Lieb--Robinson-type insertion
bound.  It must be proved from the actual interaction norm; it is not a
consequence of \eqref{eq:e8v6-local-decomposition} alone.

Choose \(R\geq vT\) and define
\begin{align}
 \mathcal V^{\rm near}&:=\sum_{d(X,B)<R}\mathcal V_X,
 &
 \mathcal V^{\rm far}&:=\sum_{d(X,B)\geq R}\mathcal V_X,
 \label{eq:e8v6-near-far}\\
 \nu_{\rm near}(B,R)&:=
 \sum_{d(X,B)<R}\|\mathcal V_X\|_\diamond,
 &
 \nu_{\mu,\rm far}(B,R)&:=
 \sum_{d(X,B)\geq R}
 e^{-\mu d(X,B)}\|\mathcal V_X\|_\diamond.
 \label{eq:e8v6-nu}
\end{align}

\begin{theorem}[Near/far generator bound]
\label{thm:e8v6-near-far}
For a common initial density \(\rho\) and local observable \(A_B\),
\begin{align}
 &\left|\Tr\bigl(
 (\mathcal S_1(T)-\mathcal S_0(T))\rho A_B\bigr)\right|
 \notag\\
 &\quad\leq T\|A_B\|\nu_{\rm near}(B,R)
 +C_{\rm LR}|B|T\|A_B\|e^{\mu vT}
 \nu_{\mu,\rm far}(B,R).
 \label{eq:e8v6-near-far-bound}
\end{align}
If every far component is supported at distance at least \(R\) and
\(\sum_X\|\mathcal V_X\|_\diamond\leq\nu_{\rm far}\), then the second
term is at most
\begin{equation}
 C_{\rm LR}|B|T\|A_B\|\nu_{\rm far}
 e^{-\mu(R-vT)}.
 \label{eq:e8v6-shell-bound}
\end{equation}
\end{theorem}

\begin{proof}
For the near part, use \eqref{eq:e8v6-global} componentwise and sum.  For
the far part, pair \eqref{eq:e8v6-Duhamel} with \(A_B\) and move all maps
to the observable side:
\begin{align*}
 &\Tr\bigl(\mathcal S_1(T-u)\mathcal V_X
                  \mathcal S_0(u)\rho\,A_B\bigr)\\
 &\qquad=\Tr\bigl(\mathcal S_0(u)\rho\,
       \mathcal V_X^*\mathcal S_1^*(T-u)A_B\bigr).
\end{align*}
The state has norm one.  Since
\(d(X,B)\geq R\geq vT\), apply \eqref{eq:e8v6-LR} with
\(t=T-u\), bound \(e^{\mu v(T-u)}\leq e^{\mu vT}\), integrate over
\(u\), and sum \(X\).  This proves
\eqref{eq:e8v6-near-far-bound}.  The last assertion uses
\(e^{-\mu d(X,B)}\leq e^{-\mu R}\).
\end{proof}

\begin{remark}[Heat-kernel alternative]
\label{rem:e8v6-heat}
If the Heisenberg insertion is dominated by the V4 heat kernel instead of
\eqref{eq:e8v6-LR}, the same proof replaces
\(e^{-\mu(R-vT)}\) by the Gaussian factor
\(2^d\exp\{-R^2/[4d(\cosh1)T]\}\).  The Duhamel time integral may only
improve that bound because every intermediate propagation time is at most
\(T\).
\end{remark}

\subsection{Explicit adjacent loader}
\label{sec:e8v6-loader}

At scale \(n\), write
\begin{align}
 \kappa_n&:=\|\rho_{n+1}-\rho_n\|_1,
 &a_n(A)&:=\|A_{n+1}-A_n\|,
 \label{eq:e8v6-kappa-a}\\
 N_n(A)&:=T_n\nu_{{\rm near},n}(B_A,R_n),
 &H_n(A)&:=C_{\rm LR}|B_A|T_n e^{\mu vT_n}
 \nu_{\mu,\rm far,n}(B_A,R_n).
 \label{eq:e8v6-NH}
\end{align}
Projector transport and recovery/leakage remainders may be included in
\(\kappa_n\) when a trace-norm state comparison is available, or kept as
separate V5 rows.

\begin{corollary}[Bounded-generator population of the V5 split]
\label{cor:e8v6-populate}
For every protected local \(A\),
\begin{equation}
 |F_{n+1}^A-F_n^A|
 \leq \|A\|\bigl(\kappa_n+N_n(A)+H_n(A)\bigr)+a_n(A),
 \label{eq:e8v6-adjacent-loader}
\end{equation}
up to the separately declared projector/contact normalization constants.
If the right side is summable for every element of a dense local core, the
V5 finite-horizon state landing theorem applies.
\end{corollary}

\begin{proof}
Combine \Cref{prop:e8v6-four,thm:e8v6-near-far}.  Summability invokes
\Cref{thm:e8v5-state}.
\end{proof}

\begin{corollary}[Pure remote-shell sufficient condition]
\label{cor:e8v6-pure-shell}
Suppose eventually \(\nu_{{\rm near},n}=0\),
\(T_n\leq T_*\), and
\(\sum_X\|\mathcal V_{n,X}\|_\diamond\leq\eta_n\).  Then the generator
row is summable if
\begin{equation}
 \sum_n\eta_n e^{-\mu(R_n-vT_*)}<\infty.
 \label{eq:e8v6-shell-sum}
\end{equation}
The initial-state and observable-contact rows remain independent.
\end{corollary}

\begin{proof}
Use \eqref{eq:e8v6-shell-bound}, absorb fixed constants, and apply the
Weierstrass test.  Duhamel's identity contains no relation between the
generator row and the other two rows, proving the last statement.
\end{proof}

\subsection{Form-domain boundary}
\label{sec:e8v6-form}

\begin{firewall}[Unbounded generators]
\label{fire:e8v6-unbounded}
The identities above are operator-norm statements for bounded finite-cutoff
generators.  If the limiting electric, Dirac, or gravitational generator is
unbounded, convergence of formal exponentials and the estimate
\(T\|\mathcal V\|\) cannot be used.  One needs a common closed quadratic
form domain, relative form bounds, strong-resolvent or Mosco convergence,
and a local observable propagation theorem compatible with that limit.
The audited YM V53--V55 programme reaches the same boundary for its
reflected transfer.
\end{firewall}

\begin{assessment}[What V6 resolves]
\label{ass:e8v6-status}
The V5 adjacent split is now an exact consequence of bounded-generator
Duhamel comparison plus one stated propagation bound.  Remote old/detail
terms acquire a distance factor, while near generator mismatch, initial
carrier mismatch, observable contact, projector rotation, and leakage
remain visible.  The missing stress row is the metric derivative of this
comparison; differentiating it will create two ordered generator
insertions and derivatives of every boundary term.
\end{assessment}

\begin{programme}[Eighth-series V7: differentiated finite-horizon
comparison]
\label{prog:e8v6-V7}
V7 will differentiate \eqref{eq:e8v6-Duhamel} in a metric direction,
derive the exact mixed adjacent/stress formula, and bound its ordered
insertions by a two-source Lieb--Robinson convolution.  It will keep the
metric derivatives of the initial state, observable, projector, and
recovery rows explicit.
\end{programme}

\begin{firewall}[Claim boundary after eighth-series V6]
\label{fire:e8v6-boundary}
V6 proves the exact bounded-generator adjacent Duhamel identity, its
four-term state/observable/generator decomposition, the near/far
Lieb--Robinson shell bound, and an explicit sufficient loader for V5
finite-horizon state projectivity.

V6 does not prove the propagation hypothesis for the interacting SM
generator, control near or unbounded generator differences, populate the
initial/projector/recovery rows, derive the stress comparison, remove the
cutoff, or construct the Standard Model--Einstein continuum.
\end{firewall}

\begin{theorem}[Eighth-series V6 result]
\label{thm:e8v6-result}
For two bounded local CPTP generators on one transported collar, the
adjacent finite-horizon expectation defect is bounded by
\eqref{eq:e8v6-adjacent-loader}.  A remote refinement shell is suppressed
by \(e^{-\mu(R-vT)}\), or by the V4 Gaussian heat factor under heat-kernel
domination.  All near, initial-state, observable-contact, projector, and
recovery/leakage terms remain separate summability obligations.
\end{theorem}

\begin{proof}
Combine \Cref{lem:e8v6-Duhamel,prop:e8v6-four,
thm:e8v6-near-far,cor:e8v6-populate,cor:e8v6-pure-shell}.
\end{proof}

\paragraph{Parent-version record.}
V6 was formed from
\texttt{Renewal\_SM\_Einstein\_Continuum\_Research\_Notes\_V5.tex}.
The V5 parent had \(1\,902\) logical source lines, \(73\,463\) bytes, and
SHA--256
\[
 \texttt{87E26FEBBB7B6157BBDD7DA0B40C3A60761B195F007894345CE7B2508F75C3A7}.
\]
The next compulsory companion audit is V10.

% =============================================================================
% END APPENDED EIGHTH-SERIES V6 MATERIAL
% =============================================================================
% =============================================================================
% BEGIN APPENDED EIGHTH-SERIES V7 MATERIAL
% =============================================================================

\section{Differentiated finite-horizon comparison and two-source locality}
\label{sec:v8v7}

V6 controls adjacent expectations.  The Einstein carrier source is the
metric derivative of that comparison.  Differentiating the generator
Duhamel identity creates one direct derivative of the refinement generator
and two ordered insertions coupling the refinement difference to the
metric generator.  These are the finite-horizon counterparts of the mixed
channel sources isolated in the preceding series.

\subsection{Exact differentiated Duhamel formula}
\label{sec:e8v7-formula}

Let all V6 objects depend \(C^1\) on a scalar metric parameter \(q\).  A
directional derivative on a Banach metric chart is obtained by restricting
to a line.  Put
\begin{equation}
 \mathcal V:=\mathcal L_1-\mathcal L_0,
 \qquad
 \mathcal V_q:=\mathcal L_{1,q}-\mathcal L_{0,q},
 \qquad
 \mathcal K(T):=\mathcal S_1(T)-\mathcal S_0(T).
 \label{eq:e8v7-VK}
\end{equation}

\begin{lemma}[Semigroup metric derivative]
\label{lem:e8v7-semigroup}
For \(j\in\{0,1\}\),
\begin{equation}
 \mathcal S_{j,q}(t)
 =\int_0^t\mathcal S_j(t-r)\mathcal L_{j,q}
                   \mathcal S_j(r)\,dr,
 \label{eq:e8v7-Sq}
\end{equation}
and
\begin{equation}
 \|\mathcal S_{j,q}(t)\|_\diamond
 \leq t\|\mathcal L_{j,q}\|_\diamond.
 \label{eq:e8v7-Sq-bound}
\end{equation}
\end{lemma}

\begin{proof}
This is the first Duhamel derivative formula from V1.  Both surrounding
semigroups are CPTP and hence have diamond norm one, giving the bound.
\end{proof}

\begin{theorem}[Exact adjacent stress operator]
\label{thm:e8v7-Kq}
The derivative of the adjacent channel difference is
\begin{align}
 \mathcal K_q(T)
 ={}&\int_0^T
 \mathcal S_1(T-u)\mathcal V_q\mathcal S_0(u)\,du
 \label{eq:e8v7-direct}\\
 &+\int_0^T\int_0^{T-u}
 \mathcal S_1(T-u-r)\mathcal L_{1,q}\mathcal S_1(r)
 \mathcal V\mathcal S_0(u)\,dr\,du
 \label{eq:e8v7-left-ordered}\\
 &+\int_0^T\int_0^u
 \mathcal S_1(T-u)\mathcal V\mathcal S_0(u-r)
 \mathcal L_{0,q}\mathcal S_0(r)\,dr\,du.
 \label{eq:e8v7-right-ordered}
\end{align}
Consequently,
\begin{equation}
 \boxed{
 \|\mathcal K_q(T)\|_\diamond
 \leq T\|\mathcal V_q\|_\diamond
 +{T^2\over2}\|\mathcal V\|_\diamond
 \bigl(\|\mathcal L_{1,q}\|_\diamond
       +\|\mathcal L_{0,q}\|_\diamond\bigr).}
 \label{eq:e8v7-global-Kq}
\end{equation}
\end{theorem}

\begin{proof}
Differentiate the V6 identity
\[
 \mathcal K(T)=\int_0^T
 \mathcal S_1(T-u)\mathcal V\mathcal S_0(u)\,du.
\]
The derivative can hit the left semigroup, \(\mathcal V\), or the right
semigroup.  Substitute \eqref{eq:e8v7-Sq} in the two semigroup terms; this
is exactly \eqref{eq:e8v7-direct}--
\eqref{eq:e8v7-right-ordered}.  In diamond norm, the direct interval has
length \(T\), while each ordered triangular domain has area \(T^2/2\).
All semigroups have norm one, proving \eqref{eq:e8v7-global-Kq}.
\end{proof}

\subsection{Complete derivative of the adjacent expectation}
\label{sec:e8v7-expectation}

Retain the notation of V6 and put
\begin{equation}
 \delta\rho:=\rho_1-\rho_0,qquad
 \delta A:=A_1-A_0.
 \label{eq:e8v7-deltas}
\end{equation}

\begin{theorem}[Global differentiated adjacent bound]
\label{thm:e8v7-global}
Let \(F_j(T,q)=\Tr(\mathcal S_j(T)\rho_jA_j)\).  Then
\begin{align}
 |\partial_q(F_1-F_0)|
 \leq{}&
 \|A_1\|\|\delta\rho_q\|_1
 +\|\delta\rho\|_1\|A_{1,q}\|
 +T\|\mathcal L_{1,q}\|_\diamond
       \|\delta\rho\|_1\|A_1\|
 \notag\\
 &+\|A_1\|\|\mathcal K_q(T)\|_\diamond
 +T\|\mathcal V\|_\diamond
   \bigl(\|\rho_{0,q}\|_1\|A_1\|+\|A_{1,q}\|\bigr)
 \notag\\
 &+T\|\mathcal L_{0,q}\|_\diamond\|\delta A\|
 +\|\rho_{0,q}\|_1\|\delta A\|
 +\|\delta A_q\|,
 \label{eq:e8v7-global-adjacent}
\end{align}
where \(\|\mathcal K_q(T)\|_\diamond\) is bounded by
\eqref{eq:e8v7-global-Kq}.
\end{theorem}

\begin{proof}
Differentiate the exact three-line decomposition in
\Cref{prop:e8v6-four}.  The initial-state line produces the first three
terms of \eqref{eq:e8v7-global-adjacent}, using
\eqref{eq:e8v7-Sq-bound}.  The generator line produces
\(\Tr(\mathcal K_q\rho_0A_1)\),
\(\Tr(\mathcal K\rho_{0,q}A_1)\), and
\(\Tr(\mathcal K\rho_0A_{1,q})\); use
\(\|\mathcal K\|_\diamond\leq T\|\mathcal V\|_\diamond\).
The contact line produces the last three terms.  Trace/diamond duality and
the triangle inequality prove the result.
\end{proof}

\begin{remark}[Direct contacts are indispensable]
\label{rem:e8v7-contact}
Even when \(\mathcal V=0\), the differentiated adjacent expectation can be
nonzero through \(\delta\rho_q\), \(\delta A_q\), or metric derivatives of
the common initial carrier.  A stress comparison which differentiates only
the semigroup therefore omits genuine contact terms.
\end{remark}

\subsection{Two-source propagation for a remote refinement shell}
\label{sec:e8v7-two-source}

Write
\begin{equation}
 \mathcal V=\sum_X\mathcal V_X,qquad
 \mathcal V_q=\sum_X\mathcal V_{q,X},qquad
 \mathcal L_{j,q}=\sum_Y\mathcal W_{j,Y}.
 \label{eq:e8v7-components}
\end{equation}
For finite sets \(B,X,Y\), let \(\mathfrak t(B,X,Y)\) be the minimum total
length of a tree in the cell graph meeting all three sets.  In particular,
\begin{equation}
 \mathfrak t(B,X,Y)\geq d(B,X).
 \label{eq:e8v7-tree-lower}
\end{equation}

Besides the V6 one-source estimate, assume the ordered two-source bound
\begin{align}
 &\|\mathcal V_X^*\mathcal S_a^*(t_1)
       \mathcal W_{j,Y}^*\mathcal S_b^*(t_2)A_B\|
 \notag\\
 &\quad+\|\mathcal W_{j,Y}^*\mathcal S_a^*(t_1)
       \mathcal V_X^*\mathcal S_b^*(t_2)A_B\|
 \notag\\
 &\qquad\leq C_2|B|\|A_B\|
 \|\mathcal V_X\|_\diamond\|\mathcal W_{j,Y}\|_\diamond
 e^{-\mu(\mathfrak t(B,X,Y)-v(t_1+t_2))}
 \label{eq:e8v7-LR2}
\end{align}
for all semigroup choices \(a,b\in\{0,1\}\) appearing in the ordered
Duhamel terms.  This is the exact extra locality input required by the
metric derivative.

Define
\begin{align}
 M_{j,q}&:=\sum_Y\|\mathcal W_{j,Y}\|_\diamond,
 \label{eq:e8v7-Mq}\\
 \nu_{q,\mu,\rm far}(B,R)
 &:=\sum_{d(X,B)\geq R}e^{-\mu d(X,B)}
          \|\mathcal V_{q,X}\|_\diamond,
 \label{eq:e8v7-nuq}\\
 \nu_{\mu,\rm far}(B,R)
 &:=\sum_{d(X,B)\geq R}e^{-\mu d(X,B)}
          \|\mathcal V_X\|_\diamond.
 \label{eq:e8v7-nuv}
\end{align}

\begin{theorem}[Remote-shell adjacent stress bound]
\label{thm:e8v7-remote}
Assume \(R\geq vT\), the V6 one-source bound for
\(\mathcal V_{q,X}\), and \eqref{eq:e8v7-LR2}.  The contribution of
components \(X\) with \(d(X,B)\geq R\) to the metric derivative of the
common-state/common-observable channel difference is bounded by
\begin{align}
 C|B|\|A_B\|e^{\mu vT}
 \left[
 T\nu_{q,\mu,\rm far}(B,R)
 +{T^2\over2}(M_{1,q}+M_{0,q})
       \nu_{\mu,\rm far}(B,R)
 \right],
 \label{eq:e8v7-remote-bound}
\end{align}
where \(C=\max\{C_{\rm LR},C_2\}\).  If all refinement components are at
distance at least \(R\), the right side is at most the same unweighted
component sums multiplied by \(e^{-\mu(R-vT)}\).
\end{theorem}

\begin{proof}
Pair the direct term \eqref{eq:e8v7-direct} with \(A_B\), move the maps to
the observable side, and apply the V6 one-source estimate to
\(\mathcal V_{q,X}\).  This gives the first term in brackets.

For each ordered term, move all maps to the observable side and apply
\eqref{eq:e8v7-LR2}.  The total propagation time is at most \(T\), and
\eqref{eq:e8v7-tree-lower} leaves the factor
\(e^{-\mu d(B,X)}e^{\mu vT}\).  Sum first over \(Y\), producing
\(M_{j,q}\), then over \(X\).  Each triangular time domain has area
\(T^2/2\), proving \eqref{eq:e8v7-remote-bound}.  The last assertion uses
\(e^{-\mu d(B,X)}\leq e^{-\mu R}\).
\end{proof}

\begin{firewall}[One-source locality does not imply two-source locality]
\label{fire:e8v7-LR2}
The V6 bound \eqref{eq:e8v6-LR} controls one insertion against a local
observable.  It does not, by itself, prove \eqref{eq:e8v7-LR2}; after the
first insertion the observable is generally quasilocal.  A proof of the
two-source estimate requires an interaction-norm convolution theorem or a
Dyson-path expansion retaining both insertion supports.
\end{firewall}

\subsection{Stress-projectivity loader}
\label{sec:e8v7-loader}

Let \(D_n^{\rm bdry}(A;v)\) denote the seven non-generator rows in
\eqref{eq:e8v7-global-adjacent}: differentiated initial state, observable,
projector, contact, and recovery/leakage terms.  Let
\(D_n^{\rm near}(A;v)\) be the global bound
\eqref{eq:e8v7-global-Kq} restricted to near generator components, and let
\(D_n^{\rm far}(A;v)\) be \eqref{eq:e8v7-remote-bound}.

\begin{corollary}[Finite-horizon stress row]
\label{cor:e8v7-stress-row}
If for every protected observable and unit metric tangent
\begin{equation}
 \sum_n\sup_{q\in Q}
 \bigl(D_n^{\rm bdry}(A;v)+D_n^{\rm near}(A;v)
       +D_n^{\rm far}(A;v)\bigr)<\infty,
 \label{eq:e8v7-stress-sum}
\end{equation}
then the V5 derivative row holds and the finite-horizon carrier
expectations land in \(C^1(Q)\).
\end{corollary}

\begin{proof}
Theorem \ref{thm:e8v7-global} bounds all boundary and near terms, while
\Cref{thm:e8v7-remote} bounds the far terms.  Their sum is an adjacent
derivative majorant.  Apply the V5 \(C^1\) landing theorem.
\end{proof}

\begin{assessment}[Progress at V7]
\label{ass:e8v7-status}
Both base and stress comparisons now have exact finite-horizon formulas.
The stress row exposes a new physical input, \eqref{eq:e8v7-LR2}, rather
than silently reusing a one-source bound.  The most economical next step is
to derive that estimate from a weighted interaction norm for a finite-range
Lindblad or Markov generator.
\end{assessment}

\begin{programme}[Eighth-series V8: Dyson-path proof of two-source
locality]
\label{prog:e8v7-V8}
V8 will introduce a finite-range interaction matrix for the Heisenberg
generator and prove a path expansion for one and two marked insertions.  It
will derive an explicit convolution constant for
\eqref{eq:e8v7-LR2}.  If the physical generator is unbounded, the theorem
will remain a finite-regulator loader and the form-domain passage will stay
open.
\end{programme}

\begin{firewall}[Claim boundary after eighth-series V7]
\label{fire:e8v7-boundary}
V7 proves the exact differentiated generator comparison, its global
diamond bound, the complete adjacent expectation-derivative estimate, and
a remote-shell stress estimate under an explicit two-source propagation
hypothesis.

V7 does not prove that two-source hypothesis for the SM generator, control
unbounded interactions, sum the boundary or near rows, establish the
Einstein residual, remove the cutoff, or construct the Standard
Model--Einstein continuum.
\end{firewall}

\begin{theorem}[Eighth-series V7 result]
\label{thm:e8v7-result}
The adjacent channel stress operator is the sum of one direct metric
refinement insertion and two ordered metric/refinement Duhamel integrals,
with global bound \eqref{eq:e8v7-global-Kq}.  If the two marked insertions
obey \eqref{eq:e8v7-LR2}, a remote refinement shell retains the factor
\(e^{-\mu(R-vT)}\), and summability of the resulting far, near, and boundary
rows yields finite-horizon \(C^1\) carrier landing.
\end{theorem}

\begin{proof}
Combine \Cref{thm:e8v7-Kq,thm:e8v7-global,
thm:e8v7-remote,cor:e8v7-stress-row}.
\end{proof}

\paragraph{Parent-version record.}
V7 was formed from
\texttt{Renewal\_SM\_Einstein\_Continuum\_Research\_Notes\_V6.tex}.
The V6 parent had \(2\,207\) logical source lines, \(84\,566\) bytes, and
SHA--256
\[
 \texttt{04D309803BFCB6D52E996AC3B569E771D50E45260AE6EC0D70616769056E35B6}.
\]
The next compulsory companion audit is V10.

% =============================================================================
% END APPENDED EIGHTH-SERIES V7 MATERIAL
% =============================================================================
% =============================================================================
% BEGIN APPENDED EIGHTH-SERIES V8 MATERIAL
% =============================================================================

\section{Marked Dyson paths and two-source locality}
\label{sec:v8v8}

V7 isolates the two-source propagation estimate needed by the adjacent
stress row.  This appendix derives it from an abstract local oscillation
packet.  The proof is finite-regulator and applies to bounded Heisenberg
generators.  Its constants are volume independent when the interaction
matrix and local support sizes are uniform.

\subsection{Oscillation propagation matrix}
\label{sec:e8v8-oscillation}

Let \(I\) be a finite cell graph and let
\(\delta_i(A)\geq0\) be local observable seminorms.  Assume a Heisenberg
semigroup \(\mathcal S^*(t)\) satisfies
\begin{equation}
 \delta_i(\mathcal S^*(t)A)
 \leq\sum_j(e^{tK})_{ij}\delta_j(A),
 \label{eq:e8v8-propagation}
\end{equation}
where \(K=(K_{ij})\) is a nonnegative matrix.  For \(\mu>0\), define
\begin{equation}
 k_\mu:=\sup_i\sum_j e^{\mu d(i,j)}K_{ij}<\infty.
 \label{eq:e8v8-kmu}
\end{equation}

\begin{lemma}[Exponential kernel of the influence semigroup]
\label{lem:e8v8-kernel}
For every \(i,j\) and \(t\geq0\),
\begin{equation}
 (e^{tK})_{ij}
 \leq e^{k_\mu t}e^{-\mu d(i,j)}.
 \label{eq:e8v8-kernel}
\end{equation}
\end{lemma}

\begin{proof}
For a matrix \(M\), put
\(\|M\|_\mu=\sup_i\sum_j e^{\mu d(i,j)}|M_{ij}|\).
The graph triangle inequality gives
\(\|MN\|_\mu\leq\|M\|_\mu\|N\|_\mu\).  Hence
\[
 \|e^{tK}\|_\mu
 \leq\sum_{n\geq0}{t^n\over n!}\|K\|_\mu^n
 =e^{k_\mu t}.
\]
Every nonnegative entry is at most its weighted row sum multiplied by
\(e^{-\mu d(i,j)}\), proving \eqref{eq:e8v8-kernel}.
\end{proof}

Assume that if \(A_B\) is supported in \(B\), then
\begin{equation}
 \delta_i(A_B)=0\quad(i\notin B),
 \qquad
 \sum_{i\in B}\delta_i(A_B)
 \leq c_0|B|\|A_B\|.
 \label{eq:e8v8-local-A}
\end{equation}
Fix an insertion radius \(r_0\) and write
\(X^+=\{i:d(i,X)\leq r_0\}\).  Suppose every local insertion
\(\mathcal J_X^*\) has a strength \(j_X\geq0\) and obeys
\begin{align}
 \delta_i(\mathcal J_X^*A)&=0\quad(i\notin X^+),
 \label{eq:e8v8-insertion-support}\\
 \sum_{i\in X^+}\delta_i(\mathcal J_X^*A)
 &\leq j_X\sum_{i\in X^+}\delta_i(A),
 \label{eq:e8v8-insertion-osc}\\
 \|\mathcal J_X^*A\|
 &\leq j_X\sum_{i\in X^+}\delta_i(A).
 \label{eq:e8v8-insertion-readout}
\end{align}
These hypotheses are satisfied by many commutator or conditional
oscillation seminorms for a bounded finite-range insertion, after enlarging
\(r_0,c_0,j_X\).  They must be verified for the actual generator terms.

Assume a uniform local cardinality bound
\begin{equation}
 |X^+|\leq q_0
 \label{eq:e8v8-cardinality}
\end{equation}
for every insertion support under consideration.

\subsection{One and two marked insertions}
\label{sec:e8v8-marked}

\begin{lemma}[One marked insertion]
\label{lem:e8v8-one}
For a local observable \(A_B\),
\begin{equation}
 \|\mathcal J_X^*\mathcal S^*(t)A_B\|
 \leq c_0q_0|B|j_X\|A_B\|
 e^{k_\mu t}e^{-\mu(d(X,B)-r_0)}.
 \label{eq:e8v8-one}
\end{equation}
\end{lemma}

\begin{proof}
Use \eqref{eq:e8v8-insertion-readout}, then
\eqref{eq:e8v8-propagation} and \eqref{eq:e8v8-kernel}:
\begin{align*}
 \|\mathcal J_X^*\mathcal S^*(t)A_B\|
 &\leq j_X\sum_{i\in X^+}\sum_{b\in B}
 (e^{tK})_{ib}\delta_b(A_B)\\
 &\leq j_Xq_0 e^{k_\mu t}
 e^{-\mu d(X^+,B)}
 \sum_{b\in B}\delta_b(A_B).
\end{align*}
Now \(d(X^+,B)\geq d(X,B)-r_0\), and
\eqref{eq:e8v8-local-A} completes the proof.
\end{proof}

For three sets define the two ordered path lengths
\begin{equation}
 \ell_{X\leftarrow Y\leftarrow B}
 :=d(X,Y)+d(Y,B),
 \qquad
 \ell_{Y\leftarrow X\leftarrow B}
 :=d(Y,X)+d(X,B).
 \label{eq:e8v8-ordered-lengths}
\end{equation}
Each is at least the minimum tree length
\(\mathfrak t(B,X,Y)\) used in V7.

\begin{theorem}[Two marked insertion bound]
\label{thm:e8v8-two}
Let \(\mathcal J_X^*,\mathcal J_Y^*\) satisfy
\eqref{eq:e8v8-insertion-support}--
\eqref{eq:e8v8-cardinality}.  Then
\begin{align}
 &\|\mathcal J_X^*\mathcal S^*(t_1)
       \mathcal J_Y^*\mathcal S^*(t_2)A_B\|
 \notag\\
 &\qquad\leq
 c_0q_0^2|B|j_Xj_Y\|A_B\|
 e^{k_\mu(t_1+t_2)}
 e^{-\mu(\ell_{X\leftarrow Y\leftarrow B}-3r_0)}.
 \label{eq:e8v8-two-ordered}
\end{align}
The same estimate with \(X,Y\) exchanged holds for the reverse order.
Consequently their sum is bounded by
\begin{equation}
 2c_0q_0^2e^{3\mu r_0}|B|j_Xj_Y\|A_B\|
 e^{k_\mu(t_1+t_2)}e^{-\mu\mathfrak t(B,X,Y)}.
 \label{eq:e8v8-two-tree}
\end{equation}
\end{theorem}

\begin{proof}
First propagate \(A_B\) for time \(t_2\) and restrict its oscillation
vector to \(Y^+\).  As in the proof of \Cref{lem:e8v8-one},
\begin{equation}
 \sum_{y\in Y^+}\delta_y(\mathcal S^*(t_2)A_B)
 \leq c_0q_0|B|\|A_B\|e^{k_\mu t_2}
 e^{-\mu(d(Y,B)-r_0)}.
 \label{eq:e8v8-first-leg}
\end{equation}
The insertion \(\mathcal J_Y^*\) multiplies this bound by \(j_Y\) and,
by \eqref{eq:e8v8-insertion-support}, leaves an oscillation vector
supported in \(Y^+\).  Propagate it for time \(t_1\).  Summing its values
over \(X^+\), using \eqref{eq:e8v8-kernel}, costs at most
\begin{equation}
 q_0e^{k_\mu t_1}e^{-\mu(d(X,Y)-2r_0)}.
 \label{eq:e8v8-second-leg}
\end{equation}
Finally \eqref{eq:e8v8-insertion-readout} for
\(\mathcal J_X^*\) contributes \(j_X\).  Multiplying
\eqref{eq:e8v8-first-leg} and \eqref{eq:e8v8-second-leg} proves
\eqref{eq:e8v8-two-ordered}.  Exchange \(X,Y\) for the reverse order.
Since each ordered length is at least \(\mathfrak t(B,X,Y)\), enlarge the
constant by \(e^{3\mu r_0}\) and sum the two orders to obtain
\eqref{eq:e8v8-two-tree}.
\end{proof}

\begin{remark}[Marked path interpretation]
\label{rem:e8v8-path}
The exponent in \eqref{eq:e8v8-two-ordered} is the length of the ordered
path from the observable support to the inner insertion and then to the
outer insertion.  Replacing it by the smaller spanning-tree length loses
order information but gives the symmetric V7 hypothesis.
\end{remark}

\subsection{Population of the V7 locality hypothesis}
\label{sec:e8v8-population}

Suppose both Heisenberg semigroups \(\mathcal S_0^*,\mathcal S_1^*\) in V7
obey \eqref{eq:e8v8-propagation} with a common matrix \(K\).  Suppose each
\(\mathcal V_X^*\) and \(\mathcal W_{j,Y}^*\) obeys the insertion packet,
with
\begin{equation}
 j_X\leq C_V\|\mathcal V_X\|_\diamond,
 \qquad
 j_{j,Y}\leq C_W\|\mathcal W_{j,Y}\|_\diamond.
 \label{eq:e8v8-strength}
\end{equation}

\begin{corollary}[Explicit V7 two-source constant]
\label{cor:e8v8-LR2}
Under these hypotheses, V7's two-source estimate holds with velocity
\begin{equation}
 v_\mu={k_\mu\over\mu}
 \label{eq:e8v8-velocity}
\end{equation}
and constant
\begin{equation}
 C_2=2c_0q_0^2C_VC_We^{3\mu r_0}.
 \label{eq:e8v8-C2}
\end{equation}
The V6 one-source estimate holds with
\(C_{\rm LR}=c_0q_0C_Ve^{\mu r_0}\) and the same velocity.
\end{corollary}

\begin{proof}
Insert \eqref{eq:e8v8-strength} into
\eqref{eq:e8v8-two-tree}.  Since
\(e^{k_\mu(t_1+t_2)}=
 e^{\mu v_\mu(t_1+t_2)}\), the result is exactly the V7 form.  The
one-source statement follows similarly from \Cref{lem:e8v8-one}.
\end{proof}

\begin{theorem}[Remote-shell stress row from the oscillation packet]
\label{thm:e8v8-stress}
Assume the V6--V7 bounded-generator decompositions and the common
oscillation packet above.  If the remote refinement and differentiated
refinement component sums obey
\begin{align}
 \sum_n T_n e^{-\mu(R_n-v_\mu T_n)}
 \sum_X\|\mathcal V_{q,n,X}\|_\diamond&<\infty,
 \label{eq:e8v8-direct-sum}\\
 \sum_n T_n^2 e^{-\mu(R_n-v_\mu T_n)}
 (M_{1,q,n}+M_{0,q,n})
 \sum_X\|\mathcal V_{n,X}\|_\diamond&<\infty,
 \label{eq:e8v8-ordered-sum}
\end{align}
then the far part of the V7 adjacent stress row is summable for every fixed
protected local observable.  The near and boundary rows remain separate.
\end{theorem}

\begin{proof}
\Cref{cor:e8v8-LR2} supplies the one- and two-source hypotheses of
\Cref{thm:e8v7-remote}.  Bound every remote distance below by \(R_n\) and
compare its two terms with
\eqref{eq:e8v8-direct-sum}--\eqref{eq:e8v8-ordered-sum}.
\end{proof}

\subsection{Deriving the propagation matrix from a generator inequality}
\label{sec:e8v8-generator}

\begin{proposition}[Dini comparison implies
\eqref{eq:e8v8-propagation}]
\label{prop:e8v8-Dini}
Suppose for every observable \(A\),
\begin{equation}
 D^+\delta_i(\mathcal S^*(t)A)
 \leq\sum_jK_{ij}\delta_j(\mathcal S^*(t)A),
 \qquad K_{ij}\geq0.
 \label{eq:e8v8-Dini}
\end{equation}
Then \eqref{eq:e8v8-propagation} holds.
\end{proposition}

\begin{proof}
Let \(z_i(t)=\delta_i(\mathcal S^*(t)A)\).  For
\(\epsilon>0\), the vector
\(y_\epsilon(t)=e^{tK}(z(0)+\epsilon\mathbf1)\) is strictly positive and
solves \(y_\epsilon'=Ky_\epsilon\).  If some component of
\(z-y_\epsilon\) first reached zero from below, the upper Dini inequality
and nonnegativity of \(K\) off that component would make its upper
derivative nonpositive, preventing a crossing.  Thus
\(z(t)\leq y_\epsilon(t)\).  Let \(\epsilon\downarrow0\) to obtain
\eqref{eq:e8v8-propagation}.
\end{proof}

\begin{firewall}[Finite regulator and bounded insertions]
\label{fire:e8v8-bounded}
The oscillation proof requires bounded local insertion maps and a finite
weighted influence constant \(k_\mu\).  It does not provide these for an
unbounded continuum Dirac, electric-field, or gravitational generator.
Nor does it prove that the local seminorms have the Ward-quotient dual
property.  Those are independent physical and functional-analytic rows.
\end{firewall}

\begin{assessment}[Progress at V8]
\label{ass:e8v8-status}
The two-source hypothesis introduced at V7 is no longer primitive.  It
follows from a common oscillation propagation matrix and three elementary
local-insertion properties.  The far adjacent stress row now has explicit
series \eqref{eq:e8v8-direct-sum}--
\eqref{eq:e8v8-ordered-sum}.  The remaining finite-regulator bottleneck is
to construct the oscillation seminorm and compute its influence matrix for
the actual gauge--Higgs--fermion carrier generator.
\end{assessment}

\begin{programme}[Eighth-series V9: tensor-sum generator assembly]
\label{prog:e8v8-V9}
V9 will assemble sector generators and their interaction terms in one
oscillation matrix.  It will prove a block-matrix weighted-norm criterion,
identify how gauge Ward, Higgs massive, fermion carrier, and metric contact
rows enter, and quantify when cross-sector couplings preserve the required
finite-horizon propagation even though they need not produce a global
contraction gap.
\end{programme}

\begin{firewall}[Claim boundary after eighth-series V8]
\label{fire:e8v8-boundary}
V8 proves exponential propagation from a weighted influence matrix, one-
and two-marked insertion locality, an explicit V7 two-source constant, and
summability criteria for the far adjacent stress row.

V8 does not construct the SM influence matrix, prove uniform constants for
unbounded fields, verify the Ward quotient duality, sum near or boundary
rows, remove the cutoff, or construct the Standard Model--Einstein
continuum.
\end{firewall}

\begin{theorem}[Eighth-series V8 result]
\label{thm:e8v8-result}
A nonnegative local influence matrix with finite exponential row norm,
together with bounded insertion support, oscillation, and readout bounds,
implies the ordered two-source decay
\eqref{eq:e8v8-two-ordered}.  It therefore supplies the missing V7
locality hypothesis and reduces far stress projectivity to the explicit
summability rows \eqref{eq:e8v8-direct-sum}--
\eqref{eq:e8v8-ordered-sum}.
\end{theorem}

\begin{proof}
Combine \Cref{lem:e8v8-kernel,lem:e8v8-one,
thm:e8v8-two,cor:e8v8-LR2,thm:e8v8-stress,
prop:e8v8-Dini}.
\end{proof}

\paragraph{Parent-version record.}
V8 was formed from
\texttt{Renewal\_SM\_Einstein\_Continuum\_Research\_Notes\_V7.tex}.
The V7 parent had \(2\,537\) logical source lines, \(96\,410\) bytes, and
SHA--256
\[
 \texttt{EC93B11558E7300F728B76B976ABBACBBD2A0386B081D87607676EF4E0652362}.
\]
The next compulsory companion audit is V10.

% =============================================================================
% END APPENDED EIGHTH-SERIES V8 MATERIAL
% =============================================================================

% =============================================================================
% BEGIN APPENDED EIGHTH-SERIES V9 MATERIAL
% =============================================================================

\section{Sectorwise influence assembly and a small-gain criterion}
\label{sec:v8v9}

V8 reduces one- and two-source locality to a nonnegative influence matrix.
The Standard Model carrier has several coupled sectors.  This appendix
assembles their influence matrices, proves a weighted small-gain criterion,
and separates finite-horizon propagation from strict contraction.  The
separation is essential: positive majorants for massive Higgs or fermion
rows cannot make a Ward-critical gauge row strictly negative.

\subsection{Combined sector--cell norm}
\label{sec:e8v9-combined}

Let sectors be indexed by \(a\in\{1,\ldots,r\}\) and cells by \(i\in I\).
For a nonnegative vector \(z=(z_{a,i})\), positive sector weights
\(\sigma_a\), a localization set \(Z\), and \(\mu>0\), define
\begin{equation}
 \|z\|_{\mu,Z;\sigma}
 :=\max_{a,i}\sigma_a e^{\mu d(i,Z)}z_{a,i}.
 \label{eq:e8v9-vector-norm}
\end{equation}
For a block matrix \(K=(K^{ab}_{ij})\), put
\begin{equation}
 \|K\|_{\mu;\sigma}
 :=\sup_{a,i}\sum_{b,j}{\sigma_a\over\sigma_b}
 e^{\mu d(i,j)}|K^{ab}_{ij}|.
 \label{eq:e8v9-matrix-norm}
\end{equation}

\begin{lemma}[Sector-weighted submultiplicativity]
\label{lem:e8v9-submultiplicative}
The norm \eqref{eq:e8v9-matrix-norm} is submultiplicative and
\begin{equation}
 \|Kz\|_{\mu,Z;\sigma}
 \leq\|K\|_{\mu;\sigma}\|z\|_{\mu,Z;\sigma}.
 \label{eq:e8v9-action}
\end{equation}
If \(K\geq0\), then
\begin{equation}
 (e^{tK})^{ab}_{ij}
 \leq {\sigma_b\over\sigma_a}
 e^{t\|K\|_{\mu;\sigma}}e^{-\mu d(i,j)}.
 \label{eq:e8v9-kernel}
\end{equation}
\end{lemma}

\begin{proof}
The factors \(\sigma_a/\sigma_b\) telescope under block multiplication,
and
\(d(i,j)\leq d(i,k)+d(k,j)\) supplies the spatial weight.  This proves
submultiplicativity.  For the vector action, use
\(e^{\mu d(i,Z)}\leq e^{\mu d(i,j)}e^{\mu d(j,Z)}\), multiply by
\(\sigma_a\), and take the maximum.  Expanding the exponential and using
submultiplicativity gives
\(\|e^{tK}\|_{\mu;\sigma}\leq e^{t\|K\|_{\mu;\sigma}}\).  A single
nonnegative entry is bounded by its weighted row sum, giving
\eqref{eq:e8v9-kernel}.
\end{proof}

Suppose sectorwise observable oscillations
\(\delta_{a,i}(A)\) satisfy a Dini comparison
\begin{equation}
 D^+\delta_{a,i}(\mathcal S^*(t)A)
 \leq\sum_{b,j}K^{ab}_{ij}
       \delta_{b,j}(\mathcal S^*(t)A),
 \qquad K^{ab}_{ij}\geq0.
 \label{eq:e8v9-Dini}
\end{equation}

\begin{theorem}[Coupled-sector finite-horizon locality]
\label{thm:e8v9-locality}
If \(\|K\|_{\mu;\sigma}<\infty\), then
\begin{equation}
 \|\delta(\mathcal S^*(t)A)\|_{\mu,Z;\sigma}
 \leq e^{t\|K\|_{\mu;\sigma}}
 \|\delta(A)\|_{\mu,Z;\sigma}.
 \label{eq:e8v9-propagation}
\end{equation}
Consequently the V8 one- and two-marked insertion theorems hold on the
combined sector--cell graph, with velocity
\begin{equation}
 v_{\mu,\sigma}:={\|K\|_{\mu;\sigma}\over\mu},
 \label{eq:e8v9-velocity}
\end{equation}
provided their local insertion support, oscillation, and readout hypotheses
hold sectorwise with uniform constants.
\end{theorem}

\begin{proof}
The finite-dimensional comparison argument of V8 gives
\(\delta(\mathcal S^*(t)A)\leq e^{tK}\delta(A)\).  Apply
\Cref{lem:e8v9-submultiplicative}.  Its pointwise kernel bound is the only
propagation input used in the V8 marked-insertion proof.  Replacing the
cell index by the pair \((a,i)\) and retaining the finite uniform support
cardinality reproduces that proof with velocity
\eqref{eq:e8v9-velocity}.
\end{proof}

\subsection{Block row estimates}
\label{sec:e8v9-blocks}

Define block weighted row constants
\begin{equation}
 k_{ab}:=sup_i\sum_j e^{\mu d(i,j)}K^{ab}_{ij}.
 \label{eq:e8v9-kab}
\end{equation}
Then
\begin{equation}
 \|K\|_{\mu;\sigma}
 \leq\max_a\sum_b{\sigma_a\over\sigma_b}k_{ab}.
 \label{eq:e8v9-block-row}
\end{equation}

\begin{corollary}[Finite cross-sector couplings preserve locality]
\label{cor:e8v9-finite}
If every \(k_{ab}\) is bounded uniformly in cutoff and volume and the
number of sectors is fixed, then any fixed positive sector weights give a
uniform finite-horizon propagation speed.  No negative contraction margin
is needed for this conclusion.
\end{corollary}

\begin{proof}
The right side of \eqref{eq:e8v9-block-row} is a finite sum of uniform
constants.  Apply \Cref{thm:e8v9-locality}.
\end{proof}

For strict contraction, let \(Q=(Q^{ab}_{ij})\) be a Metzler comparison
matrix: off-diagonal entries in the combined \((a,i)\) index are
nonnegative.  Suppose the within-sector weighted logarithmic rows obey
\begin{equation}
 \sup_i\left(Q^{aa}_{ii}
 +\sum_{j\ne i}e^{\mu d(i,j)}Q^{aa}_{ij}\right)
 \leq-m_a
 \label{eq:e8v9-diagonal-margin}
\end{equation}
with \(m_a\geq0\), and define cross-sector bounds
\begin{equation}
 c_{ab}:=sup_i\sum_j e^{\mu d(i,j)}Q^{ab}_{ij},
 \qquad a\ne b.
 \label{eq:e8v9-cross}
\end{equation}

\begin{theorem}[Sector small-gain contraction]
\label{thm:e8v9-small-gain}
The combined weighted logarithmic norm satisfies
\begin{equation}
 \ell_{\mu;\sigma}(Q)
 \leq\max_a\left(
 -m_a+\sum_{b\ne a}{\sigma_a\over\sigma_b}c_{ab}
 \right).
 \label{eq:e8v9-log-bound}
\end{equation}
Hence a sufficient condition for contraction rate \(m_*>0\) is
\begin{equation}
 \sum_{b\ne a}{\sigma_a\over\sigma_b}c_{ab}
 \leq m_a-m_*
 \qquad\text{for every }a.
 \label{eq:e8v9-small-gain}
\end{equation}
Under this condition, the combined channel semigroup contracts at rate
\(m_*\) in the sector-weighted local norm.
\end{theorem}

\begin{proof}
Conjugate the combined comparison matrix by the diagonal sector and spatial
weights used in \eqref{eq:e8v9-vector-norm}.  The logarithmic norm in the
sup norm is the largest diagonal entry plus the absolute off-diagonal row
sum.  Equation \eqref{eq:e8v9-diagonal-margin} bounds the within-sector
part, while \eqref{eq:e8v9-cross} contributes
\((\sigma_a/\sigma_b)c_{ab}\).  This proves
\eqref{eq:e8v9-log-bound}.  Under \eqref{eq:e8v9-small-gain}, the bound is
at most \(-m_*\); the V2 Metzler theorem gives contraction
\(e^{-m_*t}\).
\end{proof}

\begin{corollary}[Spectral-radius form of small gain]
\label{cor:e8v9-spectral-radius}
Assume every \(m_a>0\).  Let \(C=(c_{ab})\) with \(c_{aa}=0\) and
\(M=\operatorname{diag}(m_1,\ldots,m_r)\).  If
\begin{equation}
 \rho(M^{-1}C)<1,
 \label{eq:e8v9-rho}
\end{equation}
then there are positive sector weights \(\sigma_a\) for which
\eqref{eq:e8v9-small-gain} holds with some \(m_*>0\).
\end{corollary}

\begin{proof}
Let \(A=M^{-1}C\).  Since \(A\geq0\) and \(\rho(A)<1\),
\begin{equation}
 x:=(I-A)^{-1}\mathbf1=\sum_{n\geq0}A^n\mathbf1
 \label{eq:e8v9-x}
\end{equation}
is finite and strictly positive, and \(Ax=x-\mathbf1<x\)
componentwise.  Choose \(\sigma_a=x_a^{-1}\).  Then
\[
 \sum_{b\ne a}{\sigma_a\over\sigma_b}c_{ab}
 ={(Cx)_a\over x_a}<m_a.
\]
The finite minimum of the positive row slacks is a valid \(m_*\).
\end{proof}

\subsection{The Ward row cannot borrow a positive margin}
\label{sec:e8v9-Ward}

\begin{theorem}[Positive-majorant obstruction in a massless sector]
\label{thm:e8v9-Ward}
If some sector \(g\) has \(m_g=0\), then no positive sector weights can
make the sufficient small-gain bound
\eqref{eq:e8v9-small-gain} hold with \(m_*>0\).  Indeed,
\begin{equation}
 -m_g+\sum_{b\ne g}{\sigma_g\over\sigma_b}c_{gb}
 \geq0.
 \label{eq:e8v9-Ward-row}
\end{equation}
Thus positive absolute influence from massive sectors cannot create a
strict contraction margin in a Ward-critical row.
\end{theorem}

\begin{proof}
Every \(c_{gb}\) and every weight ratio is nonnegative, while \(m_g=0\).
The row bound is therefore nonnegative and cannot be at most
\(-m_*<0\).
\end{proof}

\begin{remark}[Scope of the obstruction]
\label{rem:e8v9-scope}
The theorem concerns a positive absolute majorant.  Signed cancellations,
a Ward quotient with observable duality, or a reflected physical-time form
can yield coercivity not visible in this matrix.  Such a mechanism must be
proved in its own norm; it cannot be inferred from the massive diagonal
rows.
\end{remark}

\subsection{SM sector interpretation}
\label{sec:e8v9-SM}

Use sectors \(g,h,f,m\) for gauge, Higgs, fermion-carrier, and metric/contact
oscillations.  The block constants have the following meanings:
\begin{equation}
 (k_{ab})_{a,b\in\{g,h,f,m\}}=
 \begin{pmatrix}
 k_{gg}&k_{gh}&k_{gf}&k_{gm}\\
 k_{hg}&k_{hh}&k_{hf}&k_{hm}\\
 k_{fg}&k_{fh}&k_{ff}&k_{fm}\\
 k_{mg}&k_{mh}&k_{mf}&k_{mm}
 \end{pmatrix}.
 \label{eq:e8v9-SM-matrix}
\end{equation}
The entries are weighted sums of actual local generator derivatives, not
coupling names.  Yukawa, covariant hopping, gauge--Higgs, and metric contact
terms contribute to off-diagonal blocks only after their oscillation
inequalities are proved.

\begin{corollary}[Conditional coupled SM finite-horizon row]
\label{cor:e8v9-SM}
If every entry of \eqref{eq:e8v9-SM-matrix} and the V8 insertion constants
is uniformly finite on the protected chart, then the coupled finite
carrier has volume-independent one- and two-source propagation for every
fixed physical horizon.  This conclusion remains valid with
\(m_g=0\).  A strict global contraction conclusion does not.
\end{corollary}

\begin{proof}
Finite-horizon propagation is \Cref{cor:e8v9-finite}; it only uses finiteness
of the nonnegative block rows.  The absence of a strict absolute
contraction conclusion when \(m_g=0\) is \Cref{thm:e8v9-Ward}.
\end{proof}

\begin{definition}[Sector influence assembly card, SIAC]
\label{def:e8v9-SIAC}
SIAC records:
\begin{enumerate}
\item one local oscillation family for each sector and a combined
observable readout inequality;
\item the within- and cross-sector generator influence kernels;
\item uniform exponential row sums \(k_{ab}\);
\item local support, oscillation, and readout bounds for refinement and
metric insertions;
\item an explicit Ward quotient or long-range row for the massless sector;
and
\item cofinal near, far, contact, projector, leakage, and recovery sums.
\end{enumerate}
\end{definition}

\begin{assessment}[Progress at V9]
\label{ass:e8v9-status}
The tensor-sector assembly is now algebraically closed.  Uniformly finite
cross rows suffice for the V8 finite-horizon locality needed by carrier
projectivity.  They do not supply a global gap in the photon row.  The
physical bottleneck is to compute the entries of
\eqref{eq:e8v9-SM-matrix} on one regulator and to prove a combined local
observable readout compatible with the Ward treatment.
\end{assessment}

\begin{programme}[Eighth-series V10: companion audit and SIAC reduction]
\label{prog:e8v9-V10}
V10 will perform the compulsory companion audit, then reduce the SIAC
matrix using exact tensor-product and bounded-interaction identities.  It
will identify which cross blocks vanish for a factorized reference carrier,
which are generated by the old/detail interaction, and which must be paid
by the refinement amplitude.
\end{programme}

\begin{firewall}[Claim boundary after eighth-series V9]
\label{fire:e8v9-boundary}
V9 proves the sector-weighted influence norm, coupled finite-horizon
locality, the block small-gain theorem and its spectral-radius form, and the
positive-majorant obstruction for a massless Ward row.

V9 does not compute the physical SM block entries, prove their cutoff
bounds, construct the Ward readout, sum the SIAC ledger, remove the cutoff,
or construct the Standard Model--Einstein continuum.
\end{firewall}

\begin{theorem}[Eighth-series V9 result]
\label{thm:e8v9-result}
A fixed finite matrix of uniformly bounded sector-to-sector exponential
influences yields volume-independent finite-horizon one- and two-source
locality.  Strict contraction follows from the weighted small-gain
condition \eqref{eq:e8v9-small-gain}, equivalently from
\(\rho(M^{-1}C)<1\) when all diagonal margins are positive.  If a Ward
sector has zero diagonal margin, nonnegative absolute cross influences
cannot make that row strictly contractive, although finite-horizon locality
remains available.
\end{theorem}

\begin{proof}
Combine \Cref{lem:e8v9-submultiplicative,
thm:e8v9-locality,thm:e8v9-small-gain,
cor:e8v9-spectral-radius,thm:e8v9-Ward,cor:e8v9-SM}.
\end{proof}

\paragraph{Parent-version record.}
V9 was formed from
\texttt{Renewal\_SM\_Einstein\_Continuum\_Research\_Notes\_V8.tex}.
The V8 parent had \(2\,886\) logical source lines, \(108\,637\) bytes, and
SHA--256
\[
 \texttt{51BC5704B1DF51863D3695275F2BF58A74FF4CDAD2887451469A85858D3A292C}.
\]
The next compulsory companion audit is V10.

% =============================================================================
% END APPENDED EIGHTH-SERIES V9 MATERIAL
% =============================================================================
% =============================================================================
% BEGIN APPENDED EIGHTH-SERIES V10 MATERIAL
% =============================================================================

\section{V10 companion audit and factorized-sector reduction}
\label{sec:v8v10}

This compulsory audit is followed by the first reduction of the V9 sector
matrix.  A tensor-product reference generator has no off-diagonal sector
influence.  Every cross-sector row is therefore owned by the interaction
generator.  If that interaction and its metric derivative carry the
refinement amplitude, the far base and stress rows inherit it through the
V8 marked-path estimates.

\subsection{Compulsory companion audit}
\label{sec:e8v10-audit}

The newest coherent companion sources inspected are
\begin{align}
 &\texttt{Renewal\_Yang\_Mills\_Mass\_Gap\_Research\_Notes\_V61.tex},
 \notag\\
 &\qquad 27\,346\ \text{logical lines},\quad962\,682\ \text{bytes},
 \notag\\
 &\qquad
 \texttt{40CA8E0F9AAD41EB60E65C07DF231F9450211C4A524FE958542F005767BADA58},
 \label{eq:e8v10-YM-hash}\\
 &\texttt{Renewal\_NS\_Stability\_Research\_Notes\_V26.tex},
 \notag\\
 &\qquad 5\,319\ \text{logical lines},\quad278\,283\ \text{bytes},
 \notag\\
 &\qquad
 \texttt{82C887F8C2C69E4F28FAD7C3DC8DE5B9099CB5A4BF431EBA1FFF3E91935978AD}.
 \label{eq:e8v10-NS-hash}
\end{align}
Each has one live document terminator and ends with
\verb|\end{document}|.  No companion file was modified.

YM V56--V61 separates local finite-horizon projectivity from a global
transfer gap, develops two-level block and Schur-form coercivity criteria,
and then constructs exact bounded gauge-invariant, centre-neutral bilocal
plaquette networks.  Their quadratic jets contain colour-singlet
opposite-momentum two-particle packets with Gaussian energy
\(O(R^{-1})\).  Thus quotienting the global colour mode does not remove an
infrared singlet sequence; a nonperturbative signed nonlinear lift would be
needed for a pure-Yang--Mills mass proof.  That lift is not proved.

For the SM programme, the transferable conclusion is narrower.  Local
state projectivity may use finite-horizon locality without waiting for a
global gap.  A global gap must not be inserted into the unbroken long-range
gauge sector merely because local or coloured modes have been quotiented.
The YM nonlinear lift target is not imported as a photon mass mechanism.
NS V26 remains unchanged and supplies no sector influence or finite-horizon
carrier estimate.

\begin{audit}[V10 companion decision]
\label{aud:e8v10-decision}
V10 retains finite-horizon propagation as the state-existence engine and
keeps every global mass or transfer gap outside SIAC.  The exact
factorized/interacting reduction below imports no YM or NS numerical
constant.  Any later mass claim must use the appropriate physical reflected
transfer and cannot be deduced from the SIAC locality matrix.
\end{audit}

\begin{proof}
The decision follows from the explicit separation theorem and singlet
infrared sequence in the audited YM source.  Those results concern its
physical reflected transfer, whereas SIAC is a finite-horizon local
propagation packet.  The unchanged NS source has no map to either object.
\end{proof}

\subsection{Tensor-product reference generator}
\label{sec:e8v10-factorized}

Let
\begin{equation}
 \mathfrak A=\bigotimes_{a=1}^r\mathfrak A^{(a)},
 \qquad
 \mathcal L_{\rm fac}
 :=\sum_{a=1}^r
 I^{(1)}\otimes\cdots\otimes\mathcal L^{(a)}
 \otimes\cdots\otimes I^{(r)}.
 \label{eq:e8v10-Lfac}
\end{equation}
The summands commute, so
\begin{equation}
 e^{t\mathcal L_{\rm fac}}
 =\bigotimes_{a=1}^r e^{t\mathcal L^{(a)}}.
 \label{eq:e8v10-product-semigroup}
\end{equation}
Assume sector oscillations \(\delta_{a,i}\) have the tensor stability
property: a channel acting only in sector \(b\ne a\) does not increase
\(\delta_{a,i}\), while the sector-\(a\) channel propagates its own row by
\(K^{(a)}\).

\begin{theorem}[Exact block diagonality of the reference influence]
\label{thm:e8v10-block-diagonal}
Under tensor stability, the product semigroup
\eqref{eq:e8v10-product-semigroup} obeys the V9 Dini comparison with
\begin{equation}
 K_{\rm fac}=\operatorname{diag}
 (K^{(1)},\ldots,K^{(r)}).
 \label{eq:e8v10-Kfac}
\end{equation}
In particular,
\begin{equation}
 (K_{\rm fac})^{ab}=0\qquad(a\ne b).
 \label{eq:e8v10-cross-zero}
\end{equation}
\end{theorem}

\begin{proof}
The commuting exponential factorization lets the sector channels be
applied in any order.  All factors with \(b\ne a\) leave the
sector-\(a\) oscillation no larger by tensor stability.  The remaining
factor propagates it by \(e^{tK^{(a)}}\).  Thus no oscillation initially in
another sector is needed to bound row \(a\), and the combined comparison
matrix is the direct block sum \eqref{eq:e8v10-Kfac}.
\end{proof}

\begin{remark}[Tensor stability must match the chosen oscillation]
\label{rem:e8v10-tensor-stability}
Block diagonality is exact for the generator, but its proof uses tensor
stability of the seminorms.  A seminorm defined through a correlated
reference state may fail that property even when the algebraic generator
factorizes.  SIAC must use either tensor-stable oscillations or charge the
resulting comparison error explicitly.
\end{remark}

\subsection{Interaction ownership of cross-sector rows}
\label{sec:e8v10-interaction}

Let the fine generator be
\begin{equation}
 \mathcal L=\mathcal L_{\rm fac}+\mathcal V,
 \qquad
 \mathcal V=\sum_X\mathcal V_X.
 \label{eq:e8v10-L}
\end{equation}
Suppose its oscillation Dini derivative is bounded by
\begin{equation}
 D^+\delta(e^{t\mathcal L^*}A)
 \leq(K_{\rm fac}+E)\delta(e^{t\mathcal L^*}A),
 \qquad E\geq0,
 \label{eq:e8v10-K}
\end{equation}
where \(E=(E^{ab})\) is obtained from the local insertion inequalities for
\(\mathcal V_X^*\).

\begin{proposition}[Cross-block ownership]
\label{prop:e8v10-cross-owner}
For \(a\ne b\), every cross-sector influence row of the fine comparison
matrix is bounded by the corresponding interaction row:
\begin{equation}
 K^{ab}\leq E^{ab}.
 \label{eq:e8v10-cross-owner}
\end{equation}
If
\begin{equation}
 \|E\|_{\mu;\sigma}\leq C_E\varepsilon,
 \label{eq:e8v10-E-small}
\end{equation}
then all cross-block exponential row sums are
\(O(\varepsilon)\), while each diagonal row differs from its factorized
reference by at most \(C_E\varepsilon\) in the combined norm.
\end{proposition}

\begin{proof}
The factorized cross blocks vanish by
\eqref{eq:e8v10-cross-zero}; the only nonnegative contribution in
\eqref{eq:e8v10-K} is therefore \(E^{ab}\).  A block row is bounded by the
full combined norm after multiplication by its fixed sector-weight ratio,
so \eqref{eq:e8v10-E-small} gives the stated estimates.
\end{proof}

\begin{theorem}[Finite-horizon influence perturbation]
\label{thm:e8v10-influence-perturbation}
Let
\(K_0=K_{\rm fac}\) and \(K_1=K_{\rm fac}+E\).  Put
\begin{equation}
 \kappa:=\max\{\|K_0\|_{\mu;\sigma},
                       \|K_1\|_{\mu;\sigma}\}.
 \label{eq:e8v10-kappa}
\end{equation}
Then
\begin{equation}
 \|e^{TK_1}-e^{TK_0}\|_{\mu;\sigma}
 \leq T e^{\kappa T}\|E\|_{\mu;\sigma}.
 \label{eq:e8v10-influence-Duhamel}
\end{equation}
Thus \eqref{eq:e8v10-E-small} gives an
\(O(T e^{\kappa T}\varepsilon)\) change of the finite-horizon propagation
matrix.
\end{theorem}

\begin{proof}
Matrix Duhamel gives
\[
 e^{TK_1}-e^{TK_0}
 =\int_0^T e^{(T-u)K_1}E e^{uK_0}\,du.
\]
Use submultiplicativity of the V9 weighted norm and
\(\|e^{tK_j}\|\leq e^{t\|K_j\|}\).  The integrand is at most
\(e^{\kappa(T-u)}\|E\|e^{\kappa u}\), and integration proves the result.
\end{proof}

\subsection{Metric derivative and refinement amplitude}
\label{sec:e8v10-metric}

Let the interaction depend on the metric chart and write
\(E_q\) for a nonnegative influence majorant of \(\mathcal V_q\).
Assume at adjacent scale \(n\)
\begin{equation}
 \|E_n\|_{\mu;\sigma}\leq C_E\varepsilon_n,
 \qquad
 \|E_{q,n}\|_{\mu;\sigma}\leq C_{E,q}\varepsilon_n
 \label{eq:e8v10-Eq-small}
\end{equation}
for unit metric tangents.  Let the factorized metric generator have bounded
combined insertion row \(M_q\).

\begin{theorem}[Interaction-tagged base and stress propagation]
\label{thm:e8v10-tagged}
For \(0\leq T\leq T_*\), assumptions
\eqref{eq:e8v10-Eq-small} imply:
\begin{enumerate}
\item the interaction correction to the propagation matrix is
\(O(\varepsilon_n)\);
\item the direct differentiated interaction insertion is
\(O(T\varepsilon_n)\); and
\item the two ordered interaction/metric insertion terms are
\(O(T^2M_q\varepsilon_n)\).
\end{enumerate}
If the interaction components are supported in a remote refinement shell,
the last two rows acquire the V8 factor
\(e^{-\mu(R_n-vT_*)}\).  Therefore their sums are finite whenever
\begin{equation}
 \sum_n\varepsilon_n e^{-\mu(R_n-vT_*)}<\infty
 \label{eq:e8v10-tagged-sum}
\end{equation}
and all uniform insertion/readout constants are populated.
\end{theorem}

\begin{proof}
Item 1 is \Cref{thm:e8v10-influence-perturbation} with bounded \(T\) and
\(\kappa\).  Item 2 is the direct interval in the V7 differentiated
Duhamel formula, using \(E_{q,n}=O(\varepsilon_n)\).  Item 3 consists of
the two triangular domains, each with one interaction insertion and one
factorized metric insertion; their total area is \(T^2\), so the bound is
\(O(T^2M_q\varepsilon_n)\).  V8 supplies the marked-path shell factor.
Comparison with \eqref{eq:e8v10-tagged-sum} gives summability.
\end{proof}

\begin{corollary}[Exact factorized reference reduction of SIAC]
\label{cor:e8v10-SIAC}
For a tensor-stable factorized reference, SIAC cross blocks need only be
proved for the interaction generator.  If
\eqref{eq:e8v10-Eq-small}, the V8 insertion packet, the remote-shell
support, and \eqref{eq:e8v10-tagged-sum} hold, then the generated
cross-sector parts of the finite-horizon base and stress rows are summable.
The diagonal factorized rows and all near, state, contact, projector,
leakage, and recovery rows remain independent.
\end{corollary}

\begin{proof}
Cross-block ownership is \Cref{prop:e8v10-cross-owner}; base and stress
summability are \Cref{thm:e8v10-tagged}.  None of those estimates contains
the listed independent terms.
\end{proof}

\begin{firewall}[Small interaction influence is not yet a physical
estimate]
\label{fire:e8v10-small}
The hypothesis \(\|E_n\|=O(\varepsilon_n)\) must be derived from the
actual old/detail Hamiltonian or channel generator after the factorized
reference has been fixed.  A small coupling constant alone does not imply a
summable cofinal row: support multiplicity, ultraviolet scaling, inverse
fermion floors, and metric derivatives all enter the influence norm.
\end{firewall}

\begin{firewall}[Locality is not a mass gap]
\label{fire:e8v10-mass}
SIAC controls propagation for fixed physical horizons and supports
projective state construction.  It neither excludes the infrared singlet
packets found in YM V61 nor proves a lower bound for a reflected Hamiltonian.
The unbroken long-range SM gauge sector must remain in the state without a
fabricated global contraction gap.
\end{firewall}

\begin{assessment}[Progress at V10]
\label{ass:e8v10-status}
The sector matrix now has an ownership rule.  The factorized reference is
block diagonal, and every generated cross block is charged to the
old/detail interaction.  If both the interaction influence and its metric
derivative carry a summable refinement amplitude, V8 turns them into
summable remote base and stress rows.  The next task is to derive those
influence bounds from a concrete finite Gibbs or Lindblad interaction,
including the inverse fermion and Higgs-chart constants.
\end{assessment}

\begin{programme}[Eighth-series V11: local interaction-to-influence
compiler]
\label{prog:e8v10-V11}
V11 will define tensor-stable conditional oscillations on a finite product
algebra and bound the influence block generated by one local Hamiltonian or
Lindblad interaction term.  It will express \(E^{ab}_{ij}\) in commutator
or completely bounded norms, then identify the ultraviolet and support
multiplicity factors which must be included in \(\varepsilon_n\).
\end{programme}

\begin{firewall}[Claim boundary after eighth-series V10]
\label{fire:e8v10-boundary}
V10 performs the compulsory companion audit, proves exact block diagonality
of a tensor-stable factorized reference, assigns every cross-sector row to
the interaction generator, proves finite-horizon influence perturbation,
and derives interaction-tagged base and stress summability.

V10 does not prove tensor stability or the interaction influence bounds for
the physical SM carrier, populate the refinement amplitude, control the
independent diagonal and boundary rows, establish a reflected transfer gap,
remove the cutoff, or construct the Standard Model--Einstein continuum.
\end{firewall}

\begin{theorem}[Eighth-series V10 result]
\label{thm:e8v10-result}
A tensor-stable product generator has an exactly block-diagonal sector
influence matrix.  Adding an interaction with weighted influence
\(O(\varepsilon_n)\) creates only
\(O(\varepsilon_n)\) cross-sector finite-horizon propagation, and if its
metric derivative has the same size, both direct and ordered adjacent
stress insertions retain that amplitude.  Remote support and
\eqref{eq:e8v10-tagged-sum} then make the generated cross-sector base and
stress rows summable, without asserting a global mass gap.
\end{theorem}

\begin{proof}
Combine \Cref{thm:e8v10-block-diagonal,
prop:e8v10-cross-owner,thm:e8v10-influence-perturbation,
thm:e8v10-tagged,cor:e8v10-SIAC}.
\end{proof}

\paragraph{Parent-version record.}
V10 was formed from
\texttt{Renewal\_SM\_Einstein\_Continuum\_Research\_Notes\_V9.tex}.
The V9 parent had \(3\,236\) logical source lines, \(121\,440\) bytes, and
SHA--256
\[
 \texttt{805BFC0389AE9BBC2CD46AF9418829D2E008DA0DBADC2F29869DCD6855599CD4}.
\]
The next compulsory companion audit is V15.

% =============================================================================
% END APPENDED EIGHTH-SERIES V10 MATERIAL
% =============================================================================

% =============================================================================
% BEGIN APPENDED EIGHTH-SERIES V11 MATERIAL
% =============================================================================

\section{Conditional oscillations and local interaction readout}
\label{sec:v8v11}

V10 reduces every generated cross-sector row to the interaction generator.
We now bound one local Hamiltonian or Lindblad interaction in a
Tensor-product conditional oscillation.  This also identifies a restriction
in V8: its support-reset insertion hypothesis applies to conditional
replacement maps, but not to an arbitrary local commutator acting on a
quasilocal observable.

\subsection{Tensor-product conditional oscillations}
\label{sec:e8v11-oscillations}

Let \(\mathfrak A=\bigotimes_{i\in I}M_{d_i}(\mathbb C)\).  Let
\(E_i:\mathfrak A\to\mathfrak A_{i^c}\) be normalized partial trace over
site \(i\), tensored with the identity on the complement, and define
\begin{equation}
 \delta_i(A):=\|A-E_iA\|.
 \label{eq:e8v11-delta}
\end{equation}
The expectations commute.  For \(X=\{i_1,\ldots,i_m\}\), put
\begin{equation}
 E_{X^c}:=E_{i_1}\cdots E_{i_m},
 \label{eq:e8v11-EX}
\end{equation}
which removes the dependence on every site of \(X\).

\begin{lemma}[Conditional telescoping]
\label{lem:e8v11-telescope}
For every finite \(X\subset I\),
\begin{equation}
 \|A-E_{X^c}A\|
 \leq\sum_{i\in X}\delta_i(A).
 \label{eq:e8v11-telescope}
\end{equation}
If a unital channel acts only on tensor factors disjoint from \(i\), it
commutes with \(E_i\) and does not increase \(\delta_i\).
\end{lemma}

\begin{proof}
Order \(X=\{i_1,\ldots,i_m\}\) and telescope
\[
 A-E_{i_1}\cdots E_{i_m}A
 =\sum_{r=1}^m E_{i_1}\cdots E_{i_{r-1}}
 (I-E_{i_r})A.
\]
Every conditional expectation is unital completely positive and hence an
operator-norm contraction.  Since it commutes with \(E_{i_r}\), the norm
of the \(r\)-th term is at most \(\delta_{i_r}(A)\).  For the last
statement, commute the disjoint channel through \(E_i\) and use unital
contractivity.
\end{proof}

\subsection{Local annihilating superoperators}
\label{sec:e8v11-annihilating}

Call a bounded Heisenberg superoperator \(\mathcal J_X^*\) locally
annihilating on \(X\) if
\begin{equation}
 \mathcal J_X^*E_{X^c}=0.
 \label{eq:e8v11-annihilating}
\end{equation}

\begin{theorem}[Local interaction readout]
\label{thm:e8v11-readout}
If \(\mathcal J_X^*\) satisfies \eqref{eq:e8v11-annihilating}, then
\begin{equation}
 \|\mathcal J_X^*A\|
 \leq\|\mathcal J_X^*\|_{\rm cb}
       \sum_{i\in X}\delta_i(A).
 \label{eq:e8v11-readout}
\end{equation}
The same estimate holds at every matrix amplification.
\end{theorem}

\begin{proof}
By local annihilation,
\[
 \mathcal J_X^*A
 =\mathcal J_X^*(A-E_{X^c}A).
\]
Apply the completely bounded norm and
\Cref{lem:e8v11-telescope}.  At an amplification, use
\(E_i\otimes I_k\); the same telescoping proof applies, and taking the
supremum over \(k\) gives the cb statement.
\end{proof}

\begin{corollary}[Hamiltonian interaction constant]
\label{cor:e8v11-Hamiltonian}
Let \(h_X=h_X^*\) be supported in \(X\) and
\begin{equation}
 \mathcal J_{h_X}^*(A):=i[h_X,A].
 \label{eq:e8v11-Hamiltonian}
\end{equation}
Then \(\mathcal J_{h_X}^*E_{X^c}=0\),
\(\|\mathcal J_{h_X}^*\|_{\rm cb}\leq2\|h_X\|\), and
\begin{equation}
 \|[h_X,A]\|
 \leq2\|h_X\|\sum_{i\in X}\delta_i(A).
 \label{eq:e8v11-H-bound}
\end{equation}
For a metric derivative, the same bound holds with \(h_X\) replaced by
\(Dh_X[v]\).
\end{corollary}

\begin{proof}
An operator in \(\mathfrak A_{X^c}\) commutes with \(h_X\), proving local
annihilation.  Left and right multiplication by \(h_X\) have cb norm
\(\|h_X\|\), so their commutator has cb norm at most
\(2\|h_X\|\).  Apply \Cref{thm:e8v11-readout}; differentiation is linear.
\end{proof}

\begin{corollary}[Lindblad interaction constant]
\label{cor:e8v11-Lindblad}
Let
\begin{equation}
 \mathcal D_X^*(A):=
 \sum_\alpha\left(
 L_{X,\alpha}^*AL_{X,\alpha}
 -{1\over2}\{L_{X,\alpha}^*L_{X,\alpha},A\}\right)
 \label{eq:e8v11-Lindblad}
\end{equation}
with every \(L_{X,\alpha}\) supported in \(X\).  Then
\begin{equation}
 \mathcal D_X^*E_{X^c}=0,
 \qquad
 \|\mathcal D_X^*\|_{\rm cb}
 \leq2\sum_\alpha\|L_{X,\alpha}\|^2,
 \label{eq:e8v11-D-cb}
\end{equation}
and
\begin{equation}
 \|\mathcal D_X^*A\|
 \leq2\left(\sum_\alpha\|L_{X,\alpha}\|^2\right)
       \sum_{i\in X}\delta_i(A).
 \label{eq:e8v11-D-bound}
\end{equation}
\end{corollary}

\begin{proof}
If \(B\) is supported outside \(X\), it commutes with every
\(L_{X,\alpha}\), and the three terms for each \(\alpha\) cancel.  The cb
norm of \(A\mapsto L^*AL\) is \(\|L\|^2\); each anticommutator half has
combined cb norm at most \(\|L\|^2\).  Sum over \(\alpha\) and apply
\Cref{thm:e8v11-readout}.
\end{proof}

\subsection{Interaction influence ledger}
\label{sec:e8v11-ledger}

For a local interaction term define
\begin{equation}
 c_X:=2\|h_X\|+2\sum_\alpha\|L_{X,\alpha}\|^2.
 \label{eq:e8v11-cX}
\end{equation}
If it couples sectors \(a\) and \(b\), assign its readout constant to the
corresponding V10 cross block.  A sufficient exponential interaction row
is
\begin{equation}
 \mathcal I_{\mu;\sigma}:=
 \sup_{a,i}\sum_{\substack{X\ni i,\ b:\ X\text{ couples }a\text{ to }b}}
 {\sigma_a\over\sigma_b}|X|e^{\mu\operatorname{diam}X}c_X<\infty.
 \label{eq:e8v11-interaction-row}
\end{equation}
The factor \(|X|\) comes from the telescoping sum in
\eqref{eq:e8v11-readout}.

\begin{proposition}[Interaction row bounds the generated readout]
\label{prop:e8v11-interaction}
For every local observable, the sum of Hamiltonian and Lindblad interaction
readouts meeting a fixed output cell is bounded by
\(\mathcal I_{\mu;\sigma}\) times the sector-weighted exponential
oscillation norm.  If, at adjacent scale \(n\),
\begin{equation}
 \mathcal I_{\mu;\sigma,n}\leq C\varepsilon_n,
 \qquad
 \mathcal I^{(q)}_{\mu;\sigma,n}\leq C_q\varepsilon_n,
 \label{eq:e8v11-refinement-row}
\end{equation}
where the second row uses \(Dh_X[v]\) and the differentiated Lindblad
operators, then the direct interaction and metric-interaction readout
constants required by V10 are \(O(\varepsilon_n)\).
\end{proposition}

\begin{proof}
Apply \Cref{cor:e8v11-Hamiltonian,cor:e8v11-Lindblad} to each term.  For
\(j\in X\),
\(e^{\mu d(i,Z)}\leq
 e^{\mu\operatorname{diam}X}e^{\mu d(j,Z)}\) whenever \(i\in X\).
The sector weights give \(\sigma_a/\sigma_b\), and summing the at most
\(|X|\) input oscillations yields \eqref{eq:e8v11-interaction-row}.
The two bounds in \eqref{eq:e8v11-refinement-row} then give the stated
orders.
\end{proof}

\begin{firewall}[Correction to the scope of V8]
\label{fire:e8v11-V8-scope}
V8 assumes
\(\delta_i(\mathcal J_X^*A)=0\) outside a fixed neighbourhood of \(X\).
A replacement or conditional-reset insertion can satisfy this.  A generic
Hamiltonian commutator does not: if
\(A=A_X\otimes B_Y\) with \(Y\cap X=\varnothing\), then
\([h_X,A]=[h_X,A_X]\otimes B_Y\) still depends on \(Y\).
Therefore V8's two-source proof applies directly only to support-reset
insertions.  For Hamiltonian or Lindblad interactions,
\eqref{eq:e8v11-readout} is valid, but the two-source tree bound requires a
separate nested-commutator or influence-path proof which does not assume
support reset.  V10's stress conclusion remains conditional on that proof.
\end{firewall}

\begin{proof}
The displayed tensor-product example proves failure of support reset.
The local readout theorem proves the replacement bound.  Since the V8 proof
uses support reset at its first marked insertion, it cannot be invoked for
the example without an additional path expansion.
\end{proof}

\begin{assessment}[Revised finite-regulator frontier]
\label{ass:e8v11-status}
The ultraviolet and multiplicity content of one interaction row is now
explicit in \eqref{eq:e8v11-interaction-row}.  Hamiltonian and Lindblad
readouts are controlled by tensor-product conditional oscillations, and
the metric row is obtained by differentiating their local coefficients.
The remote two-source stress estimate is proved for reset channels but
remains open for generic Hamiltonian interactions; the next attack must use
a genuine nested-commutator Dyson path.
\end{assessment}

\begin{programme}[Eighth-series V12: nested-commutator marked paths]
\label{prog:e8v11-V12}
V12 will replace V8 support reset by the conditional-oscillation readout
\eqref{eq:e8v11-readout}.  It will expand the evolved observable into
interaction paths, mark the refinement and metric terms, and prove that a
nonzero nested commutator requires a connected support chain.  The weighted
sum over connected chains will provide the two-source tree factor for
Hamiltonian interactions.
\end{programme}

\begin{firewall}[Claim boundary after eighth-series V11]
\label{fire:e8v11-boundary}
V11 proves conditional telescoping, cb local readout for locally
annihilating superoperators, explicit Hamiltonian and Lindblad constants,
and a sector-weighted interaction ledger.  It also corrects the scope of
the V8 support-reset hypothesis.

V11 does not yet prove the two-source marked-path estimate for generic
Hamiltonian or Lindblad interactions, populate the physical SM interaction
row, control unbounded fields, remove the cutoff, or construct the Standard
Model--Einstein continuum.
\end{firewall}

\begin{theorem}[Eighth-series V11 result]
\label{thm:e8v11-result}
For tensor-product conditional oscillations, every bounded local
Hamiltonian or Lindblad term annihilates the outside algebra and obeys the
readout estimate \eqref{eq:e8v11-readout}, with constants
\(2\|h_X\|\) and \(2\sum_\alpha\|L_{X,\alpha}\|^2\).  The weighted sum
\eqref{eq:e8v11-interaction-row} therefore controls the generated sector
readout and its metric derivative.  Support-reset two-source locality does
not automatically extend to these interactions and requires a connected
nested-commutator proof.
\end{theorem}

\begin{proof}
Combine \Cref{lem:e8v11-telescope,thm:e8v11-readout,
cor:e8v11-Hamiltonian,cor:e8v11-Lindblad,
prop:e8v11-interaction,fire:e8v11-V8-scope}.
\end{proof}

\paragraph{Parent-version record.}
V11 was formed from
\texttt{Renewal\_SM\_Einstein\_Continuum\_Research\_Notes\_V10.tex}.
The V10 parent had \(3\,590\) logical source lines, \(135\,524\) bytes,
and SHA--256
\[
 \texttt{F15B61E1237F51F4057395BB40C3EAC8C122145CDEF5BD5BE0ADEFCBF75F76DC}.
\]
The next compulsory companion audit is V15.

% =============================================================================
% END APPENDED EIGHTH-SERIES V11 MATERIAL
% =============================================================================
% =============================================================================
% BEGIN APPENDED EIGHTH-SERIES V12 MATERIAL
% =============================================================================

\section{Nested Hamiltonian commutators and two-marked locality}
\label{sec:v8v12}

V11 shows that a local Hamiltonian commutator does not reset the support of
a quasilocal observable.  The correct proof of two-source locality uses the
derivation property instead.  We first derive a one-commutator
Lieb--Robinson estimate from conditional oscillations.  A Jacobi identity
then gives two different bounds on a double commutator; taking their
minimum forces a connected three-set geometry.

\subsection{Conditional-oscillation propagation for a Hamiltonian}
\label{sec:e8v12-propagation}

On the finite tensor-product algebra of V11, let
\begin{equation}
 H=\sum_{Z\subset I}h_Z,
 \qquad h_Z=h_Z^*,
 \qquad
 \tau_t(A):=e^{itH}Ae^{-itH}.
 \label{eq:e8v12-H}
\end{equation}
Assume the interaction is finite at every regulator.  Define the
nonnegative influence matrix
\begin{equation}
 K_{ij}:=4\sum_{Z\ni i,j}\|h_Z\|.
 \label{eq:e8v12-K}
\end{equation}
The sum includes \(i=j\).  Let
\begin{equation}
 k_\mu:=\sup_i\sum_j e^{\mu d(i,j)}K_{ij}<\infty.
 \label{eq:e8v12-kmu}
\end{equation}

\begin{lemma}[Hamiltonian Dini inequality]
\label{lem:e8v12-Dini}
For the conditional oscillations
\(\delta_i(A)=\|A-E_iA\|\),
\begin{equation}
 D^+\delta_i(\tau_tA)
 \leq\sum_jK_{ij}\delta_j(\tau_tA).
 \label{eq:e8v12-Dini}
\end{equation}
Consequently,
\begin{equation}
 \delta_i(\tau_tA)
 \leq\sum_j(e^{|t|K})_{ij}\delta_j(A)
 \leq e^{k_\mu|t|}\sum_j e^{-\mu d(i,j)}\delta_j(A).
 \label{eq:e8v12-osc-propagation}
\end{equation}
\end{lemma}

\begin{proof}
Fix \(i\) and split
\[
 H=H_{i^c}+H_i,
 \qquad
 H_i:=\sum_{Z\ni i}h_Z.
\]
The automorphism generated by \(H_{i^c}\) acts outside site \(i\),
commutes with \(E_i\), and preserves the operator norm.  Passing to its
interaction picture therefore contributes zero to the upper Dini growth of
\(\delta_i\).  For one term \(h_Z\) with \(i\in Z\), the norm derivative
is bounded by
\[
 2\|[h_Z,\tau_tA]\|.
\]
By V11 conditional telescoping, the part of \(\tau_tA\) commuting with
\(h_Z\) may be subtracted, and hence
\[
 \|[h_Z,\tau_tA]\|
 \leq2\|h_Z\|\sum_{j\in Z}\delta_j(\tau_tA).
\]
Summing over \(Z\ni i\) gives
\[
 D^+\delta_i(\tau_tA)
 \leq4\sum_{Z\ni i}\|h_Z\|
                 \sum_{j\in Z}\delta_j(\tau_tA)
 =\sum_jK_{ij}\delta_j(\tau_tA).
\]
The comparison theorem from V8 yields the first inequality in
\eqref{eq:e8v12-osc-propagation}.  The V8 weighted-matrix exponential
bound gives the second.  Negative time follows by replacing \(H\) by
\(-H\), which leaves \(K\) unchanged.
\end{proof}

\begin{theorem}[Pair commutator bound]
\label{thm:e8v12-pair}
Let \(V_X\) and \(A_B\) be supported in finite sets \(X,B\).  Then
\begin{equation}
 \|[V_X,\tau_t(A_B)]\|
 \leq 2c_0|X||B|\|V_X\|\|A_B\|
 e^{-\mu(d(X,B)-v_\mu|t|)},
 \qquad
 v_\mu:={k_\mu\over\mu},
 \label{eq:e8v12-pair}
\end{equation}
where \(c_0\) is the local-observable oscillation constant in V8.
\end{theorem}

\begin{proof}
The V11 Hamiltonian readout and
\eqref{eq:e8v12-osc-propagation} give
\begin{align*}
 \|[V_X,\tau_t(A_B)]\|
 &\leq2\|V_X\|\sum_{i\in X}\delta_i(\tau_tA_B)\\
 &\leq2\|V_X\||X|e^{k_\mu|t|}
 e^{-\mu d(X,B)}
 \sum_{b\in B}\delta_b(A_B).
\end{align*}
Use
\(\sum_{b\in B}\delta_b(A_B)\leq c_0|B|\|A_B\|\) and
\(k_\mu=\mu v_\mu\).
\end{proof}

\subsection{The double-commutator geometry}
\label{sec:e8v12-double}

Let \(V_X,W_Y,A_B\) be local observables.  For
\(t_1,t_2\geq0\), define
\begin{equation}
 \mathcal D_{X,Y;B}(t_1,t_2)
 :=[V_X,\tau_{t_1}([W_Y,\tau_{t_2}(A_B)])].
 \label{eq:e8v12-D}
\end{equation}
Put
\begin{equation}
 a=d(Y,B),\qquad b=d(X,Y),\qquad c=d(X,B),
 \qquad
 m(X,Y;B):=\max\{a,\min\{b,c\}\}.
 \label{eq:e8v12-m}
\end{equation}

\begin{lemma}[Minimum of two Jacobi bounds]
\label{lem:e8v12-two-bounds}
There is a constant
\begin{equation}
 C_{X,Y,B}:=4c_0\max\{|Y||B|,|X||Y|+|X||B|\}
 \label{eq:e8v12-C}
\end{equation}
such that, with \(T=t_1+t_2\),
\begin{equation}
 \|\mathcal D_{X,Y;B}(t_1,t_2)\|
 \leq2C_{X,Y,B}|V_X\|\|W_Y\|\|A_B\|
 e^{\mu v_\mu T}e^{-\mu m(X,Y;B)}.
 \label{eq:e8v12-min-bound}
\end{equation}
\end{lemma}

\begin{proof}
The outer commutator and isometry of \(\tau_{t_1}\) give
\begin{align}
 \|\mathcal D_{X,Y;B}\|
 &\leq2\|V_X\|\|[W_Y,\tau_{t_2}(A_B)]\| 
 \notag\\
 &\leq4c_0|Y||B|\|V_X\|\|W_Y\|\|A_B\|
 e^{\mu v_\mu T}e^{-\mu a}
 \label{eq:e8v12-bound-a}
\end{align}
by \Cref{thm:e8v12-pair}; we enlarged the time factor from \(t_2\) to
\(T\).

Conjugate the outer commutator by \(\tau_{-t_1}\).  Its norm becomes
\[
 \|[\tau_{-t_1}(V_X),[W_Y,\tau_{t_2}(A_B)]]\|.
\]
Jacobi's identity writes this as the sum of
\[
 [[\tau_{-t_1}(V_X),W_Y],\tau_{t_2}(A_B)]
 \quad\text{and}\quad
 [W_Y,[\tau_{-t_1}(V_X),\tau_{t_2}(A_B)]].
\]
Bound the first outer commutator by twice \(\|A_B\|\), and the second by
twice \(\|W_Y\|\).  Pair locality gives
\begin{align}
 \|\mathcal D_{X,Y;B}\|
 \leq{}&4c_0|X||Y|\|V_X\|\|W_Y\|\|A_B\|
 e^{\mu v_\mu T}e^{-\mu b}
 \notag\\
 &+4c_0|X||B|\|V_X\|\|W_Y\|\|A_B\|
 e^{\mu v_\mu T}e^{-\mu c}.
 \label{eq:e8v12-bound-bc}
\end{align}
Here
\(\|[\tau_{-t_1}V_X,\tau_{t_2}A_B]\|
 =\|[V_X,\tau_TA_B]\|\).

Take the smaller of \eqref{eq:e8v12-bound-a} and
\eqref{eq:e8v12-bound-bc}, dominate coefficients by
\eqref{eq:e8v12-C}, and use
\[
 e^{-\mu b}+e^{-\mu c}
 \leq2e^{-\mu\min\{b,c\}}.
\]
For positive numbers \(u,v\),
\(\min\{e^{-u},2e^{-v}\}\leq2e^{-\max\{u,v\}}\).
This proves \eqref{eq:e8v12-min-bound}.
\end{proof}

For bounded-diameter marked supports, the separation in
\eqref{eq:e8v12-m} controls a genuine three-set tree.  Define
\begin{equation}
 \mathfrak t(B,X,Y):=
 \inf_{\substack{x\in X,\ y\in Y,\ b_0\in B}}
 \bigl(d(x,y)+d(y,b_0)+d(b_0,x)
 -\max\{d(x,y),d(y,b_0),d(b_0,x)\}\bigr).
 \label{eq:e8v12-tree}
\end{equation}
For three points the expression in parentheses is the length of a minimum
spanning tree.

\begin{lemma}[Pair separation controls the tree]
\label{lem:e8v12-tree}
If
\begin{equation}
 \operatorname{diam}X, \operatorname{diam}Y,
 \ \operatorname{diam}B\leq r_0,
 \label{eq:e8v12-diameter}
\end{equation}
then
\begin{equation}
 m(X,Y;B)\geq{1\over2}\mathfrak t(B,X,Y)-2r_0.
 \label{eq:e8v12-tree-comparison}
\end{equation}
\end{lemma}

\begin{proof}
Choose \(x,y,b_0\) within an arbitrarily small error of the infimum in
\eqref{eq:e8v12-tree}.  Each distance between these chosen points is at
most the corresponding set distance plus the two set diameters, hence at
most that set distance plus \(2r_0\).  For three point distances, if
\(d_1\leq d_2\leq d_3\), the tree length is \(d_1+d_2\), while
\(\max\{a,\min(b,c)\}\), with any one distance designated as \(a\), is
at least \(d_2\geq(d_1+d_2)/2\).  Replacing point distances by set
distances loses at most \(2r_0\).  Let the approximation error tend to
zero.
\end{proof}

\begin{theorem}[Hamiltonian two-marked tree bound]
\label{thm:e8v12-tree-bound}
Under \eqref{eq:e8v12-diameter},
\begin{align}
 \|\mathcal D_{X,Y;B}(t_1,t_2)\|
 \leq{}&2C_{X,Y,B}e^{2\mu r_0}
 \|V_X\|\|W_Y\|\|A_B\|
 \notag\\
 &\times
 \exp\left\{-{\mu\over2}
 \bigl(\mathfrak t(B,X,Y)-2v_\mu(t_1+t_2)\bigr)\right\}.
 \label{eq:e8v12-tree-bound}
\end{align}
Thus V7's two-source locality hypothesis holds for bounded Hamiltonian
commutator insertions, with exponent \(\mu/2\), velocity \(2v_\mu\), and
the displayed support constant.
\end{theorem}

\begin{proof}
Insert \eqref{eq:e8v12-tree-comparison} into
\eqref{eq:e8v12-min-bound}.  Rewrite
\(e^{\mu v_\mu T}e^{-\mu\mathfrak t/2}\) as the exponential in
\eqref{eq:e8v12-tree-bound}.  This is exactly the geometry and time form
of the V7 hypothesis.
\end{proof}

\subsection{Stress-row consequence}
\label{sec:e8v12-stress}

Let the refinement and metric Hamiltonian derivatives decompose as
\begin{equation}
 V^{\rm ref}=\sum_XV_X,qquad
 W^{\rm met}=\sum_YW_Y,
 \label{eq:e8v12-insertions}
\end{equation}
with uniformly bounded support diameter and cardinality.  Suppose their
weighted marked sums are finite in the V7 tree norm.

\begin{corollary}[Hamiltonian population of the V7 ordered row]
\label{cor:e8v12-V7}
For a bounded finite-regulator Hamiltonian interaction satisfying
\eqref{eq:e8v12-kmu}, the ordered refinement/metric Duhamel terms in V7
obey its two-source locality estimate.  If the resulting component sums
satisfy the V8 far-row series with \(\mu\) replaced by \(\mu/2\) and
velocity by \(2v_\mu\), then the Hamiltonian far adjacent stress row is
summable.
\end{corollary}

\begin{proof}
Apply \Cref{thm:e8v12-tree-bound} to every pair \((X,Y)\) inside the two
ordered time simplexes and sum their nonnegative majorants.  The V7 remote
stress theorem and V8 summability argument then apply with the degraded
exponent and velocity.
\end{proof}

\begin{firewall}[Lindblad insertions require their own multicommutator
argument]
\label{fire:e8v12-Lindblad}
The Jacobi step uses that a Hamiltonian insertion is a derivation.  A
Lindblad dissipator obeys no ordinary Leibniz or Jacobi identity, so
\Cref{thm:e8v12-tree-bound} does not prove the corresponding two-source
bound for generic dissipative insertions.  Reset/conditional Lindblad maps
remain covered by V8; other dissipators need a completely bounded
influence-path expansion.
\end{firewall}

\begin{firewall}[Finite-regulator boundedness]
\label{fire:e8v12-finite}
The proof assumes bounded local Hamiltonian terms and a finite uniform
interaction row \(k_\mu\).  It does not pass an unbounded Dirac, electric,
or gravitational generator to the continuum.  Such a passage still needs
form-domain convergence and uniform control of the local interaction
majorants.
\end{firewall}

\begin{assessment}[Progress at V12]
\label{ass:e8v12-status}
The principal logical gap identified at V11 is closed for bounded
Hamiltonian interactions.  Support reset is unnecessary: ordinary pair
locality plus Jacobi gives a connected two-marked tree bound, at the cost of
halving the spatial exponent.  The far Hamiltonian stress row is now a
proved implication of one finite interaction norm and the explicit
component summability series.
\end{assessment}

\begin{programme}[Eighth-series V13: completely bounded dissipative paths]
\label{prog:e8v12-V13}
V13 will treat a local Lindblad generator without assuming reset.  It will
use the carré-du-champ/product-defect identity to replace Jacobi, derive a
marked completely bounded influence recursion, and determine whether the
same tree geometry holds or whether an additional local branching factor
must enter the SIAC ledger.
\end{programme}

\begin{firewall}[Claim boundary after eighth-series V12]
\label{fire:e8v12-boundary}
V12 derives Hamiltonian conditional-oscillation propagation, a pair
Lieb--Robinson bound, two complementary Jacobi bounds for the double
commutator, their three-set tree consequence, and population of the V7
ordered far stress row for bounded Hamiltonian interactions.

V12 does not cover generic non-reset Lindblad terms, prove uniform physical
SM interaction norms, control unbounded continuum generators, sum the near
or boundary rows, remove the cutoff, or construct the Standard
Model--Einstein continuum.
\end{firewall}

\begin{theorem}[Eighth-series V12 result]
\label{thm:e8v12-result}
A bounded Hamiltonian interaction with finite exponential influence row
\eqref{eq:e8v12-kmu} obeys the pair bound
\eqref{eq:e8v12-pair}.  Two Hamiltonian refinement/metric insertions obey
the marked tree estimate \eqref{eq:e8v12-tree-bound} without any
support-reset assumption.  Consequently the Hamiltonian part of the V7 far
adjacent stress row is summable under the explicit V8 component series with
the stated exponent and velocity degradation.
\end{theorem}

\begin{proof}
Combine \Cref{lem:e8v12-Dini,thm:e8v12-pair,
lem:e8v12-two-bounds,lem:e8v12-tree,
thm:e8v12-tree-bound,cor:e8v12-V7}.
\end{proof}

\paragraph{Parent-version record.}
V12 was formed from
\texttt{Renewal\_SM\_Einstein\_Continuum\_Research\_Notes\_V11.tex}.
The V11 parent had \(3\,876\) logical source lines, \(146\,117\) bytes,
and SHA--256
\[
 \texttt{BFD7BA8FF3DF159F1840EE414E19B6110944DA228DC97E4A77CFEDD0FEEAE4EB}.
\]
The next compulsory companion audit is V15.

% =============================================================================
% END APPENDED EIGHTH-SERIES V12 MATERIAL
% =============================================================================

% =============================================================================
% BEGIN APPENDED EIGHTH-SERIES V13 MATERIAL
% =============================================================================

\section{Completely bounded dissipative insertions without support reset}
\label{sec:v8v13}

V12 proves two-marked locality for Hamiltonian derivations by Jacobi's
identity.  A Lindblad dissipator is not a derivation.  We replace Jacobi by
a conditional-oscillation decomposition: away from the insertion support,
the pre-existing oscillation is propagated; on the support, the insertion
creates one localized bump.  Two complementary estimates again force a
connected three-set decay.

\subsection{Local superoperator oscillation decomposition}
\label{sec:e8v13-insertion}

Let \(\mathcal J_Y^*\) be a bounded Heisenberg superoperator localized in
\(Y\), satisfying
\begin{align}
 \mathcal J_Y^*E_{Y^c}&=0,
 \label{eq:e8v13-annihilate}\\
 E_i\mathcal J_Y^*&=\mathcal J_Y^*E_i
 \qquad(i\notin Y),
 \label{eq:e8v13-commute}\\
 \|\mathcal J_Y^*\|_{\rm cb}&\leq j_Y.
 \label{eq:e8v13-cb}
\end{align}
Hamiltonian commutators and Lindblad terms supported in \(Y\) have these
properties, with the constants computed in V11.

\begin{lemma}[Persistent tail plus localized bump]
\label{lem:e8v13-decomposition}
For every observable \(A\),
\begin{equation}
 \delta_i(\mathcal J_Y^*A)
 \leq j_Y\delta_i(A)
 +2j_Y\mathbf1_{\{i\in Y\}}
       \sum_{y\in Y}\delta_y(A).
 \label{eq:e8v13-decomposition}
\end{equation}
Moreover,
\begin{equation}
 \|\mathcal J_Y^*A\|
 \leq j_Y\sum_{y\in Y}\delta_y(A).
 \label{eq:e8v13-readout}
\end{equation}
\end{lemma}

\begin{proof}
If \(i\notin Y\), use \eqref{eq:e8v13-commute}:
\[
 \delta_i(\mathcal J_Y^*A)
 =\|\mathcal J_Y^*(A-E_iA)\|
 \leq j_Y\delta_i(A).
\]
If \(i\in Y\), contractivity of \(E_i\) gives
\[
 \delta_i(\mathcal J_Y^*A)
 \leq2\|\mathcal J_Y^*A\|.
\]
Local annihilation and the V11 telescoping lemma give
\[
 \|\mathcal J_Y^*A\|
 =\|\mathcal J_Y^*(A-E_{Y^c}A)\|
 \leq j_Y\sum_{y\in Y}\delta_y(A).
\]
These estimates prove both claims; adding the harmless first term on
\(Y\) yields the uniform formula \eqref{eq:e8v13-decomposition}.
\end{proof}

\begin{remark}[Why this repairs the V8 restriction]
\label{rem:e8v13-repair}
Equation \eqref{eq:e8v13-decomposition} does not claim that the insertion
output is supported near \(Y\).  Its first term retains every old
oscillation, exactly as in the V11 tensor-product counterexample.  Only the
second term is newly localized at \(Y\).
\end{remark}

\subsection{Propagation for a local Lindblad generator}
\label{sec:e8v13-propagation}

Let
\begin{equation}
 \mathcal L^*=\sum_Z\mathcal J_Z^*
 \label{eq:e8v13-L}
\end{equation}
be a bounded finite-regulator Heisenberg Lindblad generator, decomposed into
local Lindblad and Hamiltonian terms.  Let \(j_Z\) dominate the cb norm of
each term and define
\begin{equation}
 K_{ik}:=2\sum_{Z\ni i,k}j_Z,
 \qquad
 k_\mu:=\sup_i\sum_k e^{\mu d(i,k)}K_{ik}.
 \label{eq:e8v13-K}
\end{equation}

\begin{lemma}[Lindblad conditional-oscillation propagation]
\label{lem:e8v13-propagation}
For the unital completely positive semigroup
\(\mathcal S^*(t)=e^{t\mathcal L^*}\),
\begin{equation}
 \delta_i(\mathcal S^*(t)A)
 \leq\sum_k(e^{tK})_{ik}\delta_k(A)
 \leq e^{k_\mu t}
       \sum_k e^{-\mu d(i,k)}\delta_k(A).
 \label{eq:e8v13-propagation}
\end{equation}
\end{lemma}

\begin{proof}
Fix \(i\).  The sum of generator terms with \(Z\not\ni i\) commutes with
\(E_i\) and generates a unital completely positive contraction outside
site \(i\); in its interaction picture it does not increase
\(\delta_i\).  For a term with \(i\in Z\), the upper norm derivative is at
most
\[
 2\|\mathcal J_Z^*A\|
 \leq2j_Z\sum_{k\in Z}\delta_k(A)
\]
by \eqref{eq:e8v13-readout}.  Therefore
\[
 D^+\delta_i(\mathcal S^*(t)A)
 \leq\sum_kK_{ik}\delta_k(\mathcal S^*(t)A).
\]
The V8 comparison and weighted kernel lemmas give
\eqref{eq:e8v13-propagation}.
\end{proof}

\subsection{Two arbitrary local insertions}
\label{sec:e8v13-two}

Let \(A_B\) be local and let
\(\mathcal J_X^*,\mathcal J_Y^*\) obey
\eqref{eq:e8v13-annihilate}--\eqref{eq:e8v13-cb}.  Put
\begin{equation}
 \mathcal M_{X,Y;B}(t_1,t_2)
 :=\mathcal J_X^*\mathcal S^*(t_1)
   \mathcal J_Y^*\mathcal S^*(t_2)A_B.
 \label{eq:e8v13-M}
\end{equation}
Let \(T=t_1+t_2\),
\(a=d(Y,B)\), \(b=d(X,Y)\), and \(c=d(X,B)\).

\begin{lemma}[Two complementary dissipative bounds]
\label{lem:e8v13-complementary}
There are explicit support constants \(C_1,C_2\), depending only on
\(c_0,|X|,|Y|,|B|\), such that
\begin{align}
 \|\mathcal M_{X,Y;B}\|
 &\leq C_1j_Xj_Y\|A_B\|
 e^{k_\mu T}e^{-\mu a},
 \label{eq:e8v13-bound-a}\\
 \|\mathcal M_{X,Y;B}\|
 &\leq C_2j_Xj_Y\|A_B\|e^{k_\mu T}
 \left(e^{-\mu c}+e^{-\mu(a+b)}\right).
 \label{eq:e8v13-bound-c}
\end{align}
One may take
\begin{equation}
 C_1=c_0|Y||B|,qquad
 C_2=c_0|X||B|\max\{1,2|Y|^2\}.
 \label{eq:e8v13-constants}
\end{equation}
\end{lemma}

\begin{proof}
For the first bound, use the cb norm of the outer insertion, contractivity
of the semigroup, and the readout estimate for the inner insertion:
\begin{align*}
 \|\mathcal M_{X,Y;B}\|
 &\leq j_X\|\mathcal J_Y^*\mathcal S^*(t_2)A_B\|\\
 &\leq j_Xj_Y\sum_{y\in Y}
        \delta_y(\mathcal S^*(t_2)A_B).
\end{align*}
Apply \eqref{eq:e8v13-propagation} and the V8 local-observable bound,
then enlarge \(e^{k_\mu t_2}\) to \(e^{k_\mu T}\).  This gives
\eqref{eq:e8v13-bound-a} with \(C_1\).

For the second bound, let
\(A'=\mathcal S^*(t_2)A_B\) and
\(C=\mathcal J_Y^*A'\).  Apply the outer readout and propagate its input:
\[
 \|\mathcal M_{X,Y;B}\|
 \leq j_X\sum_{x\in X}\sum_i(e^{t_1K})_{xi}\delta_i(C).
\]
Insert \eqref{eq:e8v13-decomposition}.  The persistent-tail part is bounded
by
\[
 c_0j_Xj_Y|X||B|\|A_B\|
 e^{k_\mu T}e^{-\mu c},
\]
because the two propagation matrices multiply to
\(e^{(t_1+t_2)K}\).  The localized bump is first bounded at \(Y\) after
time \(t_2\), producing
\(c_0|Y||B|e^{k_\mu t_2}e^{-\mu a}\|A_B\|\), and is then propagated from
\(Y\) to \(X\) for time \(t_1\).  Summing its at most \(|Y|\) entries and
the \(|X|\) readout sites gives the second term with coefficient at most
\(2c_0|X||Y|^2|B|\).  Combining the two terms gives
\eqref{eq:e8v13-bound-c}--\eqref{eq:e8v13-constants}.
\end{proof}

\begin{theorem}[Non-reset dissipative tree bound]
\label{thm:e8v13-tree}
Assume
\begin{equation}
 \operatorname{diam}X,\operatorname{diam}Y,
 \operatorname{diam}B\leq r_0.
 \label{eq:e8v13-diameter}
\end{equation}
Let \(\mathfrak t(B,X,Y)\) be the three-set tree length defined in V12 and
put \(C=2\max\{C_1,C_2e^{2\mu r_0}\}\).  Then
\begin{align}
 \|\mathcal M_{X,Y;B}(t_1,t_2)\|
 \leq{}&Cj_Xj_Y\|A_B\|e^{2\mu r_0}
 \notag\\
 &\times\exp\left\{-{\mu\over2}
 \bigl(\mathfrak t(B,X,Y)-2v_\mu T\bigr)\right\},
 \qquad v_\mu={k_\mu\over\mu}.
 \label{eq:e8v13-tree}
\end{align}
This applies to generic bounded local Lindblad insertions and does not
assume support reset.
\end{theorem}

\begin{proof}
For point supports, the triangle inequality gives \(c\leq a+b\), so the
parentheses in \eqref{eq:e8v13-bound-c} are at most \(2e^{-\mu c}\).
For sets of diameter at most \(r_0\), the same argument loses at most
\(2r_0\), hence
\[
 e^{-\mu c}+e^{-\mu(a+b)}
 \leq2e^{2\mu r_0}e^{-\mu c}.
\]
Take the smaller of this bound and \eqref{eq:e8v13-bound-a}.  It is at most
a fixed constant times
\(e^{-\mu\max\{a,c\}}\).  For any three point distances, the maximum of
two distances issuing from the point in \(B\) is at least half the minimum
spanning-tree length.  Replacing points by bounded-diameter sets loses at
most \(2r_0\), as in V12.  Therefore
\[
 \max\{a,c\}\geq{1\over2}\mathfrak t(B,X,Y)-2r_0.
\]
Insert this inequality and rewrite
\(e^{k_\mu T}=e^{\mu v_\mu T}\) to obtain
\eqref{eq:e8v13-tree}.
\end{proof}

\subsection{Application to Lindblad parameter derivatives}
\label{sec:e8v13-application}

For the local dissipator in V11,
\begin{equation}
 j_X=2\sum_\alpha\|L_{X,\alpha}\|^2
 \label{eq:e8v13-j-Lindblad}
\end{equation}
is valid.  A parameter derivative is a sum of terms
\begin{align}
 D\mathcal D_X^*[v](A)
 =\sum_\alpha\bigl(&DL_{X,\alpha}[v]^*AL_{X,\alpha}
 +L_{X,\alpha}^*A DL_{X,\alpha}[v]
 \notag\\
 &-{1\over2}\{DL_{X,\alpha}[v]^*L_{X,\alpha}
 +L_{X,\alpha}^*DL_{X,\alpha}[v],A\}\bigr),
 \label{eq:e8v13-DL}
\end{align}
with cb norm bounded by
\begin{equation}
 j_X^{(1)}leq
 4\sum_\alpha\|L_{X,\alpha}\|
                    \|DL_{X,\alpha}[v]\|.
 \label{eq:e8v13-DL-bound}
\end{equation}

\begin{corollary}[Dissipative population of the V7 ordered row]
\label{cor:e8v13-V7}
If the local Lindblad coefficients and their metric/refinement derivatives
have uniformly bounded support diameter and finite weighted interaction
rows formed from \eqref{eq:e8v13-j-Lindblad}--
\eqref{eq:e8v13-DL-bound}, then their two ordered V7 insertions obey
\eqref{eq:e8v13-tree}.  The far dissipative adjacent stress row is
summable under the V8 component series with exponent \(\mu/2\) and velocity
\(2v_\mu\).
\end{corollary}

\begin{proof}
The differentiated local dissipator retains
\eqref{eq:e8v13-annihilate}--\eqref{eq:e8v13-commute}; the four products in
\eqref{eq:e8v13-DL} give the cb estimate
\eqref{eq:e8v13-DL-bound}.  Apply \Cref{thm:e8v13-tree} to the metric and
refinement insertions and sum their component majorants exactly as in the
V12 stress corollary.
\end{proof}

\begin{firewall}[No equilibrium gap follows]
\label{fire:e8v13-gap}
The influence matrix \eqref{eq:e8v13-K} controls finite-time propagation.
It need not have a negative logarithmic norm and yields no global mixing or
Hamiltonian gap.  This is compatible with the V10 audit and with a massless
long-range sector.
\end{firewall}

\begin{assessment}[Progress at V13]
\label{ass:e8v13-status}
The V11 support-reset limitation is now removed for bounded local Lindblad
terms as well as Hamiltonian terms.  The persistent tail is not discarded;
it is paired with a second global readout estimate, and the minimum of the
two bounds restores three-set locality.  Both Hamiltonian and dissipative
far stress rows now reduce to explicit weighted interaction sums.
\end{assessment}

\begin{programme}[Eighth-series V14: unbounded-form approximation]
\label{prog:e8v13-V14}
V14 will address the remaining finite-cutoff boundedness firewall.  It will
truncate an unbounded local generator on a common form core, prove a
Duhamel estimate uniform in the truncation under a relative form bound, and
state the additional local-energy moment needed to pass the marked
insertion bounds.  If norm convergence fails, only strong local expectation
convergence will be claimed.
\end{programme}

\begin{firewall}[Claim boundary after eighth-series V13]
\label{fire:e8v13-boundary}
V13 proves the persistent-tail/local-bump decomposition for arbitrary local
cb insertions, conditional-oscillation propagation for bounded local
Lindblad generators, a non-reset dissipative two-marked tree bound, and
explicit differentiated Lindblad constants.

V13 does not control unbounded continuum generators, populate the physical
SM interaction rows, sum the near or boundary ledgers, prove a transfer
gap, remove the cutoff, or construct the Standard Model--Einstein
continuum.
\end{firewall}

\begin{theorem}[Eighth-series V13 result]
\label{thm:e8v13-result}
A bounded local Lindblad insertion preserves old conditional oscillations
outside its support and creates one controlled local bump on its support.
For a local Lindblad semigroup with finite exponential interaction row,
this decomposition yields the non-reset two-source tree estimate
\eqref{eq:e8v13-tree}.  Hence the dissipative part of the V7 far stress row
has the same explicit summability reduction as the Hamiltonian part, with
no support-reset or global-gap assumption.
\end{theorem}

\begin{proof}
Combine \Cref{lem:e8v13-decomposition,lem:e8v13-propagation,
lem:e8v13-complementary,thm:e8v13-tree,cor:e8v13-V7}.
\end{proof}

\paragraph{Parent-version record.}
V13 was formed from
\texttt{Renewal\_SM\_Einstein\_Continuum\_Research\_Notes\_V12.tex}.
The V12 parent had \(4\,255\) logical source lines, \(158\,613\) bytes,
and SHA--256
\[
 \texttt{7C0C384E6060ACBD62DC32EA2F264E7D0E5AD5F8A5958C2998417FAAFED5B69E}.
\]
The next compulsory companion audit is V15.

% =============================================================================
% END APPENDED EIGHTH-SERIES V13 MATERIAL
% =============================================================================


% =============================================================================
% BEGIN APPENDED EIGHTH-SERIES V14 MATERIAL
% =============================================================================

\section{Weak Duhamel comparison for unbounded form generators}
\label{sec:v8v14}

V6--V13 use bounded finite-regulator generators.  Their operator norms may
diverge as the cutoff is removed even when the associated closed forms
converge.  This appendix proves a weak Duhamel estimate on a common form
domain.  The clock factor is retained, but it is multiplied by a local
form-energy moment.  That moment is the new uniform-integrability row.

\subsection{Equivalent common form norms}
\label{sec:e8v14-form-norm}

Let \(\mathcal H\) be a Hilbert space.  Let \(q_0,q_1\) be densely defined,
closed, nonnegative symmetric forms with the same domain \(\mathcal Q\),
and let \(H_0,H_1\geq0\) be their self-adjoint operators.  Put
\begin{equation}
 w:=q_1-q_0.
 \label{eq:e8v14-w}
\end{equation}
Assume
\begin{equation}
 |w[u,u]|\leq a q_0[u,u]+b\|u\|^2,
 \qquad u\in\mathcal Q,qquad 0\leq a<1,quad b\geq0.
 \label{eq:e8v14-relative}
\end{equation}
Choose \(\lambda>b\) and define
\begin{align}
 \|u\|_{\mathcal Q,0}^2&:=q_0[u,u]+\lambda\|u\|^2,
 \label{eq:e8v14-Q0}\\
 c_\lambda&:=\min\{1-a,1-b/\lambda\},
 &C_\lambda&:=\max\{1+a,1+b/\lambda\},
 \label{eq:e8v14-cC}\\
 \gamma_\lambda&:=\max\{a,b/\lambda\},
 &M_\lambda&:=\sqrt{C_\lambda/c_\lambda}.
 \label{eq:e8v14-gammaM}
\end{align}

\begin{lemma}[Form equivalence and bilinear perturbation bound]
\label{lem:e8v14-equivalence}
For every \(u\in\mathcal Q\),
\begin{equation}
 c_\lambda\|u\|_{\mathcal Q,0}^2
 \leq q_1[u,u]+\lambda\|u\|^2
 \leq C_\lambda\|u\|_{\mathcal Q,0}^2.
 \label{eq:e8v14-equivalence}
\end{equation}
Moreover, for \(u,v\in\mathcal Q\),
\begin{equation}
 |w[u,v]|
 \leq\gamma_\lambda
       \|u\|_{\mathcal Q,0}\|v\|_{\mathcal Q,0}.
 \label{eq:e8v14-bilinear}
\end{equation}
\end{lemma}

\begin{proof}
Add \(\lambda\|u\|^2\) to
\(q_1=q_0+w\) and use \eqref{eq:e8v14-relative} with both signs.  Comparing
the coefficients of \(q_0[u,u]\) and \(\lambda\|u\|^2\) gives
\eqref{eq:e8v14-equivalence}.  The quadratic estimate
\[
 |w[u,u]|\leq\gamma_\lambda\|u\|_{\mathcal Q,0}^2
\]
represents \(w\) as a bounded self-adjoint operator on the Hilbert space
\((\mathcal Q,\|\cdot\|_{\mathcal Q,0})\).  Cauchy--Schwarz for its polar
decomposition gives \eqref{eq:e8v14-bilinear}.
\end{proof}

\begin{lemma}[Semigroup bounds on the common form domain]
\label{lem:e8v14-semigroup-Q}
For \(t\geq0\),
\begin{align}
 \|e^{-tH_0}u\|_{\mathcal Q,0}
 &\leq\|u\|_{\mathcal Q,0},
 \label{eq:e8v14-S0Q}\\
 \|e^{-tH_1}u\|_{\mathcal Q,0}
 &\leq M_\lambda\|u\|_{\mathcal Q,0}.
 \label{eq:e8v14-S1Q}
\end{align}
\end{lemma}

\begin{proof}
Spectral calculus gives
\[
 \|(H_j+\lambda)^{1/2}e^{-tH_j}u\|
 \leq\|(H_j+\lambda)^{1/2}u\|.
\]
For \(j=0\), this is \eqref{eq:e8v14-S0Q}.  For \(j=1\), compare the two
form norms before and after applying the semigroup using
\eqref{eq:e8v14-equivalence}; the ratio is at most
\(C_\lambda/c_\lambda=M_\lambda^2\).
\end{proof}

\subsection{Weak form Duhamel identity}
\label{sec:e8v14-Duhamel}

\begin{theorem}[Common-domain weak Duhamel formula]
\label{thm:e8v14-Duhamel}
For \(u,v\in\mathcal Q\) and \(t\geq0\),
\begin{equation}
 \langle u,(e^{-tH_1}-e^{-tH_0})v\rangle
 =-\int_0^t
 w[e^{-(t-s)H_1}u,e^{-sH_0}v],ds.
 \label{eq:e8v14-Duhamel}
\end{equation}
Consequently,
\begin{equation}
 \boxed{
 |\langle u,(e^{-tH_1}-e^{-tH_0})v\rangle|
 \leq t\gamma_\lambda M_\lambda
 \|u\|_{\mathcal Q,0}\|v\|_{\mathcal Q,0}.}
 \label{eq:e8v14-Duhamel-bound}
\end{equation}
\end{theorem}

\begin{proof}
First take \(u,v\) in a common operator core and differentiate
\[
 F(s):=\langle e^{-(t-s)H_1}u,e^{-sH_0}v\rangle.
\]
The weak derivative is
\[
 F'(s)=q_1[e^{-(t-s)H_1}u,e^{-sH_0}v]
       -q_0[e^{-(t-s)H_1}u,e^{-sH_0}v]
       =w[e^{-(t-s)H_1}u,e^{-sH_0}v].
\]
Since
\(F(t)-F(0)=
\langle u,e^{-tH_0}v\rangle-
\langle u,e^{-tH_1}v\rangle\), integration gives
\eqref{eq:e8v14-Duhamel}.  Approximation in the common form norm extends
it to all \(u,v\in\mathcal Q\).  Finally apply
\eqref{eq:e8v14-bilinear},
\eqref{eq:e8v14-S0Q}, and
\eqref{eq:e8v14-S1Q} inside the integral.
\end{proof}

\begin{remark}[Weak rather than operator-norm control]
\label{rem:e8v14-weak}
The estimate is uniform on bounded sets of the form norm.  It does not give
\(\|e^{-tH_1}-e^{-tH_0}\|=O(t)\) on the Hilbert unit ball, because vectors
of unbounded form energy are not controlled by
\eqref{eq:e8v14-Duhamel-bound}.
\end{remark}

\subsection{Local energy moments}
\label{sec:e8v14-moments}

Let \(\rho\) be a positive trace-class operator with spectral decomposition
\(\rho=\sum_r p_r|\psi_r\rangle\langle\psi_r|\).  For an operator \(A\)
which maps every \(\psi_r\) into \(\mathcal Q\), define
\begin{equation}
 \mathcal E_{0,\lambda}(\rho;A)
 :=\sum_rp_r\|A\psi_r\|_{\mathcal Q,0}^2.
 \label{eq:e8v14-energy-moment}
\end{equation}
When the trace is well defined, this is
\(\Tr(\rho A^*(H_0+\lambda)A)\).

\begin{corollary}[Transfer expectation bound]
\label{cor:e8v14-expectation}
If \(\mathcal E_{0,\lambda}(\rho;A)<\infty\), then
\begin{equation}
 \left|\Tr\left(
 \rho A^*(e^{-tH_1}-e^{-tH_0})A\right)\right|
 \leq t\gamma_\lambda M_\lambda
       \mathcal E_{0,\lambda}(\rho;A).
 \label{eq:e8v14-expectation}
\end{equation}
\end{corollary}

\begin{proof}
Apply \eqref{eq:e8v14-Duhamel-bound} with
\(u=v=A\psi_r\), multiply by \(p_r\), and sum.  Absolute convergence
follows from the assumed energy moment.
\end{proof}

\begin{firewall}[The energy moment is a new row]
\label{fire:e8v14-moment}
A bounded local observable need not have a cutoff-uniform value of
\(\mathcal E_{0,\lambda}(\rho;A)\) when the generator is unbounded.  The
clock factor in \eqref{eq:e8v14-expectation} is useful cofinally only after
this local energy moment, or a renormalized version of it, is bounded and
shown uniformly integrable.
\end{firewall}

\subsection{Uniform bounded truncations and Mosco passage}
\label{sec:e8v14-truncation}

Let \(q_{j,N}\) be closed nonnegative forms on Hilbert spaces transported
into \(\mathcal H\), with common domains \(\mathcal Q_N\), and let
\(H_{j,N}\) be their operators.  Suppose:
\begin{enumerate}
\item the constants \(a,b,\lambda\) in
\eqref{eq:e8v14-relative} are independent of \(N\);
\item \(q_{j,N}\) Mosco-converges to \(q_j\) for \(j=0,1\); and
\item test vectors \(u_N,v_N\) converge strongly to \(u,v\) and have
uniformly bounded transported \(\mathcal Q_N\)-norms.
\end{enumerate}

\begin{theorem}[Uniform weak truncation passage]
\label{thm:e8v14-truncation}
Every truncation obeys
\begin{equation}
 |\langle u_N,(e^{-tH_{1,N}}-e^{-tH_{0,N}})v_N\rangle|
 \leq t\gamma_\lambda M_\lambda
 \|u_N\|_{\mathcal Q_N,0}\|v_N\|_{\mathcal Q_N,0}.
 \label{eq:e8v14-uniform-truncation}
\end{equation}
The two matrix elements converge to their limiting semigroup matrix
elements, and the same upper bound passes to the limit with the liminf of
the right side.
\end{theorem}

\begin{proof}
Apply \Cref{thm:e8v14-Duhamel} at each \(N\); uniformity of the relative
form constants gives the displayed common coefficient.  Mosco convergence
of closed nonnegative forms implies strong convergence of their semigroups
at every \(t>0\).  Combine this with strong convergence of the test vectors
to pass the two matrix elements to the limit.  Taking the liminf in the
uniform inequality proves the last statement.
\end{proof}

\begin{corollary}[Bounded truncation norms may diverge harmlessly]
\label{cor:e8v14-norm-divergence}
The operator norms of bounded cutoff perturbations may diverge with
\(N\) while \eqref{eq:e8v14-uniform-truncation} remains uniform, provided
the relative form constants and the test-vector energy moments remain
uniform.  Thus an operator-norm Duhamel failure does not by itself prevent a
weak local transfer comparison.
\end{corollary}

\begin{proof}
The coefficient in \eqref{eq:e8v14-uniform-truncation} contains no cutoff
operator norm.  It contains only \(a,b,\lambda\) and the form norms of the
test vectors.
\end{proof}

\subsection{Two marked form insertions}
\label{sec:e8v14-two}

A symmetric form \(w\) defines a bounded map
\(\widehat w:\mathcal Q\to\mathcal Q^*\).  Two ordered insertions require
the semigroup to return a form-dual vector to the form domain.  Assume
\begin{equation}
 \|e^{-tH}\|_{\mathcal Q^*\to\mathcal Q}\leq g(t),
 \qquad 0<t\leq T,
 \label{eq:e8v14-smoothing}
\end{equation}
where
\begin{equation}
 G_T:=\int_0^T(T-r)g(r)\,dr<\infty.
 \label{eq:e8v14-GT}
\end{equation}

\begin{lemma}[Ordered form-insertion bound]
\label{lem:e8v14-two}
If \(w_1,w_2\) have
\(\|\widehat w_j\|_{\mathcal Q\to\mathcal Q^*}\leq\eta_j\) and the
outer semigroup factors are bounded by \(M\) on \(\mathcal Q\), then the
weak ordered simplex term satisfies
\begin{equation}
 \left|\int_{\substack{r,s\geq0\\r+s\leq T}}
 \langle u,S(T-r-s)\widehat w_1S(r)
              \widehat w_2S(s)v\rangle\,dr\,ds\right|
 \leq M^2\eta_1\eta_2G_T
 \|u\|_{\mathcal Q}\|v\|_{\mathcal Q}.
 \label{eq:e8v14-two}
\end{equation}
\end{lemma}

\begin{proof}
Read the pairing successively from the right.  The first outer semigroup is
bounded on \(\mathcal Q\), \(\widehat w_2\) maps to \(\mathcal Q^*\), the
middle semigroup costs \(g(r)\), \(\widehat w_1\) maps back to the dual,
and the left pairing costs the second factor \(M\).  Thus the integrand is
bounded by
\(M^2\eta_1\eta_2g(r)\|u\|_{\mathcal Q}\|v\|_{\mathcal Q}\).
For fixed \(r\), the remaining simplex coordinate has length \(T-r\);
integration gives \eqref{eq:e8v14-GT}.
\end{proof}

\begin{firewall}[Smoothing is not supplied by Mosco convergence]
\label{fire:e8v14-smoothing}
Strong semigroup convergence does not imply the form-dual smoothing row
\eqref{eq:e8v14-smoothing}, much less its spatially local marked version.
The two-insertion stress passage therefore needs both the energy moment and
a uniform integrable smoothing kernel.  V14 proves only the consequence of
those hypotheses.
\end{firewall}

\begin{assessment}[Progress at V14]
\label{ass:e8v14-status}
The bounded-generator clock factor has a rigorous unbounded analogue at the
level of local transfer matrix elements.  Uniform relative form bounds and
local form-energy moments replace a divergent operator norm, and Mosco
convergence passes the weak comparison to the limit.  Two stress insertions
add a form-dual smoothing obligation which cannot be inferred from Mosco
convergence alone.
\end{assessment}

\begin{programme}[Eighth-series V15: companion audit and form-locality
ledger]
\label{prog:e8v14-V15}
V15 will perform the compulsory companion audit, then combine the weak form
Duhamel estimate, local energy moment, smoothing kernel, and finite-horizon
spatial rows into one form-locality acquisition card.  It will compare that
card with the companion reflected-transfer work and retain the distinction
between state projectivity and a physical spectral gap.
\end{programme}

\begin{firewall}[Claim boundary after eighth-series V14]
\label{fire:e8v14-boundary}
V14 proves common form-norm equivalence, a weak Duhamel formula with a
linear clock factor, a local transfer expectation bound, uniform passage
through Mosco-convergent truncations, and a conditional two-form-insertion
smoothing estimate.

V14 does not prove the required SM relative form bounds, local energy
moments, form-dual smoothing or spatial locality, identify the physical
reflected transfer, remove the cutoff, or construct the Standard
Model--Einstein continuum.
\end{firewall}

\begin{theorem}[Eighth-series V14 result]
\label{thm:e8v14-result}
For two nonnegative self-adjoint generators with a common form domain and
uniform relative form bound \eqref{eq:e8v14-relative}, their transfer
matrix elements differ by at most the clock-scaled form-energy bound
\eqref{eq:e8v14-Duhamel-bound}.  The estimate is uniform under
Mosco-convergent bounded truncations and yields local expectation control
when \eqref{eq:e8v14-energy-moment} is finite.  Two marked form insertions
are controlled only after the additional integrable form-dual smoothing
row \eqref{eq:e8v14-smoothing} is proved.
\end{theorem}

\begin{proof}
Combine \Cref{lem:e8v14-equivalence,lem:e8v14-semigroup-Q,
thm:e8v14-Duhamel,cor:e8v14-expectation,
thm:e8v14-truncation,lem:e8v14-two}.
\end{proof}

\paragraph{Parent-version record.}
V14 was formed from
\texttt{Renewal\_SM\_Einstein\_Continuum\_Research\_Notes\_V13.tex}.
The V13 parent had \(4\,620\) logical source lines, \(171\,099\) bytes,
and SHA--256
\[
 \texttt{DBA26066FAD8ADB6F4B27736B9E6DE152ECF18DE150A1F8ABC7FF762298FE693}.
\]
The next compulsory companion audit is V15.

% =============================================================================
% END APPENDED EIGHTH-SERIES V14 MATERIAL
% =============================================================================

% =============================================================================
% BEGIN APPENDED EIGHTH-SERIES V15 MATERIAL
% =============================================================================

\section{Companion audit and the form-locality acquisition card}
\label{sec:v8v15}

V14 replaced cutoff operator norms by common-domain form estimates.  The
present appendix performs the compulsory companion audit and assembles the
remaining hypotheses into one checkable card.  Its main result is a
conditional (C^1) landing theorem for local transfer matrix elements.  A
separate signed-coercivity firewall prevents these absolute estimates from
being misread as a spectral-gap proof.

\subsection{Compulsory V15 companion audit}
\label{sec:e8v15-audit}

The newest coherent companion files read for this audit were
\begin{align*}
 &\texttt{Renewal\_Yang\_Mills\_Mass\_Gap\_Research\_Notes\_V67.tex},\\
 &\texttt{Renewal\_NS\_Stability\_Research\_Notes\_V26.tex}.
\end{align*}
At the completed read, the Yang--Mills file had (30\,301) logical source
lines, (1\,064\,104) bytes, one live document terminator, and SHA--256
\[
 \texttt{734FCAA1214F0575D0A66ECA6AA6EB7AC7A62A99A9001EBE4B5C29B0A56005D7}.
\]
The Navier--Stokes file had (5\,319) logical source lines,
(278\,283) bytes, one live document terminator, and SHA--256
\[
 \texttt{82C887F8C2C69E4F28FAD7C3DC8DE5B9099CB5A4BF431EBA1FFF3E91935978AD}.
\]
The Yang--Mills file was updated concurrently during the first read; the
metadata above belong to the completed V67 appendix, not to its transient
V66 prefix.

\begin{assessment}[Yang--Mills V67 reading]
\label{ass:e8v15-YM}
The companion note identifies normalized transfer loss with a conditional
variance and polarizes its signed derivative into a centered
conditional-mean multiplier plus a within-fibre covariance.  Both pieces
can change sign.  It proves that a conditional-mean range inequality,
together with a conditional variance of clock order, would imply a direct
signed G\aa rding inequality.  It also gives a dense-range negative-direction
obstruction for output-only insertions.  Thus absolute marked-path or
locality majorants can control the size of the response but cannot supply
its positive sign.
\end{assessment}

\begin{assessment}[Navier--Stokes V26 reading]
\label{ass:e8v15-NS}
The companion note computes rank-one norms of the first two derivatives of
the Oseen kernel and quantifies a modest improvement in one radial range.
Those constants belong to its ancient-solution kernel ledger.  They do not
populate a Standard Model form bound, local energy moment, smoothing
kernel, recovery estimate, or reflected-transfer incidence row, so no
numerical or analytic input is imported here.
\end{assessment}

\begin{theorem}[V15 audit decision]
\label{thm:e8v15-audit}
The only imported conclusion is the following logical separation:
\begin{equation}
 \boxed{\text{absolute form locality and Mosco landing}
 \quad\not\Longrightarrow\quad
 \text{signed vacuum-orthogonal coercivity}.}
 \label{eq:e8v15-separation}
\end{equation}
No companion constant or model-specific estimate is transferred into the
SM--Einstein ledger.
\end{theorem}

\begin{proof}
The Yang--Mills covariance formula and its negative-direction theorem show
that the sign of a centered insertion is information not contained in an
absolute majorant.  The Navier--Stokes estimates concern a different kernel
and a different function space.  Hence only the stated proof obligation is
common to the programmes; their quantitative rows are not.
\end{proof}

\subsection{A common-form comparison with one parameter derivative}
\label{sec:e8v15-C1-form}

Let (U) be an open subset of a real Banach space (X).  For
(j\in\{0,1\}) and (x\in U), let (h_j(x)) be a densely defined,
closed, nonnegative symmetric form on a Hilbert space \(\mathcal H\), with
one common form domain \(\mathcal Q\).  Equip \(\mathcal Q\) with a fixed
Hilbert norm \(\|\cdot\|_{\mathcal Q}\), and write
\[
 S_j(x,t)=e^{-tH_j(x)}.
\]
Assume that (x\mapsto h_j(x)) is (C^1) as a map into the bounded
symmetric forms on \(\mathcal Q\).  For a unit direction \(\xi\in X\), put
\begin{align}
 w(x)&=h_1(x)-h_0(x),
 &\dot w(x)[\xi]&=Dh_1(x)[\xi]-Dh_0(x)[\xi].
 \label{eq:e8v15-wdot}
\end{align}
On a compact convex chart (K\Subset U), suppose
\begin{align}
 \|S_j(x,t)\|_{\mathcal Q\to\mathcal Q}&\leq M,
 &&0\leq t\leq T,
 \label{eq:e8v15-Q-bound}\\
 \|S_j(x,t)\|_{\mathcal Q^*\to\mathcal Q}&\leq g(t),
 &&0<t\leq T,
 \label{eq:e8v15-smoothing}\\
 \|\widehat w(x)\|_{\mathcal Q\to\mathcal Q^*}&\leq\eta_0,
 &\|\widehat{\dot w}(x)[\xi]\|_{\mathcal Q\to\mathcal Q^*}
 &\leq\eta_1,
 \label{eq:e8v15-eta}\\
 \|\widehat{Dh_j(x)[\xi]}\|_{\mathcal Q\to\mathcal Q^*}
 &\leq\zeta.
 \label{eq:e8v15-zeta}
\end{align}
All constants are uniform in \(x\in K\), \(j\), and unit \(\xi\).  Define
\begin{equation}
 G_T(g):=\int_0^T(T-r)g(r)\,dr.
 \label{eq:e8v15-G}
\end{equation}

\begin{lemma}[Adjacent derivative comparison]
\label{lem:e8v15-adjacent-derivative}
If \(G_T(g)<\infty\), then for \(u,v\in\mathcal Q\), \(0\leq t\leq T\),
and \(\|\xi\|_X\leq1\),
\begin{align}
 &\left|D_x\langle u,(S_1(x,t)-S_0(x,t))v\rangle[\xi]\right|
 \notag\\
 &\qquad\leq
 M^2\left(t\eta_1+2\eta_0\zeta G_T(g)\right)
 \|u\|_{\mathcal Q}\|v\|_{\mathcal Q}.
 \label{eq:e8v15-adjacent-derivative}
\end{align}
Here \(u,v\) are fixed in the common transported Hilbert space.  Derivatives
of carriers or observables must be added as contact terms.
\end{lemma}

\begin{proof}
The weak form derivative is
\begin{equation}
 D_x\langle u,S_j(x,t)v\rangle[\xi]
 =-\int_0^t
 Dh_j(x)[\xi]
 [S_j(x,t-s)u,S_j(x,s)v]\,ds.
 \label{eq:e8v15-form-derivative}
\end{equation}
It follows first for vectors in an operator core by differentiating the
semigroup equation, and then on \(\mathcal Q\) by the uniform form bounds.
Subtract the \(j=0\) identity from the \(j=1\) identity and telescope each
integrand into three terms: the derivative-form difference with the two
outer semigroups fixed, the difference of the left semigroup, and the
difference of the right semigroup.  The first term is at most
\(M^2\eta_1\|u\|_{\mathcal Q}\|v\|_{\mathcal Q}\), and integration gives
\(tM^2\eta_1\).

Insert the weak Duhamel formula for \(S_1-S_0\) in each of the other two
terms.  Each resulting ordered term contains one copy of
\(\widehat w\), one copy of \(\widehat{Dh_j[\xi]}\), and a semigroup between
them acting from \(\mathcal Q^*\) to \(\mathcal Q\).  The two outer
semigroups cost \(M^2\), the insertions cost \(\eta_0\zeta\), and the
ordered simplex integral costs \(G_T(g)\), exactly as in the V14 ordered
form-insertion lemma.  There are two orders.  Adding all three estimates
gives \eqref{eq:e8v15-adjacent-derivative}.
\end{proof}

\begin{remark}[Why the smoothing row occurs only in the mixed term]
\label{rem:e8v15-smoothing-role}
The single derivative-form difference is paired directly between two form
vectors.  A base-form insertion followed by a derivative-form insertion
creates a vector in \(\mathcal Q^*\) between the two marks; returning it to
\(\mathcal Q\) is precisely the role of \eqref{eq:e8v15-smoothing}.
\end{remark}

\subsection{The form-locality acquisition card}
\label{sec:e8v15-card}

Fix a countable local observable core \(\mathfrak A_{\rm loc}^{(0)}\), a
compact parameter chart \(K\), and a time horizon \(T\).  Adjacent scales
\(n,n+1\) are transported to one comparison Hilbert space before the card
is evaluated.

\begin{definition}[Form-locality acquisition card, FLAC]
\label{def:e8v15-FLAC}
The FLAC at adjacent scale \(n\) consists of the following verified rows.
\begin{enumerate}[label=\textup{(F\arabic*)},leftmargin=4.5em]
\item \emph{Common domain.}  The two transported nonnegative forms and
their parameter derivatives act on a common Hilbert form domain
\(\mathcal Q_n\).
\item \emph{Relative base row.}  Their difference obeys the V14 relative
form estimate with uniform \(a_n<1,b_n,\lambda_n>b_n\); write
\(\gamma_n,M_n\) for the resulting constants.
\item \emph{Local energy row.}  For every
\(A\in\mathfrak A_{\rm loc}^{(0)}\), the transported carrier has
\begin{equation}
 E_{n,A}:=\sup_{x\in K}\sum_rp_{n,r}(x)
 \|A_n(x)\psi_{n,r}(x)\|_{\mathcal Q_n}^2<\infty.
 \label{eq:e8v15-energy-row}
\end{equation}
\item \emph{One- and two-mark rows.}  The quantities
\(\eta_{n,0},\eta_{n,1},\zeta_n,M_n^{Q}\) and \(g_n\) satisfy
\eqref{eq:e8v15-Q-bound}--\eqref{eq:e8v15-zeta}, with
\(G_T(g_n)<\infty\).
\item \emph{Spatial row.}  The base and derivative forms admit local
component decompositions whose one- and two-mark matrix elements obey
summable near, far, and boundary majorants.  In the bounded Hamiltonian and
Lindblad subrows these are the V12 and V13 non-reset tree bounds; in an
unbounded sector their form analogues must be proved rather than inferred.
\item \emph{Contact row.}  Changes of carrier, observable transport,
recovery map, boundary prescription, and parameter chart have scalar
bounds (c_{n,A}^{(0)}) and (c_{n,A}^{(1)}).
\item \emph{Cofinal summability.}  The scalar defects
\begin{align}
 d_{n,A}^{(0)}(T)
 &:=T\gamma_nM_nE_{n,A}+c_{n,A}^{(0)},
 \label{eq:e8v15-d0}\\
 d_{n,A}^{(1)}(T)
 &:=(M_n^{Q})^2
 \left(T\eta_{n,1}
 +2\eta_{n,0}\zeta_nG_T(g_n)\right)E_{n,A}
 +c_{n,A}^{(1)}
 \label{eq:e8v15-d1}
\end{align}
obey
\begin{equation}
 \sum_n d_{n,A}^{(0)}(T)<\infty,
 \qquad
 \sum_n d_{n,A}^{(1)}(T)<\infty.
 \label{eq:e8v15-sums}
\end{equation}
The spatial component sums from (F5) are included in the displayed scalar
rows, not counted a second time.
\item \emph{Landing and incidence.}  The common-refinement bridges are
compatible on the local core.  If a spectral or Lorentzian conclusion is
sought, the limiting form must additionally be identified with the
physical reflected-time transfer or Hamiltonian incidence.
\end{enumerate}
\end{definition}

\begin{remark}[The card is an obligation list]
\label{rem:e8v15-card-obligation}
Rows (F1)--(F8) are not asserted for the full SM--Einstein regulator
family.  The purpose of the card is to prevent a bounded finite-cutoff
locality estimate, a Mosco limit, or a recovery estimate from silently
standing in for the other rows.
\end{remark}

\subsection{A conditional (C^1) landing theorem}
\label{sec:e8v15-landing}

For \(A\in\mathfrak A_{\rm loc}^{(0)}\), let
\begin{equation}
 F_{n,A}(x,t)
 :=\sum_rp_{n,r}(x)
 \langle A_n(x)\psi_{n,r}(x),
 S_n(x,t)A_n(x)\psi_{n,r}(x)\rangle.
 \label{eq:e8v15-F}
\end{equation}
The comparison bridges are understood in this notation.  Suppose their
contact rows have been chosen so that
\begin{align}
 \sup_{(x,t)\in K\times[0,T]}
 |F_{n+1,A}(x,t)-F_{n,A}(x,t)|
 &\leq d_{n,A}^{(0)}(T),
 \label{eq:e8v15-F-defect}\\
 \sup_{(x,t)\in K\times[0,T]}
 \|D_xF_{n+1,A}(x,t)-D_xF_{n,A}(x,t)\|_{X^*}
 &\leq d_{n,A}^{(1)}(T).
 \label{eq:e8v15-DF-defect}
\end{align}

\begin{theorem}[FLAC (C^1) landing]
\label{thm:e8v15-C1-landing}
If the FLAC is verified from some scale onward, then for every local core
observable (A) there is a unique function
\[
 F_A\in C^1(K;C([0,T]))
\]
such that
\begin{align}
 \|F_A-F_{N,A}\|_{C(K\times[0,T])}
 &\leq\sum_{n\geq N}d_{n,A}^{(0)}(T),
 \label{eq:e8v15-tail0}\\
 \|D_xF_A-D_xF_{N,A}\|_{C(K\times[0,T];X^*)}
 &\leq\sum_{n\geq N}d_{n,A}^{(1)}(T).
 \label{eq:e8v15-tail1}
\end{align}
The limit is independent of a cofinal subsequence and of a comparison
route whose additional bridge defects have summable tails.
\end{theorem}

\begin{proof}
By \eqref{eq:e8v15-F-defect} and the first series in
\eqref{eq:e8v15-sums}, (F_{n,A}) is uniformly Cauchy on
(K\times[0,T]); call its limit (F_A).  By
\eqref{eq:e8v15-DF-defect} and the second series, (D_xF_{n,A}) is
uniformly Cauchy there; call its limit (G_A).

For (x,y\in K), convexity and the Banach-space fundamental theorem of
calculus give
\[
 F_{n,A}(y,t)-F_{n,A}(x,t)
 =\int_0^1D_xF_{n,A}(x+s(y-x),t)[y-x],ds.
\]
Uniform convergence permits passage to the limit inside the integral:
\[
 F_A(y,t)-F_A(x,t)
 =\int_0^1G_A(x+s(y-x),t)[y-x],ds.
\]
Continuity of (G_A) follows from its uniform convergence.  The last
identity implies that (F_A) is Fr\'echet differentiable with derivative
(G_A): subtract (G_A(x,t)[y-x]), bound the remainder by the modulus of
continuity of \(G_A\) along the segment, and divide by \(\|y-x\|\).
The estimates \eqref{eq:e8v15-tail0}--\eqref{eq:e8v15-tail1} are the
corresponding telescoping tails.

Every cofinal subsequence of a Cauchy sequence has the same limit.  A
second comparison route differs by its bridge telescoping series; a
summable tail tends to zero, so the two limits agree.
\end{proof}

\begin{corollary}[Positive local functional landing]
\label{cor:e8v15-positive-landing}
Suppose in addition that (F_{n,A}(x,0)) comes from normalized positive
finite-regulator states and that the comparison maps preserve the local
unit and involution up to the summable contact row.  Then the (t=0)
limit extends by continuity from the local core to a normalized positive
functional on the inductive local (C^*)-algebra whenever the local core
is norm dense and the finite functionals are uniformly contractive.
\end{corollary}

\begin{proof}
At (t=0), uniform convergence gives a well-defined compatible functional
on the local core.  For every core element (B), positivity follows from
\(\omega(B^*B)=\lim_n\omega_n(B_n^*B_n)\geq0\), while normalization follows
from the unit contact tail.  Uniform contractivity gives
\(|\omega(B)|\leq\|B\|\), so the functional extends uniquely and
continuously to the norm completion.  Positivity and normalization persist
under that extension.
\end{proof}

\begin{corollary}[How V14 populates the scalar rows]
\label{cor:e8v15-V14-population}
Under rows (F1)--(F4), V14 gives the generator contribution
(T\gamma_nM_nE_{n,A}) in \eqref{eq:e8v15-d0};
\Cref{lem:e8v15-adjacent-derivative} gives the displayed generator
contribution in \eqref{eq:e8v15-d1}.  Once rows (F5) and (F6) bound all
component and contact terms, the summability test \eqref{eq:e8v15-sums}
is sufficient for \Cref{thm:e8v15-C1-landing}.
\end{corollary}

\begin{proof}
Apply the V14 transfer expectation bound to each carrier vector and sum
with weights (p_{n,r}), which produces the energy moment
\eqref{eq:e8v15-energy-row}.  Apply
\Cref{lem:e8v15-adjacent-derivative} in the same way for the parameter
derivative.  The remaining differentiated carriers, observables,
boundaries, and bridges are exactly the contact terms by definition.
Equations \eqref{eq:e8v15-F-defect}--\eqref{eq:e8v15-DF-defect} follow,
so the landing theorem applies when their tails are summable.
\end{proof}

\subsection{Signed-coercivity firewall and next upstream row}
\label{sec:e8v15-sign}

\begin{proposition}[Absolute rows do not determine a signed response]
\label{prop:e8v15-sign-counterexample}
For every \(\varepsilon>0\) there are bounded symmetric perturbation forms
(w_+) and (w_-) on the same one-dimensional form domain such that
\[
 |w_+[u,v]|=|w_-[u,v]|
 \leq\varepsilon\|u\|_{\mathcal Q}\|v\|_{\mathcal Q}
\]
for all (u,v), while (w_+\succeq0) and (w_-\preceq0).  Consequently,
no theorem whose hypotheses use only the absolute FLAC rows can infer a
positive spectral derivative.
\end{proposition}

\begin{proof}
Take \(\mathcal Q=\mathbb C\) with its usual norm and set
\[
 w_+[u,v]=\varepsilon\overline u v,
 \qquad
 w_-[u,v]=-\varepsilon\overline u v.
\]
Their absolute bilinear bounds are identical and their quadratic signs are
opposite.  Any conclusion determined solely by those absolute bounds would
have to be the same for both forms, so it cannot assert positivity.
\end{proof}

\begin{firewall}[What a gap claim would additionally require]
\label{fire:e8v15-gap}
The FLAC proves a conditional continuum landing statement for local
transfer data.  A physical mass or transfer gap additionally requires:
\begin{enumerate}[label=\textup{(\roman*)},leftmargin=3.5em]
\item a signed vacuum-orthogonal form inequality with a cofinal positive
stock;
\item control of the vacuum projector or Ward quotient used in that
inequality;
\item identification of the form clock with the physical reflected-time
incidence; and
\item passage of the signed lower bound to the limiting operator.
\end{enumerate}
Neither spatial decay nor Mosco convergence supplies item (i).
\end{firewall}

\begin{assessment}[Progress at V15]
\label{ass:e8v15-status}
The unbounded-form route is now organized as a finite acquisition card.
The new adjacent-derivative estimate shows exactly where form-dual
smoothing enters, and the (C^1) landing theorem converts summable base and
stress defects into route-independent continuum local transfer data.  The
companion audit confirms that signed spectral coercivity remains a separate
row.  The first upstream unresolved FLAC entry is the cutoff-uniform local
energy moment \eqref{eq:e8v15-energy-row}; without it, even the one-mark V14
bound cannot be summed.
\end{assessment}

\begin{programme}[Eighth-series V16: acquire the local energy row]
\label{prog:e8v15-V16}
V16 will attack \eqref{eq:e8v15-energy-row} upstream.  It will derive
sufficient local energy-moment criteria for Gibbs or Euclidean-slab
carriers from a global energy moment plus commutator control, identify the
volume term that must be subtracted or localized, and test which criterion
is stable under the SM gauge, Higgs, fermion, and gravitational cutoff
bridges.  Only estimates actually uniform in volume and refinement will be
entered into the FLAC.
\end{programme}

\begin{firewall}[Claim boundary after eighth-series V15]
\label{fire:e8v15-boundary}
V15 records the compulsory YM/NS audit; proves an adjacent (C^1) weak
form comparison; defines the form-locality acquisition card; proves its
conditional (C^1), cofinal, route-independent landing theorem; and proves
that its absolute rows cannot imply signed spectral positivity.

V15 does not verify the FLAC for the Standard Model--Einstein regulator
family, prove a cutoff-uniform local energy moment or smoothing kernel,
construct the limiting state or Hamiltonian, establish reflection
positivity, prove a signed gap inequality, or construct the full
Standard Model--Einstein continuum.
\end{firewall}

\begin{theorem}[Eighth-series V15 result]
\label{thm:e8v15-result}
For (C^1) adjacent common-domain form families, the difference of their
semigroup derivatives obeys
\eqref{eq:e8v15-adjacent-derivative}.  If all FLAC rows hold with summable
base and derivative defects, local transfer matrix elements converge in
(C^1), with the explicit tail bounds
\eqref{eq:e8v15-tail0}--\eqref{eq:e8v15-tail1}.  This conclusion contains
no spectral-sign assertion.
\end{theorem}

\begin{proof}
Combine \Cref{lem:e8v15-adjacent-derivative,
thm:e8v15-C1-landing,cor:e8v15-V14-population,
prop:e8v15-sign-counterexample}.
\end{proof}

\paragraph{Parent-version record.}
V15 was formed from
\texttt{Renewal\_SM\_Einstein\_Continuum\_Research\_Notes\_V14.tex}.
The V14 parent had (4\,979) logical source lines, (184\,018) bytes,
and SHA--256
\[
 \texttt{0EB5B4E34244F97E4A080624BBF80FC602FE0983DFF5289EDCC8A686F753B4A3}.
\]
The next compulsory companion audit is V20.

% =============================================================================
% END APPENDED EIGHTH-SERIES V15 MATERIAL
% =============================================================================


% =============================================================================
% BEGIN APPENDED EIGHTH-SERIES V16 MATERIAL
% =============================================================================

\section{Acquiring the local form-energy row from a vacuum commutator}
\label{sec:v8v16}

The FLAC introduced in V15 requires a cutoff-uniform energy moment for every
fixed local source.  The Grand Tensor manuscript already proves that a finite
source bank satisfying (K\preceq E_*G_0) has a delayed semigroup floor and a
quantitative collar modulus.  Repeating that theorem would not advance the
programme.  This appendix supplies its upstream input: it derives the
source-local energy matrix from the commutator of the physical positive
Hamiltonian with the local source.  It also identifies an extensive background
term which makes the raw mixed-state trace moment unsuitable for a
thermodynamic limit.

\subsection{The exact vacuum-commutator identity}
\label{sec:e8v16-commutator}

Let (H\geq0) be self-adjoint on a Hilbert space \(\mathcal H\), and let
\(\Omega\in\operatorname{Dom}H\) be normalized with
\begin{equation}
 H\Omega=0.
 \label{eq:e8v16-vacuum}
\end{equation}
For a bounded operator (A), say that the vacuum commutator
\(\delta_H(A)\Omega\) exists if \(A\Omega\in\operatorname{Dom}H\) and
\begin{equation}
 \delta_H(A)\Omega:=HA\Omega-AH\Omega=HA\Omega.
 \label{eq:e8v16-delta}
\end{equation}

\begin{lemma}[Exact local energy identity]
\label{lem:e8v16-energy-identity}
Under \eqref{eq:e8v16-vacuum}--\eqref{eq:e8v16-delta},
\begin{equation}
 \boxed{
 \|H^{1/2}A\Omega\|^2
 =\langle A\Omega,\delta_H(A)\Omega\rangle.}
 \label{eq:e8v16-energy-identity}
\end{equation}
The right side is real and nonnegative, although that need not be apparent
from a termwise commutator expansion.
\end{lemma}

\begin{proof}
Since \(A\Omega\in\operatorname{Dom}H\),
\[
 \|H^{1/2}A\Omega\|^2
 =\langle A\Omega,HA\Omega\rangle.
\]
Equation \eqref{eq:e8v16-delta} identifies the second vector with
\(\delta_H(A)\Omega\).  Positivity follows from \(H\geq0\).
\end{proof}

\begin{corollary}[Scalar local energy bound]
\label{cor:e8v16-scalar}
If
\begin{equation}
 \|\delta_H(A)\Omega\|\leq D_H(A),
 \label{eq:e8v16-D}
\end{equation}
then for every \(\lambda>0\),
\begin{equation}
 \boxed{
 \|A\Omega\|_{\operatorname{Dom}(H+\lambda)^{1/2}}^2
 \leq \|A\|D_H(A)+\lambda\|A\|^2.}
 \label{eq:e8v16-scalar-bound}
\end{equation}
\end{corollary}

\begin{proof}
Apply Cauchy--Schwarz to \eqref{eq:e8v16-energy-identity} and use
\(\|A\Omega\|\leq\|A\|\).  Add
\(\lambda\|A\Omega\|^2\).
\end{proof}

\begin{remark}[Form-core version]
\label{rem:e8v16-form-core}
The operator-domain hypothesis is sufficient, not necessary.  If there are
bounded approximants \(A_m\) with \(A_m\Omega\in\operatorname{Dom}H\),
\(A_m\Omega\to A\Omega\), and uniformly bounded right sides in
\eqref{eq:e8v16-energy-identity}, closedness of (H^{1/2}) places
\(A\Omega\) in the form domain and preserves the bound.  This is the route
available for unbounded local field generators after a common-core
truncation.
\end{remark}

\subsection{A finite source-bank compiler}
\label{sec:e8v16-bank}

Let \(E\) be finite dimensional and let
\(\Psi:E\to\mathcal H\) have the form
\begin{equation}
 \Psi c=A(c)\Omega,
 \label{eq:e8v16-Psi}
\end{equation}
where \(A:E\to\mathcal B(\mathcal H)\) is linear and every
\(A(c)\Omega\) belongs to \(\operatorname{Dom}H\).  Define
\begin{align}
 G_0&:=\Psi^*\Psi,
 &K&:=\Psi^*H\Psi,
 \label{eq:e8v16-GK}\\
 \Delta c&:=\delta_H(A(c))\Omega.
 \label{eq:e8v16-Delta}
\end{align}

\begin{theorem}[Vacuum-commutator source-energy compiler]
\label{thm:e8v16-bank-compiler}
Suppose
\begin{equation}
 \|\Delta c\|\leq D\|\Psi c\|,
 \qquad c\in E.
 \label{eq:e8v16-relative-commutator}
\end{equation}
Then
\begin{equation}
 \boxed{0\preceq K\preceq D G_0.}
 \label{eq:e8v16-bank-energy}
\end{equation}
Consequently this source bank populates the source-local energy row with
\(E_*=D\), and the previously proved source-local semigroup theorem applies
to it.
\end{theorem}

\begin{proof}
For (c\in E), \Cref{lem:e8v16-energy-identity} gives
\[
 K[c]=\langle\Psi c,\Delta c\rangle.
\]
It is nonnegative, and \eqref{eq:e8v16-relative-commutator} gives
\[
 K[c]\leq\|\Psi c\|\|\Delta c\|
 \leq D\|\Psi c\|^2=D G_0[c].
\]
This is the matrix inequality \eqref{eq:e8v16-bank-energy}.
\end{proof}

\begin{corollary}[Coefficient estimate and its conditioning debit]
\label{cor:e8v16-conditioning}
Fix a norm \(\|\cdot\|_E\) on (E).  If
\begin{equation}
 \|\Delta c\|\leq d\|c\|_E,
 \qquad
 G_0[c]\geq m\|c\|_E^2
 \label{eq:e8v16-coefficient}
\end{equation}
with (m>0), then \eqref{eq:e8v16-bank-energy} holds with
\(D=d/\sqrt m\).  A coefficientwise commutator estimate therefore gives a
cofinal source-energy stock only if the source Gram floor (m) does not
collapse.
\end{corollary}

\begin{proof}
The Gram floor gives
\(\|c\|_E\leq m^{-1/2}\|\Psi c\|\).  Substitute this in the first
inequality of \eqref{eq:e8v16-coefficient} and apply
\Cref{thm:e8v16-bank-compiler}.
\end{proof}

\begin{countertheorem}[Separate column bounds do not control the bank]
\label{cth:e8v16-columns}
Uniform bounds on each column energy and column norm do not imply a uniform
matrix inequality \(K\preceq E_*G_0\).
\end{countertheorem}

\begin{proof}
On \(\mathbb C^2\), let \(H=\operatorname{diag}(0,1)\),
\(\psi_1=e_1\), and
\(\psi_2=e_1+\varepsilon e_2\).  Both column norms are at most \(\sqrt2\)
and both column energies are at most one.  For \(c=(1,-1)\),
\[
 G_0[c]=\varepsilon^2,
 \qquad K[c]=\varepsilon^2.
\]
This example alone has ratio one.  Replace \(H\) by
\(H_\varepsilon=\operatorname{diag}(0,\varepsilon^{-1})\).  The individual
column energy remains \(\varepsilon\leq1\), whereas
\[
 {K[c]\over G_0[c]}={1\over\varepsilon}\longrightarrow\infty.
\]
Thus near cancellation of source columns can expose arbitrarily high energy;
a relative bank estimate or a stable Gram floor is essential.
\end{proof}

\subsection{Local interactions give volume-independent commutators}
\label{sec:e8v16-local-interaction}

Let \(\Lambda\) be a finite regulator region and suppose
\begin{equation}
 H_\Lambda=\sum_{Z\subset\Lambda}\Phi_\Lambda(Z)-E_{0,\Lambda}I,
 \qquad H_\Lambda\geq0,
 \qquad H_\Lambda\Omega_\Lambda=0.
 \label{eq:e8v16-local-H}
\end{equation}
The interaction terms may depend on the cutoff.  Let (A_B) be supported in
a fixed finite set (B), so
\begin{equation}
 [\Phi_\Lambda(Z),A_B]=0
 \quad\hbox{when }Z\cap B=\varnothing.
 \label{eq:e8v16-disjoint}
\end{equation}

\begin{theorem}[Bounded-interaction local energy row]
\label{thm:e8v16-bounded-local}
Assume all relevant \(\Phi_\Lambda(Z)\) are bounded and
\begin{equation}
 J_B:=\sup_\Lambda
 \sum_{Z\subset\Lambda:\,Z\cap B\neq\varnothing}
 \|\Phi_\Lambda(Z)\|<\infty.
 \label{eq:e8v16-JB}
\end{equation}
Then
\begin{equation}
 \boxed{
 \langle A_B\Omega_\Lambda,H_\Lambda A_B\Omega_\Lambda\rangle
 \leq2J_B\|A_B\|^2}
 \label{eq:e8v16-bounded-energy}
\end{equation}
uniformly in volume and cutoff.
\end{theorem}

\begin{proof}
The scalar ground-energy shift commutes with (A_B), and
\eqref{eq:e8v16-disjoint} gives
\[
 [H_\Lambda,A_B]
 =\sum_{Z:\,Z\cap B\neq\varnothing}
 [\Phi_\Lambda(Z),A_B].
\]
Hence
\[
 \|[H_\Lambda,A_B]\Omega_\Lambda\|
 \leq2\|A_B\|
 \sum_{Z:\,Z\cap B\neq\varnothing}\|\Phi_\Lambda(Z)\|
 \leq2J_B\|A_B\|.
\]
Apply \Cref{cor:e8v16-scalar} with \(\lambda=0\).
\end{proof}

\begin{corollary}[Exponentially summable interaction row]
\label{cor:e8v16-exponential}
If for some \(\mu>0\)
\begin{equation}
 \sup_{\Lambda,i\in\Lambda}
 \sum_{Z\ni i}e^{\mu\operatorname{diam}Z}
 |Z|\,\|\Phi_\Lambda(Z)\|\leq J_\mu,
 \label{eq:e8v16-Jmu}
\end{equation}
then \(J_B\leq |B|J_\mu\), and
\eqref{eq:e8v16-bounded-energy} is uniform for every fixed local support.
\end{corollary}

\begin{proof}
Every set counted in \eqref{eq:e8v16-JB} contains at least one (i\in B).
Sum first over that site and discard the weights, which are at least one.
The possible multiple counting only enlarges the bound.
\end{proof}

For unbounded local interactions, operator norms in \eqref{eq:e8v16-JB}
are unavailable.  The exact replacement is a vacuum commutator row.

\begin{theorem}[Unbounded form-local commutator criterion]
\label{thm:e8v16-unbounded-local}
For every fixed local (A_B), suppose there are vectors
\(d_{\Lambda,Z}(A_B)\in\mathcal H_\Lambda\) such that, on a common core,
\begin{align}
 [H_\Lambda,A_B]\Omega_\Lambda
 &=\sum_{Z:\,Z\cap B\neq\varnothing}d_{\Lambda,Z}(A_B),
 \label{eq:e8v16-form-decomposition}\\
 \sup_\Lambda\sum_{Z:\,Z\cap B\neq\varnothing}
 \|d_{\Lambda,Z}(A_B)\|&=:D_B(A_B)<\infty.
 \label{eq:e8v16-form-sum}
\end{align}
Assume the truncated vectors converge to (A_B\Omega_\Lambda) in the graph
closure needed in \Cref{rem:e8v16-form-core}.  Then
\begin{equation}
 \boxed{
 \|H_\Lambda^{1/2}A_B\Omega_\Lambda\|^2
 \leq\|A_B\|D_B(A_B)}
 \label{eq:e8v16-unbounded-energy}
\end{equation}
uniformly in \(\Lambda\).
\end{theorem}

\begin{proof}
The absolutely convergent vector sum gives
\(\|[H_\Lambda,A_B]\Omega_\Lambda\|\leq D_B(A_B)\).
Apply \Cref{cor:e8v16-scalar} to the common-core truncations and then use the
closed form passage of \Cref{rem:e8v16-form-core}.
\end{proof}

\begin{remark}[Sector reading]
\label{rem:e8v16-sectors}
For finite-cutoff gauge-link and local CAR sectors, the bounded criterion is
the natural row.  Weyl observables in bosonic sectors can instead satisfy the
unbounded commutator criterion because their commutator with a polynomial
local Hamiltonian is supported near (B), although its vacuum norm still
needs a moment estimate.  The gravitational and interacting scalar sectors
therefore require the vector sum \eqref{eq:e8v16-form-sum}; locality of the
formal expression alone does not bound its norm.
\end{remark}

\subsection{Why a raw Gibbs trace is not the cofinal local row}
\label{sec:e8v16-Gibbs}

Let (H_\Lambda\geq0) be finite dimensional and
\(\rho_\Lambda=Z_\Lambda^{-1}e^{-\beta H_\Lambda}\).  For a bounded local
operator (A_B), cyclicity and \([\rho_\Lambda,H_\Lambda]=0\) give
\begin{align}
 \Tr(\rho_\Lambda A_B^*H_\Lambda A_B)
 &=\Tr(\rho_\Lambda H_\Lambda A_B^*A_B)
 +\Tr(\rho_\Lambda A_B^*[H_\Lambda,A_B]).
 \label{eq:e8v16-Gibbs-split}
\end{align}

\begin{proposition}[Extensive background obstruction]
\label{prop:e8v16-Gibbs-obstruction}
The first term on the right of \eqref{eq:e8v16-Gibbs-split} obeys only
\begin{equation}
 0\leq\Tr(\rho_\Lambda H_\Lambda A_B^*A_B)
 \leq\|A_B\|^2\Tr(\rho_\Lambda H_\Lambda).
 \label{eq:e8v16-background}
\end{equation}
For \(A_B=I\), equality holds.  Hence a uniform bound on the raw V14 moment
\(\Tr\rho_\Lambda A_B^*(H_\Lambda+\lambda)A_B\) would already require a
uniform bound on the total Gibbs energy.  In a thermodynamic family with
positive energy density this is impossible.
\end{proposition}

\begin{proof}
Because \(\rho_\Lambda H_\Lambda\geq0\) and commutes internally,
\(0\leq A_B^*A_B\leq\|A_B\|^2I\) gives
\eqref{eq:e8v16-background} after taking the trace.  Setting (A_B=I)
shows sharpness.  If the mean energy grows with volume, the asserted raw
moment grows as well.
\end{proof}

\begin{assessment}[Correct carrier for the FLAC energy row]
\label{ass:e8v16-carrier}
The cofinal local row must be evaluated in the reconstructed positive-energy
vacuum representation, where (H\Omega=0), or in an exactly equivalent
source-local OS form.  A finite-temperature density matrix on the entire
finite-volume Hilbert space inserts the background energy
\eqref{eq:e8v16-background}.  Subtracting that number informally is
insufficient: the renormalized quadratic form, its positivity, and its
incidence with the reflected Hamiltonian must be specified.
\end{assessment}

\subsection{Insertion into the FLAC}
\label{sec:e8v16-FLAC}

\begin{theorem}[Local-energy population theorem]
\label{thm:e8v16-FLAC-energy}
Consider a cofinal family with normalized vectors \(\Omega_n\), positive
Hamiltonians \(H_n\Omega_n=0\), and transported local observables \(A_n\).
If either \Cref{thm:e8v16-bounded-local} or
\Cref{thm:e8v16-unbounded-local} holds uniformly for every member of a fixed
local core, then FLAC row (F3) holds with
\begin{equation}
 E_{n,A}\leq
 \begin{cases}
 2J_B\|A\|^2+\lambda\|A\|^2,
 &\text{in the bounded interaction branch},\\
 \|A\|D_B(A)+\lambda\|A\|^2,
 &\text{in the form-commutator branch}.
 \end{cases}
 \label{eq:e8v16-FLAC-energy}
\end{equation}
For a finite source bank, the stronger relative estimate
\eqref{eq:e8v16-relative-commutator} also supplies the source-local energy
matrix required by the existing semigroup floor theorem.
\end{theorem}

\begin{proof}
The FLAC energy norm is
\(\|H_n^{1/2}A_n\Omega_n\|^2+\lambda\|A_n\Omega_n\|^2\).
Use \eqref{eq:e8v16-bounded-energy} or
\eqref{eq:e8v16-unbounded-energy}, and bound the second term by
\(\lambda\|A\|^2\).  The source-bank assertion is
\Cref{thm:e8v16-bank-compiler}.
\end{proof}

\begin{firewall}[What has and has not been acquired]
\label{fire:e8v16-energy}
V16 reduces the local energy row to a concrete commutator norm.  It does not
prove the uniform vector sum \eqref{eq:e8v16-form-sum} for the interacting
SM--Einstein Hamiltonian, construct its vacuum representation, or identify a
finite-temperature Gibbs trace with that vacuum form.  In particular, the
formal locality of a gravitational Hamiltonian density does not establish the
required domain or vacuum-moment bound.
\end{firewall}

\begin{assessment}[Progress at V16]
\label{ass:e8v16-status}
The first FLAC row now has an upstream acquisition route.  Bounded local
interactions give a volume-independent energy estimate from their overlap
sum.  Unbounded sectors require an explicit, absolutely summable vacuum
commutator vector.  The raw Gibbs trace route has been rejected because it
retains the extensive background even for the identity observable.  The next
unresolved FLAC row is the form-dual smoothing integral used by two marked
insertions.
\end{assessment}

\begin{programme}[Eighth-series V17: audit the smoothing exponent]
\label{prog:e8v16-V17}
V17 will compute the exact spectral bound for
\(e^{-tH}:\mathcal Q^*\to\mathcal Q\).  It will determine whether the
integral \(G_T\) in V14 is finite for the natural energy form domain.  If the
endpoint exponent is critical, the programme will go upstream to fractional
form scales or a renormalized/contact cancellation rather than assume an
integrable kernel.
\end{programme}

\begin{firewall}[Claim boundary after eighth-series V16]
\label{fire:e8v16-boundary}
V16 proves the exact vacuum-commutator energy identity, a finite source-bank
compiler, bounded and unbounded local-interaction sufficient criteria, and
the extensive-background obstruction for a raw Gibbs trace.

V16 does not verify the unbounded SM--Einstein commutator row, construct the
physical vacuum Hamiltonian, populate the remaining FLAC rows, remove the
cutoff, or construct the Standard Model--Einstein continuum.
\end{firewall}

\begin{theorem}[Eighth-series V16 result]
\label{thm:e8v16-result}
A local observable acting on a zero-energy vector has form energy exactly
given by \eqref{eq:e8v16-energy-identity}.  Uniform local commutator control
therefore populates the FLAC local-energy row by
\eqref{eq:e8v16-FLAC-energy}, whereas the corresponding raw Gibbs trace
contains the extensive background \eqref{eq:e8v16-background}.
\end{theorem}

\begin{proof}
Combine \Cref{lem:e8v16-energy-identity,
thm:e8v16-bounded-local,thm:e8v16-unbounded-local,
prop:e8v16-Gibbs-obstruction,thm:e8v16-FLAC-energy}.
\end{proof}

\paragraph{Parent-version record.}
V16 was formed from
\texttt{Renewal\_SM\_Einstein\_Continuum\_Research\_Notes\_V15.tex}.
The V15 parent had (5\,452) logical source lines, (203\,330) bytes,
and SHA--256
\[
 \texttt{33A98501D70AFF4F87924CDE3121C8F335D2E5D0E89B275D88DAFEFAFFF5C93F}.
\]
The next compulsory companion audit is V20.

% =============================================================================
% END APPENDED EIGHTH-SERIES V16 MATERIAL
% =============================================================================


% =============================================================================
% BEGIN APPENDED EIGHTH-SERIES V17 MATERIAL
% =============================================================================

\section{The critical form-smoothing exponent and its logarithmic obstruction}
\label{sec:v8v17}

V14 and V15 isolated an integrable map
\(e^{-tH}:\mathcal Q^*\to\mathcal Q\) as the bridge between two form
insertions.  This appendix computes that norm exactly on the Hilbert scale of
(H).  At the natural form exponent the singularity is critical and its
ordered integral diverges for every unbounded (H).  A fractional extension
of the insertions makes the integral finite; otherwise a contact
renormalization or a structural cancellation is required.

\subsection{The spectral Hilbert scale}
\label{sec:e8v17-scale}

Let (H\geq0) be self-adjoint on \(\mathcal H\), fix \(\lambda>0\), and set
\(A=H+\lambda\).  For \(s\geq0\), let
\begin{equation}
 \mathcal H_s:=\operatorname{Dom}(A^{s/2}),
 \qquad
 \|u\|_s:=\|A^{s/2}u\|.
 \label{eq:e8v17-Hs}
\end{equation}
Let \(\mathcal H_{-s}\) be the completion of \(\mathcal H\) in the norm
\(\|f\|_{-s}=\|A^{-s/2}f\|\).  The anti-dual of \(\mathcal H_s\) is
canonically \(\mathcal H_{-s}\).  The natural form domain is
\(\mathcal Q=\mathcal H_1\).

\begin{theorem}[Exact Hilbert-scale smoothing norm]
\label{thm:e8v17-exact-smoothing}
For (s\geq0) and (t>0),
\begin{equation}
 \boxed{
 \|e^{-tH}\|_{\mathcal H_{-s}\to\mathcal H_s}
 =\sup_{E\in\sigma(H)}(E+\lambda)^s e^{-tE}.}
 \label{eq:e8v17-exact}
\end{equation}
In particular,
\begin{equation}
 \|e^{-tH}\|_{\mathcal H_{-s}\to\mathcal H_s}
 \leq \sup_{E\geq0}(E+\lambda)^s e^{-tE}.
 \label{eq:e8v17-envelope}
\end{equation}
For (s>0), the last supremum is
\begin{equation}
 \mathfrak g_{s,\lambda}(t)=
 \begin{cases}
 \displaystyle
 \left({s\over t}\right)^s e^{-s+\lambda t},
 &0<t\leq s/\lambda,\\
 \lambda^s,&t\geq s/\lambda.
 \end{cases}
 \label{eq:e8v17-g-exact}
\end{equation}
\end{theorem}

\begin{proof}
The isometry \(A^{-s/2}:\mathcal H_{-s}\to\mathcal H\) identifies an input
\(f\) with \(g=A^{-s/2}f\).  Functional calculus gives
\[
 \|e^{-tH}f\|_s
 =\|A^{s/2}e^{-tH}A^{s/2}g\|
 =\|A^se^{-tH}g\|.
\]
The operator norm of the last spectral multiplier is the essential supremum
of \((E+\lambda)^se^{-tE}\) on \(\sigma(H)\), proving
\eqref{eq:e8v17-exact}.  To maximize over the entire half-line, differentiate
its logarithm.  The interior stationary point is
\(E=s/t-\lambda\); it is nonnegative exactly when \(t\leq s/\lambda\).
Substitution gives the first branch of \eqref{eq:e8v17-g-exact}; otherwise
the maximum occurs at \(E=0\).
\end{proof}

\begin{corollary}[A convenient fractional envelope]
\label{cor:e8v17-fractional-envelope}
For (0<s<1),
\begin{equation}
 \|e^{-tH}\|_{\mathcal H_{-s}\to\mathcal H_s}
 \leq \lambda^s+\left({s\over et}\right)^s.
 \label{eq:e8v17-fractional-envelope}
\end{equation}
Consequently,
\begin{equation}
 \boxed{
 \int_0^T(T-t)
 \|e^{-tH}\|_{\mathcal H_{-s}\to\mathcal H_s}\,dt
 \leq {\lambda^sT^2\over2}
 +{(s/e)^sT^{2-s}\over(1-s)(2-s)}<\infty.}
 \label{eq:e8v17-fractional-G}
\end{equation}
\end{corollary}

\begin{proof}
For (0<s<1), concavity gives
\((E+\lambda)^s\leq E^s+\lambda^s\).  The maximum of
\(E^se^{-tE}\) is \((s/(et))^s\).  Integrate the resulting bound and use
\[
 \int_0^T(T-t)t^{-s}\,dt
 ={T^{2-s}\over(1-s)(2-s)}.
\]
\end{proof}

\subsection{Endpoint failure on the natural form domain}
\label{sec:e8v17-endpoint}

For (s=1), \eqref{eq:e8v17-g-exact} has the small-time behavior
\begin{equation}
 \mathfrak g_{1,\lambda}(t)
 ={e^{-1+\lambda t}\over t}\sim{1\over et}.
 \label{eq:e8v17-critical}
\end{equation}
The upper envelope is not merely too crude: unbounded spectrum forces the
actual operator-norm integral to diverge.

\begin{theorem}[Universal endpoint divergence for unbounded generators]
\label{thm:e8v17-endpoint-divergence}
If (H\geq0) is unbounded, then for every (T>0),
\begin{align}
 \int_0^T
 \|e^{-tH}\|_{\mathcal H_{-1}\to\mathcal H_1}\,dt
 &=\infty,
 \label{eq:e8v17-L1-divergence}\\
 \int_0^T(T-t)
 \|e^{-tH}\|_{\mathcal H_{-1}\to\mathcal H_1}\,dt
 &=\infty.
 \label{eq:e8v17-G-divergence}
\end{align}
Thus the V14 smoothing hypothesis with
\(\mathcal Q=\operatorname{Dom}(H+\lambda)^{1/2}\) fails for every
unbounded (H).
\end{theorem}

\begin{proof}
Choose spectral points (E_k\in\sigma(H)) with
\(E_{k+1}\geq2E_k\) and \(E_1>2/T\).  For
(t\in[1/E_{k+1},1/E_k]), formula \eqref{eq:e8v17-exact} and the spectral
point \(E_k\) give
\[
 \|e^{-tH}\|_{\mathcal H_{-1}\to\mathcal H_1}
 \geq(E_k+\lambda)e^{-tE_k}
 \geq E_ke^{-1}.
\]
Therefore the integral on this interval is at least
\[
 {E_k\over e}\left({1\over E_k}-{1\over E_{k+1}}\right)
 \geq{1\over2e}.
\]
The intervals are disjoint and accumulate at zero, so summing over (k)
proves \eqref{eq:e8v17-L1-divergence}.  On \((0,T/2]\),
\(T-t\geq T/2\), which proves \eqref{eq:e8v17-G-divergence}.
\end{proof}

\begin{corollary}[Cutoff norms hide a cofinal logarithm]
\label{cor:e8v17-cutoff-log}
Let \(H_N\) be bounded nonnegative truncations with spectral ceiling
\(\Lambda_N\to\infty\).  The universal endpoint estimate gives, for
\(0<\tau<T\),
\begin{equation}
 \int_\tau^T(T-t)
 \|e^{-tH_N}\|_{\mathcal H_{-1}\to\mathcal H_1}\,dt
 \leq {\lambda T^2\over2}+{T\over e}\log{T\over\tau}.
 \label{eq:e8v17-collar}
\end{equation}
If the collar is taken at the ultraviolet scale
\(\tau\asymp\Lambda_N^{-1}\), this bound grows like
\(\log\Lambda_N\).  Finite-cutoff boundedness alone therefore does not give
a cofinal two-form estimate.
\end{corollary}

\begin{proof}
For (s=1), the elementary envelope following
\eqref{eq:e8v17-envelope} is
\[
 \|e^{-tH_N}\|_{-1\to1}\leq\lambda+{1\over et}.
\]
Multiply by \(T-t\leq T\), integrate the singular term on
\([\tau,T]\), and use
\(\int_0^T(T-t)\lambda\,dt=\lambda T^2/2\).
\end{proof}

\begin{remark}[The divergence is an absolute-value statement]
\label{rem:e8v17-cancellation}
The theorem rules out the V14 proof based on the absolute product of two
arbitrary order-one form norms.  It does not rule out a particular second
derivative whose two time orders cancel a contact singularity, nor a
renormalized coincident insertion.  Such a cancellation must be exhibited
before integration; it cannot be inferred from existence of every bounded
cutoff integral.
\end{remark}

\subsection{The fractional two-insertion repair}
\label{sec:e8v17-fractional}

Let \(0<s<1\).  A symmetric insertion form \(b\) is called
\emph{(s)-subcritical} relative to (H) if it extends to a bounded map
\begin{equation}
 \widehat b:\mathcal H_s\longrightarrow\mathcal H_{-s},
 \qquad \|\widehat b\|_{s\to-s}\leq\eta_b.
 \label{eq:e8v17-subcritical}
\end{equation}

\begin{theorem}[Subcritical ordered form estimate]
\label{thm:e8v17-subcritical-two}
Let \(b_1,b_2\) be \(s\)-subcritical.  Suppose every outer semigroup factor
in an ordered two-insertion term is bounded by \(M_s\) on
\(\mathcal H_s\).  Then
\begin{align}
 &\left|\int_{\substack{r,q\geq0\\r+q\leq T}}
 \langle u,S(T-r-q)\widehat b_1S(r)
 \widehat b_2S(q)v\rangle\,dr\,dq\right|
 \notag\\
 &\qquad\leq
 M_s^2\eta_{b_1}\eta_{b_2}
 \left[{\lambda^sT^2\over2}
 +{(s/e)^sT^{2-s}\over(1-s)(2-s)}\right]
 \|u\|_s\|v\|_s.
 \label{eq:e8v17-subcritical-two}
\end{align}
\end{theorem}

\begin{proof}
Read the ordered pairing from the right.  The first insertion maps
\(\mathcal H_s\) to \(\mathcal H_{-s}\); the middle semigroup returns it
to \(\mathcal H_s\) with the norm in
\eqref{eq:e8v17-fractional-envelope}; the second insertion maps back to the
anti-dual, and the left pairing costs the two outer bounds.  The simplex
integration is exactly the integral in
\eqref{eq:e8v17-fractional-G}.
\end{proof}

\begin{corollary}[Repaired V15 derivative row]
\label{cor:e8v17-repaired-V15}
In the V15 adjacent derivative comparison, replace the natural-form norms of
both the adjacent base insertion and the parameter-derivative insertions by
uniform \(s\)-subcritical norms for one fixed \(s<1\).  Then its mixed term is
finite with
\begin{equation}
 G_T(g_n)\leq
 {\lambda^sT^2\over2}
 +{(s/e)^sT^{2-s}\over(1-s)(2-s)}.
 \label{eq:e8v17-repaired-G}
\end{equation}
Consequently the FLAC derivative defect remains valid after replacing its
formal \(G_T(g_n)\) by the right side of
\eqref{eq:e8v17-repaired-G}.
\end{corollary}

\begin{proof}
Apply \Cref{thm:e8v17-subcritical-two} to the two time orders in the proof of
the V15 adjacent derivative lemma.  The single insertion-difference term
requires only the corresponding (s)-scale pairing and is unchanged in
structure.
\end{proof}

\begin{countertheorem}[Relative form boundedness does not imply
subcriticality]
\label{cth:e8v17-relative-not-fractional}
There are form-bounded insertions on \(\mathcal H_1\) which do not extend
boundedly from \(\mathcal H_s\) to \(\mathcal H_{-s}\) for any \(s<1\).
\end{countertheorem}

\begin{proof}
Take \(\mathcal H=\ell^2(\mathbb N)\), \(He_k=k e_k\), and
\(b[u,v]=\langle H^{1/2}u,H^{1/2}v\rangle\).  This is the closed form of
(H), hence bounded on \(\mathcal H_1\).  If it extended from
\(\mathcal H_s\) to \(\mathcal H_{-s}\), testing on \(e_k\) would give
\[
 k=|b[e_k,e_k]|
 \leq C\|e_k\|_s^2=C(k+\lambda)^s.
\]
This is impossible as \(k\to\infty\) when \(s<1\).
\end{proof}

\subsection{Revised smoothing branch of the FLAC}
\label{sec:e8v17-FLAC}

\begin{definition}[Three admissible two-mark branches]
\label{def:e8v17-branches}
The FLAC one- and two-mark row is populated only after one of the following is
proved uniformly along the cofinal family.
\begin{enumerate}[label=\textup{(S\arabic*)},leftmargin=4em]
\item \emph{Fractional branch:} both marked insertions have the same
\(s\)-subcritical extension for some fixed \(s<1\).
\item \emph{Collared branch:} the marks are physically separated by a fixed
positive collar.  Any removal of the collar includes the explicit logarithmic
debit \eqref{eq:e8v17-collar}.
\item \emph{Cancellation branch:} the ordered pair plus its coincident contact
form has a combined integrand with an integrable norm majorant.  The combined
identity is proved before absolute values are taken.
\end{enumerate}
\end{definition}

\begin{theorem}[Endpoint audit of FLAC row (F4)]
\label{thm:e8v17-FLAC-audit}
For an unbounded limiting Hamiltonian, the natural common-form estimate of
V14 does not by itself populate FLAC row (F4).  The fractional branch is
sufficient by \Cref{thm:e8v17-subcritical-two}; the collared branch is
sufficient only at fixed collar; and the cancellation branch remains a
model-specific proof obligation.
\end{theorem}

\begin{proof}
Failure of the unqualified natural-form branch is
\Cref{thm:e8v17-endpoint-divergence}.  Sufficiency of the fractional branch
is \Cref{thm:e8v17-subcritical-two}.  The fixed-collar estimate is
\Cref{cor:e8v17-cutoff-log}; its logarithm prevents removal without another
small factor or subtraction.  The last branch is a definition of the missing
combined estimate and has no consequence until that estimate is proved.
\end{proof}

\begin{assessment}[Upstream interpretation for SM--Einstein variations]
\label{ass:e8v17-SM}
Metric variation of a principal kinetic form is generally of the same
differential order as the form itself, so relative form boundedness naturally
lands at \(s=1\), not below it.  The fractional branch must therefore be
proved from additional regularity or from a transport which fixes the
principal kinetic form and leaves a lower-order adjacent insertion.  If that
cannot be done, the programme must retain the metric Hessian/contact term and
seek a cancellation branch.  Assuming \(G_T<\infty\) at exponent one would
be circular.
\end{assessment}

\begin{programme}[Eighth-series V18: principal-form transport]
\label{prog:e8v17-V18}
V18 will go upstream to a common kinetic-frame transport.  It will determine
when two adjacent uniformly elliptic quadratic forms can be conjugated to one
fixed principal form, compute the connection/contact terms created by the
transport, and test whether the remaining metric and refinement insertions
are (s)-subcritical.  Failure will be recorded as a precise principal-symbol
obstruction rather than hidden in the cutoff.
\end{programme}

\begin{firewall}[Claim boundary after eighth-series V17]
\label{fire:e8v17-boundary}
V17 computes the exact Hilbert-scale semigroup smoothing norm; proves
universal divergence of the natural form-dual integral for unbounded
Hamiltonians; proves an integrable fractional two-insertion estimate; and
states the three admissible FLAC smoothing branches.

V17 does not prove fractional regularity or a contact cancellation for the
SM--Einstein insertions, populate FLAC row (F4), remove the cutoff, or
construct the Standard Model--Einstein continuum.
\end{firewall}

\begin{theorem}[Eighth-series V17 result]
\label{thm:e8v17-result}
For an unbounded nonnegative Hamiltonian, the absolute smoothing bridge from
its natural form anti-dual to its form domain is never integrable at zero.
For every fixed \(0<s<1\), the fractional bridge is integrable with the
explicit bound \eqref{eq:e8v17-fractional-G}.  Hence two arbitrary order-one
forms do not satisfy the V14 two-mark hypothesis, whereas uniformly
(s)-subcritical insertions do.
\end{theorem}

\begin{proof}
Combine \Cref{thm:e8v17-exact-smoothing,
thm:e8v17-endpoint-divergence,thm:e8v17-subcritical-two,
cth:e8v17-relative-not-fractional}.
\end{proof}

\paragraph{Parent-version record.}
V17 was formed from
\texttt{Renewal\_SM\_Einstein\_Continuum\_Research\_Notes\_V16.tex}.
The V16 parent had (5\,908) logical source lines, (219\,440) bytes,
and SHA--256
\[
 \texttt{953BBA79DBE7F75049944E9E1AF18D2BC1F84C5FBD6D3DE73A81521327188E86}.
\]
The next compulsory companion audit is V20.

% =============================================================================
% END APPENDED EIGHTH-SERIES V17 MATERIAL
% =============================================================================



% =============================================================================
% BEGIN APPENDED EIGHTH-SERIES V18 MATERIAL
% =============================================================================

\section{Principal-form transport: congruence, unitary incidence, and obstruction}
\label{sec:v8v18}

V17 suggested freezing the principal kinetic form in order to make metric or
refinement insertions subcritical.  There are two different operations which
can be called such a freezing.  An energy congruence exists very generally,
but changes the Hilbert metric and does not conjugate the heat semigroup.  A
physical unitary conjugacy preserves the semigroup, but it can freeze a
kinetic operator only along an isospectral, and locally an isometric, orbit.
This appendix proves the distinction and states the exact conditional branch
which can feed the fractional estimate of V17.

\subsection{Energy congruence is not physical conjugacy}
\label{sec:e8v18-congruence}

Let (H_0,H_1\geq0) be self-adjoint on \(\mathcal H\), fix
\(\lambda>0\), and write \(A_j=H_j+\lambda\).  Suppose their form domains
coincide and their shifted form norms are equivalent.  Define on the common
form domain
\begin{equation}
 R_{10}:=A_1^{-1/2}A_0^{1/2}.
 \label{eq:e8v18-R}
\end{equation}

\begin{lemma}[Canonical energy isometry]
\label{lem:e8v18-energy-isometry}
The map \(R_{10}\) is an isometry from the (A_0)-form Hilbert space to the
(A_1)-form Hilbert space, with inverse
\(R_{01}=A_0^{-1/2}A_1^{1/2}\).  In particular,
\begin{equation}
 \boxed{
 \langle A_1^{1/2}R_{10}u,A_1^{1/2}R_{10}v\rangle
 =\langle A_0^{1/2}u,A_0^{1/2}v\rangle.}
 \label{eq:e8v18-congruence}
\end{equation}
At a bounded finite cutoff this is the matrix congruence
\(R_{10}^*A_1R_{10}=A_0\).
\end{lemma}

\begin{proof}
For \(u\) in the common form domain,
\(A_1^{1/2}R_{10}u=A_0^{1/2}u\).  This proves
\eqref{eq:e8v18-congruence}.  Interchanging the indices gives the inverse
on the form completions.  At finite cutoff all factors are bounded matrices,
so the same identity is the displayed congruence.
\end{proof}

\begin{remark}[No automatic base-Hilbert boundedness]
\label{rem:e8v18-H-boundedness}
The energy isometry acts between form Hilbert spaces.  It need not extend to
a uniformly bounded operator on the underlying \(\mathcal H\) as the cutoff
is removed.  Even when it does, it need not be unitary in \(\mathcal H\).
Both properties are separate carrier/contact rows.
\end{remark}

\begin{proposition}[Congruence does not freeze the semigroup]
\label{prop:e8v18-not-semigroup}
There are positive one-dimensional Hamiltonians for which
\eqref{eq:e8v18-congruence} holds but
\begin{equation}
 R_{10}^*e^{-tH_1}R_{10}\neq e^{-tH_0}
 \label{eq:e8v18-not-semigroup}
\end{equation}
for every \(t>0\).
\end{proposition}

\begin{proof}
Take \(\mathcal H=\mathbb C\), \(H_j=h_jI\) with
\(0\leq h_0\neq h_1\), and
\[
 R_{10}=\sqrt{h_0+\lambda\over h_1+\lambda}.
\]
Then the shifted quadratic forms are congruent, while the left side of
\eqref{eq:e8v18-not-semigroup} is
\[
 {h_0+\lambda\over h_1+\lambda}e^{-th_1},
\]
which cannot equal \(e^{-th_0}\) for every positive \(t\).  In particular,
the values at \(t=0\) already differ unless \(h_0=h_1\).
\end{proof}

\begin{assessment}[Where the dynamics moves]
\label{ass:e8v18-metric-moves}
Pulling back the shifted energy by \(R_{10}\) fixes the energy form but pulls
the ambient scalar product back to
\begin{equation}
 m_{10}[u,v]:=\langle R_{10}u,R_{10}v\rangle.
 \label{eq:e8v18-pulled-metric}
\end{equation}
The physical generator becomes a generalized operator relative to this
moving metric.  Thus the principal variation has moved into the clock and
carrier incidence; it has not disappeared.
\end{assessment}

\subsection{What a physical unitary can freeze}
\label{sec:e8v18-unitary}

\begin{theorem}[Unitary freezing requires spectral equivalence]
\label{thm:e8v18-unitary-spectrum}
If a unitary \(U\) satisfies
\begin{equation}
 U^*H_1U=H_0,
 \label{eq:e8v18-unitary-conjugacy}
\end{equation}
then the projection-valued spectral measures obey
\(E_{H_1}(B)=UE_{H_0}(B)U^*\) for every Borel set \(B\).  Consequently
\(H_0\) and \(H_1\) have the same spectrum, spectral multiplicities, heat
trace whenever defined, and all unitary spectral invariants.  A generic
metric variation which changes any of these data cannot be frozen by a
physical unitary.
\end{theorem}

\begin{proof}
Apply the Borel functional calculus to
\eqref{eq:e8v18-unitary-conjugacy}.  For every bounded Borel function \(f\),
\(f(H_1)=Uf(H_0)U^*\).  Indicator functions give the spectral projections,
and the stated invariants follow.
\end{proof}

A more local obstruction is visible at the principal symbol.  Let \(M\) be a
closed smooth manifold.  Consider scalar Laplace-type operators whose
principal symbols are
\begin{equation}
 \sigma_2(H_g)(x,\xi)=g^{ab}(x)\xi_a\xi_b.
 \label{eq:e8v18-symbol}
\end{equation}
A local geometric unitary has the form
\begin{equation}
 (U_{\phi,\theta}u)(x)
 =e^{\mathrm i\theta(x)}J_\phi(x)^{1/2}u(\phi(x)),
 \label{eq:e8v18-geometric-U}
\end{equation}
with the density factor chosen for the two (L^2) measures.

\begin{theorem}[Principal-symbol obstruction for local unitary transport]
\label{thm:e8v18-symbol-obstruction}
Let \(g_0,g_1\) be smooth Riemannian metrics.  If
\begin{equation}
 U_{\phi,\theta}^*H_{g_1}U_{\phi,\theta}-H_{g_0}
 \label{eq:e8v18-lower-order-difference}
\end{equation}
is a differential operator of order at most one, then
\begin{equation}
 g_0=\phi^*g_1.
 \label{eq:e8v18-isometry}
\end{equation}
Thus multiplication phases, density changes, and coordinate transport remove
the second-order metric variation only when it is a diffeomorphism direction.
\end{theorem}

\begin{proof}
Multiplication by the phase and by a nonvanishing scalar density changes only
terms of order at most one under conjugation.  The chain rule transforms the
quadratic principal symbol of \(H_{g_1}\) into the pullback quadratic form
\((\phi^*g_1)^{ab}\xi_a\xi_b\).  If
\eqref{eq:e8v18-lower-order-difference} has order at most one, its degree-two
symbol vanishes.  Hence
\((\phi^*g_1)^{ab}\xi_a\xi_b=g_0^{ab}\xi_a\xi_b\) for every cotangent vector
\(\xi\), which is equivalent to \eqref{eq:e8v18-isometry}.
\end{proof}

\begin{corollary}[Einstein directions cannot all be gauge-fixed away]
\label{cor:e8v18-Einstein-direction}
A metric tangent which is not a Lie derivative
\(\dot g=\mathcal L_Xg\) cannot be made lower order by differentiating a
family of local geometric unitaries of the form
\eqref{eq:e8v18-geometric-U}.  Physical transverse metric variations
therefore remain principal-form insertions.
\end{corollary}

\begin{proof}
Differentiate \eqref{eq:e8v18-isometry} along a smooth family with
\(\phi_0=\operatorname{id}\).  The tangent to a diffeomorphism orbit is
\(\mathcal L_Xg\).  A tangent outside that orbit contradicts the necessary
condition in \Cref{thm:e8v18-symbol-obstruction}.
\end{proof}

\subsection{The connection term of a moving unitary}
\label{sec:e8v18-connection}

Let \(x\mapsto U(x)\) be a strongly \(C^1\) unitary family and let
\(x\mapsto H(x)\) be differentiable on a common operator core.  Put
\begin{equation}
 \widetilde H(x)=U(x)^*H(x)U(x),
 \qquad
 K_\xi(x)=U(x)^*D_\xi U(x).
 \label{eq:e8v18-Htilde-K}
\end{equation}
Then \(K_\xi^*=-K_\xi\).

\begin{lemma}[Exact moving-frame derivative]
\label{lem:e8v18-moving-frame}
On the common core,
\begin{equation}
 \boxed{
 D_\xi\widetilde H
 =U^*(D_\xi H)U+[\widetilde H,K_\xi].}
 \label{eq:e8v18-moving-derivative}
\end{equation}
The second term is a required connection/contact term.
\end{lemma}

\begin{proof}
Differentiate \(\widetilde H=U^*HU\).  From \(U^*U=I\) one obtains
\(D_\xi U=UK_\xi\) and
\(D_\xi U^*=-K_\xi U^*\).  Therefore
\[
 D_\xi\widetilde H
 =-K_\xi\widetilde H+U^*(D_\xi H)U+\widetilde H K_\xi,
\]
which is \eqref{eq:e8v18-moving-derivative}.
\end{proof}

\begin{countertheorem}[A bounded connection need not be lower order]
\label{cth:e8v18-bounded-K}
Boundedness of \(K\) on \(\mathcal H\) does not imply that
\([H,K]\) is \(s\)-subcritical for any \(s<1\).
\end{countertheorem}

\begin{proof}
On \(\ell^2(\mathbb N_0)\), let \(He_n=2^ne_n\) and let \(S\) be the
unilateral shift.  The skew-adjoint bounded operator \(K=S-S^*\) has
matrix elements of \([H,K]\) between \(e_n\) and \(e_{n+1}\) of magnitude
\(2^n\).  If the commutator extended from \(\mathcal H_s\) to
\(\mathcal H_{-s}\), testing those two basis vectors would give
\[
 2^n\leq C(2^n+\lambda)^{s/2}
               (2^{n+1}+\lambda)^{s/2}=O(2^{ns}),
\]
which is impossible for \(s<1\).
\end{proof}

\subsection{The exact conditional transport branch}
\label{sec:e8v18-conditional}

\begin{definition}[Subcritical physical kinetic frame]
\label{def:e8v18-SPKF}
A cofinal family has a subcritical physical kinetic frame, abbreviated SPKF,
on a parameter chart if there are physical unitaries \(U_n(x)\), a fixed
reference principal form \(h_{n,\mathrm{pr}}\), and remainders \(b_n(x)\)
such that
\begin{equation}
 U_n(x)^*H_n(x)U_n(x)=H_{n,\mathrm{pr}}+B_n(x),
 \label{eq:e8v18-SPKF}
\end{equation}
and, for one fixed \(s<1\), all of the following are uniformly bounded maps
\(\mathcal H_{n,s}\to\mathcal H_{n,-s}\):
\begin{enumerate}[label=\textup{(K\arabic*)},leftmargin=4em]
\item adjacent differences \(B_{n+1}-B_n\) after the declared bridge;
\item parameter derivatives \(D_\xi B_n\);
\item the complete connection contribution
\([H_{n,\mathrm{pr}}+B_n,K_{n,\xi}]\); and
\item adjacent differences of the last two rows.
\end{enumerate}
Carrier derivatives \(D_\xi U_n\), local observable transports, and their
projective defects are retained in the FLAC contact row.
\end{definition}

\begin{theorem}[SPKF populates the fractional smoothing branch]
\label{thm:e8v18-SPKF}
If SPKF holds, its fractional norms and contact bounds obey the cofinal
summability requirements of V15, and its local source vectors have the
fractional energy moments required in V17, then the two-mark generator part
of the FLAC derivative defect is finite and is bounded by the V17
subcritical formula.  The resulting (C^1) landing statement is unitarily
identical to the original physical matrix elements.
\end{theorem}

\begin{proof}
By \Cref{lem:e8v18-moving-frame}, the complete transformed derivative
insertion is the sum of rows (K2) and (K3); its adjacent difference is
controlled by rows (K1) and (K4), with the bridge terms included.  Every
marked insertion therefore maps \(\mathcal H_s\) to
\(\mathcal H_{-s}\).  Apply the V17 subcritical ordered-form theorem to the
two time orders and add the declared contact tails.  The V15 landing theorem
then applies.  Finally,
\[
 \langle u,e^{-tH_n(x)}v\rangle
 =\langle U_n(x)^*u,e^{-t\widetilde H_n(x)}U_n(x)^*v\rangle,
\]
so the transported and original matrix elements agree exactly when all
carrier derivatives are included.
\end{proof}

\begin{firewall}[SPKF is not implied by ellipticity]
\label{fire:e8v18-SPKF}
Uniform ellipticity gives equivalent form norms and therefore the energy
congruence of \Cref{lem:e8v18-energy-isometry}.  It does not give a physical
unitary, an isospectral orbit, or the subcritical connection bounds in SPKF.
For transverse Einstein variations,
\Cref{cor:e8v18-Einstein-direction} leaves the principal insertion in place.
The general programme must therefore use an aligned/cancellation mechanism
or prove extra source regularity; it cannot declare the kinetic term fixed by
congruence.
\end{firewall}

\begin{assessment}[Progress at V18]
\label{ass:e8v18-status}
The proposed kinetic-frame shortcut has been separated into a valid special
branch and an invalid general inference.  Static energy congruence is always
available under common-form equivalence, but it moves the physical Hilbert
metric.  Local unitary transport removes a principal metric variation only
on a diffeomorphism orbit, and its moving connection can itself retain full
form order.  SPKF is sufficient, but it is a new row rather than a consequence
of the existing SM--Einstein regularity assumptions.
\end{assessment}

\begin{programme}[Eighth-series V19: aligned principal insertions]
\label{prog:e8v18-V19}
V19 will test the cancellation branch on the exact model
\(H(x)=c(x)H_0+V(x)\).  When the principal insertion is a scalar multiple of
the generator, all heat factors commute and the apparent (t^{-1}) middle
singularity can be combined before taking norms.  The resulting theorem will
identify which aligned principal variations are finite and which genuinely
noncommuting metric variations remain open.
\end{programme}

\begin{firewall}[Claim boundary after eighth-series V18]
\label{fire:e8v18-boundary}
V18 proves the canonical energy congruence, proves that it does not conjugate
the physical semigroup, derives spectral and principal-symbol obstructions
to unitary freezing, computes the moving-unitary connection term, and gives
the conditional SPKF handoff to the V17 fractional branch.

V18 does not construct SPKF for transverse metric variations, prove a
principal contact cancellation, populate the complete FLAC, remove the
cutoff, or construct the Standard Model--Einstein continuum.
\end{firewall}

\begin{theorem}[Eighth-series V18 result]
\label{thm:e8v18-result}
Equivalent kinetic forms can be made identical by the energy isometry
\eqref{eq:e8v18-R}, but this operation does not identify their semigroups.
A physical local unitary can remove the principal metric difference only
along an isometric orbit and must retain the connection term
\eqref{eq:e8v18-moving-derivative}.  If the stronger SPKF rows hold, V17
supplies the desired finite two-mark bound.
\end{theorem}

\begin{proof}
Combine \Cref{lem:e8v18-energy-isometry,
prop:e8v18-not-semigroup,thm:e8v18-unitary-spectrum,
thm:e8v18-symbol-obstruction,lem:e8v18-moving-frame,
thm:e8v18-SPKF}.
\end{proof}

\paragraph{Parent-version record.}
V18 was formed from
\texttt{Renewal\_SM\_Einstein\_Continuum\_Research\_Notes\_V17.tex}.
The V17 parent had (6\,287) logical source lines, (233\,407) bytes,
and SHA--256
\[
 \texttt{3ED9F8701C5FE0304260DC386ADD01AA0A01E5CED70C839E23117A94AD57949F}.
\]
The next compulsory companion audit is V20.

% =============================================================================
% END APPENDED EIGHTH-SERIES V18 MATERIAL
% =============================================================================



% =============================================================================
% BEGIN APPENDED EIGHTH-SERIES V19 MATERIAL
% =============================================================================

\section{Aligned principal insertions and exact heat-simplex cancellation}
\label{sec:v8v19}

V17 proved that two arbitrary order-one forms cannot be bounded by separately
smoothing the interval between their marks.  The failure occurs after the
semigroup product has been split and absolute values have been taken.  When
the insertions are spectral multipliers of the same Hamiltonian, the complete
ordered product can instead be combined first.  The simplex area then
compensates the remaining endpoint singularity.  This appendix proves that
aligned branch and computes an explicit adjacent (C^1) estimate for a
scalar physical-clock variation.

\subsection{Two commuting order-one marks}
\label{sec:e8v19-two}

Let \(H\geq0\) be self-adjoint, fix \(\lambda>0\), and put
\(A=H+\lambda\).  Let \(B_1=b_1(H)\) and \(B_2=b_2(H)\) be self-adjoint
spectral multipliers satisfying
\begin{equation}
 |b_j(E)|\leq\eta_j(E+\lambda),
 \qquad E\in\sigma(H),\quad j=1,2.
 \label{eq:e8v19-relative-multipliers}
\end{equation}
Their quadratic forms are bounded on
\(\mathcal H_1=\operatorname{Dom}A^{1/2}\), even when the operators are
unbounded.

\begin{theorem}[Exact aligned two-mark reduction]
\label{thm:e8v19-aligned-two}
For \(u,v\in\mathcal H_1\) and \(T>0\), the ordered form integral exists and
obeys
\begin{align}
 &\int_{\substack{r,q\geq0\\r+q\leq T}}
 \langle u,e^{-(T-r-q)H}B_1e^{-rH}B_2e^{-qH}v\rangle
 \,dr\,dq
 \notag\\
 &\qquad={T^2\over2}\langle u,B_1B_2e^{-TH}v\rangle,
 \label{eq:e8v19-exact-simplex}\\
 &\left|{T^2\over2}\langle u,B_1B_2e^{-TH}v\rangle\right|
 \leq {T^2\over2}\eta_1\eta_2
 \mathfrak g_{1,\lambda}(T)\|u\|_1\|v\|_1,
 \label{eq:e8v19-aligned-bound}
\end{align}
where \(\mathfrak g_{1,\lambda}\) is the exact envelope from V17.  In
particular, for \(0<T\leq1/\lambda\),
\begin{equation}
 \boxed{
 \left|\textup{the ordered term}\right|
 \leq {\eta_1\eta_2\over2}
 T e^{-1+\lambda T}\|u\|_1\|v\|_1.}
 \label{eq:e8v19-linear-clock}
\end{equation}
\end{theorem}

\begin{proof}
All spectral multipliers in the integrand strongly commute.  On their common
spectral core their product is
\[
 e^{-(T-r-q)H}B_1e^{-rH}B_2e^{-qH}
 =B_1B_2e^{-TH},
\]
independent of \(r,q\).  The simplex has area \(T^2/2\), which gives
\eqref{eq:e8v19-exact-simplex}.  Move one factor \(A^{1/2}\) to each test
vector.  The intervening multiplier is
\[
 {b_1(E)b_2(E)\over E+\lambda}e^{-TE}.
\]
Its absolute value is at most
\(\eta_1\eta_2(E+\lambda)e^{-TE}\).  The V17 spectral maximum proves
\eqref{eq:e8v19-aligned-bound}; its small-time branch gives
\eqref{eq:e8v19-linear-clock}.  Approximation by bounded spectral
projections extends the identity from the core to all form vectors.
\end{proof}

\begin{remark}[Why this does not contradict V17]
\label{rem:e8v19-no-contradiction}
V17 bounded each marked form separately and paid
\(\|e^{-rH}\|_{-1\to1}\sim r^{-1}\) before integrating \(r\) near zero.
Here commutation combines the entire heat time into \(e^{-TH}\); the
integrand no longer depends on the separation \(r\), and its remaining
\(T^{-1}\) norm is multiplied by the area \(T^2/2\).  The order of
operations is the mathematical difference.
\end{remark}

\begin{corollary}[Any finite aligned word]
\label{cor:e8v19-word}
Let \(B_j=b_j(H)\) satisfy
\(|b_j(E)|\leq\eta_j(E+\lambda)^{\alpha_j}\), with
\(\alpha_j\geq0\), and let \(m\) such marks be integrated over the ordered
\(m\)-simplex of total time \(T\).  For test vectors in
\(\mathcal H_1\), the word equals
\begin{equation}
 {T^m\over m!}\langle u,B_1\cdots B_m e^{-TH}v\rangle
 \label{eq:e8v19-word}
\end{equation}
and is bounded by
\begin{equation}
 {T^m\over m!}\left(\prod_{j=1}^m\eta_j\right)
 \sup_{E\in\sigma(H)}
 (E+\lambda)^{\sum_j\alpha_j-1}e^{-TE}
 \|u\|_1\|v\|_1.
 \label{eq:e8v19-word-bound}
\end{equation}
\end{corollary}

\begin{proof}
Strong commutation makes the integrand constant on the ordered simplex,
whose volume is \(T^m/m!\).  Moving \(A^{1/2}\) to each test vector divides
the multiplier by \(E+\lambda\).  The assumed pointwise bounds give
\eqref{eq:e8v19-word-bound}.
\end{proof}

\begin{assessment}[Power count]
\label{ass:e8v19-power}
If all \(m\) marks have full form order \(\alpha_j=1\), the spectral
supremum grows like \(T^{-(m-1)}\) and the simplex volume leaves a linear
clock factor \(O(T)\).  Thus arbitrarily many perfectly aligned principal
marks remain finite on form-energy vectors.  Alignment, not the number of
marks, is doing the work.
\end{assessment}

\subsection{A scalar clock family}
\label{sec:e8v19-clock}

Let \(H_0\geq0\), let \(K\) be a compact parameter set, and let
\(c_j\in C^1(K;(0,\infty))\) for adjacent scales \(j=0,1\).  Assume
\begin{equation}
 c_j(x)\geq c_->0.
 \label{eq:e8v19-cminus}
\end{equation}
Put
\begin{equation}
 S_j(x,T)=e^{-Tc_j(x)H_0}.
 \label{eq:e8v19-Sc}
\end{equation}
For a unit parameter direction \(\xi\), write
\(\dot c_j=Dc_j[\xi]\),
\(\Delta c=c_1-c_0\), and
\(\Delta\dot c=\dot c_1-\dot c_0\).

\begin{lemma}[Exact scalar-clock derivative]
\label{lem:e8v19-clock-derivative}
On the form domain of \(H_0\),
\begin{equation}
 D_xS_j(x,T)[\xi]
 =-T\dot c_j(x)H_0e^{-Tc_j(x)H_0}.
 \label{eq:e8v19-clock-derivative}
\end{equation}
\end{lemma}

\begin{proof}
This is scalar differentiation of the spectral multiplier
\(e^{-Tc_j(x)E}\), followed by dominated convergence on a spectral core and
closure in the weak form topology.
\end{proof}

\begin{theorem}[Adjacent aligned (C^1) bound]
\label{thm:e8v19-adjacent}
For \(u,v\in\operatorname{Dom}(H_0+\lambda)^{1/2}\),
\begin{align}
 &\left|D_x\langle u,(S_1(x,T)-S_0(x,T))v\rangle[\xi]\right|
 \notag\\
 &\quad\leq
 \left[
 T|\Delta\dot c|
 +{T\over ec_-}|\Delta c|\max_{j=0,1}|\dot c_j|
 \right]
 \|u\|_{1,0}\|v\|_{1,0},
 \label{eq:e8v19-adjacent}
\end{align}
where \(\|u\|_{1,0}=\|(H_0+\lambda)^{1/2}u\|\).  The estimate has no
form-dual smoothing integral.
\end{theorem}

\begin{proof}
Using \Cref{lem:e8v19-clock-derivative}, telescope
\begin{align*}
 &\dot c_1H_0e^{-Tc_1H_0}-\dot c_0H_0e^{-Tc_0H_0}\\
 &\quad=(\Delta\dot c)H_0e^{-Tc_1H_0}
 +\dot c_0H_0(e^{-Tc_1H_0}-e^{-Tc_0H_0}).
\end{align*}
After moving one factor \((H_0+\lambda)^{1/2}\) to each test vector, the
first multiplier \(H_0/(H_0+\lambda)\) has norm at most one.  For the second
term, the scalar mean-value theorem gives, at every energy,
\[
 |e^{-Tc_1E}-e^{-Tc_0E}|
 \leq T|\Delta c|E e^{-Tc_-E}.
\]
After the two form weights are removed, its multiplier is bounded by
\[
 T|\Delta c|{E^2\over E+\lambda}e^{-Tc_-E}
 \leq T|\Delta c|E e^{-Tc_-E}
 \leq {|\Delta c|\over ec_-}.
\]
Multiply the telescoped expression by the outer \(T\) in
\eqref{eq:e8v19-clock-derivative} and use either
\(|\dot c_0|\) or the larger symmetric maximum.  This proves
\eqref{eq:e8v19-adjacent}.
\end{proof}

\begin{corollary}[Second parameter derivative]
\label{cor:e8v19-second}
For a scalar \(C^2\) curve \(c(x)\geq c_->0\),
\begin{equation}
 D_x^2e^{-Tc(x)H_0}
 =-Tc''(x)H_0e^{-Tc(x)H_0}
 +T^2c'(x)^2H_0^2e^{-Tc(x)H_0},
 \label{eq:e8v19-second}
\end{equation}
and for form vectors
\begin{align}
 |\langle u,D_x^2e^{-Tc(x)H_0}v\rangle|
 \leq
 \left(T|c''(x)|+{T\over ec_-}|c'(x)|^2\right)
 \|u\|_{1,0}\|v\|_{1,0}.
 \label{eq:e8v19-second-bound}
\end{align}
\end{corollary}

\begin{proof}
Differentiate \eqref{eq:e8v19-clock-derivative}.  The first term is bounded
as in the first part of \Cref{thm:e8v19-adjacent}.  For the second, remove
the two form weights and use
\[
 {E^2\over E+\lambda}e^{-TcE}
 \leq E e^{-Tc_-E}\leq{1\over eTc_-}.
\]
Multiplication by \(T^2|c'|^2\) gives the result.
\end{proof}

\subsection{Aligned versus noncommuting principal variations}
\label{sec:e8v19-boundary}

\begin{definition}[Aligned principal branch]
\label{def:e8v19-APB}
A marked finite-horizon family belongs to the aligned principal branch,
abbreviated APB, if after its physical unitary transport:
\begin{enumerate}[label=\textup{(A\arabic*)},leftmargin=4em]
\item every full-order marked insertion is a spectral multiplier of the same
nonnegative Hamiltonian used by every heat factor in the word;
\item its multiplier has a declared polynomial energy order;
\item lower-order noncommuting insertions are estimated separately by the
fractional branch of V17; and
\item parameter-dependent carrier and unitary derivatives are included in
the contact ledger.
\end{enumerate}
\end{definition}

\begin{theorem}[APB population of the FLAC two-mark row]
\label{thm:e8v19-APB}
For an APB word, every purely aligned subword is bounded by
\eqref{eq:e8v19-word-bound}.  In particular two full-order aligned marks
satisfy the linear clock estimate \eqref{eq:e8v19-linear-clock} and require
no \(G_T\) row.  If the remaining fractional and contact defects are
cofinally summable, the V15 (C^1) landing theorem applies.
\end{theorem}

\begin{proof}
Apply \Cref{cor:e8v19-word} to each maximal aligned subword, the V17
fractional estimate to the declared lower-order words, and then add the
contact defects.  These are precisely the adjacent bounds required by V15.
\end{proof}

\begin{firewall}[Local metric insertions are not aligned]
\label{fire:e8v19-local-metric}
A scalar global rescaling \(H\mapsto cH\) is aligned.  A local metric
variation changes the principal symbol by
\(\dot g^{ab}(x)\xi_a\xi_b\), which is generally not a Borel function of
the original scalar symbol \(g^{ab}(x)\xi_a\xi_b\) on the full phase space.
It therefore need not commute with \(H\), and
\Cref{thm:e8v19-aligned-two} does not apply.  Such transverse Einstein marks
still require SPKF, a genuine combined heat-kernel cancellation, or stronger
source regularity.
\end{firewall}

\begin{assessment}[Progress at V19]
\label{ass:e8v19-status}
The endpoint obstruction of V17 is not universal for structured form
families.  Full-order marks aligned with the physical Hamiltonian combine
exactly, and every finite aligned word retains a linear short-time clock on
form vectors.  The scalar-clock adjacent derivative now has the explicit
summable target \eqref{eq:e8v19-adjacent}.  The unresolved SM--Einstein row
has narrowed to noncommuting local principal variations and their contact
completion.
\end{assessment}

\begin{programme}[Eighth-series V20: audit and noncommuting contact target]
\label{prog:e8v19-V20}
V20 will perform the compulsory companion audit.  It will then assemble the
three two-mark branches---fractional, aligned, and combined-contact---into a
single decision theorem and select the next upstream target for local metric
variations.  No signed spectral conclusion will be imported from a companion
programme without its physical incidence.
\end{programme}

\begin{firewall}[Claim boundary after eighth-series V19]
\label{fire:e8v19-boundary}
V19 proves exact simplex reduction and energy bounds for commuting spectral
marks, proves adjacent first- and second-derivative estimates for a scalar
clock family, and defines the aligned principal branch of the FLAC.

V19 does not align local Einstein metric variations, control their
noncommuting contact terms, populate the full FLAC, remove the cutoff, or
construct the Standard Model--Einstein continuum.
\end{firewall}

\begin{theorem}[Eighth-series V19 result]
\label{thm:e8v19-result}
Two full-order insertions which are spectral multipliers of the physical
Hamiltonian have a finite ordered form integral with the linear short-time
bound \eqref{eq:e8v19-linear-clock}.  A scalar adjacent clock variation
satisfies the explicit (C^1) estimate \eqref{eq:e8v19-adjacent}.  These
results populate only the aligned branch and do not extend to generic local
metric variations.
\end{theorem}

\begin{proof}
Combine \Cref{thm:e8v19-aligned-two,
cor:e8v19-word,thm:e8v19-adjacent,cor:e8v19-second,
thm:e8v19-APB}.
\end{proof}

\paragraph{Parent-version record.}
V19 was formed from
\texttt{Renewal\_SM\_Einstein\_Continuum\_Research\_Notes\_V18.tex}.
The V18 parent had (6\,653) logical source lines, (247\,801) bytes,
and SHA--256
\[
 \texttt{107747203322B94FB7795A55433DAA21430659840F777B7D7ECBE5469FDEBBB7}.
\]
The next compulsory companion audit is V20.

% =============================================================================
% END APPENDED EIGHTH-SERIES V19 MATERIAL
% =============================================================================



% =============================================================================
% BEGIN APPENDED EIGHTH-SERIES V20 MATERIAL
% =============================================================================

\section{Companion audit and holomorphic combined-contact control}
\label{sec:v8v20}

This appendix performs the compulsory companion audit and adds a fourth
rigorous branch to the V17--V19 two-mark decision.  A common-domain
holomorphic sectorial family permits Cauchy estimates to be applied to the
complete semigroup variation.  The argument never assigns an absolute norm
to the singular product of two order-one forms.  It therefore controls a
combined noncommuting second variation when a cutoff-uniform complex form
disk exists.

\subsection{Compulsory V20 companion audit}
\label{sec:e8v20-audit}

The newest completed Yang--Mills appendix at the audit time was
\begin{equation}
 \texttt{Renewal\_Yang\_Mills\_Mass\_Gap\_Research\_Notes\_V75.tex}.
 \label{eq:e8v20-YM-file}
\end{equation}
It had (33\,789) logical source lines, (1\,186\,484) bytes, one live
document terminator, and SHA--256
\[
 \texttt{2A3D9ED012EC77D3FAFE136CAD6DAFC6C610A9CC288E7777332029C23782A5E7}.
\]
A V76 filename was present but contained no V76 section and was therefore
treated as a concurrent staging copy, not as a completed mathematical
appendix.  The newest Navier--Stokes companion remained
\begin{equation}
 \texttt{Renewal\_NS\_Stability\_Research\_Notes\_V26.tex},
 \label{eq:e8v20-NS-file}
\end{equation}
with (5\,319) logical source lines, (278\,283) bytes, one live terminator,
and SHA--256
\[
 \texttt{82C887F8C2C69E4F28FAD7C3DC8DE5B9099CB5A4BF431EBA1FFF3E91935978AD}.
\]

\begin{assessment}[Yang--Mills V75 reading]
\label{ass:e8v20-YM}
V75 proves that the central one-plaquette Wilson convolution has
Casimir-heat character ratios on every fixed representation bank at large
coupling.  It then gives an exact low/high representation-band G\aa rding
compiler: a positive low-band response and a high-band transfer-loss floor
can absorb a bounded negative high-band response.  Uniform control on the
growing representation bank and transfer to the actual half-slab remain
open.  This is a signed spectral argument on a central gauge convolution;
it supplies no complex metric-form disk or SM contact estimate.
\end{assessment}

\begin{assessment}[Navier--Stokes V26 reading]
\label{ass:e8v20-NS}
V26 is unchanged and remains a pointwise rank-one Oseen-kernel norm
calculation.  It has no common-domain sectorial family, reflected transfer,
or SM--Einstein metric insertion to import.
\end{assessment}

\begin{theorem}[V20 audit decision]
\label{thm:e8v20-audit}
No model-specific companion estimate is imported.  The Yang--Mills band
compiler confirms that signed coercivity may use a loss term on a high
sector, but that conclusion remains separate from the absolute continuum
landing problem treated here.
\end{theorem}

\begin{proof}
The V75 hypotheses are character-ratio inequalities for one central
\(\mathrm{SU}(3)\) kernel and a same-owner transfer.  The present missing row
is a noncommuting complexified metric-form estimate.  Neither V75 nor NS V26
contains that datum.  Importing only the logical separation changes no
constant or hypothesis in the SM ledger.
\end{proof}

\subsection{Common-domain holomorphic sectorial forms}
\label{sec:e8v20-holomorphic}

Let \(\mathcal H\) be a Hilbert space and \(\mathcal Q\) a dense Hilbert
subspace continuously embedded in \(\mathcal H\).  Let
\(D_R=\{z\in\mathbb C:|z|<R\}\).  For \(z\in D_R\), let \(a_z\) be a closed
sectorial form on the common domain \(\mathcal Q\), with associated
\(m\)-sectorial operator \(H_z\) and semigroup
\begin{equation}
 S_z(t)=e^{-tH_z}.
 \label{eq:e8v20-Sz}
\end{equation}
Assume \(z\mapsto a_z[u,v]\) is holomorphic for every \(u,v\in\mathcal Q\),
and that the sectorial vertex and semi-angle are uniform on compact subdisks.

\begin{lemma}[Weak semigroup holomorphy]
\label{lem:e8v20-semigroup-holomorphy}
For \(t>0\) and \(u,v\in\mathcal H\),
\begin{equation}
 z\longmapsto\langle u,S_z(t)v\rangle
 \label{eq:e8v20-holomorphic-matrix}
\end{equation}
is holomorphic on \(D_R\).
\end{lemma}

\begin{proof}
For a common-domain holomorphic sectorial family, the variational resolvent
equation
\[
 a_z[u,w]+\zeta\langle u,w\rangle=\langle f,w\rangle
 \qquad(w\in\mathcal Q)
\]
has a locally uniform coercivity constant when \(\zeta\) lies on a contour
outside the common sector.  The bounded inverse theorem and the resolvent
identity make \((H_z+\zeta)^{-1}\) holomorphic in \(z\).  The Dunford
integral for \(e^{-tH_z}\), on one contour valid on a compact subdisk,
converges locally uniformly.  Its matrix elements are therefore
holomorphic.
\end{proof}

Fix \(0<r<R\).  Suppose for \(z\in\overline D_r\) and \\(0\leq t\leq T\),
\begin{equation}
 \|S_z(t)\|_{\mathcal Q\to\mathcal Q}
 +\|S_z(t)^*\|_{\mathcal Q\to\mathcal Q}
 \leq2M_r,
 \label{eq:e8v20-Q-semigroup}
\end{equation}
and
\begin{equation}
 \|\widehat{a_z-a_0}\|_{\mathcal Q\to\mathcal Q^*}
 \leq L_r|z|.
 \label{eq:e8v20-Lipschitz-form}
\end{equation}

\begin{lemma}[One-combined-variation bound on the complex circle]
\label{lem:e8v20-circle-bound}
For \(u,v\in\mathcal Q\), \(z\in\overline D_r\), and
\(0\leq t\leq T\),
\begin{equation}
 \boxed{
 |\langle u,(S_z(t)-S_0(t))v\rangle|
 \leq tM_r^2L_r|z|\|u\|_{\mathcal Q}\|v\|_{\mathcal Q}.}
 \label{eq:e8v20-circle-bound}
\end{equation}
\end{lemma}

\begin{proof}
The weak Duhamel identity for two sectorial forms on a common domain is
\begin{align}
 \langle u,(S_z(t)-S_0(t))v\rangle
 =-\int_0^t(a_z-a_0)
 [S_z(t-s)^*u,S_0(s)v]\,ds.
 \label{eq:e8v20-sectorial-Duhamel}
\end{align}
It follows first on operator cores by differentiating
\(\langle S_z(t-s)^*u,S_0(s)v\rangle\) and then by form closure.  Apply
\eqref{eq:e8v20-Q-semigroup} and
\eqref{eq:e8v20-Lipschitz-form} to the integrand and integrate its constant
majorant over an interval of length \(t\).
\end{proof}

\begin{theorem}[Cauchy combined-contact estimate]
\label{thm:e8v20-Cauchy-contact}
Under the preceding hypotheses, for every integer \(m\geq1\),
\begin{equation}
 \boxed{
 \left|{d^m\over dz^m}\Big|_{z=0}
 \langle u,S_z(t)v\rangle\right|
 \leq m!\,tM_r^2L_r r^{1-m}
 \|u\|_{\mathcal Q}\|v\|_{\mathcal Q}.}
 \label{eq:e8v20-Cauchy-contact}
\end{equation}
In particular, the complete second variation satisfies
\begin{equation}
 \left|{d^2\over dz^2}\Big|_{z=0}
 \langle u,S_z(t)v\rangle\right|
 \leq {2tM_r^2L_r\over r}
 \|u\|_{\mathcal Q}\|v\|_{\mathcal Q}.
 \label{eq:e8v20-second-contact}
\end{equation}
No form-dual smoothing integral occurs.
\end{theorem}

\begin{proof}
Let
\(f(z)=\langle u,(S_z(t)-S_0(t))v\rangle\).  It is holomorphic by
\Cref{lem:e8v20-semigroup-holomorphy}, vanishes at zero, and on
\(|z|=r\) obeys
\(|f(z)|\leq tM_r^2L_rr\|u\|_{\mathcal Q}\|v\|_{\mathcal Q}\) by
\Cref{lem:e8v20-circle-bound}.  Cauchy's formula gives
\[
 f^{(m)}(0)={m!\over2\pi\mathrm i}
 \int_{|z|=r}{f(z)\over z^{m+1}}\,dz.
\]
Taking absolute values proves \eqref{eq:e8v20-Cauchy-contact}; \(m=2\)
gives \eqref{eq:e8v20-second-contact}.
\end{proof}

\begin{remark}[The contact is already combined]
\label{rem:e8v20-combined}
If \(a_z=a_0+zb_1+z^2b_2/2+\cdots\), the second derivative in
\eqref{eq:e8v20-second-contact} contains both time orders of \(b_1\) and
the direct contact \(b_2\).  Cauchy's formula bounds their sum as the
coefficient of one holomorphic semigroup.  It makes no claim that the three
pieces are separately absolutely integrable on the natural form scale.
\end{remark}

\begin{corollary}[Relative-radius power count]
\label{cor:e8v20-radius}
Suppose \(a_z=a_0+zb\) and
\(\|\widehat b\|_{\mathcal Q\to\mathcal Q^*}\leq\eta\).  If uniform
sectoriality and \eqref{eq:e8v20-Q-semigroup} hold on
\(|z|\leq\theta/\eta\) for a fixed \(0<\theta<1\), then
\begin{equation}
 \left|{d^2\over dz^2}\Big|_{z=0}
 \langle u,S_z(t)v\rangle\right|
 \leq {2tM_r^2\eta^2\over\theta}
 \|u\|_{\mathcal Q}\|v\|_{\mathcal Q}.
 \label{eq:e8v20-eta2}
\end{equation}
Thus the expected quadratic insertion size is recovered without the
nonintegrable endpoint factorization.
\end{corollary}

\begin{proof}
Use \(L_r=\eta\) and \(r=\theta/\eta\) in
\eqref{eq:e8v20-second-contact}.
\end{proof}

\subsection{Adjacent-scale derivative without a Dyson split}
\label{sec:e8v20-adjacent}

Let \(a_{n,z}\) and \(a_{n+1,z}\) be two common-domain holomorphic sectorial
families on the same comparison space for \(|z|<R\).  Assume the semigroup
bounds \eqref{eq:e8v20-Q-semigroup} hold for both families with \(M_r\), and
suppose
\begin{equation}
 \sup_{|z|\leq r}
 \|\widehat{a_{n+1,z}-a_{n,z}}\|_{\mathcal Q\to\mathcal Q^*}
 \leq\varepsilon_n(r).
 \label{eq:e8v20-adjacent-form}
\end{equation}

\begin{theorem}[Holomorphic adjacent (C^1) estimate]
\label{thm:e8v20-adjacent-C1}
For \(u,v\in\mathcal Q\) and \(0\leq t\leq T\),
\begin{align}
 |\langle u,(S_{n+1,0}(t)-S_{n,0}(t))v\rangle|
 &\leq tM_r^2\varepsilon_n(r)
 \|u\|_{\mathcal Q}\|v\|_{\mathcal Q},
 \label{eq:e8v20-adjacent-base}\\
 \left|{d\over dz}\Big|_{z=0}
 \langle u,(S_{n+1,z}(t)-S_{n,z}(t))v\rangle\right|
 &\leq {tM_r^2\varepsilon_n(r)\over r}
 \|u\|_{\mathcal Q}\|v\|_{\mathcal Q}.
 \label{eq:e8v20-adjacent-derivative}
\end{align}
These estimates allow noncommuting order-one base and metric forms.
\end{theorem}

\begin{proof}
Apply the weak sectorial Duhamel identity between \(a_{n+1,z}\) and
\(a_{n,z}\).  Uniformly on \(|z|\leq r\), its matrix element is bounded by
\(tM_r^2\varepsilon_n(r)\|u\|_{\mathcal Q}\|v\|_{\mathcal Q}\).
At \(z=0\) this is \eqref{eq:e8v20-adjacent-base}.  The difference of the
two matrix elements is holomorphic in \(z\); Cauchy's first derivative
estimate on the circle of radius \(r\) gives
\eqref{eq:e8v20-adjacent-derivative}.
\end{proof}

\begin{corollary}[Holomorphic FLAC branch]
\label{cor:e8v20-HFLAC}
For a finite carrier bank with energy moment \(E_{n,A}\), the generator
parts of the adjacent FLAC defects may be replaced by
\begin{align}
 d_{n,A}^{(0),\mathrm{hol}}(T)
 &=TM_r^2\varepsilon_n(r)E_{n,A},
 \label{eq:e8v20-hol-d0}\\
 d_{n,A}^{(1),\mathrm{hol}}(T)
 &={TM_r^2\varepsilon_n(r)\over r}E_{n,A}.
 \label{eq:e8v20-hol-d1}
\end{align}
If their sums and the carrier/contact rows are finite, the V15 (C^1)
landing theorem applies without a separate \(G_T\) estimate.
\end{corollary}

\begin{proof}
Apply \Cref{thm:e8v20-adjacent-C1} to every carrier vector, sum with its
positive weight, and use the definition of \(E_{n,A}\).  The resulting
bounds are the hypotheses of the V15 telescoping theorem.
\end{proof}

\begin{remark}[Spatial form locality still matters]
\label{rem:e8v20-spatial}
The scalar \(\varepsilon_n(r)\) must be obtained after summing the local
near, far, and boundary form components uniformly on the complex disk.  A
real-axis locality bound does not automatically survive complexification,
and the Cauchy theorem does not manufacture spatial decay.
\end{remark}

\subsection{The completed two-mark decision}
\label{sec:e8v20-decision}

\begin{theorem}[Four-branch two-mark decision]
\label{thm:e8v20-four-branch}
For an unbounded physical Hamiltonian, the FLAC two-mark row is rigorously
populated by any one of the following proved alternatives:
\begin{enumerate}[label=\textup{(D\arabic*)},leftmargin=4em]
\item a fixed \(s<1\) fractional extension, using V17;
\item a physically collared word, retaining the V17 logarithmic collar
debit when the collar changes;
\item a common spectral alignment, using V19; or
\item a cutoff-uniform common-domain holomorphic sectorial disk, using
\Cref{thm:e8v20-Cauchy-contact} or
\Cref{thm:e8v20-adjacent-C1}.
\end{enumerate}
Absent one of these rows, bounded cutoff Dyson integrals do not supply a
cofinal two-mark estimate.
\end{theorem}

\begin{proof}
Sufficiency of (D1) is the V17 subcritical theorem.  At fixed collar, (D2)
is the collared estimate of V17.  Sufficiency of (D3) is the V19 exact
simplex theorem.  Sufficiency of (D4) is
\Cref{thm:e8v20-Cauchy-contact,thm:e8v20-adjacent-C1}.  V17 proves failure of
the remaining generic absolute natural-form route.
\end{proof}

\begin{assessment}[Best current branch for local metric variations]
\label{ass:e8v20-metric}
V18 rules out unitary freezing of transverse Einstein directions, and V19
rules out treating them as aligned spectral multipliers.  The holomorphic
branch is therefore the first viable general target.  It requires a complex
metric neighbourhood with a regulator-independent sectorial radius,
uniform form-semigroup bounds, and a spatially summable adjacent form defect.
Gauge fixing, ghosts, chiral determinants, moving measures, and boundary
contacts must be included in the same complexified form; omitting them would
change the derivative being bounded.
\end{assessment}

\begin{programme}[Eighth-series V21: sectorial radius compiler]
\label{prog:e8v20-V21}
V21 will construct the complex sectorial disk for a uniformly elliptic
bosonic kinetic form with a bounded relative potential.  It will give an
explicit radius and semigroup form-norm bound in terms of ellipticity and
relative-form margins, then list which SM--Einstein sectors fit that theorem
and which still require separate determinant, gauge, or domain input.
\end{programme}

\begin{firewall}[Claim boundary after eighth-series V20]
\label{fire:e8v20-boundary}
V20 records the compulsory companion audit; proves weak holomorphy and a
complex-circle Duhamel bound for common-domain sectorial forms; proves
combined second-variation and adjacent (C^1) Cauchy estimates; and completes
the four-branch FLAC decision.

V20 does not construct a cutoff-uniform complex sectorial disk for the
SM--Einstein action, control its spatial complexified defects or contacts,
populate the full FLAC, prove signed coercivity, remove the cutoff, or
construct the Standard Model--Einstein continuum.
\end{firewall}

\begin{theorem}[Eighth-series V20 result]
\label{thm:e8v20-result}
A cutoff-uniform holomorphic sectorial disk controls the complete
noncommuting form variation by \eqref{eq:e8v20-Cauchy-contact} and controls
adjacent parameter derivatives by
\eqref{eq:e8v20-adjacent-derivative}.  These estimates replace the
nonintegrable split Dyson bound and, with summable local energy and contact
rows, feed directly into the V15 (C^1) landing theorem.
\end{theorem}

\begin{proof}
Combine \Cref{lem:e8v20-semigroup-holomorphy,
lem:e8v20-circle-bound,thm:e8v20-Cauchy-contact,
thm:e8v20-adjacent-C1,cor:e8v20-HFLAC,
thm:e8v20-four-branch}.
\end{proof}

\paragraph{Parent-version record.}
V20 was formed from
\texttt{Renewal\_SM\_Einstein\_Continuum\_Research\_Notes\_V19.tex}.
The V19 parent had (6\,995) logical source lines, (260\,258) bytes,
and SHA--256
\[
 \texttt{991CAC1A59A90DD15F7B6AF8330636058206D1C85D3058DD33A92299F36879DB}.
\]
The next compulsory companion audit is V25.

% =============================================================================
% END APPENDED EIGHTH-SERIES V20 MATERIAL
% =============================================================================



% =============================================================================
% BEGIN APPENDED EIGHTH-SERIES V21 MATERIAL
% =============================================================================

\section{A sectorial-radius compiler for complexified kinetic forms}
\label{sec:v8v21}

V20 reduced noncommuting two-mark control to a cutoff-uniform holomorphic
sectorial disk with semigroup bounds on one common form domain.  This appendix
constructs the disk from a relative perturbation margin and separates the
additional square-root estimate needed to obtain the form-domain semigroup
bound.  An elliptic coefficient corollary gives an explicit complex radius
for the bosonic kinetic sectors.

\subsection{Relative complex perturbations}
\label{sec:e8v21-relative}

Let \(\mathcal Q\) be densely and continuously embedded in a Hilbert space
\(\mathcal H\).  Let \(a_0\) be a densely defined closed nonnegative
symmetric form on \(\mathcal Q\), fix \(\lambda>0\), and put
\begin{equation}
 \|u\|_{Q}^2:=a_0[u,u]+\lambda\|u\|^2.
 \label{eq:e8v21-Q}
\end{equation}
Let \(p_z\) be forms on \(\mathcal Q\), holomorphic in
\(z\in D_R\), with \(p_0=0\).  Assume that for some \(\kappa>0\),
\begin{equation}
 |p_z[u,u]|\leq\kappa|z|\|u\|_Q^2,
 \qquad u\in\mathcal Q.
 \label{eq:e8v21-relative}
\end{equation}
Define \(a_z=a_0+p_z\).

\begin{theorem}[Explicit sectorial disk]
\label{thm:e8v21-sectorial-disk}
Fix \(0<\theta<1\) and set
\begin{equation}
 r_\theta:=\min\{R,\theta/\kappa\}.
 \label{eq:e8v21-radius}
\end{equation}
Assume additionally, when the bilinear estimate below is used, that
\[
 |p_z[u,v]|\leq\kappa|z|\|u\|_Q\|v\|_Q.
\]
For \(|z|\leq r_\theta\), the shifted form
\begin{equation}
 a_{z,\lambda}[u,v]:=a_z[u,v]+\lambda\langle u,v\rangle
 \label{eq:e8v21-shifted}
\end{equation}
is closed, sectorial, and satisfies
\begin{align}
 \operatorname{Re}a_{z,\lambda}[u,u]
 &\geq(1-\theta)\|u\|_Q^2,
 \label{eq:e8v21-coercive}\\
 |\operatorname{Im}a_{z,\lambda}[u,u]|
 &\leq{\theta\over1-\theta}
 \operatorname{Re}a_{z,\lambda}[u,u],
 \label{eq:e8v21-angle}\\
 |a_{z,\lambda}[u,v]|
 &\leq(1+\theta)\|u\|_Q\|v\|_Q
 \quad\text{under the additional direct bilinear hypothesis}.
 \label{eq:e8v21-bilinear}
\end{align}
The sectorial vertex and angle are uniform in the regulator whenever
\(\kappa,R,\lambda\), and \(\theta\) are.
\end{theorem}

\begin{proof}
The real and imaginary parts satisfy
\begin{align*}
 \operatorname{Re}a_{z,\lambda}[u,u]
 &=\|u\|_Q^2+\operatorname{Re}p_z[u,u]
 \geq(1-\kappa|z|)\|u\|_Q^2,\\
 |\operatorname{Im}a_{z,\lambda}[u,u]|
 &\leq\kappa|z|\|u\|_Q^2.
\end{align*}
Use \(\kappa|z|\leq\theta\) to obtain
\eqref{eq:e8v21-coercive}--\eqref{eq:e8v21-angle}.  The shifted form norm is
therefore equivalent to the complete norm \(\|\cdot\|_Q\), so the form is
closed.  Under the additional direct bilinear hypothesis, add the reference inner product
to prove \eqref{eq:e8v21-bilinear}.
\end{proof}

\begin{remark}[Nonsymmetric polarization constant]
\label{rem:e8v21-polarization}
For a general sectorial perturbation, a diagonal numerical-radius bound does
not give \eqref{eq:e8v21-bilinear} with constant one.  The V20 complex-circle
Duhamel estimate therefore requires either a direct bilinear bound or the
corresponding numerical-radius polarization constant.  This constant must be
included in \(L_r\).
\end{remark}

\subsection{The square-root row}
\label{sec:e8v21-square-root}

Let \(H_z\) be the \(m\)-sectorial operator associated with \(a_z\).  The
common form domain need not equal
\(\operatorname{Dom}(H_z+\lambda)^{1/2}\) for an arbitrary sectorial form.
Assume the uniform square-root estimates
\begin{align}
 c_{\rm sq}\|u\|_Q
 &\leq\|(H_z+\lambda)^{1/2}u\|
 \leq C_{\rm sq}\|u\|_Q,
 \label{eq:e8v21-square-root}\\
 c_{\rm sq}\|u\|_Q
 &\leq\|(H_z^*+\lambda)^{1/2}u\|
 \leq C_{\rm sq}\|u\|_Q
 \label{eq:e8v21-adjoint-square-root}
\end{align}
for \(|z|\leq r_\theta\).

\begin{theorem}[Square-root-to-semigroup compiler]
\label{thm:e8v21-Q-semigroup}
Under \eqref{eq:e8v21-square-root}--
\eqref{eq:e8v21-adjoint-square-root}, for \(0\leq t\leq T\),
\begin{align}
 \|e^{-tH_z}\|_{\mathcal Q\to\mathcal Q}
 &\leq {C_{\rm sq}\over c_{\rm sq}}e^{\lambda t},
 \label{eq:e8v21-Q-bound}\\
 \|e^{-tH_z^*}\|_{\mathcal Q\to\mathcal Q}
 &\leq {C_{\rm sq}\over c_{\rm sq}}e^{\lambda t}.
 \label{eq:e8v21-Qstar-bound}
\end{align}
Hence the V20 constant may be taken as
\begin{equation}
 M_r={C_{\rm sq}\over c_{\rm sq}}e^{\lambda T}.
 \label{eq:e8v21-Mr}
\end{equation}
\end{theorem}

\begin{proof}
The shifted operator \(H_z+\lambda\) is \(m\)-accretive by
\eqref{eq:e8v21-coercive}; hence
\(e^{-t(H_z+\lambda)}\) is a contraction.  Fractional powers commute with
the holomorphic functional calculus of \(H_z\), so for \(u\in\mathcal Q\),
\begin{align*}
 \|e^{-tH_z}u\|_Q
 &\leq c_{\rm sq}^{-1}
 \|(H_z+\lambda)^{1/2}e^{-tH_z}u\|\\
 &=c_{\rm sq}^{-1}
 \|e^{-tH_z}(H_z+\lambda)^{1/2}u\|\\
 &\leq c_{\rm sq}^{-1}e^{\lambda t}
 \|(H_z+\lambda)^{1/2}u\|\\
 &\leq{C_{\rm sq}\over c_{\rm sq}}e^{\lambda t}\|u\|_Q.
\end{align*}
Apply the same argument to \(H_z^*\).
\end{proof}

\begin{firewall}[Sectoriality is not the square-root property]
\label{fire:e8v21-square-root}
\Cref{thm:e8v21-sectorial-disk} constructs a common closed sectorial form
disk.  It does not prove
\eqref{eq:e8v21-square-root}--\eqref{eq:e8v21-adjoint-square-root}.  The
latter is automatic on each finite-dimensional cutoff but its constants may
deteriorate with refinement.  A cofinal use of V20 requires uniform constants,
obtained from a square-root theorem for the actual discrete or continuum
operator family.
\end{firewall}

\subsection{Uniformly elliptic bosonic coefficients}
\label{sec:e8v21-elliptic}

Let \(M\) carry a fixed positive measure \(d\mu\), and let
\(\mathcal Q=H^1(M;\mathbb C^N)\) with the boundary condition declared by the
regulator.  Consider
\begin{equation}
 a_0[u,u]
 =\int_M\langle\nabla u,G_0(x)\nabla u\rangle\,d\mu
 +v_0[u,u],
 \label{eq:e8v21-a0}
\end{equation}
where \(v_0\) is nonnegative and
\begin{equation}
 \nu|\xi|^2\leq\langle\xi,G_0(x)\xi\rangle
 \leq\Lambda|\xi|^2
 \label{eq:e8v21-ellipticity}
\end{equation}
with \(0<\nu\leq\Lambda<\infty\).  Let the complexified perturbation be
\begin{equation}
 p_z[u,u]
 =\int_M\langle\nabla u,\Delta G_z(x)\nabla u\rangle\,d\mu
 +\Delta v_z[u,u].
 \label{eq:e8v21-pz}
\end{equation}
Assume
\begin{align}
 \|\Delta G_z\|_{L^\infty}&\leq L_G|z|,
 \label{eq:e8v21-G-bound}\\
 |\Delta v_z[u,u]|
 &\leq |z|\{\alpha_v a_0[u,u]+\beta_v\|u\|^2\}.
 \label{eq:e8v21-v-bound}
\end{align}

\begin{corollary}[Explicit elliptic complex radius]
\label{cor:e8v21-elliptic-radius}
For the shifted norm \eqref{eq:e8v21-Q},
\begin{equation}
 |p_z[u,u]|
 \leq\kappa_{\rm ell}|z|\|u\|_Q^2,
 \qquad
 \kappa_{\rm ell}:={L_G\over\nu}+\alpha_v+{\beta_v\over\lambda}.
 \label{eq:e8v21-kappa-ell}
\end{equation}
Consequently the disk
\begin{equation}
 |z|\leq r_\theta^{\rm ell}
 :=\min\left\{R,{\theta\over\kappa_{\rm ell}}\right\}
 \label{eq:e8v21-r-ell}
\end{equation}
has the uniform coercivity and angle bounds
\eqref{eq:e8v21-coercive}--\eqref{eq:e8v21-angle}.  If the ellipticity,
relative-potential, and square-root constants are regulator independent,
then this disk satisfies the analytic and form-semigroup parts of the V20
holomorphic branch.
\end{corollary}

\begin{proof}
The principal perturbation is bounded by
\[
 L_G|z|\|\nabla u\|_2^2
 \leq{L_G\over\nu}|z|a_0[u,u]
\]
after retaining the nonnegative reference lower-order part, or after charging
its negative part to the declared relative constants.  Equation
\eqref{eq:e8v21-v-bound} supplies the other two terms, and
\(\beta_v\|u\|^2\leq(\beta_v/\lambda)\|u\|_Q^2\).  Add them and apply
\Cref{thm:e8v21-sectorial-disk,thm:e8v21-Q-semigroup}.
\end{proof}

\begin{remark}[Moving measures]
\label{rem:e8v21-measure}
A metric-dependent volume form must first be transported to the fixed
\(L^2(d\mu)\) space by its positive density square root.  Differentiating
that transport creates lower-order connection and contact terms.  They belong
in \(\Delta v_z\) and in the carrier contact ledger.  Treating the measure as
fixed without this transport would omit part of the Einstein stress.
\end{remark}

\subsection{Local component and adjacent-scale bounds}
\label{sec:e8v21-local}

Suppose the complex perturbation decomposes as
\begin{equation}
 p_{n,z}=\sum_Xp_{n,z;X}
 \label{eq:e8v21-component}
\end{equation}
in the form topology, and that
\begin{equation}
 |p_{n,z;X}[u,v]|
 \leq |z|\kappa_{n,X}\|u\|_{Q,X}\|v\|_{Q,X}
 \label{eq:e8v21-component-bound}
\end{equation}
with localized energy seminorms.  Assume the embedding into the global form
norm has constants \(c_X\) and
\begin{equation}
 \sum_Xc_X^2\kappa_{n,X}\leq\kappa_*<\infty.
 \label{eq:e8v21-component-sum}
\end{equation}

\begin{lemma}[Spatial summation of the complex radius]
\label{lem:e8v21-spatial-sum}
Under \eqref{eq:e8v21-component-bound}--
\eqref{eq:e8v21-component-sum},
\begin{equation}
 |p_{n,z}[u,v]|
 \leq |z|\kappa_*\|u\|_Q\|v\|_Q.
 \label{eq:e8v21-global-component}
\end{equation}
Thus a weighted local component ledger gives a regulator-independent complex
radius whenever \(\kappa_*\) is uniform.
\end{lemma}

\begin{proof}
Absolute convergence in the form topology permits termwise estimation.
Use \(\|u\|_{Q,X}\leq c_X\|u\|_Q\) and the same bound for \(v\), then sum
\eqref{eq:e8v21-component-bound} and apply
\eqref{eq:e8v21-component-sum}.
\end{proof}

\begin{corollary}[Adjacent complex form defect]
\label{cor:e8v21-adjacent}
If, on the same disk,
\begin{equation}
 a_{n+1,z}-a_{n,z}=\sum_Xd_{n,z;X},
 \qquad
 |d_{n,z;X}[u,v]|
 \leq\varepsilon_{n,X}(r)\|u\|_{Q,X}\|v\|_{Q,X},
 \label{eq:e8v21-adjacent-components}
\end{equation}
then the V20 adjacent defect may be taken as
\begin{equation}
 \varepsilon_n(r)
 =\sum_Xc_X^2\varepsilon_{n,X}(r).
 \label{eq:e8v21-epsilon}
\end{equation}
In particular, summability of the weighted component tails populates the
generator part of the holomorphic FLAC branch.
\end{corollary}

\begin{proof}
Apply the proof of \Cref{lem:e8v21-spatial-sum} to the adjacent difference,
then insert \eqref{eq:e8v21-epsilon} in the V20 adjacent (C^1) theorem.
\end{proof}

\subsection{Sector applicability audit}
\label{sec:e8v21-sectors}

\begin{assessment}[Bosonic gauge and Higgs squares]
\label{ass:e8v21-bosonic}
After gauge fixing and transport to a fixed measure, positive
Laplace-type gauge and Higgs quadratic forms fit
\Cref{cor:e8v21-elliptic-radius} on a uniformly nondegenerate metric chart,
provided their interaction potentials have a uniform relative-form margin.
Polynomial scalar interactions require the common domain and relative bound
uniformly under refinement; formal positivity at each cutoff is not enough.
A lattice-to-continuum square-root estimate remains required.
\end{assessment}

\begin{assessment}[Fermions and determinant lines]
\label{ass:e8v21-fermions}
A first-order chiral Dirac operator is not a nonnegative form generator.
Its positive square may enter an elliptic heat estimate, but the phase,
Pfaffian or determinant line, chiral projector, and anomaly contacts are not
reconstructed from that square.  Those rows remain separate and must be
complexified on the same physical branch before a full stress derivative is
claimed.
\end{assessment}

\begin{assessment}[Gravity]
\label{ass:e8v21-gravity}
The Euclidean Einstein--Hilbert conformal direction is not a nonnegative
uniformly elliptic matter form.  The present theorem applies to matter
Hamiltonians on a prescribed uniformly nondegenerate metric chart.  It does
not construct a positive gravitational OS Hamiltonian or cure the conformal
factor problem.  The gravity carrier, constraint quotient, and reflected
positive clock remain independent continuum gates.
\end{assessment}

\begin{definition}[Holomorphic elliptic form card, HEFC]
\label{def:e8v21-HEFC}
For each cofinal regulator and physical parameter chart, HEFC consists of:
\begin{enumerate}[label=\textup{(H\arabic*)},leftmargin=4em]
\item a fixed comparison Hilbert space and common form domain;
\item regulator-uniform \(\nu,\Lambda,\alpha_v,\beta_v,\lambda\), and hence
a positive radius \eqref{eq:e8v21-r-ell};
\item uniform square-root estimates for the operator and its adjoint;
\item the complexified local component sum
\eqref{eq:e8v21-component-sum};
\item the summable adjacent complex defect
\eqref{eq:e8v21-epsilon}; and
\item every moving-measure, gauge-fixing, boundary, carrier, and physical
incidence contact.
\end{enumerate}
\end{definition}

\begin{theorem}[HEFC-to-holomorphic-FLAC handoff]
\label{thm:e8v21-HEFC}
If HEFC holds and the V16 local energy rows are summable with the contact
rows, then the matter-sector local transfer matrix elements land in \(C^1\)
by the V20 holomorphic adjacent theorem and the V15 landing theorem.
\end{theorem}

\begin{proof}
Rows (H1), (H2), and (H4) give the common holomorphic sectorial disk by
\Cref{thm:e8v21-sectorial-disk,lem:e8v21-spatial-sum}.  Row (H3) gives the
form-semigroup bounds by \Cref{thm:e8v21-Q-semigroup}.  Row (H5) gives the
V20 adjacent complex defect.  Multiply its scalar estimates by the V16
energy moment, add row (H6), and invoke the V15 telescoping theorem.
\end{proof}

\begin{assessment}[Progress at V21]
\label{ass:e8v21-status}
The complex-radius part of the V20 branch is now explicit for uniformly
elliptic positive bosonic forms.  The remaining analytic operator row is no
longer hidden: it is the uniform square-root property for both the
complexified operator and its adjoint.  The physical matter landing also
needs a summable adjacent complex component ledger.  Fermion-line and gravity
positivity gates remain outside this bosonic theorem.
\end{assessment}

\begin{programme}[Eighth-series V22: uniform discrete square-root row]
\label{prog:e8v21-V22}
V22 will test the square-root row first on finite weighted graphs with a
uniformly elliptic divergence-form kinetic operator.  It will seek constants
independent of graph size and mesh by an exact gradient-factorization
argument.  If arbitrary complex coefficients defeat that argument, it will
state the additional normality or accretive-factor condition needed for the
HEFC rather than importing a continuum theorem without its hypotheses.
\end{programme}

\begin{firewall}[Claim boundary after eighth-series V21]
\label{fire:e8v21-boundary}
V21 proves an explicit sectorial radius from a relative complex form margin,
derives form-domain semigroup bounds from a uniform operator-and-adjoint
square-root row, gives an elliptic bosonic radius, and compiles spatial
component and adjacent defects into HEFC.

V21 does not prove the uniform square-root row for the SM lattice operators,
complexify the determinant/Pfaffian and gravitational sectors, populate all
HEFC or FLAC rows, remove the cutoff, or construct the Standard
Model--Einstein continuum.
\end{firewall}

\begin{theorem}[Eighth-series V21 result]
\label{thm:e8v21-result}
Uniform ellipticity and a relative complex potential margin give the explicit
sectorial radius \eqref{eq:e8v21-r-ell}.  Uniform square-root estimates then
give the V20 form-semigroup constant \eqref{eq:e8v21-Mr}, while a summable
local complex component ledger gives the adjacent defect
\eqref{eq:e8v21-epsilon}.  These rows are sufficient for matter-sector
(C^1) landing but do not include fermion-line or gravitational positivity.
\end{theorem}

\begin{proof}
Combine \Cref{thm:e8v21-sectorial-disk,
thm:e8v21-Q-semigroup,cor:e8v21-elliptic-radius,
lem:e8v21-spatial-sum,cor:e8v21-adjacent,
thm:e8v21-HEFC}.
\end{proof}

\paragraph{Parent-version record.}
V21 was formed from
\texttt{Renewal\_SM\_Einstein\_Continuum\_Research\_Notes\_V20.tex}.
The V20 parent had \(7\,386\) logical source lines, \(275\,323\) bytes,
and SHA--256
\[
 \texttt{FA0653A6FB902AA277351F5A50BC5707B6753797F577913555C40F6596D43439}.
\]
The next compulsory companion audit is V25.

% =============================================================================
% END APPENDED EIGHTH-SERIES V21 MATERIAL
% =============================================================================




% =============================================================================
% BEGIN APPENDED EIGHTH-SERIES V22 MATERIAL
% =============================================================================

\section{Discrete square-root rows: exact branches and the nonnormal gap}
\label{sec:v8v22}

V21 showed that a common sectorial disk feeds V20 only after uniform
square-root estimates for the complexified operator and its adjoint are
proved.  This appendix treats finite weighted graph kinetic forms.  The
self-adjoint real branch and an energy-commuting complex branch have exact
cutoff-independent constants.  For the general nonnormal complex branch, a
finite singular-value card identifies the cofinal quantity which remains to
be bounded.

\subsection{Weighted graph kinetic forms}
\label{sec:e8v22-graph}

For a finite graph regulator \(n\), let \(\mathcal H_{v,n}\) and
\(\mathcal H_{e,n}\) be the weighted vertex and oriented-edge Hilbert spaces.
Let
\begin{equation}
 D_n:\mathcal H_{v,n}\longrightarrow\mathcal H_{e,n}
 \label{eq:e8v22-D}
\end{equation}
be the covariant incidence operator, including the declared gauge and frame
transport.  Let \(G_{n,x}=G_{n,x}^*\) act on edge fibres and satisfy
\begin{equation}
 \nu I\preceq G_{n,x}\preceq\Lambda I
 \label{eq:e8v22-ellipticity}
\end{equation}
with \(0<\nu\leq\Lambda\) independent of \(n,x\).  For \(\lambda>0\), put
\begin{align}
 a_{n,x,\lambda}[u,v]
 &=\langle D_nu,G_{n,x}D_nv\rangle
 +\lambda\langle u,v\rangle,
 \label{eq:e8v22-form}\\
 A_{n,x}&=D_n^*G_{n,x}D_n+\lambda I.
 \label{eq:e8v22-A}
\end{align}

\begin{theorem}[Exact real-axis square-root row]
\label{thm:e8v22-real-square-root}
For every finite graph and every real parameter \(x\),
\begin{equation}
 \boxed{
 \|A_{n,x}^{1/2}u\|^2=a_{n,x,\lambda}[u,u].}
 \label{eq:e8v22-exact-real}
\end{equation}
If \(x_0\) is a reference point and
\(\|u\|_{Q,n}^2=a_{n,x_0,\lambda}[u,u]\), then
\begin{equation}
 c_{\rm r}\|u\|_{Q,n}
 \leq\|A_{n,x}^{1/2}u\|
 \leq C_{\rm r}\|u\|_{Q,n},
 \label{eq:e8v22-real-comparison}
\end{equation}
where one may take
\begin{equation}
 c_{\rm r}=\min\{1,\sqrt{\nu/\Lambda}\},
 \qquad
 C_{\rm r}=\max\{1,\sqrt{\Lambda/\nu}\}.
 \label{eq:e8v22-real-constants}
\end{equation}
The same estimates hold for the adjoint because \(A_{n,x}\) is
self-adjoint.
\end{theorem}

\begin{proof}
The finite matrix \(A_{n,x}\) is positive definite, and its quadratic form
is exactly \eqref{eq:e8v22-form}; spectral calculus proves
\eqref{eq:e8v22-exact-real}.  Both the reference and target gradient terms
lie between \(\nu\|D_nu\|^2\) and \(\Lambda\|D_nu\|^2\), while their
\(\lambda\|u\|^2\) terms coincide.  Comparing the two nonnegative summands
gives \eqref{eq:e8v22-real-comparison}--
\eqref{eq:e8v22-real-constants}.
\end{proof}

\begin{remark}[Why the real row is insufficient for V20]
\label{rem:e8v22-real-insufficient}
Cauchy's formula samples a complex circle.  The exact self-adjoint identity
\eqref{eq:e8v22-exact-real} holds only on the real parameter axis and does
not control the nonnormal square roots at nonreal parameter values.
\end{remark}

\subsection{An exact energy-commuting complex branch}
\label{sec:e8v22-commuting}

Let \(A_{n,0}\) be the positive reference operator and suppose a complexified
shifted form can be written
\begin{equation}
 a_{n,z,\lambda}[u,v]
 =\langle A_{n,0}^{1/2}u,
 C_{n,z}A_{n,0}^{1/2}v\rangle,
 \label{eq:e8v22-C-form}
\end{equation}
where \(C_{n,z}\) is bounded and normal,
\begin{equation}
 [C_{n,z},A_{n,0}]=0.
 \label{eq:e8v22-commute}
\end{equation}
Assume its spectrum lies in a simply connected sector avoiding the negative
real axis, so the principal square root is defined.

\begin{theorem}[Energy-commuting complex square root]
\label{thm:e8v22-commuting-square-root}
The operator associated with \eqref{eq:e8v22-C-form} is
\begin{equation}
 A_{n,z}=A_{n,0}C_{n,z},
 \label{eq:e8v22-product}
\end{equation}
and
\begin{equation}
 A_{n,z}^{1/2}=A_{n,0}^{1/2}C_{n,z}^{1/2}.
 \label{eq:e8v22-product-root}
\end{equation}
If
\begin{equation}
 0<m_C\leq\inf_{\zeta\in\sigma(C_{n,z})}|\zeta|
 \leq\sup_{\zeta\in\sigma(C_{n,z})}|\zeta|\leq M_C,
 \label{eq:e8v22-annulus}
\end{equation}
then
\begin{equation}
 \boxed{
 \sqrt{m_C}\|u\|_{Q,n}
 \leq\|A_{n,z}^{1/2}u\|
 \leq\sqrt{M_C}\|u\|_{Q,n}.}
 \label{eq:e8v22-commuting-comparison}
\end{equation}
The adjoint has the same constants.
\end{theorem}

\begin{proof}
Strong commutation of the normal operators gives the joint functional
calculus.  Thus
\(A_{n,0}^{1/2}C_{n,z}A_{n,0}^{1/2}=A_{n,0}C_{n,z}\), proving
\eqref{eq:e8v22-product}; the principal joint functional calculus gives
\eqref{eq:e8v22-product-root}.  Normality gives
\[
 \sqrt{m_C}\|w\|
 \leq\|C_{n,z}^{1/2}w\|
 \leq\sqrt{M_C}\|w\|.
\]
Use \(w=A_{n,0}^{1/2}u\).  Replacing \(C_{n,z}\) by \(C_{n,z}^*\) proves
the adjoint statement.
\end{proof}

\begin{corollary}[Linear commuting disk]
\label{cor:e8v22-linear-commuting}
If \(C_{n,z}=I+zB_n\), where \(B_n=B_n^*\),
\([B_n,A_{n,0}]=0\), and \(\|B_n\|\leq\kappa\), then on
\(|z|\kappa\leq\theta<1\),
\begin{equation}
 \sqrt{1-\theta}\|u\|_{Q,n}
 \leq\|A_{n,z}^{1/2}u\|
 \leq\sqrt{1+\theta}\|u\|_{Q,n},
 \label{eq:e8v22-linear-root}
\end{equation}
with the same bounds for the adjoint.
\end{corollary}

\begin{proof}
Every spectral value has the form \(1+zb\), with \(|b|\leq\kappa\), and
therefore lies in the annulus
\(1-\theta\leq|1+zb|\leq1+\theta\).  Apply
\Cref{thm:e8v22-commuting-square-root}.
\end{proof}

\begin{firewall}[Energy commutation is a strong geometric restriction]
\label{fire:e8v22-commuting}
A spatially varying metric coefficient acts on edge gradients, whereas
\(A_{n,0}\) also contains incidence and divergence.  These operators do not
generally commute.  The theorem applies to global aligned rescalings and
other energy-diagonal variations; it does not cover a generic local Einstein
metric insertion.
\end{firewall}

\subsection{The exact finite-cutoff certificate}
\label{sec:e8v22-certificate}

At finite cutoff, let \(A_{n,z}\) be the shifted complex operator and let
\(A_{n,0}>0\) be the reference.  Define
\begin{align}
 R_{n,z}&:=A_{n,z}^{1/2}A_{n,0}^{-1/2},
 \label{eq:e8v22-R}\\
 R_{n,z}^{\dagger}&:=(A_{n,z}^*)^{1/2}A_{n,0}^{-1/2}.
 \label{eq:e8v22-Rdag}
\end{align}

\begin{theorem}[Finite square-root singular-value certificate]
\label{thm:e8v22-finite-certificate}
The V21 square-root and adjoint square-root constants on a compact complex
parameter set \(K\) are exactly
\begin{align}
 c_{n,K}&=\inf_{z\in K}
 \min\{s_{\min}(R_{n,z}),s_{\min}(R_{n,z}^{\dagger})\},
 \label{eq:e8v22-c}\\
 C_{n,K}&=\sup_{z\in K}
 \max\{s_{\max}(R_{n,z}),s_{\max}(R_{n,z}^{\dagger})\}.
 \label{eq:e8v22-C}
\end{align}
They are positive and finite at every fixed finite regulator whenever the
principal square roots exist continuously and are invertible on \(K\).  The HEFC square-root row
holds cofinally if and only if
\begin{equation}
 \boxed{
 \inf_n c_{n,K}>0,
 \qquad
 \sup_n C_{n,K}<\infty.}
 \label{eq:e8v22-cofinal-certificate}
\end{equation}
\end{theorem}

\begin{proof}
For \(w=A_{n,0}^{1/2}u\),
\[
 \|A_{n,z}^{1/2}u\|=\|R_{n,z}w\|,
 \qquad
 \|u\|_{Q,n}=\|w\|.
\]
The sharp lower and upper constants are therefore the smallest and largest
singular values of \(R_{n,z}\).  Repeat for the adjoint, then take the
infimum and supremum over \(K\).  Continuity and compactness give positive
finite constants at fixed \(n\).  Uniformity in \(n\) is exactly
\eqref{eq:e8v22-cofinal-certificate}.
\end{proof}

\begin{countertheorem}[Finite-dimensional automaticity is not cofinal]
\label{cth:e8v22-finite-not-cofinal}
The statement \(0<c_{n,K}\leq C_{n,K}<\infty\) for every \(n\) does not imply
\eqref{eq:e8v22-cofinal-certificate}.
\end{countertheorem}

\begin{proof}
Take \(K\) to be a singleton, \(A_{n,0}=I_2\), and
\(A_{n,z}=\operatorname{diag}(n^{-2},n^2)\).  Then the principal square root
is invertible at every finite \(n\), while
\(R_{n,z}=\operatorname{diag}(n^{-1},n)\).  Thus
\(c_{n,K}=n^{-1}\) and \(C_{n,K}=n\).  This proves both failures and shows
why a graph-uniform analytic theorem or a certified cofinal estimate is
required.
\end{proof}

\subsection{The discrete quadratic square-root row}
\label{sec:e8v22-DQSR}

\begin{definition}[Discrete quadratic square-root card, DQSRC]
\label{def:e8v22-DQSRC}
For the complex graph family on the V21 disk, DQSRC consists of:
\begin{enumerate}[label=\textup{(Q\arabic*)},leftmargin=4em]
\item a graph-uniform proof that
\(\operatorname{Dom}A_{n,z}^{1/2}=\operatorname{Dom}D_n\) and likewise for
the adjoint;
\item constants in those two norm equivalences depending only on the
ellipticity, sector angle, bounded degree, and declared geometric regularity,
not on vertex count or mesh;
\item stability under the gauge, frame, boundary, and moving-measure contacts
used in the physical form; and
\item compatibility of the square-root comparison with the adjacent-scale
bridge.
\end{enumerate}
The finite certificate \eqref{eq:e8v22-cofinal-certificate} is an equivalent
matrix target on each cutoff, but not a proof of its uniformity.
\end{definition}

\begin{theorem}[DQSRC closes the V21 square-root row]
\label{thm:e8v22-DQSRC}
If DQSRC holds on a fixed complex disk, then V21 gives a regulator-uniform
form-domain semigroup constant \(M_r\), and the V20 holomorphic combined
contact estimate applies to the graph kinetic sector.
\end{theorem}

\begin{proof}
Rows (Q1)--(Q2) are precisely the operator and adjoint estimates
\eqref{eq:e8v21-square-root}--
\eqref{eq:e8v21-adjoint-square-root}.  Rows (Q3)--(Q4) ensure they hold for
the physical transported family.  Apply the V21 square-root-to-semigroup
theorem and then V20.
\end{proof}

\begin{assessment}[Progress at V22]
\label{ass:e8v22-status}
The square-root row is closed exactly on the real axis and on the complex
energy-commuting branch.  Generic local metric coefficients are noncommuting,
so the unresolved analytic target is DQSRC or, equivalently, uniform control
of the two singular-value families in
\eqref{eq:e8v22-cofinal-certificate}.  The next independent HEFC row is the
summability of adjacent metric coefficients; it can be attacked without
waiting for DQSRC.
\end{assessment}

\begin{programme}[Eighth-series V23: adjacent metric approximation]
\label{prog:e8v22-V23}
V23 will take a uniformly elliptic \(C^{0,\alpha}\) metric family on a fixed
physical chart and its dyadic piecewise approximation.  It will prove that
the adjacent principal-form defect is \(O(2^{-\alpha n})\), including one
parameter derivative, and hence is summable.  This will populate the HEFC
adjacent generator row conditionally on DQSRC.
\end{programme}

\begin{firewall}[Claim boundary after eighth-series V22]
\label{fire:e8v22-boundary}
V22 proves cutoff-independent square-root constants for real self-adjoint
graph forms and energy-commuting complex forms; gives the exact finite
singular-value certificate; and isolates DQSRC as the nonnormal graph-uniform
row.

V22 does not prove DQSRC for generic complex metric coefficients, populate
the remaining HEFC or FLAC rows, control determinant or gravitational
positivity, remove the cutoff, or construct the Standard Model--Einstein
continuum.
\end{firewall}

\begin{theorem}[Eighth-series V22 result]
\label{thm:e8v22-result}
The discrete square-root row is automatic with uniform constants on the real
self-adjoint branch and on the complex energy-commuting branch
\eqref{eq:e8v22-commuting-comparison}.  In the general finite-cutoff branch,
its exact cofinal content is the singular-value condition
\eqref{eq:e8v22-cofinal-certificate}; finite-dimensional existence alone
does not prove that condition.
\end{theorem}

\begin{proof}
Combine \Cref{thm:e8v22-real-square-root,
thm:e8v22-commuting-square-root,cor:e8v22-linear-commuting,
thm:e8v22-finite-certificate,cth:e8v22-finite-not-cofinal,
thm:e8v22-DQSRC}.
\end{proof}

\paragraph{Parent-version record.}
V22 was formed from
\texttt{Renewal\_SM\_Einstein\_Continuum\_Research\_Notes\_V21.tex}.
The V21 parent had (7\,823) logical source lines, (291\,914) bytes,
and SHA--256
\[
 \texttt{2A7061A4EB082502476DA3178B74AD90D452A6653288609AA7639DC180084D90}.
\]
The next compulsory companion audit is V25.

% =============================================================================
% END APPENDED EIGHTH-SERIES V22 MATERIAL
% =============================================================================





% =============================================================================
% BEGIN APPENDED EIGHTH-SERIES V23 MATERIAL
% =============================================================================

\section{Summable adjacent metric forms from Hölder approximation}
\label{sec:v8v23}

The nonnormal square-root row remains open, but the adjacent complex form
defect in HEFC can be settled independently.  This appendix proves that
piecewise dyadic approximation of a uniformly elliptic
\(C^{0,\alpha}\) metric coefficient has a geometrically summable adjacent
form tail.  The same estimate holds on a fixed complex parameter disk and
therefore feeds the V20 derivative estimate by Cauchy's formula.

\subsection{A stable approximation operator}
\label{sec:e8v23-approximation}

Let \(M\) be a compact manifold covered by a fixed finite atlas, or one
bounded Lipschitz chart.  Let \(h_n=h_0 2^{-n}\).  Suppose
\(P_n:C^{0,\alpha}(M;\mathbb C^{N\times N})\to L^\infty(M)\) is a
piecewise approximation satisfying
\begin{equation}
 \|P_nF-F\|_{L^\infty}
 \leq C_{\rm app}h_n^\alpha[F]_{C^{0,\alpha}},
 \qquad 0<\alpha\leq1.
 \label{eq:e8v23-approximation}
\end{equation}
Cell averages, nodal interpolation on a shape-regular mesh, and cell-center
sampling have this property with their usual chart-dependent constants.

\begin{lemma}[Adjacent coefficient estimate]
\label{lem:e8v23-adjacent-coefficient}
For every \(F\in C^{0,\alpha}\),
\begin{equation}
 \boxed{
 \|P_{n+1}F-P_nF\|_{L^\infty}
 \leq C_{\rm app}(1+2^{-\alpha})h_n^\alpha
 [F]_{C^{0,\alpha}}.}
 \label{eq:e8v23-adjacent-coefficient}
\end{equation}
Consequently,
\begin{equation}
 \sum_{n\geq N}\|P_{n+1}F-P_nF\|_{L^\infty}
 \leq {C_{\rm app}(1+2^{-\alpha})
 h_N^\alpha\over1-2^{-\alpha}}
 [F]_{C^{0,\alpha}}.
 \label{eq:e8v23-coefficient-tail}
\end{equation}
\end{lemma}

\begin{proof}
The triangle inequality and \eqref{eq:e8v23-approximation} give
\[
 \|P_{n+1}F-P_nF\|_\infty
 \leq C_{\rm app}(h_{n+1}^\alpha+h_n^\alpha)[F]_{C^{0,\alpha}}.
\]
Use \(h_{n+1}=h_n/2\), then sum the geometric series.
\end{proof}

\subsection{Complexified uniformly elliptic metrics}
\label{sec:e8v23-complex-metric}

Let \(K\) be a compact real parameter chart and let
\(z\mapsto G(z,\cdot)\) be holomorphic on \(|z|<R\), with values in
\(C^{0,\alpha}(M;\mathbb C^{dN\times dN})\).  On a fixed disk
\(|z|\leq r<R\), assume
\begin{align}
 \sup_{|z|\leq r}[G(z)]_{C^{0,\alpha}}&\leq L_{G,\alpha}(r),
 \label{eq:e8v23-G-holder}\\
 \operatorname{Re}\langle\xi,G(z,x)\xi\rangle
 &\geq\nu_r|\xi|^2,
 \qquad \nu_r>0.
 \label{eq:e8v23-G-elliptic}
\end{align}
Put \(G_n(z)=P_nG(z)\).  Let
\begin{equation}
 a_{n,z,\lambda}[u,v]
 =\int_M\langle\nabla u,G_n(z)\nabla v\rangle\,d\mu
 +v_{n,z}[u,v]+\lambda\langle u,v\rangle
 \label{eq:e8v23-an}
\end{equation}
on one common domain \(\mathcal Q=H^1(M;\mathbb C^N)\).  Use a reference
norm satisfying
\begin{equation}
 \|u\|_Q^2\geq\nu_0\|\nabla u\|_2^2+\lambda_0\|u\|_2^2
 \label{eq:e8v23-Q-lower}
\end{equation}
with \(\nu_0,\lambda_0>0\).

Assume the lower-order adjacent forms obey, uniformly on the complex disk,
\begin{equation}
 |(v_{n+1,z}-v_{n,z})[u,v]|
 \leq C_v(r)h_n^\alpha\|u\|_Q\|v\|_Q.
 \label{eq:e8v23-v-tail}
\end{equation}

\begin{theorem}[Summable adjacent complex metric-form defect]
\label{thm:e8v23-complex-defect}
Uniformly for \(|z|\leq r\),
\begin{equation}
 \boxed{
 |(a_{n+1,z,\lambda}-a_{n,z,\lambda})[u,v]|
 \leq C_{\rm met}(r)h_n^\alpha
 \|u\|_Q\|v\|_Q,}
 \label{eq:e8v23-complex-defect}
\end{equation}
where
\begin{equation}
 C_{\rm met}(r)
 ={C_{\rm app}(1+2^{-\alpha})L_{G,\alpha}(r)\over\nu_0}
 +C_v(r).
 \label{eq:e8v23-Cmet}
\end{equation}
Hence the V20 adjacent complex-form quantity may be chosen as
\begin{equation}
 \varepsilon_n(r)=C_{\rm met}(r)h_n^\alpha,
 \qquad
 \sum_{n\geq N}\varepsilon_n(r)
 \leq {C_{\rm met}(r)h_N^\alpha\over1-2^{-\alpha}}.
 \label{eq:e8v23-epsilon-tail}
\end{equation}
\end{theorem}

\begin{proof}
The principal difference is
\[
 \int_M\langle\nabla u,(G_{n+1}(z)-G_n(z))\nabla v\rangle\,d\mu.
\]
Cauchy--Schwarz, \Cref{lem:e8v23-adjacent-coefficient}, and
\eqref{eq:e8v23-Q-lower} bound its absolute value by
\[
 {C_{\rm app}(1+2^{-\alpha})L_{G,\alpha}(r)\over\nu_0}
 h_n^\alpha\|u\|_Q\|v\|_Q.
\]
Add \eqref{eq:e8v23-v-tail}.  The geometric-series estimate is immediate.
\end{proof}

\begin{corollary}[Adjacent parameter derivative]
\label{cor:e8v23-parameter}
Let \(z\) complexify one real unit parameter direction.  Under the hypotheses
of \Cref{thm:e8v23-complex-defect}, the local transfer derivative difference
at \(z=0\) obeys the V20 estimate with
\begin{equation}
 {\varepsilon_n(r)\over r}
 ={C_{\rm met}(r)\over r}h_n^\alpha.
 \label{eq:e8v23-parameter-defect}
\end{equation}
Its cofinal tail is bounded by
\begin{equation}
 {C_{\rm met}(r)h_N^\alpha\over r(1-2^{-\alpha})}
 \label{eq:e8v23-parameter-tail}
\end{equation}
before multiplication by the local energy moment and the V20 semigroup
constant.
\end{corollary}

\begin{proof}
The complex form difference is holomorphic and has the circle bound
\eqref{eq:e8v23-complex-defect}.  Apply the V20 holomorphic adjacent
\(C^1\) theorem and sum \eqref{eq:e8v23-epsilon-tail}.
\end{proof}

\begin{corollary}[Direct real derivative coefficient estimate]
\label{cor:e8v23-real-derivative}
If \(G(x,\cdot)\) is \(C^1\) in a real parameter and
\begin{equation}
 \sup_{x\in K}[D_xG(x)[\xi]]_{C^{0,\alpha}}
 \leq L_{DG,\alpha}
 \label{eq:e8v23-DG}
\end{equation}
for unit directions \(\xi\), then
\begin{equation}
 \|P_{n+1}D_xG[\xi]-P_nD_xG[\xi]\|_\infty
 \leq C_{\rm app}(1+2^{-\alpha})L_{DG,\alpha}h_n^\alpha.
 \label{eq:e8v23-DG-tail}
\end{equation}
This is the direct derivative-form row.  It does not replace the combined
semigroup estimate when two full-order noncommuting marks are split.
\end{corollary}

\begin{proof}
Apply \Cref{lem:e8v23-adjacent-coefficient} to
\(D_xG(x)[\xi]\), uniformly in \(x,\xi\).
\end{proof}

\subsection{From a Hölder metric to its densitized inverse}
\label{sec:e8v23-densitized}

For a positive metric matrix \(g\), the scalar kinetic coefficient is
\begin{equation}
 \mathcal G(g)=\sqrt{\det g}\,g^{-1}.
 \label{eq:e8v23-G-of-g}
\end{equation}

\begin{lemma}[Hölder stability of the densitized inverse]
\label{lem:e8v23-densitized}
Let \(\mathcal K_{\nu,\Lambda}\) be the compact set of real symmetric
matrices with eigenvalues in \([\nu,\Lambda]\).  There is a finite constant
\(C_{\nu,\Lambda,d}\) such that for
\(g,h\in\mathcal K_{\nu,\Lambda}\),
\begin{equation}
 \|\mathcal G(g)-\mathcal G(h)\|
 \leq C_{\nu,\Lambda,d}\|g-h\|.
 \label{eq:e8v23-G-Lipschitz}
\end{equation}
Consequently a (C^{0,\alpha}) uniformly nondegenerate metric has a
(C^{0,\alpha}) densitized inverse, with seminorm multiplied by at most
\(C_{\nu,\Lambda,d}\).
\end{lemma}

\begin{proof}
The map \(g\mapsto\sqrt{\det g}\,g^{-1}\) is smooth on the open positive
cone.  Its derivative is continuous and therefore bounded on the compact
convex hull of \(\mathcal K_{\nu,\Lambda}\), which remains inside that cone.
Integrate the derivative along the line segment from \(g\) to \(h\).  Apply
the resulting Lipschitz estimate pointwise to obtain the Hölder seminorm.
\end{proof}

\begin{corollary}[Dyadic metric tail]
\label{cor:e8v23-metric-tail}
Let \(g(z,\cdot)\) remain in a fixed complex neighbourhood of
\(\mathcal K_{\nu,\Lambda}\) on which determinant square root and inverse use
the declared analytic branch.  If its (C^{0,\alpha}) norm is uniform there,
then \(\mathcal G(g(z))\) satisfies
\eqref{eq:e8v23-G-holder}, and its dyadic adjacent kinetic-form tail is
\(O(h_n^\alpha)\) uniformly on that complex neighbourhood.
\end{corollary}

\begin{proof}
The determinant and inverse are holomorphic away from their singular sets.
On a compact complex neighbourhood of the real nondegenerate chart, their
first derivatives are uniformly bounded.  Repeat the proof of
\Cref{lem:e8v23-densitized} and then apply
\Cref{thm:e8v23-complex-defect}.
\end{proof}

\subsection{Bridge and quadrature debits}
\label{sec:e8v23-bridge}

The preceding forms live on one fixed continuum domain.  A lattice regulator
has different finite spaces.  Let \(J_n\) transport the \(n\)-th space into a
common refinement space, and suppose
\begin{equation}
 |a_{n,z}^{\rm disc}[u,v]
 -a_{n,z}^{\rm cont}[J_nu,J_nv]|
 \leq q_n(r)\|u\|_{Q,n}\|v\|_{Q,n}
 \label{eq:e8v23-quadrature}
\end{equation}
uniformly on \(|z|\leq r\).

\begin{theorem}[Complete adjacent metric defect with bridge debit]
\label{thm:e8v23-bridge}
If the bridge norms are uniformly stable, then the adjacent discrete form
defect is bounded by
\begin{equation}
 \boxed{
 \varepsilon_n^{\rm disc}(r)
 \leq C_J C_{\rm met}(r)h_n^\alpha
 +q_n(r)+q_{n+1}(r)+b_n(r),}
 \label{eq:e8v23-discrete-defect}
\end{equation}
where \(b_n(r)\) is the declared noncommutation defect between the two
transport routes.  It is summable if \(q_n(r)\), \(b_n(r)\), and
\(h_n^\alpha\) are summable.
\end{theorem}

\begin{proof}
Insert and subtract the two transported continuum forms on the common
refinement.  Their difference is controlled by
\Cref{thm:e8v23-complex-defect} and the bridge norm \(C_J\).  The two
discrete-to-continuum replacements cost \(q_n,q_{n+1}\); comparing the
direct and staged routes costs \(b_n\).  Add the four bounds.  Summability is
termwise.
\end{proof}

\begin{countertheorem}[Ellipticity without spatial regularity is not enough]
\label{cth:e8v23-ellipticity-only}
Uniform ellipticity alone does not imply a vanishing, much less summable,
adjacent \(L^\infty\) coefficient defect.
\end{countertheorem}

\begin{proof}
On \([0,1]\), let \(r_n\) be the \(n\)-th Rademacher function, constant
\(\pm1\) on dyadic intervals of length \(2^{-n}\), and set
\(G_n=1+\varepsilon r_n\) with \(0<\varepsilon<1\).  Then
\(1-\varepsilon\leq G_n\leq1+\varepsilon\) for every \(n\), but
\begin{equation}
 \|G_{n+1}-G_n\|_\infty=2\varepsilon.
 \label{eq:e8v23-rademacher}
\end{equation}
The adjacent defects do not even tend to zero.
\end{proof}

\begin{assessment}[Progress at V23]
\label{ass:e8v23-status}
For a fixed uniformly nondegenerate (C^{0,\alpha}) metric chart, the
principal complex form and its parameter derivative now have an explicit
geometrically summable adjacent tail.  The result remains conditional on a
stable common-refinement bridge, summable quadrature/contact debits, and the
DQSRC square-root row.  It applies to matter kinetic forms on the metric
chart and does not construct the gravitational carrier.
\end{assessment}

\begin{programme}[Eighth-series V24: graph bridge and quadrature]
\label{prog:e8v23-V24}
V24 will populate the bridge debit in
\eqref{eq:e8v23-discrete-defect} for a shape-regular nested finite-element or
finite-volume scalar kinetic form.  It will prove the exact direct/staged
commutation on nested trial spaces and isolate the quadrature order needed
for summability.  Gauge holonomy and spin transport will remain separately
typed contacts unless the same proof covers them.
\end{programme}

\begin{firewall}[Claim boundary after eighth-series V23]
\label{fire:e8v23-boundary}
V23 proves geometric adjacent coefficient and form tails for uniformly
Hölder metric families on a complex disk, transports the result through the
densitized inverse, and gives the complete bridge/quadrature defect formula.

V23 does not prove DQSRC, the lattice bridge or quadrature rates, gauge and
spin contacts, determinant or gravitational positivity, removal of the
cutoff, or construction of the Standard Model--Einstein continuum.
\end{firewall}

\begin{theorem}[Eighth-series V23 result]
\label{thm:e8v23-result}
Dyadic approximation of a uniformly elliptic
(C^{0,\alpha}) complex metric coefficient produces the summable adjacent
form defect \eqref{eq:e8v23-epsilon-tail} and parameter derivative tail
\eqref{eq:e8v23-parameter-tail}.  After a stable discrete bridge, the only
additional debits are exactly those in \eqref{eq:e8v23-discrete-defect}.
\end{theorem}

\begin{proof}
Combine \Cref{lem:e8v23-adjacent-coefficient,
thm:e8v23-complex-defect,cor:e8v23-parameter,
lem:e8v23-densitized,cor:e8v23-metric-tail,
thm:e8v23-bridge}.
\end{proof}

\paragraph{Parent-version record.}
V23 was formed from
\texttt{Renewal\_SM\_Einstein\_Continuum\_Research\_Notes\_V22.tex}.
The V22 parent had (8\,172) logical source lines, (304\,337) bytes,
and SHA--256
\[
 \texttt{D71AF0680EB006E82EFFDFD364A7C596998E5724B24D70676FB3D26BBF545E6D}.
\]
The next compulsory companion audit is V25.

% =============================================================================
% END APPENDED EIGHTH-SERIES V23 MATERIAL
% =============================================================================



% =============================================================================
% BEGIN APPENDED EIGHTH-SERIES V24 MATERIAL
% =============================================================================

\section{Nested Galerkin bridges and Hölder kinetic quadrature}
\label{sec:v8v24}

V23 expressed the complete adjacent metric defect as a coefficient tail plus
quadrature and route-commutation debits.  This appendix closes those latter
debits for a scalar or vector-valued kinetic form on nested conforming
piecewise-affine spaces.  Exact Galerkin restriction has zero bridge defect.
Cell sampling of a \(C^{0,\alpha}\) coefficient has the same summable
\(h_n^\alpha\) error as the coefficient approximation, while cell averaging
is exact for the kinetic integrand on each affine element.

\subsection{Nested trial spaces}
\label{sec:e8v24-nested}

Let \(\{\mathcal T_n\}_{n\geq0}\) be nested, shape-regular simplicial meshes
of a bounded Lipschitz domain \(M\), with maximal element diameter
\(h_n=h_0 2^{-n}\).  Let \(V_n\subset H^1(M;\mathbb C^N)\) be the conforming
continuous \(P_1\) space with a fixed boundary condition.  Refinement gives
literal inclusions
\begin{equation}
 I_n:V_n\hookrightarrow V_{n+1}.
 \label{eq:e8v24-inclusion}
\end{equation}
For a bounded coefficient \(G(z,x)\), define the exact kinetic form
\begin{equation}
 k_z[u,v]=\int_M\langle\nabla u,G(z,x)\nabla v\rangle\,dx.
 \label{eq:e8v24-k}
\end{equation}
Its Galerkin restriction is \(k_{n,z}=k_z|_{V_n\times V_n}\).

\begin{theorem}[Exact direct/staged Galerkin commutation]
\label{thm:e8v24-exact-bridge}
For \(u,v\in V_n\),
\begin{equation}
 \boxed{
 k_{n+1,z}[I_nu,I_nv]=k_{n,z}[u,v].}
 \label{eq:e8v24-exact-bridge}
\end{equation}
Moreover, the direct inclusion \(V_n\hookrightarrow V_{n+2}\) equals
\(I_{n+1}I_n\) as a function, so every direct/staged kinetic-form bridge
defect is exactly zero.
\end{theorem}

\begin{proof}
The inclusion does not interpolate or alter a \(P_1\) function; it regards
the same function as an element of the refined space.  Both sides of
\eqref{eq:e8v24-exact-bridge} are therefore the same integral
\eqref{eq:e8v24-k}.  Literal inclusions compose, proving the route statement.
\end{proof}

\begin{remark}[Ritz projection is a different map]
\label{rem:e8v24-Ritz}
A projection from a fine space to a coarse space need not compose exactly and
can create a bridge defect.  The theorem concerns the cofinal inductive
inclusions.  Recovery maps and coarse marginalization remain separate FLAC
rows.
\end{remark}

\subsection{Point-sampled and cell-averaged coefficients}
\label{sec:e8v24-quadrature}

Assume \(G(z,\cdot)\in C^{0,\alpha}\) uniformly for \(|z|\leq r\), with
\begin{equation}
 \sup_{|z|\leq r}[G(z)]_{C^{0,\alpha}}\leq L_\alpha(r).
 \label{eq:e8v24-holder}
\end{equation}
For each element \(K\in\mathcal T_n\), choose \(x_K\in K\) and define
\begin{equation}
 G_n^{\rm pt}(z,x)=G(z,x_K),
 \qquad x\in K.
 \label{eq:e8v24-point}
\end{equation}
Let \(k_{n,z}^{\rm pt}\) be \eqref{eq:e8v24-k} with this piecewise constant
coefficient.  Also define the cell average
\begin{equation}
 G_{n,K}^{\rm av}(z)={1\over|K|}\int_KG(z,x)\,dx.
 \label{eq:e8v24-average}
\end{equation}

\begin{lemma}[Point-sampling coefficient error]
\label{lem:e8v24-point-error}
Uniformly on the complex disk,
\begin{equation}
 \|G_n^{\rm pt}(z)-G(z)\|_{L^\infty}
 \leq L_\alpha(r)h_n^\alpha.
 \label{eq:e8v24-point-error}
\end{equation}
\end{lemma}

\begin{proof}
For \(x\in K\), Hölder continuity gives
\(\|G(z,x_K)-G(z,x)\|\leq L_\alpha(r)|x-x_K|^\alpha\), and
\(|x-x_K|\leq h_n\).
\end{proof}

\begin{theorem}[Summable kinetic quadrature error]
\label{thm:e8v24-quadrature}
Let the common reference form norm satisfy
\begin{equation}
 \|u\|_Q^2\geq\nu_0\|\nabla u\|_2^2.
 \label{eq:e8v24-Q}
\end{equation}
Then for \(u,v\in V_n\), uniformly on \(|z|\leq r\),
\begin{equation}
 \boxed{
 |(k_{n,z}^{\rm pt}-k_{n,z})[u,v]|
 \leq {L_\alpha(r)\over\nu_0}h_n^\alpha
 \|u\|_Q\|v\|_Q.}
 \label{eq:e8v24-quadrature-error}
\end{equation}
The quadrature debit is geometrically summable:
\begin{equation}
 \sum_{n\geq N}q_n(r)
 \leq {L_\alpha(r)h_N^\alpha
 \over\nu_0(1-2^{-\alpha})}.
 \label{eq:e8v24-quadrature-tail}
\end{equation}
For the cell-average coefficient and \(P_1\) functions,
\begin{equation}
 k_{n,z}^{\rm av}[u,v]=k_{n,z}[u,v]
 \label{eq:e8v24-average-exact}
\end{equation}
exactly on every element.
\end{theorem}

\begin{proof}
The difference between the point-sampled and exact forms is bounded by
\[
 \|G_n^{\rm pt}-G\|_\infty
 \|\nabla u\|_2\|\nabla v\|_2.
\]
Apply \Cref{lem:e8v24-point-error} and \eqref{eq:e8v24-Q}; summing the
geometric sequence proves \eqref{eq:e8v24-quadrature-tail}.  On an affine
simplex, \(\nabla u\) and \(\nabla v\) are constant.  Hence
\begin{align*}
 \int_K\langle\nabla u,G(z,x)\nabla v\rangle\,dx
 &=|K|\langle\nabla u,G_{n,K}^{\rm av}(z)\nabla v\rangle,
\end{align*}
which proves \eqref{eq:e8v24-average-exact} after summing the cells.
\end{proof}

\begin{corollary}[Parameter derivative quadrature]
\label{cor:e8v24-derivative}
If the complex parameter derivative has the uniform spatial bound
\begin{equation}
 \sup_{|z|\leq r}[\partial_zG(z)]_{C^{0,\alpha}}
 \leq L_{\alpha}^{(1)}(r),
 \label{eq:e8v24-DG}
\end{equation}
then the direct derivative-form quadrature error satisfies
\begin{equation}
 |\partial_z(k_{n,z}^{\rm pt}-k_{n,z})[u,v]|
 \leq {L_{\alpha}^{(1)}(r)\over\nu_0}h_n^\alpha
 \|u\|_Q\|v\|_Q.
 \label{eq:e8v24-Dquadrature}
\end{equation}
The same derivative is recovered from the V20 circle estimate when the full
complex disk is retained.
\end{corollary}

\begin{proof}
Differentiate the coefficient at fixed trial functions and repeat
\Cref{lem:e8v24-point-error,thm:e8v24-quadrature}.  Holomorphy justifies
termwise differentiation and Cauchy's formula.
\end{proof}

\subsection{Coefficient refinement on the nested bridge}
\label{sec:e8v24-coefficient-bridge}

Let \(G_n(z)\) be any piecewise coefficient on \(\mathcal T_n\), and define
\begin{equation}
 k_{n,z}^{G_n}[u,v]
 =\int_M\langle\nabla u,G_n(z)\nabla v\rangle\,dx.
 \label{eq:e8v24-Gn-form}
\end{equation}

\begin{theorem}[Exact location of the adjacent defect]
\label{thm:e8v24-location}
For \(u,v\in V_n\),
\begin{equation}
 \boxed{
 k_{n+1,z}^{G_{n+1}}[I_nu,I_nv]-k_{n,z}^{G_n}[u,v]
 =\int_M\langle\nabla u,(G_{n+1}(z)-G_n(z))\nabla v\rangle\,dx.}
 \label{eq:e8v24-location}
\end{equation}
There is no trial-space or route term.  If
\(\|G_{n+1}-G_n\|_\infty\leq C_G(r)h_n^\alpha\), the form defect is at most
\begin{equation}
 {C_G(r)\over\nu_0}h_n^\alpha\|u\|_Q\|v\|_Q.
 \label{eq:e8v24-location-bound}
\end{equation}
\end{theorem}

\begin{proof}
Use the literal inclusion from \Cref{thm:e8v24-exact-bridge} and subtract the
two integrals.  The norm bound follows exactly as in
\Cref{thm:e8v24-quadrature}.
\end{proof}

\begin{corollary}[Population of the V23 bridge formula]
\label{cor:e8v24-V23}
For nested conforming \(P_1\) spaces with exact integration, the V23 bridge
and quadrature debits satisfy
\begin{equation}
 b_n(r)=q_n(r)=0.
 \label{eq:e8v24-zero-debits}
\end{equation}
For point sampling, \(b_n(r)=0\) and
\begin{equation}
 q_n(r)\leq {L_\alpha(r)\over\nu_0}h_n^\alpha.
 \label{eq:e8v24-point-debit}
\end{equation}
Thus the complete adjacent kinetic-form defect is geometrically summable.
\end{corollary}

\begin{proof}
Exact restriction is \Cref{thm:e8v24-exact-bridge}; exact and point-sampled
quadrature are \Cref{thm:e8v24-quadrature}.  Insert their values into the V23
bridge theorem.
\end{proof}

\subsection{Lower-order and geometric contacts}
\label{sec:e8v24-contacts}

\begin{remark}[Mass and interaction quadrature]
\label{rem:e8v24-lower-order}
The cell-average exactness in \eqref{eq:e8v24-average-exact} uses constancy of
\(P_1\) gradients.  Products \(u\overline v\), Yukawa terms, and nonlinear
potentials are not constant on a cell.  Their quadrature errors require their
own estimate
\begin{equation}
 |q_{n,z}^{\rm low}[u,v]|
 \leq C_{\rm low}(r)h_n^\beta\|u\|_Q\|v\|_Q,
 \qquad \beta>0,
 \label{eq:e8v24-low-order}
\end{equation}
and are summable only after such a positive order is proved.  Positivity of a
mass-lumped rule does not by itself supply \eqref{eq:e8v24-low-order}.
\end{remark}

Let \(\mathcal A\) be a fixed smooth gauge or spin connection and let
\(U_\mathcal A(e)\) be its path-ordered parallel transport along an oriented
edge \(e\).

\begin{lemma}[Exact holonomy subdivision]
\label{lem:e8v24-holonomy}
If a coarse edge is the oriented concatenation
\(e=e_m\circ\cdots\circ e_1\), then
\begin{equation}
 \boxed{
 U_\mathcal A(e)=U_\mathcal A(e_m)\cdots U_\mathcal A(e_1).}
 \label{eq:e8v24-holonomy}
\end{equation}
Thus direct and staged subdivision of a continuum connection has zero
holonomy-route defect.
\end{lemma}

\begin{proof}
Parallel transport is the fundamental solution of a first-order linear
ordinary differential equation along the path.  Uniqueness of that solution
makes the propagator multiplicative under concatenation.  Repeated
application proves \eqref{eq:e8v24-holonomy}.
\end{proof}

\begin{firewall}[Independent lattice links are different data]
\label{fire:e8v24-links}
The identity \eqref{eq:e8v24-holonomy} applies when every link is sampled
from one continuum connection on the same path.  Independently distributed
fine links or a nonlinear blocking rule need not multiply to the coarse
link.  Their discrepancy is a physical gauge-bridge contact and cannot be
set to zero by the finite-element inclusion theorem.  The same qualification
applies to a spin lift and its choice of frame orientation.
\end{firewall}

\begin{theorem}[Kinetic bridge acquisition theorem]
\label{thm:e8v24-acquisition}
For a uniformly Hölder complex metric coefficient on nested conforming
\(P_1\) spaces, exact Galerkin integration populates the HEFC kinetic bridge
with zero route and quadrature debits.  Point sampling populates it with the
summable debit \eqref{eq:e8v24-point-debit}.  If gauge and spin transports
come from one continuum connection and subdivide the same paths, their
holonomy route debit is also zero.  All lower-order, recovery, and independent
link-blocking contacts remain separate.
\end{theorem}

\begin{proof}
Combine \Cref{thm:e8v24-exact-bridge,
thm:e8v24-quadrature,thm:e8v24-location,
cor:e8v24-V23,lem:e8v24-holonomy}.
\end{proof}

\begin{assessment}[Progress at V24]
\label{ass:e8v24-status}
The V23 bridge formula is now populated for one concrete discretization of
the bosonic kinetic form.  The trial-space bridge contributes no error, and
the kinetic quadrature error is either zero or \(O(h_n^\alpha)\).  This does
not close DQSRC, nonlinear interaction quadrature, recovery/projectivity, or
the independently generated gauge and spin link routes.
\end{assessment}

\begin{programme}[Eighth-series V25: companion audit and HEFC ledger]
\label{prog:e8v24-V25}
V25 will perform the compulsory companion audit, then consolidate V20--V24
into a populated/unpopulated HEFC table.  The next mathematical attack will
be selected from DQSRC, nonlinear lower-order quadrature, or source-local
energy according to which has an upstream estimate in the attached
SM--Einstein construction.
\end{programme}

\begin{firewall}[Claim boundary after eighth-series V24]
\label{fire:e8v24-boundary}
V24 proves exact nested Galerkin bridge commutation, summable Hölder
point-sampling error, exact cell-average kinetic integration, the location of
the coefficient defect, and exact continuum-holonomy subdivision.

V24 does not prove the nonnormal square-root row, lower-order interaction
quadrature, projective recovery, independent lattice-link consistency,
determinant or gravitational positivity, cutoff removal, or construction of
the Standard Model--Einstein continuum.
\end{firewall}

\begin{theorem}[Eighth-series V24 result]
\label{thm:e8v24-result}
For nested conforming \(P_1\) spaces, the kinetic trial-space bridge is exact.
A uniformly \(C^{0,\alpha}\) complex coefficient has either zero cell-average
kinetic quadrature error or the summable point-sampling error
\eqref{eq:e8v24-quadrature-tail}.  Consequently the full kinetic adjacent
form row is populated, conditional only on the separately declared
square-root and contact rows.
\end{theorem}

\begin{proof}
This is \Cref{thm:e8v24-acquisition} with the V23 coefficient tail.
\end{proof}

\paragraph{Parent-version record.}
V24 was formed from
\texttt{Renewal\_SM\_Einstein\_Continuum\_Research\_Notes\_V23.tex}.
The V23 parent had (8\,528) logical source lines, (317\,031) bytes,
and SHA--256
\[
 \texttt{DAE0E04B77DB6EB151D02BA5083A236A71A88FCA405F7E6FAFA543083A7550CF}.
\]
The next compulsory companion audit is V25.

% =============================================================================
% END APPENDED EIGHTH-SERIES V24 MATERIAL
% =============================================================================



% =============================================================================
% BEGIN APPENDED EIGHTH-SERIES V25 MATERIAL
% =============================================================================

\section{Companion audit and the populated HEFC ledger}
\label{sec:v8v25}

This appendix performs the compulsory companion audit, records exactly which
holomorphic elliptic form rows were populated in V20--V24, and adjusts the
next direction.  The newest Yang--Mills obstruction shows why a theorem
uniform over every exterior configuration may be too strong even when an
averaged physical theorem remains possible.  The analogous SM--Einstein
question is whether the metric carrier can be restricted to a uniformly
nondegenerate good chart with a controlled bad-geometry capacity debit.

\subsection{Compulsory V25 companion audit}
\label{sec:e8v25-audit}

The newest completed companion files were
\begin{align*}
 &\texttt{Renewal\_Yang\_Mills\_Mass\_Gap\_Research\_Notes\_V83.tex},\\
 &\texttt{Renewal\_NS\_Stability\_Research\_Notes\_V26.tex}.
\end{align*}
At the audit read, the Yang--Mills file had (37\,684) logical source lines,
(1\,320\,999) bytes, one live document terminator, and SHA--256
\[
 \texttt{663720A4F5CEF4C66658A35F562D7024FE9DAAD53B7C4E910AE4FCCD46E6A607}.
\]
The Navier--Stokes file had (5\,319) logical source lines,
(278\,283) bytes, one live terminator, and SHA--256
\[
 \texttt{82C887F8C2C69E4F28FAD7C3DC8DE5B9099CB5A4BF431EBA1FFF3E91935978AD}.
\]

\begin{assessment}[Yang--Mills V76--V83 reading]
\label{ass:e8v25-YM}
The intervening Yang--Mills work moves from growing representation control to
an intrinsic ground-state Poincar\'e route.  V82 gives a quotient-compatible
block approximate-tensorization compiler and distinguishes its auxiliary
heat-bath clock from the physical half-slab.  V83 then constructs an
admissible center-staple exterior whose one-link conditional law has two
separated gauge-visible wells and an exponentially collapsing Poincar\'e gap.
This disproves the worst-exterior block-gap row.  It does not disprove an
averaged physical inequality; V83 replaces the failed row by a good-exterior
estimate plus a bad-exterior variance/capacity debit.
\end{assessment}

\begin{assessment}[Navier--Stokes V26 reading]
\label{ass:e8v25-NS}
The NS file is unchanged and contains no new continuum form, metric bridge,
or bad-geometry capacity input relevant to HEFC.
\end{assessment}

\begin{theorem}[V25 audit decision]
\label{thm:e8v25-audit}
No Yang--Mills constant or spectral conclusion is imported.  The transferable
proof strategy is to distinguish:
\begin{equation}
 \boxed{
 \text{uniform control on a good physical chart}
 \quad+\quad
 \text{a quantified bad-set form debit}}
 \label{eq:e8v25-good-bad}
\end{equation}
from a supremum over every admissible exterior or metric history.  V26 will
formulate and prove that averaged compiler for the SM local transfer rows.
\end{theorem}

\begin{proof}
The V83 counterexample is specific to an \(\mathrm{SU}(3)\) matrix-Fisher
conditional law, so its constants and gap statement do not transfer.  Its
logic is general: small probability alone does not control an unbounded
\(L^2\) or form-domain contribution, whereas a relative variance or capacity
bound does.  The HEFC currently assumes uniform nondegeneracy on a chosen
metric chart, so the same logical distinction identifies a missing carrier
row without asserting that bad geometries are rare.
\end{proof}

\subsection{HEFC population ledger after V24}
\label{sec:e8v25-ledger}

Recall the six HEFC rows of V21.  The following status refers only to the
bosonic kinetic sector on a prescribed uniformly nondegenerate metric chart.

\begin{center}
\begin{tabular}{p{0.09\textwidth}p{0.27\textwidth}p{0.51\textwidth}}
\toprule
row & status & exact content \\
\midrule
(H1) & partial & V18 separates physical unitary transport from form
congruence; V24 gives one fixed continuum \(H^1\) domain and nested \(P_1\)
trial spaces.  The full SM gauge, fermion, and gravity domain is open. \\
(H2) & populated on the good chart & V21 gives the explicit complex radius
\(r_\theta^{\rm ell}=\min\{R,\theta/\kappa_{\rm ell}\}\) from ellipticity and
relative-potential margins. \\
(H3) & partial & V22 proves uniform square-root constants on the real and
energy-commuting complex branches.  Generic nonnormal complex metric
coefficients require DQSRC. \\
(H4) & populated for the kinetic coefficient & V23 gives the local complex
component sum for a uniformly \(C^{0,\alpha}\) densitized inverse metric. \\
(H5) & populated for nested \(P_1\) kinetics & V23 gives an
\(O(h_n^\alpha)\) adjacent complex form defect; V24 gives zero exact-Galerkin
bridge error and zero or \(O(h_n^\alpha)\) kinetic quadrature error. \\
(H6) & partial & V24 gives exact path-subdivision holonomy for links sampled
from one continuum connection.  Moving measures, nonlinear interactions,
independent link blocking, spin lifts, recovery, determinant/Pfaffian, and
physical clock incidence remain open. \\
\bottomrule
\end{tabular}
\end{center}

\begin{theorem}[Restricted kinetic-sector handoff]
\label{thm:e8v25-kinetic-handoff}
On a uniformly nondegenerate \(C^{0,\alpha}\) metric chart, suppose DQSRC,
the V16 source-local energy row, and all remaining contact tails hold with
summable constants.  Then the nested bosonic \(P_1\) kinetic transfer matrix
elements and their metric derivatives converge in \(C^1\) on compact
parameter charts, with a geometric generator tail \(O(h_N^\alpha)\).
\end{theorem}

\begin{proof}
V21 supplies a fixed complex sectorial radius.  DQSRC and the V21
square-root compiler supply the uniform form-semigroup constant.  V23 and V24
supply an adjacent complex form tail bounded by \(Ch_n^\alpha\), whose sum
from \(N\) is \(O(h_N^\alpha)\).  Multiply by the V16 local energy moment,
add the assumed contact tails, and apply the V20 holomorphic adjacent theorem
and the V15 \(C^1\) landing theorem.
\end{proof}

\begin{firewall}[Scope of the populated rows]
\label{fire:e8v25-scope}
The phrase ``populated'' in the table means proved for the stated nested
bosonic kinetic discretization.  It does not mean proved for the interacting
Standard Model, for a fluctuating gravitational measure, for a chiral
fermion line, or for the physical reflected Hamiltonian.  DQSRC and the
source-local energy moment are still hypotheses even in that restricted
handoff.
\end{firewall}

\subsection{Why worst-history uniformity is a separate gate}
\label{sec:e8v25-worst}

Let \(\mathfrak G_n\) be the regulated metric-history space and let
\(\mu_n\) be its proposed physical carrier.  For fixed
\(0<\nu<\Lambda<\infty\), define the good chart
\begin{equation}
 \mathfrak G_n^{\rm good}(\nu,\Lambda)
 =\{g:\nu I\preceq g\preceq\Lambda I
 \text{ on the declared region, with the HEFC regularity bounds}\}.
 \label{eq:e8v25-good-chart}
\end{equation}
Its complement is the bad-geometry event
\(\mathfrak G_n^{\rm bad}\).

\begin{proposition}[Small bad probability is not a form bound]
\label{prop:e8v25-small-probability}
For every \(0<p<1\) and every \(M>0\), there are a probability space, an
event \(B\) of probability \(p\), and a nonnegative random variable \(X\)
supported on \(B\) such that \(\mathbb E X=M\).  Therefore a statement
\(\mu_n(\mathfrak G_n^{\rm bad})\to0\) alone cannot control an unbounded
energy, square-root, or stress contribution on the bad set.
\end{proposition}

\begin{proof}
Take \(X=(M/p)1_B\).  Then \(X\geq0\), it is supported on \(B\), and
\(\mathbb EX=M\), independently of how small \(p\) is.
\end{proof}

\begin{definition}[Bad-geometry form debit, BGFD]
\label{def:e8v25-BGFD}
A BGFD for a local carrier vector \(\Psi_{n,A}\) is a bound of the form
\begin{equation}
 \int_{\mathfrak G_n^{\rm bad}}
 \|\Psi_{n,A}(g)\|_{Q_n(g)}^2\,d\mu_n(g)
 \leq\varepsilon_{n,A}^{\rm bad}
 +\kappa_n^{\rm bad}
 \int_{\mathfrak G_n}
 \|\Psi_{n,A}(g)\|_{Q_n(g)}^2\,d\mu_n(g),
 \label{eq:e8v25-BGFD}
\end{equation}
with \(\sup_n\kappa_n^{\rm bad}<1\) when the term must be absorbed, and with
the remaining \(\varepsilon_{n,A}^{\rm bad}\) summable in an adjacent-scale
application.  An analogous row is required for the complexified adjacent
form and contact integrands.
\end{definition}

\begin{assessment}[Direction selected by the audit]
\label{ass:e8v25-direction}
Demanding HEFC and DQSRC uniformly over every metric history may fail for the
same structural reason that the Yang--Mills worst-exterior conditional gap
fails: admissible histories can approach degeneracy or develop separated
geometric sectors.  The correct next question is not whether their
probability is small.  It is whether the bad-history contribution satisfies a
relative form/capacity inequality such as \eqref{eq:e8v25-BGFD}, allowing it
to be absorbed into a good-chart estimate.
\end{assessment}

\begin{programme}[Eighth-series V26: good-geometry averaged compiler]
\label{prog:e8v25-V26}
V26 will prove an abstract good/bad compiler for local transfer and stress
matrix elements.  It will identify the absorption threshold, show how a
capacity or weighted-energy bound implies BGFD, and state the exact bad-set
row required to extend the restricted kinetic handoff beyond a worst-history
supremum.  It will not assume that the proposed gravitational carrier already
satisfies that row.
\end{programme}

\begin{firewall}[Claim boundary after eighth-series V25]
\label{fire:e8v25-boundary}
V25 performs the compulsory companion audit, records the populated HEFC
ledger, proves the restricted conditional kinetic handoff, and isolates the
bad-geometry form debit suggested by the Yang--Mills worst-exterior
obstruction.

V25 does not prove BGFD for the gravitational carrier, DQSRC, the complete
SM contact rows, signed coercivity, reflection-positive gravity, cutoff
removal, or construction of the Standard Model--Einstein continuum.
\end{firewall}

\begin{theorem}[Eighth-series V25 result]
\label{thm:e8v25-result}
The bosonic nested kinetic rows (H2), (H4), and (H5) of HEFC are populated on
a uniformly nondegenerate Hölder metric chart; (H1), (H3), and (H6) remain
partial.  Conditional on DQSRC, local energy, and contact tails, these rows
give \(C^1\) kinetic-sector landing.  Extension from a good metric chart to
the physical carrier requires a bad-geometry form debit, not merely a small
bad-set probability.
\end{theorem}

\begin{proof}
The table is the direct ledger of V20--V24.  The conditional landing is
\Cref{thm:e8v25-kinetic-handoff}; the necessity of a weighted bad-set row is
\Cref{prop:e8v25-small-probability}.
\end{proof}

\paragraph{Parent-version record.}
V25 was formed from
\texttt{Renewal\_SM\_Einstein\_Continuum\_Research\_Notes\_V24.tex}.
The V24 parent had (8\,883) logical source lines, (329\,948) bytes,
and SHA--256
\[
 \texttt{E65598CE2CB2653F6ED9C8696A7DDDA2F03D0169F47989AAF442100A9D80BF3D}.
\]
The next compulsory companion audit is V30.

% =============================================================================
% END APPENDED EIGHTH-SERIES V25 MATERIAL
% =============================================================================



% =============================================================================
% BEGIN APPENDED EIGHTH-SERIES V26 MATERIAL
% =============================================================================

\section{A good-geometry compiler with bad-set form absorption}
\label{sec:v8v26}

The V25 audit replaced a worst-history target by a good-geometry chart plus a
quantified bad-set debit.  This appendix proves the abstract absorption
principle for local energy moments and adjacent transfer defects.  It then
derives a bad-set form bound from a uniform Sobolev inequality.  The result is
conditional but operational: it states the exact summable rows needed to
extend the good-chart HEFC landing theorem to an averaged physical carrier.

\subsection{Energy absorption}
\label{sec:e8v26-energy}

Let \((\Omega_n,\mu_n)\) be a probability space of regulated histories,
split into a measurable good event \(G_n\) and bad event \(B_n=G_n^c\).  Let
\(e_{n,A}\geq0\) be the historywise form-energy density of a fixed local
source \(A\), and put
\begin{align}
 E_{n,A}&=\int_{\Omega_n}e_{n,A}\,d\mu_n,
 \label{eq:e8v26-E}\\
 E_{n,A}^{G}&=\int_{G_n}e_{n,A}\,d\mu_n,
 &E_{n,A}^{B}&=\int_{B_n}e_{n,A}\,d\mu_n.
 \label{eq:e8v26-EGB}
\end{align}

\begin{theorem}[Good/bad energy absorption]
\label{thm:e8v26-energy-absorption}
Suppose
\begin{align}
 E_{n,A}^{G}&\leq C_A,
 \label{eq:e8v26-good-energy}\\
 E_{n,A}^{B}&\leq\varepsilon_{n,A}^{E}
 +\kappa_n^E E_{n,A},
 \qquad 0\leq\kappa_n^E<1.
 \label{eq:e8v26-bad-energy}
\end{align}
Then
\begin{equation}
 \boxed{
 E_{n,A}\leq{C_A+\varepsilon_{n,A}^{E}\over1-\kappa_n^E}.}
 \label{eq:e8v26-energy-absorbed}
\end{equation}
In particular, if \(\sup_n\kappa_n^E<1\) and
\(\sup_n\varepsilon_{n,A}^{E}<\infty\), the global local-energy moment is
uniformly bounded.
\end{theorem}

\begin{proof}
Since \(G_n\) and \(B_n\) partition the history space,
\[
 E_{n,A}=E_{n,A}^{G}+E_{n,A}^{B}
 \leq C_A+\varepsilon_{n,A}^{E}+\kappa_n^E E_{n,A}.
\]
Move the last term to the left and divide by \(1-\kappa_n^E\).
\end{proof}

\begin{remark}[Why the relative term is useful]
\label{rem:e8v26-relative}
The bad energy need not be small absolutely.  It may be controlled by a
strict fraction of the total energy plus a finite remainder.  The strict
inequality \(\kappa_n^E<1\) is the exact absorption threshold.
\end{remark}

\subsection{Adjacent-defect absorption}
\label{sec:e8v26-defect}

After transporting adjacent regulators and their measures to one comparison
space, let \(\Delta_{n,A}^{(k)}(g;x,t)\) be the scalar adjacent defect for
\(k=0\) or its parameter derivative for \(k=1\).  Define
\begin{equation}
 D_{n,A}^{(k)}
 :=\sup_{(x,t)\in K\times[0,T]}
 \int_{\Omega_n}|\Delta_{n,A}^{(k)}(g;x,t)|\,d\mu_n(g).
 \label{eq:e8v26-D}
\end{equation}
Let \(D_{n,A}^{(k),G}\) and \(D_{n,A}^{(k),B}\) be the corresponding
integrals over \(G_n\) and \(B_n\), with the supremum retained.

\begin{theorem}[Good/bad adjacent-defect compiler]
\label{thm:e8v26-defect-compiler}
Suppose for \(k=0,1\),
\begin{align}
 D_{n,A}^{(k),G}&\leq d_{n,A}^{(k),G},
 \label{eq:e8v26-good-defect}\\
 D_{n,A}^{(k),B}&\leq\varepsilon_{n,A}^{(k),B}
 +\kappa_n^{(k)}D_{n,A}^{(k)},
 \qquad0\leq\kappa_n^{(k)}<1.
 \label{eq:e8v26-bad-defect}
\end{align}
Then
\begin{equation}
 \boxed{
 D_{n,A}^{(k)}
 \leq{d_{n,A}^{(k),G}+\varepsilon_{n,A}^{(k),B}
 \over1-\kappa_n^{(k)}}.}
 \label{eq:e8v26-defect-absorbed}
\end{equation}
If the right sides are summable in \(n\), the averaged local transfer
functions converge in \(C^1\) by the V15 landing theorem.
\end{theorem}

\begin{proof}
For every \(x,t\), split the absolute integral over \(G_n\cup B_n\), then
take the supremum.  Equations \eqref{eq:e8v26-good-defect}--
\eqref{eq:e8v26-bad-defect} give
\[
 D_{n,A}^{(k)}
 \leq d_{n,A}^{(k),G}+\varepsilon_{n,A}^{(k),B}
 +\kappa_n^{(k)}D_{n,A}^{(k)}.
\]
Absorb the last term.  Summability gives exactly the two telescoping series
required in V15.
\end{proof}

\begin{corollary}[Geometric good tail plus summable bad debit]
\label{cor:e8v26-geometric}
If
\begin{align}
 d_{n,A}^{(k),G}&\leq C_A^{(k)}2^{-\alpha n},
 \qquad \alpha>0,
 \label{eq:e8v26-good-geometric}\\
 \sum_n\varepsilon_{n,A}^{(k),B}&<\infty,
 &\sup_n\kappa_n^{(k)}&\leq\kappa_*^{(k)}<1,
 \label{eq:e8v26-bad-summable}
\end{align}
then
\begin{equation}
 \sum_{n\geq N}D_{n,A}^{(k)}
 \leq{1\over1-\kappa_*^{(k)}}
 \left{
 {C_A^{(k)}2^{-\alpha N}\over1-2^{-\alpha}}
 +\sum_{n\geq N}\varepsilon_{n,A}^{(k),B}
 \right}.
 \label{eq:e8v26-total-tail}
\end{equation}
\end{corollary}

\begin{proof}
Insert \eqref{eq:e8v26-good-geometric}--
\eqref{eq:e8v26-bad-summable} in
\eqref{eq:e8v26-defect-absorbed} and sum.
\end{proof}

\subsection{A Sobolev route to bad-set form smallness}
\label{sec:e8v26-Sobolev}

Let \((\Omega,\mu)\) be a probability space with a nonnegative closed form
\(\mathcal E\) and form norm
\begin{equation}
 \|f\|_{\mathcal Q}^2=\|f\|_2^2+\mathcal E[f].
 \label{eq:e8v26-Q}
\end{equation}
Assume that for some \(q>1\),
\begin{equation}
 \|f\|_{2q}^2\leq S_q\|f\|_{\mathcal Q}^2.
 \label{eq:e8v26-Sobolev}
\end{equation}

\begin{lemma}[Small set gives a form-small multiplier]
\label{lem:e8v26-small-set}
For every measurable \(B\subset\Omega\),
\begin{equation}
 \boxed{
 \int_B|f|^2\,d\mu
 \leq S_q\mu(B)^{1-1/q}\|f\|_{\mathcal Q}^2.}
 \label{eq:e8v26-small-set}
\end{equation}
Thus multiplication by \(1_B\) has relative form coefficient
\begin{equation}
 \kappa_B=S_q\mu(B)^{1-1/q}.
 \label{eq:e8v26-kappaB}
\end{equation}
\end{lemma}

\begin{proof}
H\"older's inequality with exponents \(q\) and \(q/(q-1)\) gives
\[
 \int_B|f|^2\,d\mu
 \leq\mu(B)^{1-1/q}\left(\int|f|^{2q}\,d\mu\right)^{1/q}.
\]
The last factor is \(\|f\|_{2q}^2\); apply
\eqref{eq:e8v26-Sobolev}.
\end{proof}

\begin{corollary}[Strict absorption threshold]
\label{cor:e8v26-threshold}
If
\begin{equation}
 S_q\mu(B)^{1-1/q}<1,
 \label{eq:e8v26-threshold}
\end{equation}
then every bad-set \(L^2\) contribution controlled by the same form norm has
a strict relative coefficient and may be absorbed by
\Cref{thm:e8v26-energy-absorption} or
\Cref{thm:e8v26-defect-compiler}, after its non-form remainder is added to
\(\varepsilon^B\).
\end{corollary}

\begin{proof}
Use \eqref{eq:e8v26-kappaB} as the relative coefficient.  The strict
inequality is exactly the denominator condition in the two absorption
theorems.
\end{proof}

\begin{remark}[Capacity formulation]
\label{rem:e8v26-capacity}
More generally, it is enough to prove directly that the bad-set multiplier is
form-small:
\begin{equation}
 \|1_Bf\|_2^2
 \leq\varepsilon_B\|f\|_2^2+\kappa_B\mathcal E[f],
 \qquad \kappa_B<1.
 \label{eq:e8v26-capacity}
\end{equation}
Such an estimate may follow from a capacity, Hardy, or repair-map inequality
even when no useful \(L^{2q}\) Sobolev estimate is available.  The name
``capacity debit'' refers to \eqref{eq:e8v26-capacity}, not to probability
alone.
\end{remark}

\subsection{Dimension-growth obstruction}
\label{sec:e8v26-dimension}

\begin{proposition}[A deteriorating Sobolev exponent can erase rarity]
\label{prop:e8v26-dimension}
Let \(q_n>1\) with \(q_n\downarrow1\).  There are probabilities
\(p_n\downarrow0\) for which
\begin{equation}
 p_n^{1-1/q_n}\longrightarrow1.
 \label{eq:e8v26-rarity-erased}
\end{equation}
Thus a bad event may become rare while the coefficient in
\eqref{eq:e8v26-kappaB} does not become small.
\end{proposition}

\begin{proof}
Take \(q_n=1+1/n\) and \(p_n=e^{-\sqrt n}\).  Then
\[
 1-{1\over q_n}={1\over n+1},
 \qquad
 p_n^{1-1/q_n}=e^{-\sqrt n/(n+1)}\longrightarrow1.
\]
\end{proof}

\begin{assessment}[Meaning for regulator history spaces]
\label{ass:e8v26-dimension}
The number of metric, gauge, and matter variables grows with the regulator.
A finite-dimensional Sobolev inequality whose exponent or constant
collapses with dimension cannot prove a cofinal BGFD merely from an
exponentially small bad probability.  The required input is dimension-free:
a uniform log-Sobolev/hypercontractive estimate, a local repair-map capacity
bound, or another form-smallness theorem on the projective carrier.
\end{assessment}

\subsection{Good-geometry averaged landing}
\label{sec:e8v26-landing}

\begin{definition}[Good-geometry averaged form card, GGAFC]
\label{def:e8v26-GGAFC}
For every local source and compact parameter chart, GGAFC consists of:
\begin{enumerate}[label=\textup{(G\arabic*)},leftmargin=4em]
\item good events on which HEFC, DQSRC, local energy, and contact estimates
hold with cofinal constants;
\item the energy absorption row \eqref{eq:e8v26-bad-energy};
\item the base and derivative defect rows \eqref{eq:e8v26-bad-defect};
\item uniform strict absorption margins for all three rows;
\item summable non-form remainders and bridge changes of the good/bad split;
and
\item compatibility of the averaged carriers with the projective local
algebra and physical reflected-time incidence.
\end{enumerate}
\end{definition}

\begin{theorem}[GGAFC (C^1) landing]
\label{thm:e8v26-GGAFC}
If GGAFC holds, then the averaged local transfer matrix elements converge in
\(C^1\) on compact parameter charts.  Their base and derivative tails are
bounded by \eqref{eq:e8v26-total-tail} with the corresponding contact and
measure-bridge remainders included.  Positivity and normalization at zero
time pass to the inductive local algebra under the V15 hypotheses.
\end{theorem}

\begin{proof}
Use \Cref{thm:e8v26-energy-absorption} to obtain the global energy moment
needed to weight the good-chart V20 estimate.  Use
\Cref{thm:e8v26-defect-compiler} for the global adjacent base and derivative
defects.  GGAFC rows (G4)--(G5) make their tails summable, so V15 gives the
\(C^1\) limit.  Its positive-local-functional corollary proves the last
statement.
\end{proof}

\begin{firewall}[The carrier inequality remains open]
\label{fire:e8v26-open}
V26 proves what follows from a bad-geometry form debit; it does not prove that
the proposed SM--Einstein carrier satisfies one.  In particular, no
probability estimate for nearly degenerate metrics has been supplied, and no
dimension-free Sobolev, capacity, or repair-map inequality has been proved on
the gravitational history space.
\end{firewall}

\begin{assessment}[Progress at V26]
\label{ass:e8v26-status}
The worst-history requirement has been replaced by a precise averaged
compiler.  A strict relative bad-set coefficient is enough to absorb local
energy and adjacent defects, while a summable remainder preserves the V15
landing series.  A Sobolev estimate converts rarity into form smallness only
when its exponent and constant remain uniform.  The next upstream target is
therefore a dimension-free local repair or capacity inequality for the
metric carrier, not another pointwise ellipticity estimate.
\end{assessment}

\begin{programme}[Eighth-series V27: local metric repair map]
\label{prog:e8v26-V27}
V27 will formulate a finite-cell metric repair that projects eigenvalues into
a fixed nondegenerate interval while changing only the cells in a declared
bad cluster.  It will prove the measure-change and energy-cost identities
needed for a repair injection, then state the exact action-cost lower bound
which would imply \eqref{eq:e8v26-capacity}.  If the gravitational action has
a conformal direction of the wrong sign, that obstruction will be recorded
before any probability conclusion.
\end{programme}

\begin{firewall}[Claim boundary after eighth-series V26]
\label{fire:e8v26-boundary}
V26 proves good/bad absorption for local energy and adjacent \(C^1\) defects,
derives a bad-set form bound from a uniform Sobolev inequality, proves the
dimension-growth obstruction, and gives the conditional GGAFC landing
theorem.

V26 does not prove a bad-geometry capacity estimate for the gravitational
carrier, DQSRC, the complete SM contacts, reflection-positive gravity,
signed coercivity, cutoff removal, or construction of the Standard
Model--Einstein continuum.
\end{firewall}

\begin{theorem}[Eighth-series V26 result]
\label{thm:e8v26-result}
Good-chart HEFC estimates extend to averaged (C^1) local transfer data when
the bad-history energy and defect contributions obey strict relative
form bounds with summable remainders.  A uniform Sobolev inequality gives the
explicit relative coefficient \eqref{eq:e8v26-kappaB}; small probability
without such a dimension-free functional inequality is insufficient.
\end{theorem}

\begin{proof}
Combine \Cref{thm:e8v26-energy-absorption,
thm:e8v26-defect-compiler,lem:e8v26-small-set,
prop:e8v26-dimension,thm:e8v26-GGAFC}.
\end{proof}

\paragraph{Parent-version record.}
V26 was formed from
\texttt{Renewal\_SM\_Einstein\_Continuum\_Research\_Notes\_V25.tex}.
The V25 parent had (9\,134) logical source lines, (341\,078) bytes,
and SHA--256
\[
 \texttt{85687530D16CCA4378A21DD8882FFFDE47F36BED727756BDC6667F418D65C29C}.
\]
The next compulsory companion audit is V30.

% =============================================================================
% END APPENDED EIGHTH-SERIES V26 MATERIAL
% =============================================================================

% =============================================================================
% BEGIN APPENDED EIGHTH-SERIES V27 MATERIAL
% =============================================================================

\section{Metric repair injections and their exact entropy debit}
\label{sec:v8v27}

V26 reduced extension off a uniformly elliptic metric chart to a
bad-geometry form debit.  The present note constructs two finite-cell
repairs and distinguishes their uses.  Spectral clipping is the nearest
uniformly elliptic tensor, but it is many-to-one and therefore cannot by
itself support a Gibbs change of variables.  A rational spectral compression
is an exact diffeomorphism from the positive cone into the interior of the
ellipticity box.  Its Jacobian yields the missing entropy term in any repair
argument.

The frozen seventh-series V59 already proves that the physical-sign
Euclidean Einstein action is unbounded below on a uniformly elliptic
conformal sequence.  That result is used here rather than repeated.  It
shows that an eigenvalue repair can control a degeneracy event only; it
cannot control the independent high-frequency conformal escape.

\subsection{Nearest-point clipping and its lost fibres}
\label{sec:e8v27-clipping}

Let
\(
 \mathsf E_d=\{A\in M_d(\mathbb R):A^{\mathsf T}=A\}
\)
with Frobenius inner product, and for
\(0<\nu<\Lambda\) set
\begin{equation}
 \mathsf K_{\nu,\Lambda}
 :=\{A\in\mathsf E_d:\nu I\preceq A\preceq\Lambda I\}.
 \label{eq:e8v27-K}
\end{equation}
For a scalar \(s\), write
\begin{equation}
 c_{\nu,\Lambda}(s)=\min\{\Lambda,\max\{\nu,s\}\}.
 \label{eq:e8v27-scalar-clip}
\end{equation}
If \(A=Q\operatorname{diag}(\lambda_1,\ldots,\lambda_d)Q^{\mathsf T}\),
define
\begin{equation}
 \Pi_{\nu,\Lambda}(A)
 =Q\operatorname{diag}
  \bigl(c_{\nu,\Lambda}(\lambda_1),\ldots,
        c_{\nu,\Lambda}(\lambda_d)\bigr)Q^{\mathsf T}.
 \label{eq:e8v27-Pi}
\end{equation}

\begin{theorem}[Spectral clipping is the exact metric projection]
\label{thm:e8v27-projection}
The set \(\mathsf K_{\nu,\Lambda}\) is closed and convex, and
\(\Pi_{\nu,\Lambda}\) is its Frobenius nearest-point projection.  Hence
\begin{align}
 \|\Pi(A)-\Pi(B)\|_F&\leq\|A-B\|_F,
 \label{eq:e8v27-nonexpansive}\\
 \langle\Pi(A)-\Pi(B),A-B\rangle_F
 &\geq\|\Pi(A)-\Pi(B)\|_F^2.
 \label{eq:e8v27-firm}
\end{align}
\end{theorem}

\begin{proof}
The two Loewner inequalities in \eqref{eq:e8v27-K} define intersections of
closed convex half-spaces, so the set is closed and convex.  Let
\(B\in\mathsf K_{\nu,\Lambda}\).  Von Neumann's trace inequality gives,
after ordering the eigenvalues of \(A\) and \(B\),
\[
 \|A-B\|_F^2
 =\|A\|_F^2+\|B\|_F^2-2\operatorname{Tr}(AB)
 \geq\sum_{i=1}^d(\lambda_i(A)-\lambda_i(B))^2.
\]
Each \(\lambda_i(B)\in[\nu,\Lambda]\), and the scalar nearest point to
\(\lambda_i(A)\) in that interval is
\(c_{\nu,\Lambda}(\lambda_i(A))\).  Equality is attained by choosing the
same eigenvectors as \(A\).  Thus \eqref{eq:e8v27-Pi} is the unique metric
projection.

For completeness, if \(P\) is the metric projection onto any closed convex
set, the variational inequalities
\(
 \langle A-P A,P B-P A\rangle_F\leq0
\)
and
\(
 \langle B-P B,P A-P B\rangle_F\leq0
\)
sum to \eqref{eq:e8v27-firm}.  Cauchy--Schwarz then gives
\eqref{eq:e8v27-nonexpansive}.
\end{proof}

\begin{proposition}[Exact clipping has no repair Jacobian]
\label{prop:e8v27-clipping-collapse}
The map \(\Pi_{\nu,\Lambda}\) is not injective.  On every open spectral
region on which at least one eigenvalue lies below \(\nu\) or above
\(\Lambda\), its differential loses at least one direction.  In particular,
there is no positive cofinal lower bound on an ordinary change-of-variables
Jacobian for exact clipping.
\end{proposition}

\begin{proof}
Already on scalar matrices,
\(
 \Pi(sI)=\nu I
\)
for every \(s<\nu\), and
\(
 \Pi(sI)=\Lambda I
\)
for every \(s>\Lambda\).  On a region with a simple eigenvalue below
\(\nu\), varying that eigenvalue while holding its eigenspace and all other
eigenvalues fixed leaves its clipped value constant.  The corresponding
differential direction is therefore in the kernel.  The same argument
applies above \(\Lambda\).  Repeated eigenvalues follow by approximation;
noninjectivity alone already rules out a diffeomorphic change of variables.
\end{proof}

\begin{remark}[Defect coordinates must be retained]
\label{rem:e8v27-defect-coordinate}
The pair
\(
 (\Pi(A),A-\Pi(A))
\)
does determine \(A\).  Thus a proof based on exact clipping must integrate
the normal defect as an additional fibre variable.  Dropping that variable
is exactly the loss detected in
\Cref{prop:e8v27-clipping-collapse}.
\end{remark}

\subsection{An injective spectral repair}
\label{sec:e8v27-injective}

Put \(a=\Lambda-\nu\).  On the positive cone
\(\mathsf P_d=\{A\in\mathsf E_d:A\succ0\}\), define
\begin{equation}
 \mathcal R_{\nu,\Lambda}(A)
 :=\nu I+aA(I+A)^{-1}.
 \label{eq:e8v27-R}
\end{equation}
All matrix entries are dimensionless here; for metric tensors one first
uses a declared positive cell reference tensor to form its self-adjoint
relative endomorphism.

\begin{theorem}[Rational repair diffeomorphism and Jacobian]
\label{thm:e8v27-repair-diffeomorphism}
Let \(m=d(d+1)/2\).  The map
\begin{equation}
 \mathcal R_{\nu,\Lambda}:\mathsf P_d
 \longrightarrow
 \mathsf K_{\nu,\Lambda}^{\circ}
 :=\{B:\nu I\prec B\prec\Lambda I\}
 \label{eq:e8v27-diffeomorphism}
\end{equation}
is a real-analytic diffeomorphism.  Its inverse and its Jacobian with
respect to Lebesgue measure on \(\mathsf E_d\) are
\begin{align}
 \mathcal R_{\nu,\Lambda}^{-1}(B)
 &=(B-\nu I)(\Lambda I-B)^{-1},
 \label{eq:e8v27-inverse}\\
 J_{\nu,\Lambda}(A)
 &=a^m\det(I+A)^{-(d+1)}.
 \label{eq:e8v27-J}
\end{align}
\end{theorem}

\begin{proof}
For an eigenvalue \(\lambda>0\), the scalar function
\[
 r(\lambda)=\nu+a{\lambda\over1+\lambda}
\]
is analytic, strictly increasing, and maps \((0,\infty)\) bijectively to
\((\nu,\Lambda)\).  Functional calculus therefore proves the range and
injectivity.  Since
\(
 B-\nu I=aA(I+A)^{-1}
\)
and
\(
 \Lambda I-B=a(I+A)^{-1}
\), the two factors commute and their quotient is \(A\), proving
\eqref{eq:e8v27-inverse}.  The rational formulae show analyticity in both
directions.

Diagonalize \(A\) with eigenvalues \(\lambda_i\).  In an orthonormal basis
of symmetric matrices, the Fr\'echet derivative of a spectral function has
eigenvalues
\[
 r'(\lambda_i)={a\over(1+\lambda_i)^2}
 \quad(1\leq i\leq d),
 \qquad
 {r(\lambda_i)-r(\lambda_j)\over\lambda_i-\lambda_j}
 ={a\over(1+\lambda_i)(1+\lambda_j)}
 \quad(i<j).
\]
The divided difference is interpreted by continuity when eigenvalues
coincide.  Multiplying the \(d\) diagonal and \(d(d-1)/2\) off-diagonal
factors gives \(a^m\); each \(1+\lambda_i\) occurs with exponent two in its
diagonal factor and once in each of \(d-1\) off-diagonal factors.  The total
exponent is \(d+1\), which proves \eqref{eq:e8v27-J}.
\end{proof}

\begin{corollary}[Cellwise repair of a declared bad cluster]
\label{cor:e8v27-cellwise}
For \(N\) positive cell matrices \(g=(A_1,\ldots,A_N)\), let
\begin{equation}
 C(g)=\{i:A_i\notin\mathsf K_{\nu,\Lambda}\}.
 \label{eq:e8v27-bad-cluster}
\end{equation}
On the stratum \(\Omega_C=\{g:C(g)=C\}\), apply
\(\mathcal R_{\nu,\Lambda}\) to precisely the cells in \(C\) and leave the
others fixed.  The resulting map \(T_C\) is injective, sends
\(\Omega_C\) into the all-good chart, and has Jacobian
\begin{equation}
 J_C(g)=\prod_{i\in C}
 a^m\det(I+A_i)^{-(d+1)}.
 \label{eq:e8v27-JC}
\end{equation}
The spectral boundary of the strata is Lebesgue-null, so it does not affect
the change-of-variables formula for an absolutely continuous carrier.
\end{corollary}

\begin{proof}
This is the product of the diffeomorphisms in
\Cref{thm:e8v27-repair-diffeomorphism} on the declared cells and identity
maps on the remaining cells.  Products multiply the Jacobians.  A spectral
boundary is contained in the zero set of
\(
 \det(A_i-\nu I)\det(A_i-\Lambda I)
\), a nonzero polynomial, and is therefore Lebesgue-null.
\end{proof}

\begin{firewall}[Covariance of the repair]
\label{fire:e8v27-covariance}
The repair is equivariant under orthogonal changes of the declared cell
frame.  A gravitational regulator must additionally prove compatibility of
the reference tensors, cell frames, gauge fixing, and inter-scale bridges
with diffeomorphism Ward identities.  The finite-dimensional map alone does
not supply that covariance theorem.
\end{firewall}

\subsection{The exact Gibbs entropy debit}
\label{sec:e8v27-Gibbs}

Let \(\Omega\subset\mathbb R^M\) be a finite-dimensional regulator chart,
with reference density \(w>0\), action \(S\), partition function
\begin{equation}
 Z=\int_\Omega w(x)e^{-S(x)}\,dx<\infty,
 \qquad
 d\mu(x)=Z^{-1}w(x)e^{-S(x)}\,dx.
 \label{eq:e8v27-Gibbs}
\end{equation}

\begin{lemma}[Repair change of variables]
\label{lem:e8v27-change}
Let \(E\subset\Omega\) and let \(T:E\to\Omega\) be an injective
\(C^1\) map with positive Jacobian \(J_T\) almost everywhere.  Define the
effective repair cost
\begin{equation}
 \Gamma_T(x)
 :=S(x)-S(Tx)+\log J_T(x)
   +\log w(Tx)-\log w(x).
 \label{eq:e8v27-Gamma}
\end{equation}
Then
\begin{equation}
 \int_E w(x)e^{-S(x)}\,dx
 =\int_{T(E)}w(y)e^{-S(y)}
   e^{-\Gamma_T(T^{-1}y)}\,dy.
 \label{eq:e8v27-change}
\end{equation}
Consequently, if \(\Gamma_T\geq\delta\) almost everywhere on \(E\), then
\begin{equation}
 \mu(E)\leq e^{-\delta}.
 \label{eq:e8v27-probability}
\end{equation}
\end{lemma}

\begin{proof}
Set \(y=Tx\), so \(dy=J_T(x)dx\).  The ratio of the transformed density to
the Gibbs density at \(y\) is
\[
 {w(x)e^{-S(x)}\over J_T(x)w(Tx)e^{-S(Tx)}}
 =e^{-\Gamma_T(x)}.
\]
This proves \eqref{eq:e8v27-change}.  Under the lower bound, its right side
is at most
\(
 e^{-\delta}\int_{T(E)}w(y)e^{-S(y)}dy\leq e^{-\delta}Z
\).  Divide by \(Z\).
\end{proof}

\begin{remark}[Why the Jacobian sign matters]
\label{rem:e8v27-J-sign}
For the cell repair \eqref{eq:e8v27-R}, \(J_C\) may be very small.  The
quantity entering \(\Gamma_{T_C}\) is \(+\log J_C\), which is then a
negative entropy debit.  An action decrease under repair is useful only if
it also pays this contraction and the reference-density ratio.  Omitting
either term reverses the correct comparison.
\end{remark}

\begin{theorem}[Cluster repair bound]
\label{thm:e8v27-cluster}
Partition an \(N\)-cell chart, modulo null boundaries, into the strata
\(\Omega_C\), \(C\subset\{1,\ldots,N\}\).  Suppose the cellwise repair of
\Cref{cor:e8v27-cellwise} obeys
\begin{equation}
 \Gamma_{T_C}(g)\geq\delta |C|
 \quad\text{on }\Omega_C
 \label{eq:e8v27-linear-cost}
\end{equation}
for one \(\delta>0\).  Then, for \(1\leq r\leq N\),
\begin{align}
 \mu\{|C(g)|\geq r\}
 &\leq\sum_{k=r}^{N}{N\choose k}e^{-\delta k},
 \label{eq:e8v27-cluster-tail}\\
 \mu\{C(g)\ne\varnothing\}
 &\leq(1+e^{-\delta})^N-1
 \leq e^{Ne^{-\delta}}-1.
 \label{eq:e8v27-any-bad}
\end{align}
\end{theorem}

\begin{proof}
Apply \Cref{lem:e8v27-change} to every stratum.  There are
\({N\choose k}\) strata with \(k\) bad cells, which proves the first
bound.  Sum from \(k=1\) and use \(1+x\leq e^x\) for the second.
\end{proof}

\begin{corollary}[Repair cost plus a uniform Sobolev row implies BGFD]
\label{cor:e8v27-repair-BGFD}
Assume in addition that the Gibbs form obeys the V26 Sobolev row
\begin{equation}
 \|f\|_{2q}^2\leq S_q\|f\|_{\mathcal Q}^2,
 \qquad q>1.
 \label{eq:e8v27-Sobolev}
\end{equation}
For the event \(B=\{C(g)\ne\varnothing\}\), the bad-set multiplier has
relative coefficient
\begin{equation}
 \kappa_B
 \leq S_q\bigl(e^{Ne^{-\delta}}-1\bigr)^{1-1/q}.
 \label{eq:e8v27-kappa}
\end{equation}
If the right side is strictly below one, the V26 energy and adjacent-defect
absorption theorems apply.  Along a regulator sequence it is summably small
whenever
\begin{equation}
 \sum_n S_{q_n}
 \bigl(e^{N_ne^{-\delta_n}}-1\bigr)^{1-1/q_n}<\infty.
 \label{eq:e8v27-summable-repair}
\end{equation}
\end{corollary}

\begin{proof}
Insert \eqref{eq:e8v27-any-bad} into the small-set form estimate of V26.
The strict and summable conclusions are precisely the hypotheses of the
V26 compilers.
\end{proof}

\begin{corollary}[The volume entropy that a per-cell cost must beat]
\label{cor:e8v27-entropy-threshold}
If \(N_n\to\infty\), a scale-independent positive \(\delta\) does not make
the all-good complement rare through \eqref{eq:e8v27-any-bad}.  A sufficient
probability condition is
\begin{equation}
 N_ne^{-\delta_n}\longrightarrow0;
 \label{eq:e8v27-logN}
\end{equation}
in particular, 
\(\delta_n-\log N_n\to+\infty\).  Summable BGFD requires the stronger
quantitative row \eqref{eq:e8v27-summable-repair}.
\end{corollary}

\begin{proof}
The upper bound in \eqref{eq:e8v27-any-bad} tends to zero when
\eqref{eq:e8v27-logN} holds.  With fixed \(\delta\), its exponent grows like
\(N_ne^{-\delta}\), so this estimate gives no rarity.  The last assertion is
\Cref{cor:e8v27-repair-BGFD}.
\end{proof}

\subsection{Why ellipticity repair does not repair Einstein gravity}
\label{sec:e8v27-Einstein}

\begin{proposition}[Derivative escape lies inside the all-good stratum]
\label{prop:e8v27-conformal-inside}
There is a fixed box \(\mathsf K_{\nu,\Lambda}\) and a sequence of smooth
metrics \(g_N\) for which every cell sample lies in that box, while the
physical-sign Euclidean Einstein--Hilbert action plus arbitrary
cosmological and Weyl-square terms tends to \(-\infty\).  Hence a repair that
acts only on \(C(g)\) cannot yield a lower bound or tightness for that action.
\end{proposition}

\begin{proof}
This is the consequence of the uniformly elliptic conformal-instability
theorem proved in the frozen file
\texttt{Renewal\_SM\_Einstein\_Continuum\_Research\_Notes\_V100\_f7.tex},
seventh-series V59, label
\texttt{thm:newS7V59-instability}.  Its sequence is
\(g_N=e^{2\varepsilon\sin(Nx^1)}\delta\) with fixed
\(0<\varepsilon<1\).  Thus all eigenvalues lie in the fixed interval
\([e^{-2\varepsilon},e^{2\varepsilon}]\), so \(C(g_N)=\varnothing\), while
the action tends to \(-\infty\).  Since an empty-cluster repair is the
identity, it assigns no positive repair cost to this sequence.
\end{proof}

\begin{definition}[Metric repair injection card, MRIC]
\label{def:e8v27-MRIC}
A finite-cutoff gravitational carrier satisfies MRIC on a declared local
region if it provides:
\begin{enumerate}[label=\textup{(M\arabic*)},leftmargin=4em]
\item a measurable bad-cell stratification and an injective repair on every
stratum, with its exact Jacobian and reference-density ratio;
\item a cofinal lower bound on the effective cost \eqref{eq:e8v27-Gamma}
strong enough to beat the number of eligible bad clusters;
\item either a dimension-free functional inequality or a direct capacity
bound converting the resulting rarity into the V26 bad-set form debit;
\item compatible repairs for the complex metric chart and its first
parameter derivative, including measure and bridge contacts;
\item covariance and Ward compatibility of the cell reference tensors; and
\item a separate coercive mechanism for derivative escapes occurring
inside the uniformly elliptic chart, together with its physical positivity
and continuum interpretation.
\end{enumerate}
\end{definition}

\begin{firewall}[The action-cost row is not automatic]
\label{fire:e8v27-action-open}
Neither spectral clipping nor the rational diffeomorphism proves
\eqref{eq:e8v27-linear-cost}.  Curvature terms couple neighbouring cells,
the Jacobian debit grows along large eigenvalue fibres, and the
physical-sign Einstein term has the separate conformal escape in
\Cref{prop:e8v27-conformal-inside}.  A proposed gravitational regulator
must prove the full effective cost for its actual action and reference
measure; deterministic ellipticity or curvature coercivity alone is not
that proof.
\end{firewall}

\begin{assessment}[Progress at V27]
\label{ass:e8v27-status}
The finite-cell geometry of repair is now exact.  Nearest-point clipping is
stable but collapses normal fibres.  Rational spectral compression retains
those fibres and has the explicit determinant
\eqref{eq:e8v27-J}.  A lower bound on action change alone is insufficient:
the correct quantity is the action change plus the logarithmic Jacobian and
reference-density terms.  Even a linear per-cell effective cost must beat
the growing cluster entropy, and no eigenvalue repair sees the archived
high-frequency conformal instability inside the good chart.
\end{assessment}

\begin{programme}[Eighth-series V28: a direct repair-capacity inequality]
\label{prog:e8v27-V28}
V28 will go one step upstream and test whether a one-dimensional normal-fibre
Hardy estimate can convert repair coordinates directly into a capacity
debit, avoiding a dimension-dependent global Sobolev constant.  It will
separate the boundary trace term from the normal energy, identify the
uniform conditional density ratio required on each bad fibre, and record a
counterexample if an action gap without a fibre Poincar\'e estimate is
insufficient.
\end{programme}

\begin{firewall}[Claim boundary after eighth-series V27]
\label{fire:e8v27-boundary}
V27 proves the clipping projection theorem, its noninjectivity, an analytic
spectral repair diffeomorphism with exact Jacobian, the Gibbs repair identity,
the cluster entropy bound, and a conditional route from repair cost plus a
uniform Sobolev row to BGFD.  It also shows, using the archived V59 theorem,
that ellipticity repair cannot control the conformal derivative escape.

V27 does not prove MRIC for a gravitational regulator, a direct fibre
capacity estimate, a conformal positivity resolution, DQSRC, the full SM
contact ledger, reflection-positive gravity, cutoff removal, or construction
of the Standard Model--Einstein continuum.
\end{firewall}

\begin{theorem}[Eighth-series V27 result]
\label{thm:e8v27-result}
Bad metric cells admit an exact injective spectral compression into a fixed
ellipticity box, with Jacobian \eqref{eq:e8v27-J}.  For a Gibbs carrier the
effective repair cost is exactly \eqref{eq:e8v27-Gamma}; a linear cost gives
the cluster bounds \eqref{eq:e8v27-cluster-tail}--
\eqref{eq:e8v27-any-bad}, and together with a uniform Sobolev row gives the
explicit BGFD coefficient \eqref{eq:e8v27-kappa}.  This controls only
eigenvalue degeneracy.  It cannot control the physical-sign conformal escape,
which occurs entirely within one uniformly elliptic chart.
\end{theorem}

\begin{proof}
Combine \Cref{thm:e8v27-repair-diffeomorphism,
cor:e8v27-cellwise,lem:e8v27-change,thm:e8v27-cluster,
cor:e8v27-repair-BGFD,prop:e8v27-conformal-inside}.
\end{proof}

\paragraph{Parent-version record.}
V27 was formed from
\texttt{Renewal\_SM\_Einstein\_Continuum\_Research\_Notes\_V26.tex}.
The V26 parent had (9\,503) logical source lines, (354\,153) bytes,
and SHA--256
\[
 \texttt{B6A753BB1E9A361996D0BCDD7653CEFD45ED87E0EE7F804CE8DD1F8074CF46FD}.
\]
The next compulsory companion audit is V30.

% =============================================================================
% END APPENDED EIGHTH-SERIES V27 MATERIAL
% =============================================================================

% =============================================================================
% BEGIN APPENDED EIGHTH-SERIES V28 MATERIAL
% =============================================================================

\section{Normal-fibre Hardy control for bad metric cells}
\label{sec:v8v28}

V27 obtained a probability estimate from an injective repair only after the
action change paid the exact Jacobian and reference-density debits.  V26
then converted probability into form smallness by a global Sobolev
inequality.  On a regulator space whose dimension grows, that Sobolev
constant or exponent may collapse.  This note proves a different route:
disintegrate the carrier along a one-dimensional repair fibre and use an
anchored Hardy estimate on that fibre.  The resulting constant depends on
the conditional tail profile and on a collar trace, not on the number of
spectator variables.

\subsection{An anchored exponential-tail inequality}
\label{sec:e8v28-Hardy}

\begin{lemma}[Anchored exponential Hardy inequality]
\label{lem:e8v28-exp-Hardy}
Let \(a>0\) and let \(g\) be locally absolutely continuous on
\([0,\infty)\), with \(g(0)=0\) and
\(\int_0^\infty(|g|^2+|g'|^2)e^{-ax}\,dx<\infty\).  Then
\begin{equation}
 \int_0^\infty |g(x)|^2e^{-ax}\,dx
 \leq {4\over a^2}
 \int_0^\infty |g'(x)|^2e^{-ax}\,dx.
 \label{eq:e8v28-exp-Hardy}
\end{equation}
\end{lemma}

\begin{proof}
First take \(g\) smooth with compact support and \(g(0)=0\).  Integration
by parts gives
\[
 \int_0^\infty g^2e^{-ax}\,dx
 ={2\over a}\int_0^\infty gg'e^{-ax}\,dx.
\]
Cauchy--Schwarz implies that the left side is at most
\[
 {2\over a}
 \left(\int_0^\infty g^2e^{-ax}\,dx\right)^{1/2}
 \left(\int_0^\infty (g')^2e^{-ax}\,dx\right)^{1/2}.
\]
If the first factor vanishes there is nothing to prove; otherwise divide
and square.  General \(g\) follows by truncation and density in the weighted
Sobolev space.
\end{proof}

\begin{lemma}[Bounded perturbations of an exponential tail]
\label{lem:e8v28-perturbed-Hardy}
Let
\begin{equation}
 \rho(x)=c\,e^{-ax-u(x)},\qquad x\geq0,
 \label{eq:e8v28-rho}
\end{equation}
where \(c>0\), \(a\geq a_0>0\), and
\(\operatorname*{ess\,osc}_{x\geq0}u(x)\leq K\).  For every locally
absolutely continuous \(g\) with \(g(0)=0\),
\begin{equation}
 \int_0^\infty |g|^2\rho\,dx
 \leq {4e^K\over a_0^2}
 \int_0^\infty |g'|^2\rho\,dx.
 \label{eq:e8v28-perturbed-Hardy}
\end{equation}
\end{lemma}

\begin{proof}
Let \(u_-=\operatorname*{ess\,inf}u\) and
\(u_+=\operatorname*{ess\,sup}u\).  The left side is at most
\(ce^{-u_-}\int g^2e^{-ax}dx\).  Apply
\Cref{lem:e8v28-exp-Hardy}, use \(a\geq a_0\), and bound the unperturbed
gradient integral by
\[
 \int_0^\infty |g'|^2e^{-ax}dx
 \leq c^{-1}e^{u_+}\int_0^\infty |g'|^2\rho\,dx.
\]
The quotient is \(e^{u_+-u_-}\leq e^K\).
\end{proof}

\subsection{A collar controls the fibre anchor}
\label{sec:e8v28-collar}

\begin{lemma}[One-dimensional collar trace]
\label{lem:e8v28-trace}
Let \(\ell>0\), let \(f\in H^1(-\ell,0)\), and suppose a density
\(\rho\) satisfies \(\rho(x)\geq m>0\) almost everywhere on
\([-\ell,0]\).  Then
\begin{equation}
 |f(0)|^2
 \leq {2\over \ell m}\int_{-\ell}^{0}|f|^2\rho\,dx
      {2\ell\over m}\int_{-\ell}^{0}|f'|^2\rho\,dx.
 \label{eq:e8v28-trace}
\end{equation}
\end{lemma}

\begin{proof}
For \(-\ell\leq t\leq0\),
\[
 |f(0)|^2
 \leq2|f(t)|^2+2\left|\int_t^0f'(s)\,ds\right|^2
 \leq2|f(t)|^2+2\ell\int_{-\ell}^0|f'(s)|^2\,ds.
\]
Average in \(t\) over the collar, then use
\(\int h\,dx\leq m^{-1}\int h\rho\,dx\) for \(h\geq0\).
\end{proof}

Let \((Y,\eta)\) be a probability space of all spectator variables.  On a
normal repair coordinate \(r\in[-\ell,\infty)\), suppose the carrier
disintegrates as
\begin{equation}
 d\mu(y,r)=d\eta(y)\rho_y(r)\,dr,
 \qquad
 \int_{-\ell}^\infty\rho_y(r)\,dr=1
 \quad\text{for \(\eta\)-almost every \(y\)}.
 \label{eq:e8v28-disintegration}
\end{equation}
The collar is \(G_\ell=\{-\ell\leq r\leq0\}\), and the bad half-fibre is
\(B=\{r\geq0\}\).  Set
\begin{equation}
 p_y=\int_0^\infty\rho_y(r)\,dr.
 \label{eq:e8v28-py}
\end{equation}

\begin{theorem}[Direct normal-fibre form debit]
\label{thm:e8v28-fibre-debit}
Assume there are constants
\[
 p_y\leq p_*,
 \qquad
 \rho_y(r)\geq m_*>0\quad(-\ell\leq r\leq0),
 \label{eq:e8v28-collar-data}
\]
and, for \(r\geq0\),
\begin{equation}
 \rho_y(r)=c_y e^{-a_yr-u_y(r)},\qquad
 a_y\geq a_0>0,\qquad
 \operatorname*{ess\,osc}u_y\leq K,
 \label{eq:e8v28-tail-data}
\end{equation}
uniformly for \(\eta\)-almost every \(y\).  Then every function in the
corresponding fibre Sobolev space obeys
\begin{align}
 \int_B|f|^2\,d\mu
 &\leq \varepsilon_{\rm fib}\int_{G_\ell}|f|^2\,d\mu
       +\kappa_{\rm fib}
        \int_{G_\ell\cup B}|\partial_rf|^2\,d\mu,
 \label{eq:e8v28-fibre-debit}\\
 \varepsilon_{\rm fib}
 &:={4p_*\over\ell m_*},
 \label{eq:e8v28-epsilon}\\
 \kappa_{\rm fib}
 &:={4p_*\ell\over m_*}+{8e^K\over a_0^2}.
 \label{eq:e8v28-kappa}
\end{align}
In particular, if \(\kappa_{\rm fib}<1\), the bad fibre is relatively
form-small with a strict, spectator-dimension-free energy coefficient.
\end{theorem}

\begin{proof}
Fix \(y\), put \(b=f(y,0)\), and set \(g(r)=f(y,r)-b\) for \(r\geq0\).
Then \(g(0)=0\), and
\[
 |f(y,r)|^2\leq2|b|^2+2|g(r)|^2.
\]
Integrating over the bad half-fibre and applying
\Cref{lem:e8v28-perturbed-Hardy} gives
\begin{equation}
 \int_0^\infty|f(y,r)|^2\rho_y(r)\,dr
 \leq2p_y|b|^2
 {8e^K\over a_0^2}
 \int_0^\infty|\partial_rf(y,r)|^2\rho_y(r)\,dr.
 \label{eq:e8v28-fixed-y}
\end{equation}
Apply \Cref{lem:e8v28-trace} to \(b\), use
\(p_y\leq p_*\), and integrate in \(y\).  The two trace terms give
\eqref{eq:e8v28-epsilon} and the first term in
\eqref{eq:e8v28-kappa}; the anchored Hardy term gives the second.
\end{proof}

\begin{corollary}[No global Sobolev exponent is used]
\label{cor:e8v28-no-global-Sobolev}
Suppose the full regulator form \(\mathcal E_n\) dominates the normal energy
in \eqref{eq:e8v28-fibre-debit}, and the constants
\((p_*,m_*,\ell,a_0,K)\) are uniform in \(n\).  Then
\begin{equation}
 \int_B|f|^2\,d\mu_n
 \leq\varepsilon_{\rm fib}\|f\|_{L^2(\mu_n)}^2
     \kappa_{\rm fib}\mathcal E_n[f].
 \label{eq:e8v28-global-form}
\end{equation}
This implication is independent of the dimension of \(Y\).
\end{corollary}

\begin{proof}
Enlarge the collar \(L^2\) integral to the full \(L^2\) norm and the normal
energy to \(\mathcal E_n[f]\).  No estimate is applied in the \(Y\)
variables.
\end{proof}

\begin{remark}[Several bad normal coordinates]
\label{rem:e8v28-overlap}
If a local source meets bad-fibre events \(B_j\) and every history belongs
to at most \(L\) of the chosen fibre charts, summing
\eqref{eq:e8v28-global-form} gives coefficients at most
\(L\varepsilon_{\rm fib}\) and \(L\kappa_{\rm fib}\).  This remains
dimension-free only when the overlap \(L\) is cofinally bounded.  Applying
the estimate separately to all \(N_n\) cells produces an \(N_n\) debit and
does not solve the global all-good problem.
\end{remark}

\subsection{What an action comparison must prove on a fibre}
\label{sec:e8v28-action}

In repair coordinates the conditional density has the form
\begin{equation}
 \rho_y(r)=Z_y^{-1}
 \exp\{-S_y(r)\}\,w_y(r)J_y(r).
 \label{eq:e8v28-conditional-density}
\end{equation}
Here \(J_y\) includes the coordinate Jacobian and \(w_y\) the declared
reference density.  Define the conditional effective potential
\begin{equation}
 V_y(r)=S_y(r)-\log w_y(r)-\log J_y(r).
 \label{eq:e8v28-effective-potential}
\end{equation}

\begin{proposition}[A sufficient affine-tail row]
\label{prop:e8v28-affine-tail}
Suppose, on every bad half-fibre,
\begin{equation}
 V_y(r)=V_y(0)+a_yr+u_y(r),
 \qquad
 a_y\geq a_0>0,
 \qquad
 \operatorname*{ess\,osc}_{r\geq0}u_y(r)\leq K.
 \label{eq:e8v28-affine-tail}
\end{equation}
Then the conditional density satisfies
\eqref{eq:e8v28-tail-data}.  If the good collar also supplies
\eqref{eq:e8v28-collar-data}, the direct capacity estimate
\eqref{eq:e8v28-fibre-debit} follows.
\end{proposition}

\begin{proof}
Substitute \eqref{eq:e8v28-affine-tail} into
\eqref{eq:e8v28-conditional-density}.  The factor
\(Z_y^{-1}e^{-V_y(0)}\) is the constant \(c_y\), leaving
\(e^{-a_yr-u_y(r)}\).  Apply \Cref{thm:e8v28-fibre-debit}.
\end{proof}

\begin{assessment}[The new requirement is a conditional slope, not a mass
gap]
\label{ass:e8v28-slope}
The V27 effective repair cost compares two points.  The affine-tail row
\eqref{eq:e8v28-affine-tail} controls the complete conditional profile
between and beyond them.  It excludes a long flat tail and a thin
low-conductance neck.  This is stronger than a probability estimate, but it
is local in one repair coordinate and therefore avoids the global
dimension-growth obstruction of V26.
\end{assessment}

\subsection{Rarity alone does not give fibre capacity}
\label{sec:e8v28-counterexample}

\begin{proposition}[A small bad well can have arbitrarily small conductance]
\label{prop:e8v28-small-well}
Fix \(0<p<1\), \(0\leq\varepsilon<1\), and
\(\kappa<\infty\).  There are a smooth strictly positive probability
density \(\rho\) on \([-2,2]\), a set \(B=[1,2]\) with
\(\int_B\rho=p\), and a smooth \(f\) such that
\begin{equation}
 \int_B f^2\rho\,dx
 >
 \varepsilon\int_{-2}^{2}f^2\rho\,dx
 +\kappa\int_{-2}^{2}(f')^2\rho\,dx.
 \label{eq:e8v28-capacity-failure}
\end{equation}
Thus the bad probability, even when prescribed and small, supplies no
uniform relative form bound.
\end{proposition}

\begin{proof}
Choose nonnegative smooth probability bumps \(\phi_L,\phi_R\), supported
respectively in \((-2,-1)\) and \((1,2)\).  For
\(0<\eta<\min\{4p,4(1-p)/3\}\), set
\[
 \rho_\eta(x)
 =\left(1-p-{3\eta\over4}\right)\phi_L(x)
  +\left(p-{\eta\over4}\right)\phi_R(x)
  +{\eta\over4}.
\]
It is smooth, strictly positive, integrates to one, and its mass on
\(B=[1,2]\) is exactly
\((p-\eta/4)+\eta/4=p\).

Choose a fixed smooth \(f\) with \(0\leq f\leq1\), equal to zero on
\([-2,-1/2]\) and to one on \([1/2,2]\).  The two bumps do not meet the
transition region, so
\[
 \int_{-2}^{2}(f')^2\rho_\eta\,dx
 ={\eta\over4}\int_{-2}^{2}(f')^2\,dx
 =C_1\eta.
\]
Moreover,
\[
 \int_Bf^2\rho_\eta\,dx=p,
 \qquad
 \int_{-2}^{2}f^2\rho_\eta\,dx\leq p+C_2\eta
\]
for a constant \(C_2\) depending only on \(f\).  Since
\((1-\varepsilon)p>0\), choose \(\eta\) so small that
\(\varepsilon C_2\eta+\kappa C_1\eta<(1-\varepsilon)p\).
This is exactly \eqref{eq:e8v28-capacity-failure}.
\end{proof}

\begin{corollary}[Repair probability does not populate BGFD]
\label{cor:e8v28-probability-not-BGFD}
The action-cost probability estimate of V27 cannot by itself populate the
V26 bad-geometry form debit.  A conditional Hardy/Poincar\'e row, a
dimension-free Sobolev row, or a direct capacity estimate is logically
independent and must be proved.
\end{corollary}

\begin{proof}
\Cref{prop:e8v28-small-well} holds for every prescribed \(p>0\) and every
candidate strict form coefficient.  It therefore rules out an implication
from bad mass alone to a uniform capacity inequality.
\end{proof}

\subsection{The fibre-capacity card}
\label{sec:e8v28-card}

\begin{definition}[Fibre repair capacity card, FRCC]
\label{def:e8v28-FRCC}
For each bad metric mode seen by a fixed local source, FRCC consists of:
\begin{enumerate}[label=\textup{(F\arabic*)},leftmargin=4em]
\item a coarea or disintegration theorem producing repair coordinates
\((y,r)\) and identifying the normal derivative inside the regulated
Dirichlet form;
\item a good collar with uniform width and conditional density lower bound;
\item a uniform affine-tail representation
\eqref{eq:e8v28-affine-tail}, including the action, reference density, and
repair Jacobian;
\item a cofinally bounded overlap number for the fibre charts met by the
local source and its refinement collar;
\item the same estimates on one fixed complex parameter circle, with first
metric and refinement contacts; and
\item summable changes of the disintegration, collar, and bad-mode
assignment under adjacent-scale bridges.
\end{enumerate}
\end{definition}

\begin{theorem}[FRCC implies the bad-set analytic rows of GGAFC]
\label{thm:e8v28-FRCC-GGAFC}
Suppose FRCC holds and the summed normal-energy coefficient satisfies
\begin{equation}
 L\left({4p_*\ell\over m_*}+{8e^K\over a_0^2}\right)<1.
 \label{eq:e8v28-strict}
\end{equation}
Suppose also that its collar \(L^2\) terms and adjacent bridge changes give
the summable non-form remainders required in V26.  Then the energy, base
defect, and first-derivative bad-set rows of GGAFC hold for that local
source.  Combined with the good-chart rows, they give the conditional
\(C^1\) landing theorem of V26.
\end{theorem}

\begin{proof}
Apply \Cref{thm:e8v28-fibre-debit} in each FRCC chart and sum using the
overlap bound.  Equation \eqref{eq:e8v28-strict} is the strict absorption
margin.  FRCC rows (F5)--(F6) provide the same estimate for the complexified
base and derivative defects and make the remaining terms summable.  These
are exactly GGAFC rows (G2)--(G5), so the V26 landing theorem applies.
\end{proof}

\begin{firewall}[No gravitational fibre law has been proved]
\label{fire:e8v28-open}
V28 proves a functional-analytic compiler.  It does not construct repair
coordinates for the full metric quotient, prove the affine-tail row for an
Einstein or stabilized gravitational action, identify the normal derivative
inside a reflection-positive gravitational form, or control the number of
bad fibres met by the physical transfer observable.
\end{firewall}

\begin{assessment}[Progress at V28]
\label{ass:e8v28-status}
A global Sobolev inequality is no longer the only route from repair to
BGFD.  Uniform conditional exponential tails, a nondegenerate collar, and a
normal Dirichlet energy give the explicit direct debit
\eqref{eq:e8v28-fibre-debit}.  The counterexample proves that the missing
conditional conductance cannot be inferred from the V27 probability bound.
The next task is to test whether the actual finite metric action can provide
the affine-tail slope after the repair Jacobian is included.
\end{assessment}

\begin{programme}[Eighth-series V29: testing the repair slope on a discrete
curvature action]
\label{prog:e8v28-V29}
V29 will analyze the simplest conformal finite-difference fibre.  It will
compute the action as one cell conformal variable tends to either
degeneracy or infinity, include the metric-coordinate Jacobian, and decide
which signs of cosmological, Einstein, and curvature-square owners can
supply the affine-tail row.  The computation will be kept separate from
reflection positivity and from the already proved smooth conformal
instability.
\end{programme}

\begin{firewall}[Claim boundary after eighth-series V28]
\label{fire:e8v28-boundary}
V28 proves an anchored Hardy inequality, its bounded-perturbation form, a
collar trace, a spectator-dimension-free fibre capacity estimate, and the
conditional FRCC-to-GGAFC implication.  It also proves by a smooth
two-well example that arbitrarily small bad probability does not imply a
relative form bound.

V28 does not prove FRCC for gravity, the conformal repair slope, physical
reflection positivity, DQSRC, the full SM contact ledger, cutoff removal,
or construction of the Standard Model--Einstein continuum.
\end{firewall}

\begin{theorem}[Eighth-series V28 result]
\label{thm:e8v28-result}
An affine conditional effective potential with slope \(a_0>0\), bounded
tail oscillation, and a nondegenerate good collar gives the explicit
bad-fibre form coefficient \eqref{eq:e8v28-kappa}, independently of all
spectator dimensions.  With bounded chart overlap and a strict summed
coefficient, FRCC populates the analytic bad-set rows of GGAFC.  A repair
probability estimate without this conditional conductance is insufficient.
\end{theorem}

\begin{proof}
Combine \Cref{lem:e8v28-perturbed-Hardy,
lem:e8v28-trace,thm:e8v28-fibre-debit,
prop:e8v28-small-well,thm:e8v28-FRCC-GGAFC}.
\end{proof}

\paragraph{Parent-version record.}
V28 was formed from
\texttt{Renewal\_SM\_Einstein\_Continuum\_Research\_Notes\_V27.tex}.
The V27 parent had (10\,011) logical source lines, (373\,253) bytes,
and SHA--256
\[
 \texttt{5D33CC954A312AFA6E55FB08997F97C7CB6A5C67AB5B1AEFBAB586AF81DEA2EC}.
\]
The next compulsory companion audit is V30.

% =============================================================================
% END APPENDED EIGHTH-SERIES V28 MATERIAL
% =============================================================================

% =============================================================================
% BEGIN APPENDED EIGHTH-SERIES V29 MATERIAL
% =============================================================================

\section{Tail slopes of a discrete conformal repair fibre}
\label{sec:v8v29}

V28 showed that a direct bad-fibre capacity estimate follows from a positive
outward slope of the conditional effective potential.  The present note
tests that slope on a declared finite-difference conformal model.  It does
not identify the model with a completed gravitational regulator.  Its role
is diagnostic: it determines which local owners can suppress a single
conformal amplitude after the metric-coordinate density is included, and
which escapes remain invisible to that calculation.

\subsection{Hardy control from a one-sided potential slope}
\label{sec:e8v29-slope-Hardy}

\begin{lemma}[Outward-slope Hardy inequality]
\label{lem:e8v29-slope-Hardy}
Let \(V\) be locally absolutely continuous on \([0,\infty)\), suppose
\(V'(r)\geq a_0>0\) almost everywhere, and let
\(\rho(r)=ce^{-V(r)}\), \(c>0\).  If \(g(0)=0\), then
\begin{equation}
 \int_0^\infty |g|^2\rho\,dr
 \leq {4\over a_0^2}
 \int_0^\infty |g'|^2\rho\,dr
 \label{eq:e8v29-slope-Hardy}
\end{equation}
for every function in the weighted Sobolev space.
\end{lemma}

\begin{proof}
It is enough first to take a compactly supported smooth \(g\).  Since
\(V'\geq a_0\),
\[
 I:=\int_0^\infty g^2e^{-V}\,dr
 \leq {1\over a_0}\int_0^\infty g^2V'e^{-V}\,dr.
\]
Because \(V'e^{-V}=-(e^{-V})'\), integration by parts and \(g(0)=0\)
give
\[
 I\leq {2\over a_0}\int_0^\infty gg'e^{-V}\,dr
 \leq {2\over a_0}I^{1/2}
 \left(\int_0^\infty(g')^2e^{-V}\,dr\right)^{1/2}.
\]
Divide and square.  The constant \(c\) cancels, and truncation plus density
gives the general case.
\end{proof}

\begin{corollary}[Slope version of the V28 fibre debit]
\label{cor:e8v29-slope-debit}
In the setting of the V28 collar theorem, replace the affine-tail hypothesis
by
\begin{equation}
 {d\over dr}V_y(r)\geq a_0>0
 \quad\text{on every bad half-fibre}.
 \label{eq:e8v29-slope-row}
\end{equation}
Then
\begin{equation}
 \int_B|f|^2\,d\mu
 \leq {4p_*\over\ell m_*}\int_{G_\ell}|f|^2\,d\mu
 +\left({4p_*\ell\over m_*}+{8\over a_0^2}\right)
 \int_{G_\ell\cup B}|\partial_rf|^2\,d\mu.
 \label{eq:e8v29-slope-debit}
\end{equation}
\end{corollary}

\begin{proof}
Repeat the proof of the V28 direct normal-fibre theorem, applying
\Cref{lem:e8v29-slope-Hardy} to
\(g(r)=f(y,r)-f(y,0)\).  The collar trace is unchanged.
\end{proof}

\subsection{A declared conformal star action}
\label{sec:e8v29-star}

Let one lattice vertex carry a conformal variable
\(\phi\in\mathbb R\), let its \(q\) neighbouring values be
\(\psi=(\psi_1,\ldots,\psi_q)\), and let \(w_j>0\).  Put
\begin{align}
 W&=\sum_{j=1}^q w_j,
 &B_\psi&=\sum_{j=1}^q w_j\psi_j,
 &C_\psi&=\sum_{j=1}^q w_j\psi_j^2,
 \label{eq:e8v29-WBC}\\
 D_\psi(\phi)
 &=\sum_{j=1}^q w_j(\phi-\psi_j)^2
 =W\phi^2-2B_\psi\phi+C_\psi,
 \label{eq:e8v29-D}\\
 L_\psi(\phi)
 &=\sum_{j=1}^q w_j(\psi_j-\phi)
 =B_\psi-W\phi,
 \label{eq:e8v29-L}\\
 Q_\psi(\phi)&=L_\psi(\phi)+D_\psi(\phi).
 \label{eq:e8v29-Q}
\end{align}
The following one-fibre action retains the exact conformal scaling of the
cosmological term, the wrong-sign two-derivative term, and a positive
curvature-square leading owner:
\begin{equation}
 S_\psi(\phi)
 =36\alpha Q_\psi(\phi)^2
  -6\kappa e^{2\phi}D_\psi(\phi)
  +\Lambda e^{4\phi}+U_\psi(\phi).
 \label{eq:e8v29-action}
\end{equation}
Here \(\kappa>0\) is the physical Einstein sign in the Euclidean conformal
sector.  The expression is a declared scalar star model, not a
diffeomorphism-invariant lattice curvature formula.

Suppose the reference measure on this coordinate is
\begin{equation}
 d\rho_\beta(\phi)=e^{\beta\phi}\,d\phi.
 \label{eq:e8v29-reference}
\end{equation}
For example, Lebesgue measure in the scalar metric eigenvalue
\(\lambda=e^{2\phi}\) gives \(d\lambda=2e^{2\phi}d\phi\) and hence
\(\beta=2\), up to a constant.  The conditional effective potential is
\begin{equation}
 V_\psi(\phi)=S_\psi(\phi)-\beta\phi.
 \label{eq:e8v29-V}
\end{equation}

\begin{lemma}[Polynomial and exponential asymptotics]
\label{lem:e8v29-asymptotics}
Fix a compact neighbour set \(|\psi_j|\leq M\).  Uniformly on that set,
\begin{align}
 D_\psi(\phi)&=W\phi^2+O(|\phi|+1),
 \label{eq:e8v29-D-asymp}\\
 Q_\psi(\phi)&=W\phi^2+O(|\phi|+1),
 \label{eq:e8v29-Q-asymp}\\
 Q_\psi(\phi)Q_\psi'(\phi)
 &=2W^2\phi^3+O(\phi^2+1).
 \label{eq:e8v29-QQ-asymp}
\end{align}
\end{lemma}

\begin{proof}
The bounds \(|B_\psi|\leq WM\) and \(0\leq C_\psi\leq WM^2\) are uniform.
Insert them into \eqref{eq:e8v29-D}--\eqref{eq:e8v29-Q}.  The leading
coefficient of \(Q_\psi\) is \(W\), and that of \(Q_\psi'\) is \(2W\);
multiplication gives the last formula.
\end{proof}

\begin{theorem}[Tail classification for the conformal star]
\label{thm:e8v29-tail-classification}
Assume \(\kappa>0\), \(|\psi_j|\leq M\), and the remainder is \(C^1\) with
\begin{align}
 U_\psi'(\phi)&=o(e^{4\phi})
 &&(\phi\to+\infty),
 \label{eq:e8v29-U-plus}\\
 U_\psi'(\phi)&=o(|\phi|^3)
 &&(\phi\to-\infty),
 \label{eq:e8v29-U-minus}
\end{align}
uniformly in the neighbour set.

\begin{enumerate}[label=\textup{(\roman*)},leftmargin=3.8em]
\item If \(\Lambda>0\), then
\(V_\psi'(\phi)\to+\infty\) as \(\phi\to+\infty\), uniformly in
\(\psi\).
\item If \(\Lambda<0\), then
\(V_\psi(\phi)\to-\infty\) and
\(V_\psi'(\phi)\to-\infty\) as \(\phi\to+\infty\).
\item If \(\Lambda=0\), then the physical Einstein term dominates every
curvature-square polynomial:
\[
 V_\psi(\phi)\to-\infty,
 \qquad
 V_\psi'(\phi)\to-\infty
 \quad(\phi\to+\infty).
 \label{eq:e8v29-zero-Lambda-escape}
\]
\item If \(\alpha>0\), the lower outward potential
\[
 V_{\psi,-}(r):=V_\psi(-r)
 \label{eq:e8v29-lower-potential}
\]
satisfies \(V_{\psi,-}'(r)\to+\infty\) as \(r\to+\infty\), uniformly in
\(\psi\).
\item If \(\alpha=0\) and
\(U_\psi'(\phi)=o(1)\) as \(\phi\to-\infty\), then
\[
 V_{\psi,-}'(r)\longrightarrow\beta.
 \label{eq:e8v29-lower-beta}
\]
Thus the displayed action gives no lower-tail slope in this case unless
the reference density itself has \(\beta>0\).
\end{enumerate}
\end{theorem}

\begin{proof}
Differentiate \eqref{eq:e8v29-action}:
\begin{align}
 S_\psi'(\phi)
 ={}&72\alpha Q_\psi Q_\psi'
 -6\kappa e^{2\phi}\bigl(2D_\psi+D_\psi'\bigr)
 +4\Lambda e^{4\phi}+U_\psi'(\phi).
 \label{eq:e8v29-Sprime}
\end{align}
The derivative of \(-\beta\phi\) is the fixed constant \(-\beta\).

As \(\phi\to+\infty\), the cosmological derivative has order
\(e^{4\phi}\), the Einstein derivative has negative order
\(e^{2\phi}\phi^2\), and the curvature-square derivative has order
\(\phi^3\).  Under \eqref{eq:e8v29-U-plus}, a positive cosmological
coefficient proves (i), while a negative one proves (ii).  When
\(\Lambda=0\), the negative Einstein term dominates all polynomial terms,
which proves (iii), including the potential statement by the same
asymptotic comparison before differentiation.

As \(\phi\to-\infty\), both exponential terms and their derivatives vanish
faster than every inverse power.  By
\Cref{lem:e8v29-asymptotics},
\[
 72\alpha Q_\psi(\phi)Q_\psi'(\phi)
 =144\alpha W^2\phi^3+O(\phi^2+1).
\]
For \(\alpha>0\) this tends to \(-\infty\).  Since
\(V_{\psi,-}'(r)=-V_\psi'(-r)\), (iv) follows.  If \(\alpha=0\), all action
derivatives tend to zero under the stronger remainder hypothesis, and
\(-V_\psi'(-r)\to\beta\), proving (v).
\end{proof}

\begin{corollary}[When the star supplies two fibre slopes]
\label{cor:e8v29-two-slopes}
Under the hypotheses of \Cref{thm:e8v29-tail-classification}, if
\(\Lambda>0\) and \(\alpha>0\), then for every chosen \(a_0>0\) there are
uniform thresholds \(R_+,R_->0\) such that
\begin{align}
 V_\psi'(\phi)&\geq a_0
 &&(\phi\geq R_+),\label{eq:e8v29-upper-slope}\\
 {d\over dr}V_\psi(-r)&\geq a_0
 &&(r\geq R_-).\label{eq:e8v29-lower-slope}
\end{align}
If the intervening collars have uniform conditional density lower bounds,
\Cref{cor:e8v29-slope-debit} gives direct upper- and lower-bad-fibre
capacity debits.
\end{corollary}

\begin{proof}
Uniform divergence of the two derivatives gives uniform thresholds.
Apply the slope debit separately to the two tails.
\end{proof}

\begin{corollary}[The coordinate Jacobian cannot cure the physical upper
escape]
\label{cor:e8v29-Jacobian-no-cure}
For every fixed \(\beta\), the linear term \(-\beta\phi\) cannot change
\eqref{eq:e8v29-zero-Lambda-escape}.  In particular, the metric-coordinate
Jacobian supplies no positive upper repair slope when
\(\Lambda=0\), \(\kappa>0\), and only polynomial curvature-square owners
oppose the escape.
\end{corollary}

\begin{proof}
The term \(-\beta\phi\) and its derivative are lower order than
\(e^{2\phi}\phi^2\).
\end{proof}

\subsection{Limits of the one-cell classification}
\label{sec:e8v29-limits}

\begin{proposition}[Amplitude-tail control does not address the archived
oscillatory instability]
\label{prop:e8v29-amplitude-vs-frequency}
The two slope bounds in \Cref{cor:e8v29-two-slopes}, when available, do not
control the uniformly elliptic high-frequency conformal sequence of the
frozen seventh-series V59 theorem.
\end{proposition}

\begin{proof}
The star bounds concern \(|\phi|\to\infty\) with neighbours in a fixed
compact set.  In the archived sequence
\(\phi_N=\varepsilon\sin(Nx^1)\), the amplitude remains bounded by the fixed
\(\varepsilon\) while its frequency diverges.  It therefore remains in the
intermediate conformal collar of every amplitude-tail decomposition, even
though the physical-sign Einstein action tends to \(-\infty\).
\end{proof}

\begin{proposition}[Neighbour compactness is an unpaid cluster row]
\label{prop:e8v29-neighbour-row}
The thresholds in \Cref{cor:e8v29-two-slopes} need not be uniform when the
neighbour values are unrestricted.  Hence a one-cell tail proof must be
combined with either a bounded-neighbour conditional estimate or a
simultaneous bad-cluster repair.
\end{proposition}

\begin{proof}
The coefficients \(B_\psi\) and \(C_\psi\) in
\eqref{eq:e8v29-D}--\eqref{eq:e8v29-Q} enter the lower-order terms.  If they
are unbounded, the constants hidden in
\eqref{eq:e8v29-D-asymp}--\eqref{eq:e8v29-QQ-asymp} are unbounded, so the
point at which the leading terms dominate may diverge with \(\psi\).
\end{proof}

\begin{assessment}[Interpretation of the signs]
\label{ass:e8v29-signs}
For the declared star model, a positive cosmological owner controls very
large positive conformal amplitude, and a positive curvature-square owner
controls degeneracy.  The physical Einstein term defeats every polynomial
curvature-square owner on the positive tail when the cosmological
coefficient vanishes.  These are conditional amplitude-tail facts.  A
positive curvature-square owner still faces the frozen V61 elementary
reflection-positivity obstruction, and none of these tails controls the
bounded-amplitude high-frequency conformal escape.
\end{assessment}

\begin{definition}[Conformal fibre slope card, CFSC]
\label{def:e8v29-CFSC}
A gravitational regulator populates CFSC only if:
\begin{enumerate}[label=\textup{(C\arabic*)},leftmargin=4em]
\item its actual gauge-fixed or quotient action admits repair coordinates
with a declared reference Jacobian;
\item every outward amplitude tail has a uniform effective-potential slope
as in \eqref{eq:e8v29-slope-row}, including unrestricted neighbouring
fields or a proved cluster decomposition;
\item the intermediate collar has a uniform trace density and covers the
bounded-amplitude high-frequency sector with a separate coercive estimate;
\item the normal derivatives are dominated by the physical regulated
Dirichlet form;
\item the estimates persist on the complex metric circle required by HEFC;
and
\item the resulting carrier satisfies the selected conformal positivity
resolution branch from the frozen V61 analysis.
\end{enumerate}
\end{definition}

\begin{firewall}[The star model is diagnostic]
\label{fire:e8v29-model}
Equation \eqref{eq:e8v29-action} has not been derived from a
reflection-positive, diffeomorphism-compatible lattice gravitational
measure.  Its tail theorem therefore identifies necessary sign and
uniformity issues for that model; it does not populate CFSC for the physical
carrier.
\end{firewall}

\begin{programme}[Eighth-series V30: companion audit and conformal route
selection]
\label{prog:e8v29-V30}
V30 will perform the compulsory audit of the newest coherent Yang--Mills
and Navier--Stokes notes.  It will compare their most recent good/bad,
cluster, and capacity mechanisms with MRIC, FRCC, and CFSC, then select the
next noncircular gravitational or Standard Model interface.  No numerical
constant will be imported without matching hypotheses.
\end{programme}

\begin{firewall}[Claim boundary after eighth-series V29]
\label{fire:e8v29-boundary}
V29 proves a Hardy estimate from a one-sided potential slope, computes the
effective conformal star derivative including the reference Jacobian,
classifies both amplitude tails, and isolates neighbour compactness and
bounded-amplitude frequency as separate rows.

V29 does not derive the star from a physical gravitational regulator, prove
CFSC, solve the conformal positivity problem, establish DQSRC or the full SM
contacts, remove the cutoff, or construct the Standard Model--Einstein
continuum.
\end{firewall}

\begin{theorem}[Eighth-series V29 result]
\label{thm:e8v29-result}
For the declared physical-sign conformal star, positive
\(\Lambda\) and positive \(\alpha\) give uniform conditional potential
slopes at respectively large positive and negative conformal amplitude,
provided the neighbours remain in a compact set.  With
\(\Lambda=0\), the negative Einstein term overwhelms every polynomial
curvature-square owner on the positive tail, and no fixed coordinate
Jacobian changes that conclusion.  These amplitude slopes neither control
unbounded neighbour clusters nor the archived bounded-amplitude,
high-frequency conformal instability.
\end{theorem}

\begin{proof}
Combine \Cref{lem:e8v29-slope-Hardy,
thm:e8v29-tail-classification,cor:e8v29-two-slopes,
cor:e8v29-Jacobian-no-cure,prop:e8v29-amplitude-vs-frequency,
prop:e8v29-neighbour-row}.
\end{proof}

\paragraph{Parent-version record.}
V29 was formed from
\texttt{Renewal\_SM\_Einstein\_Continuum\_Research\_Notes\_V28.tex}.
The V28 parent had (10\,463) logical source lines, (389\,657) bytes,
and SHA--256
\[
 \texttt{CD080393A74868AEF4EF5E28A440F366510FD638155F562B57D3D63F9317F1D7}.
\]
The next compulsory companion audit is V30.

% =============================================================================
% END APPENDED EIGHTH-SERIES V29 MATERIAL
% =============================================================================

% =============================================================================
% BEGIN APPENDED EIGHTH-SERIES V30 MATERIAL
% =============================================================================

\section{V30 companion audit: restricted smoothing must move geometry}
\label{sec:v8v30}

This is the compulsory five-version audit following V25.  The newest
coherent Yang--Mills note has moved from bad-event probability to a genuine
form-capacity inequality through a restricted physical transfer.  The
newest Navier--Stokes note remains on its certified rank-one kernel
refinement.  The audit changes the SM--Einstein direction: FRCC is retained
as a coordinate route, but a restricted-smoothing route is now added.
Before using it, one must distinguish a Markov transfer on the metric
carrier from the fixed-metric quantum semigroup studied in V17--V24.

\subsection{Exact audit record}
\label{sec:e8v30-audit}

\begin{audit}[Companion files read at V30]
\label{audit:e8v30-files}
The audit used:
\begin{enumerate}[label=\textup{(\roman*)},leftmargin=3.8em]
\item
\texttt{Renewal\_Yang\_Mills\_Mass\_Gap\_Research\_Notes\_V89.tex},
with \(41\,160\) logical source lines, \(1\,440\,324\) bytes, one live
terminator, and SHA--256
\[
 \texttt{3421D3CC92B63006BFF54915F19C66E21213674C237EB0B7BD09320C13D01F2E};
\]
\item
\texttt{Renewal\_NS\_Stability\_Research\_Notes\_V26.tex},
with \(5\,319\) logical source lines, \(278\,283\) bytes, one live
terminator, and SHA--256
\[
 \texttt{82C887F8C2C69E4F28FAD7C3DC8DE5B9099CB5A4BF431EBA1FFF3E91935978AD}.
\]
\end{enumerate}
Both files end in \(\texttt{\string\end\{document\}}\).
\end{audit}

\begin{theorem}[Yang--Mills input selected by the audit]
\label{thm:e8v30-YM-input}
The audited YM V89 proves the following finite-cutoff facts.
\begin{enumerate}[label=\textup{(\roman*)},leftmargin=3.8em]
\item If \(P\) is a reversible Markov contraction and \(B\) an event, then
the restricted norm
\[
 \eta_P(B)=\|M_BP\|_{2\to2}^2
\]
controls the bad-set \(L^2\) term by a variance term plus a restricted jump
energy.
\item For \(P=e^{-\delta H}\), the enlarged jump term gives the capacity
pair
\[
 \int_B|f|^2\,d\pi
 \leq2\eta_P(B)\|f\|_2^2
      +4\delta\,\mathcal E_P(f),
 \qquad
 \mathcal E_P(f)=\delta^{-1}\langle f,(I-P)f\rangle.
\]
\item For a positive half-slab amplitude, the singular ground function
cancels exactly from the restricted Doob Hilbert--Schmidt norm.  Under its
stated Wilson action and stable-rank hypotheses this gives an exponentially
small restricted smoothing coefficient.
\item For connected clusters, the remaining open rows are the norm of the
assembled positive operator sum and the incidence of the restricted jump
energies in the global physical Dirichlet form.
\end{enumerate}
The source-specific boundary owner, response-matched half-slab card, and
cluster assembly remain hypotheses in that programme.
\end{theorem}

\begin{proof}
These are respectively the statements labelled
\texttt{thm:s6v89-smoothing-capacity},
\texttt{lem:s6v89-Doob-HS},
\texttt{thm:s6v89-action-smoothing}, and
\texttt{thm:s6v89-CRSC} in the audited file, with its firewalls retained.
\end{proof}

\begin{theorem}[Navier--Stokes input selected by the audit]
\label{thm:e8v30-NS-input}
The audited NS V26 computes the first- and second-derivative Oseen kernel
norms on rank-one traceless inputs.  The first-derivative refinement gives
no improvement over the full symmetric-traceless norm; the
second-derivative pointwise improvement is confined to an intermediate
radial range and is expected to change the integrated constant only
slightly.  No probability, repair, restricted-transfer, or form-capacity
estimate for the gravitational carrier is supplied, and the full
Navier--Stokes regularity theorem remains unproved.
\end{theorem}

\begin{proof}
This is the content of the propositions labelled
\texttt{prop:c4v26-first} and \texttt{prop:c4v26-second}, together with the
V26 conclusion and firewall in the audited file.
\end{proof}

\begin{assessment}[Audit decision]
\label{ass:e8v30-decision}
The YM restricted-smoothing inequality is directly relevant to the V26
bad-geometry form debit and avoids a global Sobolev exponent.  It is not
equivalent to the V28 normal-fibre Hardy route: one uses a physical Markov
transfer and the other uses conditional repair coordinates.  The next
question is whether the SM--Einstein regulator possesses a
geometry-moving reversible transfer with a restricted norm and a
loader-energy incidence estimate.  The NS rank-one calculation does not
alter that acquisition order.
\end{assessment}

\subsection{Hilbert-valued restricted smoothing}
\label{sec:e8v30-vector}

Let \((\Omega,\mu)\) be a probability space, let \(P\) be a self-adjoint
Markov-kernel contraction on \(L^2(\mu)\), and let \(\mathcal K\) be a separable
Hilbert space.  The operator \(P\) acts componentwise on
\(L^2(\mu;\mathcal K)\).  For \(\delta>0\), put
\begin{equation}
 \mathcal E_{P,\mathcal K}(F)
 :=\delta^{-1}\langle F,(I-P)F\rangle_{L^2(\mu;\mathcal K)}.
 \label{eq:e8v30-vector-energy}
\end{equation}

\begin{theorem}[Hilbert-valued restricted-smoothing capacity]
\label{thm:e8v30-vector-capacity}
Let \(B\subset\Omega\) be measurable and
\(\eta=\|M_BP\|_{2\to2}^2\).  Then every
\(F\in L^2(\mu;\mathcal K)\) satisfies
\begin{equation}
 \boxed{
 \int_B\|F(\omega)\|_{\mathcal K}^2\,d\mu(\omega)
 \leq2\eta\|F\|_{L^2(\mu;\mathcal K)}^2
     +4\delta\mathcal E_{P,\mathcal K}(F).}
 \label{eq:e8v30-vector-capacity}
\end{equation}
\end{theorem}

\begin{proof}
The tensor-product operator
\((M_BP)\otimes I_{\mathcal K}\) has norm
\(\|M_BP\|_{2\to2}\).  Decompose
\[
 M_BF=M_BPF+M_B(I-P)F.
\]
The first square is at most
\(\eta\|F\|_2^2\).  If \(P(\omega,d\omega')\) is its Markov kernel,
Hilbert-space Jensen gives
\[
 \|(I-P)F(\omega)\|_{\mathcal K}^2
 \leq\int
 \|F(\omega)-F(\omega')\|_{\mathcal K}^2
 P(\omega,d\omega').
\]
After integration over \(B\), enlarge to \(\Omega\).  Reversibility gives
\[
 \int\!\!\int\|F(\omega)-F(\omega')\|_{\mathcal K}^2
 P(\omega,d\omega')d\mu(\omega)
 =2\langle F,(I-P)F\rangle.
\]
Finally use \(\|a+b\|^2\leq2\|a\|^2+2\|b\|^2\) and
\eqref{eq:e8v30-vector-energy}.
\end{proof}

\begin{corollary}[Loader incidence produces a BGFD pair]
\label{cor:e8v30-loader}
Suppose a local energy, transfer-defect, or metric-derivative loader
\(F_{n,A}\) obeys
\begin{equation}
 \mathcal E_{P_n,\mathcal K_n}(F_{n,A})
 \leq a_{n,A}+b_{n,A}\|F_{n,A}\|_2^2.
 \label{eq:e8v30-incidence}
\end{equation}
For a bad event \(B_n\), set
\begin{equation}
 \eta_{n,A}
 :=\|M_{B_n}P_n\|_{2\to2}^2.
 \label{eq:e8v30-eta}
\end{equation}
Then
\begin{equation}
 \int_{B_n}\|F_{n,A}\|^2\,d\mu_n
 \leq
 4\delta_na_{n,A}
 +\bigl(2\eta_{n,A}+4\delta_nb_{n,A}\bigr)
  \int_{\Omega_n}\|F_{n,A}\|^2\,d\mu_n.
 \label{eq:e8v30-BGFD-pair}
\end{equation}
Thus the V26 bad-set absorption applies whenever
\begin{equation}
 \sup_n\bigl(2\eta_{n,A}+4\delta_nb_{n,A}\bigr)<1
 \label{eq:e8v30-strict-window}
\end{equation}
and \(4\delta_na_{n,A}\), together with its bridge contacts, is summable.
\end{corollary}

\begin{proof}
Insert \eqref{eq:e8v30-incidence} into
\eqref{eq:e8v30-vector-capacity}.  The constant and relative terms are
exactly the two terms in \eqref{eq:e8v30-BGFD-pair}.
\end{proof}

\begin{remark}[Why Hilbert-valued loaders matter]
\label{rem:e8v30-loaders}
The form energy in V16--V24 is a norm of a source vector after applying an
energy or square-root loader.  It is not generally the square of one scalar
function of the metric.  The vector theorem permits this loader to be used
directly, provided the varying cutoff Hilbert fibres have a declared
measurable isometric trivialization or a compatible bundle transfer.
\end{remark}

\subsection{The semigroup must update the bad variable}
\label{sec:e8v30-move-geometry}

The HEFC semigroup has the direct-integral form
\begin{equation}
 (\mathsf T_tF)(g)=e^{-tH(g)}F(g).
 \label{eq:e8v30-fibrewise}
\end{equation}
It evolves quantum fields at a fixed metric \(g\); it does not move \(g\)
in the carrier.

\begin{proposition}[A fibrewise quantum semigroup cannot smooth a metric
event]
\label{prop:e8v30-fibrewise-no}
Let
\(\mathscr H=\int_\Omega^\oplus\mathcal K_g\,d\mu(g)\) and
\(\mathsf T=\int_\Omega^\oplus T(g)\,d\mu(g)\).  If
\(\|T(g)\|=1\) for almost every \(g\) in a positive-measure event \(B\),
then
\begin{equation}
 \|M_B\mathsf T\|_{\mathscr H\to\mathscr H}=1.
 \label{eq:e8v30-fibrewise-norm}
\end{equation}
In particular, after ground-energy normalization, the fixed-metric
semigroup \eqref{eq:e8v30-fibrewise} has no small restricted-smoothing
coefficient for a bad-geometry event.
\end{proposition}

\begin{proof}
The norm of a decomposable operator is the essential supremum of its fibre
norms.  Since \(M_B\mathsf T\) has fibre \(T(g)\) on \(B\) and zero fibre
off \(B\),
\[
 \|M_B\mathsf T\|
 =\operatorname*{ess\,sup}_{g\in B}\|T(g)\|=1.
\]
For \(T(g)=e^{-tH(g)}\) with \(H(g)\geq0\) and zero in its spectrum after
normalization, \(\|T(g)\|=1\).
\end{proof}

\begin{corollary}[HEFC and restricted carrier smoothing are distinct rows]
\label{cor:e8v30-distinct}
The holomorphic fixed-metric estimates of V20--V24 cannot populate
\(\eta_{n,A}\) in \eqref{eq:e8v30-eta}.  A successful restricted-smoothing
argument must use a Markov or transfer operator which changes the metric
history, or a joint block variable containing it.
\end{corollary}

\begin{proof}
For a nonnull bad event,
\Cref{prop:e8v30-fibrewise-no} gives \(\eta=1\), so even the term
\(2\eta\) in \eqref{eq:e8v30-strict-window} exceeds the absorption
threshold.
\end{proof}

\subsection{The smoothing-time window}
\label{sec:e8v30-time-window}

\begin{proposition}[Short time makes the restricted norm tend to one]
\label{prop:e8v30-short-time}
Let \(P_\delta\) be strongly continuous on \(L^2(\mu)\) with
\(P_0=I\), and let \(\mu(B)>0\).  Then
\begin{equation}
 \lim_{\delta\downarrow0}\|M_BP_\delta\|_{2\to2}=1.
 \label{eq:e8v30-short-time}
\end{equation}
\end{proposition}

\begin{proof}
Choose \(f\) supported in \(B\) with \(\|f\|_2=1\).  Strong continuity
gives \(P_\delta f\to f\), hence
\[
 \|M_BP_\delta\|
 \geq\|M_BP_\delta f\|_2\longrightarrow1.
\]
The reverse bound is one because both factors are contractions.
\end{proof}

\begin{assessment}[A genuine two-sided time constraint]
\label{ass:e8v30-time-window}
Taking \(\delta_n\) small reduces the energy term
\(4\delta_nb_{n,A}\), but by
\Cref{prop:e8v30-short-time} it drives the variance coefficient
\(2\eta_{n,A}\) toward two for every nonnull bad event.  Taking
\(\delta_n\) large may smooth the event but can make the energy debit too
large.  The physical carrier must provide an intermediate window satisfying
\eqref{eq:e8v30-strict-window}; neither rarity nor formal semigroup
existence proves that window.
\end{assessment}

\subsection{Restricted-smoothing form card}
\label{sec:e8v30-card}

\begin{definition}[Restricted-smoothing form card, RSFC]
\label{def:e8v30-RSFC}
For every local source and each loader needed for its base and first metric
derivative, RSFC consists of:
\begin{enumerate}[label=\textup{(R\arabic*)},leftmargin=4em]
\item a positive normalized carrier and a reversible local transfer
\(P_n\) that updates the metric or joint bad variable;
\item a compatible trivialization of the cutoff Hilbert bundle and a
well-defined Hilbert-valued loader;
\item a restricted smoothing estimate for the bad event, or an assembled
cluster operator estimate without a volume factor;
\item the loader-incidence bound \eqref{eq:e8v30-incidence} in the physical
carrier Dirichlet form;
\item a smoothing time satisfying the strict window
\eqref{eq:e8v30-strict-window};
\item summable constant terms, complex-circle contacts, and adjacent bridge
changes; and
\item compatibility with reflection, gauge/diffeomorphism Ward identities,
and the projective local algebra.
\end{enumerate}
\end{definition}

\begin{theorem}[RSFC populates the analytic bad rows of GGAFC]
\label{thm:e8v30-RSFC-GGAFC}
If RSFC holds with cofinal strict margins and summable remainders for the
energy, base-defect, and first-derivative loaders, then GGAFC rows
\textup{(G2)--(G5)} hold.  Combined with the good-chart HEFC rows, this gives
the V26 conditional \(C^1\) landing theorem.
\end{theorem}

\begin{proof}
Apply \Cref{cor:e8v30-loader} to every loader.  RSFC rows (R3)--(R5) give a
strict relative coefficient; row (R6) supplies the summable non-form
remainder for both base and derivative defects.  These are precisely the
V26 absorption hypotheses.  The final compatibility row supplies the
remaining GGAFC incidence data.
\end{proof}

\begin{remark}[FRCC and RSFC are alternative local mechanisms]
\label{rem:e8v30-two-routes}
FRCC proves form smallness by a conditional normal-fibre Hardy inequality.
RSFC proves it by physical transfer smoothing.  A regulator may use either
route on a given bad sector, or combine them, but their constants and
Dirichlet forms cannot be interchanged without an identification theorem.
\end{remark}

\begin{firewall}[The geometry-moving transfer is open]
\label{fire:e8v30-open}
The current SM--Einstein notes have a fixed-metric quantum semigroup and
abstract carrier channels, but no constructed positive reversible
geometry-moving transfer with the restricted norm, loader incidence,
cluster assembly, and reflection properties required by RSFC.  The chiral
determinant phase and the unresolved gravitational positivity problem also
prevent importing the positive-amplitude Yang--Mills Doob construction
unchanged.
\end{firewall}

\begin{assessment}[Progress at V30]
\label{ass:e8v30-status}
The audit replaces an exclusively coordinate-based capacity programme by
two precise alternatives.  The new RSFC compiler handles the
Hilbert-valued energy and response loaders used by the SM programme and
identifies a strict smoothing-time window.  It also closes a possible
logical loophole: HEFC's fixed-geometry heat semigroup cannot be reused as
the carrier smoother.  The next upstream target is the simplest genuine
geometry-updating kernel, a conditional heat-bath step, before attempting a
joint chiral gravitational transfer.
\end{assessment}

\begin{programme}[Eighth-series V31: exact heat-bath restricted norm]
\label{prog:e8v30-V31}
V31 will compute \(\|M_BP\|\) for a one-cell conditional heat-bath
projection.  It will identify the exact worst-exterior conditional bad
probability, derive the corresponding capacity pair and incidence count,
and test whether a good-exterior/bad-exterior split can replace a uniform
worst-history bound.  This will determine whether the YM transfer route can
be made local for metric amplitudes without assuming a global Sobolev
inequality.
\end{programme}

\begin{firewall}[Claim boundary after eighth-series V30]
\label{fire:e8v30-boundary}
V30 performs the compulsory companion audit; imports the proved YM V89
restricted-smoothing mechanism with its open assembly rows; proves its
Hilbert-valued extension and BGFD loader compiler; proves that a
fixed-geometry quantum semigroup cannot smooth a metric event; and
identifies the two-sided smoothing-time window.

V30 does not construct RSFC, a positive geometry-moving SM--Einstein
transfer, a chiral ground transform, a cluster assembly, DQSRC, complete SM
contacts, cutoff removal, or the Standard Model--Einstein continuum.
\end{firewall}

\begin{theorem}[Eighth-series V30 result]
\label{thm:e8v30-result}
A reversible geometry-moving transfer with restricted norm \(\eta_n\) and
loader incidence constants \(a_{n,A},b_{n,A}\) gives the exact BGFD pair
\[
 \varepsilon_{n,A}=4\delta_na_{n,A},
 \qquad
 \kappa_{n,A}=2\eta_n+4\delta_nb_{n,A}.
\]
The pair is useful only in the strict window \(\kappa_{n,A}<1\).
Fixed-metric HEFC semigroups have restricted norm one on every nonnull
bad-geometry event and cannot supply this row.  RSFC is therefore a
separate carrier-level acquisition card.
\end{theorem}

\begin{proof}
Use \Cref{thm:e8v30-vector-capacity,cor:e8v30-loader,
prop:e8v30-fibrewise-no,prop:e8v30-short-time,
thm:e8v30-RSFC-GGAFC}.
\end{proof}

\paragraph{Parent-version record.}
V30 was formed from
\texttt{Renewal\_SM\_Einstein\_Continuum\_Research\_Notes\_V29.tex}.
The V29 parent had (10\,858) logical source lines, (404\,421) bytes,
and SHA--256
\[
 \texttt{CD90A29DF499AC632F47C8059BE74DB3809E002C284CAD14F4D910D238304377}.
\]
The next compulsory companion audit is V35.

% =============================================================================
% END APPENDED EIGHTH-SERIES V30 MATERIAL
% =============================================================================

% =============================================================================
% BEGIN APPENDED EIGHTH-SERIES V31 MATERIAL
% =============================================================================

\section{Exact heat-bath smoothing of a bad metric fibre}
\label{sec:v8v31}

V30 showed that the capacity transfer must move the metric variable.  The
simplest such operator is exact conditional resampling of one metric cell.
This note computes its restricted norm, including the continuous-time
interpolation, and derives the optimal elementary absorption window.  The
calculation also locates the worst-exterior obstruction: the restricted norm
is governed by the essential supremum of the conditional bad probability,
not by its global average.

\subsection{Conditional resampling}
\label{sec:e8v31-heat-bath}

Let a probability space disintegrate as
\begin{equation}
 d\mu(y,x)=d\eta(y)\,d\mu_y(x),
 \label{eq:e8v31-disintegration}
\end{equation}
where \(x\) is one cell or block variable and \(y\) denotes its exterior.
The heat-bath projection is
\begin{equation}
 (Qf)(y,x)=\int f(y,z)\,d\mu_y(z).
 \label{eq:e8v31-Q}
\end{equation}
It is the orthogonal projection onto the exterior-measurable subspace.
Let \(B\) be a bad event with fibre section \(B_y\), and put
\begin{equation}
 p(y)=\mu_y(B_y),
 \qquad
 p_*=\operatorname*{ess\,sup}_{y}p(y).
 \label{eq:e8v31-pstar}
\end{equation}

\begin{theorem}[Exact restricted norm of a heat-bath projection]
\label{thm:e8v31-Q-norm}
If \(p_*>0\), then
\begin{equation}
 \boxed{\|M_BQ\|_{2\to2}^2=p_*.}
 \label{eq:e8v31-Q-norm}
\end{equation}
If \(p_*=0\), the norm is zero.
\end{theorem}

\begin{proof}
Write \(\bar f(y)=\int f(y,z)d\mu_y(z)\).  Conditional Jensen gives
\[
 \|f\|_2^2
 \geq\int|\bar f(y)|^2\,d\eta(y),
\]
whereas
\[
 \|M_BQf\|_2^2
 =\int p(y)|\bar f(y)|^2\,d\eta(y)
 \leq p_*\|f\|_2^2.
\]
For the reverse inequality, take \(f(y,x)=h(y)\), with \(h\) supported
where \(p(y)>p_*-\varepsilon\).  Then \(Qf=f\), and the norm quotient is at
least \(p_*-\varepsilon\).  Let \(\varepsilon\downarrow0\).  When \(p_*=0\),
\(B\) is null by disintegration.
\end{proof}

\begin{corollary}[The global bad probability is insufficient]
\label{cor:e8v31-average-no}
The averaged mass
\begin{equation}
 \mu(B)=\int p(y)\,d\eta(y)
 \label{eq:e8v31-average}
\end{equation}
does not control \(\|M_BQ\|\).  For every \(0<\varepsilon<1\), there are
events with \(\mu(B)=\varepsilon\) and
\(\|M_BQ\|=1\).
\end{corollary}

\begin{proof}
Take an exterior event \(E\) with \(\eta(E)=\varepsilon\), set
\(B_y\) equal to the full \(x\)-fibre on \(E\), and set it empty off \(E\).
Then \(p=1_E\), so \(\mu(B)=\varepsilon\) while \(p_*=1\).  Apply
\Cref{thm:e8v31-Q-norm}.
\end{proof}

\subsection{Continuous-time heat-bath interpolation}
\label{sec:e8v31-continuous}

For \(t\geq0\), let
\begin{equation}
 P_t=e^{-t(I-Q)}
 =Q+e^{-t}(I-Q).
 \label{eq:e8v31-Pt}
\end{equation}
This is the reversible Markov semigroup which refreshes the block at rate
one.

\begin{theorem}[Exact finite-time restricted norm]
\label{thm:e8v31-Pt-norm}
For a fibre with \(0<p(y)<1\),
\begin{equation}
 \|M_{B_y}P_t\|_{L^2(\mu_y)\to L^2(\mu_y)}^2
 =p(y)+(1-p(y))e^{-2t}.
 \label{eq:e8v31-fibre-Pt}
\end{equation}
Consequently,
\begin{equation}
 \boxed{
 \|M_BP_t\|_{2\to2}^2
 \leq p_*+(1-p_*)e^{-2t},}
 \label{eq:e8v31-Pt-upper}
\end{equation}
with equality whenever nontrivial fibre probabilities approach \(p_*\) on
sets of positive exterior measure.  The right side tends to \(p_*\) as
\(t\to\infty\) and to one as \(t\downarrow0\).
\end{theorem}

\begin{proof}
Fix \(y\), write \(p=p(y)\), \(c=e^{-t}\), and use the orthonormal basis
\[
 u={1_{B_y}\over\sqrt p},
 \qquad
 v={1_{B_y^c}\over\sqrt{1-p}}
\]
of the two-dimensional subspace of functions constant on \(B_y\) and its
complement.  On that subspace,
\[
 Q=
 \begin{pmatrix}
 p&\sqrt{p(1-p)}\\
 \sqrt{p(1-p)}&1-p
 \end{pmatrix},
 \qquad
 P_t=cI+(1-c)Q.
\]
The operator \(M_{B_y}P_t\) retains the first row.  Its squared row norm is
\begin{align*}
 \bigl(c+(1-c)p\bigr)^2
 +(1-c)^2p(1-p)
 &=c^2+p(1-c^2)\\
 &=p+(1-p)e^{-2t}.
\end{align*}
On the orthogonal complement, \(Q=0\).  A zero-mean fluctuation supported
inside \(B_y\) therefore has squared norm ratio \(c^2\), no larger than the
displayed value; fluctuations supported off \(B_y\) are killed by
\(M_{B_y}\).  This proves the fibre formula.

The full operator is decomposable over \(y\), so its norm is the essential
supremum of the fibre norms.  The displayed function is increasing in
\(p\), giving \eqref{eq:e8v31-Pt-upper} and the equality criterion.  The
limits are immediate.
\end{proof}

\begin{remark}[Null conditional fibres]
\label{rem:e8v31-null-fibres}
If \(p(y)=0\), the fibre restriction is zero rather than \(e^{-2t}\).
This does not affect \eqref{eq:e8v31-Pt-upper}.  The equality statement is
therefore phrased using nontrivial probabilities approaching \(p_*\).
\end{remark}

\subsection{The optimized heat-bath capacity window}
\label{sec:e8v31-window}

For a Hilbert-valued loader \(F\), let
\begin{equation}
 \operatorname{Var}_x(F)
 :=\langle F,(I-Q)F\rangle
 =\int\|F-QF\|_{\mathcal K}^2\,d\mu.
 \label{eq:e8v31-conditional-variance}
\end{equation}
The jump energy of \(P_t\) is
\begin{equation}
 \langle F,(I-P_t)F\rangle
 =(1-e^{-t})\operatorname{Var}_x(F).
 \label{eq:e8v31-Pt-energy}
\end{equation}

\begin{theorem}[Tunable heat-bath capacity inequality]
\label{thm:e8v31-tunable}
Let
\begin{equation}
 \eta_t=\|M_BP_t\|_{2\to2}^2.
 \label{eq:e8v31-etat}
\end{equation}
For every \(\theta>0\),
\begin{align}
 \int_B\|F\|_{\mathcal K}^2\,d\mu
 \leq{}&(1+\theta)\eta_t\|F\|_2^2
 \notag\\
 &+2(1+\theta^{-1})(1-e^{-t})
   \operatorname{Var}_x(F).
 \label{eq:e8v31-tunable}
\end{align}
\end{theorem}

\begin{proof}
Write
\[
 M_BF=M_BP_tF+M_B(I-P_t)F
\]
and use
\[
 \|a+b\|^2
 \leq(1+\theta)\|a\|^2+(1+\theta^{-1})\|b\|^2.
\]
The first term is bounded by \(\eta_t\|F\|_2^2\).  Hilbert-space Jensen for
the Markov kernel of \(P_t\), followed by enlargement from \(B\) to the
whole space and reversibility, bounds the second square by
\[
 2\langle F,(I-P_t)F\rangle.
\]
Use \eqref{eq:e8v31-Pt-energy}.
\end{proof}

\begin{corollary}[Optimal elementary absorption coefficient]
\label{cor:e8v31-optimal}
Suppose the loader incidence estimate is
\begin{equation}
 \operatorname{Var}_x(F)
 \leq a+b\|F\|_2^2,
 \qquad a,b\geq0.
 \label{eq:e8v31-incidence}
\end{equation}
Put
\begin{equation}
 A=\eta_t,
 \qquad
 C=2(1-e^{-t})b.
 \label{eq:e8v31-AC}
\end{equation}
Then \eqref{eq:e8v31-tunable} gives the BGFD pair
\begin{align}
 \varepsilon_B(\theta)
 &=2(1+\theta^{-1})(1-e^{-t})a,
 \label{eq:e8v31-epsilon}\\
 \kappa_B(\theta)
 &=(1+\theta)A+(1+\theta^{-1})C.
 \label{eq:e8v31-kappa-theta}
\end{align}
If \(A,C>0\), its minimum over \(\theta>0\) is
\begin{equation}
 \boxed{
 \inf_{\theta>0}\kappa_B(\theta)
 =(\sqrt A+\sqrt C)^2,}
 \qquad
 \theta_{\rm opt}=\sqrt{C/A}.
 \label{eq:e8v31-kappa-opt}
\end{equation}
Thus a strict heat-bath absorption window exists exactly when
\begin{equation}
 \sqrt{\eta_t}
 +\sqrt{2(1-e^{-t})b}<1.
 \label{eq:e8v31-window}
\end{equation}
\end{corollary}

\begin{proof}
Insert \eqref{eq:e8v31-incidence} in
\eqref{eq:e8v31-tunable}.  The relative coefficient is
\[
 A+C+A\theta+{C\over\theta}.
\]
The arithmetic--geometric mean inequality gives
\(A\theta+C/\theta\geq2\sqrt{AC}\), with equality at the displayed
\(\theta\).  Hence the minimum is
\(A+C+2\sqrt{AC}=(\sqrt A+\sqrt C)^2\).
\end{proof}

\begin{corollary}[Exact-resampling limit]
\label{cor:e8v31-Q-window}
As \(t\to\infty\), the sufficient strict window becomes
\begin{equation}
 \sqrt{p_*}+\sqrt{2b}<1.
 \label{eq:e8v31-Q-window}
\end{equation}
Thus a small worst-exterior bad probability is still insufficient unless
the loader's conditional-variance incidence coefficient is also small.
\end{corollary}

\begin{proof}
Use \(\eta_t\to p_*\) and \(1-e^{-t}\to1\) in
\eqref{eq:e8v31-window}.
\end{proof}

\subsection{Good exteriors and the recursive bad-exterior debit}
\label{sec:e8v31-good-exterior}

Let \(G\) be an exterior-measurable event and assume
\begin{equation}
 p(y)\leq p_{\rm good}<1
 \quad\text{for \(y\in G\)}.
 \label{eq:e8v31-good-p}
\end{equation}

\begin{proposition}[Exact good-exterior restricted norm]
\label{prop:e8v31-good-norm}
For the exact heat-bath projection,
\begin{equation}
 \|M_{B\cap G}Q\|_{2\to2}^2
 =\operatorname*{ess\,sup}_{y\in G}p(y)
 \leq p_{\rm good}.
 \label{eq:e8v31-good-norm}
\end{equation}
Moreover,
\begin{equation}
 1_B\leq1_{B\cap G}+1_{G^c},
 \label{eq:e8v31-split}
\end{equation}
so the uncontrolled part is precisely a bad-exterior form debit.
\end{proposition}

\begin{proof}
Since \(G\) is exterior measurable, the fibre sections of \(B\cap G\) are
\(B_y\) on \(G\) and empty off \(G\).  Apply
\Cref{thm:e8v31-Q-norm}.  The indicator inequality is pointwise.
\end{proof}

\begin{theorem}[Good-exterior heat-bath compiler]
\label{thm:e8v31-good-compiler}
Suppose the good-exterior loader incidence satisfies
\[
 \operatorname{Var}_x(F)
 \leq a_G+b_G\|F\|_2^2,
\]
and choose \(t,\theta\) so that the coefficient obtained from
\eqref{eq:e8v31-kappa-theta}, with
\[
 \eta_t\leq
 p_{\rm good}+(1-p_{\rm good})e^{-2t},
\]
is strictly below one.  If \(G^c\) has an independent strict form debit
with a summable remainder, then \(B\) has a strict form debit after the two
relative coefficients and remainders are added.  The conclusion is useful
only when their sum remains below one.
\end{theorem}

\begin{proof}
Apply \Cref{thm:e8v31-tunable} to \(B\cap G\), use
\Cref{prop:e8v31-good-norm} and the finite-time calculation, and then use
\eqref{eq:e8v31-split}.  Add the assumed debit for \(G^c\).  Absorption
requires the sum of relative coefficients to be strictly less than one;
the constant terms form the summed remainder.
\end{proof}

\begin{firewall}[The recursion must terminate]
\label{fire:e8v31-recursion}
Calling \(G^c\) rare does not prove its form debit.  The good-exterior
compiler closes only if the bad exterior is controlled by a different
capacity mechanism, by a cluster transfer with an assembled operator row,
or by a finite recursion whose final remainder is proved.  Reapplying the
same unproved conditional bound to \(G^c\) would be circular.
\end{firewall}

\subsection{Heat-bath repair card}
\label{sec:e8v31-card}

\begin{definition}[Heat-bath repair card, HBRC]
\label{def:e8v31-HBRC}
For each metric cell or block relevant to a local loader, HBRC consists of:
\begin{enumerate}[label=\textup{(H\arabic*)},leftmargin=4em]
\item a positive conditional carrier \(\mu_y\) and its exact heat-bath
projection on the metric or joint block variable;
\item a good-exterior conditional bad-probability bound, preferably derived
from the V27 effective repair cost including its Jacobian;
\item a loader conditional-variance incidence bound
\eqref{eq:e8v31-incidence};
\item a time and tuning parameter satisfying the strict window, after all
bad-exterior debits are added;
\item a bounded-overlap or assembled-operator theorem for the family of
updated cells and connected clusters;
\item complex-circle, metric-derivative, and adjacent-bridge versions with
summable remainders; and
\item compatibility with reflection, gauge/diffeomorphism incidence, and
the physical local algebra.
\end{enumerate}
\end{definition}

\begin{theorem}[HBRC implies the corresponding RSFC row]
\label{thm:e8v31-HBRC-RSFC}
If HBRC holds for every energy, base-defect, and first-derivative loader
with cofinal strict margins and summable remainders, then the
restricted-smoothing form card of V30 holds for those loaders.  Hence the
analytic bad rows of GGAFC and the conditional \(C^1\) landing theorem
follow.
\end{theorem}

\begin{proof}
Rows (H1)--(H4) and
\Cref{thm:e8v31-good-compiler} give the restricted norm, incidence, and
strict absorption rows of RSFC.  Row (H5) prevents a volume factor in their
assembly, while rows (H6)--(H7) give the remaining analytic and physical
compatibilities.  Apply the V30 RSFC-to-GGAFC theorem and then V26.
\end{proof}

\begin{assessment}[Progress at V31]
\label{ass:e8v31-status}
The restricted norm of a one-block heat bath is exactly the worst-exterior
conditional bad probability.  At finite time it is
\(p+(1-p)e^{-2t}\), and the sharp tuning of the elementary decomposition
gives the strict window \eqref{eq:e8v31-window}.  This supplies an explicit
bridge from V27 conditional repair costs to V30 restricted smoothing, but
only on good exteriors.  The next unresolved row is the loader incidence
and cluster assembly for many metric updates.
\end{assessment}

\begin{programme}[Eighth-series V32: heat-bath loader incidence]
\label{prog:e8v31-V32}
V32 will derive conditional-variance incidence for a local
Hilbert-valued loader from a bounded one-cell replacement difference.  It
will prove the Efron--Stein identity for heat-bath projections, retain the
metric dependence of the energy square root, and count the overlap of cell
updates with a fixed source collar.  If the number of relevant cells grows,
the note will isolate the exact weighted summability needed instead of
assuming a uniform finite support.
\end{programme}

\begin{firewall}[Claim boundary after eighth-series V31]
\label{fire:e8v31-boundary}
V31 computes the exact restricted norm for a conditional heat bath and its
finite-time semigroup; optimizes the resulting capacity coefficient; proves
the good-exterior split; and formulates HBRC.

V31 does not prove a conditional bad probability for the gravitational
carrier, loader incidence, cluster assembly, positivity of a joint chiral
heat bath, DQSRC, complete SM contacts, cutoff removal, or construction of
the Standard Model--Einstein continuum.
\end{firewall}

\begin{theorem}[Eighth-series V31 result]
\label{thm:e8v31-result}
For a one-block heat bath, restricted smoothing is quantified exactly by
\[
 \eta_t=p_*+(1-p_*)e^{-2t}
\]
when the worst conditional bad probability is attained through nontrivial
fibres.  With loader incidence coefficient \(b\), the optimal elementary
strict window is
\[
 \sqrt{\eta_t}+\sqrt{2(1-e^{-t})b}<1.
\]
An averaged bad probability cannot replace \(p_*\); a good-exterior split
is valid only when the complementary exterior has an independent
terminating form debit.
\end{theorem}

\begin{proof}
Use \Cref{thm:e8v31-Q-norm,cor:e8v31-average-no,
thm:e8v31-Pt-norm,cor:e8v31-optimal,
prop:e8v31-good-norm,thm:e8v31-good-compiler}.
\end{proof}

\paragraph{Parent-version record.}
V31 was formed from
\texttt{Renewal\_SM\_Einstein\_Continuum\_Research\_Notes\_V30.tex}.
The V30 parent had (11\,290) logical source lines, (421\,040) bytes,
and SHA--256
\[
 \texttt{A9DC38185C191B57A986811E4533F9579834AD7D01B7FB6A02FBD1FA332D81FD}.
\]
The next compulsory companion audit is V35.

% =============================================================================
% END APPENDED EIGHTH-SERIES V31 MATERIAL
% =============================================================================

% =============================================================================
% BEGIN APPENDED EIGHTH-SERIES V32 MATERIAL
% =============================================================================

\section{Heat-bath incidence for energy and response loaders}
\label{sec:v8v32}

V31 computed the restricted smoothing coefficient of a one-cell heat bath.
The other term in its strict window is the conditional variance of the
Hilbert-valued loader.  This note expresses that variance exactly as a
replacement square, then separates a local source replacement from the
unbounded energy-square-root contact.  The result identifies a weighted
influence row which remains meaningful when a source collar grows with the
regulator.

\subsection{Hilbert-valued conditional Efron--Stein identity}
\label{sec:e8v32-Efron-Stein}

Retain the disintegration
\[
 d\mu(y,x)=d\eta(y)d\mu_y(x)
\]
and heat-bath projection \(Q\) of V31.  Let \(x'\) be an independent draw
from \(\mu_y\), conditional on the same exterior \(y\).

\begin{theorem}[Conditional replacement identity]
\label{thm:e8v32-replacement}
For every \(F\in L^2(\mu;\mathcal K)\),
\begin{align}
 \operatorname{Var}_x(F)
 &:=\int\|F-QF\|_{\mathcal K}^2\,d\mu
 \label{eq:e8v32-variance}\\
 &= {1\over2}
 \int_Y\int\!\!\int
 \|F(y,x)-F(y,x')\|_{\mathcal K}^2
 d\mu_y(x)d\mu_y(x')d\eta(y).
 \label{eq:e8v32-replacement}
\end{align}
\end{theorem}

\begin{proof}
Fix \(y\) and write \(m_y=\int F(y,x)d\mu_y(x)\).  Expanding the squared
difference and using conditional independence gives
\begin{align*}
 {1\over2}\int\!\!\int\|F(x)-F(x')\|^2d\mu_y(x)d\mu_y(x')
 &=\int\|F(x)\|^2d\mu_y(x)-\|m_y\|^2\\
 &=\int\|F(x)-m_y\|^2d\mu_y(x).
\end{align*}
Integrate in \(y\).  Bochner integrability follows from \(F\in L^2\).
\end{proof}

\begin{corollary}[Relative replacement bound gives incidence]
\label{cor:e8v32-relative}
Suppose
\begin{equation}
 \|F(y,x)-F(y,x')\|^2
 \leq A(y,x,x')
 +\beta\bigl(\|F(y,x)\|^2+\|F(y,x')\|^2\bigr),
 \label{eq:e8v32-relative-replacement}
\end{equation}
where \(A\geq0\) is integrable.  Then
\begin{equation}
 \operatorname{Var}_x(F)
 \leq {1\over2}\mathbb E A+\beta\|F\|_2^2.
 \label{eq:e8v32-incidence}
\end{equation}
\end{corollary}

\begin{proof}
Insert \eqref{eq:e8v32-relative-replacement} in
\eqref{eq:e8v32-replacement}.  The two identically distributed loader
norms contribute \(2\|F\|_2^2\), which cancels the factor \(1/2\).
\end{proof}

\begin{remark}[The relative coefficient is the V31 datum]
\label{rem:e8v32-beta}
In the notation of V31, the incidence constants are
\[
 a={1\over2}\mathbb EA,\qquad b=\beta.
\]
Thus the exact-resampling strict window becomes
\(\sqrt{p_*}+\sqrt{2\beta}<1\), while the constant term contributes the
non-form remainder chosen by the tuning parameter.
\end{remark}

\subsection{Replacing an unbounded energy loader}
\label{sec:e8v32-energy-loader}

For a history \(g\), let \(\mathcal H_g\) be a Hilbert fibre,
\(H_g\geq0\) a self-adjoint Hamiltonian, and
\(\psi_g\in D(H_g^{1/2})\) a local source vector.  For a one-cell
replacement \(g\rightsquigarrow g'\), let
\[
 U_{g'\to g}:\mathcal H_{g'}\longrightarrow\mathcal H_g
\]
be a declared isometry.  The aligned energy-loader difference is
\begin{equation}
 \Delta F(g,g')
 :=H_g^{1/2}\psi_g
   -U_{g'\to g}H_{g'}^{1/2}\psi_{g'}.
 \label{eq:e8v32-DeltaF}
\end{equation}

\begin{lemma}[Exact source/contact split]
\label{lem:e8v32-source-contact}
On the common form domain,
\begin{align}
 \Delta F(g,g')
 ={}&H_g^{1/2}
 \bigl(\psi_g-U_{g'\to g}\psi_{g'}\bigr)
 \notag\\
 &+\bigl(H_g^{1/2}U_{g'\to g}
         -U_{g'\to g}H_{g'}^{1/2}\bigr)\psi_{g'}.
 \label{eq:e8v32-source-contact}
\end{align}
For every \(\theta>0\),
\begin{align}
 \|\Delta F(g,g')\|^2
 \leq{}&(1+\theta)
 \left\|H_g^{1/2}
 \bigl(\psi_g-U_{g'\to g}\psi_{g'}\bigr)\right\|^2
 \notag\\
 &+(1+\theta^{-1})
 \left\|
 \bigl(H_g^{1/2}U_{g'\to g}
       -U_{g'\to g}H_{g'}^{1/2}\bigr)\psi_{g'}
 \right\|^2.
 \label{eq:e8v32-split-bound}
\end{align}
\end{lemma}

\begin{proof}
Add and subtract
\(H_g^{1/2}U_{g'\to g}\psi_{g'}\) in
\eqref{eq:e8v32-DeltaF}.  The bound is
\(\|a+b\|^2\leq(1+\theta)\|a\|^2+
(1+\theta^{-1})\|b\|^2\).
\end{proof}

\begin{firewall}[The square-root contact is an unbounded-form row]
\label{fire:e8v32-square-root}
The second term in \eqref{eq:e8v32-source-contact} is not controlled by
\(\|H_g-H_{g'}\|\), which may be infinite or undefined for unbounded
Hamiltonians.  DQSRC supplies bounded congruence maps under its own
hypotheses, but a replacement estimate for this commutator still has to be
proved.  Treating \(H^{1/2}\) as a bounded loader would invalidate the
incidence argument.
\end{firewall}

\begin{definition}[One-cell energy replacement rows]
\label{def:e8v32-energy-rows}
For a chosen cell \(i\), the source and square-root contact rows are
\begin{align}
 &\left\|H_g^{1/2}
 \bigl(\psi_g-U_{g'\to g}\psi_{g'}\bigr)\right\|^2
 \leq
 A_i^{\rm src}(g,g')
 +\beta_i^{\rm src}
  \bigl(e(g)+e(g')\bigr),
 \label{eq:e8v32-source-row}\\
 &\left\|
 \bigl(H_g^{1/2}U_{g'\to g}
       -U_{g'\to g}H_{g'}^{1/2}\bigr)\psi_{g'}
 \right\|^2
 \leq
 A_i^{\rm sq}(g,g')
 +\beta_i^{\rm sq}
  \bigl(e(g)+e(g')\bigr),
 \label{eq:e8v32-square-root-row}\\
 &e(g):=\|H_g^{1/2}\psi_g\|^2.
 \label{eq:e8v32-e}
\end{align}
\end{definition}

\begin{theorem}[Energy-loader incidence compiler]
\label{thm:e8v32-energy-incidence}
If \eqref{eq:e8v32-source-row}--\eqref{eq:e8v32-square-root-row} hold and
\(F(g)=H_g^{1/2}\psi_g\) is aligned by the isometries, then
\begin{equation}
 \operatorname{Var}_i(F)
 \leq a_i(\theta)+b_i(\theta)\|F\|_2^2,
 \label{eq:e8v32-energy-incidence}
\end{equation}
where
\begin{align}
 a_i(\theta)
 &={1\over2}\mathbb E\left[
 (1+\theta)A_i^{\rm src}
 +(1+\theta^{-1})A_i^{\rm sq}\right],
 \label{eq:e8v32-ai}\\
 b_i(\theta)
 &=(1+\theta)\beta_i^{\rm src}
 +(1+\theta^{-1})\beta_i^{\rm sq}.
 \label{eq:e8v32-bi}
\end{align}
\end{theorem}

\begin{proof}
Insert the two replacement rows in
\eqref{eq:e8v32-split-bound}.  The result has the form
\eqref{eq:e8v32-relative-replacement}, with the displayed \(A\) and
\(\beta\).  Apply \Cref{cor:e8v32-relative}.
\end{proof}

\begin{corollary}[Best source/contact split]
\label{cor:e8v32-best-split}
If both relative coefficients are positive, then
\begin{equation}
 \inf_{\theta>0}b_i(\theta)
 =\left(\sqrt{\beta_i^{\rm src}}
        +\sqrt{\beta_i^{\rm sq}}\right)^2,
 \qquad
 \theta_{\rm opt}
 =\sqrt{\beta_i^{\rm sq}/\beta_i^{\rm src}}.
 \label{eq:e8v32-best-split}
\end{equation}
\end{corollary}

\begin{proof}
Expand
\[
 b_i=\beta_i^{\rm src}+\beta_i^{\rm sq}
 +\theta\beta_i^{\rm src}
 +\theta^{-1}\beta_i^{\rm sq}
\]
and minimize the last two terms by arithmetic--geometric mean.
\end{proof}

\subsection{Several updates and weighted influence}
\label{sec:e8v32-many}

Let \(Q_i\) be the conditional heat-bath projection for cell or block \(i\),
and let \(c_i>0\) be its update rate.  The reversible heat-bath generator
and its Hilbert-valued Dirichlet form are
\begin{align}
 \mathcal L_{\rm HB}
 &=\sum_i c_i(I-Q_i),
 \label{eq:e8v32-LHB}\\
 \mathcal E_{\rm HB}(F)
 &=\langle F,\mathcal L_{\rm HB}F\rangle
 =\sum_i c_i\operatorname{Var}_i(F).
 \label{eq:e8v32-EHB}
\end{align}

\begin{theorem}[Weighted loader-incidence sum]
\label{thm:e8v32-weighted-sum}
Suppose every one-cell replacement obeys
\[
 \operatorname{Var}_i(F)\leq a_i+b_i\|F\|_2^2
\]
and the sums below are finite.  Then
\begin{equation}
 \boxed{
 \mathcal E_{\rm HB}(F)
 \leq
 \sum_i c_i a_i
 +\left(\sum_i c_i b_i\right)\|F\|_2^2.}
 \label{eq:e8v32-weighted-sum}
\end{equation}
\end{theorem}

\begin{proof}
Multiply each incidence estimate by \(c_i\), sum, and use
\eqref{eq:e8v32-EHB}.
\end{proof}

\begin{corollary}[Finite collar count]
\label{cor:e8v32-finite-collar}
If \(F\) is independent of cells outside a set \(S\), then
\(\operatorname{Var}_i(F)=0\) for \(i\notin S\).  If
\(|S|\leq L\), \(c_i\leq c_*\), \(a_i\leq a_*\), and
\(b_i\leq b_*\) on \(S\), then
\begin{equation}
 \mathcal E_{\rm HB}(F)
 \leq Lc_*a_*+Lc_*b_*\|F\|_2^2.
 \label{eq:e8v32-finite-collar}
\end{equation}
\end{corollary}

\begin{proof}
Conditional resampling of a variable on which \(F\) does not depend fixes
\(F\).  The bound is then
\Cref{thm:e8v32-weighted-sum} with at most \(L\) nonzero terms.
\end{proof}

\begin{definition}[Cofinal loader influence row, CLIR]
\label{def:e8v32-CLIR}
For a regulator sequence and local source \(A\), CLIR is the pair of bounds
\begin{align}
 \alpha_{n,A}
 &:=\sum_i c_{n,i}a_{n,A,i},
 \label{eq:e8v32-alpha}\\
 \beta_{n,A}
 &:=\sum_i c_{n,i}b_{n,A,i},
 \label{eq:e8v32-beta}
\end{align}
with:
\begin{enumerate}[label=\textup{(\roman*)},leftmargin=3.8em]
\item \(\beta_{n,A}\) small enough for the V31 restricted-smoothing window;
\item the tuned constant contribution built from \(\alpha_{n,A}\) summable
after adjacent differences and contact remainders are included;
\item the same rows for the base, first metric derivative, and every
complex-circle loader used by HEFC; and
\item tails outside a fixed source collar summable uniformly in \(n\) when
the loader is quasilocal rather than exactly local.
\end{enumerate}
\end{definition}

\begin{proposition}[A growing collar need not cause a volume loss]
\label{prop:e8v32-growing-collar}
Let \(d(i,S_A)\) be a cell distance from the source support.  Suppose
\begin{equation}
 c_{n,i}a_{n,A,i}\leq C_A\varepsilon_n e^{-m d(i,S_A)},
 \qquad
 c_{n,i}b_{n,A,i}\leq C_A\varepsilon_n e^{-m d(i,S_A)}
 \label{eq:e8v32-decay}
\end{equation}
and the cell graph has volume growth
\[
 \#\{i:d(i,S_A)=r\}\leq C_A'(1+r)^D.
\]
Then
\begin{equation}
 \alpha_{n,A}+\beta_{n,A}
 \leq C_A C_A'\varepsilon_n
 \sum_{r=0}^\infty(1+r)^D e^{-mr}
 =O_A(\varepsilon_n).
 \label{eq:e8v32-decay-sum}
\end{equation}
In particular, a growing regulator volume does not enter the incidence
constant.
\end{proposition}

\begin{proof}
Group the cell sum by distance shells and use
\eqref{eq:e8v32-decay}.  The exponential dominates the polynomial, so the
displayed series converges.
\end{proof}

\begin{proposition}[Uniform per-cell influence does cause a volume loss]
\label{prop:e8v32-volume-loss}
If a loader depends on \(N_n\) cells and
\(c_{n,i}b_{n,A,i}\geq b_0>0\) on all of them, then
\begin{equation}
 \beta_{n,A}\geq b_0N_n.
 \label{eq:e8v32-volume-loss}
\end{equation}
Hence CLIR fails when \(N_n\to\infty\).
\end{proposition}

\begin{proof}
Sum the lower bound in \eqref{eq:e8v32-beta}.
\end{proof}

\subsection{Response loaders and the noncommuting contact ledger}
\label{sec:e8v32-response}

For \(k=0,1\), let \(F_{n,A}^{(k)}\) denote respectively a base transfer
loader and its first metric derivative after the V20 complex-circle
representation.  Each replacement difference has three logically distinct
sources:
\begin{enumerate}[label=\textup{(\roman*)},leftmargin=3.8em]
\item replacement of the source vector and local observable;
\item replacement of the form square root or semigroup;
\item replacement of the metric derivative insertion, measure, and bridge.
\end{enumerate}

\begin{definition}[Response replacement card, RRC]
\label{def:e8v32-RRC}
RRC requires a relative replacement bound of the form
\begin{equation}
 \|\Delta_iF_{n,A}^{(k)}\|^2
 \leq A_{n,A,i}^{(k)}
 +b_{n,A,i}^{(k)}
 \left(\|F_{n,A}^{(k)}(g)\|^2+
       \|F_{n,A}^{(k)}(g')\|^2\right),
 \label{eq:e8v32-RRC}
\end{equation}
for \(k=0,1\), on one fixed complex parameter circle, with:
\begin{enumerate}[label=\textup{(R\arabic*)},leftmargin=4em]
\item each of the three sources above estimated before their norms are
combined;
\item the square-root term proved in common form coordinates, not by an
operator-norm difference of unbounded generators;
\item the weighted sums of constant and relative coefficients satisfying
CLIR;
\item adjacent bridge changes summable; and
\item all isometries compatible with the local algebra and reflection
incidence.
\end{enumerate}
\end{definition}

\begin{theorem}[RRC plus conditional activity populates HBRC incidence]
\label{thm:e8v32-RRC-HBRC}
Suppose RRC holds for the energy, base, and first-derivative loaders.  Then
the conditional heat-bath incidence rows of HBRC hold with
\[
 a_{n,A}^{(k)}
 ={1\over2}\sum_i c_{n,i}\mathbb E A_{n,A,i}^{(k)},
 \qquad
 b_{n,A}^{(k)}
 =\sum_i c_{n,i}b_{n,A,i}^{(k)}.
\]
If the V31 conditional bad-activity and strict-window rows also hold, with
the tuned constant remainders summable, then the corresponding RSFC and
GGAFC rows follow.
\end{theorem}

\begin{proof}
Apply \Cref{cor:e8v32-relative} cell by cell and then
\Cref{thm:e8v32-weighted-sum}.  RRC row (R3) makes the sums finite and
small.  Insert these constants in the V31 heat-bath compiler, then apply
the V31 HBRC-to-RSFC and V30 RSFC-to-GGAFC theorems.
\end{proof}

\begin{firewall}[RRC is not yet populated]
\label{fire:e8v32-open}
The existing HEFC and DQSRC results do not prove
\eqref{eq:e8v32-square-root-row} for an arbitrary one-cell metric
replacement, nor do they prove exponential cell-distance decay for the
metric derivative and measure contacts.  Long-range constraints can also
replace exponential influence by multipole tails.  These are acquisition
rows, not consequences of formal locality.
\end{firewall}

\begin{assessment}[Progress at V32]
\label{ass:e8v32-status}
The V31 incidence coefficient is now reduced to replacement estimates with
an exact factor \(1/2\).  For energy loaders the source replacement and the
unbounded square-root contact are separated before using a norm inequality.
Many-cell incidence is governed by a weighted influence sum; exponential
quasilocality with polynomial graph growth avoids a volume factor, while a
uniform per-cell debit does not.  The next upstream target is the
square-root replacement row on a common sectorial form domain.
\end{assessment}

\begin{programme}[Eighth-series V33: square-root replacement from resolvents]
\label{prog:e8v32-V33}
V33 will derive a fractional-power replacement formula from the Balakrishnan
resolvent representation for two positive operators on a common form
domain.  It will state the relative resolvent defect needed to bound
\[
 H_g^{1/2}U-UH_{g'}^{1/2}
\]
on a local source, track the low- and high-spectral integrals separately,
and determine whether the endpoint again requires a fractional, collared,
or cancellation branch.
\end{programme}

\begin{firewall}[Claim boundary after eighth-series V32]
\label{fire:e8v32-boundary}
V32 proves the Hilbert-valued conditional replacement identity, the
relative replacement incidence compiler, the exact energy-loader
source/contact split, the weighted influence theorem, and the conditional
RRC-to-HBRC handoff.

V32 does not prove the square-root replacement row, CLIR for the physical
loaders, conditional gravitational activity, cluster assembly, a positive
joint chiral heat bath, complete SM contacts, cutoff removal, or the
Standard Model--Einstein continuum.
\end{firewall}

\begin{theorem}[Eighth-series V32 result]
\label{thm:e8v32-result}
The conditional variance of a Hilbert-valued loader is one half of its
independent replacement square.  A relative replacement coefficient
\(\beta_i\) is therefore exactly the one-cell incidence coefficient in the
V31 heat-bath window.  For the energy loader, that coefficient splits into
source replacement and square-root contact contributions, optimally as
\[
 \left(\sqrt{\beta_i^{\rm src}}+
       \sqrt{\beta_i^{\rm sq}}\right)^2.
\]
Across many updates the correct cofinal datum is the weighted influence sum
\(\sum_i c_i b_i\), which remains local under a summable distance tail but
grows with volume under a uniform per-cell debit.
\end{theorem}

\begin{proof}
Use \Cref{thm:e8v32-replacement,cor:e8v32-relative,
thm:e8v32-energy-incidence,cor:e8v32-best-split,
thm:e8v32-weighted-sum,prop:e8v32-growing-collar,
prop:e8v32-volume-loss}.
\end{proof}

\paragraph{Parent-version record.}
V32 was formed from
\texttt{Renewal\_SM\_Einstein\_Continuum\_Research\_Notes\_V31.tex}.
The V31 parent had (11\,755) logical source lines, (436\,033) bytes,
and SHA--256
\[
 \texttt{800DC15C6C20610302A0974022E2A6BE56FEF4794BD33A924289C533F02ACEE2}.
\]
The next compulsory companion audit is V35.

% =============================================================================
% END APPENDED EIGHTH-SERIES V32 MATERIAL
% =============================================================================

% =============================================================================
% BEGIN APPENDED EIGHTH-SERIES V33 MATERIAL
% =============================================================================

\section{Square-root replacement: resolvent and polar-form branches}
\label{sec:v8v33}

V32 isolated the contact
\[
 H_g^{1/2}U-UH_{g'}^{1/2}.
\]
For unbounded Hamiltonians, estimating it by a bounded operator difference
is generally impossible.  This note first records the exact resolvent
formula for the special bounded-perturbation branch.  It then proves a
form-level alternative: two-sided form comparison gives a polar isometry
between the energy fibres, and the square-root replacement is small
relative to the source energy.  The price is a new compatibility row,
because independently chosen pairwise polar maps need not define one global
Hilbert-bundle trivialization.

\subsection{The bounded-perturbation resolvent branch}
\label{sec:e8v33-resolvent}

\begin{theorem}[Balakrishnan square-root difference]
\label{thm:e8v33-Balakrishnan}
Let \(A,B\) be positive self-adjoint operators on one Hilbert space with
\(A,B\geq\lambda I\), \(\lambda>0\).  Suppose \(D=A-B\) extends to a
bounded operator.  Then
\begin{equation}
 A^{1/2}-B^{1/2}
 ={1\over\pi}\int_0^\infty
 t^{1/2}(A+t)^{-1}D(B+t)^{-1}\,dt
 \label{eq:e8v33-Balakrishnan}
\end{equation}
as a norm-convergent Bochner integral, and
\begin{equation}
 \boxed{
 \|A^{1/2}-B^{1/2}\|
 \leq{\|A-B\|\over2\sqrt\lambda}.}
 \label{eq:e8v33-Lipschitz}
\end{equation}
\end{theorem}

\begin{proof}
The Balakrishnan formula is
\[
 A^{1/2}
 ={1\over\pi}\int_0^\infty
 t^{-1/2}A(A+t)^{-1}\,dt
\]
in the strong sense on \(D(A^{1/2})\), and similarly for \(B\).  Since
\[
 A(A+t)^{-1}-B(B+t)^{-1}
 =t(A+t)^{-1}(A-B)(B+t)^{-1},
\]
subtraction gives \eqref{eq:e8v33-Balakrishnan}.  The lower spectral bound
implies
\[
 \|t^{1/2}(A+t)^{-1}D(B+t)^{-1}\|
 \leq {t^{1/2}\|D\|\over(t+\lambda)^2}.
\]
This is integrable at zero and infinity, so the difference formula
converges in norm.  Finally,
\[
 {1\over\pi}\int_0^\infty
 {t^{1/2}\over(t+\lambda)^2}\,dt
 ={1\over2\sqrt\lambda},
\]
by the substitution \(t=\lambda s\) and
\(\mathrm B(3/2,1/2)=\pi/2\).
\end{proof}

\begin{corollary}[A finite-cutoff replacement row]
\label{cor:e8v33-bounded-row}
If a regulated pair satisfies \(A,B\geq\lambda I\) and
\(\|A-B\|\leq d_i\), then for every source vector \(\psi\),
\begin{equation}
 \|(A^{1/2}-B^{1/2})\psi\|^2
 \leq {d_i^2\over4\lambda}\|\psi\|^2.
 \label{eq:e8v33-bounded-row}
\end{equation}
This contributes only a non-form replacement remainder when
\(\|\psi\|^2\) is already uniformly controlled.
\end{corollary}

\begin{proof}
Apply \eqref{eq:e8v33-Lipschitz} to \(\psi\) and square.
\end{proof}

\begin{firewall}[Why this branch is usually not cofinal]
\label{fire:e8v33-bounded}
For differential or lattice Hamiltonians whose ultraviolet norm grows with
the cutoff, a metric-coefficient replacement changes the principal part.
Then \(A-B\) need not be bounded uniformly, even though the quadratic forms
are uniformly comparable.  The bounded branch is valid when its displayed
hypothesis holds; finite dimensionality alone does not make its constant
cofinal.
\end{firewall}

\subsection{Relative forms and polar energy alignment}
\label{sec:e8v33-polar}

Let \(A,B\geq\lambda I\), and suppose their closed forms have the same
domain.  Assume that for constants \(0<m\leq M<\infty\),
\begin{equation}
 m\,\langle u,Bu\rangle
 \leq\langle u,Au\rangle
 \leq M\,\langle u,Bu\rangle
 \quad\text{for every form vector \(u\)}.
 \label{eq:e8v33-form-equivalence}
\end{equation}

\begin{theorem}[Polar energy alignment]
\label{thm:e8v33-polar}
The closure of
\begin{equation}
 C=A^{1/2}B^{-1/2}
 \label{eq:e8v33-C}
\end{equation}
is bounded and boundedly invertible.  In its polar decomposition
\begin{equation}
 C=V|C|,
 \label{eq:e8v33-polar}
\end{equation}
\(V\) is unitary and
\begin{equation}
 \boxed{
 \|C-V\|
 \leq\rho(m,M)
 :=\max\{\,|1-\sqrt m|,\ |1-\sqrt M|\,\}.}
 \label{eq:e8v33-rho}
\end{equation}
Consequently, for every \(\psi\in D(B^{1/2})\),
\begin{equation}
 \boxed{
 \|A^{1/2}\psi-VB^{1/2}\psi\|
 \leq\rho(m,M)\|B^{1/2}\psi\|.}
 \label{eq:e8v33-polar-contact}
\end{equation}
\end{theorem}

\begin{proof}
For \(\xi\) in the Hilbert space, put \(u=B^{-1/2}\xi\).  The form
comparison gives
\[
 m\|\xi\|^2
 \leq\|A^{1/2}B^{-1/2}\xi\|^2
 \leq M\|\xi\|^2.
\]
Thus \(C\) extends boundedly, is bounded below by \(\sqrt m\), and has
closed range.  The same form-domain equivalence with \(A,B\) interchanged
shows that the range is all of the Hilbert space, so the partial isometry in
the polar decomposition is unitary.

In form sense,
\[
 mI\leq C^*C\leq MI.
\]
Therefore the spectrum of \(|C|=(C^*C)^{1/2}\) lies in
\([\sqrt m,\sqrt M]\).  Functional calculus gives
\[
 \||C|-I\|\leq\rho(m,M).
\]
Since \(C-V=V(|C|-I)\), this proves \eqref{eq:e8v33-rho}.  Finally set
\(\xi=B^{1/2}\psi\).  Then
\[
 A^{1/2}\psi-VB^{1/2}\psi=(C-V)\xi,
\]
which proves \eqref{eq:e8v33-polar-contact}.
\end{proof}

\begin{corollary}[Small relative form defect]
\label{cor:e8v33-small-defect}
If \(0\leq\varepsilon<1\) and
\begin{equation}
 (1-\varepsilon)B\leq A\leq(1+\varepsilon)B
 \label{eq:e8v33-epsilon}
\end{equation}
in form sense, then
\begin{equation}
 \rho(1-\varepsilon,1+\varepsilon)
 =1-\sqrt{1-\varepsilon}
 ={\varepsilon\over1+\sqrt{1-\varepsilon}}.
 \label{eq:e8v33-rho-epsilon}
\end{equation}
In particular,
\begin{equation}
 \|A^{1/2}\psi-VB^{1/2}\psi\|^2
 \leq
 {\varepsilon^2\over
  (1+\sqrt{1-\varepsilon})^2}
 \|B^{1/2}\psi\|^2.
 \label{eq:e8v33-quadratic-small}
\end{equation}
\end{corollary}

\begin{proof}
The lower endpoint gives the larger deviation because
\[
 1-\sqrt{1-\varepsilon}
 -(\sqrt{1+\varepsilon}-1)
 =2-\sqrt{1-\varepsilon}-\sqrt{1+\varepsilon}\geq0,
\]
the last inequality following from concavity of the square root.  Rationalize
the first expression and apply
\Cref{thm:e8v33-polar}.
\end{proof}

\begin{remark}[Different Hilbert fibres]
\label{rem:e8v33-different-fibres}
If \(J:\mathcal H_{g'}\to\mathcal H_g\) is a baseline unitary, apply the
theorem to \(A=Q_g\) and
\(\widetilde B=JQ_{g'}J^*\).  The energy alignment from
\(\mathcal H_{g'}\) to \(\mathcal H_g\) is then \(VJ\), and
\[
 \|Q_g^{1/2}J\psi'-(VJ)Q_{g'}^{1/2}\psi'\|
 \leq\rho(m,M)\|Q_{g'}^{1/2}\psi'\|.
\]
Here \(Q_g=H_g+\lambda I\) is the shifted form operator used in the V17 and
V26 form norms.
\end{remark}

\subsection{Why form comparison does not give a bounded difference}
\label{sec:e8v33-obstruction}

\begin{proposition}[Relative form closeness with unbounded square-root
difference]
\label{prop:e8v33-unbounded}
For every \(\varepsilon>0\), there are positive self-adjoint
\(A,B\geq I\) with common form domain and
\[
 B\leq A\leq(1+\varepsilon)B,
\]
but \(A^{1/2}-B^{1/2}\) is unbounded.
\end{proposition}

\begin{proof}
On \(\ell^2(\mathbb N)\), let
\[
 Be_n=n^2e_n,\qquad A=(1+\varepsilon)B.
\]
The displayed form comparison is immediate and the form domains coincide.
However,
\[
 (A^{1/2}-B^{1/2})e_n
 =(\sqrt{1+\varepsilon}-1)n e_n,
\]
whose norm is unbounded in \(n\).
\end{proof}

\begin{corollary}[The correct endpoint is relative energy]
\label{cor:e8v33-relative-endpoint}
No argument using only relative form comparison can replace
\eqref{eq:e8v32-square-root-row} by a uniform operator-norm bound on the
square-root difference.  The polar estimate
\eqref{eq:e8v33-polar-contact}, which is relative to the source form
energy, has the correct ultraviolet scaling.
\end{corollary}

\begin{proof}
\Cref{prop:e8v33-unbounded} excludes the operator-norm conclusion, while
\Cref{thm:e8v33-polar} gives the relative-energy conclusion under the same
kind of form comparison.
\end{proof}

\subsection{Insertion in the V32 replacement compiler}
\label{sec:e8v33-insertion}

Let \(g,g'\) differ in one updated cell, let
\(J_{g'\to g}\) be a baseline source-space unitary, and assume
\begin{equation}
 m_iQ_{g'}\preceq
 J_{g'\to g}^*Q_gJ_{g'\to g}
 \preceq M_iQ_{g'}
 \label{eq:e8v33-cell-form}
\end{equation}
in form sense.  Let \(V_iJ_{g'\to g}\) be the polar energy alignment from
\Cref{rem:e8v33-different-fibres}.

\begin{theorem}[Polar-aligned energy replacement]
\label{thm:e8v33-energy-replacement}
For source vectors \(\psi_g,\psi_{g'}\),
\begin{align}
 &\|Q_g^{1/2}\psi_g
       -(V_iJ_{g'\to g})Q_{g'}^{1/2}\psi_{g'}\|^2
 \notag\\
 &\quad\leq
 (1+\theta)
 \|Q_g^{1/2}(\psi_g-J_{g'\to g}\psi_{g'})\|^2
 \notag\\
 &\qquad+
 (1+\theta^{-1})\rho(m_i,M_i)^2
 \|Q_{g'}^{1/2}\psi_{g'}\|^2
 \label{eq:e8v33-energy-replacement}
\end{align}
for every \(\theta>0\).
\end{theorem}

\begin{proof}
Add and subtract \(Q_g^{1/2}J_{g'\to g}\psi_{g'}\).  The first difference
is the source replacement.  The second is bounded by the different-fibre
polar estimate.  Apply the tunable two-term square inequality.
\end{proof}

\begin{corollary}[Explicit square-root incidence coefficient]
\label{cor:e8v33-incidence}
If the source replacement has relative coefficient
\(\beta_i^{\rm src}\) in the sense of V32, then the energy replacement has
relative coefficient
\begin{equation}
 b_i(\theta)
 =(1+\theta)\beta_i^{\rm src}
 +(1+\theta^{-1})\rho(m_i,M_i)^2.
 \label{eq:e8v33-bi}
\end{equation}
Its optimal value is
\begin{equation}
 \boxed{
 b_i^{\rm opt}
 =\left(\sqrt{\beta_i^{\rm src}}+
        \rho(m_i,M_i)\right)^2.}
 \label{eq:e8v33-bi-opt}
\end{equation}
\end{corollary}

\begin{proof}
Use \Cref{thm:e8v33-energy-replacement} in the V32 conditional
replacement identity.  The polar term is already proportional to one of
the two source energies and therefore fits the symmetric relative row.
Minimize in \(\theta\) as in the V32 source/contact split.
\end{proof}

\subsection{The polar-cocycle obstruction}
\label{sec:e8v33-cocycle}

\begin{proposition}[Pairwise polar maps need not be a trivialization]
\label{prop:e8v33-no-cocycle}
Let \(V(g,g')\) be the polar unitary obtained independently from every
pair of form operators.  Form comparison alone does not imply
\begin{equation}
 V(g_0,g_2)=V(g_0,g_1)V(g_1,g_2)
 \label{eq:e8v33-cocycle}
\end{equation}
around a three-history path.  Hence the pairwise estimate does not by
itself define a single Hilbert-valued random variable whose conditional
variance is given by the V32 replacement identity.
\end{proposition}

\begin{proof}
Take one finite-dimensional Hilbert space and three positive matrices
\(Q_0,Q_1,Q_2\) which do not commute.  The polar factor of
\(Q_i^{1/2}Q_j^{-1/2}\) is the unitary part of that particular product.
Polar decomposition is not multiplicative:
\[
 \operatorname{polar}(XY)
 \ne\operatorname{polar}(X)\operatorname{polar}(Y)
\]
in general when the positive factors in the two products do not commute.
For example, choose two noncommuting positive \(2\times2\) matrices \(X,Y\);
their individual polar factors are the identity, while \(XY\) is generally
not self-adjoint and has a nontrivial polar factor.  Realize
\(X=Q_0^{1/2}Q_1^{-1/2}\) and
\(Y=Q_1^{1/2}Q_2^{-1/2}\) after choosing the corresponding positive
matrices.  Thus no cocycle identity follows merely from pairwise positivity
and form equivalence.
\end{proof}

\begin{remark}[Scope of the example]
\label{rem:e8v33-cocycle-example}
The example establishes the logical nonimplication.  A particular local
form family may possess a canonical coefficient-space embedding or a flat
connection that does satisfy \eqref{eq:e8v33-cocycle}; that additional
structure is exactly what must be proved.
\end{remark}

\begin{definition}[Energy polar alignment card, EPAC]
\label{def:e8v33-EPAC}
EPAC for a regulator family consists of:
\begin{enumerate}[label=\textup{(E\arabic*)},leftmargin=4em]
\item shifted positive form operators \(Q_g\) on a common measurable field
of form domains;
\item cell-replacement form ratios \(m_i,M_i\) whose polar coefficients
\(\rho(m_i,M_i)^2\) satisfy the V32 weighted influence row;
\item a measurable global energy-fibre trivialization whose pairwise maps
obey the required cocycle, or a canonical coefficient-space embedding
which makes the replacement square literal;
\item source replacement estimates in that same trivialization;
\item complex-circle and first metric-derivative versions; and
\item compatibility with adjacent bridges, reflection, and the physical
local algebra.
\end{enumerate}
\end{definition}

\begin{theorem}[EPAC populates the V32 square-root replacement row]
\label{thm:e8v33-EPAC}
If EPAC holds, then the energy-loader square-root contact in V32 has
relative coefficient \(\rho(m_i,M_i)^2\), and the full one-cell incidence
coefficient is bounded by \eqref{eq:e8v33-bi-opt}.  If its weighted sum and
constant remainders satisfy CLIR, the V32 RRC-to-HBRC handoff applies.
\end{theorem}

\begin{proof}
EPAC row (E3) turns the pairwise aligned energy vectors into values of one
measurable Hilbert-valued loader.  Apply
\Cref{thm:e8v33-energy-replacement,cor:e8v33-incidence} cell by cell.
Rows (E2), (E4), and (E5) give CLIR for every required loader, while row
(E6) supplies the bridge and physical compatibility.  Invoke the V32
handoff.
\end{proof}

\begin{firewall}[EPAC is open beyond the coefficient-local kinetic branch]
\label{fire:e8v33-open}
Uniform ellipticity gives finite form ratios, but not necessarily ratios
close enough for the V31 strict window.  The current notes also do not
construct the global polar cocycle for the interacting SM Hamiltonian,
its chiral fermion sector, or the metric derivative loaders.  In the nested
bosonic kinetic sector a common gradient/coefficient representation is
available from V23--V24; extending it to all contacts is open.
\end{firewall}

\begin{assessment}[Progress at V33]
\label{ass:e8v33-status}
The square-root contact now has a rigorous form-level branch.  A relative
form change of size \(\varepsilon\) gives a quadratic energy-incidence
coefficient of order \(\varepsilon^2\) after polar alignment, even though
the square-root difference may be unbounded.  The remaining issue is
geometric rather than algebraic: the energy alignments must come from one
measurable cocycle, and the one-cell form ratios must be small enough after
their weighted influence sum is taken.
\end{assessment}

\begin{programme}[Eighth-series V34: coefficient-space realization]
\label{prog:e8v33-V34}
V34 will populate EPAC for the nested bosonic \(P_1\) kinetic form already
used in V23--V24.  It will embed every energy loader into a fixed direct sum
of cell-gradient spaces using the positive coefficient square root, prove
the literal cocycle and one-cell replacement bound, and compute the
weighted influence in terms of the source energy fraction in the replaced
cell.  Lower-order, gauge-link, fermionic, and interaction contacts will
remain separate.
\end{programme}

\begin{firewall}[Claim boundary after eighth-series V33]
\label{fire:e8v33-boundary}
V33 proves the norm-convergent Balakrishnan difference formula for bounded
perturbations, the polar energy-alignment theorem for relatively comparable
forms, the explicit relative incidence coefficient, and the pairwise
polar-cocycle obstruction.

V33 does not prove EPAC for the interacting Hamiltonian, small form ratios
for arbitrary heat-bath replacements, the chiral and interaction
replacement contacts, conditional gravitational activity, cluster
assembly, cutoff removal, or the Standard Model--Einstein continuum.
\end{firewall}

\begin{theorem}[Eighth-series V33 result]
\label{thm:e8v33-result}
Bounded operator perturbations give the square-root Lipschitz estimate
\(\|A^{1/2}-B^{1/2}\|\leq\|A-B\|/(2\sqrt\lambda)\), but relative form
comparison alone cannot make that difference bounded.  It instead gives a
polar-aligned energy estimate with coefficient
\[
 \rho(m,M)^2
 =\max\{|1-\sqrt m|,|1-\sqrt M|\}^2.
\]
Combined with a source replacement coefficient, the optimal one-cell
incidence is \eqref{eq:e8v33-bi-opt}.  Using this estimate in a conditional
variance requires the additional global trivialization or cocycle row of
EPAC.
\end{theorem}

\begin{proof}
Use \Cref{thm:e8v33-Balakrishnan,thm:e8v33-polar,
prop:e8v33-unbounded,thm:e8v33-energy-replacement,
cor:e8v33-incidence,prop:e8v33-no-cocycle,
thm:e8v33-EPAC}.
\end{proof}

\paragraph{Parent-version record.}
V33 was formed from
\texttt{Renewal\_SM\_Einstein\_Continuum\_Research\_Notes\_V32.tex}.
The V32 parent had (12\,254) logical source lines, (452\,399) bytes,
and SHA--256
\[
 \texttt{261A21B1E7EE2A6340EB6F6142F90BE69F283547B9203868A004E3E76CABB6ED}.
\]
The next compulsory companion audit is V35.

% =============================================================================
% END APPENDED EIGHTH-SERIES V33 MATERIAL
% =============================================================================

% =============================================================================
% BEGIN APPENDED EIGHTH-SERIES V34 MATERIAL
% =============================================================================

\section{A fixed coefficient-space loader for nested kinetic forms}
\label{sec:v8v34}

V33 showed that abstract polar alignment gives the correct relative
square-root scaling but may fail to define a cocycle.  For the real bosonic
kinetic form of V23--V24, no pairwise polar choice is needed.  The form has a
canonical factorization through cell gradients and positive coefficient
square roots.  All metric histories therefore load into one fixed
coefficient space, and one-cell replacements occupy orthogonal summands.
This yields a heat-bath incidence bound with no factor proportional to the
number of cells.

\subsection{Coefficient-space factorization}
\label{sec:e8v34-factorization}

Fix a finite mesh \(\mathcal T_n\) and its conforming space
\(V_n\subset H^1(M;\mathbb C^N)\).  For each real metric history \(g\), let
\(G_g(x)\) be a measurable Hermitian coefficient satisfying
\begin{equation}
 \nu_G I\preceq G_g(x)\preceq\Lambda_G I
 \quad\text{almost everywhere},
 \qquad 0<\nu_G\leq\Lambda_G<\infty.
 \label{eq:e8v34-ellipticity}
\end{equation}
The densitized inverse metric of V23 is an example on its protected chart.
Define
\begin{equation}
 k_g[u,v]
 =\sum_{K\in\mathcal T_n}
 \int_K\langle\nabla u,G_g\nabla v\rangle\,dx.
 \label{eq:e8v34-form}
\end{equation}
Let
\begin{equation}
 \mathscr K_n
 :=\bigoplus_{K\in\mathcal T_n}
 L^2(K;\mathbb C^{dN})
 \label{eq:e8v34-K}
\end{equation}
with its standard orthogonal direct-sum norm.

\begin{definition}[Canonical kinetic loader]
\label{def:e8v34-loader}
For \(u\in V_n\), set
\begin{equation}
 \mathcal J_g u
 :=\bigl(G_g^{1/2}\nabla u|_K\bigr)_{K\in\mathcal T_n}
 \in\mathscr K_n.
 \label{eq:e8v34-J}
\end{equation}
\end{definition}

\begin{theorem}[Literal common-space energy identity]
\label{thm:e8v34-energy-identity}
For every real protected metric history and \(u,v\in V_n\),
\begin{align}
 \langle\mathcal J_gu,\mathcal J_gv\rangle_{\mathscr K_n}
 &=k_g[u,v],
 \label{eq:e8v34-polarization}\\
 \|\mathcal J_gu\|_{\mathscr K_n}^2
 &=k_g[u,u].
 \label{eq:e8v34-energy}
\end{align}
All histories use the same codomain \(\mathscr K_n\); hence comparison of
the loaders uses the identity map and satisfies every pairwise and
three-history cocycle exactly.
\end{theorem}

\begin{proof}
On each cell, the positive coefficient square root is Hermitian and
\[
 \langle G_g^{1/2}\nabla u,G_g^{1/2}\nabla v\rangle
 =\langle\nabla u,G_g\nabla v\rangle.
\]
Integrate and sum the orthogonal cell components.  Since
\(\mathscr K_n\) is defined from the mesh and not from \(g\), all
comparisons are literal identities in one Hilbert space.
\end{proof}

\begin{remark}[Relation to the operator square root]
\label{rem:e8v34-operator-root}
The factor \(\mathcal J_g\) satisfies
\(\mathcal J_g^*\mathcal J_g=A_g\) in form sense, where \(A_g\) is the
kinetic operator.  It need not equal the positive operator square root
\(A_g^{1/2}\) as a map into the original field Hilbert space.  It is a valid
common-space loader for form-energy and capacity estimates because
\eqref{eq:e8v34-energy} is exact.  DQSRC and fixed-Hilbert semigroup
contacts remain separate.
\end{remark}

\subsection{A noncommuting matrix square-root bound}
\label{sec:e8v34-matrix-root}

\begin{lemma}[Coefficient square roots are Lipschitz on an elliptic box]
\label{lem:e8v34-matrix-root}
If \(G,H\) are Hermitian matrices with
\(G,H\succeq\nu_G I\), then
\begin{equation}
 \boxed{
 \|G^{1/2}-H^{1/2}\|
 \leq{\|G-H\|\over
 \lambda_{\min}(G)^{1/2}+\lambda_{\min}(H)^{1/2}}
 \leq{\|G-H\|\over2\sqrt{\nu_G}}.}
 \label{eq:e8v34-matrix-root}
\end{equation}
\end{lemma}

\begin{proof}
Put \(X=G^{1/2}-H^{1/2}\).  It solves the Sylvester equation
\[
 G^{1/2}X+XH^{1/2}=G-H.
\]
The solution has the convergent representation
\[
 X=\int_0^\infty
 e^{-tG^{1/2}}(G-H)e^{-tH^{1/2}}\,dt.
\]
Indeed, differentiating the integrand and integrating from zero to infinity
recovers the Sylvester equation.  Taking norms gives
\[
 \|X\|
 \leq\|G-H\|
 \int_0^\infty
 e^{-t(\sqrt{\lambda_{\min}(G)}
       +\sqrt{\lambda_{\min}(H)})}\,dt,
\]
which is the first bound; the ellipticity floor gives the second.
\end{proof}

\begin{corollary}[Cellwise coefficient contact]
\label{cor:e8v34-cell-contact}
Suppose \(G\) and \(H\) differ only on a cell \(K_i\), and
\begin{equation}
 \operatorname*{ess\,sup}_{x\in K_i}\|G(x)-H(x)\|
 \leq\Delta_i.
 \label{eq:e8v34-Delta}
\end{equation}
For a fixed \(u\in V_n\), define the local energy
\begin{equation}
 e_i(H;u)
 :=\int_{K_i}\langle\nabla u,H\nabla u\rangle\,dx.
 \label{eq:e8v34-local-energy}
\end{equation}
Then
\begin{equation}
 \boxed{
 \|\mathcal J_Gu-\mathcal J_Hu\|_{\mathscr K_n}^2
 \leq\omega_i^2 e_i(H;u),\qquad
 \omega_i:={\Delta_i\over2\nu_G}.}
 \label{eq:e8v34-contact}
\end{equation}
\end{corollary}

\begin{proof}
The loader difference is supported only in the \(K_i\) summand.  By
\Cref{lem:e8v34-matrix-root},
\[
 \int_{K_i}
 \|(G^{1/2}-H^{1/2})\nabla u\|^2dx
 \leq{\Delta_i^2\over4\nu_G}
 \int_{K_i}|\nabla u|^2dx.
\]
Ellipticity of \(H\) gives
\(\nu_G\int_{K_i}|\nabla u|^2dx\leq e_i(H;u)\).
Combine the two estimates.
\end{proof}

\subsection{Exact heat-bath incidence without a volume factor}
\label{sec:e8v34-incidence}

Let \(Q_i\) conditionally resample the coefficient on cell \(K_i\) while
holding its exterior fixed.  Given the exterior, let \(G_i,G_i'\) be two
independent coefficient draws.  Assume their conditional support diameter
obeys
\begin{equation}
 \|G_i-G_i'\|_{L^\infty(K_i)}\leq\Delta_i
 \quad\text{almost surely}.
 \label{eq:e8v34-conditional-diameter}
\end{equation}

\begin{theorem}[Fixed-source one-cell incidence]
\label{thm:e8v34-one-cell}
If \(u\in V_n\) is independent of the resampled metric coefficient and
\(F(g)=\mathcal J_gu\), then
\begin{equation}
 \operatorname{Var}_i(F)
 \leq{\omega_i^2\over2}
 \int e_i(G_g;u)\,d\mu(g),
 \qquad
 \omega_i={\Delta_i\over2\nu_G}.
 \label{eq:e8v34-one-cell}
\end{equation}
\end{theorem}

\begin{proof}
The V32 conditional replacement identity gives
\[
 \operatorname{Var}_i(F)
 ={1\over2}\mathbb E
 \|\mathcal J_Gu-\mathcal J_{G'}u\|^2.
\]
Apply \Cref{cor:e8v34-cell-contact} with \(H=G'\).  Conditional
exchangeability makes
\(\mathbb E e_i(G';u)=\mathbb E e_i(G;u)\).
\end{proof}

\begin{theorem}[Orthogonal-cell weighted incidence]
\label{thm:e8v34-orthogonal-incidence}
Let the heat-bath rates be \(c_i>0\).  Under the hypotheses of
\Cref{thm:e8v34-one-cell},
\begin{equation}
 \boxed{
 \sum_i c_i\operatorname{Var}_i(\mathcal J_gu)
 \leq {1\over2}
 \left(\sup_i c_i\omega_i^2\right)
 \int k_g[u,u]\,d\mu(g).}
 \label{eq:e8v34-orthogonal-incidence}
\end{equation}
The relative incidence coefficient contains no factor
\(\#\mathcal T_n\).
\end{theorem}

\begin{proof}
Multiply \eqref{eq:e8v34-one-cell} by \(c_i\) and sum:
\[
 \sum_i c_i\operatorname{Var}_i(F)
 \leq {1\over2}\int
 \sum_i c_i\omega_i^2 e_i(G_g;u)\,d\mu(g).
\]
All \(e_i\) are nonnegative, and their sum is
\(k_g[u,u]\).  Pull out the largest weight.
\end{proof}

\begin{remark}[Why V32's crude sum is improved]
\label{rem:e8v34-improvement}
The general V32 bound sums relative coefficients because replacement
effects may occupy the same loader directions.  Here cell coefficient
contacts lie in orthogonal summands of \(\mathscr K_n\), and their local
energies add exactly.  This converts a sum of coefficients into the
supremum in \eqref{eq:e8v34-orthogonal-incidence}.
\end{remark}

\begin{corollary}[Insertion in a one-cell V31 heat-bath window]
\label{cor:e8v34-window}
For the exact resampling projection \(Q_i\), the fixed-source kinetic loader
has incidence coefficient
\begin{equation}
 b_{{\rm kin},i}
 ={1\over2}\left({\Delta_i\over2\nu_G}\right)^2.
 \label{eq:e8v34-bkin}
\end{equation}
If \(p_{*,i}\) is the corresponding worst-exterior conditional bad
probability, a sufficient one-cell window is
\begin{equation}
 \boxed{
 \sqrt{p_{*,i}}+{\Delta_i\over2\nu_G}<1,}
 \label{eq:e8v34-window}
\end{equation}
before source-replacement, bad-exterior, and bridge debits are added.
\end{corollary}

\begin{proof}
Bound the local energy in \eqref{eq:e8v34-one-cell} by the total kinetic
energy, so \(b_i=\omega_i^2/2\).  Insert it in the V31 exact-resampling
condition \(\sqrt{p_{*,i}}+\sqrt{2b_i}<1\).
\end{proof}

\begin{remark}[The assembled restricted norm remains separate]
\label{rem:e8v34-assembled}
For the full weighted heat-bath Dirichlet form,
\eqref{eq:e8v34-orthogonal-incidence} gives the incidence coefficient
\[
 b_{\rm HB}
 ={1\over2}\sup_i c_i
 \left({\Delta_i\over2\nu_G}\right)^2.
\]
Its partner in an RSFC strict window must be the restricted norm of the
same assembled physical transfer.  The one-cell identity
\(\|M_{B_i}Q_i\|^2=p_{*,i}\) does not prove that assembled operator row;
this is precisely the cluster-assembly distinction retained from the V30
audit.
\end{remark}

\subsection{From metric diameter to coefficient diameter}
\label{sec:e8v34-metric}

Let
\[
 \mathcal G(g)=\sqrt{\det g}\,g^{-1}
\]
be the densitized inverse from V23.

\begin{corollary}[Metric-cell replacement bound]
\label{cor:e8v34-metric-diameter}
On a fixed metric box
\(\nu I\preceq g,h\preceq\Lambda I\), let
\[
 \|\mathcal G(g)-\mathcal G(h)\|
 \leq C_{\nu,\Lambda,d}\|g-h\|
\]
be the V23 Lipschitz bound.  If a conditional metric update has diameter
\begin{equation}
 \|g_i-g_i'\|_{L^\infty(K_i)}\leq d_i,
 \label{eq:e8v34-metric-diameter}
\end{equation}
then
\begin{equation}
 \Delta_i\leq C_{\nu,\Lambda,d}d_i,
 \qquad
 \omega_i\leq
 {C_{\nu,\Lambda,d}d_i\over2\nu_G}.
 \label{eq:e8v34-metric-to-coefficient}
\end{equation}
\end{corollary}

\begin{proof}
Apply the V23 densitized-inverse Lipschitz estimate pointwise on \(K_i\)
and take the essential supremum.
\end{proof}

\begin{assessment}[Full heat baths versus narrow updates]
\label{ass:e8v34-diameter}
An exact heat bath over the entire protected metric box has a coefficient
diameter controlled by the box but not necessarily small enough for
\eqref{eq:e8v34-window}.  The estimate becomes perturbative for a narrow
conditional block, a small-step reversible update, or a carrier whose
conditional coefficient diameter is concentrated.  Replacing worst-case
diameter by a conditional second moment is possible directly in the proof,
but it must remain weighted by the corresponding local energy rather than
factorized without a correlation estimate.
\end{assessment}

\begin{theorem}[Second-moment refinement]
\label{thm:e8v34-second-moment}
Without the uniform diameter assumption,
\begin{equation}
 \operatorname{Var}_i(\mathcal J_gu)
 \leq {1\over8\nu_G^2}
 \mathbb E\!\left[
 \|G_i-G_i'\|_{L^\infty(K_i)}^2
 e_i(G_i';u)\right].
 \label{eq:e8v34-second-moment}
\end{equation}
If the conditional correlation bound
\begin{equation}
 \mathbb E\!\left[
 \|G_i-G_i'\|_\infty^2e_i(G_i';u)\mid y\right]
 \leq d_{2,i}^2\,
 \mathbb E[e_i(G_i';u)\mid y]
 \label{eq:e8v34-correlation}
\end{equation}
holds, then \eqref{eq:e8v34-one-cell} holds with
\(\omega_i=d_{2,i}/(2\nu_G)\).
\end{theorem}

\begin{proof}
Keep the actual coefficient difference in the proof of
\Cref{cor:e8v34-cell-contact}, then apply the factor \(1/2\) in the V32
replacement identity.  This gives \eqref{eq:e8v34-second-moment}.  Insert
\eqref{eq:e8v34-correlation} and use conditional exchangeability.
\end{proof}

\subsection{Metric-dependent sources and response loaders}
\label{sec:e8v34-sources}

For a source \(u_g\) which changes under the heat-bath update,
\begin{align}
 \mathcal J_gu_g-\mathcal J_{g'}u_{g'}
 ={}&\mathcal J_g(u_g-u_{g'})
 \notag\\
 &+\bigl(\mathcal J_g-\mathcal J_{g'}\bigr)u_{g'}.
 \label{eq:e8v34-source-split}
\end{align}
The second term is the coefficient contact just proved.  The first may
occupy many cell summands and needs its own local or quasilocal influence
estimate.

\begin{definition}[Kinetic source replacement card, KSRC]
\label{def:e8v34-KSRC}
KSRC for a metric-dependent kinetic loader consists of:
\begin{enumerate}[label=\textup{(K\arabic*)},leftmargin=4em]
\item the canonical coefficient-space factorization
\eqref{eq:e8v34-energy};
\item a coefficient diameter or correlated second-moment row giving
\eqref{eq:e8v34-one-cell};
\item a relative replacement estimate for
\(\mathcal J_g(u_g-u_{g'})\), with a summable distance influence if it is
not strictly local;
\item a combined strict V31 heat-bath window after good/bad exterior
splitting;
\item first metric-derivative and adjacent-bridge versions; and
\item compatible mass-density, boundary, gauge, spin, and physical
reflection contacts.
\end{enumerate}
\end{definition}

\begin{theorem}[Real kinetic EPAC is populated for fixed trial sources]
\label{thm:e8v34-EPAC-populated}
For the real nested \(P_1\) kinetic forms of V23--V24 and a fixed trial
source \(u\), the common-space/cocycle row of EPAC and the coefficient
square-root replacement row are populated by
\Cref{thm:e8v34-energy-identity,thm:e8v34-one-cell}.  The assembled
kinetic incidence row is \eqref{eq:e8v34-orthogonal-incidence}.
\end{theorem}

\begin{proof}
The loader \(\mathcal J_g u\) takes values in the fixed
\(\mathscr K_n\), so its comparison maps are identities and obey the
cocycle.  The coefficient replacement is
\Cref{thm:e8v34-one-cell}; orthogonal assembly is
\Cref{thm:e8v34-orthogonal-incidence}.  Nested inclusion and coefficient
quadrature compatibility are the exact V24 results.
\end{proof}

\begin{firewall}[Scope of the populated branch]
\label{fire:e8v34-scope}
The theorem concerns the real bosonic principal kinetic form and fixed
trial sources.  It does not identify the coefficient loader with a
nonnormal complex operator square root, control a metric-dependent
interacting Gibbs or transfer source, or include mass-density, gauge-link,
spin, Yukawa, determinant, and derivative-insertion contacts.  Those are
the remaining KSRC rows.
\end{firewall}

\begin{assessment}[Progress at V34]
\label{ass:e8v34-status}
The polar-cocycle obstruction is removed on the real nested kinetic branch
by a canonical coefficient-space loader.  One-cell coefficient contacts
are local, and orthogonality of cell summands replaces the dangerous
volume sum by a supremum of weighted conditional diameters.  This gives an
explicit kinetic incidence coefficient for the V31 window.  The next audit
must determine whether the newest YM or NS work supplies an assembled
operator or influence mechanism for the remaining metric-dependent source
and bad-exterior rows.
\end{assessment}

\begin{programme}[Eighth-series V35: companion audit and next contact]
\label{prog:e8v34-V35}
V35 will perform the compulsory companion audit.  It will compare any new
YM cluster-operator assembly with the orthogonal coefficient-space
mechanism above and check whether the NS rank-one refinement has produced a
useful summable nonlocal influence row.  The next mathematical target will
be selected from the metric-dependent source replacement, the
mass-density contact, and the bad-exterior cluster assembly.
\end{programme}

\begin{firewall}[Claim boundary after eighth-series V34]
\label{fire:e8v34-boundary}
V34 constructs the canonical common coefficient-space kinetic loader;
proves a noncommuting coefficient square-root Lipschitz bound; derives
one-cell and orthogonally assembled heat-bath incidence estimates; and
populates the real fixed-source kinetic branch of EPAC.

V34 does not populate KSRC for metric-dependent interacting sources, the
complex DQSRC row, lower-order and chiral contacts, conditional
gravitational activity, bad-exterior cluster assembly, cutoff removal, or
the Standard Model--Einstein continuum.
\end{firewall}

\begin{theorem}[Eighth-series V34 result]
\label{thm:e8v34-result}
For a real uniformly elliptic nested \(P_1\) kinetic form, the canonical
loader
\[
 \mathcal J_gu=(G_g^{1/2}\nabla u|_K)_K
\]
lies in one fixed cell-gradient Hilbert space and has norm exactly equal to
the kinetic energy.  A one-cell coefficient heat bath with conditional
diameter \(\Delta_i\) has incidence
\[
 \operatorname{Var}_i(\mathcal J_gu)
 \leq{1\over2}
 \left({\Delta_i\over2\nu_G}\right)^2
 \mathbb E e_i(G_g;u),
\]
and orthogonal cell assembly gives the global relative coefficient
\(
 \frac12\sup_i c_i(\Delta_i/(2\nu_G))^2
\)
without a volume factor; its assembled restricted norm remains a separate row.
\end{theorem}

\begin{proof}
Use \Cref{thm:e8v34-energy-identity,
lem:e8v34-matrix-root,thm:e8v34-one-cell,
thm:e8v34-orthogonal-incidence,cor:e8v34-window,
thm:e8v34-EPAC-populated}.
\end{proof}

\paragraph{Parent-version record.}
V34 was formed from
\texttt{Renewal\_SM\_Einstein\_Continuum\_Research\_Notes\_V33.tex}.
The V33 parent had (12\,746) logical source lines, (469\,142) bytes,
and SHA--256
\[
 \texttt{8638958C60649E07B8429B510463091205CF6EF201422FA2768EA99E7D885C22}.
\]
The next compulsory companion audit is V35.

% =============================================================================
% END APPENDED EIGHTH-SERIES V34 MATERIAL
% =============================================================================

% =============================================================================
% BEGIN APPENDED EIGHTH-SERIES V35 MATERIAL
% =============================================================================

\section{V35 companion audit: internal blocks and commuting smoothing}
\label{sec:v8v35}

This is the compulsory audit following V30.  Yang--Mills V90--V94 has
replaced the uncentered cluster-operator route by an enlarged conditional
block on which the bad event is internal.  The conditional transfer then
commutes exactly with conditional centering, which removes the ground-score
and leakage commutators from that proof path.  The Navier--Stokes active
file is unchanged.  The present note proves the corresponding
Hilbert-valued compiler for SM kinetic and response loaders and records the
extra conditional-mean term which does not occur in a pure conditional
variance estimate.

\subsection{Exact audit record}
\label{sec:e8v35-audit}

\begin{audit}[Companion files read at V35]
\label{audit:e8v35-files}
The audit used:
\begin{enumerate}[label=\textup{(\roman*)},leftmargin=3.8em]
\item
\texttt{Renewal\_Yang\_Mills\_Mass\_Gap\_Research\_Notes\_V94.tex},
with \(44\,139\) logical source lines, \(1\,537\,433\) bytes, one live
terminator, and SHA--256
\[
 \texttt{7C45CA13CCB5556D35BF6FC7E5FCDC5E0CBF3614398C11C5F33DD932C1021506};
\]
\item
\texttt{Renewal\_NS\_Stability\_Research\_Notes\_V26.tex},
with \(5\,319\) logical source lines, \(278\,283\) bytes, one live
terminator, and SHA--256
\[
 \texttt{82C887F8C2C69E4F28FAD7C3DC8DE5B9099CB5A4BF431EBA1FFF3E91935978AD}.
\]
\end{enumerate}
Both files are structurally complete and end with one document terminator.
\end{audit}

\begin{theorem}[Yang--Mills V94 input selected by the audit]
\label{thm:e8v35-YM-input}
The audited YM continuation proves the following conditional mechanism.
\begin{enumerate}[label=\textup{(\roman*)},leftmargin=3.8em]
\item A bounded enlarged star contains every link entering the declared
gauge-invariant frustration event, so the event is internal to the
conditional fibre.
\item A direct integral of conditional positive self-adjoint Markov
transfers commutes exactly with the far-exterior conditional expectation
and its fluctuation projection.
\item The restricted direct-integral norm is the essential supremum of the
conditional restricted norms, and the singular ground function cancels
fibrewise for the stated positive half-slab Doob transform.
\item Commuting projected smoothing gives a gap compiler without an
uncentered global smoothing-operator sum, a conditional ground-score
derivative, or a transfer/projection leakage term.
\item Four source-specific rows remain open: a uniform conditional action
margin, good-sector coercivity, enlarged-block approximate tensorization,
and physical generator/form incidence.
\end{enumerate}
\end{theorem}

\begin{proof}
These are the results labelled
\texttt{lem:s6v94-star-size},
\texttt{lem:s6v94-internal},
\texttt{thm:s6v94-commute},
\texttt{lem:s6v94-direct-norm},
\texttt{cor:s6v94-fibre-cancel},
\texttt{thm:s6v94-smoothing}, and
\texttt{thm:s6v94-bootstrap}, with the acquisition firewall of V94
retained.
\end{proof}

\begin{theorem}[Navier--Stokes status at the audit]
\label{thm:e8v35-NS-input}
The NS active file remains V26 with the same rank-one Oseen kernel
calculation audited at V30.  It supplies no new conditional block,
restricted-transfer, loader-incidence, or gravitational capacity theorem.
The full Navier--Stokes regularity result remains open.
\end{theorem}

\begin{proof}
The file digest and terminal V26 programme are unchanged from the V30
audit.
\end{proof}

\begin{assessment}[Audit-guided direction]
\label{ass:e8v35-direction}
The V34 coefficient-space loader already makes cell-gradient contacts
orthogonal.  The V94 enlargement principle addresses the matching
restricted-norm side: choose a bounded metric patch containing the bad
event and all local coefficient contacts, then use a conditional transfer
on that same patch.  This makes the fluctuation projection commute
algebraically.  The SM energy debit still contains a conditional mean,
and metric differentiation may change the conditional carrier; both terms
must remain visible.
\end{assessment}

\subsection{Internal-block conditional transfers}
\label{sec:e8v35-block}

Let a regulated carrier disintegrate over a bounded block \(\Lambda\):
\begin{equation}
 d\mu(x_\Lambda,\eta)
 =d\mu_\Lambda^\eta(x_\Lambda)\,d\bar\mu(\eta).
 \label{eq:e8v35-disintegration}
\end{equation}
Let \(E\) be conditional expectation onto the far exterior \(\eta\), and
let
\begin{equation}
 D=I-E.
 \label{eq:e8v35-D}
\end{equation}
For almost every \(\eta\), let \(P^\eta\) be a positive self-adjoint Markov
contraction on \(L^2(\mu_\Lambda^\eta)\), and let \(P\) be its direct
integral.  Let \(B\) be internal to the block, with fibre section
\(B^\eta\).

\begin{theorem}[Exact conditional commutation and norm]
\label{thm:e8v35-commutation}
The direct-integral transfer satisfies
\begin{equation}
 EP=PE=E,\qquad [D,P]=0.
 \label{eq:e8v35-commutation}
\end{equation}
Moreover,
\begin{equation}
 \boxed{
 \|M_BP\|_{2\to2}^2
 =\operatorname*{ess\,sup}_{\eta}
 \|M_{B^\eta}P^\eta\|_{2\to2}^2.}
 \label{eq:e8v35-direct-norm}
\end{equation}
The same statements hold componentwise on
\(L^2(\mu;\mathcal K)\) for every separable Hilbert space \(\mathcal K\).
\end{theorem}

\begin{proof}
Each \(P^\eta\) preserves constants and \(\mu_\Lambda^\eta\).  Therefore it
fixes exterior-measurable functions and its conditional fibre average is
unchanged:
\[
 PE=E,\qquad EP=E.
\]
Subtract from the identity to obtain \([D,P]=0\).  The norm of a
decomposable operator is the essential supremum of its fibre norms, which
gives \eqref{eq:e8v35-direct-norm}.  Tensoring every operator with
\(I_{\mathcal K}\) preserves its norm and all algebraic identities.
\end{proof}

\subsection{Hilbert-valued commuting projected smoothing}
\label{sec:e8v35-projected}

\begin{theorem}[Projected loader smoothing]
\label{thm:e8v35-projected}
In the setting above, put
\begin{equation}
 \eta_B=\|M_BP\|_{2\to2}^2.
 \label{eq:e8v35-eta}
\end{equation}
Then for every \(F\in L^2(\mu;\mathcal K)\) and every \(\theta>0\),
\begin{equation}
 \boxed{
 \|M_BDF\|_2^2
 \leq
 (1+\theta)\eta_B\|DF\|_2^2
 +(1+\theta^{-1})
 \langle F,(I-P)F\rangle.}
 \label{eq:e8v35-projected}
\end{equation}
\end{theorem}

\begin{proof}
Split
\[
 M_BDF=M_BPDF+M_B(I-P)DF
\]
and apply the tunable square inequality.  The first term is at most
\(\eta_B\|DF\|_2^2\).  Since \(0\leq P\leq I\),
\[
 (I-P)^2\leq I-P.
\]
Thus the second square is at most
\(\langle DF,(I-P)DF\rangle\).  By
\eqref{eq:e8v35-commutation}, \(D\) commutes with \(I-P\); since
\(0\leq D\leq I\),
\[
 D(I-P)D=D(I-P)\leq I-P.
\]
This proves the energy term in \eqref{eq:e8v35-projected}.
\end{proof}

\begin{corollary}[Conditionally centered loader window]
\label{cor:e8v35-centered}
Suppose \(EF=0\) and
\begin{equation}
 \langle F,(I-P)F\rangle
 \leq a+b\|F\|_2^2.
 \label{eq:e8v35-incidence}
\end{equation}
Then the bad-set relative coefficient obtained from
\eqref{eq:e8v35-projected} is
\begin{equation}
 \kappa(\theta)
 =(1+\theta)\eta_B+(1+\theta^{-1})b.
 \label{eq:e8v35-centered-kappa}
\end{equation}
Its optimum is
\begin{equation}
 \inf_{\theta>0}\kappa(\theta)
 =(\sqrt{\eta_B}+\sqrt b)^2.
 \label{eq:e8v35-centered-opt}
\end{equation}
Hence the strict centered window is
\begin{equation}
 \boxed{\sqrt{\eta_B}+\sqrt b<1.}
 \label{eq:e8v35-centered-window}
\end{equation}
\end{corollary}

\begin{proof}
If \(EF=0\), then \(DF=F\).  Insert
\eqref{eq:e8v35-incidence} in
\eqref{eq:e8v35-projected}.  Minimize
\(\eta_B+b+\theta\eta_B+b/\theta\) by arithmetic--geometric mean.
\end{proof}

\begin{remark}[Gain over the unprojected split]
\label{rem:e8v35-gain}
The V30 unprojected Markov-kernel estimate paid twice the global jump form.
Positivity and exact commutation give only one copy in
\eqref{eq:e8v35-projected}.  This changes the centered optimized condition
from a square root of \(2b\) to a square root of \(b\).  The gain applies
only to \(DF\), not automatically to the conditional mean.
\end{remark}

\subsection{The conditional mean debit}
\label{sec:e8v35-mean}

Define the worst conditional bad mass
\begin{equation}
 p_*=\operatorname*{ess\,sup}_{\eta}
 \mu_\Lambda^\eta(B^\eta).
 \label{eq:e8v35-pstar}
\end{equation}

\begin{lemma}[Bad conditional mean]
\label{lem:e8v35-mean}
For every Hilbert-valued \(F\),
\begin{equation}
 \|M_BEF\|_2^2\leq p_*\|EF\|_2^2\leq p_*\|F\|_2^2.
 \label{eq:e8v35-mean}
\end{equation}
\end{lemma}

\begin{proof}
The function \(EF\) is constant on each conditional fibre.  Therefore
\[
 \|M_BEF\|_2^2
 =\int
 \mu_\Lambda^\eta(B^\eta)\|EF(\eta)\|_{\mathcal K}^2
 d\bar\mu(\eta).
\]
Use \eqref{eq:e8v35-pstar} and contractivity of \(E\).
\end{proof}

\begin{theorem}[Full internal-block loader debit]
\label{thm:e8v35-full}
Assume \eqref{eq:e8v35-incidence}.  For \(\theta>0\), set
\begin{equation}
 K_\theta
 :=(1+\theta)\eta_B+(1+\theta^{-1})b.
 \label{eq:e8v35-Ktheta}
\end{equation}
For every \(\alpha>0\),
\begin{align}
 \|M_BF\|_2^2
 \leq{}&
 \left\{(1+\alpha)p_*
 +(1+\alpha^{-1})K_\theta\right\}\|F\|_2^2
 \notag\\
 &+(1+\alpha^{-1})(1+\theta^{-1})a.
 \label{eq:e8v35-full}
\end{align}
If \(p_*,K_\theta>0\), the best relative coefficient in \(\alpha\) is
\begin{equation}
 \boxed{
 \left(\sqrt{p_*}+\sqrt{K_\theta}\right)^2,}
 \qquad
 \alpha_{\rm opt}=\sqrt{K_\theta/p_*}.
 \label{eq:e8v35-full-opt}
\end{equation}
\end{theorem}

\begin{proof}
Since \(F=EF+DF\),
\[
 \|M_BF\|^2
 \leq(1+\alpha)\|M_BEF\|^2
 +(1+\alpha^{-1})\|M_BDF\|^2.
\]
Apply \Cref{lem:e8v35-mean,thm:e8v35-projected}, use
\(\|DF\|\leq\|F\|\), and insert
\eqref{eq:e8v35-incidence}.  The relative coefficient is
\[
 p_*+K_\theta+\alpha p_*+{K_\theta\over\alpha},
\]
whose minimum is \eqref{eq:e8v35-full-opt}.
\end{proof}

\begin{corollary}[Full strict window]
\label{cor:e8v35-full-window}
A sufficient strict form-absorption condition for the full loader is
\begin{equation}
 \sqrt{p_*}
 +\sqrt{(1+\theta)\eta_B+(1+\theta^{-1})b}<1
 \label{eq:e8v35-full-window}
\end{equation}
for some \(\theta>0\), with the constant remainder in
\eqref{eq:e8v35-full} summable along the regulator sequence.
\end{corollary}

\begin{proof}
Under the displayed condition,
\eqref{eq:e8v35-full-opt} is below one.  The remaining term is the
non-form remainder.
\end{proof}

\begin{remark}[The conditional mean can be removed only honestly]
\label{rem:e8v35-centering}
For a conditional variance problem the loader is already \(DF\), as in the
audited YM route.  A Standard Model energy loader generally has
\(EF\ne0\).  It may be centered only if the resulting conditional mean is
handled elsewhere by an exact identity, a good-sector estimate, or a
summable source term.  Dropping it would remove the \(p_*\) term without
proof.
\end{remark}

\subsection{Internal metric blocks}
\label{sec:e8v35-metric-block}

Let a local kinetic source meet a set of cells \(S_A\), and let its
coefficient replacement, bad-ellipticity event, and finite-range
neighbouring action terms depend on a bounded enlargement
\(\Lambda_A\supset S_A\).

\begin{definition}[Internal metric-block smoothing card, IMBSC]
\label{def:e8v35-IMBSC}
IMBSC consists of:
\begin{enumerate}[label=\textup{(I\arabic*)},leftmargin=4em]
\item a cofinally bounded block \(\Lambda_A\) containing the bad event and
every coefficient/source contact assigned to the update;
\item a disintegration over its far exterior and a positive self-adjoint
conditional transfer \(P_A^\eta\);
\item a uniform essential-supremum restricted norm \(\eta_{B,A}\);
\item a worst-exterior conditional bad-mass bound \(p_{*,A}\), unless the
loader is conditionally centered;
\item physical generator/form incidence, using the V34 orthogonal
coefficient-space estimate where applicable;
\item a strict centered or full window from
\eqref{eq:e8v35-centered-window} or
\eqref{eq:e8v35-full-window};
\item bounded overlap or approximate tensorization for the family of
enlarged blocks; and
\item complex-circle, metric-derivative, bridge, reflection, and Ward
compatibility with summable remainders.
\end{enumerate}
\end{definition}

\begin{theorem}[IMBSC populates the internal kinetic RSFC row]
\label{thm:e8v35-IMBSC}
If IMBSC holds for the real kinetic energy loader of V34, then its bad-set
form debit has a strict relative coefficient and the explicit remainder
in \eqref{eq:e8v35-full}.  If the corresponding base and first-derivative
response loaders satisfy the same card with summable remainders, the V30
RSFC-to-GGAFC handoff applies.
\end{theorem}

\begin{proof}
Rows (I1)--(I2) give \eqref{eq:e8v35-commutation}; rows (I3)--(I5) give
the constants in \Cref{thm:e8v35-full}; row (I6) makes the relative
coefficient strict.  Row (I7) prevents a volume loss, and row (I8) supplies
the remaining RSFC analytic and physical rows.  Apply the V30 handoff.
\end{proof}

\begin{firewall}[Metric differentiation can move the disintegration]
\label{fire:e8v35-moving-carrier}
Exact commutation in \eqref{eq:e8v35-commutation} holds at a fixed carrier.
If the metric response parameter changes
\(\mu_\Lambda^\eta\), then differentiating \(E\) produces a conditional
score covariance, just as in the audited YM V92 calculation.  A
complex-circle estimate may avoid differentiating \(E\) directly only if
the carrier comparison and Radon--Nikodym contacts are included uniformly
on that circle.  IMBSC row (I8) is therefore substantive.
\end{firewall}

\begin{firewall}[Positivity remains prior]
\label{fire:e8v35-positivity}
The conditional Markov transfer requires a positive carrier.  A complex
chiral determinant phase cannot be inserted into a conditional probability
without a separate positive representation or a controlled reweighting
theorem.  The enlarged-block algebra does not solve the determinant-line or
gravitational reflection-positivity gates.
\end{firewall}

\begin{assessment}[Progress at V35]
\label{ass:e8v35-status}
The audit yields a cleaner local route.  Internalizing the metric event and
all finite-range contacts in one bounded block makes the conditional
transfer commute exactly with conditional centering.  For centered loaders
the strict window is \(\sqrt{\eta_B}+\sqrt b<1\); for full energy loaders
the conditional mean adds the explicit \(p_*\) debit in
\eqref{eq:e8v35-full-window}.  The next target is to construct the minimal
metric block for the kinetic and mass-density components and to prove its
bounded overlap.
\end{assessment}

\begin{programme}[Eighth-series V36: minimal metric block and mass density]
\label{prog:e8v35-V36}
V36 will add the metric-dependent \(L^2\) mass-density component to the V34
coefficient-space loader.  It will prove its one-cell square-root contact,
construct a bounded patch containing both gradient and mass terms for a
fixed \(P_1\) source, and compute the exact overlap number on a
shape-regular mesh.  Gauge, spin, Yukawa, and chiral determinant contacts
will remain separately typed.
\end{programme}

\begin{firewall}[Claim boundary after eighth-series V35]
\label{fire:e8v35-boundary}
V35 performs the compulsory companion audit; imports the YM V94
internal-block route with its open physical rows; proves the
Hilbert-valued commuting projected-smoothing inequality; derives centered
and full loader windows; and formulates IMBSC.

V35 does not prove IMBSC for a gravitational carrier, the uniform
conditional action margin, conditional good-sector coercivity, the moving
carrier score, chiral positivity, complete SM contacts, cutoff removal, or
the Standard Model--Einstein continuum.
\end{firewall}

\begin{theorem}[Eighth-series V35 result]
\label{thm:e8v35-result}
When a bad metric event is internal to a bounded conditional block, any
positive self-adjoint fibre transfer commutes with conditional centering and
has restricted norm equal to the essential supremum of its fibre norms.
For a conditionally centered Hilbert loader this gives the optimal
elementary strict window
\[
 \sqrt{\eta_B}+\sqrt b<1.
\]
For a full loader the conditional mean is controlled by the worst
conditional bad mass and gives the sufficient window
\eqref{eq:e8v35-full-window}.  These statements populate the analytic
compiler only after the physical restricted norm, generator incidence,
block overlap, and carrier-contact rows of IMBSC are proved.
\end{theorem}

\begin{proof}
Use \Cref{thm:e8v35-commutation,thm:e8v35-projected,
cor:e8v35-centered,lem:e8v35-mean,thm:e8v35-full,
thm:e8v35-IMBSC}.
\end{proof}

\paragraph{Parent-version record.}
V35 was formed from
\texttt{Renewal\_SM\_Einstein\_Continuum\_Research\_Notes\_V34.tex}.
The V34 parent had (13\,253) logical source lines, (486\,331) bytes,
and SHA--256
\[
 \texttt{4B4AC9F76CA2600FAB3DC1B6C89394124D155BF5DE08306D9B7F4F4E868BA2E5}.
\]
The next compulsory companion audit is V40.

% =============================================================================
% END APPENDED EIGHTH-SERIES V35 MATERIAL
% =============================================================================

% =============================================================================
% BEGIN APPENDED EIGHTH-SERIES V36 MATERIAL
% =============================================================================

\section{Mass-density loaders and bounded internal metric patches}
\label{sec:v8v36}

V34 treated the densitized inverse metric in the principal gradient form.
The \(L^2\) and mass terms also move with the metric volume density.  This
note adds that component to the same fixed cellwise Hilbert space, proves
its one-cell replacement estimate, and constructs an internal block with a
mesh-independent overlap bound for any declared finite-range metric action.

\subsection{The fourth-root volume loader}
\label{sec:e8v36-volume}

For a real positive \(d\times d\) metric matrix, put
\begin{equation}
 W(g)=\sqrt{\det g},
 \qquad
 w(g)=W(g)^{1/2}=(\det g)^{1/4}.
 \label{eq:e8v36-w}
\end{equation}
On the metric box
\begin{equation}
 \mathsf K_{\nu,\Lambda}
 =\{g:\nu I\preceq g\preceq\Lambda I\},
 \label{eq:e8v36-box}
\end{equation}
one has
\begin{equation}
 \nu^{d/4}\leq w(g)\leq\Lambda^{d/4}.
 \label{eq:e8v36-w-bounds}
\end{equation}

\begin{lemma}[Lipschitz fourth-root determinant]
\label{lem:e8v36-w-Lipschitz}
For \(g,h\in\mathsf K_{\nu,\Lambda}\),
\begin{equation}
 |w(g)-w(h)|
 \leq L_w\|g-h\|_{\rm op},
 \qquad
 L_w={d\over4}\Lambda^{d/4}\nu^{-1}.
 \label{eq:e8v36-w-Lipschitz}
\end{equation}
\end{lemma}

\begin{proof}
The derivative on the positive cone is
\begin{equation}
 Dw_g[k]
 ={1\over4}w(g)\operatorname{Tr}(g^{-1}k).
 \label{eq:e8v36-Dw}
\end{equation}
On the box,
\[
 |\operatorname{Tr}(g^{-1}k)|
 \leq d\|g^{-1}\|_{\rm op}\|k\|_{\rm op}
 \leq d\nu^{-1}\|k\|_{\rm op},
\]
and \(w(g)\leq\Lambda^{d/4}\).  The box is convex, so integrate the
derivative along the segment from \(g\) to \(h\).
\end{proof}

\begin{corollary}[Cell mass-density contact]
\label{cor:e8v36-mass-contact}
Suppose two metric histories differ only on a cell \(K_i\) and
\begin{equation}
 \|g-h\|_{L^\infty(K_i;{\rm op})}\leq d_i.
 \label{eq:e8v36-di}
\end{equation}
For a fixed trial field \(u\), define
\begin{equation}
 e_i^{\rm mass}(h;u)
 :=m_0^2\int_{K_i}W(h)|u|^2\,dx,
 \qquad m_0\geq0.
 \label{eq:e8v36-mass-energy}
\end{equation}
Then
\begin{equation}
 \|m_0(w(g)-w(h))u\|_{L^2(K_i)}^2
 \leq\zeta_i^2 e_i^{\rm mass}(h;u),
 \qquad
 \zeta_i:={L_wd_i\over\nu^{d/4}}.
 \label{eq:e8v36-mass-contact}
\end{equation}
\end{corollary}

\begin{proof}
\Cref{lem:e8v36-w-Lipschitz} bounds the left side by
\[
 m_0^2L_w^2d_i^2\int_{K_i}|u|^2dx.
\]
Since \(W(h)\geq\nu^{d/2}\), the last integral is at most
\(\nu^{-d/2}m_0^{-2}e_i^{\rm mass}(h;u)\) when \(m_0>0\).  This gives
\eqref{eq:e8v36-mass-contact}.  If \(m_0=0\), both sides vanish.
\end{proof}

\subsection{Combined gradient and mass coefficient space}
\label{sec:e8v36-combined}

Let \(G_g=\sqrt{\det g}\,g^{-1}\) be the real densitized inverse coefficient
and retain the V34 gradient space.  Define
\begin{equation}
 \mathscr K_n^{\rm km}
 :=
 \bigoplus_{K\in\mathcal T_n}
 \left[
 L^2(K;\mathbb C^{dN})
 \oplus L^2(K;\mathbb C^N)
 \right]
 \label{eq:e8v36-Kkm}
\end{equation}
and the kinetic--mass loader
\begin{equation}
 \mathcal J_g^{\rm km}u
 :=
 \bigl(G_g^{1/2}\nabla u,\ m_0w(g)u\bigr)_{K\in\mathcal T_n}.
 \label{eq:e8v36-Jkm}
\end{equation}
Its cell energy is
\begin{equation}
 e_i^{\rm km}(g;u)
 =\int_{K_i}\langle\nabla u,G_g\nabla u\rangle dx
  +m_0^2\int_{K_i}W(g)|u|^2dx.
 \label{eq:e8v36-ekm}
\end{equation}

\begin{theorem}[Exact combined energy identity]
\label{thm:e8v36-combined-energy}
For every fixed \(u\in V_n\),
\begin{equation}
 \|\mathcal J_g^{\rm km}u\|_{\mathscr K_n^{\rm km}}^2
 =\sum_i e_i^{\rm km}(g;u).
 \label{eq:e8v36-combined-energy}
\end{equation}
Every metric history uses the same direct-sum space and the same cell
decomposition.
\end{theorem}

\begin{proof}
The gradient identity is the V34 coefficient-space factorization.  The
square of the second component is
\(m_0^2w(g)^2|u|^2=m_0^2W(g)|u|^2\).  Orthogonality of the two components
and of distinct cells gives the sum.
\end{proof}

\begin{theorem}[Combined one-cell heat-bath incidence]
\label{thm:e8v36-one-cell}
Let cell \(i\) be conditionally resampled, let \(g_i,g_i'\) be independent
conditional draws, and suppose
\[
 \|g_i-g_i'\|_{L^\infty(K_i;{\rm op})}\leq d_i.
\]
Let the V34 gradient coefficient have ellipticity floor \(\nu_G\) and
Lipschitz constant \(C_{\nu,\Lambda,d}\).  Put
\begin{equation}
 \omega_i
 ={C_{\nu,\Lambda,d}d_i\over2\nu_G},
 \qquad
 \chi_i=\max\{\omega_i,\zeta_i\}.
 \label{eq:e8v36-chi}
\end{equation}
For a metric-independent trial source \(u\),
\begin{equation}
 \boxed{
 \operatorname{Var}_i(\mathcal J_g^{\rm km}u)
 \leq{\chi_i^2\over2}
 \int e_i^{\rm km}(g;u)\,d\mu(g).}
 \label{eq:e8v36-one-cell}
\end{equation}
\end{theorem}

\begin{proof}
The squared replacement difference is the sum of its gradient and mass
components.  The V34 cell contact bounds the former by
\(\omega_i^2e_i^{\rm kin}(g';u)\), and
\Cref{cor:e8v36-mass-contact} bounds the latter by
\(\zeta_i^2e_i^{\rm mass}(g';u)\).  Both are at most
\(\chi_i^2e_i^{\rm km}(g';u)\).  Apply the V32 replacement identity and
conditional exchangeability.
\end{proof}

\begin{corollary}[Orthogonal weighted incidence for kinetic plus mass]
\label{cor:e8v36-weighted}
For heat-bath rates \(c_i>0\),
\begin{equation}
 \sum_i c_i
 \operatorname{Var}_i(\mathcal J_g^{\rm km}u)
 \leq{1\over2}
 \left(\sup_i c_i\chi_i^2\right)
 \int\sum_i e_i^{\rm km}(g;u)\,d\mu(g).
 \label{eq:e8v36-weighted}
\end{equation}
\end{corollary}

\begin{proof}
Sum \eqref{eq:e8v36-one-cell}, pull out the supremum, and use
\eqref{eq:e8v36-combined-energy}.
\end{proof}

\begin{corollary}[One-cell internal-block window]
\label{cor:e8v36-window}
If \(p_{*,i}\) is the worst-exterior conditional bad probability for the
cell update, a sufficient V31 exact-resampling window for the fixed
kinetic--mass loader is
\begin{equation}
 \boxed{\sqrt{p_{*,i}}+\chi_i<1,}
 \label{eq:e8v36-window}
\end{equation}
before bad-exterior, conditional-mean, and response-source debits are
added.
\end{corollary}

\begin{proof}
Equation \eqref{eq:e8v36-one-cell} gives incidence
\(b_i=\chi_i^2/2\).  Insert it in the V31 window
\(\sqrt{p_{*,i}}+\sqrt{2b_i}<1\).
\end{proof}

\subsection{Finite-range internal patches}
\label{sec:e8v36-patches}

Let \(\mathcal I_n\) be an interaction graph whose vertices are mesh cells.
Join two cells when one declared local action term or contact uses variables
from both.  Assume
\begin{equation}
 \deg_{\mathcal I_n}(K)\leq D_*
 \label{eq:e8v36-degree}
\end{equation}
uniformly in \(n\).  For \(R\geq0\), let
\(\mathbb B_R(K)\) be the graph ball of radius \(R\).

\begin{lemma}[Uniform graph-ball count]
\label{lem:e8v36-ball}
If \(D_*\geq2\), then
\begin{equation}
 |\mathbb B_R(K)|
 \leq b_R
 :=1+D_*\sum_{r=0}^{R-1}(D_*-1)^r.
 \label{eq:e8v36-ball}
\end{equation}
The same \(b_R\) bounds the number of radius-\(R\) cell balls containing
any fixed cell.
\end{lemma}

\begin{proof}
There are at most \(D_*\) vertices at distance one.  After the first step,
each nonbacktracking extension has at most \(D_*-1\) choices, giving at most
\(D_*(D_*-1)^{r-1}\) vertices at distance \(r\geq1\).  Sum the shells.
A ball \(\mathbb B_R(K_0)\) contains \(K\) exactly when
\(\operatorname{dist}(K_0,K)\leq R\), so the number of possible centers is
bounded by the same ball count.
\end{proof}

\begin{lemma}[Shape regularity gives bounded incidence degree]
\label{lem:e8v36-shape-degree}
For a conforming shape-regular simplicial mesh in fixed dimension, the
number of simplices incident on one vertex is bounded by a constant
\(q_*(d,\sigma)\), where \(\sigma\) is the shape-regularity constant.
Consequently every interaction graph obtained by connecting cells which
share a vertex has a degree bound depending only on \(d,\sigma\).
\end{lemma}

\begin{proof}
Shape regularity gives a positive lower bound
\(\omega_*(d,\sigma)\) on the solid angle of a simplex at each vertex.
The interiors of the solid-angle sectors of distinct incident simplices
are disjoint, while their union lies in the unit sphere of total measure
\(\omega_{d-1}\).  Hence at most
\(\omega_{d-1}/\omega_*(d,\sigma)\) simplices meet a vertex.  A simplex has
\(d+1\) vertices, so the number of other cells sharing at least one of them
is bounded in terms of \(d\) and \(q_*\).
\end{proof}

\begin{definition}[Minimal finite-range metric patch]
\label{def:e8v36-patch}
For a trial source \(u\), let
\begin{equation}
 S(u)=\{K\in\mathcal T_n:
       u|_K\not\equiv0\ \text{or}\ \nabla u|_K\not\equiv0\}.
 \label{eq:e8v36-support}
\end{equation}
If the conditional density and all assigned local action contacts have
interaction range \(R\), define
\begin{equation}
 \Lambda_R(u)
 :=\bigcup_{K\in S(u)}\mathbb B_R(K).
 \label{eq:e8v36-patch}
\end{equation}
\end{definition}

\begin{theorem}[Internality and patch size]
\label{thm:e8v36-patch}
Suppose every coefficient in
\(\mathcal J_g^{\rm km}u\) is cell-local and every action/contact term
meeting a cell in \(S(u)\) is supported within interaction distance \(R\).
Then the kinetic--mass loader, its bad coefficient event, and all assigned
conditional-density terms are measurable inside \(\Lambda_R(u)\).  Moreover,
\begin{equation}
 |\Lambda_R(u)|\leq b_R|S(u)|.
 \label{eq:e8v36-patch-size}
\end{equation}
If the source family consists of single nodal \(P_1\) basis functions, then
\(|S(u)|\leq q_*(d,\sigma)\), so the patch size is uniform in the mesh.
\end{theorem}

\begin{proof}
The definition of \(S(u)\) contains every cell on which either component of
\eqref{eq:e8v36-Jkm} is nonzero.  The range assumption places every action
or contact term touching those cells in their radius-\(R\) union, proving
internality.  The union bound and
\Cref{lem:e8v36-ball} give \eqref{eq:e8v36-patch-size}.  A nodal basis
function is supported precisely on the incident simplices of its vertex;
apply \Cref{lem:e8v36-shape-degree}.
\end{proof}

\begin{theorem}[Uniform overlap for nodal internal patches]
\label{thm:e8v36-overlap}
Let \(\{\varphi_v\}\) be all nodal \(P_1\) basis functions, and suppose the
cell interaction graph has degree at most \(D_*\).  Every cell belongs to at
most
\begin{equation}
 L_R:=(d+1)b_Rq_*(d,\sigma)
 \label{eq:e8v36-overlap}
\end{equation}
of the patches \(\Lambda_R(\varphi_v)\).  In particular, the overlap is
independent of the number and size of mesh cells.
\end{theorem}

\begin{proof}
If a cell \(K\) lies in \(\Lambda_R(\varphi_v)\), then some cell
\(K_0\) incident on \(v\) lies within graph distance \(R\) of \(K\).
There are at most \(b_R\) choices for \(K_0\).  Each such simplex has
\(d+1\) vertices; allowing the incidence bound \(q_*\) gives the stated
safe upper bound.  All constants depend only on the fixed geometry of the
mesh family and the declared interaction range.
\end{proof}

\begin{remark}[The overlap bound is deliberately conservative]
\label{rem:e8v36-overlap}
For \(R=0\), a cell belongs to exactly the \(d+1\) vertex stars associated
with its vertices, so \eqref{eq:e8v36-overlap} is not sharp.  Its purpose is
a cofinal uniform constant.  A sharper incidence count may be inserted
without changing the compiler.
\end{remark}

\subsection{Finite range is a substantive gravitational hypothesis}
\label{sec:e8v36-finite-range}

\begin{firewall}[Constraint reduction can destroy finite range]
\label{fire:e8v36-constraints}
A local pre-constraint lattice action may have fixed interaction range, but
solving elliptic gravitational or gauge constraints can produce Coulomb or
multipole tails.  In that case \(\Lambda_R(u)\) does not contain the full
conditional-density contact.  The missing exterior must be controlled by a
summable influence tail of the type isolated earlier in the programme; it
cannot be discarded by the mesh overlap count.
\end{firewall}

\begin{firewall}[Curvature and frame contacts]
\label{fire:e8v36-curvature}
A curvature discretization may use more than the cell metric, and a
fermionic mass or kinetic term uses frames, spin lifts, and the chiral
measure.  V36 applies once their interaction ranges and reference
Jacobians are explicitly assigned.  It does not prove that the physical
gravitational and fermionic constructions have the cell-local form assumed
above.
\end{firewall}

\begin{definition}[Kinetic--mass internal block card, KMIBC]
\label{def:e8v36-KMIBC}
KMIBC consists of:
\begin{enumerate}[label=\textup{(M\arabic*)},leftmargin=4em]
\item the common loader \eqref{eq:e8v36-Jkm} on a protected real metric
chart;
\item conditional diameter or correlated second-moment bounds for both
\(G_g\) and \(w(g)\);
\item a finite-range interaction graph, or a finite patch plus a summable
exterior influence tail;
\item the bounded size and overlap estimates
\eqref{eq:e8v36-patch-size}--\eqref{eq:e8v36-overlap};
\item the V35 conditional restricted norm, mean debit, and generator
incidence for that same patch;
\item metric-dependent source and first-derivative replacement rows; and
\item bridge, reflection, Ward, boundary, gauge, spin, and chiral contacts.
\end{enumerate}
\end{definition}

\begin{theorem}[KMIBC populates the real fixed-source IMBSC geometry rows]
\label{thm:e8v36-KMIBC}
For metric-independent nodal trial sources, KMIBC rows (M1)--(M5) populate
the real kinetic--mass loader, one-cell incidence, internality, bounded
patch, bounded overlap, and strict-window geometry rows of IMBSC.  If rows
(M6)--(M7) also hold with summable remainders, the corresponding response
loaders enter the V35 IMBSC-to-RSFC handoff.
\end{theorem}

\begin{proof}
The loader and incidence are
\Cref{thm:e8v36-combined-energy,thm:e8v36-one-cell,
cor:e8v36-weighted}.  Internality and uniform overlap are
\Cref{thm:e8v36-patch,thm:e8v36-overlap}.  Row (M5) supplies the remaining
V35 analytic constants.  Rows (M6)--(M7) are exactly the response and
physical compatibility rows of IMBSC.
\end{proof}

\begin{assessment}[Progress at V36]
\label{ass:e8v36-status}
The real protected kinetic branch now includes the moving volume density.
Both gradient and mass contacts load into orthogonal cell summands, so
their many-cell incidence is controlled by a supremum rather than the mesh
volume.  A finite-range action produces a uniformly bounded internal patch
and overlap on shape-regular meshes.  The remaining local analytic issue is
the metric dependence of the physical source and response insertion; the
remaining physical issue is a positive conditional gravitational carrier
with the required restricted norm.
\end{assessment}

\begin{programme}[Eighth-series V37: metric-dependent source replacement]
\label{prog:e8v36-V37}
V37 will treat sources of the form
\(u_g=R_gA\Omega_g\), where \(R_g\) is a local recovery or transfer loader.
It will derive a resolvent identity for the source replacement, separate
vacuum motion from recovery motion, and state the local gap or reduced
resolvent row needed to keep
\(\mathcal J_g^{\rm km}(u_g-u_{g'})\) inside a summable influence collar.
\end{programme}

\begin{firewall}[Claim boundary after eighth-series V36]
\label{fire:e8v36-boundary}
V36 proves the fourth-root determinant contact, constructs the combined
kinetic--mass coefficient loader, derives its heat-bath incidence, and
proves uniform internal-patch size and overlap for declared finite-range
actions on shape-regular meshes.

V36 does not prove finite range after constraint reduction, a positive
conditional gravitational transfer, metric-dependent source replacement,
complex and chiral contacts, the physical restricted norm, cutoff removal,
or the Standard Model--Einstein continuum.
\end{firewall}

\begin{theorem}[Eighth-series V36 result]
\label{thm:e8v36-result}
On a protected real metric box, the loader
\[
 \mathcal J_g^{\rm km}u
 =(G_g^{1/2}\nabla u,\,
   m_0(\det g)^{1/4}u)_K
\]
has norm exactly equal to the kinetic plus metric-volume mass energy.  Its
one-cell replacement incidence is bounded by
\(\chi_i^2/2\) times the local energy, with \(\chi_i\) in
\eqref{eq:e8v36-chi}.  For a fixed-range action and nodal \(P_1\) sources,
the internal metric patches have cofinally bounded size and overlap.
\end{theorem}

\begin{proof}
Use \Cref{lem:e8v36-w-Lipschitz,
thm:e8v36-combined-energy,thm:e8v36-one-cell,
cor:e8v36-weighted,thm:e8v36-patch,
thm:e8v36-overlap,thm:e8v36-KMIBC}.
\end{proof}

\paragraph{Parent-version record.}
V36 was formed from
\texttt{Renewal\_SM\_Einstein\_Continuum\_Research\_Notes\_V35.tex}.
The V35 parent had (13\,733) logical source lines, (503\,431) bytes,
and SHA--256
\[
 \texttt{16B1FFD2ADDE38A5233AFCCF19226525DCA26EDAE3E2107F4186D9F0D0B58752}.
\]
The next compulsory companion audit is V40.

% =============================================================================
% END APPENDED EIGHTH-SERIES V36 MATERIAL
% =============================================================================

% =============================================================================
% BEGIN APPENDED EIGHTH-SERIES V37 MATERIAL
% =============================================================================

\section{Moving vacua and metric-dependent recovered sources}
\label{sec:v8v37}

V36 left the source itself fixed while the kinetic--mass coefficient
loader moved with the metric.  A physical insertion generally has the
form
\begin{equation}
 u_g=R_gA_g\Omega_g,
 \label{eq:e8v37-source}
\end{equation}
where the vacuum, the observable, and the recovery map may all move.  This
note separates those four derivatives exactly.  It first proves the
finite-cutoff reduced-resolvent formula, then replaces the inadmissible
infinite-volume gap by a source-restricted inverse row, and finally turns
spatial decay of that row into a volume-independent replacement incidence.

\subsection{The finite-cutoff vacuum derivative}
\label{sec:e8v37-vacuum}

All Hilbert spaces in this section are complex.  The scalar product is
linear in its second entry.  Let \(I\subset\mathbb R\) be an interval and
let \(s\mapsto H_s\) be a \(C^1\) family of nonnegative self-adjoint
operators on a fixed Hilbert space \(\mathcal H\).  For the theorem below
it is enough that the family have a common domain, that \(P_s\) be \(C^1\)
in operator norm, and that \(\dot H_sP_s\) extend to a bounded operator.

\begin{theorem}[Exact derivative of an isolated vacuum projector]
\label{thm:e8v37-projector-derivative}
Suppose
\begin{equation}
 \ker H_s=\operatorname{Ran}P_s,
 \qquad
 H_s\big|_{\operatorname{Ran}Q_s}\geq\gamma_s Q_s,
 \qquad Q_s=I-P_s,
 \label{eq:e8v37-gap}
\end{equation}
where \(P_s\) has fixed finite rank and \(\gamma_s>0\).  Define the
reduced resolvent
\begin{equation}
 S_s=Q_s(H_s|_{\operatorname{Ran}Q_s})^{-1}Q_s.
 \label{eq:e8v37-S}
\end{equation}
Then
\begin{equation}
 \dot P_s=-S_s\dot H_sP_s-P_s\dot H_sS_s
 \label{eq:e8v37-Pdot}
\end{equation}
and
\begin{equation}
 \|\dot P_s\|
 \leq {\|Q_s\dot H_sP_s\|\over\gamma_s}.
 \label{eq:e8v37-Pdot-bound}
\end{equation}
Consequently,
\begin{equation}
 \|P_t-P_r\|
 \leq\int_r^t {\|Q_s\dot H_sP_s\|\over\gamma_s}\,ds
 \qquad(r<t).
 \label{eq:e8v37-P-finite}
\end{equation}
\end{theorem}

\begin{proof}
Differentiate \(H_sP_s=0\) and multiply on the left by \(S_s\).  Since
\(S_sH_s=Q_s\),
\begin{equation}
 Q_s\dot P_s=-S_s\dot H_sP_s.
 \label{eq:e8v37-QPdot}
\end{equation}
Differentiating \(P_s^2=P_s\) shows that the diagonal blocks of
\(\dot P_s\) vanish:
\begin{equation}
 P_s\dot P_sP_s=0,
 \qquad Q_s\dot P_sQ_s=0.
 \label{eq:e8v37-off-diagonal}
\end{equation}
The adjoint of \eqref{eq:e8v37-QPdot} supplies the other off-diagonal
block and proves \eqref{eq:e8v37-Pdot}.  Relative to
\(\operatorname{Ran}P_s\oplus\operatorname{Ran}Q_s\), the derivative has
block form
\(
 \left(\begin{smallmatrix}0&B_s^*\\B_s&0\end{smallmatrix}\right)
\)
with \(B_s=-S_s\dot H_sP_s\); hence \(\|\dot P_s\|=\|B_s\|\).
The gap gives \(\|S_s\|\leq\gamma_s^{-1}\), proving
\eqref{eq:e8v37-Pdot-bound}.  Integrate \(\dot P_s\) in operator norm to
obtain \eqref{eq:e8v37-P-finite}.
\end{proof}

\begin{corollary}[Parallel-gauge vacuum response]
\label{cor:e8v37-vacuum-vector}
If \(P_s\) has rank one, choose a normalized \(C^1\) vector
\(\Omega_s\) with
\begin{equation}
 P_s=|\Omega_s\rangle\langle\Omega_s|,
 \qquad \langle\Omega_s,\dot\Omega_s\rangle=0.
 \label{eq:e8v37-parallel-gauge}
\end{equation}
Then
\begin{equation}
 \dot\Omega_s=-S_sQ_s\dot H_s\Omega_s,
 \qquad
 \|\dot\Omega_s\|
 \leq {\|Q_s\dot H_s\Omega_s\|\over\gamma_s}.
 \label{eq:e8v37-Omegadot}
\end{equation}
In this gauge,
\begin{equation}
 \|\Omega_t-\Omega_r\|
 \leq\int_r^t
 {\|Q_s\dot H_s\Omega_s\|\over\gamma_s}\,ds.
 \label{eq:e8v37-Omega-finite}
\end{equation}
\end{corollary}

\begin{proof}
Differentiate \(H_s\Omega_s=0\) and project with \(Q_s\):
\(
 H_sQ_s\dot\Omega_s=-Q_s\dot H_s\Omega_s
\).
The parallel-gauge condition removes the \(P_s\)-component of
\(\dot\Omega_s\), so inversion on \(\operatorname{Ran}Q_s\) gives
\eqref{eq:e8v37-Omegadot}.  The last assertion follows by integration.
\end{proof}

\begin{remark}[Moving eigenvalue]
\label{rem:e8v37-moving-energy}
For \(H_s\Omega_s=E_s\Omega_s\), apply the result to
\(H_s-E_sI\).  The Feynman--Hellmann identity gives
\(\dot E_s=\langle\Omega_s,\dot H_s\Omega_s\rangle\), while
\(Q_s\dot E_s\Omega_s=0\).  Thus the right side of
\eqref{eq:e8v37-Omegadot} is unchanged.
\end{remark}

\begin{firewall}[A global gap is not a photon--graviton continuum row]
\label{fire:e8v37-global-gap}
The estimate above is a valid finite-cutoff or genuinely massive-sector
lemma.  It does not supply a volume-uniform gap for an unbroken photon or
a massless gravitational mode.  The exact torus calculation
\Cref{thm:e8v4-gap,ass:e8v4-order} shows why quotienting the constant mode
alone does not repair the infinite-volume limit.  Hence
\(\inf_s\gamma_s>0\) may not be inserted into the full
Standard Model--Einstein ledger without an infrared regulator, a protected
massive sector, or a different local response estimate.
\end{firewall}

\subsection{A source-restricted reduced inverse}
\label{sec:e8v37-restricted}

The derivative equation only asks for the inverse of \(H_s\) on the
particular vector created by the local metric perturbation.  This is
strictly weaker than bounded invertibility on the whole vacuum
complement.

\begin{definition}[Source-restricted reduced inverse, SRRI]
\label{def:e8v37-SRRI}
Let \(P_s=|\Omega_s\rangle\langle\Omega_s|\), \(Q_s=I-P_s\), and let
\(\mathcal B_s\subset Q_s\mathcal H\) be a linear space of declared
perturbation sources.  An SRRI is a linear map
\begin{equation}
 T_s:\mathcal B_s\longrightarrow
 \mathcal D(H_s)\cap Q_s\mathcal H
 \label{eq:e8v37-T-domain}
\end{equation}
such that
\begin{equation}
 H_sT_sb=b\quad(b\in\mathcal B_s),
 \qquad
 \ker(H_s|_{\mathcal D(H_s)\cap Q_s\mathcal H})=\{0\}.
 \label{eq:e8v37-T-inverse}
\end{equation}
No bound for \(T_s\) on all of \(Q_s\mathcal H\) is part of the
definition.
\end{definition}

\begin{theorem}[Vacuum derivative from an SRRI]
\label{thm:e8v37-SRRI-response}
Suppose \(s\mapsto\Omega_s\) is \(C^1\), satisfies the parallel gauge,
and
\begin{equation}
 b_s:=Q_s\dot H_s\Omega_s\in\mathcal B_s.
 \label{eq:e8v37-b}
\end{equation}
If an SRRI exists, then
\begin{equation}
 \dot\Omega_s=-T_sb_s.
 \label{eq:e8v37-SRRI-Omega}
\end{equation}
In particular, any sourcewise estimate for \(T_sb_s\) is directly a
vacuum-response estimate.
\end{theorem}

\begin{proof}
Differentiating \(H_s\Omega_s=0\) gives
\(H_s\dot\Omega_s=-\dot H_s\Omega_s\).  Projection onto the vacuum
complement and the parallel gauge give
\begin{equation}
 H_s\dot\Omega_s=-b_s,
 \qquad \dot\Omega_s\in Q_s\mathcal H.
\end{equation}
The vector \(-T_sb_s\) solves the same equation in the same subspace.
Their difference lies in the kernel displayed in
\eqref{eq:e8v37-T-inverse}, so it vanishes.
\end{proof}

\begin{definition}[Local SRRI row]
\label{def:e8v37-local-SRRI}
Let a bounded-degree cell graph carry orthogonal local projections
\(\chi_X\) on \(\mathcal H\).  A replacement source \(b_{s,i}\) has core
\(Y_i\) if \(\chi_{Y_i}b_{s,i}=b_{s,i}\).  The local SRRI row with
constants \(C_T,\mu>0\) is
\begin{equation}
 \|\chi_XT_sb_{s,i}\|
 \leq C_Te^{-\mu d(X,Y_i)}\|b_{s,i}\|
 \label{eq:e8v37-local-SRRI}
\end{equation}
for every cell set \(X\), with constants uniform in the cutoff and the
replacement path.
\end{definition}

\begin{remark}[What the local row does and does not assert]
\label{rem:e8v37-local-SRRI}
Equation \eqref{eq:e8v37-local-SRRI} is a hypothesis to be proved from a
finite-horizon, a Ward cancellation, a massive channel, or a suitable
weighted resolvent estimate.  It is not a consequence of source locality:
the inverse of a massless elliptic operator normally has a long-range
kernel.  The purpose of SRRI is to expose precisely this missing physical
input without disguising it as a global gap.
\end{remark}

\subsection{Recovery motion and the four-term source derivative}
\label{sec:e8v37-recovery}

At a fixed cutoff all spaces below are finite dimensional, so the products
are bounded.  Uniformity of their norms is a separate cofinal row.

\begin{lemma}[Exact recovery resolvent identity]
\label{lem:e8v37-recovery-resolvent}
Let \(K_s\) be a norm-\(C^1\) family of bounded operators and suppose
\(K_s+\alpha I\) is invertible.  Put
\begin{equation}
 R_s=(K_s+\alpha I)^{-1}.
 \label{eq:e8v37-R}
\end{equation}
Then
\begin{align}
 R_t-R_s&=-R_t(K_t-K_s)R_s,
 \label{eq:e8v37-resolvent-identity}\\
 \dot R_s&=-R_s\dot K_sR_s.
 \label{eq:e8v37-Rdot}
\end{align}
\end{lemma}

\begin{proof}
Insert the identity

\[
 A^{-1}-B^{-1}=A^{-1}(B-A)B^{-1}
\]

with \(A=K_t+\alpha I\) and \(B=K_s+\alpha I\).  Divide by \(t-s\) and
take the norm limit to obtain \eqref{eq:e8v37-Rdot}.
\end{proof}

\begin{theorem}[Exact coefficient-space derivative of a recovered source]
\label{thm:e8v37-four-term}
Let \(\mathcal C\) be the fixed coefficient Hilbert space used by the V36
loader.  Along a one-cell replacement path, suppose \(J_s:\mathcal H\to
\mathcal C\), \(A_s:\mathcal H\to\mathcal H\), and \(K_s\) are
norm-\(C^1\), let \(R_s\) be \eqref{eq:e8v37-R}, and assume the hypotheses
of \Cref{thm:e8v37-SRRI-response}.  Define
\begin{equation}
 j_s=J_su_s,
 \qquad u_s=R_sA_s\Omega_s.
 \label{eq:e8v37-j}
\end{equation}
Then
\begin{equation}
 \dot j_s=D_s^J+D_s^R+D_s^A+D_s^\Omega,
 \label{eq:e8v37-four-split}
\end{equation}
where
\begin{align}
 D_s^J&=\dot J_sR_sA_s\Omega_s,
 \label{eq:e8v37-DJ}\\
 D_s^R&=-J_sR_s\dot K_sR_sA_s\Omega_s,
 \label{eq:e8v37-DR}\\
 D_s^A&=J_sR_s\dot A_s\Omega_s,
 \label{eq:e8v37-DA}\\
 D_s^\Omega&=-J_sR_sA_sT_sb_s.
 \label{eq:e8v37-DOmega}
\end{align}
Consequently,
\begin{equation}
 \|j_1-j_0\|_{\mathcal C}
 \leq\int_0^1
 \bigl(\|D_s^J\|+\|D_s^R\|+
       \|D_s^A\|+\|D_s^\Omega\|\bigr)\,ds.
 \label{eq:e8v37-four-bound}
\end{equation}
\end{theorem}

\begin{proof}
Apply the product rule to \(J_sR_sA_s\Omega_s\).  Substitute
\eqref{eq:e8v37-Rdot} and \eqref{eq:e8v37-SRRI-Omega}; this gives the four
displayed terms.  The fundamental theorem of calculus in
\(\mathcal C\), followed by the triangle inequality, gives
\eqref{eq:e8v37-four-bound}.
\end{proof}

\begin{corollary}[Four independent replacement budgets]
\label{cor:e8v37-four-budgets}
If measurable functions \(c_J,c_R,c_A,c_\Omega\) satisfy
\begin{equation}
 \|D_s^X\|\leq c_X(s),
 \qquad X\in\{J,R,A,\Omega\},
 \label{eq:e8v37-cX}
\end{equation}
then the squared replacement debit obeys
\begin{equation}
 \|j_1-j_0\|^2
 \leq
 \left(\sum_X\int_0^1c_X(s)\,ds\right)^2.
 \label{eq:e8v37-four-budget}
\end{equation}
Each row can therefore be estimated by a different mechanism; in
particular, a local coefficient Lipschitz bound for \(J_s\) does not pay
for vacuum motion.
\end{corollary}

\begin{proof}
Insert \eqref{eq:e8v37-cX} into \eqref{eq:e8v37-four-bound} and square.
\end{proof}

\begin{firewall}[Domain passage is still required]
\label{fire:e8v37-domains}
The finite-cutoff product rule is exact.  In the continuum, the kinetic
loader and the Hamiltonian are unbounded and their domains move with the
metric.  A common-form-domain theorem or a graph-norm trivialization must
be proved before \eqref{eq:e8v37-four-split} can be used after cutoff
removal.  Finite-dimensional differentiability alone does not provide
that theorem.
\end{firewall}

\subsection{From a local response tail to uniform incidence}
\label{sec:e8v37-incidence}

Let \(\mathcal G=(\mathcal V,\mathcal E)\) be a graph of maximum degree
\(D\geq2\), and write
\(\mathcal C=\bigoplus_{K\in\mathcal V}\mathcal C_K\), with \(M_K\) the
orthogonal projection onto the \(K\)-summand.  Let \(Y_i\) be the core of
replacement \(i\), and suppose each cell belongs to at most \(L\) cores.

\begin{theorem}[Exponential response gives volume-independent incidence]
\label{thm:e8v37-tail-incidence}
Suppose the coefficient replacement
\(\delta_i=j_{1,i}-j_{0,i}\) satisfies
\begin{equation}
 \|M_K\delta_i\|
 \leq C_0e^{-\mu d(K,Y_i)}
 \label{eq:e8v37-cell-tail}
\end{equation}
for all cells and replacements.  If
\begin{equation}
 (D-1)e^{-2\mu}<1,
 \label{eq:e8v37-decay-threshold}
\end{equation}
then, for every \(K\),
\begin{equation}
 \sum_i\|M_K\delta_i\|^2
 \leq C_0^2L
 \left(1+{De^{-2\mu}\over1-(D-1)e^{-2\mu}}\right).
 \label{eq:e8v37-uniform-incidence}
\end{equation}
The bound is independent of the number of cells.
\end{theorem}

\begin{proof}
For every \(i\), choose a nearest cell \(y_i\in Y_i\) to \(K\).  At a
fixed cell \(y\), at most \(L\) indices can have \(y_i=y\), because each
such core contains \(y\).  A graph of maximum degree \(D\) has at most one
vertex at radius zero and at most \(D(D-1)^{r-1}\) vertices at radius
\(r\geq1\).  Squaring \eqref{eq:e8v37-cell-tail}, grouping the chosen
cells by radius, and summing the geometric series gives
\begin{align*}
 \sum_i\|M_K\delta_i\|^2
 &\leq C_0^2L\left(
 1+\sum_{r\geq1}D(D-1)^{r-1}e^{-2\mu r}\right)\\
 &=C_0^2L
 \left(1+{De^{-2\mu}\over1-(D-1)e^{-2\mu}}\right).
\end{align*}
Condition \eqref{eq:e8v37-decay-threshold} is exactly convergence of this
series.
\end{proof}

\begin{corollary}[Derivative-tail compiler]
\label{cor:e8v37-derivative-tail}
If, along every replacement path,
\begin{equation}
 \|M_KD_{s,i}^X\|
 \leq c_X(s)e^{-\mu d(K,Y_i)},
 \qquad X\in\{J,R,A,\Omega\},
 \label{eq:e8v37-derivative-tail}
\end{equation}
then \eqref{eq:e8v37-cell-tail} holds with
\begin{equation}
 C_0=\sum_X\int_0^1c_X(s)\,ds.
 \label{eq:e8v37-C0}
\end{equation}
Thus the V36 bounded core overlap and the SRRI vacuum tail together give a
cofinal Hilbert-valued replacement incidence whenever
\eqref{eq:e8v37-decay-threshold} holds.
\end{corollary}

\begin{proof}
Project \eqref{eq:e8v37-four-split} with \(M_K\), integrate in \(s\), and
use \eqref{eq:e8v37-derivative-tail}.  Apply
\Cref{thm:e8v37-tail-incidence}.
\end{proof}

\begin{remark}[Polynomial-growth improvement]
\label{rem:e8v37-polynomial}
On a graph with sphere growth \(N(r)\leq C(1+r)^q\), every
\(\mu>0\) makes the analogous series finite.  The threshold
\eqref{eq:e8v37-decay-threshold} is only the safe bounded-degree estimate;
it is not intrinsic to Euclidean shape-regular meshes.
\end{remark}

\subsection{The protected response card and the next bottleneck}
\label{sec:e8v37-card}

\begin{definition}[Protected source-response card, PSRC]
\label{def:e8v37-PSRC}
For every metric replacement path, PSRC consists of:
\begin{enumerate}[label=\textup{(P\arabic*)},leftmargin=4em]
\item a fixed Hilbert-space trivialization and a parallel vacuum gauge;
\item either the finite-cutoff gap row \eqref{eq:e8v37-gap} or an SRRI on
the actual perturbation sources \eqref{eq:e8v37-b};
\item a local response estimate strong enough for the graph-growth sum;
\item the recovery resolvent identity with cofinal bounds on \(R_s\) and
its differentiated insertion;
\item local derivative rows for \(J_s\) and \(A_s\);
\item the V36 bounded core size and overlap, with exterior constraint tails
retained; and
\item Ward, reflection, boundary, spin, chiral, density, and conditional
mean contacts for the same protected block.
\end{enumerate}
\end{definition}

\begin{theorem}[PSRC closes the metric-dependent source row conditionally]
\label{thm:e8v37-PSRC}
At a fixed cutoff, PSRC rows (P1)--(P6) imply the exact four-term source
replacement bound \eqref{eq:e8v37-four-bound}.  If their local tails obey
\eqref{eq:e8v37-derivative-tail} and the graph-growth sum is finite, the
resulting coefficient replacement has a cutoff-independent incidence
bound.  With (P7), this supplies KMIBC row (M6) and the corresponding
metric-dependent source entry in the V35 IMBSC handoff.
\end{theorem}

\begin{proof}
Rows (P1)--(P5) give \Cref{thm:e8v37-four-term}.  Row (P6) supplies \(L\)
in \Cref{thm:e8v37-tail-incidence}; the local rows supply \(C_0\).
Therefore \eqref{eq:e8v37-uniform-incidence} is cofinal.  Row (P7) lists
the remaining compatibility contacts required by KMIBC and IMBSC.
\end{proof}

\begin{assessment}[Progress at V37]
\label{ass:e8v37-status}
Metric-dependent source motion is no longer a single opaque debit.  The
loader, recovery, observable, and vacuum contributions are separated by
an exact identity.  Vacuum motion is controlled either by an honest
finite-cutoff reduced resolvent or by an explicitly source-restricted
inverse.  Exponential local response, together with the V36 overlap row,
now has a proved volume-independent incidence compiler.  What remains is
to construct such a local response without assuming a false global gap in
the massless sectors.
\end{assessment}

\begin{programme}[Eighth-series V38: finite-horizon reduced response]
\label{prog:e8v37-V38}
V38 will go upstream of the missing SRRI row.  It will replace the full
inverse by the truncated heat response
\(
 T_T=\int_0^T e^{-tH}Q\,dt
\), compute its exact residual, and derive a local source estimate with an
explicit finite-horizon tail.  The residual and the propagation tail will
remain separate so that no equilibrium gap is smuggled into the massless
branch.
\end{programme}

\begin{firewall}[Claim boundary after eighth-series V37]
\label{fire:e8v37-boundary}
V37 proves the isolated-projector derivative, the parallel-gauge vacuum
formula, the exact recovery resolvent identity, the four-term recovered
source derivative, and the graph-growth incidence theorem.

V37 does not prove an SRRI or its decay for the interacting
Standard Model--Einstein Hamiltonian, a cofinal common-domain theorem, a
positive gravitational conditional carrier, the chiral and reflection
contacts, cutoff removal, or the full continuum.
\end{firewall}

\begin{theorem}[Eighth-series V37 result]
\label{thm:e8v37-result}
For a differentiable finite-cutoff family with a simple vacuum, the
coefficient source \(j_s=J_sR_sA_s\Omega_s\) satisfies the exact split
\eqref{eq:e8v37-four-split}.  A source-restricted reduced inverse is
sufficient for its vacuum term.  If all four derivatives decay from
bounded-overlap replacement cores as in
\eqref{eq:e8v37-derivative-tail}, then their squared replacement incidence
is volume independent under the graph-growth condition
\eqref{eq:e8v37-decay-threshold}.
\end{theorem}

\begin{proof}
Use \Cref{thm:e8v37-SRRI-response,lem:e8v37-recovery-resolvent,
thm:e8v37-four-term,cor:e8v37-derivative-tail}.
\end{proof}

\paragraph{Parent-version record.}
V37 was formed from
\texttt{Renewal\_SM\_Einstein\_Continuum\_Research\_Notes\_V36.tex}.
The V36 parent had (14\,211) logical source lines, (520\,330) bytes,
and SHA--256
\[
 \texttt{DB0E614A25AA7C71A40189AF6E818F11BA726456B587BFC8D0EA770B4F0C61D3}.
\]
The next compulsory companion audit is V40.

% =============================================================================
% END APPENDED EIGHTH-SERIES V37 MATERIAL
% =============================================================================

% =============================================================================
% BEGIN APPENDED EIGHTH-SERIES V38 MATERIAL
% =============================================================================

\section{Finite-horizon reduced response without a global gap}
\label{sec:v8v38}

V37 isolated the vacuum term
\(-JRA\,T(Q\dot H\Omega)\) and left the source-restricted inverse \(T\)
as a physical input.  In a massless sector the full inverse can be
nonlocal or unbounded.  This note replaces it by a finite-horizon heat
response.  The price is an exact terminal residual.  An infrared
cancellation row makes that residual small, while a Davies--Gaffney row
keeps the response local for every finite horizon.

\subsection{Truncated heat inversion}
\label{sec:e8v38-truncated}

Let \(H\geq0\) be self-adjoint on a Hilbert space \(\mathcal H\), let
\(P\) be its vacuum projection, and put \(Q=I-P\).  Since \(P\) is a
spectral projection of \(H\), it commutes with the heat semigroup.

\begin{definition}[Finite-horizon reduced response]
\label{def:e8v38-TT}
For \(T>0\), define the bounded operator
\begin{equation}
 \mathcal T_T
 :=\int_0^T e^{-tH}Q\,dt,
 \label{eq:e8v38-TT}
\end{equation}
where the integral is a strong Bochner integral.
\end{definition}

\begin{theorem}[Exact response and terminal residual]
\label{thm:e8v38-residual}
The operator \(\mathcal T_T\) is self-adjoint, nonnegative, commutes with
\(H,P,Q\), and satisfies
\begin{align}
 \|\mathcal T_T\|&\leq T,
 \label{eq:e8v38-TT-norm}\\
 H\mathcal T_Tb
 &=Qb-e^{-TH}Qb
 \qquad(b\in\mathcal H),
 \label{eq:e8v38-exact-residual}\\
 \mathcal T_THb
 &=Qb-e^{-TH}Qb
 \qquad(b\in\mathcal D(H)).
 \label{eq:e8v38-other-residual}
\end{align}
Moreover, \(\mathcal T_Tb\in\mathcal D(H)\) for every \(b\in\mathcal H\).
\end{theorem}

\begin{proof}
By the spectral theorem,
\begin{equation}
 \mathcal T_T
 =Q f_T(H),
 \qquad
 f_T(\lambda)=
 \begin{cases}
 (1-e^{-T\lambda})/\lambda,&\lambda>0,\\
 T,&\lambda=0.
 \end{cases}
 \label{eq:e8v38-fT}
\end{equation}
The factor \(Q\) removes the value at zero.  Since
\(0\leq f_T(\lambda)\leq T\), the first assertions and
\eqref{eq:e8v38-TT-norm} follow.  Also
\(\lambda f_T(\lambda)=1-e^{-T\lambda}\), which is bounded.  Thus the
range of \(\mathcal T_T\) lies in \(\mathcal D(H)\) and spectral
calculus gives \eqref{eq:e8v38-exact-residual}.  Commutation with \(H\)
gives \eqref{eq:e8v38-other-residual}.
\end{proof}

\begin{corollary}[Energy bound for the truncated response]
\label{cor:e8v38-energy}
For every \(b\in\mathcal H\),
\begin{align}
 \|H^{1/2}\mathcal T_Tb\|^2
 &\leq T\|Qb\|^2,
 \label{eq:e8v38-energy}\\
 \langle b,\mathcal T_Tb\rangle
 &=\int_0^T\|e^{-tH/2}Qb\|^2\,dt.
 \label{eq:e8v38-positive}
\end{align}
\end{corollary}

\begin{proof}
For \(x\geq0\),
\begin{equation}
 (1-e^{-x})^2\leq x.
 \label{eq:e8v38-scalar-energy}
\end{equation}
Indeed, for \(0\leq x\leq1\) the left side is at most \(x^2\), and for
\(x\geq1\) it is at most \(1\).  With \(x=T\lambda\),
\[
 \lambda f_T(\lambda)^2
 ={(1-e^{-T\lambda})^2\over\lambda}
 \leq T.
\]
Integrating against the spectral measure of \(Qb\) proves
\eqref{eq:e8v38-energy}.  The semigroup identity
\(\langle Qb,e^{-tH}Qb\rangle=\|e^{-tH/2}Qb\|^2\), followed by integration,
proves \eqref{eq:e8v38-positive}.
\end{proof}

\begin{remark}[The response is not an inverse]
\label{rem:e8v38-not-inverse}
Equation \eqref{eq:e8v38-exact-residual} is the exact statement.  Replacing
its right side by \(Qb\) discards \(e^{-TH}Qb\) and silently takes an
infinite-time limit.  No such limit is justified in a massless sector.
\end{remark}

\subsection{Infrared cancellation pays the terminal residual}
\label{sec:e8v38-cancellation}

\begin{theorem}[Fractional-range residual bound]
\label{thm:e8v38-fractional}
Let \(\alpha>0\), let \(c\in\mathcal D(H^\alpha)\cap Q\mathcal H\), and
set
\begin{equation}
 b=H^\alpha c.
 \label{eq:e8v38-fractional-source}
\end{equation}
Then
\begin{equation}
 \|e^{-TH}b\|
 \leq
 \left({\alpha\over eT}\right)^\alpha\|c\|.
 \label{eq:e8v38-fractional-residual}
\end{equation}
For \(\alpha=1\),
\begin{equation}
 \mathcal T_THc=(I-e^{-TH})c,
 \qquad
 \|\mathcal T_THc-c\|\leq {1\over eT}\|c\|.
 \label{eq:e8v38-range-one}
\end{equation}
\end{theorem}

\begin{proof}
Spectral calculus gives
\[
 \|e^{-TH}H^\alpha c\|
 \leq\sup_{\lambda\geq0}\lambda^\alpha e^{-T\lambda}\,\|c\|.
\]
The scalar function has its maximum at \(\lambda=\alpha/T\), with value
\((\alpha/(eT))^\alpha\), proving
\eqref{eq:e8v38-fractional-residual}.  When \(\alpha=1\),
\eqref{eq:e8v38-other-residual} and \(Qc=c\) give the identity; apply the
same maximum with \(\alpha=1\).
\end{proof}

\begin{definition}[Infrared source order]
\label{def:e8v38-infrared-order}
A family of perturbation sources \(b_n\in Q_n\mathcal H_n\) has uniform
infrared order \(\alpha>0\) if
\begin{equation}
 b_n=H_n^\alpha c_n,
 \qquad
 \sup_n\|c_n\|<\infty.
 \label{eq:e8v38-infrared-order}
\end{equation}
The vectors \(c_n\), their norms, and the quotient defining \(Q_n\) are
part of the row.
\end{definition}

\begin{corollary}[Cofinal terminal control]
\label{cor:e8v38-cofinal-residual}
If \eqref{eq:e8v38-infrared-order} holds and \(T_n\to\infty\), then
\begin{equation}
 \|e^{-T_nH_n}b_n\|=O(T_n^{-\alpha}).
 \label{eq:e8v38-cofinal-terminal}
\end{equation}
If the scale-\(n\) terminal debit is multiplied by \(a_n\geq0\), it is
summable whenever
\begin{equation}
 \sum_n a_nT_n^{-\alpha}<\infty.
 \label{eq:e8v38-terminal-sum}
\end{equation}
\end{corollary}

\begin{proof}
Apply \Cref{thm:e8v38-fractional} uniformly and then the comparison test.
\end{proof}

\begin{firewall}[Cancellation must be structural]
\label{fire:e8v38-cancellation}
The representation \(b=H^\alpha c\) is not implied by
\(\langle\Omega,b\rangle=0\).  Orthogonality removes an exact zero mode,
whereas infrared order controls the spectral density near zero.  A Ward
identity, a commutator formula, or a direct spectral estimate must produce
\eqref{eq:e8v38-infrared-order}.
\end{firewall}

\subsection{Finite-horizon locality}
\label{sec:e8v38-locality}

Let the Hilbert space carry local projections \(\chi_X\) indexed by cell
sets.  The distance between \(X\) and \(Y\) is the cell-graph distance.

\begin{definition}[Davies--Gaffney heat row]
\label{def:e8v38-DG}
The heat family has propagation constant \(c_{\rm h}>0\) if, for disjoint
cell sets \(X,Y\) and \(t>0\),
\begin{equation}
 \|\chi_Xe^{-tH}Q\chi_Y\|
 \leq
 \exp\left\{-{d(X,Y)^2\over4c_{\rm h}t}\right\}.
 \label{eq:e8v38-DG}
\end{equation}
\end{definition}

\begin{theorem}[Gaussian locality of finite-horizon response]
\label{thm:e8v38-local-response}
Under \eqref{eq:e8v38-DG}, if \(b\in Q\mathcal H\) and
\(\chi_Yb=b\), then for every \(X\) at distance \(R=d(X,Y)>0\),
\begin{equation}
 \|\chi_X\mathcal T_Tb\|
 \leq
 T\exp\left\{-{R^2\over4c_{\rm h}T}\right\}\|b\|.
 \label{eq:e8v38-local-response}
\end{equation}
\end{theorem}

\begin{proof}
For \(0<t\leq T\), one has
\[
 \exp\left\{-{R^2\over4c_{\rm h}t}\right\}
 \leq
 \exp\left\{-{R^2\over4c_{\rm h}T}\right\}.
\]
Insert \(\chi_Yb=b\) in \eqref{eq:e8v38-TT}, apply
\eqref{eq:e8v38-DG}, and integrate this uniform bound over \(0<t<T\).
\end{proof}

\begin{remark}[The short-time singular endpoint is harmless]
\label{rem:e8v38-short-time}
The integrand in \eqref{eq:e8v38-DG} tends to zero faster than any power
as \(t\downarrow0\) when \(R>0\).  The simpler upper bound used in
\eqref{eq:e8v38-local-response} loses sharpness but introduces no
endpoint divergence.
\end{remark}

\begin{theorem}[Gaussian-tail incidence on bounded-degree graphs]
\label{thm:e8v38-Gaussian-incidence}
Let the replacement cores \(Y_i\) have overlap at most \(L\) on a graph
of maximum degree \(D\geq2\).  Suppose coefficient replacements satisfy
\begin{equation}
 \|M_K\delta_i\|
 \leq C_0\exp\{-a\,d(K,Y_i)^2\},
 \qquad a>0.
 \label{eq:e8v38-Gaussian-cell}
\end{equation}
Then, for every cell \(K\),
\begin{equation}
 \sum_i\|M_K\delta_i\|^2
 \leq C_0^2L\,\mathfrak G_D(a),
 \label{eq:e8v38-Gaussian-incidence}
\end{equation}
where
\begin{equation}
 \mathfrak G_D(a)
 :=
 1+D\sum_{r=1}^\infty(D-1)^{r-1}e^{-2ar^2}<\infty.
 \label{eq:e8v38-GD}
\end{equation}
Thus no lower threshold on \(a\) is needed at a fixed horizon.
\end{theorem}

\begin{proof}
Choose a nearest \(y_i\in Y_i\) to \(K\), as in
\Cref{thm:e8v37-tail-incidence}.  At radius \(r\) there are at most
\(LD(D-1)^{r-1}\) choices for \(r\geq1\), and at most \(L\) at radius
zero.  Squaring \eqref{eq:e8v38-Gaussian-cell} gives the displayed
series.  Its ratio satisfies
\begin{equation}
 {(D-1)^{r}e^{-2a(r+1)^2}
  \over (D-1)^{r-1}e^{-2ar^2}}
 =(D-1)e^{-2a(2r+1)}\longrightarrow0,
\end{equation}
so the series converges for every \(a>0\).
\end{proof}

\begin{corollary}[Vacuum-response incidence at time \(T\)]
\label{cor:e8v38-vacuum-incidence}
Assume the vacuum source for replacement \(i\) is supported in \(Y_i\),
has norm at most \(B_*\), and all bounded maps following
\(\mathcal T_T\) in the vacuum term of
\eqref{eq:e8v37-DOmega} have local norm product at most \(C_*\).
Under the Davies--Gaffney row,
\begin{equation}
 \sum_i\|M_KD_i^\Omega\|^2
 \leq
 C_*^2B_*^2T^2L\,
 \mathfrak G_D\left({1\over4c_{\rm h}T}\right).
 \label{eq:e8v38-vacuum-incidence}
\end{equation}
The constant is independent of the total volume, though it need not be
uniform as \(T\to\infty\).
\end{corollary}

\begin{proof}
Apply \Cref{thm:e8v38-local-response} cellwise.  This gives
\eqref{eq:e8v38-Gaussian-cell} with
\(C_0=C_*B_*T\) and \(a=(4c_{\rm h}T)^{-1}\).  Invoke
\Cref{thm:e8v38-Gaussian-incidence}.
\end{proof}

\begin{firewall}[The vacuum projection can be nonlocal]
\label{fire:e8v38-Q-locality}
The Davies--Gaffney row in \eqref{eq:e8v38-DG} contains \(Q\).
Even when \(\dot H_i\) is local, the vector
\(Q\dot H_i\Omega\) need not be supported in the same cell patch.
Pinned vacuum energy gives
\(\langle\Omega,\dot H_i\Omega\rangle=0\), but entanglement still prevents
one from identifying Hilbert-space support with operator support.  The
source-support row in \Cref{thm:e8v38-local-response} must therefore be
proved in the chosen excitation representation.
\end{firewall}

\subsection{Balancing terminal accuracy and propagation length}
\label{sec:e8v38-balance}

Longer horizons reduce the terminal residual but enlarge the response
collar.  The two effects must be scheduled together.

\begin{theorem}[Finite-horizon cofinal schedule]
\label{thm:e8v38-schedule}
At scale \(n\), suppose:
\begin{enumerate}[label=\textup{(\roman*)}]
\item the source has infrared order \(\alpha>0\) with
\(\|c_n\|\leq C_{\rm ir}\);
\item the protected observable is separated from the replacement core by
\(R_n\);
\item the source norm is at most \(B_n\); and
\item the relevant heat row has the same \(c_{\rm h}\).
\end{enumerate}
Then the terminal and separated-propagation debits are bounded by
\begin{align}
 \varepsilon_n^{\rm term}
 &\leq
 C_{\rm ir}\left({\alpha\over eT_n}\right)^\alpha,
 \label{eq:e8v38-schedule-terminal}\\
 \varepsilon_n^{\rm prop}
 &\leq
 B_nT_n\exp\left\{-{R_n^2\over4c_{\rm h}T_n}\right\}.
 \label{eq:e8v38-schedule-prop}
\end{align}
In particular, if
\begin{equation}
 T_n=2^{\theta n},
 \qquad R_n=2^{\rho n},
 \qquad 0<\theta<2\rho,
 \label{eq:e8v38-dyadic-schedule}
\end{equation}
and \(B_n\) grows at most exponentially, then both series
\(\sum_n\varepsilon_n^{\rm term}\) and
\(\sum_n\varepsilon_n^{\rm prop}\) converge.
\end{theorem}

\begin{proof}
The first estimate is \Cref{thm:e8v38-fractional}; the second is
\Cref{thm:e8v38-local-response}.  Under
\eqref{eq:e8v38-dyadic-schedule}, the terminal sequence is a geometric
multiple of \(2^{-\alpha\theta n}\).  Also
\[
 {R_n^2\over T_n}=2^{(2\rho-\theta)n}.
\]
The exponential in \eqref{eq:e8v38-schedule-prop} therefore decays faster
than \(\exp(-c\,2^{\eta n})\) for some \(c,\eta>0\), which dominates any
exponential growth of \(B_nT_n\).
\end{proof}

\begin{remark}[Fixed physical volume comes first]
\label{rem:e8v38-order}
The schedule controls a sequence of separated local corrections.  It does
not make the equilibrium heat-block time uniform in the infinite physical
volume.  The order-of-limits restriction in
\Cref{ass:e8v4-order,fire:e8v4-equilibrium} remains in force.
\end{remark}

\subsection{Finite-horizon response card}
\label{sec:e8v38-card}

\begin{definition}[Finite-horizon response card, FHRC]
\label{def:e8v38-FHRC}
FHRC consists of:
\begin{enumerate}[label=\textup{(H\arabic*)},leftmargin=4em]
\item a nonnegative self-adjoint cutoff Hamiltonian and its declared
vacuum or Ward projection;
\item the exact residual identity \eqref{eq:e8v38-exact-residual};
\item a structural infrared-order representation for every moving-vacuum
source;
\item a Davies--Gaffney or stronger finite-horizon locality row in the
same excitation representation;
\item a scale schedule making both
\eqref{eq:e8v38-schedule-terminal} and
\eqref{eq:e8v38-schedule-prop} summable;
\item the V36 core-overlap and V37 four-derivative budgets; and
\item common-domain, conditional-density, Ward, bridge, boundary, spin,
chiral, reflection, and gravitational-carrier contacts.
\end{enumerate}
\end{definition}

\begin{theorem}[FHRC supplies a gap-free approximate PSRC row]
\label{thm:e8v38-FHRC}
FHRC rows (H1)--(H6) replace the SRRI in PSRC by
\(\mathcal T_{T_n}\).  The inverse defect is exactly the terminal source
\(e^{-T_nH_n}b_n\), and the response exterior is bounded by the Gaussian
tail in \eqref{eq:e8v38-schedule-prop}.  If the two debit series are
summable, their cumulative source error converges without a
volume-uniform spectral gap.
\end{theorem}

\begin{proof}
The residual statement is \Cref{thm:e8v38-residual}; its summability is
\Cref{cor:e8v38-cofinal-residual}.  The exterior estimate and its
summability are
\Cref{thm:e8v38-local-response,thm:e8v38-schedule}.  The internal
replacement incidence is
\Cref{thm:e8v38-Gaussian-incidence,cor:e8v38-vacuum-incidence}.
Summability makes the partial corrected sources Cauchy by the triangle
inequality.
\end{proof}

\begin{assessment}[Progress at V38]
\label{ass:e8v38-status}
The missing reduced inverse now has a gap-free finite-horizon substitute.
Its failure to invert is an exact terminal heat source, not an unspecified
remainder.  Fractional infrared order makes that source decay
polynomially in the horizon, while finite-time heat locality gives a
Gaussian response collar and a volume-independent incidence constant.
The new bottleneck is structural: deriving positive infrared order for
the actual metric perturbation of the vacuum.
\end{assessment}

\begin{programme}[Eighth-series V39: commutator and gauge-tangent sources]
\label{prog:e8v38-V39}
V39 will test whether a perturbation is a pure transport tangent.  For a
unitary conjugation \(H_s=U_sHU_s^*\), differentiating the vacuum equation
should put \(Q\dot H\Omega\) in \(\operatorname{Ran}H\) exactly.  The note
will separate this gauge or frame motion from a genuine physical metric
deformation and identify the remainder that still requires the
finite-horizon response.
\end{programme}

\begin{firewall}[Claim boundary after eighth-series V38]
\label{fire:e8v38-boundary}
V38 proves exact finite-horizon inversion, fractional residual decay,
Gaussian response locality conditional on a Davies--Gaffney row, its
bounded-degree incidence compiler, and a compatible cofinal schedule.

V38 does not prove infrared order or excitation support for the
interacting metric source, the Davies--Gaffney row for the constrained
Standard Model--Einstein Hamiltonian, a cofinal common-domain theorem, the
physical conditional carrier, cutoff removal, or the full continuum.
\end{firewall}

\begin{theorem}[Eighth-series V38 result]
\label{thm:e8v38-result}
For every nonnegative self-adjoint \(H\), the truncated response
\(\mathcal T_T=\int_0^Te^{-tH}Q\,dt\) obeys
\[
 H\mathcal T_T=Q-e^{-TH}Q.
\]
If a local source has uniform infrared order \(\alpha>0\), the residual is
\(O(T^{-\alpha})\).  Under a Davies--Gaffney heat row, the response at
distance \(R\) is at most
\(T\exp\{-R^2/(4c_{\rm h}T)\}\) times the source norm, and its squared
replacement incidence is volume independent for every finite \(T\).
\end{theorem}

\begin{proof}
Use \Cref{thm:e8v38-residual,thm:e8v38-fractional,
thm:e8v38-local-response,thm:e8v38-Gaussian-incidence}.
\end{proof}

\paragraph{Parent-version record.}
V38 was formed from
\texttt{Renewal\_SM\_Einstein\_Continuum\_Research\_Notes\_V37.tex}.
The V37 parent had (14\,748) logical source lines, (539\,790) bytes,
and SHA--256
\[
 \texttt{02DD87D14AB47B822468E61AD05A1CFBE1E81E8ECAB4F33EC80C0CA65FAD89A9}.
\]
The next compulsory companion audit is V40.

% =============================================================================
% END APPENDED EIGHTH-SERIES V38 MATERIAL
% =============================================================================

% =============================================================================
% BEGIN APPENDED EIGHTH-SERIES V39 MATERIAL
% =============================================================================

\section{Gauge-tangent subtraction and the physical vacuum source}
\label{sec:v8v39}

The V38 infrared row becomes automatic with exponent one when a moving
Hamiltonian is obtained by unitary conjugation.  Such motion includes
changes of representation and may include regulated gauge, frame, or
diffeomorphism tangents.  It must be separated from a genuine deformation
of the physical geometry.  This note proves the exact separation and
shows that a covariantly transported recovered source is constant along a
pure conjugation orbit.

\subsection{Correction to the V38 response-vector estimate}
\label{sec:e8v39-correction}

\begin{firewall}[Erratum: equation residual versus vector convergence]
\label{fire:e8v39-V38-correction}
The identity in \eqref{eq:e8v38-range-one},
\[
 \mathcal T_THc=(I-e^{-TH})c,
\]
is correct.  The vector-norm bound stated immediately after that identity
in V38 is not valid without additional infrared control on \(c\).  The
correct general statements are
\begin{align}
 \|\mathcal T_THc-c\|
 &=\|e^{-TH}c\|\leq\|c\|,
 \label{eq:e8v39-correct-vector}\\
 \|H(\mathcal T_THc-c)\|
 &=\|e^{-TH}Hc\|
 \leq {1\over eT}\|c\|.
 \label{eq:e8v39-correct-equation}
\end{align}
The second line controls the differentiated equation residual; it does
not imply uniform convergence in the Hilbert norm because spectral mass
of \(c\) may accumulate at zero.  Accordingly, the V38 FHRC produces an
approximate solution of the vacuum derivative equation with a summable
equation defect.  Convergence of the response vectors themselves still
requires a reduced stability estimate, an observable seminorm that kills
the infrared sector, or stronger spectral order for \(c\).
\end{firewall}
\subsection{Exact unitary tangent}
\label{sec:e8v39-unitary}

At a fixed cutoff, let \(s\mapsto U_s\) be a norm-\(C^1\) unitary family
on \(\mathcal H\), and define its right generator
\begin{equation}
 G_s=\dot U_sU_s^*.
 \label{eq:e8v39-G}
\end{equation}
Then \(G_s^*=-G_s\).  Let \(H_0\geq0\) have normalized vacuum
\(\Omega_0\), and set
\begin{equation}
 H_s=U_sH_0U_s^*,
 \qquad
 \Omega_s=U_s\Omega_0,
 \qquad
 P_s=|\Omega_s\rangle\langle\Omega_s|.
 \label{eq:e8v39-orbit}
\end{equation}

\begin{lemma}[Commutator derivative of a conjugation orbit]
\label{lem:e8v39-commutator}
The family in \eqref{eq:e8v39-orbit} satisfies
\begin{align}
 \dot H_s&=[G_s,H_s],
 \label{eq:e8v39-Hdot-commutator}\\
 \dot\Omega_s&=G_s\Omega_s,
 \label{eq:e8v39-Omegadot-generator}\\
 \dot P_s&=[G_s,P_s].
 \label{eq:e8v39-Pdot-generator}
\end{align}
\end{lemma}

\begin{proof}
Differentiate \(U_sU_s^*=I\) to obtain
\(\dot U_s^*=-U_s^*\dot U_sU_s^*=-U_s^*G_s\).  Hence
\[
 \dot H_s
 =G_sU_sH_0U_s^*-U_sH_0U_s^*G_s
 =[G_s,H_s].
\]
The other two identities follow by differentiating
\(\Omega_s=U_s\Omega_0\) and
\(P_s=U_sP_0U_s^*\).
\end{proof}

\begin{remark}[Parallel gauge]
\label{rem:e8v39-parallel}
The scalar \(\langle\Omega_s,G_s\Omega_s\rangle\) is purely imaginary.
Multiplying \(U_s\) by a scalar phase replaces \(G_s\) by
\(G_s-i\theta_s'I\).  Along an interval the phase may be chosen so that
\begin{equation}
 \langle\Omega_s,G_s\Omega_s\rangle=0,
 \qquad Q_sG_s\Omega_s=G_s\Omega_s.
 \label{eq:e8v39-parallel}
\end{equation}
This is exactly the parallel gauge in
\eqref{eq:e8v37-parallel-gauge}.  Holonomy around a closed loop is not
removed by this local choice.
\end{remark}

\begin{theorem}[A pure unitary tangent has infrared order one]
\label{thm:e8v39-order-one}
Under the parallel gauge, the moving-vacuum source
\begin{equation}
 b_s:=Q_s\dot H_s\Omega_s
 \label{eq:e8v39-b}
\end{equation}
satisfies
\begin{equation}
 b_s=-H_sG_s\Omega_s.
 \label{eq:e8v39-b-range}
\end{equation}
Thus it has infrared order one with \(c_s=-G_s\Omega_s\).  Moreover, for
the V38 truncated response,
\begin{align}
 -\mathcal T_{T,s}b_s
 &=(I-e^{-TH_s})G_s\Omega_s,
 \label{eq:e8v39-truncated-gauge}\\
 \|H_s(-\mathcal T_{T,s}b_s)+b_s\|
 &\leq {1\over eT}\|G_s\Omega_s\|.
 \label{eq:e8v39-gauge-residual}
\end{align}
\end{theorem}

\begin{proof}
Use \(H_s\Omega_s=0\) in
\eqref{eq:e8v39-Hdot-commutator}:
\[
 Q_s\dot H_s\Omega_s
 =Q_s(G_sH_s-H_sG_s)\Omega_s
 =-H_sQ_sG_s\Omega_s.
\]
Equation \eqref{eq:e8v39-parallel} proves
\eqref{eq:e8v39-b-range}.  The first response identity follows from
\eqref{eq:e8v38-range-one}.  By \Cref{thm:e8v38-residual},
\[
 H_s(-\mathcal T_{T,s}b_s)+b_s=e^{-TH_s}b_s
 =-e^{-TH_s}H_sG_s\Omega_s.
\]
The scalar bound in \Cref{thm:e8v38-fractional} with \(\alpha=1\) gives
\eqref{eq:e8v39-gauge-residual}.
\end{proof}

\begin{corollary}[No inverse is needed on an identified pure tangent]
\label{cor:e8v39-no-inverse}
If \(G_s\Omega_s\) is known directly, the exact vacuum derivative is
\(\dot\Omega_s=G_s\Omega_s\), so the reduced resolvent is unnecessary.
The truncated response solves the same differentiated vacuum equation up
to the terminal equation defect in \eqref{eq:e8v39-gauge-residual}; it
need not converge to the derivative in Hilbert norm without an additional
infrared stability row.
\end{corollary}

\begin{proof}
This is \eqref{eq:e8v39-Omegadot-generator}; the reconstruction statement
is \eqref{eq:e8v39-truncated-gauge}.
\end{proof}

\subsection{Gauge tangent plus physical deformation}
\label{sec:e8v39-split}

For a general differentiable family, choose a declared anti-self-adjoint
connection \(G_s\) and define the covariant Hamiltonian derivative
\begin{equation}
 V_s:=\dot H_s-[G_s,H_s].
 \label{eq:e8v39-V}
\end{equation}
This is an identity after \(G_s\) has been chosen; its value depends on
that choice.

\begin{theorem}[Exact gauge--physical source split]
\label{thm:e8v39-source-split}
Let \(H_s\Omega_s=0\), use the parallel phase
\(\langle\Omega_s,\dot\Omega_s\rangle=0\), and put
\begin{equation}
 q_s=Q_sG_s\Omega_s,
 \qquad
 v_s=Q_sV_s\Omega_s.
 \label{eq:e8v39-qv}
\end{equation}
Then
\begin{align}
 b_s:=Q_s\dot H_s\Omega_s
 &=-H_sq_s+v_s,
 \label{eq:e8v39-b-split}\\
 H_s(\dot\Omega_s-q_s)
 &=-v_s.
 \label{eq:e8v39-covariant-vacuum}
\end{align}
For every finite horizon \(T\), the approximate derivative
\begin{equation}
 \dot\Omega_s^{(T)}
 :=q_s-\mathcal T_{T,s}v_s
 \label{eq:e8v39-approx-derivative}
\end{equation}
has the exact equation defect
\begin{equation}
 H_s\dot\Omega_s^{(T)}+b_s
 =e^{-TH_s}v_s.
 \label{eq:e8v39-approx-defect}
\end{equation}
\end{theorem}

\begin{proof}
From \eqref{eq:e8v39-V},
\[
 Q_s\dot H_s\Omega_s
 =Q_s(G_sH_s-H_sG_s+V_s)\Omega_s
 =-H_sQ_sG_s\Omega_s+Q_sV_s\Omega_s,
\]
which is \eqref{eq:e8v39-b-split}.  Differentiating
\(H_s\Omega_s=0\) gives \(H_s\dot\Omega_s=-b_s\); subtract
\(H_sq_s\) and use the split to obtain
\eqref{eq:e8v39-covariant-vacuum}.  Finally,
\Cref{thm:e8v38-residual} gives
\(H_s\mathcal T_{T,s}v_s=v_s-e^{-TH_s}v_s\).
Substitute this and \eqref{eq:e8v39-b-split} into the left side of
\eqref{eq:e8v39-approx-defect}.
\end{proof}

\begin{corollary}[Only the physical source pays the infrared residual]
\label{cor:e8v39-physical-residual}
If
\begin{equation}
 v_s=H_s^\alpha c_s,
 \qquad \alpha>0,
 \label{eq:e8v39-v-order}
\end{equation}
then
\begin{equation}
 \|H_s\dot\Omega_s^{(T)}+b_s\|
 \leq
 \left({\alpha\over eT}\right)^\alpha\|c_s\|.
 \label{eq:e8v39-physical-bound}
\end{equation}
No norm of the gauge tangent \(q_s\) enters this terminal debit.
\end{corollary}

\begin{proof}
Combine \eqref{eq:e8v39-approx-defect} with
\Cref{thm:e8v38-fractional}.
\end{proof}

\begin{corollary}[Pinned vacuum energy removes the scalar physical part]
\label{cor:e8v39-FH}
If the vacuum energy is identically zero, then
\begin{equation}
 \langle\Omega_s,V_s\Omega_s\rangle=0,
 \qquad v_s=V_s\Omega_s.
 \label{eq:e8v39-V-orthogonal}
\end{equation}
\end{corollary}

\begin{proof}
Feynman--Hellmann gives
\[
 0={d\over ds}\langle\Omega_s,H_s\Omega_s\rangle
 =\langle\Omega_s,\dot H_s\Omega_s\rangle.
\]
The expectation of \([G_s,H_s]\) in a zero-energy vector vanishes.
Using \eqref{eq:e8v39-V} proves the first identity, and hence projection
by \(Q_s\) leaves \(V_s\Omega_s\) unchanged.
\end{proof}

\begin{firewall}[The split is connection dependent]
\label{fire:e8v39-connection}
Equation \eqref{eq:e8v39-V} does not canonically decide which part of a
metric perturbation is gauge.  A gauge fixing, a constraint connection,
or an explicit unitary implementation must determine \(G_s\).  Choosing
\(G_s\) merely to make \(V_s\) small would be circular unless the choice
also respects the Ward quotient, locality, boundary conditions, and the
physical observable algebra.
\end{firewall}

\subsection{Cancellation of a covariantly transported recovered source}
\label{sec:e8v39-covariance}

The four terms in \eqref{eq:e8v37-four-split} can cancel exactly along a
pure representation change.  Estimating them separately would then
create an artificial debit.

\begin{theorem}[Exact covariance cancellation]
\label{thm:e8v39-covariance}
Let \(U_s\) be as above and let
\begin{align}
 \Omega_s&=U_s\Omega_0,
 &
 A_s&=U_sA_0U_s^*,
 \label{eq:e8v39-covariant-A}\\
 K_s&=U_sK_0U_s^*,
 &
 R_s&=(K_s+\alpha I)^{-1}=U_sR_0U_s^*.
 \label{eq:e8v39-covariant-R}
\end{align}
If the coefficient loader is contragredient,
\begin{equation}
 J_s=J_0U_s^*,
 \label{eq:e8v39-covariant-J}
\end{equation}
then the recovered coefficient source is constant:
\begin{equation}
 j_s=J_sR_sA_s\Omega_s
 =J_0R_0A_0\Omega_0.
 \label{eq:e8v39-constant-j}
\end{equation}
Consequently, the four derivatives in
\eqref{eq:e8v37-four-split} sum to zero.
\end{theorem}

\begin{proof}
Insert \eqref{eq:e8v39-covariant-A}--\eqref{eq:e8v39-covariant-J} and
cancel adjacent factors \(U_s^*U_s\).  The result has no \(s\)-dependence.
Differentiating the constant identity gives the final assertion.
\end{proof}

\begin{corollary}[Gauge-orbit replacements have zero coefficient debit]
\label{cor:e8v39-zero-debit}
If the two endpoints of a replacement path differ only by the unitary
transport in \Cref{thm:e8v39-covariance}, then
\begin{equation}
 \|j_1-j_0\|_{\mathcal C}=0.
 \label{eq:e8v39-zero-debit}
\end{equation}
The heat-bath incidence must therefore be computed after gauge-covariant
identification of the endpoint coefficient spaces.
\end{corollary}

\begin{proof}
Evaluate \eqref{eq:e8v39-constant-j} at the two endpoints.
\end{proof}

\begin{remark}[Partial covariance]
\label{rem:e8v39-partial-covariance}
If any of \(J,R,A,\Omega\) fails its covariance identity, the remaining
derivative is a genuine contact term.  The theorem licenses cancellation
only for the explicitly transported factors; it does not remove regulator
or anomaly contacts.
\end{remark}

\subsection{Locality and constraint meaning}
\label{sec:e8v39-meaning}

\begin{definition}[Gauge-tangent decomposition card, GTDC]
\label{def:e8v39-GTDC}
For each metric replacement direction, GTDC consists of:
\begin{enumerate}[label=\textup{(G\arabic*)},leftmargin=4em]
\item a declared anti-self-adjoint connection \(G_s\) implementing the
regulated gauge, frame, or diffeomorphism tangent;
\item the exact decomposition
\(\dot H_s=[G_s,H_s]+V_s\);
\item a parallel vacuum gauge and a local or summable bound for
\(q_s=Q_sG_s\Omega_s\);
\item a positive infrared-order or other finite-horizon terminal estimate
for \(v_s=Q_sV_s\Omega_s\);
\item covariance identities for every transported \(J,R,A,\Omega\)
factor, with unmatched contacts retained;
\item compatibility with the Ward quotient, constraints, boundary
conditions, and the conditional metric block; and
\item the V38 heat locality and cofinal horizon schedule for the physical
source \(v_s\).
\end{enumerate}
\end{definition}

\begin{theorem}[GTDC removes the pure-transport inverse debit]
\label{thm:e8v39-GTDC}
GTDC rows (G1)--(G5) replace the V37 vacuum derivative by
\begin{equation}
 \dot\Omega_s
 =q_s+\nabla_s\Omega_s,
 \qquad
 H_s\nabla_s\Omega_s=-v_s,
 \label{eq:e8v39-covariant-derivative}
\end{equation}
where only \(\nabla_s\Omega_s\) requires inversion or finite-horizon
response.  If (G6)--(G7) also hold with summable constants, the physical
vacuum-motion row enters FHRC without charging a resolvent to the pure
gauge tangent.
\end{theorem}

\begin{proof}
Define \(\nabla_s\Omega_s=\dot\Omega_s-q_s\) and apply
\eqref{eq:e8v39-covariant-vacuum}.  Exact covariance cancellation is
\Cref{thm:e8v39-covariance}.  The remaining physical response and its
summability are the V38 FHRC rows.
\end{proof}

\begin{firewall}[A diffeomorphism tangent is not every metric tangent]
\label{fire:e8v39-not-all-metric}
Longitudinal regulated directions may be generated by constraints, but
transverse metric and curvature variations change physical data.  They
belong in \(V_s\).  Treating the full Einstein metric replacement as a
unitary conjugation would erase the physical stress response and cannot
serve as a proof of the coupled continuum.
\end{firewall}

\begin{firewall}[Anomalies and boundaries obstruct exact covariance]
\label{fire:e8v39-anomaly}
A formal classical gauge generator does not prove
\eqref{eq:e8v39-covariant-A}--\eqref{eq:e8v39-covariant-J} for the
regulated chiral measure or a boundary sector.  Jacobians, determinant
lines, edge modes, and anomaly cancellation are separate contacts.  Any
failure of covariance must remain in \(V_s\) or in an unmatched
coefficient derivative.
\end{firewall}

\begin{assessment}[Progress at V39]
\label{ass:e8v39-status}
Pure unitary transport now has an exact algebraic treatment.  Its vacuum
source lies in \(\operatorname{Ran}H\), its derivative is known without an
inverse, and a fully covariant recovered coefficient source has zero
replacement debit.  A general metric tangent splits into this transport
part and a physical source \(v_s\); only the latter pays the V38 infrared
and propagation estimates.  The unresolved task is to derive the
physical source order from the actual coupled constraint and stress
identities.
\end{assessment}

\begin{programme}[Eighth-series V40: companion audit and physical-source ledger]
\label{prog:e8v39-V40}
V40 is the compulsory companion audit.  It will inspect the newest
Yang--Mills and Navier--Stokes notes for a non-circular source factorization
or boundary-sector estimate.  It will then compare those rows with GTDC
and state the smallest physical metric-source lemma still missing.
\end{programme}

\begin{firewall}[Claim boundary after eighth-series V39]
\label{fire:e8v39-boundary}
V39 proves the order-one identity for a unitary tangent, the exact
gauge--physical source split, the finite-horizon equation defect, and the
covariant cancellation of a recovered coefficient source.

V39 does not construct the connection \(G_s\) for interacting Einstein
constraints, prove infrared order for \(V_s\Omega_s\), establish chiral or
boundary covariance, construct the positive gravitational carrier,
remove the cutoff, or prove the full continuum.
\end{firewall}

\begin{theorem}[Eighth-series V39 result]
\label{thm:e8v39-result}
For
\(\dot H_s=[G_s,H_s]+V_s\), the moving-vacuum source splits exactly as
\[
 Q_s\dot H_s\Omega_s
 =-H_sQ_sG_s\Omega_s+Q_sV_s\Omega_s.
\]
The first term has infrared order one and is recovered directly by
unitary transport.  Only the second term requires the V38 finite-horizon
response.  When \(J,R,A,\Omega\) all transform covariantly along the
unitary orbit, their recovered coefficient source is exactly constant.
\end{theorem}

\begin{proof}
Use \Cref{thm:e8v39-order-one,thm:e8v39-source-split,
thm:e8v39-covariance}.
\end{proof}

\paragraph{Parent-version record.}
V39 was formed from
\texttt{Renewal\_SM\_Einstein\_Continuum\_Research\_Notes\_V38.tex}.
The V38 parent had (15\,240) logical source lines, (556\,682) bytes,
and SHA--256
\[
 \texttt{B71753090EF9EE981647AD279A2220672BBCC9AB1DB1564C40AE5BC329B622D8}.
\]
The next compulsory companion audit is V40.

% =============================================================================
% END APPENDED EIGHTH-SERIES V39 MATERIAL
% =============================================================================

% =============================================================================
% BEGIN APPENDED EIGHTH-SERIES V40 MATERIAL
% =============================================================================

\section{V40 companion audit and the curvature-to-spectrum bottleneck}
\label{sec:v8v40}

This is the compulsory audit following V35.  The Yang--Mills programme has
advanced from V94 to V99: it now has an exact buffered-collar census,
full-tree frame conditioning, a rank-\(60\) curvature map with no
non-gauge linear kernel, and local nonlinear chord reconstruction.  The
Navier--Stokes active note remains V26.  The audit suggests a precise SM
use: curvature controls the transverse size of a regulated tangent after
boundary frames are conditioned.  This note proves that Hodge-type split,
then shows that it does not by itself give the infrared order required in
V38.

\subsection{Compulsory companion audit}
\label{sec:e8v40-audit}

\begin{audit}[Companion files read at V40]
\label{audit:e8v40-files}
The audit used:
\begin{enumerate}[label=\textup{(\roman*)},leftmargin=3.8em]
\item
\texttt{Renewal\_Yang\_Mills\_Mass\_Gap\_Research\_Notes\_V99.tex},
with \(46\,872\) logical source lines, \(1\,630\,579\) bytes, one live
terminator, and SHA--256
\[
 \texttt{FC4EBC96016C0AEFE6D85AF3146B5876637700CBA6540DB499BB37134E43DC10};
\]
\item
\texttt{Renewal\_NS\_Stability\_Research\_Notes\_V26.tex},
with \(5\,319\) logical source lines, \(278\,283\) bytes, one live
terminator, and SHA--256
\[
 \texttt{82C887F8C2C69E4F28FAD7C3DC8DE5B9099CB5A4BF431EBA1FFF3E91935978AD}.
\]
\end{enumerate}
Both audited files are structurally complete.  YM V98 is also present as
the immediate predecessor, but V99 is the newest complete YM continuation
and contains V98 as its append-only prefix.
\end{audit}

\begin{theorem}[New Yang--Mills input selected by the V40 audit]
\label{thm:e8v40-YM-input}
Relative to the V35 audit of YM V94, YM V95--V99 proves the following
finite-block facts.
\begin{enumerate}[label=\textup{(\roman*)},leftmargin=3.8em]
\item Small Wilson frustration gives an aligned one-link well, but
arbitrary enlarged-fibre fluctuations still require an aligned corridor.
\item Tree gauge and common-relative coordinates reduce the original
enlarged star to one common and five relative chord modes, with
gauge-covariant transport between conditional fibres.
\item The principal Wilson region must include a one-layer crossing
collar; treating those crossing plaquettes only by crude bounded
oscillation is incompatible with the desired action margin.
\item On the resulting \(64\)-vertex, \(123\)-link,
\(72\)-plaquette collar, full-tree conditioning leaves \(60\) curvature
chords.  The linear curvature map has rank \(60\), its entire kernel
before conditioning is gauge, and a certified singular floor is at least
\(102^{-59}\).
\item Local nonlinear chord reconstruction follows near every flat chord
configuration.
\item Global classification of nonlinear flat sectors, physical
response-ground stability, complete-action smoothing, and form incidence
remain open.
\end{enumerate}
\end{theorem}

\begin{proof}
The statements are the results and firewalls labelled in the YM note as
\texttt{thm:s6v95-one-well},
\texttt{thm:s6v96-fibre-transport},
\texttt{thm:s6v96-variance},
\texttt{prop:s6v98-incompatibility},
\texttt{lem:s6v98-buffer-gauge},
\texttt{thm:s6v99-rank},
\texttt{thm:s6v99-chord-frame},
\texttt{prop:s6v99-floor}, and
\texttt{thm:s6v99-local-reconstruction}, with
\texttt{fire:s6v99-flat-sectors} and
\texttt{fire:s6v99-boundary} retained.
\end{proof}

\begin{theorem}[Navier--Stokes status at V40]
\label{thm:e8v40-NS-input}
The NS file and digest are unchanged from the V35 audit.  V26 computes
pointwise rank-one norms for first and second derivatives of the Oseen
kernel and reports only a modest second-derivative improvement.  It
supplies no new gauge-frame decomposition, moving-vacuum factorization,
conditional gravitational carrier, or infrared source-order theorem.
The full Navier--Stokes regularity statement remains open.
\end{theorem}

\begin{proof}
The digest agrees with \Cref{audit:e8v35-files}; the terminal V26 theorem
and programme are unchanged.
\end{proof}

\begin{assessment}[Audit-guided use and limit]
\label{ass:e8v40-direction}
The YM collar result licenses a finite-dimensional separation of gauge
frames from curvature chords.  Its boundary qualification is essential:
outer frame variables are conditioned parameters, not directions that may
be gauge-fixed while the far exterior is held unchanged.  The result does
not identify the spectral location of
\(Q\,\dot H\Omega\).  The next subsection proves both statements in an
abstract form suitable for the SM ledger.
\end{assessment}

\subsection{Finite-complex gauge--curvature decomposition}
\label{sec:e8v40-Hodge}

Let \(\mathsf C^0,\mathsf C^1,\mathsf C^2\) be finite-dimensional Hilbert
spaces and let
\begin{equation}
 \mathsf C^0\xrightarrow{\,d_0\,}\mathsf C^1
 \xrightarrow{\,d_1\,}\mathsf C^2,
 \qquad d_1d_0=0.
 \label{eq:e8v40-complex}
\end{equation}
Write
\begin{equation}
 \mathsf G=\operatorname{Ran}d_0,
 \qquad
 \Pi_{\mathsf G}\text{ for the orthogonal projection onto }\mathsf G,
 \qquad
 \Pi_\perp=I-\Pi_{\mathsf G}.
 \label{eq:e8v40-projections}
\end{equation}

\begin{theorem}[Quantitative finite-complex Hodge split]
\label{thm:e8v40-Hodge}
Assume
\begin{equation}
 \ker d_1=\operatorname{Ran}d_0.
 \label{eq:e8v40-H1zero}
\end{equation}
Then \(d_1|_{\mathsf G^\perp}\) is injective and
\begin{equation}
 \sigma_1:=
 \inf_{\substack{x\in\mathsf G^\perp\\\|x\|=1}}\|d_1x\|>0.
 \label{eq:e8v40-sigma}
\end{equation}
Every \(x\in\mathsf C^1\) has the unique orthogonal decomposition
\begin{equation}
 x=x_{\mathsf g}+x_\perp,
 \qquad
 x_{\mathsf g}=\Pi_{\mathsf G}x,
 \qquad
 x_\perp=\Pi_\perp x,
 \label{eq:e8v40-Hodge-split}
\end{equation}
and
\begin{equation}
 d_1x=d_1x_\perp,
 \qquad
 \|x_\perp\|\leq\sigma_1^{-1}\|d_1x\|.
 \label{eq:e8v40-curvature-control}
\end{equation}
\end{theorem}

\begin{proof}
The relation \(d_1d_0=0\) gives
\(\mathsf G\subseteq\ker d_1\); equality is
\eqref{eq:e8v40-H1zero}.  Hence a vector in
\(\mathsf G^\perp\cap\ker d_1\) must vanish, proving injectivity.  The unit
sphere of the finite-dimensional space \(\mathsf G^\perp\) is compact,
and the continuous function \(x\mapsto\|d_1x\|\) has no zero there.
Its minimum \(\sigma_1\) is therefore positive.  Orthogonal projection
gives \eqref{eq:e8v40-Hodge-split}.  Since
\(d_1x_{\mathsf g}=0\),
\[
 \|d_1x\|=\|d_1x_\perp\|\geq\sigma_1\|x_\perp\|,
\]
which proves \eqref{eq:e8v40-curvature-control}.
\end{proof}

\begin{remark}[Meaning of the YM floor]
\label{rem:e8v40-YM-floor}
For the fully frame-conditioned YM V99 collar, the audited integer-minor
certificate supplies one explicit, deliberately coarse value
\(\sigma_1\geq102^{-59}\).  V40 does not transplant that numerical value
to an Einstein or Standard Model complex.  It imports the proof pattern:
condition boundary frames, identify the whole linear kernel, and only then
declare a transverse curvature floor.
\end{remark}

\subsection{Hamiltonian source induced by the Hodge split}
\label{sec:e8v40-source}

Let \(H\geq0\) have normalized vacuum \(\Omega\), with
\(P=|\Omega\rangle\langle\Omega|\) and \(Q=I-P\).  Let
\(\mathscr L:\mathsf C^1\to\mathcal B_{\rm sa}(\mathcal H)\) be the
linearized Hamiltonian response at a fixed cutoff.  Suppose gauge
directions have declared anti-self-adjoint implementers:
\begin{equation}
 \mathscr L(d_0\xi)=[G(\xi),H].
 \label{eq:e8v40-gauge-implementation}
\end{equation}
For \(x_{\mathsf g}=d_0\xi_x\), choose one fixed linear right inverse on
\(\mathsf G\) and put
\begin{equation}
 q(x)=QG(\xi_x)\Omega,
 \qquad
 v(x)=Q\mathscr L(x_\perp)\Omega.
 \label{eq:e8v40-qv}
\end{equation}

\begin{theorem}[Curvature-controlled gauge--physical source split]
\label{thm:e8v40-source-split}
Under \eqref{eq:e8v40-H1zero} and
\eqref{eq:e8v40-gauge-implementation},
\begin{equation}
 Q\mathscr L(x)\Omega
 =-Hq(x)+v(x).
 \label{eq:e8v40-source-split}
\end{equation}
If
\begin{equation}
 \|Q\mathscr L(y)\Omega\|
 \leq C_{\mathscr L}\|y\|
 \qquad(y\in\mathsf G^\perp),
 \label{eq:e8v40-L-bound}
\end{equation}
then
\begin{equation}
 \|v(x)\|
 \leq {C_{\mathscr L}\over\sigma_1}\|d_1x\|.
 \label{eq:e8v40-v-curvature}
\end{equation}
\end{theorem}

\begin{proof}
By linearity,
\[
 Q\mathscr L(x)\Omega
 =Q\mathscr L(d_0\xi_x)\Omega
  +Q\mathscr L(x_\perp)\Omega.
\]
Using \eqref{eq:e8v40-gauge-implementation} and \(H\Omega=0\), the first
term is
\[
 Q(G(\xi_x)H-HG(\xi_x))\Omega
 =-HQG(\xi_x)\Omega=-Hq(x).
\]
This proves \eqref{eq:e8v40-source-split}.  Apply
\eqref{eq:e8v40-L-bound} and then
\eqref{eq:e8v40-curvature-control} to obtain
\eqref{eq:e8v40-v-curvature}.
\end{proof}

\begin{corollary}[Compatibility with the V39 tangent split]
\label{cor:e8v40-V39}
The source in \eqref{eq:e8v40-source-split} has exactly the form
\eqref{eq:e8v39-b-split}: the cochain gauge component is an
order-one range-\(H\) source, while \(v(x)\) is the physical curvature
source.  The finite-complex theorem bounds the norm of \(v(x)\) but does
not assign it positive infrared order.
\end{corollary}

\begin{proof}
Compare \eqref{eq:e8v40-source-split} with
\eqref{eq:e8v39-b-split}.  The last sentence follows because
\eqref{eq:e8v40-v-curvature} contains no power of \(H\).
\end{proof}

\subsection{Curvature control does not imply infrared order}
\label{sec:e8v40-counterexample}

\begin{theorem}[Finite curvature floor cannot supply spectral order]
\label{thm:e8v40-no-infrared}
Fix \(\alpha>0\).  There is a family of finite-cutoff systems for which:
\begin{enumerate}[label=\textup{(\roman*)}]
\item \(\ker d_1=\operatorname{Ran}d_0\) and
\(\sigma_1=1\);
\item the physical tangent and curvature have unit norm;
\item \(\|Q\mathscr L(x)\Omega\|=1\); but
\item every solution of
\begin{equation}
 Q\mathscr L(x)\Omega=H_\varepsilon^\alpha c_\varepsilon
 \label{eq:e8v40-factor-test}
\end{equation}
satisfies
\(\|Q c_\varepsilon\|=\varepsilon^{-\alpha}\).
\end{enumerate}
Thus no cutoff-uniform infrared-order constant follows from the
finite-complex curvature floor and
\eqref{eq:e8v40-L-bound} alone.
\end{theorem}

\begin{proof}
Take
\(\mathsf C^0=\{0\}\), \(\mathsf C^1=\mathsf C^2=\mathbb R\),
\(d_0=0\), and \(d_1=1\).  Then
\(\ker d_1=\{0\}=\operatorname{Ran}d_0\) and \(\sigma_1=1\).
Let \(\mathcal H=\mathbb C^2\), with orthonormal basis
\(\{\Omega,e\}\), and set
\begin{equation}
 H_\varepsilon\Omega=0,
 \qquad
 H_\varepsilon e=\varepsilon e,
 \qquad 0<\varepsilon\leq1.
 \label{eq:e8v40-Heps}
\end{equation}
Define the self-adjoint response
\[
 \mathscr L(1)=|e\rangle\langle\Omega|
              +|\Omega\rangle\langle e|.
\]
For the unit tangent \(x=1\), one has
\(d_1x=1\) and
\(Q\mathscr L(x)\Omega=e\).  If
\eqref{eq:e8v40-factor-test} holds, projection onto \(e\) gives
\(\varepsilon^\alpha\langle e,c_\varepsilon\rangle=1\).
Therefore
\(\|Qc_\varepsilon\|\geq
|\langle e,c_\varepsilon\rangle|=\varepsilon^{-\alpha}\).
Equality is achieved by
\(c_\varepsilon=\varepsilon^{-\alpha}e\).
\end{proof}

\begin{consequence}[Smallest missing physical-source lemma]
\label{cons:e8v40-missing}
The Hodge split and a positive curvature floor reduce the V39--V38
bottleneck to a cross-operator statement.  One must prove, for the actual
physical transverse metric source, a factorization or estimate that
contains both the curvature differential and a positive power of the
Hamiltonian.  Neither factor may be inferred from the other.
\end{consequence}

\begin{proof}
\Cref{thm:e8v40-source-split} supplies curvature control;
\Cref{thm:e8v40-no-infrared} rules out the missing implication to
Hamiltonian order.
\end{proof}

\subsection{Curvature-to-spectral factorization}
\label{sec:e8v40-factorization}

\begin{definition}[Physical curvature--spectral factorization, PCSF]
\label{def:e8v40-PCSF}
At cutoff \(n\), PCSF with order \(\alpha>0\) consists of a bounded map
\begin{equation}
 C_n:\mathsf C_n^2\longrightarrow Q_n\mathcal H_n
 \label{eq:e8v40-Cn}
\end{equation}
such that every declared transverse replacement tangent \(x_n\) obeys
\begin{equation}
 Q_n\mathscr L_n(\Pi_{\perp,n}x_n)\Omega_n
 =H_n^\alpha C_nd_{1,n}x_n,
 \qquad
 \sup_n\|C_n\|\leq C_*.
 \label{eq:e8v40-PCSF}
\end{equation}
Locality or a summable influence tail for \(C_n\) is a separate part of
the row.
\end{definition}

\begin{theorem}[PCSF closes the finite-horizon terminal row]
\label{thm:e8v40-PCSF}
If PCSF holds, then the physical source
\(v_n=Q_n\mathscr L_n(\Pi_{\perp,n}x_n)\Omega_n\) has uniform infrared
order \(\alpha\), and
\begin{equation}
 \|e^{-T_nH_n}v_n\|
 \leq
 C_*\left({\alpha\over eT_n}\right)^\alpha
 \|d_{1,n}x_n\|.
 \label{eq:e8v40-PCSF-residual}
\end{equation}
Consequently, the weighted terminal debits are summable whenever
\begin{equation}
 \sum_n a_nT_n^{-\alpha}\|d_{1,n}x_n\|<\infty.
 \label{eq:e8v40-PCSF-sum}
\end{equation}
\end{theorem}

\begin{proof}
In \eqref{eq:e8v38-infrared-order}, take
\(c_n=C_nd_{1,n}x_n\).  Its norm is at most
\(C_*\|d_{1,n}x_n\|\).  Apply
\Cref{thm:e8v38-fractional} and then the comparison test.
\end{proof}

\begin{firewall}[PCSF is much stronger than a Ward quotient]
\label{fire:e8v40-PCSF}
PCSF says that the physical transverse response vanishes at low energy
with a quantitative Hamiltonian power.  Gauge orthogonality, anomaly
cancellation, a finite curvature floor, and a local stress tensor do not
individually imply it.  A proof must use the coupled dynamics or replace
PCSF by a weaker observable-level cancellation with an equally explicit
terminal estimate.
\end{firewall}

\begin{definition}[Audited physical-source card, APSC]
\label{def:e8v40-APSC}
APSC consists of:
\begin{enumerate}[label=\textup{(A\arabic*)},leftmargin=4em]
\item a finite gauge--curvature complex with all boundary frames either
conditioned or explicitly owned;
\item exactness \(\ker d_1=\operatorname{Ran}d_0\) and a transverse
singular floor;
\item the Hamiltonian gauge implementation
\eqref{eq:e8v40-gauge-implementation};
\item PCSF, or an observable-level substitute, for the transverse source;
\item a local or summable influence row for the factor \(C_n\);
\item the V38 finite-horizon schedule and V37 coefficient replacement
incidence; and
\item nonlinear flat-sector, conditional-density, response-ground,
Ward, boundary, reflection, spin, chiral, and gravitational contacts.
\end{enumerate}
\end{definition}

\begin{theorem}[APSC handoff]
\label{thm:e8v40-APSC}
APSC rows (A1)--(A5) decompose every declared regulated tangent into a
directly transported gauge source and a curvature-controlled physical
source with the terminal estimate
\eqref{eq:e8v40-PCSF-residual}.  With (A6), the moving-vacuum contribution
enters the FHRC and PSRC compilers.  Row (A7) remains necessary for the
physical Standard Model--Einstein handoff.
\end{theorem}

\begin{proof}
The decomposition and norm control are
\Cref{thm:e8v40-Hodge,thm:e8v40-source-split}.  PCSF gives
\Cref{thm:e8v40-PCSF}.  V38 and V37 supply the finite-horizon and
replacement-incidence compilers.  The final row lists contacts not proved
by those abstract results.
\end{proof}

\begin{assessment}[Progress at V40]
\label{ass:e8v40-status}
The compulsory audit has produced a precise advance rather than a formal
analogy.  A finite curvature complex with no non-gauge linear kernel gives
a quantitative gauge--physical tangent split, and its Hamiltonian source
inherits the V39 range-\(H\) gauge term.  An explicit two-level
counterexample proves that this curvature floor says nothing uniform about
the physical source near zero Hamiltonian energy.  PCSF is now the
smallest named cross-operator row that would close the V38 terminal
estimate.
\end{assessment}

\begin{programme}[Eighth-series V41: quadratic-ground-state test of PCSF]
\label{prog:e8v40-V41}
V41 will test PCSF in the exactly solvable quadratic branch.  It will
differentiate a Gaussian vacuum under a positive quadratic form, compute
the physical source and reduced response mode by mode, and determine which
Hamiltonian power is actually available uniformly as the lowest
frequency tends to zero.  If PCSF fails even there, the programme will
retreat to an observable-level seminorm rather than assume the row.
\end{programme}

\begin{firewall}[Claim boundary after eighth-series V40]
\label{fire:e8v40-boundary}
V40 records the compulsory YM/NS audit, proves the quantitative
finite-complex Hodge split and induced Hamiltonian source split, proves
that curvature coercivity alone cannot imply infrared source order, and
formulates the sufficient PCSF row.

V40 does not prove PCSF for a quadratic or interacting field, classify
the YM nonlinear flat sectors, establish the Einstein constraint complex,
construct the conditional gravitational carrier, remove any cutoff, or
prove the full Standard Model--Einstein continuum.
\end{firewall}

\begin{theorem}[Eighth-series V40 result]
\label{thm:e8v40-result}
If a finite gauge--curvature complex is exact at degree one, every tangent
splits orthogonally into gauge and curvature-controlled parts.  Under
Hamiltonian gauge covariance, the induced vacuum source is
\[
 Q\mathscr L(x)\Omega=-Hq(x)+v(x),
 \qquad
 \|v(x)\|\leq C_{\mathscr L}\sigma_1^{-1}\|d_1x\|.
\]
This estimate does not imply positive infrared order.  PCSF is a
sufficient additional row and yields
\eqref{eq:e8v40-PCSF-residual}.
\end{theorem}

\begin{proof}
Use \Cref{thm:e8v40-Hodge,thm:e8v40-source-split,
thm:e8v40-no-infrared,thm:e8v40-PCSF}.
\end{proof}

\paragraph{Parent-version record.}
V40 was formed from
\texttt{Renewal\_SM\_Einstein\_Continuum\_Research\_Notes\_V39.tex}.
The V39 parent had (15\,715) logical source lines, (572\,813) bytes,
and SHA--256
\[
 \texttt{9BDA6D4E1112E20FB770F08533CDBEB5B8A8FA1D8DB28F847651378CE658D7C3}.
\]
The next compulsory companion audit is V45.

% =============================================================================
% END APPENDED EIGHTH-SERIES V40 MATERIAL
% =============================================================================

% =============================================================================
% BEGIN APPENDED EIGHTH-SERIES V41 MATERIAL
% =============================================================================

\section{Quadratic vacuum test of curvature--spectral factorization}
\label{sec:v8v41}

V40 reduced the moving-vacuum problem to PCSF and required that the
physical source contain a positive Hamiltonian power with a cofinally
bounded coefficient.  The first honest test is a family of harmonic
modes.  It is exactly solvable and includes the free bosonic principal
symbol of a metric-dependent field.  The calculation below shows two
different facts.  Relative frequency variations have infrared order one,
but their global vacuum derivative need not have a cutoff-independent
Hilbert norm because every ultraviolet mode contributes.

\subsection{One oscillator}
\label{sec:e8v41-one}

On \(L^2(\mathbb R,dq)\), let \(p=-i\,d/dq\) and, for \(\omega>0\), set
\begin{equation}
 H_\omega
 ={1\over2}\bigl(p^2+\omega^2q^2-\omega\bigr).
 \label{eq:e8v41-H}
\end{equation}
Its normalized vacuum and annihilation operator are
\begin{align}
 \Omega_\omega(q)
 &=\left({\omega\over\pi}\right)^{1/4}
   e^{-\omega q^2/2},
 \label{eq:e8v41-vacuum}\\
 a_\omega
 &=\sqrt{\omega\over2}\,q
   +{i\over\sqrt{2\omega}}\,p,
 \qquad
 H_\omega=\omega a_\omega^*a_\omega.
 \label{eq:e8v41-a}
\end{align}

\begin{lemma}[Exact one-mode vacuum and source derivatives]
\label{lem:e8v41-one-derivative}
Let \(s\mapsto\omega_s>0\) be \(C^1\).  At any fixed \(s\), write
\(a=a_{\omega_s}\), \(\Omega=\Omega_{\omega_s}\), and
\(\eta_s=\dot\omega_s/\omega_s\).  Then
\begin{align}
 \dot\Omega_s
 &=-{\eta_s\over4}(a^*)^2\Omega,
 \label{eq:e8v41-Omegadot}\\
 Q_s\dot H_s\Omega_s
 &={\dot\omega_s\over2}(a^*)^2\Omega,
 \label{eq:e8v41-b}\\
 Q_s\dot H_s\Omega_s
 &=H_s c_s,
 \qquad
 c_s={\eta_s\over4}(a^*)^2\Omega=-\dot\Omega_s.
 \label{eq:e8v41-factor}
\end{align}
Moreover,
\begin{equation}
 \|c_s\|^2=\|\dot\Omega_s\|^2={|\eta_s|^2\over8},
 \qquad
 \|Q_s\dot H_s\Omega_s\|^2={|\dot\omega_s|^2\over2}.
 \label{eq:e8v41-one-norms}
\end{equation}
\end{lemma}

\begin{proof}
Direct differentiation of \eqref{eq:e8v41-vacuum} gives
\begin{equation}
 \partial_\omega\Omega_\omega
 =\left({1\over4\omega}-{q^2\over2}\right)\Omega_\omega.
 \label{eq:e8v41-direct-vacuum}
\end{equation}
The identity
\begin{equation}
 q^2\Omega_\omega
 ={1\over2\omega}\bigl(\Omega_\omega+(a_\omega^*)^2\Omega_\omega\bigr)
 \label{eq:e8v41-q2}
\end{equation}
turns \eqref{eq:e8v41-direct-vacuum} into
\eqref{eq:e8v41-Omegadot}.  Also
\[
 \partial_\omega H_\omega=\omega q^2-\frac12,
\]
so \eqref{eq:e8v41-q2} gives
\((\partial_\omega H_\omega)\Omega_\omega
 =\frac12(a_\omega^*)^2\Omega_\omega\), proving
\eqref{eq:e8v41-b}.  Since
\[
 H_\omega(a_\omega^*)^2\Omega_\omega
 =2\omega(a_\omega^*)^2\Omega_\omega,
\]
\eqref{eq:e8v41-factor} follows.  Finally,
\(\|(a^*)^2\Omega\|^2=2\), which gives
\eqref{eq:e8v41-one-norms}.
\end{proof}

\begin{corollary}[Exact fractional source norm]
\label{cor:e8v41-alpha}
For any \(\alpha>0\), the unique vacuum-orthogonal vector satisfying
\begin{equation}
 Q_s\dot H_s\Omega_s=H_s^\alpha c_{\alpha,s}
 \label{eq:e8v41-alpha-factor}
\end{equation}
is
\begin{equation}
 c_{\alpha,s}
 ={\dot\omega_s\over2(2\omega_s)^\alpha}
   (a^*)^2\Omega,
 \qquad
 \|c_{\alpha,s}\|^2
 =2^{-2\alpha-1}|\dot\omega_s|^2\omega_s^{-2\alpha}.
 \label{eq:e8v41-alpha-norm}
\end{equation}
If \(|\dot\omega_s|\asymp\omega_s^\beta\) as
\(\omega_s\downarrow0\), a uniform one-mode infrared coefficient can
hold only when \(\alpha\leq\beta\).
\end{corollary}

\begin{proof}
The source in \eqref{eq:e8v41-b} lies in the two-particle eigenspace of
energy \(2\omega_s\).  Division by \((2\omega_s)^\alpha\) gives the only
vacuum-orthogonal solution and its norm.  The last assertion is the power
\(\omega_s^{2(\beta-\alpha)}\) in \eqref{eq:e8v41-alpha-norm}.
\end{proof}

\begin{firewall}[Absolute and relative variations differ]
\label{fire:e8v41-relative}
If \(\dot\omega_s\) remains nonzero while \(\omega_s\downarrow0\), no
positive \(\alpha\) gives a uniform source coefficient.  If instead
\(\dot\omega_s=\eta_s\omega_s\) with bounded \(\eta_s\), order
\(\alpha=1\) is uniform.  The physical content of PCSF is therefore the
relative low-frequency behavior of the metric perturbation, not merely
smooth dependence at each fixed cutoff.
\end{firewall}

\subsection{A finite collection of independent modes}
\label{sec:e8v41-many}

Let \(\mathcal K\) be a finite mode set and let
\begin{equation}
 H_s=\sum_{k\in\mathcal K}\omega_k(s)a_{k,s}^*a_{k,s},
 \qquad
 \Omega_s=\bigotimes_{k\in\mathcal K}\Omega_{\omega_k(s)}.
 \label{eq:e8v41-many-H}
\end{equation}
Write \(\eta_k=\dot\omega_k/\omega_k\).

\begin{theorem}[Exact multimode order-one factorization]
\label{thm:e8v41-many-factor}
The parallel vacuum derivative and physical source are
\begin{align}
 \dot\Omega_s
 &=-{1\over4}\sum_{k\in\mathcal K}
   \eta_k(a_{k,s}^*)^2\Omega_s,
 \label{eq:e8v41-many-Omegadot}\\
 b_s:=Q_s\dot H_s\Omega_s
 &={1\over2}\sum_{k\in\mathcal K}
   \dot\omega_k(a_{k,s}^*)^2\Omega_s
 =H_sc_s,
 \label{eq:e8v41-many-b}\\
 c_s
 &={1\over4}\sum_{k\in\mathcal K}
   \eta_k(a_{k,s}^*)^2\Omega_s=-\dot\Omega_s.
 \label{eq:e8v41-many-c}
\end{align}
The two-particle vectors for distinct modes are orthogonal, and hence
\begin{align}
 \|c_s\|^2
 &=\|\dot\Omega_s\|^2
 ={1\over8}\sum_{k\in\mathcal K}|\eta_k|^2,
 \label{eq:e8v41-many-c-norm}\\
 \|b_s\|^2
 &={1\over2}\sum_{k\in\mathcal K}|\dot\omega_k|^2.
 \label{eq:e8v41-many-b-norm}
\end{align}
\end{theorem}

\begin{proof}
Differentiate the tensor product one factor at a time and apply
\Cref{lem:e8v41-one-derivative}.  Each summand has two quanta in mode
\(k\) and vacuum in every other mode, so distinct summands are
orthogonal.  The Hamiltonian multiplies the \(k\)-th summand by
\(2\omega_k\), giving \eqref{eq:e8v41-many-b}.  Summing the squared
one-mode norms proves the two norm identities.
\end{proof}

\begin{corollary}[A precise quadratic PCSF map]
\label{cor:e8v41-PCSF}
Give the relative-frequency tangent
\(\eta=(\eta_k)_{k\in\mathcal K}\) its \(\ell^2(\mathcal K)\) norm and
define
\begin{equation}
 C_s\eta
 ={1\over4}\sum_{k\in\mathcal K}
 \eta_k(a_{k,s}^*)^2\Omega_s.
 \label{eq:e8v41-C}
\end{equation}
Then
\begin{equation}
 \|C_s\|_{\ell^2\to\mathcal H}={1\over2\sqrt2},
 \qquad
 b_s=H_sC_s\eta.
 \label{eq:e8v41-C-bound}
\end{equation}
Thus the diagonal quadratic model satisfies PCSF with
\(\alpha=1\) uniformly in all frequencies and in the number of modes,
provided the curvature-side norm controls the
\(\ell^2\) norm of the relative frequency tangent.
\end{corollary}

\begin{proof}
Equation \eqref{eq:e8v41-many-c-norm} gives
\(\|C_s\eta\|^2=\|\eta\|_{\ell^2}^2/8\), including equality for every
\(\eta\).  The factorization is \eqref{eq:e8v41-many-b}.
\end{proof}

\begin{remark}[What the uniform operator norm means]
\label{rem:e8v41-input-norm}
The conclusion does not say that a fixed geometric parameter has bounded
\(\ell^2\) relative-frequency tangent.  It says that once that tangent is
measured in \(\ell^2\), the map from it to the source coefficient is
uniform.  The embedding of a local or global metric variation into this
mode norm is a separate ultraviolet and volume question.
\end{remark}

\subsection{Metric principal symbols have the correct infrared scaling}
\label{sec:e8v41-metric}

Let \(A_s\) be a \(C^1\) family of real symmetric matrices with
\begin{equation}
 \nu I\preceq A_s\preceq\Lambda I,
 \qquad \nu>0,
 \label{eq:e8v41-A-box}
\end{equation}
and define
\begin{equation}
 \omega_s(k)
 =\bigl(k^\top A_sk+m^2\bigr)^{1/2}.
 \label{eq:e8v41-dispersion}
\end{equation}

\begin{theorem}[Uniform relative-frequency bound for a metric tangent]
\label{thm:e8v41-relative-metric}
For every \(k\) with \(\omega_s(k)>0\),
\begin{equation}
 \eta_s(k):={\dot\omega_s(k)\over\omega_s(k)}
 ={k^\top\dot A_sk\over
   2(k^\top A_sk+m^2)}
 \label{eq:e8v41-eta-metric}
\end{equation}
and
\begin{equation}
 |\eta_s(k)|
 \leq{\|\dot A_s\|_{\rm op}\over2\nu}.
 \label{eq:e8v41-eta-bound}
\end{equation}
In particular, for \(m=0\) and nonzero momentum,
\(|\dot\omega_s(k)|\leq C\omega_s(k)\) uniformly as \(k\to0\).
\end{theorem}

\begin{proof}
Differentiate \eqref{eq:e8v41-dispersion}.  Since
\[
 |k^\top\dot A_sk|
 \leq\|\dot A_s\|_{\rm op}|k|^2
\]
and
\(k^\top A_sk+m^2\geq\nu|k|^2\), division gives
\eqref{eq:e8v41-eta-bound}.  The massless conclusion is
\(\dot\omega=\eta\omega\).
\end{proof}

\begin{corollary}[Modewise PCSF for the free metric branch]
\label{cor:e8v41-metric-PCSF}
At every finite cutoff, a metric variation of the quadratic principal
symbol \eqref{eq:e8v41-dispersion} has modewise infrared order one.  For a
finite mode set,
\begin{equation}
 \|c_s\|^2
 \leq {|\mathcal K|\over32\nu^2}
       \|\dot A_s\|_{\rm op}^2.
 \label{eq:e8v41-mode-count-bound}
\end{equation}
\end{corollary}

\begin{proof}
Use \eqref{eq:e8v41-eta-bound} in
\eqref{eq:e8v41-many-c-norm}.
\end{proof}

\subsection{The ultraviolet and volume obstruction}
\label{sec:e8v41-UV}

\begin{theorem}[Relative infrared scaling does not give a cofinal vacuum norm]
\label{thm:e8v41-mode-count-obstruction}
For each \(N\), take \(N\) massless modes and the homogeneous metric path
\(A_s=(1+s)I\) at \(s=0\).  Then every nonzero mode has
\begin{equation}
 \eta_k={1\over2},
 \label{eq:e8v41-constant-eta}
\end{equation}
while
\begin{equation}
 \|\dot\Omega_N\|^2=\|c_N\|^2={N\over32}.
 \label{eq:e8v41-sqrtN}
\end{equation}
Therefore the PCSF coefficient from a single unnormalized homogeneous
metric parameter to the global vacuum Hilbert space grows like
\(\sqrt N\), despite perfect order-one infrared behavior in every mode.
\end{theorem}

\begin{proof}
For \(m=0\),
\(\omega_s(k)=\sqrt{1+s}\,|k|\), so
\(\dot\omega_k/\omega_k=1/2\) at \(s=0\).  Substitute this value in
\eqref{eq:e8v41-many-c-norm}.
\end{proof}

\begin{consequence}[PCSF needs a local or renormalized input norm]
\label{cons:e8v41-local-PCSF}
A cofinal PCSF statement cannot measure a metric perturbation only by a
fixed finite-dimensional amplitude and then demand the global Fock norm
of its vacuum derivative.  It must instead use at least one of:
\begin{enumerate}[label=\textup{(\roman*)}]
\item a local coefficient or observable seminorm;
\item a curvature norm whose mode weights make the relative-frequency
tangent square summable;
\item a renormalized density estimate; or
\item an implementable Bogoliubov difference after cancellation of the
common ultraviolet principal symbol.
\end{enumerate}
\end{consequence}

\begin{proof}
\Cref{thm:e8v41-mode-count-obstruction} violates a uniform global bound
with none of the listed modifications.  Each item changes the norm or
subtracts the common high-frequency contribution rather than contradicting
the exact finite-mode formula.
\end{proof}

\begin{firewall}[Finite-mode unitarity is not cofinal implementability]
\label{fire:e8v41-Shale}
Every finite collection of oscillators has a differentiable normalized
vacuum in one Hilbert space.  In the infinite-mode limit, the
square-summability condition in \eqref{eq:e8v41-many-c-norm} is the
infinitesimal form of a Bogoliubov implementability requirement.  The
bounded pointwise estimate \eqref{eq:e8v41-eta-bound} does not imply that
condition.  V41 therefore proves no unitary equivalence between continuum
vacua at different metrics.
\end{firewall}

\subsection{A terminal residual density}
\label{sec:e8v41-density}

Although the global coefficient norm can diverge, the finite-horizon
terminal source has additional low-energy damping.  For the diagonal
model, \eqref{eq:e8v41-many-b} gives exactly
\begin{equation}
 \|e^{-TH_s}b_s\|^2
 ={1\over2}\sum_{k\in\mathcal K}
 |\dot\omega_k|^2e^{-4T\omega_k}.
 \label{eq:e8v41-terminal-sum}
\end{equation}

\begin{theorem}[Massless continuum terminal density]
\label{thm:e8v41-terminal-density}
Let \(m=0\), let \(A\) obey \eqref{eq:e8v41-A-box}, and suppose
\(|\eta(k)|\leq M\).  With momentum measure
\((2\pi)^{-d}dk\), the squared terminal-source norm per translation
volume is bounded by
\begin{equation}
 \mathfrak r_T^2
 :={1\over2}\int_{\mathbb R^d}
 |\dot\omega(k)|^2e^{-4T\omega(k)}
 {dk\over(2\pi)^d}
 \leq
 {M^2|\mathbb S^{d-1}|\Gamma(d+2)\over
  2(2\pi)^d\nu^{d/2}(4T)^{d+2}}.
 \label{eq:e8v41-density-bound}
\end{equation}
Thus
\begin{equation}
 \mathfrak r_T=O(T^{-(d+2)/2}).
 \label{eq:e8v41-density-rate}
\end{equation}
\end{theorem}

\begin{proof}
Since \(\dot\omega=\eta\omega\),
\[
 |\dot\omega(k)|^2e^{-4T\omega(k)}
 \leq M^2\omega(k)^2e^{-4T\omega(k)}.
\]
Set \(y=A^{1/2}k\).  Then
\(\omega(k)=|y|\) and
\(dk=(\det A)^{-1/2}dy\leq\nu^{-d/2}dy\).
Polar coordinates give
\[
 \int_{\mathbb R^d}|y|^2e^{-4T|y|}\,dy
 =|\mathbb S^{d-1}|
   \int_0^\infty r^{d+1}e^{-4Tr}\,dr
 ={|\mathbb S^{d-1}|\Gamma(d+2)\over(4T)^{d+2}}.
\]
Insert the measure and the factor \(1/2\).
\end{proof}

\begin{remark}[Density is not a state-vector norm]
\label{rem:e8v41-density}
Equation \eqref{eq:e8v41-density-bound} is a translation-volume density
for the quadratic terminal source.  It does not construct a normalized
infinite-volume vacuum vector or prove that a local interacting source has
the diagonal form used here.  Its role is to identify an observable
scale on which the massless residual improves with \(T\).
\end{remark}

\subsection{Quadratic PCSF card}
\label{sec:e8v41-card}

\begin{definition}[Quadratic curvature--spectral card, QCSC]
\label{def:e8v41-QCSC}
QCSC consists of:
\begin{enumerate}[label=\textup{(Q\arabic*)},leftmargin=4em]
\item a fixed-cutoff oscillator or Bogoliubov normal form for the protected
quadratic sector;
\item relative frequency derivatives controlled by the declared metric
curvature norm;
\item the exact order-one factorization
\eqref{eq:e8v41-many-b};
\item a cutoff-independent local, weighted, or density norm replacing
the divergent unnormalized global metric tangent;
\item a locality estimate for the change of normal form and for
off-diagonal mode mixing;
\item the V38 finite-horizon residual and propagation schedule; and
\item constraint, Ward, boundary, reflection, chiral, interaction,
renormalization, and gravitational-carrier contacts.
\end{enumerate}
\end{definition}

\begin{theorem}[What QCSC would supply]
\label{thm:e8v41-QCSC}
Rows (Q1)--(Q3) prove finite-cutoff PCSF with \(\alpha=1\) in the
relative-frequency \(\ell^2\) norm.  Row (Q4) is necessary for cofinality:
without it, \Cref{thm:e8v41-mode-count-obstruction} makes the global
coefficient diverge.  If (Q4)--(Q6) hold with summable constants, the
quadratic physical source enters APSC and FHRC.  Row (Q7) remains outside
the quadratic calculation.
\end{theorem}

\begin{proof}
\Cref{cor:e8v41-PCSF,thm:e8v41-relative-metric} give the order-one
factorization.  The obstruction is
\Cref{thm:e8v41-mode-count-obstruction}; the density alternative is
\Cref{thm:e8v41-terminal-density}.  V38 supplies the finite-horizon
compiler once the stated norm and locality rows are available.
\end{proof}

\begin{assessment}[Progress at V41]
\label{ass:e8v41-status}
The quadratic test does not reject PCSF at low energy.  A metric principal
symbol changes every massless frequency relatively, and the Gaussian
vacuum source is exactly \(H\) times a two-particle coefficient.  It does
reject the unqualified global Hilbert-norm version: the coefficient norm
counts modes and diverges for a homogeneous metric tangent.  The programme
must now formulate the response in a local, weighted, or renormalized
observable topology before interaction estimates are attempted.
\end{assessment}

\begin{programme}[Eighth-series V42: local quadratic stress and implementability]
\label{prog:e8v41-V42}
V42 will leave the diagonal homogeneous tangent and treat a spatially
localized quadratic coefficient perturbation.  It will compute the
two-particle Bogoliubov kernel, state its exact Hilbert--Schmidt criterion,
and test which Sobolev regularity or principal-symbol subtraction makes
the local vacuum response finite.  If the criterion fails at the physical
dimension, the source compiler will be moved to a local observable
seminorm rather than a vacuum-vector norm.
\end{programme}

\begin{firewall}[Claim boundary after eighth-series V41]
\label{fire:e8v41-boundary}
V41 proves the exact one-mode and finite-mode Gaussian vacuum derivatives,
the order-one relative-frequency factorization, the uniform infrared
scaling of a metric principal symbol, the mode-count obstruction, and a
massless terminal-source density estimate.

V41 does not prove a cofinal Hilbert-space vacuum derivative, a
Bogoliubov implementability theorem for local metric changes, the
interacting PCSF row, renormalized stress control, the conditional
gravitational carrier, cutoff removal, or the full continuum.
\end{firewall}

\begin{theorem}[Eighth-series V41 result]
\label{thm:e8v41-result}
For finitely many quadratic modes,
\[
 Q\dot H\Omega
 =H\left({1\over4}\sum_k
 {\dot\omega_k\over\omega_k}(a_k^*)^2\Omega\right),
 \qquad
 \left\|H^{-1}Q\dot H\Omega\right\|^2
 ={1\over8}\sum_k
 \left|{\dot\omega_k\over\omega_k}\right|^2.
\]
Metric principal symbols make each relative derivative uniformly bounded
at zero momentum, but the sum may diverge with the number of modes.  The
finite-horizon terminal source nevertheless has the density decay
\eqref{eq:e8v41-density-bound}.
\end{theorem}

\begin{proof}
Use \Cref{thm:e8v41-many-factor,thm:e8v41-relative-metric,
thm:e8v41-mode-count-obstruction,thm:e8v41-terminal-density}.
\end{proof}

\paragraph{Parent-version record.}
V41 was formed from
\texttt{Renewal\_SM\_Einstein\_Continuum\_Research\_Notes\_V40.tex}.
The V40 parent had (16\,217) logical source lines, (590\,906) bytes,
and SHA--256
\[
 \texttt{771D84ED705CB3AF8C0CB2B0C710A81D55DCA6B4534AB999F808166EB6766D82}.
\]
The next compulsory companion audit is V45.

% =============================================================================
% END APPENDED EIGHTH-SERIES V41 MATERIAL
% =============================================================================

% =============================================================================
% BEGIN APPENDED EIGHTH-SERIES V42 MATERIAL
% =============================================================================

\section{Local quadratic perturbations and the implementability barrier}
\label{sec:v8v42}

V41 diagonalized a homogeneous metric tangent and found a mode-count
divergence.  A localized perturbation mixes modes, so the relevant object
is a two-particle kernel rather than a sequence of diagonal relative
frequencies.  This note derives that kernel from the Sylvester derivative
of the one-particle frequency.  It then proves that a nonzero perturbation
of the principal second-order coefficient is not Hilbert--Schmidt on the
free massless torus, even when its coefficient is smooth.  Lower-order
perturbations behave differently in spatial dimension at most three.

\subsection{Finite-dimensional Gaussian field}
\label{sec:e8v42-Gaussian}

Let \(\mathfrak h=\mathbb R^N\) with its Euclidean inner product, let
\(A_s\) be a \(C^1\) family of positive symmetric matrices, and put
\begin{equation}
 \omega_s=A_s^{1/2}.
 \label{eq:e8v42-omega}
\end{equation}
On \(L^2(\mathbb R^N,dq)\), define the vacuum-normalized quadratic
Hamiltonian
\begin{equation}
 H_s={1\over2}\bigl(p^\top p+q^\top A_sq-\operatorname{Tr}\omega_s\bigr).
 \label{eq:e8v42-H}
\end{equation}
Its normalized vacuum is
\begin{equation}
 \Omega_s(q)
 ={\det(\omega_s)^{1/4}\over\pi^{N/4}}
   \exp\left\{-{1\over2}q^\top\omega_sq\right\}.
 \label{eq:e8v42-Omega}
\end{equation}

\begin{lemma}[Sylvester derivative of the frequency]
\label{lem:e8v42-Sylvester}
The derivative \(X_s=\dot\omega_s\) is the unique symmetric solution of
\begin{equation}
 \omega_sX_s+X_s\omega_s=\dot A_s.
 \label{eq:e8v42-Sylvester}
\end{equation}
In an eigenbasis of \(\omega_s\), with eigenvalues
\(\omega_i>0\),
\begin{equation}
 (X_s)_{ij}={(\dot A_s)_{ij}\over\omega_i+\omega_j}.
 \label{eq:e8v42-Xij}
\end{equation}
\end{lemma}

\begin{proof}
Differentiate \(\omega_s^2=A_s\) to obtain
\eqref{eq:e8v42-Sylvester}.  In an eigenbasis the equation reads
\((\omega_i+\omega_j)X_{ij}=\dot A_{ij}\).  Every denominator is positive,
which gives existence, uniqueness, and \eqref{eq:e8v42-Xij}.
\end{proof}

\begin{definition}[Infinitesimal Bogoliubov kernel]
\label{def:e8v42-K}
Define
\begin{equation}
 \mathcal K_s
 =\omega_s^{-1/2}X_s\omega_s^{-1/2}.
 \label{eq:e8v42-K}
\end{equation}
In the eigenbasis of \(\omega_s\),
\begin{equation}
 (\mathcal K_s)_{ij}
 ={(\dot A_s)_{ij}\over
   \sqrt{\omega_i\omega_j}(\omega_i+\omega_j)}.
 \label{eq:e8v42-Kij}
\end{equation}
\end{definition}

\begin{theorem}[Exact Gaussian vacuum derivative and source factorization]
\label{thm:e8v42-Gaussian-factor}
Let \(a_i^*\) be the creation operators associated with the eigenbasis of
\(\omega_s\).  Then
\begin{align}
 \dot\Omega_s
 &=-{1\over4}\sum_{i,j=1}^N
   (\mathcal K_s)_{ij}a_i^*a_j^*\Omega_s,
 \label{eq:e8v42-Omegadot}\\
 b_s:=Q_s\dot H_s\Omega_s
 &={1\over4}\sum_{i,j=1}^N
 {(\dot A_s)_{ij}\over\sqrt{\omega_i\omega_j}}\,
 a_i^*a_j^*\Omega_s,
 \label{eq:e8v42-b}\\
 b_s&=-H_s\dot\Omega_s.
 \label{eq:e8v42-order-one}
\end{align}
Furthermore,
\begin{equation}
 \|\dot\Omega_s\|^2
 ={1\over8}\|\mathcal K_s\|_{\rm HS}^2.
 \label{eq:e8v42-HS-norm}
\end{equation}
\end{theorem}

\begin{proof}
Differentiate \eqref{eq:e8v42-Omega}:
\begin{equation}
 \dot\Omega_s
 =\left({1\over4}\operatorname{Tr}(\omega_s^{-1}X_s)
       -{1\over2}q^\top X_sq\right)\Omega_s.
 \label{eq:e8v42-direct}
\end{equation}
In the eigenbasis,
\[
 q_i={a_i+a_i^*\over\sqrt{2\omega_i}}.
\]
When \(q^\top X_sq\) acts on the vacuum, its vacuum component is
\(\frac12\operatorname{Tr}(\omega_s^{-1}X_s)\Omega_s\), which cancels the
first term of \eqref{eq:e8v42-direct}.  The remaining two-particle
component is \eqref{eq:e8v42-Omegadot}.

Next,
\[
 \dot H_s={1\over2}q^\top\dot A_sq
           -{1\over2}\operatorname{Tr}X_s.
\]
Its vacuum component cancels by the diagonal part of
\eqref{eq:e8v42-Sylvester}, leaving \eqref{eq:e8v42-b}.  The vector
\(a_i^*a_j^*\Omega_s\) has energy \(\omega_i+\omega_j\).  Using
\eqref{eq:e8v42-Xij} in \eqref{eq:e8v42-Omegadot} therefore gives
\eqref{eq:e8v42-order-one}.

Finally, for a symmetric matrix \(K\), the canonical commutation
relations give
\begin{equation}
 \left\|\sum_{i,j}K_{ij}a_i^*a_j^*\Omega_s\right\|^2
 =2\sum_{i,j}|K_{ij}|^2.
 \label{eq:e8v42-pair-norm}
\end{equation}
Apply this with \(K=\mathcal K_s\) and the coefficient \(1/4\).
\end{proof}

\begin{corollary}[Exact finite-cutoff PCSF criterion]
\label{cor:e8v42-PCSF}
The quadratic physical source has order-one coefficient
\(c_s=-\dot\Omega_s\).  A family of cutoffs has a uniform global
Hilbert-space PCSF coefficient exactly when
\begin{equation}
 \sup_s\|\mathcal K_s\|_{\rm HS}<\infty.
 \label{eq:e8v42-HS-criterion}
\end{equation}
\end{corollary}

\begin{proof}
Use \eqref{eq:e8v42-order-one} and
\eqref{eq:e8v42-HS-norm}.  The vacuum-orthogonal solution of
\(H_sc_s=b_s\) is unique when the vacuum is the only zero-energy vector.
\end{proof}

\begin{remark}[Consistency with V41]
\label{rem:e8v42-diagonal}
When \(A_s\) is diagonal,
\(\dot A_{ii}=2\omega_i\dot\omega_i\), so
\((\mathcal K_s)_{ii}=\dot\omega_i/\omega_i\) and all off-diagonal
entries vanish.  Equation \eqref{eq:e8v42-HS-norm} then becomes
\eqref{eq:e8v41-many-c-norm}.
\end{remark}

\subsection{Free massless scalar on a torus}
\label{sec:e8v42-torus}

Use the flat torus \(\mathbb T^d=(\mathbb R/2\pi\mathbb Z)^d\) and remove
the constant mode.  In the Fourier basis \(e_p(x)=(2\pi)^{-d/2}e^{ipx}\),
\(p\in\mathbb Z^d\setminus\{0\}\), let
\begin{equation}
 A_0=-\Delta,
 \qquad
 \omega_0=|D|,
 \qquad
 \omega_p=|p|.
 \label{eq:e8v42-flat}
\end{equation}
For a real smooth scalar coefficient \(h\), consider the symmetric
principal perturbation
\begin{equation}
 B_h=-\nabla\cdot(h\nabla).
 \label{eq:e8v42-Bh}
\end{equation}
With the Fourier convention
\(h(x)=\sum_r\widehat h(r)e^{irx}\), its matrix elements are
\begin{equation}
 (B_h)_{pq}=p\cdot q\,\widehat h(p-q).
 \label{eq:e8v42-Bhpq}
\end{equation}

\begin{theorem}[Principal-coefficient perturbations fail the HS criterion]
\label{thm:e8v42-principal-obstruction}
If \(h\not\equiv0\), then the infinitesimal kernel
\begin{equation}
 (\mathcal K_h)_{pq}
 ={p\cdot q\,\widehat h(p-q)\over
   \sqrt{|p||q|}(|p|+|q|)}
 \label{eq:e8v42-Kh}
\end{equation}
is not Hilbert--Schmidt on the nonzero Fourier modes.  This holds in every
spatial dimension \(d\geq1\), even when \(h\) is \(C^\infty\).
\end{theorem}

\begin{proof}
Choose \(r\in\mathbb Z^d\) with \(\widehat h(r)\neq0\).  Choose a lattice
direction \(e\in\mathbb Z^d\setminus\{0\}\) and set
\(p_n=ne\), \(q_n=ne-r\).  For all sufficiently large \(n\), both modes
are nonzero and
\begin{equation}
 {p_n\cdot q_n\over
  \sqrt{|p_n||q_n|}(|p_n|+|q_n|)}
 \longrightarrow {1\over2}.
 \label{eq:e8v42-limit}
\end{equation}
Hence
\[
 |(\mathcal K_h)_{p_nq_n}|
 \longrightarrow {|\widehat h(r)|\over2}>0.
\]
The Hilbert--Schmidt sum contains the infinite subseries
\(\sum_n|(\mathcal K_h)_{p_nq_n}|^2\), which diverges.
\end{proof}

\begin{corollary}[Localization and smoothness do not repair the principal symbol]
\label{cor:e8v42-smooth}
Rapid decay of \(\widehat h(r)\) as \(|r|\to\infty\) does not affect the
obstruction: a single nonzero Fourier coefficient produces infinitely
many asymptotically constant matrix entries along \(p-q=r\).
\end{corollary}

\begin{proof}
This is the fixed-\(r\) subsequence in
\Cref{thm:e8v42-principal-obstruction}.
\end{proof}

\begin{firewall}[A local metric perturbation need not move the Fock vacuum]
\label{fire:e8v42-Fock}
The failure of \eqref{eq:e8v42-HS-criterion} means that the derivative of
the cutoff vacuum does not converge as a vector in the original global
Fock representation.  It does not mean that local observables have no
response, or that algebraic states on two backgrounds cannot be compared.
It rules out the specific global-vacuum-vector PCSF route for a
second-order metric perturbation unless its common principal response is
subtracted or the topology is weakened.
\end{firewall}

\subsection{Lower-order perturbations}
\label{sec:e8v42-lower}

For comparison, let \(V\) be a real trigonometric polynomial and use the
zeroth-order perturbation
\begin{equation}
 B_Vf=Vf,
 \qquad
 (B_V)_{pq}=\widehat V(p-q).
 \label{eq:e8v42-BV}
\end{equation}
Its infinitesimal kernel is
\begin{equation}
 (\mathcal K_V)_{pq}
 ={\widehat V(p-q)\over
   \sqrt{|p||q|}(|p|+|q|)}.
 \label{eq:e8v42-KV}
\end{equation}

\begin{lemma}[A fixed Fourier diagonal is summable below dimension four]
\label{lem:e8v42-fixed-r}
Fix \(r\in\mathbb Z^d\).  Then
\begin{equation}
 \sum_{\substack{p\neq0\\p-r\neq0}}
 {1\over |p||p-r|(|p|+|p-r|)^2}<\infty
 \label{eq:e8v42-fixed-r-sum}
\end{equation}
whenever \(d\leq3\).
\end{lemma}

\begin{proof}
Split the sum into \(|p|\leq2|r|+2\) and its complement.  The first part
is finite and has no zero denominator because the excluded modes were
removed.  On the complement,
\[
 |p-r|\geq|p|/2,
 \qquad
 |p|+|p-r|\geq|p|,
\]
so the summand is at most \(2|p|^{-4}\).  The lattice series
\(\sum_{p\neq0}|p|^{-4}\) converges exactly when \(d<4\).
\end{proof}

\begin{theorem}[Finite-Fourier lower-order perturbations are HS in \(d\leq3\)]
\label{thm:e8v42-lower-HS}
If \(d\leq3\) and \(V\) has finite Fourier support, then
\begin{equation}
 \|\mathcal K_V\|_{\rm HS}^2
 =\sum_{r\in\operatorname{supp}\widehat V}
 |\widehat V(r)|^2
 \sum_{\substack{p\neq0\\p-r\neq0}}
 {1\over |p||p-r|(|p|+|p-r|)^2}
 <\infty.
 \label{eq:e8v42-lower-HS}
\end{equation}
\end{theorem}

\begin{proof}
Set \(q=p-r\) in the Hilbert--Schmidt sum of
\eqref{eq:e8v42-KV}.  Only finitely many \(r\) occur, and each inner sum
is finite by \Cref{lem:e8v42-fixed-r}.
\end{proof}

\begin{remark}[Infrared and volume qualifications]
\label{rem:e8v42-IR}
The torus has a positive least nonzero momentum.  The constants in
\eqref{eq:e8v42-fixed-r-sum} need not remain uniform as the torus side
tends to infinity, and a massless zero mode still requires its Ward or
shift quotient.  The theorem isolates the ultraviolet distinction between
orders zero and two; it does not settle the infinite-volume infrared
limit.
\end{remark}

\subsection{Principal-symbol subtraction}
\label{sec:e8v42-subtraction}

The obstruction is the order-zero high-frequency limit of
\(\mathcal K_h\).  A viable comparison must remove this common local
principal response before asking for a Hilbert--Schmidt remainder.

\begin{definition}[Subtracted implementability row]
\label{def:e8v42-SIR}
For a cutoff family of infinitesimal kernels \(\mathcal K_n\), a
subtracted implementability row consists of:
\begin{equation}
 \mathcal K_n=\mathcal K_n^{\rm loc}+\mathcal R_n,
 \label{eq:e8v42-subtraction}
\end{equation}
where:
\begin{enumerate}[label=\textup{(\roman*)}]
\item \(\mathcal K_n^{\rm loc}\) is an explicitly local principal-symbol
contact absorbed into the coefficient-space identification or local
algebra;
\item \(\mathcal R_n\) is Hilbert--Schmidt with
\(\sup_n\|\mathcal R_n\|_{\rm HS}<\infty\); and
\item changes of the subtraction on overlaps satisfy a summable cocycle
or anomaly contact.
\end{enumerate}
\end{definition}

\begin{theorem}[The HS remainder gives a renormalized order-one source]
\label{thm:e8v42-remainder}
At a finite cutoff, define
\begin{equation}
 c_n^{\rm ren}
 ={1\over4}\sum_{i,j}(\mathcal R_n)_{ij}
 a_{i,n}^*a_{j,n}^*\Omega_n.
 \label{eq:e8v42-cren}
\end{equation}
Then
\begin{equation}
 \|c_n^{\rm ren}\|^2
 ={1\over8}\|\mathcal R_n\|_{\rm HS}^2.
 \label{eq:e8v42-cren-norm}
\end{equation}
Consequently, row (ii) of
\Cref{def:e8v42-SIR} yields a cofinally bounded renormalized source
coefficient.
\end{theorem}

\begin{proof}
Apply the pair norm identity \eqref{eq:e8v42-pair-norm} to
\(\mathcal R_n\).
\end{proof}

\begin{firewall}[Subtraction is not deletion]
\label{fire:e8v42-subtraction}
The local term \(\mathcal K_n^{\rm loc}\) cannot simply be discarded.
It changes the comparison of fields, observables, measures, and stress
tensors between metrics.  It must reappear as a local counterterm,
connection, or relative-algebra contact, with Ward and anomaly
compatibility.  Only the Hilbert--Schmidt remainder may be represented by
the global two-particle vector in \eqref{eq:e8v42-cren}.
\end{firewall}

\subsection{Quadratic implementability card}
\label{sec:e8v42-card}

\begin{definition}[Quadratic local-response card, QLRC]
\label{def:e8v42-QLRC}
QLRC consists of:
\begin{enumerate}[label=\textup{(L\arabic*)},leftmargin=4em]
\item the one-particle positive operator \(A_g\), frequency
\(\omega_g=A_g^{1/2}\), and common finite-cutoff phase space;
\item the Sylvester derivative and kernel
\eqref{eq:e8v42-Kij};
\item a subtracted implementability row
\eqref{eq:e8v42-subtraction};
\item a local-algebra realization of the principal term, including its
stress and Jacobian contacts;
\item local or summable dependence of the HS remainder on the metric
replacement;
\item the V38 finite-horizon and V37 incidence compilers for the
renormalized remainder; and
\item interaction, constraint, boundary, reflection, spin, chiral,
anomaly, density, and gravitational-carrier contacts.
\end{enumerate}
\end{definition}

\begin{theorem}[QLRC supplies the corrected quadratic APSC row]
\label{thm:e8v42-QLRC}
QLRC rows (L1)--(L3) give a cofinally bounded order-one vector source for
the Hilbert--Schmidt remainder, with norm
\eqref{eq:e8v42-cren-norm}.  Row (L4) retains the nonimplementable
principal response as a local contact.  If (L5)--(L6) hold, the remainder
enters APSC and FHRC without charging the divergent principal kernel to a
global vacuum vector.
\end{theorem}

\begin{proof}
The exact finite-cutoff split is \eqref{eq:e8v42-subtraction}.
\Cref{thm:e8v42-remainder} controls the remainder.  The principal part is
kept by (L4), and V38--V37 apply to the remainder under (L5)--(L6).
\end{proof}

\begin{assessment}[Progress at V42]
\label{ass:e8v42-status}
The global vacuum-vector bottleneck is now classified in the free
quadratic branch.  The derivative exists with norm proportional to the
Hilbert--Schmidt norm of a weighted Sylvester kernel.  Every nonzero
second-order coefficient perturbation of the free massless torus fails
that criterion, regardless of smooth localization, whereas a
finite-Fourier zeroth-order perturbation passes it in \(d\leq3\) at fixed
volume.  The SM--Einstein route must therefore keep the local principal
metric response outside the global Fock-vector debit and control only a
subtracted remainder there.
\end{assessment}

\begin{programme}[Eighth-series V43: local observable response topology]
\label{prog:e8v42-V43}
V43 will construct a topology that retains the principal response rather
than representing it by a Fock vector.  It will derive a Duhamel formula
for the derivative of a local expectation, identify the commutator that
kills exterior principal contacts, and state the weighted trace or
Hilbert--Schmidt conditions needed for a cutoff-independent local
observable seminorm.
\end{programme}

\begin{firewall}[Claim boundary after eighth-series V42]
\label{fire:e8v42-boundary}
V42 proves the exact weighted Sylvester kernel, its Gaussian
order-one source factorization, the Hilbert--Schmidt criterion, the
principal-coefficient nonimplementability theorem, and the fixed-volume
lower-order comparison in \(d\leq3\).

V42 does not construct the principal-symbol subtraction, prove local
quasi-equivalence, control an interacting or constrained field, establish
the stress and anomaly contacts, construct the gravitational carrier,
remove the cutoff, or prove the full continuum.
\end{firewall}

\begin{theorem}[Eighth-series V42 result]
\label{thm:e8v42-result}
For a finite quadratic field with one-particle frequency
\(\omega=A^{1/2}\), the moving-vacuum coefficient has squared norm
\[
 \|H^{-1}Q\dot H\Omega\|^2
 ={1\over8}
 \left\|\omega^{-1/2}\dot\omega\,\omega^{-1/2}\right\|_{\rm HS}^2.
\]
For \(A=-\Delta\) on a torus, every nonzero perturbation
\(-\nabla\cdot(h\nabla)\) makes this Hilbert--Schmidt norm infinite.
A cofinal construction must isolate that local principal response and
apply vector estimates only to a Hilbert--Schmidt remainder.
\end{theorem}

\begin{proof}
Use \Cref{thm:e8v42-Gaussian-factor,
thm:e8v42-principal-obstruction,thm:e8v42-remainder}.
\end{proof}

\paragraph{Parent-version record.}
V42 was formed from
\texttt{Renewal\_SM\_Einstein\_Continuum\_Research\_Notes\_V41.tex}.
The V41 parent had (16\,746) logical source lines, (609\,008) bytes,
and SHA--256
\[
 \texttt{09FA15B0C17E94C9416A785B7D4122B4D39C1A9577B31E1E7F5F299133221BAB}.
\]
The next compulsory companion audit is V45.

% =============================================================================
% END APPENDED EIGHTH-SERIES V42 MATERIAL
% =============================================================================

% =============================================================================
% BEGIN APPENDED EIGHTH-SERIES V43 MATERIAL
% =============================================================================

\section{Observable response beyond the global vacuum-vector topology}
\label{sec:v8v43}

V42 proved that the principal metric response is generally not
Hilbert--Schmidt and therefore cannot be represented by a cofinal global
vacuum derivative in the original Fock space.  The same weighted
Sylvester kernel can nevertheless act boundedly on one-particle space.
This note pairs that bounded kernel with trace-class two-particle tests.
It also gives a gap-free Duhamel criterion in which the source and
observable contribute complementary spectral orders.

\subsection{Finite-cutoff expectation derivative}
\label{sec:e8v43-expectation}

Let \(H_s\geq0\) have a normalized simple vacuum \(\Omega_s\), let
\(Q_s=I-|\Omega_s\rangle\langle\Omega_s|\), and use the parallel gauge.
For a norm-\(C^1\) family of bounded self-adjoint observables \(O_s\), put
\begin{equation}
 \varphi_s(O_s)=\langle\Omega_s,O_s\Omega_s\rangle,
 \qquad
 f_s=Q_sO_s\Omega_s.
 \label{eq:e8v43-state}
\end{equation}

\begin{theorem}[Exact observable derivative]
\label{thm:e8v43-observable-derivative}
At a finite cutoff, if
\[
 b_s=Q_s\dot H_s\Omega_s,
 \qquad
 \dot\Omega_s=-S_sb_s
\]
with \(S_s\) the reduced inverse, then
\begin{equation}
 {d\over ds}\varphi_s(O_s)
 =
 \langle\Omega_s,\dot O_s\Omega_s\rangle
 -2\operatorname{Re}\langle S_sb_s,f_s\rangle.
 \label{eq:e8v43-exact-response}
\end{equation}
Only the vacuum-orthogonal excitation \(f_s\) of the observable enters
the response term.
\end{theorem}

\begin{proof}
Differentiate \(\langle\Omega_s,O_s\Omega_s\rangle\).  Self-adjointness of
\(O_s\) makes the two vacuum-derivative terms conjugate, so their sum is
\(2\operatorname{Re}\langle\dot\Omega_s,O_s\Omega_s\rangle\).
The parallel gauge gives
\(\dot\Omega_s\perp\Omega_s\), hence \(O_s\Omega_s\) may be replaced by
\(f_s\).  Insert \(\dot\Omega_s=-S_sb_s\).
\end{proof}

\begin{corollary}[Observables annihilating the source channel]
\label{cor:e8v43-annihilator}
If \(Q_sO_s\Omega_s=0\), then
\begin{equation}
 {d\over ds}\varphi_s(O_s)
 =\langle\Omega_s,\dot O_s\Omega_s\rangle.
 \label{eq:e8v43-annihilator}
\end{equation}
In particular, a metric-independent observable with
\(O_s\Omega_s\in\mathbb C\Omega_s\) has zero vacuum-motion response.
\end{corollary}

\begin{proof}
Set \(f_s=0\) in \eqref{eq:e8v43-exact-response}.
\end{proof}

\subsection{Finite-horizon response as a correlation integral}
\label{sec:e8v43-Duhamel}

For fixed \(s\), suppress the subscript and let
\(\mathcal T_T=\int_0^Te^{-tH}Q\,dt\) as in V38.  Define
\begin{equation}
 \mathscr D_T(O)
 :=
 \langle\Omega,\dot O\Omega\rangle
 -2\operatorname{Re}\langle\mathcal T_Tb,f\rangle.
 \label{eq:e8v43-DT}
\end{equation}

\begin{theorem}[Exact finite-horizon Duhamel response]
\label{thm:e8v43-Duhamel}
For every \(T>0\),
\begin{equation}
 \mathscr D_T(O)
 =
 \langle\Omega,\dot O\Omega\rangle
 -2\operatorname{Re}\int_0^T
   \langle e^{-tH}b,f\rangle\,dt.
 \label{eq:e8v43-correlation}
\end{equation}
For \(T'>T\),
\begin{equation}
 |\mathscr D_{T'}(O)-\mathscr D_T(O)|
 \leq2\int_T^{T'}
 |\langle e^{-tH}b,f\rangle|\,dt.
 \label{eq:e8v43-Cauchy-tail}
\end{equation}
Consequently, if
\begin{equation}
 \int_0^\infty|\langle e^{-tH}b,f\rangle|\,dt<\infty,
 \label{eq:e8v43-integrable-correlation}
\end{equation}
then \(\mathscr D_T(O)\) has a finite limit as \(T\to\infty\), whether or
not \(\mathcal T_Tb\) converges in Hilbert norm.
\end{theorem}

\begin{proof}
Insert the Bochner integral defining \(\mathcal T_T\) into
\eqref{eq:e8v43-DT}.  Boundedness of the finite-time semigroup justifies
moving the scalar product through the integral.  Subtract the formulas at
\(T'\) and \(T\), then apply the triangle inequality.  Absolute
integrability makes the scalar integrals Cauchy.
\end{proof}

\begin{theorem}[Complementary spectral orders]
\label{thm:e8v43-complementary-orders}
Suppose
\begin{equation}
 b=H^\alpha c,
 \qquad
 f=H^\beta d,
 \qquad
 \alpha,\beta\geq0,
 \qquad
 \rho:=\alpha+\beta>1,
 \label{eq:e8v43-orders}
\end{equation}
with \(c,d\in\mathcal H\) in the required domains.  Then, for \(t>0\),
\begin{equation}
 |\langle e^{-tH}b,f\rangle|
 \leq
 \left({\rho\over et}\right)^\rho
 \|c\|\,\|d\|,
 \label{eq:e8v43-correlation-decay}
\end{equation}
and
\begin{equation}
 |\mathscr D_\infty(O)-\mathscr D_T(O)|
 \leq
 {2\over\rho-1}
 \left({\rho\over e}\right)^\rho
 T^{1-\rho}\|c\|\,\|d\|.
 \label{eq:e8v43-Duhamel-tail}
\end{equation}
\end{theorem}

\begin{proof}
By spectral calculus and commutation of functions of \(H\),
\[
 |\langle e^{-tH}H^\alpha c,H^\beta d\rangle|
 \leq
 \|H^\rho e^{-tH}\|\,\|c\|\,\|d\|.
\]
The maximum of \(\lambda^\rho e^{-t\lambda}\) on
\([0,\infty)\) is \((\rho/(et))^\rho\).  This gives
\eqref{eq:e8v43-correlation-decay}.  Integrate it from \(T\) to infinity
and use \eqref{eq:e8v43-Cauchy-tail}.
\end{proof}

\begin{remark}[The observable can supply the missing order]
\label{rem:e8v43-test-order}
V38 placed all infrared order on the source.  The present theorem only
requires the sum \(\alpha+\beta>1\).  A conserved or derivative observable
may therefore suppress low energy strongly enough to make its response
finite even when the global vacuum-response vector is unavailable.
\end{remark}

\begin{firewall}[Equation residual still does not imply response convergence]
\label{fire:e8v43-residual}
The correction in \Cref{fire:e8v39-V38-correction} remains essential.
Small \(e^{-TH}b\) controls the differentiated equation, while
\eqref{eq:e8v43-integrable-correlation} controls an observable response.
One condition does not imply the other without a spectral relation
between \(b\) and \(f\).
\end{firewall}

\subsection{Pair-kernel test topology}
\label{sec:e8v43-pair}

Let \(\mathfrak h_N\) be an \(N\)-dimensional one-particle space and let
\(\mathfrak S_2^{\rm sym}(\mathfrak h_N)\) be its symmetric
Hilbert--Schmidt matrices.  Define
\begin{equation}
 \mathfrak p_N(K)
 ={1\over\sqrt2}\sum_{i,j}K_{ij}a_i^*a_j^*\Omega_N.
 \label{eq:e8v43-pair-map}
\end{equation}

\begin{lemma}[Pair map isometry]
\label{lem:e8v43-pair-isometry}
For symmetric \(K,L\),
\begin{equation}
 \langle\mathfrak p_N(K),\mathfrak p_N(L)\rangle
 =\operatorname{Tr}(K^*L).
 \label{eq:e8v43-pair-isometry}
\end{equation}
\end{lemma}

\begin{proof}
The canonical commutation relations give
\[
 \langle\Omega_N,a_ja_i a_k^*a_l^*\Omega_N\rangle
 =\delta_{ik}\delta_{jl}+\delta_{il}\delta_{jk}.
\]
Insert this in \eqref{eq:e8v43-pair-map}.  Symmetry makes the two
contractions equal, and the factor \(1/2\) cancels their sum.
\end{proof}

For the Gaussian metric response of V42,
\begin{equation}
 \dot\Omega_N=-{1\over2\sqrt2}\mathfrak p_N(\mathcal K_N).
 \label{eq:e8v43-Gaussian-dot}
\end{equation}
For a bounded self-adjoint test observable \(O_N\), suppose its
two-particle vacuum component has kernel \(L_{O,N}\):
\begin{equation}
 P_{2,N}Q_NO_N\Omega_N=\mathfrak p_N(L_{O,N}).
 \label{eq:e8v43-test-kernel}
\end{equation}

\begin{theorem}[Exact pair-sector expectation response]
\label{thm:e8v43-pair-response}
The contribution of the two-particle sector to the vacuum-motion
derivative is
\begin{equation}
 2\operatorname{Re}
 \langle\dot\Omega_N,P_{2,N}Q_NO_N\Omega_N\rangle
 =
 -{1\over\sqrt2}\operatorname{Re}
 \operatorname{Tr}(\mathcal K_N^*L_{O,N}).
 \label{eq:e8v43-trace-pairing}
\end{equation}
If \(\mathcal K_N\) is bounded and \(L_{O,N}\) is trace class, then
\begin{equation}
 \left|
 2\operatorname{Re}
 \langle\dot\Omega_N,P_{2,N}Q_NO_N\Omega_N\rangle
 \right|
 \leq
 {1\over\sqrt2}
 \|\mathcal K_N\|_{\rm op}\|L_{O,N}\|_1.
 \label{eq:e8v43-trace-bound}
\end{equation}
\end{theorem}

\begin{proof}
Insert \eqref{eq:e8v43-Gaussian-dot} and
\eqref{eq:e8v43-test-kernel} into the left side, then use
\Cref{lem:e8v43-pair-isometry}.  The second assertion is the duality
\(|\operatorname{Tr}(K^*L)|\leq\|K\|_{\rm op}\|L\|_1\).
\end{proof}

\begin{definition}[Pair-response seminorm]
\label{def:e8v43-seminorm}
On tests satisfying \eqref{eq:e8v43-test-kernel}, define
\begin{equation}
 \|O_N\|_{\rm pair,1}:=\|L_{O,N}\|_1.
 \label{eq:e8v43-seminorm}
\end{equation}
Components of \(O_N\Omega_N\) outside the two-particle sector belong to
other response rows and are not measured by this seminorm.
\end{definition}

\begin{corollary}[Principal response is continuous in the pair-test topology]
\label{cor:e8v43-continuity}
If
\(\sup_N\|\mathcal K_N\|_{\rm op}\leq C_{\mathcal K}\), then the
two-particle response functionals are uniformly bounded on the
pair-response test space:
\begin{equation}
 |\delta_N^{(2)}(O_N)|
 \leq {C_{\mathcal K}\over\sqrt2}
 \|O_N\|_{\rm pair,1}.
 \label{eq:e8v43-uniform-functional}
\end{equation}
No Hilbert--Schmidt bound on \(\mathcal K_N\) is required.
\end{corollary}

\begin{proof}
Apply \eqref{eq:e8v43-trace-bound}.
\end{proof}

\subsection{The principal metric kernel is operator bounded}
\label{sec:e8v43-principal}

Return to the massless torus and the kernel
\(\mathcal K_h\) in \eqref{eq:e8v42-Kh}.

\begin{lemma}[Pointwise Schur majorant]
\label{lem:e8v43-majorant}
For nonzero \(p,q\),
\begin{equation}
 \left|
 {p\cdot q\over\sqrt{|p||q|}(|p|+|q|)}
 \right|
 \leq {1\over2}.
 \label{eq:e8v43-majorant}
\end{equation}
Consequently,
\begin{equation}
 |(\mathcal K_h)_{pq}|
 \leq {1\over2}|\widehat h(p-q)|.
 \label{eq:e8v43-convolution-majorant}
\end{equation}
\end{lemma}

\begin{proof}
Cauchy--Schwarz gives \(|p\cdot q|\leq|p||q|\), and the arithmetic--geometric
mean inequality gives
\(\sqrt{|p||q|}\leq(|p|+|q|)/2\).  Combine them in
\eqref{eq:e8v42-Kh}.
\end{proof}

\begin{theorem}[Uniform operator bound for the non-HS principal kernel]
\label{thm:e8v43-principal-op}
If
\begin{equation}
 \|\widehat h\|_{\ell^1}
 :=\sum_{r\in\mathbb Z^d}|\widehat h(r)|<\infty,
 \label{eq:e8v43-h-l1}
\end{equation}
then every Fourier cutoff of \(\mathcal K_h\) satisfies
\begin{equation}
 \|\mathcal K_{h,N}\|_{\ell^2\to\ell^2}
 \leq {1\over2}\|\widehat h\|_{\ell^1}.
 \label{eq:e8v43-principal-op}
\end{equation}
The constant is independent of the number and size of retained modes.
\end{theorem}

\begin{proof}
For a sequence \(u\), \eqref{eq:e8v43-convolution-majorant} gives
\[
 |(\mathcal K_hu)(p)|
 \leq {1\over2}\sum_q|\widehat h(p-q)|\,|u(q)|.
\]
Young's convolution inequality on \(\mathbb Z^d\) bounds the
\(\ell^2\) norm by
\(\frac12\|\widehat h\|_{\ell^1}\|u\|_{\ell^2}\).
Compression to a finite Fourier set cannot increase this bound.
\end{proof}

\begin{corollary}[A cutoff-independent principal observable response]
\label{cor:e8v43-principal-response}
For a real coefficient \(h\) satisfying
\eqref{eq:e8v43-h-l1}, the principal Gaussian response obeys
\begin{equation}
 |\delta_N^{(2)}(O_N)|
 \leq {1\over2\sqrt2}\|\widehat h\|_{\ell^1}
 \|O_N\|_{\rm pair,1}.
 \label{eq:e8v43-principal-response}
\end{equation}
This holds although
\(\|\mathcal K_{h,N}\|_{\rm HS}\to\infty\) for every nonzero \(h\).
\end{corollary}

\begin{proof}
Combine \Cref{thm:e8v43-principal-op,cor:e8v43-continuity}; the divergence
statement is \Cref{thm:e8v42-principal-obstruction}.
\end{proof}

\subsection{Cofinal trace-dual limit}
\label{sec:e8v43-limit}

\begin{theorem}[Compression limit in bounded--trace duality]
\label{thm:e8v43-compression}
Let \(\mathfrak h\) be separable, let \(P_N\to I\) strongly be increasing
finite-rank projections, let \(K\in\mathcal B(\mathfrak h)\), and let
\(L\in\mathfrak S_1(\mathfrak h)\).  Put
\begin{equation}
 K_N=P_NKP_N,
 \qquad L_N=P_NLP_N.
 \label{eq:e8v43-compressions}
\end{equation}
Then
\begin{equation}
 \operatorname{Tr}(K_N^*L_N)
 \longrightarrow
 \operatorname{Tr}(K^*L).
 \label{eq:e8v43-trace-limit}
\end{equation}
\end{theorem}

\begin{proof}
Strong convergence of uniformly bounded \(P_N\) implies
\begin{equation}
 \|P_NLP_N-L\|_1\longrightarrow0
 \label{eq:e8v43-trace-compression}
\end{equation}
for every trace-class \(L\).  To see this, approximate \(L\) in trace norm
by a finite-rank operator and use uniform boundedness; convergence is
norm convergence on the finite-dimensional range of the approximant.
Now
\[
 \operatorname{Tr}(K_N^*L_N)
 =\operatorname{Tr}(K^*P_NL_NP_N)
 =\operatorname{Tr}(K^*L_N).
\]
Trace duality and \eqref{eq:e8v43-trace-compression} give the limit.
\end{proof}

\begin{corollary}[Cofinal principal response on a fixed test]
\label{cor:e8v43-cofinal}
If the cutoff kernels are compressions of the bounded operator
\(\mathcal K_h\) and a test has a fixed symmetric trace-class kernel
\(L_O\), then its two-particle response converges to
\begin{equation}
 \delta^{(2)}_h(O)
 =-{1\over\sqrt2}\operatorname{Re}
   \operatorname{Tr}(\mathcal K_h^*L_O).
 \label{eq:e8v43-limit-functional}
\end{equation}
\end{corollary}

\begin{proof}
Use \Cref{thm:e8v43-pair-response,thm:e8v43-compression}.
\end{proof}

\begin{firewall}[Trace-class tests are not yet the physical local algebra]
\label{fire:e8v43-local-tests}
V43 proves a sufficient functional topology.  It does not prove that the
gauge-invariant Standard Model observables, renormalized stress tensor, or
Einstein response insertions have uniformly trace-class pair kernels.
That membership must be checked after smearing, constraints, spin and
gauge projection, and renormalization.  Describing an observable as local
does not by itself give \(\|L_O\|_1<\infty\).
\end{firewall}

\subsection{Observable-response card}
\label{sec:e8v43-card}

\begin{definition}[Local observable response card, LORC]
\label{def:e8v43-LORC}
LORC consists of:
\begin{enumerate}[label=\textup{(O\arabic*)},leftmargin=4em]
\item a protected test algebra and its vacuum excitations
\(f_O=QO\Omega\);
\item either the integrable correlation row
\eqref{eq:e8v43-integrable-correlation} or complementary spectral orders
with sum greater than one;
\item for a quadratic principal response, a uniformly bounded Sylvester
kernel and a uniformly trace-class test kernel;
\item compatibility of the cutoff kernels as compressions, or an
equivalent weak-star/trace-dual convergence row;
\item local replacement and exterior-tail bounds for the response
functional;
\item the V35 conditional-centering and V37 coefficient-incidence rows;
and
\item Ward, constraint, boundary, reflection, spin, chiral, anomaly,
stress-renormalization, and gravitational-carrier contacts.
\end{enumerate}
\end{definition}

\begin{theorem}[LORC replaces the divergent vector debit]
\label{thm:e8v43-LORC}
Rows (O1)--(O4) define a cutoff-independent derivative of every declared
test expectation.  In the quadratic principal branch it is the
trace-dual functional \eqref{eq:e8v43-limit-functional}; in the
finite-horizon branch it is the limit of
\eqref{eq:e8v43-correlation}.  If (O5)--(O6) hold with summable constants,
the local observable response enters the IMBSC handoff without requiring
a cofinal global vacuum derivative.
\end{theorem}

\begin{proof}
The finite-horizon conclusion is
\Cref{thm:e8v43-Duhamel,thm:e8v43-complementary-orders}.
The quadratic conclusion is
\Cref{cor:e8v43-principal-response,cor:e8v43-cofinal}.  The remaining
rows are exactly the local incidence and conditional-centering inputs
needed by the earlier compiler.
\end{proof}

\begin{assessment}[Progress at V43]
\label{ass:e8v43-status}
The non-Hilbert--Schmidt principal metric response now has a rigorous
observable topology.  Its kernel is uniformly operator bounded for a
smooth coefficient with summable Fourier series, and hence acts
continuously on trace-class pair tests with a cutoff-independent limit.
Independently, a heat-correlation response converges whenever the source
and observable spectral orders sum to more than one.  The next task is to
show that physically useful smeared energy--momentum tests belong to one
of these classes after their local counterterms are fixed.
\end{assessment}

\begin{programme}[Eighth-series V44: smeared quadratic stress tests]
\label{prog:e8v43-V44}
V44 will compute the pair kernel of a spatially smeared, normal-ordered
quadratic observable in the free scalar branch.  It will determine the
derivative count needed for trace-class membership and test whether
stress-tensor smearing supplies it.  Any remaining diagonal ultraviolet
term will be retained as a local counterterm rather than hidden in the
response functional.
\end{programme}

\begin{firewall}[Claim boundary after eighth-series V43]
\label{fire:e8v43-boundary}
V43 proves the exact expectation derivative, the finite-horizon
correlation criterion, the complementary-order tail, bounded--trace
pair-kernel duality, the uniform operator bound for the free principal
metric kernel, and its compression limit.

V43 does not prove trace-class membership for the physical Standard Model
or stress tests, fix the required local counterterms, handle interactions
or constraints, construct the conditional gravitational carrier, remove
the cutoff for the full theory, or prove the continuum.
\end{firewall}

\begin{theorem}[Eighth-series V43 result]
\label{thm:e8v43-result}
The principal free metric-response kernel is generally not
Hilbert--Schmidt but obeys
\[
 \|\mathcal K_h\|_{\rm op}
 \leq {1\over2}\|\widehat h\|_{\ell^1}.
\]
It therefore defines the cofinal response
\[
 \delta_h^{(2)}(O)
 =-{1\over\sqrt2}\operatorname{Re}
   \operatorname{Tr}(\mathcal K_h^*L_O)
\]
for every fixed trace-class pair test \(L_O\).  More generally, the
finite-horizon response converges when the source and test spectral orders
have sum greater than one.
\end{theorem}

\begin{proof}
Use \Cref{thm:e8v43-complementary-orders,
thm:e8v43-principal-op,thm:e8v43-compression}.
\end{proof}

\paragraph{Parent-version record.}
V43 was formed from
\texttt{Renewal\_SM\_Einstein\_Continuum\_Research\_Notes\_V42.tex}.
The V42 parent had (17\,241) logical source lines, (625\,969) bytes,
and SHA--256
\[
 \texttt{2CF24AF8A71040934D0D8CBB0A37DF73A75D71C1134DE4AA6135369D408177AF}.
\]
The next compulsory companion audit is V45.

% =============================================================================
% END APPENDED EIGHTH-SERIES V43 MATERIAL
% =============================================================================

% =============================================================================
% BEGIN APPENDED EIGHTH-SERIES V44 MATERIAL
% =============================================================================

\section{Spacetime-smeared quadratic tests and Wick contacts}
\label{sec:v8v44}

V43 made the principal metric response continuous on observables whose
two-particle vacuum kernels are trace class.  This note tests that
condition for free quadratic fields.  Pure spatial smearing at one time
does not in general suppress pairs with large opposite momenta.
Smooth spacetime smearing does: the positive pair energy
\(\omega_p+\omega_q\) supplies decay in the missing relative direction.
The derivative of the normal-ordering subtraction remains as a separate
local contact.

\subsection{A general quadratic pair symbol}
\label{sec:e8v44-symbol}

On the massless torus of V42, continue to omit the zero mode and write
\(\omega_p=|p|\).  Let a quadratic observable, after normal ordering at a
fixed cutoff, have two-particle kernel
\begin{equation}
 L_g(p,q)
 =a(p,q)\,
 \widehat g(\omega_p+\omega_q,p+q),
 \label{eq:e8v44-L}
\end{equation}
where \(a(p,q)=a(q,p)\) is its pair symbol and
\(\widehat g(\tau,k)\) is the time Fourier transform and spatial Fourier
coefficient of a test \(g(t,x)\).  Assume
\begin{equation}
 |a(p,q)|
 \leq C_a(1+|p|+|q|)^r
 \label{eq:e8v44-symbol-order}
\end{equation}
for some \(r\geq0\).  This conservative class includes stress-type
quadratic symbols with \(r=1\).

\begin{lemma}[Spacetime Fourier decay]
\label{lem:e8v44-Fourier-decay}
If \(g\in C_c^\infty(\mathbb R\times\mathbb T^d)\), then for every
integer \(M\geq0\) there is \(C_{g,M}<\infty\) such that
\begin{equation}
 |\widehat g(\tau,k)|
 \leq C_{g,M}(1+|\tau|+|k|)^{-M}.
 \label{eq:e8v44-Fourier-decay}
\end{equation}
\end{lemma}

\begin{proof}
Integrate by parts in \(t\) and in each periodic spatial coordinate.
Every time derivative produces a factor \(\tau\), and every spatial
derivative produces the corresponding component of \(k\).  A finite sum
of derivatives of total order at most \(M+d+1\), together with
compact time support, bounds all regions in frequency space and yields
\eqref{eq:e8v44-Fourier-decay}.
\end{proof}

\begin{lemma}[Absolute matrix summability implies trace class]
\label{lem:e8v44-l1-trace}
Let \(L=(L_{pq})\) be a matrix on \(\ell^2\) and suppose
\(\sum_{p,q}|L_{pq}|<\infty\).  Then \(L\) is trace class and
\begin{equation}
 \|L\|_1\leq\sum_{p,q}|L_{pq}|.
 \label{eq:e8v44-l1-trace}
\end{equation}
\end{lemma}

\begin{proof}
Write
\[
 L=\sum_{p,q}L_{pq}|e_p\rangle\langle e_q|.
\]
Each matrix unit has trace norm one.  The assumed absolute convergence
makes this series converge in trace norm, and the triangle inequality
gives \eqref{eq:e8v44-l1-trace}.
\end{proof}

\begin{theorem}[Smooth spacetime smearing gives a trace-class pair test]
\label{thm:e8v44-trace}
Under \eqref{eq:e8v44-symbol-order}, if
\begin{equation}
 M>2d+r,
 \label{eq:e8v44-M-condition}
\end{equation}
then
\begin{equation}
 \|L_g\|_1
 \leq
 C_aC_{g,M}
 \sum_{p,q\neq0}
 (1+|p|+|q|)^{r-M}
 <\infty.
 \label{eq:e8v44-trace-bound}
\end{equation}
In particular, a stress-type symbol of order at most one is trace class
once a Fourier-decay order \(M>2d+1\) is used.
\end{theorem}

\begin{proof}
Since
\[
 \omega_p+\omega_q=|p|+|q|,
\]
\eqref{eq:e8v44-Fourier-decay} and
\eqref{eq:e8v44-symbol-order} give the pointwise majorant in
\eqref{eq:e8v44-trace-bound}.  The number of pairs
\((p,q)\in\mathbb Z^{2d}\) with
\(n\leq|p|+|q|<n+1\) is at most \(C_d(1+n)^{2d-1}\).
The shell series is therefore bounded by a constant times
\[
 \sum_{n\geq0}(1+n)^{2d-1+r-M},
\]
which converges precisely under the sufficient condition
\eqref{eq:e8v44-M-condition}.  Apply
\Cref{lem:e8v44-l1-trace}.
\end{proof}

\begin{corollary}[Cutoff-independent principal metric response]
\label{cor:e8v44-response}
Let \(h\) satisfy \(\widehat h\in\ell^1\), and let \(O_g\) have pair
kernel \eqref{eq:e8v44-L} with \(g\) and \(M\) as in
\Cref{thm:e8v44-trace}.  Then
\begin{equation}
 |\delta_h^{(2)}(O_g)|
 \leq {1\over2\sqrt2}\|\widehat h\|_{\ell^1}
 \|L_g\|_1,
 \label{eq:e8v44-response}
\end{equation}
and the compressed cutoff responses converge to the trace pairing
\eqref{eq:e8v43-limit-functional}.
\end{corollary}

\begin{proof}
Use \Cref{cor:e8v43-principal-response,thm:e8v43-compression} and
\Cref{thm:e8v44-trace}.
\end{proof}

\subsection{Why fixed-time spatial smearing fails}
\label{sec:e8v44-fixed-time}

For the normal-ordered square of the free field at one time, spatially
smeared by \(f\), the pair kernel has the form
\begin{equation}
 L_f^{\phi^2}(p,q)
 ={c_d\,\widehat f(p+q)\over\sqrt{|p||q|}},
 \label{eq:e8v44-phi2}
\end{equation}
where \(c_d\neq0\) depends only on Fourier normalization.

\begin{theorem}[Fixed-time \(\phi^2\) pair kernel is not HS in \(d\geq2\)]
\label{thm:e8v44-fixed-time-obstruction}
If \(f\not\equiv0\) and \(d\geq2\), then
\begin{equation}
 L_f^{\phi^2}\notin\mathfrak S_2,
 \qquad\text{hence}\qquad
 L_f^{\phi^2}\notin\mathfrak S_1.
 \label{eq:e8v44-not-HS}
\end{equation}
Smoothness of \(f\) does not change this conclusion.
\end{theorem}

\begin{proof}
Choose \(r\in\mathbb Z^d\) with \(\widehat f(r)\neq0\), and set
\(q=r-p\).  Along this affine diagonal,
\[
 |L_f^{\phi^2}(p,r-p)|^2
 ={|c_d\widehat f(r)|^2\over |p||r-p|}
 \asymp |p|^{-2}
\]
as \(|p|\to\infty\).  The lattice sum
\(\sum_{p\in\mathbb Z^d\setminus\{0,r\}}|p|^{-2}\) diverges for
\(d\geq2\).  This is a subseries of the Hilbert--Schmidt norm.  Trace
class is contained in Hilbert--Schmidt, so the second conclusion follows.
\end{proof}

\begin{corollary}[Fixed-time stress tests are no better by power counting]
\label{cor:e8v44-stress-fixed}
A fixed-time quadratic stress component has a pair symbol with additional
momentum powers relative to \eqref{eq:e8v44-phi2}.  Unless an exact tensor
cancellation removes its leading affine-diagonal symbol, spatial
smearing alone cannot place its pair kernel in trace class in
\(d\geq2\).
\end{corollary}

\begin{proof}
On \(q=r-p\), every uncancelled derivative pair contributes a
nonnegative power of \(|p|\) multiplying the \(|p|^{-1}\) amplitude in
\eqref{eq:e8v44-phi2}.  Its squared affine-diagonal subseries therefore
cannot converge when the \(\phi^2\) subseries already diverges.  The
qualification is necessary because a particular linear combination of
stress components may cancel its leading symbol.
\end{proof}

\begin{firewall}[Spatial locality is not energy smoothing]
\label{fire:e8v44-spatial}
Rapid decay of \(\widehat f(p+q)\) controls total spatial momentum.  It
does not control the relative momentum along \(q=r-p\).  Spacetime
smearing evaluates the time transform at the positive pair energy
\(|p|+|q|\), which does grow along that diagonal.  Replacing spacetime
smearing by fixed-time locality would lose the mechanism used in
\Cref{thm:e8v44-trace}.
\end{firewall}

\subsection{A positive Euclidean-time buffer}
\label{sec:e8v44-buffer}

An equivalent source of pair-energy damping is a positive Euclidean-time
separation.  Let \(O\) have finite-cutoff pair kernel \(L_O(p,q)\), and
define
\begin{equation}
 O_\tau=e^{-\tau H}Oe^{-\tau H},
 \qquad \tau>0.
 \label{eq:e8v44-buffer}
\end{equation}

\begin{theorem}[Heat buffering makes polynomial pair symbols trace class]
\label{thm:e8v44-buffer}
If
\begin{equation}
 |L_O(p,q)|\leq C(1+|p|+|q|)^r,
 \label{eq:e8v44-polynomial-L}
\end{equation}
then the two-particle kernel of \(O_\tau\) is
\begin{equation}
 L_{O_\tau}(p,q)
 =e^{-\tau(|p|+|q|)}L_O(p,q)
 \label{eq:e8v44-buffer-kernel}
\end{equation}
and is trace class for every \(\tau>0\).
\end{theorem}

\begin{proof}
The vacuum is fixed by \(e^{-\tau H}\), while the two-particle vector
\(a_p^*a_q^*\Omega\) is multiplied by
\(e^{-\tau(|p|+|q|)}\).  This proves
\eqref{eq:e8v44-buffer-kernel}.  Its absolute matrix sum is bounded by
\[
 C\sum_{p,q}
 e^{-\tau(|p|+|q|)}(1+|p|+|q|)^r<\infty.
\]
Apply \Cref{lem:e8v44-l1-trace}.
\end{proof}

\begin{remark}[The buffer is a declared resolution scale]
\label{rem:e8v44-buffer}
The trace norm generally diverges as \(\tau\downarrow0\).  A cofinal
programme may use \(\tau_n>0\) only with an explicit schedule that balances
this divergence against scale amplitudes and separation tails.  A fixed
positive \(\tau\) defines a smeared observable, not a sharp-time stress
tensor.
\end{remark}

\subsection{The Wick-ordering contact}
\label{sec:e8v44-Wick}

At cutoff \(N\), let \(B_N(g)\) be a bare quadratic expression and define
its vacuum-normal-ordered version by
\begin{equation}
 :B_N(g):_s
 =B_N(g)-C_{N,s}(g)I,
 \qquad
 C_{N,s}(g)
 =\langle\Omega_{N,s},B_N(g)\Omega_{N,s}\rangle.
 \label{eq:e8v44-Wick}
\end{equation}

\begin{theorem}[Exact cancellation of the vacuum Wick contact]
\label{thm:e8v44-Wick-contact}
For every finite cutoff,
\begin{equation}
 \langle\Omega_{N,s},:B_N(g):_s\Omega_{N,s}\rangle=0
 \label{eq:e8v44-Wick-zero}
\end{equation}
and
\begin{equation}
 \dot C_{N,s}(g)
 =
 \langle\Omega_{N,s},\dot B_N(g)\Omega_{N,s}\rangle
 +2\operatorname{Re}
 \langle\dot\Omega_{N,s},
        Q_{N,s}B_N(g)\Omega_{N,s}\rangle.
 \label{eq:e8v44-Wick-contact}
\end{equation}
Thus the explicit derivative of the subtraction cancels the vacuum-motion
and bare-observable derivatives in the expectation of the moving
normal-ordered observable.
\end{theorem}

\begin{proof}
Equation \eqref{eq:e8v44-Wick-zero} is the definition of
\(C_{N,s}(g)\).  Differentiate that definition and use self-adjointness as
in \Cref{thm:e8v43-observable-derivative}.  The component of
\(B_N(g)\Omega_{N,s}\) parallel to the vacuum does not contribute because
the vacuum is in parallel gauge.  This gives
\eqref{eq:e8v44-Wick-contact}.
\end{proof}

\begin{firewall}[A divergent contact still needs renormalization]
\label{fire:e8v44-renormalization}
The exact finite-cutoff cancellation does not prove that
\(C_{N,s}(g)\) or its derivative converges.  A local counterterm
prescription must identify and subtract the divergent part, and its
finite remainder must satisfy covariance, conservation, Ward, and anomaly
conditions.  The trace-dual response controls the nonlocal pair term; it
does not choose the stress-tensor renormalization.
\end{firewall}

\subsection{Smeared stress-test card}
\label{sec:e8v44-card}

\begin{definition}[Smeared stress-test card, SSTC]
\label{def:e8v44-SSTC}
SSTC consists of:
\begin{enumerate}[label=\textup{(S\arabic*)},leftmargin=4em]
\item a gauge-invariant quadratic or renormalized stress insertion and
its pair symbol order \(r\);
\item smooth spacetime smearing with decay order \(M>2d+r\), or a positive
Euclidean-time buffer;
\item the trace-norm estimate
\eqref{eq:e8v44-trace-bound} and compatibility of cutoff compressions;
\item an explicit local Wick or stress counterterm and its metric
derivative;
\item a scale schedule if the time buffer tends to zero;
\item V43 trace-dual response, V37 replacement incidence, and V35
conditional centering; and
\item constraint, boundary, reflection, spin, chiral, anomaly,
conservation, interaction, and gravitational-carrier contacts.
\end{enumerate}
\end{definition}

\begin{theorem}[SSTC populates the free smeared LORC row]
\label{thm:e8v44-SSTC}
Rows (S1)--(S3) place the quadratic test in the V43 pair-response space
with a cutoff-independent principal response.  Row (S4) retains the local
renormalization contact required by
\Cref{thm:e8v44-Wick-contact}.  If (S5)--(S6) have summable constants, the
smeared free response enters LORC.  Row (S7) is not supplied by the free
calculation.
\end{theorem}

\begin{proof}
Trace-class membership is \Cref{thm:e8v44-trace} or
\Cref{thm:e8v44-buffer}.  The response bound is
\Cref{cor:e8v44-response}; the explicit contact is
\Cref{thm:e8v44-Wick-contact}.  The remaining compiler rows are those
named in the definition.
\end{proof}

\begin{assessment}[Progress at V44]
\label{ass:e8v44-status}
The V43 test topology is nonempty for physically recognizable quadratic
symbols.  Smooth spacetime smearing, or any fixed positive Euclidean-time
buffer, makes every polynomial pair symbol trace class.  Pure spatial
smearing of \(:\!\phi^2\!:\) fails even the Hilbert--Schmidt test in
dimensions at least two, and a generic fixed-time stress symbol is worse.
The local derivative of Wick ordering is now explicit and cannot be
absorbed into the nonlocal response bound.
\end{assessment}

\begin{programme}[Eighth-series V45: companion audit and response topology]
\label{prog:e8v44-V45}
V45 will perform the compulsory YM/NS audit.  It will check whether the
newest YM collar work supplies a boundary-compatible local response
topology or counterterm ownership rule, then consolidate the smallest
free-to-interacting handoff left after V42--V44.
\end{programme}

\begin{firewall}[Claim boundary after eighth-series V44]
\label{fire:e8v44-boundary}
V44 proves a sufficient trace-class criterion for spacetime-smeared
quadratic pair symbols, the fixed-time \(\phi^2\) obstruction, the
positive-time buffer theorem, and the exact finite-cutoff Wick-contact
identity.

V44 does not construct a renormalized interacting stress tensor, prove
sharp-time convergence, establish gauge or diffeomorphism Ward identities,
control the chiral and boundary sectors, construct the gravitational
carrier, remove the full cutoff, or prove the continuum.
\end{firewall}

\begin{theorem}[Eighth-series V44 result]
\label{thm:e8v44-result}
If a quadratic pair symbol grows at most like
\((1+|p|+|q|)^r\), smooth spacetime smearing with Fourier decay order
\(M>2d+r\) makes its pair kernel trace class and hence admissible for the
V43 principal-response functional.  Spatial smearing alone does not do
so for \(\phi^2\) in \(d\geq2\).  Differentiating vacuum normal ordering
produces the exact local contact \eqref{eq:e8v44-Wick-contact}.
\end{theorem}

\begin{proof}
Use \Cref{thm:e8v44-trace,thm:e8v44-fixed-time-obstruction,
thm:e8v44-Wick-contact}.
\end{proof}

\paragraph{Parent-version record.}
V44 was formed from
\texttt{Renewal\_SM\_Einstein\_Continuum\_Research\_Notes\_V43.tex}.
The V43 parent had (17\,794) logical source lines, (644\,216) bytes,
and SHA--256
\[
 \texttt{7B1157B7F0D312C701A685D138E3BC2B5351A7773290F3D576795579E760EC9D}.
\]
The next compulsory companion audit is V45.

% =============================================================================
% END APPENDED EIGHTH-SERIES V44 MATERIAL
% =============================================================================

% =============================================================================
% BEGIN APPENDED EIGHTH-SERIES V45 MATERIAL
% =============================================================================

\section{V45 companion audit and boundary-influence resolvents}
\label{sec:v8v45}

This compulsory audit follows V40.  Yang--Mills has frozen its sixth
series at V100 and restarted at a compact V1.  Its new result is directly
relevant to the SM block compiler: no nonempty finite link set is closed
under recursive Wilson-action incidence.  A finite first contact collar
is useful, but boundary defects must be owned and propagated by a strict
influence mechanism.  This note proves that abstract resolvent compiler.

\subsection{Compulsory companion audit}
\label{sec:e8v45-audit}

\begin{audit}[Companion files read at V45]
\label{audit:e8v45-files}
The newest active files are:
\begin{enumerate}[label=\textup{(\roman*)},leftmargin=3.8em]
\item
\texttt{Renewal\_Yang\_Mills\_Mass\_Gap\_Research\_Notes\_V1.tex},
seventh series, with \(703\) lines, \(28\,473\) bytes, one terminator, and
SHA--256
\[
 \texttt{73E8194BC4415A65ED8B5A2F9C50FBE1EFBC4D916AD62A28EF72B1CD07DEEBAF};
\]
\item
\texttt{Renewal\_NS\_Stability\_Research\_Notes\_V26.tex},
with \(5\,319\) lines, \(278\,283\) bytes, one terminator, and SHA--256
\[
 \texttt{82C887F8C2C69E4F28FAD7C3DC8DE5B9099CB5A4BF431EBA1FFF3E91935978AD}.
\]
\end{enumerate}
The YM V1 records its parent
\texttt{Renewal\_Yang\_Mills\_Mass\_Gap\_Research\_Notes\_V100\_f6.tex}
with \(47\,291\) lines, \(1\,647\,389\) bytes, and digest
\[
 \texttt{63A6EE472D2127FF2D928791F1C720ACFF1DBB5193047C2D3CEA14F7DB759A4A}.
\]
\end{audit}

\begin{theorem}[New Yang--Mills input selected by V45]
\label{thm:e8v45-YM-input}
The audited YM V100 and compact V1 prove:
\begin{enumerate}[label=\textup{(\roman*)},leftmargin=3.8em]
\item triviality of the buffered collar presentation, uniqueness of flat
chord completion, and global small-defect reconstruction;
\item nonexistence of a nonempty finite Wilson-action-closed link set on
\(\mathbb Z^d\), \(d\geq2\);
\item an exact first incident shell for the \(60\)-chord update with
\(250\) plaquettes, \(178\) beyond the principal \(72\), and \(259\) new
links beyond the principal \(123\);
\item incompatibility of a termwise bounded-oscillation repair with the
needed weak-coupling action scale;
\item uniform Hessian coercivity when both principal and crossing Wilson
defects are small; and
\item an open boundary bad-event capacity and shell-influence problem.
\end{enumerate}
\end{theorem}

\begin{proof}
These are the companion results
\texttt{thm:s6v100-pi1},
\texttt{thm:s6v100-global},
\texttt{thm:s7v1-no-finite-closure},
\texttt{prop:s7v1-crossing-count},
\texttt{prop:s7v1-Holley-no-go}, and
\texttt{thm:s7v1-complete-Hessian}, with their BWIC firewalls retained.
\end{proof}

\begin{theorem}[Navier--Stokes status at V45]
\label{thm:e8v45-NS-input}
The NS note and digest are unchanged.  Its rank-one Oseen-kernel
calculation supplies no new shell ownership, observable response,
gravitational capacity, or contraction theorem.
\end{theorem}

\begin{proof}
Compare \Cref{audit:e8v45-files,audit:e8v40-files}.
\end{proof}

\begin{firewall}[Meaning of a finite internal patch]
\label{fire:e8v45-patch}
The V36 finite patch theorem remains correct for the finite list of terms
touching a fixed source support.  It cannot assert that a nonempty updated
block contains every interaction touching every variable recursively
added to that block.  The YM theorem disproves that stronger reading for
Wilson plaquettes.  Hereafter an internal patch means a finite owned first
collar plus an explicit exterior influence tail.
\end{firewall}

\subsection{Boundary influence equation}
\label{sec:e8v45-resolvent}

Let \(\mathcal I\) index translated blocks and let
\(\ell^2(\mathcal I;\mathcal X)\) contain response debits in a Hilbert
space \(\mathcal X\).  An owner rule assigns each crossing contact once.
Let \(a\) be the owned boundary source and let \(K\) propagate it.

\begin{theorem}[Boundary influence resolvent]
\label{thm:e8v45-resolvent}
If
\begin{equation}
 \|K\|\leq\kappa<1,
 \label{eq:e8v45-kappa}
\end{equation}
then
\begin{equation}
 (I-K)^{-1}=\sum_{m=0}^\infty K^m,
 \qquad
 \|(I-K)^{-1}\|\leq{1\over1-\kappa}.
 \label{eq:e8v45-Neumann}
\end{equation}
The equation \(e=a+Ke\) has the unique solution
\begin{equation}
 e=\sum_{m=0}^\infty K^ma,
 \qquad
 \|e\|\leq{1\over1-\kappa}\|a\|.
 \label{eq:e8v45-solution}
\end{equation}
Its \(M\)-shell truncation has error at most
\begin{equation}
 {\kappa^M\over1-\kappa}\|a\|.
 \label{eq:e8v45-truncation}
\end{equation}
\end{theorem}

\begin{proof}
The Neumann series converges in operator norm.  Its partial sums multiply
with \(I-K\) to \(I-K^{M+1}\), which tends to \(I\).  Apply the inverse to
the boundary equation.  The geometric tail gives
\eqref{eq:e8v45-truncation}.
\end{proof}

\subsection{Owned-overlap Schur criterion}
\label{sec:e8v45-Schur}

For an operator-valued influence matrix \(K=(K_{ij})\), put
\begin{equation}
 r_*=\sup_i\sum_j\|K_{ij}\|,
 \qquad
 c_*=\sup_j\sum_i\|K_{ij}\|.
 \label{eq:e8v45-Schur-constants}
\end{equation}

\begin{theorem}[Operator-valued Schur bound]
\label{thm:e8v45-Schur}
If \(r_*,c_*<\infty\), then
\begin{equation}
 \|K\|_{\ell^2(\mathcal I;\mathcal X)\to
          \ell^2(\mathcal I;\mathcal X)}
 \leq\sqrt{r_*c_*}.
 \label{eq:e8v45-Schur}
\end{equation}
Thus \(\sqrt{r_*c_*}<1\) is sufficient for
\Cref{thm:e8v45-resolvent}.
\end{theorem}

\begin{proof}
For \(x=(x_j)\), weighted Cauchy--Schwarz gives
\[
 \|(Kx)_i\|^2
 \leq
 \left(\sum_j\|K_{ij}\|\right)
 \left(\sum_j\|K_{ij}\|\|x_j\|^2\right).
\]
Sum in \(i\), use the row bound, interchange the nonnegative sums, and
use the column bound.  This gives
\(\|Kx\|^2\leq r_*c_*\|x\|^2\).
\end{proof}

\begin{remark}[Ownership controls multiplicity]
\label{rem:e8v45-owner}
A translation-covariant owner map charges each crossing contact once.
Other incidences remain as entries of its influence column.  Ownership
does not make the contact local, but makes both Schur sums auditable.
\end{remark}

\begin{theorem}[Explicit exponential graph criterion]
\label{thm:e8v45-exponential}
Suppose the block graph has maximum degree \(D\geq2\) and
\begin{equation}
 \|K_{ij}\|\leq\kappa_0e^{-\mu d(i,j)}.
 \label{eq:e8v45-exponential-kernel}
\end{equation}
If \((D-1)e^{-\mu}<1\), then
\begin{equation}
 r_*,c_*
 \leq\kappa_0\mathfrak S_D(\mu),
 \qquad
 \mathfrak S_D(\mu)
 =1+{De^{-\mu}\over1-(D-1)e^{-\mu}}.
 \label{eq:e8v45-SD}
\end{equation}
Hence \(\kappa_0\mathfrak S_D(\mu)<1\) implies
\[
 \|(I-K)^{-1}\|
 \leq{1\over1-\kappa_0\mathfrak S_D(\mu)}.
\]
\end{theorem}

\begin{proof}
There are at most \(D(D-1)^{m-1}\) blocks at distance \(m\geq1\).
Sum \eqref{eq:e8v45-exponential-kernel} over shells to obtain
\eqref{eq:e8v45-SD} for rows and columns.  Apply
\Cref{thm:e8v45-Schur,thm:e8v45-resolvent}.
\end{proof}

\begin{corollary}[Gaussian finite-horizon criterion]
\label{cor:e8v45-Gaussian}
If
\[
 \|K_{ij}\|\leq\kappa_0e^{-a d(i,j)^2},
 \qquad a>0,
\]
then
\begin{equation}
 r_*,c_*
 \leq\kappa_0\mathfrak Q_D(a),
 \quad
 \mathfrak Q_D(a)
 :=1+D\sum_{m\geq1}(D-1)^{m-1}e^{-am^2}<\infty.
 \label{eq:e8v45-QD}
\end{equation}
The strict condition
\(\kappa_0\mathfrak Q_D(a)<1\) yields the boundary resolvent.
\end{corollary}

\begin{proof}
Use the same sphere count.  The ratio of successive summands tends to
zero, and \Cref{thm:e8v45-Schur} applies.
\end{proof}

\subsection{Cofinal response handoff}
\label{sec:e8v45-SM}

\begin{theorem}[Owned shells preserve the response topology]
\label{thm:e8v45-response}
Let \(\mathcal X\) be either the V43 trace-dual test-response space or the
V37 coefficient space for a renormalized Hilbert--Schmidt remainder.
If the direct source is \(a_n\) and
\(\|K_n\|\leq\kappa_n<1\), then the total cutoff-\(n\) boundary response
obeys
\begin{equation}
 \|e_n\|\leq{\|a_n\|\over1-\kappa_n}.
 \label{eq:e8v45-en}
\end{equation}
The cumulative response converges absolutely whenever
\begin{equation}
 \sum_n{\|a_n\|\over1-\kappa_n}<\infty.
 \label{eq:e8v45-cofinal-sum}
\end{equation}
\end{theorem}

\begin{proof}
Apply \Cref{thm:e8v45-resolvent} at each cutoff and then the Banach-space
comparison test.
\end{proof}

\begin{firewall}[Finite census does not prove strictness]
\label{fire:e8v45-strict}
The exact YM first-shell count proves finite multiplicity, not
\(\kappa<1\).  Strictness must come from a bad-boundary capacity, a heat
or spacetime separation factor, reflection positivity, or another
quantitative mechanism.  Enlarging the shell once cannot make the
Neumann series contract.
\end{firewall}

\begin{definition}[Boundary influence resolvent card, BIRC]
\label{def:e8v45-BIRC}
BIRC consists of:
\begin{enumerate}[label=\textup{(B\arabic*)},leftmargin=4em]
\item a finite first collar and an owner for every crossing interaction;
\item a response norm, direct source \(a\), and influence kernel \(K\);
\item audited row and column Schur sums with
\(\sqrt{r_*c_*}<1\);
\item cofinal summability with loss \((1-\|K\|)^{-1}\);
\item the V44 counterterm, V43 observable, V38 horizon, and V35
conditional-centering contacts; and
\item Ward, constraint, boundary, reflection, chiral, spin, anomaly,
interaction, and gravitational-carrier compatibility.
\end{enumerate}
\end{definition}

\begin{theorem}[BIRC replaces recursive action closure]
\label{thm:e8v45-BIRC}
Rows (B1)--(B4) give the unique summed boundary response
\eqref{eq:e8v45-solution} without a finite recursively closed block.
With (B5), it enters LORC or PSRC in the declared norm.  Row (B6) remains
necessary for the physical SM--Einstein handoff.
\end{theorem}

\begin{proof}
Use \Cref{thm:e8v45-Schur,thm:e8v45-resolvent,
thm:e8v45-response}; the remaining rows identify the earlier compilers.
\end{proof}

\begin{assessment}[Progress at V45]
\label{ass:e8v45-status}
The audit closes the nonlinear flat-sector issue in the YM collar and
exposes a general block-closure obstruction.  A rigorous
owner-plus-resolvent compiler now replaces recursive enlargement.
Finite first-shell multiplicity enters Schur sums, while a genuinely
analytic small factor must make propagation strict.
\end{assessment}

\begin{programme}[Eighth-series V46: strictness from buffered response]
\label{prog:e8v45-V46}
V46 will combine the positive Euclidean-time buffer of V44 with the
Gaussian propagation row of V38.  It will derive an explicit sufficient
inequality for \(\kappa_0\mathfrak Q_D(a)<1\), including bad-boundary
probability, overlap, and source amplitudes, and identify the factor still
missing from the physical conditional gravitational measure.
\end{programme}

\begin{firewall}[Claim boundary after eighth-series V45]
\label{fire:e8v45-boundary}
V45 records the compulsory YM/NS audit, corrects the scope of a finite
internal patch, and proves the boundary Neumann-resolvent, operator Schur,
graph-decay, and cofinal response compilers.

V45 does not prove a strict influence constant for YM or SM--Einstein, a
bad-boundary capacity, the physical counterterm and Ward contacts, the
gravitational carrier, cutoff removal, or the full continuum.
\end{firewall}

\begin{theorem}[Eighth-series V45 result]
\label{thm:e8v45-result}
If owned first-shell influences form \(K\) with
\(\sqrt{r_*c_*}<1\), their total response is
\[
 e=(I-K)^{-1}a,
 \qquad
 \|e\|\leq{1\over1-\sqrt{r_*c_*}}\|a\|.
\]
Exponential and Gaussian graph tails give the explicit sufficient
constants \eqref{eq:e8v45-SD} and \eqref{eq:e8v45-QD}.
\end{theorem}

\begin{proof}
Use \Cref{thm:e8v45-resolvent,thm:e8v45-Schur,
thm:e8v45-exponential,cor:e8v45-Gaussian}.
\end{proof}

\paragraph{Parent-version record.}
V45 was formed from
\texttt{Renewal\_SM\_Einstein\_Continuum\_Research\_Notes\_V44.tex}.
The V44 parent had (18\,205) logical source lines, (658\,721) bytes,
and SHA--256
\[
 \texttt{43B44F104D472CF44497BE6333EF7D0B0F5C9BC0FB86761426953612C3FCAE68}.
\]
The next compulsory companion audit is V50.

% =============================================================================
% END APPENDED EIGHTH-SERIES V45 MATERIAL
% =============================================================================

% =============================================================================
% BEGIN APPENDED EIGHTH-SERIES V46 MATERIAL
% =============================================================================

\section{Strict boundary influence from buffered local response}
\label{sec:v8v46}

V45 reduced recursive boundary propagation to the strict inequality
\(\|K\|<1\).  V38 and V44 provide a Gaussian finite-horizon or
positive-time response tail, while V31 supplies restricted bad-event
norms.  This note composes those factors, improves the graph count on
polynomial-growth meshes, and proves a scale window in which terminal
accuracy, spatial separation, and boundary strictness can hold
simultaneously.

\subsection{One-step propagated bad-boundary map}
\label{sec:e8v46-one-step}

Let blocks be indexed by \(\mathcal I\).  For a path from block \(j\) to
block \(i\), suppose the propagated contact factors as
\begin{equation}
 K_{ij}=M_{B_i}P_iR_{ij}(T)S_j.
 \label{eq:e8v46-factor}
\end{equation}
Here \(S_j\) creates an owned boundary source, \(R_{ij}(T)\) propagates it
for horizon \(T\), \(P_i\) is the chosen conditional smoothing transfer,
and \(M_{B_i}\) restricts to the bad boundary event at the receiving
block.

\begin{theorem}[One-step Gaussian influence bound]
\label{thm:e8v46-one-step}
Assume
\begin{align}
 \|S_j\|&\leq A_*,
 \label{eq:e8v46-source-bound}\\
 \|M_{B_i}P_i\|&\leq\sqrt{\eta_*},
 \qquad 0\leq\eta_*\leq1,
 \label{eq:e8v46-restricted-bound}\\
 \|R_{ij}(T)\|
 &\leq C_RT
 \exp\left\{-{d(i,j)^2\over4c_{\rm h}T}\right\}.
 \label{eq:e8v46-response-bound}
\end{align}
Then
\begin{equation}
 \|K_{ij}\|
 \leq
 \kappa_0(T)
 \exp\left\{-{d(i,j)^2\over4c_{\rm h}T}\right\},
 \qquad
 \kappa_0(T)=\sqrt{\eta_*}\,C_RA_*T.
 \label{eq:e8v46-kernel}
\end{equation}
\end{theorem}

\begin{proof}
Take the operator norm of \eqref{eq:e8v46-factor} and apply the three
displayed bounds in their order.
\end{proof}

\begin{remark}[Sources included in \(A_*\)]
\label{rem:e8v46-A}
The constant \(A_*\) includes the number and norms of contacts assigned
to one owner, including any propagated local Wick or stress counterterm.
A purely local counterterm that does not seed another block belongs in
the direct source \(a\) of V45 instead.  This distinction prevents the
same contact from being charged both to \(a\) and to \(K\).
\end{remark}

\begin{corollary}[Explicit bounded-degree strictness condition]
\label{cor:e8v46-bounded-degree}
On a graph of maximum degree \(D\), \Cref{thm:e8v46-one-step} and
\Cref{cor:e8v45-Gaussian} give
\begin{equation}
 \|K\|
 \leq
 \sqrt{\eta_*}\,C_RA_*T\,
 \mathfrak Q_D\left({1\over4c_{\rm h}T}\right).
 \label{eq:e8v46-D-bound}
\end{equation}
The sufficient strictness condition is
\begin{equation}
 \sqrt{\eta_*}\,C_RA_*T\,
 \mathfrak Q_D\left({1\over4c_{\rm h}T}\right)<1.
 \label{eq:e8v46-D-strict}
\end{equation}
\end{corollary}

\begin{proof}
Use \eqref{eq:e8v46-kernel} in the Gaussian criterion
\Cref{cor:e8v45-Gaussian} with
\(a=(4c_{\rm h}T)^{-1}\).
\end{proof}

\subsection{Polynomial-growth improvement}
\label{sec:e8v46-polynomial}

Shape-regular Euclidean mesh graphs have polynomial, rather than maximal
tree, growth.  Assume their spheres satisfy
\begin{equation}
 N_i(r):=\#\{j:d(i,j)=r\}
 \leq C_{\rm g}(1+r)^{q-1}
 \label{eq:e8v46-growth}
\end{equation}
for some \(q\geq1\), uniformly in \(i\) and the cutoff.

\begin{lemma}[Gaussian sum on a polynomial-growth graph]
\label{lem:e8v46-Gaussian-sum}
For \(a>0\),
\begin{equation}
 \sup_i\sum_j e^{-a d(i,j)^2}
 \leq
 C_{\rm g}C_q(1+a^{-q/2}),
 \label{eq:e8v46-Gaussian-sum}
\end{equation}
where \(C_q<\infty\) depends only on \(q\).
\end{lemma}

\begin{proof}
Use \eqref{eq:e8v46-growth} and split the radius-zero term:
\[
 \sum_j e^{-ad(i,j)^2}
 \leq C_{\rm g}\left[
 1+\sum_{r\geq1}(1+r)^{q-1}e^{-ar^2}\right].
\]
For \(r\geq1\), \((1+r)^{q-1}\leq2^{q-1}r^{q-1}\).
Comparison with the integral, after enlarging the constant over the
bounded interval where the summand is increasing, gives
\[
 \sum_{r\geq1}r^{q-1}e^{-ar^2}
 \leq C_q\left(1+\int_0^\infty
 x^{q-1}e^{-ax^2}\,dx\right).
\]
The integral equals
\(\frac12\Gamma(q/2)a^{-q/2}\).  Absorb all numerical constants into
\(C_q\).
\end{proof}

\begin{theorem}[Polynomial-growth strictness]
\label{thm:e8v46-polynomial}
Under \eqref{eq:e8v46-growth} and
\eqref{eq:e8v46-kernel}, both Schur sums obey
\begin{equation}
 r_*,c_*
 \leq
 \sqrt{\eta_*}\,C_RA_*T\,
 C_{\rm g}C_q
 \left(1+(4c_{\rm h}T)^{q/2}\right).
 \label{eq:e8v46-poly-Schur}
\end{equation}
Hence
\begin{equation}
 \boxed{\;
 \sqrt{\eta_*}\,C_RA_*T\,
 C_{\rm g}C_q
 \left(1+(4c_{\rm h}T)^{q/2}\right)<1
 \;}
 \label{eq:e8v46-poly-strict}
\end{equation}
is sufficient for the V45 boundary resolvent.
\end{theorem}

\begin{proof}
Sum \eqref{eq:e8v46-kernel} in \(j\) and apply
\Cref{lem:e8v46-Gaussian-sum} with
\(a=(4c_{\rm h}T)^{-1}\).  The same uniform sphere bound with the roles of
\(i,j\) reversed gives the column sum.  Apply
\Cref{thm:e8v45-Schur,thm:e8v45-resolvent}.
\end{proof}

\begin{remark}[Large-horizon cost]
\label{rem:e8v46-large-T}
For \(T\geq1\), the left side of
\eqref{eq:e8v46-poly-strict} is bounded by a constant times
\begin{equation}
 \sqrt{\eta_*}\,A_*T^{1+q/2}.
 \label{eq:e8v46-large-T-cost}
\end{equation}
The finite-horizon residual improves as \(T\) grows, but the response
collar becomes more expensive.  A cofinal proof must schedule both.
\end{remark}

\subsection{A simultaneous cofinal schedule}
\label{sec:e8v46-schedule}

At scale \(n\), let the response horizon be \(T_n\), the source-observable
separation be \(R_n\), and abbreviate
\begin{equation}
 E_n:=\sqrt{\eta_n}\,C_{R,n}A_n.
 \label{eq:e8v46-En}
\end{equation}
Assume \(C_{R,n}\) is uniformly bounded and the graph constants are
uniform.

\begin{theorem}[Nonempty horizon window]
\label{thm:e8v46-window}
Suppose, for constants \(\alpha,\beta,\rho>0\),
\begin{equation}
 E_n\leq C_E2^{-\beta n},
 \qquad
 T_n=2^{\theta n},
 \qquad
 R_n=2^{\rho n}.
 \label{eq:e8v46-scales}
\end{equation}
If
\begin{equation}
 0<\theta<
 \min\left\{2\rho,\;{\beta\over1+q/2}\right\},
 \label{eq:e8v46-window}
\end{equation}
then:
\begin{enumerate}[label=\textup{(\roman*)}]
\item the polynomial-growth Schur bound tends to zero, so it is strictly
below one for all sufficiently large \(n\);
\item an infrared-order-\(\alpha\) terminal residual is summable:
\begin{equation}
 \sum_nT_n^{-\alpha}<\infty;
 \label{eq:e8v46-terminal-sum}
\end{equation}
\item if the unpropagated source grows at most exponentially, its separated
Gaussian tail
\begin{equation}
 T_n\exp\left\{-{R_n^2\over4c_{\rm h}T_n}\right\}
 \label{eq:e8v46-propagation-tail}
\end{equation}
is summable.
\end{enumerate}
\end{theorem}

\begin{proof}
By \eqref{eq:e8v46-large-T-cost}, the Schur bound is at most a constant
times
\[
 2^{[-\beta+\theta(1+q/2)]n},
\]
which tends to zero by the second upper bound in
\eqref{eq:e8v46-window}.  The terminal series is geometric with ratio
\(2^{-\alpha\theta}<1\).  Finally,
\[
 {R_n^2\over T_n}=2^{(2\rho-\theta)n},
\]
which tends superexponentially inside the negative exponential because
\(\theta<2\rho\).  It dominates the stated exponential prefactor and any
exponential source growth.
\end{proof}

\begin{corollary}[Uniform resolvent margin after finitely many scales]
\label{cor:e8v46-margin}
Under \Cref{thm:e8v46-window}, there are \(n_0\) and
\(\bar\kappa<1\) such that
\begin{equation}
 \|K_n\|\leq\bar\kappa
 \qquad(n\geq n_0).
 \label{eq:e8v46-margin}
\end{equation}
If the finitely many earlier scales also have invertible \(I-K_n\), their
maximum resolvent norm and \((1-\bar\kappa)^{-1}\) give one uniform
constant for the whole sequence.
\end{corollary}

\begin{proof}
The Schur bounds tend to zero.  Choose \(n_0\) after they fall below any
fixed \(\bar\kappa\in(0,1)\).  A finite maximum handles the earlier
scales.
\end{proof}

\subsection{Where a bad-event probability enters}
\label{sec:e8v46-probability}

For an exact heat-bath projection, V31 proves that the restricted norm
squared is the worst conditional bad-event probability.  Thus one may
take
\begin{equation}
 \eta_n=p_{*,n}
 :=\operatorname*{ess\,sup}_{\rm exterior}
 \mathbb P(B_n\mid{\rm exterior}).
 \label{eq:e8v46-pstar}
\end{equation}
For a finite-time heat-bath transfer, the corresponding coefficient is
\begin{equation}
 \eta_n(t)
 =p_{*,n}+(1-p_{*,n})e^{-2t}.
 \label{eq:e8v46-eta-t}
\end{equation}

\begin{corollary}[Required conditional capacity scale]
\label{cor:e8v46-capacity}
If \(A_n\) and \(C_{R,n}\) are both bounded below by positive
constants, then the window hypothesis
\(E_n\lesssim2^{-\beta n}\) requires
\begin{equation}
 \sqrt{\eta_n}\lesssim2^{-\beta n}.
 \label{eq:e8v46-capacity}
\end{equation}
For exact heat bath this means
\(p_{*,n}\lesssim2^{-2\beta n}\).  An averaged or unconditional bad-event
probability cannot be substituted for \(p_{*,n}\).
\end{corollary}

\begin{proof}
Use \eqref{eq:e8v46-En} and the stated bounds.  Square the exact
heat-bath conclusion.  The essential-supremum qualification is part of
\eqref{eq:e8v46-pstar}.
\end{proof}

\begin{firewall}[The present gravitational missing factor]
\label{fire:e8v46-gravity}
The metric-box and loader estimates of V34--V36 bound \(A_n\) and the
local overlap.  The heat and spacetime rows bound propagation.  They do
not prove the decay \eqref{eq:e8v46-capacity} for the conditional
Einstein bad-boundary event.  A positive conditional gravitational
carrier with an essential-supremum capacity or restricted-transfer bound
is still the missing strictness factor.  Unconditional concentration
cannot fill this row.
\end{firewall}

\subsection{Buffered shell strictness card}
\label{sec:e8v46-card}

\begin{definition}[Buffered shell strictness card, BSSC]
\label{def:e8v46-BSSC}
BSSC consists of:
\begin{enumerate}[label=\textup{(T\arabic*)},leftmargin=4em]
\item the V45 owner map and first-collar census;
\item a restricted conditional bad-event norm \(\eta_n\);
\item owned source/contact amplitude \(A_n\) without double counting;
\item a Gaussian response with \(C_{R,n},c_{\rm h}\);
\item a uniform polynomial graph-growth row;
\item a horizon and separation schedule satisfying
\eqref{eq:e8v46-window};
\item cofinal summability of direct sources and local counterterms after
the V45 resolvent loss; and
\item Ward, constraint, reflection, boundary, spin, chiral, anomaly,
interaction, renormalization, and gravitational-carrier contacts.
\end{enumerate}
\end{definition}

\begin{theorem}[BSSC closes BIRC strictness conditionally]
\label{thm:e8v46-BSSC}
Rows (T1)--(T6) imply a uniform strict influence margin after finitely many
scales and make the terminal and separated-propagation remainders
summable.  With (T7), the total boundary response is cofinally summable by
BIRC.  Row (T8) remains required for the physical SM--Einstein handoff.
\end{theorem}

\begin{proof}
Strictness and the uniform margin are
\Cref{thm:e8v46-polynomial,thm:e8v46-window,
cor:e8v46-margin}.  Terminal and propagation summability are the other
parts of \Cref{thm:e8v46-window}.  Apply
\Cref{thm:e8v45-response} using (T7).
\end{proof}

\begin{assessment}[Progress at V46]
\label{ass:e8v46-status}
Boundary strictness is now an explicit inequality rather than a generic
smallness request.  On a \(q\)-dimensional polynomial-growth mesh, the
buffered influence costs \(E_nT_n^{1+q/2}\).  If
\(E_n\lesssim2^{-\beta n}\), there is a nonempty horizon window that also
makes infrared terminal and separated propagation errors summable.  The
remaining physical obstruction is isolated as decay of the worst
conditional gravitational bad-boundary norm.
\end{assessment}

\begin{programme}[Eighth-series V47: conditional gravitational capacity]
\label{prog:e8v46-V47}
V47 will attack the missing factor in
\eqref{eq:e8v46-capacity}.  It will test a one-cell conformal conditional
density against the V29 potential-slope Hardy estimate, retain the
bounded-amplitude high-frequency conformal instability, and determine
whether a restricted good-metric chart can yield an essential-supremum
capacity or whether the action needs a new positive higher-derivative
carrier.
\end{programme}

\begin{firewall}[Claim boundary after eighth-series V46]
\label{fire:e8v46-boundary}
V46 proves the one-step factorization bound, polynomial-growth Gaussian
Schur estimate, explicit strictness inequality, simultaneous cofinal
horizon window, and the required scale of a conditional bad-event
capacity.

V46 does not prove that capacity for the Einstein measure, positivity of
the required gravitational carrier, interacting heat locality, the full
counterterm and Ward contacts, cutoff removal, or the Standard
Model--Einstein continuum.
\end{firewall}

\begin{theorem}[Eighth-series V46 result]
\label{thm:e8v46-result}
If a one-step boundary response has restricted norm
\(\sqrt{\eta_n}\), source amplitude \(A_n\), and Gaussian propagation
constant \(C_{R,n}T_n\), then on a polynomial-growth graph its Schur norm
is bounded by \eqref{eq:e8v46-poly-Schur}.  If
\(E_n=\sqrt{\eta_n}C_{R,n}A_n\lesssim2^{-\beta n}\), any
\[
 0<\theta<
 \min\left\{2\rho,{\beta\over1+q/2}\right\}
\]
gives horizons \(T_n=2^{\theta n}\) and separations
\(R_n=2^{\rho n}\) with strict boundary influence and summable terminal
and propagation errors.
\end{theorem}

\begin{proof}
Use \Cref{thm:e8v46-one-step,thm:e8v46-polynomial,
thm:e8v46-window}.
\end{proof}

\paragraph{Parent-version record.}
V46 was formed from
\texttt{Renewal\_SM\_Einstein\_Continuum\_Research\_Notes\_V45.tex}.
The V45 parent had (18\,565) logical source lines, (670\,955) bytes,
and SHA--256
\[
 \texttt{F7F9460140E8E73F011879CAA09D08D2A3799D0F586EAE2FB29333D858E38AE7}.
\]
The next compulsory companion audit is V50.

% =============================================================================
% END APPENDED EIGHTH-SERIES V46 MATERIAL
% =============================================================================

% =============================================================================
% BEGIN APPENDED EIGHTH-SERIES V47 MATERIAL
% =============================================================================

\section{Conditional conformal capacity: a one-cell no-go and block repair}
\label{sec:v8v47}

V46 requires decay of a worst conditional gravitational bad-event norm.
The V29 conformal star has good amplitude slopes when neighboring values
remain in a compact set.  That qualification is decisive.  This note
proves that a fixed one-cell absolute-amplitude chart has worst
conditional bad probability one when all neighboring conformal values
move together toward degeneracy.  The V28 collar density then loses every
uniform lower bound.  A finite block split into its common mode and
mean-zero relative modes avoids this particular failure and gives uniform
common-mode tail slopes.

\subsection{The translated-neighbour family}
\label{sec:e8v47-translation}

Use the declared conformal star action of V29 with
\(\alpha,\Lambda,\kappa>0\), fixed weights \(w_j>0\), reference exponent
\(\beta\in\mathbb R\), and set \(U_\psi=0\).  Let
\begin{equation}
 \psi_j^{(M)}=-M
 \qquad(j=1,\ldots,q),
 \qquad M>0.
 \label{eq:e8v47-neighbours}
\end{equation}
Write \(W=\sum_jw_j\) and use the translated coordinate
\begin{equation}
 \phi=-M+z.
 \label{eq:e8v47-z}
\end{equation}

\begin{lemma}[Exact translated star polynomials]
\label{lem:e8v47-translated}
For \eqref{eq:e8v47-neighbours}--\eqref{eq:e8v47-z},
\begin{equation}
 D_{\psi^{(M)}}(-M+z)=Wz^2,
 \qquad
 L_{\psi^{(M)}}(-M+z)=-Wz,
 \qquad
 Q_{\psi^{(M)}}(-M+z)=W(z^2-z).
 \label{eq:e8v47-translated}
\end{equation}
Consequently the effective potential is
\begin{align}
 V_M(-M+z)
 ={}&36\alpha W^2(z^2-z)^2
 -6\kappa We^{-2M+2z}z^2
 +\Lambda e^{-4M+4z}
 +\beta M-\beta z.
 \label{eq:e8v47-VM}
\end{align}
\end{lemma}

\begin{proof}
Every difference \(\phi-\psi_j^{(M)}\) equals \(z\), and every
\(\psi_j^{(M)}-\phi\) equals \(-z\).  Substitute in the definitions
\eqref{eq:e8v29-D}--\eqref{eq:e8v29-Q}, then in
\eqref{eq:e8v29-action} and
\eqref{eq:e8v29-V}.
\end{proof}

\begin{lemma}[Uniform mass near the translated neighbour well]
\label{lem:e8v47-denominator}
There is \(C_0>0\), independent of all sufficiently large \(M\), such
that
\begin{equation}
 \int_{-M-1/2}^{-M+1/2}e^{-V_M(\phi)}\,d\phi
 \geq C_0e^{-\beta M}.
 \label{eq:e8v47-denominator}
\end{equation}
\end{lemma}

\begin{proof}
On \(|z|\leq1/2\), the quartic polynomial in
\eqref{eq:e8v47-VM} and the term \(-\beta z\) are uniformly bounded.
The cosmological term is nonnegative and uniformly bounded above for
large \(M\).  The Einstein term is nonpositive and therefore only
increases \(e^{-V_M}\).  Hence
\(V_M(-M+z)\leq\beta M+C\) on this interval for a fixed \(C\).
Integrate \(e^{-\beta M-C}\) over its unit length.
\end{proof}

\begin{lemma}[Superquartic exclusion of a fixed absolute chart]
\label{lem:e8v47-good-numerator}
Fix \(R>0\).  There are \(c_R,C_R>0\) such that, for all sufficiently
large \(M\) and every \(|\phi|\leq R\),
\begin{equation}
 V_M(\phi)\geq c_RM^4-C_RM^2.
 \label{eq:e8v47-good-lower}
\end{equation}
Therefore
\begin{equation}
 \int_{-R}^{R}e^{-V_M(\phi)}\,d\phi
 \leq2R\,e^{-c_RM^4+C_RM^2}.
 \label{eq:e8v47-good-numerator}
\end{equation}
\end{lemma}

\begin{proof}
For \(|\phi|\leq R\), set \(z=M+\phi\).  Then
\(|z|\asymp M\), and
\(|z^2-z|\geq M^2/4\) for large \(M\).  The positive curvature-square
term in \eqref{eq:e8v47-VM} is therefore at least
\(c_RM^4\).  The negative Einstein term has magnitude at most
\(C_Re^{2R}(M+R)^2\), hence at most \(C_RM^2\).  The cosmological term is
nonnegative and the reference term is bounded on the fixed interval.
This proves \eqref{eq:e8v47-good-lower}; integration gives the second
display.
\end{proof}

\begin{theorem}[Worst one-cell absolute bad probability is one]
\label{thm:e8v47-pstar-one}
Let \(\rho_M(\phi)=Z_M^{-1}e^{-V_M(\phi)}\), and let
\begin{equation}
 G_R=\{|\phi|\leq R\},
 \qquad
 B_R=G_R^c.
 \label{eq:e8v47-good-bad}
\end{equation}
Then
\begin{equation}
 \rho_M(G_R)\longrightarrow0,
 \qquad
 \rho_M(B_R)\longrightarrow1
 \qquad(M\to\infty).
 \label{eq:e8v47-probability-limit}
\end{equation}
In particular,
\begin{equation}
 \sup_\psi
 \mathbb P_\psi(B_R)=1
 \label{eq:e8v47-pstar-one}
\end{equation}
for the one-cell conditional family.
\end{theorem}

\begin{proof}
The partition function is at least the translated-well integral in
\eqref{eq:e8v47-denominator}.  Hence
\[
 \rho_M(G_R)
 \leq {2R\over C_0}
 \exp\{-c_RM^4+C_RM^2+\beta M\}\longrightarrow0
\]
by \Cref{lem:e8v47-good-numerator}.  The complement has probability
\(1-\rho_M(G_R)\).  The family
\(\psi^{(M)}\) is among the exteriors over which the supremum is taken.
\end{proof}

\begin{corollary}[The V46 exact heat-bath factor cannot decay]
\label{cor:e8v47-no-local-p}
For the bad event \(B_R\), the exact one-cell heat-bath restricted norm
has
\begin{equation}
 p_*=\operatorname*{ess\,sup}_\psi
 \mathbb P_\psi(B_R)=1
 \label{eq:e8v47-no-local-p}
\end{equation}
whenever the exterior law gives positive measure to arbitrarily large
neighbour translations or the conditional kernel is required pointwise
for all exteriors.  It cannot satisfy
\eqref{eq:e8v46-capacity}.
\end{corollary}

\begin{proof}
Use \Cref{thm:e8v47-pstar-one} and the exact heat-bath identity recorded
in V31 and \eqref{eq:e8v46-pstar}.  Essential supremum equals one if every
neighbourhood of the displayed exterior sequence has positive exterior
measure; the pointwise version is immediate.
\end{proof}

\subsection{Failure of the fixed good collar}
\label{sec:e8v47-collar}

\begin{theorem}[No uniform V28 collar density on a fixed absolute chart]
\label{thm:e8v47-collar-failure}
Fix a collar interval \(I=[-R,-R+\ell]\subset[-R,R]\), with
\(\ell>0\).  Then
\begin{equation}
 \inf_{\phi\in I}\rho_M(\phi)\longrightarrow0.
 \label{eq:e8v47-collar-failure}
\end{equation}
Thus the uniform lower-density constant \(m_*\) required in
\eqref{eq:e8v28-collar-data} is zero after the supremum over exterior
neighbours, and the V28 coefficient
\(4p_*\ell/m_*\) cannot be finite.
\end{theorem}

\begin{proof}
By \eqref{eq:e8v47-good-lower},
\[
 \sup_{\phi\in I}e^{-V_M(\phi)}
 \leq e^{-c_RM^4+C_RM^2}.
\]
Divide by the partition-function lower bound
\eqref{eq:e8v47-denominator}.  The resulting upper bound tends to zero
uniformly on \(I\), which proves the infimum statement and in fact the
stronger uniform convergence of the density to zero there.
\end{proof}

\begin{consequence}[Neighbour compactness is structurally necessary]
\label{cons:e8v47-neighbours}
The compact-neighbour hypothesis in
\Cref{thm:e8v29-tail-classification,prop:e8v29-neighbour-row} cannot be
removed by increasing the one-cell amplitude threshold.  A common
translation of the cell and all its neighbours moves the conditional well
outside every fixed absolute chart without paying the local curvature
terms.
\end{consequence}

\begin{proof}
The translated differences in
\eqref{eq:e8v47-translated} remain \(O(1)\) near \(z=0\), independently
of \(M\), while \(\phi=-M+z\) leaves every fixed chart.
\end{proof}

\subsection{Block common and relative conformal modes}
\label{sec:e8v47-block}

Let a finite block contain \(N\) conformal values and write
\begin{equation}
 \phi_i=m+\xi_i,
 \qquad
 \sum_{i=1}^N\xi_i=0.
 \label{eq:e8v47-common-relative}
\end{equation}
On a bounded relative chart \(|\xi_i|\leq R_\xi\), suppose the
common-mode effective potential has the form
\begin{equation}
 \mathcal V_\xi(m)
 =\Lambda A_\xi e^{4m}
 -\kappa B_\xi e^{2m}
 -\beta Nm+C_\xi,
 \label{eq:e8v47-common-potential}
\end{equation}
where
\begin{equation}
 0<A_0\leq A_\xi\leq A_1,
 \qquad
 0\leq B_\xi\leq B_1.
 \label{eq:e8v47-AB}
\end{equation}
For the pure cosmological sum,
\(A_\xi=\sum_i e^{4\xi_i}\geq N\) by Jensen's inequality.

\begin{theorem}[Uniform two-sided slopes for the block common mode]
\label{thm:e8v47-common-slopes}
Assume \(\Lambda>0\) and \(\beta>0\).  For every
\begin{equation}
 0<a_0<\beta N
 \label{eq:e8v47-a0}
\end{equation}
there are \(R_+,R_->0\), depending only on the constants in
\eqref{eq:e8v47-AB}, such that uniformly on the relative chart,
\begin{align}
 \mathcal V_\xi'(m)&\geq a_0
 &&(m\geq R_+),
 \label{eq:e8v47-upper-common}\\
 {d\over dr}\mathcal V_\xi(-r)&\geq a_0
 &&(r\geq R_-).
 \label{eq:e8v47-lower-common}
\end{align}
\end{theorem}

\begin{proof}
Differentiate:
\begin{equation}
 \mathcal V_\xi'(m)
 =4\Lambda A_\xi e^{4m}
 -2\kappa B_\xi e^{2m}
 -\beta N.
 \label{eq:e8v47-common-derivative}
\end{equation}
Using \(A_\xi\geq A_0\) and \(B_\xi\leq B_1\), the right side is bounded
below by
\[
 4\Lambda A_0e^{4m}
 -2\kappa B_1e^{2m}
 -\beta N,
\]
which tends to \(+\infty\); choose \(R_+\) uniformly.
For the lower outward coordinate,
\[
 {d\over dr}\mathcal V_\xi(-r)
 =\beta N+2\kappa B_\xi e^{-2r}
 -4\Lambda A_\xi e^{-4r}
 \geq\beta N-4\Lambda A_1e^{-4r}.
\]
Choose \(R_-\) so that the last expression is at least \(a_0\).
\end{proof}

\begin{corollary}[Common-mode Hardy debit]
\label{cor:e8v47-common-Hardy}
If the block common-mode collars also have a uniform conditional density
lower bound, then \Cref{cor:e8v29-slope-debit} applies at both tails with
any \(a_0<\beta N\).  The translated-neighbour counterexample is absent
because a common translation changes \(m\), which is updated inside the
block rather than frozen as exterior data.
\end{corollary}

\begin{proof}
Use \Cref{thm:e8v47-common-slopes} in the V29 slope debit.  The coordinate
identity \eqref{eq:e8v47-common-relative} gives the last assertion.
\end{proof}

\begin{firewall}[The relative sector retains the conformal instability]
\label{fire:e8v47-relative}
The common-mode theorem does not control mean-zero high-frequency
relative fields \(\xi\) of bounded amplitude.  The instability in
\Cref{prop:e8v27-conformal-inside,prop:e8v29-amplitude-vs-frequency}
lies in that sector.  A positive higher-derivative owner can dominate its
large frequency, but the reflection-positivity and continuum-theory
changes identified earlier remain unpaid.
\end{firewall}

\subsection{A quadratic relative-mode diagnostic}
\label{sec:e8v47-symbol}

Let \(L\geq0\) be a finite graph Laplacian on the mean-zero block
subspace.  The quadratic conformal symbol with a higher-derivative owner
is
\begin{equation}
 \mathcal Q_\alpha(\xi)
 =\alpha\|L\xi\|^2-\kappa\langle\xi,L\xi\rangle,
 \qquad \alpha,\kappa>0.
 \label{eq:e8v47-Qalpha}
\end{equation}

\begin{theorem}[High-frequency stabilization threshold]
\label{thm:e8v47-high-frequency}
Assume \(\alpha\lambda_0\geq\kappa\).  If \(\xi\) belongs to the
spectral subspace of \(L\) for \(\lambda\geq\lambda_0>0\), then
\begin{equation}
 \mathcal Q_\alpha(\xi)
 \geq
 \lambda_0(\alpha\lambda_0-\kappa)\|\xi\|^2.
 \label{eq:e8v47-high-frequency}
\end{equation}
In particular, the quadratic form is nonnegative on that subspace, and
strictly positive when \(\alpha\lambda_0>\kappa\).

\end{theorem}

\begin{proof}
In an orthonormal eigenbasis of \(L\), a mode with eigenvalue
\(\lambda\) contributes
\((\alpha\lambda^2-\kappa\lambda)|\xi_\lambda|^2\).
For \(\lambda\geq\lambda_0\) and
\(\alpha\lambda_0\geq\kappa\),
\[
 \alpha\lambda^2-\kappa\lambda
 =\lambda(\alpha\lambda-\kappa)
 \geq\lambda_0(\alpha\lambda_0-\kappa).
\]
Sum over modes.
\end{proof}

\begin{remark}[The low relative modes remain separate]
\label{rem:e8v47-low}
For \(0<\lambda<\kappa/\alpha\), the symbol in
\eqref{eq:e8v47-Qalpha} is negative.  The higher-derivative term cures the
ultraviolet conformal escape but does not by itself make the full
relative form positive.  Low modes require a mass/cosmological Hessian,
a finite-volume treatment, a contour choice, or another declared
mechanism.
\end{remark}

\begin{definition}[Conformal block capacity card, CBCC]
\label{def:e8v47-CBCC}
CBCC consists of:
\begin{enumerate}[label=\textup{(C\arabic*)},leftmargin=4em]
\item a finite conformal block containing its common translation and a
mean-zero relative decomposition;
\item uniform common-mode coefficients as in
\eqref{eq:e8v47-AB}, two tail slopes, and collar density;
\item a positive relative-mode form on the high-frequency sector and an
explicit treatment of the remaining low modes;
\item an owner and V45 influence kernel for crossing metric contacts;
\item a worst-conditional restricted norm for the resulting block bad
event, rather than the impossible one-cell absolute event;
\item compatibility with V46's horizon window and cofinal capacity scale;
and
\item reflection, diffeomorphism, Ward, boundary, chiral, spin,
renormalization, and gravitational-carrier contacts.
\end{enumerate}
\end{definition}

\begin{theorem}[CBCC is the viable replacement for one-cell CFSC]
\label{thm:e8v47-CBCC}
Rows (C1)--(C3) give conditional common-mode Hardy control and a declared
relative-mode coercivity split.  Rows (C4)--(C6), if proved with strict
constants, place the block bad event in BIRC and BSSC.  The one-cell
absolute chart cannot substitute for these rows because its worst
conditional probability is one by
\Cref{thm:e8v47-pstar-one}.
\end{theorem}

\begin{proof}
The common-mode result is
\Cref{thm:e8v47-common-slopes,cor:e8v47-common-Hardy}; the relative
high-frequency result is \Cref{thm:e8v47-high-frequency}.  V45 and V46
give the shell and horizon compilers under (C4)--(C6).  The no-go theorem
proves the last assertion.
\end{proof}

\begin{assessment}[Progress at V47]
\label{ass:e8v47-status}
The missing V46 capacity cannot come from a fixed one-cell
absolute-amplitude chart: its worst conditional bad probability is
exactly one, and its good-collar density collapses.  Updating the block
common mode removes the translated-neighbour obstruction and gives honest
two-sided slopes when \(\Lambda,\beta>0\).  The remaining gravitational
problem is now concentrated in the mean-zero relative sector, where a
higher-derivative owner controls high frequencies but introduces the
known reflection and low-mode questions.
\end{assessment}

\begin{programme}[Eighth-series V48: low relative conformal modes]
\label{prog:e8v47-V48}
V48 will analyze the finite block low-mode symbol.  It will add the exact
common-mode cosmological Hessian and a declared relative mass/Hessian
term, derive the necessary and sufficient spectral positivity threshold,
and test its scaling under refinement.  If the threshold cannot be
cofinal without changing the Einstein continuum, the note will record
that obstruction rather than promoting CBCC.
\end{programme}

\begin{firewall}[Claim boundary after eighth-series V47]
\label{fire:e8v47-boundary}
V47 proves the translated-neighbour one-cell probability and collar
no-go, uniform block common-mode slopes, and the high-frequency quadratic
stabilization threshold.

V47 does not prove the block relative-mode capacity, a
reflection-positive higher-derivative gravitational carrier, the
essential-supremum block bad-event norm, the low-mode conformal treatment,
cutoff removal, or the full continuum.
\end{firewall}

\begin{theorem}[Eighth-series V47 result]
\label{thm:e8v47-result}
For the V29 star with positive \(\alpha,\Lambda,\kappa\), conditioning all
neighbours at \(-M\) makes the central conditional law concentrate near
\(-M\).  Hence every fixed absolute good chart has worst conditional bad
probability one and no uniform collar density.  A block
common/relative split repairs this translation defect: its common-mode
potential has uniform two-sided slopes for \(\Lambda,\beta>0\), while the
relative high-frequency form is positive above
\(\lambda_0>\kappa/\alpha\).
\end{theorem}

\begin{proof}
Use \Cref{thm:e8v47-pstar-one,thm:e8v47-collar-failure,
thm:e8v47-common-slopes,thm:e8v47-high-frequency}.
\end{proof}

\paragraph{Parent-version record.}
V47 was formed from
\texttt{Renewal\_SM\_Einstein\_Continuum\_Research\_Notes\_V46.tex}.
The V46 parent had (18\,989) logical source lines, (685\,323) bytes,
and SHA--256
\[
 \texttt{F99DD3A359FE2AF2E43A3E94F018A7A66A4E2763AC995D009178EF45C4738698}.
\]
The next compulsory companion audit is V50.

% =============================================================================
% END APPENDED EIGHTH-SERIES V47 MATERIAL
% =============================================================================

% =============================================================================
% BEGIN APPENDED EIGHTH-SERIES V48 MATERIAL
% =============================================================================

\section{Low relative conformal modes and the cosmological-Hessian conflict}
\label{sec:v8v48}

V47 separated the conformal block into a common amplitude and mean-zero
relative modes.  A positive curvature-square owner stabilizes sufficiently
high graph frequencies, but leaves an intermediate negative band.  This
note computes the exact finite-block mass threshold for that band.  It
then evaluates the cosmological Hessian and proves a conflict: extending
the good common-amplitude set far enough to make its lower tail cofinally
rare drives the cosmological relative mass to zero and destroys uniform
relative coercivity.

\subsection{Exact spectral threshold}
\label{sec:e8v48-threshold}

Let \(L\geq0\) be a graph Laplacian on a finite connected block and
restrict it to the mean-zero subspace.  Its eigenvalues are
\begin{equation}
 0<\lambda_1\leq\lambda_2\leq\cdots\leq\lambda_M.
 \label{eq:e8v48-spectrum}
\end{equation}
For \(\alpha,\kappa>0\) and a relative mass \(\mu\geq0\), define
\begin{equation}
 \mathcal Q_{\alpha,\mu}(\xi)
 =\alpha\|L\xi\|^2
 -\kappa\langle\xi,L\xi\rangle
 +\mu\|\xi\|^2.
 \label{eq:e8v48-Q}
\end{equation}

\begin{theorem}[Necessary and sufficient finite-block positivity]
\label{thm:e8v48-exact-threshold}
Put
\begin{equation}
 \mu_{\rm crit}(L;\alpha,\kappa)
 :=
 \max_{1\leq j\leq M}
 \bigl(\kappa\lambda_j-\alpha\lambda_j^2\bigr).
 \label{eq:e8v48-mucrit}
\end{equation}
Then \(\mathcal Q_{\alpha,\mu}\geq0\) on the mean-zero subspace if and
only if
\begin{equation}
 \mu\geq\mu_{\rm crit}(L;\alpha,\kappa).
 \label{eq:e8v48-iff}
\end{equation}
It has gap
\begin{equation}
 \delta_{\alpha,\mu}
 =\min_j\bigl(\alpha\lambda_j^2-\kappa\lambda_j+\mu\bigr)
 =\mu-\mu_{\rm crit}.
 \label{eq:e8v48-gap}
\end{equation}
\end{theorem}

\begin{proof}
Expand \(\xi=\sum_j\xi_je_j\) in an orthonormal eigenbasis of \(L\).
Then
\[
 \mathcal Q_{\alpha,\mu}(\xi)
 =\sum_j
 \bigl(\alpha\lambda_j^2-\kappa\lambda_j+\mu\bigr)|\xi_j|^2.
\]
Nonnegativity is equivalent to nonnegativity of every coefficient,
which is precisely \eqref{eq:e8v48-iff}.  Their minimum gives
\eqref{eq:e8v48-gap}.
\end{proof}

\begin{corollary}[Completed-square formula]
\label{cor:e8v48-square}
The critical mass has the exact representation
\begin{equation}
 \mu_{\rm crit}
 ={\,\kappa^2\over4\alpha}
 -\alpha\,
 \operatorname{dist}\left(
 {\kappa\over2\alpha},
 \{\lambda_1,\ldots,\lambda_M\}
 \right)^2.
 \label{eq:e8v48-square}
\end{equation}
In particular,
\begin{equation}
 \mu_{\rm crit}\leq{\kappa^2\over4\alpha}.
 \label{eq:e8v48-universal-sufficient}
\end{equation}
\end{corollary}

\begin{proof}
Complete the square:
\[
 \kappa\lambda-\alpha\lambda^2
 ={\kappa^2\over4\alpha}
 -\alpha\left(\lambda-{\kappa\over2\alpha}\right)^2.
\]
Maximize over the finite spectrum.
\end{proof}

\begin{corollary}[Mass-free alternative]
\label{cor:e8v48-mass-free}
The form with \(\mu=0\) is nonnegative if and only if
\begin{equation}
 \alpha\lambda_1\geq\kappa.
 \label{eq:e8v48-alpha-condition}
\end{equation}
It is strictly positive if the inequality is strict.
\end{corollary}

\begin{proof}
The scalar factor is
\(\lambda_j(\alpha\lambda_j-\kappa)\).  It is nonnegative for every
\(j\) exactly when it is nonnegative at \(\lambda_1\).
\end{proof}

\subsection{Cosmological Hessian in relative directions}
\label{sec:e8v48-cosmological}

Use the block coordinates
\(\phi_i=m+\xi_i\), \(\sum_i\xi_i=0\), from V47.  The positive
cosmological action is
\begin{equation}
 S_\Lambda(m,\xi)
 =\Lambda e^{4m}\sum_{i=1}^Ne^{4\xi_i},
 \qquad \Lambda>0.
 \label{eq:e8v48-SLambda}
\end{equation}

\begin{lemma}[Exact cosmological relative Hessian]
\label{lem:e8v48-cosmological-Hessian}
At \(\xi=0\), for every mean-zero direction \(h\),
\begin{equation}
 {d^2\over dt^2}S_\Lambda(m,th)\bigg|_{t=0}
 =16\Lambda e^{4m}\|h\|^2.
 \label{eq:e8v48-Hessian}
\end{equation}
Thus the quadratic Taylor coefficient in
\eqref{eq:e8v48-Q} is
\begin{equation}
 \mu_\Lambda(m)=8\Lambda e^{4m}.
 \label{eq:e8v48-muLambda}
\end{equation}
\end{lemma}

\begin{proof}
Differentiate each \(e^{4th_i}\) twice.  The first derivative is
\(4\Lambda e^{4m}\sum_ih_i=0\) by the mean-zero condition, and the second
is \(16\Lambda e^{4m}\sum_ih_i^2\).  The Taylor quadratic form is one half
of the Hessian.
\end{proof}

\begin{corollary}[Exact common-amplitude coercivity threshold]
\label{cor:e8v48-m-threshold}
The quadratic relative form is nonnegative at common amplitude \(m\) if
and only if
\begin{equation}
 m\geq m_{\rm crit}
 :={1\over4}\log
 \left({\mu_{\rm crit}\over8\Lambda}\right)
 \label{eq:e8v48-mcrit}
\end{equation}
when \(\mu_{\rm crit}>0\).  If \(\mu_{\rm crit}=0\), no lower bound on
\(m\) is needed at quadratic order.
\end{corollary}

\begin{proof}
Insert \eqref{eq:e8v48-muLambda} into
\eqref{eq:e8v48-iff} and solve the monotone exponential inequality.
\end{proof}

\subsection{Rarity versus coercivity}
\label{sec:e8v48-conflict}

To use a lower common-amplitude tail as a bad event, let
\begin{equation}
 G_R^{\rm com}=\{m\geq-R\},
 \qquad
 B_R^{\rm com}=\{m<-R\}.
 \label{eq:e8v48-common-event}
\end{equation}
The V47 common-mode slope gives exponentially improving tail estimates as
\(R\to\infty\), provided its collar data hold.

\begin{theorem}[Cosmological-mass conflict]
\label{thm:e8v48-conflict}
Assume \(\mu_{\rm crit}>0\).  Uniform nonnegativity of the quadratic
relative form on all of \(G_R^{\rm com}\) requires
\begin{equation}
 8\Lambda e^{-4R}\geq\mu_{\rm crit},
 \qquad\text{equivalently}\qquad
 R\leq R_{\max}
 :={1\over4}\log{8\Lambda\over\mu_{\rm crit}}.
 \label{eq:e8v48-Rmax}
\end{equation}
Consequently, no sequence \(R_n\to\infty\) can simultaneously:
\begin{enumerate}[label=\textup{(\roman*)}]
\item push the lower common-mode bad event arbitrarily far into its
exponential tail; and
\item retain uniform relative-mode nonnegativity on the complementary
good set using only the cosmological Hessian.
\end{enumerate}
\end{theorem}

\begin{proof}
The smallest cosmological mass on \(m\geq-R\) is
\(8\Lambda e^{-4R}\).  Apply the necessary and sufficient condition
\eqref{eq:e8v48-iff}.  Solving it gives
\eqref{eq:e8v48-Rmax}, which is a fixed finite upper bound when
\(\mu_{\rm crit}>0\).
\end{proof}

\begin{corollary}[No cofinal V46 capacity from this split alone]
\label{cor:e8v48-no-capacity}
Suppose the only mechanism making
\(\mathbb P(B_R^{\rm com}\mid{\rm exterior})\) small is sending
\(R\to\infty\).  If \(\mu_{\rm crit}>0\), the cosmological Hessian cannot
provide the good-set relative coercivity and the decaying conditional
capacity required in \eqref{eq:e8v46-capacity} at the same time.
\end{corollary}

\begin{proof}
Capacity decay requires an unbounded sequence of tail thresholds by the
stated hypothesis; \Cref{thm:e8v48-conflict} bounds every coercive
threshold above by \(R_{\max}\).
\end{proof}

\begin{firewall}[Probability scale is an additional input]
\label{fire:e8v48-probability}
The corollary is conditional on how the tail probability is reduced.  A
coupling-dependent slope, an increasing block size, or another action
owner might make a fixed threshold cofinally rare.  Those mechanisms
must be proved in the actual conditional measure.  V48 rules out only the
route that sends the lower threshold to infinity while relying on
\(\mu_\Lambda(m)\) for relative positivity.
\end{firewall}

\subsection{Scaling of the mass-free alternative}
\label{sec:e8v48-scaling}

\begin{theorem}[Large-volume cost of mass-free stabilization]
\label{thm:e8v48-large-volume}
Let \(L_{\rm phys}\) be the side length of a periodic physical box and let
the relative Laplacian have least nonzero eigenvalue
\begin{equation}
 \lambda_1(L_{\rm phys})
 \leq {C\over L_{\rm phys}^2}.
 \label{eq:e8v48-lambda1}
\end{equation}
If \(\mu=0\), positivity
\eqref{eq:e8v48-alpha-condition} requires
\begin{equation}
 \alpha\geq{\kappa\over\lambda_1(L_{\rm phys})}
 \geq{\kappa\over C}L_{\rm phys}^2.
 \label{eq:e8v48-alpha-growth}
\end{equation}
Thus a fixed higher-derivative coefficient cannot stabilize all
mean-zero relative modes uniformly in the infinite-volume limit.
\end{theorem}

\begin{proof}
Rearrange \eqref{eq:e8v48-alpha-condition} and use
\eqref{eq:e8v48-lambda1}.
\end{proof}

\begin{corollary}[Dense-spectrum critical mass]
\label{cor:e8v48-dense}
If a cofinal family of spectra approaches
\(\kappa/(2\alpha)\), then
\begin{equation}
 \mu_{\rm crit}\longrightarrow{\kappa^2\over4\alpha}.
 \label{eq:e8v48-dense}
\end{equation}
Hence the universal sufficient mass
\(\kappa^2/(4\alpha)\) is asymptotically necessary on such a family.
\end{corollary}

\begin{proof}
Use the distance formula \eqref{eq:e8v48-square}.
\end{proof}

\begin{firewall}[A uniform conformal mass changes the physical question]
\label{fire:e8v48-mass}
Adding an \(m\)-independent relative mass
\(\mu_0\geq\kappa^2/(4\alpha)\) makes the quadratic symbol nonnegative,
but no such term has been derived from the target Einstein action.  It is
a new relevant coupling unless removed by constraints or a proved
renormalization condition.  It cannot be inserted merely to close the
capacity estimate.
\end{firewall}

\subsection{Branches left by the quadratic test}
\label{sec:e8v48-branches}

\begin{definition}[Low conformal-mode resolution card, LCMRC]
\label{def:e8v48-LCMRC}
At least one of the following physical branches must be populated:
\begin{enumerate}[label=\textup{(L\arabic*)},leftmargin=4em]
\item a constraint or quotient theorem removing the negative conformal
relative band from the physical conditional carrier;
\item a reflection-compatible contour or signed construction with an
explicit positive reconstructed sector;
\item a derived, renormalized stabilizing Hessian satisfying
\(\mu\geq\mu_{\rm crit}\) without changing the target continuum;
\item a finite-volume treatment whose constants survive the subsequent
volume limit by a separate cancellation; or
\item an observable-level construction that never requires a positive
probability carrier in the conformal coordinate.
\end{enumerate}
Every branch must also retain the V47 common-mode tail and V45 boundary
influence contacts.
\end{definition}

\begin{theorem}[Quadratic data required by LCMRC]
\label{thm:e8v48-LCMRC}
For branches (L1)--(L4), the remaining physical relative spectrum must
obey either
\begin{equation}
 \mu\geq
 \max_{\lambda\in\sigma(L)\setminus\{0\}}
 (\kappa\lambda-\alpha\lambda^2)
 \label{eq:e8v48-LCMRC-threshold}
\end{equation}
or a stated replacement inequality implying the same good-set
coercivity.  Branch (L5) must instead supply the local response
positivity and Ward properties required by LORC.  None of these
conclusions follows from the V47 common-mode slopes.
\end{theorem}

\begin{proof}
The threshold is necessary and sufficient by
\Cref{thm:e8v48-exact-threshold}.  LORC states the requirements for the
observable branch.  Common-mode slopes contain no relative spectral
information.
\end{proof}

\begin{assessment}[Progress at V48]
\label{ass:e8v48-status}
The low relative conformal bottleneck now has an exact scalar threshold.
The cosmological Hessian meets it only above a fixed common amplitude.
Pushing the lower common bad set outward to gain capacity does the
opposite: it removes that Hessian.  A fixed curvature-square coefficient
also fails in infinite physical volume because the least relative
frequency closes.  The programme must now use a genuine constraint,
contour, renormalized stabilizer, finite-volume cancellation, or
observable-algebra branch.
\end{assessment}

\begin{programme}[Eighth-series V49: constraint elimination test]
\label{prog:e8v48-V49}
V49 will test branch (L1) in a finite-dimensional linearized ADM model.
It will distinguish the lapse constraint, longitudinal gauge directions,
the trace/conformal coordinate, and transverse-traceless modes, and will
compute whether solving the linearized Hamiltonian constraint removes the
negative scalar mode or merely transfers it into a nonlocal matter source.
\end{programme}

\begin{firewall}[Claim boundary after eighth-series V48]
\label{fire:e8v48-boundary}
V48 proves the exact finite-spectrum positivity threshold, the
cosmological relative Hessian, the rarity/coercivity conflict, and the
large-volume growth required by mass-free higher-derivative stabilization.

V48 does not prove any LCMRC branch for Einstein gravity, a
reflection-positive carrier, the conditional capacity, nonlinear
constraint control, cutoff removal, or the full continuum.
\end{firewall}

\begin{theorem}[Eighth-series V48 result]
\label{thm:e8v48-result}
For the quadratic mean-zero conformal form
\(\alpha L^2-\kappa L+\mu I\), positivity is equivalent to
\[
 \mu\geq
 \mu_{\rm crit}
 =\max_{\lambda\in\sigma(L)\setminus\{0\}}
 (\kappa\lambda-\alpha\lambda^2).
\]
The cosmological mass is \(8\Lambda e^{4m}\).  If
\(\mu_{\rm crit}>0\), uniform positivity on \(m\geq-R\) bounds \(R\)
above and is incompatible with sending \(R\to\infty\) to suppress the
lower common tail.  With \(\mu=0\), the required
\(\alpha\) grows at least like the square of physical volume length.
\end{theorem}

\begin{proof}
Use \Cref{thm:e8v48-exact-threshold,
lem:e8v48-cosmological-Hessian,thm:e8v48-conflict,
thm:e8v48-large-volume}.
\end{proof}

\paragraph{Parent-version record.}
V48 was formed from
\texttt{Renewal\_SM\_Einstein\_Continuum\_Research\_Notes\_V47.tex}.
The V47 parent had (19\,463) logical source lines, (702\,075) bytes,
and SHA--256
\[
 \texttt{91D06E56E3F876347AB4CC0C2894A2BE89636BBE2C9CACA68D1801AE9ECD5EF7}.
\]
The next compulsory companion audit is V50.

% =============================================================================
% END APPENDED EIGHTH-SERIES V48 MATERIAL
% =============================================================================

% =============================================================================
% BEGIN APPENDED EIGHTH-SERIES V49 MATERIAL
% =============================================================================

\section{Linearized ADM reduction and the transferred scalar source}
\label{sec:v8v49}

V48 left constraint elimination as the first physically conservative
resolution of the negative conformal band.  This note performs the exact
nonzero-momentum linearized calculation.  In source-free gravity the
scalar configuration modes are removed: one is fixed by the Hamiltonian
constraint and one is longitudinal gauge.  With matter, the constrained
scalar becomes an inverse-Laplacian response to energy density.  Thus the
conformal instability is not a physical scalar oscillator at this order,
but its analytic cost reappears as a nonlocal infrared source.

\subsection{Fourier decomposition of a symmetric metric tangent}
\label{sec:e8v49-decomposition}

Fix spatial dimension \(d\geq3\) and a nonzero Fourier momentum
\(k\in\mathbb R^d\).  Put
\begin{equation}
 n={k\over|k|},
 \qquad
 P=I-n\otimes n,
 \qquad
 T={P\over\sqrt{d-1}},
 \qquad
 N=n\otimes n.
 \label{eq:e8v49-PTN}
\end{equation}
The space \(\operatorname{Sym}^2(\mathbb R^d)\) has its Frobenius inner
product.  Define
\begin{align}
 \mathsf{TT}_k
 &:=\{h:h\,n=0,\ \operatorname{Tr}h=0\},
 \label{eq:e8v49-TT}\\
 \mathsf V_k
 &:=\{n\otimes v+v\otimes n:v\perp n\}.
 \label{eq:e8v49-vector}
\end{align}

\begin{lemma}[Orthogonal scalar--vector--tensor decomposition]
\label{lem:e8v49-decomposition}
Every symmetric tensor has a unique orthogonal decomposition
\begin{equation}
 h=h^{\rm TT}+h^{\rm V}+tT+bN,
 \label{eq:e8v49-decomposition}
\end{equation}
with
\(h^{\rm TT}\in\mathsf{TT}_k\),
\(h^{\rm V}\in\mathsf V_k\), and \(t,b\in\mathbb R\).
\end{lemma}

\begin{proof}
Relative to
\(\mathbb R^d=n^\perp\oplus\mathbb Rn\), a symmetric matrix consists of
a symmetric \((d-1)\)-dimensional transverse block, one transverse vector
in its mixed block, and one longitudinal scalar.  Split the transverse
block into its trace-free part and its multiple of \(P\).  The four
summands are mutually orthogonal, and their dimensions add to
\(d(d+1)/2\).
\end{proof}

\subsection{Constraint and spatial gauge}
\label{sec:e8v49-constraint}

Use the linearized scalar-curvature constraint
\begin{equation}
 \mathcal C_k(h)
 :=|k|^2\operatorname{Tr}h-k_ik_jh_{ij}.
 \label{eq:e8v49-C}
\end{equation}
The image of a spatial infinitesimal diffeomorphism is
\begin{equation}
 \mathcal G_k(\xi)
 :=i(k\otimes\xi+\xi\otimes k).
 \label{eq:e8v49-G}
\end{equation}

\begin{lemma}[The constraint selects the transverse scalar]
\label{lem:e8v49-C-scalar}
For the decomposition \eqref{eq:e8v49-decomposition},
\begin{equation}
 \mathcal C_k(h)
 =|k|^2\sqrt{d-1}\,t.
 \label{eq:e8v49-C-scalar}
\end{equation}
Moreover,
\begin{equation}
 \operatorname{Ran}\mathcal G_k
 =\mathsf V_k\oplus\operatorname{span}\{N\}.
 \label{eq:e8v49-gauge-image}
\end{equation}
\end{lemma}

\begin{proof}
The TT and vector pieces have zero trace and zero double contraction with
\(k\).  For \(T\), the trace is \(\sqrt{d-1}\) and
\(n_i n_jT_{ij}=0\).  For \(N\), both quantities equal one and cancel in
\eqref{eq:e8v49-C}.  This proves the first identity.

Write \(\xi=\xi_\perp+\xi_\parallel n\).  The transverse part generates
\(i|k|(n\otimes\xi_\perp+\xi_\perp\otimes n)\), spanning
\(\mathsf V_k\), and the longitudinal part generates
\(2i|k|\xi_\parallel N\).
\end{proof}

\begin{theorem}[Source-free scalar elimination at nonzero momentum]
\label{thm:e8v49-vacuum-reduction}
At \(k\neq0\),
\begin{equation}
 \{h:\mathcal C_k(h)=0\}/\operatorname{Ran}\mathcal G_k
 \cong\mathsf{TT}_k.
 \label{eq:e8v49-reduced}
\end{equation}
Thus no scalar metric configuration mode remains in the linearized
source-free reduced space.
\end{theorem}

\begin{proof}
The constraint sets \(t=0\) by
\eqref{eq:e8v49-C-scalar}.  The vector part and \(bN\) are exactly the
gauge image by \eqref{eq:e8v49-gauge-image}.  The TT component is
orthogonal to both and remains free.
\end{proof}

\begin{remark}[Canonical partner]
\label{rem:e8v49-canonical}
The momentum constraint and lapse-generated momentum gauge remove the
corresponding non-TT momentum variables.  Hence the usual linearized
reduced phase space consists of TT canonical pairs.  V49 uses only the
configuration calculation above; a nonlinear symplectic reduction and
its measure Jacobian remain separate.
\end{remark}

\begin{firewall}[The zero mode is not covered]
\label{fire:e8v49-zero}
At \(k=0\), both \(\mathcal C_k\) and \(\mathcal G_k\) vanish.  Global
volume, ADM boundary data, total energy, and topology determine that
sector.  It cannot be removed by the nonzero-mode calculation.
\end{firewall}

\subsection{Matter transfers the scalar into an inverse Laplacian}
\label{sec:e8v49-matter}

Let the sourced linearized Hamiltonian constraint be
\begin{equation}
 \mathcal C_k(h)=\gamma_G\rho(k),
 \label{eq:e8v49-sourced}
\end{equation}
where \(\rho\) is the matter energy source and \(\gamma_G\) contains the
declared gravitational normalization.

\begin{theorem}[Exact sourced scalar coefficient]
\label{thm:e8v49-sourced}
Every solution of \eqref{eq:e8v49-sourced} has
\begin{equation}
 t(k)
 ={\gamma_G\over\sqrt{d-1}}\,
 {\rho(k)\over|k|^2}.
 \label{eq:e8v49-t}
\end{equation}
Modulo spatial gauge, this coefficient is fixed and cannot be absorbed
into the TT component.
\end{theorem}

\begin{proof}
Insert \eqref{eq:e8v49-C-scalar} in
\eqref{eq:e8v49-sourced}.  Gauge directions have zero constraint and do
not change \(t\); TT tensors are orthogonal to \(T\).
\end{proof}

\begin{corollary}[Infrared amplification]
\label{cor:e8v49-amplification}
On a periodic box of side \(L\), if the least nonzero momentum satisfies
\(|k_{\min}|\asymp L^{-1}\), then a source with
\(|\rho(k_{\min})|\geq\rho_0>0\) produces
\begin{equation}
 |t(k_{\min})|\gtrsim L^2\rho_0.
 \label{eq:e8v49-L2}
\end{equation}
Therefore eliminating the conformal scalar does not give a
volume-uniform local metric response to a generic matter source.
\end{corollary}

\begin{proof}
Use \eqref{eq:e8v49-t} and
\(|k_{\min}|^{-2}\asymp L^2\).
\end{proof}

\begin{firewall}[Constraint solution destroys finite range]
\label{fire:e8v49-finite-range}
The local pre-constraint equation becomes the nonlocal multiplier
\(|k|^{-2}\) after solving it.  This is the concrete linearized instance
of \Cref{fire:e8v36-constraints}.  A post-constraint SM compiler must
retain the resulting long-range influence rather than apply the finite
first-patch theorem as if the source were still local.
\end{firewall}

\subsection{The energy topology and a divergence cancellation}
\label{sec:e8v49-energy}

The potential itself is not the correct local norm for the constraint
response.  Its gradient has better spatial decay.  In
\(\mathbb R^d\), \(d>2\), let
\begin{equation}
 G_d(x)=c_d|x|^{2-d}
 \label{eq:e8v49-Green}
\end{equation}
be the Laplace fundamental solution away from the origin.

\begin{lemma}[Exterior gradient energy of a point source]
\label{lem:e8v49-gradient-tail}
For \(R>0\),
\begin{equation}
 \int_{|x|\geq R}|\nabla G_d(x)|^2\,dx
 =C_dR^{2-d}<\infty.
 \label{eq:e8v49-gradient-tail}
\end{equation}
By contrast,
\begin{equation}
 \int_{|x|\geq R}|G_d(x)|^2\,dx=\infty
 \label{eq:e8v49-potential-tail}
\end{equation}
for \(d=3,4\).
\end{lemma}

\begin{proof}
Since
\(|\nabla G_d(x)|=|c_d|(d-2)|x|^{1-d}\), polar coordinates give
\[
 C\int_R^\infty r^{2-2d}r^{d-1}\,dr
 =C\int_R^\infty r^{1-d}\,dr
 =C_dR^{2-d}.
\]
For the potential the radial integral is
\(\int_R^\infty r^{4-2d}r^{d-1}dr
=\int_R^\infty r^{3-d}dr\), which diverges for \(d\leq4\).
\end{proof}

\begin{theorem}[Divergence sources have a uniform constrained energy response]
\label{thm:e8v49-divergence}
On a periodic box, omit the zero mode.  Suppose
\begin{equation}
 \rho(k)=ik\cdot J(k)
 \label{eq:e8v49-divergence-source}
\end{equation}
for a vector source \(J\).  Let
\begin{equation}
 E(k)=ik\,|k|^{-2}\rho(k)
 \label{eq:e8v49-E}
\end{equation}
be the gradient of the constrained scalar potential.  Then
\begin{equation}
 |E(k)|\leq|J(k)|,
 \qquad
 \|E\|_{\ell^2_k}\leq\|J\|_{\ell^2_k},
 \label{eq:e8v49-Riesz}
\end{equation}
with constant one, independently of the box side and lattice refinement.
\end{theorem}

\begin{proof}
Substitute \eqref{eq:e8v49-divergence-source}:
\[
 E(k)=-{k\otimes k\over|k|^2}J(k).
\]
The matrix \(k\otimes k/|k|^2\) is the orthogonal projection onto
\(\mathbb Ck\) and has norm one.  Square and sum over nonzero modes.
\end{proof}

\begin{corollary}[The exact missing constraint-source row]
\label{cor:e8v49-missing-row}
The gradient-energy loader can absorb the post-constraint scalar response
uniformly if the relevant matter source has the divergence factorization
\eqref{eq:e8v49-divergence-source} with a cofinally bounded \(J\).
Without that factorization, the multiplier has norm
\(|k_{\min}|^{-1}\) from the source \(L^2\) norm to the gradient-energy
norm and diverges with physical volume.
\end{corollary}

\begin{proof}
The positive statement is \Cref{thm:e8v49-divergence}.  For a source
supported on \(k_{\min}\), \eqref{eq:e8v49-E} has magnitude
\(|\rho(k_{\min})|/|k_{\min}|\), which attains the stated operator norm.
\end{proof}

\begin{remark}[Conservation is not automatically the needed factorization]
\label{rem:e8v49-conservation}
Stress-energy conservation relates time and spatial derivatives of
several components.  It does not by itself write the energy density on a
fixed time slice as the spatial divergence of an \(L^2\) vector with
uniform norm.  Total energy and boundary flux determine the zero and
near-zero modes.  The row in \eqref{eq:e8v49-divergence-source} must be
proved for the particular fluctuation or replaced by a charge-sector
decomposition.
\end{remark}

\subsection{Constraint-reduced carrier card}
\label{sec:e8v49-card}

\begin{definition}[Linearized constraint transfer card, LCTC]
\label{def:e8v49-LCTC}
LCTC consists of:
\begin{enumerate}[label=\textup{(A\arabic*)},leftmargin=4em]
\item a regulated scalar--vector--TT decomposition and all constraint
Jacobians;
\item source-free removal of nonzero-momentum scalar and vector gauge
directions;
\item an explicit treatment of the zero mode, total charges, and boundary
data;
\item for matter sources, a divergence/charge-sector factorization giving
a uniform energy-topology response;
\item V45 ownership and long-range influence bounds after constraint
solution;
\item nonlinear stability of the constraint chart and compatibility with
V43--V44 local observables; and
\item reflection, Ward, diffeomorphism, chiral, spin, anomaly,
renormalization, and positive-carrier contacts.
\end{enumerate}
\end{definition}

\begin{theorem}[What LCTC resolves]
\label{thm:e8v49-LCTC}
Rows (A1)--(A2) remove the V48 negative conformal band from the
source-free linearized physical configuration space at nonzero momentum.
Rows (A3)--(A5) are necessary to prevent its sourced replacement
\eqref{eq:e8v49-t} from violating the cofinal response bounds.  If the
divergence row holds, the constrained gradient response has the
cutoff-independent estimate \eqref{eq:e8v49-Riesz}.
\end{theorem}

\begin{proof}
Use \Cref{thm:e8v49-vacuum-reduction,thm:e8v49-sourced,
thm:e8v49-divergence}.  The zero-mode and boundary qualifications are
\Cref{fire:e8v49-zero,fire:e8v49-finite-range}.
\end{proof}

\begin{assessment}[Progress at V49]
\label{ass:e8v49-status}
At linear order and nonzero momentum, pure Einstein gravity has no
physical conformal oscillator: the transverse scalar is constrained and
the longitudinal scalar is gauge.  Matter does not leave this result
cost-free.  It produces an inverse-Laplacian scalar response with
volume-divergent low modes.  In the gradient-energy topology, a precise
divergence factorization converts that response into an orthogonal Riesz
projection with norm one.  This factorization, plus zero-mode and
nonlinear constraint control, is the next physical interface.
\end{assessment}

\begin{programme}[Eighth-series V50: companion audit and constraint source]
\label{prog:e8v49-V50}
V50 will perform the compulsory YM/NS audit.  It will compare any new
boundary-shell or owner construction with the post-constraint long-range
tail above, then state whether the next attack should target the matter
energy source factorization, the nonlinear constraint Jacobian, or the
zero/charge sector.
\end{programme}

\begin{firewall}[Claim boundary after eighth-series V49]
\label{fire:e8v49-boundary}
V49 proves the nonzero-mode linearized scalar--vector--TT decomposition,
source-free scalar elimination, exact sourced inverse-Laplacian
coefficient, exterior energy-tail comparison, and the norm-one divergence
source estimate.

V49 does not prove nonlinear ADM reduction, the matter divergence row,
zero-mode or ADM boundary control, a positive reflection-compatible
carrier, the conditional capacity, cutoff removal, or the full continuum.
\end{firewall}

\begin{theorem}[Eighth-series V49 result]
\label{thm:e8v49-result}
For every nonzero momentum, the source-free linearized ADM configuration
space modulo spatial gauge is the TT tensor space.  With matter,
\[
 t(k)={\gamma_G\over\sqrt{d-1}}{\rho(k)\over|k|^2}.
\]
Its gradient response is uniform in volume if
\(\rho(k)=ik\cdot J(k)\), because it is then the orthogonal Fourier
projection
\(-k\otimes k\,J(k)/|k|^2\).
\end{theorem}

\begin{proof}
Use \Cref{thm:e8v49-vacuum-reduction,thm:e8v49-sourced,
thm:e8v49-divergence}.
\end{proof}

\paragraph{Parent-version record.}
V49 was formed from
\texttt{Renewal\_SM\_Einstein\_Continuum\_Research\_Notes\_V48.tex}.
The V48 parent had (19\,870) logical source lines, (716\,025) bytes,
and SHA--256
\[
 \texttt{6D88E16BAD2FD32A54629D301666347D3153E4F0CE00616EBCBB7F10F0024E24}.
\]
The next compulsory companion audit is V50.

% =============================================================================
% END APPENDED EIGHTH-SERIES V49 MATERIAL
% =============================================================================

% =============================================================================
% BEGIN APPENDED EIGHTH-SERIES V50 MATERIAL
% =============================================================================

\clearpage
\section{Eighth-series V50: companion audit and partial-polar constraint cancellation}
\label{sec:v8v50}

\subsection{Purpose and audit record}
\label{sec:e8v50-audit}

V49 isolated a precise infrared interface.  Solving the scalar Einstein
constraint costs an inverse Laplacian, whereas a source which is already a
spatial divergence has a norm-one response in the gradient-energy topology.
The compulsory V50 audit asks whether the parallel Yang--Mills and
Navier--Stokes programmes supply a non-circular way to exploit that
interface.

The newest companion files inspected at this audit are
\begin{align}
 &\texttt{Renewal\_Yang\_Mills\_Mass\_Gap\_Research\_Notes\_V8.tex},
 \notag\\
 &\hspace{2em}3136\ \text{logical lines},\quad116402\ \text{bytes},
 \quad\operatorname{SHA256}
 =\texttt{576D70B7E5E70D0424A72E0BE7DF7142BEAD73354A015AE9AE2A565EFF57D6B4},
 \label{eq:e8v50-YM-audit}\\
 &\texttt{Renewal\_NS\_Stability\_Research\_Notes\_V26.tex},
 \notag\\
 &\hspace{2em}5319\ \text{logical lines},\quad278283\ \text{bytes},
 \quad\operatorname{SHA256}
 =\texttt{82C887F8C2C69E4F28FAD7C3DC8DE5B9099CB5A4BF431EBA1FFF3E91935978AD}.
 \label{eq:e8v50-NS-audit}
\end{align}
The NS file is unchanged since the V45 audit.  The YM file has advanced
from the seventh-series V1 to V8.

\begin{audit}[New Yang--Mills information at V50]
\label{audit:e8v50-YM}
The new YM chain proves six facts relevant to the present programme.
First, a translation-covariant plaquette owner and normalized shell debit
give a strict local factor only after the debit is fixed independently of
the source.  Second, source and owner conditional projections need not
commute, so their leakage has to be compiled by an orthogonal martingale
decomposition.  Third, the Gaussian infrared test requires cancellation
in the source symbol; locality alone does not control the inverse response.
Fourth, a curvature-adjoint source cancels the square root of the Hodge
operator.  Fifth, the cancellation persists nonlinearly through the polar
factor of an injective collar curvature derivative.  Sixth, writing a
source in the adjoint range by first applying the inverse is explicitly
circular: the adjoint primitive must instead be supplied by a Ward,
curvature, conservation, or martingale identity with a uniform norm.
\end{audit}

\begin{audit}[Effect on the Einstein direction]
\label{audit:e8v50-direction}
The V49 Fourier projection is the scalar-constraint analogue of the YM
Hodge cancellation, but the Einstein constraint derivative has a kernel
coming from constant modes, gauge directions, and boundary charges.  Thus
the useful new task is not to repeat the injective YM theorem.  It is to
prove its partial-isometry version on the physical quotient, identify the
surviving charge sector, and then construct the matter adjoint primitive
directly from stress--energy balance.  The nonlinear constraint Jacobian
is postponed until that source row is non-circular.
\end{audit}

\subsection{Partial polar cancellation with a kernel}
\label{sec:e8v50-polar}

Let \(X,Y\) be finite-dimensional Hilbert spaces at a fixed regulator and
let \(B:X\to Y\).  Put
\begin{equation}
 K:=\ker B,
 \qquad X_0:=K^\perp,
 \qquad H_0:=\left.B^*B\right|_{X_0},
 \qquad Y_0:=\operatorname{Ran}B.
 \label{eq:e8v50-spaces}
\end{equation}
The restriction \(H_0\) is positive definite.  Its least eigenvalue will
be denoted by \(\lambda_0(B)>0\).

\begin{theorem}[Partial-polar square-root cancellation]
\label{thm:e8v50-partial-polar}
Define
\begin{equation}
 V:=BH_0^{-1/2}:X_0\longrightarrow Y.
 \label{eq:e8v50-V}
\end{equation}
Then
\begin{equation}
 V^*V=I_{X_0},
 \qquad
 VV^*=P_{Y_0},
 \qquad
 H_0^{-1/2}B^*=V^*.
 \label{eq:e8v50-polar-identities}
\end{equation}
Consequently, for \(Z\in Y\) and \(R\in X_0\), the source
\begin{equation}
 S=B^*Z+R
 \label{eq:e8v50-source-split}
\end{equation}
satisfies
\begin{equation}
 \left\|H_0^{-1/2}S\right\|_X
 \leq
 \left\|P_{Y_0}Z\right\|_Y
 +\lambda_0(B)^{-1/2}\|R\|_X.
 \label{eq:e8v50-source-bound}
\end{equation}
The coefficient of the exact adjoint part is one and has no dependence on
the least nonzero singular value of \(B\).
\end{theorem}

\begin{proof}
For \(x\in X_0\),
\[
 \|Vx\|_Y^2
 =\langle H_0^{-1/2}x,B^*BH_0^{-1/2}x\rangle_X
 =\|x\|_X^2.
\]
Thus \(V^*V=I_{X_0}\), \(V\) is an isometry, and its range equals
\(Y_0\).  Hence \(VV^*\) is the orthogonal projection onto \(Y_0\).
Taking the adjoint of \(V=BH_0^{-1/2}\) gives the third identity in
\eqref{eq:e8v50-polar-identities}.  Therefore
\[
 H_0^{-1/2}S=V^*Z+H_0^{-1/2}R.
\]
Now \(V^*Z=V^*P_{Y_0}Z\), and \(V^*\) is an isometry from \(Y_0\)
onto \(X_0\).  The spectral theorem gives
\(\|H_0^{-1/2}R\|\leq\lambda_0(B)^{-1/2}\|R\|\).
The triangle inequality proves \eqref{eq:e8v50-source-bound}.
\end{proof}

\begin{corollary}[Pseudoinverse energy response]
\label{cor:e8v50-pseudoinverse}
Let \(H^\dagger\) be the Moore--Penrose inverse of \(B^*B\), extended by
zero on \(K\).  If \(S=B^*Z\in X_0\), then
\begin{equation}
 \left\|B H^\dagger S\right\|_Y
 =\left\|H_0^{-1/2}S\right\|_X
 =\left\|P_{Y_0}Z\right\|_Y
 \leq\|Z\|_Y.
 \label{eq:e8v50-pseudoinverse-energy}
\end{equation}
\end{corollary}

\begin{proof}
On \(X_0\), \(BH_0^{-1}=VH_0^{-1/2}\).  Since \(V\) is an
isometry,
\[
 \|BH^\dagger S\|
 =\|H_0^{-1/2}S\|
 =\|V^*Z\|
 =\|P_{Y_0}Z\|.
\]
\end{proof}

\begin{remark}[What is new relative to the audited YM result]
\label{rem:e8v50-new}
The audited YM theorem treats an injective collar derivative and obtains
an isometric polar factor.  \Cref{thm:e8v50-partial-polar} permits a
nontrivial kernel, identifies the exact quotient \(X_0\), uses the
pseudoinverse rather than a fictitious inverse, and displays the residual
term which still pays the least positive singular value.  Those additions
are essential for gravitational gauge and charge modes.
\end{remark}

\subsection{The scalar Einstein constraint as an adjoint complex}
\label{sec:e8v50-Einstein}

On a periodic spatial regulator, let \(X\) be scalar functions and \(Y\)
vector fields with their regulated \(L^2\) inner products.  Take
\begin{equation}
 B=\nabla,
 \qquad B^*=-\operatorname{div},
 \qquad B^*B=-\Delta.
 \label{eq:e8v50-grad-complex}
\end{equation}
Then \(K\) is the constant scalar sector, \(X_0\) is the mean-zero scalar
sector, and \(Y_0\) is the space of exact vector fields.

\begin{theorem}[Constraint-adjoint cancellation on the mean-zero sector]
\label{thm:e8v50-constraint-adjoint}
Suppose the mean-zero part of a regulated matter source has a directly
constructed representation
\begin{equation}
 \rho_0=-\operatorname{div}J+R,
 \qquad R\in X_0.
 \label{eq:e8v50-rho-split}
\end{equation}
Let \(u=(-\Delta)^\dagger\rho_0\).  Then
\begin{equation}
 \|\nabla u\|_2
 \leq
 \|P_{\mathrm{grad}}J\|_2
 +|k_{\min}|^{-1}\|R\|_2.
 \label{eq:e8v50-grad-bound}
\end{equation}
If \(R=0\), the constant is one, independently of volume and refinement.
The average
\begin{equation}
 Q(\rho):=|\Omega|^{-1}\int_\Omega\rho
 \label{eq:e8v50-charge}
\end{equation}
is not inverted by this construction and remains a separate charge row.
\end{theorem}

\begin{proof}
Apply \Cref{thm:e8v50-partial-polar,cor:e8v50-pseudoinverse} to
\eqref{eq:e8v50-grad-complex}.  The least eigenvalue of \(-\Delta\) on
\(X_0\) is \(|k_{\min}|^2\), which gives the residual coefficient in
\eqref{eq:e8v50-grad-bound}.  The orthogonal complement of \(X_0\) is
the constant mode, so the pseudoinverse annihilates it rather than
solving the integrated constraint.
\end{proof}

\begin{corollary}[Recovery of the V49 Fourier projection]
\label{cor:e8v50-V49}
If \(R=0\), then at each nonzero momentum
\begin{equation}
 \widehat{\nabla u}(k)
 =-{k\otimes k\over |k|^2}\widehat J(k).
 \label{eq:e8v50-Riesz}
\end{equation}
Thus \Cref{thm:e8v49-divergence} is the Fourier realization of the
partial-polar identity.
\end{corollary}

\begin{proof}
Fourier transform \(u=(-\Delta)^\dagger(-\operatorname{div}J)\) and
multiply by \(ik\).  The resulting multiplier is the orthogonal
projection onto the span of \(k\).
\end{proof}

\begin{firewall}[Formal range membership is circular]
\label{fire:e8v50-circular}
For every mean-zero \(\rho_0\) one can formally set
\(J=-\nabla(-\Delta)^\dagger\rho_0\).  This does not prove the desired
factorization: its norm is exactly the inverse response one is trying to
bound.  An admissible \(J\) must be produced before constraint inversion
from a local conservation, Ward, flux, ownership, or martingale identity,
and its norm and locality must be controlled uniformly in the regulator
and volume.
\end{firewall}

\begin{firewall}[The charge row cannot be hidden in the quotient]
\label{fire:e8v50-charge}
The equation \(-\Delta u=\rho\) on a periodic regulator requires
\(Q(\rho)=0\).  Projecting away \(Q(\rho)\) is a compatibility operation,
not a solution of the total-energy constraint.  In asymptotically flat or
bounded geometry the corresponding datum is carried by boundary flux and
ADM charge.  It must be matched by a boundary sector rather than erased.
\end{firewall}

\subsection{Constraint-adjoint source factorization card}
\label{sec:e8v50-card}

\begin{definition}[Constraint-adjoint source factorization card, CASFC]
\label{def:e8v50-CASFC}
For each regulator and physical patch, CASFC consists of:
\begin{enumerate}[label=\textup{(C\arabic*)},leftmargin=4em]
\item a constraint derivative \(B\), its gauge/charge kernel \(K\), and
the physical complement \(X_0=K^\perp\);
\item a directly constructed decomposition
\(S=B^*Z+R+S_{\mathrm{charge}}\);
\item a uniform bound on \(P_{\operatorname{Ran}B}Z\) in the local
energy norm;
\item a residual estimate strong enough to absorb
\(\lambda_0(B)^{-1/2}\|R\|\), including its volume scaling;
\item an explicit boundary or global treatment of
\(S_{\mathrm{charge}}\);
\item stability of (C1)--(C5) under nonlinear constraint charts and
owner/martingale conditioning; and
\item compatibility with the V43--V44 observable tests and the V45
long-range boundary compiler.
\end{enumerate}
\end{definition}

\begin{theorem}[What CASFC would close]
\label{thm:e8v50-CASFC}
If CASFC holds with \(R=0\), then the square-root constraint response on
the physical complement has norm at most
\(\|P_{\operatorname{Ran}B}Z\|\), with no dependence on the least
positive constraint singular value.  If \(R\neq0\), all infrared loss is
isolated in the single explicit factor
\(\lambda_0(B)^{-1/2}\|R\|\).  Neither statement controls the charge
sector without (C5).
\end{theorem}

\begin{proof}
The first two conclusions are \eqref{eq:e8v50-source-bound}.  Since
\(B^*Z+R\in X_0\), it is orthogonal to the charge/kernel sector; hence the
last sector cannot be estimated by the same pseudoinverse identity.
\end{proof}

\begin{assessment}[Progress at V50]
\label{ass:e8v50-status}
The companion audit selects matter energy factorization as the immediate
upstream target.  The YM polar mechanism extends exactly to the
non-injective Einstein constraint after the kernel is separated.  It
turns the V49 Fourier observation into a regulator-independent Hilbert
complex theorem and exposes only two remaining losses: a non-adjoint
residual pays the spectral gap, while the total charge requires a global
or boundary equation.  No inverse-defined primitive is accepted.
\end{assessment}

\begin{programme}[Eighth-series V51: direct conservation primitive]
\label{prog:e8v50-V51}
V51 will use the local stress--energy balance law on a time slab to build
an adjoint primitive for time-window differences of the energy density.
It will separate the exact spatial-divergence term, endpoint charge, and
metric-work residual, with norm estimates before any constraint inverse
is applied.
\end{programme}

\begin{firewall}[Claim boundary after eighth-series V50]
\label{fire:e8v50-boundary}
V50 proves the partial-polar theorem, the pseudoinverse energy identity,
and the scalar constraint-adjoint estimate with explicit residual and
charge sectors.  It does not prove that the Standard Model energy source
has the required direct factorization, that the residual is cofinally
small, that the nonlinear ADM constraint chart is stable, or that the
global charge is matched.  It does not prove a positive carrier, cutoff
removal, or the full SM--Einstein continuum.
\end{firewall}

\begin{theorem}[Eighth-series V50 result]
\label{thm:e8v50-result}
On the physical complement of the kernel of a regulated constraint
operator \(B\), an exact adjoint source \(B^*Z\) obeys
\[
 \|(B^*B)^{-1/2}B^*Z\|
 =\|P_{\operatorname{Ran}B}Z\|
 \leq\|Z\|.
\]
For the scalar Einstein constraint this is the norm-one gradient response
to a directly constructed divergence source.  Constants and total
charges lie outside this inversion and remain separate data.
\end{theorem}

\begin{proof}
Use \Cref{thm:e8v50-partial-polar,thm:e8v50-constraint-adjoint}.
\end{proof}

\paragraph{Parent-version record.}
V50 was formed from
\texttt{Renewal\_SM\_Einstein\_Continuum\_Research\_Notes\_V49.tex}.
The V49 parent had (20\,274) logical source lines, (730\,129) bytes,
and SHA--256
\[
 \texttt{345B10AA597C15891384565DE31DF0ABCC3E942D856F8C20B7F4A5CCFD91CCFC}.
\]
The next compulsory companion audit is V55.

% =============================================================================
% END APPENDED EIGHTH-SERIES V50 MATERIAL
% =============================================================================
% =============================================================================
% BEGIN APPENDED EIGHTH-SERIES V51 MATERIAL
% =============================================================================

\clearpage
\section{Eighth-series V51: stress--energy balance as a direct constraint primitive}
\label{sec:v8v51}

\subsection{Why a time slab is the correct upstream object}
\label{sec:e8v51-purpose}

V50 forbids defining an adjoint primitive by applying the inverse
constraint operator to the source.  A local balance law supplies a
different construction: the physical spatial flux exists before the
Einstein constraint is solved.  This section proves exactly which
energy-density combinations inherit that flux factorization.

Let \(\Omega\) be a regulated periodic spatial box or a bounded spatial
cell, let \(I=[a,b]\), and suppose
\begin{equation}
 \partial_t q+\operatorname{div}J=W
 \label{eq:e8v51-balance}
\end{equation}
in distributions on \(I\times\Omega\).  Here \(q\) is a densitized
energy source, \(J\) is its physical spatial flux, and \(W\) is the work
term.  It is important that \(J\) and \(W\) are defined by the local
matter identity, before any elliptic constraint inverse is introduced.

\begin{theorem}[Endpoint transfer factorization]
\label{thm:e8v51-endpoint}
Assume
\[
 q\in C(I;L^2(\Omega)),\qquad
 J,W\in L^1(I;L^2(\Omega)).
\]
Then in distributions on \(\Omega\),
\begin{equation}
 q(b)-q(a)
 =-\operatorname{div}Z_{a,b}+R_{a,b},
 \qquad
 Z_{a,b}:=\int_a^bJ(t)\,dt,
 \qquad
 R_{a,b}:=\int_a^bW(t)\,dt.
 \label{eq:e8v51-endpoint-split}
\end{equation}
Moreover, with \(T=b-a\),
\begin{align}
 \|Z_{a,b}\|_{L^2_x}
 &\leq T^{1/2}\|J\|_{L^2_{t,x}},
 \label{eq:e8v51-Z-endpoint}\\
 \|R_{a,b}\|_{L^2_x}
 &\leq T^{1/2}\|W\|_{L^2_{t,x}}
 \label{eq:e8v51-R-endpoint}
\end{align}
whenever the right-hand sides are finite.
\end{theorem}

\begin{proof}
Integrate \eqref{eq:e8v51-balance} from \(a\) to \(b\):
\[
 q(b)-q(a)
 =-\operatorname{div}\int_a^bJ(t)\,dt+\int_a^bW(t)\,dt.
\]
This identity is valid after pairing with any spatial test function, so
it holds distributionally under the stated hypotheses.  Pointwise in
\(x\), Cauchy--Schwarz gives
\[
 \left|\int_a^bJ(t,x)\,dt\right|^2
 \leq T\int_a^b|J(t,x)|^2\,dt.
\]
Integration in \(x\) proves \eqref{eq:e8v51-Z-endpoint}; the same
argument proves \eqref{eq:e8v51-R-endpoint}.
\end{proof}

\begin{theorem}[Compact time-window factorization]
\label{thm:e8v51-window}
Let \(\chi\in H^1_0(a,b)\), and define
\begin{equation}
 S_\chi:=\int_a^b\chi'(t)q(t)\,dt.
 \label{eq:e8v51-Schi}
\end{equation}
Then
\begin{equation}
 S_\chi=-\operatorname{div}Z_\chi+R_\chi,
 \qquad
 Z_\chi:=-\int_a^b\chi(t)J(t)\,dt,
 \qquad
 R_\chi:=-\int_a^b\chi(t)W(t)\,dt.
 \label{eq:e8v51-window-split}
\end{equation}
If \(J,W\in L^2(I\times\Omega)\), then
\begin{align}
 \|Z_\chi\|_{L^2_x}
 &\leq\|\chi\|_{L^2_t}\|J\|_{L^2_{t,x}},
 \label{eq:e8v51-Z-window}\\
 \|R_\chi\|_{L^2_x}
 &\leq\|\chi\|_{L^2_t}\|W\|_{L^2_{t,x}}.
 \label{eq:e8v51-R-window}
\end{align}
\end{theorem}

\begin{proof}
The vanishing trace of \(\chi\) and integration by parts give
\[
 S_\chi
 =-\int_a^b\chi\,\partial_tq\,dt
 =\operatorname{div}\int_a^b\chi J\,dt-\int_a^b\chi W\,dt,
\]
which is \eqref{eq:e8v51-window-split}.  For almost every \(x\),
\[
 \left|\int_a^b\chi(t)J(t,x)\,dt\right|^2
 \leq\|\chi\|_{L^2_t}^2\int_a^b|J(t,x)|^2\,dt.
\]
Integrate over \(x\).  The proof for \(R_\chi\) is identical.
\end{proof}

\subsection{Constraint-energy estimates obtained before inversion}
\label{sec:e8v51-response}

On the periodic box, write \(P_0\) for mean-zero projection and
\(\lambda_1=|k_{\min}|^2\) for the least positive eigenvalue of
\(-\Delta\).

\begin{corollary}[Endpoint constrained response]
\label{cor:e8v51-endpoint-response}
Let
\[
 u_{a,b}:=(-\Delta)^\dagger P_0\bigl(q(b)-q(a)\bigr).
\]
Under the hypotheses of \Cref{thm:e8v51-endpoint},
\begin{equation}
 \|\nabla u_{a,b}\|_2
 \leq
 T^{1/2}\|J\|_{L^2_{t,x}}
 +\lambda_1^{-1/2}T^{1/2}\|P_0W\|_{L^2_{t,x}}.
 \label{eq:e8v51-endpoint-response}
\end{equation}
In the exactly conserved case \(W=0\), the bound has no inverse spectral
factor.
\end{corollary}

\begin{proof}
Project \eqref{eq:e8v51-endpoint-split} onto mean-zero functions and
apply \Cref{thm:e8v50-constraint-adjoint}, followed by
\eqref{eq:e8v51-Z-endpoint}--\eqref{eq:e8v51-R-endpoint}.
\end{proof}

\begin{corollary}[Time-window constrained response]
\label{cor:e8v51-window-response}
Let \(u_\chi=(-\Delta)^\dagger P_0S_\chi\).  Then
\begin{equation}
 \|\nabla u_\chi\|_2
 \leq
 \|\chi\|_{L^2_t}
 \left(
 \|J\|_{L^2_{t,x}}
 +\lambda_1^{-1/2}\|P_0W\|_{L^2_{t,x}}
 \right).
 \label{eq:e8v51-window-response}
\end{equation}
\end{corollary}

\begin{proof}
Use \Cref{thm:e8v51-window,thm:e8v50-constraint-adjoint}.
\end{proof}

\begin{remark}[Why the exact term is cofinally preferable]
\label{rem:e8v51-cofinal}
The flux terms in \eqref{eq:e8v51-endpoint-response} and
\eqref{eq:e8v51-window-response} are controlled in spacetime \(L^2\)
with constant one up to the elementary time-window factor.  They do not
pay \(|k_{\min}|^{-1}\).  Only the geometric or external work residual
does.  Thus the next cofinal estimate should be placed on \(W\), rather
than on the full density difference.
\end{remark}

\subsection{Charge and boundary flux}
\label{sec:e8v51-charge}

For a bounded \(\Omega\), define
\begin{equation}
 Q(t):=\int_\Omega q(t,x)\,dx.
 \label{eq:e8v51-Q}
\end{equation}

\begin{theorem}[Integrated charge balance]
\label{thm:e8v51-charge}
If the normal trace is defined, then
\begin{equation}
 Q(b)-Q(a)
 =-\int_a^b\int_{\partial\Omega}J\cdot n\,d\sigma\,dt
 +\int_a^b\int_\Omega W\,dx\,dt.
 \label{eq:e8v51-charge-balance}
\end{equation}
On a periodic box with \(W=0\), \(Q(t)\) is constant.  Consequently,
\[
 \int_\Omega \bigl(q(b)-q(a)\bigr)\,dx=0,
 \qquad
 \int_\Omega S_\chi\,dx=0.
\]
\end{theorem}

\begin{proof}
Integrate \eqref{eq:e8v51-balance} over spacetime and use the divergence
theorem.  In the periodic case the boundary flux cancels.  For
\(S_\chi\), either integrate its definition and use constancy of \(Q\),
or integrate \eqref{eq:e8v51-window-split}.
\end{proof}

\begin{firewall}[Conservation does not factor a fixed-time density]
\label{fire:e8v51-fixed-time}
Equation \eqref{eq:e8v51-balance} factors endpoint differences and
zero-boundary time smearings.  It does not imply that \(q(t_0)\) itself
is a spatial divergence.  Indeed, on a periodic box,
\[
 q(t,x)=q_*\neq0,\qquad J(t,x)=0,\qquad W(t,x)=0
\]
satisfies the balance law, while \(q_*\) has nonzero spatial mean and
cannot be a periodic divergence.  A fixed-time source therefore needs a
charge subtraction, a reference-state difference, or a boundary carrier.
\end{firewall}

\subsection{Exact covariant origin of the work term}
\label{sec:e8v51-covariant}

Let \(T^{\mu\nu}\) be a symmetric stress tensor satisfying
\(\nabla_\mu T^{\mu\nu}=0\), and let \(\xi\) be a smooth vector field.
Define the energy current
\begin{equation}
 {\cal J}^\mu:=T^\mu{}_\nu\xi^\nu.
 \label{eq:e8v51-current}
\end{equation}

\begin{theorem}[Densitized covariant balance]
\label{thm:e8v51-covariant}
In any coordinate chart,
\begin{equation}
 \partial_\mu\!\left(\sqrt{|g|}\,{\cal J}^\mu\right)
 =
 {1\over2}\sqrt{|g|}\,T^{\mu\nu}
 ({\cal L}_\xi g)_{\mu\nu}.
 \label{eq:e8v51-covariant-balance}
\end{equation}
Hence \eqref{eq:e8v51-balance} holds with
\begin{equation}
 q=\sqrt{|g|}\,{\cal J}^0,\qquad
 J^i=\sqrt{|g|}\,{\cal J}^i,\qquad
 W={1\over2}\sqrt{|g|}\,T^{\mu\nu}
 ({\cal L}_\xi g)_{\mu\nu}.
 \label{eq:e8v51-qJW}
\end{equation}
If \(\xi\) is Killing, then \(W=0\).
\end{theorem}

\begin{proof}
Stress-energy conservation gives
\[
 \nabla_\mu{\cal J}^\mu
 =T^\mu{}_\nu\nabla_\mu\xi^\nu.
\]
Symmetry of \(T^{\mu\nu}\) removes the antisymmetric part of
\(\nabla_\mu\xi_\nu\), so
\[
 T^\mu{}_\nu\nabla_\mu\xi^\nu
 ={1\over2}T^{\mu\nu}
 (\nabla_\mu\xi_\nu+\nabla_\nu\xi_\mu)
 ={1\over2}T^{\mu\nu}({\cal L}_\xi g)_{\mu\nu}.
\]
Finally,
\(\sqrt{|g|}\nabla_\mu{\cal J}^\mu
=\partial_\mu(\sqrt{|g|}{\cal J}^\mu)\).
Splitting the time and spatial derivatives proves
\eqref{eq:e8v51-qJW}.
\end{proof}

\begin{corollary}[Geometric work is the sole non-adjoint residual]
\label{cor:e8v51-work}
For the covariant Standard Model stress tensor, the endpoint and
time-window sources have the CASFC decomposition with directly
constructed primitive equal to the corresponding integral of the
physical current \({\cal J}^i\).  The residual is precisely the smeared
deformation work
\[
 {1\over2}\sqrt{|g|}T^{\mu\nu}({\cal L}_\xi g)_{\mu\nu}.
\]
\end{corollary}

\begin{proof}
Insert \eqref{eq:e8v51-qJW} into
\Cref{thm:e8v51-endpoint,thm:e8v51-window}.
\end{proof}

\begin{firewall}[Dynamical spacetime is not stationary]
\label{fire:e8v51-stationary}
The choice \(W=0\) is exact for a Killing evolution, but a general
Einstein--Standard Model spacetime need not possess a timelike Killing
field.  The full programme must estimate the deformation tensor
\({\cal L}_\xi g\) together with the stress tensor, or absorb their
product into the nonlinear gravitational energy balance.  Stationarity
cannot be imposed as a substitute for this row.
\end{firewall}

\subsection{Stress--energy transfer card}
\label{sec:e8v51-card}

\begin{definition}[Stress--energy transfer card, SETC]
\label{def:e8v51-SETC}
SETC consists of:
\begin{enumerate}[label=\textup{(S\arabic*)},leftmargin=4em]
\item a regulated local Ward identity yielding
\(\partial_tq+\operatorname{div}J=W\);
\item uniform spacetime \(L^2\) control of the directly defined flux
\(J\) for the selected density difference or time window;
\item an estimate on the deformation-work residual \(W\) strong enough
to absorb \(\lambda_1^{-1/2}\|P_0W\|_2\);
\item a boundary-flux and total-charge identity;
\item compatibility with the constraint quotient and pseudoinverse of
V50;
\item stability under the V45 owner/martingale conditioning and the
V43--V44 observable tests; and
\item nonlinear, anomaly, renormalization, reflection, spin, and chiral
contacts for the full Standard Model stress tensor.
\end{enumerate}
\end{definition}

\begin{assessment}[Progress at V51]
\label{ass:e8v51-status}
The desired adjoint primitive is now constructed without an elliptic
inverse for two physically natural source classes: endpoint differences
and compact time-window derivatives.  Exact conservation gives a
volume-uniform constraint-energy response controlled by the physical
flux.  On a dynamical geometry the only residual is the covariant
deformation work.  Fixed-time absolute energy and total charge remain
outside this result.
\end{assessment}

\begin{programme}[Eighth-series V52: reference subtraction and work absorption]
\label{prog:e8v51-V52}
V52 will determine whether a reference-subtracted fixed-time Standard
Model source can be represented as a flux transfer plus a controlled
deformation-work remainder.  It will first prove a quantitative
comparison identity along an interpolation of states or geometries and
will reject any construction that defines its flux through the
constraint inverse.
\end{programme}

\begin{firewall}[Claim boundary after eighth-series V51]
\label{fire:e8v51-boundary}
V51 proves exact balance-law factorizations and their \(L^2\) bounds for
endpoint differences and compact time-window derivatives.  It identifies
the covariant metric-deformation residual and the integrated charge
equation.  It does not factor an arbitrary fixed-time energy density,
bound the interacting Standard Model flux uniformly, control the
deformation work cofinally, solve the nonlinear constraints, construct a
positive carrier, remove the cutoffs, or prove the full continuum.
\end{firewall}

\begin{theorem}[Eighth-series V51 result]
\label{thm:e8v51-result}
Every local balance law
\(\partial_tq+\operatorname{div}J=W\) gives the direct decompositions
\[
 q(b)-q(a)
 =-\operatorname{div}\int_a^bJ\,dt+\int_a^bW\,dt
\]
and
\[
 \int_a^b\chi' q\,dt
 =-\operatorname{div}\!\left(-\int_a^b\chi J\,dt\right)
 -\int_a^b\chi W\,dt.
\]
When \(W=0\), V50 converts either source into a constrained
gradient-energy response with no infrared spectral loss.
\end{theorem}

\begin{proof}
Use \Cref{thm:e8v51-endpoint,thm:e8v51-window,
cor:e8v51-endpoint-response,cor:e8v51-window-response}.
\end{proof}

\paragraph{Parent-version record.}
V51 was formed from
\texttt{Renewal\_SM\_Einstein\_Continuum\_Research\_Notes\_V50.tex}.
The V50 parent had (20\,631) logical source lines, (743\,719) bytes,
and SHA--256
\[
 \texttt{14FF04CF6DC934EA68C091D0B36AC74243205C40BF5DB0E30D65FC46F36E64C4}.
\]
The next compulsory companion audit is V55.

% =============================================================================
% END APPENDED EIGHTH-SERIES V51 MATERIAL
% =============================================================================
% =============================================================================
% BEGIN APPENDED EIGHTH-SERIES V52 MATERIAL
% =============================================================================

\clearpage
\section{Eighth-series V52: fixed-time reference subtraction by a local variation current}
\label{sec:v8v52}

\subsection{The preparation-cylinder identity}
\label{sec:e8v52-cylinder}

V51 factors density differences taken along physical time.  A
fixed-time comparison with a reference configuration requires another
local parameter.  Let \(s\in[0,1]\) label a preparation path and suppose
\begin{equation}
 \partial_s q_s+\operatorname{div}K_s=U_s
 \label{eq:e8v52-cylinder-balance}
\end{equation}
on a fixed spatial regulator.  The fields \(K_s\) and \(U_s\) must be
constructed locally from the path; they may not be defined using an
inverse Laplacian.

\begin{theorem}[Preparation-cylinder transfer]
\label{thm:e8v52-cylinder}
If \(q_s\in C([0,1];L^2(\Omega))\) and
\(K,U\in L^1_sL^2_x\), then
\begin{equation}
 q_1-q_0=-\operatorname{div}Z_{\rm prep}+R_{\rm prep},
 \qquad
 Z_{\rm prep}:=\int_0^1K_s\,ds,
 \qquad
 R_{\rm prep}:=\int_0^1U_s\,ds.
 \label{eq:e8v52-cylinder-split}
\end{equation}
Furthermore,
\begin{equation}
 \|Z_{\rm prep}\|_2\leq\int_0^1\|K_s\|_2\,ds,
 \qquad
 \|R_{\rm prep}\|_2\leq\int_0^1\|U_s\|_2\,ds.
 \label{eq:e8v52-cylinder-norm}
\end{equation}
\end{theorem}

\begin{proof}
Integrate \eqref{eq:e8v52-cylinder-balance} in \(s\).  The norm
estimates are Minkowski's integral inequality.
\end{proof}

\begin{firewall}[An arbitrary interpolation has no content]
\label{fire:e8v52-arbitrary}
The linear path \(q_s=(1-s)q_0+sq_1\), with \(K_s=0\) and
\(U_s=q_1-q_0\), satisfies
\eqref{eq:e8v52-cylinder-balance} but gives no infrared improvement.
Likewise,
\(K_s=\nabla(-\Delta)^\dagger\partial_sq_s\) merely hides the desired
inverse estimate inside the primitive.  A useful preparation identity
must arise from a local first variation or Ward law.
\end{firewall}

\subsection{A local Hamiltonian first-variation current}
\label{sec:e8v52-variation}

Let \(z=(z^1,\ldots,z^N)\) denote regulated canonical matter variables
on \(\Omega\), including momenta when present.  Consider a first-order
local density
\begin{equation}
 h(z,\partial z)
 \label{eq:e8v52-density}
\end{equation}
with no explicit dependence on a constraint inverse.  For a \(C^1\)
path \(z_s\), set \(\dot z_s=\partial_sz_s\), and define
\begin{align}
 {\cal E}_A(h)(z_s)
 &:={\partial h\over\partial z^A}(z_s,\partial z_s)
 -\partial_i\!\left(
 {\partial h\over\partial(\partial_i z^A)}(z_s,\partial z_s)
 \right),
 \label{eq:e8v52-Euler}\\
 \Theta_s^i
 &:=
 \sum_{A=1}^N
 {\partial h\over\partial(\partial_i z^A)}
 (z_s,\partial z_s)\,\dot z_s^A.
 \label{eq:e8v52-Theta}
\end{align}

\begin{theorem}[Local fixed-time comparison identity]
\label{thm:e8v52-first-variation}
Pointwise in the preparation parameter and distributionally in space,
\begin{equation}
 \partial_sh(z_s,\partial z_s)
 =
 \sum_A{\cal E}_A(h)(z_s)\dot z_s^A
 +\operatorname{div}\Theta_s.
 \label{eq:e8v52-first-variation}
\end{equation}
Consequently,
\begin{equation}
 h(z_1)-h(z_0)
 =-\operatorname{div}Z_h+R_h,
 \label{eq:e8v52-h-split}
\end{equation}
where
\begin{equation}
 Z_h:=-\int_0^1\Theta_s\,ds,
 \qquad
 R_h:=\int_0^1
 \sum_A{\cal E}_A(h)(z_s)\dot z_s^A\,ds.
 \label{eq:e8v52-ZR}
\end{equation}
Both terms are local functions of the path and the density.
\end{theorem}

\begin{proof}
The chain rule gives
\[
 \partial_sh
 =\sum_A{\partial h\over\partial z^A}\dot z_s^A
 +\sum_{A,i}
 {\partial h\over\partial(\partial_i z^A)}
 \partial_i\dot z_s^A.
\]
Apply the product rule to the second term:
\[
 {\partial h\over\partial(\partial_i z^A)}
 \partial_i\dot z_s^A
 =
 \partial_i\!\left(
 {\partial h\over\partial(\partial_i z^A)}\dot z_s^A
 \right)
 -
 \partial_i\!\left(
 {\partial h\over\partial(\partial_i z^A)}
 \right)\dot z_s^A.
\]
Summing proves \eqref{eq:e8v52-first-variation}.  Integrating in \(s\)
and choosing the sign in \eqref{eq:e8v52-ZR} proves
\eqref{eq:e8v52-h-split}.  Only differentiation and multiplication at
the same regulated point were used.
\end{proof}

\begin{corollary}[Quantitative local comparison]
\label{cor:e8v52-quantitative}
For exponents \(p_A,r_A\in[2,\infty]\) satisfying
\(p_A^{-1}+r_A^{-1}=1/2\), suppose
\[
 P^i_{A,s}:={\partial h\over\partial(\partial_i z^A)}(z_s,\partial z_s)
 \in L^{p_A}(\Omega),
 \qquad
 \dot z_s^A\in L^{r_A}(\Omega).
\]
Then
\begin{equation}
 \|Z_h\|_2
 \leq
 \int_0^1
 \sum_{A,i}\|P^i_{A,s}\|_{p_A}\|\dot z_s^A\|_{r_A}\,ds.
 \label{eq:e8v52-Z-bound}
\end{equation}
If \(a_A^{-1}+b_A^{-1}=1/2\), then also
\begin{equation}
 \|R_h\|_2
 \leq
 \int_0^1
 \sum_A
 \|{\cal E}_A(h)(z_s)\|_{a_A}
 \|\dot z_s^A\|_{b_A}\,ds.
 \label{eq:e8v52-R-bound}
\end{equation}
\end{corollary}

\begin{proof}
Use \eqref{eq:e8v52-ZR}, Minkowski's inequality in \(s\), and
Hölder's inequality in \(x\) for each product.
\end{proof}

\begin{corollary}[Constraint response of a reference-subtracted density]
\label{cor:e8v52-response}
On the periodic regulator, set
\[
 u_h=(-\Delta)^\dagger P_0\bigl(h(z_1)-h(z_0)\bigr).
\]
Then
\begin{equation}
 \|\nabla u_h\|_2
 \leq
 \|Z_h\|_2+\lambda_1^{-1/2}\|P_0R_h\|_2.
 \label{eq:e8v52-response}
\end{equation}
In particular, a path contained in the spatial critical set
\({\cal E}_A(h)=0\) has a norm-one constraint response controlled by
the local variation current \(Z_h\).
\end{corollary}

\begin{proof}
Apply \Cref{thm:e8v50-constraint-adjoint} to
\eqref{eq:e8v52-h-split}.  If all Euler expressions vanish along the
path, then \(R_h=0\).
\end{proof}

\begin{remark}[Canonical fields and momenta]
\label{rem:e8v52-canonical}
Variables without spatial derivatives in \(h\), such as many canonical
momenta, contribute to the Euler residual but not to \(\Theta_s\).
Thus the theorem does not falsely convert an arbitrary change of local
kinetic energy into a spatial flux.  It records exactly which part comes
from integration by parts and which part remains a physical bulk source.
\end{remark}

\begin{example}[Scalar-field density]
\label{ex:e8v52-scalar}
For
\[
 h(\phi,\pi,\nabla\phi)
 ={1\over2}\pi^2+{1\over2}|\nabla\phi|^2+V(\phi),
\]
one has
\[
 \Theta_s=\dot\phi_s\nabla\phi_s,\qquad
 {\cal E}_\phi(h)=-\Delta\phi+V'(\phi),\qquad
 {\cal E}_\pi(h)=\pi.
\]
Hence
\begin{equation}
 h(z_1)-h(z_0)
 =
 \operatorname{div}\int_0^1\dot\phi_s\nabla\phi_s\,ds
 +\int_0^1
 \left[
 \bigl(-\Delta\phi_s+V'(\phi_s)\bigr)\dot\phi_s
 +\pi_s\dot\pi_s
 \right]ds.
 \label{eq:e8v52-scalar-split}
\end{equation}
The kinetic change \(\int\pi_s\dot\pi_sds\) is a genuine residual unless
the preparation path supplies an additional identity.
\end{example}

\begin{proof}
Differentiate the displayed density and apply
\Cref{thm:e8v52-first-variation}.  The formula also follows directly
from
\(\nabla\phi_s\cdot\nabla\dot\phi_s
=\operatorname{div}(\dot\phi_s\nabla\phi_s)
-\dot\phi_s\Delta\phi_s\).
\end{proof}

\subsection{Quantifying the covariant deformation work}
\label{sec:e8v52-work}

\begin{lemma}[Hölder estimate for geometric work]
\label{lem:e8v52-work}
Let
\[
 W={1\over2}\sqrt{|g|}\,T^{\mu\nu}
 ({\cal L}_\xi g)_{\mu\nu}
\]
as in V51.  With tensor norms taken in a uniformly nondegenerate
regulated frame, if \(p^{-1}+r^{-1}=1/2\), then
\begin{equation}
 \|W\|_{L^2_{t,x}}
 \leq {1\over2}
 \|\sqrt{|g|}\,T\|_{L^p_{t,x}}
 \|{\cal L}_\xi g\|_{L^r_{t,x}}.
 \label{eq:e8v52-work-holder}
\end{equation}
\end{lemma}

\begin{proof}
The pointwise Cauchy--Schwarz contraction bound gives
\[
 |W|\leq{1\over2}\sqrt{|g|}\,|T|\,|{\cal L}_\xi g|.
\]
Hölder's inequality with \(1/2=1/p+1/r\) proves the claim.
\end{proof}

\begin{definition}[Work-absorption quantity]
\label{def:e8v52-work-quantity}
For regulator \(n\), define
\begin{equation}
 {\mathfrak W}_n(p,r)
 :=
 {1\over2}\lambda_{1,n}^{-1/2}
 \|\sqrt{|g_n|}\,T_n\|_{L^p}
 \|{\cal L}_{\xi_n}g_n\|_{L^r}.
 \label{eq:e8v52-work-quantity}
\end{equation}
\end{definition}

\begin{theorem}[Sufficient work-absorption condition]
\label{thm:e8v52-work-absorption}
For a cofinal family of endpoint or compact-window sources of V51,
assume the factors \(T_n^{1/2}\) or \(\|\chi_n\|_{L^2_t}\) are uniformly
bounded.  The residual part of the constraint-energy bound is uniformly
bounded if
\begin{equation}
 \sup_n{\mathfrak W}_n(p,r)<\infty.
 \label{eq:e8v52-work-uniform}
\end{equation}
It vanishes cofinally if
\(\mathfrak W_n(p,r)\to0\).
\end{theorem}

\begin{proof}
Insert \eqref{eq:e8v52-work-holder} into
\eqref{eq:e8v51-endpoint-response} or
\eqref{eq:e8v51-window-response}, including the corresponding fixed
time-window prefactor.  The asserted uniformity and convergence follow
directly from \eqref{eq:e8v52-work-quantity}.
\end{proof}

\begin{firewall}[Small metric deformation is an assumption, not a theorem]
\label{fire:e8v52-smallness}
The factor \(\lambda_{1,n}^{-1/2}\) grows with physical box size.
Bounded stress energy and pointwise small deformation do not alone imply
\eqref{eq:e8v52-work-uniform}.  One needs a quantitative product estimate
at the same cofinal scales, or an exact combined
matter--gravitational current which converts \(W\) into another
divergence.
\end{firewall}

\subsection{Reference-comparison card}
\label{sec:e8v52-card}

\begin{definition}[Local reference-comparison card, LRCC]
\label{def:e8v52-LRCC}
LRCC consists of:
\begin{enumerate}[label=\textup{(R\arabic*)},leftmargin=4em]
\item a reference configuration \(z_0\) with its charge and boundary
data;
\item a regulator-compatible local preparation path \(z_s\) to \(z_1\);
\item the local variation current \(\Theta_s\) and the exact Euler
residual in \eqref{eq:e8v52-first-variation};
\item uniform bounds of the form
\eqref{eq:e8v52-Z-bound}--\eqref{eq:e8v52-R-bound};
\item equality or explicit transfer of total charge between the two
endpoints;
\item the work-absorption condition
\eqref{eq:e8v52-work-uniform}, or an exact combined current replacing
it; and
\item compatibility with constraint reduction, conditioning,
renormalization, reflection, and the local observable class.
\end{enumerate}
\end{definition}

\begin{assessment}[Progress at V52]
\label{ass:e8v52-status}
A fixed-time reference subtraction now has an exact local decomposition.
Its spatial-derivative variation is a directly constructed divergence,
while its spatial Euler variation is an explicit residual.  The latter,
together with the V51 deformation work, is the only part that pays the
infrared gap in the scalar constraint-energy estimate.  The construction
does not use an inverse operator and therefore passes the V50
non-circularity test.
\end{assessment}

\begin{programme}[Eighth-series V53: diffeomorphism directions and the combined current]
\label{prog:e8v52-V53}
V53 will specialize the local variation identity to infinitesimal
spatial diffeomorphisms.  It will prove the associated Noether
factorization, determine which Euler residual is a constraint term, and
test whether matter deformation work cancels against the gravitational
current rather than being estimated through the infrared gap.
\end{programme}

\begin{firewall}[Claim boundary after eighth-series V52]
\label{fire:e8v52-boundary}
V52 proves a local first-variation decomposition for first-order
regulated densities, quantitative product estimates, and a sufficient
work-absorption criterion.  It does not construct a uniformly controlled
preparation path for the interacting Standard Model, prove the Euler
residual small, establish the combined Einstein--matter Noether current,
control composite-operator renormalization, remove cutoffs, or prove the
full continuum.
\end{firewall}

\begin{theorem}[Eighth-series V52 result]
\label{thm:e8v52-result}
For a first-order local density and a local preparation path,
\[
 h(z_1)-h(z_0)
 =
 \operatorname{div}\int_0^1
 \sum_A{\partial h\over\partial(\partial_i z^A)}
 \dot z_s^A\,ds
 +\int_0^1\sum_A{\cal E}_A(h)(z_s)\dot z_s^A\,ds.
\]
The first term receives the norm-one V50 constraint response.  Only the
explicit Euler residual pays the inverse spectral factor.
\end{theorem}

\begin{proof}
Use \Cref{thm:e8v52-first-variation,cor:e8v52-response}.
\end{proof}

\paragraph{Parent-version record.}
V52 was formed from
\texttt{Renewal\_SM\_Einstein\_Continuum\_Research\_Notes\_V51.tex}.
The V51 parent had (21\,024) logical source lines, (756\,548) bytes,
and SHA--256
\[
 \texttt{E2662DE17703520E7C64127942D718D9F66EF170E298B0428A31CB4426ECC687}.
\]
The next compulsory companion audit is V55.

% =============================================================================
% END APPENDED EIGHTH-SERIES V52 MATERIAL
% =============================================================================
% =============================================================================
% BEGIN APPENDED EIGHTH-SERIES V53 MATERIAL
% =============================================================================

\clearpage
\section{Eighth-series V53: diffeomorphism transfer and the combined Einstein--matter current}
\label{sec:v8v53}

\subsection{Spatial gauge orbits give exact preparation currents}
\label{sec:e8v53-spatial}

The V52 preparation residual vanishes for a pure spatial coordinate
change for a reason stronger than criticality.  The energy source is a
spatial density, and the Lie derivative of a density is a divergence.

\begin{theorem}[Finite spatial-diffeomorphism factorization]
\label{thm:e8v53-spatial}
Let \(F_s:\Omega\to\Omega\), \(s\in[0,1]\), be a \(C^1\) path of
orientation-preserving diffeomorphisms.  Put
\[
 v_s:=\dot F_s\circ F_s^{-1},
 \qquad
 \eta_s:=F_s^*v_s,
\]
where the second expression denotes the pullback of a vector field.
Let \(q_0\,d^dx\) be an integrable density and set
\[
 q_s\,d^dx=F_s^*(q_0\,d^dx).
\]
Then
\begin{equation}
 \partial_sq_s=\partial_i(q_s\eta_s^i)
 \label{eq:e8v53-density-Lie}
\end{equation}
with the pullback sign convention above, and
\begin{equation}
 q_1-q_0=-\operatorname{div}Z_{\rm diff},
 \qquad
 Z_{\rm diff}:=-\int_0^1q_s\eta_s\,ds.
 \label{eq:e8v53-diff-split}
\end{equation}
If \(q_s\eta_s\in L^1_sL^2_x\), then
\begin{equation}
 \|Z_{\rm diff}\|_2
 \leq\int_0^1\|q_s\eta_s\|_2\,ds.
 \label{eq:e8v53-diff-bound}
\end{equation}
On a periodic box or for a boundary-preserving flow,
\(\int_\Omega q_s\) is independent of \(s\).
\end{theorem}

\begin{proof}
The derivative of a pullback is
\[
 {d\over ds}F_s^*(q_0\,d^dx)
 =F_s^*{\cal L}_{v_s}(q_0\,d^dx)
 ={\cal L}_{\eta_s}F_s^*(q_0\,d^dx).
\]
The Lie derivative of a weight-one density is
\[
 {\cal L}_{\eta_s}(q_s\,d^dx)
 =\partial_i(q_s\eta_s^i)\,d^dx.
\]
Integrate in \(s\) and use the definition of \(Z_{\rm diff}\).
Minkowski's inequality proves \eqref{eq:e8v53-diff-bound}.  Integration
of \eqref{eq:e8v53-density-Lie} gives only a boundary flux, which
vanishes in the two stated cases.
\end{proof}

\begin{corollary}[No infrared loss along a controlled gauge orbit]
\label{cor:e8v53-gauge-response}
On the periodic regulator, let
\[
 u_{\rm diff}=(-\Delta)^\dagger(q_1-q_0).
\]
Then
\begin{equation}
 \|\nabla u_{\rm diff}\|_2
 \leq\int_0^1\|q_s\eta_s\|_2\,ds.
 \label{eq:e8v53-gauge-response}
\end{equation}
\end{corollary}

\begin{proof}
The charge equality in \Cref{thm:e8v53-spatial} makes \(q_1-q_0\)
mean zero.  Apply \Cref{thm:e8v50-constraint-adjoint} to
\eqref{eq:e8v53-diff-split}.
\end{proof}

\begin{firewall}[Gauge transfer does not compare physical states]
\label{fire:e8v53-gauge}
\Cref{thm:e8v53-spatial} compares two representatives on the same
spatial diffeomorphism orbit.  A path changing gauge-invariant matter
content is not covered.  Treating a physical preparation as a coordinate
flow would erase precisely the energy source that the Einstein
constraint must carry.
\end{firewall}

\subsection{A local diffeomorphism Noether identity}
\label{sec:e8v53-Noether}

Let \(\Phi=(g,\psi)\) denote all regulated metric and Standard Model
fields.  On a coordinate patch, suppose a local total Lagrangian density
\({\cal L}_{\rm tot}\) has first variation
\begin{equation}
 \delta{\cal L}_{\rm tot}
 =
 {\cal E}_A(\Phi)\,\delta\Phi^A
 +\partial_\mu\Theta^\mu(\Phi;\delta\Phi).
 \label{eq:e8v53-first-variation}
\end{equation}
Boundary improvements may be included in \(\Theta\).  Diffeomorphism
covariance means that for a compactly supported vector field \(\xi\),
\begin{equation}
 \delta_\xi\Phi={\cal L}_\xi\Phi,
 \qquad
 \delta_\xi{\cal L}_{\rm tot}
 =\partial_\mu(\xi^\mu{\cal L}_{\rm tot}).
 \label{eq:e8v53-covariance}
\end{equation}

\begin{theorem}[Off-shell local Noether identity]
\label{thm:e8v53-Noether}
Define the densitized current
\begin{equation}
 {\cal N}_\xi^\mu
 :=
 \xi^\mu{\cal L}_{\rm tot}
 -\Theta^\mu(\Phi;{\cal L}_\xi\Phi).
 \label{eq:e8v53-Noether-current}
\end{equation}
Then
\begin{equation}
 \partial_\mu{\cal N}_\xi^\mu
 =
 {\cal E}_A(\Phi)\,{\cal L}_\xi\Phi^A.
 \label{eq:e8v53-Noether-identity}
\end{equation}
In particular, \({\cal N}_\xi\) is conserved on every coupled
Einstein--Standard Model solution.
\end{theorem}

\begin{proof}
Insert \(\delta\Phi={\cal L}_\xi\Phi\) in
\eqref{eq:e8v53-first-variation} and equate the result with
\eqref{eq:e8v53-covariance}.  Moving the boundary variation to the
other side gives \eqref{eq:e8v53-Noether-identity}.  On a coupled
solution all Euler expressions vanish.
\end{proof}

\begin{remark}[Second derivatives and boundary improvements]
\label{rem:e8v53-boundary-improvement}
The Einstein--Hilbert density contains second derivatives in its raw
coordinate form.  One may either use the corresponding higher-order
first-variation formula or subtract the standard local boundary
divergence on a patch.  Both choices give
\eqref{eq:e8v53-Noether-identity}; they change
\({\cal N}_\xi\) by a local superpotential.  V54 will quantify that
ambiguity rather than assume it is harmless.
\end{remark}

\subsection{Exact cancellation of matter deformation work}
\label{sec:e8v53-work}

Let the matter equations hold and define the densitized Hilbert
stress-energy current from V51,
\begin{equation}
 {\cal C}_{m,\xi}^\mu
 :=
 \sqrt{|g|}\,T^\mu{}_\nu\xi^\nu.
 \label{eq:e8v53-matter-current}
\end{equation}
By \Cref{thm:e8v51-covariant},
\begin{equation}
 \partial_\mu{\cal C}_{m,\xi}^\mu=W_\xi,
 \qquad
 W_\xi:={1\over2}\sqrt{|g|}\,T^{\mu\nu}
 ({\cal L}_\xi g)_{\mu\nu}.
 \label{eq:e8v53-matter-work}
\end{equation}

\begin{theorem}[Combined-current work cancellation]
\label{thm:e8v53-combined}
On a coupled solution, define the local gravitational remainder current
\begin{equation}
 {\cal C}_{g,\xi}^\mu
 :={\cal N}_\xi^\mu-{\cal C}_{m,\xi}^\mu.
 \label{eq:e8v53-gravity-current}
\end{equation}
Then
\begin{equation}
 \partial_\mu{\cal C}_{g,\xi}^\mu=-W_\xi,
 \qquad
 \partial_\mu\bigl(
 {\cal C}_{m,\xi}^\mu+{\cal C}_{g,\xi}^\mu
 \bigr)=0.
 \label{eq:e8v53-work-cancel}
\end{equation}
Thus the deformation work is an internal transfer between the matter
and gravitational currents; it is not a residual of the combined local
balance.
\end{theorem}

\begin{proof}
By \Cref{thm:e8v53-Noether},
\(\partial_\mu{\cal N}_\xi^\mu=0\) on the coupled equations.
Subtract \eqref{eq:e8v53-matter-work} and use
\eqref{eq:e8v53-gravity-current}.
\end{proof}

\begin{corollary}[Exact combined endpoint factorization]
\label{cor:e8v53-combined-endpoint}
Write
\[
 q_{\rm tot}:={\cal N}_\xi^0,\qquad
 J_{\rm tot}^i:={\cal N}_\xi^i.
\]
On every coordinate time slab on which the coupled equations and current
are defined,
\begin{equation}
 q_{\rm tot}(b)-q_{\rm tot}(a)
 =-\operatorname{div}\int_a^bJ_{\rm tot}(t)\,dt.
 \label{eq:e8v53-combined-endpoint}
\end{equation}
Consequently, on a periodic regulator,
\begin{equation}
 \left\|
 \nabla(-\Delta)^\dagger
 \bigl(q_{\rm tot}(b)-q_{\rm tot}(a)\bigr)
 \right\|_2
 \leq
 T^{1/2}\|J_{\rm tot}\|_{L^2_{t,x}}.
 \label{eq:e8v53-combined-response}
\end{equation}
\end{corollary}

\begin{proof}
Split the second identity in \eqref{eq:e8v53-work-cancel} into time and
space.  Then apply
\Cref{thm:e8v51-endpoint,thm:e8v50-constraint-adjoint}.
\end{proof}

\begin{corollary}[Matter work need not pay the infrared gap]
\label{cor:e8v53-no-work-gap}
If the source used by the scalar constraint is the combined density
difference in \eqref{eq:e8v53-combined-endpoint}, the work term
\(W_\xi\) cancels before constraint inversion and the sufficient
smallness condition \(\mathfrak W_n\to0\) from V52 is unnecessary for
that source.
\end{corollary}

\begin{proof}
The combined balance has zero residual by
\eqref{eq:e8v53-work-cancel}.  Apply the exact-adjoint part of
\Cref{thm:e8v50-constraint-adjoint}.
\end{proof}

\begin{firewall}[Current alignment is still missing]
\label{fire:e8v53-alignment}
The local Noether density \(q_{\rm tot}={\cal N}_\xi^0\) depends on the
choice of boundary improvement, evolution field \(\xi\), and constraint
representative.  V53 has not proved that the source appearing in the
regulated scalar Einstein constraint is exactly this density, nor that
their difference is a uniformly bounded divergence.  The cancellation
of \(W_\xi\) may be used in the compiler only after that alignment is
established.
\end{firewall}

\begin{firewall}[Conservation does not imply positivity]
\label{fire:e8v53-positivity}
The combined Noether current is conserved, but its local gravitational
part is generally gauge- and improvement-dependent and need not have a
positive density.  Therefore
\eqref{eq:e8v53-combined-response} is a source-response estimate, not a
construction of the positive reflection-compatible carrier required by
the continuum programme.
\end{firewall}

\subsection{Combined-current alignment card}
\label{sec:e8v53-card}

\begin{definition}[Combined-current alignment card, CCAC]
\label{def:e8v53-CCAC}
CCAC consists of:
\begin{enumerate}[label=\textup{(N\arabic*)},leftmargin=4em]
\item a regulated diffeomorphism-covariant total action and its local
first-variation current \(\Theta\);
\item the off-shell identity
\eqref{eq:e8v53-Noether-identity}, including gauge-fixing and ghost
terms when present;
\item an exact comparison between the scalar-constraint source and
\({\cal N}_\xi^0\), modulo a directly controlled divergence, equations
of motion, and charge terms;
\item a uniform \(L^2\) bound on the aligned spatial current;
\item control of superpotential and boundary-improvement ambiguities;
\item a separate positive-carrier or reflection argument; and
\item compatibility with anomalies, chiral matter, renormalization,
owner conditioning, and cutoff removal.
\end{enumerate}
\end{definition}

\begin{assessment}[Progress at V53]
\label{ass:e8v53-status}
Pure spatial diffeomorphism comparisons now have an exact local
constraint primitive.  More substantially, the covariant matter work
residual identified in V51 cancels identically against a local
gravitational Noether current on the coupled equations.  This removes
the V52 large-volume smallness demand for an aligned combined density.
The next bottleneck is precise current alignment and control of the
superpotential ambiguity, not the deformation-work product itself.
\end{assessment}

\begin{programme}[Eighth-series V54: superpotential invariance and alignment defect]
\label{prog:e8v53-V54}
V54 will quantify changes
\({\cal N}^\mu\mapsto{\cal N}^\mu+\partial_\nu
U^{[\mu\nu]}\), prove their norm-one effect on the scalar
constraint-energy response, and isolate the remaining equations-of-motion
and charge defects in the comparison with the ADM constraint source.
\end{programme}

\begin{firewall}[Claim boundary after eighth-series V53]
\label{fire:e8v53-boundary}
V53 proves the exact gauge-orbit density factorization, the off-shell
diffeomorphism Noether identity, and on-shell cancellation of matter
deformation work by a local gravitational remainder current.  It does
not prove current-to-constraint alignment, uniform current norms,
positive gravitational energy density, quantum Ward identities,
anomaly cancellation at the regulator, cutoff removal, or the full
continuum.
\end{firewall}

\begin{theorem}[Eighth-series V53 result]
\label{thm:e8v53-result}
For a local diffeomorphism-covariant regulated Einstein--Standard Model
action, the current
\[
 {\cal N}_\xi^\mu
 =\xi^\mu{\cal L}_{\rm tot}
 -\Theta^\mu(\Phi;{\cal L}_\xi\Phi)
\]
satisfies
\(\partial_\mu{\cal N}_\xi^\mu
={\cal E}_A{\cal L}_\xi\Phi^A\).
On coupled solutions it is conserved, and its gravitational remainder
has divergence exactly opposite to the matter deformation work.
\end{theorem}

\begin{proof}
Use \Cref{thm:e8v53-Noether,thm:e8v53-combined}.
\end{proof}

\paragraph{Parent-version record.}
V53 was formed from
\texttt{Renewal\_SM\_Einstein\_Continuum\_Research\_Notes\_V52.tex}.
The V52 parent had (21\,432) logical source lines, (769\,270) bytes,
and SHA--256
\[
 \texttt{D1F3E2BED6AA2FD9A4423478CB460C663F9DDC13711BAED69786FE0023329853}.
\]
The next compulsory companion audit is V55.

% =============================================================================
% END APPENDED EIGHTH-SERIES V53 MATERIAL
% =============================================================================
% =============================================================================
% BEGIN APPENDED EIGHTH-SERIES V54 MATERIAL
% =============================================================================

\clearpage
\section{Eighth-series V54: superpotential stability and the constraint-alignment defect}
\label{sec:v8v54}

\subsection{Local current improvements}
\label{sec:e8v54-improvement}

Let \({\cal N}^\mu\) and \(\widetilde{\cal N}^\mu\) be two local
representatives of an Einstein--matter current.  The standard
improvement freedom is
\begin{equation}
 \widetilde{\cal N}^\mu-{\cal N}^\mu
 =
 \partial_\nu U^{[\mu\nu]}+F^\mu,
 \qquad
 U^{[\mu\nu]}=-U^{[\nu\mu]}.
 \label{eq:e8v54-improvement}
\end{equation}
The term \(F^\mu\) records equations-of-motion, gauge-fixing, regulator,
or Ward defects rather than hiding them in the superpotential.

\begin{theorem}[Density and charge change under improvement]
\label{thm:e8v54-improvement}
Let \(q={\cal N}^0\), \(\widetilde q=\widetilde{\cal N}^0\), and
\(Z_U^i:=-U^{[0i]}\).  Then
\begin{equation}
 \widetilde q-q
 =-\operatorname{div}Z_U+F^0.
 \label{eq:e8v54-density-change}
\end{equation}
On a bounded spatial cell,
\begin{equation}
 \int_\Omega(\widetilde q-q)\,dx
 =
 \int_{\partial\Omega}U^{[0i]}n_i\,d\sigma
 +\int_\Omega F^0\,dx.
 \label{eq:e8v54-charge-change}
\end{equation}
In particular, a pure superpotential preserves the integrated charge on
a periodic regulator.
\end{theorem}

\begin{proof}
Set \(\mu=0\) in \eqref{eq:e8v54-improvement}.  Antisymmetry gives
\(U^{[00]}=0\), so
\[
 \widetilde q-q=\partial_iU^{[0i]}+F^0
 =-\operatorname{div}Z_U+F^0.
\]
The divergence theorem proves \eqref{eq:e8v54-charge-change}.
\end{proof}

\begin{theorem}[Norm-one stability of the constraint-energy response]
\label{thm:e8v54-response}
On the periodic regulator, let
\[
 u_U=(-\Delta)^\dagger P_0(\widetilde q-q).
\]
Then
\begin{equation}
 \|\nabla u_U\|_2
 \leq
 \|P_{\rm grad}Z_U\|_2
 +\lambda_1^{-1/2}\|P_0F^0\|_2
 \leq
 \|U^{[0\cdot]}\|_2
 +\lambda_1^{-1/2}\|P_0F^0\|_2.
 \label{eq:e8v54-response}
\end{equation}
For a pure superpotential the estimate has no inverse spectral factor.
\end{theorem}

\begin{proof}
Apply \Cref{thm:e8v50-constraint-adjoint} to
\eqref{eq:e8v54-density-change}.  Since \(Z_U=-U^{[0\cdot]}\),
orthogonal projection is contractive.
\end{proof}

\begin{corollary}[Two improved endpoint transfers remain equivalent]
\label{cor:e8v54-endpoint}
Suppose \(F^\mu=0\) and both currents are conserved.  Their endpoint
density transfers differ by
\begin{equation}
 \bigl(\widetilde q(b)-\widetilde q(a)\bigr)
 -\bigl(q(b)-q(a)\bigr)
 =
 \operatorname{div}\!\left(
 U^{[0\cdot]}(b)-U^{[0\cdot]}(a)
 \right).
 \label{eq:e8v54-endpoint}
\end{equation}
The associated scalar constraint-energy responses differ by at most
\begin{equation}
 \|U^{[0\cdot]}(b)-U^{[0\cdot]}(a)\|_2.
 \label{eq:e8v54-endpoint-bound}
\end{equation}
\end{corollary}

\begin{proof}
Apply \eqref{eq:e8v54-density-change} at \(a\) and \(b\), subtract, and
use \Cref{thm:e8v54-response} with zero residual.
\end{proof}

\subsection{The exact alignment defect}
\label{sec:e8v54-alignment}

Let \(S_{\rm ADM}\) denote the regulated source appearing in the scalar
Einstein constraint after the V49 scalar--vector--TT split.  Let
\(q_{\cal N}\) be the density of a combined current from V53.  Define an
alignment decomposition to be an identity
\begin{equation}
 S_{\rm ADM}-q_{\cal N}
 =
 -\operatorname{div}Z_{\rm al}
 +R_{\rm EOM}+R_{\rm reg}+S_{\rm ch},
 \label{eq:e8v54-alignment}
\end{equation}
where \(S_{\rm ch}\) lies in the constant, harmonic, boundary-charge, or
ADM-charge sector.

\begin{theorem}[Alignment-defect response estimate]
\label{thm:e8v54-alignment}
After projecting away the explicitly retained charge sector, put
\[
 u_{\rm al}
 :=
 (-\Delta)^\dagger P_0
 \bigl(S_{\rm ADM}-q_{\cal N}-S_{\rm ch}\bigr).
\]
Then
\begin{equation}
 \|\nabla u_{\rm al}\|_2
 \leq
 \|P_{\rm grad}Z_{\rm al}\|_2
 +\lambda_1^{-1/2}
 \|P_0(R_{\rm EOM}+R_{\rm reg})\|_2.
 \label{eq:e8v54-alignment-bound}
\end{equation}
Thus superpotential and local boundary improvements enter with constant
one, whereas equation and regulator defects still pay the infrared
factor.
\end{theorem}

\begin{proof}
Subtract \(S_{\rm ch}\) in \eqref{eq:e8v54-alignment}, apply \(P_0\),
and use \Cref{thm:e8v50-constraint-adjoint}.
\end{proof}

\begin{corollary}[Cofinal alignment criterion]
\label{cor:e8v54-cofinal}
For regulators \(n\), assume
\begin{align}
 \sup_n\|Z_{{\rm al},n}\|_2&<\infty,
 \label{eq:e8v54-Z-cofinal}\\
 \lambda_{1,n}^{-1/2}
 \|P_0(R_{{\rm EOM},n}+R_{{\rm reg},n})\|_2&\longrightarrow0.
 \label{eq:e8v54-R-cofinal}
\end{align}
Then the aligned constraint-energy discrepancy is uniformly bounded,
and its non-adjoint part vanishes cofinally.
\end{corollary}

\begin{proof}
Insert \eqref{eq:e8v54-Z-cofinal}--\eqref{eq:e8v54-R-cofinal} in
\eqref{eq:e8v54-alignment-bound}.
\end{proof}

\begin{remark}[Canonical target identity]
\label{rem:e8v54-canonical}
For a diffeomorphism-covariant Hamiltonian formulation, the target is
the local version of
\[
 H[\xi]
 =
 \int_\Omega
 \bigl(\xi^\perp{\cal C}_\perp+\xi^i{\cal C}_i\bigr)\,dx
 +B_{\partial\Omega}[\xi].
\]
Its density identifies the Noether generator with the Hamiltonian and
momentum constraints modulo a boundary divergence.  V54 does not assume
this identity survives a particular lattice, gauge-fixing, chiral, and
renormalized construction; its failure is exactly \(R_{\rm reg}\).
\end{remark}

\subsection{Classical equation terms and quantum contacts}
\label{sec:e8v54-contact}

An equation-of-motion factor vanishes on a classical solution, but its
insertion in a quantum expectation need not vanish.  The elementary
finite-regulator identity already shows the contact term.

\begin{theorem}[Finite-dimensional Schwinger--Dyson contact identity]
\label{thm:e8v54-SD}
Let \(x\in\mathbb R^N\), let \(S\in C^1\), and suppose
\(e^{-S}F\) decays sufficiently fast that integration by parts has no
boundary term.  With
\[
 d\mu=Z^{-1}e^{-S(x)}\,dx,
\]
one has for every \(C^1\) test \(F\),
\begin{equation}
 \left\langle
 {\partial S\over\partial x^a}F
 \right\rangle_\mu
 =
 \left\langle{\partial F\over\partial x^a}\right\rangle_\mu.
 \label{eq:e8v54-SD}
\end{equation}
\end{theorem}

\begin{proof}
Integrate
\[
 0=\int_{\mathbb R^N}
 {\partial\over\partial x^a}\bigl(F(x)e^{-S(x)}\bigr)\,dx
 =\int
 \left({\partial F\over\partial x^a}
 -F{\partial S\over\partial x^a}\right)e^{-S}\,dx
\]
and divide by \(Z\).
\end{proof}

\begin{corollary}[Equation defects become differentiated tests]
\label{cor:e8v54-contact}
If a regulated alignment defect contains
\[
 R_0=\sum_a A^a(x){\partial S\over\partial x^a},
\]
then against a test \(F\),
\begin{equation}
 \langle R_0F\rangle_\mu
 =
 \sum_a
 \left\langle
 {\partial A^a\over\partial x^a}F
 \right\rangle_\mu
 +
 \sum_a\langle A^a\partial_aF\rangle_\mu.
 \label{eq:e8v54-contact-expanded}
\end{equation}
For the divergence-corrected equation defect
\[
 R_{\rm EOM}
 :=
 \sum_a\left(
 A^a{\partial S\over\partial x^a}
 -{\partial A^a\over\partial x^a}
 \right),
\]
one has the pure contact identity
\begin{equation}
 \langle R_{\rm EOM}F\rangle_\mu
 =
 \sum_a\langle A^a\partial_aF\rangle_\mu.
 \label{eq:e8v54-contact-simple}
\end{equation}
\end{corollary}

\begin{proof}
Apply \eqref{eq:e8v54-SD} to \(A^aF\):
\[
 \left\langle
 A^a{\partial S\over\partial x^a}F
 \right\rangle
 =
 \langle\partial_a(A^aF)\rangle.
\]
Expanding the derivative proves the first form.  Subtracting the
coefficient-divergence term gives the second.
\end{proof}

\begin{firewall}[On-shell equality is insufficient for quantum alignment]
\label{fire:e8v54-on-shell}
One may discard \(R_{\rm EOM}\) in a classical solution identity.
Inside regulated correlation functions it differentiates the observable
and creates contact terms as in
\eqref{eq:e8v54-contact-simple}.  These contacts must be paired with the
V44 spacetime-smoothed observable class and bounded uniformly; an
on-shell classical formula alone does not close CCAC.
\end{firewall}

\subsection{Positivity and boundary limitations}
\label{sec:e8v54-limitations}

\begin{proposition}[Improvement freedom does not prove local positivity]
\label{prop:e8v54-positivity}
Even when \(F^\mu=0\), the sign of a local current density is not
invariant under superpotential improvement.
\end{proposition}

\begin{proof}
On a one-dimensional periodic spatial regulator, start from \(q=0\).
Choose a nonconstant smooth periodic \(f\) and set \(U^{[01]}=f\).
Then \(\widetilde q=f'\).  Its integral is zero, and because it is
nonzero it must take both positive and negative values.  Thus a pure
improvement changes local sign while preserving total charge.
\end{proof}

\begin{firewall}[Boundary norms remain physical data]
\label{fire:e8v54-boundary-norm}
On a bounded domain, the superpotential changes the charge by
\eqref{eq:e8v54-charge-change}.  A bulk \(L^2\) bound on
\(U^{[0\cdot]}\) does not by itself control its normal trace.  The V45
owner/shell compiler must therefore receive a trace or collar estimate,
not only the periodic norm-one response theorem.
\end{firewall}

\subsection{Alignment-defect card}
\label{sec:e8v54-card}

\begin{definition}[Constraint-current alignment defect card, CCADC]
\label{def:e8v54-CCADC}
CCADC consists of:
\begin{enumerate}[label=\textup{(D\arabic*)},leftmargin=4em]
\item the exact regulated identity
\eqref{eq:e8v54-alignment};
\item uniform \(L^2\) control of its local superpotential primitive;
\item the cofinal residual estimate \eqref{eq:e8v54-R-cofinal};
\item explicit constant, harmonic, boundary, and ADM charge matching;
\item trace/collar control for every boundary improvement;
\item Schwinger--Dyson contact estimates for the chosen smoothed
observable class;
\item regulator Ward restoration, anomaly cancellation, and BRST
compatibility; and
\item a separate positive-carrier construction unaffected by local
improvement sign changes.
\end{enumerate}
\end{definition}

\begin{assessment}[Progress at V54]
\label{ass:e8v54-status}
The current ambiguity left by V53 is now quantitatively harmless in the
bulk energy topology whenever its superpotential has a uniform
\(L^2\) bound: it enters with constant one.  All dangerous infrared
dependence is isolated in the equation and regulator defect.  At the
quantum level the equation defect becomes a differentiated-test contact,
so alignment must be proved in the V44 observable topology.  Boundary
trace and positivity remain independent rows.
\end{assessment}

\begin{programme}[Eighth-series V55: companion audit and alignment route]
\label{prog:e8v54-V55}
V55 will perform the compulsory YM/NS audit.  It will compare the newest
curvature-adjoint and owner-martingale constructions with CCADC, then
choose between a regulated Ward derivation of
\eqref{eq:e8v54-alignment}, a boundary trace estimate, and a positive
combined carrier as the next upstream attack.
\end{programme}

\begin{firewall}[Claim boundary after eighth-series V54]
\label{fire:e8v54-boundary}
V54 proves superpotential stability of the scalar constraint-energy
response, the alignment-defect bound, the finite-regulator
Schwinger--Dyson contact identity, and the non-invariance of local
positivity.  It does not derive the ADM--Noether alignment for a concrete
regulator, bound its contacts or boundary trace cofinally, construct the
positive carrier, remove cutoffs, or prove the full continuum.
\end{firewall}

\begin{theorem}[Eighth-series V54 result]
\label{thm:e8v54-result}
If two regulated current densities differ by
\[
 \widetilde q-q
 =\partial_iU^{[0i]}+F^0,
\]
then their scalar constraint-energy responses differ by at most
\[
 \|U^{[0\cdot]}\|_2
 +\lambda_1^{-1/2}\|P_0F^0\|_2.
\]
Classical equation terms within \(F^0\) become differentiated-observable
contacts in regulated quantum expectations and may not simply be
discarded.
\end{theorem}

\begin{proof}
Use \Cref{thm:e8v54-response,thm:e8v54-SD,
cor:e8v54-contact}.
\end{proof}

\paragraph{Parent-version record.}
V54 was formed from
\texttt{Renewal\_SM\_Einstein\_Continuum\_Research\_Notes\_V53.tex}.
The V53 parent had (21\,797) logical source lines, (781\,610) bytes,
and SHA--256
\[
 \texttt{47AA6D2E37540B5E4520353EC562B32CFFB2313D686638AE8BAC99B8B5F5671E}.
\]
The next compulsory companion audit is V55.

% =============================================================================
% END APPENDED EIGHTH-SERIES V54 MATERIAL
% =============================================================================
% =============================================================================
% BEGIN APPENDED EIGHTH-SERIES V55 MATERIAL
% =============================================================================

\clearpage
\section{Eighth-series V55: compulsory companion audit and the Ward-score route}
\label{sec:v8v55}

\subsection{Files and results inspected}
\label{sec:e8v55-audit}

The compulsory fifth-version audit inspected
\begin{align}
 &\texttt{Renewal\_Yang\_Mills\_Mass\_Gap\_Research\_Notes\_V13.tex},
 \notag\\
 &\hspace{2em}5256\ \text{logical lines},\quad193564\ \text{bytes},
 \quad\operatorname{SHA256}
 =\texttt{4FAE700A2316CADE772844637EA4A006134668DCE78F46E7A628C08009078BB9},
 \label{eq:e8v55-YM-file}\\
 &\texttt{Renewal\_NS\_Stability\_Research\_Notes\_V26.tex},
 \notag\\
 &\hspace{2em}5319\ \text{logical lines},\quad278283\ \text{bytes},
 \quad\operatorname{SHA256}
 =\texttt{82C887F8C2C69E4F28FAD7C3DC8DE5B9099CB5A4BF431EBA1FFF3E91935978AD}.
 \label{eq:e8v55-NS-file}
\end{align}
Each file has exactly one live document terminator.  The NS file is
unchanged since the V50 audit.  The YM programme has advanced from V8
to V13.

\begin{audit}[Yang--Mills V9--V11]
\label{audit:e8v55-YM-score}
V9 differentiates a normalized conditional heat bath and obtains the
exact centered physical score.  The Wilson action gradient is proved to
be a curvature adjoint, so the V8 polar cancellation controls that part;
conditional normalization remains a genuine centered residual.  Haar
integration by parts converts an internal score to a generator
derivative on the test.  V10 then separates source, owner, and ground
conditional projections, proves the exact nested-projection Pythagoras
identity, and identifies the required lower angle only on the active
score space.  V11 derives the moving-ground Poisson equation, proves the
maximal-tree Haar Jacobian is one, and isolates the physical/ground
measure discrepancy.
\end{audit}

\begin{audit}[Yang--Mills V12--V13]
\label{audit:e8v55-YM-gap}
V12 shows that compactness gives only fixed-parameter positivity:
oscillation comparison, min--max, and Doeblin minorization require
constants that may degenerate.  A smooth globally strictly convex
potential is impossible on the closed owner fibre, so the available
Wilson Hessian remains local.  V13 proves an overlap-cover Poincaré
gluing theorem and a three-state obstruction: small bad probability and
a uniform good-core gap do not imply a global gap without interface
conductance or a second local inequality.  The YM programme therefore
goes upstream to block dynamics and an influence matrix.
\end{audit}

\begin{audit}[Navier--Stokes V26]
\label{audit:e8v55-NS}
No new NS result has appeared since V50.  Its terminal source,
concentration, ancestry, and singular-thread closure rows therefore do
not change the present Einstein constraint decision.
\end{audit}

\subsection{Consequences for the SM--Einstein route}
\label{sec:e8v55-consequences}

\begin{theorem}[The good-chart route cannot close CCADC from rarity]
\label{thm:e8v55-rarity}
No regulator-uniform positive alignment or carrier gap can be inferred
solely from:
\begin{enumerate}[label=\textup{(\roman*)},leftmargin=3.5em]
\item an upper bound on the probability of leaving a good
Einstein--matter chart; and
\item a Poincaré or Hessian constant inside that chart.
\end{enumerate}
One must additionally prove interface conductance, overlap mass plus a
second local inequality, or a strict block-influence estimate.
\end{theorem}

\begin{proof}
The audited YM V13 three-state family has fixed bad probability and a
fixed Poincaré constant on its two-state good restriction, while its
global gap tends to zero with the single interface conductance.  That
family is already a finite-dimensional regulated conditional fibre.
Relabeling its good and bad states as chart sectors leaves both
hypotheses unchanged.  Hence no conclusion depending only on the two
listed data can give a positive uniform global gap.
\end{proof}

\begin{corollary}[Application to the conformal chart]
\label{cor:e8v55-conformal}
Even if the V47 conformal relative-mode chart acquired a small
bad-probability estimate and the V48 quadratic form were positive
inside it, those facts would not by themselves construct the global
positive carrier or its conditional capacity.
\end{corollary}

\begin{proof}
Apply \Cref{thm:e8v55-rarity}.  V47 already gives the stronger warning
that a fixed absolute chart can have conditional probability tending to
zero under a translated well; the present corollary shows that reversing
that probability estimate would still not supply interface mixing.
\end{proof}

\begin{theorem}[Audit-selected order of attack]
\label{thm:e8v55-order}
Among the three V54 candidates,
\begin{enumerate}[label=\textup{(\alph*)},leftmargin=3.5em]
\item a finite-regulator Ward derivation of the alignment defect,
\item a boundary trace estimate, and
\item a positive combined carrier,
\end{enumerate}
the Ward derivation is logically first.  It determines the actual
equation, Jacobian, gauge-fixing, ghost, anomaly, and normalization
residuals that the boundary and positivity estimates would have to
control.
\end{theorem}

\begin{proof}
The V54 response estimate applies only after the decomposition
\[
 S_{\rm ADM}-q_{\cal N}
 =-\operatorname{div}Z_{\rm al}
 +R_{\rm EOM}+R_{\rm reg}+S_{\rm ch}
\]
has been derived.  A boundary estimate addresses \(Z_{\rm al}\) and
\(S_{\rm ch}\), while a positive carrier concerns the sign and spectral
properties of a chosen representative.  Neither identifies
\(R_{\rm reg}\).  The audited YM V9--V11 chain obtains its corresponding
residual by first differentiating the normalized finite-regulator
measure and applying Haar integration by parts.  Therefore the
change-of-variables Ward score must be computed before the downstream
two rows can be stated without circular assumptions.
\end{proof}

\subsection{The gravitational Ward-score target}
\label{sec:e8v55-target}

Let \(x\) denote all finite-regulator fields, let \(R_\xi(x)\) be the
infinitesimal regulated diffeomorphism vector field, and let
\(\operatorname{div}_{m_n}R_\xi\) denote its divergence relative to the
chosen reference density \(m_n\).  The target score is
\begin{equation}
 \sigma_{\xi,n}
 :=
 R_\xi S_n-\operatorname{div}_{m_n}R_\xi.
 \label{eq:e8v55-score}
\end{equation}
Gauge fixing, ghosts, fermion measure, chiral phase, and boundary
conditions are part of \(S_n\) and \(m_n\), not informal corrections
added after the identity.

\begin{definition}[Diffeomorphism Ward-score card, DWSC]
\label{def:e8v55-DWSC}
DWSC consists of:
\begin{enumerate}[label=\textup{(W\arabic*)},leftmargin=4em]
\item a finite-regulator field space, reference density, and
infinitesimal diffeomorphism vector field \(R_\xi\);
\item the exact score \(\sigma_{\xi,n}\) in
\eqref{eq:e8v55-score};
\item its decomposition into constraint adjoint, gauge-fixing/ghost,
Jacobian/anomaly, boundary, and normalization pieces;
\item exact integration by parts moving the centered score to a
derivative of the V44 smoothed test;
\item an active-score projection and lower-angle estimate, rather than a
full-space norm demand;
\item uniform contact, current, boundary trace, and charge bounds;
\item restoration of the continuum diffeomorphism Ward identity; and
\item compatibility with Standard Model gauge and chiral Ward identities.
\end{enumerate}
\end{definition}

\begin{firewall}[Centering is not smallness]
\label{fire:e8v55-centering}
Normalized change of variables will imply
\(\langle\sigma_{\xi,n}\rangle=0\).  It does not imply
\(\|\sigma_{\xi,n}\|_2\to0\), a spectral gap, or a strict
active-score angle.  The audited YM chain keeps these as separate rows,
and DWSC does the same.
\end{firewall}

\begin{assessment}[Progress at V55]
\label{ass:e8v55-status}
The audit rejects a tempting return to good-chart rarity and selects an
exact finite-regulator Ward-score ledger as the next upstream step.  It
also imports the active-score lesson: only the component paired with the
observable and constraint response should enter the angle estimate.
This direction connects the classical V53 Noether identity to the V54
quantum contact row without assuming regulator covariance.
\end{assessment}

\begin{programme}[Eighth-series V56: finite-regulator Ward score]
\label{prog:e8v55-V56}
V56 will prove the change-of-variables score identity for a general
finite regulator, including a non-volume-preserving reference density.
It will derive centering, the differentiated-test formula, conditional
normalization, and an active-score projection.  It will then state the
exact additional identification needed for that score to populate the
ADM alignment defect.
\end{programme}

\begin{firewall}[Claim boundary after eighth-series V55]
\label{fire:e8v55-boundary}
V55 records the mandatory companion audit, transfers the V13
rarity/interface obstruction to the Einstein chart problem, and proves
the logical order of the next alignment attack.  It does not yet derive
the gravitational score, identify its anomaly or ADM component, prove
active-score strictness, control boundary traces, construct a positive
carrier, remove cutoffs, or prove the full continuum.
\end{firewall}

\begin{theorem}[Eighth-series V55 result]
\label{thm:e8v55-result}
Good-chart probability and local coercivity cannot supply the missing
global Einstein--matter carrier without an interface estimate.  The
next non-circular alignment object is instead the exact finite-regulator
diffeomorphism score
\[
 \sigma_{\xi,n}
 =R_\xi S_n-\operatorname{div}_{m_n}R_\xi,
\]
whose constraint, Jacobian, gauge, boundary, normalization, and anomaly
pieces must be separated before any inverse response or positivity
argument is used.
\end{theorem}

\begin{proof}
Use \Cref{thm:e8v55-rarity,thm:e8v55-order}.
\end{proof}

\paragraph{Parent-version record.}
V55 was formed from
\texttt{Renewal\_SM\_Einstein\_Continuum\_Research\_Notes\_V54.tex}.
The V54 parent had (22\,195) logical source lines, (794\,162) bytes,
and SHA--256
\[
 \texttt{05F4C614CFA684CC5B39AE7B323A0986BC1746FB1BB51A95377B073551E4E0EE}.
\]
The next compulsory companion audit is V60.

% =============================================================================
% END APPENDED EIGHTH-SERIES V55 MATERIAL
% =============================================================================
% =============================================================================
% BEGIN APPENDED EIGHTH-SERIES V56 MATERIAL
% =============================================================================

\clearpage
\section{Eighth-series V56: finite-regulator Ward scores and active observables}
\label{sec:v8v56}

\subsection{Change of variables with a nontrivial reference density}
\label{sec:e8v56-Ward}

Let \(M\) be a finite-dimensional oriented regulated field manifold with
a smooth positive density \(dm\).  For a \(C^1\) vector field \(R\),
define its density divergence by
\begin{equation}
 {\cal L}_Rdm=(\operatorname{div}_mR)\,dm.
 \label{eq:e8v56-divergence}
\end{equation}
Let
\begin{equation}
 d\mu=Z^{-1}e^{-S}\,dm,
 \qquad
 Z=\int_Me^{-S}\,dm.
 \label{eq:e8v56-measure}
\end{equation}
The boundaryless case includes compact lattice gauge fibres.  On a
manifold with boundary, write \(d\iota_Rm\) for the induced oriented
flux density.

\begin{theorem}[Exact finite-regulator Ward-score identity]
\label{thm:e8v56-Ward}
Define
\begin{equation}
 \sigma_R:=RS-\operatorname{div}_mR.
 \label{eq:e8v56-score}
\end{equation}
If \(F\) is \(C^1\) and the terms are integrable, then
\begin{equation}
 \langle F\sigma_R\rangle_\mu
 =
 \langle RF\rangle_\mu
 -
 {1\over Z}\int_{\partial M}Fe^{-S}\,d\iota_Rm.
 \label{eq:e8v56-Ward-boundary}
\end{equation}
If \(M\) has no boundary, \(R\) is tangent to the boundary, or the flux
vanishes, then
\begin{equation}
 \langle F\sigma_R\rangle_\mu=\langle RF\rangle_\mu,
 \qquad
 \langle\sigma_R\rangle_\mu=0.
 \label{eq:e8v56-Ward-closed}
\end{equation}
\end{theorem}

\begin{proof}
The product and density-divergence rules give
\[
 \operatorname{div}_m(Fe^{-S}R)
 =
 e^{-S}\bigl(RF-FRS+F\operatorname{div}_mR\bigr)
 =
 e^{-S}(RF-F\sigma_R).
\]
Integrate over \(M\) and apply the divergence theorem.  Rearranging and
dividing by \(Z\) proves
\eqref{eq:e8v56-Ward-boundary}.  Under zero boundary flux it reduces to
the first identity in \eqref{eq:e8v56-Ward-closed}; take \(F=1\) for
centering.
\end{proof}

\begin{corollary}[V54 Schwinger--Dyson identity as a Ward score]
\label{cor:e8v56-SD}
For \(M=\mathbb R^N\), Lebesgue reference density, and
\[
 R=\sum_aA^a(x){\partial\over\partial x^a},
\]
one has
\begin{equation}
 \sigma_R
 =
 \sum_a\left(
 A^a{\partial S\over\partial x^a}
 -{\partial A^a\over\partial x^a}
 \right),
 \qquad
 \langle F\sigma_R\rangle
 =\sum_a\langle A^a\partial_aF\rangle.
 \label{eq:e8v56-SD}
\end{equation}
\end{corollary}

\begin{proof}
For Lebesgue density,
\(\operatorname{div}_mR=\sum_a\partial_aA^a\).
Use \Cref{thm:e8v56-Ward} and the assumed vanishing flux at infinity.
\end{proof}

\begin{firewall}[The Jacobian is part of the score]
\label{fire:e8v56-Jacobian}
The action derivative \(RS\) is not the complete Ward source unless the
chosen reference density is invariant under \(R\).  The term
\(-\operatorname{div}_mR\) is the infinitesimal Jacobian.  Gauge fixing,
coordinate charts, fermion determinants, and a chiral regulator can
contribute here.  Dropping it would assume away the regulator anomaly
that DWSC is required to measure.
\end{firewall}

\subsection{Moving normalized measures and conditional centering}
\label{sec:e8v56-moving}

Let \(S_\theta\) and \(dm_\theta\) be differentiable in a parameter
\(\theta\).  Relative to a fixed local reference density \(dm_*\), write
\begin{equation}
 dm_\theta=w_\theta\,dm_*,
 \qquad
 \tau_\theta
 :=
 \partial_\theta S_\theta-\partial_\theta\log w_\theta.
 \label{eq:e8v56-tau}
\end{equation}

\begin{theorem}[Derivative of a normalized regulated expectation]
\label{thm:e8v56-moving}
For
\[
 d\mu_\theta=Z_\theta^{-1}e^{-S_\theta}dm_\theta,
\]
and a differentiable test \(F_\theta\),
\begin{equation}
 \partial_\theta\langle F_\theta\rangle_{\mu_\theta}
 =
 \langle\partial_\theta F_\theta\rangle_{\mu_\theta}
 -
 \left\langle
 \bigl(F_\theta-\langle F_\theta\rangle_{\mu_\theta}\bigr)
 \tau_\theta
 \right\rangle_{\mu_\theta}.
 \label{eq:e8v56-moving}
\end{equation}
Equivalently, the normalized parameter score
\begin{equation}
 \widehat\tau_\theta
 :=
 \tau_\theta-\langle\tau_\theta\rangle_{\mu_\theta}
 \label{eq:e8v56-centered-tau}
\end{equation}
satisfies
\begin{equation}
 \partial_\theta\langle F_\theta\rangle_{\mu_\theta}
 =
 \langle\partial_\theta F_\theta\rangle_{\mu_\theta}
 -\langle F_\theta\widehat\tau_\theta\rangle_{\mu_\theta}.
 \label{eq:e8v56-moving-centered}
\end{equation}
\end{theorem}

\begin{proof}
The logarithmic derivative of the unnormalized density is
\(-\tau_\theta\).  Thus
\[
 \partial_\theta Z_\theta
 =-Z_\theta\langle\tau_\theta\rangle_{\mu_\theta}.
\]
Differentiate
\[
 Z_\theta^{-1}
 \int F_\theta e^{-S_\theta}w_\theta\,dm_*
\]
and insert the displayed derivative of \(Z_\theta\).  Collecting terms
gives \eqref{eq:e8v56-moving}; subtracting the mean of \(\tau_\theta\)
gives \eqref{eq:e8v56-moving-centered}.
\end{proof}

\begin{corollary}[Conditional normalized score]
\label{cor:e8v56-conditional}
Let \(x\) be an active block and \(y\) the exterior field, with
conditional law
\[
 d\mu_\theta(x\mid y)
 =
 Z_\theta(y)^{-1}e^{-S_\theta(x,y)}
 w_\theta(x,y)\,dm_*(x).
\]
For almost every fixed \(y\), define
\begin{equation}
 \widehat\tau_\theta(x\mid y)
 :=
 \tau_\theta(x,y)
 -\mathbb E_\theta[\tau_\theta\mid y].
 \label{eq:e8v56-conditional-score}
\end{equation}
Then
\begin{equation}
 \mathbb E_\theta[\widehat\tau_\theta\mid y]=0
 \label{eq:e8v56-conditional-centered}
\end{equation}
and
\begin{equation}
 \partial_\theta\mathbb E_\theta[F_\theta\mid y]
 =
 \mathbb E_\theta[\partial_\theta F_\theta\mid y]
 -
 \mathbb E_\theta[F_\theta\widehat\tau_\theta\mid y].
 \label{eq:e8v56-conditional-derivative}
\end{equation}
\end{corollary}

\begin{proof}
Apply \Cref{thm:e8v56-moving} to the normalized \(x\)-integral with
\(y\) fixed.  Equation \eqref{eq:e8v56-conditional-centered} follows
from the definition.
\end{proof}

\begin{remark}[Extension beyond the audited YM formula]
\label{rem:e8v56-extension}
The audited YM V9 calculation supplies the conditional-score template
for a Haar fibre.  V56 retains a moving reference density, its Jacobian,
an explicit field-space boundary flux, and the ADM-alignment interface.
These are the pieces needed for a gravitational regulator and are not
provided by conditional centering alone.
\end{remark}

\subsection{Only the observable-active score is visible}
\label{sec:e8v56-active}

Let \({\cal H}=L^2_0(\mu)\) be the centered real \(L^2\) space, and let
\({\cal T}\subset{\cal H}\) be the closed span of the centered,
spacetime-smoothed tests admitted by V44.  Denote its orthogonal
projection by \(P_{\cal T}\).

\begin{theorem}[Exact active-score dual norm]
\label{thm:e8v56-active}
For every centered score \(\sigma\in{\cal H}\),
\begin{equation}
 \sup_{\substack{F\in{\cal T}\\\|F\|_2\leq1}}
 |\langle F\sigma\rangle_\mu|
 =
 \|P_{\cal T}\sigma\|_2.
 \label{eq:e8v56-active}
\end{equation}
If
\(\sigma=\sigma_{\rm adj}+\sigma_{\rm rem}\), then
\begin{equation}
 \|P_{\cal T}\sigma\|_2
 \leq
 \|P_{\cal T}\sigma_{\rm adj}\|_2
 +\|P_{\cal T}\sigma_{\rm rem}\|_2.
 \label{eq:e8v56-active-split}
\end{equation}
No norm of the component in \({\cal T}^\perp\) enters either formula.
\end{theorem}

\begin{proof}
For \(F\in{\cal T}\),
\(\langle F,\sigma\rangle=\langle F,P_{\cal T}\sigma\rangle\).
Cauchy--Schwarz gives the upper bound in
\eqref{eq:e8v56-active}.  If \(P_{\cal T}\sigma\neq0\), equality is
attained by
\(F=P_{\cal T}\sigma/\|P_{\cal T}\sigma\|_2\); otherwise both sides
vanish.  The triangle inequality proves
\eqref{eq:e8v56-active-split}.
\end{proof}

\begin{corollary}[Ward control by differentiated tests]
\label{cor:e8v56-test-derivative}
Under the zero-flux hypotheses of \Cref{thm:e8v56-Ward},
\begin{equation}
 \|P_{\cal T}\sigma_R\|_2
 =
 \sup_{\substack{F\in{\cal T}\\\|F\|_2\leq1}}
 |\langle RF\rangle_\mu|.
 \label{eq:e8v56-Ward-active}
\end{equation}
Thus an \(L^2\) operator bound
\begin{equation}
 |\langle RF\rangle_\mu|
 \leq C_R\|F\|_2,
 \qquad F\in{\cal T},
 \label{eq:e8v56-test-bound}
\end{equation}
controls the entire observable-active Ward score with the same constant,
even if the full score norm is large.
\end{corollary}

\begin{proof}
Combine \eqref{eq:e8v56-Ward-closed} with
\Cref{thm:e8v56-active}.
\end{proof}

\begin{firewall}[The full score may be large]
\label{fire:e8v56-full-score}
Centering and the active projection do not bound
\(\|\sigma_R\|_2\).  The identity only removes the component that no
admitted observable can detect.  To use a larger observable algebra one
must enlarge \({\cal T}\) and reprove
\eqref{eq:e8v56-test-bound}; the projection cannot be fixed after seeing
the source.
\end{firewall}

\subsection{From a Ward score to a local ADM source}
\label{sec:e8v56-ADM}

At a finite regulator, let \(\Xi_n\) be the space of local
diffeomorphism parameters.  Assume the transformation is linear in the
parameter,
\[
 R_\xi=\sum_{\alpha}\xi^\alpha R_\alpha.
\]
Then the score is linear:
\begin{equation}
 \sigma_{\xi,n}
 =\sum_\alpha\xi^\alpha\sigma_{\alpha,n},
 \qquad
 \sigma_{\alpha,n}
 :=R_\alpha S_n-\operatorname{div}_{m_n}R_\alpha.
 \label{eq:e8v56-local-score}
\end{equation}
This coefficient ledger is defined without an elliptic inverse.

\begin{proposition}[Finite-regulator coefficient Ward identity]
\label{prop:e8v56-coefficient}
Under zero field-space boundary flux, every local parameter coefficient
satisfies
\begin{equation}
 \langle F\sigma_{\alpha,n}\rangle_{\mu_n}
 =
 \langle R_\alpha F\rangle_{\mu_n}.
 \label{eq:e8v56-coefficient-Ward}
\end{equation}
If the field-space boundary flux is nonzero, its coefficient must be
subtracted on the right-hand side exactly as in
\eqref{eq:e8v56-Ward-boundary}.
\end{proposition}

\begin{proof}
Apply \Cref{thm:e8v56-Ward} with \(R=R_\alpha\).  Linearity in \(\xi\)
permits coefficient comparison.
\end{proof}

\begin{definition}[ADM Ward-alignment identity]
\label{def:e8v56-AWAI}
AWAI is the regulator-local identity
\begin{equation}
 S_{{\rm ADM},n}-\sigma_{\perp,n}
 =
 -\operatorname{div}Z_{{\rm W},n}
 +R_{{\rm gf},n}+R_{{\rm gh},n}
 +R_{{\rm an},n}+S_{{\rm ch},n},
 \label{eq:e8v56-AWAI}
\end{equation}
where \(\sigma_{\perp,n}\) is the coefficient of the normal
diffeomorphism parameter, \(R_{\rm gf}\) and \(R_{\rm gh}\) are
gauge-fixing and ghost terms, \(R_{\rm an}\) is the measure or anomaly
defect, and \(S_{\rm ch}\) is the retained charge sector.
\end{definition}

\begin{theorem}[What AWAI would give]
\label{thm:e8v56-AWAI}
If \eqref{eq:e8v56-AWAI} holds, then after charge projection
\begin{equation}
 \left\|
 \nabla(-\Delta)^\dagger P_0
 (S_{{\rm ADM},n}-\sigma_{\perp,n}-S_{{\rm ch},n})
 \right\|_2
 \leq
 \|Z_{{\rm W},n}\|_2
 +\lambda_{1,n}^{-1/2}
 \|P_0(R_{{\rm gf},n}+R_{{\rm gh},n}+R_{{\rm an},n})\|_2.
 \label{eq:e8v56-AWAI-bound}
\end{equation}
Against admitted tests, the score coefficient is controlled by
\eqref{eq:e8v56-coefficient-Ward} and the active norm
\eqref{eq:e8v56-active}.
\end{theorem}

\begin{proof}
Apply \Cref{thm:e8v50-constraint-adjoint} to
\eqref{eq:e8v56-AWAI}.  The test statement follows from
\Cref{prop:e8v56-coefficient,thm:e8v56-active}.
\end{proof}

\begin{firewall}[AWAI is not supplied by linear algebra]
\label{fire:e8v56-AWAI}
The coefficient ledger \eqref{eq:e8v56-local-score} exists for every
finite parameter space, but it does not identify
\(\sigma_{\perp,n}\) with the ADM scalar constraint.  That equality is
a Noether or canonical statement about the concrete regulator.  Nor
does centering make \(R_{\rm an}\) vanish.  AWAI must be derived from
the regulated action, reference measure, gauge fixing, and boundary
terms.
\end{firewall}

\subsection{Ward-score assessment}
\label{sec:e8v56-assessment}

\begin{assessment}[Progress at V56]
\label{ass:e8v56-status}
The gravitational Ward target now has an exact finite-regulator
identity.  It includes the reference-density Jacobian, conditional
normalization, and field-space boundary flux.  The full residual is
replaced by its projection onto the fixed V44 observable class, with an
exact dual characterization through differentiated tests.  The
remaining substantive step is AWAI: identifying the normal
diffeomorphism score with the ADM scalar source and exposing its
gauge-fixing, ghost, anomaly, boundary, and charge defects.
\end{assessment}

\begin{programme}[Eighth-series V57: canonical normal-deformation coefficient]
\label{prog:e8v56-V57}
V57 will derive the normal-deformation score for a finite-dimensional
canonical constrained action.  It will show exactly how the lapse
derivative produces the Hamiltonian constraint, how lapse-independent
reference measures avoid a normalization Jacobian, and where
gauge-fixing or regulator dependence enters AWAI.
\end{programme}

\begin{firewall}[Claim boundary after eighth-series V56]
\label{fire:e8v56-boundary}
V56 proves the finite-regulator Ward-score identity, moving and
conditional normalized-score formulas, the exact active-score dual
norm, and the conditional consequence of AWAI.  It does not derive AWAI
for a concrete Einstein--Standard Model regulator, treat Berezin
integration and chiral anomalies, prove active-score uniformity, close
boundary or charge sectors, construct a positive carrier, remove
cutoffs, or prove the full continuum.
\end{firewall}

\begin{theorem}[Eighth-series V56 result]
\label{thm:e8v56-result}
For a finite regulator with reference density \(dm_n\), the exact
diffeomorphism score is
\[
 \sigma_{\xi,n}
 =R_\xi S_n-\operatorname{div}_{m_n}R_\xi,
\]
and, up to the explicit field-space boundary flux,
\[
 \langle F\sigma_{\xi,n}\rangle
 =\langle R_\xi F\rangle.
\]
Only its orthogonal projection onto the fixed smoothed-test space enters
these Ward pairings.  Relating its normal-deformation coefficient to the
ADM scalar source requires the separate local identity
\eqref{eq:e8v56-AWAI}.
\end{theorem}

\begin{proof}
Use \Cref{thm:e8v56-Ward,thm:e8v56-active,
prop:e8v56-coefficient}.
\end{proof}

\paragraph{Parent-version record.}
V56 was formed from
\texttt{Renewal\_SM\_Einstein\_Continuum\_Research\_Notes\_V55.tex}.
The V55 parent had (22\,440) logical source lines, (804\,549) bytes,
and SHA--256
\[
 \texttt{93D56FE924930BE60C9DE08A8320D3D86B57DBE29D7894AC6E4668B29924D93A}.
\]
The next compulsory companion audit is V60.

% =============================================================================
% END APPENDED EIGHTH-SERIES V56 MATERIAL
% =============================================================================
% =============================================================================
% BEGIN APPENDED EIGHTH-SERIES V57 MATERIAL
% =============================================================================

\clearpage
\section{Eighth-series V57: lapse-multiplier scores versus diffeomorphism Ward scores}
\label{sec:v8v57}

\subsection{The canonical multiplier ledger}
\label{sec:e8v57-multiplier}

The normal diffeomorphism Ward score of V56 and the ADM Hamiltonian
constraint are related by canonical gauge structure, but they are not
the same finite-regulator object.  The constraint first appears as the
coefficient of the lapse multiplier.

Let \(x\in M_n\) collect the finite-regulator canonical metric, matter,
gauge, and auxiliary fields.  Let \(N=(N^\alpha)\) denote fixed lapse
and shift multipliers.  Suppose
\begin{equation}
 S_n(x;N)
 =
 S_{0,n}(x)
 +\sum_\alpha N^\alpha{\cal C}_{\alpha,n}(x)
 +S_{{\rm gf},n}(x;N)
 +S_{{\rm bd},n}(x;N),
 \label{eq:e8v57-action}
\end{equation}
and let
\begin{equation}
 d\mu_{n,N}(x)
 =
 Z_n(N)^{-1}e^{-S_n(x;N)}
 w_n(x;N)\,dm_{*,n}(x).
 \label{eq:e8v57-measure}
\end{equation}
The measure is assumed normalizable for the fixed allowed multipliers
under discussion.

\begin{definition}[Effective multiplier score]
\label{def:e8v57-score}
Set
\begin{align}
 G_{\alpha,n}
 &:=
 \partial_{N^\alpha}S_{{\rm gf},n},
 &
 B_{\alpha,n}
 &:=
 \partial_{N^\alpha}S_{{\rm bd},n},
 \label{eq:e8v57-GB}\\
 J_{\alpha,n}
 &:=
 \partial_{N^\alpha}\log w_n,
 &
 \tau_{\alpha,n}
 &:=
 {\cal C}_{\alpha,n}
 +G_{\alpha,n}+B_{\alpha,n}-J_{\alpha,n}.
 \label{eq:e8v57-tau}
\end{align}
Its normalized version is
\begin{equation}
 \widehat\tau_{\alpha,n}
 :=
 \tau_{\alpha,n}
 -\langle\tau_{\alpha,n}\rangle_{n,N}.
 \label{eq:e8v57-centered}
\end{equation}
\end{definition}

\begin{theorem}[Exact lapse-response identity]
\label{thm:e8v57-lapse}
For every differentiable regulated test \(F_N\),
\begin{equation}
 \partial_{N^\alpha}\langle F_N\rangle_{n,N}
 =
 \langle\partial_{N^\alpha}F_N\rangle_{n,N}
 -
 \langle F_N\widehat\tau_{\alpha,n}\rangle_{n,N}.
 \label{eq:e8v57-lapse-response}
\end{equation}
The bare canonical constraint has the exact decomposition
\begin{equation}
 {\cal C}_{\alpha,n}
 =
 \widehat\tau_{\alpha,n}
 -G_{\alpha,n}-B_{\alpha,n}+J_{\alpha,n}
 +\langle\tau_{\alpha,n}\rangle_{n,N}.
 \label{eq:e8v57-constraint-ledger}
\end{equation}
\end{theorem}

\begin{proof}
Differentiate \eqref{eq:e8v57-action}:
\[
 \partial_{N^\alpha}S_n
 =
 {\cal C}_{\alpha,n}+G_{\alpha,n}+B_{\alpha,n}.
\]
Since
\(\partial_{N^\alpha}\log w_n=J_{\alpha,n}\), the parameter score
\(\partial_{N^\alpha}S_n-\partial_{N^\alpha}\log w_n\) is exactly
\(\tau_{\alpha,n}\).  Apply
\Cref{thm:e8v56-moving} to obtain
\eqref{eq:e8v57-lapse-response}.  Solving
\eqref{eq:e8v57-tau} for \({\cal C}_{\alpha,n}\) and inserting
\(\tau=\widehat\tau+\langle\tau\rangle\) proves
\eqref{eq:e8v57-constraint-ledger}.
\end{proof}

\begin{corollary}[Lapse-independent gauge, boundary, and measure]
\label{cor:e8v57-clean}
If \(G_{\alpha,n}=B_{\alpha,n}=J_{\alpha,n}=0\), then
\begin{equation}
 {\cal C}_{\alpha,n}
 =
 \widehat{\cal C}_{\alpha,n}
 +\langle{\cal C}_{\alpha,n}\rangle,
 \qquad
 \widehat{\cal C}_{\alpha,n}
 :={\cal C}_{\alpha,n}
 -\langle{\cal C}_{\alpha,n}\rangle.
 \label{eq:e8v57-clean}
\end{equation}
For a lapse-independent test,
\begin{equation}
 \langle F\widehat{\cal C}_{\alpha,n}\rangle
 =
 -\partial_{N^\alpha}\langle F\rangle.
 \label{eq:e8v57-clean-response}
\end{equation}
\end{corollary}

\begin{proof}
Specialize \eqref{eq:e8v57-constraint-ledger} and
\eqref{eq:e8v57-lapse-response}.
\end{proof}

\begin{firewall}[The mean constraint is a separate row]
\label{fire:e8v57-mean}
Normalization controls only the centered multiplier score.  The number
\(\langle\tau_{\alpha,n}\rangle\) in
\eqref{eq:e8v57-constraint-ledger} is not small by centering.  For a
local lapse coefficient it is a one-point constraint defect; for the
constant lapse it contributes to the integrated energy or charge row.
It must be set by the constraint measure, a boundary equation, a
renormalization condition, or a physical-state condition.
\end{firewall}

\subsection{Why the Ward score is not the constraint}
\label{sec:e8v57-distinction}

Let \(R_\xi\) be an infinitesimal field redefinition representing a
regulated diffeomorphism, with V56 Ward score
\begin{equation}
 \sigma_{\xi,n}
 =
 R_\xi S_n-\operatorname{div}_{m_n}R_\xi.
 \label{eq:e8v57-Ward-score}
\end{equation}
By contrast, the multiplier score is a parameter derivative at fixed
fields.  The following elementary model rules out their naive
identification.

\begin{proposition}[Invariant Ward score with nonzero multiplier coefficient]
\label{prop:e8v57-obstruction}
Let \(M=\mathbb R^2\) with Lebesgue density,
\[
 S_N(x)
 ={1\over2}|x|^2+N\bigl(|x|^2-1\bigr),
 \qquad N>-{1\over2},
\]
and let
\[
 R=x^1{\partial\over\partial x^2}
 -x^2{\partial\over\partial x^1}
\]
be the rotation vector field.  Then
\begin{equation}
 \sigma_R=0
 \label{eq:e8v57-zero-Ward}
\end{equation}
pointwise, whereas
\begin{equation}
 \partial_NS_N=|x|^2-1
 \label{eq:e8v57-nonzero-constraint}
\end{equation}
is not the zero function.
\end{proposition}

\begin{proof}
Both \(|x|^2\) and the Lebesgue density are rotation invariant, so
\(RS_N=0\) and \(\operatorname{div}R=0\).  This proves
\eqref{eq:e8v57-zero-Ward}.  Direct differentiation in \(N\) gives
\eqref{eq:e8v57-nonzero-constraint}.
\end{proof}

\begin{corollary}[Correction to the naive AWAI target]
\label{cor:e8v57-AWAI-correction}
There is no general identity equating the Hamiltonian constraint source
to the normal-diffeomorphism Ward score merely because both arise from
diffeomorphism symmetry.  The V56 AWAI definition can hold only after
its defect ledger includes the canonical multiplier coefficient; as a
route for deriving that coefficient from the Ward score alone, it is
insufficient.
\end{corollary}

\begin{proof}
\Cref{prop:e8v57-obstruction} has an exact invariant measure and zero
Ward score but a nonzero multiplier coefficient.  Hence symmetry and
centering alone cannot imply the proposed equality.
\end{proof}

\begin{firewall}[Two distinct derivatives]
\label{fire:e8v57-two-derivatives}
The field-space derivative \(R_\xi S\) tests invariance under a change
of integration variables.  The parameter derivative
\(\partial_NS\) inserts the constraint conjugate to the lapse.  A
canonical BRST or Noether identity can relate their correlation
functions, but replacing one by the other without deriving that identity
is circular.
\end{firewall}

\subsection{The observable-active multiplier norm}
\label{sec:e8v57-active}

Let \({\cal T}_n\subset L^2_0(\mu_{n,N})\) be the fixed centered test
space.  For each local multiplier component, define its susceptibility
functional
\begin{equation}
 {\cal D}_{\alpha,n}(F)
 :=
 -\partial_{N^\alpha}\langle F\rangle_{n,N}
 +\langle\partial_{N^\alpha}F\rangle_{n,N}.
 \label{eq:e8v57-susceptibility}
\end{equation}

\begin{theorem}[Active multiplier score equals lapse susceptibility]
\label{thm:e8v57-active}
One has
\begin{equation}
 \|P_{{\cal T}_n}\widehat\tau_{\alpha,n}\|_2
 =
 \sup_{\substack{F\in{\cal T}_n\\\|F\|_2\leq1}}
 |{\cal D}_{\alpha,n}(F)|.
 \label{eq:e8v57-active}
\end{equation}
If the gauge, boundary, and measure terms are separated as in
\eqref{eq:e8v57-constraint-ledger}, then
\begin{align}
 \|P_{{\cal T}_n}
 ({\cal C}_{\alpha,n}
 -\langle\tau_{\alpha,n}\rangle)\|_2
 &\leq
 \|P_{{\cal T}_n}\widehat\tau_{\alpha,n}\|_2
 \notag\\
 &\quad+
 \|P_{{\cal T}_n}
 (-G_{\alpha,n}-B_{\alpha,n}+J_{\alpha,n})\|_2.
 \label{eq:e8v57-active-ledger}
\end{align}
\end{theorem}

\begin{proof}
Equation \eqref{eq:e8v57-lapse-response} gives
\[
 {\cal D}_{\alpha,n}(F)
 =\langle F\widehat\tau_{\alpha,n}\rangle.
\]
Apply \Cref{thm:e8v56-active}.  Project
\eqref{eq:e8v57-constraint-ledger} onto \({\cal T}_n\), note that
centered tests are orthogonal to the scalar mean, and use the triangle
inequality.
\end{proof}

\begin{corollary}[A measurable target for the constraint insertion]
\label{cor:e8v57-measurable}
Uniform differentiability of the admitted observable expectations in
the local lapse,
\begin{equation}
 \sup_{n,\alpha}
 \sup_{\substack{F\in{\cal T}_n\\\|F\|_2\leq1}}
 |{\cal D}_{\alpha,n}(F)|<\infty,
 \label{eq:e8v57-uniform-susceptibility}
\end{equation}
is exactly a uniform bound on the observable-active normalized
constraint score, before gauge, boundary, and measure defects are added.
\end{corollary}

\begin{proof}
This is \eqref{eq:e8v57-active}.
\end{proof}

\subsection{Constraint-source ledger after the correction}
\label{sec:e8v57-ledger}

\begin{definition}[Canonical lapse-score card, CLSC]
\label{def:e8v57-CLSC}
CLSC consists of:
\begin{enumerate}[label=\textup{(L\arabic*)},leftmargin=4em]
\item a regulated canonical action with explicit lapse and shift
coefficients \({\cal C}_{\alpha,n}\);
\item a normalizable fixed-multiplier measure and its reference-density
derivative \(J_{\alpha,n}\);
\item the exact ledger \eqref{eq:e8v57-constraint-ledger};
\item a physical equation for the mean constraint or integrated charge;
\item the active susceptibility estimate
\eqref{eq:e8v57-uniform-susceptibility};
\item a separate diffeomorphism Ward identity controlling regulator
symmetry, gauge fixing, ghosts, and anomalies;
\item a canonical or BRST compatibility identity relating the two
ledgers without identifying them; and
\item the V50 spatial-adjoint factorization, boundary, positivity,
reflection, and cutoff-removal rows.
\end{enumerate}
\end{definition}

\begin{assessment}[Progress at V57]
\label{ass:e8v57-status}
The scalar constraint insertion is now derived from the correct
finite-regulator object: the normalized lapse-multiplier score.  Its
gauge-fixing, boundary, measure-Jacobian, and mean pieces are explicit,
and its observable-active norm is exactly the lapse susceptibility.
An explicit invariant model proves that the diffeomorphism Ward score
cannot replace this multiplier score.  The two ledgers must be connected
by a canonical or BRST identity, not equated.
\end{assessment}

\begin{programme}[Eighth-series V58: multiplier integration and constraint enforcement]
\label{prog:e8v57-V58}
V58 will analyze how integrating or averaging over a regulated lapse
enforces the constraint.  It will distinguish Fourier delta
enforcement, positive Laplace weights, and Gaussian penalty
regularization, and will quantify the residual constraint and positivity
cost of each choice before selecting a carrier.
\end{programme}

\begin{firewall}[Claim boundary after eighth-series V57]
\label{fire:e8v57-boundary}
V57 proves the exact lapse-score ledger, the Ward-score obstruction, and
the active lapse-susceptibility identity.  It does not construct a
normalizable gravitational lapse integral, impose the quantum
Hamiltonian constraint, relate the Ward and multiplier ledgers through
BRST, prove the V50 spatial-divergence row for the constraint insertion,
construct a positive carrier, remove cutoffs, or prove the full
continuum.
\end{firewall}

\begin{theorem}[Eighth-series V57 result]
\label{thm:e8v57-result}
For a regulated canonical action linear in the lapse before gauge and
boundary corrections, the Hamiltonian constraint obeys
\[
 {\cal C}_{\alpha,n}
 =
 \widehat\tau_{\alpha,n}
 -G_{\alpha,n}-B_{\alpha,n}+J_{\alpha,n}
 +\langle\tau_{\alpha,n}\rangle.
\]
Its active normalized part is exactly the local lapse susceptibility.
The diffeomorphism change-of-variables score is a distinct object and
can vanish identically while the multiplier coefficient is nonzero.
\end{theorem}

\begin{proof}
Use \Cref{thm:e8v57-lapse,prop:e8v57-obstruction,
thm:e8v57-active}.
\end{proof}

\paragraph{Parent-version record.}
V57 was formed from
\texttt{Renewal\_SM\_Einstein\_Continuum\_Research\_Notes\_V56.tex}.
The V56 parent had (22\,905) logical source lines, (819\,071) bytes,
and SHA--256
\[
 \texttt{8F8BEA7D3E8A11DE83DC2F17BF7F089D25EA621F91F007074764B0A17152B98C}.
\]
The next compulsory companion audit is V60.

% =============================================================================
% END APPENDED EIGHTH-SERIES V57 MATERIAL
% =============================================================================
% =============================================================================
% BEGIN APPENDED EIGHTH-SERIES V58 MATERIAL
% =============================================================================

\clearpage
\section{Eighth-series V58: regulated lapse integration and positive constraint enforcement}
\label{sec:v8v58}

\subsection{Three multiplier prescriptions}
\label{sec:e8v58-prescriptions}

V57 identifies the Hamiltonian constraint as the lapse-multiplier
coefficient.  It does not follow that integrating the lapse produces a
positive measure concentrated on that constraint.  This section compares
the elementary prescriptions before they are inserted in the
Einstein--Standard Model carrier.

\begin{proposition}[Finite Fourier enforcement is signed]
\label{prop:e8v58-Fourier}
For \(L>0\),
\begin{equation}
 K_L(c):={1\over2\pi}\int_{-L}^Le^{iNc}\,dN
 =
 \begin{cases}
 \dfrac{\sin(Lc)}{\pi c},&c\neq0,\\
 \dfrac L\pi,&c=0.
 \end{cases}
 \label{eq:e8v58-Dirichlet}
\end{equation}
The kernels \(K_L\) converge to \(\delta_0\) in the usual distributional
sense on Schwartz tests, but \(K_L\) changes sign and is not a positive
probability density.
\end{proposition}

\begin{proof}
Direct integration gives \eqref{eq:e8v58-Dirichlet}.  The function
\(\sin(Lc)\) alternates sign on successive half-periods.  Distributional
convergence is Fourier inversion applied to the expanding indicator
\(\mathbf1_{[-L,L]}\).
\end{proof}

\begin{proposition}[A real Euclidean lapse does not enforce a
two-sided constraint]
\label{prop:e8v58-Laplace}
For \(L>0\),
\begin{equation}
 \int_{-L}^Le^{-Nc}\,dN
 =
 \begin{cases}
 \dfrac{2\sinh(Lc)}c,&c\neq0,\\
 2L,&c=0.
 \end{cases}
 \label{eq:e8v58-Laplace}
\end{equation}
This positive factor is increasing as a function of \(|c|\) for
\(|c|>0\), so it favors rather than suppresses large two-sided
constraint violations.  Moreover,
\(\int_{\mathbb R}e^{-Nc}dN\) diverges for every real \(c\).
\end{proposition}

\begin{proof}
The integral is elementary.  For \(x>0\),
\[
 {d\over dx}{\sinh(Lx)\over x}
 ={Lx\cosh(Lx)-\sinh(Lx)\over x^2}>0,
\]
because the numerator vanishes at zero and has derivative
\(L^2x\sinh(Lx)>0\).  On the full real lapse line, one of the two tails
grows exponentially when \(c\neq0\), while both tails have infinite
length when \(c=0\).
\end{proof}

\begin{firewall}[Wick rotation of the lapse is substantive]
\label{fire:e8v58-Wick}
The oscillatory Fourier multiplier enforces the constraint
distributionally but is not a positive carrier.  The real Euclidean
multiplier is positive but fails to enforce a sign-indefinite
constraint.  Passing to an imaginary lapse contour or replacing it by
a penalty changes the integration cycle or the action; it must be
declared and tested for reflection, locality, and continuum equivalence.
\end{firewall}

\subsection{Gaussian lapse averaging}
\label{sec:e8v58-Gaussian}

\begin{theorem}[Gaussian multiplier produces a positive penalty]
\label{thm:e8v58-Gaussian}
For \(c\in\mathbb R^r\) and \(\kappa>0\),
\begin{equation}
 (2\pi\kappa)^{-r/2}
 \int_{\mathbb R^r}
 \exp\!\left(-{|N|^2\over2\kappa}+iN\cdot c\right)dN
 =
 \exp\!\left(-{\kappa\over2}|c|^2\right).
 \label{eq:e8v58-Gaussian-Fourier}
\end{equation}
Thus an oscillatory Gaussian lapse integral can be performed exactly,
leaving a positive constraint-squared weight.
\end{theorem}

\begin{proof}
The integral factors into \(r\) one-dimensional characteristic
functions of centered Gaussians with variance \(\kappa\):
\[
 (2\pi\kappa)^{-1/2}
 \int_{\mathbb R}
 e^{-N^2/(2\kappa)+iNc}\,dN
 =e^{-\kappa c^2/2}.
\]
Multiplication proves \eqref{eq:e8v58-Gaussian-Fourier}.
\end{proof}

\begin{remark}[Locality of the penalty]
\label{rem:e8v58-local}
If a separate multiplier \(N_\alpha\) is assigned to each local
constraint coefficient \(c_\alpha\) with product Gaussian law, then
integrating all multipliers gives
\[
 \exp\!\left(-{\kappa\over2}\sum_\alpha c_\alpha^2\right).
\]
This is regulator-local whenever each \(c_\alpha\) is local.  Its
coupling \(\kappa\) is nevertheless a new stiffness scale.
\end{remark}

\subsection{A quantitative penalty residual}
\label{sec:e8v58-residual}

The following estimate avoids assuming a spectral gap for the full
field measure.  It uses only the marginal law of the constraint
coordinate.

\begin{theorem}[Constraint residual under a Gaussian penalty]
\label{thm:e8v58-residual}
Let \(c\in\mathbb R^r\) have probability density
\begin{equation}
 d\nu_\kappa(c)
 =
 Z_\kappa^{-1}e^{-\kappa|c|^2/2}w(c)\,dc,
 \label{eq:e8v58-penalty-law}
\end{equation}
where \(w>0\) is \(C^1\), the boundary terms at infinity vanish, and
\begin{equation}
 |\nabla\log w(c)|\leq M.
 \label{eq:e8v58-log-bound}
\end{equation}
Then, with
\(X_\kappa=\langle|c|^2\rangle_{\nu_\kappa}^{1/2}\),
\begin{equation}
 \kappa X_\kappa^2\leq r+MX_\kappa
 \label{eq:e8v58-quadratic}
\end{equation}
and hence
\begin{equation}
 X_\kappa
 \leq
 {M+\sqrt{M^2+4\kappa r}\over2\kappa}
 \leq {M\over\kappa}+\sqrt{r\over\kappa}.
 \label{eq:e8v58-residual}
\end{equation}
\end{theorem}

\begin{proof}
Integrate the divergence
\[
 \operatorname{div}\!\left(
 c\,e^{-\kappa|c|^2/2}w(c)
 \right)
 =
 e^{-\kappa|c|^2/2}w(c)
 \left(
 r-\kappa|c|^2+c\cdot\nabla\log w(c)
 \right).
\]
The integral vanishes by the boundary hypothesis, so
\[
 \kappa\langle|c|^2\rangle
 =
 r+\langle c\cdot\nabla\log w(c)\rangle.
\]
Cauchy--Schwarz and \eqref{eq:e8v58-log-bound} give
\[
 \kappa X_\kappa^2
 \leq r+M\langle|c|\rangle
 \leq r+MX_\kappa,
\]
which is \eqref{eq:e8v58-quadratic}.  Solving the quadratic inequality
for its nonnegative root proves the first bound in
\eqref{eq:e8v58-residual}.  The second follows from
\(\sqrt{M^2+4\kappa r}\leq M+2\sqrt{\kappa r}\).
\end{proof}

\begin{corollary}[Scaling needed by the Einstein constraint response]
\label{cor:e8v58-scaling}
For a scalar constraint residual \(c_n\) entering the V50
gradient-energy estimate with spatial gap \(\lambda_{1,n}\), suppose the
conditional marginal hypotheses of \Cref{thm:e8v58-residual} hold with
\(r\) fixed and constants \(M_n\).  Then
\begin{equation}
 \lambda_{1,n}^{-1/2}
 \|c_n\|_{L^2(\nu_{\kappa_n})}
 \leq
 \lambda_{1,n}^{-1/2}
 \left({M_n\over\kappa_n}
 +\sqrt{r\over\kappa_n}\right).
 \label{eq:e8v58-response-residual}
\end{equation}
This residual vanishes if
\begin{equation}
 {M_n\over\kappa_n\sqrt{\lambda_{1,n}}}\longrightarrow0,
 \qquad
 \kappa_n\lambda_{1,n}\longrightarrow\infty.
 \label{eq:e8v58-scaling}
\end{equation}
\end{corollary}

\begin{proof}
Insert \eqref{eq:e8v58-residual} in the residual term of
\Cref{thm:e8v50-constraint-adjoint}.  The two limits in
\eqref{eq:e8v58-scaling} make the two terms on the right of
\eqref{eq:e8v58-response-residual} vanish.
\end{proof}

\begin{firewall}[Growing constraint rank]
\label{fire:e8v58-rank}
The global estimate \eqref{eq:e8v58-residual} contains
\(\sqrt{r/\kappa}\).  If \(r\) is the number of all lattice constraint
sites, it grows with volume and refinement.  A useful continuum argument
must apply the theorem conditionally to fixed-size local blocks and
compile them with owner or martingale estimates, or choose
\(\kappa_n\) large enough to pay the global rank.  One may not call the
bound uniform while suppressing \(r_n\).
\end{firewall}

\subsection{The hard-constraint surface}
\label{sec:e8v58-hard}

Let \(M\) be a compact Riemannian manifold with volume \(dm\), and let
\(C:M\to\mathbb R^r\) be smooth.  Write
\[
 J_C(x):=\sqrt{\det(DC_xDC_x^*)}.
\]

\begin{theorem}[Positive hard-constraint limit by coarea]
\label{thm:e8v58-coarea}
Assume \(0\) is a regular value of \(C\).  For every continuous test
\(F\) supported in a neighborhood containing no critical point of
\(C\),
\begin{equation}
 \lim_{\kappa\to\infty}
 \left({\kappa\over2\pi}\right)^{r/2}
 \int_M
 F(x)e^{-\kappa|C(x)|^2/2}\,dm(x)
 =
 \int_{C^{-1}(0)}{F(x)\over J_C(x)}\,d{\cal H}_{d-r}(x).
 \label{eq:e8v58-coarea}
\end{equation}
The limiting surface measure is positive.
\end{theorem}

\begin{proof}
The coarea formula gives
\[
 \int_MF(x)e^{-\kappa|C(x)|^2/2}dm(x)
 =
 \int_{\mathbb R^r}
 e^{-\kappa|y|^2/2}A_F(y)\,dy,
\]
where
\[
 A_F(y)
 =
 \int_{C^{-1}(y)}{F(x)\over J_C(x)}\,d{\cal H}_{d-r}(x).
\]
By the regular-value hypothesis and the support restriction,
\(A_F\) is continuous near zero and bounded with compact support after
shrinking the neighborhood if necessary.  The normalized Gaussian
\((\kappa/2\pi)^{r/2}e^{-\kappa|y|^2/2}\) is an approximate identity.
Therefore the limit is \(A_F(0)\), which is the right-hand side of
\eqref{eq:e8v58-coarea}.  Positivity follows when \(F\geq0\).
\end{proof}

\begin{corollary}[Normalized penalty measures approach the surface law]
\label{cor:e8v58-normalized}
If the denominator in
\eqref{eq:e8v58-coarea} with \(F=1\) is positive and the same regularity
holds on the relevant support, then normalized Gaussian-penalty
expectations converge to
\begin{equation}
 {\displaystyle
 \int_{C^{-1}(0)}FJ_C^{-1}\,d{\cal H}_{d-r}
 \over\displaystyle
 \int_{C^{-1}(0)}J_C^{-1}\,d{\cal H}_{d-r}}.
 \label{eq:e8v58-surface-law}
\end{equation}
\end{corollary}

\begin{proof}
Apply \Cref{thm:e8v58-coarea} to the numerator and denominator and take
their ratio.
\end{proof}

\begin{firewall}[The coarea determinant is not optional]
\label{fire:e8v58-determinant}
The positive hard-constraint law carries \(J_C^{-1}\).  At gauge
degeneracies \(DC\) can lose rank and this factor becomes singular.
Gauge quotient, Faddeev--Popov or BRST structure, Gribov sectors, and
boundary charges cannot be replaced by the formal symbol
\(\delta(C)\,dm\).
\end{firewall}

\subsection{Constraint-enforcement card}
\label{sec:e8v58-card}

\begin{definition}[Positive constraint-enforcement card, PCEC]
\label{def:e8v58-PCEC}
PCEC consists of:
\begin{enumerate}[label=\textup{(P\arabic*)},leftmargin=4em]
\item a declared lapse contour or Gaussian penalty at every regulator;
\item a local constraint rank and the corresponding multiplier
normalization;
\item conditional marginal bounds of the form
\eqref{eq:e8v58-log-bound};
\item the cofinal stiffness conditions
\eqref{eq:e8v58-scaling}, including growing-rank costs;
\item regular-value, determinant, gauge-quotient, and charge control for
the hard-constraint limit;
\item reflection and locality tests for the constraint-squared action;
\item compatibility with the V57 multiplier score and the V50
constraint-adjoint response; and
\item proof that the limiting positive surface carrier has the desired
Einstein--Standard Model observables.
\end{enumerate}
\end{definition}

\begin{assessment}[Progress at V58]
\label{ass:e8v58-status}
Finite Fourier enforcement is correctly constraint-selecting but signed;
a real Euclidean lapse is positive but anti-confining for a two-sided
constraint.  Gaussian lapse averaging yields a positive local
constraint-squared penalty with an explicit \(L^2\) residual, and the
coarea theorem identifies its positive hard-constraint limit and
Jacobian.  The price is a stiffness scale that must dominate the
large-volume spatial infrared scale without losing local conditional
control.
\end{assessment}

\begin{programme}[Eighth-series V59: penalty Hessian and local conditional stability]
\label{prog:e8v58-V59}
V59 will differentiate the constraint-squared penalty.  It will separate
the positive normal Hessian
\(\kappa\,DC^*DC\) from the off-surface term
\(\kappa\,C\cdot D^2C\), derive a quantitative tubular-core condition,
and test whether the V58 residual scale is compatible with a uniform
local conditional gap.
\end{programme}

\begin{firewall}[Claim boundary after eighth-series V58]
\label{fire:e8v58-boundary}
V58 proves the sign properties of three lapse prescriptions, the
Gaussian-penalty residual bound, its required cofinal scaling, and the
positive coarea hard-constraint limit.  It does not prove the marginal
log-derivative bound for Einstein--Standard Model constraints, handle
growing rank globally, control gauge singularities or reflection,
construct a uniform conditional gap, remove cutoffs, or prove the full
continuum.
\end{firewall}

\begin{theorem}[Eighth-series V58 result]
\label{thm:e8v58-result}
Gaussian lapse averaging replaces the oscillatory multiplier by the
positive weight \(e^{-\kappa|C|^2/2}\).  If the constraint marginal has
\(|\nabla\log w|\leq M\), then
\[
 \langle|C|^2\rangle^{1/2}
 \leq {M\over\kappa}+\sqrt{r\over\kappa}.
\]
After multiplication by the scalar constraint infrared factor, the
residual vanishes under \eqref{eq:e8v58-scaling}; the hard-constraint
limit is the positive coarea surface measure weighted by \(J_C^{-1}\).
\end{theorem}

\begin{proof}
Use \Cref{thm:e8v58-Gaussian,thm:e8v58-residual,
cor:e8v58-scaling,thm:e8v58-coarea}.
\end{proof}

\paragraph{Parent-version record.}
V58 was formed from
\texttt{Renewal\_SM\_Einstein\_Continuum\_Research\_Notes\_V57.tex}.
The V57 parent had (23\,288) logical source lines, (831\,331) bytes,
and SHA--256
\[
 \texttt{829DCD680A946BBD4F1D9D3B0991AD1CBF9BB754FD74CE2FE6B509D4C064C9F1}.
\]
The next compulsory companion audit is V60.

% =============================================================================
% END APPENDED EIGHTH-SERIES V58 MATERIAL
% =============================================================================
% =============================================================================
% BEGIN APPENDED EIGHTH-SERIES V59 MATERIAL
% =============================================================================

\clearpage
\section{Eighth-series V59: nonlinear penalty Hessians and normal--surface splitting}
\label{sec:v8v59}

\subsection{Exact Hessian of the constraint-squared penalty}
\label{sec:e8v59-Hessian}

Let \(U\subset\mathbb R^m\), let
\(C=(C^1,\ldots,C^r):U\to\mathbb R^r\) be \(C^2\), and define
\begin{equation}
 P_\kappa(x):={\kappa\over2}|C(x)|^2.
 \label{eq:e8v59-penalty}
\end{equation}

\begin{theorem}[Exact nonlinear penalty Hessian]
\label{thm:e8v59-Hessian}
For every \(x\in U\) and \(v\in\mathbb R^m\),
\begin{align}
 DP_\kappa(x)[v]
 &=\kappa\langle C(x),DC_xv\rangle,
 \label{eq:e8v59-gradient}\\
 D^2P_\kappa(x)[v,v]
 &=
 \kappa|DC_xv|^2
 +\kappa\langle C(x),D^2C_x[v,v]\rangle.
 \label{eq:e8v59-Hessian}
\end{align}
\end{theorem}

\begin{proof}
Differentiate \(|C|^2/2\) once and twice by the chain rule.
\end{proof}

At a regular point, let
\[
 T_x:=\ker DC_x,
 \qquad
 N_x:=T_x^\perp,
 \qquad
 v=v_T+v_N.
\]

\begin{corollary}[Tubular lower bound]
\label{cor:e8v59-tube}
Suppose on a set \(U_\delta=\{|C|\leq\delta\}\),
\begin{equation}
 |DC_xv_N|\geq\sigma|v_N|,
 \qquad
 |D^2C_x[v,v]|\leq K|v|^2.
 \label{eq:e8v59-geometry}
\end{equation}
Then
\begin{equation}
 D^2P_\kappa(x)[v,v]
 \geq
 \kappa\sigma^2|v_N|^2
 -\kappa\delta K\bigl(|v_T|^2+|v_N|^2\bigr).
 \label{eq:e8v59-tube-bound}
\end{equation}
\end{corollary}

\begin{proof}
The first term in \eqref{eq:e8v59-Hessian} is at least
\(\kappa\sigma^2|v_N|^2\), since \(DC_xv_T=0\).  Cauchy--Schwarz and
\eqref{eq:e8v59-geometry} bound the second term below by
\(-\kappa\delta K|v|^2\).
\end{proof}

\begin{theorem}[Combined local coercivity criterion]
\label{thm:e8v59-combined}
Let \(V\) be the remaining local regulated action and suppose on
\(U_\delta\)
\begin{equation}
 D^2V(x)[v,v]
 \geq a_T|v_T|^2-b_N|v_N|^2
 \label{eq:e8v59-base}
\end{equation}
with \(a_T>0\) and \(b_N\geq0\).  Then
\begin{align}
 D^2(V+P_\kappa)(x)[v,v]
 &\geq
 (a_T-\kappa\delta K)|v_T|^2
 \notag\\
 &\quad+
 (\kappa\sigma^2-b_N-\kappa\delta K)|v_N|^2.
 \label{eq:e8v59-combined}
\end{align}
In particular, if
\begin{equation}
 \kappa\delta K\leq{a_T\over2},
 \qquad
 \kappa\sigma^2-b_N-\kappa\delta K\geq c_N>0,
 \label{eq:e8v59-coercive-conditions}
\end{equation}
then
\begin{equation}
 D^2(V+P_\kappa)[v,v]
 \geq
 {a_T\over2}|v_T|^2+c_N|v_N|^2.
 \label{eq:e8v59-coercive}
\end{equation}
\end{theorem}

\begin{proof}
Add \eqref{eq:e8v59-base} and
\eqref{eq:e8v59-tube-bound}.  The two assumptions in
\eqref{eq:e8v59-coercive-conditions} give
\eqref{eq:e8v59-coercive}.
\end{proof}

\begin{corollary}[Affine constraints have no off-surface Hessian debit]
\label{cor:e8v59-affine}
If \(C(x)=Ax-b\), then \(D^2C=0\) and
\begin{equation}
 D^2P_\kappa=\kappa A^*A.
 \label{eq:e8v59-affine}
\end{equation}
Hence the penalty is nonnegative in every direction and is
\(\kappa\sigma^2\)-coercive on
\((\ker A)^\perp\), where \(\sigma\) is the least positive singular
value of \(A\).
\end{corollary}

\begin{proof}
Set \(D^2C=0\) in \eqref{eq:e8v59-Hessian}.
\end{proof}

\subsection{The tubular-width conflict in ambient Hessian estimates}
\label{sec:e8v59-conflict}

\begin{proposition}[The crude convexity tube shrinks faster than the
penalty fluctuation]
\label{prop:e8v59-conflict}
Assume \(a_T,K>0\) are independent of \(\kappa\).  The first condition
in \eqref{eq:e8v59-coercive-conditions} requires
\begin{equation}
 \delta_\kappa\leq{a_T\over2K\kappa}.
 \label{eq:e8v59-delta}
\end{equation}
If the constraint fluctuation has only the V58 scale
\begin{equation}
 \langle|C|^2\rangle^{1/2}\asymp\kappa^{-1/2},
 \label{eq:e8v59-typical}
\end{equation}
then Markov's estimate at the tube
\eqref{eq:e8v59-delta} is non-informative:
\begin{equation}
 \mathbb P(|C|>\delta_\kappa)
 \leq{\langle|C|^2\rangle\over\delta_\kappa^2}
 =O(\kappa).
 \label{eq:e8v59-Markov}
\end{equation}
\end{proposition}

\begin{proof}
The tangential coefficient in
\eqref{eq:e8v59-combined} is positive by the stated argument only if
\(\kappa\delta_\kappa K<a_T\); the displayed half-margin gives
\eqref{eq:e8v59-delta}.  Insert
\(\langle|C|^2\rangle=O(\kappa^{-1})\) and
\(\delta_\kappa^{-2}=O(\kappa^2)\) in Markov's inequality.
\end{proof}

\begin{firewall}[A bad upper bound is not instability]
\label{fire:e8v59-conflict}
\Cref{prop:e8v59-conflict} does not show that the penalty measure is
unstable.  It shows that ambient global convexity, using only a uniform
bound on \(D^2C\), asks for a tube much thinner than the typical Gaussian
normal fluctuation.  Adapted normal coordinates or a conditional
normal--surface Poincaré argument may remove the spurious tangential
debit.
\end{firewall}

\subsection{Exact product normal--surface model}
\label{sec:e8v59-product}

Let \(Y=\mathbb R^r\) carry the centered Gaussian
\[
 d\gamma_\kappa(y)
 =
 \left({\kappa\over2\pi}\right)^{r/2}
 e^{-\kappa|y|^2/2}\,dy.
\]
Let \((Z,\nu)\) carry a Dirichlet form
\({\cal E}_Z\) with spectral gap \(g_Z>0\).
On \(Y\times Z\), define
\begin{equation}
 {\cal E}(f,f)
 :=
 \int|\nabla_yf|^2\,d\gamma_\kappa d\nu
 +\int{\cal E}_Z(f(y,\cdot),f(y,\cdot))\,d\gamma_\kappa(y).
 \label{eq:e8v59-product-form}
\end{equation}

\begin{theorem}[Exact product penalty gap]
\label{thm:e8v59-product}
The product form \eqref{eq:e8v59-product-form} has spectral gap
\begin{equation}
 g_{\rm prod}=\min\{\kappa,g_Z\}.
 \label{eq:e8v59-product-gap}
\end{equation}
\end{theorem}

\begin{proof}
Variance decomposition gives
\begin{align*}
 \operatorname{Var}_{\gamma_\kappa\otimes\nu}(f)
 &=
 \int\operatorname{Var}_\nu(f(y,\cdot))\,d\gamma_\kappa(y)
 +
 \operatorname{Var}_{\gamma_\kappa}
 \left(\int f(y,z)\,d\nu(z)\right).
\end{align*}
The \(Z\) Poincaré inequality bounds the first term by
\(g_Z^{-1}\) times the second term of
\eqref{eq:e8v59-product-form}.  The Gaussian Poincaré inequality with
gap \(\kappa\), followed by Jensen's inequality for the \(z\)-average
of \(\nabla_yf\), bounds the second variance by
\(\kappa^{-1}\) times the first term of the form.  Hence the gap is at
least \(\min\{\kappa,g_Z\}\).

For the reverse inequality, use a function depending only on \(y\) and
approaching a Gaussian first eigenfunction to obtain gap at most
\(\kappa\); use a function depending only on \(z\) and a minimizing
sequence for the \(Z\) Rayleigh quotient to obtain gap at most \(g_Z\).
\end{proof}

\begin{corollary}[Large penalty does not damage an exact product gap]
\label{cor:e8v59-product}
If \(g_Z\geq g_*>0\) uniformly and the regulated constraint coordinates
factor exactly as above, then every
\(\kappa\geq g_*\) gives
\[
 g_{\rm prod}\geq g_*.
\]
Thus stiffness strengthens the normal gap and leaves the surface gap as
the bottleneck.
\end{corollary}

\begin{proof}
Use \eqref{eq:e8v59-product-gap}.
\end{proof}

\begin{firewall}[Coarea coordinates need not factor]
\label{fire:e8v59-nonproduct}
The coarea law in V58 generally has a \(y\)-dependent surface measure,
Jacobian \(J_C^{-1}\), gauge determinant, and interacting action.
Therefore it is not the product law of
\Cref{thm:e8v59-product}.  Passing from the exact product theorem to the
Einstein--Standard Model constraint tube requires uniform conditional
surface gaps and quantitative control of how the conditional laws vary
with the normal coordinate.
\end{firewall}

\subsection{Conditional normal--surface card}
\label{sec:e8v59-card}

\begin{definition}[Normal--surface penalty card, NSPC]
\label{def:e8v59-NSPC}
NSPC consists of:
\begin{enumerate}[label=\textup{(H\arabic*)},leftmargin=4em]
\item a uniform regular-value tube and lower singular bound for \(DC\);
\item exact coarea coordinates with their \(J_C^{-1}\), gauge, ghost,
and boundary factors;
\item a normal conditional Poincaré estimate at scale \(\kappa\);
\item a uniform Poincaré or positive-carrier estimate on the constraint
surface;
\item a strict interdependence, block-dynamics, or martingale estimate
coupling normal and surface conditionals;
\item the V58 residual scaling and growing-rank ledger;
\item reflection and locality control of the constraint-squared action;
and
\item compatibility with lapse scores, charges, observables, and cutoff
removal.
\end{enumerate}
\end{definition}

\begin{theorem}[What NSPC would resolve]
\label{thm:e8v59-NSPC}
In the exact product case, rows (H3)--(H4) give the gap
\(\min\{\kappa,g_Z\}\).  For a nonproduct tube, any valid comparison
must add row (H5); good-tube probability and the local Hessian estimate
alone cannot supply it.
\end{theorem}

\begin{proof}
The product assertion is \Cref{thm:e8v59-product}.  The necessity of an
interface or interdependence input follows from the audited V13
three-state obstruction recorded in \Cref{thm:e8v55-rarity}: local gap
and bad probability can remain fixed while global mixing vanishes.
\end{proof}

\begin{assessment}[Progress at V59]
\label{ass:e8v59-status}
The penalty contributes a strong positive normal Hessian
\(\kappa DC^*DC\), but nonlinear constraint curvature creates an
off-surface term \(\kappa C\cdot D^2C\).  A crude ambient convexity
argument demands an \(O(\kappa^{-1})\) tube, narrower than the
\(O(\kappa^{-1/2})\) penalty fluctuation.  In exact normal--surface
product coordinates the gap is instead
\(\min\{\kappa,g_{\rm surface}\}\), showing that the real missing datum
is surface mixing plus controlled normal--surface dependence.
\end{assessment}

\begin{programme}[Eighth-series V60: companion audit and conditional-gap route]
\label{prog:e8v59-V60}
V60 will perform the compulsory YM/NS audit.  It will compare the newest
YM block-dynamics influence estimate with NSPC and decide whether the
next Einstein step should derive a normal--surface conditional
interdependence matrix, a surface Ward estimate, or a boundary/charge
decomposition.
\end{programme}

\begin{firewall}[Claim boundary after eighth-series V59]
\label{fire:e8v59-boundary}
V59 proves the exact nonlinear penalty Hessian, a tubular coercivity
criterion, the ambient-width obstruction, and the exact product
normal--surface gap.  It does not construct coarea coordinates uniformly
for the Einstein constraints, prove a surface gap or strict
interdependence, control gauge singularities and reflection, remove
cutoffs, or prove the full continuum.
\end{firewall}

\begin{theorem}[Eighth-series V59 result]
\label{thm:e8v59-result}
The nonlinear penalty satisfies
\[
 D^2P_\kappa[v,v]
 =
 \kappa|DCv|^2+\kappa\langle C,D^2C[v,v]\rangle.
\]
Ambient convexity alone is incompatible with the typical penalty tube
at the available bounds.  In exact coarea product coordinates, however,
the normal--surface gap is
\(\min\{\kappa,g_{\rm surface}\}\); hence large penalty stiffness is
compatible with a uniform carrier precisely when the surface and
interdependence rows are uniform.
\end{theorem}

\begin{proof}
Use \Cref{thm:e8v59-Hessian,prop:e8v59-conflict,
thm:e8v59-product,thm:e8v59-NSPC}.
\end{proof}

\paragraph{Parent-version record.}
V59 was formed from
\texttt{Renewal\_SM\_Einstein\_Continuum\_Research\_Notes\_V58.tex}.
The V58 parent had (23\,687) logical source lines, (844\,592) bytes,
and SHA--256
\[
 \texttt{6642796F957100B86C3DBC35EED83AC6C14418D492D5CC68F561DFE9960AB0E4}.
\]
The next compulsory companion audit is V60.

% =============================================================================
% END APPENDED EIGHTH-SERIES V59 MATERIAL
% =============================================================================
% =============================================================================
% BEGIN APPENDED EIGHTH-SERIES V60 MATERIAL
% =============================================================================

\clearpage
\section{Eighth-series V60: compulsory audit and Gaussian block-frame transport}
\label{sec:v8v60}

\subsection{Companion files inspected}
\label{sec:e8v60-audit}

The V60 audit inspected the newest complete companion notes
\begin{align}
 &\texttt{Renewal\_Yang\_Mills\_Mass\_Gap\_Research\_Notes\_V16.tex},
 \notag\\
 &\hspace{2em}6472\ \text{logical lines},\quad237803\ \text{bytes},
 \quad\operatorname{SHA256}
 =\texttt{98B47DEDC8093BCAF903C83F5816A8AEB5B888A0E2373A396CE14D837041B95F},
 \label{eq:e8v60-YM-file}\\
 &\texttt{Renewal\_NS\_Stability\_Research\_Notes\_V26.tex},
 \notag\\
 &\hspace{2em}5319\ \text{logical lines},\quad278283\ \text{bytes},
 \quad\operatorname{SHA256}
 =\texttt{82C887F8C2C69E4F28FAD7C3DC8DE5B9099CB5A4BF431EBA1FFF3E91935978AD}.
 \label{eq:e8v60-NS-file}
\end{align}
Both have exactly one live terminator.  A V15 Yang--Mills predecessor
also remains in the shared folder; it was not modified or removed.  The
NS note is unchanged.

\begin{audit}[Yang--Mills V14]
\label{audit:e8v60-YM14}
V14 proves an exact block-dynamics comparison: uniform conditional
block gaps multiply a heat-bath projection gap.  It expresses the
projection gap through an operator-valued Gram matrix and supplies a
scalar diagonal-dominance certificate.  Differentiating a one-block
conditional law gives the exact source-to-block commutator and its
incidence sparsity.  The declared one-link partition fails uniformity,
so the programme enlarges blocks rather than concealing that failure in
a formal Dobrushin constant.
\end{audit}

\begin{audit}[Yang--Mills V15--V16]
\label{audit:e8v60-YM16}
V15 constructs five half-overlapping twenty-chord blocks with a
certified triangular curvature coordinate in every block and only
constants in their common invariant sigma field.  V16 replaces a crude
singular estimate by exact integer inverse norms, obtaining a full
flat-curvature floor \(1/\sqrt{28}\), certified block floors, and a
first-chaos frame bound \(1/952\).  It proves that the associated
correlated Gaussian block heat bath has the same gap on all Wiener
chaoses.  The remaining YM step is stability of those conditional
projections under physical measure deformation.
\end{audit}

\begin{audit}[Effect on the V59 normal--surface problem]
\label{audit:e8v60-effect}
The audited Gaussian theorem supplies exactly the missing method for
NSPC row (H3): whiten the linearized constraint-normal Gaussian,
identify the conditional block subspaces, and prove a first-chaos frame
bound.  It does not supply the constraint-surface gap in row (H4), and
it warns that common invariants or local Jacobian floors alone do not
give a quantitative gap.  The Einstein calculation must therefore
separate the normal frame from the physical tangent dynamics.
\end{audit}

\subsection{Transport to a linearized constraint-normal Gaussian}
\label{sec:e8v60-transport}

Let \(X\) and \(Y\) be finite-dimensional Hilbert spaces and let
\(A:X\to Y\).  Put \(N=(\ker A)^\perp\).  On \(N\), consider
\begin{equation}
 d\gamma_{\kappa,A}(x)
 =
 Z^{-1}\exp\!\left(-{\kappa\over2}\|Ax\|_Y^2\right)dx.
 \label{eq:e8v60-normal-Gaussian}
\end{equation}
Let \(\Lambda_\alpha\) be coordinate blocks in \(X\), and let
\(K_\alpha^{\kappa,A}\) be conditional expectation for
\(\gamma_{\kappa,A}\) given the coordinates outside
\(\Lambda_\alpha\), restricted to the normal fibre.

\begin{proposition}[Penalty stiffness does not change the whitened
block frame]
\label{prop:e8v60-kappa}
The map
\begin{equation}
 W_\kappa:N\longrightarrow\operatorname{Ran}A,
 \qquad
 W_\kappa x=\sqrt{\kappa}\,Ax
 \label{eq:e8v60-whitening}
\end{equation}
is a bijection carrying \(\gamma_{\kappa,A}\) to a standard Gaussian
on \(\operatorname{Ran}A\).  Let \(X_{\Lambda_\alpha}\) be the coordinate subspace supported
in \(\Lambda_\alpha\), and put
\begin{equation}
 V_\alpha:=N\cap X_{\Lambda_\alpha},
 \qquad
 U_\alpha:=A V_\alpha\subset\operatorname{Ran}A.
 \label{eq:e8v60-block-space}
\end{equation}
The first-chaos subspace changed by the conditional block update is
\(U_\alpha\), which is independent of \(\kappa\).
\end{proposition}

\begin{proof}
Since \(A|_N\) is injective and has range \(\operatorname{Ran}A\),
\(W_\kappa\) is a bijection.  Its Jacobian is constant, and
\(\kappa\|Ax\|^2=\|W_\kappa x\|^2\), proving the Gaussian claim.
Two points of \(N\) with identical coordinates outside
\(\Lambda_\alpha\) differ by an element of
\(N\cap X_{\Lambda_\alpha}=V_\alpha\).  Their whitened coordinates
therefore differ in \(\sqrt\kappa\,A V_\alpha\).  Multiplication by the
nonzero scalar \(\sqrt\kappa\) does not change that subspace, giving
\eqref{eq:e8v60-block-space}.
\end{proof}

\begin{corollary}[Audited all-chaos frame compiler]
\label{cor:e8v60-frame}
Let \(P_\alpha\) be orthogonal projection onto \(U_\alpha\).  If
\begin{equation}
 \sum_\alpha P_\alpha\succeq\delta_N I_{\operatorname{Ran}A}
 \label{eq:e8v60-frame}
\end{equation}
for some \(\delta_N>0\), then the audited YM V16 all-chaos Gaussian
theorem gives
\begin{equation}
 \operatorname{gap}
 \left(\sum_\alpha(I-K_\alpha^{\kappa,A})\right)
 \geq\delta_N
 \label{eq:e8v60-normal-gap}
\end{equation}
for every \(\kappa>0\).  Thus the V58 stiffness limit does not itself
degrade the linearized normal heat-bath gap.
\end{corollary}

\begin{proof}
\Cref{prop:e8v60-kappa} unitarily identifies the first chaos with the
standard Gaussian block subspaces \(U_\alpha\).  The audited V16 theorem
states that the complete Gaussian heat-bath gap equals the least
eigenvalue of the first-chaos frame operator.  Apply
\eqref{eq:e8v60-frame}.
\end{proof}

\begin{firewall}[The tangent sector is invisible to normal updates]
\label{fire:e8v60-tangent}
If the state space includes \(\ker A\) and every block update resamples
only the constraint-normal conditional, then every nonconstant function
of the tangent coordinate is invariant.  The full heat-bath generator
therefore has zero gap even when \(\delta_N>0\).  A surface dynamics and
its own uniform gap or positive-carrier estimate remain necessary.
\end{firewall}

\begin{proposition}[Common invariants identify the missing sector]
\label{prop:e8v60-invariants}
For conditional projections \(K_\alpha\), one has
\begin{equation}
 \ker\sum_\alpha(I-K_\alpha)
 =
 \bigcap_\alpha\operatorname{Ran}K_\alpha.
 \label{eq:e8v60-common-kernel}
\end{equation}
Consequently, proving that the common invariant space consists of
functions of the constraint-surface coordinate identifies the tangent
sector but gives no positive quantitative gap on that sector.
\end{proposition}

\begin{proof}
Since each \(I-K_\alpha\) is an orthogonal projection,
\[
 \left\langle f,\sum_\alpha(I-K_\alpha)f\right\rangle
 =\sum_\alpha\|(I-K_\alpha)f\|_2^2.
\]
The sum vanishes exactly when every \(K_\alpha f=f\), proving
\eqref{eq:e8v60-common-kernel}.  An infinite-dimensional nonnegative
operator can have the stated kernel while its positive spectrum
accumulates at zero, so the kernel identity supplies no spectral floor.
\end{proof}

\subsection{Audit decision}
\label{sec:e8v60-decision}

\begin{theorem}[V60 direction selection]
\label{thm:e8v60-direction}
The next Einstein calculation should test the frame
\eqref{eq:e8v60-frame} for the discrete linearized scalar constraint
before attempting nonlinear conditional-measure transport.  If a
fixed-size local block cover already has
\(\delta_N\to0\) with physical volume, nonlinear density comparison
cannot create a uniform margin from that cover.  If it has a strict
flat margin, the subsequent target is projection stability, following
the audited YM route.
\end{theorem}

\begin{proof}
The exact Gaussian normal gap is bounded below by the flat frame margin
through \Cref{cor:e8v60-frame}.  A zero or collapsing flat margin leaves
no perturbative reserve for nonlinear transport.  Conversely, a
positive flat margin is only the reference estimate; the audited V16
firewall shows that physical conditional projections and local
conditional gaps still require a separate perturbation theorem.
\end{proof}

\begin{assessment}[Progress at V60]
\label{ass:e8v60-status}
The audit converts the abstract V59 normal gap into a concrete
finite-matrix problem.  Penalty stiffness cancels under whitening, and
the all-chaos normal gap is controlled by a first-chaos block frame.
The surface sector remains a genuine common invariant.  The next
calculation will test whether local Einstein scalar-constraint updates
have a volume-uniform frame or inherit the inverse-Laplacian slow mode.
\end{assessment}

\begin{programme}[Eighth-series V61: scalar-constraint local-block test]
\label{prog:e8v60-V61}
V61 will compute the Gaussian heat-bath frame for a periodic
linearized scalar constraint with precision
\(\kappa(-\Delta)^2\).  It will begin with one-site blocks, derive the
exact first-chaos eigenvalues, and determine their volume scaling.  A
collapsing margin will force block enlargement or a multiscale update
rather than a nonlinear perturbation of the same cover.
\end{programme}

\begin{firewall}[Claim boundary after eighth-series V60]
\label{fire:e8v60-boundary}
V60 records the mandatory audit and rigorously transports the audited
Gaussian block-frame theorem to a linearized constraint-normal
Gaussian, including independence from penalty stiffness and the tangent
common-invariant obstruction.  It does not prove a strict Einstein
frame, a surface gap, nonlinear projection stability, boundary or
charge control, cutoff removal, or the full continuum.
\end{firewall}

\begin{theorem}[Eighth-series V60 result]
\label{thm:e8v60-result}
For the linearized normal penalty measure
\[
 d\gamma_{\kappa,A}\propto e^{-\kappa\|Ax\|^2/2}dx,
\]
the Gaussian block heat-bath frame depends on the subspaces
\(A(N\cap X_{\Lambda_\alpha})\) and is independent of
\(\kappa\).  A first-chaos frame floor \(\delta_N\) gives the same
all-chaos normal gap, while functions of \(\ker A\) require a separate
surface dynamics.
\end{theorem}

\begin{proof}
Use \Cref{prop:e8v60-kappa,cor:e8v60-frame,
fire:e8v60-tangent}.
\end{proof}

\paragraph{Parent-version record.}
V60 was formed from
\texttt{Renewal\_SM\_Einstein\_Continuum\_Research\_Notes\_V59.tex}.
The V59 parent had (24\,048) logical source lines, (856\,180) bytes,
and SHA--256
\[
 \texttt{FDE2FCB055FEE02E85DB1FD48447D330ECBB32D97D90D482FDCC387F37DA8815}.
\]
The next compulsory companion audit is V65.

% =============================================================================
% END APPENDED EIGHTH-SERIES V60 MATERIAL
% =============================================================================
% =============================================================================
% BEGIN APPENDED EIGHTH-SERIES V61 MATERIAL
% =============================================================================

\clearpage
\section{Eighth-series V61: exact local-block gap for the linear scalar constraint}
\label{sec:v8v61}

\subsection{Periodic scalar penalty model}
\label{sec:e8v61-model}

Let
\[
 \Lambda_L=(\mathbb Z/L\mathbb Z)^d,\qquad L\geq5,
\]
with counting inner product.  Let \(X_0\) be the mean-zero real scalar
fields and let \(A=-\Delta_L\) be the combinatorial periodic Laplacian
\begin{equation}
 (A\phi)_x
 =
 \sum_{j=1}^d
 \bigl(2\phi_x-\phi_{x+e_j}-\phi_{x-e_j}\bigr).
 \label{eq:e8v61-Laplacian}
\end{equation}
The linear scalar-constraint penalty measure is
\begin{equation}
 d\gamma_{\kappa,L}(\phi)
 =
 Z_{\kappa,L}^{-1}
 \exp\!\left(-{\kappa\over2}\|A\phi\|_2^2\right)d\phi,
 \qquad \phi\in X_0.
 \label{eq:e8v61-measure}
\end{equation}
This is the normal Gaussian for the inverse-Laplacian scalar response
isolated in V49, with its constant charge mode removed.

\subsection{One-site blocks are frozen by the charge condition}
\label{sec:e8v61-one-site}

\begin{proposition}[One-site normal update is degenerate]
\label{prop:e8v61-one-site}
For every site \(x\),
\begin{equation}
 X_0\cap\operatorname{span}\{e_x\}=\{0\}.
 \label{eq:e8v61-one-site-intersection}
\end{equation}
Consequently, conditioning a field in \(X_0\) on all sites except \(x\)
determines \(\phi_x\) exactly.  The one-site conditional expectation is
the identity and the one-site heat-bath generator has gap zero.
\end{proposition}

\begin{proof}
A vector \(ae_x\) has spatial sum \(a\), so it is mean zero only when
\(a=0\).  If every coordinate except \(\phi_x\) is fixed, the equation
\(\sum_y\phi_y=0\) fixes
\(\phi_x=-\sum_{y\neq x}\phi_y\).  Hence the conditional law is a point
mass, its conditional-expectation operator is \(I\), and
\(\sum_x(I-K_x)=0\).
\end{proof}

\begin{firewall}[Updating the charge is a different model]
\label{fire:e8v61-charge}
Allowing a one-site update by varying the constant mode would remove the
degeneracy, but it would also resample the charge sector that V49--V50
keeps separate.  A gap obtained that way is not a gap on the physical
mean-zero constraint-normal fibre.
\end{firewall}

\subsection{Nearest-neighbor pair blocks}
\label{sec:e8v61-pairs}

Let \(E_L\) be the set of positively oriented nearest-neighbor edges
\(e=(x,x+e_j)\).  Put
\begin{equation}
 d_e:=e_x-e_{x+e_j}\in X_0,
 \qquad
 u_e:=Ad_e.
 \label{eq:e8v61-edge-vectors}
\end{equation}
For the two-site block \(\{x,x+e_j\}\), the allowed variation at fixed
complement and fixed mean is precisely \(\operatorname{span}\{d_e\}\).
After the V60 whitening map its first-chaos update subspace is therefore
\(\operatorname{span}\{u_e\}\).

\begin{lemma}[Exact edge-vector norm]
\label{lem:e8v61-edge-norm}
For every \(e\in E_L\),
\begin{equation}
 \|u_e\|_2^2=4d(2d+3).
 \label{eq:e8v61-edge-norm}
\end{equation}
\end{lemma}

\begin{proof}
Translation and coordinate symmetry make the norm independent of the
edge.  The Laplacian matrix has
\[
 (A^2)_{00}=(2d)^2+2d=4d^2+2d.
\]
For a nearest neighbor \(e_1\), and \(L\geq5\), the only nonzero
two-step contributions shared by the matrix product are through the two
endpoints, giving
\[
 (A^2)_{0,e_1}=-2d-2d=-4d.
\]
Therefore
\[
 \|A(e_0-e_{e_1})\|^2
 =
 \langle e_0-e_{e_1},A^2(e_0-e_{e_1})\rangle
 =
 2\bigl((A^2)_{00}-(A^2)_{0,e_1}\bigr)
 =8d^2+12d.
\]
\end{proof}

\begin{lemma}[Edge incidence identity]
\label{lem:e8v61-incidence}
On all scalar fields,
\begin{equation}
 \sum_{e\in E_L}d_ed_e^*=A.
 \label{eq:e8v61-incidence}
\end{equation}
\end{lemma}

\begin{proof}
For any \(\phi\),
\[
 \left\langle\phi,\sum_ed_ed_e^*\phi\right\rangle
 =
 \sum_{x,j}|\phi_x-\phi_{x+e_j}|^2
 =
 \langle\phi,A\phi\rangle.
\]
Both operators are self-adjoint, so equality of their quadratic forms
proves \eqref{eq:e8v61-incidence}.
\end{proof}

\begin{theorem}[Exact first-chaos pair-frame operator]
\label{thm:e8v61-frame}
Let \(P_e\) be orthogonal projection in the whitened first chaos onto
\(\operatorname{span}\{u_e\}\).  Then on \(X_0=\operatorname{Ran}A\),
\begin{equation}
 {\mathbb F}_{\rm pair}
 :=
 \sum_{e\in E_L}P_e
 =
 {A^3\over4d(2d+3)}.
 \label{eq:e8v61-frame}
\end{equation}
\end{theorem}

\begin{proof}
By \Cref{lem:e8v61-edge-norm},
\[
 P_e={u_eu_e^*\over4d(2d+3)}
 ={Ad_ed_e^*A\over4d(2d+3)}.
\]
Sum over edges and apply \Cref{lem:e8v61-incidence}:
\[
 \sum_eP_e
 ={A(\sum_ed_ed_e^*)A\over4d(2d+3)}
 ={A^3\over4d(2d+3)}.
\]
\end{proof}

\subsection{Fourier spectrum and volume collapse}
\label{sec:e8v61-spectrum}

For \(k\in\Lambda_L\), the Fourier eigenvalue of \(A\) is
\begin{equation}
 \lambda(k)
 =
 4\sum_{j=1}^d\sin^2\!\left({\pi k_j\over L}\right).
 \label{eq:e8v61-lambda}
\end{equation}

\begin{corollary}[Exact pair-frame margin]
\label{cor:e8v61-margin}
The least pair-frame eigenvalue on \(X_0\) is
\begin{equation}
 \delta_{\rm pair}(L,d)
 =
 { \bigl(4\sin^2(\pi/L)\bigr)^3
 \over4d(2d+3)}.
 \label{eq:e8v61-margin}
\end{equation}
In particular,
\begin{equation}
 \delta_{\rm pair}(L,d)
 \sim
 {16\pi^6\over d(2d+3)}\,L^{-6}
 \qquad(L\to\infty).
 \label{eq:e8v61-asymptotic}
\end{equation}
\end{corollary}

\begin{proof}
Diagonalize \eqref{eq:e8v61-frame} in Fourier space.  The least nonzero
Laplacian eigenvalue is
\(\lambda_{\min}=4\sin^2(\pi/L)\), attained at a unit coordinate
momentum.  Cubing and dividing gives
\eqref{eq:e8v61-margin}.  Since
\(\sin(\pi/L)\sim\pi/L\),
\[
 {64(\pi/L)^6\over4d(2d+3)}
 ={16\pi^6\over d(2d+3)}L^{-6},
\]
which proves \eqref{eq:e8v61-asymptotic}.
\end{proof}

\begin{theorem}[All-chaos pair heat-bath gap]
\label{thm:e8v61-all-chaos}
Let \(K_e\) be conditional expectation for
\(\gamma_{\kappa,L}\) given all sites outside the pair \(e\).  Then
\begin{equation}
 \operatorname{gap}
 \left(\sum_{e\in E_L}(I-K_e)\right)
 =
 \delta_{\rm pair}(L,d).
 \label{eq:e8v61-all-chaos}
\end{equation}
The value is independent of \(\kappa\) and collapses as \(L^{-6}\).
\end{theorem}

\begin{proof}
V60 whitening identifies the pair first-chaos subspaces with
\(\operatorname{span}\{u_e\}\), independently of \(\kappa\).
The audited YM V16 all-chaos Gaussian theorem states that the Gaussian
block heat-bath gap equals the least eigenvalue of the first-chaos frame
operator.  Apply \Cref{cor:e8v61-margin}.
\end{proof}

\begin{firewall}[Penalty stiffness cannot cure local mixing]
\label{fire:e8v61-kappa}
Increasing \(\kappa\) suppresses the constraint residual as in V58, but
it does not improve \eqref{eq:e8v61-all-chaos}.  The scalar whitening
removes \(\kappa\) from every conditional projection.  Thus the
conditions \(\kappa_n\lambda_{1,n}\to\infty\) and a uniform local-block
mixing gap are logically independent.
\end{firewall}

\begin{firewall}[Pair locality has a sixth-order infrared cost]
\label{fire:e8v61-locality}
The \(L^{-6}\) collapse is already present in the exact linear Gaussian
normal fibre.  A bounded perturbation or nonlinear conditional-density
comparison can preserve only a fraction of an existing margin; it
cannot turn this collapsing pair cover into a volume-uniform one.
Larger or multiscale blocks must act directly on the low Fourier modes.
\end{firewall}

\subsection{Scalar block-frame card}
\label{sec:e8v61-card}

\begin{definition}[Multiscale scalar constraint card, MSCC]
\label{def:e8v61-MSCC}
MSCC consists of:
\begin{enumerate}[label=\textup{(M\arabic*)},leftmargin=4em]
\item exact removal or separate treatment of the constant charge mode;
\item a family of charge-preserving blocks, including coarse scales;
\item a first-chaos frame lower bound uniform in physical volume and
refinement;
\item the audited all-chaos Gaussian lift;
\item stability of conditional projections under nonlinear constraint,
coarea-Jacobian, gauge, ghost, and Standard Model density deformations;
\item a separate constraint-surface dynamics and gap;
\item compatibility with the V58 stiffness scaling and V50
constraint-energy response; and
\item boundary, reflection, positivity, observable, and cutoff-removal
contacts.
\end{enumerate}
\end{definition}

\begin{assessment}[Progress at V61]
\label{ass:e8v61-status}
The first local Einstein scalar block test is decisive.  One-site
updates are frozen by the fixed charge.  The minimal pair update is
nontrivial, but its exact first-chaos and all-chaos gap is
\[
 { \bigl(4\sin^2(\pi/L)\bigr)^3\over4d(2d+3)}
 \asymp L^{-6}.
\]
This failure is linear, exact, and independent of penalty stiffness.
The programme must introduce coarse or multiscale charge-preserving
updates before nonlinear projection transport.
\end{assessment}

\begin{programme}[Eighth-series V62: Fourier-shell and multiscale frames]
\label{prog:e8v61-V62}
V62 will first compute the ideal Fourier-mode update, which gives a
unit frame but is nonlocal.  It will then formulate a dyadic
charge-preserving block frame and determine the scale weights required
to remove the \(L^{-6}\) pair-mode collapse without hiding volume
dependence in the update rates.
\end{programme}

\begin{firewall}[Claim boundary after eighth-series V61]
\label{fire:e8v61-boundary}
V61 proves one-site degeneracy, the exact pair-frame operator
\(A^3/[4d(2d+3)]\), and its exact all-chaos Gaussian gap and
large-volume asymptotics.  It does not construct a uniform multiscale
frame, prove nonlinear transport or a surface gap, close charge and
boundary sectors, remove cutoffs, or prove the full continuum.
\end{firewall}

\begin{theorem}[Eighth-series V61 result]
\label{thm:e8v61-result}
For the periodic mean-zero scalar constraint penalty
\(e^{-\kappa\|(-\Delta)\phi\|^2/2}\), one-site normal heat-bath updates
are trivial.  Nearest-neighbor charge-preserving pair updates have exact
all-chaos gap
\[
 \operatorname{gap}
 =
 { \bigl(4\sin^2(\pi/L)\bigr)^3\over4d(2d+3)}
 \asymp L^{-6},
\]
independently of \(\kappa\).
\end{theorem}

\begin{proof}
Use \Cref{prop:e8v61-one-site,thm:e8v61-frame,
cor:e8v61-margin,thm:e8v61-all-chaos}.
\end{proof}

\paragraph{Parent-version record.}
V61 was formed from
\texttt{Renewal\_SM\_Einstein\_Continuum\_Research\_Notes\_V60.tex}.
The V60 parent had (24\,312) logical source lines, (866\,974) bytes,
and SHA--256
\[
 \texttt{6F4DE971C5AC97CABDD23B0E655D56EA396686B0EAFF9517A28F871348684AAF}.
\]
The next compulsory companion audit is V65.

% =============================================================================
% END APPENDED EIGHTH-SERIES V61 MATERIAL
% =============================================================================
% =============================================================================
% BEGIN APPENDED EIGHTH-SERIES V62 MATERIAL
% =============================================================================

\clearpage
\section{Eighth-series V62: finite-range obstruction and the ideal Fourier-shell frame}
\label{sec:v8v62}

\subsection{Translation-orbit block directions}
\label{sec:e8v62-templates}

Continue with the periodic scalar penalty model of V61.  Let
\({\cal V}=\{v^1,\ldots,v^m\}\) be a fixed finite collection of real
templates on \(\mathbb Z^d\), embedded periodically for all sufficiently
large \(L\).  Assume
\begin{equation}
 \operatorname{supp}v^a\subset\{x:|x|\leq R\},
 \qquad
 \sum_xv^a(x)=0,
 \qquad
 v^a\neq0.
 \label{eq:e8v62-template}
\end{equation}
The zero-sum condition makes every translated update preserve the
constant charge.  Put
\begin{equation}
 u^a:=Av^a,
 \qquad
 u^a_z:=\tau_zu^a,
 \qquad
 {\mathbb F}_{\cal V}
 :=
 \sum_{a=1}^m\sum_{z\in\Lambda_L}
 {u^a_z(u^a_z)^*\over\|u^a\|_2^2}.
 \label{eq:e8v62-frame}
\end{equation}
This is the first-chaos frame for rank-one conditional updates along
all translated template directions after constraint whitening.

\begin{lemma}[Exact translation-frame multiplier]
\label{lem:e8v62-multiplier}
Let
\[
 \widetilde v^a(k)
 :=
 \sum_xv^a(x)e^{-2\pi i k\cdot x/L}.
\]
Then every Fourier mode \(e_k\) is an eigenvector of
\({\mathbb F}_{\cal V}\), with eigenvalue
\begin{equation}
 f_{\cal V}(k)
 =
 \sum_{a=1}^m
 {\lambda(k)^2|\widetilde v^a(k)|^2\over\|Av^a\|_2^2},
 \label{eq:e8v62-multiplier}
\end{equation}
where \(\lambda(k)\) is given by
\eqref{eq:e8v61-lambda}.
\end{lemma}

\begin{proof}
Translation invariance makes the frame diagonal in Fourier space.  With
\(e_k(x)=|\Lambda_L|^{-1/2}e^{2\pi ik\cdot x/L}\),
\[
 \sum_z
 |\langle e_k,\tau_zu^a\rangle|^2
 =
 |\widetilde u^a(k)|^2.
\]
Since \(u^a=Av^a\),
\(\widetilde u^a(k)=\lambda(k)\widetilde v^a(k)\).
Divide by \(\|u^a\|_2^2\), sum over \(a\), and obtain
\eqref{eq:e8v62-multiplier}.
\end{proof}

\begin{theorem}[Universal sixth-order obstruction for fixed local
charge-preserving shapes]
\label{thm:e8v62-obstruction}
For templates satisfying \eqref{eq:e8v62-template}, there is a constant
\(C_{\cal V}<\infty\), independent of \(L\), such that
\begin{equation}
 \lambda_{\min}
 \left({\mathbb F}_{\cal V}\big|_{X_0}\right)
 \leq C_{\cal V}L^{-6}.
 \label{eq:e8v62-obstruction}
\end{equation}
Consequently no fixed finite family of bounded-range rank-one
charge-preserving update shapes has a volume-uniform Gaussian
heat-bath gap for the scalar constraint penalty.
\end{theorem}

\begin{proof}
Take \(k_*=(1,0,\ldots,0)\).  Since each template has zero sum,
\begin{align*}
 |\widetilde v^a(k_*)|
 &=
 \left|
 \sum_xv^a(x)
 \left(e^{-2\pi ix_1/L}-1\right)
 \right|\\
 &\leq
 {2\pi\over L}\sum_x|x_1v^a(x)|
 =:{M_a\over L}.
\end{align*}
Also
\[
 \lambda(k_*)=4\sin^2(\pi/L)\leq{4\pi^2\over L^2}.
\]
Insert both bounds in \eqref{eq:e8v62-multiplier}:
\[
 f_{\cal V}(k_*)
 \leq
 {16\pi^4\over L^4}
 \sum_a{M_a^2\over L^2\|Av^a\|_2^2}
 =C_{\cal V}L^{-6}.
\]
The least eigenvalue on \(X_0\) is at most this Rayleigh eigenvalue.
By the audited all-chaos Gaussian block theorem, the heat-bath gap
equals the least first-chaos frame eigenvalue.
\end{proof}

\begin{corollary}[The V61 pair exponent is structural]
\label{cor:e8v62-pair}
The \(L^{-6}\) pair-gap collapse in V61 is the generic low-frequency
order forced by three facts: the constraint symbol contributes
\(|k|^2\), the frame squares it, and charge preservation makes every
fixed local update template vanish at least linearly at \(k=0\).
\end{corollary}

\begin{proof}
The proof of \Cref{thm:e8v62-obstruction} gives
\[
 |\lambda(k)\widetilde v(k)|^2
 =O(|k|^4)|O(|k|^2)=O(|k|^6).
\]
V61 computes a pair family for which the resulting multiplier is
exactly proportional to \(\lambda(k)^3\), attaining this order.
\end{proof}

\begin{remark}[Higher moment cancellation is worse]
\label{rem:e8v62-moments}
If a template also has vanishing first moments, then
\(\widetilde v(k)=O(|k|^2)\), and its contribution begins at order
\(|k|^8\).  Additional local cancellations cannot repair the infrared
frame; they suppress it further.
\end{remark}

\subsection{The ideal nonlocal Fourier-shell update}
\label{sec:e8v62-Fourier}

Let \({\cal K}_L\) be a choice of real Fourier sine/cosine subspaces
covering every nonzero momentum exactly once.  Equivalently, work in
the complexification and pair \(k\) with \(-k\).  Partition the nonzero
modes into disjoint shells
\begin{equation}
 \Lambda_L\setminus\{0\}
 =
 {\cal S}_1\mathbin{\dot\cup}\cdots\mathbin{\dot\cup}{\cal S}_J.
 \label{eq:e8v62-shells}
\end{equation}
Let \(P_j\) be the orthogonal Fourier projection onto shell
\({\cal S}_j\).

\begin{theorem}[Unit Fourier-shell frame]
\label{thm:e8v62-Fourier}
For the Gaussian measure \(\gamma_{\kappa,L}\), conditional resampling
of the Fourier coefficients in shell \({\cal S}_j\), holding all other
coefficients fixed, has first-chaos difference projection \(P_j\).
Therefore
\begin{equation}
 \sum_{j=1}^JP_j=I_{X_0}
 \label{eq:e8v62-unit-frame}
\end{equation}
and the full Gaussian shell heat-bath generator has all-chaos gap
\begin{equation}
 \operatorname{gap}
 \left(\sum_{j=1}^J(I-K_j^{\rm shell})\right)=1,
 \label{eq:e8v62-unit-gap}
\end{equation}
independently of \(L\) and \(\kappa\).
\end{theorem}

\begin{proof}
The Fourier transform diagonalizes \(A\), and
\eqref{eq:e8v61-measure} becomes the product
\[
 \prod_{k\neq0}
 \exp\!\left(-{\kappa\over2}\lambda(k)^2|\widehat\phi(k)|^2\right)
 d\widehat\phi(k),
\]
with the usual reality pairing.  Conditional expectation given the
complement of a shell integrates precisely its independent Gaussian
coordinates.  On first chaos, \(I-K_j^{\rm shell}=P_j\).
The shells partition \(X_0\), proving
\eqref{eq:e8v62-unit-frame}.  The audited Gaussian all-chaos theorem
then gives a gap equal to the least eigenvalue of the frame, which is
one.
\end{proof}

\begin{firewall}[The unit frame is spatially nonlocal]
\label{fire:e8v62-nonlocal}
A Fourier shell has support across the entire spatial box.  Its exact
unit gap therefore does not satisfy a finite-range owner or Markov
locality requirement.  It is a reference preconditioner and a target
for multiscale approximation, not by itself a local continuum
construction.
\end{firewall}

\subsection{What a multiscale construction must change}
\label{sec:e8v62-multiscale}

\begin{theorem}[Necessary growth of spatial scale or update intensity]
\label{thm:e8v62-necessary}
Consider a cofinal family built from finitely many translation orbits
of charge-preserving templates.  If their support radii and normalized
first moments remain uniformly bounded and their update intensities
remain uniformly bounded, then the frame margin tends to zero at least
as \(L^{-6}\).  Hence a volume-uniform margin requires at least one of:
\begin{enumerate}[label=\textup{(\roman*)},leftmargin=3.5em]
\item template support or first moment growing with \(L\);
\item scale-dependent update intensity growing at low frequency; or
\item a nonlocal or auxiliary coarse variable which directly resamples
the low modes.
\end{enumerate}
\end{theorem}

\begin{proof}
Uniform boundedness of the number of template orbits, their normalized
moment constants \(M_a/\|Av^a\|_2\), and their intensities makes the
constant \(C_{\cal V}\) in
\eqref{eq:e8v62-obstruction} uniform.  The margin then tends to zero.
The contrapositive gives the three possible ways to escape one of the
hypotheses used in that estimate.
\end{proof}

\begin{firewall}[Accelerated rates must enter the physical clock]
\label{fire:e8v62-rates}
Multiplying a slow coarse projection by an \(L\)-dependent rate can
produce a formal generator gap, but that rate changes the dynamics and
its time normalization.  It cannot be hidden in the definition of the
Dirichlet form when the programme later compares spectral scales,
Duhamel time, reflection, or physical mass.
\end{firewall}

\begin{definition}[Multiscale constraint-frame card, MCFC]
\label{def:e8v62-MCFC}
MCFC consists of:
\begin{enumerate}[label=\textup{(F\arabic*)},leftmargin=4em]
\item a dyadic or coarse-variable family preserving the constant charge;
\item an exact first-chaos frame multiplier and a uniform lower bound;
\item explicit scale-dependent support, overlap, and update intensity;
\item a comparison with the ideal unit Fourier-shell frame;
\item an all-chaos Gaussian lift and stability under physical
conditional-density deformation;
\item a local implementation or an explicit controlled nonlocal tail;
\item a separate constraint-surface dynamics; and
\item compatibility with V45 ownership, V50 response, V58 stiffness,
reflection, boundaries, charges, and cutoff removal.
\end{enumerate}
\end{definition}

\begin{assessment}[Progress at V62]
\label{ass:e8v62-status}
The pair failure is now a finite-range no-go theorem: every fixed
bounded-range charge-preserving rank-one template family has a frame
mode of order \(L^{-6}\).  The exact Fourier-shell update has unit gap
on all chaoses but is spatially global.  A successful scalar constraint
compiler must therefore be genuinely multiscale, with all support and
rate growth declared and reconciled with the physical clock and
long-range ownership ledger.
\end{assessment}

\begin{programme}[Eighth-series V63: dyadic block-average frame]
\label{prog:e8v62-V63}
V63 will construct charge-preserving differences of neighboring dyadic
block averages.  It will compute their translation-averaged Fourier
multiplier, determine the scale weights needed to cover every nonzero
mode, and test whether those weights can be implemented with bounded
overlap per scale and an explicit physical-time normalization.
\end{programme}

\begin{firewall}[Claim boundary after eighth-series V62]
\label{fire:e8v62-boundary}
V62 proves the fixed finite-range \(L^{-6}\) obstruction and the exact
unit-gap Fourier-shell reference.  It does not yet construct a local
uniform dyadic frame, choose physically admissible scale rates, prove
nonlinear conditional transport or a surface gap, close boundary and
charge sectors, remove cutoffs, or prove the full continuum.
\end{firewall}

\begin{theorem}[Eighth-series V62 result]
\label{thm:e8v62-result}
For the linear scalar constraint penalty, every fixed finite family of
bounded-range charge-preserving translated update shapes has Gaussian
block gap at most \(C L^{-6}\).  Disjoint Fourier-shell updates instead
give the identity first-chaos frame and all-chaos gap one, exposing the
precise locality--mixing tradeoff that a multiscale construction must
resolve.
\end{theorem}

\begin{proof}
Use \Cref{thm:e8v62-obstruction,thm:e8v62-Fourier,
thm:e8v62-necessary}.
\end{proof}

\paragraph{Parent-version record.}
V62 was formed from
\texttt{Renewal\_SM\_Einstein\_Continuum\_Research\_Notes\_V61.tex}.
The V61 parent had (24\,650) logical source lines, (877\,573) bytes,
and SHA--256
\[
 \texttt{7AD356BFCC988B801706452D6C6109D17A6E5B22B266F77EDFBD6DDE0D272EF8}.
\]
The next compulsory companion audit is V65.

% =============================================================================
% END APPENDED EIGHTH-SERIES V62 MATERIAL
% =============================================================================
% =============================================================================
% BEGIN APPENDED EIGHTH-SERIES V63 MATERIAL
% =============================================================================

\clearpage
\section{Eighth-series V63: an exact dyadic block-difference frame and its clock cost}
\label{sec:v8v63}

\subsection{One-dimensional dyadic templates}
\label{sec:e8v63-templates}

Work first on the one-dimensional torus
\(\Lambda_L=\mathbb Z/L\mathbb Z\), where \(L=2^m\), \(m\geq3\).
Let
\begin{equation}
 {\cal S}_L=\{1,2,4,\ldots,L/4\}.
 \label{eq:e8v63-scales}
\end{equation}
For \(s\in{\cal S}_L\), define the charge-preserving difference of two
neighboring block sums
\begin{equation}
 v_s(x)
 =
 \mathbf1_{\{0,\ldots,s-1\}}(x)
 -
 \mathbf1_{\{s,\ldots,2s-1\}}(x).
 \label{eq:e8v63-template}
\end{equation}
Normalization by \(s^{-1}\) would not change any rank-one projection,
so block averages give the same frame.  Put \(u_s=Av_s\), where
\(A=-\Delta_L\), and sum the projections over all translations:
\begin{equation}
 {\mathbb F}_s
 :=
 \sum_{z\in\Lambda_L}
 {(\tau_zu_s)(\tau_zu_s)^*\over\|u_s\|_2^2}.
 \label{eq:e8v63-scale-frame}
\end{equation}

\begin{lemma}[Exact dyadic symbols and norms]
\label{lem:e8v63-symbol}
For \(\theta=2\pi k/L\),
\begin{align}
 \widetilde v_s(\theta)
 &=
 { (1-e^{-is\theta})^2\over1-e^{-i\theta}},
 \label{eq:e8v63-v-symbol}\\
 |\widetilde v_s(\theta)|^2
 &=
 {4\sin^4(s\theta/2)\over\sin^2(\theta/2)}.
 \label{eq:e8v63-v-absolute}
\end{align}
Moreover,
\begin{equation}
 \|Av_1\|_2^2=20,
 \qquad
 \|Av_s\|_2^2=12\quad(2\leq s\leq L/4).
 \label{eq:e8v63-norms}
\end{equation}
\end{lemma}

\begin{proof}
The geometric sum gives
\[
 \widetilde v_s(\theta)
 =
 (1-e^{-is\theta})\sum_{j=0}^{s-1}e^{-ij\theta}
 ={(1-e^{-is\theta})^2\over1-e^{-i\theta}},
\]
and taking absolute values proves
\eqref{eq:e8v63-v-absolute}.  For \(s=1\), V61 computes the norm
\(20\).  For \(s\geq2\), the vector \(Av_s\) is supported at the three
jump interfaces, with six nonzero values
\[
 -1,\ 1,\ 2,\ -2,\ -1,\ 1.
\]
The supports do not wrap or collide because \(2s\leq L/2\).
Their squared sum is \(12\).
\end{proof}

\begin{corollary}[Exact scale-frame multiplier]
\label{cor:e8v63-multiplier}
The Fourier multiplier of \({\mathbb F}_s\) is
\begin{equation}
 f_1(\theta)
 ={16\over5}\sin^6(\theta/2)
 \label{eq:e8v63-f1}
\end{equation}
for \(s=1\), and
\begin{equation}
 f_s(\theta)
 ={16\over3}\sin^2(\theta/2)\sin^4(s\theta/2)
 \label{eq:e8v63-fs}
\end{equation}
for \(2\leq s\leq L/4\).
\end{corollary}

\begin{proof}
Use the translation-frame formula
\eqref{eq:e8v62-multiplier}, the Laplacian symbol
\(\lambda(\theta)=4\sin^2(\theta/2)\), and
\Cref{lem:e8v63-symbol}.
\end{proof}

\subsection{A uniform weighted frame}
\label{sec:e8v63-weighted}

Assign scale intensity
\begin{equation}
 a_s=s^2
 \label{eq:e8v63-weights}
\end{equation}
and define
\begin{equation}
 {\mathbb F}_{\rm dyad}
 :=
 \sum_{s\in{\cal S}_L}a_s{\mathbb F}_s.
 \label{eq:e8v63-dyadic-frame}
\end{equation}

\begin{theorem}[Uniform dyadic first-chaos frame]
\label{thm:e8v63-frame}
On the mean-zero scalar space,
\begin{equation}
 {\mathbb F}_{\rm dyad}\succeq{1\over3}I.
 \label{eq:e8v63-frame}
\end{equation}
Consequently the Gaussian all-chaos heat-bath generator with translated
dyadic direction updates at rates \(a_s=s^2\) has gap at least \(1/3\),
independently of \(L\) and the penalty stiffness \(\kappa\).
\end{theorem}

\begin{proof}
By reflection, take \(\theta\in[2\pi/L,\pi]\).  If
\(\theta\geq\pi/2\), use \(s=1\).  Then
\(\sin(\theta/2)\geq2^{-1/2}\), and
\[
 a_1f_1(\theta)
 \geq {16\over5}{1\over8}={2\over5}>{1\over3}.
\]
If \(2\pi/L\leq\theta<\pi/2\), dyadic doubling supplies
\(s\in{\cal S}_L\), \(s\geq2\), such that
\begin{equation}
 {\pi\over2}\leq s\theta\leq\pi.
 \label{eq:e8v63-matched-scale}
\end{equation}
Indeed, the largest allowed scale \(L/4\) reaches the lower endpoint
when \(\theta=2\pi/L\).  Now
\[
 \sin^4(s\theta/2)\geq{1\over4},
 \qquad
 \sin(\theta/2)\geq{\theta\over\pi},
 \qquad
 s^2\theta^2\geq{\pi^2\over4}.
\]
Using \eqref{eq:e8v63-fs},
\[
 a_sf_s(\theta)
 \geq
 {16\over3}
 {s^2\theta^2\over\pi^2}{1\over4}
 \geq{1\over3}.
\]
Every other scale contributes a nonnegative operator, proving
\eqref{eq:e8v63-frame}.  The all-chaos conclusion follows from the
audited Gaussian block-frame theorem.
\end{proof}

\begin{remark}[A concrete multiscale escape from V62]
\label{rem:e8v63-escape}
The support diameter \(2s\) grows through all dyadic scales up to
\(L/2\).  Thus \Cref{thm:e8v63-frame} escapes the V62 finite-range
obstruction by the first allowed mechanism: it lets the update range
grow with the box.
\end{remark}

\subsection{The intensity cost}
\label{sec:e8v63-clock}

\begin{proposition}[Quadratic total scale intensity]
\label{prop:e8v63-intensity}
The sum of dyadic rate coefficients is
\begin{equation}
 W_L:=\sum_{s\in{\cal S}_L}a_s
 =\sum_{\ell=0}^{m-2}4^\ell
 ={L^2\over12}-{1\over3}.
 \label{eq:e8v63-intensity}
\end{equation}
\end{proposition}

\begin{proof}
Sum the finite geometric series and use \(L=2^m\).
\end{proof}

\begin{theorem}[Unit-total-intensity normalization restores a slow mode]
\label{thm:e8v63-clock}
Let
\begin{equation}
 \overline{\mathbb F}_{\rm dyad}
 :=
 W_L^{-1}{\mathbb F}_{\rm dyad}.
 \label{eq:e8v63-normalized-frame}
\end{equation}
There is a numerical constant \(C<\infty\), independent of \(L\), such
that
\begin{equation}
 \lambda_{\min}
 (\overline{\mathbb F}_{\rm dyad}|_{X_0})
 \leq {C\over L^2}.
 \label{eq:e8v63-clock}
\end{equation}
\end{theorem}

\begin{proof}
Evaluate the frame at the least mode
\(\theta_*=2\pi/L\).  For \(s\geq2\), use
\(\sin x\leq x\) in \eqref{eq:e8v63-fs}:
\[
 a_sf_s(\theta_*)
 \leq
 {16\over3}s^2{\theta_*^2\over4}
 {s^4\theta_*^4\over16}
 ={1\over12}s^6\theta_*^6.
\]
The dyadic sum of \(s^6\) through \(L/4\) is at most a fixed constant
times \(L^6\), so these terms have a total bounded independently of
\(L\).  The \(s=1\) term is also bounded.  Hence
\[
 \langle e_{k_*},{\mathbb F}_{\rm dyad}e_{k_*}\rangle
 \leq C_0.
\]
By \eqref{eq:e8v63-intensity},
\(W_L\geq cL^2\) for \(L\geq8\).  Dividing gives
\eqref{eq:e8v63-clock}.
\end{proof}

\begin{firewall}[A uniform gap was purchased with a faster clock]
\label{fire:e8v63-clock}
The unnormalized lower bound \eqref{eq:e8v63-frame} uses rates
\(s^2\), whose total scale intensity is \(O(L^2)\).  If the physical
comparison fixes a unit aggregate scale clock, the least mode collapses
as in \eqref{eq:e8v63-clock}.  Therefore the uniform frame cannot be
reported without also declaring which clock, per-site exposure, and
block cost enter the Duhamel and mass scales.
\end{firewall}

\begin{firewall}[The coarsest update is macroscopic]
\label{fire:e8v63-macroscopic}
At \(s=L/4\), the direction \(v_s\) occupies half the spatial torus.
Although each template is a difference of adjacent block averages, the
coarsest update is macroscopic.  Nonlinear physical transport must
control its long-range conditional density and its compatibility with
the V45 boundary ownership ledger.
\end{firewall}

\subsection{Dyadic constraint-frame card}
\label{sec:e8v63-card}

\begin{definition}[Dyadic clock and locality card, DCLC]
\label{def:e8v63-DCLC}
DCLC consists of:
\begin{enumerate}[label=\textup{(Y\arabic*)},leftmargin=4em]
\item the exact dyadic multiplier
\eqref{eq:e8v63-f1}--\eqref{eq:e8v63-fs};
\item a declared scale-rate sequence and its physical-time
normalization;
\item per-site incidence and computational cost at every scale;
\item a long-range owner/boundary estimate for macroscopic blocks;
\item nonlinear conditional-projection transport with a loss below the
frame margin;
\item a constraint-surface dynamics and positive carrier;
\item extension from the scalar model to tensor, gauge, ghost, and
Standard Model sectors; and
\item compatibility with reflection, charges, observables, and cutoff
removal.
\end{enumerate}
\end{definition}

\begin{assessment}[Progress at V63]
\label{ass:e8v63-status}
An explicit real-space multiscale family now closes the Gaussian
first-chaos and all-chaos frame with margin \(1/3\).  The exact
calculation also exposes its price: the required dyadic intensity grows
as \(L^2/12\), and unit-total-intensity normalization leaves an
\(O(L^{-2})\) slow mode.  The remaining issue is therefore physical
clock and long-range implementation, not existence of a multiscale
frame in abstract Gaussian \(L^2\).
\end{assessment}

\begin{programme}[Eighth-series V64: per-site clock and randomized tilings]
\label{prog:e8v63-V64}
V64 will replace all translations at every scale by randomized
nonoverlapping dyadic tilings.  It will compute the averaged frame and
per-site incidence exactly, then determine whether any rate choice can
retain a uniform low-mode margin with bounded per-site clock.
\end{programme}

\begin{firewall}[Claim boundary after eighth-series V63]
\label{fire:e8v63-boundary}
V63 proves the exact dyadic multiplier, an unnormalized uniform
Gaussian gap, and the quadratic scale-clock cost with its normalized
slow mode.  It does not construct a bounded-clock local implementation,
transport the frame to the nonlinear physical measure, prove the
surface gap or positive carrier, close boundaries and charges, remove
cutoffs, or prove the full continuum.
\end{firewall}

\begin{theorem}[Eighth-series V63 result]
\label{thm:e8v63-result}
Translated differences of neighboring dyadic block averages, weighted
at scale \(s\) by \(s^2\), give a Gaussian scalar-constraint frame
\[
 {\mathbb F}_{\rm dyad}\succeq{1\over3}I
\]
uniformly in volume and penalty stiffness.  Their total scale intensity
is \(L^2/12-1/3\), and normalizing that intensity to one produces a
least frame mode bounded above by \(CL^{-2}\).
\end{theorem}

\begin{proof}
Use \Cref{thm:e8v63-frame,prop:e8v63-intensity,
thm:e8v63-clock}.
\end{proof}

\paragraph{Parent-version record.}
V63 was formed from
\texttt{Renewal\_SM\_Einstein\_Continuum\_Research\_Notes\_V62.tex}.
The V62 parent had (24\,968) logical source lines, (888\,886) bytes,
and SHA--256
\[
 \texttt{7B7CEF151513FECF2655911C0D423B8443053F4DFC1AD8D96AB02DD1DC6EF0D9}.
\]
The next compulsory companion audit is V65.

% =============================================================================
% END APPENDED EIGHTH-SERIES V63 MATERIAL
% =============================================================================
% =============================================================================
% BEGIN APPENDED EIGHTH-SERIES V64 MATERIAL
% =============================================================================

\clearpage
\section{Eighth-series V64: randomized dyadic tilings and the per-site clock}
\label{sec:v8v64}

\subsection{Nonoverlapping tilings}
\label{sec:e8v64-tilings}

Continue in the one-dimensional scalar model of V63 with
\(L=2^m\) and \(s\in{\cal S}_L\).  For a shift
\(r\in\{0,\ldots,2s-1\}\), define the tiling
\begin{equation}
 {\cal B}_{s,r}
 :=
 \left\{
 \{r+2sj,\ldots,r+2sj+2s-1\}:
 j=0,\ldots,{L\over2s}-1
 \right\},
 \label{eq:e8v64-tiling}
\end{equation}
with positions read modulo \(L\).  Each block carries the translated
charge-preserving direction \(\tau_{r+2sj}v_s\).  In whitened
coordinates put
\begin{equation}
 {\mathbb F}_{s,r}^{\rm tile}
 :=
 \sum_{j=0}^{L/(2s)-1}
 {(\tau_{r+2sj}u_s)(\tau_{r+2sj}u_s)^*
 \over\|u_s\|_2^2}.
 \label{eq:e8v64-tile-frame}
\end{equation}

\begin{proposition}[Exact randomized-tiling average]
\label{prop:e8v64-average}
Uniform averaging over the \(2s\) shifts gives
\begin{equation}
 {1\over2s}\sum_{r=0}^{2s-1}
 {\mathbb F}_{s,r}^{\rm tile}
 =
 {1\over2s}{\mathbb F}_s,
 \label{eq:e8v64-average}
\end{equation}
where \({\mathbb F}_s\) is the all-translation frame of V63.
\end{proposition}

\begin{proof}
Every \(z\in\Lambda_L\) has a unique representation
\[
 z=r+2sj\pmod L,
 \qquad
 0\leq r<2s,\quad0\leq j<{L\over2s}.
\]
Therefore the double sum over \(r,j\) contains every translated
rank-one projection in \({\mathbb F}_s\) exactly once.  Divide by
\(2s\).
\end{proof}

\begin{proposition}[Exact per-site incidence]
\label{prop:e8v64-incidence}
For fixed \(s,r\), the supports of the directions in
\eqref{eq:e8v64-tile-frame} partition \(\Lambda_L\).  Thus every site
belongs to exactly one scale-\(s\) block.  If the tiling generator at
scale \(s\) is run with coefficient \(b_s\), its support-incidence rate
at each site is \(b_s\).
\end{proposition}

\begin{proof}
The intervals in \eqref{eq:e8v64-tiling} are consecutive, disjoint, and
their total cardinality is
\[
 {L\over2s}\,2s=L.
\]
Multiplying every block term at that scale by \(b_s\) gives the stated
incidence.
\end{proof}

\subsection{Weighted randomized generator}
\label{sec:e8v64-generator}

For nonnegative scale rates \(b_s\), define the averaged first-chaos
frame
\begin{equation}
 {\mathbb G}_{\boldsymbol b}
 :=
 \sum_{s\in{\cal S}_L}
 {b_s\over2s}{\mathbb F}_s.
 \label{eq:e8v64-weighted-frame}
\end{equation}
Its declared per-site clock is
\begin{equation}
 B_L(\boldsymbol b):=\sum_{s\in{\cal S}_L}b_s.
 \label{eq:e8v64-site-clock}
\end{equation}

\begin{lemma}[Weighted Gaussian all-chaos lift]
\label{lem:e8v64-all-chaos}
For nonnegative coefficients, the all-chaos gap of the averaged
Gaussian heat-bath generator equals
\begin{equation}
 \lambda_{\min}
 \left({\mathbb G}_{\boldsymbol b}\big|_{X_0}\right).
 \label{eq:e8v64-all-chaos}
\end{equation}
\end{lemma}

\begin{proof}
On first chaos the generator is exactly
\({\mathbb G}_{\boldsymbol b}\), so the full gap is no larger than its
least eigenvalue.  On the \(n\)-th Wiener chaos, each conditional
expectation is the symmetric restriction of \(Q^{\otimes n}\) for a
first-chaos orthogonal projection \(Q\).  The audited YM V16 inequality
\[
 I-Q^{\otimes n}
 \succeq
 {1\over n}\sum_{\ell=1}^n
 I^{\otimes(\ell-1)}\otimes(I-Q)
 \otimes I^{\otimes(n-\ell)}
\]
remains valid after multiplication by any nonnegative coefficient.
Sum all weighted block terms.  On each tensor position the resulting
first-chaos frame is at least its least eigenvalue, so every nonzero
chaos has the same lower bound.  Equality follows from first chaos.
\end{proof}

\subsection{A uniform frame and its exact cubic clock}
\label{sec:e8v64-uniform}

\begin{theorem}[Randomized tiling rates that recover the V63 frame]
\label{thm:e8v64-uniform}
Choose
\begin{equation}
 b_s=2s^3.
 \label{eq:e8v64-rates}
\end{equation}
Then
\begin{equation}
 {\mathbb G}_{\boldsymbol b}
 =
 \sum_s s^2{\mathbb F}_s
 ={\mathbb F}_{\rm dyad}
 \succeq{1\over3}I_{X_0}.
 \label{eq:e8v64-uniform}
\end{equation}
The corresponding Gaussian all-chaos gap is at least \(1/3\), while
the per-site clock is exactly
\begin{equation}
 B_L
 =
 \sum_{\ell=0}^{m-2}2(2^\ell)^3
 ={L^3\over28}-{2\over7}.
 \label{eq:e8v64-cubic-clock}
\end{equation}
\end{theorem}

\begin{proof}
Equation \eqref{eq:e8v64-rates} gives
\(b_s/(2s)=s^2\), so
\eqref{eq:e8v64-uniform} is exactly
\Cref{thm:e8v63-frame}.  The gap follows from
\Cref{lem:e8v64-all-chaos}.  Finally,
\[
 \sum_{\ell=0}^{m-2}2\,8^\ell
 ={2(8^{m-1}-1)\over7}
 ={L^3\over28}-{2\over7}.
\]
\end{proof}

\begin{remark}[Why random tiling costs an additional factor \(s\)]
\label{rem:e8v64-factor}
At scale \(s\), a particular translated direction appears in exactly
one of the \(2s\) tiling shifts.  Randomizing the shift therefore
reduces its frame intensity by \(1/(2s)\).  Recovering the V63
all-translation coefficient \(s^2\) requires tiling rate \(2s^3\).
\end{remark}

\subsection{A matching necessary clock order}
\label{sec:e8v64-necessary}

\begin{theorem}[Bounded per-site clock forces a cubic slow mode]
\label{thm:e8v64-necessary}
There is a numerical constant \(C<\infty\) such that, for every
nonnegative rate sequence \(\boldsymbol b\),
\begin{equation}
 \lambda_{\min}
 \left({\mathbb G}_{\boldsymbol b}\big|_{X_0}\right)
 \leq
 {C\,B_L(\boldsymbol b)\over L^3}.
 \label{eq:e8v64-necessary}
\end{equation}
Consequently a frame margin at least \(\delta>0\) requires
\begin{equation}
 B_L(\boldsymbol b)\geq{\delta\over C}L^3.
 \label{eq:e8v64-clock-lower}
\end{equation}
The cubic clock order in \eqref{eq:e8v64-cubic-clock} is therefore
optimal for this randomized-tiling family up to constants.
\end{theorem}

\begin{proof}
Evaluate the Fourier multiplier at
\(\theta_*=2\pi/L\).  For \(s\geq2\),
\eqref{eq:e8v63-fs} and \(\sin x\leq x\) give
\[
 f_s(\theta_*)
 \leq {1\over12}s^4\theta_*^6.
\]
Hence
\[
 {b_s\over2s}f_s(\theta_*)
 \leq {b_s\over24}s^3\theta_*^6
 \leq {b_s\over24}\left({L\over4}\right)^3
 \left({2\pi\over L}\right)^6
 \leq {C_1b_s\over L^3}.
\]
For \(s=1\),
\[
 {b_1\over2}f_1(\theta_*)
 \leq C_2b_1L^{-6}
 \leq C_2b_1L^{-3}.
\]
Sum over scales:
\[
 \langle e_{k_*},{\mathbb G}_{\boldsymbol b}e_{k_*}\rangle
 \leq {C\over L^3}\sum_sb_s.
\]
The least eigenvalue is at most this eigenvalue, proving
\eqref{eq:e8v64-necessary}.  Rearrangement gives
\eqref{eq:e8v64-clock-lower}; compare with
\eqref{eq:e8v64-cubic-clock}.
\end{proof}

\begin{corollary}[Unit per-site clock has vanishing gap]
\label{cor:e8v64-unit-clock}
If \(B_L(\boldsymbol b)=1\), then the all-chaos Gaussian gap is at most
\begin{equation}
 CL^{-3}.
 \label{eq:e8v64-unit-clock}
\end{equation}
\end{corollary}

\begin{proof}
Combine \Cref{thm:e8v64-necessary,lem:e8v64-all-chaos}.
\end{proof}

\begin{firewall}[Randomization does not remove the clock cost]
\label{fire:e8v64-randomization}
Random tilings reduce simultaneous overlap: every site lies in one
block per scale.  They do not create low-frequency frame strength for
free.  The probability \(1/(2s)\) of selecting a given translate moves
the required compensation into the scale rate, producing the cubic
per-site clock proved above.
\end{firewall}

\begin{firewall}[Generator time is not physical mass]
\label{fire:e8v64-physical}
The lower bound \(1/3\) in
\eqref{eq:e8v64-uniform} is measured in a generator whose per-site clock
grows as \(L^3\).  It cannot be identified with a volume-uniform
physical mass or Duhamel decay rate.  With a fixed per-site time
normalization, \Cref{cor:e8v64-unit-clock} is the relevant conclusion.
\end{firewall}

\subsection{Clock-aware multiscale card}
\label{sec:e8v64-card}

\begin{definition}[Randomized tiling clock card, RTCC]
\label{def:e8v64-RTCC}
RTCC consists of:
\begin{enumerate}[label=\textup{(T\arabic*)},leftmargin=4em]
\item the exact shift average \eqref{eq:e8v64-average};
\item per-site incidence and the clock
\eqref{eq:e8v64-site-clock};
\item a frame margin stated in the same time normalization used by the
Duhamel and spectral comparisons;
\item an implementation cost for macroscopic conditional updates;
\item nonlinear physical-measure transport below the normalized frame
margin;
\item a constraint-surface dynamics with its own clock and gap;
\item boundary ownership and long-range influence estimates; and
\item reflection, charge, observable, and cutoff-removal contacts.
\end{enumerate}
\end{definition}

\begin{assessment}[Progress at V64]
\label{ass:e8v64-status}
Randomized nonoverlapping tilings have now been computed exactly.  They
recover the uniform dyadic Gaussian frame only at a per-site rate of
order \(L^3\), and every rate choice in this family obeys the matching
upper bound
\(\operatorname{gap}\leq CB_L/L^3\).  Thus bounded-overlap
randomization does not solve the physical-clock problem.  A viable next
route must change the dynamics, the topology being gapped, or the role
of the scalar constraint response.
\end{assessment}

\begin{programme}[Eighth-series V65: companion audit and clock-aware route]
\label{prog:e8v64-V65}
V65 will perform the compulsory YM/NS audit.  It will compare any new
physical-measure projection transport with the cubic clock obstruction
and choose whether to pursue a preconditioned nonlocal scalar response,
a gap only in the V44 observable topology, or direct elimination of the
constraint-normal field.
\end{programme}

\begin{firewall}[Claim boundary after eighth-series V64]
\label{fire:e8v64-boundary}
V64 proves the exact randomized-tiling average, per-site incidence,
uniform unnormalized frame, cubic clock cost, and matching necessary
clock order.  It does not construct a bounded-clock uniform scalar
dynamics, transport to the nonlinear physical measure, prove the
surface carrier, close boundaries or charges, remove cutoffs, or prove
the full continuum.
\end{firewall}

\begin{theorem}[Eighth-series V64 result]
\label{thm:e8v64-result}
Randomized nonoverlapping dyadic tilings at rates \(b_s=2s^3\) give the
uniform Gaussian scalar frame
\({\mathbb G}_{\boldsymbol b}\succeq I/3\), but their exact per-site
clock is \(L^3/28-2/7\).  For every nonnegative rate choice,
\[
 \operatorname{gap}
 \leq {C\sum_sb_s\over L^3}.
\]
Hence a bounded per-site clock forces the Gaussian scalar normal gap to
vanish.
\end{theorem}

\begin{proof}
Use \Cref{thm:e8v64-uniform,thm:e8v64-necessary,
lem:e8v64-all-chaos}.
\end{proof}

\paragraph{Parent-version record.}
V64 was formed from
\texttt{Renewal\_SM\_Einstein\_Continuum\_Research\_Notes\_V63.tex}.
The V63 parent had (25\,305) logical source lines, (899\,300) bytes,
and SHA--256
\[
 \texttt{666801AFFF6EC6D0D155B016B52FA06F825BF85B4A5A4A98343C4C2009042CD2}.
\]
The next compulsory companion audit is V65.

% =============================================================================
% END APPENDED EIGHTH-SERIES V64 MATERIAL
% =============================================================================
% =============================================================================
% BEGIN APPENDED EIGHTH-SERIES V65 MATERIAL
% =============================================================================

\clearpage
\section{Eighth-series V65: compulsory audit and the observable-restricted route}
\label{sec:v8v65}

\subsection{Companion audit record}
\label{sec:e8v65-audit}

The compulsory V65 audit inspected
\begin{align}
 &\texttt{Renewal\_Yang\_Mills\_Mass\_Gap\_Research\_Notes\_V20.tex},
 \notag\\
 &\hspace{2em}7679\ \text{logical lines},\quad279890\ \text{bytes},
 \quad\operatorname{SHA256}
 =\texttt{FB5903451D8A22637A60651ED87D8DD047A0B8E7CCD3D2C8745AB8D2DCCA4F31},
 \label{eq:e8v65-YM-file}\\
 &\texttt{Renewal\_NS\_Stability\_Research\_Notes\_V26.tex},
 \notag\\
 &\hspace{2em}5319\ \text{logical lines},\quad278283\ \text{bytes},
 \quad\operatorname{SHA256}
 =\texttt{82C887F8C2C69E4F28FAD7C3DC8DE5B9099CB5A4BF431EBA1FFF3E91935978AD}.
 \label{eq:e8v65-NS-file}
\end{align}
Each file has one live terminator.  The V20 YM file is byte-for-byte
identical to the V19 file also present in the shared folder, so the
newest mathematical material is the V19 chapter.  Neither companion
file was modified or removed.

\begin{audit}[Yang--Mills V17--V18]
\label{audit:e8v65-YM-flow}
V17 bounds conditional-projection perturbations by conditional density
ratios and gives an explicit gap-loss compiler, but its full
\(L^\infty\) oscillation condition is too strong in the physical
scaling.  V18 replaces the global ratio by an intrinsic unitary spectral
flow.  The derivative of each transported conditional projection is an
off-diagonal centered-score operator.  A square-root eigenbranch bound
uses only the active score matrix element, while spectral-edge isolation
remains an explicit hypothesis.
\end{audit}

\begin{audit}[Yang--Mills V19 type correction]
\label{audit:e8v65-YM-type}
V19 proves that a coupling interpolation produces a scalar action-value
fluctuation, whereas curvature-adjoint polar cancellation controls a
coordinate derivative.  Those objects have different type and
different Taylor order.  A configuration-space continuity equation
creates the type-correct directional derivative and moves its score to
the test.  However, constructing a generic transport velocity by a
weighted Poisson equation has optimal norm equal to the inverse square
root of the missing Poincaré gap.  Generic Moser transport is therefore
gap-equivalent and cannot prove that gap.
\end{audit}

\begin{audit}[Navier--Stokes status]
\label{audit:e8v65-NS}
The NS V26 programme remains unchanged.  It supplies no new mechanism
that alters the scalar clock obstruction or the constraint-source
factorization established in V49--V64.
\end{audit}

\subsection{Consequences for the scalar constraint programme}
\label{sec:e8v65-consequences}

\begin{theorem}[A generic inverse transport is circular here]
\label{thm:e8v65-circular}
Let \({\cal A}_{L}\) be a bounded-per-site-clock scalar block generator
from the randomized dyadic family, with gap \(g_L\).  A minimum-norm
solution of
\begin{equation}
 -\operatorname{div}_{\mu_L}X=h,
 \qquad \langle h\rangle_{\mu_L}=0,
 \label{eq:e8v65-divergence}
\end{equation}
has solution-operator norm \(g_L^{-1/2}\).  Since V64 gives
\(g_L\leq CB_LL^{-3}\), bounded \(B_L\) permits the inverse norm to grow
at least as
\begin{equation}
 g_L^{-1/2}\geq cL^{3/2}.
 \label{eq:e8v65-inverse-growth}
\end{equation}
Thus a generic Poisson/Moser preconditioner cannot establish a
volume-uniform scalar compiler.
\end{theorem}

\begin{proof}
The audited YM V19 minimum-norm divergence theorem is a Hilbert-space
identity: the squared solution norm is
\(\langle h,{\cal A}_L^{-1}h\rangle\), so the optimal operator norm is
\(g_L^{-1/2}\).  Insert the V64 upper bound
\(g_L\leq CB_LL^{-3}\).  If \(B_L\leq B_*\), then
\[
 g_L^{-1/2}\geq C^{-1/2}B_*^{-1/2}L^{3/2}.
\]
\end{proof}

\begin{firewall}[Type-correct polar cancellation]
\label{fire:e8v65-type}
The V50 constraint-adjoint cancellation acts on a spatial divergence
source or a directional derivative after integration by parts.  It
does not turn an arbitrary scalar action-value fluctuation into an
adjoint vector.  Any future spectral flow must derive a
configuration-space continuity equation first, as in the audited YM
correction.
\end{firewall}

\begin{theorem}[Audit-selected route]
\label{thm:e8v65-route}
The next attack should restrict the response estimate to the fixed
V44 spacetime-smoothed observable space, rather than seek a
bounded-clock full-space scalar gap by generic inversion.  The required
estimate must:
\begin{enumerate}[label=\textup{(\roman*)},leftmargin=3.5em]
\item keep the slow spectral projection explicit;
\item obtain cancellation from the source or differentiated test before
applying an inverse;
\item retain the constant and ADM charge sector separately; and
\item use direct constraint elimination with its V45 long-range
ownership ledger.
\end{enumerate}
\end{theorem}

\begin{proof}
\Cref{thm:e8v65-circular} excludes the generic inverse-transport route
at bounded physical clock.  V43 proves that response pairings can be
bounded by complementary source and test spectral orders even when a
full response norm diverges, and V44 fixes a spacetime-smoothed test
class for which contact operators are controlled.  V49--V51 provide
direct spatial-divergence factorizations and explicit charge
separation.  Combining these existing interfaces yields the four
requirements without assuming a full scalar mixing gap.
\end{proof}

\begin{firewall}[Restricted observables do not prove a full gap]
\label{fire:e8v65-restricted}
Uniform response on a declared test space is weaker than a spectral gap
on the full carrier Hilbert space.  It can support convergence of the
selected observable algebra, but it cannot be reported as a positive
Hamiltonian mass gap or as uniform mixing of all constraint-normal
fluctuations.
\end{firewall}

\subsection{Observable-restricted scalar card}
\label{sec:e8v65-card}

\begin{definition}[Observable-restricted constraint card, ORCC]
\label{def:e8v65-ORCC}
ORCC consists of:
\begin{enumerate}[label=\textup{(O\arabic*)},leftmargin=4em]
\item the bounded-clock scalar generator and its explicit slow spectral
projection;
\item the fixed V44 spacetime-smoothed test space, chosen independently
of the source;
\item a source-side divergence, Ward, conservation, or multiplier
factorization obtained without inverse construction;
\item a differentiated-test estimate on the slow subspace;
\item direct constraint elimination and the V45 long-range
owner/boundary compiler;
\item separate constant, harmonic, boundary, and ADM charge sectors;
\item quantum contact, gauge, ghost, anomaly, and reflection control;
and
\item proof that the resulting restricted Schwinger functions form the
desired continuum observable algebra.
\end{enumerate}
\end{definition}

\begin{assessment}[Progress at V65]
\label{ass:e8v65-status}
The audit prevents two circular moves: inserting a scalar coupling score
into a vector polar identity and constructing a transport through the
inverse missing gap.  In light of the exact V64 clock obstruction, the
programme now targets observable-restricted scalar response with direct
source/test cancellation.  Full scalar mixing, surface positivity, and
the global charge remain separate claims.
\end{assessment}

\begin{programme}[Eighth-series V66: exact slow-mode pairing]
\label{prog:e8v65-V66}
V66 will introduce a low-mode projection for the periodic scalar
constraint and prove exact source--test pairing bounds.  It will compute
the moment cancellation required of a spatially localized smoothed test
to offset the inverse-Laplacian response, and will distinguish fixed
physical support from volume-growing test families.
\end{programme}

\begin{firewall}[Claim boundary after eighth-series V65]
\label{fire:e8v65-boundary}
V65 records the mandatory audit, transfers the generic-transport
circularity theorem to the scalar clock obstruction, and selects the
observable-restricted route.  It does not yet prove the slow-mode test
bound, build the continuum observable algebra, establish a full scalar
or surface gap, construct the positive carrier, remove cutoffs, or prove
the full continuum.
\end{firewall}

\begin{theorem}[Eighth-series V65 result]
\label{thm:e8v65-result}
At bounded per-site clock, the randomized dyadic scalar gap permits a
generic inverse-transport norm growing at least as \(L^{3/2}\).  Such a
transport is circular.  The next non-circular target is the scalar
response paired only with the fixed spacetime-smoothed observable
space, using source or test cancellation before inversion and retaining
all charge sectors.
\end{theorem}

\begin{proof}
Use \Cref{thm:e8v65-circular,thm:e8v65-route}.
\end{proof}

\paragraph{Parent-version record.}
V65 was formed from
\texttt{Renewal\_SM\_Einstein\_Continuum\_Research\_Notes\_V64.tex}.
The V64 parent had (25\,654) logical source lines, (910\,345) bytes,
and SHA--256
\[
 \texttt{DF66771137BAFCAB113A1B53CFDC1713CEB042B69314AD596D7644B48FF78806}.
\]
The next compulsory companion audit is V70.

% =============================================================================
% END APPENDED EIGHTH-SERIES V65 MATERIAL
% =============================================================================
% =============================================================================
% BEGIN APPENDED EIGHTH-SERIES V66 MATERIAL
% =============================================================================

\clearpage
\section{Eighth-series V66: local observable pairings across the scalar infrared tail}
\label{sec:v8v66}

\subsection{Periodic negative Sobolev norm}
\label{sec:e8v66-Hminus}

Let \(\Omega_L=(\mathbb R/L\mathbb Z)^d\), with physical Fourier
momenta
\[
 {\cal K}_L={2\pi\over L}\mathbb Z^d.
\]
For \(f\in L^2(\Omega_L)\), use
\[
 \widehat f(k)=\int_{\Omega_L}f(x)e^{-ik\cdot x}\,dx,
 \qquad
 \|f\|_2^2={1\over L^d}\sum_{k\in{\cal K}_L}|\widehat f(k)|^2.
\]
Define
\begin{equation}
 \|f\|_{\dot H^{-1}_L}^2
 :=
 {1\over L^d}\sum_{\substack{k\in{\cal K}_L\\k\neq0}}
 {|\widehat f(k)|^2\over|k|^2}.
 \label{eq:e8v66-Hminus}
\end{equation}
The zero mode is omitted and remains in the charge sector.

\begin{lemma}[Uniform low-momentum lattice sum in \(d>2\)]
\label{lem:e8v66-lattice-sum}
For every \(d>2\), there is \(C_d<\infty\) such that for all \(L\geq1\),
\begin{equation}
 {1\over L^d}
 \sum_{\substack{k\in{\cal K}_L\\0<|k|\leq1}}
 {1\over|k|^2}
 \leq C_d.
 \label{eq:e8v66-lattice-sum}
\end{equation}
\end{lemma}

\begin{proof}
Write \(k=2\pi n/L\).  The left side is
\[
 {L^{2-d}\over4\pi^2}
 \sum_{0<|n|\leq L/(2\pi)}{|n|^{-2}}.
\]
The number of integer points in
\(\{r\leq|n|<r+1\}\) is at most
\(C'_d(r+1)^{d-1}\).  Hence the sum through radius \(R\) is at most
\[
 C''_d\sum_{r=1}^{\lceil R\rceil}r^{d-3}
 \leq C'''_d(1+R^{d-2}).
\]
With \(R=L/(2\pi)\), multiplication by \(L^{2-d}\) is uniformly
bounded.  The finitely many \(L\) below \(2\pi\) are absorbed in the
constant.
\end{proof}

\begin{theorem}[Fixed local tests are uniformly \(H^{-1}\) in
physical dimension]
\label{thm:e8v66-local-test}
Let \(d>2\).  Suppose \(f_L\) is supported in a fixed cube embedded in
\(\Omega_L\), and
\begin{equation}
 \sup_L\bigl(\|f_L\|_1+\|f_L\|_2\bigr)\leq M_f.
 \label{eq:e8v66-local-hyp}
\end{equation}
Then
\begin{equation}
 \sup_L\|f_L\|_{\dot H^{-1}_L}
 \leq C_dM_f.
 \label{eq:e8v66-local-test}
\end{equation}
No zero-mean condition on \(f_L\) is needed when \(d>2\).
\end{theorem}

\begin{proof}
For \(0<|k|\leq1\),
\(|\widehat f_L(k)|\leq\|f_L\|_1\).  Therefore the low-momentum part of
\eqref{eq:e8v66-Hminus} is at most
\(C_d\|f_L\|_1^2\) by
\Cref{lem:e8v66-lattice-sum}.  For \(|k|>1\), discard the factor
\(|k|^{-2}\leq1\) and use Parseval:
\[
 {1\over L^d}\sum_{|k|>1}
 {|\widehat f_L(k)|^2\over|k|^2}
 \leq\|f_L\|_2^2.
\]
Add both estimates and take a square root.
\end{proof}

\begin{theorem}[Moment cancellation covers \(d=1,2\)]
\label{thm:e8v66-moment}
Let \(d\geq1\), assume \(f_L\) is supported in a fixed coordinate cube,
and suppose
\begin{equation}
 \int_{\Omega_L}f_L(x)\,dx=0,
 \qquad
 \int_{\Omega_L}|x|\,|f_L(x)|\,dx\leq M_1,
 \qquad
 \|f_L\|_2\leq M_2.
 \label{eq:e8v66-moment-hyp}
\end{equation}
Then
\begin{equation}
 \sup_L\|f_L\|_{\dot H^{-1}_L}
 \leq C_d(M_1+M_2).
 \label{eq:e8v66-moment}
\end{equation}
\end{theorem}

\begin{proof}
Zero mean gives
\[
 \widehat f_L(k)
 =
 \int f_L(x)(e^{-ik\cdot x}-1)\,dx,
\]
so
\[
 |\widehat f_L(k)|\leq|k|M_1.
\]
For \(0<|k|\leq1\), every summand in
\eqref{eq:e8v66-Hminus} is therefore at most \(M_1^2\).  The normalized
number \(L^{-d}\#\{k:|k|\leq1\}\) is bounded uniformly in \(L\).  The
high-momentum part is at most \(M_2^2\) by Parseval, as in
\Cref{thm:e8v66-local-test}.
\end{proof}

\subsection{Exact pairing with a divergence source}
\label{sec:e8v66-pairing}

\begin{theorem}[Uniform local potential pairing in \(d>2\)]
\label{thm:e8v66-pairing}
Let the mean-zero scalar source have the direct factorization
\begin{equation}
 \rho=-\operatorname{div}J
 \label{eq:e8v66-source}
\end{equation}
on \(\Omega_L\), and let
\(u=(-\Delta)^\dagger\rho\).  Then for every test \(f\),
\begin{equation}
 |\langle f,u\rangle|
 \leq
 \|f\|_{\dot H^{-1}_L}\|J\|_2.
 \label{eq:e8v66-pairing}
\end{equation}
Consequently, the local tests of
\Cref{thm:e8v66-local-test} obey
\begin{equation}
 |\langle f_L,u_L\rangle|
 \leq C_dM_f\|J_L\|_2
 \label{eq:e8v66-local-pairing}
\end{equation}
uniformly in physical volume.
\end{theorem}

\begin{proof}
On nonzero Fourier modes,
\[
 \widehat\rho(k)=-ik\cdot\widehat J(k).
\]
Hence
\[
 \|\rho\|_{\dot H^{-1}_L}^2
 ={1\over L^d}\sum_{k\neq0}
 {|\widehat\rho(k)|^2\over|k|^2}
 \leq
 {1\over L^d}\sum_k|\widehat J(k)|^2
 =\|J\|_2^2.
\]
Self-adjoint spectral calculus gives
\[
 \langle f,(-\Delta)^\dagger\rho\rangle
 =
 \left\langle
 (-\Delta)^{-1/2}P_0f,
 (-\Delta)^{-1/2}\rho
 \right\rangle.
\]
Cauchy--Schwarz proves \eqref{eq:e8v66-pairing}; then apply
\Cref{thm:e8v66-local-test}.
\end{proof}

\begin{corollary}[Double-divergence bound in every dimension]
\label{cor:e8v66-double}
If in addition
\begin{equation}
 f=-\operatorname{div}K
 \label{eq:e8v66-test-divergence}
\end{equation}
with directly constructed \(K\in L^2\), then
\begin{equation}
 |\langle f,(-\Delta)^\dagger\rho\rangle|
 \leq\|K\|_2\|J\|_2
 \label{eq:e8v66-double}
\end{equation}
in every spatial dimension, with constant one.
\end{corollary}

\begin{proof}
Both
\((-\Delta)^{-1/2}(-\operatorname{div})K\) and the corresponding
operator applied to \(J\) are adjoints of the gradient polar isometry
from V50 and have norm at most one.  Pair them and use
Cauchy--Schwarz.
\end{proof}

\subsection{A local divergence primitive for zero-mean tests}
\label{sec:e8v66-primitive}

\begin{theorem}[Rectangular local primitive]
\label{thm:e8v66-primitive}
Let \(Q=(0,R)^d\) and \(f\in L^2(Q)\) satisfy \(\int_Qf=0\).
There is \(K\in L^2(Q;\mathbb R^d)\) with vanishing normal trace,
\begin{equation}
 \operatorname{div}K=f,
 \qquad
 \|K\|_2\leq C_dR\|f\|_2.
 \label{eq:e8v66-primitive}
\end{equation}
Extending \(K\) by zero gives the same distributional identity in any
larger periodic box containing \(Q\).
\end{theorem}

\begin{proof}
Proceed by induction on \(d\).  In one dimension set
\[
 K_1(x)=\int_0^xf(t)\,dt.
\]
It vanishes at both endpoints and
\(\|K_1\|_2\leq R\|f\|_2\).

For \(d>1\), write \(x=(x_1,x')\) and put
\[
 \bar f(x')=\int_0^Rf(t,x')\,dt,
\qquad
 K_1(x)=\int_0^{x_1}f(t,x')\,dt-{x_1\over R}\bar f(x').
\]
Then \(K_1=0\) on the two \(x_1\)-faces and
\[
 \partial_1K_1=f-g,
 \qquad
 g(x')={1\over R}\bar f(x').
\]
The function \(g\) has zero integral on \((0,R)^{d-1}\), and
\[
 \|g\|_{L^2_{x'}}\leq R^{-1/2}\|f\|_{L^2_x},
 \qquad
 \|K_1\|_2\leq C R\|f\|_2
\]
by Cauchy--Schwarz.  By induction choose
\(K'(x')\) with
\(\operatorname{div}_{x'}K'=g\),
zero normal trace, and
\(\|K'\|_{L^2_{x'}}\leq C_{d-1}R\|g\|_2\).
Extend \(K'\) independently of \(x_1\).  Its \(L^2(Q)\) norm is at most
\(C_{d-1}R\|f\|_2\).  Then
\[
 K=(K_1,K'),\qquad\operatorname{div}K=f,
\]
with the stated norm and normal trace.  Zero extension creates no
boundary distribution because the normal trace vanishes.
\end{proof}

\begin{corollary}[Volume-uniform zero-mean local test pairing]
\label{cor:e8v66-local-primitive}
If \(f_L\) is supported in a fixed cube \(Q\), has zero integral, and
\(\sup_L\|f_L\|_2<\infty\), then
\begin{equation}
 |\langle f_L,(-\Delta)^\dagger\rho_L\rangle|
 \leq C_dR\|f_L\|_2\|J_L\|_2
 \label{eq:e8v66-local-primitive-pairing}
\end{equation}
uniformly in the ambient volume and in every dimension.
\end{corollary}

\begin{proof}
Use \Cref{thm:e8v66-primitive,cor:e8v66-double}.
\end{proof}

\subsection{The excluded volume-growing tests}
\label{sec:e8v66-obstruction}

\begin{proposition}[Lowest-mode test has divergent response norm]
\label{prop:e8v66-lowest}
Let \(e_{k_{\min}}\) be an \(L^2\)-normalized lowest nonzero Fourier
mode.  Then
\begin{equation}
 \|e_{k_{\min}}\|_{\dot H^{-1}_L}
 ={1\over|k_{\min}|}
 ={L\over2\pi}.
 \label{eq:e8v66-lowest}
\end{equation}
\end{proposition}

\begin{proof}
Only one Fourier coefficient is nonzero in
\eqref{eq:e8v66-Hminus}, and
\(|k_{\min}|=2\pi/L\).
\end{proof}

\begin{firewall}[Local and volume-growing tests are different claims]
\label{fire:e8v66-test-growth}
\Cref{thm:e8v66-local-test} assumes fixed physical support and uniform
\(L^1\cap L^2\) norms.  The lowest Fourier mode spreads across the
whole box and violates that locality class.  Uniform local observable
pairings therefore do not imply uniform response to macroscopic
low-mode probes.
\end{firewall}

\begin{firewall}[The source primitive remains substantive]
\label{fire:e8v66-source}
The bound \eqref{eq:e8v66-pairing} requires a directly constructed
\(J\) in \eqref{eq:e8v66-source}.  Defining
\(J=\nabla(-\Delta)^\dagger\rho\) would insert the desired inverse
bound into the hypothesis and repeat the V50 circularity.  V51 provides
such a primitive only for endpoint differences and compact time-window
derivatives, not for every fixed-time density.
\end{firewall}

\subsection{Local observable response card}
\label{sec:e8v66-card}

\begin{definition}[Local scalar pairing card, LSPC]
\label{def:e8v66-LSPC}
LSPC consists of:
\begin{enumerate}[label=\textup{(S\arabic*)},leftmargin=4em]
\item a fixed physical-support smearing with uniform \(L^1\cap L^2\)
norms, or a local zero-mean divergence primitive;
\item a direct matter-source primitive from conservation, Ward, or
reference transfer;
\item the pairing bounds
\eqref{eq:e8v66-local-pairing} or
\eqref{eq:e8v66-local-primitive-pairing};
\item separate zero, harmonic, ADM-charge, and macroscopic-test sectors;
\item V44 composite-operator and spacetime contact control;
\item V45 long-range ownership after direct constraint elimination;
\item nonlinear constraint-chart and regulator stability; and
\item reflection, positivity, gauge, anomaly, and cutoff-removal
contacts for the selected observable algebra.
\end{enumerate}
\end{definition}

\begin{assessment}[Progress at V66]
\label{ass:e8v66-status}
The full bounded-clock scalar gap is unnecessary for a sharply declared
class of local observables.  In physical spatial dimension \(d=3\), a
fixed \(L^1\cap L^2\) smearing has uniformly bounded
\(\dot H^{-1}\) norm, and a direct divergence source has norm-one
\(\dot H^{-1}\) control.  Their inverse-Laplacian pairing is therefore
uniform in volume.  Zero-mean local tests admit an explicit local
divergence primitive in every dimension.  Macroscopic lowest-mode tests
remain divergent and are not hidden.
\end{assessment}

\begin{programme}[Eighth-series V67: spacetime test and endpoint source synthesis]
\label{prog:e8v66-V67}
V67 will combine the V44 spacetime-smoothed test kernels with the V51
endpoint and compact-window source primitives.  It will prove the
mixed spacetime norm bound, track the time-window constants, and state
the exact class of linearized Einstein scalar Schwinger pairings now
uniform without a full scalar mixing gap.
\end{programme}

\begin{firewall}[Claim boundary after eighth-series V66]
\label{fire:e8v66-boundary}
V66 proves uniform negative-Sobolev bounds for fixed local tests in
\(d>2\), moment-cancelled tests in every dimension, exact
source--test pairings, and a local rectangular divergence primitive.
It does not supply the source primitive for arbitrary fixed-time energy,
control nonlinear or quantum composite constraints, construct the
positive surface carrier, remove cutoffs, or prove the full continuum.
\end{firewall}

\begin{theorem}[Eighth-series V66 result]
\label{thm:e8v66-result}
In three spatial dimensions, if \(f_L\) has fixed support and uniformly
bounded \(L^1\cap L^2\) norms and
\(\rho_L=-\operatorname{div}J_L\), then
\[
 |\langle f_L,(-\Delta)^\dagger\rho_L\rangle|
 \leq C\bigl(\|f_L\|_1+\|f_L\|_2\bigr)\|J_L\|_2
\]
uniformly in physical volume.  The corresponding norm for a normalized
lowest Fourier-mode test grows as \(L/(2\pi)\).
\end{theorem}

\begin{proof}
Use \Cref{thm:e8v66-local-test,thm:e8v66-pairing,
prop:e8v66-lowest}.
\end{proof}

\paragraph{Parent-version record.}
V66 was formed from
\texttt{Renewal\_SM\_Einstein\_Continuum\_Research\_Notes\_V65.tex}.
The V65 parent had (25\,878) logical source lines, (919\,696) bytes,
and SHA--256
\[
 \texttt{1EB7C6EED8DB87ABE55B476144E7CB6E295DB475C175B9579EA8F670E06B031D}.
\]
The next compulsory companion audit is V70.

% =============================================================================
% END APPENDED EIGHTH-SERIES V66 MATERIAL
% =============================================================================
% =============================================================================
% BEGIN APPENDED EIGHTH-SERIES V67 MATERIAL
% =============================================================================

\clearpage
\section{Eighth-series V67: spacetime-smoothed conservation sources and local scalar response}
\label{sec:v8v67}

\subsection{Mixed spacetime response norm}
\label{sec:e8v67-mixed}

Let \(\Omega_L\) be the periodic spatial box and write
\[
 R_L:=(-\Delta_L)^\dagger P_0.
\]
For a time-dependent mean-zero source \(\rho(t)\), set
\(u(t)=R_L\rho(t)\).

\begin{theorem}[Mixed spacetime divergence-source bound]
\label{thm:e8v67-mixed}
Suppose
\begin{equation}
 \rho(t)=-\operatorname{div}J(t)
 \label{eq:e8v67-source}
\end{equation}
for \(J\in L^2(I\times\Omega_L)\).  If
\[
 f\in L^2\bigl(I;\dot H^{-1}_L(\Omega_L)\bigr),
\]
then
\begin{equation}
 \left|
 \int_I\langle f(t),u(t)\rangle_x\,dt
 \right|
 \leq
 \|f\|_{L^2_t\dot H^{-1}_x}
 \|J\|_{L^2_{t,x}}.
 \label{eq:e8v67-mixed}
\end{equation}
The constant is one and is independent of physical volume, lattice
refinement, and any scalar block-dynamics gap.
\end{theorem}

\begin{proof}
For almost every \(t\), \Cref{thm:e8v66-pairing} gives
\[
 |\langle f(t),R_L\rho(t)\rangle_x|
 \leq\|f(t)\|_{\dot H^{-1}_x}\|J(t)\|_{L^2_x}.
\]
Integrate in \(t\) and apply Cauchy--Schwarz.
\end{proof}

\begin{corollary}[Fixed-support spacetime tests in \(d=3\)]
\label{cor:e8v67-local}
In three spatial dimensions, suppose every \(f_L(t,\cdot)\) is
supported in one fixed physical cube and
\begin{equation}
 M_f^2:=
 \int_I
 \bigl(\|f_L(t)\|_{L^1_x}+\|f_L(t)\|_{L^2_x}\bigr)^2dt
 \label{eq:e8v67-test-norm}
\end{equation}
is bounded uniformly in \(L\).  Then
\begin{equation}
 \left|
 \int_I\langle f_L(t),R_L\rho_L(t)\rangle_xdt
 \right|
 \leq C M_f\|J_L\|_{L^2_{t,x}}
 \label{eq:e8v67-local}
\end{equation}
uniformly in volume.
\end{corollary}

\begin{proof}
\Cref{thm:e8v66-local-test} applied at almost every time gives
\[
 \|f_L(t)\|_{\dot H^{-1}_L}
 \leq C\bigl(\|f_L(t)\|_1+\|f_L(t)\|_2\bigr).
\]
Square, integrate in time, and use \Cref{thm:e8v67-mixed}.
\end{proof}

\subsection{Compact time-window sources}
\label{sec:e8v67-window}

Assume the local balance law from V51,
\begin{equation}
 \partial_tq+\operatorname{div}J=W
 \label{eq:e8v67-balance}
\end{equation}
on \(I=[a,b]\).  For \(\chi\in H^1_0(a,b)\), define
\begin{equation}
 S_\chi:=\int_a^b\chi'(t)q(t)\,dt,
 \qquad
 U_\chi:=R_LS_\chi.
 \label{eq:e8v67-window}
\end{equation}

\begin{theorem}[Exact compact-window scalar response]
\label{thm:e8v67-window}
For every spatial test \(f\in\dot H^{-1}_L\),
\begin{equation}
 |\langle f,U_\chi\rangle|
 \leq
 \|f\|_{\dot H^{-1}_L}\|\chi\|_{L^2_t}
 \left(
 \|J\|_{L^2_{t,x}}
 +\lambda_{1,L}^{-1/2}\|P_0W\|_{L^2_{t,x}}
 \right).
 \label{eq:e8v67-window-bound}
\end{equation}
If \(W=0\), the estimate is volume-uniform whenever the test and
physical current norms are uniform.
\end{theorem}

\begin{proof}
V51 gives the direct decomposition
\[
 S_\chi=-\operatorname{div}Z_\chi+R_\chi,
 \qquad
 Z_\chi=-\int_a^b\chi J\,dt,
 \qquad
 R_\chi=-\int_a^b\chi W\,dt,
\]
with
\[
 \|Z_\chi\|_2\leq\|\chi\|_2\|J\|_{L^2_{t,x}},
 \qquad
 \|R_\chi\|_2\leq\|\chi\|_2\|W\|_{L^2_{t,x}}.
\]
By V50,
\[
 \|S_\chi\|_{\dot H^{-1}_L}
 \leq
 \|Z_\chi\|_2+
 \lambda_{1,L}^{-1/2}\|P_0R_\chi\|_2.
\]
Finally,
\[
 |\langle f,R_LS_\chi\rangle|
 \leq
 \|f\|_{\dot H^{-1}_L}
 \|S_\chi\|_{\dot H^{-1}_L}.
\]
Combining the three displays proves
\eqref{eq:e8v67-window-bound}.
\end{proof}

\begin{corollary}[Combined Einstein--matter window has no work debit]
\label{cor:e8v67-combined}
Let \(q_{\rm tot},J_{\rm tot}\) be the aligned combined current of V53,
so that
\[
 \partial_tq_{\rm tot}+\operatorname{div}J_{\rm tot}=0.
\]
Then
\begin{equation}
 \left|
 \left\langle
 f,
 R_L\int_a^b\chi'(t)q_{\rm tot}(t)\,dt
 \right\rangle
 \right|
 \leq
 \|f\|_{\dot H^{-1}_L}
 \|\chi\|_{L^2_t}
 \|J_{\rm tot}\|_{L^2_{t,x}}.
 \label{eq:e8v67-combined}
\end{equation}
\end{corollary}

\begin{proof}
Set \(W=0\) in \Cref{thm:e8v67-window}.
\end{proof}

\begin{firewall}[Alignment remains a hypothesis]
\label{fire:e8v67-alignment}
The combined-current estimate applies to the scalar constraint source
only after the current-to-ADM alignment defect of V54 is controlled.
Conservation of a convenient improved current does not by itself prove
that its density is the source in the regulated Hamiltonian constraint.
\end{firewall}

\subsection{Endpoint transfer}
\label{sec:e8v67-endpoint}

\begin{theorem}[Endpoint scalar response]
\label{thm:e8v67-endpoint}
Let \(T=b-a\) and
\[
 S_{a,b}:=q(b)-q(a),\qquad U_{a,b}:=R_LS_{a,b}.
\]
Then
\begin{equation}
 |\langle f,U_{a,b}\rangle|
 \leq
 T^{1/2}\|f\|_{\dot H^{-1}_L}
 \left(
 \|J\|_{L^2_{t,x}}
 +\lambda_{1,L}^{-1/2}\|P_0W\|_{L^2_{t,x}}
 \right).
 \label{eq:e8v67-endpoint}
\end{equation}
\end{theorem}

\begin{proof}
Use the V51 endpoint decomposition
\[
 S_{a,b}
 =-\operatorname{div}\int_a^bJ\,dt+\int_a^bW\,dt
\]
and its \(T^{1/2}\) norm estimates, then repeat the last two steps in
the proof of \Cref{thm:e8v67-window}.
\end{proof}

\begin{firewall}[Long time windows carry a real cost]
\label{fire:e8v67-time}
The factors \(\|\chi\|_{L^2_t}\) and \(T^{1/2}\) are explicit.
Uniformity for fixed compact time support does not imply a bound for
time windows growing with the infrared scale.  Such growth must be
matched against decay, a spectral estimate, or a normalized temporal
test.
\end{firewall}

\subsection{Relation to the V44 observable class}
\label{sec:e8v67-V44}

\begin{theorem}[A rigorous linearized scalar Schwinger test class]
\label{thm:e8v67-class}
The following tests have volume-uniform linearized scalar response in
three spatial dimensions:
\begin{enumerate}[label=\textup{(\roman*)},leftmargin=3.5em]
\item spatial smearings supported in a fixed physical cube with the
uniform norm \eqref{eq:e8v67-test-norm};
\item compact \(H^1_0\) time windows;
\item endpoint or time-window energy sources with a uniformly bounded
direct physical flux primitive; and
\item vanishing work residual, or a residual satisfying the explicit
weighted bound in \eqref{eq:e8v67-window-bound}.
\end{enumerate}
For zero-mean spatial tests, the \(L^1\) hypothesis may be replaced by
the local divergence primitive of V66.
\end{theorem}

\begin{proof}
Items (i)--(iv) are precisely the hypotheses of
\Cref{cor:e8v67-local,thm:e8v67-window,thm:e8v67-endpoint}.
For the last statement use
\Cref{thm:e8v66-primitive,cor:e8v66-double}.
\end{proof}

\begin{remark}[What V44 adds]
\label{rem:e8v67-V44}
V44 supplies spacetime smoothing strong enough to make the declared
quadratic composite tests trace class and to control their Wick
contacts.  V67 supplies the separate infrared condition on the scalar
Einstein response.  A test used in the continuum compiler must satisfy
both interfaces; neither theorem silently implies the other.
\end{remark}

\begin{firewall}[This is a pairing theorem, not carrier positivity]
\label{fire:e8v67-pairing}
The estimates above control selected linearized Schwinger pairings
without a full scalar mixing gap.  They do not construct a positive
Einstein--matter measure, prove reflection positivity, control the
nonlinear constraint Jacobian, or give a Hamiltonian mass gap.
\end{firewall}

\subsection{Spacetime scalar response card}
\label{sec:e8v67-card}

\begin{definition}[Spacetime scalar response card, SSRC]
\label{def:e8v67-SSRC}
SSRC consists of:
\begin{enumerate}[label=\textup{(R\arabic*)},leftmargin=4em]
\item the fixed-support spatial test norm
\eqref{eq:e8v67-test-norm};
\item a compact time window or controlled endpoint interval;
\item the direct balance-law primitive and its \(L^2_{t,x}\) bound;
\item aligned combined-current cancellation or the explicit work
residual in \eqref{eq:e8v67-window-bound};
\item V44 trace-class composite smoothing and contact estimates;
\item V54 ADM-current alignment and quantum Ward contacts;
\item V45 long-range ownership and separate charge/boundary sectors;
and
\item nonlinear stability, reflection, positivity, and cutoff removal
for the declared test algebra.
\end{enumerate}
\end{definition}

\begin{assessment}[Progress at V67]
\label{ass:e8v67-status}
The spatial and temporal interfaces now close at the linearized pairing
level.  Compact time-window derivatives and endpoint energy changes
have direct flux primitives, and fixed local three-dimensional tests
pair uniformly with their inverse-Laplacian response.  Combined
Einstein--matter conservation removes the deformation-work debit when
the current is aligned.  The result bypasses the V64 full-space clock
obstruction without misreporting a full scalar gap.
\end{assessment}

\begin{programme}[Eighth-series V68: cofinal convergence of local scalar pairings]
\label{prog:e8v67-V68}
V68 will place the boxes in a common \(\mathbb R^3\) space and prove
convergence of the local scalar pairings under strong convergence of
tests and weak convergence of the physical flux primitives.  It will
separate compactness needed for existence from uniqueness of the
limiting constrained field.
\end{programme}

\begin{firewall}[Claim boundary after eighth-series V67]
\label{fire:e8v67-boundary}
V67 proves mixed spacetime, compact-window, endpoint, and local
fixed-support scalar pairing bounds by synthesizing V44, V50--V51, and
V66.  It does not prove convergence of the flux primitives, nonlinear
alignment, composite constraint renormalization, carrier positivity,
cutoff removal, or the full continuum.
\end{firewall}

\begin{theorem}[Eighth-series V67 result]
\label{thm:e8v67-result}
For fixed-support spacetime-smoothed tests in three spatial dimensions,
compact time-window or endpoint sources satisfying
\(\partial_tq+\operatorname{div}J=W\) obey
\[
 |\langle f,R_LS\rangle|
 \leq
 C_fC_t
 \left(
 \|J\|_{L^2_{t,x}}
 +\lambda_{1,L}^{-1/2}\|P_0W\|_{L^2_{t,x}}
 \right).
\]
When an aligned combined Einstein--matter current makes \(W=0\), the
bound is uniform in volume without a full scalar mixing gap.
\end{theorem}

\begin{proof}
Use \Cref{thm:e8v67-window,thm:e8v67-endpoint,
thm:e8v67-class}.
\end{proof}

\paragraph{Parent-version record.}
V67 was formed from
\texttt{Renewal\_SM\_Einstein\_Continuum\_Research\_Notes\_V66.tex}.
The V66 parent had (26\,288) logical source lines, (932\,046) bytes,
and SHA--256
\[
 \texttt{BE4CD86231DF94677A30E00D18BCEFAB3A7E68F45360AC05E193552CBA21B6A7}.
\]
The next compulsory companion audit is V70.

% =============================================================================
% END APPENDED EIGHTH-SERIES V67 MATERIAL
% =============================================================================
% =============================================================================
% BEGIN APPENDED EIGHTH-SERIES V68 MATERIAL
% =============================================================================

\clearpage
\section{Eighth-series V68: cofinal convergence of the linearized local scalar response}
\label{sec:v8v68}

\subsection{The continuum energy solution}
\label{sec:e8v68-continuum}

Use the Fourier convention
\[
 \widehat f(\xi)=\int_{\mathbb R^3}f(x)e^{-i\xi\cdot x}\,dx,
 \qquad
 \|f\|_2^2=(2\pi)^{-3}\int_{\mathbb R^3}|\widehat f(\xi)|^2d\xi.
\]

\begin{theorem}[Continuum scalar constraint solution in the energy
topology]
\label{thm:e8v68-energy}
Let \(J\in L^2(\mathbb R^3;\mathbb R^3)\) and
\(\rho=-\operatorname{div}J\in\dot H^{-1}(\mathbb R^3)\).  There is a
unique \(u\in\dot H^1(\mathbb R^3)\), modulo constants, satisfying
\begin{equation}
 -\Delta u=\rho
 \label{eq:e8v68-Poisson}
\end{equation}
weakly, and
\begin{equation}
 \widehat{\nabla u}(\xi)
 =
 {\xi\otimes\xi\over|\xi|^2}\widehat J(\xi)
 \quad\hbox{for a.e. }\xi\neq0.
 \label{eq:e8v68-Riesz}
\end{equation}
Consequently,
\begin{equation}
 \|\nabla u\|_2\leq\|J\|_2.
 \label{eq:e8v68-energy}
\end{equation}
\end{theorem}

\begin{proof}
Define the gradient by the Fourier multiplier in
\eqref{eq:e8v68-Riesz}.  The matrix
\(\xi\otimes\xi/|\xi|^2\) is an orthogonal projection, so the field lies
in \(L^2\) and obeys \eqref{eq:e8v68-energy}.  It is curl free and hence
is the distributional gradient of an element of the homogeneous Sobolev
space, unique modulo constants.  Moreover,
\[
 |\xi|^2\widehat u(\xi)
 =-i\xi\cdot\widehat J(\xi)
 =\widehat\rho(\xi)
\]
with the sign fixed by
\(\rho=-\operatorname{div}J\), proving
\eqref{eq:e8v68-Poisson}.  If two energy solutions exist, their
difference has zero gradient norm and is constant.
\end{proof}

\begin{corollary}[The primitive is unique only longitudinally]
\label{cor:e8v68-longitudinal}
If \(J'\in L^2\) also satisfies
\(-\operatorname{div}J'=\rho\), then
\begin{equation}
 P_{\rm grad}J'=P_{\rm grad}J=\nabla u.
 \label{eq:e8v68-longitudinal}
\end{equation}
The solenoidal difference \(J'-J\) does not affect any scalar response
pairing.
\end{corollary}

\begin{proof}
The Fourier transform of \(J'-J\) is orthogonal to \(\xi\) almost
everywhere.  The multiplier in \eqref{eq:e8v68-Riesz} annihilates it.
\end{proof}

\subsection{A common Fourier space for expanding tori}
\label{sec:e8v68-embedding}

Let \(L_n\to\infty\), set \(h_n=2\pi/L_n\), and let
\({\cal K}_n=h_n\mathbb Z^3\).  For each \(k\in{\cal K}_n\), let
\[
 Q_{n,k}=k+[-h_n/2,h_n/2)^3.
\]
For a periodic field \(g_n\), define the Fourier-step embedding
\begin{equation}
 ({\cal E}_ng_n)(\xi):=\widehat g_n(k),
 \qquad \xi\in Q_{n,k}.
 \label{eq:e8v68-embedding}
\end{equation}

\begin{lemma}[The Fourier-step embedding is isometric]
\label{lem:e8v68-isometry}
One has
\begin{equation}
 (2\pi)^{-3}\|{\cal E}_ng_n\|_{L^2_\xi}^2
 =
 \|g_n\|_{L^2(\Omega_{L_n})}^2.
 \label{eq:e8v68-isometry}
\end{equation}
\end{lemma}

\begin{proof}
Every cell has volume
\(h_n^3=(2\pi/L_n)^3\).  Therefore
\[
 (2\pi)^{-3}\sum_kh_n^3|\widehat g_n(k)|^2
 =
 L_n^{-3}\sum_k|\widehat g_n(k)|^2,
\]
which is periodic Parseval.
\end{proof}

Let \(f\in C_c^\infty(\mathbb R^3)\).  For all sufficiently large
\(n\), regard \(f\) as supported inside the centered fundamental box
and define
\begin{equation}
 K_{n,f}:=\nabla(-\Delta_{L_n})^\dagger P_0f.
 \label{eq:e8v68-periodic-test}
\end{equation}
Its embedded Fourier field is the step function with values
\begin{equation}
 ({\cal E}_nK_{n,f})(\xi)
 =
 {ik\over|k|^2}\widehat f(k),
 \qquad \xi\in Q_{n,k},\quad k\neq0,
 \label{eq:e8v68-periodic-symbol}
\end{equation}
and value zero on \(Q_{n,0}\).

\begin{theorem}[Strong convergence of periodic Riesz test fields]
\label{thm:e8v68-test-convergence}
For every \(f\in C_c^\infty(\mathbb R^3)\),
\begin{equation}
 {\cal E}_nK_{n,f}
 \longrightarrow
 K_f(\xi):={i\xi\over|\xi|^2}\widehat f(\xi)
 \quad\hbox{strongly in }L^2(\mathbb R^3_\xi).
 \label{eq:e8v68-test-convergence}
\end{equation}
\end{theorem}

\begin{proof}
Fix \(\epsilon>0\).  On an annulus
\(\epsilon\leq|\xi|\leq\epsilon^{-1}\), the multiplier times
\(\widehat f\) is uniformly continuous, and its cellwise samples
converge uniformly, hence in \(L^2\).

Near zero,
\[
 |K_f(\xi)|\leq{\|f\|_1\over|\xi|},
\]
whose square is integrable in three dimensions.  The continuum norm on
\(|\xi|<\epsilon\) tends to zero with \(\epsilon\).  The corresponding
step-function norms are uniformly small by the lattice-shell estimate
in \Cref{lem:e8v66-lattice-sum}, with the same proof restricted to a
ball of radius \(\epsilon\).

At infinity, \(\widehat f\) is Schwartz.  Hence
\(|\xi|^{-1}\widehat f(\xi)\) has an arbitrarily small \(L^2\) tail, and
its lattice Riemann-sum tails are uniformly small for sufficiently
large cutoff radius.  First choose the inner and outer radii, then let
\(n\to\infty\) on the intervening annulus.  This proves strong
convergence.
\end{proof}

\subsection{Weak flux compactness gives convergent pairings}
\label{sec:e8v68-compactness}

\begin{theorem}[Weak--strong cofinal pairing theorem]
\label{thm:e8v68-pairing}
Let \(J_n\in L^2(\Omega_{L_n};\mathbb R^3)\) satisfy
\begin{equation}
 \sup_n\|J_n\|_2<\infty
 \label{eq:e8v68-J-bound}
\end{equation}
and suppose, along a subsequence,
\begin{equation}
 {\cal E}_nJ_n\rightharpoonup\widehat J
 \quad\hbox{weakly in }L^2(\mathbb R^3_\xi).
 \label{eq:e8v68-J-weak}
\end{equation}
Set
\[
 \rho_n=-\operatorname{div}J_n,
 \qquad
 u_n=(-\Delta_{L_n})^\dagger\rho_n.
\]
Then for every \(f\in C_c^\infty(\mathbb R^3)\),
\begin{equation}
 \lim_{n\to\infty}\langle f,u_n\rangle_{\Omega_{L_n}}
 =
 (2\pi)^{-3}
 \int_{\mathbb R^3}
 \overline{K_f(\xi)}\cdot\widehat J(\xi)\,d\xi.
 \label{eq:e8v68-pairing-limit}
\end{equation}
The right-hand side is the continuum pairing
\(\langle f,u\rangle\) for the energy solution generated by the
longitudinal part of \(J\).
\end{theorem}

\begin{proof}
Periodic integration by parts and self-adjointness give
\[
 \langle f,u_n\rangle
 =
 \langle K_{n,f},J_n\rangle.
\]
By \Cref{lem:e8v68-isometry}, this equals
\[
 (2\pi)^{-3}
 \int
 \overline{{\cal E}_nK_{n,f}}\cdot{\cal E}_nJ_n\,d\xi.
\]
The first factor converges strongly by
\Cref{thm:e8v68-test-convergence}; the second converges weakly by
\eqref{eq:e8v68-J-weak}.  Their pairings therefore converge to
\eqref{eq:e8v68-pairing-limit}.  The identification with the continuum
energy solution is \Cref{thm:e8v68-energy}.
\end{proof}

\begin{corollary}[Subsequential existence from the flux bound]
\label{cor:e8v68-subsequence}
The uniform bound \eqref{eq:e8v68-J-bound} alone gives a subsequence on
which all pairings in \eqref{eq:e8v68-pairing-limit} converge
simultaneously for a countable dense family of compactly supported
smooth tests.
\end{corollary}

\begin{proof}
\Cref{lem:e8v68-isometry} makes
\({\cal E}_nJ_n\) a bounded sequence in one Hilbert space.  Weak
sequential compactness gives a weakly convergent subsequence.  Apply
\Cref{thm:e8v68-pairing}; a diagonal extraction handles a chosen
countable dense test family, although the single weakly convergent
subsequence already works for every fixed test.
\end{proof}

\begin{firewall}[Compactness is not uniqueness]
\label{fire:e8v68-uniqueness}
Different cofinal subsequences may have different weak longitudinal
flux limits.  The uniform bound proves subsequential existence of local
scalar pairings, not uniqueness of the continuum state.  Uniqueness
requires convergence of the longitudinal physical current, a limiting
Ward/balance identity determining it, or a separate uniqueness theorem
for the full interacting measure.
\end{firewall}

\subsection{Spacetime extension}
\label{sec:e8v68-spacetime}

\begin{corollary}[Spacetime weak--strong limit]
\label{cor:e8v68-spacetime}
Let \(I\) be bounded, let
\(f\in C_c^\infty(I\times\mathbb R^3)\), and assume
\[
 {\cal E}_nJ_n\rightharpoonup\widehat J
 \quad\hbox{in }L^2(I\times\mathbb R^3_\xi).
\]
Then
\begin{equation}
 \int_I\langle f(t),u_n(t)\rangle\,dt
 \longrightarrow
 (2\pi)^{-3}\int_I\int_{\mathbb R^3}
 \overline{
 {i\xi\over|\xi|^2}\widehat f(t,\xi)}
 \cdot\widehat J(t,\xi)\,d\xi\,dt.
 \label{eq:e8v68-spacetime}
\end{equation}
\end{corollary}

\begin{proof}
The proof of \Cref{thm:e8v68-test-convergence} applies in the Bochner
space \(L^2(I\times\mathbb R^3_\xi)\), using the compact time support
and smoothness of \(f\).  Pair this strong convergence with the stated
weak convergence.
\end{proof}

\begin{firewall}[The embedding proves the linear scalar sector only]
\label{fire:e8v68-scope}
The Fourier-step construction uses the flat periodic Laplacian and a
linear divergence source.  A nonlinear random metric changes the
elliptic operator, volume density, and constraint chart.  Extending the
argument requires resolvent or form convergence of those random
operators together with uniform local test and flux bounds.
\end{firewall}

\subsection{Cofinal scalar compactness card}
\label{sec:e8v68-card}

\begin{definition}[Cofinal scalar compactness card, CSCC]
\label{def:e8v68-CSCC}
CSCC consists of:
\begin{enumerate}[label=\textup{(C\arabic*)},leftmargin=4em]
\item a common Fourier or geometric Hilbert-space embedding;
\item strong convergence of the local Riesz test fields;
\item a cofinal \(L^2\) bound and weak compactness for directly
constructed flux primitives;
\item uniqueness of the longitudinal weak limit;
\item aligned balance and Ward identities in the limit;
\item form/resolvent stability under nonlinear random metrics;
\item V44 composite tests, V45 ownership, and charge/boundary control;
and
\item reflection, positivity, tightness, and reconstruction of the
continuum observable algebra.
\end{enumerate}
\end{definition}

\begin{assessment}[Progress at V68]
\label{ass:e8v68-status}
The local scalar estimate now yields a genuine cofinal compactness
statement.  Periodic Riesz test fields converge strongly in a single
Fourier \(L^2(\mathbb R^3)\) space, while a uniform physical-flux bound
gives weakly convergent subsequences.  Weak--strong pairing then
constructs the limiting local linearized scalar response.  Existence is
proved subsequentially; uniqueness and nonlinear metric stability remain
separate.
\end{assessment}

\begin{programme}[Eighth-series V69: variable-metric form stability]
\label{prog:e8v68-V69}
V69 will replace the flat Dirichlet form by uniformly elliptic regulated
metric forms on fixed local patches.  It will prove a quantitative
resolvent-energy comparison and determine the coefficient convergence
needed for strong convergence of local test solutions.
\end{programme}

\begin{firewall}[Claim boundary after eighth-series V68]
\label{fire:e8v68-boundary}
V68 proves continuum energy solvability for divergence sources, strong
periodic Riesz-test convergence, weak--strong convergence of local
pairings, and subsequential compactness from a flux bound.  It does not
prove uniqueness of the flux limit, variable-metric resolvent
convergence, nonlinear constraint control, positive carrier
reconstruction, full cutoff removal, or the complete continuum.
\end{firewall}

\begin{theorem}[Eighth-series V68 result]
\label{thm:e8v68-result}
If the directly constructed periodic flux primitives are uniformly
bounded in \(L^2\), then along a subsequence their linearized scalar
constraint responses converge against every compactly supported smooth
three-dimensional test.  The limit is the unique homogeneous-energy
solution generated by the weak longitudinal flux limit.
\end{theorem}

\begin{proof}
Use \Cref{thm:e8v68-energy,thm:e8v68-test-convergence,
thm:e8v68-pairing,cor:e8v68-subsequence}.
\end{proof}

\paragraph{Parent-version record.}
V68 was formed from
\texttt{Renewal\_SM\_Einstein\_Continuum\_Research\_Notes\_V67.tex}.
The V67 parent had (26\,643) logical source lines, (942\,716) bytes,
and SHA--256
\[
 \texttt{413D4A33C74D537A24FA7ADFF9367A73222C5A206DA3741755320A28C9C2B9B6}.
\]
The next compulsory companion audit is V70.

% =============================================================================
% END APPENDED EIGHTH-SERIES V68 MATERIAL
% =============================================================================
% =============================================================================
% BEGIN APPENDED EIGHTH-SERIES V69 MATERIAL
% =============================================================================

\clearpage
\section{Eighth-series V69: variable-metric energy forms and strong local test resolvents}
\label{sec:v8v69}

\subsection{Uniformly elliptic local forms}
\label{sec:e8v69-forms}

Let \(D\subset\mathbb R^3\) be a bounded Lipschitz patch and set
\(V=H^1_0(D)\), with norm \(\|v\|_V=\|\nabla v\|_2\).
Let \(A(x)\) be a measurable symmetric matrix satisfying
\begin{equation}
 m|\zeta|^2\leq\zeta\cdot A(x)\zeta\leq M|\zeta|^2
 \label{eq:e8v69-ellipticity}
\end{equation}
for almost every \(x\), every \(\zeta\), and fixed
\(0<m\leq M<\infty\).  Define
\begin{equation}
 a_A(u,v):=\int_D\nabla u\cdot A\nabla v\,dx.
 \label{eq:e8v69-form}
\end{equation}

\begin{theorem}[Metric energy solution for a flux source]
\label{thm:e8v69-flux}
For every \(J\in L^2(D;\mathbb R^3)\), there is a unique
\(u_A\in V\) such that
\begin{equation}
 a_A(u_A,v)=\int_DJ\cdot\nabla v\,dx
 \qquad(v\in V).
 \label{eq:e8v69-flux-equation}
\end{equation}
It satisfies
\begin{equation}
 \|\nabla u_A\|_2\leq m^{-1}\|J\|_2.
 \label{eq:e8v69-flux-bound}
\end{equation}
\end{theorem}

\begin{proof}
The form is bounded and coercive by
\eqref{eq:e8v69-ellipticity}.  The right-hand side has norm at most
\(\|J\|_2\) in \(V^*\).  Lax--Milgram gives existence and uniqueness.
Taking \(v=u_A\) yields
\[
 m\|\nabla u_A\|_2^2
 \leq\|J\|_2\|\nabla u_A\|_2,
\]
which proves \eqref{eq:e8v69-flux-bound}.
\end{proof}

\begin{theorem}[Test resolvent and pairing identity]
\label{thm:e8v69-test}
For every \(f\in V^*\), let \(w_{A,f}\in V\) solve
\begin{equation}
 a_A(w_{A,f},v)=\langle f,v\rangle
 \qquad(v\in V).
 \label{eq:e8v69-test-equation}
\end{equation}
Then
\begin{equation}
 \|\nabla w_{A,f}\|_2\leq m^{-1}\|f\|_{V^*}
 \label{eq:e8v69-test-bound}
\end{equation}
and, for the flux solution in \Cref{thm:e8v69-flux},
\begin{equation}
 \langle f,u_A\rangle
 =
 \int_DJ\cdot\nabla w_{A,f}\,dx.
 \label{eq:e8v69-pairing}
\end{equation}
Consequently,
\begin{equation}
 |\langle f,u_A\rangle|
 \leq m^{-1}\|f\|_{V^*}\|J\|_2.
 \label{eq:e8v69-pairing-bound}
\end{equation}
\end{theorem}

\begin{proof}
Lax--Milgram gives \(w_{A,f}\), and testing
\eqref{eq:e8v69-test-equation} with \(w_{A,f}\) proves
\eqref{eq:e8v69-test-bound}.  Symmetry of \(a_A\) gives
\[
 \langle f,u_A\rangle
 =a_A(w_{A,f},u_A)
 =a_A(u_A,w_{A,f})
 =\int_DJ\cdot\nabla w_{A,f}.
\]
Cauchy--Schwarz and \eqref{eq:e8v69-test-bound} prove the last bound.
\end{proof}

\subsection{Quantitative coefficient stability}
\label{sec:e8v69-stability}

\begin{theorem}[Strong test-resolvent comparison]
\label{thm:e8v69-comparison}
Let \(A,B\) both satisfy \eqref{eq:e8v69-ellipticity} with the same
\(m,M\).  For fixed \(f\in V^*\),
\begin{equation}
 \|\nabla(w_{A,f}-w_{B,f})\|_2
 \leq
 {\|A-B\|_{L^\infty}\over m^2}\|f\|_{V^*}.
 \label{eq:e8v69-comparison}
\end{equation}
More generally, for \(f_A,f_B\in V^*\),
\begin{equation}
 \|\nabla(w_{A,f_A}-w_{B,f_B})\|_2
 \leq
 {1\over m}\|f_A-f_B\|_{V^*}
 +{\|A-B\|_\infty\over m^2}\|f_B\|_{V^*}.
 \label{eq:e8v69-comparison-general}
\end{equation}
\end{theorem}

\begin{proof}
Put \(z=w_{A,f_A}-w_{B,f_B}\).  Subtract the two weak equations and
test with \(z\):
\[
 a_A(z,z)
 =
 \langle f_A-f_B,z\rangle
 +\int_D\nabla z\cdot(B-A)\nabla w_{B,f_B}.
\]
Thus
\[
 m\|\nabla z\|_2^2
 \leq
 \left(
 \|f_A-f_B\|_{V^*}
 +\|A-B\|_\infty\|\nabla w_{B,f_B}\|_2
 \right)\|\nabla z\|_2.
\]
Use \eqref{eq:e8v69-test-bound} for \(B,f_B\) and divide by
\(m\|\nabla z\|_2\) when the latter is nonzero.  This proves
\eqref{eq:e8v69-comparison-general}; take \(f_A=f_B=f\) for
\eqref{eq:e8v69-comparison}.
\end{proof}

\begin{corollary}[Weak flux plus strong metric convergence]
\label{cor:e8v69-weak-strong}
Suppose
\begin{equation}
 A_n\longrightarrow A\quad\hbox{in }L^\infty(D),
 \qquad
 f_n\longrightarrow f\quad\hbox{in }V^*,
 \qquad
 J_n\rightharpoonup J\quad\hbox{in }L^2(D).
 \label{eq:e8v69-convergence}
\end{equation}
Let \(u_n\) solve \eqref{eq:e8v69-flux-equation} for \(A_n,J_n\).
Then
\begin{equation}
 \langle f_n,u_n\rangle
 \longrightarrow
 \int_DJ\cdot\nabla w_{A,f}\,dx.
 \label{eq:e8v69-pairing-limit}
\end{equation}
\end{corollary}

\begin{proof}
\Cref{thm:e8v69-comparison} gives
\[
 \nabla w_{A_n,f_n}\longrightarrow\nabla w_{A,f}
 \quad\hbox{strongly in }L^2(D).
\]
By \eqref{eq:e8v69-pairing},
\[
 \langle f_n,u_n\rangle
 =\int_DJ_n\cdot\nabla w_{A_n,f_n}.
\]
Pair weak convergence of \(J_n\) with strong convergence of the test
gradients.
\end{proof}

\begin{corollary}[Quantitative pairing perturbation]
\label{cor:e8v69-pairing-perturbation}
For fixed \(f,J\),
\begin{equation}
 |\langle f,u_A-u_B\rangle|
 \leq
 {\|A-B\|_\infty\over m^2}
 \|f\|_{V^*}\|J\|_2.
 \label{eq:e8v69-pairing-perturbation}
\end{equation}
\end{corollary}

\begin{proof}
Use the pairing representation
\[
 \langle f,u_A-u_B\rangle
 =
 \int_DJ\cdot\nabla(w_{A,f}-w_{B,f})
\]
and \Cref{thm:e8v69-comparison}.
\end{proof}

\subsection{Metric coefficients}
\label{sec:e8v69-metric}

For a Riemannian spatial metric \(g\), the scalar energy form uses
\begin{equation}
 A(g):=\sqrt{\det g}\,g^{-1}.
 \label{eq:e8v69-metric-coefficient}
\end{equation}

\begin{lemma}[Lipschitz dependence on a uniformly nondegenerate metric]
\label{lem:e8v69-metric}
Fix \(0<m_g\leq M_g<\infty\).  On the set of symmetric matrices
satisfying
\begin{equation}
 m_gI\preceq g\preceq M_gI,
 \label{eq:e8v69-metric-window}
\end{equation}
there is \(C=C(m_g,M_g,d)\) such that
\begin{equation}
 \|A(g)-A(h)\|
 \leq C\|g-h\|.
 \label{eq:e8v69-metric-Lipschitz}
\end{equation}
The corresponding matrices \(A(g)\) have uniform ellipticity constants
depending only on \(m_g,M_g,d\).
\end{lemma}

\begin{proof}
The map
\[
 g\longmapsto\sqrt{\det g}\,g^{-1}
\]
is smooth on the positive-definite cone.  The line segment between
two matrices satisfying \eqref{eq:e8v69-metric-window} remains in that
compact spectral window.  The derivative of the map is uniformly
bounded there, so the mean-value formula gives
\eqref{eq:e8v69-metric-Lipschitz}.  Its eigenvalues are
\(\sqrt{\det g}/\lambda_j(g)\), which lie between positive constants
determined by the same window.
\end{proof}

\begin{corollary}[Strong metric convergence gives form-resolvent
stability]
\label{cor:e8v69-metric}
If regulated metrics \(g_n\) satisfy a common window
\eqref{eq:e8v69-metric-window} and
\[
 g_n\to g\quad\hbox{in }L^\infty(D),
\]
then \(A(g_n)\to A(g)\) in \(L^\infty(D)\), and
\Cref{cor:e8v69-weak-strong} applies.
\end{corollary}

\begin{proof}
Apply \Cref{lem:e8v69-metric} pointwise and take the essential supremum.
\end{proof}

\begin{firewall}[Uniform ellipticity is not automatic in the conformal
sector]
\label{fire:e8v69-ellipticity}
The V47 translated conformal wells show that a fixed absolute metric
chart need not have positive worst-conditional mass.  V69 assumes a
uniform nondegeneracy window on the local patch; it does not prove that
the Einstein--Standard Model measure concentrates there.  Degenerate
metrics and changing signature remain outside the theorem.
\end{firewall}

\begin{firewall}[Principal scalar form is not the full ADM constraint]
\label{fire:e8v69-ADM}
The coefficient \(A(g)\) controls the Laplace--Beltrami principal form.
The nonlinear ADM Hamiltonian constraint also contains curvature,
extrinsic-curvature, matter, gauge, boundary, and lower-order terms.
Those terms require form-bounded perturbation estimates and the
alignment ledger; they cannot be discarded because the principal
resolvent is stable.
\end{firewall}

\subsection{Patch-to-global interface}
\label{sec:e8v69-global}

\begin{theorem}[Local deterministic continuum interface]
\label{thm:e8v69-interface}
On a fixed patch \(D\), the following data suffice for convergence of
the scalar response pairing:
\begin{enumerate}[label=\textup{(\roman*)},leftmargin=3.5em]
\item common uniform ellipticity;
\item strong \(L^\infty\) convergence of the metric coefficients;
\item strong convergence of the test in \(H^{-1}(D)\); and
\item weak \(L^2\) convergence of the direct flux primitive.
\end{enumerate}
No spectral gap for a scalar sampling dynamics enters this conclusion.
\end{theorem}

\begin{proof}
This is exactly \Cref{cor:e8v69-weak-strong}, with the metric
coefficient supplied by \Cref{cor:e8v69-metric}.
\end{proof}

\begin{firewall}[Boundary influence is a separate limit]
\label{fire:e8v69-boundary-influence}
Dirichlet data on \(D\) were fixed to place all forms in one Hilbert
space.  The physical periodic or asymptotic solution restricted to
\(D\) differs by a homogeneous boundary response.  V45's long-range
owner/boundary compiler must show that this difference is controlled or
convergent before the local Dirichlet limit is identified with the
global Einstein response.
\end{firewall}

\subsection{Variable-metric scalar card}
\label{sec:e8v69-card}

\begin{definition}[Variable-metric scalar compactness card, VMSC]
\label{def:e8v69-VMSC}
VMSC consists of:
\begin{enumerate}[label=\textup{(V\arabic*)},leftmargin=4em]
\item a local metric nondegeneracy window;
\item strong coefficient convergence or a weaker proved Mosco-form
convergence;
\item strong local test-resolvent convergence;
\item weak compactness and longitudinal uniqueness of flux primitives;
\item form bounds for all lower-order ADM and Standard Model terms;
\item boundary-influence convergence linking the patch to the global
solution;
\item quantum composite, Ward, charge, and reflection control; and
\item assembly of the local limits into a positive continuum state.
\end{enumerate}
\end{definition}

\begin{assessment}[Progress at V69]
\label{ass:e8v69-status}
The flat scalar compactness theorem now has a variable-metric local
extension.  Uniform ellipticity and strong coefficient convergence make
the local test resolvents converge strongly in energy, so weak physical
flux convergence suffices for convergence of every response pairing.
This is independent of the obstructed scalar sampling gap.  Metric
chart probability, lower-order ADM terms, and boundary influence remain
explicit interfaces.
\end{assessment}

\begin{programme}[Eighth-series V70: companion audit and lower-order
form route]
\label{prog:e8v69-V70}
V70 will perform the compulsory YM/NS audit.  It will compare new
configuration-transport or projection-flow results with VMSC, then
choose between lower-order ADM form bounds, boundary influence, and
metric-chart tightness as the next upstream target.
\end{programme}

\begin{firewall}[Claim boundary after eighth-series V69]
\label{fire:e8v69-boundary}
V69 proves variable-coefficient flux solvability, quantitative
test-resolvent stability, weak--strong pairing convergence, and
Lipschitz dependence of the principal metric coefficient.  It does not
prove metric-chart tightness, lower-order ADM form bounds, global
boundary convergence, carrier positivity, full cutoff removal, or the
complete continuum.
\end{firewall}

\begin{theorem}[Eighth-series V69 result]
\label{thm:e8v69-result}
On a fixed patch, if
\(A_n\to A\) in \(L^\infty\) with common ellipticity,
\(f_n\to f\) in \(H^{-1}\), and \(J_n\rightharpoonup J\) in \(L^2\),
then the variable-metric scalar response pairings converge:
\[
 \langle f_n,u_n\rangle
 =
 \int_DJ_n\cdot\nabla w_{A_n,f_n}
 \longrightarrow
 \int_DJ\cdot\nabla w_{A,f}.
\]
\end{theorem}

\begin{proof}
Use \Cref{thm:e8v69-comparison,cor:e8v69-weak-strong}.
\end{proof}

\paragraph{Parent-version record.}
V69 was formed from
\texttt{Renewal\_SM\_Einstein\_Continuum\_Research\_Notes\_V68.tex}.
The V68 parent had (27\,010) logical source lines, (954\,955) bytes,
and SHA--256
\[
 \texttt{FF96EC8924548F5B4E237493757919E6C75322ED250FBFA576AA703157291189}.
\]
The next compulsory companion audit is V70.

% =============================================================================
% END APPENDED EIGHTH-SERIES V69 MATERIAL
% =============================================================================
% =============================================================================
% BEGIN APPENDED EIGHTH-SERIES V70 MATERIAL
% =============================================================================

\clearpage
\section{Eighth-series V70: compulsory audit and the lower-order form ledger}
\label{sec:v8v70}

\subsection{Companion files inspected}
\label{sec:e8v70-audit}

The V70 audit inspected
\begin{align}
 &\texttt{Renewal\_Yang\_Mills\_Mass\_Gap\_Research\_Notes\_V22.tex},
 \notag\\
 &\hspace{2em}8875\ \text{logical lines},\quad322409\ \text{bytes},
 \quad\operatorname{SHA256}
 =\texttt{EB8E96B0378FBA49A52C6F3CC13C67D48EE4859145E6FC8B9826363FF533EE04},
 \label{eq:e8v70-YM-file}\\
 &\texttt{Renewal\_NS\_Stability\_Research\_Notes\_V26.tex},
 \notag\\
 &\hspace{2em}5319\ \text{logical lines},\quad278283\ \text{bytes},
 \quad\operatorname{SHA256}
 =\texttt{82C887F8C2C69E4F28FAD7C3DC8DE5B9099CB5A4BF431EBA1FFF3E91935978AD}.
 \label{eq:e8v70-NS-file}
\end{align}
Both files have one live terminator.  The NS programme remains
unchanged.

\begin{audit}[Yang--Mills V20]
\label{audit:e8v70-YM20}
The V20 companion audit reads the SM scalar clock calculations through
V63 and does not import their \(L^{-6}\) exponent into Yang--Mills.
It constructs an explicit selected-curvature dilation, computes its
continuity-equation density defect, and keeps material and unselected
plaquette remainders explicit.  The construction is type-correct
because the transport score contains a directional action derivative.
\end{audit}

\begin{audit}[Yang--Mills V21]
\label{audit:e8v70-YM21}
For the global linearized curvature complex, one-link heat-bath updates
have exact gap of order \(L^{-2}\), and every fixed bounded-range
curvature-image block family has the same upper-order obstruction.
Fourier shells have gap one but are nonlocal.  The exponent differs
from the SM scalar exponent because the curvature symbol has one
derivative and no scalar charge-preserving difference is required.
\end{audit}

\begin{audit}[Yang--Mills V22]
\label{audit:e8v70-YM22}
Translated dyadic cube link fields, weighted by
\(s^{-(d-1)}\), give a uniform all-chaos curvature-Gaussian frame while
their total per-translate scale intensity is summable for \(d>1\).
The proof exploits face support of the curvature and explicitly retains
macroscopic block range, torons, nonlinear group realization, and
physical conditional scores as open rows.
\end{audit}

\begin{theorem}[Derivative order prevents a direct clock transfer]
\label{thm:e8v70-no-transfer}
The V22 bounded-clock rate cannot be transferred to the SM scalar
constraint frame by analogy alone.  The audited Yang--Mills update
symbol has one derivative and a cube amplitude, whereas the SM scalar
normal frame contains two Laplacian factors and a charge-preserving
template zero.  Their low-momentum frame orders are respectively
\(O(|k|^2)\) and \(O(|k|^6)\) for fixed local shapes.
\end{theorem}

\begin{proof}
The audited YM V21 factorization is
\(\widehat u=d_1(k)\widehat v\), with
\(\|d_1(k)\|=O(|k|)\), so the squared rank-one frame is
\(O(|k|^2)\).  In V62 the SM whitened template is
\[
 \widehat u(k)=\lambda(k)\widehat v(k),
\]
where \(\lambda(k)=O(|k|^2)\) and charge preservation gives
\(\widehat v(k)=O(|k|)\).  Squaring yields \(O(|k|^6)\).
\end{proof}

\subsection{Why lower-order forms come next}
\label{sec:e8v70-lower}

\begin{proposition}[Principal convergence alone does not determine the
resolvent]
\label{prop:e8v70-counterexample}
Let \(D\) be bounded and let \(A_n=I\) for all \(n\).  Choose constants
\(c_0\neq c_1\), both nonnegative, and set
\[
 c_n=
 \begin{cases}
 c_0,&n\ \hbox{even},\\
 c_1,&n\ \hbox{odd}.
 \end{cases}
\]
For a nonzero Dirichlet Laplacian eigenfunction \(f\), the solutions of
\begin{equation}
 (-\Delta+c_n)u_n=f
 \label{eq:e8v70-counterexample}
\end{equation}
do not converge, although the principal coefficients converge
identically.
\end{proposition}

\begin{proof}
If \(-\Delta f=\lambda f\), then
\[
 u_n={1\over\lambda+c_n}f.
\]
The even and odd subsequences have distinct limits because
\(c_0\neq c_1\).
\end{proof}

\begin{theorem}[Audit-selected order at V70]
\label{thm:e8v70-direction}
Before using V69 to identify the nonlinear local scalar limit, the
programme must:
\begin{enumerate}[label=\textup{(\roman*)},leftmargin=3.5em]
\item write every lower-order scalar ADM contribution as a bilinear
form on the same energy space;
\item prove a relative form bound strictly below the principal
ellipticity constant;
\item prove convergence of those forms in \(V\to V^*\) norm or a
declared weaker form topology; and
\item retain boundary and charge terms outside the bulk form.
\end{enumerate}
Metric-chart tightness and boundary influence remain necessary
downstream, but they cannot repair an unidentified bulk form.
\end{theorem}

\begin{proof}
\Cref{prop:e8v70-counterexample} shows that the V69 principal
coefficient hypotheses do not determine the full response.  A relative
bound below ellipticity preserves coercivity, while convergence in
operator-form norm is exactly what permits the V69 subtraction argument
to extend.  Boundary and charge data act outside \(H^1_0(D)\) and must
remain separate by V50 and V69.
\end{proof}

\begin{firewall}[Local form control does not create chart probability]
\label{fire:e8v70-chart}
Even perfect lower-order form estimates on a uniformly elliptic patch
do not prove that the gravitational measure remains in that patch.
The V47 translated-well obstruction and metric degeneracy row remain.
\end{firewall}

\subsection{Lower-order target card}
\label{sec:e8v70-card}

\begin{definition}[Lower-order ADM form card, LOAFC]
\label{def:e8v70-LOAFC}
LOAFC consists of:
\begin{enumerate}[label=\textup{(A\arabic*)},leftmargin=4em]
\item the principal metric form \(a_{A_n}\);
\item a bulk lower-order form \(q_n\) containing the declared curvature,
extrinsic, matter, gauge-fixing, and renormalized terms;
\item a uniform relative bound
\(|q_n(u,v)|\leq\alpha\|u\|_V\|v\|_V\) with \(\alpha<m\);
\item convergence \(q_n\to q\) in the chosen form topology;
\item separate boundary, zero-mode, harmonic, and ADM-charge data;
\item compatibility with direct flux sources and local smoothed tests;
\item random-metric probability and uniform-integrability estimates;
and
\item reflection, positivity, Ward, anomaly, and cutoff-removal
contacts.
\end{enumerate}
\end{definition}

\begin{assessment}[Progress at V70]
\label{ass:e8v70-status}
The audit confirms that derivative order controls multiscale clock
cost and prevents importing the Yang--Mills repair into the scalar
constraint.  It also selects the missing deterministic row in V69:
convergence and coercivity of the complete lower-order ADM form.
Principal metric convergence without that row is demonstrably
insufficient.
\end{assessment}

\begin{programme}[Eighth-series V71: relatively bounded ADM
perturbations]
\label{prog:e8v70-V71}
V71 will prove existence, quantitative resolvent stability, and
weak--strong flux pairing convergence for
\[
 b_n(u,v)=a_{A_n}(u,v)+q_n(u,v)
\]
under a uniform relative bound \(\alpha<m\).  It will then specialize
the abstract bound to bounded drift and potential terms on a fixed
patch.
\end{programme}

\begin{firewall}[Claim boundary after eighth-series V70]
\label{fire:e8v70-boundary}
V70 records the compulsory audit, proves the derivative-order
non-transfer and a counterexample to principal-only resolvent
convergence, and selects LOAFC.  It does not yet prove the perturbed
resolvent theorem, establish the actual ADM relative bounds, control
metric-chart probability or boundary influence, construct the positive
carrier, remove cutoffs, or prove the full continuum.
\end{firewall}

\begin{theorem}[Eighth-series V70 result]
\label{thm:e8v70-result}
The bounded-clock Yang--Mills dyadic frame relies on one-derivative
curvature scaling and does not repair the two-derivative,
charge-preserving scalar constraint frame.  For the observable
continuum route, convergence of principal metric coefficients is also
insufficient: the full lower-order ADM form must be uniformly
form-bounded below ellipticity and must converge.
\end{theorem}

\begin{proof}
Use \Cref{thm:e8v70-no-transfer,prop:e8v70-counterexample,
thm:e8v70-direction}.
\end{proof}

\paragraph{Parent-version record.}
V70 was formed from
\texttt{Renewal\_SM\_Einstein\_Continuum\_Research\_Notes\_V69.tex}.
The V69 parent had (27\,399) logical source lines, (967\,041) bytes,
and SHA--256
\[
 \texttt{D113BBB85584F202AC0E101B9A58674C2FA940376460492CB8F2C4B8FBD10584}.
\]
The next compulsory companion audit is V75.

% =============================================================================
% END APPENDED EIGHTH-SERIES V70 MATERIAL
% =============================================================================
% =============================================================================
% BEGIN APPENDED EIGHTH-SERIES V71 MATERIAL
% =============================================================================

\clearpage
\section{Eighth-series V71: relatively bounded lower-order forms and adjoint test resolvents}
\label{sec:v8v71}

\subsection{Abstract perturbed form}
\label{sec:e8v71-abstract}

Let \(D\subset\mathbb R^3\) be a fixed bounded Lipschitz patch and
\(V=H^1_0(D)\), with
\(\|v\|_V=\|\nabla v\|_2\).  Let \(A\) satisfy the V69 ellipticity
window
\[
 mI\preceq A(x)\preceq MI.
\]
Let \(q:V\times V\to\mathbb R\) be a bounded bilinear form satisfying
\begin{equation}
 |q(u,v)|\leq\alpha\|u\|_V\|v\|_V,
 \qquad
 0\leq\alpha<m.
 \label{eq:e8v71-relative}
\end{equation}
Define
\begin{equation}
 b(u,v)
 :=
 \int_D\nabla u\cdot A\nabla v\,dx+q(u,v),
 \qquad
 \gamma:=m-\alpha>0.
 \label{eq:e8v71-form}
\end{equation}
The form need not be symmetric.

\begin{theorem}[Coercivity and flux solvability]
\label{thm:e8v71-flux}
The form \(b\) is bounded and coercive:
\begin{equation}
 b(v,v)\geq\gamma\|v\|_V^2.
 \label{eq:e8v71-coercivity}
\end{equation}
For every \(J\in L^2(D;\mathbb R^3)\), there is a unique \(u\in V\)
such that
\begin{equation}
 b(u,v)=\int_DJ\cdot\nabla v\,dx
 \qquad(v\in V),
 \label{eq:e8v71-flux-equation}
\end{equation}
and
\begin{equation}
 \|u\|_V\leq\gamma^{-1}\|J\|_2.
 \label{eq:e8v71-flux-bound}
\end{equation}
\end{theorem}

\begin{proof}
Ellipticity and \eqref{eq:e8v71-relative} give
\[
 b(v,v)
 \geq m\|v\|_V^2-\alpha\|v\|_V^2
 =\gamma\|v\|_V^2.
\]
Boundedness is immediate from the upper ellipticity bound and
\eqref{eq:e8v71-relative}.  Lax--Milgram gives the unique solution.
Testing with \(v=u\) and applying Cauchy--Schwarz proves
\eqref{eq:e8v71-flux-bound}.
\end{proof}

\subsection{Adjoint test resolvent}
\label{sec:e8v71-adjoint}

\begin{theorem}[Adjoint test representation]
\label{thm:e8v71-adjoint}
For every \(f\in V^*\), there is a unique \(w_f\in V\) satisfying
\begin{equation}
 b(v,w_f)=\langle f,v\rangle
 \qquad(v\in V).
 \label{eq:e8v71-adjoint}
\end{equation}
It obeys
\begin{equation}
 \|w_f\|_V\leq\gamma^{-1}\|f\|_{V^*}.
 \label{eq:e8v71-adjoint-bound}
\end{equation}
For the flux solution \(u\),
\begin{equation}
 \langle f,u\rangle
 =
 \int_DJ\cdot\nabla w_f\,dx,
 \label{eq:e8v71-pairing}
\end{equation}
and hence
\begin{equation}
 |\langle f,u\rangle|
 \leq\gamma^{-1}\|f\|_{V^*}\|J\|_2.
 \label{eq:e8v71-pairing-bound}
\end{equation}
\end{theorem}

\begin{proof}
The adjoint form \(b^*(w,v)=b(v,w)\) is bounded and satisfies
\(b^*(v,v)=b(v,v)\geq\gamma\|v\|_V^2\).  Lax--Milgram gives \(w_f\);
testing with \(v=w_f\) proves
\eqref{eq:e8v71-adjoint-bound}.  Taking \(v=u\) in
\eqref{eq:e8v71-adjoint} and \(v=w_f\) in
\eqref{eq:e8v71-flux-equation} gives
\[
 \langle f,u\rangle=b(u,w_f)
 =\int_DJ\cdot\nabla w_f.
\]
Cauchy--Schwarz proves \eqref{eq:e8v71-pairing-bound}.
\end{proof}

\begin{remark}[Why the adjoint is required]
\label{rem:e8v71-adjoint}
For a drift term, \(b(u,v)\) need not equal \(b(v,u)\).  The symmetric
V69 test equation would then give the wrong pairing identity.  Defining
the test through the second form argument as in
\eqref{eq:e8v71-adjoint} preserves the exact weak--strong interface.
\end{remark}

\subsection{Quantitative convergence of complete forms}
\label{sec:e8v71-convergence}

Let \(b_n\) be forms of the type \eqref{eq:e8v71-form}, all with the
same coercivity floor \(\gamma>0\).  Define the form distance
\begin{equation}
 \eta_n
 :=
 \sup_{\substack{u,v\in V\\\|u\|_V,\|v\|_V\leq1}}
 |b_n(u,v)-b(u,v)|.
 \label{eq:e8v71-form-distance}
\end{equation}

\begin{theorem}[Strong adjoint-resolvent stability]
\label{thm:e8v71-stability}
Let \(f_n,f\in V^*\), and let \(w_n,w\) solve the adjoint equations for
\((b_n,f_n)\) and \((b,f)\).  Then
\begin{equation}
 \|w_n-w\|_V
 \leq
 {1\over\gamma}\|f_n-f\|_{V^*}
 +{\eta_n\over\gamma^2}\|f\|_{V^*}.
 \label{eq:e8v71-stability}
\end{equation}
\end{theorem}

\begin{proof}
Put \(z=w_n-w\).  For every \(v\in V\),
\[
 b_n(v,z)
 =
 \langle f_n-f,v\rangle
 +(b-b_n)(v,w).
\]
Take \(v=z\).  Coercivity, the definition of \(\eta_n\), and
\eqref{eq:e8v71-adjoint-bound} give
\[
 \gamma\|z\|_V^2
 \leq
 \left(
 \|f_n-f\|_{V^*}
 +\eta_n\gamma^{-1}\|f\|_{V^*}
 \right)\|z\|_V.
\]
Divide by \(\gamma\|z\|_V\) when \(z\neq0\).
\end{proof}

\begin{corollary}[Weak flux and complete-form convergence]
\label{cor:e8v71-weak-strong}
Assume
\begin{equation}
 \eta_n\to0,\qquad
 f_n\to f\ \hbox{in }V^*,\qquad
 J_n\rightharpoonup J\ \hbox{in }L^2(D).
 \label{eq:e8v71-data-convergence}
\end{equation}
Let \(u_n\) solve
\[
 b_n(u_n,v)=\int_DJ_n\cdot\nabla v.
\]
Then
\begin{equation}
 \langle f_n,u_n\rangle
 \longrightarrow
 \int_DJ\cdot\nabla w_f\,dx.
 \label{eq:e8v71-limit}
\end{equation}
\end{corollary}

\begin{proof}
\Cref{thm:e8v71-stability} gives
\(w_n\to w_f\) strongly in \(V\).  The adjoint pairing identity gives
\[
 \langle f_n,u_n\rangle
 =\int_DJ_n\cdot\nabla w_n.
\]
Pair weak convergence of \(J_n\) with strong convergence of
\(\nabla w_n\).
\end{proof}

\begin{corollary}[Coefficient and lower-form distance]
\label{cor:e8v71-distance}
If
\[
 b_n=a_{A_n}+q_n,\qquad b=a_A+q,
\]
then
\begin{equation}
 \eta_n
 \leq
 \|A_n-A\|_{L^\infty}
 +\|q_n-q\|_{{\cal B}(V,V^*)}.
 \label{eq:e8v71-distance}
\end{equation}
\end{corollary}

\begin{proof}
For unit \(V\)-norm arguments,
\[
 |a_{A_n}(u,v)-a_A(u,v)|
 \leq\|A_n-A\|_\infty.
\]
Add the operator-form difference of \(q_n\) and \(q\).
\end{proof}

\subsection{Bounded drift and potential}
\label{sec:e8v71-drift}

Let \(C_P(D)\) be the Poincaré constant
\(\|v\|_2\leq C_P\|\nabla v\|_2\).  Consider
\begin{equation}
 q_{\beta,c}(u,v)
 :=
 \int_D\beta\cdot\nabla u\,v\,dx
 +\int_Dcuv\,dx,
 \label{eq:e8v71-drift-form}
\end{equation}
with \(\beta\in L^\infty(D;\mathbb R^3)\) and
\(c\in L^\infty(D)\).

\begin{theorem}[Explicit sufficient lower-order bound]
\label{thm:e8v71-drift}
The form \eqref{eq:e8v71-drift-form} satisfies
\begin{equation}
 |q_{\beta,c}(u,v)|
 \leq
 \alpha_{\beta,c}\|u\|_V\|v\|_V,
 \qquad
 \alpha_{\beta,c}
 :=
 C_P\|\beta\|_\infty+C_P^2\|c\|_\infty.
 \label{eq:e8v71-drift-bound}
\end{equation}
If \(\alpha_{\beta,c}<m\), then
\Cref{thm:e8v71-flux,thm:e8v71-adjoint,
thm:e8v71-stability} apply with
\(\gamma=m-\alpha_{\beta,c}\).
\end{theorem}

\begin{proof}
Hölder and Poincaré give
\[
 \left|\int\beta\cdot\nabla u\,v\right|
 \leq
 \|\beta\|_\infty\|\nabla u\|_2\|v\|_2
 \leq C_P\|\beta\|_\infty\|u\|_V\|v\|_V
\]
and
\[
 \left|\int cuv\right|
 \leq
 \|c\|_\infty\|u\|_2\|v\|_2
 \leq C_P^2\|c\|_\infty\|u\|_V\|v\|_V.
\]
Add the bounds.
\end{proof}

\begin{corollary}[Convergence of bounded lower-order coefficients]
\label{cor:e8v71-drift-convergence}
If
\[
 \beta_n\to\beta,\qquad c_n\to c
 \quad\hbox{in }L^\infty(D),
\]
then
\begin{equation}
 \|q_{\beta_n,c_n}-q_{\beta,c}\|_{{\cal B}(V,V^*)}
 \leq
 C_P\|\beta_n-\beta\|_\infty
 +C_P^2\|c_n-c\|_\infty
 \longrightarrow0.
 \label{eq:e8v71-drift-convergence}
\end{equation}
\end{corollary}

\begin{proof}
Repeat the proof of \Cref{thm:e8v71-drift} for the coefficient
differences.
\end{proof}

\begin{firewall}[The absolute bound is sufficient, not necessary]
\label{fire:e8v71-sufficient}
A divergence-free drift with suitable boundary behavior may be
skew-symmetric and contribute nothing to \(q(v,v)\), while
\eqref{eq:e8v71-drift-bound} still charges its full norm.  Failure of
\(\alpha_{\beta,c}<m\) does not prove loss of coercivity; it means the
coarse absolute-value estimate is inconclusive and the symmetric,
antisymmetric, and negative-potential parts must be analyzed separately.
\end{firewall}

\begin{firewall}[Patch size enters through Poincaré]
\label{fire:e8v71-patch}
The constant \(C_P(D)\) grows with the diameter of a dilated patch.
V71 is a fixed-local-patch theorem.  Applying
\eqref{eq:e8v71-drift-bound} on volume-growing domains would reintroduce
infrared factors and must not be called uniform.
\end{firewall}

\subsection{Lower-order continuum card}
\label{sec:e8v71-card}

\begin{definition}[Perturbed local resolvent card, PLRC]
\label{def:e8v71-PLRC}
PLRC consists of:
\begin{enumerate}[label=\textup{(P\arabic*)},leftmargin=4em]
\item a common local energy space and principal ellipticity floor \(m\);
\item a complete, possibly nonsymmetric lower-order form;
\item a coercivity proof with floor \(\gamma>0\), using structure rather
than absolute bounds when necessary;
\item convergence in form norm or a proved Mosco/sectorial substitute;
\item the adjoint test-resolvent representation
\eqref{eq:e8v71-pairing};
\item weak compactness of direct flux primitives;
\item boundary, charge, metric-chart, Ward, and composite-operator
control; and
\item assembly into a positive reflected continuum state.
\end{enumerate}
\end{definition}

\begin{assessment}[Progress at V71]
\label{ass:e8v71-status}
The V69 weak--strong continuum interface now survives a complete
relatively bounded lower-order form, including nonsymmetric drift.
The adjoint test resolvent restores the exact flux pairing, and a
quantitative form-distance estimate gives strong test convergence.
Bounded drift and potential coefficients provide an explicit
sufficient criterion on fixed patches.  Identifying and bounding the
actual ADM coefficients remains the next physical step.
\end{assessment}

\begin{programme}[Eighth-series V72: ADM scalar lower-order
decomposition]
\label{prog:e8v71-V72}
V72 will linearize the scalar Hamiltonian constraint about a general
uniformly elliptic background patch and separate its Laplace--Beltrami
principal part, curvature potential, extrinsic-curvature coupling, and
matter variation.  Each coefficient will be assigned to the symmetric,
skew, source, or boundary ledger before any form-smallness claim is
made.
\end{programme}

\begin{firewall}[Claim boundary after eighth-series V71]
\label{fire:e8v71-boundary}
V71 proves coercive solvability, adjoint test representation,
quantitative complete-form convergence, and explicit bounded
drift/potential estimates.  It does not derive the actual nonlinear ADM
lower-order coefficients, prove their form smallness, control
metric-chart probability or global boundary influence, construct the
positive carrier, remove cutoffs, or prove the full continuum.
\end{firewall}

\begin{theorem}[Eighth-series V71 result]
\label{thm:e8v71-result}
If the complete local forms
\[
 b_n=a_{A_n}+q_n
\]
share coercivity floor \(\gamma>0\), converge in
\({\cal B}(V,V^*)\), the tests converge strongly in \(V^*\), and the
direct flux primitives converge weakly in \(L^2\), then all local scalar
response pairings converge.  Nonsymmetric lower-order terms are handled
by the adjoint test resolvent.
\end{theorem}

\begin{proof}
Use \Cref{thm:e8v71-adjoint,thm:e8v71-stability,
cor:e8v71-weak-strong}.
\end{proof}

\paragraph{Parent-version record.}
V71 was formed from
\texttt{Renewal\_SM\_Einstein\_Continuum\_Research\_Notes\_V70.tex}.
The V70 parent had (27\,629) logical source lines, (975\,872) bytes,
and SHA--256
\[
 \texttt{8A03BEB49B2986478A3771EABE2C73B368340F05D6DFC4F1600049E966215555}.
\]
The next compulsory companion audit is V75.

% =============================================================================
% END APPENDED EIGHTH-SERIES V71 MATERIAL
% =============================================================================
% =============================================================================
% BEGIN APPENDED EIGHTH-SERIES V72 MATERIAL
% =============================================================================

\clearpage
\section{Eighth-series V72: conformal linearization of the ADM scalar constraint}
\label{sec:v8v72}

\subsection{Declared convention}
\label{sec:e8v72-convention}

On a \(d\)-dimensional spatial patch, use the Hamiltonian-constraint
convention
\begin{equation}
 {\cal H}(g,K,\Psi)
 :=
 R_g+\tau^2-|K|_g^2-2\kappa_G\rho(g,\Psi)=0,
 \qquad
 \tau:=\operatorname{tr}_gK.
 \label{eq:e8v72-constraint}
\end{equation}
Here \(K\) is treated as a covariant two-tensor.  The variation below
holds \(K\) fixed in the metric derivative and then adds an independent
tensor variation \(\dot K\).  A canonical calculation holding ADM
momentum fixed gives a different algebraic potential and must be
derived separately.

\begin{firewall}[Convention dependence is explicit]
\label{fire:e8v72-convention}
The signs in \eqref{eq:e8v72-constraint}, the choice of \(K\) rather
than canonical momentum as independent variable, and the definition of
\(\rho\) are part of V72.  The resulting lower-order coefficient may
not be transplanted into a different ADM convention without repeating
the variation.
\end{firewall}

\subsection{Metric and extrinsic-curvature variations}
\label{sec:e8v72-variation}

\begin{lemma}[Scalar-curvature linearization]
\label{lem:e8v72-curvature}
For a symmetric metric variation \(h=\dot g\),
\begin{equation}
 DR_g[h]
 =
 \nabla^i\nabla^jh_{ij}
 -\Delta_g(\operatorname{tr}_gh)
 -\langle\operatorname{Ric}_g,h\rangle_g.
 \label{eq:e8v72-DR}
\end{equation}
For the conformal direction \(h=2\phi g\),
\begin{equation}
 DR_g[2\phi g]
 =
 -2(d-1)\Delta_g\phi-2R_g\phi.
 \label{eq:e8v72-conformal-R}
\end{equation}
\end{lemma}

\begin{proof}
The connection variation is
\[
 \dot\Gamma^k{}_{ij}
 ={1\over2}g^{k\ell}
 \bigl(\nabla_ih_{j\ell}+\nabla_jh_{i\ell}
 -\nabla_\ell h_{ij}\bigr).
\]
Insert it in
\(\dot R_{ij}=\nabla_k\dot\Gamma^k{}_{ij}
-\nabla_j\dot\Gamma^k{}_{ik}\), and include
\(\dot g^{ij}=-h^{ij}\) when tracing.  Contraction gives
\eqref{eq:e8v72-DR}.  If \(h=2\phi g\), then
\[
 \operatorname{tr}_gh=2d\phi,\qquad
 \nabla^i\nabla^jh_{ij}=2\Delta_g\phi,\qquad
 \langle\operatorname{Ric},h\rangle=2R_g\phi.
\]
Substitution proves \eqref{eq:e8v72-conformal-R}.
\end{proof}

\begin{lemma}[Extrinsic scalar linearization]
\label{lem:e8v72-K}
Let
\[
 {\cal S}(g,K):=\tau^2-|K|_g^2.
\]
For a simultaneous variation \((h,\dot K)\),
\begin{align}
 D{\cal S}_{(g,K)}[h,\dot K]
 &=
 2\langle h,K\circ K-\tau K\rangle_g
 \notag\\
 &\quad
 -2\langle K,\dot K\rangle_g
 +2\tau\operatorname{tr}_g\dot K,
 \label{eq:e8v72-DS}
\end{align}
where
\((K\circ K)_{ij}=K_i{}^kK_{kj}\).
In particular, for \(h=2\phi g\),
\begin{equation}
 D{\cal S}[2\phi g,\dot K]
 =
 4\phi(|K|^2-\tau^2)
 -2\langle K,\dot K\rangle
 +2\tau\operatorname{tr}\dot K.
 \label{eq:e8v72-conformal-S}
\end{equation}
\end{lemma}

\begin{proof}
At fixed \(K\),
\[
 D\tau[h]=-\langle h,K\rangle,\qquad
 D|K|^2[h]=-2\langle h,K\circ K\rangle.
\]
At fixed metric,
\[
 D\tau[\dot K]=\operatorname{tr}\dot K,\qquad
 D|K|^2[\dot K]=2\langle K,\dot K\rangle.
\]
Differentiate \(\tau^2-|K|^2\) and add the two independent
contributions.  With \(h=2\phi g\),
\[
 \langle g,K\circ K\rangle=|K|^2,\qquad
 \langle g,K\rangle=\tau,
\]
which gives \eqref{eq:e8v72-conformal-S}.
\end{proof}

\subsection{The complete conformal scalar equation}
\label{sec:e8v72-scalar}

Write the matter variation as
\begin{equation}
 D\rho[2\phi g,\dot\Psi]
 =
 D_g\rho[2\phi g]+D_\Psi\rho[\dot\Psi].
 \label{eq:e8v72-matter-variation}
\end{equation}

\begin{theorem}[Exact conformal ADM linearization]
\label{thm:e8v72-linearization}
The first variation of \eqref{eq:e8v72-constraint} in the direction
\((2\phi g,\dot K,\dot\Psi)\) is
\begin{align}
 D{\cal H}[2\phi g,\dot K,\dot\Psi]
 &=
 -2(d-1)\Delta_g\phi
 -2R_g\phi
 +4(|K|^2-\tau^2)\phi
 \notag\\
 &\quad
 -2\langle K,\dot K\rangle
 +2\tau\operatorname{tr}\dot K
 -2\kappa_GD_g\rho[2\phi g]
 -2\kappa_GD_\Psi\rho[\dot\Psi].
 \label{eq:e8v72-linearization}
\end{align}
\end{theorem}

\begin{proof}
Add \Cref{lem:e8v72-curvature,lem:e8v72-K} and the variation of the
last term in \eqref{eq:e8v72-constraint}.
\end{proof}

Assume on the declared matter chart that the pure conformal metric
variation is algebraic,
\begin{equation}
 D_g\rho[2\phi g]=m_\rho\,\phi.
 \label{eq:e8v72-matter-multiplier}
\end{equation}
Define
\begin{align}
 c_{\rm ADM}
 &:=
 -2R_g+4(|K|^2-\tau^2)-2\kappa_Gm_\rho,
 \label{eq:e8v72-potential}\\
 F_{\rm phys}
 &:=
 -2\langle K,\dot K\rangle
 +2\tau\operatorname{tr}\dot K
 -2\kappa_GD_\Psi\rho[\dot\Psi].
 \label{eq:e8v72-source}
\end{align}

\begin{corollary}[Principal, potential, and source ledger]
\label{cor:e8v72-ledger}
Under \eqref{eq:e8v72-matter-multiplier}, the linearized constraint is
\begin{equation}
 \bigl[-2(d-1)\Delta_g+c_{\rm ADM}\bigr]\phi
 +F_{\rm phys}=0.
 \label{eq:e8v72-operator}
\end{equation}
Thus:
\begin{enumerate}[label=\textup{(\roman*)},leftmargin=3.5em]
\item \(-2(d-1)\Delta_g\) is the principal scalar form;
\item \(c_{\rm ADM}\) is a symmetric multiplication form; and
\item \(\dot K\) and the independent matter variation enter the source
ledger.
\end{enumerate}
\end{corollary}

\begin{proof}
Regroup \eqref{eq:e8v72-linearization} using
\eqref{eq:e8v72-potential}--\eqref{eq:e8v72-source}.
\end{proof}

\begin{firewall}[Matter metric variation may contain more structure]
\label{fire:e8v72-matter}
Equation \eqref{eq:e8v72-matter-multiplier} is a declared chart
hypothesis.  If a regulator, spin connection, improvement term, or
renormalized composite stress tensor makes \(D_g\rho\) contain
derivatives of \(\phi\), those contributions must be assigned to the
principal, drift, or boundary form.  They may not be hidden in
\(m_\rho\).
\end{firewall}

\subsection{Coercivity of the conformal scalar form}
\label{sec:e8v72-coercivity}

On a fixed patch \(D\), define
\begin{equation}
 b_{\rm ADM}(u,v)
 =
 2(d-1)\int_D\langle\nabla u,\nabla v\rangle_g\,dV_g
 +\int_Dc_{\rm ADM}uv\,dV_g.
 \label{eq:e8v72-form}
\end{equation}
Let \(C_{P,g}\) satisfy
\[
 \|v\|_{L^2_g}\leq C_{P,g}\|\nabla v\|_{L^2_g}.
\]

\begin{theorem}[Negative-potential coercivity criterion]
\label{thm:e8v72-coercivity}
Let
\[
 c_{\rm ADM,-}:=\max\{-c_{\rm ADM},0\}.
\]
If
\begin{equation}
 \|c_{\rm ADM,-}\|_\infty C_{P,g}^2
 <2(d-1),
 \label{eq:e8v72-coercivity-condition}
\end{equation}
then
\begin{equation}
 b_{\rm ADM}(v,v)
 \geq
 \left(
 2(d-1)-\|c_{\rm ADM,-}\|_\infty C_{P,g}^2
 \right)
 \|\nabla v\|_{L^2_g}^2.
 \label{eq:e8v72-coercivity}
\end{equation}
\end{theorem}

\begin{proof}
The positive part of \(c_{\rm ADM}\) can be discarded in a lower bound.
For the negative part,
\[
 \int c_{\rm ADM}v^2dV_g
 \geq
 -\|c_{\rm ADM,-}\|_\infty\|v\|_{L^2_g}^2
 \geq
 -\|c_{\rm ADM,-}\|_\infty C_{P,g}^2
 \|\nabla v\|_{L^2_g}^2.
\]
Add the principal term.
\end{proof}

\begin{corollary}[Fixed-patch insertion into V71]
\label{cor:e8v72-V71}
If \eqref{eq:e8v72-coercivity-condition} holds uniformly along a
regulator sequence and \(c_{{\rm ADM},n}\to c_{\rm ADM}\) in
\(L^\infty\), then the potential forms converge in
\({\cal B}(H^1_0,H^{-1})\), and the V71 weak--strong pairing theorem
applies together with convergence of the principal metric forms.
\end{corollary}

\begin{proof}
For \(q_n(u,v)=\int c_{{\rm ADM},n}uv\,dV_{g_n}\), uniform metric
control and Poincaré give
\[
 |q_n(u,v)-q(u,v)|
 \leq C\|c_{{\rm ADM},n}-c_{\rm ADM}\|_\infty
 \|u\|_{H^1_0}\|v\|_{H^1_0},
\]
with the volume-density difference included in the principal metric
convergence.  The coercivity floor is
\eqref{eq:e8v72-coercivity}; apply
\Cref{cor:e8v71-weak-strong}.
\end{proof}

\begin{firewall}[Curvature convergence is stronger than metric
convergence]
\label{fire:e8v72-curvature-convergence}
The coefficient \(c_{\rm ADM}\) contains \(R_g\).
Strong \(L^\infty\) convergence of \(g_n\) does not imply convergence,
or even boundedness, of \(R_{g_n}\), which contains two derivatives.
V69 metric-coefficient convergence therefore does not prove the
potential-form convergence required in
\Cref{cor:e8v72-V71}.
\end{firewall}

\begin{firewall}[The potential need not be small]
\label{fire:e8v72-smallness}
Condition \eqref{eq:e8v72-coercivity-condition} is sufficient on a
fixed patch.  Its failure does not prove that the scalar operator has a
negative mode; positive parts, boundary conditions, constraint
couplings, or a sharper spectral estimate may restore coercivity.
Conversely, the background constraint
\(R+\tau^2-|K|^2=2\kappa_G\rho\) does not by itself force
\(c_{\rm ADM}\geq0\).
\end{firewall}

\subsection{ADM scalar form card}
\label{sec:e8v72-card}

\begin{definition}[Conformal ADM form card, CAFC]
\label{def:e8v72-CAFC}
CAFC consists of:
\begin{enumerate}[label=\textup{(C\arabic*)},leftmargin=4em]
\item a declared Hamiltonian-constraint sign and canonical-variable
convention;
\item the exact scalar-curvature and extrinsic variations
\eqref{eq:e8v72-conformal-R}--\eqref{eq:e8v72-conformal-S};
\item a complete matter metric-variation ledger;
\item a principal, drift, potential, source, and boundary decomposition;
\item uniform coercivity of the resulting complete form;
\item convergence of \(R_g\), \(K\), \(\rho\), and renormalized
coefficients in a form-compatible topology;
\item direct flux-source, current-alignment, charge, and boundary
control; and
\item metric-chart tightness, reflection, positivity, and continuum
reconstruction.
\end{enumerate}
\end{definition}

\begin{assessment}[Progress at V72]
\label{ass:e8v72-status}
The abstract V71 lower-order form now has a concrete ADM scalar
decomposition under an explicit convention.  The conformal principal
operator is \(-2(d-1)\Delta_g\); curvature, extrinsic curvature, and the
metric derivative of matter form the multiplication potential; and
independent \(K\) and matter variations remain sources.  A sharp
negative-part condition gives fixed-patch coercivity.  The newly exposed
bottleneck is curvature and matter-coefficient convergence, which is
not supplied by uniform metric convergence.
\end{assessment}

\begin{programme}[Eighth-series V73: curvature potential in weak form]
\label{prog:e8v72-V73}
V73 will avoid demanding \(L^\infty\) convergence of scalar curvature.
It will rewrite the curvature-potential pairing using the connection
difference and derive an \(H^{-1}\) or form-norm estimate from stronger
metric Sobolev convergence, with all boundary terms explicit.
\end{programme}

\begin{firewall}[Claim boundary after eighth-series V72]
\label{fire:e8v72-boundary}
V72 proves the conformal linearization of the declared ADM constraint,
its principal/potential/source ledger, and a fixed-patch
negative-potential coercivity criterion.  It does not prove the matter
multiplier hypothesis, curvature-potential convergence, canonical
momentum version, metric-chart tightness, boundary convergence,
positive carrier construction, cutoff removal, or the full continuum.
\end{firewall}

\begin{theorem}[Eighth-series V72 result]
\label{thm:e8v72-result}
Under the declared \(K\)-variable convention and algebraic matter
metric variation, the conformal scalar linearization is
\[
 \left[
 -2(d-1)\Delta_g
 -2R_g+4(|K|^2-\tau^2)-2\kappa_Gm_\rho
 \right]\phi
 +F_{\rm phys}=0.
\]
On a fixed patch its form is coercive if
\(\|c_{\rm ADM,-}\|_\infty C_{P,g}^2<2(d-1)\).
\end{theorem}

\begin{proof}
Use \Cref{thm:e8v72-linearization,cor:e8v72-ledger,
thm:e8v72-coercivity}.
\end{proof}

\paragraph{Parent-version record.}
V72 was formed from
\texttt{Renewal\_SM\_Einstein\_Continuum\_Research\_Notes\_V71.tex}.
The V71 parent had (28\,034) logical source lines, (987\,270) bytes,
and SHA--256
\[
 \texttt{4203B5C94551A019B8E6A7251F3E0EF11F8637DC6484F5DA31A09A89BFE38C5A}.
\]
The next compulsory companion audit is V75.

% =============================================================================
% END APPENDED EIGHTH-SERIES V72 MATERIAL
% =============================================================================
% =============================================================================
% BEGIN APPENDED EIGHTH-SERIES V73 MATERIAL
% =============================================================================

\clearpage
\section{Eighth-series V73: scalar curvature as a first-derivative energy form}
\label{sec:v8v73}

\subsection{The curvature-density decomposition}
\label{sec:e8v73-density}

Let \(D\subset\mathbb R^3\) be a bounded Lipschitz coordinate patch.
For a smooth uniformly positive metric \(g\), define
\begin{align}
 Q_g
 &:=
 \sqrt{\det g}\,g^{ij}
 \left(
 \Gamma^k{}_{i\ell}\Gamma^\ell{}_{jk}
 -\Gamma^k{}_{ij}\Gamma^\ell{}_{k\ell}
 \right),
 \label{eq:e8v73-Q}\\
 V_g^k
 &:=
 \sqrt{\det g}
 \left(
 g^{ij}\Gamma^k{}_{ij}
 -g^{ik}\Gamma^j{}_{ij}
 \right).
 \label{eq:e8v73-V}
\end{align}

\begin{lemma}[Einstein density identity]
\label{lem:e8v73-density}
In the declared coordinate chart,
\begin{equation}
 \sqrt{\det g}\,R_g=Q_g+\partial_kV_g^k.
 \label{eq:e8v73-density}
\end{equation}
\end{lemma}

\begin{proof}
Start from
\[
 R_g
 =
 g^{ij}\left(
 \partial_k\Gamma^k{}_{ij}
 -\partial_j\Gamma^k{}_{ik}
 +\Gamma^k{}_{k\ell}\Gamma^\ell{}_{ij}
 -\Gamma^k{}_{j\ell}\Gamma^\ell{}_{ik}
 \right).
\]
Apply the product rule to the two derivative terms after multiplication
by \(\sqrt{\det g}\).  Using metric compatibility to replace
derivatives of \(\sqrt{\det g}\,g^{ij}\), the remaining quadratic
connection terms combine into \eqref{eq:e8v73-Q}, while the total
derivative is exactly \(\partial_kV_g^k\).
\end{proof}

For \(u,v\in H^1_0(D)\), define the weak curvature form
\begin{equation}
 {\mathfrak r}_g(u,v)
 :=
 \int_DQ_guv\,dx
 -\int_DV_g\cdot\nabla(uv)\,dx.
 \label{eq:e8v73-curvature-form}
\end{equation}

\begin{corollary}[Agreement with smooth scalar curvature]
\label{cor:e8v73-agreement}
For smooth \(g\) and \(u,v\in C_c^\infty(D)\),
\begin{equation}
 {\mathfrak r}_g(u,v)
 =
 \int_DR_guv\,dV_g.
 \label{eq:e8v73-agreement}
\end{equation}
\end{corollary}

\begin{proof}
Multiply \eqref{eq:e8v73-density} by \(uv\) and integrate by parts.
The compact support removes the boundary term.
\end{proof}

\subsection{Critical three-dimensional form bound}
\label{sec:e8v73-bound}

Assume
\begin{equation}
 m_gI\preceq g(x)\preceq M_gI,
 \qquad
 g\in W^{1,3}(D)\cap L^\infty(D).
 \label{eq:e8v73-metric-class}
\end{equation}
The Christoffel symbols are then in \(L^3\), so
\(V_g\in L^3\) and \(Q_g\in L^{3/2}\).

\begin{lemma}[Coefficient estimates]
\label{lem:e8v73-coefficients}
There is \(C=C(m_g,M_g)\) such that
\begin{align}
 \|V_g\|_{L^3}
 &\leq C\|\partial g\|_{L^3},
 \label{eq:e8v73-V-bound}\\
 \|Q_g\|_{L^{3/2}}
 &\leq C\|\partial g\|_{L^3}^2.
 \label{eq:e8v73-Q-bound}
\end{align}
\end{lemma}

\begin{proof}
The formula
\[
 \Gamma^k{}_{ij}
 ={1\over2}g^{k\ell}
 (\partial_ig_{j\ell}+\partial_jg_{i\ell}
 -\partial_\ell g_{ij})
\]
and the uniform inverse-metric bound give
\(|\Gamma|\leq C|\partial g|\).  The determinant, inverse metric, and
their products are uniformly bounded on the spectral window
\eqref{eq:e8v73-metric-class}.  Equation \eqref{eq:e8v73-V} is linear
in \(\Gamma\), while \eqref{eq:e8v73-Q} is quadratic.  Apply Hölder.
\end{proof}

\begin{theorem}[Weak scalar-curvature form bound]
\label{thm:e8v73-form-bound}
Let \(C_S(D)\) be the Sobolev constant
\(\|u\|_6\leq C_S\|\nabla u\|_2\).  Under
\eqref{eq:e8v73-metric-class},
\begin{align}
 |{\mathfrak r}_g(u,v)|
 &\leq
 C_S^2
 \left(
 \|Q_g\|_{3/2}+2\|V_g\|_3
 \right)
 \|u\|_{H^1_0}\|v\|_{H^1_0}
 \notag\\
 &\leq
 C
 \left(
 \|\partial g\|_3+\|\partial g\|_3^2
 \right)
 \|u\|_{H^1_0}\|v\|_{H^1_0}.
 \label{eq:e8v73-form-bound}
\end{align}
Thus scalar curvature defines an element of
\({\cal B}(H^1_0,H^{-1})\) using only first weak derivatives of the
metric.
\end{theorem}

\begin{proof}
Hölder and Sobolev give
\[
 \left|\int Q_guv\right|
 \leq\|Q_g\|_{3/2}\|u\|_6\|v\|_6.
\]
Since
\(\nabla(uv)=u\nabla v+v\nabla u\),
\[
 \left|\int V_g\cdot\nabla(uv)\right|
 \leq
 \|V_g\|_3
 \left(
 \|u\|_6\|\nabla v\|_2
 +\|v\|_6\|\nabla u\|_2
 \right).
\]
Apply the Sobolev inequality and then
\Cref{lem:e8v73-coefficients}.
\end{proof}

\begin{remark}[Why the exponent three is natural]
\label{rem:e8v73-critical}
In three dimensions, \(H^1_0\hookrightarrow L^6\).
The quadratic connection coefficient must therefore lie in
\(L^{3/2}\), while the divergence coefficient paired with
\(u\nabla v\) must lie in \(L^3\).  The assumption
\(\partial g\in L^3\) produces exactly these two spaces.
\end{remark}

\subsection{Strong form convergence}
\label{sec:e8v73-convergence}

\begin{theorem}[Curvature-form convergence from \(W^{1,3}\)]
\label{thm:e8v73-convergence}
Let \(g_n,g\) satisfy a common spectral window
\eqref{eq:e8v73-metric-class} and assume
\begin{equation}
 g_n\to g\quad\hbox{in }L^\infty(D),
 \qquad
 \partial g_n\to\partial g\quad\hbox{in }L^3(D).
 \label{eq:e8v73-metric-convergence}
\end{equation}
Then
\begin{align}
 V_{g_n}&\to V_g\quad\hbox{in }L^3(D),
 \label{eq:e8v73-V-convergence}\\
 Q_{g_n}&\to Q_g\quad\hbox{in }L^{3/2}(D),
 \label{eq:e8v73-Q-convergence}
\end{align}
and
\begin{equation}
 \|{\mathfrak r}_{g_n}-{\mathfrak r}_g\|
 _{{\cal B}(H^1_0,H^{-1})}
 \longrightarrow0.
 \label{eq:e8v73-form-convergence}
\end{equation}
\end{theorem}

\begin{proof}
Uniform nondegeneracy and \(L^\infty\) convergence give convergence of
\(g_n^{-1}\) and \(\sqrt{\det g_n}\) in \(L^\infty\).
The Christoffel formula then gives
\[
 \Gamma(g_n)\to\Gamma(g)\quad\hbox{in }L^3.
\]
The expression \(V_g\) is an \(L^\infty\) coefficient times one
Christoffel symbol, proving \eqref{eq:e8v73-V-convergence}.  For the
quadratic term,
\[
 \|\Gamma_n\Gamma_n-\Gamma\Gamma\|_{3/2}
 \leq
 \|\Gamma_n-\Gamma\|_3
 (\|\Gamma_n\|_3+\|\Gamma\|_3),
\]
and the remaining coefficients converge in \(L^\infty\), proving
\eqref{eq:e8v73-Q-convergence}.  Repeat the Hölder--Sobolev estimate in
\Cref{thm:e8v73-form-bound} for the coefficient differences to obtain
\eqref{eq:e8v73-form-convergence}.
\end{proof}

\begin{firewall}[No second-derivative convergence was proved]
\label{fire:e8v73-second}
The theorem defines and controls the integrated curvature form.  It
does not assert that \(R_{g_n}\) converges pointwise, in \(L^p\), or as
a bounded multiplication operator.  The divergence part has been moved
onto the test product, which is precisely why only first metric
derivatives are required.
\end{firewall}

\subsection{Extrinsic and matter multiplication forms}
\label{sec:e8v73-other}

\begin{theorem}[Quadratic extrinsic-curvature form]
\label{thm:e8v73-K}
If \(K\in L^3(D)\) and \(g\) obeys the uniform spectral window, then
\[
 c_K:=4(|K|_g^2-\tau^2)\in L^{3/2}(D)
\]
and
\begin{equation}
 \left|\int_Dc_Kuv\,dV_g\right|
 \leq
 C\|K\|_3^2
 \|u\|_{H^1_0}\|v\|_{H^1_0}.
 \label{eq:e8v73-K-bound}
\end{equation}
If \(K_n\to K\) in \(L^3\) and \(g_n\to g\) in \(L^\infty\), the
corresponding forms converge in
\({\cal B}(H^1_0,H^{-1})\).
\end{theorem}

\begin{proof}
Uniform metric equivalence gives
\(|c_K|\leq C|K|^2\), so
\(\|c_K\|_{3/2}\leq C\|K\|_3^2\).
Hölder with \(u,v\in L^6\) proves
\eqref{eq:e8v73-K-bound}.  For convergence use
\[
 \||K_n|^2-|K|^2\|_{3/2}
 \leq(\|K_n\|_3+\|K\|_3)\|K_n-K\|_3
\]
together with the \(L^\infty\) metric convergence.
\end{proof}

\begin{corollary}[Matter multiplier form]
\label{cor:e8v73-matter}
If the V72 coefficient \(m_\rho\in L^{3/2}(D)\), then
\begin{equation}
 \left|
 \int_Dm_\rho uv\,dV_g
 \right|
 \leq
 C\|m_\rho\|_{3/2}
 \|u\|_{H^1_0}\|v\|_{H^1_0}.
 \label{eq:e8v73-matter}
\end{equation}
Strong \(L^{3/2}\) convergence of \(m_{\rho,n}\), together with strong
metric \(L^\infty\) convergence, gives form-norm convergence.
\end{corollary}

\begin{proof}
Apply Hölder and \(H^1_0\hookrightarrow L^6\); repeat for differences.
\end{proof}

\subsection{The complete weak ADM potential}
\label{sec:e8v73-complete}

Define the V72 lower-order weak form by
\begin{equation}
 q_{{\rm ADM},g,K,m_\rho}(u,v)
 :=
 -2{\mathfrak r}_g(u,v)
 +4\int_D(|K|^2-\tau^2)uv\,dV_g
 -2\kappa_G\int_Dm_\rho uv\,dV_g.
 \label{eq:e8v73-complete-form}
\end{equation}

\begin{theorem}[Complete first-derivative ADM form compiler]
\label{thm:e8v73-complete}
Under the hypotheses above,
\begin{equation}
 |q_{{\rm ADM}}(u,v)|
 \leq
 \alpha_{\rm ADM}
 \|u\|_{H^1_0}\|v\|_{H^1_0},
 \label{eq:e8v73-complete-bound}
\end{equation}
where
\begin{equation}
 \alpha_{\rm ADM}
 \leq
 C\left(
 \|\partial g\|_3+\|\partial g\|_3^2
 +\|K\|_3^2+\kappa_G\|m_\rho\|_{3/2}
 \right).
 \label{eq:e8v73-alpha}
\end{equation}
If \(g_n\to g\) as in
\eqref{eq:e8v73-metric-convergence},
\(K_n\to K\) in \(L^3\), and
\(m_{\rho,n}\to m_\rho\) in \(L^{3/2}\), then
\begin{equation}
 q_{{\rm ADM},n}\to q_{\rm ADM}
 \quad\hbox{in }{\cal B}(H^1_0,H^{-1}).
 \label{eq:e8v73-complete-convergence}
\end{equation}
If the principal form has ellipticity floor \(m_0\) and
\(\sup_n\alpha_{{\rm ADM},n}<m_0\), the V71 perturbed-resolvent theorem
applies.
\end{theorem}

\begin{proof}
Add \Cref{thm:e8v73-form-bound,thm:e8v73-K,
cor:e8v73-matter}.  Their convergence statements prove
\eqref{eq:e8v73-complete-convergence}.  The last assertion is
\Cref{thm:e8v71-flux,cor:e8v71-weak-strong}.
\end{proof}

\begin{firewall}[Small total form norm is only a sufficient route]
\label{fire:e8v73-smallness}
The bound \(\alpha_{\rm ADM}<m_0\) charges positive and antisymmetric
pieces by absolute value.  Its failure does not prove noncoercivity.
It signals that the negative part, constraint coupling, or a sharper
spectral estimate must be separated before V71 is applied.
\end{firewall}

\begin{firewall}[Coordinate and boundary ledger]
\label{fire:e8v73-coordinate}
For smooth metrics the sum
\eqref{eq:e8v73-curvature-form} is coordinate invariant, although
\(Q_g\) and \(V_g\) separately are not.  At weak regularity the chart
and transition regularity are part of the definition.  For tests with
nonzero trace, integration by parts also produces
\(\int_{\partial D}uv\,V_g\cdot n\), which belongs to the boundary
compiler rather than the bulk form.
\end{firewall}

\subsection{Weak curvature card}
\label{sec:e8v73-card}

\begin{definition}[Weak ADM curvature-form card, WACFC]
\label{def:e8v73-WACFC}
WACFC consists of:
\begin{enumerate}[label=\textup{(W\arabic*)},leftmargin=4em]
\item uniformly nondegenerate local metrics in
\(W^{1,3}\cap L^\infty\);
\item the exact density decomposition
\eqref{eq:e8v73-density};
\item uniform \(L^3\) connection and \(L^{3/2}\) quadratic bounds;
\item strong form convergence or a proved weak/Mosco replacement;
\item \(L^3\) extrinsic curvature and \(L^{3/2}\) matter multiplier
control;
\item coercivity of the complete form, not merely its principal part;
\item coordinate-transition and boundary-flux compatibility; and
\item metric-chart probability, Ward, reflection, positivity, and
cutoff-removal contacts.
\end{enumerate}
\end{definition}

\begin{assessment}[Progress at V73]
\label{ass:e8v73-status}
The scalar-curvature potential no longer requires pointwise curvature
control.  In three dimensions its integrated action on scalar tests is
a bounded energy form for uniformly nondegenerate
\(W^{1,3}\cap L^\infty\) metrics, and strong
\(W^{1,3}\) convergence yields form-norm convergence.  Extrinsic
curvature in \(L^3\) and the matter multiplier in \(L^{3/2}\) fit the
same V71 compiler.  Metric-chart tightness, boundary traces, and
coercivity beyond the small-form regime remain open.
\end{assessment}

\begin{programme}[Eighth-series V74: boundary curvature flux and collar
trace]
\label{prog:e8v73-V74}
V74 will retain the boundary term
\(\int_{\partial D}uv\,V_g\cdot n\), prove the trace-space norm needed
to control it, and compare that norm with the V45 owner/collar
compiler.  If \(W^{1,3}\) does not define the required normal trace,
the programme will identify the stronger metric or divergence-measure
hypothesis explicitly.
\end{programme}

\begin{firewall}[Claim boundary after eighth-series V73]
\label{fire:e8v73-boundary}
V73 proves the first-derivative weak scalar-curvature form, its critical
three-dimensional bound and convergence, and compatible extrinsic and
matter form estimates.  It does not prove the metric regularity under
the physical measure, boundary normal-trace control, complete
coercivity outside a small-form regime, positive carrier
reconstruction, cutoff removal, or the full continuum.
\end{firewall}

\begin{theorem}[Eighth-series V73 result]
\label{thm:e8v73-result}
Uniformly nondegenerate metrics converging strongly in
\(L^\infty\cap W^{1,3}\), extrinsic curvatures converging in \(L^3\),
and matter metric multipliers converging in \(L^{3/2}\) produce
convergent complete ADM conformal lower-order forms on every fixed
three-dimensional Dirichlet patch, without any \(L^\infty\) scalar
curvature hypothesis.
\end{theorem}

\begin{proof}
Use \Cref{thm:e8v73-convergence,thm:e8v73-K,
cor:e8v73-matter,thm:e8v73-complete}.
\end{proof}

\paragraph{Parent-version record.}
V73 was formed from
\texttt{Renewal\_SM\_Einstein\_Continuum\_Research\_Notes\_V72.tex}.
The V72 parent had (28\,437) logical source lines, (999\,419) bytes,
and SHA--256
\[
 \texttt{1AED43FFCDA4D5CF70C62AA715BA34B25392C1357D16FBA8C737EA2A20DDE5DC}.
\]
The next compulsory companion audit is V75.

% =============================================================================
% END APPENDED EIGHTH-SERIES V73 MATERIAL
% =============================================================================
% =============================================================================
% BEGIN APPENDED EIGHTH-SERIES V74 MATERIAL
% =============================================================================

\clearpage
\section{Eighth-series V74: boundary curvature flux and the exact trace debit}
\label{sec:v8v74}

V73 rewrote the scalar-curvature density as a first-derivative interior
form.  That result is sufficient for compactly supported tests and for
Dirichlet tests, but it does not license discarding the divergence term at a
physical patch boundary.  This note proves two complementary statements.
First, the V73 class \(W^{1,3}\cap L^\infty\) does not carry a continuous
normal trace of the curvature flux.  Second, the stronger class
\(W^{2,3/2}\cap W^{1,3}\cap L^\infty\) does: the flux trace lies in
\(L^2(\partial D)\), exactly dual to the product of two \(H^1(D)\) traces.
This supplies a checkable boundary row for the V45 owner/resolvent compiler.

\subsection{The smooth Green identity and its critical exponents}
\label{sec:e8v74-green}

Retain the V73 definitions
\begin{align}
 Q_g
 &:=\sqrt{\det g}\,g^{ij}
 \left(\Gamma^k{}_{i\ell}\Gamma^\ell{}_{jk}
       -\Gamma^k{}_{ij}\Gamma^\ell{}_{k\ell}\right),
 \label{eq:e8v74-Q}\\
 V_g^k
 &:=\sqrt{\det g}
 \left(g^{ij}\Gamma^k{}_{ij}-g^{ik}\Gamma^j{}_{ij}\right).
 \label{eq:e8v74-V}
\end{align}
For smooth \(g,u,v\) on a bounded smooth domain \(D\subset\mathbb R^3\),
the density identity of V73 and the divergence theorem give
\begin{equation}
 \int_D\sqrt{\det g}\,R_g,uv
 =\int_DQ_guv-\int_DV_g\cdot\nabla(uv)
  +\int_{\partial D}uv\,V_g\cdot n.
 \label{eq:e8v74-green}
\end{equation}

The interior expression
\begin{equation}
 r_g^{\rm int}(u,v)
 :=\int_DQ_guv-\int_DV_g\cdot\nabla(uv)
 \label{eq:e8v74-interior}
\end{equation}
is bounded on \(H^1(D)\times H^1(D)\) whenever
\(V_g\in L^3(D)\) and \(Q_g\in L^{3/2}(D)\).  Indeed,
\begin{align}
 \left|\int_DQ_guv\right|
 &\leq \|Q_g\|_{3/2}\|u\|_6\|v\|_6,
 \label{eq:e8v74-Q-bound}\\
 \left|\int_DV_g\cdot\nabla(uv)\right|
 &\leq \|V_g\|_3
 \bigl(\|u\|_6\|\nabla v\|_2
       +\|v\|_6\|\nabla u\|_2\bigr).
 \label{eq:e8v74-V-bound}
\end{align}
The inhomogeneous Sobolev embedding on a bounded Lipschitz domain therefore
extends \eqref{eq:e8v74-interior} continuously to all of \(H^1(D)\).

\begin{firewall}[An interior distribution is not a boundary flux]
\label{fire:e8v74-interior-boundary}
The distributional identity
\(\sqrt{\det g}R_g=Q_g+\partial_kV_g^k\) on the open set \(D\)
determines \(r_g^{\rm int}\), but it does not determine a value of
\(V_g\cdot n\) on \(\partial D\).  The latter requires a trace theorem or
separate boundary data.  Equation \eqref{eq:e8v74-green} may not be used at
V73 regularity merely by approximation.
\end{firewall}

\subsection{A conformal obstruction at the V73 threshold}
\label{sec:e8v74-obstruction}

The missing trace is not a technical omission.  It already fails within a
bounded conformal family.

\begin{lemma}[Curvature flux of a three-dimensional conformal metric]
\label{lem:e8v74-conformal-flux}
Let \(g=e^{2\phi}\delta\) on a domain in \(\mathbb R^3\).  Then
\begin{equation}
 V_g=-4e^\phi\nabla\phi=-4\nabla(e^\phi).
 \label{eq:e8v74-conformal-flux}
\end{equation}
\end{lemma}

\begin{proof}
The Christoffel symbols are
\[
 \Gamma^k{}_{ij}=\delta^k_j\phi_i+\delta^k_i\phi_j
                  -\delta_{ij}\phi^k.
\]
Consequently
\[
 \delta^{ij}\Gamma^k{}_{ij}=-\phi^k,
 \qquad
 \delta^{ik}\Gamma^j{}_{ij}=3\phi^k.
\]
Since \(\sqrt{\det g}\,g^{ij}=e^\phi\delta^{ij}\), substitution in
\eqref{eq:e8v74-V} gives \eqref{eq:e8v74-conformal-flux}.
\end{proof}

\begin{theorem}[No continuous curvature-flux trace at \(W^{1,3}\)]
\label{thm:e8v74-no-trace}
Let \(D=(0,1)^3\), choose \(2/3<\alpha<1\), and set
\begin{equation}
 \phi(x)=(x^1)^\alpha,
 \qquad g=e^{2\phi}\delta.
 \label{eq:e8v74-singular-metric}
\end{equation}
Then \(g\in W^{1,3}(D)\cap L^\infty(D)\) and
\(e^{-2}\delta\preceq g\preceq e^2\delta\), while the classical inward
flux along the faces \(x^1=t\) diverges as \(t\downarrow0\):
\begin{equation}
 \left.V_g\cdot(-e_1)\right|_{x^1=t}
 =4\alpha e^{t^\alpha}t^{\alpha-1}\longrightarrow+\infty.
 \label{eq:e8v74-flux-blowup}
\end{equation}
Moreover, there are smooth uniformly elliptic metrics \(g_\varepsilon\)
converging to \(g\) in \(W^{1,3}\cap L^\infty\) for which
\begin{equation}
 \|V_{g_\varepsilon}\cdot n\|_{L^2(\{x^1=0\})}
 \longrightarrow\infty.
 \label{eq:e8v74-trace-blowup}
\end{equation}
Thus no locally bounded map from the V73 metric topology to
\(L^2(\partial D)\) can agree with the smooth curvature-flux trace.
\end{theorem}

\begin{proof}
We have
\(\partial_1\phi=\alpha(x^1)^{\alpha-1}\).  Its cube is integrable at
zero precisely when \(3(\alpha-1)>-1\), which is the assumed
\(\alpha>2/3\).  The metric bounds follow from \(0\leq\phi\leq1\).
Formula \eqref{eq:e8v74-flux-blowup} follows from
\Cref{lem:e8v74-conformal-flux}.

For \(\varepsilon>0\), let
\[
 \phi_\varepsilon(x)
 :=(x^1+\varepsilon)^\alpha-\varepsilon^\alpha,
 \qquad g_\varepsilon=e^{2\phi_\varepsilon}\delta.
\]
These metrics are smooth up to the closed cube and have the same uniform
ellipticity window.  Pointwise convergence, boundedness, and dominated
convergence give
\(\phi_\varepsilon\to\phi\) in \(L^\infty\) and
\(\nabla\phi_\varepsilon\to\nabla\phi\) in \(L^3\).  The smooth chain rule
therefore gives \(g_\varepsilon\to g\) in
\(W^{1,3}\cap L^\infty\).  On \(x^1=0\), however,
\[
 V_{g_\varepsilon}\cdot(-e_1)
 =4\alpha\varepsilon^{\alpha-1},
\]
which is constant on the unit face and tends to infinity.  This proves
\eqref{eq:e8v74-trace-blowup} and the asserted nonexistence.
\end{proof}

\begin{corollary}[Approximation does not select a bounded V73 trace]
\label{cor:e8v74-approximation}
Strong convergence of the V73 interior curvature forms does not imply
boundedness, convergence, or even existence of the associated boundary
curvature forms on \(H^1(D)\times H^1(D)\).
\end{corollary}

\begin{proof}
The metrics in \Cref{thm:e8v74-no-trace} converge in exactly the topology of
\Cref{thm:e8v73-convergence}, so their interior forms converge.  Taking
\(u=v=1\) in the smooth boundary expression on the face \(x^1=0\) yields
the divergent quantity in \eqref{eq:e8v74-trace-blowup}.
\end{proof}

\subsection{A sufficient second-derivative trace class}
\label{sec:e8v74-sufficient}

We now give a concrete sufficient condition.  It is deliberately stated as
a boundary debit rather than folded invisibly into the V73 bulk topology.

\begin{theorem}[Critical \(L^2\) curvature-flux trace]
\label{thm:e8v74-L2-trace}
Let \(D\subset\mathbb R^3\) be a bounded Lipschitz domain and suppose
\begin{equation}
 mI\preceq g\preceq MI,
 \qquad
 g\in W^{2,3/2}(D)\cap W^{1,3}(D)\cap L^\infty(D).
 \label{eq:e8v74-strong-class}
\end{equation}
Then
\begin{equation}
 V_g\in W^{1,3/2}(D;\mathbb R^3),
 \qquad
 \operatorname{Tr}(V_g\cdot n)\in L^2(\partial D),
 \label{eq:e8v74-trace-class}
\end{equation}
and
\begin{align}
 \|V_g\|_{W^{1,3/2}}
 &\leq C_{D,m,M}
 \left(\|\nabla g\|_3+\|\nabla^2g\|_{3/2}
                  +\|\nabla g\|_3^2\right),
 \label{eq:e8v74-V-W1}\\
 \|\operatorname{Tr}(V_g\cdot n)\|_{L^2(\partial D)}
 &\leq C_{D,m,M}
 \left(\|\nabla g\|_3+\|\nabla^2g\|_{3/2}
                  +\|\nabla g\|_3^2\right).
 \label{eq:e8v74-normal-bound}
\end{align}
\end{theorem}

\begin{proof}
On the fixed ellipticity window, the maps
\(g\mapsto g^{-1}\) and \(g\mapsto\sqrt{\det g}\) have bounded first and
second derivatives.  Formula \eqref{eq:e8v74-V} therefore has the schematic
form
\[
 V_g=A(g)\nabla g
\]
with \(A\) smooth and uniformly bounded on that window.  The Sobolev chain
and product rules give
\[
 \nabla V_g=A(g)\nabla^2g+DA(g)[\nabla g]\nabla g.
\]
The first term lies in \(L^{3/2}\), and the second does by
\(L^3\cdot L^3\subset L^{3/2}\).  The \(L^{3/2}\) norm of \(V_g\) is
controlled by its \(L^3\) norm and the finite measure of \(D\), proving
\eqref{eq:e8v74-V-W1}.

The trace theorem gives
\[
 W^{1,3/2}(D)\longrightarrow W^{1/3,3/2}(\partial D).
\]
On the two-dimensional boundary, the fractional Sobolev embedding gives
\(W^{1/3,3/2}(\partial D)\hookrightarrow L^2(\partial D)\).
Multiplication by the essentially bounded unit normal preserves \(L^2\).
These three bounds prove \eqref{eq:e8v74-normal-bound}.
\end{proof}

\begin{theorem}[Bounded full curvature form]
\label{thm:e8v74-full-form}
Under \eqref{eq:e8v74-strong-class}, define on \(H^1(D)\)
\begin{equation}
 r_g^{\rm full}(u,v)
 :=r_g^{\rm int}(u,v)
 +\int_{\partial D}(\operatorname{Tr}u)(\operatorname{Tr}v)
       \operatorname{Tr}(V_g\cdot n).
 \label{eq:e8v74-full-form}
\end{equation}
Then
\begin{align}
 |r_g^{\rm full}(u,v)|
 \leq C_D\bigl(&\|Q_g\|_{3/2}+\|V_g\|_3
 +\|\operatorname{Tr}(V_g\cdot n)\|_{L^2(\partial D)}\bigr)
 \|u\|_{H^1}\|v\|_{H^1}.
 \label{eq:e8v74-full-bound}
\end{align}
For smooth data it equals the left side of \eqref{eq:e8v74-green}.
\end{theorem}

\begin{proof}
The interior terms obey \eqref{eq:e8v74-Q-bound}--
\eqref{eq:e8v74-V-bound}.  The trace theorem and the boundary Sobolev
embedding give
\[
 \|\operatorname{Tr}u\|_{L^4(\partial D)}
 \leq C_D\|u\|_{H^1(D)}.
\]
Hence Hölder on \(\partial D\) yields
\[
 \left|\int_{\partial D}(\operatorname{Tr}u)(\operatorname{Tr}v)
                 \operatorname{Tr}(V_g\cdot n)\right|
 \leq C_D\|V_g\cdot n\|_{L^2(\partial D)}
       \|u\|_{H^1}\|v\|_{H^1}.
\]
The smooth equality is \eqref{eq:e8v74-green}; density then defines its
unique continuous extension.
\end{proof}

\begin{remark}[Why the exponent is matched]
\label{rem:e8v74-critical}
In boundary dimension two, the two \(H^1(D)\) traces lie in \(L^4\), so
their product lies in \(L^2\).  The \(W^{1,3/2}(D)\) trace of \(V_g\) lies
in \(L^2\).  This is the endpoint Hölder triple
\(L^4\cdot L^4\cdot L^2\subset L^1\); no hidden pointwise boundary value is
used.
\end{remark}

\subsection{Convergence of the boundary row}
\label{sec:e8v74-convergence}

\begin{theorem}[Strong boundary-form convergence]
\label{thm:e8v74-boundary-convergence}
Let \(g_j,g\) share the ellipticity window in
\eqref{eq:e8v74-strong-class} and suppose
\begin{equation}
 \|g_j-g\|_\infty+\|\nabla(g_j-g)\|_3
 +\|\nabla^2(g_j-g)\|_{3/2}\longrightarrow0.
 \label{eq:e8v74-strong-convergence}
\end{equation}
Then
\begin{align}
 \|V_{g_j}-V_g\|_{W^{1,3/2}}&\longrightarrow0,
 \label{eq:e8v74-V-convergence}\\
 \|\operatorname{Tr}((V_{g_j}-V_g)\cdot n)\|_{L^2(\partial D)}
 &\longrightarrow0,
 \label{eq:e8v74-trace-convergence}\\
 \|r_{g_j}^{\rm full}-r_g^{\rm full}\|_{
       \mathcal B(H^1\times H^1)}&\longrightarrow0.
 \label{eq:e8v74-full-convergence}
\end{align}
\end{theorem}

\begin{proof}
Write \(V_g=A(g)\nabla g\) as in
\Cref{thm:e8v74-L2-trace}.  Uniform convergence of \(g_j\), strong
\(L^3\) convergence of their first derivatives, and strong
\(L^{3/2}\) convergence of their second derivatives show, term by term,
that
\[
 A(g_j)\nabla g_j\to A(g)\nabla g
 \quad\hbox{in }L^3
\]
and
\[
 A(g_j)\nabla^2g_j
 +DA(g_j)[\nabla g_j]\nabla g_j
 \to
 A(g)\nabla^2g+DA(g)[\nabla g]\nabla g
 \quad\hbox{in }L^{3/2}.
\]
This proves \eqref{eq:e8v74-V-convergence}; the continuous trace and
boundary embedding prove \eqref{eq:e8v74-trace-convergence}.  V73 already
gives convergence of the interior form.  Applying the boundary estimate in
the proof of \Cref{thm:e8v74-full-form} to \(V_{g_j}-V_g\) proves
\eqref{eq:e8v74-full-convergence}.
\end{proof}

\begin{corollary}[Boundary-compatible V71 resolvent input]
\label{cor:e8v74-resolvent}
Suppose the principal elliptic forms and every other lower-order ADM form
converge as in V71--V73, and include \(-2r_g^{\rm full}\) in place of the
interior curvature form.  If their summed relative bound remains strictly
below the principal coercivity margin, the V71 primal and adjoint
resolvents exist and converge in form norm, including the curvature boundary
contribution.
\end{corollary}

\begin{proof}
The new boundary summand is a bounded form by
\Cref{thm:e8v74-full-form} and converges in form norm by
\Cref{thm:e8v74-boundary-convergence}.  These are exactly the hypotheses of
the V71 complete-form stability theorem.
\end{proof}

\subsection{Interface with the V45 owner and influence compiler}
\label{sec:e8v74-owner}

\begin{definition}[Boundary curvature trace row, BCTR]
\label{def:e8v74-BCTR}
For every owned coordinate patch \(D_a\), BCTR records
\begin{enumerate}[label=\textup{(T\arabic*)},leftmargin=4em]
\item an orientation and owner for each boundary face, so a shared face is
counted once or is cancelled by a proved opposite-orientation identity;
\item the trace amplitude
\begin{equation}
 B_a:=\|\operatorname{Tr}(V_{g,a}\cdot n_a)\|_{L^2(\partial D_a)};
 \label{eq:e8v74-Ba}
\end{equation}
\item a uniform \(H^1(D_a)\to L^4(\partial D_a)\) trace constant;
\item form-norm convergence of the owned boundary term under adjacent
refinement; and
\item a finite first-collar incidence count before the V45 propagation
kernel is applied.
\end{enumerate}
\end{definition}

\begin{proposition}[BCTR supplies the direct V45 boundary amplitude]
\label{prop:e8v74-BCTR-amplitude}
If a patch obeys BCTR with trace constant \(C_{{\rm tr},a}\), then its
owned curvature-boundary form satisfies
\begin{equation}
 |\beta_a(u,v)|
 \leq C_{{\rm tr},a}^2B_a
       \|u\|_{H^1(D_a)}\|v\|_{H^1(D_a)}.
 \label{eq:e8v74-direct-amplitude}
\end{equation}
If each source meets at most \(N_\partial\) owned faces, its unpropagated
boundary debit is at most the sum of at most \(N_\partial\) such amplitudes.
This debit may be inserted as the direct source vector in the V45 boundary
influence equation.
\end{proposition}

\begin{proof}
Equation \eqref{eq:e8v74-direct-amplitude} is the boundary Hölder estimate
from \Cref{thm:e8v74-full-form}.  The finite-incidence assertion is the
triangle inequality after the owner map removes double counting.  V45 then
applies to this direct source exactly when its independently required Schur
strictness holds.
\end{proof}

\begin{firewall}[A trace bound is not an influence contraction]
\label{fire:e8v74-not-contraction}
BCTR bounds the direct boundary form and its adjacent change.  It does not
make the V45 influence kernel strict, prove decay through successive
collars, or identify an asymptotic ADM charge.  Those are separate rows.
Likewise, imposing Dirichlet data deletes the displayed trace product but
does not prove that a physical Einstein boundary condition is Dirichlet.
\end{firewall}

\begin{assessment}[Progress at V74]
\label{ass:e8v74-status}
The boundary gap left by V73 is now exact.  The V73
\(W^{1,3}\cap L^\infty\) topology cannot control curvature flux, as shown by
an explicit conformal sequence.  Strong
\(W^{2,3/2}\cap W^{1,3}\cap L^\infty\) convergence does control it in
\(L^2(\partial D)\), produces a convergent full curvature form on \(H^1\),
and supplies the direct boundary amplitude required by the V45 owner
compiler.  Whether the companion YM and NS programmes supply a weaker
collar or divergence-structure route must now be audited.
\end{assessment}

\begin{programme}[Eighth-series V75: compulsory companion audit]
\label{prog:e8v74-V75}
V75 will inspect the newest Yang--Mills and Navier--Stokes active notes.  It
will test their newest shell, flux, or contraction results against BCTR and
the V71--V74 weak ADM resolvent.  Any imported result will retain its exact
hypotheses; a vector or curvature estimate will not be transferred to the
Einstein scalar channel without an explicit operator identity.
\end{programme}

\begin{firewall}[Claim boundary after eighth-series V74]
\label{fire:e8v74-boundary}
V74 proves the conformal no-trace example, the sufficient
\(W^{2,3/2}\) curvature-flux theorem, convergence of the full curvature
form, and the BCTR-to-V45 direct-amplitude interface.  It does not prove
metric tightness in the stronger topology, the V45 contraction row,
physical boundary matching, asymptotic flatness, nonlinear constraint
solvability, the continuum measure, or the full SM--Einstein continuum.
\end{firewall}

\begin{theorem}[Eighth-series V74 result]
\label{thm:e8v74-result}
For three-dimensional uniformly elliptic metrics, the V73 interior
curvature form is stable at \(W^{1,3}\), but its boundary flux is not.
The stronger topology
\(W^{2,3/2}\cap W^{1,3}\cap L^\infty\) yields an \(L^2\) normal curvature
flux, a bounded and convergent full \(H^1\) curvature form, and an explicit
owned boundary amplitude for the V45 compiler.
\end{theorem}

\begin{proof}
Use \Cref{thm:e8v74-no-trace,thm:e8v74-L2-trace,
thm:e8v74-full-form,thm:e8v74-boundary-convergence,
prop:e8v74-BCTR-amplitude}.
\end{proof}

\paragraph{Parent-version record.}
V74 was formed from
\texttt{Renewal\_SM\_Einstein\_Continuum\_Research\_Notes\_V73.tex}.
The V73 parent had (28\,891) logical source lines, (1\,012\,871) bytes,
and SHA--256
\[
 \texttt{5B8A97842039485CE8E60ACA93CE481A09512C99E128F90DA3687A2C5BAAF704}.
\]
The next version is the compulsory V75 companion audit.

% =============================================================================
% END APPENDED EIGHTH-SERIES V74 MATERIAL
% =============================================================================
% =============================================================================
% BEGIN APPENDED EIGHTH-SERIES V75 MATERIAL
% =============================================================================

\clearpage
\section{Eighth-series V75: companion audit and the global weak-curvature pivot}
\label{sec:v8v75}

This is the compulsory audit following V70.  During the audit the active
Yang--Mills note advanced from V26 to V27, so the final V27 file was inspected.
The Navier--Stokes active note remains V26.  The audit changes the immediate
SM direction in one precise way: coordinate-patch boundaries will no longer be
treated as physical boundary sources.  V76 will construct one global,
background-covariant weak scalar-curvature density, for which a partition of
unity creates no artificial interface debit.  The V74 trace theorem remains
necessary only at an actual boundary or asymptotic end.

\subsection{Files inspected}
\label{sec:e8v75-files}

The final companion checkpoints used here are
\begin{align}
 &\texttt{Renewal\_Yang\_Mills\_Mass\_Gap\_Research\_Notes\_V27.tex},
 \notag\\
 &\qquad 11\,068\ \text{logical lines},\quad
 404\,167\ \text{bytes},\quad
 \operatorname{SHA256}
 =\texttt{42EB3B0FAFFEE7F8437C3854B85F0C7B7233169656BD6BF15A232D568D92A29F},
 \label{eq:e8v75-YM-file}\\
 &\texttt{Renewal\_NS\_Stability\_Research\_Notes\_V26.tex},
 \notag\\
 &\qquad 5\,319\ \text{logical lines},\quad
 278\,283\ \text{bytes},\quad
 \operatorname{SHA256}
 =\texttt{82C887F8C2C69E4F28FAD7C3DC8DE5B9099CB5A4BF431EBA1FFF3E91935978AD}.
 \label{eq:e8v75-NS-file}
\end{align}
The NS file is byte-for-byte unchanged from the V70 audit.  The YM file has
advanced from V22 to V27.

\subsection{Yang--Mills V23--V27}
\label{sec:e8v75-YM}

\begin{audit}[YM V23: nonlinear boundary localization and its clock]
\label{audit:e8v75-YM23}
The naive nonabelian parallel-link block has a volume-order Wilson score even
though its flat linearization is boundary supported.  A compact cut-gauge
block repairs the nonlinear support exactly: only mixed plaquettes at the box
boundary vary, and its flat tangent contains the desired dyadic curvature
direction.  The repair does not make the physical clock uniform.  With the
frame rates used there, the exact per-link participation grows linearly in the
lattice diameter.
\end{audit}

\begin{audit}[YM V24: boundary support cannot carry a bounded-clock frame]
\label{audit:e8v75-YM24}
Every flat cut-gauge curvature tangent is boundary supported and annihilates
constant curvature.  Testing its projection on a lowest Fourier mode proves
\[
 \hbox{frame floor}\leq C_d
 {\hbox{maximum local participation}\over L}.
\]
This holds for arbitrary selected boxes and nonnegative weights.  Therefore a
uniform frame inside that boundary-block class forces at least linear local
participation.  Sparse ownership does not remove the obstruction.
\end{audit}

\begin{audit}[YM V25--V26: direct physical dense-core criterion]
\label{audit:e8v75-YM25-26}
After closing the bounded-clock branch, the YM programme proves that a common
exponential Euclidean-time decay rate on a dense gauge-invariant
Osterwalder--Schrader core implies a Hamiltonian gap.  At fixed cutoff it then
constructs the positive gauge-projected transfer kernel, proves simplicity of
its positive vacuum, proves density of finite spin networks in the physical
space, and identifies the exact connected two-slab correlator.  A global
operator-norm perturbation about the electric point is volume extensive.
\end{audit}

\begin{audit}[YM V27: connected expansion removes the volume debit only at
strong coupling]
\label{audit:e8v75-YM27}
For the Wilson regulator, a tilted plaquette-polymer expansion converges under
the explicit small-activity condition
\[
 e^{2\beta}-1<(2e\Delta_4^2)^{-1}.
\]
It yields volume-uniform exponential clustering of every fixed local
spin-network observable and hence an infinite-volume lattice Hamiltonian gap
in that strong-coupling interval.  The note explicitly leaves open the bridge
to the asymptotically free weak-bare-coupling continuum trajectory.
\end{audit}

\subsection{Navier--Stokes V26}
\label{sec:e8v75-NS}

\begin{audit}[NS V26: rank-one Oseen-kernel norms]
\label{audit:e8v75-NS26}
The active NS note computes exact pointwise rank-one norms of the first two
Oseen-kernel derivatives.  The first-derivative rank-one norm equals the full
symmetric-traceless norm.  The second-derivative norm is smaller on an
intermediate radial interval, by at most about eleven percent pointwise; the
expected integrated gain is only a few percent.  These are sharp constants
for the NS ancient-form input, not curvature-flux or metric compactness
estimates.
\end{audit}

\subsection{Type audit for the SM--Einstein programme}
\label{sec:e8v75-types}

\begin{theorem}[Companion non-transfer at V75]
\label{thm:e8v75-no-transfer}
None of the new companion results proves a missing SM--Einstein row.  In
particular, the following inferences are invalid:
\begin{enumerate}[label=\textup{(\roman*)},leftmargin=3.8em]
\item treating an owned codimension-one curvature boundary form as a
volume-uniform bulk scalar-constraint frame;
\item importing the YM compact-group cut-gauge action into the noncompact
metric conformal sector;
\item importing the positive small-activity Wilson polymer expansion into an
indefinite Einstein action without a positive normalized activity majorant;
\item reading a strong-coupling lattice mass as a cutoff-uniform physical mass
on the weak-coupling continuum trajectory; or
\item using the NS rank-one Oseen constants as a trace, coercivity, or
tightness estimate for the ADM form.
\end{enumerate}
\end{theorem}

\begin{proof}
For (i), the YM lowest-mode theorem shows the opposite mechanism within its
boundary-supported class; the SM scalar symbol and state space differ, so even
that no-go is guidance rather than a numerical transfer.  Items (ii) and (iii)
fail because the field spaces, reference measures, signs, and conditional
activities differ.  Item (iv) is the explicit regime firewall of the audited
YM V27 theorem.  Item (v) fails by operator type: an Oseen convolution tensor
norm is neither a metric Sobolev trace nor a lower-form bound.
\end{proof}

\begin{corollary}[What remains valid from V74]
\label{cor:e8v75-V74}
The V74 BCTR is a correct direct-amplitude row for an actual boundary.  It is
not, by itself, a bulk coercivity or continuum-tightness argument, and summing
it over an artificial coordinate tessellation is not licensed by the audit.
\end{corollary}

\begin{proof}
V74 proves only the local trace and owner estimate.  The first non-transfer in
\Cref{thm:e8v75-no-transfer} excludes the stronger reading.
\end{proof}

\subsection{The reusable structural lessons}
\label{sec:e8v75-lessons}

Two proof patterns are nevertheless type-correct after their inputs are
reconstructed in the Einstein setting.

\begin{proposition}[Localization identity for a single global weak density]
\label{prop:e8v75-localization}
Let \(M\) be a compact three-manifold without boundary with a fixed smooth
measure \(d\mu_0\).  Suppose
\[
 Q\in L^{3/2}(M),\qquad V\in L^3(TM),
\]
and define
\begin{equation}
 \mathfrak r(f):=int_MQf\,d\mu_0-int_MV\cdot\nabla f\,d\mu_0
 \label{eq:e8v75-global-density}
\end{equation}
for \(f\in W^{1,3/2}(M)\).  If
\(\{\chi_a\}_{a=1}^N\subset C^\infty(M)\) is a partition of unity, then
\begin{equation}
 \mathfrak r(f)
 =\sum_{a=1}^N
 \left[
  \int_MQ\chi_af\,d\mu_0
  -\int_MV\cdot\nabla(\chi_af)\,d\mu_0
 \right].
 \label{eq:e8v75-partition}
\end{equation}
There is no internal-face trace term.
\end{proposition}

\begin{proof}
All integrals are finite by Hölder and the Sobolev embedding
\(W^{1,3/2}\hookrightarrow L^3\) in dimension three.  Sum the right side.
Since \(\sum_a\chi_a=1\) and \(\sum_a\nabla\chi_a=0\), it is exactly
\eqref{eq:e8v75-global-density}.
\end{proof}

\begin{firewall}[The V73 coordinate flux is not yet the global vector]
\label{fire:e8v75-not-global}
The Christoffel symbols in the V73 coordinate formula are not tensors, and
the corresponding \(V_g\) is not a globally defined vector under arbitrary
coordinate changes.  Consequently
\Cref{prop:e8v75-localization} cannot yet be applied by declaring the local
V73 expressions equal on overlaps.  V76 must replace Christoffel symbols by
the tensorial difference between the metric connection and one fixed smooth
reference connection.
\end{firewall}

\begin{proposition}[Connected-local rather than extensive bookkeeping]
\label{prop:e8v75-connected-principle}
Let a normalized expectation have an absolutely convergent connected-cluster
expansion.  The mixed derivative with respect to two local source parameters
contains only connected clusters meeting both marked supports; disconnected
vacuum components cancel against normalization.  Therefore a bound on the
sum of clusters joining the marked supports can be volume uniform even when a
global norm of the complete interaction is extensive.
\end{proposition}

\begin{proof}
Write the sourced logarithm of the partition function as its absolutely
convergent sum over connected clusters.  A mixed derivative in the two
sources kills every cluster meeting neither or only one source.  Absolute
convergence permits termwise differentiation, and the remaining rooted sum
depends on the two fixed supports and their connecting clusters, not on the
number of disconnected vacuum clusters.
\end{proof}

\begin{firewall}[The structural lemma is not an Einstein cluster expansion]
\label{fire:e8v75-no-Einstein-polymer}
\Cref{prop:e8v75-connected-principle} is an implication once an absolutely
convergent normalized cluster expansion exists.  No such expansion has been
proved for the coupled SM--Einstein measure, and the conformal-sign problem
prevents importing the positive YM activity estimate.  The lemma only fixes
the correct later target: local normalized observables, rather than the global
interaction norm.
\end{firewall}

\subsection{Audit-guided direction}
\label{sec:e8v75-direction}

\begin{decision}[Global weak curvature before boundary propagation]
\label{dec:e8v75-direction}
The next step goes upstream of BCTR.  On a smooth reference manifold
\((M,\bar g)\), V76 will form the tensor
\(C=\Gamma(g)-\Gamma(\bar g)\), derive a global identity
\[
 {d\mu_g\over d\mu_{\bar g}}R_g
 =Q_{g,\bar g}+\operatorname{div}_{\bar g}V_{g,\bar g},
\]
and prove that \(g\in W^{1,3}\cap L^\infty\) gives
\(V_{g,\bar g}\in L^3\) and \(Q_{g,\bar g}\in L^{3/2}\).  The partition
identity \eqref{eq:e8v75-partition} will then cancel every artificial chart
interface.  V74 will be invoked only for a genuine boundary.
\end{decision}

\begin{assessment}[Progress at V75]
\label{ass:e8v75-status}
The companion audit blocks a tempting misuse of V74: boundary ownership is
not a substitute for bulk low-frequency control.  It also selects a sharper
geometric route.  A single global reference-connection weak curvature density
will eliminate coordinate-interface traces before any V45 propagation
estimate is attempted.  The connected-cluster lesson is retained as a later
measure-level target, with its Einstein hypotheses explicitly unpaid.
\end{assessment}

\begin{programme}[Eighth-series V76: reference-connection weak scalar
curvature]
\label{prog:e8v75-V76}
V76 will execute \Cref{dec:e8v75-direction}: derive the exact relative
connection identity, prove its critical Sobolev bounds and convergence, prove
partition-of-unity invariance, and separate true geometric boundary flux from
coordinate artifacts.
\end{programme}

\begin{firewall}[Claim boundary after eighth-series V75]
\label{fire:e8v75-boundary}
V75 records the current YM V27 and NS V26 audit, proves the type non-transfer,
the abstract global localization identity, and the connected-local
bookkeeping lemma.  It does not construct the global weak curvature density,
an Einstein cluster expansion, metric tightness, constraint solvability,
continuum reconstruction, or the full SM--Einstein continuum.
\end{firewall}

\begin{theorem}[Eighth-series V75 result]
\label{thm:e8v75-result}
The new YM boundary-block analysis confirms that codimension-one ownership
cannot be promoted to a uniform bulk frame, while its connected strong-coupling
expansion shows how normalized local observables can avoid an extensive global
norm.  For the SM--Einstein programme the immediate rigorous consequence is
to replace coordinatewise curvature fluxes by one global
reference-connection weak density before gluing patches.
\end{theorem}

\begin{proof}
Use \Cref{audit:e8v75-YM23,audit:e8v75-YM24,
audit:e8v75-YM25-26,audit:e8v75-YM27,thm:e8v75-no-transfer,
prop:e8v75-localization,prop:e8v75-connected-principle,
dec:e8v75-direction}.
\end{proof}

\paragraph{Parent-version record.}
V75 was formed from
\texttt{Renewal\_SM\_Einstein\_Continuum\_Research\_Notes\_V74.tex}.
The V74 parent had (29\,361) logical source lines, (1\,030\,075) bytes,
and SHA--256
\[
 \texttt{E15EFC53DDB9BEEC0A1E8AD39641ADDEFD3E7AFBD735C07259A3914EDCACAF04}.
\]
The next compulsory companion audit is V80.

% =============================================================================
% END APPENDED EIGHTH-SERIES V75 MATERIAL
% =============================================================================
% =============================================================================
% BEGIN APPENDED EIGHTH-SERIES V76 MATERIAL
% =============================================================================

\clearpage
\section{Eighth-series V76: global reference-connection weak scalar
curvature}
\label{sec:v8v76}

V75 selected an upstream repair of the V74 boundary ledger.  The local
coordinate vector \(V_g\) of V73 is not tensorial, so separate chart faces
cannot be assigned physical fluxes and then summed.  This note replaces it by
the difference of two Levi--Civita connections, which is a tensor.  The result
is one global weak scalar-curvature density at the critical regularity
\(g\in W^{1,3}\cap L^\infty\).  A partition of unity then creates no internal
boundary term.  The construction is independent of the chosen smooth
reference metric along every strongly convergent smooth cofinal sequence.

\subsection{Relative connection tensors}
\label{sec:e8v76-relative}

Let \(M\) be a compact oriented smooth three-manifold, possibly with
boundary, and fix a smooth reference metric \(\bar g\), its connection
\(\bar\nabla\), curvature \(\bar R_{ij}\), and measure \(d\mu_{\bar g}\).
Let \(g\) be a measurable metric satisfying, almost everywhere,
\begin{equation}
 m\bar g\preceq g\preceq M\bar g,
 \qquad
 g\in W^{1,3}(M)\cap L^\infty(M).
 \label{eq:e8v76-class}
\end{equation}
All Sobolev spaces and norms in this note may be defined with \(\bar g\);
compactness makes the choice of smooth background equivalent.

Define
\begin{align}
 J_g&:={d\mu_g\over d\mu_{\bar g}},
 &A_g^{ij}&:=J_gg^{ij},                                  \label{eq:e8v76-JA}\\
 C^k{}_{ij}
 &:=\frac12g^{k\ell}
 \left(\bar\nabla_i g_{j\ell}+\bar\nabla_j g_{i\ell}
       -\bar\nabla_\ell g_{ij}\right).                 \label{eq:e8v76-C}
\end{align}

\begin{lemma}[Tensor type and critical regularity]
\label{lem:e8v76-tensor}
The field \(C\) is the tensorial connection difference
\begin{equation}
 C(X,Y)=\nabla^g_XY-\bar\nabla_XY                         \label{eq:e8v76-difference}
\end{equation}
whenever \(g\) is smooth, and formula \eqref{eq:e8v76-C} defines an
\(L^3\) \((1,2)\)-tensor for every \(g\) in
\eqref{eq:e8v76-class}.  Moreover
\begin{equation}
 J_g,A_g\in W^{1,3}\cap L^\infty,
 \qquad
 \|C\|_3+\|\bar\nabla A_g\|_3
 \leq C_{\bar g,m,M}\|\bar\nabla g\|_3.
 \label{eq:e8v76-basic-bound}
\end{equation}
\end{lemma}

\begin{proof}
The difference of two connections is tensorial, and the Koszul formula gives
\eqref{eq:e8v76-C}.  On the compact set of positive matrices with spectrum in
\([m,M]\), inversion, determinant, and square root are smooth with bounded
derivatives.  The Sobolev chain rule therefore gives
\(g^{-1},J_g,A_g\in W^{1,3}\cap L^\infty\), with their first derivatives
bounded pointwise by a constant times \(|\bar\nabla g|\).  Formula
\eqref{eq:e8v76-C} then proves the \(L^3\) assertion and
\eqref{eq:e8v76-basic-bound}.
\end{proof}

\subsection{The exact relative curvature identity}
\label{sec:e8v76-identity}

Set
\begin{equation}
 V_{g,\bar g}^{,k}
 :=A_g^{ij}C^k{}_{ij}-A_g^{ik}C^j{}_{ij}                 \label{eq:e8v76-V}
\end{equation}
and
\begin{align}
 Q_{g,\bar g}
 :=\;&A_g^{ij}\bar R_{ij}
 +A_g^{ij}
 \left(C^k{}_{ij}C^\ell{}_{k\ell}
       -C^k{}_{i\ell}C^\ell{}_{jk}\right)              \notag\\
 &-(\bar\nabla_kA_g^{ij})C^k{}_{ij}
 +(\bar\nabla_jA_g^{ij})C^k{}_{ik}.                    \label{eq:e8v76-Q}
\end{align}

\begin{theorem}[Global relative-connection curvature identity]
\label{thm:e8v76-identity}
If \(g\) is smooth, then
\begin{equation}
 J_gR_g
 =Q_{g,\bar g}+\bar\nabla_kV_{g,\bar g}^{,k}.
 \label{eq:e8v76-density}
\end{equation}
Both sides are scalar functions relative to the fixed measure
\(d\mu_{\bar g}\); in particular, \(V_{g,\bar g}\) is a genuine vector field,
not a coordinate vector density.
\end{theorem}

\begin{proof}
The curvature-difference formula in the convention of V73 is
\begin{align}
 R_{ij}(g)
 ={}&\bar R_{ij}
 +\bar\nabla_kC^k{}_{ij}-\bar\nabla_jC^k{}_{ik}           \notag\\
 &+C^k{}_{ij}C^\ell{}_{k\ell}
  -C^k{}_{i\ell}C^\ell{}_{jk}.                         \label{eq:e8v76-Ricci}
\end{align}
Contract with \(A_g^{ij}=J_gg^{ij}\).  The two derivative terms obey
\begin{align*}
 A_g^{ij}\bar\nabla_kC^k{}_{ij}
 &=\bar\nabla_k(A_g^{ij}C^k{}_{ij})
   -(\bar\nabla_kA_g^{ij})C^k{}_{ij},\\
 -A_g^{ij}\bar\nabla_jC^k{}_{ik}
 &=-\bar\nabla_j(A_g^{ij}C^k{}_{ik})
   +(\bar\nabla_jA_g^{ij})C^k{}_{ik}.
\end{align*}
After renaming dummy indices, the two divergences combine into
\(\bar\nabla_kV_{g,\bar g}^{,k}\); all remaining terms are exactly
\eqref{eq:e8v76-Q}.  The left side is \(J_gR_g\), proving
\eqref{eq:e8v76-density}.  Tensoriality follows from
\Cref{lem:e8v76-tensor} and from contraction of tensors.
\end{proof}

\begin{corollary}[Recovery of the V73 coordinate formula]
\label{cor:e8v76-V73}
In a Euclidean coordinate patch with \(\bar g=\delta\), one has
\(C^k{}_{ij}=\Gamma^k{}_{ij}(g)\),
\(d\mu_{\bar g}=dx\), and \eqref{eq:e8v76-V} is exactly the V73 vector
density.  Formula \eqref{eq:e8v76-density} therefore reduces to the V73
\(\Gamma\Gamma+\partial V\) identity.
\end{corollary}

\begin{proof}
The Euclidean reference connection and curvature vanish, while
\(J_g=\sqrt{\det g}\).  Substitute these identities into
\eqref{eq:e8v76-C}--\eqref{eq:e8v76-Q}.
\end{proof}

\subsection{Weak curvature at \(W^{1,3}\)}
\label{sec:e8v76-weak}

\begin{theorem}[Critical coefficient bounds]
\label{thm:e8v76-coefficients}
Under \eqref{eq:e8v76-class},
\begin{equation}
 V_{g,\bar g}\in L^3(TM),
 \qquad
 Q_{g,\bar g}\in L^{3/2}(M),                            \label{eq:e8v76-spaces}
\end{equation}
and
\begin{align}
 \|V_{g,\bar g}\|_3
 &\leq C_{\bar g,m,M}\|\bar\nabla g\|_3,              \label{eq:e8v76-V-bound}\\
 \|Q_{g,\bar g}\|_{3/2}
 &\leq C_{\bar g,m,M}
 \left(1+\|\bar\nabla g\|_3^2\right).                 \label{eq:e8v76-Q-bound}
\end{align}
\end{theorem}

\begin{proof}
The coefficients \(A_g\) are bounded, and \(C\in L^3\), so
\eqref{eq:e8v76-V-bound} follows from \eqref{eq:e8v76-V}.  In
\eqref{eq:e8v76-Q}, the reference-curvature term is bounded, both quadratic
connection terms lie in \(L^{3/2}\), and the last two terms are products of
two \(L^3\) fields.  Hölder and \Cref{lem:e8v76-tensor} prove
\eqref{eq:e8v76-Q-bound}.
\end{proof}

\begin{definition}[Global weak scalar-curvature density]
\label{def:e8v76-weak-curvature}
For \(f\in W^{1,3/2}(M)\), define
\begin{equation}
 \mathscr R_{g,\bar g}(f)
 :=\int_MQ_{g,\bar g}f\,d\mu_{\bar g}
   -\int_MV_{g,\bar g}^{,k}\bar\nabla_kf\,d\mu_{\bar g}.
 \label{eq:e8v76-weak-curvature}
\end{equation}
This is the interior distribution associated with
\eqref{eq:e8v76-density}.  On a manifold with boundary it does not include a
normal trace.
\end{definition}

\begin{theorem}[Weak curvature is a bounded critical functional]
\label{thm:e8v76-functional}
There is \(C=C(M,\bar g,m,M)\) such that
\begin{equation}
 |\mathscr R_{g,\bar g}(f)|
 \leq C
 \left(1+\|\bar\nabla g\|_3+
             \|\bar\nabla g\|_3^2\right)
 \|f\|_{W^{1,3/2}}.                                    \label{eq:e8v76-functional-bound}
\end{equation}
If \(M\) has no boundary and \(g\) is smooth, then
\begin{equation}
 \mathscr R_{g,\bar g}(f)=\int_MR_gf\,d\mu_g.           \label{eq:e8v76-smooth-closed}
\end{equation}
\end{theorem}

\begin{proof}
The critical Sobolev embedding in dimension three gives
\(W^{1,3/2}(M)\hookrightarrow L^3(M)\).  Hence
\begin{align*}
 |\mathscr R_{g,\bar g}(f)|
 &\leq \|Q_{g,\bar g}\|_{3/2}\|f\|_3
 +\|V_{g,\bar g}\|_3\|\bar\nabla f\|_{3/2}.
\end{align*}
Apply \Cref{thm:e8v76-coefficients}.  In the smooth closed case, integrate
the divergence in \eqref{eq:e8v76-density}; there is no boundary term.
\end{proof}

\begin{corollary}[Global curvature form on \(H^1\)]
\label{cor:e8v76-H1-form}
For \(u,v\in H^1(M)\), the product \(uv\) belongs to
\(W^{1,3/2}(M)\), and
\begin{equation}
 \|uv\|_{W^{1,3/2}}
 \leq C_{M,\bar g}\|u\|_{H^1}\|v\|_{H^1}.             \label{eq:e8v76-product}
\end{equation}
Consequently
\begin{equation}
 r_{g,\bar g}^{\rm glob}(u,v):=\mathscr R_{g,\bar g}(uv)
 \label{eq:e8v76-global-form}
\end{equation}
is a bounded bilinear form on \(H^1(M)\times H^1(M)\).
\end{corollary}

\begin{proof}
The embedding \(H^1\hookrightarrow L^6\) gives
\(uv\in L^3\).  The weak product rule gives
\(\bar\nabla(uv)=u\bar\nabla v+v\bar\nabla u\), and
\(L^6\cdot L^2\subset L^{3/2}\).  This proves
\eqref{eq:e8v76-product}; apply
\Cref{thm:e8v76-functional}.
\end{proof}

\subsection{Strong cofinal stability}
\label{sec:e8v76-stability}

\begin{theorem}[Strong convergence of the global weak density]
\label{thm:e8v76-convergence}
Let \(g_n,g\) share the ellipticity window in
\eqref{eq:e8v76-class} and suppose
\begin{equation}
 \|g_n-g\|_\infty+\|\bar\nabla(g_n-g)\|_3\longrightarrow0.
 \label{eq:e8v76-metric-convergence}
\end{equation}
Then
\begin{align}
 \|C_n-C\|_3+
 \|V_{g_n,\bar g}-V_{g,\bar g}\|_3
 &\longrightarrow0,                                    \label{eq:e8v76-VC-convergence}\\
 \|Q_{g_n,\bar g}-Q_{g,\bar g}\|_{3/2}
 &\longrightarrow0,                                    \label{eq:e8v76-Q-convergence}\\
 \|r_{g_n,\bar g}^{\rm glob}-r_{g,\bar g}^{\rm glob}\|_{
     \mathcal B(H^1\times H^1)}
 &\longrightarrow0.                                    \label{eq:e8v76-form-convergence}
\end{align}
\end{theorem}

\begin{proof}
Smoothness of the matrix maps on the common ellipticity window gives
\[
 \|A_{g_n}-A_g\|_\infty\to0,
 \qquad
 \|\bar\nabla A_{g_n}-\bar\nabla A_g\|_3\to0.
\]
Using \eqref{eq:e8v76-C},
\[
 \|C_n-C\|_3
 \leq C\bigl(
 \|\bar\nabla(g_n-g)\|_3
 +\|g_n-g\|_\infty\|\bar\nabla g\|_3
 \bigr).
\]
This proves the first connection convergence.  Formula
\eqref{eq:e8v76-V} and Hölder then give the \(V\) convergence.  Every term in
\eqref{eq:e8v76-Q} is either a bounded coefficient, an \(L^3\cdot L^3\)
product, or a smooth reference-curvature term; expanding differences proves
\eqref{eq:e8v76-Q-convergence}.  Finally apply the estimate in
\Cref{thm:e8v76-functional} to the coefficient differences and use
\eqref{eq:e8v76-product}.
\end{proof}

\subsection{Reference independence along the continuum sequence}
\label{sec:e8v76-reference}

\begin{theorem}[Reference independence for strongly approximable metrics]
\label{thm:e8v76-reference}
Let \(\bar g\) and \(\widehat g\) be two smooth reference metrics.  Suppose
there are smooth metrics \(g_n\), sharing one uniform ellipticity window
relative to both references, such that
\begin{equation}
 g_n\to g\quad\hbox{in }L^\infty\cap W^{1,3}.            \label{eq:e8v76-approximable}
\end{equation}
If \(M\) is closed, then for every \(f\in W^{1,3/2}(M)\),
\begin{equation}
 \mathscr R_{g,\bar g}(f)=\mathscr R_{g,\widehat g}(f).
 \label{eq:e8v76-reference-independence}
\end{equation}
Thus the common distribution may be denoted simply \(\mathscr R_g\) on this
cofinal closure class.
\end{theorem}

\begin{proof}
For every smooth \(g_n\), \Cref{thm:e8v76-functional} gives
\[
 \mathscr R_{g_n,\bar g}(f)
 =\int_MR_{g_n}f\,d\mu_{g_n}
 =\mathscr R_{g_n,\widehat g}(f).
\]
Apply \Cref{thm:e8v76-convergence} once with background \(\bar g\) and once
with background \(\widehat g\), then pass to the limit.
\end{proof}

\begin{firewall}[Scope of reference independence]
\label{fire:e8v76-reference-scope}
V76 proves reference independence on the strong
\(L^\infty\cap W^{1,3}\) closure of smooth uniformly elliptic metrics, which
is the topology used in V73 and the present cofinal form theorem.  It does not
assert density of smooth metrics in that strong topology for every abstract
bounded \(W^{1,3}\) metric.  A weaker approximation topology would require a
separate compensated-product argument.
\end{firewall}

\subsection{Exact cancellation of artificial interfaces}
\label{sec:e8v76-gluing}

\begin{theorem}[Partition-of-unity gluing without chart flux]
\label{thm:e8v76-gluing}
Let \(M\) be closed, let \(g\) obey \eqref{eq:e8v76-class}, and let
\(\{\chi_a\}_{a=1}^N\) be a smooth partition of unity subordinate to any
finite atlas.  For \(f\in W^{1,3/2}(M)\),
\begin{align}
 \mathscr R_{g,\bar g}(f)
 =\sum_{a=1}^N\bigg\{
 &\int_MQ_{g,\bar g}\chi_af\,d\mu_{\bar g}              \notag\\
 &-\int_MV_{g,\bar g}\cdot
       \bar\nabla(\chi_af)\,d\mu_{\bar g}\bigg\}.      \label{eq:e8v76-gluing}
\end{align}
No normal trace on the artificial boundary of a coordinate patch occurs.
The identity remains true for \(f=uv\), \(u,v\in H^1(M)\).
\end{theorem}

\begin{proof}
All terms are integrable by \Cref{thm:e8v76-functional}.  Sum the right side
of \eqref{eq:e8v76-gluing}.  The terms in which the derivative hits the
partition cancel because \(\sum_a\bar\nabla\chi_a=0\); the remaining terms
use \(\sum_a\chi_a=1\) and give
\eqref{eq:e8v76-weak-curvature}.  The final assertion follows from
\Cref{cor:e8v76-H1-form}.
\end{proof}

\begin{corollary}[Artificial-face BCTR debits are deleted]
\label{cor:e8v76-delete-artificial}
For a closed manifold, an atlas localization of the weak scalar-curvature
form requires no V74 BCTR entry on chart faces.  Only cutoff-gradient terms in
\eqref{eq:e8v76-gluing} appear, and their sum cancels exactly.  Consequently
these chart interfaces must not be inserted into the V45 boundary-influence
kernel.
\end{corollary}

\begin{proof}
This is the last assertion of \Cref{thm:e8v76-gluing}.  BCTR concerns an
actual normal trace; no such trace occurs in \eqref{eq:e8v76-gluing}.
\end{proof}

\subsection{True boundaries and the V74 row}
\label{sec:e8v76-boundary}

\begin{theorem}[Relative Green identity at a true boundary]
\label{thm:e8v76-boundary}
Assume \(M\) has Lipschitz boundary and strengthen the metric regularity to
\begin{equation}
 g\in W^{2,3/2}(M)\cap W^{1,3}(M)\cap L^\infty(M).       \label{eq:e8v76-boundary-class}
\end{equation}
Then
\begin{equation}
 V_{g,\bar g}\in W^{1,3/2}(TM),
 \qquad
 \operatorname{Tr}\langle V_{g,\bar g},\bar n\rangle_{\bar g}
 \in L^2(\partial M).                                  \label{eq:e8v76-relative-trace}
\end{equation}
For \(u,v\in H^1(M)\), the full smooth curvature form extends as
\begin{align}
 \int_MR_guv\,d\mu_g
 ={}&\mathscr R_{g,\bar g}(uv)                           \notag\\
 &+\int_{\partial M}
   (\operatorname{Tr}u)(\operatorname{Tr}v)
   \operatorname{Tr}\langle V_{g,\bar g},\bar n\rangle_{\bar g}
   \,d\sigma_{\bar g}.                                \label{eq:e8v76-relative-Green}
\end{align}
Here equality means equality with the unique continuous extension of the
smooth expression.
\end{theorem}

\begin{proof}
The formula \(V_{g,\bar g}=A(g)\bar\nabla g\) has smooth bounded
coefficients on the ellipticity window.  Its derivative is a sum of
\(A(g)\bar\nabla^2g\), quadratic first-derivative terms, and smooth
background-connection terms.  Hence it lies in \(L^{3/2}\).  The trace chain
\[
 W^{1,3/2}(M)\longrightarrow W^{1/3,3/2}(\partial M)
 \longrightarrow L^2(\partial M)
\]
proves \eqref{eq:e8v76-relative-trace}.  Integrate
\eqref{eq:e8v76-density} against \(uv\), apply the divergence theorem, and
use the V74 endpoint product
\(L^4(\partial M)\cdot L^4(\partial M)\cdot L^2(\partial M)\subset L^1\).
Density of smooth \(u,v\) gives the extension.
\end{proof}

\begin{firewall}[The boundary pieces depend on the convention]
\label{fire:e8v76-boundary-convention}
On a manifold with boundary, changing the smooth reference connection changes
the separate interior \(Q\), divergence vector \(V\), and displayed boundary
flux.  Their complete Green sum is the geometric scalar-curvature integral.
Therefore BCTR must record one declared reference convention and the complete
paired expression; an isolated flux term is not reference invariant.
\end{firewall}

\subsection{Insertion into the ADM lower form}
\label{sec:e8v76-ADM}

\begin{theorem}[Global weak ADM curvature row]
\label{thm:e8v76-ADM}
On a closed three-manifold, replace the coordinate-patch curvature summand of
V73 by
\begin{equation}
 q_{R,g}(u,v):=-2\mathscr R_g(uv).                       \label{eq:e8v76-ADM-row}
\end{equation}
For strongly approximable uniformly elliptic metrics, this is a
reference-independent bounded form on \(H^1(M)\), and strong
\(L^\infty\cap W^{1,3}\) metric convergence implies its form-norm
convergence.  Adding the V73 \(L^3\) extrinsic-curvature and
\(L^{3/2}\) matter-multiplier rows preserves complete lower-form convergence.
\end{theorem}

\begin{proof}
Boundedness is \Cref{cor:e8v76-H1-form}, reference independence is
\Cref{thm:e8v76-reference}, and convergence is
\Cref{thm:e8v76-convergence}.  The remaining terms are the V73 multiplication
forms; addition is continuous in the bilinear-form norm.
\end{proof}

\begin{firewall}[Global definition does not supply global coercivity]
\label{fire:e8v76-no-coercivity}
The form \eqref{eq:e8v76-ADM-row} is globally defined, but V76 does not prove
that its negative part is relatively smaller than the principal elliptic
form.  On a closed manifold the gradient form also has a constant zero mode,
so the Dirichlet Poincar\'e estimate used on the V71 fixed patch cannot be
copied globally.  The constant/volume mode and the mean-zero sector must be
separated before applying Lax--Milgram.
\end{firewall}

\begin{assessment}[Progress at V76]
\label{ass:e8v76-status}
Scalar curvature is now a global critical weak form rather than a collection
of coordinate expressions.  The relative connection tensor gives
\(V\in L^3\), \(Q\in L^{3/2}\), strong cofinal form convergence, and exact
partition-of-unity cancellation at artificial interfaces.  A genuine
boundary still needs the V74 \(W^{2,3/2}\) trace row.  The next analytic
bottleneck is no longer chart gluing; it is the constant mode and coercivity
of the global scalar constraint.
\end{assessment}

\begin{programme}[Eighth-series V77: constant-mode Schur complement]
\label{prog:e8v76-V77}
V77 will split \(H^1(M)=\mathbb R\mathbf1\oplus H^1_\perp(M)\), use the
mean-zero Poincar\'e inequality on the second summand, and derive the exact
two-by-two Schur complement that determines whether the complete weak ADM
form is invertible.  It will identify which averaged Hamiltonian-constraint
coefficient controls the constant mode and will not hide that condition in a
patchwise Dirichlet constant.
\end{programme}

\begin{firewall}[Claim boundary after eighth-series V76]
\label{fire:e8v76-claim}
V76 proves the global reference-connection curvature identity, its critical
weak extension, cofinal stability, reference independence on the declared
smooth closure class, exact cancellation of artificial chart interfaces, and
the true-boundary trace formula.  It does not prove the stronger metric
tightness, global ADM coercivity, nonlinear constraint solvability, a coupled
continuum measure, Osterwalder--Schrader reconstruction, or the full
SM--Einstein continuum.
\end{firewall}

\begin{theorem}[Eighth-series V76 result]
\label{thm:e8v76-result}
For uniformly elliptic three-metrics in the strong
\(L^\infty\cap W^{1,3}\) closure of smooth metrics, scalar curvature defines
a global, reference-independent weak bilinear form on \(H^1(M)\), stable under
cofinal convergence.  Smooth partitions of unity glue it without internal
normal traces.  Only a true geometric boundary requires the separate V74
\(L^2\) curvature-flux row.
\end{theorem}

\begin{proof}
Use \Cref{thm:e8v76-identity,thm:e8v76-functional,
cor:e8v76-H1-form,thm:e8v76-convergence,thm:e8v76-reference,
thm:e8v76-gluing,thm:e8v76-boundary}.
\end{proof}

\paragraph{Parent-version record.}
V76 was formed from
\texttt{Renewal\_SM\_Einstein\_Continuum\_Research\_Notes\_V75.tex}.
The V75 parent had (29\,665) logical source lines, (1\,043\,350) bytes,
and SHA--256
\[
 \texttt{51661B6F7DE0DD5E6FB04A20FB63AE443E71DE6395B9ED0B8B0DDB856B33C46D}.
\]
The next compulsory companion audit is V80.

% =============================================================================
% END APPENDED EIGHTH-SERIES V76 MATERIAL
% =============================================================================
% =============================================================================
% BEGIN APPENDED EIGHTH-SERIES V77 MATERIAL
% =============================================================================

\clearpage
\section{Eighth-series V77: the global constant mode and its Schur
complement}
\label{sec:v8v77}

V76 constructs the global weak ADM scalar form on a closed spatial
manifold.  The passage from a Dirichlet patch to a closed manifold exposes a
mode that no patchwise Poincar\'e estimate can see: constants lie in the
kernel of the gradient principal part.  This note separates that volume mode,
proves the exact Schur-complement criterion for invertibility and positivity,
and identifies its ADM coefficient.  The result replaces an implicit global
coercivity assumption by one scalar test and one mean-zero coercivity row.

\subsection{Mean-zero decomposition}
\label{sec:e8v77-decomposition}

Let \((M,\bar g)\) be a connected closed smooth three-manifold.  Put
\begin{equation}
 H:=H^1(M),\qquad
 e_0:=\operatorname{Vol}_{\bar g}(M)^{-1/2},\qquad
 H_\perp:=\left\{w\in H:\int_Mw\,d\mu_{\bar g}=0\right\}.
 \label{eq:e8v77-splitting}
\end{equation}
Then \(H=\mathbb Re_0\oplus H_\perp\).  Let
\begin{equation}
 p_g(u,v):=4\int_M\langle\nabla u,\nabla v\rangle_g\,d\mu_g
 \label{eq:e8v77-principal}
\end{equation}
be the three-dimensional conformal principal form.  Uniform ellipticity and
the mean-zero Poincar\'e inequality give
\begin{equation}
 p_g(w,w)\geq\gamma_0\|w\|_{H^1}^2
 \qquad(w\in H_\perp)                                   \label{eq:e8v77-Poincare}
\end{equation}
for a positive \(\gamma_0\) depending on the fixed ellipticity window and
\((M,\bar g)\).

Let \(q:H\times H\to\mathbb R\) be a bounded symmetric lower-order form and
set
\begin{equation}
 b:=p_g+q.                                               \label{eq:e8v77-b}
\end{equation}
The V76 weak curvature row and the V73 extrinsic and matter multiplication
rows have this symmetry because each depends on the product \(uv\).

\begin{assumption}[Mean-zero coercivity row]
\label{ass:e8v77-perp}
There is \(\gamma>0\) such that
\begin{equation}
 b(w,w)\geq\gamma\|w\|_{H^1}^2
 \qquad(w\in H_\perp).                                  \label{eq:e8v77-perp-coercivity}
\end{equation}
\end{assumption}

\begin{remark}[A sufficient inherited condition]
\label{rem:e8v77-relative}
If
\[
 |q(w,w)|\leq\alpha\|w\|_{H^1}^2
 \quad(w\in H_\perp),
 \qquad \alpha<\gamma_0,
\]
then \Cref{ass:e8v77-perp} holds with
\(\gamma=\gamma_0-\alpha\).  Unlike the V71 Dirichlet estimate, this uses the
mean-zero Poincar\'e inequality and makes no claim on constants.
\end{remark}

\subsection{The exact scalar Schur complement}
\label{sec:e8v77-Schur}

Define
\begin{equation}
 a:=b(e_0,e_0),
 \qquad
 \ell(w):=b(e_0,w)quad(w\in H_\perp).                  \label{eq:e8v77-aell}
\end{equation}
By \Cref{ass:e8v77-perp} and Lax--Milgram, there is a unique
\(z\in H_\perp\) such that
\begin{equation}
 b(z,w)=\ell(w)qquad(w\in H_\perp).                    \label{eq:e8v77-z}
\end{equation}
Set
\begin{equation}
 s:=a-b(z,z)=a-\ell(z).                                 \label{eq:e8v77-Schur}
\end{equation}

\begin{theorem}[Exact completion of the constant mode]
\label{thm:e8v77-completion}
For every \(c\in\mathbb R\) and \(w\in H_\perp\),
\begin{equation}
 b(ce_0+w,ce_0+w)
 =b(w+cz,w+cz)+s c^2.                                   \label{eq:e8v77-completion}
\end{equation}
\end{theorem}

\begin{proof}
By symmetry and \eqref{eq:e8v77-z},
\[
 b(w+cz,w+cz)
 =b(w,w)+2c\ell(w)+c^2b(z,z).
\]
On the other hand,
\[
 b(ce_0+w,ce_0+w)=b(w,w)+2c\ell(w)+a c^2.
\]
Subtract and use \eqref{eq:e8v77-Schur}.
\end{proof}

\begin{theorem}[Global invertibility criterion]
\label{thm:e8v77-invertibility}
Assume \Cref{ass:e8v77-perp}.  The operator
\(B:H\to H^*\) defined by \(\langle Bu,v\rangle=b(u,v)\) is an isomorphism
if and only if
\begin{equation}
 s\ne0.                                                  \label{eq:e8v77-nonzero}
\end{equation}
More explicitly, for \(F\in H^*\), let \(y\in H_\perp\) be the unique
solution of
\begin{equation}
 b(y,w)=F(w)qquad(w\in H_\perp).                       \label{eq:e8v77-y}
\end{equation}
Then the solution \(u=ce_0+w\) of \(Bu=F\) is
\begin{align}
 c&={F(e_0)-\ell(y)\over s},                            \label{eq:e8v77-c}\\
 w&=y-cz.                                               \label{eq:e8v77-w}
\end{align}
\end{theorem}

\begin{proof}
Testing \(b(ce_0+w,\cdot)=F\) on \(H_\perp\) gives
\[
 b(w+cz,v)=F(v)qquad(v\in H_\perp),
\]
so uniqueness in Lax--Milgram gives \(w+cz=y\).  Testing on \(e_0\) and
using symmetry yields
\[
 ca+\ell(w)=F(e_0).
\]
Substitute \(w=y-cz\) and use
\(a-\ell(z)=s\) to obtain \eqref{eq:e8v77-c} and
\eqref{eq:e8v77-w}.  Hence \(s\ne0\) implies existence and uniqueness for
every \(F\), with a bounded inverse by the displayed formulas.

If \(s=0\), set \(h=e_0-z\).  Equation \eqref{eq:e8v77-z} gives
\(b(h,v)=0\) for \(v\in H_\perp\), while
\(b(h,e_0)=a-\ell(z)=s=0\).  Thus \(Bh=0\), and
\(h\ne0\) because its constant component is \(e_0\).  The operator is not
injective.
\end{proof}

\begin{corollary}[Compatibility at a zero Schur complement]
\label{cor:e8v77-compatibility}
If \(s=0\), then \(h=e_0-z\) spans the kernel of \(B\).  The equation
\(Bu=F\) can be solvable only if
\begin{equation}
 F(h)=0.                                                 \label{eq:e8v77-compatibility}
\end{equation}
If this condition holds, solutions are unique only modulo \(h\).
\end{corollary}

\begin{proof}
The proof of \Cref{thm:e8v77-invertibility} shows that every kernel element
has mean-zero component \(-cz\), hence is a multiple of \(h\).  Symmetry
implies \(F(h)=b(u,h)=0\) for every solution.  The remaining assertion follows
by adding multiples of the kernel vector.
\end{proof}

\begin{theorem}[Global positivity criterion]
\label{thm:e8v77-positivity}
Under \Cref{ass:e8v77-perp}, the full form \(b\) is coercive on \(H\) if and
only if
\begin{equation}
 s>0.                                                    \label{eq:e8v77-positive-Schur}
\end{equation}
If \(s>0\), a coercivity constant depends only on
\(\gamma,s,\|z\|_{H^1}\), and the norm of the fixed splitting
\(H=\mathbb Re_0\oplus H_\perp\).
\end{theorem}

\begin{proof}
If \(s>0\), \Cref{thm:e8v77-completion} gives
\[
 b(ce_0+w,ce_0+w)
 \geq \gamma\|w+cz\|_{H^1}^2+s c^2.
\]
The triangular map
\((c,w)\mapsto(c,w+cz)\) is a bounded isomorphism of
\(\mathbb R\oplus H_\perp\), so the right side controls
\(\|ce_0+w\|_{H^1}^2\).

Conversely, evaluate \eqref{eq:e8v77-completion} at \(w=-cz\).  Then
\[
 b(c(e_0-z),c(e_0-z))=s c^2.
\]
Coercivity for nonzero \(c\) forces \(s>0\).
\end{proof}

\begin{corollary}[A checkable sufficient margin]
\label{cor:e8v77-margin}
Let \(\|\ell\|_{H_\perp^*}\) be computed with the \(H^1\) norm.  Then
\begin{equation}
 0\leq b(z,z)=\ell(z)
 \leq {\|\ell\|_{H_\perp^*}^2\over\gamma}.             \label{eq:e8v77-mixing-bound}
\end{equation}
Consequently
\begin{equation}
 a>{\|\ell\|_{H_\perp^*}^2\over\gamma}                \label{eq:e8v77-sufficient}
\end{equation}
is sufficient for global coercivity.
\end{corollary}

\begin{proof}
Coercivity and \eqref{eq:e8v77-z} give
\[
 \gamma\|z\|_{H^1}^2
 \leq b(z,z)=\ell(z)
 \leq\|\ell\|_{H_\perp^*}\|z\|_{H^1}.
\]
Thus \(\|z\|_{H^1}\leq\|\ell\|/\gamma\), which proves
\eqref{eq:e8v77-mixing-bound}.  Combine it with
\(s=a-b(z,z)\) and apply \Cref{thm:e8v77-positivity}.
\end{proof}

\subsection{The ADM value of the constant block}
\label{sec:e8v77-ADM}

Let the complete global lower form use the V72 convention and V76 curvature
distribution:
\begin{align}
 q_{\rm ADM}(u,v)
 :={}&-2\mathscr R_g(uv)                                  \notag\\
 &+\int_M
 \left[
 4\bigl(|K|_g^2-\tau^2\bigr)-2\kappa_Gm_\rho
 \right]uv\,d\mu_g.                                    \label{eq:e8v77-qADM}
\end{align}
Here \(m_\rho\) is the declared algebraic matter metric multiplier from V72.

\begin{theorem}[Exact averaged constant-mode coefficient]
\label{thm:e8v77-aADM}
For \(b=p_g+q_{\rm ADM}\),
\begin{equation}
 a={1\over\operatorname{Vol}_{\bar g}(M)}
 \left{
 -2\mathscr R_g(1)
 +\int_M
 \left[4\bigl(|K|_g^2-\tau^2\bigr)-2\kappa_Gm_\rho\right]
 d\mu_g
 \right}.                                              \label{eq:e8v77-aADM}
\end{equation}
If \(g\) is smooth, this is
\begin{equation}
 a={1\over\operatorname{Vol}_{\bar g}(M)}
 \int_M
 \left[-2R_g+4\bigl(|K|_g^2-\tau^2\bigr)
             -2\kappa_Gm_\rho\right]d\mu_g.            \label{eq:e8v77-aADM-smooth}
\end{equation}
If the background data also obey the Hamiltonian constraint
\begin{equation}
 R_g+\tau^2-|K|_g^2-2\kappa_G\rho=0,                   \label{eq:e8v77-constraint}
\end{equation}
then
\begin{equation}
 a={1\over\operatorname{Vol}_{\bar g}(M)}
 \int_M
 \left[
 2\bigl(|K|_g^2-\tau^2\bigr)-4\kappa_G\rho
 -2\kappa_Gm_\rho
 \right]d\mu_g.                                        \label{eq:e8v77-a-constraint}
\end{equation}
\end{theorem}

\begin{proof}
The principal gradient form vanishes on constants and
\(e_0^2=\operatorname{Vol}_{\bar g}(M)^{-1}\).  Substitution in
\eqref{eq:e8v77-qADM} gives \eqref{eq:e8v77-aADM}.  For smooth \(g\), V76
identifies \(\mathscr R_g(1)=\int_MR_g\,d\mu_g\), giving
\eqref{eq:e8v77-aADM-smooth}.  Finally
\eqref{eq:e8v77-constraint} gives
\(R_g=|K|_g^2-\tau^2+2\kappa_G\rho\); insert this into
\eqref{eq:e8v77-aADM-smooth} and simplify.
\end{proof}

\begin{corollary}[A necessary average sign for positive coercivity]
\label{cor:e8v77-average-sign}
If \(b=p_g+q_{\rm ADM}\) is positive coercive on all of \(H^1(M)\), then
\begin{equation}
 a>0.                                                     \label{eq:e8v77-a-positive}
\end{equation}
Thus the right side of \eqref{eq:e8v77-aADM}, or of
\eqref{eq:e8v77-a-constraint} on a smooth constraint background, must be
strictly positive.  Positivity of this average is not sufficient unless the
mixing debit \(b(z,z)\) is also smaller than \(a\).
\end{corollary}

\begin{proof}
Test coercivity on \(e_0\) to obtain \(a=b(e_0,e_0)>0\).  The last statement
is the exact formula \(s=a-b(z,z)\) and
\Cref{thm:e8v77-positivity}.
\end{proof}

\begin{firewall}[The constant mode is not silently gauge]
\label{fire:e8v77-not-gauge}
A constant conformal variation changes the spatial volume.  V77 does not
identify it with a spatial diffeomorphism or remove it by gauge.  A
fixed-volume ensemble could impose the mean-zero condition, but that is a
different declared state space and requires a corresponding treatment of the
integrated Hamiltonian constraint.
\end{firewall}

\subsection{Cofinal stability of the Schur test}
\label{sec:e8v77-stability}

\begin{theorem}[Schur-complement convergence]
\label{thm:e8v77-Schur-convergence}
Let \(b_n,b\) be bounded symmetric forms on the fixed space \(H\) such that
\begin{equation}
 \|b_n-b\|_{\mathcal B(H\times H)}\longrightarrow0,     \label{eq:e8v77-form-convergence}
\end{equation}
and suppose their restrictions to \(H_\perp\) share the coercivity constant
\(\gamma>0\).  Let \(a_n,\ell_n,z_n,s_n\) be defined by
\eqref{eq:e8v77-aell}--\eqref{eq:e8v77-Schur}.  Then
\begin{equation}
 a_n\to a,qquad
 \ell_n\to\ell\text{ in }H_\perp^*,qquad
 z_n\to z\text{ in }H^1,qquad
 s_n\to s.                                              \label{eq:e8v77-Schur-convergence}
\end{equation}
If \(s\ne0\), then \(B_n^{-1}\to B^{-1}\) in
\(\mathcal B(H^*,H)\).
\end{theorem}

\begin{proof}
Form-norm convergence immediately gives \(a_n\to a\) and
\(\ell_n\to\ell\).  Subtract the equations for \(z_n\) and \(z\):
\[
 b_n(z_n-z,w)
 = (\ell_n-\ell)(w)+(b-b_n)(z,w).
\]
Set \(w=z_n-z\) and use the common coercivity constant to obtain
\[
 \|z_n-z\|_{H^1}
 \leq\gamma^{-1}
 \left(\|\ell_n-\ell\|_{H_\perp^*}
       +\|b_n-b\|\|z\|_{H^1}\right).
\]
Thus \(z_n\to z\), and the definition of \(s_n\) gives \(s_n\to s\).
If \(s\ne0\), then \(|s_n|\geq|s|/2\) eventually.  The explicit solution
formulas \eqref{eq:e8v77-c}--\eqref{eq:e8v77-w}, together with the standard
Lax--Milgram resolvent identity on \(H_\perp\), give operator-norm convergence
of the full inverses.
\end{proof}

\begin{corollary}[Global extension of the V71 stability compiler]
\label{cor:e8v77-V71}
The V71 form-norm convergence theorem extends from a Dirichlet patch to a
closed spatial manifold provided that:
\begin{enumerate}[label=\textup{(\roman*)},leftmargin=3.8em]
\item the complete forms are symmetric and converge in form norm;
\item their mean-zero restrictions share a positive coercivity margin; and
\item the limiting Schur complement satisfies \(s\ne0\), or \(s>0\) when
positive coercivity is required.
\end{enumerate}
\end{corollary}

\begin{proof}
Apply \Cref{thm:e8v77-Schur-convergence}; V76 and V73 supply the form-norm
convergence of the declared ADM coefficients.
\end{proof}

\begin{assessment}[Progress at V77]
\label{ass:e8v77-status}
The global scalar resolvent no longer depends on an unspoken Poincar\'e
inequality for constants.  Its full obstruction is the exact scalar
\(s=a-b(z,z)\): \(a\) is the averaged conformal ADM coefficient, while
\(b(z,z)\) is the mixing debit into mean-zero modes.  The global operator is
invertible exactly when \(s\ne0\), positive coercive exactly when \(s>0\),
and stable under cofinal form convergence when this margin persists.
\end{assessment}

\begin{programme}[Eighth-series V78: physical source compatibility and
volume fixing]
\label{prog:e8v77-V78}
V78 will apply the kernel vector \(h=e_0-z\) to the physical ADM source of
V72.  It will derive the exact compatibility condition when \(s=0\), compare
the unconstrained and fixed-volume source spaces, and determine whether the
integrated Hamiltonian constraint or the V51 balance law supplies that
condition without applying an inverse circularly.
\end{programme}

\begin{firewall}[Claim boundary after eighth-series V77]
\label{fire:e8v77-boundary}
V77 proves the constant/mean-zero Schur decomposition, exact invertibility and
positivity criteria, the averaged ADM constant coefficient, and cofinal
stability.  It does not prove that the mean-zero coercivity or Schur sign holds
for the physical continuum sequence, justify fixed-volume reduction, solve
the nonlinear constraints, construct the coupled measure, or establish the
full SM--Einstein continuum.
\end{firewall}

\begin{theorem}[Eighth-series V77 result]
\label{thm:e8v77-result}
On a closed spatial manifold, mean-zero coercivity plus the single scalar
Schur complement \(s=a-b(z,z)\) completely determines the global conformal
ADM linearized form: it is invertible iff \(s\ne0\), positive coercive iff
\(s>0\), and cofinally stable away from zero.  The coefficient \(a\) is the
explicit averaged ADM expression in \eqref{eq:e8v77-aADM}.
\end{theorem}

\begin{proof}
Use \Cref{thm:e8v77-invertibility,thm:e8v77-positivity,
thm:e8v77-aADM,thm:e8v77-Schur-convergence}.
\end{proof}

\paragraph{Parent-version record.}
V77 was formed from
\texttt{Renewal\_SM\_Einstein\_Continuum\_Research\_Notes\_V76.tex}.
The V76 parent had (30\,187) logical source lines, (1\,063\,387) bytes,
and SHA--256
\[
 \texttt{ABEB67CCF3E90349CD440AD20D6358D806EF9FAEE5DD1078637DFD32A11B4C33}.
\]
The next compulsory companion audit is V80.

% =============================================================================
% END APPENDED EIGHTH-SERIES V77 MATERIAL
% =============================================================================
% =============================================================================
% BEGIN APPENDED EIGHTH-SERIES V78 MATERIAL
% =============================================================================

\clearpage
\section{Eighth-series V78: physical source compatibility and the
fixed-volume alternative}
\label{sec:v8v78}

V77 shows that a degenerate global conformal ADM operator has the one-
dimensional kernel generated by \(h=e_0-z\).  This note computes the exact
condition imposed on the physical \((\dot K,\dot\Psi)\) source, tests the V51
balance law against it, and proves the fixed-volume saddle-point alternative.
The main distinction is sharp: conservation can remove the constant part of
a source, but it removes the full resonant component only when the constant
mode does not mix with the mean-zero sector.

\subsection{The physical source functional}
\label{sec:e8v78-source}

Retain the closed-manifold setting and notation of V77.  Let
\(b:H^1(M)\times H^1(M)\to\mathbb R\) be the complete symmetric conformal ADM
form, coercive on \(H_\perp\), with Schur data
\begin{equation}
 a,\quad \ell,\quad z,\quad s=a-b(z,z),\quad h=e_0-z.
 \label{eq:e8v78-Schur-data}
\end{equation}
For a canonical and matter perturbation \((\dot K,\dot\Psi)\), V72 gives
\begin{equation}
 F_{\rm phys}
 =-2\langle K,\dot K\rangle_g
  +2\tau\operatorname{tr}_g\dot K
  -2\kappa_GD_\Psi\rho[\dot\Psi].                      \label{eq:e8v78-F-density}
\end{equation}
Whenever this density belongs to \(L^{6/5}(M,d\mu_g)\), it defines
\begin{equation}
 F(v):=\int_MF_{\rm phys}v\,d\mu_g\in H^1(M)^*.         \label{eq:e8v78-F-functional}
\end{equation}
The same argument below applies directly to any \(F\in H^1(M)^*\).

\begin{lemma}[Source continuity]
\label{lem:e8v78-source-continuity}
Under the common ellipticity window,
\begin{equation}
 \|F\|_{H^1{}^*}
 \leq C_{M,g}\|F_{\rm phys}\|_{L^{6/5}(d\mu_g)}.       \label{eq:e8v78-source-bound}
\end{equation}
In particular, \(K,\dot K\in L^{12/5}\) and
\(D_\Psi\rho[\dot\Psi]\in L^{6/5}\) are a sufficient source class.
\end{lemma}

\begin{proof}
H\"older and \(H^1\hookrightarrow L^6\) give
\[
 |F(v)|\leq\|F_{\rm phys}\|_{6/5}\|v\|_6
 \leq C\|F_{\rm phys}\|_{6/5}\|v\|_{H^1}.
\]
The product of two \(L^{12/5}\) tensors lies in \(L^{6/5}\); uniform
ellipticity controls the contractions in \eqref{eq:e8v78-F-density}.
\end{proof}

\subsection{The exact resonant condition}
\label{sec:e8v78-resonance}

The linearized constraint is the weak equation
\begin{equation}
 b(\phi,v)+F(v)=0
 \qquad(v\in H^1(M)).                                   \label{eq:e8v78-equation}
\end{equation}

\begin{theorem}[Physical Fredholm compatibility]
\label{thm:e8v78-compatibility}
Assume the mean-zero coercivity row of V77 and suppose \(s=0\).  Equation
\eqref{eq:e8v78-equation} is solvable if and only if
\begin{equation}
 F(h)=0.                                                 \label{eq:e8v78-compatibility}
\end{equation}
For the density \eqref{eq:e8v78-F-density}, this condition is
\begin{align}
 0=\int_M\bigl[&-2\langle K,\dot K\rangle_g
  +2\tau\operatorname{tr}_g\dot K                         \notag\\
 &-2\kappa_GD_\Psi\rho[\dot\Psi]\bigr](e_0-z)\,d\mu_g.
 \label{eq:e8v78-physical-compatibility}
\end{align}
When it holds, solutions are unique modulo \(h\).
\end{theorem}

\begin{proof}
V77 proves that \(\ker B=\mathbb Rh\) when \(s=0\).  Since \(B\) is
symmetric, its range is the annihilator of this kernel: necessity follows by
testing \eqref{eq:e8v78-equation} on \(h\), and sufficiency follows from the
explicit mean-zero solution and constant-mode formula in V77.  Substituting
\eqref{eq:e8v78-F-density} gives
\eqref{eq:e8v78-physical-compatibility}.
\end{proof}

\begin{corollary}[The decoupled constant-mode case]
\label{cor:e8v78-decoupled}
If \(\ell=0\), then \(z=0\) and \(h=e_0\).  At \(s=a=0\), compatibility is
exactly the vanishing of the spatial mean:
\begin{equation}
 \int_MF_{\rm phys}\,d\mu_g=0.                          \label{eq:e8v78-zero-mean}
\end{equation}
\end{corollary}

\begin{proof}
The defining equation \(b(z,w)=\ell(w)\) and coercivity on \(H_\perp\) give
\(z=0\).  Since \(e_0\) is a nonzero constant, \(F(e_0)=0\) is equivalent to
\eqref{eq:e8v78-zero-mean}.
\end{proof}

\begin{theorem}[Zero mean is insufficient when the blocks mix]
\label{thm:e8v78-mean-insufficient}
Suppose \(s=0\) and \(z\ne0\).  There exists \(F\in H^1(M)^*\) such that
\begin{equation}
 F(e_0)=0,qquad F(h)\ne0.                               \label{eq:e8v78-counter-source}
\end{equation}
Thus an integrated source balance alone does not imply the Fredholm
compatibility condition.
\end{theorem}

\begin{proof}
Use the Hilbert \(H^1\) inner product on \(H_\perp\).  Define
\(F(ce_0+w):=(w,z)_{H^1}\).  This is a bounded functional and
\(F(e_0)=0\).  Since \(h=e_0-z\),
\[
 F(h)=-(z,z)_{H^1}\ne0.
\]
\end{proof}

\subsection{What the V51 balance law does and does not close}
\label{sec:e8v78-balance}

Suppose a source has the distributional balance form on the closed manifold
\begin{equation}
 F=-\operatorname{div}_gJ+W,
 \qquad J\in L^2(TM),\quad W\in H^1(M)^*.               \label{eq:e8v78-balance}
\end{equation}
Here \(W\) denotes the geometric-work or anomaly residual in the sense of
V51, interpreted as a functional when necessary.

\begin{theorem}[Balance-law projection onto the resonant mode]
\label{thm:e8v78-balance-projection}
For \(h=e_0-z\),
\begin{equation}
 F(h)
 =-\int_M\langle J,\nabla z\rangle_g\,d\mu_g+W(e_0-z).
 \label{eq:e8v78-balance-resonance}
\end{equation}
In contrast, the integrated balance controls only
\begin{equation}
 F(e_0)=W(e_0).                                         \label{eq:e8v78-balance-mean}
\end{equation}
Consequently, exact conservation \(W=0\) implies Fredholm compatibility when
\(z=0\), but not in general when \(z\ne0\).
\end{theorem}

\begin{proof}
The weak divergence convention gives
\[
 F(v)=\int_M\langle J,\nabla v\rangle_g\,d\mu_g+W(v).
\]
Constants have zero gradient.  Substitute first \(v=e_0-z\) and then
\(v=e_0\) to obtain \eqref{eq:e8v78-balance-resonance} and
\eqref{eq:e8v78-balance-mean}.
\end{proof}

\begin{corollary}[Quantitative residual supplied by V51 data]
\label{cor:e8v78-balance-bound}
If \(W\) is bounded on \(H^1\), then
\begin{equation}
 |F(h)|
 \leq \|J\|_2\|\nabla z\|_2
 +\|W\|_{H^1{}^*}\bigl(\|e_0\|_{H^1}+\|z\|_{H^1}\bigr),
 \label{eq:e8v78-balance-bound}
\end{equation}
and the V77 estimate gives
\begin{equation}
 \|z\|_{H^1}\leq{\|\ell\|_{H_\perp^*}\over\gamma}.   \label{eq:e8v78-z-bound}
\end{equation}
These are size bounds, not exact compatibility when \(s=0\).
\end{corollary}

\begin{proof}
Apply Cauchy--Schwarz and duality to
\eqref{eq:e8v78-balance-resonance}, then use the V77 mixing bound.
\end{proof}

\begin{firewall}[Conservation is not a resonant Ward identity]
\label{fire:e8v78-conservation}
The V51 conservation law annihilates constants on a closed stationary
background.  The kernel vector of the full ADM form is \(h=e_0-z\), which is
not constant when lower-order coefficients mix the two blocks.  Closing
\eqref{eq:e8v78-compatibility} therefore requires either \(\ell=0\), an exact
identity annihilating \(h\), or a nonzero Schur complement.  A small bound on
\(F(h)\) cannot replace equality at \(s=0\).
\end{firewall}

\subsection{Near-resonant cofinal sequences}
\label{sec:e8v78-near-resonance}

For \(s\ne0\), let \(y_F\in H_\perp\) solve
\begin{equation}
 b(y_F,w)=-F(w)qquad(w\in H_\perp).                    \label{eq:e8v78-yF}
\end{equation}

\begin{theorem}[Exact pole numerator]
\label{thm:e8v78-pole}
The solution of \eqref{eq:e8v78-equation} has constant coefficient
\begin{equation}
 c=-{F(h)\over s},                                      \label{eq:e8v78-pole}
\end{equation}
and mean-zero part
\begin{equation}
 w=y_F-cz.                                               \label{eq:e8v78-pole-w}
\end{equation}
Thus the only constant-mode pole is the ratio of the resonant source
projection to the Schur complement.
\end{theorem}

\begin{proof}
Apply the V77 inverse formula to the right side \(-F\).  It gives
\[
 c={-F(e_0)-\ell(y_F)\over s}.
\]
The equations for \(y_F\) and \(z\), plus symmetry, imply
\[
 F(z)=-b(y_F,z)=-b(z,y_F)=-\ell(y_F).
\]
Hence \(-F(e_0)-\ell(y_F)=-F(e_0)+F(z)=-F(h)\), proving
\eqref{eq:e8v78-pole}.  The mean-zero formula is the same block elimination.
\end{proof}

\begin{corollary}[Necessary scaling at a closing Schur gap]
\label{cor:e8v78-scaling}
Let \(b_n,F_n\) be a cofinal sequence satisfying the V77 uniform mean-zero
coercivity hypotheses, and let \(s_n\to0\).  If the corresponding solutions
\(\phi_n\) remain bounded in \(H^1\), then necessarily
\begin{equation}
 |F_n(h_n)|\leq C|s_n|.                                 \label{eq:e8v78-scaling}
\end{equation}
If \(F_n\to F\) in \(H^*\) and \(h_n\to h\) in \(H^1\), this implies the
limiting compatibility \(F(h)=0\).
\end{corollary}

\begin{proof}
The continuous projection onto constants bounds \(|c_n|\) by
\(C\|\phi_n\|_{H^1}\).  Equation \eqref{eq:e8v78-pole} then gives
\eqref{eq:e8v78-scaling}.  Strong convergence of \(F_n,h_n\) passes the
pairing to the limit.
\end{proof}

\begin{firewall}[Limit compatibility alone does not prevent blow-up]
\label{fire:e8v78-rate}
The condition \(F(h)=0\) at the limiting degenerate operator does not control
\(F_n(h_n)/s_n\).  A cofinal construction approaching \(s=0\) must prove the
rate \eqref{eq:e8v78-scaling}, or choose a normalized state space in which the
constant coefficient is removed with a declared multiplier.
\end{firewall}

\subsection{Fixed volume as an augmented equation}
\label{sec:e8v78-volume}

At a fixed metric \(g\), define
\begin{equation}
 H_{\rm vol}:=left\{w\in H^1(M):\int_Mw\,d\mu_g=0\right\}.
 \label{eq:e8v78-Hvol}
\end{equation}
Assume the restriction of \(b\) to \(H_{\rm vol}\) is coercive.

\begin{theorem}[Fixed-volume saddle-point solution]
\label{thm:e8v78-fixed-volume}
For every \(F\in H^1(M)^*\), there is a unique pair
\((\phi,\lambda)\in H_{\rm vol}\times\mathbb R\) satisfying
\begin{equation}
 b(\phi,v)+\lambda\int_Mv\,d\mu_g=-F(v)
 \qquad(v\in H^1(M)).                                   \label{eq:e8v78-augmented}
\end{equation}
It is given by the unique mean-zero solution
\begin{equation}
 b(\phi,w)=-F(w)qquad(w\in H_{\rm vol})                \label{eq:e8v78-volume-solve}
\end{equation}
and
\begin{equation}
 \lambda=-{F(1)+b(\phi,1)\over\operatorname{Vol}_g(M)}.
 \label{eq:e8v78-lambda}
\end{equation}
This construction does not require a nonzero full-space Schur complement.
\end{theorem}

\begin{proof}
Lax--Milgram on \(H_{\rm vol}\) gives the unique \(\phi\) in
\eqref{eq:e8v78-volume-solve}.  Every \(v\in H^1\) decomposes uniquely as
\[
 v=w+{1\over\operatorname{Vol}_g(M)}
       \left(\int_Mv\,d\mu_g\right)1,
 \qquad w\in H_{\rm vol}.
\]
Equation \eqref{eq:e8v78-volume-solve} handles \(w\), while testing the
remaining constant direction gives \eqref{eq:e8v78-lambda}.  Thus
\eqref{eq:e8v78-augmented} holds for every \(v\).  The same decomposition
proves uniqueness.
\end{proof}

\begin{lemma}[Linearized meaning of fixed volume]
\label{lem:e8v78-volume-variation}
For a three-dimensional conformal variation \(\dot g=2\phi g\),
\begin{equation}
 {d\over dt}\operatorname{Vol}_{g(t)}(M)\bigg|_{t=0}
 =3\int_M\phi\,d\mu_g.                                 \label{eq:e8v78-volume-variation}
\end{equation}
Thus \(\phi\in H_{\rm vol}\) is exactly the infinitesimal fixed-volume
condition.
\end{lemma}

\begin{proof}
The standard density variation is
\(D(d\mu_g)[\dot g]=\tfrac12\operatorname{tr}_g(\dot g)d\mu_g\).
For \(\dot g=2\phi g\) in dimension three, the trace is \(6\phi\), giving
\eqref{eq:e8v78-volume-variation}.
\end{proof}

\begin{corollary}[When the augmented solution solves the original constraint]
\label{cor:e8v78-lambda-zero}
The fixed-volume solution satisfies the unmodified linearized constraint
\eqref{eq:e8v78-equation} if and only if \(\lambda=0\).  When the full
operator has \(s=0\), this is equivalent to \(F(h)=0\).
\end{corollary}

\begin{proof}
The first equivalence is immediate from \eqref{eq:e8v78-augmented}.  If
\(s=0\), \Cref{thm:e8v78-compatibility} is the necessary and sufficient
condition for the unmodified equation.
\end{proof}

\begin{firewall}[Volume fixing changes the equation when \(\lambda\ne0\)]
\label{fire:e8v78-volume}
The multiplier in \eqref{eq:e8v78-augmented} is a spatially constant source.
It may be interpreted as a volume constraint reaction, and in a declared
action it may correspond to varying a cosmological parameter.  V78 does not
identify it with an already fixed physical cosmological constant.  Therefore
fixed-volume solvability with \(\lambda\ne0\) is not solvability of the
original Hamiltonian constraint.
\end{firewall}

\subsection{Non-circular source card}
\label{sec:e8v78-card}

\begin{definition}[Global conformal source compatibility card, GCSCC]
\label{def:e8v78-GCSCC}
For a cofinal family of complete global ADM forms and sources, GCSCC records:
\begin{enumerate}[label=\textup{(S\arabic*)},leftmargin=4em]
\item uniform coercivity on the declared mean-zero or fixed-volume space;
\item the Schur data \(z_n,s_n,h_n\), constructed from the form without using
the physical source;
\item either a uniform nonzero bound on \(|s_n|\), or the resonant estimate
\(|F_n(h_n)|\leq C|s_n|\);
\item at a degenerate limit, the exact Ward or constraint identity
\(F(h)=0\);
\item if V51 balance is used, separate estimates for
\(-\int J\cdot\nabla z\) and \(W(e_0-z)\);
\item if volume is fixed, the multiplier \(\lambda_n\) and proof that its
limit has the intended physical interpretation; and
\item stability of all pairings under the V73--V77 cofinal topologies.
\end{enumerate}
\end{definition}

\begin{firewall}[No inverse-defined compatibility]
\label{fire:e8v78-circularity}
The identity \(F(h)=0\) must come from the physical linearized constraints,
a Ward identity, or a declared restricted source space.  Defining admissible
sources by first solving \(B\phi=-F\), or projecting them with an inverse of
the nearly singular full operator, is circular.  The mean-zero inverse used
to construct \(z\) is legitimate only after its independent coercivity row
has been proved.
\end{firewall}

\begin{assessment}[Progress at V78]
\label{ass:e8v78-status}
The global source obstruction is now explicit.  At a zero Schur complement,
the V72 canonical--matter source must annihilate \(h=e_0-z\).  V51
conservation annihilates only the constant component; the remaining flux and
work pairing with \(z\) is the exact unpaid term.  Fixed volume gives a unique
augmented solution, but its multiplier must vanish or receive an independent
physical interpretation before it solves the original constraint.
\end{assessment}

\begin{programme}[Eighth-series V79: nonlinear augmented constraint chart]
\label{prog:e8v78-V79}
V79 will formulate a Banach-space implicit-function theorem for the nonlinear
Hamiltonian constraint together with the volume normalization.  Its
linearization will be the saddle-point operator of
\Cref{thm:e8v78-fixed-volume}.  The note will state the precise differentiable
mapping spaces and prove local existence only under verified mean-zero
coercivity and coefficient regularity, keeping the multiplier and momentum
constraint as explicit rows.
\end{programme}

\begin{firewall}[Claim boundary after eighth-series V78]
\label{fire:e8v78-boundary}
V78 proves the exact physical source compatibility, its balance-law
decomposition, the near-resonant scaling condition, and the fixed-volume
saddle-point theorem.  It does not prove the required physical Ward identity,
vanishing of the volume multiplier, nonlinear constraint solvability, the
momentum constraint, measure tightness, or the full SM--Einstein continuum.
\end{firewall}

\begin{theorem}[Eighth-series V78 result]
\label{thm:e8v78-result}
When the global conformal ADM Schur complement vanishes, the physical source
is solvable exactly when it annihilates \(h=e_0-z\).  Conservation supplies
this only in the decoupled case or after the explicit \(z\)-weighted flux and
work terms cancel.  A fixed-volume formulation is always solvable under
mean-zero coercivity, but generally introduces the nonzero multiplier
\eqref{eq:e8v78-lambda}.
\end{theorem}

\begin{proof}
Use \Cref{thm:e8v78-compatibility,thm:e8v78-balance-projection,
thm:e8v78-pole,thm:e8v78-fixed-volume,cor:e8v78-lambda-zero}.
\end{proof}

\paragraph{Parent-version record.}
V78 was formed from
\texttt{Renewal\_SM\_Einstein\_Continuum\_Research\_Notes\_V77.tex}.
The V77 parent had (30\,618) logical source lines, (1\,078\,811) bytes,
and SHA--256
\[
 \texttt{B629A356F4013019E01531A194690A2EFC13DC868097F641C2F3858CF8768078}.
\]
The next compulsory companion audit is V80.

% =============================================================================
% END APPENDED EIGHTH-SERIES V78 MATERIAL
% =============================================================================
% =============================================================================
% BEGIN APPENDED EIGHTH-SERIES V79 MATERIAL
% =============================================================================

\clearpage
\section{Eighth-series V79: a nonlinear fixed-volume Hamiltonian-constraint
chart}
\label{sec:v8v79}

V78 proves that the volume-normalized linearized constraint is an invertible
saddle system under mean-zero coercivity, even when the unconstrained Schur
complement vanishes.  This note lifts that statement to a local nonlinear
implicit-function theorem in a deliberately strong Sobolev class.  The strong
class is an analytic chart, not a claim of continuum tightness.  The volume
multiplier remains part of the solution and the original Hamiltonian
constraint is recovered only on the locus where that multiplier vanishes.

\subsection{Strong chart and exact conformal formula}
\label{sec:e8v79-chart}

Let \(M\) be a connected closed smooth three-manifold, fix \(p>3\), and let
\(g\in W^{2,p}\) be uniformly positive.  Let \(K\in L^{2p}\) be a covariant
symmetric two-tensor.  For \(\varphi\in W^{2,p}(M)\), put
\begin{equation}
 g_\varphi:=e^{2\varphi}g.                              \label{eq:e8v79-conformal-metric}
\end{equation}
Because \(W^{2,p}\hookrightarrow C^{1,\alpha}\),
\(\alpha=1-3/p>0\), the positive-metric condition is open and all products
below lie in \(L^p\).

\begin{lemma}[Exact three-dimensional conformal identities]
\label{lem:e8v79-conformal}
Holding the covariant tensor \(K\) fixed while varying the metric,
\begin{align}
 R_{g_\varphi}
 &=e^{-2\varphi}
 \left(R_g-4\Delta_g\varphi-2|\nabla\varphi|_g^2\right),
 \label{eq:e8v79-R}\
 \tau_{g_\varphi}&=e^{-2\varphi}\tau_g,
 &|K|_{g_\varphi}^2&=e^{-4\varphi}|K|_g^2,              \label{eq:e8v79-K}
\end{align}
and hence
\begin{equation}
 \tau_{g_\varphi}^2-|K|_{g_\varphi}^2
 =e^{-4\varphi}(\tau_g^2-|K|_g^2).                     \label{eq:e8v79-extrinsic}
\end{equation}
\end{lemma}

\begin{proof}
The scalar-curvature transformation in dimension \(d\) is
\[
 R_{e^{2\varphi}g}
 =e^{-2\varphi}
 \left(R_g-2(d-1)\Delta_g\varphi
 -(d-1)(d-2)|\nabla\varphi|_g^2\right).
\]
Set \(d=3\).  Since \(g_\varphi^{-1}=e^{-2\varphi}g^{-1}\), one inverse
metric in the trace of \(K\) gives the first scaling in
\eqref{eq:e8v79-K}, and two inverse metrics in its squared norm give the
second.  Squaring the trace proves \eqref{eq:e8v79-extrinsic}.
\end{proof}

Let \({\cal P}\) be a Banach space of canonical--matter parameters and let
\(q\mapsto K(q)\in L^{2p}\) be \(C^1\).  Assume the renormalized matter energy
density has a declared map
\begin{equation}
 (\varphi,q)\longmapsto\rho(\varphi,q)\in L^p(M)         \label{eq:e8v79-matter-map}
\end{equation}
which is \(C^1\) near \((0,q_0)\).  The metric dependence of spin connections,
improvement terms, and counterterms must be included in this map.

Define the nonlinear Hamiltonian-constraint map
\begin{align}
 {\cal H}(\varphi,q)
 :={}&e^{-2\varphi}
 \left(R_g-4\Delta_g\varphi-2|\nabla\varphi|_g^2\right) \notag\\
 &+e^{-4\varphi}
 \left(\tau(q)^2-|K(q)|_g^2\right)
 -2\kappa_G\rho(\varphi,q).                             \label{eq:e8v79-H}
\end{align}
Here \(\tau(q)=\operatorname{tr}_gK(q)\); the base metric \(g\) is fixed in
this chart and all conformal metric dependence is displayed explicitly.

\begin{theorem}[Differentiability of the strong constraint chart]
\label{thm:e8v79-C1}
Under the preceding hypotheses,
\begin{equation}
 {\cal H}:W^{2,p}(M)\times{\cal P}\longrightarrow L^p(M)
 \label{eq:e8v79-map}
\end{equation}
is \(C^1\).  If \({\cal H}(0,q_0)=0\), its conformal derivative is
\begin{equation}
 L\phi
 =-4\Delta_g\phi
 +\left[-2R_g+4(|K|_g^2-\tau^2)-2\kappa_Gm_\rho\right]\phi,
 \label{eq:e8v79-L}
\end{equation}
provided
\begin{equation}
 D_\varphi\rho(0,q_0)[\phi]=m_\rho\phi.                 \label{eq:e8v79-mrho}
\end{equation}
The parameter derivative
\begin{equation}
 F_{\dot q}:=D_q{\cal H}(0,q_0)[\dot q]                 \label{eq:e8v79-Fq}
\end{equation}
is the V72 physical source when \(\dot q=(\dot K,\dot\Psi)\) and the matter
map uses the V72 canonical convention.
\end{theorem}

\begin{proof}
For \(p>3\), \(W^{2,p}\) is a Banach algebra embedded in \(C^1\).
Composition with the exponential is smooth on bounded sets,
\(\Delta_g:W^{2,p}\to L^p\) is bounded, and
\(|\nabla\varphi|^2\in W^{1,p}\subset L^p\).  The products in
\eqref{eq:e8v79-H} are therefore \(C^1\); the parameter terms are \(C^1\) by
hypothesis.  This proves \eqref{eq:e8v79-map}.

Differentiate the first line of \eqref{eq:e8v79-H} at \(\varphi=0\).  The
quadratic gradient term has zero derivative and the result is
\(-4\Delta_g\phi-2R_g\phi\).  Differentiating the extrinsic factor gives
\(-4(\tau^2-|K|^2)\phi=4(|K|^2-\tau^2)\phi\).  The matter derivative is
\(-2\kappa_Gm_\rho\phi\).  Their sum is \eqref{eq:e8v79-L}, which agrees with
V72 at \(d=3\).  Differentiation in \(q\) gives precisely the independent
canonical--matter source ledger.
\end{proof}

\begin{firewall}[Strong chart versus critical weak closure]
\label{fire:e8v79-regularity}
The \(W^{2,p}\), \(p>3\), hypothesis makes the nonlinear map classically
\(C^1\) into \(L^p\).  V76 proves that the curvature form itself extends to
the much weaker critical class \(W^{1,3}\cap L^\infty\), but V79 does not
prove that the nonlinear constraint map is differentiable or locally
invertible at that endpoint.  A strong local chart is not a tightness theorem
for the continuum measure.
\end{firewall}

\subsection{The augmented linear operator}
\label{sec:e8v79-linear}

Let
\begin{equation}
 b(\phi,v):=\int_M(L\phi)v\,d\mu_g                      \label{eq:e8v79-b}
\end{equation}
in its weak symmetric realization, and set
\begin{equation}
 H_{\rm vol}:=left\{w\in H^1(M):\int_Mw\,d\mu_g=0\right\}.
 \label{eq:e8v79-Hvol}
\end{equation}
Assume
\begin{equation}
 b(w,w)\geq\gamma\|w\|_{H^1}^2
 \qquad(w\in H_{\rm vol})                               \label{eq:e8v79-coercivity}
\end{equation}
for some \(\gamma>0\).

\begin{theorem}[Isomorphism of the strong augmented linearization]
\label{thm:e8v79-augmented-isomorphism}
Assume the coefficients of \(L\) have the regularity in
\eqref{eq:e8v79-L} and the standard elliptic estimate
\begin{equation}
 \|u\|_{W^{2,p}}
 \leq C\bigl(\|Lu\|_{L^p}+\|u\|_{L^p}\bigr).           \label{eq:e8v79-elliptic-estimate}
\end{equation}
Then
\begin{align}
 {\cal S}:W^{2,p}(M)\times\mathbb R
 &\longrightarrow L^p(M)\times\mathbb R,               \notag\\
 {\cal S}(\phi,\lambda)
 &: =\left(L\phi+\lambda,
     3\int_M\phi\,d\mu_g\right)                       \label{eq:e8v79-S}
\end{align}
is a bounded isomorphism.
\end{theorem}

\begin{proof}
Fix \((f,\eta)\in L^p\times\mathbb R\) and set
\[
 c:={\eta\over3\operatorname{Vol}_g(M)},
 \qquad \phi=c+w,quad w\in H_{\rm vol}.
\]
By Lax--Milgram and \eqref{eq:e8v79-coercivity}, there is a unique
\(w\in H_{\rm vol}\) satisfying
\begin{equation}
 b(w,v)=\int_M(f-Lc)v\,d\mu_g
 \qquad(v\in H_{\rm vol}).                              \label{eq:e8v79-wsolve}
\end{equation}
Define
\begin{equation}
 \lambda:={1\over\operatorname{Vol}_g(M)}
 \int_M(f-L\phi)\,d\mu_g.                              \label{eq:e8v79-lambda}
\end{equation}
Equation \eqref{eq:e8v79-wsolve} says that
\(L\phi+\lambda-f\) annihilates \(H_{\rm vol}\); the definition of
\(\lambda\) says that it also annihilates constants.  Hence
\(L\phi+\lambda=f\) weakly on all of \(H^1\).

The right side is in \(L^p\), so
\eqref{eq:e8v79-elliptic-estimate} and the coercive bound upgrade
\(w\) to \(W^{2,p}\) and give a bounded solution map.  If
\({\cal S}(\phi,\lambda)=0\), then \(\phi\in H_{\rm vol}\); testing the
first equation on \(\phi\) gives \(b(\phi,\phi)=0\), hence \(\phi=0\).
The first equation then gives \(\lambda=0\).  Thus the map is bijective with
bounded inverse.
\end{proof}

\begin{remark}[Why the full Schur complement is absent]
\label{rem:e8v79-Schur}
The second component of \({\cal S}\) fixes the constant coefficient of
\(\phi\), while \(\lambda\) supplies the missing constant equation.  Only
coercivity on \(H_{\rm vol}\) is required.  This is the nonlinear-chart form
of the V78 saddle theorem, and it remains valid when the unconstrained Schur
complement is zero.
\end{remark}

\subsection{Local nonlinear existence}
\label{sec:e8v79-IFT}

Fix the target volume \(V_0=\operatorname{Vol}_g(M)\) and define
\begin{align}
 {\cal F}:W^{2,p}(M)\times\mathbb R\times{\cal P}
 &\longrightarrow L^p(M)\times\mathbb R,                \notag\\
 {\cal F}(\varphi,\lambda,q)
 &: =\left(
 {\cal H}(\varphi,q)+\lambda,
 \operatorname{Vol}_{g_\varphi}(M)-V_0
 \right).                                               \label{eq:e8v79-F}
\end{align}

\begin{theorem}[Nonlinear augmented constraint chart]
\label{thm:e8v79-IFT}
Suppose
\begin{equation}
 {\cal H}(0,q_0)=0                                      \label{eq:e8v79-base-constraint}
\end{equation}
and assume \eqref{eq:e8v79-coercivity}--
\eqref{eq:e8v79-elliptic-estimate}.  Then there are neighborhoods
\(U\subset{\cal P}\) of \(q_0\) and
\({\cal O}\subset W^{2,p}\times\mathbb R\) of \((0,0)\), and unique
\(C^1\) maps
\begin{equation}
 q\longmapsto\bigl(\varphi(q),\lambda(q)\bigr)\in{\cal O}
 \label{eq:e8v79-solution-map}
\end{equation}
such that
\begin{align}
 {\cal H}(\varphi(q),q)+\lambda(q)&=0,                  \label{eq:e8v79-augmented-constraint}\\
 \operatorname{Vol}_{e^{2\varphi(q)}g}(M)&=V_0.         \label{eq:e8v79-fixed-volume}
\end{align}
For a given \(q\in U\), the original Hamiltonian constraint is solved by
this chart exactly when
\begin{equation}
 \lambda(q)=0.                                          \label{eq:e8v79-physical-locus}
\end{equation}
\end{theorem}

\begin{proof}
The first component is \(C^1\) by \Cref{thm:e8v79-C1}.  Since
\(d\mu_{g_\varphi}=e^{3\varphi}d\mu_g\), the volume map is smooth and
\begin{equation}
 D_\varphi\operatorname{Vol}_{g_\varphi}(M)\big|_0[\phi]
 =3\int_M\phi\,d\mu_g.                                 \label{eq:e8v79-volume-derivative}
\end{equation}
Consequently
\[
 D_{(\varphi,\lambda)}{\cal F}(0,0,q_0)={\cal S},
\]
which is an isomorphism by
\Cref{thm:e8v79-augmented-isomorphism}.  The Banach-space implicit-function
theorem gives the neighborhoods and the unique \(C^1\) solution map.
Equation \eqref{eq:e8v79-augmented-constraint} reduces to the original
constraint precisely on \eqref{eq:e8v79-physical-locus}.
\end{proof}

\begin{corollary}[Tangent equation and the V78 multiplier]
\label{cor:e8v79-tangent}
For \(\dot q\in{\cal P}\), let
\begin{equation}
 \dot\varphi:=D\varphi(q_0)[\dot q],
 \qquad
 \dot\lambda:=D\lambda(q_0)[\dot q].                  \label{eq:e8v79-dots}
\end{equation}
Then
\begin{align}
 L\dot\varphi+\dot\lambda+F_{\dot q}&=0,               \label{eq:e8v79-linear-tangent}\\
 \int_M\dot\varphi\,d\mu_g&=0.                        \label{eq:e8v79-linear-volume}
\end{align}
Equivalently,
\begin{equation}
 \dot\lambda
 =-{F_{\dot q}(1)+b(\dot\varphi,1)
       \over\operatorname{Vol}_g(M)},                  \label{eq:e8v79-dotlambda}
\end{equation}
which is exactly the V78 fixed-volume multiplier formula.
\end{corollary}

\begin{proof}
Differentiate \eqref{eq:e8v79-augmented-constraint} and
\eqref{eq:e8v79-fixed-volume} at \(q_0\), using
\Cref{thm:e8v79-C1} and \eqref{eq:e8v79-volume-derivative}.  Integrate the
first equation against the constant test to obtain
\eqref{eq:e8v79-dotlambda}.
\end{proof}

\begin{corollary}[First-order physical compatibility]
\label{cor:e8v79-first-order-physical}
If the unconstrained linearized operator has \(s=0\), then
\begin{equation}
 \dot\lambda=0
 \quad\Longleftrightarrow\quad
 F_{\dot q}(h)=0.                                       \label{eq:e8v79-first-order-compatibility}
\end{equation}
Thus the V78 resonant condition is exactly the tangent condition for the
augmented solution curve to remain on the physical locus
\(\lambda=0\).
\end{corollary}

\begin{proof}
If \(\dot\lambda=0\), \eqref{eq:e8v79-linear-tangent} is the unmodified
linearized equation, whose solvability is equivalent to
\(F_{\dot q}(h)=0\) by V78.  Conversely, that compatibility supplies a
fixed-volume solution of the unmodified equation; uniqueness of the augmented
saddle solution forces its multiplier to vanish.
\end{proof}

\subsection{Cosmological-parameter interpretation}
\label{sec:e8v79-cosmological}

\begin{proposition}[Declared cosmological absorption]
\label{prop:e8v79-cosmological}
Suppose the physical constraint is enlarged by a variable cosmological
parameter in the convention
\begin{equation}
 {\cal H}_\Lambda(\varphi,q,\Lambda)
 :={\cal H}(\varphi,q)-2\Lambda.                        \label{eq:e8v79-HLambda}
\end{equation}
Then every augmented solution in \Cref{thm:e8v79-IFT} is a solution of the
enlarged constraint with
\begin{equation}
 \Lambda(q)=-\frac12\lambda(q).                         \label{eq:e8v79-Lambda}
\end{equation}
\end{proposition}

\begin{proof}
Substitution gives
\[
 {\cal H}_{\Lambda(q)}
 ={\cal H}+\lambda=0.
\]
\end{proof}

\begin{firewall}[A parameter is not a counterterm unless declared]
\label{fire:e8v79-cosmological}
Equation \eqref{eq:e8v79-Lambda} is an algebraic interpretation in the
enlarged family \eqref{eq:e8v79-HLambda}.  If the cosmological constant is
fixed by the physical theory, the augmented chart solves the desired
constraint only on \(\lambda=0\).  Allowing \(\Lambda\) to depend on the
regulator is a renormalization choice whose convergence and observable
normalization must be proved separately.
\end{firewall}

\subsection{What remains outside the chart}
\label{sec:e8v79-boundaries}

\begin{firewall}[The momentum constraint remains coupled]
\label{fire:e8v79-momentum}
V79 varies only the conformal scalar and treats \(K(q)\) as supplied canonical
data.  It does not solve
\(\nabla^j(K_{ij}-\tau g_{ij})=\kappa_GJ_i\), construct the transverse vector
potential used in the conformal method, or prove that the matter current obeys
the required compatibility.  A full constraint chart must augment
\eqref{eq:e8v79-F} by that vector equation and remove conformal Killing
kernels.
\end{firewall}

\begin{firewall}[Local existence is not measure construction]
\label{fire:e8v79-measure}
The implicit-function theorem gives a deterministic local graph over a strong
parameter neighborhood.  It does not show that regulated random data remain
in this neighborhood with high probability, that its radius is cutoff
uniform, or that the induced Jacobian and counterterms define a tight positive
measure.
\end{firewall}

\begin{assessment}[Progress at V79]
\label{ass:e8v79-status}
Under a strong \(W^{2,p}\), \(p>3\), chart and mean-zero coercivity, the
nonlinear Hamiltonian constraint plus fixed volume now has a genuine local
\(C^1\) solution graph.  Its derivative is the V78 saddle system, and the
resonant compatibility condition is precisely the condition that the graph be
tangent to \(\lambda=0\).  The remaining constraint-level tasks are the
vector momentum equation, conformal-Killing kernels, and a cutoff-uniform
chart radius.
\end{assessment}

\begin{programme}[Eighth-series V80: compulsory companion audit and coupled
constraint target]
\label{prog:e8v79-V80}
V80 will perform the compulsory YM/NS audit, then test whether the newest
companion results supply a type-correct kernel removal, connected local chart,
or cutoff-uniform radius mechanism.  It will use that audit to select the
coupled scalar--vector constraint operator for V81 and will continue beyond
the audit before the next pause.
\end{programme}

\begin{firewall}[Claim boundary after eighth-series V79]
\label{fire:e8v79-claim}
V79 proves the strong nonlinear conformal formula, differentiability of the
constraint map, invertibility of its augmented linearization, the local
fixed-volume implicit-function chart, and the cosmological-parameter
interpretation.  It does not prove a cutoff-uniform chart, the momentum
constraint, physical vanishing of \(\lambda\), critical-endpoint nonlinear
solvability, a coupled continuum measure, or the full SM--Einstein continuum.
\end{firewall}

\begin{theorem}[Eighth-series V79 result]
\label{thm:e8v79-result}
At a strong three-dimensional constraint background, mean-zero coercivity is
sufficient for a local nonlinear fixed-volume solution graph
\(q\mapsto(\varphi(q),\lambda(q))\), without a nonzero full-space Schur
complement.  The graph solves the original Hamiltonian constraint exactly on
\(\lambda=0\), whose tangent condition at a degenerate background is the V78
identity \(F_{\dot q}(h)=0\).
\end{theorem}

\begin{proof}
Use \Cref{thm:e8v79-C1,thm:e8v79-augmented-isomorphism,
thm:e8v79-IFT,cor:e8v79-first-order-physical}.
\end{proof}

\paragraph{Parent-version record.}
V79 was formed from
\texttt{Renewal\_SM\_Einstein\_Continuum\_Research\_Notes\_V78.tex}.
The V78 parent had (31\,053) logical source lines, (1\,095\,795) bytes,
and SHA--256
\[
 \texttt{F3F6F4AA43C7FCEEF9EA5C96DA5F084755F5976E1229D61A49FECFDB5B233581}.
\]
The next version is the compulsory V80 companion audit.

% =============================================================================
% END APPENDED EIGHTH-SERIES V79 MATERIAL
% =============================================================================
% =============================================================================
% BEGIN APPENDED EIGHTH-SERIES V80 MATERIAL
% =============================================================================

\clearpage
\section{Eighth-series V80: companion audit and the momentum-projector
no-go}
\label{sec:v8v80}

This compulsory audit follows V75.  The Yang--Mills programme has advanced
from V27 to V34; the Navier--Stokes active note remains V26.  The relevant YM
advance is a sequence of exact free renormalization tests: sharp low/high
splitting contracts discarded ultraviolet modes but preserves the massless
retained sector; exact finite-range scalar projections do not exist; local
lifting filters preserve the cochain complex but acquire mixing after
orthogonalization; and analytic buffered filters trade polynomially small
leakage for logarithmic localization plus an explicit transition channel.

For the SM--Einstein constraint problem, this rules out one tempting next
move: the momentum constraint may not be reduced by postulating an exact
finite-range Helmholtz projector.  This note proves the corresponding
vector-projector no-go directly.  V81 will instead use the global elliptic
quotient by conformal Killing fields.

\subsection{Files inspected}
\label{sec:e8v80-files}

The companion checkpoints are
\begin{align}
 &\texttt{Renewal\_Yang\_Mills\_Mass\_Gap\_Research\_Notes\_V34.tex},
 \notag\\
 &\qquad 13\,591\ \text{logical lines},\quad
 494\,269\ \text{bytes},\quad
 \operatorname{SHA256}
 =\texttt{A6762C08333B443F1C93A44F33CC64E709A59ABF613B2E01F55E0EAFEC12261B},
 \label{eq:e8v80-YM-file}\\
 &\texttt{Renewal\_NS\_Stability\_Research\_Notes\_V26.tex},
 \notag\\
 &\qquad 5\,319\ \text{logical lines},\quad
 278\,283\ \text{bytes},\quad
 \operatorname{SHA256}
 =\texttt{82C887F8C2C69E4F28FAD7C3DC8DE5B9099CB5A4BF431EBA1FFF3E91935978AD}.
 \label{eq:e8v80-NS-file}
\end{align}
The NS file is byte-for-byte unchanged from V75.

\subsection{Yang--Mills V28--V34}
\label{sec:e8v80-YM}

\begin{audit}[YM V28: exact blocking requires a full effective action]
\label{audit:e8v80-YM28}
Gauge-covariant block pushforward preserves coarse reflected correlators
exactly.  Its effective action is generally not a one-coupling Wilson action:
higher representations, larger loops, multi-loop terms, and nonlocal
interactions are generated.  A fine observable splits into its conditional
coarse part and a fine complement; exact matching on the coarse subspace does
not control the three covariance terms involving that complement.  The note
therefore states a transfer recursion with separate coarse, residual, and
off-diagonal errors.
\end{audit}

\begin{audit}[YM V29: exact free split contracts only the discarded sector]
\label{audit:e8v80-YM29}
Sharp Fourier restriction gives an exact Gaussian Fock-space tensor
factorization, zero coarse-matching and off-diagonal errors, and a
scale-independent contraction of every state carrying a high-mode quantum.
The retained lowest Fourier mode remains massless in growing volume.  Thus
perfect ultraviolet complement control does not create the terminal
interacting mass.
\end{audit}

\begin{audit}[YM V31--V32: local chain maps versus orthogonal projectors]
\label{audit:e8v80-YM31-32}
A translation-invariant finite-range scalar convolution that is an orthogonal
projection on arbitrarily large tori is either zero or the identity.  Smooth
filters can be truncated with controlled error and still commute with the
lattice de Rham differential.  A two-channel lifting construction gives
finite-range perfect reconstruction and exact derivative intertwining, but
orthogonalization introduces exponentially quasi-local factors and a
nonzero coarse--detail derivative mixing.
\end{audit}

\begin{audit}[YM V33--V34: fixed errors fail the ultraviolet margin; buffers
repair only separated bands]
\label{audit:e8v80-YM33-34}
The exact Haar-block calculation gives uniform high-mode contraction but a
fixed off-diagonal transfer error, which is too large for a physical margin
of order one lattice spacing.  Analytic logistic spectral flattening reduces
low/high leakage to \(O(a^p)\) with exponential localization length
\(O(\log(1/a))\); literal finite range at the same accuracy costs
\(O((\log(1/a))^2)\).  A populated band edge cannot be approximated in global
operator norm below one half by a continuous filter.  The transition band
must remain an explicit third channel, though its free transfer is uniformly
contracting.
\end{audit}

\subsection{Navier--Stokes V26}
\label{sec:e8v80-NS}

\begin{audit}[No new NS input]
\label{audit:e8v80-NS26}
The NS file and its rank-one Oseen-kernel norms are unchanged.  They do not
supply a conformal-Killing quotient, a momentum-constraint inverse, or a
nonlinear Einstein chart radius.
\end{audit}

\subsection{Type audit}
\label{sec:e8v80-types}

\begin{theorem}[Companion non-transfer at V80]
\label{thm:e8v80-no-transfer}
The following transfers from the audited companion work are invalid:
\begin{enumerate}[label=\textup{(\roman*)},leftmargin=3.8em]
\item using contraction of the discarded free YM sector as coercivity of the
retained Einstein constraint modes;
\item treating an exact coarse pushforward as closure of the SM--Einstein
action in a finite list of couplings;
\item importing a scalar lattice filter as an exact transverse or
conformal-Killing projector for the continuum momentum constraint;
\item deleting the transition band of a quasi-local split because its free
transfer contracts; or
\item applying the unchanged NS kernel constants to the elliptic vector
constraint.
\end{enumerate}
\end{theorem}

\begin{proof}
Item (i) confuses the eliminated and retained spectral sectors; YM V29 proves
the retained free sector stays gapless.  Item (ii) contradicts the exact
effective-action ledger of YM V28.  Item (iii) changes both operator and field
type, and the direct no-go below rules out the intended exact local projector.
Item (iv) contradicts the global one-half obstruction in YM V34.  Item (v)
has unrelated equations and state spaces.
\end{proof}

\begin{corollary}[Counterterm closure is an unpaid row of V79]
\label{cor:e8v80-counterterms}
The local solution graph of V79 is constructed for a declared nonlinear
constraint map.  Under coarse graining, closure of that map in the chosen
finite family \((g,K,\rho,m_\rho,\Lambda)\) is not automatic.  All generated
diffeomorphism-, gauge-, reflection-, and locality-compatible interactions
must be bounded or added to the parameter space before a cofinal chart radius
can be claimed.
\end{corollary}

\begin{proof}
The statement is the type-correct implication of the YM V28 effective-action
ledger: exact pushforward preserves observables but does not preserve a chosen
finite action ansatz.  V79 contains no theorem eliminating the additional
coordinates.
\end{proof}

\subsection{Direct no-go for an exact local Helmholtz projector}
\label{sec:e8v80-projector}

Consider complex one-cochains on \(\mathbb Z^d\), \(d\ge2\), with forward
difference symbol
\begin{equation}
 q(k)=\bigl(e^{ik_1}-1,\ldots,e^{ik_d}-1\bigr)^{\mathsf T}.
 \label{eq:e8v80-q}
\end{equation}
For \(k\ne0\) modulo \(2\pi\), the orthogonal projection onto exact
one-cochains is
\begin{equation}
 P_{\rm ex}(k):={q(k)q(k)^*\over q(k)^*q(k)},            \label{eq:e8v80-Pex}
\end{equation}
and the transverse projection is \(P_{\rm tr}(k)=I-P_{\rm ex}(k)\).

\begin{theorem}[No translation-invariant finite-range exact Helmholtz
projection]
\label{thm:e8v80-Helmholtz-no-go}
There is no translation-invariant finite-range matrix convolution on
one-cochains whose multiplier equals \(P_{\rm ex}(k)\) for every nonzero
momentum on periodic lattices of arbitrarily large period.  The same is true
for \(P_{\rm tr}(k)\), regardless of the value assigned at the harmonic mode
\(k=0\).
\end{theorem}

\begin{proof}
A finite-range matrix convolution has a matrix-valued trigonometric-polynomial
multiplier \(P(k)\), hence a continuous limit on the Brillouin torus.  Periodic
momenta from arbitrarily large tori are dense, so equality at all their
nonzero momenta forces equality with \eqref{eq:e8v80-Pex} at every nonzero
\(k\) by continuity.

Approach zero along the first coordinate axis.  Since
\(e^{it}-1=it+O(t^2)\),
\begin{equation}
 \lim_{t\to0}P_{\rm ex}(t,0,\ldots,0)
 =\operatorname{diag}(1,0,\ldots,0).                   \label{eq:e8v80-axis1}
\end{equation}
Approach along the second coordinate axis instead:
\begin{equation}
 \lim_{t\to0}P_{\rm ex}(0,t,0,\ldots,0)
 =\operatorname{diag}(0,1,0,\ldots,0).                 \label{eq:e8v80-axis2}
\end{equation}
These limits differ, so \(P_{\rm ex}\) has no continuous extension at zero,
contradicting continuity of \(P(k)\).  Since
\(P_{\rm tr}=I-P_{\rm ex}\), the transverse case follows.
\end{proof}

\begin{corollary}[Harmonic assignment does not repair locality]
\label{cor:e8v80-harmonic}
Choosing the zero-momentum projector to be zero, the identity, or any fixed
matrix does not alter the obstruction: the directional limits
\eqref{eq:e8v80-axis1}--\eqref{eq:e8v80-axis2} remain incompatible.
\end{corollary}

\begin{proof}
Changing one value cannot create a common limit for a function with two
distinct directional limits.
\end{proof}

\begin{remark}[Continuum symbol]
\label{rem:e8v80-continuum}
The continuum longitudinal symbol \(kk^{\mathsf T}/|k|^2\) has the same
directional discontinuity at \(k=0\).  Its inverse Fourier transform is
long-ranged.  Thus the obstruction is not an artefact of forward differences.
\end{remark}

\begin{firewall}[Scope of the projector no-go]
\label{fire:e8v80-projector-scope}
\Cref{thm:e8v80-Helmholtz-no-go} excludes an exact, translation-invariant,
finite-range orthogonal Helmholtz projection uniformly across volumes.  It
does not exclude a global elliptic projector, a local nonorthogonal lifting
with several channels, an exponentially quasi-local buffered filter, or a
finite-volume gauge choice.  Each alternative must retain harmonic or
transition modes explicitly.
\end{firewall}

\subsection{Audit-guided constraint direction}
\label{sec:e8v80-direction}

\begin{decision}[Use the elliptic quotient before any local approximation]
\label{dec:e8v80-direction}
V81 will formulate the momentum constraint with the conformal Killing
operator
\[
 ({\cal L}_gX)_{ij}
 =\nabla_iX_j+\nabla_jX_i-\frac23(\nabla_kX^k)g_{ij}.
\]
It will prove positivity of the vector form, identify its kernel as the
finite-dimensional conformal Killing space, impose the exact source
compatibility on that kernel, and derive a coupled scalar--vector coercivity
criterion on the quotient.  No finite-range transverse projector will be
introduced.  If a later regulator requires localization, the YM V34 lesson
requires an explicit spectral buffer and transition sector.
\end{decision}

\begin{assessment}[Progress at V80]
\label{ass:e8v80-status}
The audit prevents two hidden closures: coarse graining need not preserve the
finite V79 action family, and an exact local Helmholtz projection does not
exist.  The new no-go theorem proves the latter directly from the directional
singularity at zero momentum.  The next rigorous constraint step is therefore
the global Fredholm quotient by conformal Killing fields, followed by a
coupled block estimate.
\end{assessment}

\begin{programme}[Eighth-series V81: momentum quotient and coupled block]
\label{prog:e8v80-V81}
V81 will execute \Cref{dec:e8v80-direction}, including the conformal-Killing
compatibility condition, Korn coercivity on the orthogonal complement, and a
quantitative small-coupling criterion for the scalar--vector constraint
linearization.  This is the required continuation beyond the V80 audit.
\end{programme}

\begin{firewall}[Claim boundary after eighth-series V80]
\label{fire:e8v80-boundary}
V80 records the current YM V34 and NS V26 audit, proves the type non-transfer,
identifies the generated-interaction row, and proves the exact local Helmholtz
projector no-go.  It does not solve the momentum constraint, prove Korn
coercivity, control the nonlinear coupled chart, construct the cofinal
measure, or establish the full SM--Einstein continuum.
\end{firewall}

\begin{theorem}[Eighth-series V80 result]
\label{thm:e8v80-result}
The companion audit and the direct Fourier theorem rule out an exact
finite-range Helmholtz reduction of the Einstein momentum constraint.  The
type-correct route is a global elliptic quotient retaining conformal-Killing
kernels, with any later quasi-local approximation carrying an explicit buffer
and transition channel.
\end{theorem}

\begin{proof}
Use \Cref{audit:e8v80-YM28,audit:e8v80-YM29,
audit:e8v80-YM31-32,audit:e8v80-YM33-34,
thm:e8v80-no-transfer,thm:e8v80-Helmholtz-no-go,
dec:e8v80-direction}.
\end{proof}

\paragraph{Parent-version record.}
V80 was formed from
\texttt{Renewal\_SM\_Einstein\_Continuum\_Research\_Notes\_V79.tex}.
The V79 parent had (31\,495) logical source lines, (1\,113\,043) bytes,
and SHA--256
\[
 \texttt{3EEE24069FC5CA08FE2CF21DEE74707D2EE557F063EA4BACDD400AF4E9636FA4}.
\]
The next compulsory companion audit is V85.

% =============================================================================
% END APPENDED EIGHTH-SERIES V80 MATERIAL
% =============================================================================
% =============================================================================
% BEGIN APPENDED EIGHTH-SERIES V81 MATERIAL
% =============================================================================

\clearpage
\section{Eighth-series V81: the momentum quotient and the coupled
scalar--vector constraint block}
\label{sec:v8v81}

V80 rules out an exact finite-range Helmholtz projector.  The correct first
construction is global and elliptic.  This note develops it: the conformal
Killing operator produces a positive vector form, its kernel is retained
exactly, and the momentum source must annihilate that kernel.  The scalar
fixed-volume space of V79 and the vector conformal-Killing quotient are then
combined.  A quantitative cross-block condition gives coercivity and a unique
augmented scalar--vector solution with explicit volume and kernel multipliers.

\subsection{The conformal Killing vector form}
\label{sec:e8v81-vector}

Let \((M,g)\) be a connected closed smooth three-manifold.  For
\(X\in H^1(TM)\), define
\begin{equation}
 ({\cal L}_gX)_{ij}
 :=\nabla_iX_j+\nabla_jX_i
 -{2\over3}(\nabla_kX^k)g_{ij}.                         \label{eq:e8v81-L}
\end{equation}
This tensor is symmetric and trace free.  Define
\begin{equation}
 a_V(X,Y):={1\over2}
 \int_M\langle{\cal L}_gX,{\cal L}_gY\rangle_g\,d\mu_g.
 \label{eq:e8v81-aV}
\end{equation}

\begin{lemma}[Weak vector Laplacian identity]
\label{lem:e8v81-vector-identity}
For smooth vector fields,
\begin{equation}
 a_V(X,Y)=\int_M\langle P_gX,Y\rangle_g\,d\mu_g,
 \qquad
 (P_gX)_i:=-\nabla^j({\cal L}_gX)_{ij}.                 \label{eq:e8v81-P}
\end{equation}
Thus \(P_g\) is nonnegative and formally self-adjoint.
\end{lemma}

\begin{proof}
Because \({\cal L}_gX\) is symmetric and trace free,
\[
 \frac12({\cal L}_gX)^{ij}({\cal L}_gY)_{ij}
 =({\cal L}_gX)^{ij}\nabla_iY_j.
\]
Integrate by parts on the closed manifold.  The result is
\eqref{eq:e8v81-P}; symmetry is already visible in
\eqref{eq:e8v81-aV}.
\end{proof}

Define the conformal Killing space
\begin{equation}
 {\cal K}_g:=\ker{\cal L}_g
 =\{X\in H^1(TM):{\cal L}_gX=0\}.                       \label{eq:e8v81-K}
\end{equation}

\begin{theorem}[Kernel and quotient coercivity]
\label{thm:e8v81-Korn}
The space \({\cal K}_g\) is finite dimensional and consists of smooth vector
fields.  Let
\begin{equation}
 H_V:=\left\{X\in H^1(TM):(X,Y)_{L^2_g}=0
              \text{ for every }Y\in{\cal K}_g\right\}. 
 \label{eq:e8v81-HV}
\end{equation}
There is \(C_K<\infty\) such that
\begin{equation}
 \|X\|_{H^1}\leq C_K\|{\cal L}_gX\|_2
 \qquad(X\in H_V).                                     \label{eq:e8v81-Korn}
\end{equation}
Consequently
\begin{equation}
 a_V(X,X)\geq\gamma_V\|X\|_{H^1}^2,
 \qquad \gamma_V:={1\over2C_K^2}.                      \label{eq:e8v81-vector-coercivity}
\end{equation}
\end{theorem}

\begin{proof}
The principal symbol of \({\cal L}_g\) is injective: for a nonzero covector
\(\xi\), the equation
\[
 \xi_iX_j+\xi_jX_i-\frac23(\xi\cdot X)g_{ij}=0
\]
forces \(X=0\) after contracting once with \(\xi^i\) and separating the
component parallel to \(\xi\).  Hence \(P_g={\cal L}_g^*{\cal L}_g/2\) is
elliptic.  Elliptic regularity makes its kernel finite dimensional and smooth,
and the kernel equals \({\cal K}_g\) because
\(a_V(X,X)=\|{\cal L}_gX\|_2^2/2\).

The conformal Korn inequality with an \(L^2\) remainder is
\begin{equation}
 \|X\|_{H^1}\leq C
 \bigl(\|{\cal L}_gX\|_2+\|X\|_2\bigr).                \label{eq:e8v81-Korn-remainder}
\end{equation}
It follows by local ellipticity and a finite partition of unity.  If
\eqref{eq:e8v81-Korn} failed, there would be \(X_n\in H_V\) with
\(\|X_n\|_{H^1}=1\) and \(\|{\cal L}_gX_n\|_2\to0\).  Rellich compactness
gives a subsequence converging strongly in \(L^2\) and weakly in \(H^1\) to
some \(X\).  Then \({\cal L}_gX=0\) and \(X\perp{\cal K}_g\), so \(X=0\).
Equation \eqref{eq:e8v81-Korn-remainder} applied to \(X_n\) now forces
\(\|X_n\|_{H^1}\to0\), a contradiction.  Squaring
\eqref{eq:e8v81-Korn} in \eqref{eq:e8v81-aV} proves
\eqref{eq:e8v81-vector-coercivity}.
\end{proof}

\begin{firewall}[The quotient is global]
\label{fire:e8v81-global}
The projection onto \(H_V\) is the global \(L^2\) projection away from
\({\cal K}_g\), not a finite-range transverse projector.  V80 proves that an
exact translation-invariant local Helmholtz projection is unavailable even
in the flat lattice model.  No locality claim is attached to
\eqref{eq:e8v81-HV}.
\end{firewall}

\subsection{Exact momentum compatibility}
\label{sec:e8v81-momentum}

Use the V72 sign convention for the momentum residual
\begin{equation}
 ({\cal M}_{g,K,J})_i
 :=\nabla^j(K_{ij}-\tau g_{ij})-\kappa_GJ_i,
 \qquad \tau=\operatorname{tr}_gK.                     \label{eq:e8v81-M}
\end{equation}
It defines a functional on \(H^1(TM)\) whenever
\(K\in L^2\), \(\tau\in L^2\), and \(J\in H^{-1}(T^*M)\), by weak
integration by parts.

Correct \(K\) by a longitudinal trace-free tensor:
\begin{equation}
 K_X:=K+{\cal L}_gX.                                    \label{eq:e8v81-KX}
\end{equation}
Since \(\operatorname{tr}_g{\cal L}_gX=0\), the trace \(\tau\) is unchanged,
and
\begin{equation}
 {\cal M}_{g,K_X,J}
 ={\cal M}_{g,K,J}+\operatorname{div}_g{\cal L}_gX.
 \label{eq:e8v81-MX}
\end{equation}

\begin{theorem}[Momentum Fredholm alternative]
\label{thm:e8v81-momentum-Fredholm}
There is a solution \(X\in H^1(TM)\) of
\begin{equation}
 {\cal M}_{g,K_X,J}=0                                   \label{eq:e8v81-momentum-equation}
\end{equation}
if and only if
\begin{equation}
 {\cal M}_{g,K,J}(Y)=0
 \qquad(Y\in{\cal K}_g).                               \label{eq:e8v81-CKF-compatibility}
\end{equation}
When this condition holds there is a unique solution \(X\in H_V\), and
\begin{equation}
 \|X\|_{H^1}leq\gamma_V^{-1}
 \|{\cal M}_{g,K,J}\|_{H_V^*}.                        \label{eq:e8v81-X-bound}
\end{equation}
All other solutions differ by a conformal Killing field.
\end{theorem}

\begin{proof}
Testing \eqref{eq:e8v81-MX} on a smooth \(Y\) and integrating by parts gives
\[
 {\cal M}_{g,K_X,J}(Y)
 ={\cal M}_{g,K,J}(Y)-a_V(X,Y).
\]
If \(Y\in{\cal K}_g\), then \(a_V(X,Y)=0\), proving necessity of
\eqref{eq:e8v81-CKF-compatibility}.  Conversely, Lax--Milgram on \(H_V\)
with \eqref{eq:e8v81-vector-coercivity} gives a unique \(X\in H_V\) such that
\[
 a_V(X,Y)={\cal M}_{g,K,J}(Y)qquad(Y\in H_V).
\]
The compatibility condition supplies the same identity on
\({\cal K}_g\), hence on all of \(H^1(TM)\).  This is
\eqref{eq:e8v81-momentum-equation}.  The energy estimate gives
\eqref{eq:e8v81-X-bound}, and the kernel statement follows from
\Cref{thm:e8v81-Korn}.
\end{proof}

\begin{theorem}[Explicit conformal-Killing source condition]
\label{thm:e8v81-CKF-source}
For every \(Y\in{\cal K}_g\),
\begin{equation}
 {\cal M}_{g,K,J}(Y)
 =-{2\over3}\int_M\langle\nabla\tau,Y\rangle_g\,d\mu_g
 -\kappa_GJ(Y),                                         \label{eq:e8v81-CKF-source}
\end{equation}
where the first term is interpreted weakly.  Equivalently, for smooth
\(\tau\),
\begin{equation}
 {\cal M}_{g,K,J}(Y)
 ={2\over3}\int_M\tau\,\operatorname{div}_gY\,d\mu_g
 -\kappa_G\int_M\langle J,Y\rangle_g\,d\mu_g.          \label{eq:e8v81-CKF-source2}
\end{equation}
In particular, constant mean curvature and a matter current annihilating
\({\cal K}_g\) are sufficient for compatibility.
\end{theorem}

\begin{proof}
For a conformal Killing field,
\[
 \frac12(\nabla_iY_j+\nabla_jY_i)
 ={1\over3}(\operatorname{div}Y)g_{ij}.
\]
The trace of \(K-\tau g\) in dimension three is \(-2\tau\).  Therefore weak
integration by parts in \eqref{eq:e8v81-M} gives
\[
 -\int_M(K-\tau g):\nabla Y\,d\mu_g
 ={2\over3}\int_M\tau\operatorname{div}Y\,d\mu_g.
\]
The closed-manifold divergence theorem converts this to
\(-\frac23\int\langle\nabla\tau,Y\rangle\), proving both formulas.  The last
claim is immediate.
\end{proof}

\begin{firewall}[The CKF condition is physical data]
\label{fire:e8v81-CKF}
Quotienting the unknown \(X\) by \({\cal K}_g\) removes nonuniqueness; it does
not remove the source condition \eqref{eq:e8v81-CKF-compatibility}.  Nonconstant
mean curvature or matter momentum can excite a conformal Killing direction.
Projecting that source away would change the momentum constraint unless a
separate gauge or Lagrange multiplier is declared.
\end{firewall}

\subsection{The scalar--vector cross blocks}
\label{sec:e8v81-cross}

Let
\begin{equation}
 H_S:=\left\{\phi\in H^1(M):\int_M\phi\,d\mu_g=0\right\}
 \label{eq:e8v81-HS}
\end{equation}
be the V79 fixed-volume scalar space.  Let \(b_S\) be the complete scalar ADM
form and assume
\begin{equation}
 b_S(\phi,\phi)\geq\gamma_S\|\phi\|_{H^1}^2
 \qquad(\phi\in H_S)                                    \label{eq:e8v81-scalar-coercivity}
\end{equation}
for \(\gamma_S>0\).

The Hamiltonian derivative produced by a longitudinal variation
\(\dot K={\cal L}_gX\) is
\begin{equation}
 c_H(X,\psi)
 :=-2\int_M\langle K,{\cal L}_gX\rangle_g\psi\,d\mu_g.
 \label{eq:e8v81-cH}
\end{equation}
For smooth data, the geometric part of the momentum derivative under
\(\dot g=2\phi g\), holding covariant \(K\) fixed, is
\begin{equation}
 (D_g{\cal M}_{\rm geom}[2\phi g])_i
 =-2\phi({\cal M}_{\rm geom})_i
 +(K_i{}^j+\tau\delta_i{}^j)\nabla_j\phi,             \label{eq:e8v81-DMgeom}
\end{equation}
where \({\cal M}_{\rm geom}=\operatorname{div}(K-\tau g)\).
Including the matter-current metric derivative, define
\begin{equation}
 c_M(\phi,Y)
 :=-\int_M\left\langle
 D_g{\cal M}_{\rm geom}[2\phi g]
 -\kappa_GD_gJ[2\phi g],Y\right\rangle_gd\mu_g.         \label{eq:e8v81-cM}
\end{equation}
The minus sign matches the positive vector equation
\(a_V(X,Y)={\cal M}(Y)\).

\begin{lemma}[Conformal derivative of the geometric momentum residual]
\label{lem:e8v81-DM}
Formula \eqref{eq:e8v81-DMgeom} is exact in dimension three.
\end{lemma}

\begin{proof}
For \(\dot g=2\phi g\),
\[
 \dot g^{ij}=-2\phi g^{ij},qquad
 \dot\Gamma^k{}_{ij}
 =\delta^k_i\nabla_j\phi+delta^k_j\nabla_i\phi
  -g_{ij}\nabla^k\phi,qquad
 \dot\tau=-2\phi\tau.
\]
Direct contraction gives, in dimension \(d\),
\[
 D_g(\nabla^jK_{ij})[2\phi g]
 =-2\phi\nabla^jK_{ij}-\tau\nabla_i\phi
 +(d-2)K_i{}^j\nabla_j\phi.
\]
Also
\[
 D_g(\nabla_i\tau)[2\phi g]
 =-2\tau\nabla_i\phi-2\phi\nabla_i\tau.
\]
Subtract, set \(d=3\), and collect
\(-2\phi(\nabla^jK_{ij}-\nabla_i\tau)\).  This is
\eqref{eq:e8v81-DMgeom}.
\end{proof}

Assume the cross forms extend boundedly with
\begin{align}
 |c_H(X,\psi)|&\leq C_H\|X\|_{H^1}\|\psi\|_{H^1},
 \label{eq:e8v81-cH-bound}\\
 |c_M(\phi,Y)|&\leq C_M\|\phi\|_{H^1}\|Y\|_{H^1}.
 \label{eq:e8v81-cM-bound}
\end{align}
For example, \(K,\tau,{\cal M}_{\rm geom}\in L^\infty\) and a bounded matter
current derivative from \(H^1\) to \(H^{-1}\) are sufficient.

\subsection{Coupled quotient coercivity}
\label{sec:e8v81-coupled}

On \({\cal X}:=H_S\times H_V\), define
\begin{align}
 {\mathbb b}((\phi,X),(\psi,Y))
 :={}&b_S(\phi,\psi)+a_V(X,Y)                            \notag\\
 &+c_H(X,\psi)+c_M(\phi,Y).                            \label{eq:e8v81-coupled-form}
\end{align}

\begin{theorem}[Quantitative scalar--vector coercivity]
\label{thm:e8v81-coupled-coercivity}
If
\begin{equation}
 C_H+C_M<2\sqrt{\gamma_S\gamma_V},                     \label{eq:e8v81-cross-small}
\end{equation}
then \({\mathbb b}\) is coercive on \({\cal X}\).  More precisely,
\begin{equation}
 {\mathbb b}((\phi,X),(\phi,X))
 \geq\gamma_*\left(\|\phi\|_{H^1}^2+
                         \|X\|_{H^1}^2\right),         \label{eq:e8v81-coupled-floor}
\end{equation}
where
\begin{equation}
 \gamma_*
 :={\gamma_S+\gamma_V-
 \sqrt{(\gamma_S-\gamma_V)^2+(C_H+C_M)^2}\over2}>0.
 \label{eq:e8v81-gammastar}
\end{equation}
\end{theorem}

\begin{proof}
Put \(x=\|\phi\|_{H^1}\), \(y=\|X\|_{H^1}\), and
\(C=C_H+C_M\).  The diagonal coercivities and cross bounds give
\[
 {\mathbb b}((\phi,X),(\phi,X))
 \geq\gamma_Sx^2+\gamma_Vy^2-Cxy.
\]
The smaller eigenvalue of the symmetric two-by-two matrix
\[
 \begin{pmatrix}\gamma_S&-C/2\\-C/2&\gamma_V\end{pmatrix}
\]
is \eqref{eq:e8v81-gammastar}.  It is positive exactly under
\eqref{eq:e8v81-cross-small}, proving
\eqref{eq:e8v81-coupled-floor}.
\end{proof}

\begin{corollary}[Coupled normalized linear solve]
\label{cor:e8v81-coupled-solve}
Under \eqref{eq:e8v81-cross-small}, for every
\((F_S,F_V)\in H_S^*\times H_V^*\) there is a unique
\((\phi,X)\in H_S\times H_V\) satisfying
\begin{equation}
 {\mathbb b}((\phi,X),(\psi,Y))
 =F_S(\psi)+F_V(Y)
 \quad((\psi,Y)\in{\cal X}),                            \label{eq:e8v81-coupled-solve}
\end{equation}
and
\begin{equation}
 \|\phi\|_{H^1}+\|X\|_{H^1}
 \leq {\sqrt2\over\gamma_*}
 \left(\|F_S\|_{H_S^*}^2+\|F_V\|_{H_V^*}^2\right)^{1/2}.
 \label{eq:e8v81-coupled-bound}
\end{equation}
\end{corollary}

\begin{proof}
Apply Lax--Milgram to the bounded coercive, possibly nonsymmetric, form
\({\mathbb b}\).  Its energy estimate controls the product norm; the displayed
sum is at most \(\sqrt2\) times that norm.
\end{proof}

\begin{remark}[Exact block alternative]
\label{rem:e8v81-Schur}
If the vector quotient operator \(D:H_V\to H_V^*\) induced by \(a_V\) is
inverted first, the exact scalar Schur operator is
\[
 S=A-CD^{-1}E,
\]
where \(A\) is the scalar block, \(E\phi=c_M(\phi,\cdot)\), and
\((CX)(\psi)=c_H(X,\psi)\).  The coupled operator is invertible exactly when
\(S\) is.  Condition \eqref{eq:e8v81-cross-small} is a direct sufficient
coercivity test that avoids assuming symmetry of the two cross blocks.
\end{remark}

\subsection{Augmented full-space equations}
\label{sec:e8v81-augmented}

Choose an \(L^2_g\)-orthonormal basis
\(Y_1,\ldots,Y_N\) of \({\cal K}_g\).  For scalar and vector multipliers
\(\lambda\in\mathbb R\), \(\nu=(\nu_1,\ldots,\nu_N)\in\mathbb R^N\), consider
\begin{align}
 b_S(\phi,\psi)+c_H(X,\psi)
 +\lambda\int_M\psi\,d\mu_g
 &=F_S(\psi),                                           \label{eq:e8v81-aug-scalar}\\
 a_V(X,Y)+c_M(\phi,Y)
 +\sum_{A=1}^N\nu_A(Y_A,Y)_{L^2_g}
 &=F_V(Y),                                              \label{eq:e8v81-aug-vector}\\
 \int_M\phi\,d\mu_g&=0,                               \label{eq:e8v81-aug-volume}\\
 (X,Y_A)_{L^2_g}&=0\quad(1\leq A\leq N).              \label{eq:e8v81-aug-CKF}
\end{align}

\begin{theorem}[Unique augmented scalar--vector solution]
\label{thm:e8v81-augmented}
Under \eqref{eq:e8v81-cross-small}, for every full-space source
\((F_S,F_V)\in H^1(M)^*\times H^1(TM)^*\), equations
\eqref{eq:e8v81-aug-scalar}--\eqref{eq:e8v81-aug-CKF} have a unique solution
\begin{equation}
 (\phi,X,\lambda,\nu)
 \in H_S\times H_V\times\mathbb R\times\mathbb R^N.   \label{eq:e8v81-aug-space}
\end{equation}
The same \((\phi,X)\) solves the unaugmented physical linearized constraints
if and only if
\begin{equation}
 \lambda=0,qquad \nu_1=\cdots=\nu_N=0.                 \label{eq:e8v81-physical-multipliers}
\end{equation}
\end{theorem}

\begin{proof}
Restrict the two equations to tests \((\psi,Y)\in H_S\times H_V\).  The
multiplier terms vanish, so \Cref{cor:e8v81-coupled-solve} gives a unique
normalized pair \((\phi,X)\).  The residual in the scalar equation annihilates
\(H_S\), whose annihilator is the one-dimensional span of
\(\psi\mapsto\int_M\psi\,d\mu_g\); this determines a unique \(\lambda\).
The residual in the vector equation annihilates \(H_V\), so its Riesz
representative lies in \({\cal K}_g\); expansion in the orthonormal basis
determines a unique \(\nu\).  The normalization equations already hold by
construction.  Finally the augmented equations coincide with the physical
ones exactly when every multiplier vanishes.
\end{proof}

\begin{corollary}[Momentum multiplier as compatibility defect]
\label{cor:e8v81-nu}
If the cross term \(c_M(\phi,Y_A)\) vanishes on conformal Killing tests, then
\begin{equation}
 \nu_A=F_V(Y_A).                                        \label{eq:e8v81-nu}
\end{equation}
For the uncoupled physical momentum residual this is exactly
\eqref{eq:e8v81-CKF-source}.  Hence \(\nu=0\) is the conformal-Killing source
condition, not a gauge choice.
\end{corollary}

\begin{proof}
Test \eqref{eq:e8v81-aug-vector} with \(Y=Y_A\).  The vector energy term
vanishes because \({\cal L}_gY_A=0\), the cross term vanishes by hypothesis,
and orthonormality leaves \(\nu_A=F_V(Y_A)\).
\end{proof}

\subsection{Cofinal stability}
\label{sec:e8v81-stability}

\begin{theorem}[Stability of the coupled quotient solve]
\label{thm:e8v81-stability}
Let \({\mathbb b}_n\to{\mathbb b}\) in bilinear-form norm on one identified
product quotient space, suppose the lower bound
\eqref{eq:e8v81-coupled-floor} holds with one \(\gamma_*>0\), and let
\(F_n\to F\) in the product dual.  Then the normalized solutions converge
strongly in \(H^1\times H^1\).  If the volume and conformal-Killing constraint
functionals and their chosen finite-dimensional bases also converge, the
augmented multipliers converge.
\end{theorem}

\begin{proof}
Let \(U_n,U\) denote the normalized solutions.  Subtract their equations and
test with \(U_n-U\).  Coercivity gives
\[
 \gamma_*\|U_n-U\|^2
 \leq\|F_n-F\|\|U_n-U\|
 +\|{\mathbb b}_n-{\mathbb b}\|\|U\|\|U_n-U\|.
\]
Divide by \(\|U_n-U\|\) when nonzero.  This proves strong convergence.  Each
multiplier is a finite collection of residual pairings on the constant and
kernel basis vectors, so convergence of forms, sources, and bases gives its
convergence.
\end{proof}

\begin{firewall}[Moving conformal-Killing dimension]
\label{fire:e8v81-moving-kernel}
The stability theorem assumes one identified quotient and convergence of the
kernel bases.  Along a metric sequence, \(\dim{\cal K}_{g_n}\) may jump at a
more symmetric limiting metric.  Uniform vector coercivity then fails unless
the emerging near-kernel is added to the finite-dimensional slow sector.
This is the geometric analogue of the explicit transition channel selected
by the YM V34 audit.
\end{firewall}

\begin{assessment}[Progress at V81]
\label{ass:e8v81-status}
The momentum equation now has a complete Fredholm ledger: the conformal
Killing kernel, its exact matter/mean-curvature compatibility, quotient Korn
coercivity, and a unique longitudinal correction.  Coupling this quotient to
the V79 fixed-volume scalar sector is coercive when
\(C_H+C_M<2\sqrt{\gamma_S\gamma_V}\).  Arbitrary sources have a unique
augmented solution; vanishing of the volume and conformal-Killing multipliers
is the exact physical-source condition.
\end{assessment}

\begin{programme}[Eighth-series V82: nonlinear coupled constraint chart]
\label{prog:e8v81-V82}
V82 will lift \Cref{thm:e8v81-augmented} to a strong Banach-space
implicit-function theorem for the Hamiltonian and momentum constraints
together.  It will include moving kernel coordinates, derive the derivative
of the kernel multipliers, and isolate the uniform-radius and generated-
interaction rows needed before the chart can be used in the continuum
measure.
\end{programme}

\begin{firewall}[Claim boundary after eighth-series V81]
\label{fire:e8v81-boundary}
V81 proves the vector Korn quotient, exact momentum Fredholm condition, the
conformal metric cross derivative, coupled quotient coercivity, augmented
solvability, and fixed-space cofinal stability.  It does not prove the cross
smallness for the physical sequence, control a jumping conformal-Killing
space, construct the nonlinear coupled chart, solve renormalized Ward
identities, build the continuum measure, or establish the full SM--Einstein
continuum.
\end{firewall}

\begin{theorem}[Eighth-series V81 result]
\label{thm:e8v81-result}
The global linearized Einstein constraint pair is well posed on the
fixed-volume scalar space times the conformal-Killing vector quotient whenever
its two diagonal coercivity floors dominate the summed cross-block norm as in
\eqref{eq:e8v81-cross-small}.  Its full-space augmentation has unique volume
and conformal-Killing multipliers, whose vanishing is exactly the physical
compatibility condition.
\end{theorem}

\begin{proof}
Use \Cref{thm:e8v81-Korn,thm:e8v81-momentum-Fredholm,
thm:e8v81-CKF-source,thm:e8v81-coupled-coercivity,
thm:e8v81-augmented}.
\end{proof}

\paragraph{Parent-version record.}
V81 was formed from
\texttt{Renewal\_SM\_Einstein\_Continuum\_Research\_Notes\_V80.tex}.
The V80 parent had (31\,800) logical source lines, (1\,126\,355) bytes,
and SHA--256
\[
 \texttt{07616F05E3E53883035C797B7C7AA8ECA464D0DA19A7AE1A034F37176A377B5F}.
\]
The next compulsory companion audit is V85.

% =============================================================================
% END APPENDED EIGHTH-SERIES V81 MATERIAL
% =============================================================================
% =============================================================================
% BEGIN APPENDED EIGHTH-SERIES V82 MATERIAL
% =============================================================================

\clearpage
\section{Eighth-series V82: a nonlinear coupled Einstein-constraint chart}
\label{sec:v8v82}

V81 proves invertibility of the normalized scalar--vector linear block and
identifies its volume and conformal-Killing multipliers.  This note constructs
the corresponding nonlinear chart in a strong Sobolev class.  It also resolves
one apparent moving-kernel difficulty: conformal Killing fields are invariant
under conformal changes of the metric, so the kernel is fixed inside a
conformal chart.  Kernel jumps between genuinely different background metrics
remain possible and are treated as a finite-dimensional slow spectral sector.

\subsection{The nonlinear constraint map}
\label{sec:e8v82-map}

Let \(M\) be a connected closed smooth three-manifold, fix \(p>3\), and let
\(g\in W^{2,p}\) be uniformly positive.  Let \({\cal P}\) be a Banach space of
canonical and matter data.  Assume
\begin{equation}
 q\longmapsto K(q)\in W^{1,p}(S^2T^*M)
 \label{eq:e8v82-K-map}
\end{equation}
is \(C^1\), and assume the renormalized matter maps
\begin{align}
 \rho &:W^{2,p}(M)\times W^{2,p}(TM)\times{\cal P}
       \longrightarrow L^p(M),                         \label{eq:e8v82-rho-map}\\
 J &:W^{2,p}(M)\times W^{2,p}(TM)\times{\cal P}
       \longrightarrow L^p(T^*M)                       \label{eq:e8v82-J-map}
\end{align}
are \(C^1\) near \((0,0,q_0)\).

For \(\varphi\in W^{2,p}(M)\) and \(X\in W^{2,p}(TM)\), set
\begin{align}
 g_\varphi&:=e^{2\varphi}g,                             \label{eq:e8v82-gphi}\\
 \widetilde K(\varphi,X,q)
 &:=K(q)+{\cal L}_{g_\varphi}X,                         \label{eq:e8v82-Ktilde}\\
 \widetilde\tau(\varphi,X,q)
 &:=\operatorname{tr}_{g_\varphi}\widetilde K(\varphi,X,q).
 \label{eq:e8v82-tautilde}
\end{align}
Define
\begin{align}
 {\cal H}(\varphi,X,q)
 :={}&R_{g_\varphi}+\widetilde\tau^2
 -|\widetilde K|_{g_\varphi}^2
 -2\kappa_G\rho(\varphi,X,q),                           \label{eq:e8v82-H}\\
 {\cal M}(\varphi,X,q)
 :={}&\operatorname{div}_{g_\varphi}
 \bigl(\widetilde K-\widetilde\tau g_\varphi\bigr)
 -\kappa_GJ(\varphi,X,q).                               \label{eq:e8v82-M}
\end{align}

\begin{theorem}[Strong differentiability of the coupled constraints]
\label{thm:e8v82-C1}
Under the preceding hypotheses,
\begin{equation}
 ({\cal H},{\cal M}):
 W^{2,p}(M)\times W^{2,p}(TM)\times{\cal P}
 \longrightarrow L^p(M)\times L^p(T^*M)                \label{eq:e8v82-C1-map}
\end{equation}
is \(C^1\) near \((0,0,q_0)\).
\end{theorem}

\begin{proof}
Since \(p>3\), \(W^{2,p}\hookrightarrow C^{1,\alpha}\) and \(W^{1,p}\) is a
Banach algebra.  The maps \(g_\varphi\mapsto g_\varphi^{-1}\), its
Christoffel symbols, and its volume density are smooth on a uniform
ellipticity neighborhood.  The scalar curvature contains at most two
derivatives of \(\varphi\) and products of first derivatives, hence lies in
\(L^p\).  The conformal Killing tensor of \(X\) lies in \(W^{1,p}\);
its divergence lies in \(L^p\).  The contractions defining
\(\widetilde\tau^2\), \(|\widetilde K|^2\), and the momentum expression are
finite sums of smooth coefficient maps and products in these Banach spaces.
The matter terms are \(C^1\) by hypothesis.  The Sobolev chain and product
rules therefore prove the claim.
\end{proof}

\begin{firewall}[Declared matter regularity]
\label{fire:e8v82-matter}
Equations \eqref{eq:e8v82-rho-map}--\eqref{eq:e8v82-J-map} are hypotheses on
renormalized composite matter sources.  They include spin-connection,
improvement, gauge-fixing, anomaly, and counterterm dependence.  V82 does not
derive these maps from a regulated Standard Model measure.
\end{firewall}

\subsection{Conformal invariance of the vector kernel}
\label{sec:e8v82-kernel}

\begin{lemma}[Conformal covariance of the Killing operator]
\label{lem:e8v82-conformal-Killing}
For every vector field \(X\),
\begin{equation}
 {\cal L}_{e^{2\varphi}g}X
 =e^{2\varphi}{\cal L}_gX                               \label{eq:e8v82-L-covariance}
\end{equation}
as covariant symmetric two-tensors.  Consequently
\begin{equation}
 {\cal K}_{g_\varphi}={\cal K}_g                        \label{eq:e8v82-kernel-fixed}
\end{equation}
for every \(\varphi\) for which \(g_\varphi\) is defined.
\end{lemma}

\begin{proof}
Write the conformal Killing operator as the trace-free part of the Lie
derivative:
\[
 {\cal L}_gX={\cal L}_Xg-\frac23(\operatorname{div}_gX)g.
\]
Now
\[
 {\cal L}_X(e^{2\varphi}g)
 =e^{2\varphi}\bigl({\cal L}_Xg+2X(\varphi)g\bigr),
\qquad
 \operatorname{div}_{e^{2\varphi}g}X
 =\operatorname{div}_gX+3X(\varphi).
\]
Substitution cancels the \(2X(\varphi)g\) term and gives
\eqref{eq:e8v82-L-covariance}.  Multiplication by the positive function
\(e^{2\varphi}\) does not change the kernel.
\end{proof}

\begin{corollary}[A fixed quotient inside the nonlinear chart]
\label{cor:e8v82-fixed-quotient}
Let \(Y_1,\ldots,Y_N\) be an \(L^2_g\)-orthonormal basis of
\({\cal K}_g\), and set
\begin{equation}
 H_V=\{X\in H^1(TM):(X,Y_A)_{L^2_g}=0,\ 1\leq A\leq N\}.
 \label{eq:e8v82-HV}
\end{equation}
This is a valid fixed complement for every conformal metric \(g_\varphi\).
No differentiating family of kernel vectors is needed within the chart.
\end{corollary}

\begin{proof}
The kernel itself is fixed by \Cref{lem:e8v82-conformal-Killing}; any fixed
closed complement of that finite-dimensional space remains a vector-space
complement.
\end{proof}

\subsection{The augmented nonlinear map}
\label{sec:e8v82-augmented}

Assume the base data solve both constraints:
\begin{equation}
 {\cal H}(0,0,q_0)=0,\qquad {\cal M}(0,0,q_0)=0.         \label{eq:e8v82-base}
\end{equation}
Let \(V_0=\operatorname{Vol}_g(M)\), and use
\(Y_A^\flat=g(Y_A,\cdot)\).  Define
\begin{align}
 {\mathfrak F}(\varphi,X,\lambda,\nu,q)
 :=\bigg(&{\cal H}(\varphi,X,q)+\lambda,                 \notag\\
 &-{\cal M}(\varphi,X,q)+\sum_{A=1}^N\nu_AY_A^\flat,    \notag\\
 &\operatorname{Vol}_{g_\varphi}(M)-V_0,                \notag\\
 &\bigl((X,Y_A)_{L^2_g}\bigr)_{A=1}^N\bigg).            \label{eq:e8v82-F}
\end{align}
Its domain is
\[
 W^{2,p}(M)\times W^{2,p}(TM)\times\mathbb R\times
 \mathbb R^N\times{\cal P},
\]
and its range is
\[
 L^p(M)\times L^p(T^*M)\times\mathbb R\times\mathbb R^N.
\]

Let \(b_S,a_V,c_H,c_M\) be the four base linear forms of V81.  Thus the
weak derivative of the first two components of \(\mathfrak F\) in
\((\varphi,X)\) is
\begin{align}
 b_S(\phi,\psi)+c_H(Z,\psi),                            \label{eq:e8v82-first-linear}\\
 a_V(Z,Y)+c_M(\phi,Y).                                  \label{eq:e8v82-second-linear}
\end{align}

\begin{lemma}[Exact augmented derivative]
\label{lem:e8v82-derivative}
At \((0,0,0,0,q_0)\), the derivative of \(\mathfrak F\) in
\((\varphi,X,\lambda,\nu)\), tested weakly against \((\psi,Y)\), is
\begin{align}
 b_S(\phi,\psi)+c_H(Z,\psi)
 +\lambda\int_M\psi\,d\mu_g,                            \label{eq:e8v82-D1}\\
 a_V(Z,Y)+c_M(\phi,Y)
 +\sum_{A=1}^N\nu_A(Y_A,Y)_{L^2_g},                    \label{eq:e8v82-D2}\\
 3\int_M\phi\,d\mu_g,                                   \label{eq:e8v82-D3}\\
 \bigl((Z,Y_A)_{L^2_g}\bigr)_{A=1}^N.                  \label{eq:e8v82-D4}
\end{align}
\end{lemma}

\begin{proof}
The first two rows are the definitions of the V81 derivative blocks, with
the sign in the second row fixed by using \(-{\cal M}\) in
\eqref{eq:e8v82-F}.  Differentiation in the multipliers gives the displayed
constant and kernel directions.  Since
\(d\mu_{g_\varphi}=e^{3\varphi}d\mu_g\), the volume derivative is
\eqref{eq:e8v82-D3}.  The last row is linear in \(X\).
\end{proof}

\subsection{Strong invertibility of the augmented derivative}
\label{sec:e8v82-invertibility}

Let
\[
 H_S=\left\{\phi\in H^1(M):\int_M\phi\,d\mu_g=0\right\}.
\]
Assume the V81 diagonal bounds
\begin{align}
 b_S(\phi,\phi)&\geq\gamma_S\|\phi\|_{H^1}^2
 &&(\phi\in H_S),                                      \label{eq:e8v82-scalar-floor}\\
 a_V(Z,Z)&\geq\gamma_V\|Z\|_{H^1}^2
 &&(Z\in H_V),                                         \label{eq:e8v82-vector-floor}
\end{align}
and the cross-block margin
\begin{equation}
 C_H+C_M<2\sqrt{\gamma_S\gamma_V}.                      \label{eq:e8v82-cross-margin}
\end{equation}

\begin{assumption}[Strong coupled elliptic regularity]
\label{ass:e8v82-regularity}
Every weak solution of the normalized coupled linear system with
\(L^p\) right side belongs to \(W^{2,p}(M)\times W^{2,p}(TM)\), with
\begin{equation}
 \|\phi\|_{W^{2,p}}+\|Z\|_{W^{2,p}}
 \leq C_{\rm reg}\bigl(
 \|f\|_{L^p}+\|G\|_{L^p}
 +\|\phi\|_{H^1}+\|Z\|_{H^1}\bigr).                    \label{eq:e8v82-regularity}
\end{equation}
\end{assumption}

\begin{remark}[Sufficient coefficient class]
\label{rem:e8v82-regularity}
The assumption follows from the standard \(W^{2,p}\) estimate for a strongly
elliptic diagonal scalar/vector system when its principal coefficients are in
\(W^{1,p}\), its displayed cross terms have order at most one with
\(W^{1,p}\) or \(L^\infty\) coefficients, and the remaining coefficients lie
in \(L^p\).  These conditions hold for the geometric part of the strong chart;
the matter derivatives must be checked against them.
\end{remark}

\begin{theorem}[Strong augmented isomorphism]
\label{thm:e8v82-strong-isomorphism}
Under \eqref{eq:e8v82-scalar-floor}--\eqref{eq:e8v82-cross-margin} and
\Cref{ass:e8v82-regularity}, the derivative in
\Cref{lem:e8v82-derivative} is a bounded isomorphism
\begin{align}
 W^{2,p}(M)\times W^{2,p}(TM)\times\mathbb R\times\mathbb R^N
 \longrightarrow
 L^p(M)\times L^p(T^*M)\times\mathbb R\times\mathbb R^N.
 \label{eq:e8v82-isomorphism}
\end{align}
\end{theorem}

\begin{proof}
Given the two normalization values, first split the scalar into a prescribed
constant plus an element of \(H_S\), and split the vector into prescribed
components along \(Y_A\) plus an element of \(H_V\).  Move all terms involving
the prescribed finite-dimensional components to the right side.  The
remaining normalized weak system has the V81 coupled form.  Its coercivity
floor \(\gamma_*>0\) follows from
\eqref{eq:e8v82-cross-margin}, so Lax--Milgram gives a unique
\((\phi,Z)\in H_S\times H_V\) with an \(H^1\) bound.
\Cref{ass:e8v82-regularity} upgrades this pair to \(W^{2,p}\).

The residual of the scalar equation annihilates \(H_S\), hence is a unique
multiple of the volume functional and determines \(\lambda\).  The residual
of the vector equation annihilates \(H_V\), hence has a unique expansion in
the \(L^2_g\)-orthonormal kernel basis and determines \(\nu\).  This proves
surjectivity and a bounded solution map.  Applied to zero data, the same
argument gives zero normalized fields and zero multipliers, proving
injectivity.
\end{proof}

\subsection{Nonlinear coupled implicit-function theorem}
\label{sec:e8v82-IFT}

\begin{theorem}[Local nonlinear augmented Einstein-constraint chart]
\label{thm:e8v82-IFT}
Under the hypotheses of \Cref{thm:e8v82-C1,thm:e8v82-strong-isomorphism},
there are neighborhoods \(U\subset{\cal P}\) of \(q_0\) and
\({\cal O}\) of \((0,0,0,0)\), and a unique \(C^1\) map
\begin{equation}
 q\longmapsto
 \bigl(\varphi(q),X(q),\lambda(q),\nu(q)\bigr)\in{\cal O}
 \label{eq:e8v82-solution}
\end{equation}
such that
\begin{align}
 {\cal H}(\varphi(q),X(q),q)+\lambda(q)&=0,              \label{eq:e8v82-solved-H}\\
 -{\cal M}(\varphi(q),X(q),q)
 +\sum_A\nu_A(q)Y_A^\flat&=0,                           \label{eq:e8v82-solved-M}\\
 \operatorname{Vol}_{g_{\varphi(q)}}(M)&=V_0,           \label{eq:e8v82-solved-volume}\\
 (X(q),Y_A)_{L^2_g}&=0.                                 \label{eq:e8v82-solved-gauge}
\end{align}
This chart solves the original Hamiltonian and momentum constraints exactly
on the physical multiplier locus
\begin{equation}
 \lambda(q)=0,\qquad \nu(q)=0.                          \label{eq:e8v82-physical-locus}
\end{equation}
\end{theorem}

\begin{proof}
\Cref{thm:e8v82-C1} and the smooth volume and finite-dimensional constraint
maps show that \(\mathfrak F\) is \(C^1\).  It vanishes at the base point by
\eqref{eq:e8v82-base}.  Its derivative in the four unknowns is the
isomorphism of \Cref{thm:e8v82-strong-isomorphism}.  The Banach-space
implicit-function theorem gives \eqref{eq:e8v82-solution} and
\eqref{eq:e8v82-solved-H}--\eqref{eq:e8v82-solved-gauge}.  Removing the
augmentation terms gives the original two constraints precisely when
\eqref{eq:e8v82-physical-locus} holds.
\end{proof}

\begin{firewall}[The multiplier locus is a real condition]
\label{fire:e8v82-multiplier}
The implicit-function theorem solves the augmented equations for every nearby
parameter.  It does not force \(\lambda\) or \(\nu\) to vanish.  The physical
parameter set is the zero set of \(N+1\) additional functions in
\eqref{eq:e8v82-physical-locus}, unless Ward identities, a variable
cosmological parameter, or current/mean-curvature compatibility removes those
conditions.
\end{firewall}

\subsection{Derivative of the physical multipliers}
\label{sec:e8v82-multiplier-derivative}

For \(\dot q\in{\cal P}\), define the direct parameter sources
\begin{align}
 F_H&:=D_q{\cal H}(0,0,q_0)[\dot q]\in L^p(M),           \label{eq:e8v82-FH}\\
 F_V&:=D_q(-{\cal M})(0,0,q_0)[\dot q]\in L^p(T^*M).
 \label{eq:e8v82-FV}
\end{align}

\begin{theorem}[Exact tangent multiplier ledger]
\label{thm:e8v82-multiplier-derivative}
Let
\[
 (\dot\varphi,\dot X,\dot\lambda,\dot\nu)
 =D(\varphi,X,\lambda,\nu)(q_0)[\dot q].
\]
Then
\begin{align}
 \dot\lambda
 =-{1\over\operatorname{Vol}_g(M)}
 \bigl[
 &\int_MF_H\,d\mu_g+b_S(\dot\varphi,1)
 +c_H(\dot X,1)\bigr],                                  \label{eq:e8v82-dotlambda}\\
 \dot\nu_A
 ={}-&\int_M\langle F_V,Y_A\rangle_g\,d\mu_g
 -c_M(\dot\varphi,Y_A).                                 \label{eq:e8v82-dotnu}
\end{align}
Thus a parameter direction is tangent to the physical multiplier locus if
and only if every right side in
\eqref{eq:e8v82-dotlambda}--\eqref{eq:e8v82-dotnu} vanishes.
\end{theorem}

\begin{proof}
Differentiate \(\mathfrak F=0\).  The derivative is
\eqref{eq:e8v82-D1}--\eqref{eq:e8v82-D4}, with direct sources
\(F_H,F_V\) added to the first two rows.  Test the scalar row with \(1\) to
obtain \eqref{eq:e8v82-dotlambda}.  Test the vector row with \(Y_A\).
The energy term vanishes because \({\cal L}_gY_A=0\), and orthonormality
isolates \(\dot\nu_A\), giving \eqref{eq:e8v82-dotnu}.  The last statement is
the derivative of \eqref{eq:e8v82-physical-locus}.
\end{proof}

\begin{corollary}[CMC current compatibility reappears in the tangent chart]
\label{cor:e8v82-CMC}
If the metric-to-momentum cross form annihilates conformal Killing tests,
\(c_M(\phi,Y_A)=0\), then
\[
 \dot\nu_A=0
 \quad\Longleftrightarrow\quad
 \int_M\langle F_V,Y_A\rangle_g\,d\mu_g=0.
\]
For the physical momentum source this is the V81 conformal-Killing
compatibility condition.  Constant mean curvature and a matter current
annihilating \({\cal K}_g\) are sufficient.
\end{corollary}

\begin{proof}
Use \eqref{eq:e8v82-dotnu} and the explicit V81 source identity.
\end{proof}

\subsection{Quantitative chart radius}
\label{sec:e8v82-radius}

The qualitative implicit-function theorem does not give the cutoff-uniform
neighborhood required by a continuum construction.  The following elementary
version records the exact missing constants.

\begin{theorem}[Uniform contraction radius criterion]
\label{thm:e8v82-radius}
Let \(U,Z\) be Banach spaces and let
\({\cal G}:B_U(0,r)\times{\cal Q}\to Z\) be \(C^1\).  Suppose
\[
 S:=D_u{\cal G}(0,q_0)
\]
is invertible with \(\|S^{-1}\|\leq M\), and for a parameter \(q\) one has
\begin{align}
 \sup_{\|u\|\leq r}
 \|D_u{\cal G}(u,q)-S\|&\leq{1\over2M},                 \label{eq:e8v82-radius-derivative}\\
 \|{\cal G}(0,q)\|&\leq{r\over2M}.                      \label{eq:e8v82-radius-source}
\end{align}
Then \({\cal G}(u,q)=0\) has a unique solution in \(B_U(0,r)\), and
\begin{equation}
 \|u\|\leq2M\|{\cal G}(0,q)\|.                          \label{eq:e8v82-radius-bound}
\end{equation}
\end{theorem}

\begin{proof}
Define
\[
 T_q(u):=u-S^{-1}{\cal G}(u,q).
\]
The mean-value formula and \eqref{eq:e8v82-radius-derivative} give
\(\|D T_q(u)\|\leq1/2\) on the ball.  Moreover
\[
 \|T_q(u)\|
 \leq\|T_q(u)-T_q(0)\|+\|T_q(0)\|
 \leq\frac12r+M\|{\cal G}(0,q)\|\leq r,
\]
so \(T_q\) is a contraction of the closed ball.  Its unique fixed point is
the desired zero.  At that point,
\[
 \|u\|\leq\frac12\|u\|+M\|{\cal G}(0,q)\|,
\]
which proves \eqref{eq:e8v82-radius-bound}.
\end{proof}

\begin{definition}[Uniform coupled constraint chart card, UCCCC]
\label{def:e8v82-UCCCC}
For a cofinal regulator family, UCCCC consists of:
\begin{enumerate}[label=\textup{(C\arabic*)},leftmargin=4em]
\item one strong or endpoint Banach chart for all scalar, vector, canonical,
matter, and counterterm variables;
\item uniform scalar and vector quotient coercivity and a strict cross-block
margin;
\item uniform coupled elliptic regularity and a bounded augmented inverse
\(M\);
\item a common radius \(r>0\) satisfying
\eqref{eq:e8v82-radius-derivative};
\item a cofinally controlled base residual satisfying
\eqref{eq:e8v82-radius-source};
\item all generated symmetry-allowed interactions included in \(q\), or a
proved small remainder in the derivative norm;
\item stable volume and conformal-Killing slow sectors, with every emerging
near-kernel retained explicitly;
\item Ward identities or tunable parameters placing the solution on
\(\lambda=\nu=0\); and
\item compatible Jacobian, positivity, locality, and observable estimates for
the induced constrained measure.
\end{enumerate}
\end{definition}

\begin{firewall}[Pointwise IFT constants are not cofinal constants]
\label{fire:e8v82-radius}
At each fixed regulator, smoothness and invertibility produce some local
radius.  A continuum chart requires uniform \(M\), derivative modulus, and
source radius in \Cref{thm:e8v82-radius}.  None follows merely from existence
of the fixed-regulator chart.
\end{firewall}

\subsection{Kernel changes between background metrics}
\label{sec:e8v82-slow}

Let \(P_n=-\operatorname{div}_{g_n}{\cal L}_{g_n}\) be transported unitarily
to one \(L^2\) space.  This outer family changes the base metric, rather than
only its conformal factor.

\begin{theorem}[Finite-dimensional slow-sector transport]
\label{thm:e8v82-slow-sector}
Assume \(P_n\to P\) in norm-resolvent sense and that for some
\(0<\delta_-<\delta_+\),
\begin{equation}
 \operatorname{spec}(P)\cap[\delta_-,\delta_+]=\varnothing.
 \label{eq:e8v82-spectral-buffer}
\end{equation}
Let \(E\) be the spectral projection of \(P\) below the buffer and \(E_n\)
the corresponding Riesz projections for \(P_n\).  Then, for all large \(n\),
\begin{equation}
 \operatorname{rank}E_n=\operatorname{rank}E<\infty,
 \qquad \|E_n-E\|\longrightarrow0,                      \label{eq:e8v82-projection-convergence}
\end{equation}
and \(P_n\) has a uniform positive spectral floor on
\(\operatorname{Ran}(I-E_n)\).  If
\(\dim{\cal K}_{g_n}<\operatorname{rank}E_n\), the extra directions are
lifted near-kernel modes and must remain in the slow sector.
\end{theorem}

\begin{proof}
Choose a closed contour in the resolvent set enclosing the spectrum below
\(\delta_-\) and no spectrum above \(\delta_+\).  Norm-resolvent convergence
gives uniform convergence of resolvents on the contour.  The Riesz formula
\[
 E_n={1\over2\pi i}\int_\Gamma(z-P_n)^{-1}\,dz
\]
then proves norm convergence.  Projections within operator distance less
than one have the same rank, giving the first assertion.  Spectral stability
and the buffer give a common positive lower bound on the complementary
spectrum.  A vector in \(\operatorname{Ran}E_n\) outside the exact kernel has
a positive but buffered eigenvalue, so deleting it would lose the uniform
complement bound.
\end{proof}

\begin{firewall}[A slow sector is not a gauge kernel]
\label{fire:e8v82-slow}
Only \(\ker P_n={\cal K}_{g_n}\) is a conformal-Killing redundancy.  Lifted
eigenvectors in \(\operatorname{Ran}E_n\setminus\ker P_n\) are physical slow
directions.  They require finite-dimensional equations, not gauge
identification.  This is the constraint analogue of the transition channel
retained after the V80 audit.
\end{firewall}

\begin{assessment}[Progress at V82]
\label{ass:e8v82-status}
The scalar and momentum constraints now possess a local strong nonlinear
augmented chart.  The conformal-Killing kernel is exactly fixed inside the
conformal class, the derivative multipliers give the physical tangent
conditions, and the chart-radius theorem identifies the uniform constants
needed for a cofinal construction.  Between nonconformal backgrounds, a
buffered Riesz projection transports every exact and emerging slow mode
without pretending that lifted modes are gauge.
\end{assessment}

\begin{programme}[Eighth-series V83: induced constrained measure and
Jacobian]
\label{prog:e8v82-V83}
V83 will use the local graph
\(q\mapsto(\varphi(q),X(q))\) to derive the exact finite-dimensional or
Banach-chart Jacobian of constraint reduction.  It will separate the
Faddeev--Popov/Korn determinant, the volume and conformal-Killing Gram
determinants, and the physical multiplier delta functions.  It will state
which determinant ratios can be normalized locally and which require
renormalization before any continuum measure claim.
\end{programme}

\begin{firewall}[Claim boundary after eighth-series V82]
\label{fire:e8v82-boundary}
V82 proves strong differentiability of the coupled constraint map, conformal
invariance of its Killing kernel, invertibility of the augmented derivative,
the local nonlinear chart, tangent multiplier formulas, a quantitative
uniform-radius criterion, and buffered slow-sector transport.  It does not
prove UCCCC uniformly, derive the Standard Model matter maps, make the
multipliers vanish, construct or normalize the constrained measure, remove
the regulator, or establish the full SM--Einstein continuum.
\end{firewall}

\begin{theorem}[Eighth-series V82 result]
\label{thm:e8v82-result}
Under the strong matter-map, elliptic-regularity, and V81 coercivity
hypotheses, the Hamiltonian and momentum constraints admit a unique local
fixed-volume, conformal-Killing-normalized \(C^1\) solution graph with explicit
volume and kernel multipliers.  A cutoff-uniform graph requires the concrete
inverse, derivative-modulus, source, generated-interaction, and slow-sector
rows collected in UCCCC.
\end{theorem}

\begin{proof}
Use \Cref{thm:e8v82-C1,lem:e8v82-conformal-Killing,
thm:e8v82-strong-isomorphism,thm:e8v82-IFT,
thm:e8v82-multiplier-derivative,thm:e8v82-radius,
thm:e8v82-slow-sector}.
\end{proof}

\paragraph{Parent-version record.}
V82 was formed from
\texttt{Renewal\_SM\_Einstein\_Continuum\_Research\_Notes\_V81.tex}.
The V81 parent had (32\,356) logical source lines, (1\,147\,004) bytes,
and SHA--256
\[
 \texttt{E1F7A5C7D0D10E4BC90D3F40F5B582F666AEB92FF283732FBD505E8DBDC951E4}.
\]
The next compulsory companion audit is V85.

% =============================================================================
% END APPENDED EIGHTH-SERIES V82 MATERIAL
% =============================================================================
% =============================================================================
% BEGIN APPENDED EIGHTH-SERIES V83 MATERIAL
% =============================================================================

\clearpage
\section{Eighth-series V83: the constrained-measure Jacobian and its
continuum determinant debit}
\label{sec:v8v83}

V82 constructs a local graph for the augmented Hamiltonian and momentum
constraints.  Turning that graph into a measure requires two distinct
Jacobian operations.  First, the constraint delta function eliminates the
conformal scalar, longitudinal vector, volume multiplier, and
conformal-Killing multipliers.  Second, the physical equations restrict the
remaining parameter space to the zero set of the multiplier map
\((\lambda,\nu)\).  This note proves both finite-regulator coarea formulas,
factorizes the augmented determinant, and shows why its continuum limit is
not an ordinary Fredholm determinant without further subtraction.

\subsection{Finite-regulator square constraint chart}
\label{sec:e8v83-finite}

Fix a finite regulator \(R\).  Let \(U_R\cong\mathbb R^{m_R}\) be the
discretized scalar--vector field space, let
\(Z_R\cong\mathbb R^{r_R}\) contain the volume and kernel multipliers, and let
\(Q_R\cong\mathbb R^{n_R}\) contain the remaining canonical, matter, and
counterterm parameters.  Put
\[
 Y_R:=U_R\times Z_R.
\]
Let
\begin{equation}
 {\mathfrak F}_R:Y_R\times Q_R\longrightarrow Y_R
 \label{eq:e8v83-F}
\end{equation}
be the finite-dimensional V82 augmented constraint map, after fixed Euclidean
structures and volume forms have been declared.  Suppose
\begin{equation}
 {\mathfrak F}_R(y_0,q_0)=0,\qquad
 A_R(y,q):=D_y{\mathfrak F}_R(y,q)                      \label{eq:e8v83-A}
\end{equation}
is invertible on a product neighborhood.  The implicit-function theorem gives
a unique graph
\begin{equation}
 y_R(q)=\bigl(u_R(q),\mu_R(q)\bigr),\qquad
 \mu_R(q)=\bigl(\lambda_R(q),\nu_R(q)\bigr)\in Z_R.
 \label{eq:e8v83-graph}
\end{equation}

\begin{theorem}[First coarea elimination]
\label{thm:e8v83-first-coarea}
Let \(W_R\) be a continuous nonnegative weight and let \(\Phi_R\) have compact
support inside a chart on which \(y_R(q)\) is the unique zero of
\({\mathfrak F}_R(\cdot,q)\).  Then
\begin{align}
 \int_{Q_R}\int_{Y_R}
 &\Phi_R(y,q)W_R(y,q)
 \delta_0\!\left({\mathfrak F}_R(y,q)\right)\,dy\,dq
 \notag\\
 &=
 \int_{Q_R}
 {\Phi_R(y_R(q),q)W_R(y_R(q),q)
  \over
  |\det A_R(y_R(q),q)|}\,dq.                            \label{eq:e8v83-first-coarea}
\end{align}
\end{theorem}

\begin{proof}
For each fixed \(q\), the inverse-function theorem makes
\(y\mapsto{\mathfrak F}_R(y,q)\) a diffeomorphism near its unique zero.
Use \(z={\mathfrak F}_R(y,q)\) as the integration variable.  Its Jacobian at
\(z=0\) is
\(|\det D_y{\mathfrak F}_R(y_R(q),q)|^{-1}\).  Evaluation against
\(\delta_0(z)\), followed by integration in \(q\), gives
\eqref{eq:e8v83-first-coarea}.
\end{proof}

\begin{firewall}[The determinant is part of the measure]
\label{fire:e8v83-no-drop}
Solving the constraints and substituting \(y=y_R(q)\) without the factor
\(|\det A_R|^{-1}\) does not produce the delta-constrained measure.  The
Jacobian may be absorbed into a declared effective action, but it may not be
discarded.
\end{firewall}

\subsection{Factorization of the augmented determinant}
\label{sec:e8v83-factor}

Let the field space split orthogonally as
\begin{equation}
 U_R=S_R\oplus Z_R,                                     \label{eq:e8v83-split}
\end{equation}
where \(S_R\) is the fixed-volume and kernel-orthogonal quotient and \(Z_R\)
is the finite slow subspace.  In compatible bases, the augmented derivative
has the block form
\begin{equation}
 {\cal A}_R=
 \begin{pmatrix}
 D_R&B_R&0\\
 C_R&E_R&J_R\\
 0&R_R&0
 \end{pmatrix}.
 \label{eq:e8v83-aug-matrix}
\end{equation}
Here \(D_R:S_R\to S_R^*\) is the normalized coupled scalar--vector
Jacobian, \(R_R:Z_R\to Z_R\) is the normalization Gram map, and
\(J_R:Z_R\to Z_R^*\) inserts the multipliers into the slow equations.

\begin{theorem}[Exact augmented determinant factorization]
\label{thm:e8v83-aug-factor}
If \(D_R,J_R,R_R\) are invertible, then
\begin{equation}
 |\det{\cal A}_R|
 =|\det D_R|\,|\det J_R|\,|\det R_R|.                  \label{eq:e8v83-aug-factor}
\end{equation}
The result is independent of the slow--slow block \(E_R\) and of both
field/slow cross blocks.
\end{theorem}

\begin{proof}
Use block row elimination with the invertible \(D_R\).  It changes neither
determinant nor the last two block columns and reduces
\eqref{eq:e8v83-aug-matrix} to \(D_R\) times
\[
 \begin{pmatrix}
 E_R-C_RD_R^{-1}B_R&J_R\\
 R_R&0
 \end{pmatrix}.
\]
Interchange the two block columns in this \(2r_R\)-dimensional matrix.  Up to
the sign \((-1)^{r_R}\), it becomes block triangular with diagonal blocks
\(J_R\) and \(R_R\).  Taking absolute values gives
\eqref{eq:e8v83-aug-factor}.
\end{proof}

\begin{corollary}[Orthogonal constraint convention]
\label{cor:e8v83-Gram}
If multiplier insertion is the Euclidean adjoint of the normalization map,
\(J_R=R_R^*\), then
\begin{equation}
 |\det{\cal A}_R|
 =|\det D_R|\,|\det R_R|^2.                             \label{eq:e8v83-Gram}
\end{equation}
For an orthonormal volume/kernel basis \(R_R=I\), only the normalized coupled
determinant remains.
\end{corollary}

\begin{proof}
\(|\det R_R^*|=|\det R_R|\); substitute in
\eqref{eq:e8v83-aug-factor}.
\end{proof}

\begin{remark}[Unnormalized constant mode]
\label{rem:e8v83-volume}
If the scalar slow basis uses the literal constant \(1\), its Gram entry is
\(\operatorname{Vol}(M)\).  In that coordinate convention
\eqref{eq:e8v83-Gram} displays a squared volume factor.  Rescaling the basis
to the normalized constant transfers this factor into the Lebesgue volume
form and the multiplier delta function.  A determinant number has meaning
only together with the declared measures on fields, equations, and
multipliers.
\end{remark}

\subsection{Scalar--vector factorization on the quotient}
\label{sec:e8v83-coupled-factor}

Write the normalized coupled derivative from V81 as
\begin{equation}
 D_R=
 \begin{pmatrix}
 A_{S,R}&C_{H,R}\\
 C_{M,R}&P_{V,R}
 \end{pmatrix}.
 \label{eq:e8v83-D}
\end{equation}

\begin{theorem}[Exact quotient Schur determinant]
\label{thm:e8v83-Schur-det}
If \(P_{V,R}\) is invertible on the conformal-Killing complement, then
\begin{equation}
 \det D_R
 =\det P_{V,R}\,
 \det\!\left(A_{S,R}
 -C_{H,R}P_{V,R}^{-1}C_{M,R}\right).                   \label{eq:e8v83-Schur-det}
\end{equation}
If \(A_{S,R}\) is also invertible, this can be written
\begin{align}
 \det D_R
 ={}&\det A_{S,R}\det P_{V,R}                           \notag\\
 &\times
 \det\!\left(
 I-A_{S,R}^{-1}C_{H,R}P_{V,R}^{-1}C_{M,R}
 \right).                                               \label{eq:e8v83-relative-cross}
\end{align}
\end{theorem}

\begin{proof}
Multiply \eqref{eq:e8v83-D} on the right by the determinant-one block matrix
\[
 \begin{pmatrix}
 I&0\\
 -P_{V,R}^{-1}C_{M,R}&I
 \end{pmatrix}.
\]
The result is block upper triangular with diagonal blocks
\(A_{S,R}-C_{H,R}P_{V,R}^{-1}C_{M,R}\) and \(P_{V,R}\), proving
\eqref{eq:e8v83-Schur-det}.  Factor \(A_{S,R}\) from the first determinant to
obtain \eqref{eq:e8v83-relative-cross}.
\end{proof}

\begin{firewall}[Terminology]
\label{fire:e8v83-FP}
The factor \(\det P_{V,R}\) is the determinant of the York/conformal-Killing
quotient operator in this constraint parameterization.  Calling it a
Faddeev--Popov determinant additionally requires a declared gauge group,
gauge slice, invariant base measure, and proof that the orbit derivative is
exactly \(P_{V,R}\).  V83 does not assume that identification.
\end{firewall}

\subsection{The physical multiplier zero set}
\label{sec:e8v83-physical}

The first coarea step leaves the multiplier map
\[
 \mu_R:Q_R\longrightarrow Z_R.
\]
The physical parameter set is
\begin{equation}
 \Sigma_R:=\{q\in Q_R:\mu_R(q)=0\}.                     \label{eq:e8v83-Sigma}
\end{equation}

\begin{theorem}[Second coarea reduction to the physical locus]
\label{thm:e8v83-second-coarea}
Suppose \(0\) is a regular value of \(\mu_R\), so that
\(D\mu_R(q)\) has rank \(r_R\) on \(\Sigma_R\).  Define
\begin{equation}
 J_{\mu_R}(q)
 :=\sqrt{\det\!\left(D\mu_R(q)D\mu_R(q)^*\right)}.      \label{eq:e8v83-Jmu}
\end{equation}
Then
\begin{align}
 \int_{Q_R}\int_{Y_R}
 &\Phi_R(y,q)W_R(y,q)
 \delta_0({\mathfrak F}_R(y,q))
 \delta_0(\mu)\,dy\,dq                                  \notag\\
 &=
 \int_{\Sigma_R}
 {\Phi_R(y_R(q),q)W_R(y_R(q),q)
  \over
  |\det A_R(y_R(q),q)|\,J_{\mu_R}(q)}
 \,d{\cal H}^{n_R-r_R}(q).                              \label{eq:e8v83-second-coarea}
\end{align}
On the left, \(\delta_0(\mu)\) is evaluated on the multiplier coordinates
inside \(y\).
\end{theorem}

\begin{proof}
Apply \Cref{thm:e8v83-first-coarea}.  It evaluates the multiplier coordinate
at \(\mu_R(q)\), leaving the factor \(\delta_0(\mu_R(q))\).  The Euclidean
coarea formula for the submersion \(\mu_R:Q_R\to Z_R\) gives
\[
 \int_{Q_R}f(q)\delta_0(\mu_R(q))\,dq
 =
 \int_{\Sigma_R}{f(q)\over J_{\mu_R}(q)}
 \,d{\cal H}^{n_R-r_R}(q).
\]
Use the first-coarea integrand as \(f\).
\end{proof}

\begin{corollary}[Failure of transversality is a third zero-mode problem]
\label{cor:e8v83-transversality}
If \(D\mu_R\) loses rank on \(\Sigma_R\), formula
\eqref{eq:e8v83-second-coarea} is not licensed there.  The physical locus can
be singular or have excess dimension, and \(J_{\mu_R}\) vanishes.  A further
Lyapunov--Schmidt reduction or stratification is required.
\end{corollary}

\begin{proof}
The regular-value theorem and the coarea density used above both require full
rank.  When rank fails, their Jacobian is zero and their manifold conclusion
does not follow.
\end{proof}

\begin{firewall}[First-order compatibility is not transversality]
\label{fire:e8v83-first-order}
V82 computes \(D\mu_R(q_0)\).  The statement that a selected direction lies
in its kernel means that direction is tangent to the physical locus.  It does
not prove that \(D\mu_R\) is surjective, that zero is a regular value, or that
the locus has the expected codimension.
\end{firewall}

\subsection{Basis and coordinate covariance}
\label{sec:e8v83-coordinates}

\begin{proposition}[Consistent finite-dimensional reparameterization]
\label{prop:e8v83-coordinates}
Let \(y=T\widetilde y\) and
\(\widetilde{\mathfrak F}=S{\mathfrak F}\), where \(T,S\) are invertible
linear maps.  If the Lebesgue measures and delta functions are transformed by
the ordinary change-of-variables rules, the measure in
\eqref{eq:e8v83-first-coarea} is unchanged.  In particular, changing the
conformal-Killing basis changes Gram determinants, multiplier coordinates,
and their volume forms by cancelling factors.
\end{proposition}

\begin{proof}
The transformed derivative is
\[
 D_{\widetilde y}\widetilde{\mathfrak F}
 =S(D_y{\mathfrak F})T.
\]
Thus its determinant gains
\(|\det S||\det T|\).  Meanwhile
\(dy=|\det T|\,d\widetilde y\) and
\(\delta_0(Sz)=|\det S|^{-1}\delta_0(z)\).  The factors cancel.  Applying
this to a basis change in \(Z_R\) proves the last assertion.
\end{proof}

\subsection{Relative determinants and the continuum obstruction}
\label{sec:e8v83-continuum}

Absolute determinants in \eqref{eq:e8v83-aug-factor} grow with the regulator
dimension.  A meaningful continuum candidate must be a normalized ratio or a
renormalized determinant.

\begin{lemma}[Finite-dimensional determinant transport]
\label{lem:e8v83-Jacobi}
Let \(A_t\) be a \(C^1\) path of invertible finite-dimensional matrices.  Then
\begin{equation}
 \log{\det A_1\over\det A_0}
 =\int_0^1\operatorname{Tr}(A_t^{-1}\dot A_t)\,dt,       \label{eq:e8v83-Jacobi}
\end{equation}
with a continuous logarithm chosen along the path.  For positive matrices,
the real logarithm is unambiguous.
\end{lemma}

\begin{proof}
Jacobi's formula gives
\[
 {d\over dt}\det A_t
 =\det A_t\,\operatorname{Tr}(A_t^{-1}\dot A_t).
\]
Divide by \(\det A_t\) and integrate.
\end{proof}

\begin{theorem}[Trace-class relative determinant criterion]
\label{thm:e8v83-Fredholm}
Let \(A_0\) be invertible on a separable Hilbert space and suppose
\begin{equation}
 A=A_0(I+T),\qquad T\in{\cal S}_1,\qquad \|T\|<1,        \label{eq:e8v83-traceclass}
\end{equation}
where \({\cal S}_1\) is the trace class.  Then the relative Fredholm
determinant
\begin{equation}
 \det(A,A_0):=\det(I+T)                                 \label{eq:e8v83-relative-det}
\end{equation}
is defined and
\begin{equation}
 |\log\det(I+T)|
 \leq{\|T\|_1\over1-\|T\|}.                             \label{eq:e8v83-det-bound}
\end{equation}
\end{theorem}

\begin{proof}
The trace-norm convergent series
\[
 \log(I+T)=\sum_{k\geq1}{(-1)^{k+1}\over k}T^k
\]
defines the determinant by
\(\det(I+T)=\exp(\operatorname{Tr}\log(I+T))\).
Since
\(\|T^k\|_1\leq\|T\|^{k-1}\|T\|_1\), summation and
\(1/k\leq1\) give \eqref{eq:e8v83-det-bound}.
\end{proof}

\begin{theorem}[Three-dimensional determinant diagnosis]
\label{thm:e8v83-Schatten}
Let \(A_0\) be a positive elliptic differential operator of order two on a
closed three-manifold, restricted away from its finite kernel.
\begin{enumerate}[label=\textup{(\roman*)},leftmargin=3.8em]
\item If \(\delta A\) is multiplication by a smooth potential, then
\(A_0^{-1}\delta A\) has pseudodifferential order \(-2\).  It is
Hilbert--Schmidt but need not be trace class.
\item If \(\delta A\) has order one, then \(A_0^{-1}\delta A\) has order
\(-1\) and need not be Hilbert--Schmidt.
\item If \(\delta A\) changes the principal order-two symbol, then
\(A_0^{-1}\delta A\) has order zero and is generally not compact.
\end{enumerate}
Consequently the ordinary criterion in
\Cref{thm:e8v83-Fredholm} does not cover the full scalar--vector metric
Jacobian.
\end{theorem}

\begin{proof}
Composition of elliptic pseudodifferential orders gives respectively
\(-2,-1,0\).  An operator of order \(-s\) on a closed \(d\)-manifold has
singular values \(O(j^{-s/d})\); hence it belongs to \({\cal S}_r\) when
\(sr>d\).  For \(d=3,s=2\), membership holds for every \(r>3/2\), including
\({\cal S}_2\), but the trace-class sum at \(r=1\) is not forced.  For
\(s=1\), the condition is \(r>3\), so Hilbert--Schmidt membership is not
forced.  Order zero is not generally compact.  These facts prove all three
claims.
\end{proof}

\begin{corollary}[Modified determinant for a potential-only comparison]
\label{cor:e8v83-det2}
In case (i) of \Cref{thm:e8v83-Schatten}, if
\(T=A_0^{-1}\delta A\in{\cal S}_2\) and \(-1\notin\operatorname{spec}T\), the
modified determinant
\begin{equation}
 \det{}_2(I+T):=\det\bigl((I+T)e^{-T}\bigr)              \label{eq:e8v83-det2}
\end{equation}
is defined.  For \(\|T\|<1\),
\begin{equation}
 \log\det{}_2(I+T)
 =\sum_{k\geq2}{(-1)^{k+1}\over k}\operatorname{Tr}(T^k).
 \label{eq:e8v83-det2-series}
\end{equation}
The missing \(k=1\) trace is exactly the local divergent term that an
ordinary determinant would require one to renormalize.
\end{corollary}

\begin{proof}
For Hilbert--Schmidt \(T\), every \(T^k\), \(k\geq2\), is trace class and the
series converges under \(\|T\|<1\).  Multiplication by \(e^{-T}\) cancels the
linear term in the logarithm, yielding
\eqref{eq:e8v83-det2-series}.  The omitted term would be
\(\operatorname{Tr}T\), which need not exist.
\end{proof}

\begin{firewall}[No formal determinant promotion]
\label{fire:e8v83-continuum-det}
The finite determinants in \eqref{eq:e8v83-second-coarea} are exact at each
regulator.  Their convergence is not implied by the V82 operator-norm chart
estimates.  Principal-symbol changes are not relative trace class, and even a
potential perturbation in three dimensions generally needs a modified
determinant plus a local subtraction.  Zeta, heat-kernel, Pauli--Villars, or
other determinant prescriptions must be matched to the already declared
counterterm convention before a continuum density is claimed.
\end{firewall}

\subsection{Constraint-measure card}
\label{sec:e8v83-card}

\begin{definition}[Constrained chart measure card, CCMC]
\label{def:e8v83-CCMC}
For every regulator and every chart, CCMC records:
\begin{enumerate}[label=\textup{(M\arabic*)},leftmargin=4em]
\item the finite-dimensional base volume forms and the normalized regulated
weight \(W_R\);
\item the square augmented constraint map and the domain on which its zero is
unique;
\item the exact factor
\(|\det A_R|^{-1}\), including quotient, cross, and Gram factors;
\item the multiplier map \(\mu_R\), its physical zero set, and the
transversality density \(J_{\mu_R}^{-1}\);
\item transition functions and determinant covariance on chart overlaps;
\item treatment of Gribov multiplicity or multiple constraint roots outside
the single-root chart;
\item a relative determinant prescription and all local counterterms;
\item uniform integrability of determinant ratios and observable weights;
\item compatibility with gauge, diffeomorphism, reflection, BRST, chiral, and
anomaly identities; and
\item tightness and convergence of the resulting normalized physical
measures.
\end{enumerate}
\end{definition}

\begin{assessment}[Progress at V83]
\label{ass:e8v83-status}
The finite-regulator measure induced by the V82 chart is now exact.  Constraint
elimination contributes the augmented inverse determinant, which factorizes
into the normalized coupled determinant and slow-sector Gram factors.
Restriction to the physical multiplier locus contributes a second coarea
Jacobian.  In the continuum, these determinants are not automatically
Fredholm: potential changes are generally only Hilbert--Schmidt relative to a
second-order operator, and metric principal-symbol changes are order zero.
\end{assessment}

\begin{programme}[Eighth-series V84: adjacent-chart and adjacent-regulator
density ratios]
\label{prog:e8v83-V84}
V84 will compare two overlapping V82 charts and two adjacent regulators.  It
will prove the finite-dimensional Radon--Nikodym cocycle, show the cancellation
of basis-dependent Gram factors, and formulate a uniform-integrability
criterion for local observables.  Determinant counterterms will remain
explicit rather than being hidden in normalization.
\end{programme}

\begin{firewall}[Claim boundary after eighth-series V83]
\label{fire:e8v83-boundary}
V83 proves the two finite-regulator coarea formulas, the exact augmented and
scalar--vector determinant factorizations, coordinate covariance, a
trace-class sufficient criterion, and the three-dimensional Schatten
obstruction.  It does not choose or prove a physical determinant
renormalization, establish transversality, control multiple roots, prove
uniform integrability or tightness, remove the regulator, or construct the
full SM--Einstein continuum.
\end{firewall}

\begin{theorem}[Eighth-series V83 result]
\label{thm:e8v83-result}
Inside a finite-regulator single-root chart, the physical constrained measure
is the Hausdorff measure on the multiplier zero set weighted by the reciprocal
augmented determinant and reciprocal multiplier coarea Jacobian, exactly as
in \eqref{eq:e8v83-second-coarea}.  Passing this density to the continuum
requires a declared renormalized determinant because the full relative
constraint Jacobian is not trace class in general.
\end{theorem}

\begin{proof}
Use \Cref{thm:e8v83-first-coarea,thm:e8v83-aug-factor,
thm:e8v83-Schur-det,thm:e8v83-second-coarea,
prop:e8v83-coordinates,thm:e8v83-Schatten}.
\end{proof}

\paragraph{Parent-version record.}
V83 was formed from
\texttt{Renewal\_SM\_Einstein\_Continuum\_Research\_Notes\_V82.tex}.
The V82 parent had (32\,934) logical source lines, (1\,170\,293) bytes,
and SHA--256
\[
 \texttt{1DFB2BAD454C171100F2D8AB134770F529F918DA65D138D0DB9526A78C3D62E5}.
\]
The next compulsory companion audit is V85.

% =============================================================================
% END APPENDED EIGHTH-SERIES V83 MATERIAL
% =============================================================================
% =============================================================================
% BEGIN APPENDED EIGHTH-SERIES V84 MATERIAL
% =============================================================================

\clearpage
\section{Eighth-series V84: chart cocycles and adjacent-regulator density
ratios}
\label{sec:v8v84}

V83 derives the constrained density in one finite-regulator, single-root
chart.  A continuum construction needs compatible densities on chart
overlaps and controlled changes between regulators.  This note proves the
finite-dimensional chart cocycle, records the exact adjacent-density ledger,
and gives two sufficient convergence criteria: a summable total-variation
criterion on a fixed local marginal, and a stronger summable logarithmic
density criterion.  The counterterm subtraction is required to telescope;
pairwise arbitrary renormalizations do not define a measure cocycle.

\subsection{Physical-locus chart density}
\label{sec:e8v84-chart}

At regulator \(R\), retain the V83 physical locus
\[
 \Sigma_R=\{q:\mu_R(q)=0\}\subset Q_R
\]
and assume it is a smooth \(k_R\)-manifold on the region considered.  Let
\[
 \iota_\alpha:\Omega_\alpha\subset\mathbb R^{k_R}
 \longrightarrow\Sigma_R
\]
be a \(C^1\) embedding.  Define its tangential metric Jacobian
\begin{equation}
 J_{\iota_\alpha}(x)
 :=\sqrt{\det\!\left(
 D\iota_\alpha(x)^*D\iota_\alpha(x)\right)}.            \label{eq:e8v84-Jiota}
\end{equation}
Let
\begin{equation}
 H_R(q):=
 {W_R(y_R(q),q)\over
  |\det A_R(y_R(q),q)|\,J_{\mu_R}(q)}.                  \label{eq:e8v84-H}
\end{equation}
The local coordinate density induced by the V83 second coarea formula is
\begin{equation}
 h_\alpha(x):=H_R(\iota_\alpha(x))J_{\iota_\alpha}(x).
 \label{eq:e8v84-h}
\end{equation}

\begin{theorem}[Exact chart-overlap transformation]
\label{thm:e8v84-overlap}
Suppose \(\iota_\alpha(x)=\iota_\beta(\theta_{\beta\alpha}(x))\) on an
overlap, where
\(\theta_{\beta\alpha}\) is a \(C^1\) diffeomorphism between coordinate
domains.  Then
\begin{equation}
 h_\alpha(x)
 =h_\beta(\theta_{\beta\alpha}(x))
  |\det D\theta_{\beta\alpha}(x)|.                      \label{eq:e8v84-overlap}
\end{equation}
Consequently the measures \(h_\alpha(x)\,dx\) and
\(h_\beta(y)\,dy\) agree on the physical overlap.
\end{theorem}

\begin{proof}
Differentiate
\(\iota_\alpha=\iota_\beta\circ\theta_{\beta\alpha}\).  The Euclidean
area formula gives
\[
 J_{\iota_\alpha}(x)
 =J_{\iota_\beta}(\theta_{\beta\alpha}(x))
  |\det D\theta_{\beta\alpha}(x)|.
\]
The ambient scalar \(H_R\) takes the same value at the common point of
\(\Sigma_R\).  Substitute these two identities in
\eqref{eq:e8v84-h}.
\end{proof}

\begin{corollary}[Radon--Nikodym chart cocycle]
\label{cor:e8v84-chart-cocycle}
Define on overlaps
\begin{equation}
 {\cal R}_{\alpha\beta}(x)
 :={h_\alpha(x)\over
 h_\beta(\theta_{\beta\alpha}(x))}.
 \label{eq:e8v84-Rab}
\end{equation}
Then
\begin{align}
 {\cal R}_{\alpha\beta}(x)
 &=|\det D\theta_{\beta\alpha}(x)|,                     \label{eq:e8v84-Rab-J}\\
 {\cal R}_{\alpha\beta}
 {\cal R}_{\beta\gamma}
 &={\cal R}_{\alpha\gamma}                              \label{eq:e8v84-cocycle}
\end{align}
after the usual pullbacks to one coordinate domain.
\end{corollary}

\begin{proof}
The first identity is \Cref{thm:e8v84-overlap}; the second is the chain rule
for the determinant of
\(\theta_{\gamma\alpha}
=\theta_{\gamma\beta}\circ\theta_{\beta\alpha}\).
\end{proof}

\begin{firewall}[Density values are coordinate dependent]
\label{fire:e8v84-density}
The scalar density \(h_\alpha\) is not expected to have the same numerical
value in two parameterizations.  Equation \eqref{eq:e8v84-overlap}, not
pointwise equality of \(h_\alpha\) and \(h_\beta\), is the coordinate-invariant
statement.
\end{firewall}

\subsection{Chart covers and multiple roots}
\label{sec:e8v84-cover}

Let \(\{\chi_\alpha\}\) be a locally finite partition of unity on the
regular part of \(\Sigma_R\), subordinate to single-root constraint charts.

\begin{theorem}[No-overcount chart assembly]
\label{thm:e8v84-assembly}
For every nonnegative or integrable physical observable \(\Phi\),
\begin{equation}
 \int_{\Sigma_R}\Phi(q)H_R(q)\,d{\cal H}^{k_R}(q)
 =
 \sum_\alpha\int_{\Omega_\alpha}
 \chi_\alpha(\iota_\alpha(x))
 \Phi(\iota_\alpha(x))h_\alpha(x)\,dx.                  \label{eq:e8v84-assembly}
\end{equation}
Overlaps are counted exactly once after summation.
\end{theorem}

\begin{proof}
Apply the area formula in each chart to the right side.  It becomes
\[
 \sum_\alpha\int_{\Sigma_R}
 \chi_\alpha(q)\Phi(q)H_R(q)\,d{\cal H}^{k_R}(q).
\]
Since \(\sum_\alpha\chi_\alpha=1\), this is the left side.
\end{proof}

\begin{firewall}[A chart partition does not solve Gribov multiplicity]
\label{fire:e8v84-Gribov}
Equation \eqref{eq:e8v84-assembly} assembles a known regular physical
manifold.  If the original constraint or gauge map has several inequivalent
roots over the same \(q\), each branch must be included with its own
Jacobian and orientation/multiplicity rule.  A partition of unity on one
branch does not prove that other branches are absent.
\end{firewall}

\subsection{Cancellation of slow-basis Gram factors}
\label{sec:e8v84-Gram}

\begin{proposition}[Kernel-basis covariance in the physical density]
\label{prop:e8v84-Gram}
Let \(Y'_A=\sum_BC_{BA}Y_B\) be an invertible change of the finite slow
basis.  Transform multiplier coordinates by
\(\nu'=C^{-1}\nu\) and transform the kernel normalization equations
contragrediently.  Then the complete product
\begin{equation}
 \delta_0({\mathfrak F}_R)\,\delta_0(\nu)\,dy
 \label{eq:e8v84-complete-delta}
\end{equation}
and hence the physical density in
\eqref{eq:e8v83-second-coarea} are unchanged.  The changes of the two Gram
determinants in \eqref{eq:e8v83-Gram} cancel the multiplier and normalization
coordinate Jacobians.
\end{proposition}

\begin{proof}
The multiplier insertion acquires \(C\), the normalization row acquires
\(C^*\), and their determinant product acquires
\(|\det C|^2\).  The multiplier coordinate change contributes
\(|\det C|^{-1}\), and the corresponding change of normalization coordinates
contributes the second inverse factor.  Equivalently, apply the V83
reparameterization theorem to the block-diagonal transformations.  All
factors cancel in \eqref{eq:e8v84-complete-delta}.
\end{proof}

\subsection{Adjacent-regulator density ledger}
\label{sec:e8v84-adjacent}

To compare regulators, fix a local physical marginal space
\(({\cal X}_A,{\cal B}_A)\) associated with a bounded observable support
\(A\).  Transport the regulator-\(n\) constrained measure to a probability
measure \(\mu_{n,A}\) on this one space.  Suppose
\(\mu_{n+1,A}\ll\mu_{n,A}\) and write
\begin{equation}
 r_{n,A}:={d\mu_{n+1,A}\over d\mu_{n,A}}.               \label{eq:e8v84-r}
\end{equation}

In a common regular chart before normalization, the unnormalized adjacent
log-density change has the exact ledger
\begin{align}
 \Delta_{n,A}
 :={}&\log{W_{n+1,A}\over W_{n,A}}
 -\log{|\det A_{n+1,A}|\over|\det A_{n,A}|}             \notag\\
 &-\log{J_{\mu_{n+1,A}}\over J_{\mu_{n,A}}}
 +\log J_{{\rm tr},n,A}
 -\Delta C_{n,A}.                                       \label{eq:e8v84-ledger}
\end{align}
Here \(J_{{\rm tr},n,A}\) is the coordinate/transport Jacobian and
\(\Delta C_{n,A}\) is the declared local determinant/action counterterm
increment.  The normalized density is
\begin{equation}
 r_{n,A}
 ={e^{\Delta_{n,A}}\over
   \int_{{\cal X}_A}e^{\Delta_{n,A}}\,d\mu_{n,A}}.
 \label{eq:e8v84-normalized-r}
\end{equation}

\begin{proposition}[Normalization identity]
\label{prop:e8v84-normalization}
Equation \eqref{eq:e8v84-normalized-r} is the unique normalization of the
unnormalized adjacent density \(e^{\Delta_{n,A}}\), and
\begin{equation}
 \int_{{\cal X}_A}r_{n,A}\,d\mu_{n,A}=1.                \label{eq:e8v84-r-one}
\end{equation}
An additive constant in \(\Delta_{n,A}\) cancels; a field-dependent
counterterm does not.
\end{proposition}

\begin{proof}
The Radon--Nikodym derivative of a probability measure must integrate to one,
which determines the denominator in
\eqref{eq:e8v84-normalized-r}.  Multiplying numerator and denominator by one
constant leaves the ratio unchanged.  A nonconstant factor cannot be removed
from both integrals.
\end{proof}

\begin{firewall}[Local marginal transport must be declared]
\label{fire:e8v84-transport}
Measures at different cutoffs live on different spaces.  The ratio
\eqref{eq:e8v84-r} exists only after a common local marginal and a transport
map have been specified.  Writing a ratio of two formal functional densities
without that identification is not a Radon--Nikodym derivative.
\end{firewall}

\subsection{Cocycle and counterterm integrability}
\label{sec:e8v84-counterterm}

\begin{theorem}[Adjacent Radon--Nikodym cocycle]
\label{thm:e8v84-RN-cocycle}
Suppose \(\mu_k\ll\mu_j\ll\mu_i\) are probability measures on one measurable
space.  Then
\begin{equation}
 {d\mu_k\over d\mu_i}
 =
 {d\mu_k\over d\mu_j}
 {d\mu_j\over d\mu_i}
 \quad\mu_i\text{-almost everywhere}.                  \label{eq:e8v84-RN-cocycle}
\end{equation}
Consequently a product of adjacent regulator densities represents the direct
ratio only when every factor is transported to the same measurable space
with the correct pullback.
\end{theorem}

\begin{proof}
For every bounded measurable \(f\),
\[
 \int f\,d\mu_k
 =\int f\,{d\mu_k\over d\mu_j}\,d\mu_j
 =\int f\,{d\mu_k\over d\mu_j}
          {d\mu_j\over d\mu_i}\,d\mu_i.
\]
Uniqueness of the Radon--Nikodym derivative gives
\eqref{eq:e8v84-RN-cocycle}.
\end{proof}

\begin{theorem}[Counterterm telescoping criterion]
\label{thm:e8v84-counterterm}
Let \(C_n(x)\) be a declared sequence of local counterterm functionals on the
common transported marginal.  If
\begin{equation}
 \Delta C_n=C_{n+1}-C_n,                                \label{eq:e8v84-coboundary}
\end{equation}
then along \(m>n\)
\begin{equation}
 \sum_{j=n}^{m-1}\Delta C_j=C_m-C_n.                    \label{eq:e8v84-telescope}
\end{equation}
The corresponding exponential factors obey the regulator cocycle.  A family
of arbitrary pairwise subtractions \(\Delta C_{n,m}\) defines a cocycle only
if
\begin{equation}
 \Delta C_{n,k}=\Delta C_{n,j}+\Delta C_{j,k}            \label{eq:e8v84-counter-cocycle}
\end{equation}
after transport, modulo additive normalization constants.
\end{theorem}

\begin{proof}
Equation \eqref{eq:e8v84-telescope} is the telescoping sum of
\eqref{eq:e8v84-coboundary}.  Exponentiation turns sums into products.
Conversely, taking logarithms of a multiplicative density cocycle gives
\eqref{eq:e8v84-counter-cocycle}, with constants invisible after probability
normalization.
\end{proof}

\begin{firewall}[Local subtraction must respect overlaps]
\label{fire:e8v84-counterterm}
A heat-kernel or modified-determinant subtraction chosen independently in
each chart can violate both the chart cocycle
\eqref{eq:e8v84-cocycle} and the regulator cocycle
\eqref{eq:e8v84-RN-cocycle}.  Counterterms must be restrictions of one
geometric functional, or their differences must satisfy the explicit
Čech/regulator cocycle conditions.
\end{firewall}

\subsection{Two sufficient convergence criteria}
\label{sec:e8v84-convergence}

\begin{theorem}[Summable adjacent total variation]
\label{thm:e8v84-TV}
Let \(\mu_n\) be probability measures on one measurable space with
\(\mu_{n+1}=r_n\mu_n\).  If
\begin{equation}
 \sum_{n=1}^\infty
 \|r_n-1\|_{L^1(\mu_n)}<\infty,                         \label{eq:e8v84-TV-sum}
\end{equation}
then \(\mu_n\) converges in total variation to a probability measure
\(\mu_\infty\).  For every bounded measurable \(F\),
\begin{equation}
 \left|\int F\,d\mu_m-\int F\,d\mu_n\right|
 \leq\|F\|_\infty
 \sum_{j=n}^{m-1}\|r_j-1\|_{L^1(\mu_j)}.                \label{eq:e8v84-observable-Cauchy}
\end{equation}
If the space is Polish and one \(\mu_n\) is tight, then \(\mu_\infty\) is
tight.
\end{theorem}

\begin{proof}
By definition,
\[
 \|\mu_{n+1}-\mu_n\|_{\rm TV}
 ={1\over2}\|r_n-1\|_{L^1(\mu_n)}.
\]
The triangle inequality and \eqref{eq:e8v84-TV-sum} make \(\mu_n\) Cauchy in
the Banach space of finite signed measures with total-variation norm.  Its
limit has mass one and is positive, hence is a probability measure.
The direct expectation estimate follows without the \(1/2\) convention by
\[
 \left|\int F(r_j-1)\,d\mu_j\right|
 \leq\|F\|_\infty\|r_j-1\|_1
\]
and telescoping.  Total-variation convergence preserves tightness after
enlarging a compact set to absorb the small tail of the distance.
\end{proof}

\begin{theorem}[Summable logarithmic density changes]
\label{thm:e8v84-log}
In the setting of \eqref{eq:e8v84-normalized-r}, suppose
\begin{equation}
 \|\Delta_n\|_{L^\infty(\mu_n)}\leq\varepsilon_n,
 \qquad \sum_n\varepsilon_n<\infty.                     \label{eq:e8v84-log-sum}
\end{equation}
Then
\begin{align}
 e^{-2\varepsilon_n}\leq r_n\leq e^{2\varepsilon_n},
 \label{eq:e8v84-r-bound}\\
 \|r_n-1\|_{L^1(\mu_n)}
 \leq e^{2\varepsilon_n}-1,                             \label{eq:e8v84-r-L1}
\end{align}
and the measures converge in total variation.
\end{theorem}

\begin{proof}
The hypothesis gives
\(e^{-\varepsilon_n}\leq e^{\Delta_n}\leq e^{\varepsilon_n}\), so its
\(\mu_n\)-integral obeys the same bounds.  Division proves
\eqref{eq:e8v84-r-bound}, hence
\eqref{eq:e8v84-r-L1}.  Since
\(e^{2\varepsilon}-1\leq C\varepsilon\) for all sufficiently small
\(\varepsilon\), the right sides are summable.  Apply
\Cref{thm:e8v84-TV}.
\end{proof}

\begin{remark}[Why the logarithmic criterion is strong]
\label{rem:e8v84-strong}
Uniform control of the complete log-density change includes action,
constraint determinant, multiplier coarea, transport, and counterterm rows.
Continuum measures are often mutually singular globally, so
\eqref{eq:e8v84-log-sum} is intended first for each fixed local transported
marginal.  Weaker entropy or \(L^p\) criteria may be needed globally.
\end{remark}

\subsection{Projective local continuum state}
\label{sec:e8v84-projective}

\begin{theorem}[Consistent local-marginal construction]
\label{thm:e8v84-projective}
Let \(\{{\cal X}_A\}\) be a directed family of standard Borel local
configuration spaces with restriction maps
\(\pi_{BA}:{\cal X}_B\to{\cal X}_A\) for \(A\subset B\).  Suppose:
\begin{enumerate}[label=\textup{(\roman*)},leftmargin=3.8em]
\item for every \(A\), the regulator marginals \(\mu_{n,A}\) converge in
total variation to \(\mu_A\);
\item the regulated marginals are exactly consistent,
\((\pi_{BA})_*\mu_{n,B}=\mu_{n,A}\).
\end{enumerate}
Then the limits are consistent:
\begin{equation}
 (\pi_{BA})_*\mu_B=\mu_A.                               \label{eq:e8v84-limit-consistency}
\end{equation}
For a countable projective system, they define a unique probability measure
on its cylinder \(\sigma\)-algebra.
\end{theorem}

\begin{proof}
Pushforward contracts total variation.  Therefore
\[
 \|(\pi_{BA})_*\mu_B-\mu_A\|_{\rm TV}
 \leq
 \|\mu_B-\mu_{n,B}\|_{\rm TV}
 +\|\mu_{n,A}-\mu_A\|_{\rm TV}\longrightarrow0,
\]
which proves \eqref{eq:e8v84-limit-consistency}.  The standard projective
extension theorem for countable standard Borel systems gives the cylinder
measure and uniqueness.
\end{proof}

\begin{firewall}[Cylinder measure versus continuum field topology]
\label{fire:e8v84-cylinder}
The projective measure in \Cref{thm:e8v84-projective} is defined on the
cylinder \(\sigma\)-algebra.  Concentration on a desired distribution,
connection, spinor, or metric space requires tightness in that topology.
Reflection positivity, locality, and the constraints must also pass to the
limit on their respective test algebras.
\end{firewall}

\subsection{Adjacent-density card}
\label{sec:e8v84-card}

\begin{definition}[Adjacent constrained density card, ACDC]
\label{def:e8v84-ACDC}
For each fixed local support and adjacent regulator pair, ACDC records:
\begin{enumerate}[label=\textup{(D\arabic*)},leftmargin=4em]
\item the common marginal space and explicit transport map;
\item the action/weight ratio;
\item the augmented determinant and multiplier-coarea ratios;
\item every chart-coordinate and slow-basis Gram factor;
\item a counterterm increment satisfying both chart and regulator cocycles;
\item the probability normalization in
\eqref{eq:e8v84-normalized-r};
\item a summable \(L^1\), logarithmic, entropy, or other proved
uniform-integrability bound;
\item consistency under enlargement of the local support;
\item tightness in the intended continuum topology; and
\item preservation of constraints, Ward identities, reflection positivity,
spin/chiral structure, and the observable algebra.
\end{enumerate}
\end{definition}

\begin{assessment}[Progress at V84]
\label{ass:e8v84-status}
The single-chart constrained density of V83 now has exact overlap and
regulator bookkeeping.  Chart densities transform by the tangential
Jacobian, slow-basis Gram factors cancel when measures and delta functions
are transformed consistently, and counterterms must be a telescoping
coboundary.  Summable adjacent \(L^1\) density changes give total-variation
limits for fixed local marginals and hence a consistent cylinder state.
\end{assessment}

\begin{programme}[Eighth-series V85: compulsory companion audit]
\label{prog:e8v84-V85}
V85 will inspect the newest YM and NS active notes.  It will test their
current multiscale or kernel estimates against ACDC, especially the
transition-sector, generated-interaction, and summable adjacent-density rows.
It will then select a concrete V86 route and continue beyond the audit.
\end{programme}

\begin{firewall}[Claim boundary after eighth-series V84]
\label{fire:e8v84-boundary}
V84 proves the chart and regulator cocycles, consistent Gram cancellation,
counterterm telescoping, two sufficient adjacent-measure convergence
criteria, and the projective local-marginal theorem.  It does not establish
the ACDC estimates for the SM--Einstein regulator, choose determinant
counterterms, control singular chart strata, prove continuum tightness or
reflection positivity, or construct the full SM--Einstein continuum.
\end{firewall}

\begin{theorem}[Eighth-series V84 result]
\label{thm:e8v84-result}
Finite-regulator constrained chart measures assemble invariantly exactly
when their coordinate, Gram, determinant, multiplier-coarea, and counterterm
factors obey the displayed cocycles.  On every fixed transported local
marginal, summable adjacent \(L^1\) density changes suffice for a
total-variation limit and a consistent projective cylinder state.
\end{theorem}

\begin{proof}
Use \Cref{thm:e8v84-overlap,thm:e8v84-assembly,
prop:e8v84-Gram,thm:e8v84-RN-cocycle,
thm:e8v84-counterterm,thm:e8v84-TV,
thm:e8v84-log,thm:e8v84-projective}.
\end{proof}

\paragraph{Parent-version record.}
V84 was formed from
\texttt{Renewal\_SM\_Einstein\_Continuum\_Research\_Notes\_V83.tex}.
The V83 parent had (33\,463) logical source lines, (1\,190\,230) bytes,
and SHA--256
\[
 \texttt{8D3785344E0DFE14AC057DD2D432B31BA1CF401D52C7D2F13A8B440892C739B4}.
\]
The next version is the compulsory V85 companion audit.

% =============================================================================
% END APPENDED EIGHTH-SERIES V84 MATERIAL
% =============================================================================

% =============================================================================
% BEGIN APPENDED EIGHTH-SERIES V85 MATERIAL
% =============================================================================

\clearpage
\section{Eighth-series V85: companion audit and the anchored determinant route}
\label{sec:v8v85}

\subsection{Purpose of this checkpoint}
\label{sec:e8v85-purpose}

This is the compulsory five-version comparison requested for the three
parallel programmes.  The comparison is used only to import proof
mechanisms whose hypotheses can be stated in the present constrained
SM--Einstein setting.  Statements about a Yang--Mills Hamiltonian gap or a
Navier--Stokes Oseen kernel are not silently reinterpreted as statements
about the augmented constraint determinant.

The question left by V84 is concrete.  The adjacent-density ledger contains
a determinant row, but an absolute logarithmic determinant is extensive and
need not have a volume-uniform adjacent bound.  The companion audit asks
whether the other programmes supply a local replacement for that false
global target.

\subsection{Files inspected and duplicate detection}
\label{sec:e8v85-files}

The newest files present at this checkpoint were
\begin{align*}
 &\texttt{Renewal\_Yang\_Mills\_Mass\_Gap\_Research\_Notes\_V41.tex},\\
 &\texttt{Renewal\_Yang\_Mills\_Mass\_Gap\_Research\_Notes\_V42.tex},\\
 &\texttt{Renewal\_NS\_Stability\_Research\_Notes\_V26.tex}.
\end{align*}
The two Yang--Mills files each had \(17\,311\) logical source lines,
\(630\,757\) bytes, and the common SHA--256 digest
\[
 \texttt{72E8B61385DA792CE7196CD875CC1B173D86A82D4B2614EF60F09252E67E1E93}.
\]
They are byte-for-byte identical.  Thus V42 is a duplicate filename at this
checkpoint and contributes no mathematical material beyond V41.  Neither
file is modified or removed here because the notes directory is shared.

The Navier--Stokes file had \(5\,319\) logical source lines,
\(278\,283\) bytes, and SHA--256 digest
\[
 \texttt{82C887F8C2C69E4F28FAD7C3DC8DE5B9099CB5A4BF431EBA1FFF3E91935978AD}.
\]
It is unchanged from the preceding companion audit.

\subsection{Yang--Mills developments since the preceding audit}
\label{sec:e8v85-YM}

The substantive Yang--Mills additions V35--V41 establish the following
finite-regulator lessons.
\begin{enumerate}[label=\textup{(Y\arabic*)},leftmargin=4.5em]
\item A three-channel low/transition/high transfer has an exact row-sum
bound.  Separated low/high filter estimates do not by themselves control an
adjacent low/transition block.
\item Physical-time Duhamel estimates use a residual gap of order \(a^{-1}\).
A global interaction norm is extensive, so the useful estimate must be
written in a volume-uniform incidence or connected norm.
\item Exponentially weighted row and column incidence norms give
volume-independent quasi-local bounds.  With a logarithmically growing
locality radius, the locality and finite-range scheduling debits can both
vanish under the declared asymptotic-freedom scale.
\item Exact Feshbach factorization separates open low--residual mixing from
closed residual excursions.  A residual resolvent of size \(O(a)\) turns a
closed pair of vertices of size \(O(g/a)\) into a self-energy of size
\(O(g^2/a)\).
\item Unique owners and support overlap can remove a second-order return
exactly, but this second-order fact is not an all-orders connected expansion.
\item V41 corrects an upstream error: a plaquette crossing two link blocks
has a quadratic bilinear term of order \(a^{-1}\).  It cannot be charged to
an \(O(g/a)\) nonlinear remainder.  The complete cross-block quadratic form
must remain in the exact global Gaussian reference.
\item After subtracting that complete quadratic reference, the Wilson
remainder begins at cubic order with coefficient \(g/a\).  On an
oscillator-number window it is controlled by number-operator weighted
bounds.  The complement of that window remains open.
\end{enumerate}

The item relevant to V84 is the division of labour in (Y6)--(Y7): keep the
nonlocal but exactly solvable quadratic reference intact, and apply
anchored locality estimates only to a marked nonlinear response.  This
prevents an artificial cross-block error from being called small.

\subsection{Navier--Stokes status}
\label{sec:e8v85-NS}

The Navier--Stokes note remains at its V26 rank-one Oseen-kernel and singular
thread closure analysis.  It adds no new hypothesis to the ACDC determinant
row.  In particular, a parabolic kernel estimate does not supply a
determinant trace ideal, an augmented-constraint inverse, or a regulator
Radon--Nikodym estimate.

\subsection{A type-correct transfer principle}
\label{sec:e8v85-transfer}

\begin{proposition}[Marked-response transfer]
\label{prop:e8v85-transfer}
Let \(A_0\) be an invertible finite-dimensional reference matrix, let \(K\)
be a correction such that \(A=A_0(I+K)\) is invertible, and let \(B\) be a
source insertion.  Then
\begin{equation}
 \log\det(A+B)-\log\det A
 =
 \log\det\!\left(I+(I+K)^{-1}A_0^{-1}B\right)          \label{eq:e8v85-marked}
\end{equation}
on every common continuous branch of the logarithm.  Hence the vacuum term
\(\log\det A_0+\log\det(I+K)\), which may be extensive, cancels exactly from
the marked determinant response.
\end{proposition}

\begin{proof}
Factor
\[
 A+B
 =A_0(I+K)
 \left[I+(I+K)^{-1}A_0^{-1}B\right].
\]
Take determinants and subtract \(\log\det A\) on a branch on which the
factors are joined without crossing the determinant zero set.
\end{proof}

\begin{remark}[Branch-free form]
\label{rem:e8v85-branchfree}
The determinant ratio
\[
 {\det(A+B)\over\det A}
 =
 \det\!\left(I+(I+K)^{-1}A_0^{-1}B\right)
\]
is algebraic and branch-free.  The logarithmic statement requires the
declared continuous branch.  Later estimates may therefore be proved first
for the ratio or for the logarithmic derivative.
\end{remark}

\begin{proposition}[What can and cannot be imported]
\label{prop:e8v85-import}
The following implications are type-correct.
\begin{enumerate}[label=\textup{(\roman*)},leftmargin=3.8em]
\item A weighted row/column bound for the relative augmented-constraint
kernel \(K\), with norm strictly below one, gives a Neumann inverse with the
same quasi-local weight.
\item If \(B\) is supported near a fixed observable region, the right side
of \eqref{eq:e8v85-marked} is an anchored response and can admit a
volume-independent estimate.
\item A regulator change supported far from the source can affect the
marked response only through resolvent kernels connecting the source to
that remote support.
\end{enumerate}
None of the following implications follows from the companion notes:
\begin{enumerate}[label=\textup{(\roman*)},leftmargin=3.8em,start=4]
\item the Yang--Mills residual spectral gap does not prove invertibility of
the SM--Einstein constraint linearization;
\item a fixed-time transfer estimate does not prove convergence of
constrained probability densities;
\item second-order Feshbach overlap cancellation does not prove an
all-orders determinant expansion; and
\item an Oseen kernel bound does not prove Schatten or trace-class control
of a constraint determinant.
\end{enumerate}
\end{proposition}

\begin{proof}
For (i), the weighted Schur algebra is submultiplicative, so the Neumann
series converges in that algebra.  Statement (ii) follows from the exact
factorization in \Cref{prop:e8v85-transfer}: the unmarked determinant has
cancelled.  Statement (iii) is the support content of the resolvent identity
\[
 R'-R=-R'(K'-K)R.
\]
The four non-implications compare objects with different domains,
codomains, norms, or limiting statements.  No theorem in the inspected
notes supplies the missing identification or hypothesis.
\end{proof}

\begin{firewall}[Global determinant control is the wrong next target]
\label{fire:e8v85-global}
Even when \(K\) has a volume-uniform weighted Schur norm,
\(\log\det(I+K)\) can grow like the number of regulator cells.  For example,
if \(K=tI\) on \(N\) one-dimensional cells, then
\[
 \|K\|_{\rm Schur}=|t|,
 \qquad
 \log\det(I+K)=N\log(1+t).
\]
Thus locality of the kernel does not make the absolute determinant
nonextensive.  The next estimate must concern a normalized, source-marked,
or counterterm-subtracted local determinant response.
\end{firewall}

\subsection{Decision for V86}
\label{sec:e8v85-decision}

\begin{decision}[Anchored log-determinant locality]
\label{dec:e8v85-V86}
V86 will work at finite regulator on a finite metric set with
finite-dimensional fibres.  It will:
\begin{enumerate}[label=\textup{(A\arabic*)},leftmargin=4em]
\item define a two-sided exponentially weighted Schur algebra;
\item prove its product estimate, Neumann inversion, and exponential kernel
decay;
\item differentiate the source-marked log determinant and bound the trace
by an anchored incidence norm, independently of ambient volume;
\item use the resolvent identity to prove exponential insensitivity to a
remote regulator shell;
\item identify the same result in the all-orders loop expansion, where the
source marks every contributing loop; and
\item state the exact hypotheses still needed to insert the augmented
SM--Einstein constraint operator.
\end{enumerate}
\end{decision}

\begin{theorem}[Audit conclusion]
\label{thm:e8v85-result}
The companion audit replaces the false target of a volume-uniform absolute
determinant by the type-correct target of a volume-uniform marked determinant
response.  The algebraic cancellation is
\eqref{eq:e8v85-marked}; the remaining analytic input is a weighted
two-sided Schur bound for the relative augmented-constraint kernel and its
source or remote-shell perturbations.
\end{theorem}

\begin{proof}
The obstruction to the absolute target is
\Cref{fire:e8v85-global}.  The marked factorization is
\Cref{prop:e8v85-transfer}.  The weighted inverse and remote-response
mechanism are precisely the type-correct transfers isolated in
\Cref{prop:e8v85-import}.  The detailed estimates are deferred to V86 so
that the audit itself does not claim hypotheses not present in the
SM--Einstein construction.
\end{proof}

\begin{assessment}[Progress at V85]
\label{ass:e8v85-status}
The every-five-version audit is complete.  YM V41 corrects its own
cross-block Gaussian bookkeeping and thereby exposes a useful general
principle: exact extensive reference structure should cancel before a
locality norm is applied.  V85 translates that principle into the exact
marked determinant identity and fixes the V86 target.  The unchanged NS
note supplies no additional determinant estimate.
\end{assessment}

\begin{firewall}[Claim boundary after eighth-series V85]
\label{fire:e8v85-boundary}
V85 proves only the finite-dimensional marked determinant factorization and
the logical transfer classification.  It does not prove a weighted Schur
bound for the actual augmented constraint operator, a continuum
renormalized determinant, adjacent-measure summability, singular-stratum
control, reflection positivity, or the full SM--Einstein continuum.
\end{firewall}

\paragraph{Parent-version record.}
V85 was formed from
\texttt{Renewal\_SM\_Einstein\_Continuum\_Research\_Notes\_V84.tex}.
The V84 parent had \(33\,970\) logical source lines, \(1\,209\,483\) bytes,
and SHA--256
\[
 \texttt{1344EC910ED28F67C0C011740C217E90EC7A77DF620F2707B252DBCD93C546AC}.
\]
The next compulsory companion audit is V90.  V86 now executes the anchored
determinant route selected above.

% =============================================================================
% END APPENDED EIGHTH-SERIES V85 MATERIAL
% =============================================================================


% =============================================================================
% BEGIN APPENDED EIGHTH-SERIES V86 MATERIAL
% =============================================================================

\clearpage
\section{Eighth-series V86: anchored determinant locality}
\label{sec:v8v86}

\subsection{Finite-fibre metric kernels}
\label{sec:e8v86-setting}

Let \((X,d)\) be a finite or countable metric set.  At each \(x\in X\) let
\(E_x\) be a real or complex Hilbert space with
\begin{equation}
 1\leq\dim E_x\leq r<\infty,                            \label{eq:e8v86-rank}
\end{equation}
and put
\[
 {\cal H}_X:=\ell^2\mathord{\bigoplus}_{x\in X}E_x.
\]
A block kernel \(M=(M_{xy})\) has
\(M_{xy}:E_y\to E_x\).  For \(\mu\geq0\), define the two-sided weighted
Schur norm
\begin{equation}
 \|M\|_{{\cal S}_\mu}
 :=
 \max\left\{
 \sup_x\sum_y e^{\mu d(x,y)}\|M_{xy}\|,
 \sup_y\sum_x e^{\mu d(x,y)}\|M_{xy}\|
 \right\}.                                             \label{eq:e8v86-Schur}
\end{equation}
All block norms in this section are operator norms.

\begin{lemma}[Weighted Schur algebra]
\label{lem:e8v86-algebra}
If \(M,N\in{\cal S}_\mu\), then
\begin{align}
 \|MN\|_{{\cal S}_\mu}
 &\leq
 \|M\|_{{\cal S}_\mu}\|N\|_{{\cal S}_\mu},             \label{eq:e8v86-product}\\
 \|M\|_{{\cal B}({\cal H}_X)}
 &\leq
 \|M\|_{{\cal S}_0}
 \leq\|M\|_{{\cal S}_\mu}.                             \label{eq:e8v86-Schur-test}
\end{align}
\end{lemma}

\begin{proof}
The triangle inequality for \(d\) gives
\[
 e^{\mu d(x,z)}
 \|(MN)_{xz}\|
 \leq
 \sum_y
 e^{\mu d(x,y)}\|M_{xy}\|\,
 e^{\mu d(y,z)}\|N_{yz}\|.
\]
Sum first in \(z\) and then in \(y\) to obtain the row estimate in
\eqref{eq:e8v86-product}.  Reversing the order gives the column estimate.
For \eqref{eq:e8v86-Schur-test}, the scalar matrix
\((\|M_{xy}\|)\) has row and column sums bounded by
\(\|M\|_{{\cal S}_0}\).  The Schur test therefore bounds its
\(\ell^2\)-operator norm, and the same scalar matrix dominates the block
kernel pointwise.
\end{proof}

\begin{theorem}[Neumann inversion with exponential locality]
\label{thm:e8v86-Neumann}
Suppose \(T\in{\cal S}_\mu\) and
\begin{equation}
 q:=\|T\|_{{\cal S}_\mu}<1.                            \label{eq:e8v86-q}
\end{equation}
Then \(I+T\) is invertible on \({\cal H}_X\), and
\begin{equation}
 R_T:=(I+T)^{-1}=\sum_{k=0}^\infty(-T)^k,\qquad
 \|R_T\|_{{\cal S}_\mu}\leq {1\over1-q}.               \label{eq:e8v86-R}
\end{equation}
In particular,
\begin{equation}
 \|(R_T)_{xy}\|
 \leq {e^{-\mu d(x,y)}\over1-q}.                       \label{eq:e8v86-kernel-decay}
\end{equation}
The constants do not depend on \(|X|\).
\end{theorem}

\begin{proof}
\Cref{lem:e8v86-algebra} makes \({\cal S}_\mu\) a normed algebra on its
completion.  The geometric series converges there and obeys
\[
 \sum_{k=0}^\infty\|T\|_{{\cal S}_\mu}^k=(1-q)^{-1}.
\]
Its product with \(I+T\) is \(I\) on both sides.  The operator statement
follows from \eqref{eq:e8v86-Schur-test}.  One summand in either the weighted
row or column sum bounds a fixed block, giving
\eqref{eq:e8v86-kernel-decay}.
\end{proof}

\begin{remark}[Sufficient, not necessary]
\label{rem:e8v86-sufficient}
The strict inequality \eqref{eq:e8v86-q} is a transparent sufficient
condition.  The determinant-response theorems below remain valid if one
proves directly that every required resolvent belongs to
\({\cal S}_\mu\) with a uniform bound.  This distinction matters when a
parametrix or a Combes--Thomas estimate is available although the chosen
relative kernel is not Neumann-small.
\end{remark}

\subsection{Anchored trace estimates}
\label{sec:e8v86-anchored}

For a finite set \(A\subset X\), say that \(V\) is anchored in \(A\) when
\begin{equation}
 V_{xy}=0\quad\hbox{unless }x,y\in A,                  \label{eq:e8v86-anchor}
\end{equation}
and define
\begin{equation}
 |V|_{1,A}:=\sum_{x,y\in A}\|V_{xy}\|.                 \label{eq:e8v86-Vone}
\end{equation}
Such a \(V\) has rank at most \(r|A|\), independently of the size of \(X\).

\begin{lemma}[Anchored trace]
\label{lem:e8v86-trace}
If \(L\in{\cal S}_0\) and \(V\) is anchored in \(A\), then \(LV\) is
finite rank and
\begin{equation}
 |\operatorname{Tr}(LV)|
 \leq r\|L\|_{{\cal S}_0}|V|_{1,A}.                   \label{eq:e8v86-trace}
\end{equation}
\end{lemma}

\begin{proof}
Only diagonal blocks with base point in \(A\) can contribute after the
block trace is expanded:
\[
 \operatorname{Tr}(LV)
 =
 \sum_{x,y\in A}
 \operatorname{tr}_{E_x}(L_{xy}V_{yx}).
\]
For an endomorphism of \(E_x\),
\(|\operatorname{tr}_{E_x}C|\leq r\|C\|\).  Hence
\[
 |\operatorname{Tr}(LV)|
 \leq
 r\sum_{x,y\in A}\|L_{xy}\|\|V_{yx}\|
 \leq r\|L\|_{{\cal S}_0}|V|_{1,A}.
\]
\end{proof}

For finite sets \(A,B\subset X\), write
\[
 d(A,B):=\inf\{d(x,y):x\in A,\ y\in B\}.
\]

\begin{lemma}[Remote sandwich estimate]
\label{lem:e8v86-sandwich}
Let \(L,R\in{\cal S}_\mu\).  Let \(V\) be anchored in \(A\), and suppose
\(S\) is anchored in \(B\).  If \(D=d(A,B)\), then
\begin{equation}
 |\operatorname{Tr}(LSRV)|
 \leq
 r e^{-2\mu D}
 \|L\|_{{\cal S}_\mu}
 \|S\|_{{\cal S}_0}
 \|R\|_{{\cal S}_\mu}
 |V|_{1,A}.                                            \label{eq:e8v86-sandwich}
\end{equation}
The same conclusion holds for an infinite \(B\) if \(S_{yz}=0\) unless
\(y,z\in B\), \(d(A,B)\geq D\), and \(S\in{\cal S}_0\).
\end{lemma}

\begin{proof}
Expanding the trace gives terms
\[
 \operatorname{tr}_{E_x}
 \left(L_{xy}S_{yz}R_{zw}V_{wx}\right),
 \qquad x,w\in A,\quad y,z\in B.
\]
Both \(d(x,y)\) and \(d(z,w)\) are at least \(D\).  Therefore
\begin{align*}
 &\sum_{y,z}
 \|L_{xy}\|\|S_{yz}\|\|R_{zw}\|\\
 &\quad\leq e^{-2\mu D}
 \sum_{y,z}
 e^{\mu d(x,y)}\|L_{xy}\|\,
 \|S_{yz}\|\,
 e^{\mu d(z,w)}\|R_{zw}\|\\
 &\quad\leq e^{-2\mu D}
 \|L\|_{{\cal S}_\mu}
 \|S\|_{{\cal S}_0}
 \|R\|_{{\cal S}_\mu}.
\end{align*}
For the last inequality, use the weighted row sum of \(L\), the unweighted
row sum of \(S\), and the weighted column sum of \(R\).  Multiply by
\(r\|V_{wx}\|\), sum over \(x,w\in A\), and use
\eqref{eq:e8v86-Vone}.  The proof uses only convergent Schur sums, so it
also covers the stated infinite \(B\).
\end{proof}

\begin{firewall}[Why both Schur directions are recorded]
\label{fire:e8v86-two-sided}
The remote sandwich uses a row sum on the resolvent leaving \(A\) and a
column sum on the resolvent returning to \(A\).  A row estimate alone does
not supply the displayed volume-independent trace bound for a nonnormal
kernel.
\end{firewall}

\subsection{The source-marked determinant}
\label{sec:e8v86-determinant}

Assume \(T\) satisfies \eqref{eq:e8v86-q}, let \(V\) be anchored in a finite
set \(A\), and put
\begin{equation}
 v:=\|V\|_{{\cal S}_\mu}.
 \label{eq:e8v86-v}
\end{equation}
Suppose
\begin{equation}
 q+v<1.                                                \label{eq:e8v86-source-small}
\end{equation}
Since \(R_TV\) has finite rank, its Fredholm determinant is an ordinary
finite-dimensional determinant on its range.  Define the continuous
source-marked logarithm by
\begin{equation}
 {\cal D}_T(V)
 :=
 \int_0^1
 \operatorname{Tr}\!\left[(I+T+sV)^{-1}V\right]\,ds.
 \label{eq:e8v86-D}
\end{equation}

\begin{theorem}[Local determinant ratio and volume-independent bound]
\label{thm:e8v86-localdet}
Under \eqref{eq:e8v86-q}--\eqref{eq:e8v86-source-small},
\begin{align}
 \exp{\cal D}_T(V)
 &=
 \det_{\rm F}\!\left(I+R_TV\right),                    \label{eq:e8v86-Fredholm}\\
 |{\cal D}_T(V)|
 &\leq
 {r|V|_{1,A}\over1-q-v}.                              \label{eq:e8v86-D-bound}
\end{align}
When \(X\) is finite,
\begin{equation}
 \exp{\cal D}_T(V)
 =
 {\det(I+T+V)\over\det(I+T)}.                          \label{eq:e8v86-ratio}
\end{equation}
All estimates are independent of the ambient cardinality of \(X\).
\end{theorem}

\begin{proof}
For \(0\leq s\leq1\),
\[
 \|T+sV\|_{{\cal S}_\mu}\leq q+sv\leq q+v<1.
\]
\Cref{thm:e8v86-Neumann} therefore gives
\begin{equation}
 \|(I+T+sV)^{-1}\|_{{\cal S}_\mu}
 \leq {1\over1-q-sv}
 \leq {1\over1-q-v}.                                  \label{eq:e8v86-Rs}
\end{equation}
The Fredholm determinant differentiation rule for the finite-rank family
\(I+sR_TV\) gives
\[
 {d\over ds}\log\det_{\rm F}(I+sR_TV)
 =
 \operatorname{Tr}\!\left[(I+T+sV)^{-1}V\right].
\]
The logarithm is fixed by its value zero at \(s=0\); invertibility follows
from the same Neumann estimate.  Integration proves
\eqref{eq:e8v86-Fredholm}.  Apply
\Cref{lem:e8v86-trace} and \eqref{eq:e8v86-Rs} under the
integral to prove \eqref{eq:e8v86-D-bound}.  In finite volume,
\[
 I+T+V=(I+T)(I+R_TV),
\]
so multiplicativity of the determinant proves
\eqref{eq:e8v86-ratio}.
\end{proof}

\begin{corollary}[Absolute coarea determinant]
\label{cor:e8v86-absolute}
For real matrices, along the source path of
\Cref{thm:e8v86-localdet},
\begin{equation}
 \log{|\det(I+T+V)|\over|\det(I+T)|}
 =
 \operatorname{Re}{\cal D}_T(V),                      \label{eq:e8v86-absolute}
\end{equation}
and its absolute value obeys \eqref{eq:e8v86-D-bound}.
\end{corollary}

\begin{proof}
The determinant never vanishes on the path.  Its sign is therefore
constant in the real case.  Taking real parts of the continuously chosen
complex logarithm gives the logarithm of the absolute determinant.
\end{proof}

\subsection{Remote regulator shells}
\label{sec:e8v86-remote}

\begin{theorem}[Exponential insensitivity to a remote shell]
\label{thm:e8v86-remote}
Let \(T'=T+S\), where
\begin{equation}
 \|T\|_{{\cal S}_\mu}\leq q,\qquad
 \|T'\|_{{\cal S}_\mu}\leq q',\qquad
 q+v<1,\quad q'+v<1.                                  \label{eq:e8v86-twoq}
\end{equation}
Suppose \(V\) is anchored in \(A\), \(S\) is supported in
\(B\times B\), and \(d(A,B)\geq D\).  Then
\begin{equation}
 |{\cal D}_{T'}(V)-{\cal D}_T(V)|
 \leq
 {r e^{-2\mu D}\|S\|_{{\cal S}_0}|V|_{1,A}
  \over
  (1-q'-v)(1-q-v)}.                                   \label{eq:e8v86-remote}
\end{equation}
\end{theorem}

\begin{proof}
Put
\[
 R_s=(I+T+sV)^{-1},\qquad
 R'_s=(I+T'+sV)^{-1}.
\]
The resolvent identity is
\begin{equation}
 R'_s-R_s=-R'_s S R_s.                                \label{eq:e8v86-resolvent}
\end{equation}
Subtract the two definitions \eqref{eq:e8v86-D}, insert
\eqref{eq:e8v86-resolvent}, and obtain
\[
 {\cal D}_{T'}(V)-{\cal D}_T(V)
 =
 -\int_0^1\operatorname{Tr}(R'_s S R_sV)\,ds.
\]
By \Cref{thm:e8v86-Neumann},
\[
 \|R_s\|_{{\cal S}_\mu}\leq(1-q-v)^{-1},
 \qquad
 \|R'_s\|_{{\cal S}_\mu}\leq(1-q'-v)^{-1}.
\]
Apply \Cref{lem:e8v86-sandwich} and integrate.
\end{proof}

\begin{corollary}[Summable shell changes]
\label{cor:e8v86-shellsum}
Let \(T_{n+1}=T_n+S_n\).  Suppose
\[
 \sup_n\|T_n\|_{{\cal S}_\mu}\leq q_*,
 \qquad q_*+v<1,
\]
and \(S_n\) is supported in \(B_n\times B_n\), where
\(d(A,B_n)\geq D_n\).  If
\begin{equation}
 \sum_n e^{-2\mu D_n}\|S_n\|_{{\cal S}_0}<\infty,      \label{eq:e8v86-shellsum}
\end{equation}
then \({\cal D}_{T_n}(V)\) is Cauchy, with
\begin{equation}
 |{\cal D}_{T_m}(V)-{\cal D}_{T_n}(V)|
 \leq
 {r|V|_{1,A}\over(1-q_*-v)^2}
 \sum_{j=n}^{m-1}
 e^{-2\mu D_j}\|S_j\|_{{\cal S}_0}.                   \label{eq:e8v86-Cauchy}
\end{equation}
\end{corollary}

\begin{proof}
Apply \Cref{thm:e8v86-remote} to each adjacent pair and telescope.
\end{proof}

\begin{firewall}[A shell meeting the observable is not remote]
\label{fire:e8v86-near}
When \(D_n=0\), \eqref{eq:e8v86-remote} gives no small geometric factor.
The ultraviolet cells that meet \(A\) require local matching and
counterterm estimates.  Remote-shell locality controls the thermodynamic
tail; it does not remove the local cutoff divergence.
\end{firewall}

\subsection{All-orders marked-loop interpretation}
\label{sec:e8v86-loops}

\begin{theorem}[Every determinant loop is source-marked]
\label{thm:e8v86-loops}
If
\begin{equation}
 \rho:=\|R_TV\|_{{\cal B}({\cal H}_X)}<1,              \label{eq:e8v86-rho}
\end{equation}
then
\begin{equation}
 {\cal D}_T(V)
 =
 \sum_{m=1}^\infty {(-1)^{m+1}\over m}
 \operatorname{Tr}\bigl[(R_TV)^m\bigr],               \label{eq:e8v86-loopseries}
\end{equation}
and the series is absolutely convergent.  Every term contains at least one
copy of \(V\).  If also \(q<1\), substituting
\(R_T=\sum_{k\geq0}(-T)^k\) expands each connecting resolvent into
unmarked \(T\)-walks whose endpoints lie in the finite anchor \(A\).
\end{theorem}

\begin{proof}
The operator \(R_TV\) has rank at most \(r|A|\).  Its nonzero spectrum is
finite, and \eqref{eq:e8v86-rho} places it in the unit disk.  The scalar
series for \(\log(1+z)\), applied to those eigenvalues with algebraic
multiplicity, gives \eqref{eq:e8v86-loopseries}.  Moreover,
\[
 \left|\operatorname{Tr}(R_TV)^m\right|
 \leq r|A|\rho^m,
\]
so the series is absolutely convergent.  Each power visibly contains the
source insertion.  The Neumann substitution converges in
\({\cal S}_\mu\) by \Cref{thm:e8v86-Neumann}; the intervening products
remain trace class because they contain the finite-rank factor \(V\).
\end{proof}

\begin{remark}[Vacuum loops cancel before estimation]
\label{rem:e8v86-vacuum}
In finite volume, separately expanding
\(\log\det(I+T+V)\) and \(\log\det(I+T)\) produces extensive unmarked
loops.  Equation \eqref{eq:e8v86-loopseries} is their exact reorganization:
all unmarked loops have cancelled, and only loops touching the observable
anchor remain.  In infinite volume the two absolute determinants may not
exist, while the finite-rank relative determinant in
\eqref{eq:e8v86-Fredholm} still does.
\end{remark}

\subsection{Insertion into the augmented constraint determinant}
\label{sec:e8v86-constraint}

At a fixed regulator, let \({\mathbb A}\) denote the augmented
Hamiltonian--momentum constraint derivative constructed in V82, including
the finite slow-sector multiplier rows and columns.  Choose an invertible
reference \({\mathbb A}_0\) on the same augmented space and write
\begin{equation}
 {\mathbb A}
 =
 {\mathbb A}_0(I+T),\qquad
 T={\mathbb A}_0^{-1}({\mathbb A}-{\mathbb A}_0).
 \label{eq:e8v86-relative}
\end{equation}
For a source or local chart perturbation \({\mathbb B}_A\), put
\begin{equation}
 V_A:={\mathbb A}_0^{-1}{\mathbb B}_A.                 \label{eq:e8v86-preconditioned-source}
\end{equation}

\begin{theorem}[Conditional augmented-determinant card]
\label{thm:e8v86-constraint}
Assume that, uniformly in the finite regulator volume,
\begin{enumerate}[label=\textup{(\roman*)},leftmargin=3.8em]
\item \(T\in{\cal S}_\mu\) with \(\|T\|_{{\cal S}_\mu}\leq q<1\);
\item the preconditioned source \(V_A\) is anchored in a fixed finite
region \(A\), with \(\|V_A\|_{{\cal S}_\mu}\leq v\),
\(|V_A|_{1,A}\leq c_A\), and \(q+v<1\);
\item an adjacent remote regulator change induces \(S_n=T_{n+1}-T_n\)
supported in \(B_n\times B_n\).
\end{enumerate}
Then the local constraint-determinant response
\begin{equation}
 \log
 {|\det({\mathbb A}+{\mathbb B}_A)|\over|\det{\mathbb A}|}
 \label{eq:e8v86-constraint-ratio}
\end{equation}
is bounded in absolute value by
\[
 {r c_A\over1-q-v},
\]
independently of ambient volume.  Its response to the remote change obeys
\eqref{eq:e8v86-remote}; if \eqref{eq:e8v86-shellsum} holds, these remote
increments are summable.
\end{theorem}

\begin{proof}
The factors \(\det{\mathbb A}_0\) cancel:
\[
 {\det({\mathbb A}+{\mathbb B}_A)\over\det{\mathbb A}}
 =
 {\det(I+T+V_A)\over\det(I+T)}.
\]
Apply \Cref{thm:e8v86-localdet,cor:e8v86-absolute,
thm:e8v86-remote,cor:e8v86-shellsum}.
\end{proof}

\begin{firewall}[Preconditioning can destroy locality]
\label{fire:e8v86-preconditioner}
Locality of \({\mathbb B}_A\) does not imply locality of
\({\mathbb A}_0^{-1}{\mathbb B}_A\).  The reference operator must retain
all principal and mixed quadratic couplings needed for a uniform inverse,
and that inverse must itself have a proved weighted kernel bound.  This is
the direct constraint-theory analogue of the corrected global Gaussian
bookkeeping found in the V85 YM audit.
\end{firewall}

\begin{firewall}[Trace ideal and continuum boundary]
\label{fire:e8v86-continuum}
The weighted Schur estimate is not a uniform trace-class estimate for the
absolute operator \(T\).  In three continuum dimensions, V83 already showed
that relative perturbations of order \(-2\) are generally
Hilbert--Schmidt rather than trace class.  V86 avoids an extensive trace by
using a finite-rank local source at fixed regulator.  Passing
\eqref{eq:e8v86-constraint-ratio} to the continuum still requires local
counterterms or a compatible modified determinant, convergence of the
preconditioned kernels, and control of sources that cease to be finite rank
as the lattice is refined.
\end{firewall}

\subsection{Revised ACDC determinant row}
\label{sec:e8v86-ACDC}

\begin{definition}[Anchored determinant locality card, ADLC]
\label{def:e8v86-ADLC}
For every regulator and local observable support \(A\), ADLC records:
\begin{enumerate}[label=\textup{(L\arabic*)},leftmargin=4em]
\item the augmented operator \({\mathbb A}\) and exact reference
\({\mathbb A}_0\), including all principal mixed couplings;
\item the metric block decomposition and fibre-rank bound;
\item a two-sided weighted resolvent estimate, proved either by
\eqref{eq:e8v86-q} or by a direct parametrix estimate;
\item the preconditioned local source \(V_A\) and its anchored incidence
bound;
\item the finite-rank determinant ratio
\eqref{eq:e8v86-constraint-ratio};
\item near-cell matching and its local counterterm;
\item the remote-shell supports, distances, and the summable tail
\eqref{eq:e8v86-shellsum};
\item compatibility with the chart and regulator cocycles of V84; and
\item the trace ideal or modified-determinant statement used after cutoff
removal.
\end{enumerate}
\end{definition}

\begin{assessment}[Progress at V86]
\label{ass:e8v86-status}
V86 proves the finite-regulator analytic mechanism selected by the V85
audit.  A strict weighted Schur bound yields an exponentially localized
inverse.  A finite local source then defines a relative Fredholm
determinant whose logarithm is independent of ambient volume, and a remote
shell changes it by the explicit factor \(e^{-2\mu D}\).  The marked-loop
series explains the cancellation of extensive vacuum loops at every order.
\end{assessment}

\begin{programme}[Eighth-series V87: the actual augmented-kernel test]
\label{prog:e8v86-V87}
V87 will test ADLC against the block symbol of the V82 augmented
Hamiltonian--momentum derivative.  It will separate the finite slow block,
construct a principal reference containing the scalar--vector mixing,
derive the exact Schur-complement kernel, and determine whether a Neumann
bound is available.  If it is not, the work will move upstream to a
parameter-elliptic resolvent or Combes--Thomas estimate rather than assuming
locality of the inverse.
\end{programme}

\begin{firewall}[Claim boundary after eighth-series V86]
\label{fire:e8v86-boundary}
V86 proves weighted-algebra, anchored determinant, remote-shell, and marked
loop theorems under explicit finite-regulator hypotheses.  It does not yet
verify those hypotheses for the V82 augmented SM--Einstein operator, choose
the local determinant counterterms, control nonregular constraint strata,
prove adjacent probability-density summability, establish continuum
tightness or reflection positivity, or construct the full SM--Einstein
continuum.
\end{firewall}

\begin{theorem}[Eighth-series V86 result]
\label{thm:e8v86-result}
Under uniform two-sided weighted Schur smallness and a fixed finite source
anchor, the augmented determinant has a volume-independent local relative
response.  Its change under a regulator perturbation supported a distance
\(D\) from the anchor is bounded by
\[
 C_A e^{-2\mu D}\|S\|_{{\cal S}_0},
\]
and summable remote shells give a Cauchy local determinant response.
\end{theorem}

\begin{proof}
The inverse and decay are \Cref{thm:e8v86-Neumann}; the local determinant is
\Cref{thm:e8v86-localdet}; the remote estimate and summability are
\Cref{thm:e8v86-remote,cor:e8v86-shellsum}; and the augmented-constraint
factorization is \Cref{thm:e8v86-constraint}.
\end{proof}

\paragraph{Parent-version record.}
V86 was formed from
\texttt{Renewal\_SM\_Einstein\_Continuum\_Research\_Notes\_V85.tex}.
The V85 parent had \(34\,239\) logical source lines, \(1\,221\,102\) bytes,
and SHA--256
\[
 \texttt{F6330582B5F81C3C972140A31CC04DE60C6C9A345714244DD4591A387848B992}.
\]
The next compulsory companion audit remains V90.

% =============================================================================
% END APPENDED EIGHTH-SERIES V86 MATERIAL
% =============================================================================


% =============================================================================
% BEGIN APPENDED EIGHTH-SERIES V87 MATERIAL
% =============================================================================

\clearpage
\section{Eighth-series V87: testing the augmented Einstein-constraint kernel}
\label{sec:v8v87}

\subsection{Reduction to the normalized quotient}
\label{sec:e8v87-reduction}

At finite regulator \(R\), V83 writes the augmented derivative as
\[
 {\cal A}_R=
 \begin{pmatrix}
 D_R&B_R&0\\
 C_R&E_R&J_R\\
 0&R_R&0
 \end{pmatrix}
\]
and proves
\begin{equation}
 |\det{\cal A}_R|
 =
 |\det D_R|\,|\det J_R|\,|\det R_R|.                  \label{eq:e8v87-augfactor}
\end{equation}
Here \(D_R\) acts on the fixed-volume scalar and
conformal-Killing-orthogonal vector quotient.  Thus the slow--slow block and
the field/slow cross blocks do not enter the determinant.  If a local source
changes only \(D_R\) and the same slow-coordinate convention is used in
numerator and denominator, then
\begin{equation}
 {|\det({\cal A}_R+{\cal B}_{A,R})|\over|\det{\cal A}_R|}
 =
 {|\det(D_R+B_{A,R})|\over|\det D_R|}.                 \label{eq:e8v87-quotient-ratio}
\end{equation}

\begin{proof}
Apply \eqref{eq:e8v87-augfactor} to numerator and denominator.  The
\(J_R\) and \(R_R\) factors cancel because the source leaves those maps
fixed.
\end{proof}

\begin{firewall}[Sources that move the slow coordinates]
\label{fire:e8v87-moving-slow}
Equation \eqref{eq:e8v87-quotient-ratio} does not apply when the source
changes the volume functional, the conformal-Killing basis, or the
multiplier insertion.  Those finite-dimensional changes must be kept in
the Gram and coarea ledger.  V84 proves their cancellation only after the
coordinate measures and delta functions are transformed consistently.
\end{firewall}

\subsection{Actual differential orders and principal symbol}
\label{sec:e8v87-symbol}

Write the normalized coupled derivative as
\begin{equation}
 D_R=
 \begin{pmatrix}
 A_R&C_R\\
 E_R&P_R
 \end{pmatrix}.                                        \label{eq:e8v87-D}
\end{equation}
The notation matches the V81 forms:
\[
 C_RX=c_H(X,\cdot),\qquad E_R\phi=c_M(\phi,\cdot).
\]
The scalar principal form is
\[
 4\int_M\langle\nabla\phi,\nabla\psi\rangle_g\,d\mu_g,
\]
and the vector principal operator is
\[
 P_gX=-\nabla^i({\cal L}_gX)_{ij}.
\]

\begin{theorem}[Diagonal strong principal symbol]
\label{thm:e8v87-symbol}
Assume, as in the strong regularity row of V82, that the matter derivatives
add no second-order scalar--vector cross term.  At a nonzero covector
\(\xi\), the second-order principal symbol of \(D_R\), before finite
discretization, is
\begin{equation}
 \sigma_2(D)(x,\xi)
 =
 \begin{pmatrix}
 4|\xi|_g^2&0\\
 0&|\xi|_g^2 I+\frac13\xi^\sharp\otimes\xi
 \end{pmatrix}.                                        \label{eq:e8v87-symbol}
\end{equation}
Its vector eigenvalues are \(|\xi|_g^2\) on
\(\xi^\perp\) and \(\frac43|\xi|_g^2\) on
\(\mathbb R\xi^\sharp\).  Hence the system is strongly elliptic.  Both
cross blocks \(C_R,E_R\) have differential order at most one.
\end{theorem}

\begin{proof}
The scalar statement follows from the principal form by integration by
parts.  In normal coordinates at \(x\), the conformal-Killing symbol is
\[
 \sigma_1({\cal L})(\xi)X
 =
 \mathrm i\left(
 \xi_iX_j+\xi_jX_i-\frac23(\xi\cdot X)g_{ij}
 \right).
\]
Contracting once with \(\xi^i\) in
\(-\nabla^i{\cal L}\) gives
\[
 \sigma_2(P_g)(\xi)X
 =
 |\xi|^2X+\frac13\xi^\sharp(\xi\cdot X).
\]
The displayed eigenvalues follow by decomposing \(X\) parallel and
orthogonal to \(\xi^\sharp\).

The Hamiltonian cross form
\[
 c_H(X,\psi)
 =-2\int_M\langle K,{\cal L}_gX\rangle_g\psi\,d\mu_g
\]
contains one derivative of \(X\).  The geometric part of the other cross
block is a sum of a zeroth-order multiple of \(\phi\) and
\((K_i{}^j+\tau\delta_i{}^j)\nabla_j\phi\), hence has at most one derivative
of \(\phi\).  The hypothesis gives the same order for the matter part.
Thus neither cross block contributes to \(\sigma_2(D)\).
\end{proof}

\begin{corollary}[Correct exact reference]
\label{cor:e8v87-reference}
An exact diagonal reference
\begin{equation}
 D_{0,R}:=
 \begin{pmatrix}A_R&0\\0&P_R\end{pmatrix}              \label{eq:e8v87-D0}
\end{equation}
retains every second-order term of the coupled constraint derivative.
The relative correction contains only the preconditioned first- and
zeroth-order cross blocks.  There is no missing cross-block second-order
form analogous to the obstruction found in the V85 Yang--Mills audit.
\end{corollary}

\begin{proof}
This is exactly the order decomposition in
\Cref{thm:e8v87-symbol}.  The diagonal blocks in
\eqref{eq:e8v87-D0} are the complete blocks, not merely their principal
parts, so diagonal lower-order potentials are also retained.
\end{proof}

\subsection{Exact relative cross kernel}
\label{sec:e8v87-relative}

Assume \(A_R\) and \(P_R\) are invertible on their normalized quotient
spaces and define
\begin{equation}
 U_R:=A_R^{-1}C_R,\qquad
 V_R:=P_R^{-1}E_R.                                     \label{eq:e8v87-UV}
\end{equation}
Then
\begin{equation}
 D_R=D_{0,R}(I+T_R),\qquad
 T_R=
 \begin{pmatrix}0&U_R\\V_R&0\end{pmatrix}.             \label{eq:e8v87-T}
\end{equation}

\begin{theorem}[Exact off-diagonal inverse]
\label{thm:e8v87-offdiag}
Let \({\cal B}\) be any unital Banach algebra of compatible block kernels.
If
\begin{equation}
 u:=\|U_R\|_{\cal B},\qquad
 v:=\|V_R\|_{\cal B},\qquad uv<1,                      \label{eq:e8v87-uv}
\end{equation}
then \(I+T_R\) is invertible and
\begin{equation}
 (I+T_R)^{-1}
 =
 \begin{pmatrix}
 (I-U_RV_R)^{-1}&-(I-U_RV_R)^{-1}U_R\\
 -(I-V_RU_R)^{-1}V_R&(I-V_RU_R)^{-1}
 \end{pmatrix}.                                        \label{eq:e8v87-offdiag-inverse}
\end{equation}
The four blocks obey
\begin{align}
 \|(I-U_RV_R)^{-1}\|_{\cal B},
 \|(I-V_RU_R)^{-1}\|_{\cal B}
 &\leq {1\over1-uv},                                   \label{eq:e8v87-diag-bound}\\
 \|(I-U_RV_R)^{-1}U_R\|_{\cal B}
 &\leq {u\over1-uv},                                   \label{eq:e8v87-U-bound}\\
 \|(I-V_RU_R)^{-1}V_R\|_{\cal B}
 &\leq {v\over1-uv}.                                   \label{eq:e8v87-V-bound}
\end{align}
\end{theorem}

\begin{proof}
The products \(U_RV_R\) and \(V_RU_R\) have norms at most \(uv<1\).
Their inverse geometric series therefore converge in \({\cal B}\) and give
\eqref{eq:e8v87-diag-bound}.  Direct block multiplication proves
\eqref{eq:e8v87-offdiag-inverse}; use
\[
 (I-U_RV_R)^{-1}U_R
 =
 U_R(I-V_RU_R)^{-1}
\]
for the off-diagonal entries.  The remaining bounds follow by
submultiplicativity.
\end{proof}

\begin{corollary}[Sharp weighted locality criterion]
\label{cor:e8v87-weighted}
Take \({\cal B}={\cal S}_\mu\), with the weighted Schur algebra of V86.  If
\begin{equation}
 \|A_R^{-1}C_R\|_{{\cal S}_\mu}
 \|P_R^{-1}E_R\|_{{\cal S}_\mu}<1                     \label{eq:e8v87-weighted-criterion}
\end{equation}
uniformly in \(R\), then the coupled inverse \(D_R^{-1}\) has a uniform
weighted kernel bound whenever \(A_R^{-1}\) and \(P_R^{-1}\) do.  More
explicitly,
\begin{equation}
 D_R^{-1}=(I+T_R)^{-1}
 \begin{pmatrix}A_R^{-1}&0\\0&P_R^{-1}\end{pmatrix},   \label{eq:e8v87-Dinverse}
\end{equation}
and the bounds are obtained by multiplying
\eqref{eq:e8v87-diag-bound}--\eqref{eq:e8v87-V-bound}
by the corresponding diagonal inverse norms.
\end{corollary}

\begin{proof}
V86 proves that \({\cal S}_\mu\) is a Banach algebra.  Apply
\Cref{thm:e8v87-offdiag} and then
\eqref{eq:e8v87-Dinverse}.
\end{proof}

\begin{remark}[Balanced block norm]
\label{rem:e8v87-balanced}
The product condition is sharper than demanding both \(u<1\) and \(v<1\).
Indeed, rescale the scalar and vector components by
\(\operatorname{diag}(sI,I)\).  The two off-diagonal norms become \(su\)
and \(s^{-1}v\).  Their optimal common upper bound is
\(\sqrt{uv}\), attained at \(s=\sqrt{v/u}\) when \(u,v>0\).
\end{remark}

\subsection{Relation to the V81 coercivity margin}
\label{sec:e8v87-coercivity}

\begin{proposition}[Hilbert-space product bound]
\label{prop:e8v87-Hilbert}
Measure scalar and vector fields in their \(H^1\) quotient norms and
sources in the corresponding dual norms.  Under the V81 diagonal
coercivities and cross estimates,
\begin{align}
 \|A_R^{-1}C_R\|_{{\cal B}(H_V,H_S)}
 &\leq {C_H\over\gamma_S},                             \label{eq:e8v87-uH}\\
 \|P_R^{-1}E_R\|_{{\cal B}(H_S,H_V)}
 &\leq {C_M\over\gamma_V}.                             \label{eq:e8v87-vH}
\end{align}
Moreover, the V81 condition
\[
 C_H+C_M<2\sqrt{\gamma_S\gamma_V}
\]
implies
\begin{equation}
 {C_HC_M\over\gamma_S\gamma_V}<1.                     \label{eq:e8v87-H-product}
\end{equation}
Thus it implies the off-diagonal product criterion in these global Hilbert
operator norms.
\end{proposition}

\begin{proof}
If \(f=C_RX\), coercivity and the weak equation
\(A_R\phi=f\) give
\[
 \gamma_S\|\phi\|_{H^1}
 \leq\|f\|_{H_S^*}
 \leq C_H\|X\|_{H^1},
\]
which proves \eqref{eq:e8v87-uH}.  The vector estimate is identical.
Finally,
\[
 C_HC_M\leq\left({C_H+C_M\over2}\right)^2
 <\gamma_S\gamma_V
\]
by the arithmetic--geometric mean inequality.
\end{proof}

\begin{firewall}[Coercivity is not weighted locality]
\label{fire:e8v87-coercivity-locality}
The norms in \Cref{prop:e8v87-Hilbert} are global Sobolev operator norms.
They do not bound either weighted row sums or weighted column sums of a
regulator kernel.  Therefore V81 proves invertibility of \(D_R\), but it
does not prove \eqref{eq:e8v87-weighted-criterion}.  Replacing one norm by
the other without a kernel theorem would be circular.
\end{firewall}

\subsection{What diagonal estimates would suffice}
\label{sec:e8v87-sufficient}

\begin{proposition}[A checkable Schur product card]
\label{prop:e8v87-card}
Suppose, on one regulator metric block decomposition,
\begin{align}
 \|A_R^{-1}\|_{{\cal S}_\mu}&\leq a_\mu,&
 \|P_R^{-1}\|_{{\cal S}_\mu}&\leq p_\mu,               \label{eq:e8v87-diagonal-Schur}\\
 \|C_R\|_{{\cal S}_\mu}&\leq c_\mu,&
 \|E_R\|_{{\cal S}_\mu}&\leq e_\mu.                   \label{eq:e8v87-cross-Schur}
\end{align}
If
\begin{equation}
 a_\mu p_\mu c_\mu e_\mu<1,                            \label{eq:e8v87-four-product}
\end{equation}
then \eqref{eq:e8v87-weighted-criterion} holds and the V86 determinant
locality theorem applies to sources whose preconditioned quotient kernel
has the required anchor bound.
\end{proposition}

\begin{proof}
Submultiplicativity gives
\[
 \|U_R\|_{{\cal S}_\mu}\leq a_\mu c_\mu,\qquad
 \|V_R\|_{{\cal S}_\mu}\leq p_\mu e_\mu.
\]
Their product is strictly below one by
\eqref{eq:e8v87-four-product}.  Apply
\Cref{cor:e8v87-weighted} and the V86 local determinant theorem.
\end{proof}

\begin{remark}[Cross blocks are lower order but not automatically small]
\label{rem:e8v87-lower-order}
Differential order controls ultraviolet scaling after a regulator norm and
cell normalization have been fixed.  It does not imply that
\(c_\mu e_\mu\) is numerically small at every physical background.  The
smallness in \eqref{eq:e8v87-four-product} is a quantitative background
condition, parallel to the V81 coercivity margin.
\end{remark}

\subsection{Slow-mode and refinement obstructions}
\label{sec:e8v87-obstructions}

\begin{proposition}[Finite slow projectors are not uniformly local]
\label{prop:e8v87-slow}
Let \(e_0=|X_R|^{-1/2}(1,\ldots,1)\) on a connected finite regulator graph
with counting measure, and let
\(\Pi_0=e_0e_0^*\).  If the graph diameter tends to infinity and
\(\mu>0\), then
\begin{equation}
 \|\Pi_0\|_{{\cal S}_\mu}
 =
 \sup_x {1\over|X_R|}
 \sum_y e^{\mu d(x,y)}                                 \label{eq:e8v87-projector}
\end{equation}
need not remain bounded and grows exponentially on ordinary expanding
boxes.  The same issue occurs for delocalized conformal-Killing projectors.
\end{proposition}

\begin{proof}
The kernel of \(\Pi_0\) is the constant
\(|X_R|^{-1}\).  Substitution into
\eqref{eq:e8v86-Schur} gives
\eqref{eq:e8v87-projector}.  On a path or box, choose \(x\) near one
boundary.  A fixed positive fraction of the vertices lie at distance
comparable with the diameter, so the exponential weight dominates the
polynomial volume denominator.
\end{proof}

\begin{consequence}[The slow sector must remain finite dimensional]
\label{cons:e8v87-slow}
One must not insert the volume and conformal-Killing projectors into the
weighted local kernel and then ask for a uniform Schur norm.  Their
determinants and coordinate Jacobians remain in the finite-dimensional
V83--V84 ledger, while the weighted kernel theorem is applied to the
normalized fast quotient.
\end{consequence}

\begin{proof}
\Cref{prop:e8v87-slow} obstructs the first procedure.
Equation \eqref{eq:e8v87-augfactor} gives the exact algebraic separation
used by the second.
\end{proof}

\begin{firewall}[A fixed physical anchor is not a fixed-rank anchor]
\label{fire:e8v87-refinement}
Under ultraviolet refinement, a physical region \(A\) contains more
regulator cells.  The V86 constant \(r|V|_{1,A}\) is useful only after the
matrix blocks, quadrature weights, and source scaling have been normalized
so that \(|V|_{1,A}\) remains controlled.  Fixed cardinality of the anchor
cannot be assumed for a fixed physical observable.
\end{firewall}

\subsection{ADLC verdict for the V82 operator}
\label{sec:e8v87-verdict}

\begin{theorem}[Actual augmented-kernel test]
\label{thm:e8v87-verdict}
For the V82 augmented Hamiltonian--momentum derivative:
\begin{enumerate}[label=\textup{(\roman*)},leftmargin=3.8em]
\item the finite slow determinant separates exactly by
\eqref{eq:e8v87-augfactor};
\item the fast quotient has the strongly elliptic diagonal principal symbol
\eqref{eq:e8v87-symbol};
\item its exact relative cross kernel is
\eqref{eq:e8v87-T}, and the sharp weighted sufficient condition is
\eqref{eq:e8v87-weighted-criterion};
\item V81 coercivity proves the analogous global Hilbert product bound
\eqref{eq:e8v87-H-product}; but
\item the current programme has not yet proved the diagonal weighted inverse
bounds, the weighted cross bounds, or the refinement-uniform source
incidence needed for \eqref{eq:e8v87-four-product}.
\end{enumerate}
Consequently the V86 determinant theorem is structurally compatible with
the actual constraint operator, but its principal locality hypothesis
remains unpaid.
\end{theorem}

\begin{proof}
Items (i)--(iii) are
\Cref{sec:e8v87-reduction,thm:e8v87-symbol,
thm:e8v87-offdiag,cor:e8v87-weighted}.  Item (iv) is
\Cref{prop:e8v87-Hilbert}.  The earlier V82 assumptions are Sobolev
coercivity and elliptic regularity estimates; neither contains the
two-sided weighted kernel bounds in
\eqref{eq:e8v87-diagonal-Schur}--\eqref{eq:e8v87-cross-Schur}, nor the
cell-normalized anchor estimate identified in
\Cref{fire:e8v87-refinement}.  This proves (v) and the conclusion.
\end{proof}

\begin{assessment}[Progress at V87]
\label{ass:e8v87-status}
The abstract V86 hypothesis has now been tested against the actual
Hamiltonian--momentum derivative.  No second-order cross coupling was lost:
the complete second-order symbol is diagonal and strongly elliptic.  The
remaining cross correction has the exact off-diagonal form, whose inverse
is controlled by the product of two preconditioned cross kernels.  The
existing coercivity margin closes that product only in a global Sobolev
norm.  Weighted locality and refinement normalization are the next upstream
tasks.
\end{assessment}

\begin{programme}[Eighth-series V88: upstream weighted inverse theorem]
\label{prog:e8v87-V88}
V88 will prove a discrete Combes--Thomas estimate for a local coercive
diagonal regulator operator, including the passage from pointwise
exponential decay to a two-sided Schur bound under a uniform volume-growth
condition.  It will state the cell-weight normalization explicitly and then
apply the result separately to \(A_R\) and \(P_R\).  This is the upstream
input required by \eqref{eq:e8v87-four-product}.
\end{programme}

\begin{firewall}[Claim boundary after eighth-series V87]
\label{fire:e8v87-boundary}
V87 proves the actual principal symbol, the exact preconditioned
off-diagonal inverse, its sharp product criterion, and the distinction
between global coercivity and weighted locality.  It does not yet prove the
regulator-uniform weighted diagonal resolvents or cross/source incidence,
select a determinant counterterm, control singular constraint strata,
establish adjacent-measure convergence, or construct the full
SM--Einstein continuum.
\end{firewall}

\begin{theorem}[Eighth-series V87 result]
\label{thm:e8v87-result}
After exact removal of the finite slow determinant, the normalized
SM--Einstein constraint derivative is a strongly elliptic diagonal operator
times an off-diagonal relative kernel.  Its weighted inverse exists with
uniform bounds whenever
\[
 \|A_R^{-1}C_R\|_{{\cal S}_\mu}
 \|P_R^{-1}E_R\|_{{\cal S}_\mu}<1.
\]
The V81 coercivity margin proves this inequality only in the corresponding
global Sobolev operator norm, so a weighted diagonal resolvent theorem is
the next necessary input.
\end{theorem}

\begin{proof}
Combine \Cref{thm:e8v87-symbol,thm:e8v87-offdiag,
cor:e8v87-weighted,prop:e8v87-Hilbert,thm:e8v87-verdict}.
\end{proof}

\paragraph{Parent-version record.}
V87 was formed from
\texttt{Renewal\_SM\_Einstein\_Continuum\_Research\_Notes\_V86.tex}.
The V86 parent had \(34\,838\) logical source lines, \(1\,241\,661\) bytes,
and SHA--256
\[
 \texttt{846B27B83DF6F5F91B339050E20797755E0DF23537FAF2AE69770E1859B6BC06}.
\]
The next compulsory companion audit remains V90.

% =============================================================================
% END APPENDED EIGHTH-SERIES V87 MATERIAL
% =============================================================================


% =============================================================================
% BEGIN APPENDED EIGHTH-SERIES V88 MATERIAL
% =============================================================================

\clearpage
\section{Eighth-series V88: an upstream Combes--Thomas route}
\label{sec:v8v88}

\subsection{Local regulator operators and exponential conjugation}
\label{sec:e8v88-setting}

Let \(X_R\) be a finite metric regulator set with finite-dimensional fibres
and Hilbert space
\[
 {\cal H}_R=\ell^2\mathord{\bigoplus}_{x\in X_R}E_x.
\]
The Hilbert structure may already contain cell-volume and mass-matrix
weights; all statements below are invariant under the corresponding
unitary change to orthonormal cell coordinates.  For \(A\subset X_R\), let
\(\chi_A\) denote the orthogonal projection onto the fibres over \(A\).

For a real function \(\varphi\) on \(X_R\), let
\[
 W_{\alpha,\varphi}u(x):=e^{\alpha\varphi(x)}u(x).
\]
Write \({\rm Lip}_1(X_R)\) for the functions satisfying
\(|\varphi(x)-\varphi(y)|\leq d(x,y)\).  For a block operator \(H_R\), define
the exponential conjugation debit
\begin{equation}
 \delta_R(\alpha)
 :=
 \sup_{\varphi\in{\rm Lip}_1(X_R)}
 \left\|
 W_{\alpha,\varphi}H_RW_{\alpha,\varphi}^{-1}-H_R
 \right\|_{{\cal B}({\cal H}_R)}.                      \label{eq:e8v88-delta}
\end{equation}

\begin{lemma}[Incidence bound for the conjugation debit]
\label{lem:e8v88-incidence}
In orthonormal cell coordinates, put
\begin{equation}
 \kappa_R(\alpha)
 :=
 \max\left\{
 \sup_x\sum_y
 (e^{\alpha d(x,y)}-1)\|(H_R)_{xy}\|,
 \sup_y\sum_x
 (e^{\alpha d(x,y)}-1)\|(H_R)_{xy}\|
 \right\}.                                             \label{eq:e8v88-kappa}
\end{equation}
Then
\begin{equation}
 \delta_R(\alpha)\leq\kappa_R(\alpha).                 \label{eq:e8v88-delta-kappa}
\end{equation}
If \(H_R\) has propagation at most \(\ell_R\) and its off-diagonal
two-sided incidence is at most \(h_R\), then
\begin{equation}
 \delta_R(\alpha)
 \leq (e^{\alpha\ell_R}-1)h_R.                         \label{eq:e8v88-finite-range}
\end{equation}
\end{lemma}

\begin{proof}
The \(xy\) block of the conjugation difference is
\[
 \left(e^{\alpha(\varphi(x)-\varphi(y))}-1\right)
 (H_R)_{xy}.
\]
The Lipschitz property bounds its norm by
\((e^{\alpha d(x,y)}-1)\|(H_R)_{xy}\|\).
The two-sided Schur test gives
\eqref{eq:e8v88-delta-kappa}.  Under finite propagation the exponential
factor is at most \(e^{\alpha\ell_R}-1\); factor it from the row and column
sums to obtain \eqref{eq:e8v88-finite-range}.
\end{proof}

\subsection{Accretive Combes--Thomas estimate}
\label{sec:e8v88-CT}

\begin{theorem}[Finite-regulator Combes--Thomas bound]
\label{thm:e8v88-CT}
Let \(H_R\) be a square operator on \({\cal H}_R\) satisfying
\begin{equation}
 \operatorname{Re}\langle H_Ru,u\rangle
 \geq\gamma\|u\|^2
 \qquad(u\in{\cal H}_R)                                \label{eq:e8v88-accretive}
\end{equation}
for some \(\gamma>0\).  If
\begin{equation}
 \delta_R(\alpha)<\gamma,                              \label{eq:e8v88-CT-condition}
\end{equation}
then for all \(A,B\subset X_R\),
\begin{equation}
 \|\chi_AH_R^{-1}\chi_B\|
 \leq
 {e^{-\alpha d(A,B)}\over
  \gamma-\delta_R(\alpha)}.                            \label{eq:e8v88-CT}
\end{equation}
The estimate does not contain the number of cells in \(A\), \(B\), or
\(X_R\).
\end{theorem}

\begin{proof}
Accretivity implies \(\|H_Ru\|\geq\gamma\|u\|\).
Since the space is finite dimensional and \(H_R\) is square, it is
invertible and \(\|H_R^{-1}\|\leq\gamma^{-1}\).

Fix \(B\) and set
\(\varphi(x)=d(x,B)\).  This function is one-Lipschitz and vanishes on
\(B\).  Put
\[
 H_{R,\alpha}:=W_{\alpha,\varphi}H_RW_{\alpha,\varphi}^{-1}
 =H_R+K_{\alpha,\varphi}.
\]
By definition,
\(\|K_{\alpha,\varphi}\|\leq\delta_R(\alpha)<\gamma\).
The factorization
\[
 H_{R,\alpha}
 =H_R(I+H_R^{-1}K_{\alpha,\varphi})
\]
and its Neumann series give
\begin{equation}
 \|H_{R,\alpha}^{-1}\|
 \leq {1\over\gamma-\delta_R(\alpha)}.                 \label{eq:e8v88-conjugate-inverse}
\end{equation}
Because \(W_{\alpha,\varphi}\chi_B=\chi_B\) and
\(\|\chi_AW_{\alpha,\varphi}^{-1}\|
\leq e^{-\alpha d(A,B)}\),
\begin{align*}
 \chi_AH_R^{-1}\chi_B
 &=
 \chi_AW_{\alpha,\varphi}^{-1}
 H_{R,\alpha}^{-1}
 W_{\alpha,\varphi}\chi_B,
\end{align*}
which proves \eqref{eq:e8v88-CT}.
\end{proof}

\begin{remark}[Countable regulators]
\label{rem:e8v88-countable}
On a countable locally finite regulator, use
\(\varphi_L(x)=\min\{d(x,B),L\}\).  If the conjugation debit and accretive
floor are uniform in \(L\), the same proof gives
\eqref{eq:e8v88-CT} first for finitely supported vectors and then by strong
operator closure.
\end{remark}

\begin{corollary}[Uniform finite-range criterion]
\label{cor:e8v88-finite-range}
If
\[
 \inf_R\gamma_R\geq\gamma_*>0,\qquad
 \sup_R\ell_R\leq\ell_*,\qquad
 \sup_Rh_R\leq h_*<\infty,
\]
then every \(\alpha>0\) satisfying
\begin{equation}
 (e^{\alpha\ell_*}-1)h_*<\gamma_*                     \label{eq:e8v88-alpha}
\end{equation}
gives the regulator-uniform estimate
\[
 \|\chi_AH_R^{-1}\chi_B\|
 \leq
 {e^{-\alpha d(A,B)}\over
  \gamma_*-(e^{\alpha\ell_*}-1)h_*}.
\]
\end{corollary}

\begin{proof}
Use \Cref{lem:e8v88-incidence,thm:e8v88-CT}.
\end{proof}

\begin{firewall}[A shrinking exponent is not continuum locality]
\label{fire:e8v88-shrinking}
For a nearest-neighbour elliptic matrix, the raw hopping incidence can grow
as the mesh is refined.  Then \eqref{eq:e8v88-alpha} may force
\(\alpha=\alpha_R\to0\).  The theorem remains correct, but it gives no fixed
physical decay rate.  A refinement-uniform result requires either a
form-level Agmon estimate, heat-kernel gradient bounds, or a normalization
in which the conjugation debit is uniformly smaller than the coercive
floor.
\end{firewall}

\subsection{Localized trace-class determinant response}
\label{sec:e8v88-trace}

The operator-norm localization above can replace the Schur sandwich of V86
when the source insertion has a uniform trace norm.

\begin{theorem}[Local trace-class determinant]
\label{thm:e8v88-localdet}
Let \(D_s=D_0+sV\), \(0\leq s\leq1\), be finite-regulator square operators.
Assume
\begin{align}
 \operatorname{Re}\langle D_su,u\rangle
 &\geq\gamma\|u\|^2,                                  \label{eq:e8v88-source-floor}\\
 V&=\chi_AV\chi_A,\qquad \|V\|_1<\infty.              \label{eq:e8v88-trace-source}
\end{align}
Then the continuously normalized determinant response
\begin{equation}
 {\cal F}_{D_0}(V)
 :=
 \int_0^1\operatorname{Tr}(D_s^{-1}V)\,ds              \label{eq:e8v88-F}
\end{equation}
satisfies
\begin{align}
 \exp{\cal F}_{D_0}(V)
 &={\det(D_0+V)\over\det D_0},                         \label{eq:e8v88-F-ratio}\\
 |{\cal F}_{D_0}(V)|
 &\leq{\|V\|_1\over\gamma}.                            \label{eq:e8v88-F-bound}
\end{align}
The same conclusion holds on an infinite regulator with the ordinary
determinants replaced by
\(\det_{\rm F}(I+D_0^{-1}V)\).
\end{theorem}

\begin{proof}
The accretive floor makes every \(D_s\) invertible with
\(\|D_s^{-1}\|\leq\gamma^{-1}\).  Jacobi's determinant formula gives
\[
 {d\over ds}\log\det D_s
 =\operatorname{Tr}(D_s^{-1}V).
\]
Integration on the branch fixed at \(s=0\) proves
\eqref{eq:e8v88-F-ratio}.  Trace duality gives
\[
 |\operatorname{Tr}(D_s^{-1}V)|
 \leq\|D_s^{-1}\|\|V\|_1
 \leq\gamma^{-1}\|V\|_1,
\]
and integration proves \eqref{eq:e8v88-F-bound}.  In infinite dimension,
\(D_0^{-1}V\) is trace class and the same derivative formula is the
Fredholm determinant formula.
\end{proof}

\begin{theorem}[Remote perturbation of the local determinant]
\label{thm:e8v88-remote}
Let
\[
 D'_s=D_s+S,\qquad
 V=\chi_AV\chi_A,\qquad
 S=\chi_BS\chi_B,
\]
where \(d(A,B)\geq L\).  Suppose every \(D_s,D'_s\) has the common
accretive floor \(\gamma>0\) and
\begin{equation}
 \max\{\delta_{D_s}(\alpha),\delta_{D'_s}(\alpha)\}
 \leq\delta<\gamma                                    \label{eq:e8v88-uniform-debit}
\end{equation}
for \(0\leq s\leq1\).  Then
\begin{equation}
 |{\cal F}_{D'_0}(V)-{\cal F}_{D_0}(V)|
 \leq
 {e^{-2\alpha L}\|S\|\|V\|_1\over(\gamma-\delta)^2}.
 \label{eq:e8v88-remote}
\end{equation}
\end{theorem}

\begin{proof}
The resolvent identity gives
\[
 (D'_s)^{-1}-D_s^{-1}
 =-(D'_s)^{-1}SD_s^{-1}.
\]
Using the support projections and cyclicity of the trace,
\begin{align*}
 &\left|
 \operatorname{Tr}\left(
 [(D'_s)^{-1}-D_s^{-1}]V
 \right)\right|\\
 &\quad=
 \left|
 \operatorname{Tr}\left(
 \chi_A(D'_s)^{-1}\chi_B
 S
 \chi_BD_s^{-1}\chi_A
 V
 \right)\right|\\
 &\quad\leq
 \|\chi_A(D'_s)^{-1}\chi_B\|
 \|S\|
 \|\chi_BD_s^{-1}\chi_A\|
 \|V\|_1.
\end{align*}
\Cref{thm:e8v88-CT} bounds each off-diagonal resolvent by
\(e^{-\alpha L}/(\gamma-\delta)\).  Integrate in \(s\).
\end{proof}

\begin{corollary}[Cauchy response under remote shells]
\label{cor:e8v88-shells}
For adjacent remote changes \(S_n\) at distances \(L_n\), under common
\(\gamma,\delta,\alpha\), the local determinant response is Cauchy whenever
\begin{equation}
 \sum_n e^{-2\alpha L_n}\|S_n\|<\infty.                \label{eq:e8v88-shell-sum}
\end{equation}
\end{corollary}

\begin{proof}
Apply \Cref{thm:e8v88-remote} and telescope.
\end{proof}

\subsection{Application to the coupled constraint derivative}
\label{sec:e8v88-constraint}

Let \(D_R\) be the normalized scalar--vector quotient matrix of V87,
represented on the regulator Hilbert space associated with its weak
bilinear form.

\begin{theorem}[Conditional coupled Combes--Thomas card]
\label{thm:e8v88-coupled}
Suppose:
\begin{enumerate}[label=\textup{(\roman*)},leftmargin=3.8em]
\item the discrete V81 form satisfies
\[
 \operatorname{Re}{\mathbb b}_R(U,U)
 \geq\gamma_*\|U\|_{{\cal H}_R}^2
\]
with one \(\gamma_*>0\);
\item the assembled scalar, vector, and cross stencils are local in the
regulator metric and have
\[
 \sup_R\delta_{D_R}(\alpha)<\gamma_*
\]
for one physical exponent \(\alpha>0\).
\end{enumerate}
Then
\begin{equation}
 \|\chi_AD_R^{-1}\chi_B\|
 \leq
 {e^{-\alpha d(A,B)}\over
  \gamma_*-\sup_R\delta_{D_R}(\alpha)}.                \label{eq:e8v88-coupled-CT}
\end{equation}
This estimate treats the complete first-order scalar--vector coupling and
does not require the separate weighted product condition of V87.
\end{theorem}

\begin{proof}
The first hypothesis is the accretive floor required in
\eqref{eq:e8v88-accretive}; the second is
\eqref{eq:e8v88-CT-condition}.  Apply
\Cref{thm:e8v88-CT} directly to the complete quotient operator \(D_R\).
\end{proof}

\begin{remark}[Relation to V81]
\label{rem:e8v88-V81}
The V81 bound
\[
 {\mathbb b}_R(U,U)\geq\gamma_*\|U\|_{H^1\times H^1}^2
\]
implies the first hypothesis after a uniform embedding of the discrete
\({\cal H}_R\)-norm into the discrete \(H^1\)-norm.  It says nothing by
itself about the conjugation debit in the second hypothesis.
\end{remark}

\begin{firewall}[Quotient projections are not inserted into the stencil]
\label{fire:e8v88-quotient}
The constant and conformal-Killing modes must be split by the V83
finite-dimensional determinant identity before the locality calculation.
Multiplying the local stencil by global quotient projectors would reinsert
the nonlocal kernels diagnosed in V87 and can destroy the conjugation-debit
bound.
\end{firewall}

\subsection{Cell weights and ultraviolet trace ideals}
\label{sec:e8v88-cell}

\begin{proposition}[Representation invariance of the support estimate]
\label{prop:e8v88-cell}
Let \(G_R>0\) be a block-diagonal cell mass matrix and let
\({\cal U}_R=G_R^{1/2}\) identify the weighted cell Hilbert space with
orthonormal coordinates.  If
\[
 \widehat D_R={\cal U}_RD_R{\cal U}_R^{-1},
 \qquad
 \widehat\chi_A={\cal U}_R\chi_A{\cal U}_R^{-1},
\]
then
\begin{equation}
 \|\chi_AD_R^{-1}\chi_B\|_{{\cal H}_R}
 =
 \|\widehat\chi_A\widehat D_R^{-1}
   \widehat\chi_B\|_{\ell^2}.                           \label{eq:e8v88-unitary}
\end{equation}
The trace and trace norm of a source transform in the same unitary-invariant
way.
\end{proposition}

\begin{proof}
\({\cal U}_R\) is unitary by definition of the weighted inner product.
Operator norm, trace, and singular values are invariant under unitary
conjugation.
\end{proof}

\begin{consequence}[Correct ultraviolet target]
\label{cons:e8v88-UV}
Cell-count growth does not enter
\eqref{eq:e8v88-coupled-CT}.  It enters through the two genuinely analytic
quantities
\[
 \delta_{D_R}(\alpha)
 \quad\hbox{and}\quad
 \|V_{A,R}\|_1.
\]
A refinement-uniform local determinant theorem therefore requires both
\begin{equation}
 \sup_R\delta_{D_R}(\alpha)<\gamma_*,
 \qquad
 \sup_R\|V_{A,R}\|_1<\infty.                           \label{eq:e8v88-two-targets}
\end{equation}
\end{consequence}

\begin{proof}
The first quantity is the only regulator-dependent input in
\Cref{thm:e8v88-CT}; the second is the only source-size input in
\Cref{thm:e8v88-localdet,thm:e8v88-remote}.
\Cref{prop:e8v88-cell} shows that neither can be changed by merely
renaming weighted cell coordinates.
\end{proof}

\begin{firewall}[Metric sources fail the raw trace-class row]
\label{fire:e8v88-metric}
A local metric variation changes the principal symbol.  V83 showed that in
three dimensions the relative order-zero perturbation is not compact and
the order-\(-2\) potential response is generally Hilbert--Schmidt rather
than trace class.  Hence \(\sup_R\|V_{A,R}\|_1<\infty\) is not available for
an unsmoothed physical metric source.  The trace-class theorem applies to
sufficiently smoothing local observables; the metric determinant requires
a modified determinant and local counterterm subtraction.
\end{firewall}

\subsection{Upstream result and next branch}
\label{sec:e8v88-result}

\begin{assessment}[Progress at V88]
\label{ass:e8v88-status}
The weighted locality problem has been separated into a support estimate
and a trace-ideal estimate.  A regulator-uniform accretive floor plus a
uniform exponential conjugation debit gives off-diagonal resolvent decay
for the complete coupled constraint operator, without a cell-count factor.
For a trace-class local source this yields a determinant response whose
remote-shell change is \(O(e^{-2\alpha L})\).  The physical metric source
does not satisfy the raw trace-class row, so the next step must use the
Hilbert--Schmidt modified determinant already indicated in V83.
\end{assessment}

\begin{programme}[Eighth-series V89: localized modified determinant]
\label{prog:e8v88-V89}
V89 will derive the exact differential and remote-shell identity for
\(\det_2(I+K)\) when the relative source is Hilbert--Schmidt.  It will show
where the subtracted linear trace becomes a local counterterm, identify the
products that are trace class, and state a summable localized
Hilbert--Schmidt criterion.  No ordinary determinant will be used for the
order-\(-2\) three-dimensional row.
\end{programme}

\begin{firewall}[Claim boundary after eighth-series V88]
\label{fire:e8v88-boundary}
V88 proves the finite-regulator accretive Combes--Thomas bound, its
localized trace-class determinant consequence, and the exact two
refinement-uniform targets.  It does not yet prove a fixed physical
conjugation exponent for a chosen discretization, control the physical
metric source in trace norm, construct its modified determinant and
counterterms, prove adjacent-density summability, or construct the full
SM--Einstein continuum.
\end{firewall}

\begin{theorem}[Eighth-series V88 result]
\label{thm:e8v88-result}
If the normalized coupled regulator derivative has a uniform accretive
floor \(\gamma_*\) and a uniform exponential conjugation debit
\(\delta(\alpha)<\gamma_*\), then
\[
 \|\chi_AD_R^{-1}\chi_B\|
 \leq {e^{-\alpha d(A,B)}\over\gamma_*-\delta(\alpha)}.
\]
For a uniformly trace-class source in \(A\), a perturbation in a region
\(B\) changes the local determinant response by at most
\[
 {e^{-2\alpha d(A,B)}\|S\|\|V_A\|_1
  \over(\gamma_*-\delta(\alpha))^2}.
\]
\end{theorem}

\begin{proof}
Use \Cref{thm:e8v88-CT,thm:e8v88-remote,
thm:e8v88-coupled}.
\end{proof}

\paragraph{Parent-version record.}
V88 was formed from
\texttt{Renewal\_SM\_Einstein\_Continuum\_Research\_Notes\_V87.tex}.
The V87 parent had \(35\,330\) logical source lines, \(1\,259\,435\) bytes,
and SHA--256
\[
 \texttt{F7FD3B6D5774D382D2B7ACF69996EF3E42EC7CB877767BC2185A6014A7122AFF}.
\]
The next compulsory companion audit remains V90.

% =============================================================================
% END APPENDED EIGHTH-SERIES V88 MATERIAL
% =============================================================================


% =============================================================================
% BEGIN APPENDED EIGHTH-SERIES V89 MATERIAL
% =============================================================================

\clearpage
\section{Eighth-series V89: localized Hilbert--Schmidt modified determinants}
\label{sec:v8v89}

\subsection{The correct determinant class}
\label{sec:e8v89-class}

Let \({\cal H}\) be a separable Hilbert space.  Write
\({\mathfrak S}_1\) and \({\mathfrak S}_2\) for the trace-class and
Hilbert--Schmidt ideals, with norms \(\|\cdot\|_1\) and
\(\|\cdot\|_2\).  Their ideal products obey
\begin{align}
 \|AB\|_1&\leq\|A\|_2\|B\|_2,                          \label{eq:e8v89-holder}\\
 \|CAB\|_1&\leq\|C\|\|A\|_2\|B\|_2.                   \label{eq:e8v89-holder-bounded}
\end{align}
For \(K\in{\mathfrak S}_2\), define
\begin{equation}
 \det\nolimits_2(I+K)
 :=
 \det_{\rm F}\bigl((I+K)e^{-K}\bigr).                 \label{eq:e8v89-det2}
\end{equation}
Indeed,
\((I+K)e^{-K}-I\in{\mathfrak S}_1\), because the constant and linear
terms cancel and every remaining term contains at least two
Hilbert--Schmidt factors.

\begin{firewall}[No raw trace of a Hilbert--Schmidt operator]
\label{fire:e8v89-rawtrace}
An operator in \({\mathfrak S}_2\) need not belong to
\({\mathfrak S}_1\), so \(\operatorname{Tr}K\) need not exist.
Formula \eqref{eq:e8v89-det2} defines one trace-class combination; it does
not license the separate symbols \(\det(I+K)\) or
\(\operatorname{Tr}K\).
\end{firewall}

\subsection{Differentiation without an illegal trace}
\label{sec:e8v89-differential}

\begin{theorem}[Differential of the modified determinant]
\label{thm:e8v89-differential}
Let \(t\mapsto K(t)\) be \(C^1\) in
\({\mathfrak S}_2\), and suppose \(I+K(t)\) is invertible on an interval.
Then a continuous logarithm of \(\det_2(I+K(t))\) satisfies
\begin{align}
 {d\over dt}\log\det\nolimits_2(I+K(t))
 &=
 \operatorname{Tr}\left(
 [(I+K(t))^{-1}-I]K'(t)
 \right)                                               \label{eq:e8v89-differential}\\
 &=
 -\operatorname{Tr}\left(
 (I+K(t))^{-1}K(t)K'(t)
 \right).                                              \label{eq:e8v89-differential2}
\end{align}
Both displayed trace arguments are trace class.
\end{theorem}

\begin{proof}
First suppose that \(K(t)\) has finite rank.  From
\[
 \log\det\nolimits_2(I+K)
 =
 \log\det(I+K)-\operatorname{Tr}K
\]
and Jacobi's formula,
\[
 {d\over dt}\log\det\nolimits_2(I+K)
 =
 \operatorname{Tr}\left((I+K)^{-1}K'-K'\right).
\]
The resolvent identity
\[
 (I+K)^{-1}-I=-(I+K)^{-1}K
\]
gives \eqref{eq:e8v89-differential2}.  Now
\(KK'\in{\mathfrak S}_1\) by
\eqref{eq:e8v89-holder}, so the right side is continuous in the
\({\mathfrak S}_2\) topology on every set with uniformly bounded inverse.
Approximate \(K,K'\) in \({\mathfrak S}_2\) by finite-rank paths, or
equivalently apply the standard Fredholm determinant differential to
\eqref{eq:e8v89-det2}.  Passing to the limit proves both formulas.
\end{proof}

\begin{corollary}[Quantitative path bound]
\label{cor:e8v89-path}
If
\[
 \sup_t\|(I+K(t))^{-1}\|\leq M,\qquad
 \sup_t\|K(t)\|_2\leq H,\qquad
 \sup_t\|K'(t)\|_2\leq J,
\]
on an interval of length \(L\), then
\begin{equation}
 \left|
 \log{\det_2(I+K(t_1))\over\det_2(I+K(t_0))}
 \right|
 \leq LMHJ.                                             \label{eq:e8v89-path-bound}
\end{equation}
\end{corollary}

\begin{proof}
Apply \eqref{eq:e8v89-holder-bounded} to
\eqref{eq:e8v89-differential2} and integrate.
\end{proof}

\subsection{A source-marked modified determinant}
\label{sec:e8v89-source}

Let \(K,V\in{\mathfrak S}_2\), and suppose
\[
 R_s:=(I+K+sV)^{-1}
\]
exists for \(0\leq s\leq1\).  Define
\begin{equation}
 {\cal G}^{(2)}_K(V)
 :=
 \log\det\nolimits_2(I+K+V)
 -\log\det\nolimits_2(I+K)                             \label{eq:e8v89-G}
\end{equation}
on the source path.

\begin{theorem}[Exact marked differential]
\label{thm:e8v89-marked}
The source response has the trace-class representation
\begin{align}
 {\cal G}^{(2)}_K(V)
 &=
 \int_0^1
 \operatorname{Tr}[(R_s-I)V]\,ds                       \label{eq:e8v89-marked}\\
 &=
 -\int_0^1
 \operatorname{Tr}[R_s(K+sV)V]\,ds.                    \label{eq:e8v89-marked2}
\end{align}
If \(\sup_s\|R_s\|\leq M\), then
\begin{equation}
 |{\cal G}^{(2)}_K(V)|
 \leq
 M(\|K\|_2+\|V\|_2)\|V\|_2.                           \label{eq:e8v89-marked-bound}
\end{equation}
\end{theorem}

\begin{proof}
Apply \Cref{thm:e8v89-differential} to \(K(s)=K+sV\).
The two trace-class claims follow because
\((K+sV)V\in{\mathfrak S}_1\).  The ideal H\"older inequality gives
\[
 |\operatorname{Tr}R_s(K+sV)V|
 \leq M(\|K\|_2+s\|V\|_2)\|V\|_2.
\]
Integrate and use the slightly weaker displayed bound.
\end{proof}

\begin{proposition}[Anchored Hilbert--Schmidt improvement]
\label{prop:e8v89-anchored}
Let \(\chi_A\) be an orthogonal support projection and suppose
\[
 V=\chi_AV\chi_A.
\]
Then \eqref{eq:e8v89-marked-bound} improves to
\begin{equation}
 |{\cal G}^{(2)}_K(V)|
 \leq
 M\bigl(\|K\chi_A\|_2+\|V\|_2\bigr)\|V\|_2.           \label{eq:e8v89-local-bound}
\end{equation}
Thus an extensive global Hilbert--Schmidt norm of \(K\) may be replaced by
the column-localized norm at the source anchor.
\end{proposition}

\begin{proof}
Because \(V=\chi_AV\),
\[
 (K+sV)V=(K\chi_A+sV)V.
\]
Apply \eqref{eq:e8v89-holder-bounded} and integrate as in
\Cref{thm:e8v89-marked}.
\end{proof}

\begin{firewall}[The anchor norm must survive refinement]
\label{fire:e8v89-anchor}
The term \(\emph{localized}\) in
\eqref{eq:e8v89-local-bound} is an operator-ideal statement, not a
cardinality statement.  For a fixed physical region, the number of
regulator cells grows.  One must prove uniform bounds for
\(\|K_R\chi_{A,R}\|_2\) and \(\|V_{A,R}\|_2\) in the cell mass-matrix
normalization.
\end{firewall}

\subsection{Remote-shell identity in the Hilbert--Schmidt class}
\label{sec:e8v89-remote}

\begin{theorem}[Exact remote variation identity]
\label{thm:e8v89-remote-identity}
Let \(S\in{\mathfrak S}_2\), set
\[
 R'_s=(I+K+S+sV)^{-1},
\]
and suppose both source paths are invertible.  Then
\begin{equation}
 {\cal G}^{(2)}_{K+S}(V)-{\cal G}^{(2)}_K(V)
 =
 -\int_0^1
 \operatorname{Tr}(R'_s S R_s V)\,ds.                 \label{eq:e8v89-remote-identity}
\end{equation}
If
\(\sup_s\|R_s\|\leq M\) and
\(\sup_s\|R'_s\|\leq M'\), then
\begin{equation}
 |{\cal G}^{(2)}_{K+S}(V)-{\cal G}^{(2)}_K(V)|
 \leq MM'\|S\|_2\|V\|_2.                              \label{eq:e8v89-remote-global}
\end{equation}
\end{theorem}

\begin{proof}
Use the first line of \eqref{eq:e8v89-marked} for each background.  The
identity terms cancel:
\[
 [(R'_s-I)-(R_s-I)]V=(R'_s-R_s)V.
\]
The resolvent identity
\[
 R'_s-R_s=-R'_sSR_s
\]
proves \eqref{eq:e8v89-remote-identity}.  The product
\((R'_sS)(R_sV)\) is trace class, and
\[
 |\operatorname{Tr}(R'_sSR_sV)|
 \leq M'M\|S\|_2\|V\|_2.
\]
Integrate.
\end{proof}

\begin{theorem}[Exponentially localized modified-determinant response]
\label{thm:e8v89-remote-local}
Assume
\[
 V=\chi_AV\chi_A,\qquad
 S=\chi_BS\chi_B,\qquad d(A,B)\geq L.
\]
Suppose the two source-path resolvents obey
\begin{align}
 \|\chi_AR'_s\chi_B\|&\leq C'e^{-\alpha L},&
 \|\chi_BR_s\chi_A\|&\leq Ce^{-\alpha L}               \label{eq:e8v89-offdiag}
\end{align}
uniformly for \(0\leq s\leq1\).  Then
\begin{equation}
 |{\cal G}^{(2)}_{K+S}(V)-{\cal G}^{(2)}_K(V)|
 \leq
 CC'e^{-2\alpha L}\|S\|_2\|V\|_2.                     \label{eq:e8v89-remote-local}
\end{equation}
\end{theorem}

\begin{proof}
Insert the support projections in
\eqref{eq:e8v89-remote-identity} and use cyclicity:
\[
 \operatorname{Tr}(R'_sSR_sV)
 =
 \operatorname{Tr}\left[
 (\chi_AR'_s\chi_BS)
 (\chi_BR_s\chi_AV)
 \right].
\]
The first parenthesis is Hilbert--Schmidt with norm at most
\(C'e^{-\alpha L}\|S\|_2\); the second has norm at most
\(Ce^{-\alpha L}\|V\|_2\).  Their product is trace class.  Apply
\eqref{eq:e8v89-holder} and integrate.
\end{proof}

\begin{corollary}[Summable localized Hilbert--Schmidt shells]
\label{cor:e8v89-shells}
Let \(K_{n+1}=K_n+S_n\).  Suppose the uniform off-diagonal bounds
\eqref{eq:e8v89-offdiag} hold for one anchored source \(V\), with shell
distances \(L_n\).  If
\begin{equation}
 \sum_n e^{-2\alpha L_n}\|S_n\|_2<\infty,              \label{eq:e8v89-shellsum}
\end{equation}
then \({\cal G}^{(2)}_{K_n}(V)\) is Cauchy.
\end{corollary}

\begin{proof}
Apply \Cref{thm:e8v89-remote-local} to adjacent backgrounds and telescope.
\end{proof}

\subsection{Subtracted linear trace and counterterm meaning}
\label{sec:e8v89-subtraction}

\begin{proposition}[Finite-dimensional comparison with the ordinary ratio]
\label{prop:e8v89-finite}
In finite dimension,
\begin{equation}
 {\cal G}^{(2)}_K(V)
 =
 \log{\det(I+K+V)\over\det(I+K)}
 -\operatorname{Tr}V.                                 \label{eq:e8v89-finite}
\end{equation}
The first-order loop in \(V\) is subtracted exactly.
\end{proposition}

\begin{proof}
Use
\(\det_2(I+L)=\det(I+L)e^{-\operatorname{Tr}L}\) at
\(L=K+V\) and \(L=K\), then divide.
\end{proof}

\begin{remark}[Local counterterm interpretation]
\label{rem:e8v89-counterterm}
At finite regulator, \(\operatorname{Tr}V\) is the linear determinant
counterterm removed by \(\det_2\).  In a continuum Hilbert--Schmidt limit it
may not exist separately; only the combined expression
\eqref{eq:e8v89-marked2} is defined.  Any heat-kernel realization of this
subtraction must be local and must obey the chart and regulator cocycles of
V84.
\end{remark}

\subsection{Multiplicative anomaly and the regulator cocycle}
\label{sec:e8v89-anomaly}

\begin{theorem}[Modified-determinant product formula]
\label{thm:e8v89-anomaly}
Let \(A,B\in{\mathfrak S}_2\).  Then \(AB\in{\mathfrak S}_1\), and
\begin{equation}
 \det\nolimits_2\bigl((I+A)(I+B)\bigr)
 =
 \det\nolimits_2(I+A)\det\nolimits_2(I+B)
 e^{-\operatorname{Tr}(AB)}.                          \label{eq:e8v89-anomaly}
\end{equation}
\end{theorem}

\begin{proof}
For finite-rank \(A,B\),
\begin{align*}
 \det\nolimits_2((I+A)(I+B))
 &=
 \det((I+A)(I+B))
 e^{-\operatorname{Tr}(A+B+AB)}\\
 &=
 \det\nolimits_2(I+A)\det\nolimits_2(I+B)
 e^{-\operatorname{Tr}(AB)}.
\end{align*}
Finite-rank operators are dense in \({\mathfrak S}_2\), the product map
\({\mathfrak S}_2\times{\mathfrak S}_2\to{\mathfrak S}_1\) is continuous
by \eqref{eq:e8v89-holder}, and \(\det_2\) is continuous in the
Hilbert--Schmidt norm.  Approximation proves the result.
\end{proof}

\begin{corollary}[Relative factorization anomaly]
\label{cor:e8v89-relative-anomaly}
If \(I+K\) is invertible and \(S\in{\mathfrak S}_2\), put
\[
 Q:=(I+K)^{-1}S.
\]
Then
\begin{equation}
 {\det_2(I+K+S)\over\det_2(I+K)}
 =
 \det\nolimits_2(I+Q)
 \exp\!\left[-\operatorname{Tr}(KQ)\right].            \label{eq:e8v89-relative-anomaly}
\end{equation}
\end{corollary}

\begin{proof}
The exact factorization is
\[
 I+K+S=(I+K)(I+Q).
\]
Both \(K,Q\) are Hilbert--Schmidt and \(KQ\) is trace class.  Apply
\Cref{thm:e8v89-anomaly}.
\end{proof}

\begin{firewall}[Pairwise modified ratios do not multiply naively]
\label{fire:e8v89-anomaly}
The endpoint difference
\[
 \log\det_2(I+K_m)-\log\det_2(I+K_n)
\]
telescopes exactly.  A product assembled instead from separately
factorized relative \(\det_2(I+Q_n)\) terms must include every anomaly
\(-\operatorname{Tr}(K_nQ_n)\).  Dropping these trace-class terms violates
the V84 regulator cocycle and changes the effective counterterm.
\end{firewall}

\subsection{Insertion into the ACDC determinant row}
\label{sec:e8v89-ACDC}

\begin{definition}[Localized modified-determinant card, LMDC]
\label{def:e8v89-LMDC}
For a lower-order local constraint perturbation, LMDC records:
\begin{enumerate}[label=\textup{(M\arabic*)},leftmargin=4em]
\item a common reference operator and the relative kernels
\(K_R,V_{A,R}\in{\mathfrak S}_2\);
\item a uniform inverse bound along the source path;
\item the localized ideal bounds
\(\|K_R\chi_{A,R}\|_2\) and \(\|V_{A,R}\|_2\);
\item the modified response \({\cal G}^{(2)}_{K_R}(V_{A,R})\);
\item the finite-regulator linear subtraction and its local
heat-kernel/counterterm interpretation;
\item the Combes--Thomas off-diagonal bound for remote shells;
\item the summability row \eqref{eq:e8v89-shellsum};
\item every multiplicative anomaly
\eqref{eq:e8v89-relative-anomaly}; and
\item compatibility with the chart and regulator cocycles of V84.
\end{enumerate}
\end{definition}

\begin{theorem}[Conditional lower-order determinant closure]
\label{thm:e8v89-closure}
Suppose LMDC holds uniformly along a cofinal regulator sequence and the
localized kernels converge in \({\mathfrak S}_2\) on every fixed source
anchor.  Then the lower-order source-marked modified determinant has a
well-defined local limit.  Remote-shell changes are summable under
\eqref{eq:e8v89-shellsum}, and changing the relative factorization preserves
the regulator cocycle precisely when the anomaly terms in
\eqref{eq:e8v89-relative-anomaly} are retained.
\end{theorem}

\begin{proof}
Continuity of \(\det_2\) in \({\mathfrak S}_2\), together with the uniform
inverse path, gives convergence of the fixed-anchor endpoint response.
\Cref{cor:e8v89-shells} controls the remote tail.
\Cref{thm:e8v89-anomaly,cor:e8v89-relative-anomaly} give the exact
factorization law, while endpoint differences telescope by definition.
\end{proof}

\begin{firewall}[Principal metric variations remain outside LMDC]
\label{fire:e8v89-principal}
A metric variation that changes the order-two principal symbol produces an
order-zero relative operator and is generally not compact.  Neither
\(\det_2\) nor a higher Schatten determinant repairs a noncompact relative
kernel.  Such a variation requires a matched principal reference, local
symbol subtraction, or a determinant-line construction before LMDC can be
applied to its lower-order remainder.
\end{firewall}

\begin{assessment}[Progress at V89]
\label{ass:e8v89-status}
The determinant row now has a rigorous Hilbert--Schmidt branch.  The
\(\det_2\) differential combines two individually illegal traces into the
trace-class product
\(-(I+K)^{-1}KK'\).  A local source needs only localized
Hilbert--Schmidt control, and the V88 off-diagonal resolvent estimate makes
its response to a remote Hilbert--Schmidt shell exponentially small.  The
modified determinant has a trace-class multiplicative anomaly, which must
be carried as part of the V84 counterterm cocycle.
\end{assessment}

\begin{programme}[Eighth-series V90: compulsory companion audit]
\label{prog:e8v89-V90}
V90 will inspect the newest YM and NS notes.  It will compare their current
high-occupation, connected-kernel, and concentration estimates with LMDC,
especially the localized ideal norm, remote-shell summability, and anomaly
rows.  The programme will then continue in V91 with the principal-symbol
determinant branch selected by that audit.
\end{programme}

\begin{firewall}[Claim boundary after eighth-series V89]
\label{fire:e8v89-boundary}
V89 proves the \(\det_2\) differential, localized source and remote-shell
estimates, and the multiplicative anomaly under explicit
Hilbert--Schmidt hypotheses.  It does not verify LMDC for a chosen
SM--Einstein discretization, treat principal metric variations, construct
the required local symbol counterterms, prove the full adjacent-density
bound, or construct the full SM--Einstein continuum.
\end{firewall}

\begin{theorem}[Eighth-series V89 result]
\label{thm:e8v89-result}
For Hilbert--Schmidt relative kernels, the source-marked modified
determinant is controlled by trace-class products.  Under the V88
off-diagonal resolvent estimate, a Hilbert--Schmidt shell a distance \(L\)
from the source changes the response by at most
\[
 CC'e^{-2\alpha L}\|S\|_2\|V\|_2.
\]
Relative \(\det_2\) factors obey the regulator cocycle only after inclusion
of the anomaly \(-\operatorname{Tr}(KQ)\).
\end{theorem}

\begin{proof}
Use \Cref{thm:e8v89-marked,thm:e8v89-remote-local,
thm:e8v89-anomaly,cor:e8v89-relative-anomaly}.
\end{proof}

\paragraph{Parent-version record.}
V89 was formed from
\texttt{Renewal\_SM\_Einstein\_Continuum\_Research\_Notes\_V88.tex}.
The V88 parent had \(35\,834\) logical source lines, \(1\,275\,925\) bytes,
and SHA--256
\[
 \texttt{FD3102C30BB1867820BFAB5FFEC7155C3B3039C83D14E1044664F057D2A31A50}.
\]
The next version is the compulsory V90 companion audit.

% =============================================================================
% END APPENDED EIGHTH-SERIES V89 MATERIAL
% =============================================================================


% =============================================================================
% BEGIN APPENDED EIGHTH-SERIES V90 MATERIAL
% =============================================================================

\clearpage
\section{Eighth-series V90: compulsory companion audit}
\label{sec:v8v90}

\subsection{Audit scope and files}
\label{sec:e8v90-files}

This checkpoint compares the lower-order modified-determinant branch of
V89 with the newest companion work.  The inspected files were
\begin{align*}
 &\texttt{Renewal\_Yang\_Mills\_Mass\_Gap\_Research\_Notes\_V47.tex},\\
 &\texttt{Renewal\_NS\_Stability\_Research\_Notes\_V26.tex}.
\end{align*}
The YM file had \(20\,001\) logical source lines, \(726\,337\) bytes, and
SHA--256
\[
 \texttt{79216C07B0364D36CFC04CE255AC33F6B69590FEE75E11AB4C0EB52549E5B0B3}.
\]
The NS file had \(5\,319\) logical source lines, \(278\,283\) bytes, and
SHA--256
\[
 \texttt{82C887F8C2C69E4F28FAD7C3DC8DE5B9099CB5A4BF431EBA1FFF3E91935978AD}.
\]
The NS file is unchanged from V85.  The substantive new companion material
is YM V42--V47.

\subsection{Yang--Mills V42--V47}
\label{sec:e8v90-YM}

The new YM sequence establishes or isolates the following points.
\begin{enumerate}[label=\textup{(Y\arabic*)},leftmargin=4.5em]
\item V42 controls a high-occupation complement under an explicit
relative-form hypothesis.  Its residual floor is of order \(M/a\), and the
occupation-window mixing is \(O(gM^{1/2})\).  Closed returns are retained as
self-energy terms.  Successive momentum and occupation Schur eliminations
agree with direct elimination.
\item V43 proves the required relative-form card on a fixed contractible
gauge-reduced block.  The relative coefficient is \(O(gR)\), while the IMS
and density remainder is
\(O(R^{-2}+g^2R^2)\).  Torons and boundary holonomies remain slow variables.
\item V44 rewrites connected correlations as logarithmic derivatives of a
normalized generating function.  Product-reference contributions cancel,
but a volume-uniform all-support result still requires a spatial tree
estimate for anchored activities.
\item V45 separates genuine residual tails from low effective terms.  It
constructs a finite buffered cover of compact slow-label space and retains
transition modes when spectral rank changes on overlaps.
\item V46 proves a volume-uniform Combes--Thomas estimate for a finite-range
quadratic kernel with a common spectral buffer.  Contour functional
calculus gives exponentially localized high-band projections, inverse
square roots, and slow-label derivatives.  On patch overlaps, common
refinement enlarges the slow sector and preserves the residual floor.
\item V47 proves a zero-radius obstruction for the untruncated quartic
Gaussian expansion.  It replaces that expansion by a bounded
occupation-window cumulant budget and returns the complement by a Feshbach
estimate.
\end{enumerate}

\begin{proposition}[Type-correct imports from the new YM work]
\label{prop:e8v90-import}
The following proof mechanisms transfer to the SM determinant programme.
\begin{enumerate}[label=\textup{(\roman*)},leftmargin=3.8em]
\item A spectral mode whose rank changes across parameter patches must be
retained in a finite slow sector on the overlap.
\item A high-sector inverse should be defined by a contour enclosing the
high spectrum, rather than as a difference involving a possibly nonlocal
slow projector.
\item Parameter derivatives of that contour inverse contain two resolvents
and one local operator derivative, the same product structure used in the
V89 trace-class response.
\item An asymptotic local heat or symbol expansion must not be promoted to a
convergent untruncated interaction series.
\end{enumerate}
The Yang--Mills occupation gap, Wilson relative form, and bounded-window
cumulant estimates do not prove a Schatten estimate or heat coefficient for
the augmented Einstein constraint operator.
\end{proposition}

\begin{proof}
Items (i)--(iii) are algebraic or functional-calculus statements about
buffered spectral subspaces and resolvent differentiation, so they apply
once the SM operator satisfies the corresponding spectral hypotheses.
Item (iv) is a logical restriction supplied by the explicit V47
zero-radius example.  The final sentence follows because the companion
operators, domains, and norms differ from those in the coarea determinant.
\end{proof}

\subsection{Navier--Stokes status}
\label{sec:e8v90-NS}

NS V26 remains a pointwise rank-one Oseen-kernel calculation.  It neither
changes nor closes the V89 localized Hilbert--Schmidt card.  Its certified
kernel constants concern a parabolic integral operator and do not supply
the elliptic normal-operator heat coefficients needed below.

\subsection{The principal-symbol branch}
\label{sec:e8v90-principal}

V89 handles relative kernels in \({\mathfrak S}_2\).  A metric variation of
the order-two principal symbol is order zero after preconditioning and is
generally noncompact.  The appropriate next object is therefore the exact
positive normal operator.

Let \(D_R\) be the normalized finite-regulator constraint derivative and
put
\begin{equation}
 H_R:=D_R^*D_R.                                        \label{eq:e8v90-H}
\end{equation}
If \(D_R\) is invertible, then
\begin{equation}
 \log|\det D_R|
 ={1\over2}\log\det H_R.                               \label{eq:e8v90-normal-det}
\end{equation}
Unlike a formal expansion of
\(\log(I+D_0^{-1}(D-D_0))\), the right side keeps the changed principal
symbol inside the exact heat semigroup.

\begin{lemma}[Finite-regulator normal determinant]
\label{lem:e8v90-normal}
For every invertible square matrix \(D_R\),
\eqref{eq:e8v90-normal-det} holds.
\end{lemma}

\begin{proof}
\[
 \det(D_R^*D_R)
 =\overline{\det D_R}\det D_R
 =|\det D_R|^2.
\]
Take the real logarithm.
\end{proof}

\begin{firewall}[Squaring changes the heat order]
\label{fire:e8v90-order}
The continuum normal operator \(D^*D\) has differential order four.
In spatial dimension three its small-time heat trace begins at order
\(t^{-3/4}\).  The determinant integral contains an additional factor
\(t^{-1}\), so all local heat terms through homogeneity zero must be
subtracted.  Reusing the second-order subtraction count would be wrong.
\end{firewall}

\subsection{Audit decision}
\label{sec:e8v90-decision}

\begin{decision}[V91 finite-part heat determinant]
\label{dec:e8v90-V91}
V91 will:
\begin{enumerate}[label=\textup{(H\arabic*)},leftmargin=4em]
\item define the exact relative heat trace
\[
 \Theta_{10}(t)
 =\operatorname{Tr}(e^{-tH_1}-e^{-tH_0});
\]
\item assume or derive its order-four, dimension-three local asymptotic
coefficients through homogeneity zero;
\item prove an explicit finite-part integral with power and logarithmic
subtractions;
\item prove additivity under \(H_0\to H_1\to H_2\) when one subtraction
scheme is used;
\item identify the subtraction coefficients as local endpoint
counterterms and incorporate them into the V84 cocycle; and
\item keep the heat expansion asymptotic, in accordance with the V47
zero-radius warning.
\end{enumerate}
\end{decision}

\begin{theorem}[V90 audit conclusion]
\label{thm:e8v90-result}
The companion audit selects an exact heat-semigroup treatment of principal
metric changes.  The finite slow sector must be refined on overlaps, the
high inverse may be defined by buffered contour calculus, and the
principal-symbol determinant must be renormalized by local finite-part heat
subtractions rather than by an untruncated power series.
\end{theorem}

\begin{proof}
Slow refinement and contour locality are items (Y5) and
\Cref{prop:e8v90-import}.  The noncompact principal response is the V89
claim boundary.  The exact positive determinant identity is
\Cref{lem:e8v90-normal}; the necessary order correction is
\Cref{fire:e8v90-order}; and the restriction on untruncated expansions is
item (Y6).
\end{proof}

\begin{assessment}[Progress at V90]
\label{ass:e8v90-status}
The compulsory audit is complete.  YM V46 corroborates the buffered
spectral and Combes--Thomas route while adding a precise overlap rule.  YM
V47 prevents a false analytic closure and changes the acquisition order:
the principal determinant will be defined by an exact heat integral with
only finitely many local asymptotic subtractions.
\end{assessment}

\begin{firewall}[Claim boundary after eighth-series V90]
\label{fire:e8v90-boundary}
V90 proves the normal-matrix determinant identity and records the
type-correct audit transfers.  It does not yet construct the continuum heat
parametrix, compute the SM--Einstein heat coefficients, prove their
regulator convergence, establish the complete ACDC density bound, or
construct the full SM--Einstein continuum.
\end{firewall}

\paragraph{Parent-version record.}
V90 was formed from
\texttt{Renewal\_SM\_Einstein\_Continuum\_Research\_Notes\_V89.tex}.
The V89 parent had \(36\,335\) logical source lines, \(1\,292\,463\) bytes,
and SHA--256
\[
 \texttt{04EDBB4A08013CAE8BA50A1DCE903E345AC5DB38F7BD76FF1DD5FFC2C24827B1}.
\]
The next compulsory companion audit is V95.  V91 now executes the
finite-part heat route selected above.

% =============================================================================
% END APPENDED EIGHTH-SERIES V90 MATERIAL
% =============================================================================


% =============================================================================
% BEGIN APPENDED EIGHTH-SERIES V91 MATERIAL
% =============================================================================

\clearpage
\section{Eighth-series V91: finite-part heat determinants for principal changes}
\label{sec:v8v91}

\subsection{The fourth-order normal operator}
\label{sec:e8v91-normal}

Let \(D\) be the normalized scalar--vector constraint derivative on the
fast quotient, and define
\begin{equation}
 H:=D^*D.                                               \label{eq:e8v91-H}
\end{equation}
The adjoint is taken in the declared field and equation Hilbert metrics.
At finite regulator,
\[
 \log|\det D|={1\over2}\log\det H.
\]
The continuum use of the right side requires the heat subtraction developed
below.

\begin{theorem}[Principal symbol of the normal operator]
\label{thm:e8v91-symbol}
Under the V87 matter-order hypothesis, \(H\) is strongly elliptic of order
four.  At a nonzero covector \(\xi\),
\begin{equation}
 \sigma_4(H)(x,\xi)
 =
 \begin{pmatrix}
 16|\xi|_g^4&0\\
 0&
 \left(
 |\xi|_g^2I+\frac13\xi^\sharp\otimes\xi
 \right)^2
 \end{pmatrix}.                                        \label{eq:e8v91-symbol}
\end{equation}
Its vector eigenvalues are \(|\xi|_g^4\) on
\(\xi^\perp\) and \(\frac{16}{9}|\xi|_g^4\) on
\(\mathbb R\xi^\sharp\).
\end{theorem}

\begin{proof}
V87 gives the diagonal second-order symbol
\[
 \sigma_2(D)(\xi)
 =
 \operatorname{diag}\left(
 4|\xi|^2,\,
 |\xi|^2I+\frac13\xi^\sharp\otimes\xi
 \right).
\]
For a differential operator,
\(\sigma_4(D^*D)=\sigma_2(D)^*\sigma_2(D)\).
The displayed symbol is self-adjoint, so this is its square.  Squaring the
scalar eigenvalue and the transverse/longitudinal vector eigenvalues gives
\eqref{eq:e8v91-symbol}.  All are positive for \(\xi\ne0\), proving strong
ellipticity.
\end{proof}

\begin{remark}[Lower-order cross terms remain exact]
\label{rem:e8v91-cross}
The first-order scalar--vector blocks contribute to the order-three and
lower symbols of \(H\).  They are not discarded.  They enter the exact heat
semigroup and its local coefficients, while
\eqref{eq:e8v91-symbol} determines the leading heat homogeneity.
\end{remark}

\subsection{Relative heat data in dimension three}
\label{sec:e8v91-heat}

Let \(H_0,H_1\) be positive invertible strongly elliptic operators of order
four on one closed three-manifold and one finite-rank bundle.  Define the
exact relative heat trace
\begin{equation}
 \Theta_{10}(t)
 :=
 \operatorname{Tr}(e^{-tH_1}-e^{-tH_0}),
 \qquad t>0.                                           \label{eq:e8v91-Theta}
\end{equation}
Assume the small-time expansion through homogeneity zero has the rigorous
remainder form
\begin{equation}
 \Theta_{10}(t)
 =
 \sum_{j=0}^3c_j(H_1,H_0)t^{a_j}
 +{\cal R}_{10}(t),
 \qquad
 a_j:={j-3\over4},                                     \label{eq:e8v91-expansion}
\end{equation}
where, for some \(\eta>0\) and \(t_*>0\),
\begin{equation}
 |{\cal R}_{10}(t)|\leq C t^\eta
 \qquad(0<t\leq t_*).                                  \label{eq:e8v91-remainder}
\end{equation}
Thus
\[
 a_0=-\frac34,\qquad a_1=-\frac12,\qquad
 a_2=-\frac14,\qquad a_3=0.
\]

\begin{theorem}[Local endpoint form of the coefficients]
\label{thm:e8v91-local-coefficients}
For smooth strongly elliptic differential operators of order four on a
closed three-manifold, there are heat coefficients \(A_j(H)\), each the
integral of a density determined locally by finitely many jets of the full
symbol of \(H\), such that
\begin{equation}
 \operatorname{Tr}(e^{-tH})
 \sim\sum_{j=0}^\infty A_j(H)t^{(j-3)/4}.              \label{eq:e8v91-heat-asymptotic}
\end{equation}
Consequently
\begin{equation}
 c_j(H_1,H_0)=A_j(H_1)-A_j(H_0)                       \label{eq:e8v91-endpoint}
\end{equation}
in \eqref{eq:e8v91-expansion}.  Coefficients that vanish by parity or a
special operator symmetry may be omitted only after that vanishing is
proved for the chosen constraint operator.
\end{theorem}

\begin{proof}
In a coordinate patch, introduce a parameter-dependent symbol
\(q(x,\xi,\lambda)\) solving
\[
 (H-\lambda)\# q=I
\]
recursively in decreasing symbol order, beginning with
\((\sigma_4(H)-\lambda)^{-1}\).  Strong ellipticity gives the uniform
sectorial estimates needed for the recursion.  The contour representation
\[
 e^{-tH}
 ={1\over2\pi\mathrm i}
 \int_\Gamma e^{-t\lambda}(H-\lambda)^{-1}\,d\lambda
\]
turns each homogeneous resolvent-symbol term into a power
\(t^{(j-3)/4}\) after the scaling
\(\xi=t^{-1/4}\zeta\), \(\lambda=t^{-1}\rho\).
The diagonal kernel coefficient is a universal polynomial in finitely many
symbol derivatives at the same base point.  A partition of unity and the
symbol change-of-coordinate rule assemble these diagonal terms into global
densities.  Truncating the recursive parametrix after the term \(j=N\)
leaves a trace remainder \(O(t^{(N+1-3)/4})\); choosing \(N\geq3\) gives
\eqref{eq:e8v91-remainder} after retaining enough terms.

Subtract the two expansions.  Linearity of the trace makes every relative
coefficient the endpoint difference
\eqref{eq:e8v91-endpoint}.  The final warning follows because the
parametrix construction does not force a coefficient to vanish without an
additional symmetry.
\end{proof}

\begin{firewall}[Asymptotic means finite-order remainder control]
\label{fire:e8v91-asymptotic}
Equation \eqref{eq:e8v91-heat-asymptotic} is an asymptotic expansion.  V91
uses only the four displayed coefficients and the remainder estimate
\eqref{eq:e8v91-remainder}.  It makes no convergence claim for the infinite
series, in accordance with the V90 audit.
\end{firewall}

\subsection{Explicit finite-part definition}
\label{sec:e8v91-finite-part}

Fix a positive reference time \(\tau_0\), which is the determinant
renormalization scale.  Choose any splitting time
\(0<\tau\leq t_*\), and put
\begin{equation}
 P_{10}(t):=\sum_{j=0}^3c_j(H_1,H_0)t^{a_j}.           \label{eq:e8v91-P}
\end{equation}
Define
\begin{align}
 {\cal L}_{\tau_0}(H_1,H_0)
 :={}&-{1\over2}\int_0^\tau
 {\Theta_{10}(t)-P_{10}(t)\over t}\,dt                 \notag\\
 &-{1\over2}\int_\tau^\infty{\Theta_{10}(t)\over t}\,dt
 -{1\over2}\sum_{j=0}^2
 {c_j(H_1,H_0)\tau^{a_j}\over a_j}                    \notag\\
 &-{1\over2}c_3(H_1,H_0)
 \log{\tau\over\tau_0}.                                \label{eq:e8v91-L}
\end{align}

\begin{theorem}[Existence and splitting-time independence]
\label{thm:e8v91-finite-part}
The integrals in \eqref{eq:e8v91-L} converge, and
\({\cal L}_{\tau_0}(H_1,H_0)\) is independent of the auxiliary splitting
time \(\tau\).
\end{theorem}

\begin{proof}
Near zero,
\[
 {\Theta_{10}(t)-P_{10}(t)\over t}
 =O(t^{\eta-1}),
\]
which is integrable because \(\eta>0\).  If
\(H_i\geq\gamma_iI\), then for \(t\geq1\),
\[
 \operatorname{Tr}e^{-tH_i}
 \leq
 e^{-(t-1)\gamma_i}\operatorname{Tr}e^{-H_i}.
\]
Thus the large-time integral converges.

Differentiate the right side of \eqref{eq:e8v91-L} with respect to
\(\tau\).  Twice its negative derivative is
\begin{align*}
 &{\Theta_{10}(\tau)-P_{10}(\tau)\over\tau}
 -{\Theta_{10}(\tau)\over\tau}
 +\sum_{j=0}^2c_j\tau^{a_j-1}
 +{c_3\over\tau}\\
 &\qquad=
 -{P_{10}(\tau)\over\tau}
 +{P_{10}(\tau)\over\tau}=0.
\end{align*}
Hence the value is independent of \(\tau\).
\end{proof}

\begin{theorem}[Equivalent cutoff subtraction]
\label{thm:e8v91-cutoff}
Let
\begin{equation}
 {\cal L}_\varepsilon(H_1,H_0)
 :=
 -{1\over2}\int_\varepsilon^\infty
 {\Theta_{10}(t)\over t}\,dt.                          \label{eq:e8v91-Leps}
\end{equation}
Then
\begin{align}
 {\cal L}_{\tau_0}(H_1,H_0)
 =
 \lim_{\varepsilon\downarrow0}
 \biggl[
 {\cal L}_\varepsilon(H_1,H_0)
 &-{1\over2}\sum_{j=0}^2
 {c_j(H_1,H_0)\varepsilon^{a_j}\over a_j}             \notag\\
 &+{c_3(H_1,H_0)\over2}
 \log{\tau_0\over\varepsilon}
 \biggr].                                              \label{eq:e8v91-cutoff}
\end{align}
\end{theorem}

\begin{proof}
Split the integral in \eqref{eq:e8v91-Leps} at \(\tau\), insert
\(\Theta_{10}=P_{10}+{\cal R}_{10}\) on
\((\varepsilon,\tau)\), and integrate the four model terms:
\begin{align*}
 \int_\varepsilon^\tau t^{a_j-1}\,dt
 &={\tau^{a_j}-\varepsilon^{a_j}\over a_j}
 \quad(j=0,1,2),\\
 \int_\varepsilon^\tau {dt\over t}
 &=\log{\tau\over\varepsilon}.
\end{align*}
The remainder integral converges to its integral from zero by
\eqref{eq:e8v91-remainder}.  Rearranging the power and logarithmic terms
gives exactly \eqref{eq:e8v91-L}.
\end{proof}

\begin{corollary}[Renormalization-scale dependence]
\label{cor:e8v91-scale}
For \(\tau_0,\tau'_0>0\),
\begin{equation}
 {\cal L}_{\tau'_0}(H_1,H_0)
 -{\cal L}_{\tau_0}(H_1,H_0)
 =
 {c_3(H_1,H_0)\over2}
 \log{\tau'_0\over\tau_0}.                             \label{eq:e8v91-scale}
\end{equation}
\end{corollary}

\begin{proof}
Only the final logarithm in \eqref{eq:e8v91-L} contains the reference
scale.
\end{proof}

\subsection{Exact endpoint cocycle}
\label{sec:e8v91-cocycle}

\begin{theorem}[Additivity of the renormalized relative determinant]
\label{thm:e8v91-cocycle}
Let \(H_0,H_1,H_2\) satisfy the preceding hypotheses and use the same
renormalization scale and coefficient convention for every pair.  Then
\begin{align}
 {\cal L}_{\tau_0}(H_2,H_0)
 &=
 {\cal L}_{\tau_0}(H_2,H_1)
 +{\cal L}_{\tau_0}(H_1,H_0),                          \label{eq:e8v91-cocycle}\\
 {\cal L}_{\tau_0}(H_0,H_1)
 &=-{\cal L}_{\tau_0}(H_1,H_0).                        \label{eq:e8v91-antisymmetry}
\end{align}
\end{theorem}

\begin{proof}
The exact heat traces satisfy
\[
 \Theta_{20}=\Theta_{21}+\Theta_{10}.
\]
The endpoint formula \eqref{eq:e8v91-endpoint} gives
\[
 c_j(H_2,H_0)
 =
 c_j(H_2,H_1)+c_j(H_1,H_0).
\]
Every term in \eqref{eq:e8v91-L} is linear in the relative heat trace and
its four relative coefficients.  Add the two formulas to obtain
\eqref{eq:e8v91-cocycle}.  Interchanging the endpoints changes the sign of
the heat trace and every coefficient, proving
\eqref{eq:e8v91-antisymmetry}.
\end{proof}

\begin{corollary}[Local counterterms telescope]
\label{cor:e8v91-counterterms}
For each \(j\leq3\), the subtraction coefficient is the coboundary
\[
 c_j(H_{n+1},H_n)
 =A_j(H_{n+1})-A_j(H_n).
\]
Consequently its power or logarithmic counterterm telescopes along the
regulator sequence and is compatible with the V84 ACDC cocycle.
\end{corollary}

\begin{proof}
This is \eqref{eq:e8v91-endpoint}; summing endpoint differences telescopes.
The same linearity proves the full cocycle in
\Cref{thm:e8v91-cocycle}.
\end{proof}

\begin{firewall}[Pair-dependent finite parts can break the cocycle]
\label{fire:e8v91-scheme}
If different adjacent pairs use different heat scales, discard different
coefficients, or add unrelated finite local terms, the proof of
\eqref{eq:e8v91-cocycle} no longer applies.  Scheme changes must themselves
be endpoint differences of one declared local functional, including the
scale term \eqref{eq:e8v91-scale}.
\end{firewall}

\subsection{Coordinate and slow-patch covariance}
\label{sec:e8v91-covariance}

\begin{proposition}[Unitary covariance]
\label{prop:e8v91-unitary}
If \(U\) is unitary and
\(\widetilde H_i=UH_iU^{-1}\), then
\begin{equation}
 {\cal L}_{\tau_0}(\widetilde H_1,\widetilde H_0)
 =
 {\cal L}_{\tau_0}(H_1,H_0).                           \label{eq:e8v91-unitary}
\end{equation}
\end{proposition}

\begin{proof}
Functional calculus gives
\(e^{-t\widetilde H_i}=Ue^{-tH_i}U^{-1}\).
The trace is unitary invariant, so the relative heat trace, its asymptotic
coefficients, and every term of \eqref{eq:e8v91-L} agree.
\end{proof}

\begin{proposition}[Buffered slow-sector refinement]
\label{prop:e8v91-slow}
Suppose two parameter patches retain different finite slow spectral
subspaces, and on their overlap use the common refinement that retains
their union.  If the same refined fast normal operator is used on both
sides, their finite-part fast determinants agree.  The modes added to the
slow sector contribute only the corresponding finite-dimensional
determinant and Gram factors.
\end{proposition}

\begin{proof}
On the common refined fast space the two normal operators are identical, so
\Cref{prop:e8v91-unitary} applies to any change of orthonormal fast basis.
The orthogonal finite slow complement is finite dimensional.  The V83
block determinant factorization and V84 Gram transformation law place its
contribution in the finite-dimensional ledger.
\end{proof}

\begin{firewall}[No subtraction of an extended slow eigenfunction]
\label{fire:e8v91-slow}
The heat trace in \eqref{eq:e8v91-Theta} is taken directly on the declared
refined fast sector.  It is not computed as a full heat trace minus a
separately estimated nonlocal slow projector.  This follows the buffered
functional-calculus rule imported at V90.
\end{firewall}

\subsection{A cofinal convergence criterion}
\label{sec:e8v91-convergence}

\begin{theorem}[Dominated finite-part convergence]
\label{thm:e8v91-convergence}
Let \(\Theta^{(n)}_{10}\) be a sequence of relative heat traces with
coefficients \(c_j^{(n)}\).  Suppose for one
\(\tau\leq t_*\):
\begin{enumerate}[label=\textup{(\roman*)},leftmargin=3.8em]
\item \(c_j^{(n)}\to c_j\) for \(0\leq j\leq3\);
\item for \(0<t\leq\tau\),
\[
 \left|
 \Theta^{(n)}_{10}(t)
 -\sum_{j=0}^3c_j^{(n)}t^{a_j}
 \right|
 \leq Ct^\eta
\]
with common \(C,\eta>0\);
\item the subtracted traces converge pointwise on \((0,\tau]\);
\item \(\Theta^{(n)}_{10}(t)\to\Theta_{10}(t)\) for \(t\geq\tau\), and
\[
 \int_\tau^\infty
 \sup_n|\Theta^{(n)}_{10}(t)|\,{dt\over t}<\infty.
\]
\end{enumerate}
Then
\begin{equation}
 {\cal L}^{(n)}_{\tau_0}(H_1,H_0)
 \longrightarrow
 {\cal L}_{\tau_0}(H_1,H_0).                           \label{eq:e8v91-convergence}
\end{equation}
\end{theorem}

\begin{proof}
On \((0,\tau]\), the subtracted integrands divided by \(t\) are dominated by
\(Ct^{\eta-1}\), which is integrable.  On \([\tau,\infty)\), use hypothesis
(iv).  Dominated convergence passes to both integrals in
\eqref{eq:e8v91-L}; ordinary convergence handles the four explicit
coefficient terms.
\end{proof}

\begin{firewall}[Finite lattice heat traces are not uniformly continuum-like at
zero time]
\label{fire:e8v91-lattice}
At fixed finite regulator,
\(\operatorname{Tr}e^{-tH_R}\to\dim{\cal H}_R\) as \(t\downarrow0\).
It cannot satisfy a continuum \(t^{-3/4}\) expansion uniformly below the
mesh time.  Therefore
\Cref{thm:e8v91-convergence} applies only after a correlated heat/lattice
comparison or to already constructed continuum operators.  Interchanging
the limits \(R\to\infty\) and \(t\downarrow0\) without such a theorem is not
licensed.
\end{firewall}

\subsection{Principal determinant card}
\label{sec:e8v91-card}

\begin{definition}[Normal heat finite-part card, NHFC]
\label{def:e8v91-NHFC}
For each regular constraint chart and pair of backgrounds, NHFC records:
\begin{enumerate}[label=\textup{(N\arabic*)},leftmargin=4em]
\item the fast quotient derivative \(D\), its Hilbert metrics, and
\(H=D^*D\);
\item a positive spectral floor and the strongly elliptic symbol
\eqref{eq:e8v91-symbol};
\item the exact relative heat trace \eqref{eq:e8v91-Theta};
\item the four local endpoint coefficients \(A_0,\ldots,A_3\);
\item a small-time remainder estimate
\eqref{eq:e8v91-remainder};
\item the common scale \(\tau_0\) and finite part
\eqref{eq:e8v91-L};
\item the slow-patch common refinement and finite Gram ledger;
\item the endpoint counterterm convention ensuring
\eqref{eq:e8v91-cocycle};
\item a correlated regulator/heat convergence theorem; and
\item compatibility with the V89 lower-order
\(\det_2\) factor and its multiplicative anomaly.
\end{enumerate}
\end{definition}

\begin{theorem}[Conditional principal determinant closure]
\label{thm:e8v91-closure}
If NHFC holds and the cofinal sequence satisfies
\Cref{thm:e8v91-convergence} in every fixed regular local chart, then the
principal-symbol part of the relative coarea determinant has the
well-defined value
\[
 \exp{\cal L}_{\tau_0}(H_1,H_0).
\]
These relative factors obey the exact chart/regulator cocycle
\eqref{eq:e8v91-cocycle}.
\end{theorem}

\begin{proof}
Existence is \Cref{thm:e8v91-finite-part}; cofinal convergence is
\Cref{thm:e8v91-convergence}; and the cocycle is
\Cref{thm:e8v91-cocycle}.  Since
\({\cal L}\) is one half of the normal-operator log determinant, its
exponential is the renormalized absolute coarea determinant ratio.
\end{proof}

\begin{assessment}[Progress at V91]
\label{ass:e8v91-status}
The noncompact principal metric response now has a precise renormalized
target.  The positive normal operator is fourth order with an explicit
strongly elliptic symbol.  In three dimensions its relative heat
determinant requires exactly the coefficients of
\(t^{-3/4},t^{-1/2},t^{-1/4},t^0\).  The finite-part formula is independent
of the splitting time, has one explicit logarithmic scale dependence, and
obeys an exact endpoint cocycle because every subtraction coefficient is a
local coboundary.
\end{assessment}

\begin{programme}[Eighth-series V92: correlated lattice--heat window]
\label{prog:e8v91-V92}
V92 will address \Cref{fire:e8v91-lattice}.  It will formulate a two-scale
criterion with mesh time \(t_R\), compare the discrete and continuum heat
traces for \(t\geq t_R\), isolate the finite-dimensional interval
\(0<t<t_R\) as a local lattice counterterm, and prove a diagonal choice of
\(t_R\downarrow0\) that implies NHFC convergence.  If a uniform heat-kernel
comparison is unavailable, the programme will keep that estimate as the
named discretization bottleneck rather than exchange the limits.
\end{programme}

\begin{firewall}[Claim boundary after eighth-series V91]
\label{fire:e8v91-boundary}
V91 proves the fourth-order symbol, finite-part heat formula, scale law,
endpoint cocycle, and an abstract dominated convergence theorem.  It does
not compute the actual four SM--Einstein heat densities, prove a
lattice--continuum heat comparison, control singular constraint strata,
assemble the full adjacent probability density, or construct the full
SM--Einstein continuum.
\end{firewall}

\begin{theorem}[Eighth-series V91 result]
\label{thm:e8v91-result}
For two positive invertible fourth-order normal constraint operators on a
closed three-manifold, the renormalized relative coarea log determinant is
the finite part \eqref{eq:e8v91-L}.  It subtracts precisely four local heat
coefficients, is independent of the auxiliary heat splitting time, changes
with renormalization scale by \eqref{eq:e8v91-scale}, and satisfies the
exact regulator cocycle \eqref{eq:e8v91-cocycle}.
\end{theorem}

\begin{proof}
Use \Cref{thm:e8v91-symbol,thm:e8v91-local-coefficients,
thm:e8v91-finite-part,cor:e8v91-scale,thm:e8v91-cocycle}.
\end{proof}

\paragraph{Parent-version record.}
V91 was formed from
\texttt{Renewal\_SM\_Einstein\_Continuum\_Research\_Notes\_V90.tex}.
The V90 parent had \(36\,557\) logical source lines, \(1\,301\,756\) bytes,
and SHA--256
\[
 \texttt{0CBC47F99A967FB6F4099C9DD00FA1BDF898D13260C8C3EA9D9E8B565F49C467}.
\]
The next compulsory companion audit is V95.

% =============================================================================
% END APPENDED EIGHTH-SERIES V91 MATERIAL
% =============================================================================


% =============================================================================
% BEGIN APPENDED EIGHTH-SERIES V92 MATERIAL
% =============================================================================

\clearpage
\section{Eighth-series V92: a correlated lattice--heat window}
\label{sec:v8v92}

\subsection{Exact finite-regulator heat identity}
\label{sec:e8v92-finite}

At regulator \(R\), let \(H_{0,R}\) and \(H_{1,R}\) be positive invertible
matrices on the same finite-dimensional fast quotient space
\({\cal H}_R\).  Put
\begin{equation}
 \Theta_R(t)
 :=
 \operatorname{Tr}\left(e^{-tH_{1,R}}-e^{-tH_{0,R}}\right).
 \label{eq:e8v92-ThetaR}
\end{equation}

\begin{theorem}[Finite-matrix Frullani identity]
\label{thm:e8v92-Frullani}
For every such pair,
\begin{equation}
 {1\over2}\log{\det H_{1,R}\over\det H_{0,R}}
 =
 -{1\over2}\int_0^\infty{\Theta_R(t)\over t}\,dt.
 \label{eq:e8v92-Frullani}
\end{equation}
The integral converges at both endpoints.
\end{theorem}

\begin{proof}
Diagonalize each positive matrix.  For positive numbers \(a,b\),
\begin{equation}
 \int_0^\infty{e^{-ta}-e^{-tb}\over t}\,dt
 =\log{b\over a}.                                      \label{eq:e8v92-scalar-Frullani}
\end{equation}
Indeed, differentiate the integral with respect to \(a\), obtaining
\(-1/a\), and fix the integration constant by setting \(a=b\).
The finite spectral sums may be interchanged with the integral.  Equal
matrix dimensions cancel the constant small-time terms, so
\(\Theta_R(t)=O(t)\) at zero; positivity gives exponential decay at
infinity.  Summing \eqref{eq:e8v92-scalar-Frullani} proves the result.
\end{proof}

\begin{firewall}[Comparison is made at one regulator]
\label{fire:e8v92-dimension}
\Cref{thm:e8v92-Frullani} compares two backgrounds on the same regulator
space.  Matrices at adjacent regulators generally have different
dimensions.  Their renormalized relative values may be compared only after
each has first been formed relative to its declared same-regulator
reference and the local marginals have been transported as in V84.
\end{firewall}

\subsection{Replacing the lattice ultraviolet interval}
\label{sec:e8v92-window}

Retain the exponents
\[
 a_j={j-3\over4},\qquad0\leq j\leq3
\]
from V91.  Let \(c_{j,R}\) be regulator approximations to the four relative
continuum heat coefficients.  For \(t>0\), define
\begin{equation}
 B_R(t)
 :=
 -{1\over2}\sum_{j=0}^2{c_{j,R}t^{a_j}\over a_j}
 +{c_{3,R}\over2}\log{\tau_0\over t}.                 \label{eq:e8v92-BR}
\end{equation}
For a mesh heat time \(t_R>0\), set
\begin{equation}
 \widehat{\cal L}_R(t_R)
 :=
 -{1\over2}\int_{t_R}^\infty{\Theta_R(t)\over t}\,dt
 +B_R(t_R).                                            \label{eq:e8v92-Lhat}
\end{equation}

\begin{definition}[Exact lattice ultraviolet counterterm]
\label{def:e8v92-counterterm}
Let
\begin{equation}
 {\cal L}_R
 :={1\over2}\log{\det H_{1,R}\over\det H_{0,R}}.
 \label{eq:e8v92-LR}
\end{equation}
The ultraviolet counterterm removed from the raw determinant is
\begin{align}
 {\cal C}^{\rm UV}_R(t_R)
 &:={\cal L}_R-\widehat{\cal L}_R(t_R)                 \notag\\
 &=
 -{1\over2}\int_0^{t_R}{\Theta_R(t)\over t}\,dt
 -B_R(t_R).                                            \label{eq:e8v92-CUV}
\end{align}
\end{definition}

\begin{proposition}[No small-time lattice approximation is hidden]
\label{prop:e8v92-exact}
Equation \eqref{eq:e8v92-CUV} removes the complete finite-matrix interval
\(0<t<t_R\) exactly.  Consequently convergence of
\(\widehat{\cal L}_R(t_R)\) requires comparison of lattice and continuum
heat traces only on \(t\geq t_R\).
\end{proposition}

\begin{proof}
Split the exact identity \eqref{eq:e8v92-Frullani} at \(t_R\) and subtract
\eqref{eq:e8v92-CUV}.  The first integral cancels identically and the
remaining expression is \eqref{eq:e8v92-Lhat}.
\end{proof}

\subsection{Endpoint and closed-walk form of the counterterm}
\label{sec:e8v92-local}

For a positive matrix \(H_R\) of dimension \(N_R\), define the finite
endpoint functional
\begin{equation}
 {\cal U}_R(H_R;t_R)
 :=
 -{1\over2}\int_0^{t_R}
 {\operatorname{Tr}(e^{-tH_R})-N_R\over t}\,dt.
 \label{eq:e8v92-U}
\end{equation}

\begin{lemma}[Entire closed-walk expansion]
\label{lem:e8v92-walk}
The endpoint functional has the absolutely convergent series
\begin{equation}
 {\cal U}_R(H_R;t_R)
 =
 -{1\over2}\sum_{n=1}^\infty
 {(-1)^nt_R^n\over n\,n!}\operatorname{Tr}(H_R^n).
 \label{eq:e8v92-walk}
\end{equation}
If \(H_R\) has propagation at most \(\ell_R\), every summand
\(\operatorname{Tr}(H_R^n)\) is a sum over closed regulator walks of \(n\)
steps and diameter at most \(n\ell_R\).
\end{lemma}

\begin{proof}
The matrix exponential series converges uniformly on
\([0,t_R]\):
\[
 e^{-tH_R}-I
 =\sum_{n=1}^\infty{(-t)^n\over n!}H_R^n.
\]
Divide by \(t\), take the finite trace, and integrate term by term to obtain
\eqref{eq:e8v92-walk}.  Expanding the diagonal blocks of \(H_R^n\) gives
products
\[
 (H_R)_{x_0x_1}\cdots(H_R)_{x_{n-1}x_0}.
\]
Every nonzero factor joins points at distance at most \(\ell_R\), proving
the walk statement.
\end{proof}

\begin{corollary}[Quantitative long-walk tail]
\label{cor:e8v92-walk-tail}
If \(t_R\|H_R\|\leq\rho\), the part of
\eqref{eq:e8v92-walk} with \(n>m\) is bounded by
\begin{equation}
 {N_R e^\rho\rho^{m+1}\over
  2(m+1)(m+1)!}.                                      \label{eq:e8v92-walk-tail}
\end{equation}
\end{corollary}

\begin{proof}
Use
\(|\operatorname{Tr}H_R^n|\leq N_R\|H_R\|^n\) and
\[
 \sum_{n=m+1}^\infty{\rho^n\over n\,n!}
 \leq {1\over m+1}
 \sum_{n=m+1}^\infty{\rho^n\over n!}
 \leq {e^\rho\rho^{m+1}\over(m+1)(m+1)!}.
\]
\end{proof}

Let \(A_{j,R}(H_R)\) be endpoint coefficient functionals and set
\[
 c_{j,R}(H_{1,R},H_{0,R})
 =
 A_{j,R}(H_{1,R})-A_{j,R}(H_{0,R}).
\]
Define
\begin{equation}
 {\cal B}_R(H_R;t)
 :=
 -{1\over2}\sum_{j=0}^2
 {A_{j,R}(H_R)t^{a_j}\over a_j}
 +{A_{3,R}(H_R)\over2}\log{\tau_0\over t}.
 \label{eq:e8v92-endpoint-B}
\end{equation}

\begin{theorem}[Counterterm is an endpoint coboundary]
\label{thm:e8v92-coboundary}
When \(H_{0,R},H_{1,R}\) have the same dimension,
\begin{align}
 {\cal C}^{\rm UV}_R(H_{1,R},H_{0,R};t_R)
 =
 \bigl[{\cal U}_R(H_{1,R};t_R)
      -{\cal B}_R(H_{1,R};t_R)\bigr]                  \notag\\
 -
 \bigl[{\cal U}_R(H_{0,R};t_R)
      -{\cal B}_R(H_{0,R};t_R)\bigr].                 \label{eq:e8v92-coboundary}
\end{align}
Thus the exact small-time lattice term, the four heat subtractions, and
their combination all telescope under changes of background.
\end{theorem}

\begin{proof}
The dimension terms in \eqref{eq:e8v92-U} cancel between the two endpoints,
leaving the first integral in \eqref{eq:e8v92-CUV}.  The difference of the
two functionals \eqref{eq:e8v92-endpoint-B} is exactly \(B_R(t_R)\).
Substitution proves \eqref{eq:e8v92-coboundary}.
\end{proof}

\begin{remark}[Locality content]
\label{rem:e8v92-locality}
When the endpoint matrices are finite range,
\Cref{lem:e8v92-walk} expresses their ultraviolet functionals as sums of
closed local walks.  The functional is not strictly finite range because
all walk lengths occur, but \Cref{cor:e8v92-walk-tail} gives factorial
suppression after a fixed length when \(t_R\|H_R\|\) is controlled.
\end{remark}

\subsection{The two-scale convergence theorem}
\label{sec:e8v92-convergence}

Let \(\Theta(t)\) be the continuum relative heat trace of V91, with
coefficients \(c_j\), and let
\({\cal L}_{\tau_0}\) denote its finite part.

\begin{theorem}[Correlated lattice--heat convergence]
\label{thm:e8v92-two-scale}
Let \(t_R\downarrow0\).  Suppose:
\begin{enumerate}[label=\textup{(\roman*)},leftmargin=3.8em]
\item the mesoscopic heat error satisfies
\begin{equation}
 \int_{t_R}^\infty
 {|\Theta_R(t)-\Theta(t)|\over t}\,dt
 \longrightarrow0;                                    \label{eq:e8v92-heat-error}
\end{equation}
\item for \(j=0,1,2\),
\begin{equation}
 |c_{j,R}-c_j|\,t_R^{a_j}\longrightarrow0;             \label{eq:e8v92-coeff-power}
\end{equation}
\item the logarithmic coefficient satisfies
\begin{equation}
 |c_{3,R}-c_3|\,|\log t_R|\longrightarrow0.            \label{eq:e8v92-coeff-log}
\end{equation}
\end{enumerate}
Then
\begin{equation}
 \widehat{\cal L}_R(t_R)
 \longrightarrow
 {\cal L}_{\tau_0}(H_1,H_0).                           \label{eq:e8v92-two-scale}
\end{equation}
Equivalently, the raw finite determinant minus the exact counterterm
\eqref{eq:e8v92-CUV} converges to the V91 finite-part determinant.
\end{theorem}

\begin{proof}
Define the continuum cutoff expression
\begin{equation}
 \widehat{\cal L}(t_R)
 :=
 -{1\over2}\int_{t_R}^\infty{\Theta(t)\over t}\,dt
 -{1\over2}\sum_{j=0}^2{c_jt_R^{a_j}\over a_j}
 +{c_3\over2}\log{\tau_0\over t_R}.                   \label{eq:e8v92-cont-cutoff}
\end{equation}
V91 proves
\(\widehat{\cal L}(t_R)\to{\cal L}_{\tau_0}\).
The triangle inequality gives
\begin{align*}
 |\widehat{\cal L}_R(t_R)-\widehat{\cal L}(t_R)|
 \leq{}&
 {1\over2}\int_{t_R}^\infty
 {|\Theta_R-\Theta|\over t}\,dt\\
 &+{1\over2}\sum_{j=0}^2
 {|c_{j,R}-c_j|t_R^{a_j}\over|a_j|}
 +{|c_{3,R}-c_3|\over2}
 \left|\log{\tau_0\over t_R}\right|.
\end{align*}
Every term tends to zero by the three hypotheses; the fixed
\(\log\tau_0\) part follows from \eqref{eq:e8v92-coeff-log}.
\end{proof}

\begin{firewall}[Coefficient convergence needs a rate]
\label{fire:e8v92-rate}
Because \(a_0,a_1,a_2<0\), ordinary convergence
\(c_{j,R}\to c_j\) is too weak for an arbitrarily chosen \(t_R\).
The divergent powers amplify coefficient errors.  Conditions
\eqref{eq:e8v92-coeff-power}--\eqref{eq:e8v92-coeff-log} are therefore
essential unless the continuum coefficients are inserted exactly at every
regulator.
\end{firewall}

\subsection{Existence of a diagonal heat window}
\label{sec:e8v92-diagonal}

\begin{theorem}[Diagonal selection]
\label{thm:e8v92-diagonal}
Suppose:
\begin{enumerate}[label=\textup{(\roman*)},leftmargin=3.8em]
\item for every fixed \(\delta>0\),
\begin{equation}
 \int_\delta^\infty
 {|\Theta_R(t)-\Theta(t)|\over t}\,dt\longrightarrow0;
 \label{eq:e8v92-fixed-delta}
\end{equation}
\item \(c_{j,R}\to c_j\) for \(0\leq j\leq3\).
\end{enumerate}
Then there is a nonincreasing choice \(t_R\downarrow0\), after passage to
the given cofinal sequence, for which
\eqref{eq:e8v92-heat-error}--\eqref{eq:e8v92-coeff-log} all hold.
\end{theorem}

\begin{proof}
Choose a decreasing numerical sequence
\(\delta_k\downarrow0\).  For each fixed \(k\), hypotheses (i)--(ii) permit
an index \(R_k\) so large that for \(R\geq R_k\),
\begin{align}
 \int_{\delta_k}^\infty
 {|\Theta_R-\Theta|\over t}\,dt&<k^{-1},               \label{eq:e8v92-diag-heat}\\
 |c_{j,R}-c_j|\delta_k^{a_j}&<k^{-1}
 \quad(0\leq j\leq2),                                  \label{eq:e8v92-diag-power}\\
 |c_{3,R}-c_3||\log\delta_k|&<k^{-1}.                  \label{eq:e8v92-diag-log}
\end{align}
Increase the indices so that \(R_{k+1}>R_k\).  Set
\[
 t_R:=\delta_k\quad\hbox{for }R_k\leq R<R_{k+1}.
\]
Then \(t_R\downarrow0\), and the three errors are at most \(k^{-1}\) on the
\(k\)-th block.  This proves all required limits.
\end{proof}

\begin{remark}[What the diagonal theorem buys]
\label{rem:e8v92-diagonal}
Uniform comparison down to the microscopic lattice heat time is not needed.
It suffices to prove convergence on every fixed positive heat interval and
then let the lower heat cutoff approach zero slowly enough.  The price is
the exact endpoint counterterm \eqref{eq:e8v92-CUV}.
\end{remark}

\subsection{Exact cocycle at finite regulator and in the limit}
\label{sec:e8v92-cocycle}

\begin{theorem}[Renormalized finite-regulator cocycle]
\label{thm:e8v92-cocycle}
Let \(H_{0,R},H_{1,R},H_{2,R}\) act on the same regulator space.  Suppose
all coefficient approximants are endpoint differences:
\[
 c_{j,R}(H_b,H_a)=A_{j,R}(H_b)-A_{j,R}(H_a).
\]
Using one \(t_R\) and one scale \(\tau_0\),
\begin{equation}
 \widehat{\cal L}_R(H_2,H_0)
 =
 \widehat{\cal L}_R(H_2,H_1)
 +\widehat{\cal L}_R(H_1,H_0).                         \label{eq:e8v92-cocycle}
\end{equation}
If the three terms converge by
\Cref{thm:e8v92-two-scale}, their limits obey the V91 continuum cocycle.
\end{theorem}

\begin{proof}
Both the relative heat trace and each \(c_{j,R}\) are endpoint
differences.  Every term of \eqref{eq:e8v92-Lhat} is therefore additive.
Pass to the limit in the finite identity.
\end{proof}

\begin{firewall}[The heat window is part of the scheme]
\label{fire:e8v92-window}
Choosing different \(t_R\) for different adjacent background pairs can
destroy the exact finite-regulator cocycle.  The heat window must be fixed
for the regulator, or changes of window must be recorded as endpoint
counterterm differences.
\end{firewall}

\subsection{Application to the SM--Einstein normal operator}
\label{sec:e8v92-application}

\begin{definition}[Correlated normal heat card, CNHC]
\label{def:e8v92-CNHC}
For a cofinal constraint regulator, CNHC records:
\begin{enumerate}[label=\textup{(C\arabic*)},leftmargin=4em]
\item same-regulator positive normal matrices
\(H_{i,R}=D_{i,R}^*D_{i,R}\);
\item a common finite slow-sector convention and fast quotient;
\item endpoint coefficient approximants \(A_{0,R},\ldots,A_{3,R}\);
\item fixed-positive-time relative heat convergence
\eqref{eq:e8v92-fixed-delta};
\item coefficient convergence;
\item one diagonal heat window from
\Cref{thm:e8v92-diagonal};
\item the exact closed-walk endpoint counterterm
\eqref{eq:e8v92-CUV};
\item the finite-regulator cocycle \eqref{eq:e8v92-cocycle}; and
\item compatibility with V89's lower-order modified determinant and
multiplicative anomaly.
\end{enumerate}
\end{definition}

\begin{theorem}[Conditional correlated-cutoff closure]
\label{thm:e8v92-closure}
If CNHC holds on every fixed regular local chart, then the renormalized
finite-regulator principal coarea determinant converges to the V91
finite-part determinant and obeys the endpoint regulator cocycle.
\end{theorem}

\begin{proof}
\Cref{thm:e8v92-diagonal} supplies a heat window satisfying the hypotheses
of \Cref{thm:e8v92-two-scale}; that theorem proves convergence.
\Cref{thm:e8v92-cocycle} proves compatibility of the finite and limiting
cocycles.
\end{proof}

\begin{firewall}[The discretization estimate is still unpaid]
\label{fire:e8v92-unpaid}
V92 proves that fixed-positive-time heat convergence and coefficient
convergence are sufficient after diagonal selection.  It does not prove
those two facts for a chosen lattice, finite-element, or spectral
SM--Einstein regulator.  In particular, consistency of the discrete
adjoint, the quotient projection, and the four local heat coefficient
approximants remains a concrete discretization theorem.
\end{firewall}

\begin{assessment}[Progress at V92]
\label{ass:e8v92-status}
The noncommuting lattice and heat limits no longer require a uniform
continuum expansion at zero lattice time.  The entire microscopic interval
is removed as an exact endpoint coboundary with a closed-walk expansion.
Fixed-positive-time heat convergence and ordinary coefficient convergence
then admit a diagonal heat window for which the renormalized determinants
converge.  The remaining bottleneck is now a conventional semigroup and
coefficient consistency theorem for the chosen regulator.
\end{assessment}

\begin{programme}[Eighth-series V93: leading normal heat density]
\label{prog:e8v92-V93}
V93 will compute the leading coefficient \(A_0(H)\) from the explicit
principal symbol \eqref{eq:e8v91-symbol}.  It will diagonalize the scalar,
transverse-vector, and longitudinal-vector channels, evaluate the radial
Fourier integral, and identify the exact metric-volume counterterm.  The
lower coefficients \(A_1,A_2,A_3\) will remain separate until their symbol
recursion is written.
\end{programme}

\begin{firewall}[Claim boundary after eighth-series V92]
\label{fire:e8v92-boundary}
V92 proves the exact finite heat identity, closed-walk ultraviolet
counterterm, two-scale convergence theorem, diagonal heat-window selection,
and finite-regulator cocycle.  It does not verify heat convergence for a
specific SM--Einstein discretization, compute the four normal heat
coefficients, control singular constraint strata, prove the full ACDC
density estimate, or construct the full SM--Einstein continuum.
\end{firewall}

\begin{theorem}[Eighth-series V92 result]
\label{thm:e8v92-result}
Fixed-positive-time heat convergence and convergence of the four local
coefficient approximants suffice to choose \(t_R\downarrow0\) such that
\[
 {1\over2}\log{\det H_{1,R}\over\det H_{0,R}}
 -{\cal C}^{\rm UV}_R(t_R)
 \longrightarrow{\cal L}_{\tau_0}(H_1,H_0).
\]
The counterterm is an exact endpoint coboundary and the renormalized
determinants obey the regulator cocycle at every finite \(R\).
\end{theorem}

\begin{proof}
Use \Cref{thm:e8v92-Frullani,thm:e8v92-coboundary,
thm:e8v92-two-scale,thm:e8v92-diagonal,thm:e8v92-cocycle}.
\end{proof}

\paragraph{Parent-version record.}
V92 was formed from
\texttt{Renewal\_SM\_Einstein\_Continuum\_Research\_Notes\_V91.tex}.
The V91 parent had \(37\,105\) logical source lines, \(1\,320\,988\) bytes,
and SHA--256
\[
 \texttt{A5D2B0888EFBEEB4BF2C3B9D2BC70A170AF5F5FF454E71AB8957183D9A15B307}.
\]
The next compulsory companion audit is V95.

% =============================================================================
% END APPENDED EIGHTH-SERIES V92 MATERIAL
% =============================================================================


% =============================================================================
% BEGIN APPENDED EIGHTH-SERIES V93 MATERIAL
% =============================================================================

\clearpage
\section{Eighth-series V93: the leading normal heat density}
\label{sec:v8v93}

\subsection{Channel decomposition of the principal symbol}
\label{sec:e8v93-channels}

The V91 normal symbol acts on one scalar component and one three-component
vector.  At \(\xi\ne0\), write
\[
 T_xM=\xi^\perp\oplus\mathbb R\xi^\sharp.
\]
The four fibre channels and their fourth-order eigenvalues are
\begin{center}
\begin{tabular}{c|c|c}
channel&multiplicity&principal eigenvalue\\ \hline
scalar&\(1\)&\(16|\xi|_g^4\)\\
transverse vector&\(2\)&\(|\xi|_g^4\)\\
longitudinal vector&\(1\)&\(\frac{16}{9}|\xi|_g^4\).
\end{tabular}
\end{center}

\begin{lemma}[Radial fourth-order Gaussian integral]
\label{lem:e8v93-radial}
For \(c>0\),
\begin{equation}
 \int_{\mathbb R^3}e^{-c|\xi|^4}\,d\xi
 =
 \pi\Gamma\!\left({3\over4}\right)c^{-3/4}.            \label{eq:e8v93-radial}
\end{equation}
\end{lemma}

\begin{proof}
In spherical coordinates,
\[
 \int_{\mathbb R^3}e^{-c|\xi|^4}\,d\xi
 =
 4\pi\int_0^\infty r^2e^{-cr^4}\,dr.
\]
Set \(u=cr^4\).  Then
\[
 r^2\,dr={1\over4}c^{-3/4}u^{-1/4}\,du,
\]
so the radial integral is
\[
 \pi c^{-3/4}
 \int_0^\infty u^{3/4-1}e^{-u}\,du,
\]
which is \eqref{eq:e8v93-radial}.
\end{proof}

\subsection{Exact leading coefficient}
\label{sec:e8v93-A0}

\begin{theorem}[Leading heat density of the constraint normal operator]
\label{thm:e8v93-A0}
Let \(H=D^*D\) be the fourth-order normal operator of V91 on a closed
three-manifold.  Under the V87 principal-order hypothesis, its leading heat
coefficient is
\begin{equation}
 A_0(H)
 =
 C_{\rm con}\operatorname{Vol}_g(M),                   \label{eq:e8v93-A0}
\end{equation}
where
\begin{equation}
 C_{\rm con}
 :=
 {\Gamma(3/4)(17+3\sqrt3)\over64\pi^2}.                \label{eq:e8v93-C}
\end{equation}
Equivalently,
\begin{equation}
 \operatorname{Tr}(e^{-tH})
 =
 C_{\rm con}\operatorname{Vol}_g(M)t^{-3/4}
 +O(t^{-1/2})                                          \label{eq:e8v93-leading}
\end{equation}
as \(t\downarrow0\).
\end{theorem}

\begin{proof}
The leading diagonal heat density of a positive order-four system is
\[
 (2\pi)^{-3}
 \int_{T_x^*M}
 \operatorname{tr}
 e^{-\sigma_4(H)(x,\xi)}\,d\xi
\]
in an orthonormal cotangent frame.  Use the four channels above and
\Cref{lem:e8v93-radial}:
\begin{align*}
 a_0(x)
 &=
 {\Gamma(3/4)\over8\pi^2}
 \left[
 16^{-3/4}
 +2
 +\left({16\over9}\right)^{-3/4}
 \right]\\
 &=
 {\Gamma(3/4)\over8\pi^2}
 \left[
 {1\over8}+2+{3\sqrt3\over8}
 \right]\\
 &=
 {\Gamma(3/4)(17+3\sqrt3)\over64\pi^2}.
\end{align*}
This scalar is independent of \(x\).  The invariant leading density is
\(a_0(x)d\mu_g(x)\); integration proves
\eqref{eq:e8v93-A0}.  The next possible heat homogeneity is
\(t^{-1/2}\), giving \eqref{eq:e8v93-leading}.
\end{proof}

\begin{remark}[Scope of the constant]
\label{rem:e8v93-scope}
\(C_{\rm con}\) is the leading density for the scalar--vector
Hamiltonian/momentum constraint Jacobian.  It is not the heat coefficient
of the complete Standard Model matter kinetic operator or of a
gauge-fixed four-dimensional effective action.  Those determinants require
their own field, ghost, and chirality ledgers.
\end{remark}

\subsection{Fixed-volume cancellation}
\label{sec:e8v93-volume}

\begin{corollary}[Exact cancellation of the strongest relative divergence]
\label{cor:e8v93-fixed-volume}
Let \(H_0,H_1\) be normal constraint operators at two backgrounds in the
same V82 fixed-volume chart:
\begin{equation}
 \operatorname{Vol}_{g_1}(M)
 =
 \operatorname{Vol}_{g_0}(M)=V_0.                     \label{eq:e8v93-fixed-volume}
\end{equation}
Then the leading relative coefficient vanishes:
\begin{equation}
 c_0(H_1,H_0)
 =A_0(H_1)-A_0(H_0)=0.                                \label{eq:e8v93-c0-zero}
\end{equation}
Consequently the \(t_R^{-3/4}\) term is absent from the V92 relative
counterterm.
\end{corollary}

\begin{proof}
Substitute \eqref{eq:e8v93-fixed-volume} into
\eqref{eq:e8v93-A0}.  In V92 the \(j=0\) subtraction equals
\[
 -{1\over2}{c_0t_R^{-3/4}\over-3/4}
 ={2\over3}c_0t_R^{-3/4},
\]
which vanishes when \eqref{eq:e8v93-c0-zero} holds.
\end{proof}

\begin{corollary}[Infinitesimal conformal variation]
\label{cor:e8v93-variation}
For \(g_s=e^{2s\phi}g\),
\begin{equation}
 {d\over ds}A_0(H_{g_s})\bigg|_{s=0}
 =
 3C_{\rm con}\int_M\phi\,d\mu_g.                       \label{eq:e8v93-variation}
\end{equation}
Thus a tangent satisfying the linearized volume constraint
\(\int_M\phi\,d\mu_g=0\) has zero leading heat variation.
\end{corollary}

\begin{proof}
In dimension three,
\[
 d\mu_{e^{2s\phi}g}=e^{3s\phi}d\mu_g.
\]
Differentiate \eqref{eq:e8v93-A0}.
\end{proof}

\begin{firewall}[Volume cancellation removes only one coefficient]
\label{fire:e8v93-only-A0}
The fixed-volume condition proves \(c_0=0\).  It does not force
\(c_1,c_2\), or \(c_3\) to vanish.  Those coefficients contain lower
symbols of \(D^*D\), including curvature, extrinsic-curvature,
scalar--vector cross, and declared matter-derivative terms.
\end{firewall}

\subsection{Finite slow modes and the constant heat term}
\label{sec:e8v93-slow}

The leading local coefficient is insensitive to removal of finitely many
slow modes, but the homogeneity-zero coefficient of an absolute heat trace
is not.

\begin{proposition}[Effect of a finite spectral subtraction]
\label{prop:e8v93-finite-modes}
Let \(H>0\), and let \(P\) be a rank-\(N\) spectral projection with
eigenvalues \(\lambda_1,\ldots,\lambda_N\).  If the fast heat trace is
\[
 \operatorname{Tr}_{\rm f}e^{-tH}
 =
 \operatorname{Tr}e^{-tH}
 -\sum_{\alpha=1}^Ne^{-t\lambda_\alpha},
\]
then
\begin{align}
 A_j(H_{\rm f})&=A_j(H)
 &&(j=0,1,2),                                         \label{eq:e8v93-slow-low}\\
 A_3(H_{\rm f})&=A_3(H)-N.                            \label{eq:e8v93-slow-A3}
\end{align}
\end{proposition}

\begin{proof}
For each removed eigenvalue,
\[
 e^{-t\lambda_\alpha}=1-t\lambda_\alpha+O(t^2).
\]
Its subtraction changes only the \(t^0\) coefficient among the
homogeneities retained through \(j=3\), and changes it by \(-1\).
Summing over \(\alpha\) proves the result.
\end{proof}

\begin{theorem}[Matched slow transfer preserves the relative determinant]
\label{thm:e8v93-slow-transfer}
Let \(H_0,H_1\) have matched rank-\(N\) slow spectral subspaces with positive
eigenvalues \(\lambda_{\alpha,0}\) and \(\lambda_{\alpha,1}\).
Move all these modes from the fast heat traces to the finite slow ledger.
Then:
\begin{enumerate}[label=\textup{(\roman*)},leftmargin=3.8em]
\item none of the four relative coefficients \(c_0,c_1,c_2,c_3\) changes;
\item the fast relative finite part changes by
\begin{equation}
 -{1\over2}\sum_{\alpha=1}^N
 \log{\lambda_{\alpha,1}\over\lambda_{\alpha,0}};      \label{eq:e8v93-fast-change}
\end{equation}
\item adding the finite slow determinant
\begin{equation}
 {1\over2}\sum_{\alpha=1}^N
 \log{\lambda_{\alpha,1}\over\lambda_{\alpha,0}}       \label{eq:e8v93-slow-change}
\end{equation}
restores the original relative determinant exactly.
\end{enumerate}
\end{theorem}

\begin{proof}
The removed relative heat trace is
\[
 \sum_{\alpha=1}^N
 \left(e^{-t\lambda_{\alpha,1}}
      -e^{-t\lambda_{\alpha,0}}\right)
 =
 -t\sum_{\alpha=1}^N
 (\lambda_{\alpha,1}-\lambda_{\alpha,0})+O(t^2).
\]
It has no terms of homogeneity
\(t^{-3/4},t^{-1/2},t^{-1/4},t^0\), proving (i).
The finite-matrix Frullani identity of V92 assigns this relative heat trace
the value \(\frac12\sum_\alpha\log
(\lambda_{\alpha,1}/\lambda_{\alpha,0})\).
Removing it from the fast trace gives (ii); recording it separately gives
(iii).
\end{proof}

\begin{firewall}[Unmatched slow rank]
\label{fire:e8v93-rank}
If two sides of a comparison remove different numbers of slow modes, their
relative heat trace acquires a nonzero \(t^0\) term and the logarithmic
scale dependence changes.  On a slow-rank transition one must first pass to
the common refined slow space of V90.  Comparing unmatched fast spaces
directly is not a determinant cocycle.
\end{firewall}

\subsection{Leading coefficient approximation at finite regulator}
\label{sec:e8v93-regulator}

Define the finite-regulator leading endpoint approximant by
\begin{equation}
 A_{0,R}(H_{g,R})
 :=
 C_{\rm con}V_R(g),                                    \label{eq:e8v93-A0R}
\end{equation}
where \(V_R(g)\) is the declared quadrature volume of the discrete metric.

\begin{proposition}[Rate required by the V92 diagonal window]
\label{prop:e8v93-rate}
Let \(t_R\downarrow0\).  The V92 \(j=0\) coefficient condition is equivalent
to
\begin{equation}
 |[V_R(g_1)-V_R(g_0)]
   -[\operatorname{Vol}_{g_1}(M)-\operatorname{Vol}_{g_0}(M)]|
 =
 o(t_R^{3/4}).                                         \label{eq:e8v93-volume-rate}
\end{equation}
If the discrete fixed-volume constraint imposes
\[
 V_R(g_{1,R})=V_R(g_{0,R})
\]
and the continuum endpoints have the common volume \(V_0\), then the
condition holds identically.
\end{proposition}

\begin{proof}
By \eqref{eq:e8v93-A0} and \eqref{eq:e8v93-A0R},
\[
 c_{0,R}-c_0
 =
 C_{\rm con}
 \left(
 [V_R(g_1)-V_R(g_0)]
 -[\operatorname{Vol}_{g_1}-\operatorname{Vol}_{g_0}]
 \right).
\]
The V92 condition is
\(|c_{0,R}-c_0|t_R^{-3/4}\to0\), which is exactly
\eqref{eq:e8v93-volume-rate}.  Under matched fixed volume both brackets
vanish.
\end{proof}

\begin{consequence}[One CNHC coefficient closes geometrically]
\label{cons:e8v93-CNHC}
On a regulator that enforces the same discrete fixed-volume equation used
by the nonlinear constraint chart, the strongest coefficient row of CNHC
is exact and needs no asymptotic rate estimate.
\end{consequence}

\begin{proof}
Use \Cref{cor:e8v93-fixed-volume,prop:e8v93-rate}.
\end{proof}

\subsection{Status and next coefficient}
\label{sec:e8v93-status}

\begin{assessment}[Progress at V93]
\label{ass:e8v93-status}
The leading normal heat density is now explicit:
\[
 C_{\rm con}
 ={\Gamma(3/4)(17+3\sqrt3)\over64\pi^2}.
\]
It multiplies only the spatial volume.  Because the V82 chart fixes volume,
the \(t^{-3/4}\) relative coefficient and the strongest V92 counterterm
cancel exactly.  Finite slow refinements do not alter any relative
coefficient through homogeneity zero when their ranks are matched, and
their finite determinants restore the fast-mode transfer exactly.
\end{assessment}

\begin{programme}[Eighth-series V94: parity and the \(A_1\) row]
\label{prog:e8v93-V94}
V94 will determine whether the \(t^{-1/2}\) coefficient \(A_1\) vanishes
for the closed-manifold normal operator.  It will use the parity of the
parameter-dependent resolvent symbol and will keep the order-three symbol
of \(D^*D\) explicit.  If parity does not force vanishing for the complete
coupled operator, V94 will give the local symbol integral defining \(A_1\)
rather than assume the Laplace-type pattern.
\end{programme}

\begin{firewall}[Claim boundary after eighth-series V93]
\label{fire:e8v93-boundary}
V93 computes \(A_0\), proves fixed-volume cancellation, and tracks finite
slow modes.  It does not compute \(A_1,A_2,A_3\), verify the remaining
discrete heat coefficients or semigroup convergence, control singular
constraint strata, prove the full ACDC density estimate, or construct the
full SM--Einstein continuum.
\end{firewall}

\begin{theorem}[Eighth-series V93 result]
\label{thm:e8v93-result}
The leading heat coefficient of the fourth-order scalar--vector normal
constraint operator is
\[
 A_0(H)
 =
 {\Gamma(3/4)(17+3\sqrt3)\over64\pi^2}
 \operatorname{Vol}_g(M).
\]
It cancels exactly in every relative determinant inside the fixed-volume
constraint chart, eliminating the most singular
\(t_R^{-3/4}\) counterterm.
\end{theorem}

\begin{proof}
Use \Cref{thm:e8v93-A0,cor:e8v93-fixed-volume}.
\end{proof}

\paragraph{Parent-version record.}
V93 was formed from
\texttt{Renewal\_SM\_Einstein\_Continuum\_Research\_Notes\_V92.tex}.
The V92 parent had \(37\,596\) logical source lines, \(1\,338\,393\) bytes,
and SHA--256
\[
 \texttt{313088C76E3DEB0A805025D31F95D44D964544F40F0ACD0A96926BE5C30B979E}.
\]
The next compulsory companion audit is V95.

% =============================================================================
% END APPENDED EIGHTH-SERIES V93 MATERIAL
% =============================================================================


% =============================================================================
% BEGIN APPENDED EIGHTH-SERIES V94 MATERIAL
% =============================================================================

\clearpage
\section{Eighth-series V94: parity of the normal heat coefficients}
\label{sec:v8v94}

\subsection{Parity class of an even-order differential operator}
\label{sec:e8v94-symbol}

Let \(H\) be a positive differential operator of order four on a smooth
vector bundle over a closed three-manifold.  In local coordinates and a
local bundle frame, write its complete symbol as
\begin{equation}
 p(x,\xi)=p_4(x,\xi)+p_3(x,\xi)+p_2(x,\xi)
 +p_1(x,\xi)+p_0(x).                                  \label{eq:e8v94-p}
\end{equation}
Because \(H\) is differential, \(p_{4-k}\) is polynomial homogeneous of
degree \(4-k\) in \(\xi\), and therefore
\begin{equation}
 p_{4-k}(x,-\xi)=(-1)^kp_{4-k}(x,\xi),
 \qquad0\leq k\leq4.                                  \label{eq:e8v94-p-parity}
\end{equation}
Factors of \(\mathrm i\) from the symbol convention do not change this
parity identity.

\subsection{Parameter-resolvent recursion}
\label{sec:e8v94-resolvent}

Let \(\lambda\) lie on a sectorial contour avoiding the positive principal
spectrum.  The symbol of \((H-\lambda)^{-1}\) has a parameter-homogeneous
expansion
\begin{equation}
 q(x,\xi,\lambda)
 \sim\sum_{j=0}^\infty q_{-4-j}(x,\xi,\lambda),
 \label{eq:e8v94-q}
\end{equation}
where \(\xi\) has weight one and \(\lambda\) has weight four.

\begin{lemma}[Alternating resolvent-symbol parity]
\label{lem:e8v94-q-parity}
Every homogeneous resolvent symbol satisfies
\begin{equation}
 q_{-4-j}(x,-\xi,\lambda)
 =
 (-1)^jq_{-4-j}(x,\xi,\lambda).                       \label{eq:e8v94-q-parity}
\end{equation}
\end{lemma}

\begin{proof}
The leading term is
\[
 q_{-4}=(p_4-\lambda)^{-1}.
\]
It is even in \(\xi\) by \eqref{eq:e8v94-p-parity}, proving the case
\(j=0\).

For the left Kohn--Nirenberg product, the equation
\((p-\lambda)\# q=I\) gives the recursion
\begin{equation}
 q_{-4-j}
 =
 -q_{-4}
 \sum_{\substack{k+\ell+|\alpha|=j\\ \ell<j}}
 {1\over\alpha!}
 (\partial_\xi^\alpha p_{4-k})
 (D_x^\alpha q_{-4-\ell}).                            \label{eq:e8v94-recursion}
\end{equation}
Assume the claim for every \(\ell<j\).  Differentiating
\eqref{eq:e8v94-p-parity} in \(\xi\) shows that
\(\partial_\xi^\alpha p_{4-k}\) has parity
\((-1)^{k+|\alpha|}\).  Spatial differentiation does not change parity, so
\(D_x^\alpha q_{-4-\ell}\) has parity \((-1)^\ell\).
Every summand in \eqref{eq:e8v94-recursion} therefore has parity
\[
 (-1)^{k+|\alpha|+\ell}=(-1)^j.
\]
Multiplication by the even symbol \(q_{-4}\) preserves it.  Induction proves
\eqref{eq:e8v94-q-parity}.
\end{proof}

\subsection{Vanishing of the odd heat coefficients}
\label{sec:e8v94-vanishing}

\begin{theorem}[Odd coefficient vanishing on a closed manifold]
\label{thm:e8v94-odd}
For a positive strongly elliptic differential operator of even order on a
closed manifold, the local interior heat coefficient with odd symbol index
vanishes.  In the present order-four, dimension-three case,
\begin{equation}
 A_1(H)=0,\qquad A_3(H)=0.                             \label{eq:e8v94-A13}
\end{equation}
\end{theorem}

\begin{proof}
The local coefficient of index \(j\) is obtained from
\(q_{-4-j}\) by the contour integral in \(\lambda\), followed by the
cotangent integral:
\begin{equation}
 a_j(x)
 =
 {1\over(2\pi)^3}
 \int_{\mathbb R^3}
 \operatorname{tr}
 \left[
 {1\over2\pi\mathrm i}
 \int_\Gamma e^{-\lambda}
 q_{-4-j}(x,\xi,\lambda)\,d\lambda
 \right]d\xi.                                         \label{eq:e8v94-aj}
\end{equation}
The standard parameter scaling has already extracted the factor
\(t^{(j-3)/4}\).  By
\Cref{lem:e8v94-q-parity}, the integrand in \(\xi\) is odd when \(j\) is
odd.  The regularized homogeneous integral is defined symmetrically in
\(\xi\), and its value is zero.  Thus the local coefficient density itself
vanishes for odd \(j\).  Integrating over the closed manifold gives
\eqref{eq:e8v94-A13}.

There are no boundary parametrices or boundary heat densities because the
manifold is closed.  Coordinate changes preserve the zero density, so the
local calculation assembles globally.
\end{proof}

\begin{remark}[The order-three symbol need not vanish]
\label{rem:e8v94-p3}
The normal operator \(D^*D\) generally has a nonzero order-three symbol,
coming from the first-order scalar--vector blocks and derivatives of the
order-two coefficients.  The theorem does not set \(p_3=0\).  It shows that
the complete index-one heat integrand, including every recursive correction,
is odd in \(\xi\) and hence integrates to zero.
\end{remark}

\begin{corollary}[Relative odd coefficients]
\label{cor:e8v94-relative}
For two local closed-manifold normal operators \(H_0,H_1\),
\begin{equation}
 c_1(H_1,H_0)=0,\qquad c_3(H_1,H_0)=0.                 \label{eq:e8v94-relative}
\end{equation}
The V91 relative determinant is independent of the logarithmic scale
\(\tau_0\).
\end{corollary}

\begin{proof}
Each relative coefficient is an endpoint difference.  Apply
\eqref{eq:e8v94-A13}.  The V91 scale law is proportional to \(c_3\).
\end{proof}

\subsection{Finite slow reductions}
\label{sec:e8v94-slow}

Let \(P_{\rm slow}\) be a rank-\(N\) spectral projection removed from the
local differential operator.  V93 gives
\[
 A_3(H_{\rm fast})=A_3(H)-N=-N.
\]
Thus parity applies to the local full differential expression before the
finite spectral subtraction.

\begin{proposition}[Matched slow rank preserves relative parity]
\label{prop:e8v94-slow}
If \(H_0,H_1\) remove slow spectral subspaces of the same rank \(N\), then
their fast relative coefficients still satisfy
\begin{equation}
 c_{1,{\rm fast}}=0,\qquad c_{3,{\rm fast}}=0.          \label{eq:e8v94-fast-relative}
\end{equation}
\end{proposition}

\begin{proof}
Finite spectral subtraction does not change \(A_1\).  It changes both
absolute \(A_3\) coefficients by the same number \(-N\), which cancels in
their difference.  Use \Cref{thm:e8v94-odd}.
\end{proof}

\begin{firewall}[Global projectors are bookkeeping, not differential
coefficients]
\label{fire:e8v94-projector}
The operator obtained by composing a differential stencil with a global
fast projector is nonlocal and need not satisfy
\eqref{eq:e8v94-p-parity} as a differential symbol.  The parity theorem is
applied to the local differential normal operator; the finitely many
removed modes are then handled by
\Cref{prop:e8v94-slow} and the V93 finite determinant transfer.
\end{firewall}

\subsection{The reduced finite-part formula}
\label{sec:e8v94-reduced}

For the local normal operator on a closed three-manifold, the relative heat
trace through homogeneity zero is
\begin{equation}
 \Theta_{10}(t)
 =
 c_0t^{-3/4}+c_2t^{-1/4}+{\cal R}_{10}(t),
 \qquad {\cal R}_{10}(t)=O(t^\eta).                   \label{eq:e8v94-expansion}
\end{equation}

\begin{theorem}[Two-coefficient finite part]
\label{thm:e8v94-finite-part}
The V91 finite part reduces to
\begin{align}
 {\cal L}(H_1,H_0)
 ={}&
 -{1\over2}\int_0^\tau
 {\Theta_{10}(t)-c_0t^{-3/4}-c_2t^{-1/4}\over t}\,dt  \notag\\
 &-{1\over2}\int_\tau^\infty{\Theta_{10}(t)\over t}\,dt
 +{2\over3}c_0\tau^{-3/4}
 +2c_2\tau^{-1/4}.                                    \label{eq:e8v94-finite-part}
\end{align}
It is independent of \(\tau\), and it has no renormalization-scale
parameter.  Inside a fixed-volume chart, \(c_0=0\), so
\begin{align}
 {\cal L}(H_1,H_0)
 ={}&
 -{1\over2}\int_0^\tau
 {\Theta_{10}(t)-c_2t^{-1/4}\over t}\,dt
 -{1\over2}\int_\tau^\infty{\Theta_{10}(t)\over t}\,dt
 +2c_2\tau^{-1/4}.                                    \label{eq:e8v94-one-coefficient}
\end{align}
\end{theorem}

\begin{proof}
Set \(c_1=c_3=0\) in the V91 formula.  Since
\[
 -{1\over2}{c_0\tau^{-3/4}\over-3/4}
 ={2\over3}c_0\tau^{-3/4},
 \qquad
 -{1\over2}{c_2\tau^{-1/4}\over-1/4}
 =2c_2\tau^{-1/4},
\]
one obtains \eqref{eq:e8v94-finite-part}.  Splitting-time independence is
the V91 theorem or follows by direct differentiation.
\Cref{cor:e8v94-relative} removes the scale term, and V93 fixed-volume
cancellation gives \eqref{eq:e8v94-one-coefficient}.
\end{proof}

\begin{corollary}[Revised V92 convergence rates]
\label{cor:e8v94-V92}
For a closed-manifold, matched-slow-rank, fixed-volume comparison, the V92
coefficient conditions reduce to the single requirement
\begin{equation}
 |c_{2,R}-c_2|t_R^{-1/4}\longrightarrow0.              \label{eq:e8v94-only-rate}
\end{equation}
One may set
\[
 c_{0,R}=c_{0}=0,\qquad
 c_{1,R}=c_{1}=0,\qquad
 c_{3,R}=c_{3}=0
\]
exactly in the endpoint subtraction scheme.
\end{corollary}

\begin{proof}
Use V93 for \(c_0\), \Cref{cor:e8v94-relative} for \(c_1,c_3\), and retain
only the \(j=2\) condition from V92.
\end{proof}

\subsection{Where parity can fail}
\label{sec:e8v94-failures}

\begin{firewall}[Boundary heat terms]
\label{fire:e8v94-boundary-heat}
On a manifold with boundary, the boundary parametrix produces additional
heat homogeneities and odd-index coefficients need not vanish.  The
closed-spatial-manifold result cannot be applied to a local patch with an
artificial boundary.  Patch calculations must glue before the global
parity statement is used, or must carry the boundary coefficients
explicitly.
\end{firewall}

\begin{firewall}[Pseudodifferential gauge reductions]
\label{fire:e8v94-pseudodiff}
A genuinely pseudodifferential gauge projector can have symbol terms that
do not obey the differential parity
\eqref{eq:e8v94-p-parity}.  This is another reason to retain the local
augmented differential operator and separate its finite kernel, rather
than insert a nonlocal Helmholtz projector into the determinant.
\end{firewall}

\begin{firewall}[Nonsmooth coefficients]
\label{fire:e8v94-regularity}
The symbol recursion and pointwise heat densities require smooth
coefficients.  The V82 \(W^{2,p}\) chart is sufficient for the nonlinear
constraint map but not automatically for a classical all-orders heat
parametrix.  A rough-coefficient approximation theorem or a smoother
regulator core is required before applying the coefficient calculation to
the full measure support.
\end{firewall}

\subsection{Parity card and status}
\label{sec:e8v94-status}

\begin{definition}[Closed normal parity card, CNPC]
\label{def:e8v94-CNPC}
CNPC records:
\begin{enumerate}[label=\textup{(P\arabic*)},leftmargin=4em]
\item a local even-order differential normal operator on a closed manifold;
\item smooth coefficients for the parameter-resolvent recursion;
\item separation of finite slow modes without a nonlocal projector in the
local symbol;
\item matched slow rank on every comparison;
\item the exact identities \(A_1=A_3=0\);
\item fixed-volume cancellation of the relative \(A_0\) row; and
\item the remaining \(A_2\) coefficient and regulator rate
\eqref{eq:e8v94-only-rate}.
\end{enumerate}
\end{definition}

\begin{assessment}[Progress at V94]
\label{ass:e8v94-status}
The order-three symbol does not generate a \(t^{-1/2}\) heat divergence:
parameter-resolvent parity makes its complete fibre trace odd in momentum.
The same argument kills the homogeneity-zero local coefficient.  After
matched slow-mode and fixed-volume bookkeeping, the principal determinant
has exactly one remaining divergent coefficient, \(c_2t^{-1/4}\), and no
renormalization-scale ambiguity.
\end{assessment}

\begin{programme}[Eighth-series V95: compulsory companion audit]
\label{prog:e8v94-V95}
V95 will inspect the newest YM and NS notes and test their current
finite-range, bounded-window, and kernel estimates against CNPC.  It will
then choose the V96 route to the remaining \(A_2\) density or to the rough
coefficient/discretization theorem if the audit shows that this must come
first.
\end{programme}

\begin{firewall}[Claim boundary after eighth-series V94]
\label{fire:e8v94-claim}
V94 proves the differential-symbol parity recursion, vanishing of
\(A_1,A_3\), matched slow-rank transfer, and the one-coefficient
fixed-volume finite part.  It does not compute \(A_2\), extend the heat
coefficient theorem to the rough measure support, prove discrete semigroup
convergence, control singular constraint strata, or construct the full
SM--Einstein continuum.
\end{firewall}

\begin{theorem}[Eighth-series V94 result]
\label{thm:e8v94-result}
For the local fourth-order normal constraint operator on a closed
three-manifold,
\[
 A_1(H)=A_3(H)=0.
\]
With matched slow rank and fixed volume, the renormalized relative
determinant therefore requires only the \(A_2\) subtraction
\eqref{eq:e8v94-one-coefficient}, and its V92 coefficient condition is
\eqref{eq:e8v94-only-rate}.
\end{theorem}

\begin{proof}
Use \Cref{lem:e8v94-q-parity,thm:e8v94-odd,
prop:e8v94-slow,thm:e8v94-finite-part,cor:e8v94-V92}.
\end{proof}

\paragraph{Parent-version record.}
V94 was formed from
\texttt{Renewal\_SM\_Einstein\_Continuum\_Research\_Notes\_V93.tex}.
The V93 parent had \(37\,993\) logical source lines, \(1\,350\,799\) bytes,
and SHA--256
\[
 \texttt{A712DCE0B269BDBC65E010264433D38EF7F6E007E84B9425EB33386CCB76C79F}.
\]
The next version is the compulsory V95 companion audit.

% =============================================================================
% END APPENDED EIGHTH-SERIES V94 MATERIAL
% =============================================================================


% =============================================================================
% BEGIN APPENDED EIGHTH-SERIES V95 MATERIAL
% =============================================================================

\clearpage
\section{Eighth-series V95: compulsory companion audit}
\label{sec:v8v95}

\subsection{Files and version status}
\label{sec:e8v95-files}

The files inspected at this checkpoint were
\begin{align*}
 &\texttt{Renewal\_Yang\_Mills\_Mass\_Gap\_Research\_Notes\_V52.tex},\\
 &\texttt{Renewal\_NS\_Stability\_Research\_Notes\_V26.tex}.
\end{align*}
The YM file had \(21\,638\) logical source lines, \(781\,190\) bytes, and
SHA--256
\[
 \texttt{F900BF6905254177A1B4FB5FE5DBB111E44B1E6626157F4C64F36A45FEA4710D}.
\]
Its final internal appendix is labelled V51, its parent record says that
V51 was formed from V50, and its decision assigns V52 as the next step.
No distinct V52 appendix or V52 marker is present.  Thus the filename is
ahead of the substantive internal version at this checkpoint.  The file is
not modified because the directory is shared.

The NS file remains unchanged: \(5\,319\) logical source lines,
\(278\,283\) bytes, and SHA--256
\[
 \texttt{82C887F8C2C69E4F28FAD7C3DC8DE5B9099CB5A4BF431EBA1FFF3E91935978AD}.
\]

\subsection{New Yang--Mills results}
\label{sec:e8v95-YM}

The substantive YM additions after the V90 audit are V48--V51.
\begin{enumerate}[label=\textup{(Y\arabic*)},leftmargin=4.5em]
\item V48 derives a two-cluster Gaussian mixing estimate from canonical
correlation and the high-marginal covariance floor.  It explicitly
separates that result from the still-missing multi-edge normalized tree
lift.
\item V49 identifies the local marginal floor through principal angles.
Finite-rank delocalized slow modes are harmless for fixed local marginals
only when their normalization and overlap with the local support are kept
explicit.
\item V50 proves that the sharp complement of a delocalized constant
projector cannot be uniformly exponentially Schur local on expanding
boxes.  It replaces the false pure-locality statement by
\[
 C=C^{\rm loc}+F,
\]
where \(C^{\rm loc}\) is exponentially local and \(F\) is an explicit
finite-rank \(O(N^{-1})\) kernel with bounded unweighted row sum.
\item V50 also shows that the corrected local-plus-finite-rank edge remains
tree summable.  It warns that fixed physical supports require cell weights
under ultraviolet refinement.
\item V51 proves a path-resolved convolution bound: every tree joining two
anchors has an exponentially decaying all-local contribution plus an
\(O(N^{-1})\) term from paths using a delocalized edge.  The result remains
conditional on the normalized Gaussian forest lift and on a separate
slow-label sensitivity estimate.
\end{enumerate}

\begin{proposition}[Transfer to the normal heat programme]
\label{prop:e8v95-transfer}
The following conclusions are type-correct for V94.
\begin{enumerate}[label=\textup{(\roman*)},leftmargin=3.8em]
\item A delocalized volume or conformal-Killing projector must not be
inserted into the differential symbol used to compute \(A_2\).
\item The local differential normal operator supplies the heat density;
finite slow eigenvalues and Gram factors are added through the exact
finite-dimensional V83--V94 ledger.
\item A fixed physical heat or source localization must be expressed with
the regulator mass matrix or cell weights before any ultraviolet norm
claim.
\item Conditional tree summability of Yang--Mills activities does not
compute an elliptic heat coefficient or prove constraint-measure
tightness.
\end{enumerate}
\end{proposition}

\begin{proof}
Item (i) is the same finite-rank obstruction proved abstractly in V87 and
now verified in the companion sharp-projection setting.  Item (ii) follows
from the matched slow transfer of V93 and the parity bookkeeping of V94.
Item (iii) is the shared refinement issue in the V88 cell normalization and
the YM V50 audit.  The objects in (iv) have different operators, norms, and
limiting statements, so no implication is available.
\end{proof}

\subsection{Navier--Stokes status}
\label{sec:e8v95-NS}

The NS file supplies no new input.  Its rank-one Oseen derivative constants
do not determine the order-two symbol coefficient of the fourth-order
normal constraint operator, and its numerical certification is not a
substitute for a local elliptic symbol recursion.

\subsection{Why \(A_2\) is now the correct target}
\label{sec:e8v95-A2}

V93 proves \(c_0=0\) in the fixed-volume chart.  V94 proves
\[
 c_1=c_3=0
\]
on a closed manifold with matched slow rank.  Hence the entire principal
relative counterterm is
\begin{equation}
 2c_2t_R^{-1/4}.                                       \label{eq:e8v95-only}
\end{equation}
The coefficient \(A_2(H)\) is local and depends only on the symbol
components \(p_4,p_3,p_2\) of \(H=D^*D\) and the derivatives forced by the
parameter-resolvent recursion.  It should therefore be acquired before
any finite-rank slow or forest assembly.

\begin{decision}[V96 acquisition order]
\label{dec:e8v95-V96}
V96 will:
\begin{enumerate}[label=\textup{(A\arabic*)},leftmargin=4em]
\item write the exact parameter-resolvent terms
\(q_{-4},q_{-5},q_{-6}\);
\item give the invariant local integral defining \(a_2(x,H)\);
\item prove that no symbols below \(p_2\) enter \(A_2\);
\item identify the required coefficient jets of \(D\) and the still-generic
matter derivatives;
\item prove continuity of \(A_2\) under convergence of those finite symbol
jets; and
\item translate that continuity rate into the sole V92 condition
\(|c_{2,R}-c_2|t_R^{-1/4}\to0\).
\end{enumerate}
\end{decision}

\begin{firewall}[No universal numerical \(A_2\) is yet determined]
\label{fire:e8v95-A2}
The leading symbol alone determines \(A_0\), but it does not determine
\(A_2\).  The latter also uses the order-three and order-two symbols, which
contain the scalar potential, extrinsic curvature, scalar--vector
couplings, and declared matter derivatives.  A numerical or closed
geometric \(A_2\) cannot be claimed until those data are fixed.
\end{firewall}

\begin{theorem}[V95 audit conclusion]
\label{thm:e8v95-result}
The companion audit confirms the local-symbol acquisition order.  The
remaining heat counterterm is the single coefficient
\eqref{eq:e8v95-only}; it must be computed from the local differential
normal operator, while delocalized slow modes remain in a separate
finite-dimensional ledger.
\end{theorem}

\begin{proof}
The one-coefficient reduction is V93--V94.  The necessity of separating
delocalized slow modes follows from
\Cref{prop:e8v95-transfer}.  The required lower-symbol data are identified
in \Cref{fire:e8v95-A2}.
\end{proof}

\begin{assessment}[Progress at V95]
\label{ass:e8v95-status}
The compulsory audit is complete and the next step has been corrected
against the companion finite-rank obstruction.  The YM file's internal
version mismatch is recorded without treating it as a new theorem.  The
SM programme now proceeds to the exact \(q_{-6}\) symbol rather than hiding
a global quotient projector in a local heat coefficient.
\end{assessment}

\begin{firewall}[Claim boundary after eighth-series V95]
\label{fire:e8v95-boundary}
V95 records the companion results and proves their type-correct
classification.  It does not compute \(A_2\), specify the Standard Model
matter derivative symbols, prove discrete semigroup convergence, control
singular constraint strata, establish the full ACDC density estimate, or
construct the full SM--Einstein continuum.
\end{firewall}

\paragraph{Parent-version record.}
V95 was formed from
\texttt{Renewal\_SM\_Einstein\_Continuum\_Research\_Notes\_V94.tex}.
The V94 parent had \(38\,363\) logical source lines, \(1\,364\,131\) bytes,
and SHA--256
\[
 \texttt{9306735C55683885FAE22D0A95699FE91519CB057BB28CD606A582362E6EB083}.
\]
The next compulsory companion audit is V100, after which the active note
will be frozen according to the user's hundred-version rule.

% =============================================================================
% END APPENDED EIGHTH-SERIES V95 MATERIAL
% =============================================================================


% =============================================================================
% BEGIN APPENDED EIGHTH-SERIES V96 MATERIAL
% =============================================================================

\clearpage
\section{Eighth-series V96: the exact \(A_2\) symbol and its stability}
\label{sec:v8v96}

\subsection{Symbol convention}
\label{sec:e8v96-convention}

Let \(H=D^*D\) be the local positive fourth-order normal operator, written
in one coordinate patch and bundle frame as
\[
 p=p_4+p_3+p_2+p_1+p_0.
\]
Use left Kohn--Nirenberg composition
\begin{equation}
 (a\# b)(x,\xi)
 =
 \sum_\alpha{1\over\alpha!}
 (\partial_\xi^\alpha a)(x,\xi)
 (D_x^\alpha b)(x,\xi),
 \qquad D_x={1\over\mathrm i}\partial_x.               \label{eq:e8v96-product}
\end{equation}
Let \(\lambda\) lie on a fixed sectorial contour for the positive operator,
and put
\begin{equation}
 r(x,\xi,\lambda):=(p_4(x,\xi)-\lambda)^{-1}.          \label{eq:e8v96-r}
\end{equation}

\subsection{The first three parameter-resolvent symbols}
\label{sec:e8v96-recursion}

\begin{theorem}[Explicit recursion through \(q_{-6}\)]
\label{thm:e8v96-q6}
The parameter symbol of \((H-\lambda)^{-1}\) begins
\[
 q\sim q_{-4}+q_{-5}+q_{-6}+\cdots,
\]
where
\begin{align}
 q_{-4}={}&r,                                           \label{eq:e8v96-q4}\\
 q_{-5}={}&-r\left[
 p_3r+\sum_j(\partial_{\xi_j}p_4)D_{x_j}r
 \right],                                              \label{eq:e8v96-q5}\\
 q_{-6}={}&-r\biggl[
 p_2r+p_3q_{-5}\\
 &\quad+\sum_j\left(
 (\partial_{\xi_j}p_4)D_{x_j}q_{-5}
 +(\partial_{\xi_j}p_3)D_{x_j}r
 \right)                                               \notag\\
 &\quad+{1\over2}\sum_{j,k}
 (\partial_{\xi_j}\partial_{\xi_k}p_4)
 D_{x_j}D_{x_k}r
 \biggr].                                              \label{eq:e8v96-q6}
\end{align}
The order in every matrix product is the displayed order.
\end{theorem}

\begin{proof}
Insert
\[
 p-\lambda=\sum_{k=0}^4p_{4-k}-\lambda,
 \qquad
 q\sim\sum_{\ell\geq0}q_{-4-\ell}
\]
into \((p-\lambda)\# q=I\).  At parameter degree zero,
\[
 (p_4-\lambda)q_{-4}=I,
\]
giving \eqref{eq:e8v96-q4}.

At degree \(-1\),
\[
 (p_4-\lambda)q_{-5}
 +p_3r
 +\sum_j(\partial_{\xi_j}p_4)D_{x_j}r=0.
\]
Left multiplication by \(r\) gives \eqref{eq:e8v96-q5}.

At degree \(-2\), the triples
\((k,\ell,|\alpha|)\) satisfying
\(k+\ell+|\alpha|=2\) are
\[
 (0,2,0),(1,1,0),(2,0,0),
 (0,1,1),(1,0,1),(0,0,2).
\]
Their contributions are exactly the six groups in
\eqref{eq:e8v96-q6}, together with
\((p_4-\lambda)q_{-6}\).  The multiindex coefficient for
\(|\alpha|=2\) is equivalently
\(\frac12\sum_{j,k}\), because mixed derivatives occur twice in the
ordered sum.  Move all lower terms to the right and left-multiply by \(r\).
\end{proof}

\begin{corollary}[Finite jet dependence]
\label{cor:e8v96-jets}
The symbol \(q_{-6}(x,\xi,\lambda)\) depends only on
\begin{equation}
 {\cal J}_2(p)
 :=
 \left(
 p_4,\partial_xp_4,\partial_x^2p_4;\
 p_3,\partial_xp_3;\
 p_2
 \right)                                               \label{eq:e8v96-jet}
\end{equation}
and their displayed \(\xi\)-derivatives.  It is independent of
\(p_1,p_0\).
\end{corollary}

\begin{proof}
The inverse \(r\) depends on \(p_4\).  Its first and second spatial
derivatives are obtained from
\[
 D_xr=-r(D_xp_4)r
\]
and one more differentiation, so they use at most two spatial derivatives
of \(p_4\).  Formula \eqref{eq:e8v96-q5} uses \(p_3\) and the first jet of
\(p_4\); differentiating it once uses \(\partial_xp_3\) and
\(\partial_x^2p_4\).  Formula \eqref{eq:e8v96-q6} then uses \(p_2\) but no
lower symbol.  This proves the claim.
\end{proof}

\subsection{Invariant local formula for \(A_2\)}
\label{sec:e8v96-A2}

\begin{theorem}[The remaining heat density]
\label{thm:e8v96-A2}
In the convention \eqref{eq:e8v96-product}, the local index-two heat
density is
\begin{equation}
 a_2(x;H)
 =
 {1\over(2\pi)^3}
 \int_{\mathbb R^3}
 \operatorname{tr}\left[
 {1\over2\pi\mathrm i}
 \int_\Gamma e^{-\lambda}
 q_{-6}(x,\xi,\lambda)\,d\lambda
 \right]d\xi,                                         \label{eq:e8v96-a2}
\end{equation}
and
\begin{equation}
 A_2(H)=\int_Ma_2(x;H).                                \label{eq:e8v96-A2}
\end{equation}
Although the coordinate integrand changes under a different quantization
or bundle frame, the density in \eqref{eq:e8v96-A2} is invariant.
\end{theorem}

\begin{proof}
Use the sectorial contour formula
\[
 e^{-tH}
 ={1\over2\pi\mathrm i}
 \int_\Gamma e^{-t\lambda}(H-\lambda)^{-1}\,d\lambda.
\]
The parameter-homogeneous term \(q_{-4-j}\) produces the diagonal heat
homogeneity \(t^{(j-3)/4}\) after
\(\xi=t^{-1/4}\zeta\) and \(\lambda=t^{-1}\rho\).
For \(j=2\) this is \(t^{-1/4}\), and its diagonal kernel is precisely
\eqref{eq:e8v96-a2}.  Taking the fibre trace and integrating the diagonal
gives \eqref{eq:e8v96-A2}.

The diagonal heat kernel is intrinsically defined.  Two coordinate
parametrices differ by symbol divergences and smoothing remainders whose
diagonal trace has zero contribution at the fixed homogeneous order after
partition-of-unity assembly.  Therefore the coefficient extracted from
that diagonal kernel is a global density, independent of the local
calculation used to obtain it.
\end{proof}

\begin{corollary}[Only the middle normal symbols are unpaid]
\label{cor:e8v96-middle}
The one remaining fixed-volume counterterm
\[
 2c_2t_R^{-1/4}
\]
is completely determined by the order-four, order-three, and order-two
symbols of the two endpoint normal operators.  No order-one or order-zero
coefficient of \(H\) enters this divergent row.
\end{corollary}

\begin{proof}
Combine \Cref{cor:e8v96-jets,thm:e8v96-A2} with the V94
one-coefficient finite-part formula.
\end{proof}

\subsection{Translation back to the constraint derivative}
\label{sec:e8v96-D}

Write the complete local symbol of the order-two constraint derivative as
\begin{equation}
 d=d_2+d_1+d_0.                                        \label{eq:e8v96-d}
\end{equation}

\begin{proposition}[Constraint jets needed for \(A_2\)]
\label{prop:e8v96-D-jets}
For \(H=D^*D\):
\begin{enumerate}[label=\textup{(\roman*)},leftmargin=3.8em]
\item \(p_4\) is determined by \(d_2^*d_2\);
\item \(p_3\) is determined by \(d_2,d_1\) and the first spatial jet of
\(d_2\), including the adjoint density and connection terms;
\item \(p_2\) is determined by \(d_2,d_1,d_0\), the second spatial jet of
\(d_2\), and the first spatial jet of \(d_1\).
\end{enumerate}
Consequently \(A_2(D^*D)\) is fixed once
\begin{equation}
 \left(d_2,\partial_xd_2,\partial_x^2d_2;\
 d_1,\partial_xd_1;\
 d_0\right)                                            \label{eq:e8v96-D-jet}
\end{equation}
and the field/equation Hilbert densities are fixed.
\end{proposition}

\begin{proof}
Apply the symbol product \eqref{eq:e8v96-product} to \(D^*\# D\).
Terms of total degree four use only the two degree-two symbols.  Degree
three uses one degree-one symbol or one symbol derivative.  Degree two uses
one degree-zero symbol, two degree-one symbols, one derivative of a
degree-one symbol, or two derivatives distributed among the degree-two
symbols.  The formal adjoint differentiates the coefficient and Hilbert
density, but introduces no jet beyond the displayed orders.  Insert these
three normal symbols into \Cref{cor:e8v96-jets}.
\end{proof}

\begin{firewall}[The generic matter maps now enter decisively]
\label{fire:e8v96-matter}
V82 assumes that the matter density and current maps are \(C^1\) between
the strong Sobolev spaces and that their linearizations have the required
differential order.  That is enough for the implicit-function theorem, but
it does not specify the local coefficients \(d_1,d_0\) or their jets in
\eqref{eq:e8v96-D-jet}.  Hence
\eqref{eq:e8v96-a2} is the rigorous universal answer at the current level
of matter specification; a reduced curvature formula requires the explicit
Standard Model stress/current linearization.
\end{firewall}

\subsection{Stability of the coefficient}
\label{sec:e8v96-stability}

Let \(\Sigma\) be a compact parameter cosphere consisting of
\((x,\xi,\lambda)\) with \(x\in M\), \(\lambda\) on the normalized
sectorial contour, and
\(|\xi|+|\lambda|^{1/4}=1\).  Measure the jet difference in
\eqref{eq:e8v96-jet} and its required \(\xi\)-derivatives by
\(\varepsilon_n\) on \(\Sigma\).

\begin{theorem}[Lipschitz stability of \(A_2\)]
\label{thm:e8v96-stability}
Let \(H_n,H\) be smooth order-four systems on one identified bundle.
Assume:
\begin{enumerate}[label=\textup{(\roman*)},leftmargin=3.8em]
\item their principal symbols have one common sectorial ellipticity
constant;
\item the jets \({\cal J}_2(p^{(n)})\) are uniformly bounded in the symbol
seminorms used in \eqref{eq:e8v96-q6};
\item \(\varepsilon_n\to0\).
\end{enumerate}
Then, for all sufficiently large \(n\),
\begin{align}
 \sup_x|a_2(x;H_n)-a_2(x;H)|
 &\leq C\varepsilon_n,                                 \label{eq:e8v96-a2-stability}\\
 |A_2(H_n)-A_2(H)|
 &\leq C'\varepsilon_n.                                \label{eq:e8v96-A2-stability}
\end{align}
The constants depend only on the common ellipticity margin, finitely many
uniform symbol bounds, the contour, and the compact manifold.
\end{theorem}

\begin{proof}
The resolvent identity gives
\[
 r_n-r
 =
 -r_n(p_4^{(n)}-p_4)r.
\]
Uniform sectorial ellipticity bounds both resolvents on \(\Sigma\), hence
\(\|r_n-r\|\leq C\varepsilon_n\).  Differentiate this identity in \(x\)
and \(\xi\) as many times as occur in
\eqref{eq:e8v96-q5}--\eqref{eq:e8v96-q6}.  The uniform symbol bounds give
the same \(O(\varepsilon_n)\) estimate for every required derivative.

Formulas \eqref{eq:e8v96-q5}--\eqref{eq:e8v96-q6} are finite sums of
products of these resolvents, derivatives, \(p_3,p_2\), and their displayed
jets.  A telescoping product estimate therefore gives
\[
 \|q_{-6}^{(n)}-q_{-6}\|
 \leq C\varepsilon_n
\]
on \(\Sigma\).  Parameter homogeneity extends this to an integrable
majorant in \((\xi,\lambda)\) for
\eqref{eq:e8v96-a2}.  Dominated convergence with the quantitative bound
proves \eqref{eq:e8v96-a2-stability}.  Integrate over compact \(M\), also
including the convergent density identification, to obtain
\eqref{eq:e8v96-A2-stability}.
\end{proof}

\begin{corollary}[The remaining V92 rate]
\label{cor:e8v96-rate}
Suppose the endpoint coefficient approximants are obtained from symbols
whose jet errors are at most \(\varepsilon_R\).  Then the sole coefficient
condition of V94 follows from
\begin{equation}
 \varepsilon_Rt_R^{-1/4}\longrightarrow0.             \label{eq:e8v96-rate}
\end{equation}
In particular, if \(\varepsilon_R\to0\), the choice
\begin{equation}
 t_R=\varepsilon_R^2                                  \label{eq:e8v96-example}
\end{equation}
gives a coefficient error \(O(\varepsilon_R^{1/2})\), subject also to the
mesoscopic heat comparison required by V92.
\end{corollary}

\begin{proof}
\Cref{thm:e8v96-stability} gives
\(|c_{2,R}-c_2|\leq C\varepsilon_R\) after applying the estimate at both
endpoints.  Multiply by \(t_R^{-1/4}\).  For
\eqref{eq:e8v96-example}, the product is
\(\varepsilon_R^{1/2}\).
\end{proof}

\begin{firewall}[Symbol convergence is not semigroup convergence]
\label{fire:e8v96-semigroup}
The finite-jet estimate closes the coefficient row of CNHC.  It does not
prove the fixed-positive-time trace convergence
\[
 \int_\delta^\infty
 {|\Theta_R(t)-\Theta(t)|\over t}\,dt\to0.
\]
That row also requires spectral approximation without pollution, uniform
large-time control, and correct treatment of the finite slow modes.
\end{firewall}

\subsection{The \(A_2\) acquisition card}
\label{sec:e8v96-card}

\begin{definition}[Middle-symbol heat card, MSHC]
\label{def:e8v96-MSHC}
MSHC records:
\begin{enumerate}[label=\textup{(S\arabic*)},leftmargin=4em]
\item the local constraint symbols \(d_2,d_1,d_0\);
\item the field and equation densities defining \(D^*\);
\item the normal symbols \(p_4,p_3,p_2\);
\item the exact resolvent symbols
\eqref{eq:e8v96-q4}--\eqref{eq:e8v96-q6};
\item the local invariant density \eqref{eq:e8v96-a2};
\item the explicit Standard Model matter contributions to
\eqref{eq:e8v96-D-jet};
\item a uniform sectorial ellipticity margin;
\item the symbol-jet approximation rate \(\varepsilon_R\);
\item a heat window satisfying \eqref{eq:e8v96-rate}; and
\item the separate fixed-positive-time semigroup convergence row.
\end{enumerate}
\end{definition}

\begin{theorem}[Conditional \(A_2\) closure]
\label{thm:e8v96-closure}
If MSHC holds and
\(\varepsilon_Rt_R^{-1/4}\to0\), then the last divergent coefficient in the
fixed-volume closed-manifold normal determinant satisfies the V92
coefficient convergence condition.  If the separate semigroup row also
holds, the renormalized principal determinant converges.
\end{theorem}

\begin{proof}
Coefficient convergence is
\Cref{thm:e8v96-stability,cor:e8v96-rate}.  V93--V94 remove the other
three coefficients.  The V92 two-scale theorem then proves determinant
convergence once its heat-trace hypothesis is supplied.
\end{proof}

\begin{assessment}[Progress at V96]
\label{ass:e8v96-status}
The final divergent heat coefficient is no longer an unnamed geometric
constant.  It is the explicit parameter-resolvent integral
\eqref{eq:e8v96-a2} with \(q_{-6}\) given by
\eqref{eq:e8v96-q6}.  Only the finite jet
\eqref{eq:e8v96-D-jet} of the constraint derivative enters.  Uniform
convergence of that jet closes the coefficient row at the quantitative rate
\(\varepsilon_Rt_R^{-1/4}\).  The remaining determinant bottleneck is the
fixed-positive-time semigroup trace comparison and the explicit Standard
Model matter-symbol input.
\end{assessment}

\begin{programme}[Eighth-series V97: semigroup comparison without spectral
pollution]
\label{prog:e8v96-V97}
V97 will formulate a norm-resolvent criterion for the identified fast
normal operators and use a contour/Laplace representation to prove
fixed-positive-time heat convergence.  Because ordinary operator-norm
convergence does not imply trace convergence in growing dimension, it will
add the precise localized trace or collective compactness hypothesis needed
for the V92 integral.  Finite slow eigenvalues will remain in the separate
ledger.
\end{programme}

\begin{firewall}[Claim boundary after eighth-series V96]
\label{fire:e8v96-boundary}
V96 derives the exact \(A_2\) symbol, its finite-jet dependence, and its
quantitative stability.  It does not specify every Standard Model
matter-symbol coefficient, prove the regulator semigroup trace comparison,
control singular constraint strata, establish the full ACDC probability
density, or construct the full SM--Einstein continuum.
\end{firewall}

\begin{theorem}[Eighth-series V96 result]
\label{thm:e8v96-result}
The sole remaining divergent coefficient in the fixed-volume relative
normal determinant is
\[
 A_2(H)=\int_M{1\over(2\pi)^3}
 \int_{\mathbb R^3}\operatorname{tr}\left[
 {1\over2\pi\mathrm i}\int_\Gamma e^{-\lambda}
 q_{-6}\,d\lambda\right]d\xi.
\]
It depends only on
\((d_2,\partial d_2,\partial^2d_2;d_1,\partial d_1;d_0)\), and symbol-jet
error \(\varepsilon_R\) closes the V92 coefficient row whenever
\(\varepsilon_Rt_R^{-1/4}\to0\).
\end{theorem}

\begin{proof}
Use \Cref{thm:e8v96-q6,cor:e8v96-jets,thm:e8v96-A2,
prop:e8v96-D-jets,thm:e8v96-stability,cor:e8v96-rate}.
\end{proof}

\paragraph{Parent-version record.}
V96 was formed from
\texttt{Renewal\_SM\_Einstein\_Continuum\_Research\_Notes\_V95.tex}.
The V95 parent had \(38\,555\) logical source lines, \(1\,372\,286\) bytes,
and SHA--256
\[
 \texttt{274B12CECB6194A6582D59DF637D1144E8476A39F21962D70ACB880CA9109AF3}.
\]
The next compulsory companion audit is V100.

% =============================================================================
% END APPENDED EIGHTH-SERIES V96 MATERIAL
% =============================================================================


% =============================================================================
% BEGIN APPENDED EIGHTH-SERIES V97 MATERIAL
% =============================================================================

\clearpage
\section{Eighth-series V97: semigroup traces without spectral pollution}
\label{sec:v8v97}

\subsection{Why operator-norm resolvent convergence is insufficient}
\label{sec:e8v97-obstruction}

The remaining CNHC row is, for each fixed \(\delta>0\),
\[
 \int_\delta^\infty
 {|\Theta_R(t)-\Theta(t)|\over t}\,dt\longrightarrow0.
\]
This is a trace statement.  An operator-norm resolvent estimate does not
control the number of regulator modes.

\begin{example}[Growing multiplicity invisible to a resolvent norm]
\label{ex:e8v97-multiplicity}
Let \(H\) be a positive compact-resolvent operator.  Adjoin to its
regulator approximation a spurious block
\[
 \mu_R I_{M_R},\qquad \mu_R\longrightarrow\infty.
\]
At the resolvent point \(-1\), the norm of the extra block is
\[
 \|(\mu_RI+1)^{-1}\|={1\over\mu_R+1}\longrightarrow0.
\]
Thus it is invisible to generalized norm-resolvent convergence.  Its heat
trace is
\[
 M_Re^{-t\mu_R}.
\]
For a fixed \(\delta>0\), choosing
\(M_R=\lceil e^{2\delta\mu_R}\rceil\) makes this quantity diverge at
\(t=\delta\).
\end{example}

\begin{consequence}[A counting estimate is compulsory]
\label{cons:e8v97-counting}
Norm-resolvent convergence must be supplemented by a uniform eigenvalue
counting bound, a trace-ideal resolvent estimate, or an equivalent
collective compactness condition strong enough to control multiplicity.
\end{consequence}

\subsection{Spectral hypotheses}
\label{sec:e8v97-hypotheses}

For \(i=0,1\), let \(H_i\) be positive self-adjoint operators with compact
resolvent on the continuum fast Hilbert space.  List their eigenvalues with
multiplicity:
\[
 0<\gamma\leq\lambda_{1,i}\leq\lambda_{2,i}\leq\cdots.
\]
At regulator \(R\), let \(H_{i,R}\) have dimension \(N_R\) and eigenvalues
\[
 0<\gamma\leq
 \lambda_{1,i,R}\leq\cdots\leq\lambda_{N_R,i,R}.
\]

\begin{assumption}[Spectral convergence and uniform fourth-order Weyl floor]
\label{ass:e8v97-spectral}
For both endpoints:
\begin{align}
 \lambda_{k,i,R}&\longrightarrow\lambda_{k,i}
 &&\text{for every fixed }k,                           \label{eq:e8v97-eigen-convergence}\\
 \lambda_{k,i,R}&\geq c k^{4/3}-C
 &&(1\leq k\leq N_R),                                  \label{eq:e8v97-Weyl-R}\\
 \lambda_{k,i}&\geq c k^{4/3}-C
 &&(k\geq1),                                           \label{eq:e8v97-Weyl}
\end{align}
with constants \(c>0,C\geq0,\gamma>0\) independent of \(R\) and \(i\).
\end{assumption}

\begin{remark}[Counting-function form]
\label{rem:e8v97-counting}
The lower bound \eqref{eq:e8v97-Weyl-R} is equivalent, up to changed
constants, to
\[
 N_{i,R}(\Lambda)
 :=
 \#\{k:\lambda_{k,i,R}\leq\Lambda\}
 \leq C_W(1+\Lambda)^{3/4}.
\]
The exponent \(3/4\) is dimension divided by order for a fourth-order
operator in three dimensions.
\end{remark}

\subsection{Heat-trace convergence}
\label{sec:e8v97-heat}

\begin{lemma}[Uniform high-eigenvalue tail]
\label{lem:e8v97-tail}
For every \(\delta>0\),
\begin{equation}
 \sup_{R,i}\sup_{t\geq\delta}
 \sum_{k>K}^{N_R}e^{-t\lambda_{k,i,R}}
 \longrightarrow0
 \quad(K\to\infty),                                    \label{eq:e8v97-tail-R}
\end{equation}
and the same conclusion holds for the continuum tails.
\end{lemma}

\begin{proof}
Positivity makes \(e^{-t\lambda}\leq e^{-\delta\lambda}\) for
\(t\geq\delta\).  The Weyl floor gives
\[
 \sum_{k>K}^{N_R}e^{-t\lambda_{k,i,R}}
 \leq
 e^{\delta C}\sum_{k>K}e^{-\delta c k^{4/3}}.
\]
The series converges, and its tail tends to zero independently of \(R,i\).
Use \eqref{eq:e8v97-Weyl} for the continuum statement.
\end{proof}

\begin{theorem}[Uniform fixed-positive-time heat convergence]
\label{thm:e8v97-heat}
Under \Cref{ass:e8v97-spectral}, for every \(\delta>0\),
\begin{equation}
 \sup_{t\geq\delta}
 \left|
 \operatorname{Tr}(e^{-tH_{i,R}})
 -\operatorname{Tr}(e^{-tH_i})
 \right|
 \longrightarrow0                                    \label{eq:e8v97-heat}
\end{equation}
for \(i=0,1\).
\end{theorem}

\begin{proof}
Fix \(\varepsilon>0\).  By
\Cref{lem:e8v97-tail}, choose \(K\) so that the discrete and continuum
tails beyond \(K\) are each below \(\varepsilon/3\), uniformly for
\(t\geq\delta\).

For \(1\leq k\leq K\), use the mean-value theorem:
\[
 |e^{-t\lambda_{k,i,R}}-e^{-t\lambda_{k,i}}|
 \leq
 t e^{-t\min(\lambda_{k,i,R},\lambda_{k,i})}
 |\lambda_{k,i,R}-\lambda_{k,i}|.
\]
The common floor gives
\[
 \sup_{t\geq\delta}te^{-t\gamma}<\infty.
\]
Since only finitely many \(k\) occur,
\eqref{eq:e8v97-eigen-convergence} makes their total contribution smaller
than \(\varepsilon/3\) for large \(R\).  Add the two tails.
\end{proof}

\begin{corollary}[Relative heat convergence]
\label{cor:e8v97-relative}
With
\begin{align}
 \Theta_R(t)
 &:=
 \operatorname{Tr}(e^{-tH_{1,R}}-e^{-tH_{0,R}}),\\
 \Theta(t)
 &:=
 \operatorname{Tr}(e^{-tH_1}-e^{-tH_0}),
\end{align}
one has
\begin{equation}
 \sup_{t\geq\delta}|\Theta_R(t)-\Theta(t)|
 \longrightarrow0                                    \label{eq:e8v97-relative}
\end{equation}
for every \(\delta>0\).
\end{corollary}

\begin{proof}
Apply \Cref{thm:e8v97-heat} to the two endpoints and use the triangle
inequality.
\end{proof}

\subsection{The integrated CNHC row}
\label{sec:e8v97-integrated}

\begin{lemma}[Uniform large-time domination]
\label{lem:e8v97-large}
For every \(\delta>0\), there is \(C_\delta<\infty\) such that
\begin{equation}
 \sup_{R,i}\operatorname{Tr}(e^{-tH_{i,R}})
 +\sup_i\operatorname{Tr}(e^{-tH_i})
 \leq C_\delta e^{-\gamma t/2}
 \qquad(t\geq\delta).                                  \label{eq:e8v97-large}
\end{equation}
\end{lemma}

\begin{proof}
For \(t\geq\delta\),
\[
 e^{-t\lambda}
 =
 e^{-(t-\delta/2)\lambda}e^{-(\delta/2)\lambda}
 \leq
 e^{-(t-\delta/2)\gamma}e^{-(\delta/2)\lambda}.
\]
The Weyl floor bounds the heat traces at time \(\delta/2\) uniformly by
\[
 e^{\delta C/2}\sum_{k\geq1}
 e^{-(\delta c/2)k^{4/3}}.
\]
Absorb \(e^{\delta\gamma/2}\) into \(C_\delta\).
\end{proof}

\begin{theorem}[Fixed-positive-time trace integral]
\label{thm:e8v97-integrated}
Under \Cref{ass:e8v97-spectral}, for every fixed \(\delta>0\),
\begin{equation}
 \int_\delta^\infty
 {|\Theta_R(t)-\Theta(t)|\over t}\,dt
 \longrightarrow0.                                    \label{eq:e8v97-integrated}
\end{equation}
\end{theorem}

\begin{proof}
For each fixed \(t\geq\delta\),
\Cref{cor:e8v97-relative} gives convergence to zero.  The triangle
inequality and \Cref{lem:e8v97-large} dominate the integrand by
\[
 {2C_\delta e^{-\gamma t/2}\over t},
\]
which is integrable on \([\delta,\infty)\).  Dominated convergence proves
\eqref{eq:e8v97-integrated}.
\end{proof}

\begin{corollary}[The semigroup row of CNHC]
\label{cor:e8v97-CNHC}
Assumption \ref{ass:e8v97-spectral} implies the fixed-positive-time heat
condition used in the V92 diagonal-selection theorem.
\end{corollary}

\begin{proof}
Equation \eqref{eq:e8v97-integrated} is exactly that condition.
\end{proof}

\subsection{A conforming Galerkin route}
\label{sec:e8v97-Galerkin}

Let \(h_i\) be the closed positive quadratic form of \(H_i\), with form
domain \({\cal V}_i\) compactly embedded in the continuum Hilbert space.
Let \({\cal V}_{i,R}\subset{\cal V}_i\) be conforming finite-dimensional
spaces whose union is form-dense, and let \(H_{i,R}\) be the Ritz operator
defined by restricting \(h_i\) to \({\cal V}_{i,R}\).

\begin{theorem}[No pollution for a conforming form approximation]
\label{thm:e8v97-Galerkin}
For each fixed \(k\), the Ritz eigenvalues satisfy
\begin{equation}
 \lambda_{k,i,R}\longrightarrow\lambda_{k,i}.          \label{eq:e8v97-Ritz-convergence}
\end{equation}
If the spaces are nested, the convergence is monotone from above.  For
every \(R\),
\begin{equation}
 \lambda_{k,i,R}\geq\lambda_{k,i}.                     \label{eq:e8v97-Ritz-floor}
\end{equation}
Consequently a continuum Weyl lower bound
\(\lambda_{k,i}\geq ck^{4/3}-C\) automatically supplies the uniform
regulator bound \eqref{eq:e8v97-Weyl-R}.
\end{theorem}

\begin{proof}
The min--max principles are
\begin{align*}
 \lambda_{k,i}
 &=
 \inf_{\substack{L\subset{\cal V}_i\\\dim L=k}}
 \sup_{0\ne u\in L}{h_i(u,u)\over\|u\|^2},\\
 \lambda_{k,i,R}
 &=
 \inf_{\substack{L\subset{\cal V}_{i,R}\\\dim L=k}}
 \sup_{0\ne u\in L}{h_i(u,u)\over\|u\|^2}.
\end{align*}
The second infimum is over fewer subspaces, proving
\eqref{eq:e8v97-Ritz-floor}.  Form density permits approximation of a
continuum \(k\)-dimensional spectral subspace by a \(k\)-dimensional
subspace of \({\cal V}_{i,R}\), with uniformly convergent Rayleigh
quotients.  This gives
\(\limsup_R\lambda_{k,i,R}\leq\lambda_{k,i}\), and hence convergence.
Nestedness enlarges the family over which the infimum is taken and makes
the sequence nonincreasing.
\end{proof}

\begin{corollary}[Conforming fourth-order heat convergence]
\label{cor:e8v97-Galerkin}
If the positive normal constraint operators have the V91 strong
ellipticity, a common positive floor, and conforming form-dense regulator
spaces, then the continuum fourth-order Weyl bound and
\Cref{thm:e8v97-Galerkin} imply
\eqref{eq:e8v97-integrated}.
\end{corollary}

\begin{proof}
Strong ellipticity on a closed three-manifold gives compact resolvent and
the lower Weyl estimate of order \(k^{4/3}\).  Apply
\Cref{thm:e8v97-Galerkin} and then
\Cref{thm:e8v97-integrated}.
\end{proof}

\begin{firewall}[A changing discrete operator is not an exact Ritz
restriction]
\label{fire:e8v97-Strang}
If the regulator changes coefficients, quadrature, the adjoint, or the
constraint quotient, \(H_{i,R}\) need not be the exact restriction of
\(H_i\).  Then \eqref{eq:e8v97-Ritz-floor} does not follow automatically.
One needs a uniform form-consistency estimate, stability, and a
no-pollution theorem, together with the counting bound.  Calling such a
method Galerkin does not supply those facts.
\end{firewall}

\subsection{Finite slow modes}
\label{sec:e8v97-slow}

\begin{proposition}[Matched slow removal preserves spectral convergence]
\label{prop:e8v97-slow}
Suppose a fixed number \(N\) of slow eigenvalues is removed at both
endpoints and regulators, with those eigenvalues converging individually.
After relabelling the remaining eigenvalues, assumptions
\eqref{eq:e8v97-eigen-convergence}--\eqref{eq:e8v97-Weyl-R} and conclusion
\eqref{eq:e8v97-integrated} remain valid.
\end{proposition}

\begin{proof}
Removing \(N\) entries shifts the eigenvalue index by \(N\).  Fixed-index
convergence is preserved.  The bound
\[
 c(k+N)^{4/3}-C\geq ck^{4/3}-C
\]
preserves the Weyl floor after an inessential constant adjustment.
The omitted finite heat traces converge directly from convergence of their
eigenvalues, consistently with the V93 slow determinant ledger.
\end{proof}

\begin{firewall}[Near-kernel crossings require common refinement]
\label{fire:e8v97-crossing}
If the number of removed modes changes with \(R\) or the background,
fixed-index relabelling can fail and the positive floor can collapse.
The buffered common-refinement rule of V90 must first retain every
crossing mode in the finite slow sector.
\end{firewall}

\subsection{Determinant convergence with all heat rows}
\label{sec:e8v97-determinant}

\begin{theorem}[Conditional completion of CNHC]
\label{thm:e8v97-CNHC}
Assume:
\begin{enumerate}[label=\textup{(\roman*)},leftmargin=3.8em]
\item the fixed-volume and parity reductions of V93--V94;
\item the \(A_2\) symbol-jet convergence of V96;
\item the spectral hypotheses \ref{ass:e8v97-spectral}.
\end{enumerate}
Then there is a diagonal heat window \(t_R\downarrow0\) for which
\begin{equation}
 {1\over2}\log{\det H_{1,R}\over\det H_{0,R}}
 -{\cal C}^{\rm UV}_R(t_R)
 \longrightarrow
 {\cal L}(H_1,H_0),                                    \label{eq:e8v97-determinant}
\end{equation}
provided the window is chosen also to satisfy
\(\varepsilon_Rt_R^{-1/4}\to0\).
\end{theorem}

\begin{proof}
\Cref{thm:e8v97-integrated} supplies fixed-positive-time heat convergence.
V96 supplies convergence of the sole coefficient \(c_2\).  The V92
diagonal-selection proof may choose \(t_R\) sufficiently slowly that both
the heat and coefficient errors vanish.  Apply the V92 two-scale theorem.
\end{proof}

\begin{assessment}[Progress at V97]
\label{ass:e8v97-status}
The semigroup trace row is now reduced to standard spectral data.  Fixed
eigenvalue convergence removes spectral pollution, while the uniform
fourth-order Weyl floor controls growing multiplicity and upgrades heat
convergence to the integrated CNHC norm.  A conforming form-dense
approximation supplies both facts automatically from min--max and the
continuum Weyl bound.  Nonconforming regulator operators still require an
explicit form-consistency theorem.
\end{assessment}

\begin{programme}[Eighth-series V98: return to the constrained density]
\label{prog:e8v97-V98}
V98 will insert the convergent principal heat determinant and the V89
lower-order modified determinant into the V84 adjacent constrained-density
ledger.  It will isolate the remaining action, multiplier-coarea,
transport, and probability-normalization rows and prove a composition
theorem showing exactly which summable local errors yield the ACDC
\(L^1\) condition.
\end{programme}

\begin{firewall}[Claim boundary after eighth-series V97]
\label{fire:e8v97-boundary}
V97 proves heat-trace convergence from spectral convergence and a uniform
Weyl floor, plus a conforming Galerkin route.  It does not verify that a
chosen SM--Einstein regulator is conforming or no-pollution, specify all
matter symbols, control singular constraint strata, prove the remaining
adjacent-density rows, or construct the full SM--Einstein continuum.
\end{firewall}

\begin{theorem}[Eighth-series V97 result]
\label{thm:e8v97-result}
Fixed-index spectral convergence, a common positive floor, and the uniform
bound
\[
 \lambda_{k,R}\geq ck^{4/3}-C
\]
imply
\[
 \int_\delta^\infty
 {|\Theta_R(t)-\Theta(t)|\over t}\,dt\to0
\]
for every \(\delta>0\).  Together with V96 coefficient convergence, these
hypotheses close the correlated heat-determinant limit of V92.
\end{theorem}

\begin{proof}
Use \Cref{thm:e8v97-heat,thm:e8v97-integrated,
thm:e8v97-CNHC}.
\end{proof}

\paragraph{Parent-version record.}
V97 was formed from
\texttt{Renewal\_SM\_Einstein\_Continuum\_Research\_Notes\_V96.tex}.
The V96 parent had \(38\,998\) logical source lines, \(1\,388\,117\) bytes,
and SHA--256
\[
 \texttt{A07CC1BDC652C8E8A456F920FD5480C93B49EAB8F33319AA06DE34DC33CD8120}.
\]
The next compulsory companion audit is V100.

% =============================================================================
% END APPENDED EIGHTH-SERIES V97 MATERIAL
% =============================================================================


% =============================================================================
% BEGIN APPENDED EIGHTH-SERIES V98 MATERIAL
% =============================================================================

\clearpage
\section{Eighth-series V98: composition of the constrained-density rows}
\label{sec:v8v98}

\subsection{The completed determinant entry}
\label{sec:e8v98-determinant}

For one fixed transported local support \(A\), write the exact adjacent
unnormalized log-density change from V84 as
\begin{equation}
 \Delta_{n,A}
 =
 \delta^{W}_{n,A}
 -\delta^{\det}_{n,A}
 -\delta^{\rm coa}_{n,A}
 +\delta^{\rm tr}_{n,A}
 -\delta^{C}_{n,A}.                                   \label{eq:e8v98-ledger}
\end{equation}
The five entries are respectively the action/weight, augmented determinant,
multiplier-coarea Jacobian, transport Jacobian, and declared counterterm
increments.  The determinant entry now has the internal decomposition
\begin{equation}
 \delta^{\det}_{n,A}
 =
 \delta^{\rm heat}_{n,A}
 +\delta^{(2)}_{n,A}
 +\delta^{\rm anom}_{n,A}
 +\delta^{\rm slow}_{n,A}.                            \label{eq:e8v98-detrows}
\end{equation}
Here the terms are the V91--V97 normal heat finite part, the V89
lower-order modified determinant, its multiplicative anomaly, and the
finite slow/Gram determinant.

\begin{firewall}[A convergent determinant is only one density row]
\label{fire:e8v98-one-row}
Even exact convergence of every term in
\eqref{eq:e8v98-detrows} does not prove convergence of the probability
measure.  The action, coarea, transport, counterterm, and normalization
rows in \eqref{eq:e8v98-ledger} can still fail.
\end{firewall}

\subsection{Constants disappear under probability normalization}
\label{sec:e8v98-centering}

Let \(\mu\) be a probability measure and let \(\Delta\) be a measurable
function with \(0<\int e^\Delta\,d\mu<\infty\).  Define
\begin{equation}
 r_\Delta
 :=
 {e^\Delta\over\int e^\Delta\,d\mu}.                  \label{eq:e8v98-r}
\end{equation}

\begin{lemma}[Centering invariance]
\label{lem:e8v98-centering}
For every constant \(c\),
\[
 r_{\Delta+c}=r_\Delta.
\]
If \(\Delta\in L^1(\mu)\) and
\[
 h:=\Delta-\int\Delta\,d\mu,
\]
then
\begin{equation}
 r_\Delta={e^h\over\int e^h\,d\mu}.                   \label{eq:e8v98-centered-r}
\end{equation}
\end{lemma}

\begin{proof}
The factor \(e^c\) cancels between numerator and denominator.  Use
\(c=-\int\Delta\,d\mu\) for the second statement.
\end{proof}

Define the essential oscillation seminorm
\begin{equation}
 \|f\|_{\rm osc,\infty}
 :=
 \inf_{c\in\mathbb R}\|f-c\|_{L^\infty}.              \label{eq:e8v98-osc}
\end{equation}

\subsection{First composition route: oscillation budgets}
\label{sec:e8v98-Linfty}

\begin{theorem}[Row-wise oscillation compiler]
\label{thm:e8v98-osc}
Suppose
\[
 \Delta=\sum_{\rho=1}^m\sigma_\rho\delta_\rho,
 \qquad \sigma_\rho\in\{-1,1\},
\]
and
\begin{equation}
 \|\delta_\rho\|_{\rm osc,\infty}\leq\varepsilon_\rho.
 \label{eq:e8v98-row-osc}
\end{equation}
Put \(E=\sum_\rho\varepsilon_\rho\).  Then
\begin{align}
 e^{-2E}\leq r_\Delta\leq e^{2E}
 \quad\mu\text{-almost everywhere},                   \label{eq:e8v98-r-osc}\\
 \|r_\Delta-1\|_{L^1(\mu)}
 \leq e^{2E}-1.                                       \label{eq:e8v98-L1-osc}
\end{align}
\end{theorem}

\begin{proof}
For each row choose a constant \(c_\rho\) within an arbitrarily small
error of the infimum in \eqref{eq:e8v98-osc}.  Then
\[
 \left|\Delta-\sum_\rho\sigma_\rho c_\rho\right|
 \leq E
\]
after taking the infimum.  By \Cref{lem:e8v98-centering}, the summed
constant is irrelevant.  Hence the unnormalized factor lies between
\(e^{-E}\) and \(e^E\), and its expectation obeys the same bounds.
Division gives \eqref{eq:e8v98-r-osc}.  The pointwise interval implies
\(|r_\Delta-1|\leq e^{2E}-1\); integrate.
\end{proof}

\begin{corollary}[Summable oscillation rows]
\label{cor:e8v98-osc}
For the adjacent densities \(r_{n,A}\), suppose every row
\(\rho\) in \eqref{eq:e8v98-ledger}--\eqref{eq:e8v98-detrows} has
\[
 \|\delta^\rho_{n,A}\|_{\rm osc,\infty}
 \leq\varepsilon^\rho_{n,A},
\qquad
 \sum_n\sum_\rho\varepsilon^\rho_{n,A}<\infty.
 \label{eq:e8v98-osc-sum}
\]
Then
\[
 \sum_n\|r_{n,A}-1\|_{L^1(\mu_{n,A})}<\infty,
\]
so the local marginals converge in total variation by V84.
\end{corollary}

\begin{proof}
Apply \Cref{thm:e8v98-osc}.  For all sufficiently large \(n\), the total
budget \(E_{n,A}\) is small and
\(e^{2E_{n,A}}-1\leq CE_{n,A}\).  Sum in \(n\).
\end{proof}

\subsection{Second composition route: centered \(L^2\)}
\label{sec:e8v98-L2}

\begin{theorem}[Centered \(L^2\) density bound]
\label{thm:e8v98-L2}
Let
\[
 h=\Delta-\int\Delta\,d\mu.
\]
Assume
\begin{equation}
 \int e^{2|h|}\,d\mu\leq M<\infty.                    \label{eq:e8v98-exp}
\end{equation}
Then
\begin{equation}
 \|r_\Delta-1\|_{L^1(\mu)}
 \leq2\sqrt M\,\|h\|_{L^2(\mu)}.                      \label{eq:e8v98-L2}
\end{equation}
\end{theorem}

\begin{proof}
Put \(Z=\int e^h\,d\mu\).  Jensen's inequality and
\(\int h\,d\mu=0\) give \(Z\geq1\).  Hence
\begin{align*}
 \|r_\Delta-1\|_1
 &={1\over Z}\int|e^h-Z|\,d\mu\\
 &\leq\int|e^h-1|\,d\mu+|Z-1|\\
 &\leq2\int|e^h-1|\,d\mu.
\end{align*}
The elementary inequality
\[
 |e^x-1|\leq|x|e^{|x|}
\]
and Cauchy--Schwarz give
\[
 \int|e^h-1|\,d\mu
 \leq
 \|h\|_2
 \left(\int e^{2|h|}\,d\mu\right)^{1/2}.
\]
This proves \eqref{eq:e8v98-L2}.
\end{proof}

\begin{theorem}[Row-wise \(L^2\) compiler]
\label{thm:e8v98-row-L2}
Write
\[
 h_\rho
 :=
 \delta_\rho-\int\delta_\rho\,d\mu,
\qquad
 h=\sum_\rho\sigma_\rho h_\rho.
\]
If
\begin{align}
 \|h_\rho\|_2&\leq s_\rho,                             \label{eq:e8v98-row-L2}\\
 \int e^{2|h|}\,d\mu&\leq M,                           \label{eq:e8v98-combined-exp}
\end{align}
then
\begin{equation}
 \|r_\Delta-1\|_1
 \leq2\sqrt M\sum_\rho s_\rho.                        \label{eq:e8v98-row-L2-result}
\end{equation}
\end{theorem}

\begin{proof}
The triangle inequality gives
\(\|h\|_2\leq\sum_\rho s_\rho\).  Apply
\Cref{thm:e8v98-L2}.
\end{proof}

\begin{corollary}[Summable \(L^2\) rows]
\label{cor:e8v98-L2}
If, for a fixed local support \(A\),
\[
 \sup_n\int e^{2|h_{n,A}|}\,d\mu_{n,A}\leq M_A
\]
and
\begin{equation}
 \sum_n\sum_\rho
 \left\|
 \delta^\rho_{n,A}
 -\int\delta^\rho_{n,A}\,d\mu_{n,A}
 \right\|_{L^2(\mu_{n,A})}<\infty,                    \label{eq:e8v98-L2-sum}
\end{equation}
then the local marginals converge in total variation.
\end{corollary}

\begin{proof}
Use \Cref{thm:e8v98-row-L2} at every \(n\), sum, and apply the V84
total-variation theorem.
\end{proof}

\begin{firewall}[Separate exponential moments do not combine for free]
\label{fire:e8v98-exp}
Bounds on \(\int e^{2|h_\rho|}\) for each row do not immediately give
\eqref{eq:e8v98-combined-exp} with the same exponent.  H\"older's
inequality requires higher row moments whose exponents depend on the
number of rows.  The combined exponential-integrability budget must be
proved or compiled with those explicit exponents.
\end{firewall}

\subsection{A mixed compiler}
\label{sec:e8v98-mixed}

\begin{proposition}[Bounded rows plus one integrable remainder]
\label{prop:e8v98-mixed}
Decompose
\[
 \Delta=b+u,
 \qquad
 \|b\|_{\rm osc,\infty}\leq E,
\]
and center \(u\) as
\(\widetilde u=u-\int u\,d\mu\).  If
\[
 \int e^{2|\widetilde u|}\,d\mu\leq M,
\]
then
\begin{equation}
 \|r_\Delta-1\|_1
 \leq
 2e^{2E}\sqrt M
 \left(2E+\|\widetilde u\|_2\right).
 \label{eq:e8v98-mixed}
\end{equation}
\end{proposition}

\begin{proof}
Choose \(c\) with \(\|b-c\|_\infty\leq E\), up to an arbitrarily small
error.  Since the mean of \(b-c\) lies in \([-E,E]\),
\[
 \left\|b-\int b\,d\mu\right\|_\infty\leq2E.
\]
The centered total log density is
\[
 h=
 \left(b-\int b\,d\mu\right)+\widetilde u.
\]
Therefore
\[
 \|h\|_2\leq2E+\|\widetilde u\|_2,
 \qquad
 \int e^{2|h|}\,d\mu
 \leq e^{4E}\int e^{2|\widetilde u|}\,d\mu
 \leq e^{4E}M.
\]
Apply \Cref{thm:e8v98-L2} to \(h\).
\end{proof}

\begin{remark}[Intended row assignment]
\label{rem:e8v98-assignment}
Smooth transport and finite Gram factors are natural oscillation rows.
The heat and modified-determinant responses may instead be controlled in a
centered \(L^2\) or localized ideal norm.  The action row usually carries
the main exponential-integrability burden.
\end{remark}

\subsection{From convergence to a summable cofinal subsequence}
\label{sec:e8v98-subsequence}

\begin{lemma}[Summable subsequence in a metric space]
\label{lem:e8v98-subsequence}
If \(x_R\to x\) in a metric space, there is a cofinal subsequence
\(R_n\) such that
\begin{equation}
 d(x_{R_{n+1}},x_{R_n})\leq2^{-n}.                    \label{eq:e8v98-subsequence}
\end{equation}
\end{lemma}

\begin{proof}
Choose \(R_n\) strictly increasing so that
\(d(x_{R_n},x)\leq2^{-n-2}\).  Then
\[
 d(x_{R_{n+1}},x_{R_n})
 \leq2^{-n-3}+2^{-n-2}<2^{-n}.
\]
\end{proof}

\begin{theorem}[One subsequence for countably many supports and rows]
\label{thm:e8v98-diagonal}
Let \(A_1,A_2,\ldots\) be a countable exhausting family of local supports
and let \(\rho=1,\ldots,m\) enumerate finitely many density rows.  Suppose,
after transport to fixed spaces, every row object
\(x^\rho_{R,A_j}\) converges in its declared metric.  Then there is one
cofinal subsequence \(R_n\) such that, for every fixed \(j,\rho\),
\begin{equation}
 \sum_{n\geq j}
 d\left(
 x^\rho_{R_{n+1},A_j},
 x^\rho_{R_n,A_j}
 \right)<\infty.                                      \label{eq:e8v98-diagonal}
\end{equation}
\end{theorem}

\begin{proof}
At stage \(n\), convergence permits choosing \(R_n>R_{n-1}\) so large that
for every \(j\leq n\), every \(\rho\), and every \(R\geq R_n\),
\[
 d(x^\rho_{R,A_j},x^\rho_{\infty,A_j})\leq2^{-n-2}.
\]
For fixed \(j,\rho\) and \(n\geq j\), the triangle inequality bounds the
adjacent distance by \(2^{-n-2}+2^{-n-3}<2^{-n}\).
The geometric series proves \eqref{eq:e8v98-diagonal}.
\end{proof}

\begin{consequence}[Use for the determinant rows]
\label{cons:e8v98-determinant}
The V89 and V97 convergence theorems permit a common cofinal subsequence on
which all fixed-support determinant-response increments are summable,
provided their convergence is in the oscillation, centered \(L^2\), or
other row metric used by the density compiler.
\end{consequence}

\begin{proof}
Apply \Cref{thm:e8v98-diagonal} to the heat, modified-determinant, anomaly,
and finite slow rows.  Pointwise numerical convergence alone is
insufficient when the row is a field-dependent function; convergence must
hold in the stated row metric.
\end{proof}

\subsection{Passing structural properties to the limit}
\label{sec:e8v98-properties}

\begin{proposition}[Closed support conditions]
\label{prop:e8v98-support}
Let \(\mu_n\to\mu\) in total variation on a Polish local space, and let
\(\Sigma\) be closed with \(\mu_n(\Sigma)=1\) for every \(n\).  Then
\(\mu(\Sigma)=1\).
\end{proposition}

\begin{proof}
Total variation convergence gives
\[
 |\mu_n(\Sigma)-\mu(\Sigma)|
 \leq2\|\mu_n-\mu\|_{\rm TV}
\]
under the convention of V84.  Pass to the limit.
\end{proof}

\begin{proposition}[Reflection-positive cylinder inequalities]
\label{prop:e8v98-RP}
Suppose \(\mu_n\to\mu\) in total variation and, for every bounded cylinder
function \(F\) in a fixed reflected local algebra,
\begin{equation}
 \int\overline{\vartheta F}\,F\,d\mu_n\geq0.           \label{eq:e8v98-RPn}
\end{equation}
Then the same inequality holds for \(\mu\).
\end{proposition}

\begin{proof}
The integrand is bounded.  Total variation convergence implies convergence
of its expectation, and the limit of nonnegative real numbers is
nonnegative.
\end{proof}

\begin{corollary}[Bounded Ward identities]
\label{cor:e8v98-Ward}
Any exact identity
\(\int G\,d\mu_n=0\) for one bounded transported cylinder observable \(G\)
passes to the total-variation limit.
\end{corollary}

\begin{proof}
\[
 \left|\int G\,d\mu\right|
 \leq
 \left|\int G\,d\mu-\int G\,d\mu_n\right|
 \leq C\|G\|_\infty\|\mu-\mu_n\|_{\rm TV}\to0.
\]
\end{proof}

\begin{firewall}[Differential Ward identities are not bounded observables]
\label{fire:e8v98-Ward}
Ward identities involving field derivatives, unbounded currents, or
integration by parts require uniform integrability and closability of the
generator.  \Cref{cor:e8v98-Ward} covers only bounded transported cylinder
identities.
\end{firewall}

\subsection{Revised ACDC closure theorem}
\label{sec:e8v98-closure}

\begin{theorem}[Local constrained-density closure]
\label{thm:e8v98-closure}
Fix a countable projective family of local supports.  Suppose:
\begin{enumerate}[label=\textup{(\roman*)},leftmargin=3.8em]
\item every adjacent regulator measure is transported to the declared
common local space and the exact ledger
\eqref{eq:e8v98-ledger} holds;
\item chart, slow-basis, counterterm, heat, and modified-determinant
cocycles are exact;
\item along one cofinal subsequence, either
\eqref{eq:e8v98-osc-sum}, \eqref{eq:e8v98-L2-sum}, or a proved mixed
compiler holds for every support;
\item the regulated marginals are projectively consistent.
\end{enumerate}
Then every local marginal converges in total variation, and the limits
define a unique projective cylinder probability measure.  Closed physical
support conditions, bounded cylinder Ward identities, and reflection
positivity pass to the limit when they hold at every regulator on the same
transported algebra.
\end{theorem}

\begin{proof}
\Cref{cor:e8v98-osc,cor:e8v98-L2,prop:e8v98-mixed} give the summable
adjacent \(L^1\) condition of V84.  The V84 total-variation and projective
theorems give the local limits and cylinder measure.  Apply
\Cref{prop:e8v98-support,prop:e8v98-RP,cor:e8v98-Ward} to the stated
properties.
\end{proof}

\begin{firewall}[Cylinder closure is not full continuum support]
\label{fire:e8v98-topology}
The theorem constructs a cylinder measure.  Concentration on a chosen
distributional metric, connection, spinor, and matter-field topology still
requires tightness there.  It also requires the action/coarea/transport
budgets in the hypotheses; determinant convergence alone does not supply
them.
\end{firewall}

\begin{assessment}[Progress at V98]
\label{ass:e8v98-status}
The determinant programme has been reinserted into the probability measure.
Each field-independent divergence is removed automatically by centering,
while field-dependent rows may be compiled either in oscillation norm or
in centered \(L^2\) with a combined exponential moment.  Convergent row
objects admit one summable cofinal subsequence for countably many supports.
The exact remaining analytic work is now concentrated in the action,
multiplier-coarea, transport, and tightness rows.
\end{assessment}

\begin{programme}[Eighth-series V99: uniform physical-locus coarea]
\label{prog:e8v98-V99}
V99 will attack the multiplier-coarea row.  It will prove stability of a
regular zero set under \(C^1\) perturbations with a uniform smallest
singular-value floor, construct the local normal graph, and bound the
coarea Jacobian ratio.  It will also state the precise failure when
transversality collapses.
\end{programme}

\begin{firewall}[Claim boundary after eighth-series V98]
\label{fire:e8v98-boundary}
V98 proves density-row composition, summable subsequence selection, and
passage of bounded structural properties.  It does not establish the
action exponential moment, uniform multiplier transversality, transport
Jacobian bounds, continuum tightness, unbounded Ward identities, or the
full SM--Einstein continuum.
\end{firewall}

\begin{theorem}[Eighth-series V98 result]
\label{thm:e8v98-result}
Summable centered density rows imply the V84 adjacent \(L^1\) condition.
This holds either through row oscillations or through row \(L^2\) norms plus
a uniform combined exponential moment.  Under exact projective
consistency, the resulting total-variation limits define a cylinder
measure and preserve closed support, bounded Ward, and reflection-positive
cylinder conditions.
\end{theorem}

\begin{proof}
Use \Cref{thm:e8v98-osc,thm:e8v98-row-L2,
thm:e8v98-diagonal,thm:e8v98-closure}.
\end{proof}

\paragraph{Parent-version record.}
V98 was formed from
\texttt{Renewal\_SM\_Einstein\_Continuum\_Research\_Notes\_V97.tex}.
The V97 parent had \(39\,437\) logical source lines, \(1\,402\,938\) bytes,
and SHA--256
\[
 \texttt{FB2831F129AAB6A3898FA083E183719151A39E0F3523C5A2587152569C2EC595}.
\]
The next compulsory companion audit is V100.

% =============================================================================
% END APPENDED EIGHTH-SERIES V98 MATERIAL
% =============================================================================


% =============================================================================
% BEGIN APPENDED EIGHTH-SERIES V99 MATERIAL
% =============================================================================

\clearpage
\section{Eighth-series V99: uniform physical-locus coarea}
\label{sec:v8v99}

\subsection{A flattened reference physical locus}
\label{sec:e8v99-setting}

Work in one finite-regulator parameter chart after the first constraint
elimination.  Let
\[
 x\in K\subset\mathbb R^{d-r},\qquad
 z\in B_\rho(0)\subset\mathbb R^r,
\]
where \(K\) is compact.  Let
\[
 \mu_*:K\times B_\rho(0)\longrightarrow\mathbb R^r
\]
be a \(C^2\) reference multiplier map whose physical locus has been
flattened:
\begin{equation}
 \mu_*(x,0)=0\qquad(x\in K).                           \label{eq:e8v99-flat}
\end{equation}
Put
\begin{equation}
 B(x):=D_z\mu_*(x,0),\qquad
 \sigma_{\min}(B(x))\geq\beta>0.                      \label{eq:e8v99-beta}
\end{equation}
The flattening follows locally from the constant-rank theorem when zero is
a regular value.  V99 quantifies its stability.

\subsection{Uniform contraction graph}
\label{sec:e8v99-graph}

\begin{theorem}[Quantitative normal graph]
\label{thm:e8v99-graph}
Let
\(\mu:K\times B_\rho(0)\to\mathbb R^r\) be \(C^1\).  Suppose, for some
\(0\leq\theta<1\),
\begin{align}
 \sup_{x,z}
 \|D_z\mu(x,z)-B(x)\|
 &\leq\theta\beta,                                     \label{eq:e8v99-Dz}\\
 \varepsilon
 :=
 \sup_x|\mu(x,0)|
 &\leq(1-\theta)\beta\rho.                             \label{eq:e8v99-epsilon}
\end{align}
Then for every \(x\in K\) there is a unique
\(z_\mu(x)\in\overline{B_\rho(0)}\) satisfying
\begin{equation}
 \mu(x,z_\mu(x))=0.                                    \label{eq:e8v99-zero}
\end{equation}
Moreover,
\begin{equation}
 |z_\mu(x)|
 \leq{\varepsilon\over(1-\theta)\beta}.                \label{eq:e8v99-z-bound}
\end{equation}
If the inequality in \eqref{eq:e8v99-epsilon} is strict, the graph lies in
the open ball and \(z_\mu\) is \(C^1\).
\end{theorem}

\begin{proof}
For fixed \(x\), define
\begin{equation}
 \Phi_x(z):=z-B(x)^{-1}\mu(x,z).                       \label{eq:e8v99-Phi}
\end{equation}
The inverse bound from \eqref{eq:e8v99-beta} is
\(\|B(x)^{-1}\|\leq\beta^{-1}\).  By
\eqref{eq:e8v99-Dz},
\[
 \|D_z\Phi_x(z)\|
 =
 \|I-B(x)^{-1}D_z\mu(x,z)\|
 \leq\theta.
\]
Thus \(\Phi_x\) is a contraction.  Also,
\[
 |\Phi_x(z)|
 \leq|\Phi_x(z)-\Phi_x(0)|+|\Phi_x(0)|
 \leq\theta\rho+{\varepsilon\over\beta}
 \leq\rho.
\]
It maps the closed ball to itself, so the contraction theorem gives a
unique fixed point.  The fixed-point equation is
\eqref{eq:e8v99-zero}.  Since
\[
 |z_\mu|
 \leq\theta|z_\mu|+\varepsilon/\beta,
\]
one obtains \eqref{eq:e8v99-z-bound}.  Under strict inclusion the ordinary
implicit-function theorem applies because
\[
 \sigma_{\min}(D_z\mu)\geq(1-\theta)\beta>0,
\]
and gives \(C^1\) dependence on \(x\).
\end{proof}

\begin{corollary}[Exact graph derivative]
\label{cor:e8v99-derivative}
On the open graph,
\begin{equation}
 D_xz_\mu(x)
 =
 -[D_z\mu(x,z_\mu(x))]^{-1}
 D_x\mu(x,z_\mu(x)),                                  \label{eq:e8v99-Dx-z}
\end{equation}
and hence
\begin{equation}
 \|D_xz_\mu(x)\|
 \leq
 {\|D_x\mu(x,z_\mu(x))\|\over(1-\theta)\beta}.
 \label{eq:e8v99-slope}
\end{equation}
\end{corollary}

\begin{proof}
Differentiate \eqref{eq:e8v99-zero} in \(x\) and invert the normal
derivative.  The singular-value bound in the proof of
\Cref{thm:e8v99-graph} gives \eqref{eq:e8v99-slope}.
\end{proof}

\subsection{Stability of the normal graph}
\label{sec:e8v99-graph-stability}

\begin{theorem}[\(C^1\) stability of the physical graph]
\label{thm:e8v99-graph-stability}
Let \(\mu,\widetilde\mu\) satisfy the hypotheses of
\Cref{thm:e8v99-graph} with the same \(B,\beta,\theta,\rho\).  Assume their
first derivatives are Lipschitz with common constant \(L_1\).  Put
\begin{equation}
 \delta
 :=
 \|\mu-\widetilde\mu\|_{C^1(K\times B_\rho)}.          \label{eq:e8v99-delta}
\end{equation}
If both graphs remain a fixed positive distance from the boundary of the
normal ball, then
\begin{align}
 \|z_\mu-z_{\widetilde\mu}\|_{C^0(K)}
 &\leq{\delta\over(1-\theta)\beta},                    \label{eq:e8v99-C0}\\
 \|D_xz_\mu-D_xz_{\widetilde\mu}\|_{C^0(K)}
 &\leq C_{\rm gr}\delta,                               \label{eq:e8v99-C1}
\end{align}
where \(C_{\rm gr}\) depends only on
\(\beta,\theta,L_1\) and uniform \(C^1\) bounds for the maps.
\end{theorem}

\begin{proof}
At the point \(z_{\widetilde\mu}(x)\),
\[
 |\mu(x,z_{\widetilde\mu})|
 =
 |\mu-\widetilde\mu|(x,z_{\widetilde\mu})
 \leq\delta.
\]
Apply the inverse Lipschitz estimate in the normal variable, obtained from
the contraction proof:
\[
 |\mu(x,z)-\mu(x,w)|
 \geq(1-\theta)\beta|z-w|.
\]
With \(z=z_\mu\), \(w=z_{\widetilde\mu}\), this gives
\eqref{eq:e8v99-C0}.

For the derivative, use \eqref{eq:e8v99-Dx-z}.  The inverse identity
\[
 A^{-1}-\widetilde A^{-1}
 =
 A^{-1}(\widetilde A-A)\widetilde A^{-1}
\]
and the common inverse bound control the difference of the two normal
inverses.  The \(C^1\) difference contributes \(O(\delta)\); evaluating the
same derivative at the two graph points costs at most
\(L_1\|z_\mu-z_{\widetilde\mu}\|_\infty=O(\delta)\).
The same estimates apply to the tangential derivative.  Multiplying the
bounded factors proves \eqref{eq:e8v99-C1}.
\end{proof}

\subsection{The coarea singular-value floor}
\label{sec:e8v99-J}

For a row map \(A:\mathbb R^d\to\mathbb R^r\), define
\begin{equation}
 J(A):=\sqrt{\det(AA^*)}.                              \label{eq:e8v99-JA}
\end{equation}
For the multiplier map,
\[
 J_\mu(q)=J(D\mu(q)).
\]

\begin{lemma}[Normal transversality controls the full coarea Jacobian]
\label{lem:e8v99-J-floor}
Under \eqref{eq:e8v99-Dz},
\begin{equation}
 \sigma_{\min}(D\mu(x,z))
 \geq(1-\theta)\beta,                                  \label{eq:e8v99-full-smin}
\end{equation}
and therefore
\begin{equation}
 J_\mu(x,z)\geq[(1-\theta)\beta]^r.                    \label{eq:e8v99-J-floor}
\end{equation}
\end{lemma}

\begin{proof}
Write \(D\mu=(D_x\mu,D_z\mu)\).  Then
\[
 D\mu\,D\mu^*
 =
 D_x\mu\,D_x\mu^*
 +D_z\mu\,D_z\mu^*
 \succeq D_z\mu\,D_z\mu^*.
\]
The smallest singular value of \(D_z\mu\) is at least
\((1-\theta)\beta\) by Weyl's singular-value inequality.  This proves
\eqref{eq:e8v99-full-smin}; the determinant is the product of the \(r\)
singular values, giving \eqref{eq:e8v99-J-floor}.
\end{proof}

\begin{lemma}[Lipschitz logarithmic coarea determinant]
\label{lem:e8v99-logJ}
Let \(A,\widetilde A:\mathbb R^d\to\mathbb R^r\) satisfy
\[
 \sigma_{\min}(A)\geq b,\qquad
 \|A-\widetilde A\|\leq b/2.
\]
Then
\begin{equation}
 |\log J(\widetilde A)-\log J(A)|
 \leq {2r\over b}\|A-\widetilde A\|.                  \label{eq:e8v99-logJ}
\end{equation}
\end{lemma}

\begin{proof}
Set \(A_t=A+t(\widetilde A-A)\).  Weyl's inequality gives
\(\sigma_{\min}(A_t)\geq b/2\).  Differentiation yields
\[
 {d\over dt}\log J(A_t)
 =
 \operatorname{Re}\operatorname{Tr}
 \left[(A_tA_t^*)^{-1}A_t(\widetilde A-A)^*\right].
\]
The operator
\((A_tA_t^*)^{-1}A_t\) has norm
\(\sigma_{\min}(A_t)^{-1}\leq2/b\).
The trace of an endomorphism of \(\mathbb R^r\) is bounded by \(r\) times
its operator norm.  Integrate in \(t\).
\end{proof}

\subsection{Surface density on the physical graph}
\label{sec:e8v99-surface}

The graph embedding is
\[
 T_\mu(x):=(x,z_\mu(x)).
\]
Its induced Hausdorff density relative to \(dx\) is
\begin{equation}
 G_\mu(x)
 :=
 \sqrt{\det\left(I+(D_xz_\mu)^*D_xz_\mu\right)}.
 \label{eq:e8v99-G}
\end{equation}
After pulling the second coarea measure back to \(K\), its geometric
density is
\begin{equation}
 \omega_\mu(x)
 :=
 {G_\mu(x)\over J_\mu(T_\mu(x))}.                     \label{eq:e8v99-omega}
\end{equation}

\begin{lemma}[Graph-area stability]
\label{lem:e8v99-G}
If \(\|Z\|,\|\widetilde Z\|\leq K_1\), then
\begin{equation}
 \left|
 \log\sqrt{\det(I+Z^*Z)}
 -\log\sqrt{\det(I+\widetilde Z^*\widetilde Z)}
 \right|
 \leq C(d,r,K_1)\|Z-\widetilde Z\|.                   \label{eq:e8v99-G-stability}
\end{equation}
\end{lemma}

\begin{proof}
The function
\[
 Z\longmapsto\frac12\log\det(I+Z^*Z)
\]
is smooth on the finite-dimensional ball \(\|Z\|\leq K_1\).  Its
derivative is continuous and hence bounded on that compact ball.  The
mean-value theorem gives \eqref{eq:e8v99-G-stability}.
\end{proof}

\begin{theorem}[Uniform logarithmic coarea ratio]
\label{thm:e8v99-coarea-ratio}
Under the common hypotheses of
\Cref{thm:e8v99-graph-stability}, assume
\(\delta\) is small enough that the derivative matrices on corresponding
graphs differ by at most
\((1-\theta)\beta/2\).  Then
\begin{equation}
 \|\log\omega_\mu-\log\omega_{\widetilde\mu}\|_{L^\infty(K)}
 \leq C_{\rm coa}\delta,                               \label{eq:e8v99-coarea-ratio}
\end{equation}
where \(C_{\rm coa}\) depends only on the common tubular constants and
dimensions.
\end{theorem}

\begin{proof}
\Cref{thm:e8v99-graph-stability} and
\Cref{lem:e8v99-G} give
\[
 \|\log G_\mu-\log G_{\widetilde\mu}\|_\infty
 \leq C\delta.
\]
The derivative matrices at their graph points differ by \(O(\delta)\):
the direct \(C^1\) difference contributes \(\delta\), and movement of the
evaluation point contributes
\(L_1\|z_\mu-z_{\widetilde\mu}\|_\infty\).
Use the singular-value floor
\eqref{eq:e8v99-full-smin} and
\Cref{lem:e8v99-logJ} to obtain
\[
 \|\log J_\mu(T_\mu)
 -\log J_{\widetilde\mu}(T_{\widetilde\mu})\|_\infty
 \leq C\delta.
\]
Subtract the two logarithms in \eqref{eq:e8v99-omega}.
\end{proof}

\begin{corollary}[Summable multiplier-coarea row]
\label{cor:e8v99-summable}
Let \(\mu_n\) share one uniform tubular chart and suppose
\begin{equation}
 \sum_n
 \|\mu_{n+1}-\mu_n\|_{C^1(K\times B_\rho)}<\infty.     \label{eq:e8v99-C1-sum}
\end{equation}
If their first derivatives have one Lipschitz bound and their normal
singular values have one positive floor, then
\begin{equation}
 \sum_n
 \|\log\omega_{\mu_{n+1}}
      -\log\omega_{\mu_n}\|_{L^\infty(K)}
 <\infty.                                              \label{eq:e8v99-coarea-sum}
\end{equation}
Thus the multiplier-coarea entry satisfies the oscillation-row condition
of V98.
\end{corollary}

\begin{proof}
Apply \Cref{thm:e8v99-coarea-ratio} to adjacent maps and sum.  An
\(L^\infty\) difference bounds the oscillation seminorm.
\end{proof}

\subsection{Transported observables and complete coarea density}
\label{sec:e8v99-transport}

\begin{proposition}[Graph transport is a controlled chart map]
\label{prop:e8v99-transport}
Under the hypotheses above, the map
\[
 {\cal T}_{n+1,n}
 :=
 T_{\mu_{n+1}}\circ T_{\mu_n}^{-1}
\]
between the two physical graphs is a \(C^1\) diffeomorphism over \(K\), and
\begin{equation}
 \|{\cal T}_{n+1,n}-I\|_{C^1}
 \leq C
 \|\mu_{n+1}-\mu_n\|_{C^1}.                           \label{eq:e8v99-transport}
\end{equation}
\end{proposition}

\begin{proof}
Both graphs use the same base coordinate \(x\).  The map sends
\((x,z_{\mu_n}(x))\) to
\((x,z_{\mu_{n+1}}(x))\), with inverse obtained by reversing the indices.
\Cref{thm:e8v99-graph-stability} gives the \(C^1\) estimate.
\end{proof}

\begin{remark}[Separation from the ambient transport row]
\label{rem:e8v99-transport}
The graph transport controls movement of the physical multiplier locus.
The full ACDC transport row may also include interpolation of field
variables, bundle frames, regulator cells, and base measures.  Those
additional Jacobians remain separate.
\end{remark}

\subsection{Failure of transversality}
\label{sec:e8v99-failure}

\begin{theorem}[Quantitative collapse at a singular stratum]
\label{thm:e8v99-collapse}
If the best normal singular-value floor is \(\beta_n\downarrow0\), then the
general bounds above necessarily permit
\begin{align}
 \|z_{\mu_n}\|_\infty
 &\sim O(\varepsilon_n/\beta_n),                       \label{eq:e8v99-collapse-z}\\
 J_{\mu_n}^{-1}
 &\sim O(\beta_n^{-r}),                                \label{eq:e8v99-collapse-J}\\
 |\log J_{\mu_{n+1}}-\log J_{\mu_n}|
 &\sim O(\|D\mu_{n+1}-D\mu_n\|/\beta_n).              \label{eq:e8v99-collapse-log}
\end{align}
Thus \(C^1\) convergence of multiplier maps alone cannot control the coarea
row near loss of transversality.
\end{theorem}

\begin{proof}
Equations \eqref{eq:e8v99-z-bound},
\eqref{eq:e8v99-J-floor}, and
\eqref{eq:e8v99-logJ} contain exactly the displayed inverse powers of the
singular-value floor.  No uniform bound remains when \(\beta_n\to0\)
unless the numerator errors vanish at the corresponding stronger rates.
\end{proof}

\begin{firewall}[Multiple roots and singular physical loci]
\label{fire:e8v99-singular}
At \(\beta=0\), the physical locus may branch, acquire excess dimension,
or cease to be a manifold.  The graph theorem and the second coarea formula
do not apply.  One must use a Lyapunov--Schmidt reduction, a stratified
coarea theorem, or prove that the regulated measure of the singular
neighborhood vanishes with a quantitative rate.
\end{firewall}

\subsection{Uniform physical-locus card}
\label{sec:e8v99-card}

\begin{definition}[Uniform multiplier coarea card, UMCC]
\label{def:e8v99-UMCC}
For every local constraint chart, UMCC records:
\begin{enumerate}[label=\textup{(U\arabic*)},leftmargin=4em]
\item a flattened reference physical locus and compact base \(K\);
\item the normal block \(D_z\mu\) and its singular-value floor \(\beta\);
\item a common tubular radius and contraction constant \(\theta<1\);
\item a common Lipschitz bound for first derivatives;
\item the graph \(z_{\mu_n}\) and its \(C^1\) stability;
\item the full coarea floor \eqref{eq:e8v99-J-floor};
\item the pulled-back density \(\omega_{\mu_n}\);
\item the adjacent \(C^1\) sum \eqref{eq:e8v99-C1-sum};
\item the physical-graph transport map; and
\item a separate treatment or measure-exclusion estimate for every
singular stratum.
\end{enumerate}
\end{definition}

\begin{theorem}[Conditional closure of the ACDC coarea row]
\label{thm:e8v99-closure}
If UMCC holds on a regular chart, then the multiplier-coarea density and
physical-graph transport contribute a summable \(L^\infty\) logarithmic row
to the V98 density compiler.  A countable regular cover with compatible
overlap maps gives the same conclusion on its union.  To include an
exceptional complement \(E_n\) in the adjacent \(L^1\) compiler, it is
sufficient in addition that
\begin{equation}
 \sum_n\left[
 \mu_{n,A}(E_n)+\mu_{n+1,A}(E_n)
 \right]<\infty.                                      \label{eq:e8v99-exceptional}
\end{equation}
\end{theorem}

\begin{proof}
The local summability is
\Cref{cor:e8v99-summable,prop:e8v99-transport}.  On overlaps, the V84
coarea and coordinate cocycles identify the densities.  Restricting an
adjacent signed-measure difference to \(E_n\) costs at most the sum of the
two probabilities displayed in \eqref{eq:e8v99-exceptional}.  Summation
therefore adds a finite exceptional-set debit to the V98 \(L^1\) compiler.
\end{proof}

\begin{firewall}[Uniform transversality is still a physical hypothesis]
\label{fire:e8v99-physical}
V82 computes the derivative of the volume and conformal-Killing
multipliers, but it does not prove a regulator-uniform lower singular value
for the multiplier map on the measure-relevant parameter set.  UMCC turns
that missing statement into the explicit number \(\beta\); V99 does not
assume that the physical Standard Model sources automatically supply it.
\end{firewall}

\begin{assessment}[Progress at V99]
\label{ass:e8v99-status}
The multiplier-coarea row now has a quantitative regular-chart theorem.
A normal singular-value floor produces a unique physical graph, a uniform
coarea lower bound, a controlled graph transport, and a Lipschitz
log-density response.  Summable \(C^1\) changes of the multiplier map
therefore give the summable oscillation row required by V98.  Loss of
transversality exposes the exact inverse powers that must be compensated or
stratified.
\end{assessment}

\begin{programme}[Eighth-series V100: audit, archive, and handoff]
\label{prog:e8v99-V100}
V100 will perform the compulsory YM/NS audit, summarize the determinant and
coarea block V81--V99, and select the next action/transport/tightness
acquisition.  It will then be frozen with the next archive suffix and the
active notes will restart at V1 with a compact context and result summary,
in accordance with the user's hundred-version rule.
\end{programme}

\begin{firewall}[Claim boundary after eighth-series V99]
\label{fire:e8v99-boundary}
V99 proves regular-locus graph, coarea, and transport stability under a
uniform singular-value floor.  It does not prove that floor for the
measure-relevant SM--Einstein sequence, control singular strata, establish
the action exponential moment or full transport row, prove continuum
tightness, or construct the full SM--Einstein continuum.
\end{firewall}

\begin{theorem}[Eighth-series V99 result]
\label{thm:e8v99-result}
On a uniform regular multiplier chart,
\[
 \sum_n\|\mu_{n+1}-\mu_n\|_{C^1}<\infty
\]
implies a summable logarithmic coarea-density row.  The constants depend
explicitly on the normal singular-value floor, and the conclusion fails
without further analysis when that floor collapses.
\end{theorem}

\begin{proof}
Use \Cref{thm:e8v99-graph,thm:e8v99-graph-stability,
lem:e8v99-J-floor,thm:e8v99-coarea-ratio,
cor:e8v99-summable}.
\end{proof}

\paragraph{Parent-version record.}
V99 was formed from
\texttt{Renewal\_SM\_Einstein\_Continuum\_Research\_Notes\_V98.tex}.
The V98 parent had \(39\,963\) logical source lines, \(1\,419\,500\) bytes,
and SHA--256
\[
 \texttt{4496729452F0131ECD9DA4C453F29CD6CBC9FB9CAADC303C32556FBC3A090AE2}.
\]
The next version is the compulsory V100 audit and archive boundary.

% =============================================================================
% END APPENDED EIGHTH-SERIES V99 MATERIAL
% =============================================================================


% =============================================================================
% BEGIN APPENDED EIGHTH-SERIES V100 MATERIAL
% =============================================================================

\clearpage
\section{Eighth-series V100: compulsory audit, closure ledger, and rollover}
\label{sec:v8v100}

\subsection{Scope and source integrity}
\label{sec:e8v100-scope}

This version closes the eighth active series at the hundred-version
boundary.  It has three purposes: to perform the compulsory companion
audit, to state exactly what the constraint-density block V81--V99 now
proves, and to choose the first acquisition after the compact restart.
The audit distinguishes filenames from mathematical content.

The newest Yang--Mills files inspected were
\begin{center}
\small\ttfamily
Renewal\_Yang\_Mills\_Mass\_Gap\_Research\_Notes\_V54.tex,\\
Renewal\_Yang\_Mills\_Mass\_Gap\_Research\_Notes\_V55.tex.
\end{center}
They are byte-identical: each has (22,867) logical lines, (822,329)
bytes, and SHA--256
\[
 \texttt{2405E6305086309F2E5ABEC7BABB6B29A77FD3B7E441D213099D9C584E450AAA}.
\]
Both end with the internal V54 record.  Hence no distinct V55 theorem is
counted.  The Navier--Stokes file inspected was
\begin{center}
\small\ttfamily
Renewal\_NS\_Stability\_Research\_Notes\_V26.tex,
\end{center}
with (5,319) logical lines, (278,283) bytes, and SHA--256
\[
 \texttt{82C887F8C2C69E4F28FAD7C3DC8DE5B9099CB5A4BF431EBA1FFF3E91935978AD}.
\]

\subsection{Yang--Mills V52--V54 and the type-correct lesson}
\label{sec:e8v100-YM}

The new YM block contains three relevant advances.  First, its V52 proves
the positive Gaussian forest interpolation identity and an exact connected
tree formula, then shows that the needed analytic derivative seminorm is
incompatible with a nonzero compact cutoff.  It therefore moves upstream
to an operator--Duhamel route in which bounded insertions do not pay cutoff
derivatives.  Second, V53 proves a weighted propagation estimate for a
gapped finite-range quadratic reference and reduces spatial clustering to
a spacetime tree card.  Third, V54 computes exact finite-window matrix-unit
costs and proves that logarithmic occupation windows absorb every fixed
polynomial loss, while explicitly leaving the operator forest card and
high-occupation return open.

\begin{proposition}[Type-correct consequence of the YM audit]
\label{prop:e8v100-YM-transfer}
The YM V52--V54 results do not imply an action, determinant, coarea, or
tightness estimate for the SM--Einstein constrained measure.  They do give
the following design rule for the next SM action argument:
\begin{enumerate}[label=\textup{(\roman*)},leftmargin=3.8em]
\item decompose the adjacent action into bounded or exponentially
integrable local insertions rather than differentiating a compact cutoff
to arbitrarily high order;
\item expose the dependence of the estimate on every truncation window;
and
\item require only a quadratic aggregate of local influence constants when
a martingale or bounded-difference argument permits it.
\end{enumerate}
\end{proposition}

\begin{proof}
The negative part is a type statement: the YM estimates concern a
Hamiltonian/Fock-space cluster expansion, whereas the V98 density compiler
requires an adjacent logarithmic Radon--Nikodym estimate on an already
transported constrained SM--Einstein marginal.  None of the YM theorems
identifies those measures or their action increments.  The three positive
items are precisely the structural features of the proved YM route: the
bounded-insertion replacement, the explicit occupation schedule, and the
fact that a tree or martingale estimate aggregates local costs before the
thermodynamic limit.  They are constraints on an SM proof, not imported SM
conclusions.
\end{proof}

\begin{firewall}[No hidden operator-tree import]
\label{fire:e8v100-YM}
The YM notes still leave their operator spacetime tree card conditional.
V100 neither assumes that card nor relabels it as an SM action theorem.
The logarithmic-window theorem absorbs a fixed polynomial loss only after
the required high-occupation return is proved for that slower window.
\end{firewall}

\subsection{Navier--Stokes V26}
\label{sec:e8v100-NS}

The NS V26 note computes closed-form pointwise rank-one norms for the first
two derivatives of its Oseen kernel and plans certified integral constants.
Those constants are tied to a parabolic kernel and the Navier--Stokes
ancient-profile problem.

\begin{proposition}[NS non-transfer at V100]
\label{prop:e8v100-NS}
No NS V26 theorem supplies a bound for the elliptic SM--Einstein constraint
normal operator, its heat determinant, the physical-locus coarea factor,
or the Euclidean SM action.  Consequently the V100 acquisition order is
unchanged by the NS file.
\end{proposition}

\begin{proof}
The domains, generators, scaling laws, and measures differ.  A numerical
rank-one norm for the Oseen kernel becomes usable here only after a proved
representation of an SM row by that kernel.  No such representation occurs
in either programme.
\end{proof}

\subsection{The V81--V99 closure ledger}
\label{sec:e8v100-ledger}

The recent block began by restoring the momentum quotient and ends with a
quantitative physical-locus coarea theorem.  The logical dependencies can
now be stated without merging hypotheses into conclusions.

\begin{theorem}[Constraint-density ledger at the eighth-series boundary]
\label{thm:e8v100-ledger}
The following statements have been proved in V81--V99.
\begin{enumerate}[label=\textup{(L\arabic*)},leftmargin=4em]
\item On a fixed conformal-Killing quotient, the linear scalar--vector
constraint derivative has the V81 quantitative inverse under its stated
coercivity margin, and the augmented full-space system records compatibility
defects as finite-dimensional multipliers.
\item Under the declared differentiability and uniform inverse bounds, V82
constructs a local nonlinear augmented Einstein-constraint chart and gives
the exact tangent multiplier ledger.
\item V83 applies coarea twice at finite regulator, factors the augmented
Jacobian, and identifies multiplier transversality as a separate condition.
V84 gives the exact chart and adjacent-regulator density cocycles.
\item V86--V90 identify the locality and trace-ideal requirements for
relative determinants.  V91--V97 give a finite-part heat scheme for the
fourth-order normal operator, including exact cocycles, the leading
coefficient
\begin{equation}
 C_{\rm con}
 =
 {\Gamma(3/4)(17+3\sqrt3)\over64\pi^2},               \label{eq:e8v100-Ccon}
\end{equation}
the fixed-volume cancellation of the leading relative divergence, parity
vanishing of the odd closed-manifold coefficients, an exact symbol formula
for (A_2), and a no-spectral-pollution route to positive-time heat-trace
convergence.
\item V98 proves oscillation, centered-(L^2), and mixed density compilers,
a common summable cofinal subsequence for countably many supports and rows,
and passage of closed support, bounded Ward identities, and
reflection-positive cylinder inequalities under total variation.
\item V99 proves that a uniform multiplier normal singular-value floor and
summable adjacent (C^1) changes yield a summable logarithmic coarea row on
regular charts, with an explicit exceptional-set debit.
\end{enumerate}
If, in addition, the conditional cards CNHC, UMCC, the remaining action and
ambient-transport rows, and projective consistency all hold along one
cofinal subsequence, then the V98 theorem constructs a projective cylinder
probability measure.
\end{theorem}

\begin{proof}
Items (L1)--(L3) are the results of V81--V84.  Item (L4) combines the
weighted determinant analysis of V86--V90 with
\Cref{thm:e8v91-finite-part,thm:e8v91-cocycle,
thm:e8v92-two-scale,thm:e8v93-A0,thm:e8v94-odd,
thm:e8v96-A2,thm:e8v97-CNHC}.  Item (L5) is
\Cref{thm:e8v98-osc,thm:e8v98-row-L2,
thm:e8v98-diagonal,thm:e8v98-closure}.  Item (L6) is
\Cref{thm:e8v99-closure}.  The final implication is exactly the conditional
application of \Cref{thm:e8v98-closure}; it does not assert its unpaid
hypotheses.
\end{proof}

\begin{corollary}[The determinant and regular coarea rows are no longer the
first bottleneck]
\label{cor:e8v100-first-bottleneck}
Within the regular-chart branch, the first missing density estimate is the
adjacent action row.  After it come the parts of ambient field transport
not covered by V99 and tightness in a declared continuum topology.
\end{corollary}

\begin{proof}
The determinant and regular coarea entries have conditional closure cards
that feed the V98 compiler.  V98 explicitly lists action, ambient
transport, and tightness as unpaid.  The action appears directly in the
finite-regulator logarithmic density and must be controlled before the
compiler can show adjacent total-variation convergence.  Tightness is a
later upgrade from cylinder convergence to a Radon measure on a chosen
field space.
\end{proof}

\subsection{Upstream choice for the compact restart}
\label{sec:e8v100-choice}

The next step should not assume a global constructive estimate for the
full interacting action.  It should first isolate a reusable implication
from local fluctuation data to the exact V98 density norm.

\begin{decision}[First post-rollover acquisition]
\label{dec:e8v100-action}
The compact series will begin with a conditional-sub-Gaussian martingale
compiler for the centered adjacent action increment.  It will prove a
direct (L^1) bound for the normalized exponential density, derive the
variance proxy from local bounded differences, and state the precise
summability schedule needed across regulators.  Subsequent versions will
test the actual transported Einstein, gauge, Higgs, and fermion action
blocks against that card.
\end{decision}

\begin{assessment}[Position at the archive boundary]
\label{ass:e8v100-status}
The eighth series has moved the programme from a formal constrained
Jacobian to a rigorous conditional density architecture.  The strongest
new facts are the coupled constraint chart, exact coarea and cocycle
ledger, renormalized normal determinant, density-row compilers, and regular
physical-locus stability.  The full construction is still open because
the hypotheses have not yet been proved uniformly for the physical
SM--Einstein regulator family, and no continuum tightness theorem has been
established.
\end{assessment}

\begin{firewall}[Claim boundary after eighth-series V100]
\label{fire:e8v100-boundary}
V100 is an audit and dependency theorem.  It does not prove CNHC or UMCC
for the complete physical regulator sequence, control singular multiplier
strata, establish the interacting action row, complete ambient transport,
prove tightness, construct unbounded Ward generators, perform
Osterwalder--Schrader reconstruction, or construct the full SM--Einstein
continuum.
\end{firewall}

\begin{programme}[Compact ninth-series handoff]
\label{prog:e8v100-handoff}
After validation, this V100 source will be frozen as
\texttt{Renewal\_SM\_Einstein\_Continuum\_Research\_Notes\_V100\_f8.tex}.
The active file will restart at V1 with programme context, an audited
summary of the major results, the conditional-card ledger, and the open
bottlenecks.  V2 will prove the martingale action compiler selected in
\Cref{dec:e8v100-action}.
\end{programme}

\begin{theorem}[Eighth-series V100 result]
\label{thm:e8v100-result}
The V81--V99 block supplies a conditional route from the coupled
Einstein-constraint chart through determinant and regular-coarea
renormalization to a projective cylinder measure.  The required cards are
explicit and remain hypotheses until verified for the physical regulator
family.  The audited next bottleneck is the centered adjacent action row.
\end{theorem}

\begin{proof}
Apply \Cref{thm:e8v100-ledger,cor:e8v100-first-bottleneck} and the
type-correct audit in
\Cref{prop:e8v100-YM-transfer,prop:e8v100-NS}.
\end{proof}

\paragraph{Parent-version record.}
V100 was formed from
\texttt{Renewal\_SM\_Einstein\_Continuum\_Research\_Notes\_V99.tex}.
The V99 parent had (40,504) logical source lines, (1,437,389) bytes,
and SHA--256
\[
 \texttt{B14196491E9C5D5418BD141B59F2255D9A10FF1658A6B5C6F4EB1F66822673EC}.
\]
This is the eighth-series archive boundary.

% =============================================================================
% END APPENDED EIGHTH-SERIES V100 MATERIAL
% =============================================================================

\end{document}
